% Faithful English translation of Jacques Dixmier's 1969 second French edition.
% Source: Les algèbres d'opérateurs dans l'espace hilbertien
%         (Algèbres de von Neumann), 2e édition revue et augmentée.
\RequirePackage{fix-cm}
\documentclass[
  paper=b5,
  fontsize=13pt,
  twoside,
  openany,
  headings=normal,
  chapterprefix=true,
  bibliography=totoc,
  bookmarkpackage=false
]{scrbook}

\usepackage[T1]{fontenc}
\usepackage[english]{babel}
\usepackage{amsmath,amssymb,amsthm,mathtools}
\usepackage{newpxtext}
\usepackage{newpxmath}
\usepackage{microtype}
\usepackage{enumitem}
\usepackage{needspace}
\usepackage{multicol}
\usepackage{xparse}

\setkomafont{disposition}{\normalcolor\rmfamily\bfseries}
\setkomafont{pageheadfoot}{\normalfont\normalsize}
\renewcommand{\thechapter}{\Roman{chapter}}
\renewcommand{\thesection}{\arabic{section}}
\renewcommand*{\sectionformat}{\S~\thesection.\enskip}

\usepackage[
  paper=b5paper,
  left=0.7cm,
  right=0.7cm,
  top=0.91cm,
  bottom=0.91cm,
  includehead,
  headheight=18pt,
  headsep=4pt
]{geometry}

\linespread{1.1}
\allowdisplaybreaks[4]

\usepackage[automark]{scrlayer-scrpage}
\clearpairofpagestyles
\lehead{\pagemark}
\rohead{\pagemark}
\renewcommand*{\chapterpagestyle}{scrheadings}
\pagestyle{scrheadings}

\definecolor{DixmierLink}{HTML}{1D4F91}
\definecolor{DixmierCite}{HTML}{8A2E5D}
\definecolor{DixmierURL}{HTML}{006B6B}
\usepackage[
  colorlinks=true,
  linkcolor=DixmierLink,
  citecolor=DixmierCite,
  urlcolor=DixmierURL,
  filecolor=DixmierURL,
  linktoc=all,
  bookmarksnumbered=true,
  bookmarksopen=true,
  unicode=true,
  pdfencoding=auto,
  psdextra,
  pdfstartpage=1,
  pdfstartview=FitH
]{hyperref}
\usepackage[nameinlink,noabbrev,capitalise]{cleveref}

\hypersetup{
  pdftitle={Les algèbres d'opérateurs dans l'espace hilbertien (Algèbres de von Neumann) / Operator Algebras on Hilbert Space (von Neumann Algebras)},
  pdfauthor={Jacques Dixmier},
  pdfsubject={Faithful English translation of the revised and enlarged second edition (1969)},
  pdfkeywords={operator algebras, Hilbert space, von Neumann algebras}
}

% Formal statements retain the source's independent numbering within each section.
\newtheoremstyle{dixmierformal}
  {6pt}{6pt}{\itshape}{}{\bfseries}{.}{0.5em}{}
\newtheoremstyle{dixmierupright}
  {6pt}{6pt}{\normalfont}{}{\bfseries}{.}{0.5em}{}

\theoremstyle{dixmierformal}
\newtheorem{definition}{Definition}[section]
\newtheorem{theorem}{Theorem}[section]
\newtheorem{proposition}{Proposition}[section]
\newtheorem{lemma}{Lemma}[section]
\newtheorem{corollary}{Corollary}[section]

\renewcommand{\thedefinition}{\arabic{definition}}
\renewcommand{\thetheorem}{\arabic{theorem}}
\renewcommand{\theproposition}{\arabic{proposition}}
\renewcommand{\thelemma}{\arabic{lemma}}
\renewcommand{\thecorollary}{\arabic{corollary}}

% Printed corollary numbers are reset repeatedly in the source.  A separate
% monotonically increasing counter keeps every PDF destination unique.
\newcounter{dixmierhypercorollary}
\AddToHook{env/corollary/before}{\stepcounter{dixmierhypercorollary}}
\renewcommand*{\theHcorollary}{dixmier.corollary.\arabic{dixmierhypercorollary}}

\theoremstyle{dixmierupright}
\newtheorem{remark}{Remark}[section]
\newtheorem{exercise}{Exercise}[section]
\renewcommand{\theremark}{\arabic{remark}}
\renewcommand{\theexercise}{\arabic{exercise}}

\newtheorem*{corollary*}{Corollary}
\newtheorem*{remark*}{Remark}
\newtheorem*{remarks*}{Remarks}

% amsthm's proof label is italic and non-bold, while its body remains upright.
\renewcommand{\qedsymbol}{\rule{0.55ex}{1.5ex}}

\crefname{chapter}{Chapter}{Chapters}
\crefname{section}{Section}{Sections}
\crefname{definition}{Definition}{Definitions}
\crefname{theorem}{Theorem}{Theorems}
\crefname{proposition}{Proposition}{Propositions}
\crefname{lemma}{Lemma}{Lemmas}
\crefname{corollary}{Corollary}{Corollaries}
\crefname{remark}{Remark}{Remarks}
\crefname{exercise}{Exercise}{Exercises}
\crefname{equation}{equation}{equations}

% The source uses single brackets for journal articles and double brackets
% for general treatises. Both forms become live links to the two bibliographies.
\NewDocumentCommand{\ArticleRef}{m}{\hyperref[bib:article:#1]{[#1]}}
\NewDocumentCommand{\TreatiseRef}{m}{\hyperref[bib:treatise:#1]{[[#1]]}}
\NewDocumentCommand{\SourceSectionRef}{m m m}{%
  \hyperref[sec:#1-#2]{Chapter~#1, \S~#2, no.~#3}%
}
\newcounter{translatornote}
\NewDocumentCommand{\TranslatorNote}{m}{%
  \stepcounter{translatornote}%
  \begingroup
  \renewcommand{\thefootnote}{TN\arabic{translatornote}}%
  \footnote{Translator's note: #1}%
  \endgroup
}
\NewDocumentCommand{\NumberedParagraph}{m m}{%
  \par\addvspace{0.55\baselineskip}%
  \phantomsection\noindent\textbf{#1. #2.}\nobreak\ ---\enspace
}

\setlist[enumerate]{leftmargin=*,itemsep=0.35\baselineskip}
\setlist[itemize]{leftmargin=*,itemsep=0.25\baselineskip}

\begin{document}
\pagenumbering{arabic}
\setcounter{page}{1}

\begin{titlepage}
\thispagestyle{empty}
\centering

{\large\bfseries SCIENTIFIC MONOGRAPHS\par}
\vspace{0.25\baselineskip}
{\small PUBLISHED UNDER THE DIRECTION OF M. GASTON JULIA\par}
\vspace{0.8\baselineskip}
{\large\bfseries VOLUME XXV\par}

\vfill

{\Huge\bfseries
LES ALG\`EBRES D'OP\'ERATEURS\\[0.25\baselineskip]
DANS L'ESPACE HILBERTIEN\par}
\vspace{0.45\baselineskip}
{\Large\bfseries (ALG\`EBRES DE VON NEUMANN)\par}

\vspace{0.9\baselineskip}
{\color{DixmierLink}\rule{0.72\textwidth}{0.7pt}\par}
\vspace{0.9\baselineskip}

{\LARGE\bfseries
OPERATOR ALGEBRAS ON HILBERT SPACE\par}
\vspace{0.3\baselineskip}
{\Large\bfseries (VON NEUMANN ALGEBRAS)\par}

\vfill

{\large by\par}
\vspace{0.35\baselineskip}
{\LARGE\bfseries Jacques DIXMIER\par}
\vspace{0.25\baselineskip}
{\small PROFESSOR AT THE FACULTY OF SCIENCES OF PARIS\par}

\vfill

{\large\itshape SECOND EDITION\par}
{\normalsize\itshape revised and enlarged\par}

\vfill

{\large PARIS\par}
{\large GAUTHIER--VILLARS, PUBLISHER\par}
{\large 1969\par}
\end{titlepage}

% Source PDF page 372 (printed p. 363)
\chapter*{Table of Contents}\label{back:source-contents}\label{srcpage:363}
\addcontentsline{toc}{chapter}{Table of Contents}

\begingroup
\newcommand{\SourceTOCChapter}[2]{%
  \par\addvspace{1.2\baselineskip}\begin{center}
  \hyperref[#1]{\bfseries #2}
  \end{center}\addvspace{0.2\baselineskip}}
\newcommand{\SourceTOCSection}[3]{%
  \par\addvspace{0.75\baselineskip}\begin{center}
  \hyperref[#1]{\bfseries\S~#2. --- \itshape #3}
  \end{center}\addvspace{0.1\baselineskip}}
\newcommand{\SourceTOCLine}[4]{%
  \par\begingroup\emergencystretch=2em\sloppy
  \noindent\hangindent=2em\hangafter=1
  \hyperref[#1]{#2. #3}\unskip
  \penalty50\hskip0.6em plus 0.2em\hbox{}\nobreak\dotfill\nobreak
  \hbox{\hyperref[srcpage:#4]{#4}}\par\endgroup}
\newcommand{\SourceTOCBackLine}[3]{%
  \par\noindent\hangindent=1em\hangafter=1
  \hyperref[#1]{#2}\dotfill
  \hyperref[srcpage:#3]{#3}\par}

\SourceTOCChapter{ch:I}{Chapter I.\quad Global Theory}

\SourceTOCSection{sec:I-1}{1}{Definition and First Properties of von Neumann Algebras}
\SourceTOCLine{no:I-1-1}{1}{Commutant and bicommutant}{1}
\SourceTOCLine{no:I-1-2}{2}{Hermitian operators in a von Neumann algebra}{3}
\SourceTOCLine{no:I-1-3}{3}{Unitary operators in a von Neumann algebra}{4}
\SourceTOCLine{no:I-1-4}{4}{Cyclic projections in a von Neumann algebra}{5}
\SourceTOCLine{no:I-1-5}{5}{Homomorphisms}{7}
\SourceTOCLine{no:I-1-6}{6}{Ideals in von Neumann algebras}{9}
\SourceTOCLine{no:I-1-7}{7}{Maximal Abelian von Neumann subalgebras}{12}

\SourceTOCSection{sec:I-2}{2}{Elementary Operations on von Neumann Algebras}
\SourceTOCLine{no:I-2-1}{1}{Induced and reduced von Neumann algebras}{16}
\SourceTOCLine{no:I-2-2}{2}{Product of von Neumann algebras}{19}
\SourceTOCLine{no:I-2-3}{3}{Operators in a tensor product of Hilbert spaces}{21}
\SourceTOCLine{no:I-2-4}{4}{Tensor product of von Neumann algebras}{24}

\SourceTOCSection{sec:I-3}{3}{Density Theorems}
\SourceTOCLine{no:I-3-1}{1}{Topologies on \(\mathcal L(\mathfrak H)\)}{30}
\SourceTOCLine{no:I-3-2}{2}{Comparison between the preceding topologies}{34}
\SourceTOCLine{no:I-3-3}{3}{Linear forms on \(\mathcal L(\mathfrak H)\)}{34}
\SourceTOCLine{no:I-3-4}{4}{Von Neumann's density theorem}{40}
\SourceTOCLine{no:I-3-5}{5}{Kaplansky's density theorem}{43}

\SourceTOCSection{sec:I-4}{4}{Positive Linear Forms}
\SourceTOCLine{no:I-4-1}{1}{Positive linear forms on an involutive algebra of operators}{47}
\SourceTOCLine{no:I-4-2}{2}{Normal positive linear forms on a von Neumann algebra}{50}
\SourceTOCLine{no:I-4-3}{3}{Normal positive linear mappings}{53}
\SourceTOCLine{no:I-4-4}{4}{Structure of normal homomorphisms}{55}
\SourceTOCLine{no:I-4-5}{5}{Application: Isomorphisms of tensor products}{56}
\SourceTOCLine{no:I-4-6}{6}{Support of a normal positive linear form}{57}
\SourceTOCLine{no:I-4-7}{7}{Polar decomposition of a linear form}{59}
\SourceTOCLine{no:I-4-8}{8}{Decomposition of a Hermitian form into positive and negative parts}{64}

% Source PDF page 373 (printed p. 364)
\phantomsection\label{srcpage:364}
\SourceTOCSection{sec:I-5}{5}{Hilbert Algebras}
\SourceTOCLine{no:I-5-1}{1}{Definition of Hilbert algebras}{69}
\SourceTOCLine{no:I-5-2}{2}{The commutation theorem}{70}
\SourceTOCLine{no:I-5-3}{3}{Bounded elements in Hilbert algebras}{72}
\SourceTOCLine{no:I-5-4}{4}{Central elements in Hilbert algebras}{75}
\SourceTOCLine{no:I-5-5}{5}{Elementary operations on Hilbert algebras}{76}

\SourceTOCSection{sec:I-6}{6}{Traces}
\SourceTOCLine{no:I-6-1}{1}{Definition of traces}{81}
\SourceTOCLine{no:I-6-2}{2}{Traces and Hilbert algebras}{85}
\SourceTOCLine{no:I-6-3}{3}{Trace-elements}{89}
\SourceTOCLine{no:I-6-4}{4}{Order relation in the set of traces}{91}
\SourceTOCLine{no:I-6-5}{5}{Application: Isomorphisms of standard von Neumann algebras}{92}
\SourceTOCLine{no:I-6-6}{6}{Normal traces on \(\mathcal L(\mathfrak H)\)}{94}
\SourceTOCLine{no:I-6-7}{7}{First classification of von Neumann algebras}{97}
\SourceTOCLine{no:I-6-8}{8}{Classification and elementary operations}{103}
\SourceTOCLine{no:I-6-9}{9}{Commutant of the tensor product of two semifinite von Neumann algebras}{104}
\SourceTOCLine{no:I-6-10}{10}{Space \(L^1\) defined by a trace}{105}
\SourceTOCLine{no:I-6-11}{11}{Trace and determinant}{107}

\SourceTOCSection{sec:I-7}{7}{Abelian von Neumann Algebras}
\SourceTOCLine{no:I-7-1}{1}{Basic measures}{113}
\SourceTOCLine{no:I-7-2}{2}{Existence of basic measures}{115}
\SourceTOCLine{no:I-7-3}{3}{Structure of Abelian von Neumann algebras}{118}

\SourceTOCSection{sec:I-8}{8}{Discrete von Neumann Algebras}
\SourceTOCLine{no:I-8-1}{1}{Second classification of von Neumann algebras}{121}
\SourceTOCLine{no:I-8-2}{2}{Abelian projections}{123}
\SourceTOCLine{no:I-8-3}{3}{Discrete algebras and elementary operations}{125}
\SourceTOCLine{no:I-8-4}{4}{Definition of the types}{126}
\SourceTOCLine{no:I-8-5}{5}{Complete Hilbert algebras and factors of type I}{126}

\SourceTOCSection{sec:I-9}{9}{Existence of the Different Types of Factors}
\SourceTOCLine{no:I-9-1}{1}{A lemma}{129}
\SourceTOCLine{no:I-9-2}{2}{Construction of certain von Neumann algebras}{130}
\SourceTOCLine{no:I-9-3}{3}{Examples drawn from measure theory}{133}
\SourceTOCLine{no:I-9-4}{4}{Existence of the different types of factors}{135}

\SourceTOCChapter{ch:II}{Chapter II.\quad Reduction of von Neumann Algebras}

\SourceTOCSection{sec:II-1}{1}{Fields of Hilbert Spaces}
\SourceTOCLine{no:II-1-1}{1}{Borel spaces, measures}{139}
\SourceTOCLine{no:II-1-2}{2}{Fields of vectors}{140}
\SourceTOCLine{no:II-1-3}{3}{Measurable fields of Hilbert spaces}{142}
\SourceTOCLine{no:II-1-4}{4}{First properties of measurable fields of Hilbert spaces}{144}
\SourceTOCLine{no:II-1-5}{5}{Fields of square-integrable vectors}{146}

% Source PDF page 374 (printed p. 365)
\phantomsection\label{srcpage:365}
\SourceTOCLine{no:II-1-6}{6}{First properties of Hilbert integrals}{148}
\SourceTOCLine{no:II-1-7}{7}{Measurable fields of subspaces}{150}
\SourceTOCLine{no:II-1-8}{8}{Measurable fields of tensor products}{151}

\SourceTOCSection{sec:II-2}{2}{Fields of Operators}
\SourceTOCLine{no:II-2-1}{1}{Measurable fields of linear mappings}{155}
\SourceTOCLine{no:II-2-2}{2}{Examples}{157}
\SourceTOCLine{no:II-2-3}{3}{Decomposable linear mappings}{158}
\SourceTOCLine{no:II-2-4}{4}{Diagonalizable operators}{161}
\SourceTOCLine{no:II-2-5}{5}{Characterization of decomposable mappings}{163}
\SourceTOCLine{no:II-2-6}{6}{Constant fields of operators}{165}

\SourceTOCSection{sec:II-3}{3}{Fields of von Neumann Algebras}
\SourceTOCLine{no:II-3-1}{1}{Preliminary theorem}{171}
\SourceTOCLine{no:II-3-2}{2}{Measurable fields of von Neumann algebras}{172}
\SourceTOCLine{no:II-3-3}{3}{Relations between a decomposable von Neumann algebra and its components}{174}
\SourceTOCLine{no:II-3-4}{4}{Constant fields of von Neumann algebras}{177}
\SourceTOCLine{no:II-3-5}{5}{Reduction of discrete or continuous von Neumann algebras}{180}
\SourceTOCLine{no:II-3-6}{6}{Measurable fields of homomorphisms}{183}

\SourceTOCSection{sec:II-4}{4}{Fields of Hilbert Algebras}
\SourceTOCLine{no:II-4-1}{1}{Measurable fields of Hilbert algebras}{187}
\SourceTOCLine{no:II-4-2}{2}{Decomposable Hilbert algebras}{188}
\SourceTOCLine{no:II-4-3}{3}{Involution and von Neumann algebras associated with \(\mathfrak A\)}{189}
\SourceTOCLine{no:II-4-4}{4}{Elements bounded relative to \(\mathfrak A\)}{191}
\SourceTOCLine{no:II-4-5}{5}{Elements central relative to \(\mathfrak A\)}{193}
\SourceTOCLine{no:II-4-6}{6}{Uniqueness and existence of the decomposition}{194}

\SourceTOCSection{sec:II-5}{5}{Fields of Traces}
\SourceTOCLine{no:II-5-1}{1}{Measurable fields of traces}{198}
\SourceTOCLine{no:II-5-2}{2}{Decomposition of traces}{199}
\SourceTOCLine{no:II-5-3}{3}{Uniqueness of the decomposition}{203}
\SourceTOCLine{no:II-5-4}{4}{Reduction of semifinite, finite, purely infinite, and properly infinite von Neumann algebras}{204}

\SourceTOCSection{sec:II-6}{6}{Decomposition of a Hilbert Space into a Hilbert Integral}
\SourceTOCLine{no:II-6-1}{1}{Statement of the problem}{207}
\SourceTOCLine{no:II-6-2}{2}{Existence theorems}{208}
\SourceTOCLine{no:II-6-3}{3}{Uniqueness theorems}{211}

\SourceTOCChapter{ch:III}{Chapter III.\quad Further Topics}

\SourceTOCSection{sec:III-1}{1}{Comparison of Projections}
\SourceTOCLine{no:III-1-1}{1}{Comparison of projections}{215}
\SourceTOCLine{no:III-1-2}{2}{Comparability theorem}{217}

% Source PDF page 375 (printed p. 366)
\phantomsection\label{srcpage:366}
\SourceTOCLine{no:III-1-3}{3}{Cyclic projections of \(\mathcal A\) and cyclic projections of \(\mathcal A'\)}{220}
\SourceTOCLine{no:III-1-4}{4}{Applications: I. Properties of totalizing and separating elements}{221}
\SourceTOCLine{no:III-1-5}{5}{Applications: II. Characterization of standard von Neumann algebras}{223}

\SourceTOCSection{sec:III-2}{2}{Classification of Projections}
\SourceTOCLine{no:III-2-1}{1}{Definitions}{229}
\SourceTOCLine{no:III-2-2}{2}{Cyclic projections of \(\mathcal A\) and cyclic projections of \(\mathcal A'\)}{230}
\SourceTOCLine{no:III-2-3}{3}{Finite projections}{231}
\SourceTOCLine{no:III-2-4}{4}{Semifinite projections}{232}
\SourceTOCLine{no:III-2-5}{5}{Properly infinite projections}{234}
\SourceTOCLine{no:III-2-6}{6}{Purely infinite projections}{235}
\SourceTOCLine{no:III-2-7}{7}{Comparison of projections and dimension}{235}

\SourceTOCSection{sec:III-3}{3}{Supplements on Discrete von Neumann Algebras}
\SourceTOCLine{no:III-3-1}{1}{Structure of discrete von Neumann algebras}{239}
\SourceTOCLine{no:III-3-2}{2}{Isomorphisms of discrete von Neumann algebras}{240}

\SourceTOCSection{sec:III-4}{4}{Operator-Valued Traces}
\SourceTOCLine{no:III-4-1}{1}{Definition}{243}
\SourceTOCLine{no:III-4-2}{2}{Traces on \(\widehat{\mathcal Z}_+\)}{244}
\SourceTOCLine{no:III-4-3}{3}{Relations between scalar traces and operator-valued traces}{246}
\SourceTOCLine{no:III-4-4}{4}{Existence and uniqueness theorems for operator-valued traces}{248}

\SourceTOCSection{sec:III-5}{5}{Approximation Theorem}
\SourceTOCLine{no:III-5-1}{1}{Approximation theorem}{251}
\SourceTOCLine{no:III-5-2}{2}{Application: Two-sided ideals of \(\mathcal A\) and ideals of \(\mathcal Z\)}{254}
\SourceTOCLine{no:III-5-3}{3}{Characters of finite von Neumann algebras}{257}

\SourceTOCSection{sec:III-6}{6}{Coupling Operator}
\SourceTOCLine{no:III-6-1}{1}{Coupling operator}{261}
\SourceTOCLine{no:III-6-2}{2}{Properties of the coupling operator}{263}
\SourceTOCLine{no:III-6-3}{3}{Applications: I. Comparison of the strong and ultra-strong, weak and ultra-weak topologies}{265}
\SourceTOCLine{no:III-6-4}{4}{Applications: II. Conditions for an isomorphism to be spatial}{268}

\SourceTOCSection{sec:III-7}{7}{Hyperfinite Factors}
\SourceTOCLine{no:III-7-1}{1}{Factors contained in a finite von Neumann algebra}{270}
\SourceTOCLine{no:III-7-2}{2}{Existence and uniqueness of continuous hyperfinite factors}{271}
\SourceTOCLine{no:III-7-3}{3}{Some inequalities}{273}
\SourceTOCLine{no:III-7-4}{4}{New definition of hyperfinite factors}{276}
\SourceTOCLine{no:III-7-5}{5}{Hyperfinite factors and elementary operations}{280}
\SourceTOCLine{no:III-7-6}{6}{New examples of finite factors}{280}
\SourceTOCLine{no:III-7-7}{7}{Existence of non-hyperfinite finite factors}{283}

\SourceTOCSection{sec:III-8}{8}{New Definition of Finite von Neumann Algebras}
\SourceTOCLine{no:III-8-1}{1}{Statement of the theorem}{288}
\SourceTOCLine{no:III-8-2}{2}{Fundamental projections}{289}
\SourceTOCLine{no:III-8-3}{3}{Weight on the set of fundamental projections}{290}

% Source PDF page 376 (printed p. 367)
\phantomsection\label{srcpage:367}
\SourceTOCLine{no:III-8-4}{4}{Construction of a trace within \(\varepsilon\)}{294}
\SourceTOCLine{no:III-8-5}{5}{End of the proof of the theorem}{296}
\SourceTOCLine{no:III-8-6}{6}{Consequences of the theorem}{298}
\SourceTOCLine{no:III-8-7}{7}{Supplements on tensor products}{303}

\SourceTOCSection{sec:III-9}{9}{Derivations and Automorphisms of von Neumann Algebras}
\SourceTOCLine{no:III-9-1}{1}{Derivations of algebras}{307}
\SourceTOCLine{no:III-9-2}{2}{Derivations of \(C^*\)-algebras: continuity, extension}{308}
\SourceTOCLine{no:III-9-3}{3}{Derivations of von Neumann algebras}{310}
\SourceTOCLine{no:III-9-4}{4}{Automorphisms of von Neumann algebras}{313}

\par\addvspace{0.8\baselineskip}\noindent\hyperref[app:I]{\bfseries Appendix I:}
\SourceTOCLine{no:App-I-1}{1}{Spectrum}{317}
\SourceTOCLine{no:App-I-2}{2}{Spectral measures}{317}
\SourceTOCLine{no:App-I-3}{3}{Extension of the Gelfand isomorphism}{318}

\SourceTOCBackLine{app:II}{Appendix II}{321}

\par\addvspace{0.5\baselineskip}\noindent\hyperref[app:III]{\bfseries Appendix III:}
\SourceTOCLine{no:App-III-1}{1}{Support of an operator}{323}
\SourceTOCLine{no:App-III-2}{2}{Partially isometric operators}{323}
\SourceTOCLine{no:App-III-3}{3}{Polar decomposition of an operator}{324}

\SourceTOCBackLine{app:IV}{Appendix IV}{325}

\par\addvspace{0.5\baselineskip}\noindent\hyperref[app:V]{\bfseries Appendix V:}
\SourceTOCLine{no:App-V-1}{1}{Borel sets}{329}
\SourceTOCLine{no:App-V-2}{2}{Polish spaces}{329}
\SourceTOCLine{no:App-V-3}{3}{Souslin sets}{330}
\SourceTOCLine{no:App-V-4}{4}{Measurability of Souslin sets}{331}
\SourceTOCLine{no:App-V-5}{5}{Existence of measurable mappings}{332}

\SourceTOCBackLine{back:terminology}{Remarks on Terminology}{333}

\par\addvspace{0.5\baselineskip}\noindent\hyperref[back:bibliography]{\bfseries Bibliography:}
\SourceTOCBackLine{bib:journal-articles}{Journal Articles}{337}
\SourceTOCBackLine{bib:general-treatises}{General Treatises}{355}

\SourceTOCBackLine{back:notation-index}{Index of Notation}{357}
\SourceTOCBackLine{back:terminological-index}{Terminological Index}{359}

\begin{center}\bfseries End of the Table of Contents.\end{center}
\endgroup

\chapter*{Preface to the Second Edition}
\addcontentsline{toc}{chapter}{Preface to the Second Edition}
\label{front:preface-second-edition}

In comparison with the first edition, the changes are as follows:

\begin{enumerate}[label=\arabic*\textsuperscript{o}]
\item Results of Kadison, Ringrose, and Sakai on derivations and
automorphisms have been inserted; results of Sakai and Tomita on the polar
decomposition of linear forms; a result of Niiro on the disintegration of
traces; a result of J.~Schwartz on the types of the factors obtained in a
central decomposition; results of Sakai on the noncommutative
Lebesgue--Nikodym theorem and on the type of a tensor product; and a result
of Hugenholtz and St\o rmer on the type of certain factors.

Other results have been incorporated into the exercises. It has not been
possible to take account of the many recent theorems due to theoretical
physicists, in particular those of Araki, Woods, and Powers.

\item The theory of completed Hilbert algebras, useful in group theory, has
been added.

\item In reduction theory, a Borel space has been taken as the base space,
in accordance with Mackey's point of view. This gives statements more
convenient for applications to $C^*$-algebras.

\item Various corrections and improvements of detail have been made (I
thank in this connection the colleagues who sent me comments), the
terminology has been brought up to date, and the bibliography has been
completed.
\end{enumerate}

\clearpage
\chapter*{Introduction}
\addcontentsline{toc}{chapter}{Introduction}
\label{front:introduction}

In this book we study, under the name of \emph{von Neumann algebras}, the
algebras commonly called ``rings of operators'' or ``$W^*$-algebras.'' The
new terminology, suggested by J.~Dieudonn\'e, is amply justified from the
historical point of view.

Some of the results are valid for more general algebras. We have, however,
systematically avoided this kind of generalization (except when it
facilitates the study of von Neumann algebras themselves).
\hyperref[ch:I]{Chapters~I} and \hyperref[ch:II]{II} collect the results
which seem, at the present time, to be the most
useful for applications (although we do not undertake the study of those
applications). \hyperref[ch:III]{Chapter~III}, which is more technical, is
intended chiefly for specialists; it is practically independent of
\hyperref[ch:II]{Chapter~II}.

To limit the length of this book, many interesting questions have been
left aside. Some are sketched in the exercises. The bibliographical
references given throughout the text will help the reader supplement the
documentation. These references are not intended to establish priorities.

The greater part of \hyperref[ch:II]{Chapter~II} is directly inspired by a
forthcoming book of R.~Godement on spectral theory.

\begin{center}
*\quad *\quad *
\end{center}

Elementary results from algebra and general topology are used without
reference. Reading the book also presupposes a fairly thorough knowledge
of topological vector spaces and, in particular, of Hilbert spaces; for
this subject we refer to \TreatiseRef{3}, \TreatiseRef{15}\footnote{Numbers
in single brackets refer to the bibliography of journal articles. Numbers
in double brackets refer to the bibliography of general treatises (in
which, for the reader's convenience, a list of books devoted at least in
part to Hilbert space has been inserted).}, and the
\hyperref[app:I]{Appendices}. With regard
to spectral theory, in principle only the Gelfand--Neumark functional
representation theorem is assumed, for which we refer to
\TreatiseRef{9}; it is nevertheless preferable that the reader be
somewhat familiar with the details of that theory. Reading
\hyperref[ch:II]{Chapter~II} and certain exceptional portions of
\hyperref[ch:I]{Chapters~I} and \hyperref[ch:III]{III} presupposes a
knowledge of integration theory. Finally, a small number of exercises call
on other theories.

\begin{center}
*\quad *\quad *
\end{center}

When the number of the chapter (or of the section) is not specified in a
reference, the chapter (or section) currently under consideration is
meant.

The sign \(\rule{0.55ex}{1.5ex}\) marks the end of a proof.

The reader is invited at this point to consult the
\hyperref[back:terminology]{remarks on terminology} at the end of the book.

% ==================== BEGIN MAIN TRANSLATION ====================

% Source PDF page 006 (printed p. 1)
\chapter{Global Theory}\label{ch:I}\label{srcpage:1}

\section{Definition and First Properties of von Neumann Algebras}
\label{sec:I-1}

\NumberedParagraph{1}{Commutant and bicommutant}\label{no:I-1-1}
Let \(\mathfrak H\) be a complex Hilbert space, and let
\(\mathcal L(\mathfrak H)\) be the set of continuous linear mappings of
\(\mathfrak H\) into \(\mathfrak H\). This set is endowed with the structure
of an algebra over the field of complex numbers. Let \(\mathfrak M\) be an
arbitrary subset of \(\mathcal L(\mathfrak H)\). In accordance with the usual
terminology in algebra, we shall call the \emph{commutant} of \(\mathfrak M\),
and denote by \(\mathfrak M'\), the set of elements of
\(\mathcal L(\mathfrak H)\) that commute with the elements of \(\mathfrak M\).
For example, \(\mathcal L(\mathfrak H)'\) is the set of scalar operators; for a
linear mapping of \(\mathfrak H\) into \(\mathfrak H\) that commutes with all
linear mappings of rank one must leave invariant every one-dimensional
subspace of \(\mathfrak H\), and consequently, as is well known, must be a
homothety. We shall soon be able to give another proof of this fact.

We shall put
\[
  (\mathfrak M')'=\mathfrak M''
  \quad\text{(the bicommutant of \(\mathfrak M\))},
  \qquad
  (\mathfrak M'')'=\mathfrak M''',\ldots.
\]
It is clear that \(\mathfrak M'\) is a subalgebra of
\(\mathcal L(\mathfrak H)\) containing the identity operator of
\(\mathfrak H\) (an operator which we shall denote by \(I_{\mathfrak H}\), or
simply by \(1\) when no confusion is to be feared). We have
\(\mathfrak M\subset\mathfrak M''\). The inclusion
\(\mathfrak M\subset\mathfrak N\) implies
\(\mathfrak M'\supset\mathfrak N'\), and therefore
\(\mathfrak M''\subset\mathfrak N''\); in particular,
\(\mathfrak M'\supset(\mathfrak M')''=\mathfrak M'''\). Since, on the other
hand, \(\mathfrak M'\subset(\mathfrak M')''=\mathfrak M'''\), we see that
\[
  \mathfrak M'=\mathfrak M'''=\mathfrak M^{(5)}=\cdots,
  \qquad
  \mathfrak M\subset\mathfrak M''=\mathfrak M^{(4)}=\cdots.
\]

% Source PDF page 007 (printed p. 2)
\phantomsection\label{srcpage:2}
There is an adjunction operation in \(\mathcal L(\mathfrak H)\). If
\(S\in\mathcal L(\mathfrak H)\), we shall always denote the adjoint of \(S\)
by \(S^*\). We have
\[
  (S+T)^*=S^*+T^*,
  \qquad
  (\lambda S)^*=\overline\lambda S^*,
  \qquad
  (ST)^*=T^*S^*,
  \qquad
  S^{**}=S
\]
(where \(\lambda\) denotes a complex number and \(\overline\lambda\) its
complex conjugate). In other words, \(\mathcal L(\mathfrak H)\) may be
regarded as an involutive algebra. Every subalgebra of
\(\mathcal L(\mathfrak H)\) that is stable under adjunction is called an
\emph{involutive subalgebra} of \(\mathcal L(\mathfrak H)\), or an
\emph{involutive algebra of operators}. Let
\(\mathfrak M\subset\mathcal L(\mathfrak H)\); if \(\mathfrak M\) is stable
under adjunction, \(\mathfrak M'\) is an involutive subalgebra of
\(\mathcal L(\mathfrak H)\).

\begin{definition}\label{def:I-1-1}
A \emph{von Neumann algebra} in \(\mathfrak H\) is an involutive subalgebra
\(\mathcal A\) of \(\mathcal L(\mathfrak H)\) such that
\(\mathcal A=\mathcal A''\).
\end{definition}

The algebra \(\mathcal L(\mathfrak H)\) is a von Neumann algebra. The set of
scalar operators on \(\mathfrak H\) (which we shall often denote by
\(\mathbb C_{\mathfrak H}\)) is a von Neumann algebra. Every von Neumann
algebra in \(\mathfrak H\) contains \(\mathbb C_{\mathfrak H}\).

Let \(\mathfrak M\) be a subset of \(\mathcal L(\mathfrak H)\) stable under
adjunction. The set \(\mathfrak M'\) is a von Neumann algebra. The set
\(\mathfrak M''\) is a von Neumann algebra containing \(\mathfrak M\); and if
\(\mathcal A\) is a von Neumann algebra such that
\(\mathfrak M\subset\mathcal A\), then
\(\mathfrak M''\subset\mathcal A''=\mathcal A\). Thus \(\mathfrak M''\) is
the smallest von Neumann algebra containing \(\mathfrak M\). If now
\(\mathfrak M\) is an arbitrary subset of \(\mathcal L(\mathfrak H)\), put
\(\mathfrak N=\mathfrak M\cup\mathfrak M^*\), where \(\mathfrak M^*\) is the
image of \(\mathfrak M\) under adjunction. The von Neumann algebras containing
\(\mathfrak M\) are identical with the von Neumann algebras containing
\(\mathfrak N\); hence there is a smallest von Neumann algebra containing
\(\mathfrak M\), namely \(\mathfrak N''\). It is called the \emph{von Neumann
algebra generated by \(\mathfrak M\)}.

\begin{proposition}\label{prop:I-1-1}
Let \((\mathcal A_i)_{i\in I}\) be a family of von Neumann algebras. Then
\[
  \mathcal A=\bigcap_{i\in I}\mathcal A_i
\]
is a von Neumann algebra, and \(\mathcal A'\) is the von Neumann algebra
generated by the \(\mathcal A_i'\).
\end{proposition}

\begin{proof}
To say that \(T\in\mathcal A\) is equivalent to saying that \(T\) commutes
with the \(\mathcal A_i'\), hence with \(\bigcup_i\mathcal A_i'\), and hence
with \(\bigl(\bigcup_i\mathcal A_i'\bigr)''\), that is, with the von Neumann
algebra generated by the \(\mathcal A_i'\).
\end{proof}

The \emph{centre} \(\mathcal Z\) of a von Neumann algebra \(\mathcal A\) is
\(\mathcal A\cap\mathcal A'\). It is also the centre of \(\mathcal A'\). Its
commutant \(\mathcal Z'\) is the von Neumann algebra generated by
\(\mathcal A\) and \(\mathcal A'\).

\begin{definition}\label{def:I-1-2}
A \emph{factor} is a von Neumann algebra whose centre contains only the scalar
operators.
\end{definition}

% Source PDF page 008 (printed p. 3)
\phantomsection\label{srcpage:3}
Let \(\mathcal A\) be a von Neumann algebra. To say that \(\mathcal A\) is a
factor is equivalent to saying that \(\mathcal A'\) is a factor, or again that
the von Neumann algebra generated by \(\mathcal A\) and \(\mathcal A'\) is
\(\mathcal L(\mathfrak H)\). The algebra \(\mathcal L(\mathfrak H)\) itself is
a factor.

In contrast to factors, one must consider abelian von Neumann algebras. It is
clear that every operator in an abelian von Neumann algebra is normal (that is,
commutes with its adjoint). Conversely, let \((T_i)_{i\in I}\) be a family of
normal operators; suppose that \(T_i\) commutes with \(T_k\) and with
\(T_k^*\), whatever \(i,k\in I\). Let \(\mathfrak M\) be the set of the
\(T_i\) and the \(T_i^*\). We have
\(\mathfrak M\subset\mathfrak M'\), hence
\(\mathfrak M''\subset\mathfrak M'\), so that \(\mathfrak M''\) is abelian.
In other words, the von Neumann algebra generated by the \(T_i\) is abelian.

In \hyperref[no:I-3-4]{\S~3, no.~4} we shall give topological definitions of
von Neumann algebras.

\noindent\textbf{Bibliography:} \ArticleRef{65}, \ArticleRef{73}.

\NumberedParagraph{2}{Hermitian operators in a von Neumann algebra}
\label{no:I-1-2}
Let \(\mathcal A\) be an involutive subalgebra of
\(\mathcal L(\mathfrak H)\). Every \(T\in\mathcal A\) can be written uniquely
in the form \(T=T_1+iT_2\), where \(T_1\) and \(T_2\) are Hermitian elements
of \(\mathcal A\). Indeed,
\[
  T_1=\frac12(T+T^*)\in\mathcal A,
  \qquad
  T_2=\frac{i}{2}(T^*-T)\in\mathcal A.
\]

\begin{proposition}\label{prop:I-1-2}
Let \(\mathcal A\) be a von Neumann algebra and \(T\) a Hermitian element of
\(\mathcal A\). For every continuous function \(f\) of a real variable,
\(f(T)\) belongs to \(\mathcal A\). The spectral projections of \(T\) belong
to \(\mathcal A\).
\end{proposition}

(Throughout this book the word ``projection'' is used in the sense of
``orthogonal projection.'')

\begin{proof}
It is known (cf. \hyperref[app:I]{Appendix~I}) that \(f(T)\) and the spectral
projections of \(T\) commute with every operator commuting with \(T\), and in
particular with the operators of \(\mathcal A'\). Hence \(f(T)\) and the
spectral projections of \(T\) belong to
\(\mathcal A''=\mathcal A\).
\end{proof}

Since the only projections on \(\mathfrak H\) that commute with
\(\mathcal L(\mathfrak H)\) are plainly \(0\) and \(I_{\mathfrak H}\), we
recover the fact that
\(\mathcal L(\mathfrak H)'=\mathbb C_{\mathfrak H}\).

Let us adopt once and for all the following convention: if
\(\mathfrak M\subset\mathcal L(\mathfrak H)\), then \(\mathfrak M^+\) denotes
the set of Hermitian elements \(\geq0\) of \(\mathfrak M\).

\setcounter{corollary}{0}
\begin{corollary}\label{cor:I-1-2-1}
Let \(\mathcal A\) be a von Neumann algebra and \(T\in\mathcal A\). In order
that \(T\in\mathcal A^+\), it is necessary and sufficient that \(T\) can be
written in the form \(S^*S\), with \(S\in\mathcal A\).
\end{corollary}

\begin{proof}
The condition is plainly sufficient. Conversely, if \(T\in\mathcal A^+\),
then \(T^{1/2}\in\mathcal A\), and
\(T=(T^{1/2})^*T^{1/2}\).
\end{proof}

% Source PDF page 009 (printed p. 4)
\phantomsection\label{srcpage:4}
\begin{corollary}\label{cor:I-1-2-2}
Let \(\mathcal A\) be a von Neumann algebra, and let \(\mathfrak M\) be the
set of its projections. Then \(\mathcal A\) is the von Neumann algebra
generated by \(\mathfrak M\).
\end{corollary}

\begin{proof}
We have \(\mathfrak M''\subset\mathcal A''=\mathcal A\). Let us prove that
\(\mathcal A\subset\mathfrak M''\). Let, therefore, \(S\in\mathcal A\) and
\(T\in\mathfrak M'\), and let us show that \(S\) and \(T\) commute. It is
enough to consider the case in which \(S\) is Hermitian, and it is then enough
to prove that the spectral projections of \(S\) commute with \(T\) (cf.
\hyperref[app:I]{Appendix~I}). But these spectral projections belong to
\(\mathcal A\), and hence to \(\mathfrak M\).
\end{proof}

A von Neumann algebra is said to be \emph{of countable type} if every family
of nonzero, pairwise orthogonal projections of \(\mathcal A\) is countable. In
a Hilbert space with a countable basis, every von Neumann algebra is of
countable type.

\begin{quote}\small
Von Neumann algebras of countable type are called, in English,
``\(\sigma\)-finite,'' or ``countably decomposable.'' The theory of von
Neumann algebras generalizes, to a certain extent, the theory of integration
\ArticleRef{101}. It is known that, in the latter theory, countability
hypotheses on the base space are sometimes useful; the consideration of von
Neumann algebras of countable type is useful in an analogous way.
\end{quote}

\noindent\textbf{Bibliography:} \ArticleRef{73}.

\NumberedParagraph{3}{Unitary operators in a von Neumann algebra}
\label{no:I-1-3}
\begin{proposition}\label{prop:I-1-3}
Let \(\mathcal A\) be a von Neumann algebra. Every element of \(\mathcal A\)
is a linear combination of unitary elements of \(\mathcal A\).
\end{proposition}

\begin{proof}
By \hyperref[no:I-1-2]{no.~2}, it is enough to consider the case of a
Hermitian operator \(T\in\mathcal A\) such that \(\lVert T\rVert\leq1\). We
then have \(I_{\mathfrak H}-T^2\geq0\). Put
\[
  U=T+i(I_{\mathfrak H}-T^2)^{1/2}.
\]
We have \(U\in\mathcal A\), and
\begin{align*}
  UU^*=U^*U
    &=\bigl(T+i(I_{\mathfrak H}-T^2)^{1/2}\bigr)
      \bigl(T-i(I_{\mathfrak H}-T^2)^{1/2}\bigr)\\
    &=T^2+I_{\mathfrak H}-T^2=I_{\mathfrak H},
\end{align*}
so \(U\) is unitary. On the other hand,
\(T=\tfrac12(U+U^*)\).
\end{proof}

\begin{corollary*}\phantomsection\label{cor:I-1-3}
Let \(T\in\mathcal L(\mathfrak H)\). In order that \(T\in\mathcal A\), it is
necessary and sufficient that
\[
  U'TU'^{-1}=T
\]
for every unitary operator \(U'\in\mathcal A'\).
\end{corollary*}

\begin{proof}
The condition says that \(T\) commutes with the unitary operators of
\(\mathcal A'\), and hence, by \cref{prop:I-1-3}, with all the operators of
\(\mathcal A'\).
\end{proof}

The corollary will often be used without explicit reference, in the following
form: an operator constructed from operators of \(\mathcal A\) by a
``unitarily invariant'' process lies in \(\mathcal A\). For example, let
\(\mathcal F\subset\mathcal A^+\) be an increasing directed set bounded above,
and let \(T\) be its supremum in \(\mathcal L(\mathfrak H)\) (cf.
\hyperref[app:II]{Appendix~II}); then \(T\in\mathcal A^+\). As another example,
if \(S=U\lvert S\rvert\) is the polar decomposition of an element \(S\) of
\(\mathcal A\)
% Source PDF page 010 (printed p. 5)
\phantomsection\label{srcpage:5}
(cf. \hyperref[app:III]{Appendix~III}), then \(U\) and \(\lvert S\rvert\)
belong to \(\mathcal A\). The support of \(S\) (cf.
\hyperref[app:III]{Appendix~III}) belongs to \(\mathcal A\).

\noindent\textbf{Bibliography:} \ArticleRef{65}, \ArticleRef{66}.

\NumberedParagraph{4}{Cyclic projections in a von Neumann algebra}
\label{no:I-1-4}
Let \(\mathcal A\) be an involutive subalgebra of
\(\mathcal L(\mathfrak H)\), and let \(\mathfrak X\) be a closed vector
subspace of \(\mathfrak H\). To say that \(\mathfrak X\) is invariant under
\(\mathcal A\) amounts to saying that \(P_{\mathfrak X}\in\mathcal A'\)
(here and throughout this book, \(P_{\mathfrak X}\) denotes the projection
onto \(\mathfrak X\)). Indeed, if \(P_{\mathfrak X}\in\mathcal A'\), then,
for \(T\in\mathcal A\) and \(x\in\mathfrak X\),
\(Tx=TP_{\mathfrak X}x=P_{\mathfrak X}Tx\), so that
\(Tx\in\mathfrak X\). Conversely, suppose that \(\mathfrak X\) is invariant
under \(\mathcal A\). Then, for every \(T\in\mathcal A\),
\[
  TP_{\mathfrak X}(\mathfrak H)\subset\mathfrak X,
\]
and hence
\[
  P_{\mathfrak X}TP_{\mathfrak X}=TP_{\mathfrak X};
\]
therefore
\[
  P_{\mathfrak X}T
   =(T^*P_{\mathfrak X})^*
   =(P_{\mathfrak X}T^*P_{\mathfrak X})^*
   =P_{\mathfrak X}TP_{\mathfrak X}
   =TP_{\mathfrak X}.
\]

If \(\mathfrak M\) is an arbitrary subset of \(\mathfrak H\), denote by
\(\mathfrak X_{\mathfrak M}^{\mathcal A}\) the closed vector subspace of
\(\mathfrak H\) generated by the \(Tx\), where \(T\in\mathcal A\) and
\(x\in\mathfrak M\). This subspace is plainly invariant under
\(\mathcal A\). Thus the corresponding projection, denoted by
\(E_{\mathfrak M}^{\mathcal A}\), is a projection of \(\mathcal A'\). If
\(\mathfrak M\) consists of a single element \(x\), we use the notations
\(\mathfrak X_x^{\mathcal A}\) and \(E_x^{\mathcal A}\); a projection such
as \(E_x^{\mathcal A}\) is called a \emph{cyclic projection} of
\(\mathcal A'\).

\begin{definition}\label{def:I-1-3}
Let \(\mathcal A\) be an involutive subalgebra of
\(\mathcal L(\mathfrak H)\), and let \(\mathfrak M\) be a subset of
\(\mathfrak H\). The set \(\mathfrak M\) is said to be \emph{totalizing} for
\(\mathcal A\) if the union of the \(T(\mathfrak M)\), \(T\in\mathcal A\),
is a total set in \(\mathfrak H\), that is, if
\(\mathfrak X_{\mathfrak M}^{\mathcal A}=\mathfrak H\). The set
\(\mathfrak M\) is said to be \emph{separating} for \(\mathcal A\) if the
conditions \(T\in\mathcal A\), \(T(\mathfrak M)=0\), imply \(T=0\).

If \(\mathfrak M=\{x\}\), where \(x\in\mathfrak H\), then \(x\) is said to
be totalizing (respectively, separating) for \(\mathcal A\). This means that
the \(Tx\), \(T\in\mathcal A\), are dense in \(\mathfrak H\) (respectively,
that the conditions \(T\in\mathcal A\), \(Tx=0\), imply \(T=0\)).
\end{definition}

\begin{proposition}\label{prop:I-1-4}
Let \(\mathcal A\) be an involutive subalgebra of
\(\mathcal L(\mathfrak H)\) containing \(I_{\mathfrak H}\), and let
\(\mathfrak M\) be an arbitrary subset of \(\mathfrak H\). Then
\(\mathfrak X_{\mathfrak M}^{\mathcal A}\) is the smallest of the closed
vector subspaces \(\mathfrak N\) of \(\mathfrak H\) such that
\(\mathfrak M\subset\mathfrak N\) and \(P_{\mathfrak N}\in\mathcal A'\).
\end{proposition}

\begin{proof}
We have \(\mathfrak M\subset\mathfrak X_{\mathfrak M}^{\mathcal A}\) because
\(I_{\mathfrak H}\in\mathcal A\). On the other hand, if
\(\mathfrak M\subset\mathfrak N\) and \(P_{\mathfrak N}\in\mathcal A'\),
then \(\mathfrak N\) is invariant under \(\mathcal A\), and
\(\mathfrak X_{\mathfrak M}^{\mathcal A}\subset\mathfrak N\).
\end{proof}

% Source PDF page 011 (printed p. 6)
\phantomsection\label{srcpage:6}
\begin{proposition}\label{prop:I-1-5}
Let \(\mathcal A\) be an involutive subalgebra of
\(\mathcal L(\mathfrak H)\) containing \(I_{\mathfrak H}\), and let
\(\mathfrak M\) be an arbitrary subset of \(\mathfrak H\). In order that
\(\mathfrak M\) be totalizing for \(\mathcal A\), it is necessary and
sufficient that \(\mathfrak M\) be separating for \(\mathcal A'\).
\end{proposition}

\begin{proof}
If \(\mathfrak M\) is separating for \(\mathcal A'\), the evident equality
\[
  (I_{\mathfrak H}-E_{\mathfrak M}^{\mathcal A})(\mathfrak M)=0
\]
implies \(I_{\mathfrak H}-E_{\mathfrak M}^{\mathcal A}=0\), and hence
\(\mathfrak X_{\mathfrak M}^{\mathcal A}=\mathfrak H\). Conversely, suppose
that \(\mathfrak X_{\mathfrak M}^{\mathcal A}=\mathfrak H\). Then the
conditions \(T'\in\mathcal A'\), \(T'(\mathfrak M)=0\), imply, for every
\(T\in\mathcal A\),
\[
  T'T(\mathfrak M)=TT'(\mathfrak M)=0,
\]
and hence \(T'(\mathfrak H)=0\), so that \(T'=0\). Thus \(\mathfrak M\) is
separating for \(\mathcal A'\).
\end{proof}

\begin{corollary*}\phantomsection\label{cor:I-1-5}
Let \(\mathcal A\) be a von Neumann algebra. In order that \(\mathfrak M\)
be separating for \(\mathcal A\), it is necessary and sufficient that
\(\mathfrak M\) be totalizing for \(\mathcal A'\).
\end{corollary*}

\begin{proof}
Interchange the roles of \(\mathcal A\) and \(\mathcal A'\) in
\cref{prop:I-1-5}.
\end{proof}

\begin{proposition}\label{prop:I-1-6}
Let \(\mathcal A\) be a von Neumann algebra. In order that \(\mathcal A\) be
of countable type, it is necessary and sufficient that \(\mathcal A\) possess
a countable separating set.
\end{proposition}

\begin{proof}
Let \(\mathfrak M\) be a countable set separating for \(\mathcal A\), and let
\((E_i)_{i\in I}\) be a family of pairwise orthogonal projections of
\(\mathcal A\). For every \(x\in\mathfrak M\), we have \(E_i x=0\), except
for a countable family of indices \(i\). Hence \(E_i(\mathfrak M)=0\), and
consequently \(E_i=0\), except for a countable family of indices \(i\). This
proves that \(\mathcal A\) is of countable type.

Conversely, suppose that \(\mathcal A\) is of countable type. Let
\((y_x)_{x\in K}\) be a maximal family of nonzero elements of \(\mathfrak H\)
such that the \(E_{y_x}^{\mathcal A'}\) are pairwise orthogonal. We have
\[
  \sum_{x\in K}E_{y_x}^{\mathcal A'}=I_{\mathfrak H}
\]
because the family \((y_x)_{x\in K}\) is maximal, and \(K\) is countable
because \(\mathcal A\) is of countable type. Thus the set of the \(y_x\) is a
countable separating set.
\end{proof}

Let \(\mathcal A\) be a von Neumann algebra, \(\mathcal Z\) its centre, and
\(T\) an element of \(\mathcal A\). For every projection \(G\) of
\(\mathcal Z\) such that \(TG=T\), we have \(T^*G=(TG)^*=T^*\). Thus the
infimum \(F\) of the projections of \(\mathcal Z\) majorizing the support of
\(T\) (cf. \hyperref[app:III]{Appendix~III}), which is the smallest projection
of \(\mathcal Z\) majorizing the support of \(T\), is also the smallest
projection of \(\mathcal Z\) majorizing the support of \(T^*\). We have
\(TF=T\), and every projection \(E\) of \(\mathcal Z\) such that \(TE=0\) is
orthogonal to \(F\). The projection \(F\) is called the \emph{central support}
of \(T\).

\begin{proposition}\label{prop:I-1-7}
Let \(\mathfrak M\) be a subset of \(\mathfrak H\), and put
\(\mathfrak N=\mathfrak X_{\mathfrak M}^{\mathcal A}\). Then
\[
  \mathfrak X_{\mathfrak N}^{\mathcal A'}
  =\mathfrak X_{\mathfrak M}^{\mathcal Z'}.
\]
\end{proposition}

\begin{proof}
We have
\(\mathfrak N\subset\mathfrak X_{\mathfrak M}^{\mathcal Z'}\), and hence
\(\mathfrak X_{\mathfrak N}^{\mathcal A'}
 \subset\mathfrak X_{\mathfrak M}^{\mathcal Z'}\). Since
\(\mathfrak X_{\mathfrak M}^{\mathcal Z'}\) is the smallest of the closed
vector subspaces \(\mathfrak R\) of \(\mathfrak H\) containing
\(\mathfrak M\)
% Source PDF page 012 (printed p. 7)
\phantomsection\label{srcpage:7}
and such that \(P_{\mathfrak R}\in\mathcal Z\), it is enough to show that
\(E_{\mathfrak N}^{\mathcal A'}\in\mathcal Z\), in other words that
\(\mathfrak X_{\mathfrak N}^{\mathcal A'}\) is invariant under
\(\mathcal A\) and under \(\mathcal A'\). It is clear that
\(\mathfrak X_{\mathfrak N}^{\mathcal A'}\) is invariant under
\(\mathcal A'\). Now let \(x\in\mathfrak N\), \(T'\in\mathcal A'\), and
\(T\in\mathcal A\). We have
\[
  T(T'x)=T'(Tx)\in T'(\mathfrak N)
\]
(because \(\mathfrak N\) is invariant under \(\mathcal A\)); hence
\(T(\mathcal A'x)\subset
\mathfrak X_{\mathfrak N}^{\mathcal A'}\), and therefore
\(T(\mathfrak X_{\mathfrak N}^{\mathcal A'})\subset
\mathfrak X_{\mathfrak N}^{\mathcal A'}\).
\end{proof}

\setcounter{corollary}{0}
\begin{corollary}\label{cor:I-1-7-1}
Let \(\mathfrak N\) be a closed vector subspace of \(\mathfrak H\) such that
\(P_{\mathfrak N}\in\mathcal A\). Then \(E_{\mathfrak N}^{\mathcal A}\) is
the central support of \(P_{\mathfrak N}\).
\end{corollary}

\begin{proof}
Since \(\mathfrak X_{\mathfrak N}^{\mathcal A'}=\mathfrak N\), we have
\(\mathfrak X_{\mathfrak N}^{\mathcal A}
=\mathfrak X_{\mathfrak N}^{\mathcal Z'}\) by \cref{prop:I-1-7}. Thus
\(\mathfrak X_{\mathfrak N}^{\mathcal A}\) is the smallest of the closed
vector subspaces \(\mathfrak R\) of \(\mathfrak H\) containing
\(\mathfrak N\) and such that \(P_{\mathfrak R}\in\mathcal Z\).
\end{proof}

\begin{corollary}\label{cor:I-1-7-2}
Let \(\mathfrak M\) be a subset of \(\mathfrak H\). The projections
\(E_{\mathfrak M}^{\mathcal A'}\in\mathcal A\) and
\(E_{\mathfrak M}^{\mathcal A}\in\mathcal A'\) have the same central
support, namely \(E_{\mathfrak M}^{\mathcal Z'}\).
\end{corollary}

\begin{proof}
This follows at once from \cref{prop:I-1-7} and
\cref{cor:I-1-7-1}.
\end{proof}

\begin{corollary}\label{cor:I-1-7-3}
The following conditions are equivalent:
\begin{enumerate}[label=(\roman*)]
\item \(\mathcal A\) is a factor;
\item whatever the nonzero elements \(R,S\) of \(\mathcal A\), there exists
\(T\in\mathcal A\) such that \(RTS\neq0\).
\end{enumerate}
\end{corollary}

\begin{proof}
\((ii)\Rightarrow(i)\). If \(\mathcal A\) is not a factor, there exists in
\(\mathcal Z\) a projection \(P\) such that \(P\neq0\) and
\(I_{\mathfrak H}-P\neq0\) (\hyperref[cor:I-1-2-2]{Corollary~2 of
Proposition~2}). For every \(T\in\mathcal A\),
\[
  PT(I_{\mathfrak H}-P)=P(I_{\mathfrak H}-P)T=0.
\]

\((i)\Rightarrow(ii)\). Suppose that there exist \(R,S\in\mathcal A\) such
that \(R\neq0\), \(S\neq0\), and \(RTS=0\) for every \(T\in\mathcal A\).
Let \(\mathfrak Y\) be the closure of \(S(\mathfrak H)\), and put
\(P=E_{\mathfrak Y}^{\mathcal A}\). By \cref{cor:I-1-7-1}, \(P\) is a
nonzero projection of \(\mathcal Z\). On the other hand, for every
\(T\in\mathcal A\), \(R\) vanishes on \(TS(\mathfrak H)\), and therefore on
\(T(\mathfrak Y)\); hence \(R\) vanishes on
\(\mathfrak X_{\mathfrak Y}^{\mathcal A}\), so that \(RP=0\). Thus
\(P\neq I_{\mathfrak H}\), and \(\mathcal A\) is not a factor.
\end{proof}

\begin{quote}\small
\Cref{prop:I-1-4} can plainly fail if \(I_{\mathfrak H}\notin\mathcal A\)
(take \(\mathcal A=\{0\}\)). The \hyperref[cor:I-1-5]{corollary to
Proposition~5} can fail if \(\mathcal A\) is an arbitrary involutive
subalgebra of \(\mathcal L(\mathfrak H)\) (Exercise~3). The existence of
separating or totalizing elements will be studied repeatedly (and already in
\hyperref[no:I-2-1]{\S~2, no.~1}). In English, ``cyclic'' is used for
``totalizing.''
\end{quote}

\noindent\textbf{Bibliography:} \ArticleRef{65}, \ArticleRef{73}.

\NumberedParagraph{5}{Homomorphisms}\label{no:I-1-5}
Let \(\mathcal A\) and \(\mathcal B\) be von Neumann algebras. A mapping
\(\Phi\) from \(\mathcal A\) into \(\mathcal B\) is called a homomorphism if
it is linear, if
\[
  \Phi(ST)=\Phi(S)\Phi(T)
  \qquad\text{for \(S\in\mathcal A\), \(T\in\mathcal A\)},
\]
% Source PDF page 013 (printed p. 8)
\phantomsection\label{srcpage:8}
and if
\[
  \Phi(S^*)=\Phi(S)^*
  \qquad\text{for \(S\in\mathcal A\)},
\]
in other words, if \(\Phi\) is a homomorphism for the involutive algebra
structures of \(\mathcal A\) and \(\mathcal B\). A mapping \(\Phi'\) from
\(\mathcal A\) into \(\mathcal B\) is called an antihomomorphism if it is
linear, if
\[
  \Phi'(ST)=\Phi'(T)\Phi'(S)
  \qquad\text{for \(S\in\mathcal A\), \(T\in\mathcal A\)},
\]
and if
\[
  \Phi'(S^*)=\Phi'(S)^*
  \qquad\text{for \(S\in\mathcal A\)}.
\]

\begin{proposition}\label{prop:I-1-8}
Let \(\mathcal A\) and \(\mathcal B\) be von Neumann algebras, and let
\(\Phi\) be a homomorphism or an antihomomorphism from \(\mathcal A\) into
\(\mathcal B\).
\begin{enumerate}[label=(\roman*)]
\item \(\Phi(\mathcal A^+)\subset\mathcal B^+\).
\item If \(E\) is a projection of \(\mathcal A\), then \(\Phi(E)\) is a
projection of \(\mathcal B\).
\item For every \(S\in\mathcal A\),
\(\lVert\Phi(S)\rVert\leq\lVert S\rVert\); if \(\Phi\) is injective, then in
fact \(\lVert\Phi(S)\rVert=\lVert S\rVert\).
\item If \(S\) is a Hermitian operator of \(\mathcal A\), and if \(f\) is a
continuous numerical function of a real variable such that \(f(0)=0\), then
\(\Phi(S)\) is a Hermitian operator of \(\mathcal B\), and
\(\Phi(f(S))=f(\Phi(S))\).
\end{enumerate}
\end{proposition}

\begin{proof}
We shall consider the case in which \(\Phi\) is a homomorphism. The case of an
antihomomorphism is treated similarly.

If \(S\in\mathcal A^+\), then \(S=T^*T\) for some \(T\in\mathcal A\), and
hence
\[
  \Phi(S)=\Phi(T)^*\Phi(T)\in\mathcal B^+.
\]
This proves (i).

If \(E\) is a projection of \(\mathcal A\), then \(E=E^2=E^*\), and hence
\[
  \Phi(E)=\Phi(E)^2=\Phi(E)^*,
\]
so \(\Phi(E)\) is a projection of \(\mathcal B\). This proves (ii).

For every \(S\in\mathcal A\), we have
\(S^*S\leq\lVert S\rVert^2 I\), and therefore
\[
  \Phi(S)^*\Phi(S)
  \leq\lVert S\rVert^2\Phi(I)
  \leq\lVert S\rVert^2 I
\]
[because \(\Phi(I)\) is a projection]. Hence
\[
  \lVert\Phi(S)\rVert^2\leq\lVert S\rVert^2,
\]
which proves the first assertion of (iii).

Assertion (iv) is clear if \(f\) is a polynomial (without constant term). The
general case follows by continuity (using the relation
\(\lVert\Phi(S)\rVert\leq\lVert S\rVert\)).

Finally, let us prove that, if \(\Phi\) is injective, then
\[
  \lVert\Phi(S)\rVert=\lVert S\rVert.
\]
% Source PDF page 014 (printed p. 9)
\phantomsection\label{srcpage:9}
Since \(\lVert S\rVert^2=\lVert S^*S\rVert\), it is enough to consider the
case in which \(S\in\mathcal A^+\). If we had
\[
  \lVert\Phi(S)\rVert<\lVert S\rVert
\]
for an \(S\in\mathcal A^+\), there would exist a continuous numerical
function \(f\) of a real variable such that \(f(0)=0\), with \(f(S)\neq0\)
and
\[
  \Phi(f(S))=f(\Phi(S))=0,
\]
which is a contradiction.
\end{proof}

If \(\Phi\) is a bijective homomorphism (respectively, antihomomorphism) from
\(\mathcal A\) onto \(\mathcal B\), then \(\Phi^{-1}\) is a homomorphism
(respectively, antihomomorphism). In this case \(\Phi\) is called an
\emph{isomorphism} (respectively, \emph{antiisomorphism}) of \(\mathcal A\)
onto \(\mathcal B\), or sometimes an \emph{algebraic isomorphism}
(respectively, \emph{algebraic antiisomorphism}); one also says that
\(\mathcal A\) and \(\mathcal B\) are \emph{isomorphic} (respectively,
\emph{antiisomorphic}), or sometimes \emph{algebraically isomorphic}
(respectively, \emph{algebraically antiisomorphic}).

Let \(\mathfrak H\) and \(\mathfrak K\) be two complex Hilbert spaces, let
\(U\) be an isomorphism of \(\mathfrak H\) onto \(\mathfrak K\) (that is, an
isometric linear mapping of \(\mathfrak H\) onto \(\mathfrak K\)), and let
\(\mathcal A\) be a von Neumann algebra in \(\mathfrak H\). The mapping
\[
  S\longmapsto USU^{-1}
\]
is an isomorphism of \(\mathcal A\) onto a von Neumann algebra \(\mathcal B\)
in \(\mathfrak K\). Such an isomorphism is said to be \emph{spatial}; one
says that \(\mathcal A\) and \(\mathcal B\) are \emph{spatially isomorphic}.
Now let \(V\) be a bijective mapping of \(\mathfrak H\) onto \(\mathfrak K\)
such that\TranslatorNote{The French text omits \(V\) from the right-hand side
of the first identity; the mathematically necessary factors \(Vx\) and
\(Vy\) have been restored.}
\[
  V(\lambda x+\mu y)
    =\overline\lambda Vx+\overline\mu Vy,
  \qquad
  (Vx\mid Vy)=(y\mid x)
  \quad
  \text{for \(x,y\in\mathfrak H\), \(\lambda,\mu\in\mathbb C\)}.
\]
The mapping
\[
  S\longmapsto VS^*V^{-1}
\]
is an antiisomorphism of \(\mathcal A\) onto a von Neumann algebra
\(\mathcal C\) in \(\mathfrak K\). Such an antiisomorphism is said to be
\emph{spatial}; one says that \(\mathcal A\) and \(\mathcal C\) are
\emph{spatially antiisomorphic}.

Let \(\mathfrak H^c\) be the conjugate Hilbert space of \(\mathfrak H\), that
is, the space \(\mathfrak H\) endowed with the operations
\[
  (\lambda,x)\longmapsto\overline\lambda x,
  \qquad
  (x,y)\longmapsto x+y,
\]
and with the scalar product \((x,y)\longmapsto(y\mid x)\). The identity
mapping of \(\mathfrak H\) onto \(\mathfrak H^c\) can play the role of the
preceding mapping \(V\). Thus the mapping \(S\longmapsto S^*\) is a spatial
antiisomorphism of \(\mathcal A\) onto a von Neumann algebra
\(\mathcal A^c\) in \(\mathfrak H^c\). Every antihomomorphism of
\(\mathcal A\) into a von Neumann algebra \(\mathcal B\) is the composite of
this antiisomorphism and a homomorphism of \(\mathcal A^c\) into
\(\mathcal B\). This reduces most problems concerning antihomomorphisms to
problems concerning homomorphisms.

\begin{quote}\small
The structure of isomorphisms and, to some extent, of homomorphisms will be
studied in detail on various occasions, and already in
\hyperref[sec:I-4]{\S~4}. No von Neumann
algebras are known which are antiisomorphic without at the same time being
isomorphic.
\end{quote}

\noindent\textbf{Bibliography:} \ArticleRef{65}, \ArticleRef{66},
\ArticleRef{67}, \ArticleRef{79}.

\NumberedParagraph{6}{Ideals in von Neumann algebras}\label{no:I-1-6}
Let \(\mathcal A\) be a von Neumann algebra. If \(\mathfrak m\) is a left
(respectively, right) ideal of \(\mathcal A\), then \(\mathfrak m^*\) (the set
of the \(S^*\) as \(S\) ranges over \(\mathfrak m\)) is a right
(respectively, left) ideal. Thus a self-adjoint left (respectively, right) ideal
% Source PDF page 015 (printed p. 10)
\phantomsection\label{srcpage:10}
of \(\mathcal A\) is two-sided. Conversely, a two-sided ideal \(\mathfrak m\)
of \(\mathcal A\) is self-adjoint; for let \(S\in\mathfrak m\), and let
\(S=W\lvert S\rvert\) be its polar decomposition. We have
\[
  W\in\mathcal A,
  \qquad
  S^*=\lvert S\rvert W^*=W^*SW^*\in\mathfrak m.
\]

Let \(\mathfrak m\) be a vector subspace of \(\mathcal A\). In order that
\(\mathfrak m\) be a left (respectively, right) ideal of \(\mathcal A\), it
is necessary and sufficient that, for every \(S\in\mathfrak m\) and every
unitary operator \(U\in\mathcal A\), one have \(US\in\mathfrak m\)
(respectively, \(SU\in\mathfrak m\)). Necessity is clear, and sufficiency
follows from \cref{prop:I-1-3}.

Let \(\mathfrak m\) be a left ideal of \(\mathcal A\), and let
\(S\in\mathcal A\). In order that \(S\in\mathfrak m\), it is necessary and
sufficient that \(\lvert S\rvert\in\mathfrak m\). Indeed, let
\(S=W\lvert S\rvert\) be the polar decomposition of \(S\). We have
\(\lvert S\rvert=W^*S\), and \(W\in\mathcal A\).

\begin{proposition}\label{prop:I-1-9}
Let \(\mathfrak m\) be a two-sided ideal of \(\mathcal A\). Every element of
\(\mathfrak m\) is a linear combination of elements of \(\mathfrak m^+\).
\end{proposition}

\begin{proof}
Let \(S\in\mathfrak m\). We have \(S=S_1+iS_2\), with
\(S_1\in\mathfrak m\), \(S_2\in\mathfrak m\), and \(S_1,S_2\) Hermitian. It
is therefore enough to consider the case in which \(S\) is Hermitian. There
exist spectral projections \(E,F\) of \(S\) such that
\[
  E+F=I_{\mathfrak H},
  \qquad
  ES\geq0,
  \qquad
  FS\leq0.
\]
Then \(E\in\mathcal A\), \(F\in\mathcal A\), and therefore
\[
  ES\in\mathfrak m^+,
  \qquad
  -FS\in\mathfrak m^+,
  \qquad
  S=ES+FS.
\]
\end{proof}

\begin{lemma}\label{lem:I-1-1}
Let \(\mathfrak m_0\) be a subset of \(\mathcal A^+\) possessing the following
properties:
\begin{enumerate}[label=(\roman*)]
\item If \(S\in\mathfrak m_0\), and if \(U\) is a unitary operator of
\(\mathcal A\), then \(U^{-1}SU\in\mathfrak m_0\).
\item If \(S\in\mathfrak m_0\), and if \(T\) is an operator of
\(\mathcal A^+\) majorized by \(S\), then \(T\in\mathfrak m_0\).
\item If \(S\in\mathfrak m_0\) and \(T\in\mathfrak m_0\), then
\(S+T\in\mathfrak m_0\).
\end{enumerate}
Then the set \(\mathfrak n\) of the \(S\in\mathcal A\) such that
\(SS^*\in\mathfrak m_0\) is a two-sided ideal of \(\mathcal A\), and
\[
  \mathfrak m_0=(\mathfrak n^2)^+.
\]
\end{lemma}

(Here and in what follows, the symbol \(\mathfrak n^2\) denotes the square
two-sided ideal of the ideal \(\mathfrak n\), in the usual algebraic sense.)

\begin{proof}
Let \(S\in\mathfrak n\), and let \(U\) be a unitary operator of
\(\mathcal A\). We have
\begin{align*}
  (SU)(SU)^*&=SUU^*S^*=SS^*\in\mathfrak m_0,\\
  (US)(US)^*&=U(SS^*)U^{-1}\in\mathfrak m_0,
\end{align*}
so \(SU\in\mathfrak n\) and \(US\in\mathfrak n\). On the other hand, if
\(S\in\mathfrak n\) and \(T\in\mathfrak n\), then
\(S+T\in\mathfrak n\); indeed,
\[
  (S+T)(S+T)^*\leq2SS^*+2TT^*\in\mathfrak m_0,
\]
and hence
\[
  (S+T)(S+T)^*\in\mathfrak m_0.
\]
Thus \(\mathfrak n\) is a two-sided ideal of \(\mathcal A\).
% Source PDF page 016 (printed p. 11)
\phantomsection\label{srcpage:11}
If \(S\in\mathfrak m_0\), then \(S^{1/2}\in\mathfrak n\), and hence
\(S\in\mathfrak n^2\); thus
\(\mathfrak m_0\subset(\mathfrak n^2)^+\). Conversely, let
\(T\in\mathfrak n^2\). The operator \(T\) is a sum of operators \(AB^*\),
where \(A\in\mathfrak n\) and \(B\in\mathfrak n\). If \(T\) is Hermitian,
the identity
\begin{align}
  4AB^*
    &=(A+B)(A+B)^*-(A-B)(A-B)^* \notag\\
    &\quad+i(A+iB)(A+iB)^*-i(A-iB)(A-iB)^*
      \tag{1}\label{eq:I-1-1}
\end{align}
proves that \(T\) is majorized by an operator of the form
\(\sum_{i=1}^n C_iC_i^*\), where the \(C_i\) belong to \(\mathfrak n\),
that is, by an operator of \(\mathfrak m_0\). If \(T\geq0\), it follows
that \(T\in\mathfrak m_0\); thus
\((\mathfrak n^2)^+\subset\mathfrak m_0\).
\end{proof}

\begin{lemma}\label{lem:I-1-2}
Let \(S\) and \(T\) be elements of \(\mathcal A^+\) such that \(S\leq T\).
There exists a unique operator \(A\in\mathcal L(\mathfrak H)\) such that:
\begin{enumerate}[label=\arabic*\textsuperscript{o}.]
\item \(S^{1/2}=AT^{1/2}\);
\item the support of \(A\) is majorized by that of \(T\).
\end{enumerate}
Moreover, \(A\in\mathcal A\).
\end{lemma}

\begin{proof}
For every \(x\in\mathfrak H\),
\[
  \lVert S^{1/2}x\rVert^2
    =(Sx\mid x)
    \leq(Tx\mid x)
    =\lVert T^{1/2}x\rVert^2.
\]
In particular, \(T^{1/2}x=0\) implies \(S^{1/2}x=0\). The mapping
\[
  T^{1/2}x\longmapsto S^{1/2}x
\]
is a continuous linear mapping \(C\) from \(T^{1/2}(\mathfrak H)\) into
\(\mathfrak H\). Let \(B\) be the unique continuous linear mapping from
\[
  \overline{T^{1/2}(\mathfrak H)}=\overline{T(\mathfrak H)}
\]
into \(\mathfrak H\) that extends \(C\). (Here and in what follows,
\(\overline{\mathfrak M}\) denotes the closure of \(\mathfrak M\).) We have
\(S^{1/2}=BT^{1/2}\). The existence and uniqueness of \(A\) follow at
once. Uniqueness also implies that, if \(U'\) is a unitary operator of
\(\mathcal A'\), then
\[
  U'AU'^{-1}=A.
\]
Consequently \(A\in\mathcal A\), by the
\hyperref[cor:I-1-3]{corollary to Proposition~3}.
\end{proof}

\begin{proposition}\label{prop:I-1-10}
Let \(\mathfrak m_0\) be a subset of \(\mathcal A^+\). In order that there
exist a two-sided ideal \(\mathfrak m\) of \(\mathcal A\) such that
\(\mathfrak m_0=\mathfrak m^+\), it is necessary and sufficient that
\(\mathfrak m_0\) possess properties \textup{(i)}, \textup{(ii)}, and
\textup{(iii)} of \cref{lem:I-1-1}.
\end{proposition}

\begin{proof}
The conditions \textup{(i)}, \textup{(ii)}, and \textup{(iii)} are
sufficient by \cref{lem:I-1-1}. Conditions \textup{(i)} and \textup{(iii)}
are plainly necessary. Finally, suppose that \(T\in\mathfrak m^+\),
\(S\in\mathcal A^+\), and \(S\leq T\). By \cref{lem:I-1-2}, there exists
\(A\in\mathcal A\) such that
\[
  S=S^{1/2}(S^{1/2})^*
    =AT^{1/2}T^{1/2}A^*
    =ATA^*,
\]
and therefore \(S\in\mathfrak m\).
\end{proof}

\begin{proposition}\label{prop:I-1-11}
Let \(\mathfrak m\) be a two-sided ideal of \(\mathcal A\).
\begin{enumerate}[label=(\roman*)]
\item The set of the \(S\in\mathcal A\) such that
\(SS^*\in\mathfrak m\) is a two-sided ideal \(\mathfrak n\) of
\(\mathcal A\) such that \(\mathfrak n^2=\mathfrak m\).
% Source PDF page 017 (printed p. 12)
\phantomsection\label{srcpage:12}
\item If \(\Theta\) is a linear mapping from \(\mathfrak m\) into a
complex vector space such that
\[
  \Theta(ST)=\Theta(TS)
  \qquad
  (S\in\mathfrak m,\ T\in\mathcal A),
\]
then
\[
  \Theta(ST)=\Theta(TS)
  \qquad
  (S\in\mathfrak n,\ T\in\mathfrak n).
\]
\end{enumerate}
\end{proposition}

\begin{proof}
Assertion \textup{(i)} follows at once from \cref{lem:I-1-1} and
\cref{prop:I-1-10}. Let us prove assertion \textup{(ii)}. By
\eqref{eq:I-1-1} and the identity
\begin{align}
  4B^*A
    &=(A+B)^*(A+B)-(A-B)^*(A-B) \notag\\
    &\quad+i(A+iB)^*(A+iB)-i(A-iB)^*(A-iB),
      \tag{1'}\label{eq:I-1-1-prime}
\end{align}
it is enough to prove that
\(\Theta(S^*S)=\Theta(SS^*)\) for \(S\in\mathfrak n\). Now, if
\(S=W\lvert S\rvert\) is the polar decomposition of \(S\), then
\[
  S^*S=\lvert S\rvert^2=W^*W\lvert S\rvert^2,
  \qquad
  SS^*=W\lvert S\rvert^2W^*,
  \qquad
  \lvert S\rvert^2\in\mathfrak m,
\]
and therefore
\[
  \Theta(S^*S)
    =\Theta(W^*W\lvert S\rvert^2)
    =\Theta(W\lvert S\rvert^2W^*)
    =\Theta(SS^*).
\]
\end{proof}

The ideal \(\mathfrak n\) of \cref{prop:I-1-11} will be denoted by
\(\mathfrak m^{1/2}\).

One can define \(\mathfrak m^\alpha\) for an arbitrary \(\alpha\geq0\)
\ArticleRef{14}.

\noindent\textbf{Bibliography:} \ArticleRef{4}, \ArticleRef{12},
\ArticleRef{14}, \ArticleRef{25}, \ArticleRef{97}.

\NumberedParagraph{7}{Maximal abelian von Neumann subalgebras}
\label{no:I-1-7}
Let \(\mathcal A\) be a von Neumann algebra. Order by inclusion the abelian
von Neumann subalgebras of \(\mathcal A\). We then have the notion of a
maximal abelian von Neumann subalgebra of \(\mathcal A\). The following
proposition proves the existence of such subalgebras.

\begin{proposition}\label{prop:I-1-12}
Every abelian von Neumann subalgebra of \(\mathcal A\) is contained in a
maximal abelian von Neumann subalgebra of \(\mathcal A\).
\end{proposition}

\begin{proof}
It is enough to prove that the set of abelian von Neumann subalgebras of
\(\mathcal A\) is inductive. Now, if \((\mathcal A_i)_{i\in I}\) is a
totally ordered family of such subalgebras, the von Neumann algebra
\(\mathcal B\) generated by the \(\mathcal A_i\) is abelian, by what was
seen at the end of \hyperref[no:I-1-1]{no.~1}, and is contained in
\(\mathcal A\). It is clear that \(\mathcal B\) is an upper bound for the
\(\mathcal A_i\) in the set of abelian von Neumann subalgebras of
\(\mathcal A\).

\end{proof}

\begin{proposition}\label{prop:I-1-13}
In order that a von Neumann subalgebra \(\mathcal B\) of \(\mathcal A\) be
maximal abelian in \(\mathcal A\), it is necessary and sufficient that it
possess either of the following equivalent properties:
\begin{enumerate}[label=(\roman*)]
\item \(\mathcal B'\cap\mathcal A=\mathcal B\);
\item the von Neumann algebra generated by \(\mathcal B\) and
\(\mathcal A'\) is \(\mathcal B'\).
\end{enumerate}
\end{proposition}

% Source PDF page 018 (printed p. 13)
\phantomsection\label{srcpage:13}
\begin{proof}
The equivalence of conditions \textup{(i)} and \textup{(ii)} follows from
\cref{prop:I-1-1}. Now let \(\mathcal B\) be a von Neumann subalgebra of
\(\mathcal A\). To say that \(\mathcal B\) is abelian is to say that
\(\mathcal B\subset\mathcal B'\), and hence that
\(\mathcal B\subset\mathcal B'\cap\mathcal A\). Suppose henceforth that
this condition is satisfied. To say that \(\mathcal B\) is not maximal
abelian in \(\mathcal A\) is to say that there exists
\(S\in\mathcal A\), \(S\notin\mathcal B\), commuting with \(S^*\) and
with \(\mathcal B\), and hence that
\(\mathcal B\neq\mathcal B'\cap\mathcal A\). Thus, to say that
\(\mathcal B\) is maximal abelian in \(\mathcal A\) is to say that
\(\mathcal B=\mathcal B'\cap\mathcal A\).
\end{proof}

In particular, a maximal abelian von Neumann subalgebra of \(\mathcal A\)
contains the centre of \(\mathcal A\) (which was evident \emph{a priori}).

The notion of a maximal abelian von Neumann subalgebra will be useful in
\hyperref[sec:I-9]{\S~9} and in \hyperref[sec:II-3]{Chapter~II, \S~3}.

\noindent\textbf{Bibliography:} \ArticleRef{28}, \ArticleRef{57},
\ArticleRef{78}, \ArticleRef{100}.

\par\addvspace{0.75\baselineskip}
\noindent\textit{Exercises.}\ ---\enspace
\begin{enumerate}[label=\textbf{\arabic*.},ref=\arabic*,leftmargin=*]
\item\label{ex:I-1-1}
Let \(\mathcal A\) be a von Neumann algebra on \(\mathfrak H\), let
\(\mathcal Z\) be its centre, and let \(E\) be a projection of
\(\mathcal A\) such that the relations \(T\in\mathcal A\), \(ET=0\)
imply \(TE=0\). Prove that \(E\in\mathcal Z\). [Let
\(S\in\mathcal A\). On taking \(T=SE-ESE\), prove that \(SE=ESE\).
Conclude that \(E(\mathfrak H)\) is invariant under \(\mathcal A\).]
\ArticleRef{53}.

\item\label{ex:I-1-2}
Let \(\mathcal A\) be a von Neumann algebra on \(\mathfrak H\), and let
\(\mathcal A_1\) be its unit ball.
\begin{enumerate}[label=\textit{\alph*.},ref=\theenumi\alph*,leftmargin=*]
\item\label{ex:I-1-2-a}
Prove that \(I_{\mathfrak H}\) is an extreme point of
\(\mathcal A_1\). [If
\(I_{\mathfrak H}=\tfrac12(A+B)\), with
\(A\in\mathcal A_1\) and \(B\in\mathcal A_1\), put
\[
  C=\tfrac12(A+A^*),
  \qquad
  D=\tfrac12(B+B^*),
\]
so that
\[
  I_{\mathfrak H}=\tfrac12(C+D),
  \qquad
  C\in\mathcal A_1,
  \qquad
  D\in\mathcal A_1.
\]
Conclude that \(C=D=I_{\mathfrak H}\). Then
\(A=I_{\mathfrak H}+iA_1\), with \(A_1\) Hermitian, and
\(A\in\mathcal A_1\), imply \(A_1=0\). Similarly,
\(B=I_{\mathfrak H}\).]

\item\label{ex:I-1-2-b}
An element \(U\in\mathcal A\) such that
\(U^*U=I_{\mathfrak H}\) is an extreme point of \(\mathcal A_1\). [If
\(U=\tfrac12(A+B)\), with \(A\in\mathcal A_1\) and
\(B\in\mathcal A_1\), deduce from part~\textit{a} that
\(U^*A=U^*B=I_{\mathfrak H}\). Let \(F=UU^*\). Prove that \(A\) and
\(B\) map \(\mathfrak H\) isometrically onto \(F(\mathfrak H)\), and then
that \(A=B=U\).] Similarly, an element \(V\in\mathcal A\) such that
\(VV^*=I_{\mathfrak H}\) is an extreme point of \(\mathcal A_1\). (Use
the isometry \(T\mapsto T^*\) of \(\mathcal A\) onto \(\mathcal A\).)

\item\label{ex:I-1-2-c}
Let \(T\) be an extreme point of \(\mathcal A_1\). Prove that \(T^*T\)
is a projection (cf. \hyperref[lem:I-4-8]{\S~4, Lemma~8}). [Otherwise,
there exists \(C\in\mathcal A^+\), commuting with \(T^*T\), such that
\(T^*TC\neq0\), and
\[
  1
    =\bigl\lVert T^*T(I_{\mathfrak H}+C)^2\bigr\rVert
    =\bigl\lVert T^*T(I_{\mathfrak H}-C)^2\bigr\rVert.
\]
Then \(T(I_{\mathfrak H}+C)\in\mathcal A_1\) and
\(T(I_{\mathfrak H}-C)\in\mathcal A_1\), whence
\[
  T=\tfrac12\bigl(T(I_{\mathfrak H}+C)
                   +T(I_{\mathfrak H}-C)\bigr),
\]
and hence
\[
  T=T(I_{\mathfrak H}+C)=T(I_{\mathfrak H}-C),
\]
and hence \(TC=0\).] Let \(E=T^*T\) and \(F=TT^*\). Prove that
\[
  (I_{\mathfrak H}-F)\mathcal A(I_{\mathfrak H}-E)=0.
\]
[If
\[
  A\in(I_{\mathfrak H}-F)\mathcal A(I_{\mathfrak H}-E),
  \qquad
  \lVert A\rVert\leq1,
\]
then \(\lVert T\pm A\rVert\leq1\), and hence
\(T=T+A=T-A\).] If \(\mathcal A\) is a factor, prove that
\(T^*T=I_{\mathfrak H}\) or \(TT^*=I_{\mathfrak H}\). (Use
\cref{cor:I-1-7-1}.)

\item\label{ex:I-1-2-d}
The extreme points of \(\mathcal A^+\cap\mathcal A_1\) are the projections
of \(\mathcal A\). [If \(E=\tfrac12(A+B)\) is a projection of
\(\mathcal A\), with \(A\geq0\) and \(B\geq0\), then
\(A\in E\mathcal A E\) and \(B\in E\mathcal A E\). Apply
part~\textit{a} to \(E\mathcal A E\).]
\ArticleRef{34}, \ArticleRef{174}, \ArticleRef{289}, \ArticleRef{322}.
\end{enumerate}

\item\label{ex:I-1-3}
Let \(\mathfrak H\) be the Hilbert space whose elements are the complex
numerical functions defined on the interval \([0,1]\) and square-integrable
with respect to
% Source PDF page 019 (printed p. 14)
\phantomsection\label{srcpage:14}
Lebesgue measure. For every measurable, essentially bounded complex
numerical function \(g\) on \([0,1]\), let \(T_g\) be the operator
\(f\mapsto fg\) on \(\mathfrak H\). Let \(\mathcal A\) be the set of the
\(T_g\) as \(g\) ranges over the polynomials with complex coefficients.
Let \(h\) be the characteristic function of \([\tfrac12,1]\). Prove that
\(h\) is separating for \(\mathcal A\), but not totalizing for
\(\mathcal A'\). (Prove that \(T_h\) is a strong limit of operators of
\(\mathcal A\), and hence that
\[
  T'h=T'T_hh=T_hT'h
  \qquad\text{for every \(T'\in\mathcal A'\).}
\])

\item\label{ex:I-1-4}
Let \(\mathcal A\) be an involutive algebra of operators on \(\mathfrak H\).
If \(T\) is an invertible element of \(\mathcal L(\mathfrak H)\) such that
the mapping \(A\mapsto TAT^{-1}\) is an automorphism of \(\mathcal A\)
(for its involutive-algebra structure), then there exists a unitary operator
\(U\) on \(\mathfrak H\) such that
\[
  UAU^{-1}=TAT^{-1}
  \qquad\text{for every \(A\in\mathcal A\).}
\]
[The equality \((TAT^{-1})^*=TA^*T^{-1}\) proves that
\(T^*T\in\mathcal A'\). Let \(T=U\lvert T\rvert\) be the polar
decomposition of \(T\). We have \(\lvert T\rvert\in\mathcal A'\), and
hence \(TAT^{-1}=UAU^{-1}\) for \(A\in\mathcal A\).]

\item\label{ex:I-1-5}
Let \(\mathfrak A\) be a complete normed algebra over the field
\(\mathbb C\), and let \(\mathcal B\) be an involutive algebra of
operators on the complex Hilbert space \(\mathfrak H\).
\begin{enumerate}[label=\textit{\alph*.},ref=\theenumi\alph*,leftmargin=*]
\item\label{ex:I-1-5-a}
Let \(\Phi\) be an isomorphism of the algebra \(\mathfrak A\) onto the
algebra \(\mathcal B\). For \(x\in\mathcal B\), let \(\lVert x\rVert\)
be the norm of \(x\) in \(\mathcal B\), and let \(\lVert x\rVert_1\) be
the norm obtained by transporting the norm of \(\mathfrak A\) by
\(\Phi\). Prove that
\[
  \lVert x\rVert^2\leq\lVert x^*\rVert_1\lVert x\rVert_1.
\]
(If \(\lambda>\lVert x^*x\rVert_1\), then
\(I_{\mathfrak H}-\lambda^{-1}x^*x\) is invertible, and hence
\(\lambda>\lVert x^*x\rVert\).)

\item\label{ex:I-1-5-b}
Prove that the involution of \(\mathcal B\) is continuous for the norm
\(\lVert x\rVert_1\). (Using part~\textit{a}, prove that the hypotheses
\[
  \lVert x_n-y\rVert_1\longrightarrow0,
  \qquad
  \lVert x_n^*-z\rVert_1\longrightarrow0
\]
imply \(z=y^*\); then apply the closed graph theorem.)

\item\label{ex:I-1-5-c}
Prove that \(\Phi\) is continuous. (Use parts~\textit{a} and
\textit{b}.)

\item\label{ex:I-1-5-d}
Let \(\Psi\) be a homomorphism of the algebra \(\mathfrak A\) onto the
algebra \(\mathcal B\). Prove that \(\Psi\) is continuous. [The condition
\(\Psi(x)=0\) is equivalent to \(\Psi(x)\in\mathfrak M\) for every regular
maximal right ideal \(\mathfrak M\) of the semisimple algebra
\(\mathcal B\); but \(\Psi^{-1}(\mathfrak M)\) is a regular maximal right ideal
of \(\mathfrak A\), and hence is closed. Conclude that the kernel
\(N\) of \(\Psi\) is closed. Then \(\mathfrak A/N\) is a complete normed
algebra; apply part~\textit{c} to the isomorphism obtained from \(\Psi\) by
passing to the quotient.]

\item\label{ex:I-1-5-e}
Deduce from part~\textit{c} that, if \(\mathcal B\) is an involutive algebra
of operators on \(\mathfrak H\) that is norm-closed, every automorphism of
the algebra \(\mathcal B\) is continuous.
\end{enumerate}

\noindent\textit{Bibliography:} C.~E.~Rickart,
\emph{The uniqueness of norm problem in Banach algebras}
(\emph{Ann. Math.}, vol.~51, 1950, pp.~615--628).

\item\label{ex:I-1-6}
Let \(\mathcal A\) be a von Neumann algebra. For every two-sided ideal
\(\mathfrak m\) of \(\mathcal A\), let \(\mathfrak m^\infty\) be the
two-sided ideal generated by the projections of \(\mathfrak m\), and let
\(\mathfrak m^0\) be the norm closure of \(\mathfrak m\), which is a
two-sided ideal. The ideal \(\mathfrak m\) is said to be \emph{restricted}
if \(\mathfrak m=\mathfrak m^\infty\).
\begin{enumerate}[label=\textit{\alph*.},ref=\theenumi\alph*,leftmargin=*]
\item\label{ex:I-1-6-a}
Let \(T\in\mathfrak m^+\). For \(\lambda>0\), let \(E_\lambda\) be the
largest spectral projection of \(T\) such that
\(TE_\lambda\leq\lambda E_\lambda\). Prove that
\(I_{\mathfrak H}-E_\lambda\in\mathfrak m\). (We have
\(I_{\mathfrak H}-E_\lambda=TT_\lambda\), with
\(T_\lambda\in\mathcal A\).)

\item\label{ex:I-1-6-b}
Deduce from part~\textit{a} that \(\mathfrak m^0\) is the norm closure of
\(\mathfrak m^\infty\).

\item\label{ex:I-1-6-c}
Prove that every projection \(E\) of \(\mathfrak m^0\) belongs to
\(\mathfrak m\). [Let \((T_n)\) be a sequence of elements of
\(\mathfrak m\) such that \(\lVert T_n-E\rVert\to0\). If
\[
  S_n=ET_nE,
  \qquad
  R_n=S_n^*S_n,
\]
then
\[
  \lVert S_n-E\rVert\longrightarrow0,
  \qquad
  \lVert R_n-E\rVert\longrightarrow0.
\]
Consequently
\[
  \tfrac12E\leq R_n\leq\tfrac32E
\]
for all sufficiently large \(n\), and then \(E=R_nR\) with
\(R\in\mathcal A\).]

\item\label{ex:I-1-6-d}
Deduce from parts~\textit{b} and \textit{c} that the mappings
\(\mathfrak m\mapsto\mathfrak m^\infty\) and
\(\mathfrak m\mapsto\mathfrak m^0\) define mutually inverse one-to-one
correspondences between the set of
% Source PDF page 020 (printed p. 15)
\phantomsection\label{srcpage:15}
restricted two-sided ideals of \(\mathcal A\) and the set of norm-closed
two-sided ideals of \(\mathcal A\). \ArticleRef{12}, \ArticleRef{124}.
\end{enumerate}

\item\label{ex:I-1-7}
A von Neumann algebra with no nontrivial two-sided ideal is a factor of
countable type. (If the von Neumann algebra \(\mathcal A\) is not of
countable type, the set \(\mathfrak m\) of the \(T\in\mathcal A\) whose
support is of the form \(E_{\mathfrak M}^{\mathcal A'}\), with
\(\mathfrak M\) countable, is a nontrivial two-sided ideal of
\(\mathcal A\).) \ArticleRef{4}, \ArticleRef{94}.

\item\label{ex:I-1-8}
Let \(\mathfrak H\) be an infinite-dimensional complex Hilbert space with
a countable orthonormal basis, and let \(\mathfrak m_1\) (respectively,
\(\mathfrak m_2\)) be the set of compact operators (respectively,
finite-rank operators) on \(\mathfrak H\).
\begin{enumerate}[label=\textit{\alph*.},ref=\theenumi\alph*,leftmargin=*]
\item\label{ex:I-1-8-a}
Prove that \(\mathfrak m_1\) and \(\mathfrak m_2\) are two-sided ideals of
\(\mathcal L(\mathfrak H)\), and that, if \(\mathfrak m\) is an arbitrary
two-sided ideal of \(\mathcal L(\mathfrak H)\), distinct from
\(\{0\}\) and from \(\mathcal L(\mathfrak H)\), then
\[
  \mathfrak m_2\subset\mathfrak m\subset\mathfrak m_1.
\]

\item\label{ex:I-1-8-b}
Let \(\mathfrak X\) be the set of infinite sequences of nonnegative numbers
tending to zero. A subset \(\mathfrak S\) of \(\mathfrak X\) is called an
ideal of \(\mathfrak X\) if it possesses the following properties:
\begin{enumerate}[label=(\roman*),leftmargin=*]
\item if \((\lambda_n)\in\mathfrak S\) and \(\pi\) is a permutation of
\(\{1,2,\ldots\}\), then \((\lambda_{\pi(n)})\in\mathfrak S\);
\item if \((\lambda_n)\in\mathfrak S\) and
\((\mu_n)\in\mathfrak S\), then
\((\lambda_n+\mu_n)\in\mathfrak S\);
\item if \((\lambda_n)\in\mathfrak S\),
\((\mu_n)\in\mathfrak X\), and \(\mu_n\leq\lambda_n\) for every \(n\),
then \((\mu_n)\in\mathfrak S\).
\end{enumerate}
Let \(\mathfrak m\) be a two-sided ideal of
\(\mathcal L(\mathfrak H)\) contained in \(\mathfrak m_1\). For every
\(T\in\mathfrak m^+\) and every orthonormal basis \((e_n)\) of
\(\mathfrak H\) consisting of eigenvectors of \(T\), the sequence
\((\lambda_n)\) of the corresponding eigenvalues of \(T\) is an element
of \(\mathfrak X\). Let \(\mathfrak S(T)\) be the subset of
\(\mathfrak X\) thus associated with \(T\) as \((e_n)\) varies, and put
\[
  \mathfrak S=\bigcup_{T\in\mathfrak m^+}\mathfrak S(T).
\]
Then \(\mathfrak S\) is an ideal of \(\mathfrak X\), and the mapping
\(\mathfrak m\mapsto\mathfrak S\) is a one-to-one mapping of the set of
two-sided ideals of \(\mathcal L(\mathfrak H)\) contained in
\(\mathfrak m_1\) onto the set of ideals of \(\mathfrak X\). (To prove
that every ideal of \(\mathfrak X\) comes from a two-sided ideal
\(\mathfrak m\subset\mathfrak m_1\), use \cref{prop:I-1-10} and the
maximum-minimum definition of eigenvalues.) \ArticleRef{4}.
\end{enumerate}

\item\label{ex:I-1-9}
\begin{enumerate}[label=\textit{\alph*.},ref=\theenumi\alph*,leftmargin=*]
\item\label{ex:I-1-9-a}
Let \(\mathcal A\) be a von Neumann algebra on \(\mathfrak H\), let
\(\mathfrak M\) be a total subset of \(\mathfrak H\), let \((T_n)\) be a
sequence of elements of \(\mathcal A\), and let
\(T\in\mathcal L(\mathfrak H)\) be such that \(T_nx\) and \(T_n^*x\)
converge weakly to \(Tx\) and \(T^*x\), respectively, for every
\(x\in\mathfrak M\), as \(n\to+\infty\). Prove that \(T\in\mathcal A\).
[Prove that, if \(S\in\mathcal A'\), \(x\in\mathfrak M\), and
\(y\in\mathfrak M\), then
\[
  (TSx\mid y)=(STx\mid y).
\]] \ArticleRef{28}.

\item\label{ex:I-1-9-b}
Let \((e_1,e_2,\ldots)\) be an orthonormal basis of \(\mathfrak H\). Let
\(T_n\in\mathcal L(\mathfrak H)\) be defined by
\begin{align*}
  T_ne_2&=2e_2-3e_3,\\
  T_ne_n&=-ne_2+\tfrac32ne_3,\\
  T_ne_i&=0\qquad(i\neq2,n).
\end{align*}
Let \(T\in\mathcal L(\mathfrak H)\) be defined by
\[
  Te_2=2e_2-3e_3,
  \qquad
  Te_i=0\quad(i\neq2).
\]
Prove that \(T_ne_i\) converges strongly to \(Te_i\), for every \(i\), but
that \(T\) does not belong to the von Neumann algebra \(\mathcal A\)
generated by the \(T_n\). [Let \(S\) be the Hermitian operator in
\(\mathcal L(\mathfrak H)\) defined by
\[
  Se_1=e_1+\tfrac12e_2+\tfrac13e_3+\cdots,
  \qquad
  Se_i=\tfrac1i e_1\quad(i>1).
\]
Prove that \(S\) commutes with the \(T_n\), but that
\[
  (STe_1\mid e_2)\neq(TSe_1\mid e_2).
\]] Cf. \hyperref[sec:I-3]{\S~3}.
\end{enumerate}

\item\label{ex:I-1-10}
Let \(\mathcal A\) be a von Neumann algebra on \(\mathfrak H\), and let
\(T\) be a linear mapping (not necessarily continuous) from a vector
subspace of \(\mathfrak H\) into \(\mathfrak H\). The operator \(T\) is
said to be \emph{affiliated} with \(\mathcal A\), and one writes
\(T\mathrel{\eta}\mathcal A\), if
\[
  U'TU'^{-1}=T
\]
for every unitary operator \(U'\) of \(\mathcal A'\).
\begin{enumerate}[label=\textit{\alph*.},ref=\theenumi\alph*,leftmargin=*]
\item\label{ex:I-1-10-a}
If \(T\in\mathcal L(\mathfrak H)\), the conditions
\(T\mathrel{\eta}\mathcal A\) and \(T\in\mathcal A\) are equivalent.

\item\label{ex:I-1-10-b}
If \(T\mathrel{\eta}\mathcal A\), and if \(T^*\) exists, then
\(T^*\mathrel{\eta}\mathcal A\).

% Source PDF page 023 (printed p. 16)
\phantomsection\label{srcpage:16}
\item\label{ex:I-1-10-c}
Suppose that \(T\) is closed and that its domain is dense in
\(\mathfrak H\). Let \(T=U\lvert T\rvert\) be its polar decomposition. In
order that \(T\mathrel{\eta}\mathcal A\), it is necessary and sufficient
that \(U\in\mathcal A\) and \(\lvert T\rvert\mathrel{\eta}\mathcal A\).
In order that \(\lvert T\rvert\mathrel{\eta}\mathcal A\), it is necessary
and sufficient that the spectral projections of \(\lvert T\rvert\) belong
to \(\mathcal A\). \ArticleRef{65}, \ArticleRef{101}.
\end{enumerate}
\end{enumerate}

\section{Elementary Operations on von Neumann Algebras}\label{sec:I-2}

\NumberedParagraph{1}{Induced and reduced von Neumann algebras}
\label{no:I-2-1}
There are several procedures for constructing von Neumann algebras from
given von Neumann algebras. One of these ``elementary operations'' is
already known: given a von Neumann algebra, it consists in taking its
commutant. In this section we shall define other elementary operations.

Let \(\mathfrak H\) be a complex Hilbert space, let \(E=P_{\mathfrak X}\)
be a projection on \(\mathfrak H\), and let
\(T\in\mathcal L(\mathfrak H)\). We denote by \(T_{\mathfrak X}\), or by
\(T_E\), the restriction of \(ET\) to \(\mathfrak X\), which is an element
of \(\mathcal L(\mathfrak X)\). We have
\[
  T_{\mathfrak X}=(TE)_{\mathfrak X}
    =(ET)_{\mathfrak X}=(ETE)_{\mathfrak X}.
\]
If \(\mathfrak M\) is a subset of \(\mathcal L(\mathfrak H)\), we denote
by \(\mathfrak M_{\mathfrak X}\), or by \(\mathfrak M_E\), the set of the
\(T_{\mathfrak X}\) as \(T\) ranges over \(\mathfrak M\). Consider the
following two cases:
\begin{enumerate}[label=\arabic*\textsuperscript{o}.,leftmargin=*]
\item \(\mathfrak M\) is an involutive algebra of operators
\(\mathcal A\), and \(E\in\mathcal A\). Let \(\mathcal B\) be the set of
the \(T\in\mathcal A\) such that \(TE=ET=T\), in other words the set of
the \(T\in\mathcal A\) such that
\[
  T(\mathfrak X)\subset\mathfrak X,
  \qquad
  T(\mathfrak H\ominus\mathfrak X)=0.
\]
The set \(\mathcal B\) is an involutive subalgebra of \(\mathcal A\); we
have
\[
  \mathcal B=E\mathcal B E\subset E\mathcal A E,
  \qquad
  E\mathcal A E\subset\mathcal B,
\]
and hence \(\mathcal B=E\mathcal A E\). The mapping \(T\mapsto T_E\) is
an isomorphism, for the involutive-algebra structure, of \(\mathcal B\)
onto \(\mathcal A_E\).

\item \(\mathfrak M\) is an involutive algebra of operators
\(\mathcal A\), and \(E\in\mathcal A'\). Every operator
\(T\in\mathcal A\) leaves \(\mathfrak X\) and
\(\mathfrak H\ominus\mathfrak X\) invariant. The mapping
\(T\mapsto T_E\) is a homomorphism, for the involutive-algebra structure,
of \(\mathcal A\) onto \(\mathcal A_E\).
\end{enumerate}

\begin{proposition}\label{prop:I-2-1}
Let \(\mathcal A\) be a von Neumann algebra, and let
\(E=P_{\mathfrak X}\) be a projection of \(\mathcal A\).
\begin{enumerate}[label=(\roman*)]
\item \(\mathcal A_E\) and \((\mathcal A')_E\) are von Neumann algebras,
and
\[
  (\mathcal A')_E=(\mathcal A_E)'.
\]
\item If \(\mathfrak M\) is a set stable under adjunction and
multiplication that generates the von Neumann algebra \(\mathcal A\), then
\(\mathfrak M_E\) generates the von Neumann algebra \(\mathcal A_E\).
\item If \(\mathfrak N\) is a set that generates the von Neumann algebra
\(\mathcal A'\), then \(\mathfrak N_E\) generates the von Neumann algebra
\((\mathcal A_E)'\).
\end{enumerate}
\end{proposition}

% Source PDF page 024 (printed p. 17)
\phantomsection\label{srcpage:17}
\begin{proof}
We may suppose that \(\mathfrak N\) is stable under adjunction, replacing
\(\mathfrak N\), if necessary, by
\(\mathfrak N\cup\mathfrak N^*\). We may suppose that
\(I_{\mathfrak H}\in\mathfrak N\), replacing \(\mathfrak N\), if
necessary, by \(\mathfrak N\cup\{I_{\mathfrak H}\}\).

It is evident that the involutive algebras \(\mathcal A_E\) and
\((\mathcal A')_E\) commute. Conversely, let
\(T\in\mathcal L(\mathfrak X)\) be an operator commuting with the operators
\(T'_E\), where \(T'\in\mathfrak N\). Let
\(S=T\circ E\in\mathcal L(\mathfrak H)\). It is immediate that \(S\)
commutes with \(\mathfrak N\); hence
\[
  S\in\mathfrak N'=\mathcal A''=\mathcal A,
\]
and consequently \(T=S_E\in\mathcal A_E\). Taking
\(\mathfrak N=\mathcal A'\), we see that
\[
  \mathcal A_E=((\mathcal A')_E)'
\]
is a von Neumann algebra. Returning to the case of arbitrary
\(\mathfrak N\), we also see that \(\mathfrak N_E\) generates
\((\mathcal A_E)'\).

Now let \(T'\in\mathcal L(\mathfrak X)\) be an operator commuting with
\(\mathfrak M_E\). We shall prove that \(T'\in(\mathcal A')_E\). This will
establish that \((\mathcal A')_E=(\mathcal A_E)'\), that
\((\mathcal A')_E\) is a von Neumann algebra, and that \(\mathfrak M_E\)
generates \(\mathcal A_E\). The proof will then be complete.

It is enough to consider the case in which \(T'\) is unitary
(\cref{prop:I-1-3}). For \(x_1,x_2,\ldots,x_n\in\mathfrak X\) and
\(T_1,T_2,\ldots,T_n\in\mathfrak M\), we have
\begin{align*}
  \left\lVert\sum_{i=1}^n T_iT'x_i\right\rVert^2
    &=\sum_{i,j=1}^n
      (T_iET'x_i\mid T_jET'x_j)\\
    &=\sum_{i,j=1}^n
      (ET_j^*T_iET'x_i\mid T'x_j)\\
    &=\sum_{i,j=1}^n
      (T'ET_j^*T_iEx_i\mid T'x_j)\\
    &=\sum_{i,j=1}^n
      (ET_j^*T_iEx_i\mid x_j)\\
    &=\left\lVert\sum_{i=1}^n T_ix_i\right\rVert^2.
\end{align*}
Let \(\mathfrak J\) be the closed vector subspace of \(\mathfrak H\)
generated by the elements \(\sum_{i=1}^nT_ix_i\). By
\cref{prop:I-1-4},
\[
  P_{\mathfrak J}=E_{\mathfrak X}^{\mathcal A};
\]
by \cref{cor:I-1-7-1}, \(P_{\mathfrak J}\) is therefore the central
support of \(E\). There exists a unique operator
\(S'\in\mathcal L(\mathfrak H)\) such that
\[
  S'\left(\sum_{i=1}^nT_ix_i\right)
    =\sum_{i=1}^nT_iT'x_i,
  \qquad
  S'P_{\mathfrak J}=P_{\mathfrak J}S'=S'.
\]
For every \(T\in\mathfrak M\), we have
\begin{align*}
  S'T\left(\sum_{i=1}^nT_ix_i\right)
    &=S'\left(\sum_{i=1}^nTT_ix_i\right)\\
    &=\sum_{i=1}^nTT_iT'x_i\\
    &=TS'\left(\sum_{i=1}^nT_ix_i\right),
\end{align*}
and \(P_{\mathfrak J}T=TP_{\mathfrak J}\); hence \(S'T=TS'\). Thus
\(S'\in\mathfrak M'=\mathcal A'\). Finally, for \(x\in\mathfrak X\),
we have \(S'x=T'x\); hence \(T'=S'_E\in(\mathcal A')_E\).
\end{proof}

Instead of writing \((\mathcal A')_E\) or \((\mathcal A_E)'\), we shall
henceforth write \(\mathcal A'_E\). The algebra \(\mathcal A'_E\) is
called the \emph{von Neumann algebra induced by \(\mathcal A'\) in
\(\mathfrak X\)}. The homomorphism \(T'\mapsto T'_E\) of \(\mathcal A'\)
onto \(\mathcal A'_E\) is called the \emph{induction of \(\mathcal A'\)
onto \(\mathcal A'_E\)}. Every von Neumann algebra such as
\(\mathcal A_E\) is called a \emph{reduced von Neumann algebra} of
\(\mathcal A\).

% Source PDF page 025 (printed p. 18)
\phantomsection\label{srcpage:18}
There is an evident transitivity property: if
\(P_{\mathfrak X_1}\) and \(P_{\mathfrak X_2}\) are projections of
\(\mathcal A\) such that \(\mathfrak X_1\supset\mathfrak X_2\), then
\[
  \mathcal A_{\mathfrak X_2}
    =(\mathcal A_{\mathfrak X_1})_{\mathfrak X_2},
  \qquad
  \mathcal A'_{\mathfrak X_2}
    =(\mathcal A'_{\mathfrak X_1})_{\mathfrak X_2}.
\]

Let \(\mathcal A\) and \(\mathcal B\) be von Neumann algebras, let
\(\Phi\) be a homomorphism from \(\mathcal A\) into \(\mathcal B\), let
\(E\) be a projection of \(\mathcal A\), and put \(F=\Phi(E)\). The
restriction of \(\Phi\) to \(E\mathcal A E\) is a homomorphism from the
involutive algebra \(E\mathcal A E\) into the involutive algebra
\(F\mathcal B F\); in the evident manner it defines a homomorphism
\(\Phi_E\) from \(\mathcal A_E\) into \(\mathcal B_F\), called the
\emph{reduced homomorphism} of \(\Phi\). If \(\Phi\) is an isomorphism of
\(\mathcal A\) onto \(\mathcal B\), then \(\Phi_E\) is an isomorphism of
\(\mathcal A_E\) onto \(\mathcal B_F\). There are analogous results for
antihomomorphisms and anti-isomorphisms.

\begin{proposition}\label{prop:I-2-2}
Let \(\mathcal A\) be a von Neumann algebra, let \(E\) be a projection of
\(\mathcal A\), and let \(F\) be its central support. In order that the
induction of \(\mathcal A'\) onto \(\mathcal A'_E\) be an isomorphism, it
is necessary and sufficient that \(F=I_{\mathfrak H}\).
\end{proposition}

\begin{proof}
The condition is necessary, since
\((I_{\mathfrak H}-F)_E=0\). Conversely, suppose that
\(F=I_{\mathfrak H}\). Put \(\mathfrak X=E(\mathfrak H)\). The subspaces
\(T(\mathfrak X)\), \(T\in\mathcal A\), generate \(\mathfrak H\)
(\cref{cor:I-1-7-1}). If \(T'\in\mathcal A'\) is such that \(T'_E=0\),
then \(T'(\mathfrak X)=0\), and hence
\[
  T'T(\mathfrak X)=TT'(\mathfrak X)=0
  \qquad\text{for every \(T\in\mathcal A\)};
\]
therefore \(T'=0\).
\end{proof}

\begin{corollary*}\phantomsection\label{cor:I-2-2}
Let \(\mathcal A\) be a von Neumann algebra, let \(\mathcal Z\) be its
centre, and let \(E\) be a projection of \(\mathcal A\). The centre of
\(\mathcal A_E\) is \(\mathcal Z_E\). In particular, if \(\mathcal A\)
is a factor, then \(\mathcal A_E\) and \(\mathcal A'_E\) are factors.
\end{corollary*}

\begin{proof}
It is clear that \(\mathcal Z_E\) is contained in the centre of
\(\mathcal A_E\). Conversely, let
\(T\in\mathcal A_E\cap\mathcal A'_E\). We have \(T=T'_E\), with
\(T'\in\mathcal A'\). Let \(F\) be the central support of \(E\). Replacing
\(T'\) by \(FT'\), we may suppose that \(T'=FT'\). Applying
\cref{prop:I-2-2} to \(\mathcal A_F\), we then see that \(T'_F\) belongs
to the centre of \(\mathcal A'_F\). Consequently \(T'\in\mathcal Z\), and
hence \(T\in\mathcal Z_E\).
\end{proof}

\begin{proposition}\label{prop:I-2-3}
Let \(\mathcal A\) be a von Neumann algebra of countable type, and let
\(\mathcal Z\) be its centre. There exists a projection
\(G\in\mathcal Z\) such that \(\mathcal A_G\) possesses a totalizing
element and \(\mathcal A_{I_{\mathfrak H}-G}\) possesses a separating
element.
\end{proposition}

\begin{proof}
Let \((x_i)_{i\in I}\) be a maximal family of nonzero elements of
\(\mathfrak H\) possessing the following properties:
\begin{enumerate}[label=\arabic*\textsuperscript{o}.]
\item the \(E_i=E_{x_i}^{\mathcal A'}\) are pairwise orthogonal;
\item the \(E_i'=E_{x_i}^{\mathcal A}\) are pairwise orthogonal.
\end{enumerate}
Put
\[
  E=\sum_{i\in I}E_i\in\mathcal A,
  \qquad
  E'=\sum_{i\in I}E_i'\in\mathcal A',
  \qquad
  F=I_{\mathfrak H}-E,
  \qquad
  F'=I_{\mathfrak H}-E'.
\]
Let \(F_1\) (respectively, \(F_1'\)) be the central support of \(F\)
(respectively, of \(F'\)). If \(F_1F_1'\neq0\), then
\((F')_{F_1}\neq0\), and hence \((F')_F\neq0\), by
\cref{prop:I-2-2}; in other words, \(F'F\neq0\). Let \(y\) be a
% Source PDF page 026 (printed p. 19)
\phantomsection\label{srcpage:19}
nonzero element of \(F(\mathfrak H)\cap F'(\mathfrak H)\); then
\(E_y^{\mathcal A'}\) is orthogonal to the \(E_i\), and
\(E_y^{\mathcal A}\) is orthogonal to the \(E_i'\), contradicting the
maximality of the family \((x_i)_{i\in I}\). Thus \(F_1F_1'=0\). On putting
\(G=F_1\), we have
\[
  E=I_{\mathfrak H}-F\geq I_{\mathfrak H}-G,
  \qquad
  E'=I_{\mathfrak H}-F'\geq I_{\mathfrak H}-F_1'\geq G.
\]
On the other hand, since \(\mathcal A\) is of countable type, \(I\) is
countable; by multiplying the \(x_i\) by suitable scalars, we may suppose
that
\[
  \sum_{i\in I}\lVert x_i\rVert^2<+\infty;
  \qquad
  x=\sum_{i\in I}x_i\in\mathfrak H.
\]
Since \(x_i=E_i x=E_i'x\), we see that \(E_x^{\mathcal A}\) majorizes the
\(E_i'\), and hence \(E'\), and that \(E_x^{\mathcal A'}\) majorizes the
\(E_i\), and hence \(E\). Consequently \(Gx\) is totalizing for
\(\mathcal A_G\), and \((I_{\mathfrak H}-G)x\) is totalizing for
\(\mathcal A'_{I_{\mathfrak H}-G}\).
\end{proof}

\begin{corollary*}
An abelian von Neumann algebra of countable type possesses a separating
element.
\end{corollary*}

\begin{proof}
Retaining the preceding notation, \(I_{\mathfrak H}-G\) is majorized by
\(E_x^{\mathcal A'}\), and \(G\) is majorized by
\(E_x^{\mathcal A}\), and therefore, a fortiori, by
\(E_x^{\mathcal A'}\). Thus \(E_x^{\mathcal A'}=I_{\mathfrak H}\), and
\(x\) is separating for \(\mathcal A\).
\end{proof}

\Cref{prop:I-2-1}(ii) is false if \(\mathfrak M\) is an arbitrary subset
generating the von Neumann algebra \(\mathcal A\) (Exercise~1).
\Cref{prop:I-2-3} is false without the countability hypothesis: if
\(\mathcal A\) is abelian, if \(\mathcal A_G\) possesses a totalizing
element, and if \(\mathcal A_{I_{\mathfrak H}-G}\) possesses a separating
element, then \(\mathcal A\) possesses a separating element and is therefore
of countable type.

\noindent\textbf{Bibliography:} \ArticleRef{12}, \ArticleRef{65},
\ArticleRef{100}, \ArticleRef{117}.

\NumberedParagraph{2}{Product of von Neumann algebras}\label{no:I-2-2}
Let \((\mathfrak H_i)_{i\in I}\) be a family of complex Hilbert spaces, and
let \(\mathfrak H\) be their Hilbert direct sum, namely the set of families
\((x_i)_{i\in I}\), where \(x_i\in\mathfrak H_i\) and
\(\sum_{i\in I}\lVert x_i\rVert^2<+\infty\). For every family
\((T_i)_{i\in I}\), where \(T_i\in\mathcal L(\mathfrak H_i)\) and
\(\sup_i\lVert T_i\rVert<+\infty\), one plainly defines
\(T\in\mathcal L(\mathfrak H)\) by
\[
  T\bigl((x_i)_{i\in I}\bigr)=(T_i x_i)_{i\in I}.
\]
We shall denote this operator by \((T_i)\), which will cause no confusion.
It is clear that
\[
  (T_i+S_i)=(T_i)+(S_i),\qquad
  (T_iS_i)=(T_i)(S_i),\qquad
  (T_i^*)=(T_i)^*.
\]
Let \(E_i\) be the projection \(P_{\mathfrak H_i}\) in \(\mathfrak H\).
For an operator in \(\mathcal L(\mathfrak H)\) to be of the form \((T_i)\),
it is necessary and sufficient that it commute with the \(E_i\).

This being so, let \(\mathcal A_i\), for each \(i\in I\), be a von Neumann
algebra on \(\mathfrak H_i\). Let \(\mathcal A\) (respectively,
\(\mathcal B\)) be the set of operators \((T_i)\) such that
\(T_i\in\mathcal A_i\) (respectively, \(T_i\in\mathcal A_i'\)) for every
\(i\in I\), and \(\sup_i\lVert T_i\rVert<+\infty\). The involutive
operator algebras \(\mathcal A\) and \(\mathcal B\) on \(\mathfrak H\)
commute with one another. We shall prove that
\(\mathcal A'=\mathcal B\) and \(\mathcal B'=\mathcal A\), which will in
particular imply that \(\mathcal A\) and \(\mathcal B\) are
% Source PDF page 027 (printed p. 20)
\phantomsection\label{srcpage:20}
von Neumann algebras. Let \(T\in\mathcal L(\mathfrak H)\) commute with
\(\mathcal B\). Since it is clear that \(E_i\in\mathcal B\), \(T\)
commutes with the \(E_i\), and hence \(T\) is of the form \((T_i)\). For
every operator \((T_i')\), where \(T_i'\in\mathcal A_i'\) for all
\(i\in I\), we must have
\[
  (T_iT_i')=(T_i)(T_i')=(T_i')(T_i)=(T_i'T_i);
\]
thus \(T_iT_i'=T_i'T_i\), and therefore
\(T_i\in\mathcal A_i''=\mathcal A_i\), for every \(i\in I\). Hence
\(T\in\mathcal A\), which proves that \(\mathcal B'=\mathcal A\). By
interchanging the roles of the \(\mathcal A_i\) and the
\(\mathcal A_i'\), we similarly obtain \(\mathcal A'=\mathcal B\).

The von Neumann algebra \(\mathcal A\) is called the \emph{product von
Neumann algebra} of the von Neumann algebras \(\mathcal A_i\), and is
denoted by \(\prod_{i\in I}\mathcal A_i\). We have therefore proved the
formula
\[
  \left(\prod_{i\in I}\mathcal A_i\right)'
    =\prod_{i\in I}\mathcal A_i'.
\]
(For a finite family \(\mathcal A_1,\mathcal A_2,\ldots,\mathcal A_n\),
the notation \(\mathcal A_1\times\mathcal A_2\times\cdots\times
\mathcal A_n\) is also used.) The centre \(\mathcal Z\) of
\(\mathcal A\) is the set of operators \((T_i)\) such that \(T_i\)
belongs, for every \(i\in I\), to the centre \(\mathcal Z_i\) of
\(\mathcal A_i\); in other words,
\[
  \mathcal Z=\prod_{i\in I}\mathcal Z_i.
\]
In particular, the \(E_i\) belong to the centre of \(\mathcal A\).

Conversely, let \(\mathcal A\) be a von Neumann algebra, and let
\((E_i)_{i\in I}\) be a family of pairwise orthogonal projections in the
centre of \(\mathcal A\) such that
\[
  \sum_{i\in I}E_i=I_{\mathfrak H}.
\]
We shall see that these data canonically define a family
\((\mathcal A_i)_{i\in I}\) of von Neumann algebras and a spatial
isomorphism of \(\mathcal A\) onto
\(\prod_{i\in I}\mathcal A_i\). Indeed, put
\(\mathfrak H_i=E_i(\mathfrak H)\). The space \(\mathfrak H\) is the
Hilbert direct sum of the \(\mathfrak H_i\). Put
\(\mathcal A_i=\mathcal A_{E_i}\). We have
\(\mathcal A_i'=\mathcal A'_{E_i}\). For every \(T\in\mathcal A\), put
\(T_i=T_{E_i}\), and form the operator
\((T_i)\in\prod_{i\in I}\mathcal A_i\). The mapping
\[
  T\longmapsto (T_i)
\]
is an isomorphism of \(\mathcal A\) onto
\(\prod_{i\in I}\mathcal A_i\). Indeed, if
\((T_i)\in\prod_{i\in I}\mathcal A_i\), there exists an operator
\(T\in\mathcal L(\mathfrak H)\) which induces \(T_i\) on
\(\mathfrak H_i\), for every \(i\in I\), and \(T\) commutes with every
element of \(\mathcal A'=\prod_{i\in I}\mathcal A_i'\); hence
\(T\in\mathcal A''=\mathcal A\).

The product is associative in the evident sense. On the other hand, for
each \(i\), let \(F_i\) be a projection in \(\mathcal A_i\). Put
\(F=(F_i)\), which is a projection in \(\mathcal A\). We have
\[
  \mathcal A_F=\prod_{i\in I}(\mathcal A_i)_{F_i},
  \qquad
  \mathcal A_F'=\prod_{i\in I}(\mathcal A_i')_{F_i}.
\]

% Source PDF page 028 (printed p. 21)
\phantomsection\label{srcpage:21}
Finally, let \(\Phi_i\), for every \(i\in I\), be a homomorphism of
\(\mathcal A_i\) into a von Neumann algebra \(\mathcal B_i\). If, for
every \(T=(T_i)\in\prod_{i\in I}\mathcal A_i\), we put
\[
  \Phi(T)=\bigl(\Phi_i(T_i)\bigr)_{i\in I}
  \in\prod_{i\in I}\mathcal B_i,
\]
then \(\Phi\) is a homomorphism of
\(\prod_{i\in I}\mathcal A_i\) into
\(\prod_{i\in I}\mathcal B_i\). There is an analogous property for
antihomomorphisms.

\noindent\textbf{Bibliography:} \ArticleRef{118}.

\NumberedParagraph{3}{Operators in a tensor product of Hilbert spaces}
\label{no:I-2-3}
Let \(\mathfrak H_1\) and \(\mathfrak H_2\) be two complex Hilbert spaces.
Form their algebraic tensor product \(\mathfrak H_0\), which is a complex
vector space. There is a unique pre-Hilbert-space structure on
\(\mathfrak H_0\) such that
\[
  (x_1\otimes x_2\mid y_1\otimes y_2)
    =(x_1\mid y_1)(x_2\mid y_2)
\]
for \(x_1,y_1\in\mathfrak H_1\) and
\(x_2,y_2\in\mathfrak H_2\). This structure is separated. The Hilbert
space \(\mathfrak H\) obtained by completing \(\mathfrak H_0\) is called
the Hilbert tensor product of \(\mathfrak H_1\) and \(\mathfrak H_2\),
and is denoted, by abuse of notation, by
\(\mathfrak H_1\otimes\mathfrak H_2\).

Let \(T_1\in\mathcal L(\mathfrak H_1)\) and
\(T_2\in\mathcal L(\mathfrak H_2)\). The algebraic tensor product
\(T_0\) of \(T_1\) and \(T_2\) acts on \(\mathfrak H_0\) and is
continuous. Indeed, it is enough to prove this successively when
\(T_2=I_{\mathfrak H_2}\) and when \(T_1=I_{\mathfrak H_1}\). Suppose,
for example, that \(T_2=I_{\mathfrak H_2}\). Let
\(\sum_{i=1}^{n}x_1^i\otimes x_2^i\) be an element of
\(\mathfrak H_0\); we may suppose that the \(x_2^i\) are orthonormal.
Then
\begin{align*}
  \left\|T_0\left(\sum_{i=1}^{n}x_1^i\otimes x_2^i\right)\right\|^2
    &=
      \left\|\sum_{i=1}^{n}T_1x_1^i\otimes x_2^i\right\|^2\\
    &=\sum_{i=1}^{n}\lVert T_1x_1^i\rVert^2\\
    &\leq \lVert T_1\rVert^2
      \sum_{i=1}^{n}\lVert x_1^i\rVert^2\\
    &=\lVert T_1\rVert^2
      \left\|\sum_{i=1}^{n}x_1^i\otimes x_2^i\right\|^2.
\end{align*}
The continuous extension of \(T_0\) to \(\mathfrak H\) is called, by
abuse of language, the tensor product of \(T_1\) and \(T_2\), and is
denoted by \(T_1\otimes T_2\). We see at once that
\(T_1\otimes T_2\) is a bilinear function of \(T_1\) and \(T_2\), that
\[
  (S_1T_1)\otimes(S_2T_2)
    =(S_1\otimes S_2)(T_1\otimes T_2),
\]
and that
\[
  (T_1\otimes T_2)^*=T_1^*\otimes T_2^*.
\]

% Source PDF page 029 (printed p. 22)
\phantomsection\label{srcpage:22}
The Hilbert tensor product
\(\mathfrak H_1\otimes\mathfrak H_2\otimes\cdots\otimes\mathfrak H_n\)
of Hilbert spaces \(\mathfrak H_1,\mathfrak H_2,\ldots,\mathfrak H_n\),
and the tensor product
\[
  T_1\otimes T_2\otimes\cdots\otimes T_n
  \in\mathcal L(
    \mathfrak H_1\otimes\mathfrak H_2\otimes\cdots\otimes\mathfrak H_n)
\]
of operators \(T_i\in\mathcal L(\mathfrak H_i)\), are defined
analogously.

We shall study \(\mathfrak H_1\otimes\mathfrak H_2\) asymmetrically. Let
\((e_i)_{i\in I}\) be an orthonormal basis of \(\mathfrak H_2\), which
allows us canonically to identify \(\mathfrak H_2\) with the space
\(L_{\mathbb C}^{2}(I)\) of families \((\lambda_i)_{i\in I}\) of complex
numbers such that
\[
  \sum_{i\in I}\lvert\lambda_i\rvert^2<+\infty.
\]
The mapping \(x_1\mapsto x_1\otimes e_i\) is a linear isometry \(U_i\)
of \(\mathfrak H_1\) onto a complete, and hence closed, vector subspace
\(\mathfrak H_i'\) of \(\mathfrak H\). The spaces
\(\mathfrak H_i'\) are pairwise orthogonal; the vector subspace of
\(\mathfrak H_0\) spanned by the \(\mathfrak H_i'\) is dense in
\(\mathfrak H_0\). Thus \(\mathfrak H\) is the Hilbert direct sum of the
\(\mathfrak H_i'\). It follows that every \(x\in\mathfrak H\) can be
written uniquely in the form
\[
  x=\sum_{i\in I}x_1^i\otimes e_i,
\]
where \((x_1^i)_{i\in I}\) is an arbitrary family of elements of
\(\mathfrak H_1\) such that
\(\sum_{i\in I}\lVert x_1^i\rVert^2<+\infty\), and
\[
  \lVert x\rVert^2=\sum_{i\in I}\lVert x_1^i\rVert^2.
\]

Conversely, let \(\mathfrak H\) be a complex Hilbert space which is the
Hilbert direct sum of a family \((\mathfrak H_i')_{i\in I}\) of pairwise
orthogonal closed vector subspaces; let \(\mathfrak H_1\) be another
Hilbert space, and, for every \(i\in I\), let \(U_i\) be an isomorphism
of the Hilbert space \(\mathfrak H_1\) onto the Hilbert space
\(\mathfrak H_i'\). On putting
\(\mathfrak H_2=L_{\mathbb C}^{2}(I)\), the Hilbert space having the
canonical orthonormal basis \((e_i)_{i\in I}\), the preceding data
define a canonical isomorphism of
\(\mathfrak H_1\otimes\mathfrak H_2\) onto \(\mathfrak H\), which
assigns to \(x_1\otimes e_i\) the element \(U_i(x_1)\) of
\(\mathfrak H_i'\). We shall often identify
\(\mathfrak H_1\otimes\mathfrak H_2\) with \(\mathfrak H\) by this
isomorphism.

Retaining the preceding notation, \(U_i^*\) is a linear mapping from
\(\mathfrak H\) into \(\mathfrak H_1\) which vanishes on
\(\mathfrak H\ominus\mathfrak H_i'\) and maps
\(\mathfrak H_i'\) isometrically onto \(\mathfrak H_1\);
\[
  U_i^*U_i=I_{\mathfrak H_1},
  \qquad
  U_iU_i^*=P_{\mathfrak H_i'}.
\]
For every \(T\in\mathcal L(\mathfrak H)\), we have
\(U_i^*TU_\kappa\in\mathcal L(\mathfrak H_1)\), and \(T\) is completely
determined by the \(U_i^*TU_\kappa=T_{i\kappa}\), or, equivalently, by
the matrix \((T_{i\kappa})\). By abuse of notation, we shall write
\(T=(T_{i\kappa})\).

For example, suppose that \(\mathfrak H_1\) is the one-dimensional
Hilbert space \(\mathbb C\). The mapping
\(1\otimes x_2\mapsto x_2\) then canonically identifies
\(\mathfrak H=\mathfrak H_1\otimes\mathfrak H_2\) with
\(\mathfrak H_2=L_{\mathbb C}^{2}(I)\). The operators
\(T_{i\kappa}\) on \(\mathfrak H_1\) are canonically identified with
complex numbers \(\lambda_{i\kappa}\). One readily sees that
\[
  \lambda_{i\kappa}=(Te_\kappa\mid e_i).
\]

% Source PDF page 030 (printed p. 23)
\phantomsection\label{srcpage:23}
Returning to the general case, the representation of \(T\) by the matrix
\((T_{i\kappa})\) thus extends the usual matrix representation of
operators. If \(S=(S_{i\kappa})\) and \(T=(T_{i\kappa})\), we have at once
\[
  \lambda S=(\lambda S_{i\kappa}),\qquad
  S+T=(S_{i\kappa}+T_{i\kappa}),\qquad
  S^*=(S_{\kappa i}^*),\qquad
  ST=(R_{i\kappa}),
\]
where, for every \(x_1\in\mathfrak H_1\),
\begin{align*}
  R_{i\kappa}x_1
    &=U_i^*STU_\kappa x_1\\
    &=U_i^*S
      \left(\sum_{\alpha\in I}U_\alpha U_\alpha^*\right)
      TU_\kappa x_1\\
    &=\sum_{\alpha\in I}S_{i\alpha}T_{\alpha\kappa}x_1,
\end{align*}
the series converging strongly in \(\mathfrak H_1\).

\emph{Examples.} Let \(T_1\in\mathcal L(\mathfrak H_1)\) and
\(T_2\in\mathcal L(\mathfrak H_2)\), and put
\[
  (T_2e_\kappa\mid e_i)=\lambda_{i\kappa},
\]
so that \((\lambda_{i\kappa})\) is the matrix representing \(T_2\), in
the usual sense, with respect to the basis \((e_i)\). Then the matrix
\((T_{i\kappa})\) of \(T=T_1\otimes T_2\) is given by
\(T_{i\kappa}=\lambda_{i\kappa}T_1\). Indeed, for
\(x_1\in\mathfrak H_1\),
\begin{align*}
  U_i^*(T_1\otimes T_2)U_\kappa x_1
    &=U_i^*(T_1\otimes T_2)(x_1\otimes e_\kappa)\\
    &=U_i^*(T_1x_1\otimes T_2e_\kappa)\\
    &=U_i^*\left(
        \sum_{\alpha\in I}T_1x_1\otimes
          \lambda_{\alpha\kappa}e_\alpha\right)\\
    &=\sum_{\alpha\in I}\lambda_{\alpha\kappa}
      U_i^*U_\alpha T_1x_1
     =\lambda_{i\kappa}T_1x_1.
\end{align*}
In particular, the matrix representing
\(T_1\otimes I_{\mathfrak H_2}\) is
\((\delta_{i\kappa}T_1)\), where \(\delta_{i\kappa}\) denotes the
Kronecker symbol, and the matrix representing
\(I_{\mathfrak H_1}\otimes T_2\) is
\((\lambda_{i\kappa}I_{\mathfrak H_1})\).

\begin{lemma}\label{lem:I-2-1}
If \(T\in\mathcal L(\mathfrak H)\) commutes with the
\(U_iU_\kappa^*\), then \(T\) is of the form
\(T_1\otimes I_{\mathfrak H_2}\), with
\(T_1\in\mathcal L(\mathfrak H_1)\).
\end{lemma}

\begin{proof}
Choose a particular element \(\alpha\) of \(I\). We have
\[
  T_{i\kappa}=U_i^*TU_\kappa
    =U_\alpha^*U_\alpha U_i^*TU_\kappa
    =U_\alpha^*TU_\alpha U_i^*U_\kappa.
\]
Now \(U_i^*U_\kappa=0\) if \(i\ne\kappa\), and
\(U_i^*U_i=I_{\mathfrak H_1}\). Hence
\[
  T_{i\kappa}=\delta_{i\kappa}T_1
\]
for some \(T_1\in\mathcal L(\mathfrak H_1)\).
\end{proof}

\begin{lemma}\label{lem:I-2-2}
For every subset \(\mathfrak M\) of \(\mathcal L(\mathfrak H_1)\), let
\(\mathfrak N_{\mathfrak M}\) be the set of
\(T=(T_{i\kappa})\in\mathcal L(\mathfrak H)\) such that
\(T_{i\kappa}\in\mathfrak M\) for every \(i,\kappa\in I\); and let
\(\mathfrak T_{\mathfrak M}\) be the set of
\(T=(T_{i\kappa})\in\mathcal L(\mathfrak H)\) such that:
\begin{enumerate}[label=\textup{\arabic*}\textsuperscript{\(\circ\)}]
\item \(T_{i\kappa}=0\) for \(i\ne\kappa\);
\item \(T_{ii}\) is an element of \(\mathfrak M\) independent of \(i\).
\end{enumerate}
Then
\[
  (\mathfrak T_{\mathfrak M})'=\mathfrak N_{\mathfrak M'},
  \qquad
  (\mathfrak T_{\mathfrak M})''=\mathfrak T_{\mathfrak M''};
\]
and, if \(\mathfrak M\) contains \(0\) and \(I_{\mathfrak H_1}\),
\[
  (\mathfrak N_{\mathfrak M})'=\mathfrak T_{\mathfrak M'},
  \qquad
  (\mathfrak N_{\mathfrak M})''=\mathfrak N_{\mathfrak M''}.
\]
\end{lemma}

\begin{proof}
Let \(R=(R_{i\kappa})\in\mathcal L(\mathfrak H)\) and
\(S=(S_{i\kappa})\in\mathcal L(\mathfrak H)\), and suppose that
\(S_{i\kappa}=0\) for \(i\ne\kappa\), while \(S_{ii}\) is an element
\(T\) of \(\mathcal L(\mathfrak H_1)\) independent of \(i\). The
condition \(RS=SR\) translates at once into
\[
  R_{i\kappa}T=TR_{i\kappa}
  \qquad(i,\kappa\in I).
\]
It follows that
\((\mathfrak T_{\mathfrak M})'=\mathfrak N_{\mathfrak M'}\).
Replacing \(\mathfrak M\) by \(\mathfrak M'\), we deduce that
\[
  (\mathfrak T_{\mathfrak M'})'
    =\mathfrak N_{\mathfrak M''}
    \supset\mathfrak N_{\mathfrak M},
\]
and hence
\(\mathfrak T_{\mathfrak M'}\subset
(\mathfrak N_{\mathfrak M})'\). Let us prove that
\((\mathfrak N_{\mathfrak M})'\subset\mathfrak T_{\mathfrak M'}\) if
\(0,I_{\mathfrak H_1}\in\mathfrak M\). If
\(R=(R_{i\kappa})\in(\mathfrak N_{\mathfrak M})'\), then
% Source PDF page 031 (printed p. 24)
\phantomsection\label{srcpage:24}
first
\[
  (\mathfrak N_{\mathfrak M})'
    \subset(\mathfrak T_{\mathfrak M})'
    =\mathfrak N_{\mathfrak M'},
\]
so \(R_{i\kappa}\in\mathfrak M'\) for all \(i,\kappa\in I\). Next,
\(U_iU_\kappa^*\in\mathfrak N_{\mathfrak M}\), since every entry of the
matrix representing \(U_iU_\kappa^*\) is either \(0\) or
\(I_{\mathfrak H_1}\). Thus \(R\) commutes with the
\(U_iU_\kappa^*\), and \cref{lem:I-2-1} proves that
\(R\in\mathfrak T_{\mathfrak M'}\). We have therefore proved that
\[
  (\mathfrak N_{\mathfrak M})'=\mathfrak T_{\mathfrak M'}
  \quad\text{if}\quad
  0,I_{\mathfrak H_1}\in\mathfrak M.
\]
Since \(0,I_{\mathfrak H_1}\in\mathfrak M'\) for every
\(\mathfrak M\), we have
\[
  (\mathfrak T_{\mathfrak M})''
    =(\mathfrak N_{\mathfrak M'})'
    =\mathfrak T_{\mathfrak M''}
\]
for every \(\mathfrak M\). Finally, if
\(0,I_{\mathfrak H_1}\in\mathfrak M\), then
\[
  (\mathfrak N_{\mathfrak M})''
    =(\mathfrak T_{\mathfrak M'})'
    =\mathfrak N_{\mathfrak M''}.
\]
\end{proof}

In particular, the operators in \(\mathcal L(\mathfrak H)\) which
commute with the operators of the form
\(I_{\mathfrak H_1}\otimes T_2\), where
\(T_2\in\mathcal L(\mathfrak H_2)\), are the operators of the form
\(T_1\otimes I_{\mathfrak H_2}\), where
\(T_1\in\mathcal L(\mathfrak H_1)\).

These lemmas will be useful not only in
\hyperref[no:I-2-4]{no.~4}, but also in
\hyperref[lem:I-3-6]{\S~3, Lemma~6}.

\noindent\textbf{Bibliography:} \ArticleRef{65}.

\NumberedParagraph{4}{Tensor product of von Neumann algebras}
\label{no:I-2-4}
Retain the notation of \hyperref[no:I-2-3]{no.~3}. Let, in addition,
\(\mathcal A_1\) be a von Neumann algebra on \(\mathfrak H_1\), and
\(\mathcal A_2\) a von Neumann algebra on \(\mathfrak H_2\). The
operators of the form
\[
  R_1\otimes R_2+S_1\otimes S_2+\cdots+T_1\otimes T_2,
\]
where
\[
  R_1,S_1,\ldots,T_1\in\mathcal A_1,
  \qquad
  R_2,S_2,\ldots,T_2\in\mathcal A_2,
\]
form an involutive algebra of operators \(\mathcal A_0\), containing
\(I_{\mathfrak H}\), on \(\mathfrak H\). The von Neumann algebra
generated by \(\mathcal A_0\) is called the \emph{tensor-product von
Neumann algebra} of \(\mathcal A_1\) and \(\mathcal A_2\), and is
denoted by \(\mathcal A_1\otimes\mathcal A_2\), by abuse of notation.
The tensor product
\(\mathcal A_1\otimes\mathcal A_2\otimes\cdots\otimes\mathcal A_n\)
of von Neumann algebras
\(\mathcal A_1,\mathcal A_2,\ldots,\mathcal A_n\) is defined
analogously.

If \(\mathcal A_2=\mathbb C_{\mathfrak H_2}\), then \(\mathcal A_0\)
is the set of operators
\(T_1\otimes I_{\mathfrak H_2}\), with \(T_1\in\mathcal A_1\).
By \cref{lem:I-2-2}, \(\mathcal A_0\) is a von Neumann algebra, since
\(\mathcal A_1''=\mathcal A_1\). Thus the mapping
\[
  T_1\longmapsto T_1\otimes I_{\mathfrak H_2}
\]
is an isomorphism of \(\mathcal A_1\) onto
\(\mathcal A_1\otimes\mathbb C_{\mathfrak H_2}\), called the
\emph{amplification} of \(\mathcal A_1\) to
\(\mathcal A_1\otimes\mathbb C_{\mathfrak H_2}\). Similarly, the
mapping
\[
  T_2\longmapsto I_{\mathfrak H_1}\otimes T_2
\]
is an isomorphism of \(\mathcal A_2\) onto
\(\mathbb C_{\mathfrak H_1}\otimes\mathcal A_2\). And
\(\mathcal A_1\otimes\mathcal A_2\) is the von Neumann algebra
generated by
\(\mathcal A_1\otimes\mathbb C_{\mathfrak H_2}\) and
\(\mathbb C_{\mathfrak H_1}\otimes\mathcal A_2\).

\begin{proposition}\label{prop:I-2-4}
\begin{enumerate}[label=(\roman*)]
\item The operators in
\(\mathcal A_1\otimes\mathbb C_{\mathfrak H_2}\) are the operators
representable by a matrix \((T_{i\kappa})\) such that
\(T_{i\kappa}=0\) for \(i\ne\kappa\), and such that \(T_{ii}\) is an
element of \(\mathcal A_1\) independent of \(i\).
\item The operators in
\(\mathcal A_1\otimes\mathcal L(\mathfrak H_2)\) are the operators
representable by a matrix \((T_{i\kappa})\) such that
\(T_{i\kappa}\in\mathcal A_1\) for all \(i,\kappa\in I\).
\item We have
\[
  \mathcal A_1\otimes\mathbb C_{\mathfrak H_2}
    =\bigl(\mathcal A_1'\otimes\mathcal L(\mathfrak H_2)\bigr)'.
\]
\end{enumerate}
\end{proposition}

\begin{proof}
Part (i) follows at once from what precedes. The algebra
\(\bigl(\mathcal A_1'\otimes\mathcal L(\mathfrak H_2)\bigr)'\) is the
intersection of
\(\bigl(\mathcal A_1'\otimes\mathbb C_{\mathfrak H_2}\bigr)'\) and
\[
  \bigl(\mathbb C_{\mathfrak H_1}\otimes
    \mathcal L(\mathfrak H_2)\bigr)'
  =\mathcal L(\mathfrak H_1)\otimes\mathbb C_{\mathfrak H_2}.
\]
% Source PDF page 032 (printed p. 25)
\phantomsection\label{srcpage:25}
Now, for an operator of the form
\(T_1\otimes I_{\mathfrak H_2}\) to commute with
\(\mathcal A_1'\otimes\mathbb C_{\mathfrak H_2}\), it is necessary and
sufficient that \(T_1\in\mathcal A_1\). This proves (iii). Finally, (ii)
is a consequence of (i), (iii), and \cref{lem:I-2-2}.
\end{proof}

Still using the notation of \hyperref[no:I-2-3]{no.~3}, the projections
\[
  E_i=P_{\mathfrak H_i'}
\]
belong to
\(\mathbb C_{\mathfrak H_1}\otimes\mathcal L(\mathfrak H_2)\), and
therefore, a fortiori, to
\[
  \mathcal A_1\otimes\mathcal L(\mathfrak H_2)
    =\bigl(\mathcal A_1'\otimes
      \mathbb C_{\mathfrak H_2}\bigr)'.
\]
Likewise, the partial isometries
\[
  U_\kappa U_i^*=U_{i\kappa},
\]
whose initial projection is \(E_i\) and final projection is \(E_\kappa\),
belong to
\[
  \mathcal A_1\otimes\mathcal L(\mathfrak H_2)
    =\bigl(\mathcal A_1'\otimes
      \mathbb C_{\mathfrak H_2}\bigr)'.
\]
This being so, we have the following result.

\begin{proposition}\label{prop:I-2-5}
\begin{enumerate}[label=(\roman*)]
\item
\(\bigl(\mathcal A_1\otimes\mathcal L(\mathfrak H_2)\bigr)_{E_i}\)
is spatially isomorphic to \(\mathcal A_1\), and
\(\bigl(\mathcal A_1'\otimes
\mathbb C_{\mathfrak H_2}\bigr)_{E_i}\) is spatially isomorphic to
\(\mathcal A_1'\).
\item Let \(\mathcal A\) be a von Neumann algebra on a complete Hilbert
space \(\mathfrak H\), and let \((E_i)_{i\in I}\) be a family of pairwise
orthogonal projections in \(\mathcal A\), with sum
\(I_{\mathfrak H}\). Suppose that, for every \(i,\kappa\in I\), there is
a partial isometry \(U_{i\kappa}\in\mathcal A\) with initial projection
\(E_i\) and final projection \(E_\kappa\). Let \(\alpha\) be a particular
element of \(I\), and put
\[
  \mathfrak H_1=E_\alpha(\mathfrak H),
  \qquad
  \mathfrak H_2=L_{\mathbb C}^{2}(I).
\]
Then \(\mathcal A\) is spatially isomorphic to the von Neumann algebra
\[
  \mathcal A_{E_\alpha}\otimes\mathcal L(\mathfrak H_2)
\]
on \(\mathfrak H_1\otimes\mathfrak H_2\), and \(\mathcal A'\) is
spatially isomorphic to
\[
  \mathcal A'_{E_\alpha}\otimes\mathbb C_{\mathfrak H_2}.
\]
\end{enumerate}
\end{proposition}

\begin{proof}
The proof will be more precise than the statement, since it will define
canonical isomorphisms.

Let us prove (i). The isomorphism \(U_i\) of \(\mathfrak H_1\) onto
\(\mathfrak H_i'\) transforms \(\mathcal A_1'\) into
\[
  \bigl(\mathcal A_1'\otimes
    \mathbb C_{\mathfrak H_2}\bigr)_{E_i}
\]
by \cref{prop:I-2-4}(i), and therefore transforms \(\mathcal A_1\) into
\[
  \bigl(\mathcal A_1\otimes
    \mathcal L(\mathfrak H_2)\bigr)_{E_i}.
\]

Let us prove (ii). Let \(U_i\) be the isomorphism of
\(\mathfrak H_1\) onto
\(\mathfrak H_i'=E_i(\mathfrak H)\) defined by \(U_{\alpha i}\).
The family of the \(U_i\) defines a canonical isomorphism of
\(\mathfrak H\) onto \(\mathfrak H_1\otimes\mathfrak H_2\)
(\hyperref[no:I-2-3]{no.~3}), by which we identify these two spaces. For
\(T\in\mathcal A\), the representing matrix \((T_{i\kappa})\) of \(T\)
is given by
\[
  T_{i\kappa}=U_i^*TU_\kappa
    =\bigl(U_{\alpha i}^*TU_{\alpha\kappa}\bigr)_{E_\alpha}
    \in\mathcal A_{E_\alpha};
\]
therefore
\[
  \mathcal A\subset
  \mathcal A_{E_\alpha}\otimes\mathcal L(\mathfrak H_2).
\]
For \(T'\in\mathcal A'\), the representing matrix
\((T_{i\kappa}')\) of \(T'\) is given by
\begin{align*}
  T_{i\kappa}'
    &=\bigl(U_{\alpha i}^*T'U_{\alpha\kappa}\bigr)_{E_\alpha}\\
    &=\bigl(T'U_{\alpha i}^*U_{\alpha\kappa}\bigr)_{E_\alpha}\\
    &=\delta_{i\kappa}T'_{E_\alpha}
      \in\mathcal A'_{E_\alpha};
\end{align*}
therefore
\[
  \mathcal A'\subset
  \mathcal A'_{E_\alpha}\otimes\mathbb C_{\mathfrak H_2}.
\]
Consequently
\[
  \mathcal A\supset
  \mathcal A_{E_\alpha}\otimes\mathcal L(\mathfrak H_2),
\]
and finally
\[
  \mathcal A
    =\mathcal A_{E_\alpha}\otimes\mathcal L(\mathfrak H_2),
  \qquad
  \mathcal A'
    =\mathcal A'_{E_\alpha}\otimes\mathbb C_{\mathfrak H_2}.
\]
\end{proof}

% Source PDF page 033 (printed p. 26)
\phantomsection\label{srcpage:26}
Thus the passage from \(\mathcal A_1\) to
\(\mathcal A_1\otimes\mathcal L(\mathfrak H_2)\) (respectively, to
\(\mathcal A_1\otimes\mathbb C_{\mathfrak H_2}\)), and the passage from
a von Neumann algebra to a reduced (respectively, induced) algebra are,
to a certain extent, reciprocal operations. The inverse isomorphism of
an amplification is an induction.

\begin{proposition}\label{prop:I-2-6}
For \(i=1,2,\ldots,n\), let \(\mathfrak H_i\) be a complete Hilbert
space, let \(\mathfrak M_i\) be a subset of
\(\mathcal L(\mathfrak H_i)\), and let \(\mathcal A_i\) be the von
Neumann algebra generated by \(\mathfrak M_i\). Let \(\mathcal A\) be
the von Neumann algebra generated by the operators
\begin{gather*}
  T_1\otimes I_{\mathfrak H_2}\otimes\cdots\otimes I_{\mathfrak H_n}
    \quad(T_1\in\mathfrak M_1),\\
  I_{\mathfrak H_1}\otimes T_2\otimes\cdots\otimes I_{\mathfrak H_n}
    \quad(T_2\in\mathfrak M_2),\\
  \ldots,\qquad
  I_{\mathfrak H_1}\otimes I_{\mathfrak H_2}\otimes\cdots\otimes T_n
    \quad(T_n\in\mathfrak M_n).
\end{gather*}
Then
\[
  \mathcal A
    =\mathcal A_1\otimes\mathcal A_2\otimes\cdots\otimes\mathcal A_n.
\]
\end{proposition}

\begin{proof}
If \(\mathfrak M_i\) is replaced by
\(\mathfrak M_i\cup\mathfrak M_i^*\), the algebras
\(\mathcal A_i\) and \(\mathcal A\) remain unchanged. We may therefore
suppose that \(\mathfrak M_i=\mathfrak M_i^*\). Then
\(\mathcal A_i=\mathfrak M_i''\). \Cref{lem:I-2-2} proves that
\[
  \mathcal A_1\otimes
    \mathbb C_{\mathfrak H_2}\otimes\cdots\otimes
    \mathbb C_{\mathfrak H_n}
\]
is the von Neumann algebra generated by the operators
\[
  T_1\otimes I_{\mathfrak H_2}\otimes\cdots\otimes
  I_{\mathfrak H_n},
  \qquad T_1\in\mathfrak M_1.
\]
Thus
\[
  \mathcal A_1\otimes
    \mathbb C_{\mathfrak H_2}\otimes\cdots\otimes
    \mathbb C_{\mathfrak H_n}
  \subset\mathcal A.
\]
Similarly,
\[
  \mathbb C_{\mathfrak H_1}\otimes\cdots\otimes
    \mathcal A_i\otimes\cdots\otimes
    \mathbb C_{\mathfrak H_n}
  \subset\mathcal A
  \qquad(1\leq i\leq n),
\]
and therefore
\[
  \mathcal A_1\otimes\mathcal A_2\otimes\cdots\otimes\mathcal A_n
  \subset\mathcal A.
\]
The reverse inclusion is evident.
\end{proof}

The tensor product of von Neumann algebras is \emph{associative} in the
following sense. Let
\(\mathfrak H_1,\mathfrak H_2,\ldots,\mathfrak H_n\) be complex Hilbert
spaces, and put
\[
  \mathfrak H
    =\mathfrak H_1\otimes\mathfrak H_2\otimes\cdots\otimes\mathfrak H_n,
  \qquad
  \mathfrak K
    =\mathfrak H_1\otimes\mathfrak H_2\otimes\cdots\otimes\mathfrak H_p.
\]
The Hilbert space \(\mathfrak H\) is canonically identified with
\(\mathfrak K\otimes\mathfrak H_{p+1}\otimes\cdots\otimes
\mathfrak H_n\). Let \(\mathcal A_i\) be a von Neumann algebra on
\(\mathfrak H_i\), and put
\[
  \mathcal A
    =\mathcal A_1\otimes\mathcal A_2\otimes\cdots\otimes\mathcal A_n,
  \qquad
  \mathcal B
    =\mathcal A_1\otimes\mathcal A_2\otimes\cdots\otimes\mathcal A_p.
\]
The von Neumann algebra \(\mathcal A\) is generated by the algebras
\begin{gather*}
  \mathcal A_1\otimes\mathbb C_{\mathfrak H_2}\otimes\cdots\otimes
    \mathbb C_{\mathfrak H_n},\\
  \mathbb C_{\mathfrak H_1}\otimes\mathcal A_2\otimes\cdots\otimes
    \mathbb C_{\mathfrak H_n},\qquad\ldots,\qquad
  \mathbb C_{\mathfrak H_1}\otimes\mathbb C_{\mathfrak H_2}
    \otimes\cdots\otimes\mathcal A_n.
\end{gather*}
The same is true of
\[
  \mathcal B\otimes\mathcal A_{p+1}\otimes
    \mathcal A_{p+2}\otimes\cdots\otimes\mathcal A_n
\]
by \cref{prop:I-2-6}, since \(\mathcal B\) is the von Neumann algebra
generated by
\begin{gather*}
  \mathcal A_1\otimes\mathbb C_{\mathfrak H_2}\otimes\cdots\otimes
    \mathbb C_{\mathfrak H_p},\\
  \mathbb C_{\mathfrak H_1}\otimes\mathcal A_2\otimes\cdots\otimes
    \mathbb C_{\mathfrak H_p},\qquad\ldots,\qquad
  \mathbb C_{\mathfrak H_1}\otimes\mathbb C_{\mathfrak H_2}
    \otimes\cdots\otimes\mathcal A_p.
\end{gather*}
Therefore
% Source PDF page 034 (printed p. 27)
\phantomsection\label{srcpage:27}
\[
  \mathcal A
    =\mathcal B\otimes\mathcal A_{p+1}\otimes\cdots\otimes\mathcal A_n.
\]

Let \(\mathcal A_1\) (respectively, \(\mathcal A_2\)) be a von Neumann
algebra on \(\mathfrak H_1\) (respectively, \(\mathfrak H_2\)), and let
\(E_1=P_{\mathfrak K_1}\) (respectively,
\(E_2=P_{\mathfrak K_2}\)) be a projection in \(\mathcal A_1\)
(respectively, \(\mathcal A_2\)). The operator \(E_1\otimes E_2\) is a
projection in \(\mathcal A_1\otimes\mathcal A_2\), and
\((E_1\otimes E_2)(\mathfrak H_1\otimes\mathfrak H_2)\) is identified
with \(\mathfrak K_1\otimes\mathfrak K_2\). This being so,
\[
  (\mathcal A_1\otimes\mathcal A_2)_{E_1\otimes E_2}
    =(\mathcal A_1)_{E_1}\otimes(\mathcal A_2)_{E_2}.
\]
Indeed, \(\mathcal A_1\otimes\mathcal A_2\) is generated by the
operators
\[
  R_1\otimes R_2+S_1\otimes S_2+\cdots+T_1\otimes T_2,
\]
where
\[
  R_1,S_1,\ldots,T_1\in\mathcal A_1,
  \qquad
  R_2,S_2,\ldots,T_2\in\mathcal A_2.
\]
Hence, by \cref{prop:I-2-1}(ii),
\((\mathcal A_1\otimes\mathcal A_2)_{E_1\otimes E_2}\) is generated by
the operators
\begin{align*}
&\Bigl((E_1\otimes E_2)
  (R_1\otimes R_2+S_1\otimes S_2+\cdots+T_1\otimes T_2)
  (E_1\otimes E_2)\Bigr)_{E_1\otimes E_2}\\
&\quad
  =(E_1R_1E_1)_{E_1}\otimes(E_2R_2E_2)_{E_2}
    +\cdots+
    (E_1T_1E_1)_{E_1}\otimes(E_2T_2E_2)_{E_2}\\
&\quad\in(\mathcal A_1)_{E_1}\otimes(\mathcal A_2)_{E_2}.
\end{align*}
Thus
\[
  (\mathcal A_1\otimes\mathcal A_2)_{E_1\otimes E_2}
    \subset
  (\mathcal A_1)_{E_1}\otimes(\mathcal A_2)_{E_2};
\]
the reverse inclusion is evident. Analogously, we see that
\[
  (\mathcal A_1'\otimes\mathcal A_2')_{E_1\otimes E_2}
    =(\mathcal A_1')_{E_1}\otimes(\mathcal A_2')_{E_2}.
\]

Finally, for every \(i\in I\) (respectively, \(\kappa\in K\)), let
\(\mathcal A_i\) (respectively, \(\mathcal B_\kappa\)) be a von Neumann
algebra on \(\mathfrak H_i\) (respectively, \(\mathfrak K_\kappa\)).
Let
\[
  \mathcal A=\prod_{i\in I}\mathcal A_i
  \quad\text{on}\quad
  \mathfrak H=\bigoplus_{i\in I}\mathfrak H_i,
  \qquad
  \mathcal B=\prod_{\kappa\in K}\mathcal B_\kappa
  \quad\text{on}\quad
  \mathfrak K=\bigoplus_{\kappa\in K}\mathfrak K_\kappa.
\]
The space \(\mathfrak H\otimes\mathfrak K\) is canonically identified
with
\[
  \bigoplus_{(i,\kappa)\in I\times K}
    (\mathfrak H_i\otimes\mathfrak K_\kappa).
\]
We then have
\[
  \mathcal A\otimes\mathcal B
    =\prod_{(i,\kappa)\in I\times K}
      (\mathcal A_i\otimes\mathcal B_\kappa);
\]
we leave its verification to the reader.

\Cref{prop:I-2-4} also has the following application.

\begin{proposition}\label{prop:I-2-7}
Let \(\mathfrak H\) be a Hilbert space, let \(\mathcal A\) be a von
Neumann algebra on \(\mathfrak H\), let
\(\mathcal Z=\mathcal A\cap\mathcal A'\), let
\(T_{ij}\) (\(i,j=1,\ldots,n\)) be elements of \(\mathcal A\), and let
\(T_{ij}'\) (\(i,j=1,\ldots,n\)) be elements of \(\mathcal A'\). The
following conditions are equivalent:
\begin{enumerate}[label=(\roman*)]
\item
\[
  \sum_k T_{ik}T_{kj}'=0
  \qquad(i,j=1,\ldots,n).
\]
\item There are elements \(Z_{ij}\) (\(i,j=1,\ldots,n\)) of
\(\mathcal Z\) such that
\[
  \sum_k T_{ik}Z_{kj}=0,
  \qquad
  \sum_k Z_{ik}T_{kj}'=T_{ij}'
  \qquad(i,j=1,\ldots,n).
\]
\end{enumerate}
\end{proposition}

\begin{proof}
If there are \(Z_{ij}\) satisfying (ii), then
\[
  \sum_kT_{ik}T_{kj}'
    =\sum_kT_{ik}\sum_lZ_{kl}T_{lj}'
    =\sum_l\left(\sum_kT_{ik}Z_{kl}\right)T_{lj}'
    =0.
\]

% Source PDF page 035 (printed p. 28)
\phantomsection\label{srcpage:28}
Suppose that
\[
  \sum_kT_{ik}T_{kj}'=0
  \qquad(i,j=1,\ldots,n).
\]
Let \(\mathfrak K\) be the Hilbert direct sum of \(n\) spaces equal to
\(\mathfrak H\). We identify \(\mathfrak K\) with
\(\mathfrak H\otimes\mathfrak L\), where \(\mathfrak L\) is an
\(n\)-dimensional Hilbert space. The elements of
\(\mathcal L(\mathfrak K)\) are identified with \(n\)-row,
\(n\)-column matrices having entries in \(\mathcal L(\mathfrak H)\).
Let
\[
  T=(T_{ij})_{i,j=1,\ldots,n}
    \in\mathcal A\otimes\mathcal L(\mathfrak L),
  \qquad
  T'=(T_{ij}')_{i,j=1,\ldots,n}
    \in\mathcal A'\otimes\mathcal L(\mathfrak L).
\]
We have \(TT'=0\). Let \(\mathcal E\) be the set of projections
\(E'\in\mathcal A'\otimes\mathcal L(\mathfrak L)\) such that
\(TE'=0\), and let \(E_0'\) be the supremum of the elements of
\(\mathcal E\). The projection onto \(\overline{T'(\mathfrak K)}\) is an element
of \(\mathcal E\), since \(TT'=0\) and
\(T'\in\mathcal A'\otimes\mathcal L(\mathfrak L)\); hence
\[
  E_0'T'=T'.
\]
Moreover, it is clear that \(TE_0'=0\). Thus, if we put
\[
  E_0'=(Z_{ij})_{i,j=1,\ldots,n},
\]
then
\[
  Z_{ij}\in\mathcal A',
  \qquad
  \sum_kT_{ik}Z_{kj}=0,
  \qquad
  \sum_kZ_{ik}T_{kj}'=T_{ij}'.
\]
It remains to prove that \(Z_{ij}\in\mathcal A\). It is enough to prove
that, for every Hermitian \(R'\in\mathcal A'\),
\[
  Z_{ij}R'=R'Z_{ij}.
\]
Let
\[
  S'=R'\otimes I_{\mathfrak L}
    \in\mathcal A'\otimes\mathbb C_{\mathfrak L};
\]
this operator is Hermitian. We have
\[
  TS'E_0'=S'TE_0'=0;
\]
hence the projection onto \(\overline{S'E_0'(\mathfrak K)}\) is an element of
\(\mathcal E\). Consequently
\[
  S'E_0'(\mathfrak K)\subset E_0'(\mathfrak K),
  \qquad
  E_0'S'E_0'=S'E_0'.
\]
Taking adjoints gives
\[
  E_0'S'E_0'=E_0'S',
\]
and hence \(S'E_0'=E_0'S'\), whence
\[
  Z_{ij}R'=R'Z_{ij}.
\]
\end{proof}

Some simple properties of the tensor product will be obtained only later
(\hyperref[prop:I-4-2]{\S~4, Proposition~2};
\hyperref[prop:I-6-14]{\S~6, Corollary to Proposition~14}). The
following problem arises: do we have
\[
  (\mathcal A_1\otimes\mathcal A_2)'
    =\mathcal A_1'\otimes\mathcal A_2'?
\]
(See \hyperref[ex:I-2-5]{Exercise~5},
\hyperref[prop:I-6-14]{\S~6, Proposition~14}, and \ArticleRef{432}.)

Elementary operations, applied to abelian von Neumann algebras, already
lead to all discrete von Neumann algebras
(\hyperref[sec:I-8]{\S~8}).

\noindent\textbf{Bibliography:} \ArticleRef{31}, \ArticleRef{50},
\ArticleRef{65}, \ArticleRef{77}, \ArticleRef{89}, \ArticleRef{128},
\ArticleRef{190}.

\par\medskip
\noindent\textit{Exercises.}---
\begin{enumerate}[label=\arabic*.,leftmargin=*,itemsep=0.5\baselineskip]
\item\label{ex:I-2-1}
Let \(\mathfrak H\) be a complex Hilbert space, and let
\(\mathfrak H_1\) and \(\mathfrak H_2\) be two complementary orthogonal
subspaces of \(\mathfrak H\) having the same Hilbert dimension. Let
\(\mathfrak M\) be the set of operators
\(T\in\mathcal L(\mathfrak H)\) such that
\[
  T(\mathfrak H_1)\subset\mathfrak H_2,
  \qquad
  T(\mathfrak H_2)=0.
\]
Show that the von Neumann algebra generated by \(\mathfrak M\) is
\(\mathcal L(\mathfrak H)\), but that the von Neumann algebra generated
by \(\mathfrak M_{\mathfrak H_1}\) is
\(\mathbb C_{\mathfrak H_1}\).

\item\label{ex:I-2-2}
Let \(\mathcal Z,\mathcal A,\mathcal B\) be von Neumann algebras, with
\(\mathcal Z\subset\mathcal A\subset\mathcal B\), and with
\(\mathcal Z\) contained in the centre of \(\mathcal B\).
\begin{enumerate}[label=\alph*.,leftmargin=2em]
\item Let \(T\in\mathcal B\). Show that, in the set \(\mathfrak M\) of
projections \(E\in\mathcal Z\) such that \(TE\in\mathcal A\), there is
a projection greater than all the others. (Consider a maximal element
of \(\mathfrak M\).) Deduce that, if \(T\notin\mathcal A\), there is a
nonzero projection \(F\in\mathcal Z\) such that, for every nonzero
projection \(G\in\mathcal Z\) majorized by \(F\), one has
\(TG\notin\mathcal A\).
\item Suppose that, for every nonzero projection \(E\in\mathcal Z\),
\(\mathcal A_E\ne\mathcal B_E\). Show that there is a projection
\(T\in\mathcal B\) such that, for every nonzero projection
\(E\in\mathcal Z\), one has \(T_E\notin\mathcal A_E\).

\(\bigl[\)Let \(\mathfrak N\) be the set of projections
\(F\in\mathcal Z\) for which there is a projection
\(S\in\mathcal B_F\) having the following property: for every nonzero
projection \(G\in\mathcal Z\) majorized by \(F\), one has
\(S_G\notin\mathcal A_G\). Consider a maximal family of nonzero,
pairwise orthogonal projections in \(\mathfrak N\), and show, with the
aid of (a), that their sum is \(I_{\mathfrak H}\).\(\bigr]\)
\end{enumerate}

\item\label{ex:I-2-3}
Let \(\mathcal A\) and \(\mathcal B\) be von Neumann algebras. For
\(\mathcal A\otimes\mathcal B\) to be of countable type, it is necessary
and sufficient that \(\mathcal A\) and \(\mathcal B\) be of countable
type.

\(\bigl[\)If \(\mathcal A\) and \(\mathcal B\) are of countable type,
let \((x_i)\) and \((y_j)\) be separating sequences for
\(\mathcal A\) and \(\mathcal B\). Then the double sequence
\((x_i\otimes y_j)\) is totalizing for
\(\mathcal A'\otimes\mathcal B'\), hence for
\((\mathcal A\otimes\mathcal B)'\), and therefore separating for
\(\mathcal A\otimes\mathcal B\).\(\bigr]\)
\end{enumerate}
% Source PDF page 036 (printed p. 29)
\phantomsection\label{srcpage:29}
\begin{exercise}\label{ex:I-2-4}
Let \(\mathcal A\) and \(\mathcal B\) be von Neumann algebras on
\(\mathfrak H\) and \(\mathfrak K\), respectively. For
\(\mathcal A\otimes\mathcal B\) to be a factor, it is necessary and
sufficient that \(\mathcal A\) and \(\mathcal B\) be factors. [If
\(\mathcal A\) and \(\mathcal B\) are factors, show that
\(\mathcal A\otimes\mathcal B\) and
\(\mathcal A'\otimes\mathcal B'\) generate
\(\mathcal L(\mathfrak H\otimes\mathfrak K)\).]
\end{exercise}

\begin{exercise}\label{ex:I-2-5}
Let \(\mathcal A\) (respectively, \(\mathcal B\)) be a von Neumann
algebra on \(\mathfrak H\) (respectively, \(\mathfrak K\)). Let
\(T\in(\mathcal A\otimes\mathcal B)'\).
\begin{enumerate}[label=\textit{\alph*.}]
\item Show that
\[
  (\mathcal A\otimes\mathcal B)'
  =\bigl(\mathcal A'\otimes\mathcal L(\mathfrak K)\bigr)
   \cap
   \bigl(\mathcal L(\mathfrak H)\otimes\mathcal B'\bigr).
\]
[Use the fact that \(\mathcal A\otimes\mathcal B\) is generated by
\(\mathcal A\otimes\mathbb C_{\mathfrak K}\) and
\(\mathbb C_{\mathfrak H}\otimes\mathcal B\).]

\item Let \(x\in\mathfrak H\), \(y\in\mathfrak K\). There exists an
\(S\) in the involutive operator algebra \(\mathcal C_0\) generated by
\(\mathcal A'\otimes\mathbb C_{\mathfrak K}\) and
\(\mathbb C_{\mathfrak H}\otimes\mathcal B'\) such that
\(\lVert(T-S)(x\otimes y)\rVert\) is arbitrarily small. [Show that
\[
  T(x\otimes y)
   =\bigl(E_x^{\mathcal A'}\otimes E_y^{\mathcal B'}\bigr)
      T(x\otimes y);
\]
deduce that \(T(x\otimes y)\) is a limit of elements of the form
\begin{equation*}
  x_1\otimes y_1+\cdots+x_n\otimes y_n,
  \qquad
  x_i\in\mathfrak X_x^{\mathcal A'},\quad
  y_i\in\mathfrak X_y^{\mathcal B'}.
  \tag*{\textup{]}}
\end{equation*}

\item Let \((x_m)\) be a sequence of elements of \(\mathfrak H\) such
that
\[
  \sum_m\lVert x_m\rVert^2<+\infty,
\]
and let \((y_n)\) be a sequence of elements of \(\mathfrak K\) such
that
\[
  \sum_n\lVert y_n\rVert^2<+\infty.
\]
There exists an \(R\in\mathcal C_0\) such that
\[
  \sum_{m,n}\bigl\lVert(T-R)(x_m\otimes y_n)\bigr\rVert^2
\]
is arbitrarily small. [Let \(\mathfrak H_1\) and \(\mathfrak K_1\) be
infinite-dimensional Hilbert spaces; identify \((x_m)\) with an element
\(x\) of \(\mathfrak H_1\otimes\mathfrak H\), and \((y_n)\) with an
element \(y\) of \(\mathfrak K\otimes\mathfrak K_1\); form, on the
space
\[
  \mathfrak H_1\otimes(\mathfrak H\otimes\mathfrak K)
       \otimes\mathfrak K_1
  =(\mathfrak H_1\otimes\mathfrak H)
       \otimes(\mathfrak K\otimes\mathfrak K_1),
\]
the operator
\(I_{\mathfrak H_1}\otimes T\otimes I_{\mathfrak K_1}\), which
commutes with
\[
  \bigl(\mathcal L(\mathfrak H_1)\otimes\mathcal A\bigr)
  \otimes
  \bigl(\mathcal B\otimes\mathcal L(\mathfrak K_1)\bigr),
\]
and apply part~\textit{b}.]
\end{enumerate}
This result makes it plausible that
\[
  (\mathcal A\otimes\mathcal B)'=\mathcal A'\otimes\mathcal B'
\]
(cf., however, \hyperref[sec:I-1]{\S~1}, Exercise~9\textit{b}).
\end{exercise}

\begin{exercise}\label{ex:I-2-6}
Suppose that \(\mathcal A\) is a factor on \(\mathfrak K\), and that
\(T_1,T_2,\ldots,T_n\in\mathcal A\) are linearly independent (over
\(\mathbb C\)). Let \(T_1',T_2',\ldots,T_n'\in\mathcal A'\). Then the
relation
\[
  \sum_{i=1}^n T_iT_i'=0
\]
implies that the \(T_i'\) are zero (use \cref{prop:I-1-7}). Let
\(\mathcal B\) be the involutive operator algebra generated by
\(\mathcal A\) and \(\mathcal A'\), that is, the set of operators
\(\sum_{i=1}^n R_iR_i'\), where \(R_i\in\mathcal A\) and
\(R_i'\in\mathcal A'\). Show that the mapping
\[
  \sum_{i=1}^n R_iR_i'
   \longmapsto
  \sum_{i=1}^n R_i\otimes R_i'
\]
defines an isomorphism of the involutive algebra \(\mathcal B\) onto an
involutive algebra of operators on
\(\mathfrak K\otimes\mathfrak K\) \ArticleRef{65}.
\end{exercise}

% Source PDF page 037 (printed p. 30)
\phantomsection\label{srcpage:30}
\begin{exercise}\label{ex:I-2-7}
Let \(\mathcal A\) be a von Neumann algebra on \(\mathfrak H\). If
\(T_1,\ldots,T_n\) belong to \(\mathcal A\), and
\(x_1,\ldots,x_n\) belong to \(\mathfrak H\), and are such that
\[
  T_1x_1+\cdots+T_nx_n=0,
\]
there exist \(T_{ij}\in\mathcal A\)
\((i,j=1,\ldots,n)\) such that
\[
  x_i=\sum_jT_{ij}x_j,
  \qquad
  \sum_iT_iT_{ij}=0;
\]
in other words, \(\mathfrak H\) is a flat \(\mathcal A\)-module.
[Introduce \(\mathfrak K\) and \(\mathfrak L\) as in the proof of
\cref{prop:I-2-7}. Let
\(x=(x_1,\ldots,x_n)\in\mathfrak K\). Let \(E\) be the cyclic
projection of \(\mathcal A\otimes\mathcal L(\mathfrak L)\) defined by
\(x\). Put \(E=(T_{ij})\).] \ArticleRef{333}.
\end{exercise}

\section{Density Theorems}\label{sec:I-3}

\NumberedParagraph{1}{Topologies on \(\mathcal L(\mathfrak H)\)}
\label{no:I-3-1}
\emph{Norm topology.} --- The norm defines a topology on
\(\mathcal L(\mathfrak H)\), called the norm topology. It is the
topology of uniform convergence on the bounded subsets of
\(\mathfrak H\), where \(\mathfrak H\) is endowed with the strong
topology.

\emph{Strong topology.} --- Let \(x\in\mathfrak H\). The function
\(T\mapsto\lVert Tx\rVert\) is a seminorm on
\(\mathcal L(\mathfrak H)\). The set of these seminorms defines on
\(\mathcal L(\mathfrak H)\) a separated locally convex topology called
the topology of simple strong convergence, or, for short, the strong
topology. A fundamental system of neighborhoods of zero for this
topology is obtained by considering, for every finite sequence
\((x_1,x_2,\ldots,x_n)\) of elements of \(\mathfrak H\), the set of
\(T\in\mathcal L(\mathfrak H)\) such that
\[
  \lVert Tx_i\rVert\leq1\qquad(1\leq i\leq n).
\]
The strong topology may also be defined as the coarsest topology on
\(\mathcal L(\mathfrak H)\) making the mappings
\(T\mapsto Tx\) from \(\mathcal L(\mathfrak H)\) into
\(\mathfrak H\), endowed with its strong topology, continuous.

This topology is compatible with the vector-space structure of
\(\mathcal L(\mathfrak H)\), but not, in general, with its algebra
structure; in other words, the mapping \((S,T)\mapsto ST\) is not in
general strongly continuous (Exercise~2). Nevertheless, if
\(\mathcal L_1(\mathfrak H)\) denotes the unit ball of
\(\mathcal L(\mathfrak H)\), the equality
\[
  ST-S_0T_0=S(T-T_0)+(S-S_0)T_0
\]
shows that the mapping \((S,T)\mapsto ST\) from
\(\mathcal L_1(\mathfrak H)\times\mathcal L(\mathfrak H)\) into
\(\mathcal L(\mathfrak H)\) is strongly continuous. The partial
mappings \(S\mapsto ST\) and \(T\mapsto ST\) are strongly continuous.
When \(\mathfrak H\) is infinite-dimensional, the mapping
\(T\mapsto T^*\) is not strongly continuous. Indeed, let
\((e_1,e_2,\ldots)\) be an orthonormal sequence in \(\mathfrak H\),
and let \(U_n\) be the operator
\[
  x\longmapsto(x\mid e_n)e_1
\]
on \(\mathfrak H\). Then \(U_n\) tends strongly to zero; but
\[
  (U_n^*e_1\mid x)
   =(e_1\mid(x\mid e_n)e_1)
   =(e_n\mid x)
\]
for every \(x\in\mathfrak H\), and hence \(U_n^*e_1=e_n\), so that
\(U_n^*\) does not tend strongly to zero.

The unit ball \(\mathcal L_1(\mathfrak H)\), endowed with the uniform
structure associated with the strong topology, is a complete space
(\TreatiseRef{3}, Chapter~III, \S~3, Theorem~4). If \(\mathfrak H\)
has a countable basis, this space is, in addition, metrizable with a
countable base (\TreatiseRef{3}, Chapter~III, \S~3, Proposition~6), so
that every von Neumann algebra can then be generated by a countable
family of elements.

% Source PDF page 038 (printed p. 31)
\phantomsection\label{srcpage:31}
\begin{proposition}\label{prop:I-3-1}
Let \(\mathcal A\) be a von Neumann algebra of countable type, and let
\(\mathcal A_1\) be its unit ball endowed with the strong topology.
Then \(\mathcal A_1\) is metrizable.
\end{proposition}

\begin{proof}
Let \(\mathfrak M=\{x_1,x_2,\ldots\}\subset\mathfrak H\) be a
countable separating set for \(\mathcal A\). Let \(\mathfrak N\) be
the set of elements of the form
\[
  T_1'x_1+T_2'x_2+\cdots+T_n'x_n,
\]
where the \(T_i'\) belong to \(\mathcal A'\). Then \(\mathfrak N\) is
dense in \(\mathfrak H\). Since \(\mathcal A_1\) is equicontinuous,
the strong topology on \(\mathcal A_1\) is also the topology of simple
convergence on \(\mathfrak N\); since
\[
  T(T_1'x_1+T_2'x_2+\cdots+T_n'x_n)
   =\sum_{i=1}^nT_i'Tx_i
  \qquad(T\in\mathcal A),
\]
this topology is also that of simple convergence on \(\mathfrak M\).
Consequently the mapping
\[
  T\longmapsto(Tx_1,Tx_2,\ldots)
\]
is a homeomorphism of \(\mathcal A_1\) onto a subset of the product
space \(\mathfrak H\times\mathfrak H\times\cdots\), which is
metrizable.
\end{proof}

\begin{corollary*}
Let \(\mathcal A\) be a von Neumann algebra of countable type,
\(\mathfrak M\) a bounded subset of \(\mathcal A\), and let \(T\) belong
to the strong closure of \(\mathfrak M\). There exists a sequence of
elements of \(\mathfrak M\) which converges strongly to \(T\).
\end{corollary*}

\emph{Weak topology.} --- Let \(x,y\in\mathfrak H\). The function
\(T\mapsto\lvert(Tx\mid y)\rvert\) is a seminorm on
\(\mathcal L(\mathfrak H)\). The set of these seminorms defines on
\(\mathcal L(\mathfrak H)\) a separated locally convex topology called
the topology of simple weak convergence, or, for short, the weak
topology. The seminorms
\(T\mapsto\lvert(Tx\mid x)\rvert\) suffice to define this topology,
by virtue of the identity
\begin{align*}
  4(Tx\mid y)
   ={}&(T(x+y)\mid x+y)-(T(x-y)\mid x-y)\\
     &+i(T(x+iy)\mid x+iy)-i(T(x-iy)\mid x-iy).
\end{align*}
A fundamental system of neighborhoods of zero for this topology is
obtained by considering, for every finite system
\((x_1,y_1),(x_2,y_2),\ldots,(x_n,y_n)\) of pairs of points of
\(\mathfrak H\), the set of \(T\in\mathcal L(\mathfrak H)\) such that
\[
  \lvert(Tx_i\mid y_i)\rvert\leq1
  \qquad(1\leq i\leq n).
\]
The weak topology may also be defined as the coarsest topology on
\(\mathcal L(\mathfrak H)\) making the mappings
\(T\mapsto(Tx\mid y)\) from \(\mathcal L(\mathfrak H)\) into
\(\mathbb C\) continuous, or, equivalently, the mappings
\(T\mapsto Tx\) from \(\mathcal L(\mathfrak H)\) into
\(\mathfrak H\), endowed with its weak topology, continuous.

This topology is compatible with the vector-space structure of
\(\mathcal L(\mathfrak H)\), but not, in general, with its algebra
structure (Exercise~2).

% Source PDF page 039 (printed p. 32)
\phantomsection\label{srcpage:32}
Nevertheless, the mappings \(S\mapsto ST\) and \(T\mapsto ST\) are
weakly continuous. The mapping \(T\mapsto T^*\) is weakly continuous,
by virtue of the equality
\[
  \lvert(Tx\mid y)\rvert=\lvert(T^*y\mid x)\rvert.
\]
The unit ball \(\mathcal L_1(\mathfrak H)\), endowed with the weak
topology, is a compact space (\TreatiseRef{3}, Chapter~IV, \S~2,
Theorem~1, Corollary~3). If, in addition, \(\mathfrak H\) has a
countable basis, let \((x_1,x_2,\ldots)\) be a dense sequence in
\(\mathfrak H\). Then the weak topology on
\(\mathcal L_1(\mathfrak H)\) is defined by the sequence of deviations
\[
  (S,T)\longmapsto((S-T)x_i\mid x_j),
\]
and is therefore metrizable with a countable base.

\emph{Ultra-strong topology.} --- Let \((x_1,x_2,\ldots)\) be a
sequence of elements of \(\mathfrak H\) such that
\[
  \sum_{i=1}^{\infty}\lVert x_i\rVert^2<+\infty.
\]
For every \(T\in\mathcal L(\mathfrak H)\), we have
\[
  \sum_{i=1}^{\infty}\lVert Tx_i\rVert^2
   \leq\lVert T\rVert^2
       \sum_{i=1}^{\infty}\lVert x_i\rVert^2
   <+\infty.
\]
The function
\[
  T\longmapsto
  \left[\sum_{i=1}^{\infty}\lVert Tx_i\rVert^2\right]^{1/2}
\]
is a seminorm on \(\mathcal L(\mathfrak H)\). The set of these
seminorms defines on \(\mathcal L(\mathfrak H)\) a separated locally
convex topology called the ultra-strong topology. A fundamental system
of neighborhoods of zero for this topology is obtained by considering,
for every sequence \((x_i)\) of elements of \(\mathfrak H\) such that
\(\sum_{i=1}^{\infty}\lVert x_i\rVert^2<+\infty\), the set of
\(T\in\mathcal L(\mathfrak H)\) such that
\[
  \sum_{i=1}^{\infty}\lVert Tx_i\rVert^2\leq1.
\]
This topology is also the coarsest of the topologies making the mappings
\[
  T\longmapsto(Tx_1,Tx_2,\ldots)
\]
from \(\mathcal L(\mathfrak H)\) into the Hilbert direct sum
\(\mathfrak H\oplus\mathfrak H\oplus\cdots\) continuous, where
\((x_i)\) is any sequence of elements of \(\mathfrak H\) such that
\(\sum_i\lVert x_i\rVert^2<+\infty\).

This topology, which is compatible with the vector-space structure of
\(\mathcal L(\mathfrak H)\), is not compatible, in general, with its
algebra structure (Exercise~2). Nevertheless, the mapping
\((S,T)\mapsto ST\) from
\(\mathcal L_1(\mathfrak H)\times\mathcal L(\mathfrak H)\) into
\(\mathcal L(\mathfrak H)\), and the mappings \(S\mapsto ST\) and
\(T\mapsto ST\), are ultra-strongly continuous. The mapping
\(T\mapsto T^*\) is not, in general, ultra-strongly continuous
(cf. \hyperref[no:I-3-2]{no.~2}).

\emph{Ultra-weak topology.} --- Let \((x_1,x_2,\ldots)\) and
\((y_1,y_2,\ldots)\) be two sequences of elements of \(\mathfrak H\)
such that
\[
  \sum_{i=1}^{\infty}\lVert x_i\rVert^2<+\infty,
  \qquad
  \sum_{i=1}^{\infty}\lVert y_i\rVert^2<+\infty.
\]
% Source PDF page 040 (printed p. 33)
\phantomsection\label{srcpage:33}
For every \(T\in\mathcal L(\mathfrak H)\), we have
\begin{align*}
  \sum_{i=1}^{\infty}\lvert(Tx_i\mid y_i)\rvert
  &\leq\sum_{i=1}^{\infty}\lVert Tx_i\rVert\,\lVert y_i\rVert\\
  &\leq\lVert T\rVert
       \sum_{i=1}^{\infty}\lVert x_i\rVert\,\lVert y_i\rVert\\
  &\leq\lVert T\rVert
       \left[\sum_{i=1}^{\infty}\lVert x_i\rVert^2\right]^{1/2}
       \left[\sum_{i=1}^{\infty}\lVert y_i\rVert^2\right]^{1/2}
   <+\infty.
\end{align*}
The function
\[
  T\longmapsto
  \left\lvert\sum_{i=1}^{\infty}(Tx_i\mid y_i)\right\rvert
\]
is a seminorm on \(\mathcal L(\mathfrak H)\). The set of these
seminorms defines on \(\mathcal L(\mathfrak H)\) a separated locally
convex topology called the ultra-weak topology. The seminorms
\[
  T\longmapsto
  \left\lvert\sum_{i=1}^{\infty}(Tx_i\mid x_i)\right\rvert
\]
suffice to define it. A fundamental system of neighborhoods of zero for
this topology is obtained by considering, for every finite system
\[
  (x_i^1),(y_i^1),(x_i^2),(y_i^2),\ldots,(x_i^n),(y_i^n)
\]
of sequences of elements of \(\mathfrak H\) such that
\[
  \sum_{i=1}^{\infty}\lVert x_i^k\rVert^2<+\infty,
  \qquad
  \sum_{i=1}^{\infty}\lVert y_i^k\rVert^2<+\infty
  \qquad(1\leq k\leq n),
\]
the set of \(T\in\mathcal L(\mathfrak H)\) such that
\[
  \left\lvert\sum_{i=1}^{\infty}(Tx_i^k\mid y_i^k)\right\rvert
  \leq1
  \qquad(1\leq k\leq n).
\]
This topology is also the coarsest of the topologies making continuous
the linear functionals
\[
  T\longmapsto\sum_{i=1}^{\infty}(Tx_i\mid y_i),
\]
where \((x_i)\) and \((y_i)\) are any two sequences of elements of
\(\mathfrak H\) such that
\[
  \sum_{i=1}^{\infty}\lVert x_i\rVert^2<+\infty,
  \qquad
  \sum_{i=1}^{\infty}\lVert y_i\rVert^2<+\infty.
\]
This topology is not compatible, in general, with the algebra structure
of \(\mathcal L(\mathfrak H)\) (Exercise~2). The mappings
\(S\mapsto ST\), \(T\mapsto ST\), and \(T\mapsto T^*\) are
ultra-weakly continuous.\footnote{The ultra-strong and ultra-weak
topologies are of interest because they are invariant under isomorphism
(\hyperref[sec:I-4]{\S~4}, \hyperref[thm:I-4-2]{Theorem~2},
\hyperref[cor:I-4-1]{Corollary~1}). An English term for
``ultra-strong'' is ``strongest.''}

\noindent\textbf{Bibliography:} \ArticleRef{15}, \ArticleRef{73},
\ArticleRef{74}, \ArticleRef{234}.

% Source PDF page 041 (printed p. 34)
\phantomsection\label{srcpage:34}
\NumberedParagraph{2}{Comparison of the preceding topologies}
\label{no:I-3-2}
We have at once the following diagram, in which the sign \(<\) means
``finer than'':
\[
\begin{array}{ccccc}
  \text{Norm topology} & < & \text{Ultra-strong topology}
    & < & \text{Strong topology}\\
  && \uparrow && \uparrow\\[-0.2ex]
  && \text{Ultra-weak topology} & < & \text{Weak topology}.
\end{array}
\]
When \(\mathfrak H\) is infinite-dimensional, the sign \(<\) may even
be interpreted as meaning ``strictly finer than''
(\hyperref[ex:I-3-1]{Exercise~1}, and
\hyperref[prop:III-6-7]{Chapter~III, \S~6, Propositions~7}
and~\hyperref[prop:III-6-8]{8}).

The strong and ultra-strong topologies coincide on bounded subsets
of \(\mathcal L(\mathfrak H)\). Indeed, let \(\mathcal U\) be an
ultra-strong neighborhood of zero. There exists a sequence \((x_i)\) of
elements of \(\mathfrak H\) such that
\[
  \sum_{i=1}^{\infty}\lVert x_i\rVert^2<+\infty,
\]
and such that the inequality
\[
  \sum_{i=1}^{\infty}\lVert Tx_i\rVert^2\leq1
\]
implies \(T\in\mathcal U\). Let \(n\) be an integer such that
\[
  \sum_{i=n}^{\infty}\lVert x_i\rVert^2\leq\frac12.
\]
Let \(\mathcal U_1\) be the strong neighborhood of zero defined by the
inequality
\[
  \sum_{i=1}^{n-1}\lVert Tx_i\rVert^2\leq\frac12.
\]
Then the conditions \(\lVert T\rVert\leq1\) and
\(T\in\mathcal U_1\) imply
\[
  \sum_{i=1}^{\infty}\lVert Tx_i\rVert^2
  \leq
  \sum_{i=1}^{n-1}\lVert Tx_i\rVert^2
  +\lVert T\rVert^2\sum_{i=n}^{\infty}\lVert x_i\rVert^2
  \leq\frac12+\frac12=1,
\]
and hence \(T\in\mathcal U\), which proves our assertion. One sees in
the same way that the weak and ultra-weak topologies coincide on the
bounded subsets of \(\mathcal L(\mathfrak H)\).

For \(T\) to tend strongly to zero, it is necessary and sufficient that
\(T^*T\) tend weakly to zero, as follows from the equality
\[
  \lVert Tx\rVert^2=(T^*Tx\mid x).
\]
Similarly, for \(T\) to tend ultra-strongly to zero, it is necessary and
sufficient that \(T^*T\) tend ultra-weakly to zero.

\noindent\textbf{Bibliography:} \ArticleRef{15}, \ArticleRef{73},
\ArticleRef{74}.

\NumberedParagraph{3}{Linear functionals on \(\mathcal L(\mathfrak H)\)}
\label{no:I-3-3}
The space \(\mathcal L(\mathfrak H)=\mathcal L\) is a complex Banach
space. The set of linear functionals on \(\mathcal L\) that are
continuous for the norm topology constitutes the dual
\(\mathcal L^*\) of \(\mathcal L\), which is itself a complex Banach
space. The linear functionals on \(\mathcal L\) continuous for one of
the strong, weak, ultra-strong, or ultra-weak topologies constitute
subspaces of \(\mathcal L^*\) which we shall now study.

For \(x,y\in\mathfrak H\), we shall denote by \(\omega_{x,y}\) the
weakly continuous linear functional on \(\mathcal L\) defined by
\[
  \omega_{x,y}(T)=(Tx\mid y).
\]
The set \(\mathcal L_{\sim}\) of weakly continuous linear functionals
on \(\mathcal L\) coincides with the set of strongly continuous linear
% Source PDF page 042 (printed p. 35)
\phantomsection\label{srcpage:35}
functionals, and with the set of finite sums of functionals
\(\omega_{x,y}\) (\TreatiseRef{3}, Chapter~IV, \S~2,
Proposition~11). This set is a vector subspace of \(\mathcal L^*\), not
in general closed for the norm of \(\mathcal L^*\)
(\hyperref[prop:III-6-7]{Chapter~III, \S~6, Propositions~7}
and~\hyperref[prop:III-6-8]{8}); let \(\mathcal L_*\) be its closure in
\(\mathcal L^*\) in the sense of the norm.

In what follows, it will be convenient to use the following
terminology. Let \(E\) and \(F\) be two complex normed vector spaces,
\(E^*\) and \(F^*\) their topological duals, and let
\((x,y)\mapsto\beta(x,y)\) be a continuous bilinear form on
\(E\times F\). Every \(y\in F\) defines the linear functional
\(f_y:x\mapsto\beta(x,y)\) on \(E\), and the mapping
\(y\mapsto f_y\) from \(F\) into \(E^*\) is linear and continuous. If
this mapping is bijective and isometric, we shall say that \(F\) is the
dual of \(E\) for the form \(\beta\). If this is so, then
\[
  \lVert y\rVert
   =\sup_{\lVert x\rVert\leq1}\lvert\beta(x,y)\rvert
  \quad\text{for every }y\in F,
\]
and
\[
  \lVert x\rVert
   =\sup_{\lVert y\rVert\leq1}\lvert\beta(x,y)\rvert
  \quad\text{for every }x\in E.
\]

\begin{lemma}\label{lem:I-3-1}
The canonical bilinear form on
\(\mathcal L^*\times\mathcal L\) induces on
\(\mathcal L_{\sim}\times\mathcal L\) a bilinear form for which
\(\mathcal L\) is the dual of the normed vector space
\(\mathcal L_{\sim}\).
\end{lemma}

\begin{proof}
Let \(\mathcal E\) be the dual of the normed vector space
\(\mathcal L_{\sim}\). Every element of \(\mathcal L\) defines a continuous
linear functional on \(\mathcal L_{\sim}\), and hence an element of
\(\mathcal E\). Let \(\Theta\) be the resulting linear, and plainly
injective, mapping from \(\mathcal L\) into \(\mathcal E\). This
mapping is continuous when \(\mathcal E\) is endowed with the weak
dual topology \(\sigma(\mathcal E,\mathcal L_{\sim})\), and \(\mathcal L\)
with the topology \(\sigma(\mathcal L,\mathcal L_{\sim})\), which is nothing
other than the weak operator topology. Let \(\mathcal L_1\) and
\(\mathcal E_1\) be the unit balls of \(\mathcal L\) and
\(\mathcal E\). The ball \(\mathcal L_1\) is weakly compact, and hence
\(\Theta(\mathcal L_1)\) is weakly compact. Finally,
\(\Theta(\mathcal L_1)\) and \(\mathcal E_1\) have the same polar set
in \(\mathcal L_{\sim}\), namely the unit ball of
\(\mathcal L_{\sim}\). Thus
\(\Theta(\mathcal L_1)=\mathcal E_1\), which proves that \(\Theta\) is
an isometric mapping of \(\mathcal L\) onto \(\mathcal E\).
\end{proof}

\Cref{lem:I-3-1} also allows \(\mathcal L\) to be identified with the
dual of the complex Banach space \(\mathcal L_*\).

\begin{lemma}\label{lem:I-3-2}
The ultra-strongly continuous linear functionals on \(\mathcal L\)
coincide with the ultra-weakly continuous linear functionals, that is,
with the linear functionals
\[
  \sum_{i=1}^{\infty}\omega_{x_i,y_i},
  \qquad
  \sum_{i=1}^{\infty}\lVert x_i\rVert^2<+\infty,
  \quad
  \sum_{i=1}^{\infty}\lVert y_i\rVert^2<+\infty.
\]
They belong to \(\mathcal L_*\).
\end{lemma}

% Source PDF page 043 (printed p. 36)
\phantomsection\label{srcpage:36}
\begin{proof}
Let
\[
  \varphi=\sum_{i=1}^{\infty}\omega_{x_i,y_i},
  \qquad
  \sum_{i=1}^{\infty}\lVert x_i\rVert^2<+\infty,
  \quad
  \sum_{i=1}^{\infty}\lVert y_i\rVert^2<+\infty.
\]
It is clear that \(\varphi\) is the uniform limit on the ball
\(\mathcal L_1\) of the functionals
\(\sum_{i=1}^n\omega_{x_i,y_i}\), and hence
\(\varphi\in\mathcal L_*\). The functional \(\varphi\) is ultra-weakly
continuous, and every ultra-weakly continuous functional is
ultra-strongly continuous. Finally, let \(\varphi\) be an
ultra-strongly continuous linear functional. Consider the Hilbert space
\[
  \widetilde{\mathfrak H}
   =\bigoplus_{i=1}^{\infty}\mathfrak H_i,
  \qquad
  \mathfrak H=\mathfrak H_1=\mathfrak H_2=\cdots.
\]
We have
\[
  \lvert\varphi(T)\rvert
   \leq
  \left[\sum_{i=1}^{\infty}\lVert Tx_i\rVert^2\right]^{1/2},
\]
where \((x_i)\) is a sequence of vectors of \(\mathfrak H\) such that
\[
  \sum_{i=1}^{\infty}\lVert x_i\rVert^2<+\infty.
\]
The sequence \((x_i)\) is an element \(\widetilde x\) of
\(\widetilde{\mathfrak H}\). Let \(\widetilde T\) be the operator on
\(\widetilde{\mathfrak H}\) defined by
\(\widetilde T((z_i))=(Tz_i)\). We have
\(\lvert\varphi(T)\rvert\leq\lVert\widetilde T\widetilde x\rVert\),
and therefore there exists an element
\(\widetilde y=(y_i)\) of \(\widetilde{\mathfrak H}\) such that
\[
  \varphi(T)
   =(\widetilde T\widetilde x\mid\widetilde y)
   =\sum_{i=1}^{\infty}(Tx_i\mid y_i).
\]
\end{proof}

\begin{lemma}\label{lem:I-3-3}
The linear functionals in \(\mathcal L_*\) are ultra-weakly
continuous.
\end{lemma}

\begin{proof}
We first establish the following result.

For every linear functional \(\varphi\in\mathcal L_{\sim}\), there
exist two orthonormal systems
\[
  (e_1,e_2,\ldots,e_n),
  \qquad
  (e_1',e_2',\ldots,e_n')
\]
in \(\mathfrak H\), and numbers
\(\lambda_i\geq0\) \((1\leq i\leq n)\), such that
\[
  \varphi=\sum_{i=1}^n\lambda_i\omega_{e_i,e_i'},
  \qquad
  \lVert\varphi\rVert=\sum_{i=1}^n\lambda_i.
\]

We have \(\varphi=\sum_{j=1}^p\omega_{x_j,y_j}\). Let
\(x_1',x_2',\ldots,x_q'\), \(y_1',y_2',\ldots,y_q'\) be other elements
of \(\mathfrak H\) such that the finite-rank linear mappings
\[
  y\longmapsto\sum_{j=1}^p(y\mid y_j)x_j
  \quad\text{and}\quad
  y\longmapsto\sum_{k=1}^q(y\mid y_k')x_k'
\]
from \(\mathfrak H\) into \(\mathfrak H\) are identical.
% Source PDF page 044 (printed p. 37)
\phantomsection\label{srcpage:37}
For every pair of vectors \((z_1,z_2)\) of \(\mathfrak H\), denote by
\(\{z_1,z_2\}\) the linear mapping of rank at most one which sends
\(z\) to \((z\mid z_1)z_2\). We have
\begin{align*}
  \sum_{j=1}^p(\{z_1,z_2\}x_j\mid y_j)
   &=\sum_{j=1}^p(x_j\mid z_1)(z_2\mid y_j)
    =\left(\sum_{j=1}^p(z_2\mid y_j)x_j\mathrel{\big|}z_1\right),\\
  \sum_{k=1}^q(\{z_1,z_2\}x_k'\mid y_k')
   &=\left(\sum_{k=1}^q(z_2\mid y_k')x_k'\mathrel{\big|}z_1\right),
\end{align*}
and hence, by linearity and continuity,
\[
  \sum_{j=1}^p(Tx_j\mid y_j)
   =\sum_{k=1}^q(Tx_k'\mid y_k')
  \qquad\text{for every }T\in\mathcal L,
\]
that is,
\[
  \sum_{j=1}^p\omega_{x_j,y_j}
   =\sum_{k=1}^q\omega_{x_k',y_k'}.
\]
Perform the polar decomposition of the operator
\[
  y\longmapsto\sum_{j=1}^p(y\mid y_j)x_j,
\]
and then apply the spectral theory to the finite-rank Hermitian operator
thus obtained. We see that the \((x_k')\) and the \((y_k')\) can be
chosen to be two orthogonal systems. It follows at once that there exist
orthonormal systems \((e_i)\), \((e_i')\) and scalars
\(\lambda_i\geq0\) such that
\[
  \varphi=\sum_{i=1}^n\lambda_i\omega_{e_i,e_i'}.
\]
We then have
\[
  \lvert\varphi(T)\rvert
   \leq\sum_{i=1}^n\lambda_i\lVert T\rVert
          \lVert e_i\rVert\,\lVert e_i'\rVert
   =\lVert T\rVert\sum_{i=1}^n\lambda_i,
\]
whence \(\lVert\varphi\rVert\leq\sum_{i=1}^n\lambda_i\). On the other
hand, for
\[
  T=\sum_{i=1}^n\{e_i,e_i'\},
\]
we have
\[
  \lVert T\rVert\leq1,
  \qquad
  \varphi(T)=\sum_{i=1}^n\lambda_i,
\]
and hence \(\lVert\varphi\rVert\geq\sum_{i=1}^n\lambda_i\), proving
our intermediate result.

This being established, let \(\varphi\in\mathcal L_*\). We have
\[
  \varphi=\sum_{k=1}^{\infty}\varphi_k,
  \qquad
  \varphi_k\in\mathcal L_{\sim},
  \qquad
  \lVert\varphi_k\rVert\leq2^{-k}\quad(k\geq2),
\]
and therefore
\[
  \varphi_k=\sum_{i=1}^{n_k}\lambda_i^k
             \omega_{e_i^k,e_i^{\prime k}},
\]
where
\[
  \lVert e_i^k\rVert=\lVert e_i^{\prime k}\rVert=1,
  \qquad
  \lambda_i^k\geq0,
  \qquad
  \sum_{i=1}^{n_k}\lambda_i^k\leq2^{-k}.
\]
% Source PDF page 045 (printed p. 38)
\phantomsection\label{srcpage:38}
Hence
\[
  \varphi
   =\sum_{k=1}^{\infty}\left(
      \sum_{i=1}^{n_k}
       \omega_{(\lambda_i^k)^{1/2}e_i^k,
               (\lambda_i^k)^{1/2}e_i^{\prime k}}
     \right).
\]
Finally,
\begin{align*}
  \sum_{k=1}^{\infty}\left(
    \sum_{i=1}^{n_k}
      \bigl\lVert(\lambda_i^k)^{1/2}e_i^k\bigr\rVert^2
  \right)
  &=\sum_{k=1}^{\infty}\left(
      \sum_{i=1}^{n_k}\lambda_i^k\right)
    \leq\sum_{k=1}^{\infty}2^{-k}<+\infty,
\end{align*}
and, similarly,
\[
  \sum_{k=1}^{\infty}\left(
    \sum_{i=1}^{n_k}
      \bigl\lVert(\lambda_i^k)^{1/2}e_i^{\prime k}\bigr\rVert^2
  \right)<+\infty,
\]
so that \(\varphi\) is ultra-weakly continuous.
\end{proof}

\begin{theorem}\label{thm:I-3-1}
Let \(\mathfrak M\) be an ultra-weakly closed vector subspace of
\(\mathcal L(\mathfrak H)\), let \(\mathfrak M^*\) be the dual of the
Banach space \(\mathfrak M\), let \(\mathfrak M_r\) be the ball
\(\lVert T\rVert\leq r\) of \(\mathfrak M\), and let \(\varphi\) be a
linear functional on \(\mathfrak M\).

\noindent\textup{(i)} The following conditions are equivalent:

\noindent\textup{(i 1)}\quad \(\varphi\) is weakly continuous;

\noindent\textup{(i 2)}\quad \(\varphi\) is strongly continuous;

\noindent\textup{(i 3)}\quad
\(\displaystyle\varphi=\sum_{i=1}^n\omega_{x_i,y_i}.\)

\noindent\textup{(ii)} The following conditions are equivalent:

\noindent\textup{(ii 1)}\quad \(\varphi\) is ultra-weakly continuous;

\noindent\textup{(ii 2)}\quad \(\varphi\) is ultra-strongly continuous;

\noindent\textup{(ii 3)}\quad
\(\displaystyle\varphi=\sum_{i=1}^{\infty}\omega_{x_i,y_i}\), with
\[
  \sum_{i=1}^{\infty}\lVert x_i\rVert^2<+\infty,
  \qquad
  \sum_{i=1}^{\infty}\lVert y_i\rVert^2<+\infty;
\]

\noindent\textup{(ii 4) [respectively, (ii 5)].}\quad
The restriction of \(\varphi\) to \(\mathfrak M_1\) is ultra-weakly
(respectively, weakly) continuous;

\noindent\textup{(ii 6) [respectively, (ii 7)].}\quad
The restriction of \(\varphi\) to \(\mathfrak M_1\) is ultra-strongly
(respectively, strongly) continuous.

\noindent\textup{(iii)} Let \(\mathfrak M_{\sim}\) (respectively,
\(\mathfrak M_*\)) be the set of weakly (respectively, ultra-weakly)
continuous linear functionals on \(\mathfrak M\). Then
\(\mathfrak M_*\) is the closure of \(\mathfrak M_{\sim}\) for the
norm of \(\mathfrak M^*\), and \(\mathfrak M\) is the dual of
\(\mathfrak M_*\) for the bilinear form induced on
\(\mathfrak M_*\times\mathfrak M\) by the canonical bilinear form on
\(\mathfrak M^*\times\mathfrak M\).

\noindent\textup{(iv)} Let \(\mathfrak X\) be a convex subset of
\(\mathfrak M\). The following conditions are equivalent:

\noindent\textup{(iv 1)}\quad \(\mathfrak X\) is ultra-weakly closed;

\noindent\textup{(iv 2)}\quad \(\mathfrak X\) is ultra-strongly closed;

\noindent\textup{(iv 3) [respectively, (iv 4)].}\quad
\(\mathfrak X\cap\mathfrak M_r\) is ultra-weakly (respectively,
weakly) closed for every \(r\);
% Source PDF page 046 (printed p. 39)
\phantomsection\label{srcpage:39}
\noindent\textup{(iv 5) [respectively, (iv 6)].}\quad
\(\mathfrak X\cap\mathfrak M_r\) is ultra-strongly (respectively,
strongly) closed for every \(r\).
\end{theorem}

\begin{proof}
The implications
\[
  \text{(i 3)}\Longrightarrow\text{(i 1)}
  \Longrightarrow\text{(i 2)}
\]
are clear. On the other hand, if \(\varphi\) is strongly continuous on
\(\mathfrak M\), then \(\varphi\) extends to a strongly continuous
linear functional on \(\mathcal L(\mathfrak H)\), and hence
\[
  \varphi=\sum_{i=1}^{n}\omega_{x_i,y_i},
\]
which proves that (i 2) implies (i 3). This proves part (i) of the
theorem. Similarly, the use of \cref{lem:I-3-2} proves the equivalence
of (ii 1), (ii 2), and (ii 3). The equivalence
\[
  \text{(ii 4)}\Longleftrightarrow\text{(ii 5)}
  \quad
  [\text{respectively, }\,
  \text{(ii 6)}\Longleftrightarrow\text{(ii 7)}]
\]
follows from the identity of the topologies induced on
\(\mathfrak M_1\) by the weak and ultra-weak topologies (respectively,
by the strong and ultra-strong topologies). The implications
\[
  \text{(ii 3)}\Longrightarrow\text{(ii 4)}
  \Longrightarrow\text{(ii 6)}
\]
are clear. Thus, to prove part (ii) of the theorem, only the implication
(ii 6) \(\Longrightarrow\) (ii 1) remains.

On the other hand, \cref{lem:I-3-2} implies that
\(\mathfrak M_{\sim}\) is dense in \(\mathfrak M_*\) for the norm of
\(\mathfrak M^*\). Let \(\mathfrak M^\perp\) be the set of elements of
\(\mathcal L_*\) orthogonal to \(\mathfrak M\); this is a norm-closed
vector subspace of \(\mathcal L_*\). Since \(\mathfrak M\) is
ultra-weakly closed (a hypothesis that had not yet been used),
\(\mathfrak M\) is precisely the vector subspace of \(\mathcal L\)
orthogonal to \(\mathfrak M^\perp\). The Banach space \(\mathcal L\) is
the dual of the Banach space \(\mathcal L_*\) for the bilinear form
\[
  (\psi,T)\longmapsto\psi(T)
\]
on \(\mathcal L_*\times\mathcal L\). If \(\psi\) belongs to the Banach
space \(\mathcal L_*/\mathfrak M^\perp\), and if \(T\in\mathfrak M\),
denote by \(\langle\psi,T\rangle\) the value at \(T\) of any
representative of \(\psi\) in \(\mathcal L_*\). By duality theory, the
Banach space \(\mathfrak M\) is the dual of the Banach space
\(\mathcal L_*/\mathfrak M^\perp\) for the form
\[
  (\psi,T)\longmapsto\langle\psi,T\rangle.
\]
On the other hand, the restriction mapping
\[
  \mathcal L_*\longrightarrow\mathfrak M_*
\]
has kernel \(\mathfrak M^\perp\); its image is the set of functionals
on \(\mathfrak M\) which can be extended to ultra-weakly continuous
functionals on \(\mathcal L\) (\cref{lem:I-3-2,lem:I-3-3}), that is,
\(\mathfrak M_*\). We therefore have a canonical bijection
\[
  \theta:\mathcal L_*/\mathfrak M^\perp\longrightarrow\mathfrak M_*.
\]
If \(\omega\in\mathfrak M_*\) and
\(\psi=\theta^{-1}(\omega)\), then
\(\langle\psi,T\rangle=\omega(T)\) for every \(T\in\mathfrak M\), and
therefore
\[
  \lVert\psi\rVert
    =\sup_{\substack{T\in\mathfrak M\\ \lVert T\rVert\leq1}}
      \lvert\langle\psi,T\rangle\rvert
    =\sup_{\substack{T\in\mathfrak M\\ \lVert T\rVert\leq1}}
      \lvert\omega(T)\rvert
    =\lVert\omega\rVert.
\]
Thus \(\mathfrak M_*\) is complete for the norm, and
\(\mathfrak M\) is the dual of \(\mathfrak M_*\) for the bilinear form
\[
  (\omega,T)\longmapsto\omega(T).
\]
This proves (iii).

The equivalences
\[
  \text{(iv 1)}\Longleftrightarrow\text{(iv 2)},
  \qquad
  \text{(iv 3)}\Longleftrightarrow\text{(iv 5)}
\]
follow from the equivalence (ii 1) \(\Longleftrightarrow\) (ii 2). The
equivalence
\[
  \text{(iv 3)}\Longleftrightarrow\text{(iv 4)}
  \quad
  [\text{respectively, }\,
  \text{(iv 5)}\Longleftrightarrow\text{(iv 6)}]
\]
follows from the identity of the topologies induced on
\(\mathfrak M_r\) by the weak and ultra-weak topologies (respectively,
by the strong and ultra-strong topologies). The equivalence
\[
  \text{(iv 1)}\Longleftrightarrow\text{(iv 3)}
\]
follows from (iii) and a property of Banach spaces
(\TreatiseRef{3}, Chapter~IV, \S~2, Theorem~5). Thus (iv) is proved.

We can now finally establish the implication
(ii 6) \(\Longrightarrow\) (ii 1). Let \(\varphi\) be a linear
functional on \(\mathfrak M\) whose restriction to
\(\mathfrak M_1\) is ultra-strongly continuous, and let
\(\mathfrak K\subset\mathfrak M\) be the hyperplane with equation
\(\varphi(T)=0\). The set
\(\mathfrak K\cap\mathfrak M_1\) is ultra-strongly closed; hence,
% Source PDF page 047 (printed p. 40)
\phantomsection\label{srcpage:40}
by (iv), \(\mathfrak K\) is ultra-weakly closed, so that \(\varphi\) is
ultra-weakly continuous.
\end{proof}

The space \(\mathfrak M_*\) is called the \emph{predual} of
\(\mathfrak M\). Thus \(\mathfrak M\) is canonically identified with
the Banach-space dual of its predual. The ultra-weak topology on
\(\mathfrak M\) is the weakest topology for which the linear
functionals belonging to \(\mathfrak M_*\) are continuous.

\noindent\textbf{Bibliography:} \ArticleRef{7}, \ArticleRef{15},
\ArticleRef{67}, \ArticleRef{68}, \ArticleRef{73}, \ArticleRef{74}.

\NumberedParagraph{4}{von Neumann's density theorem}
\label{no:I-3-4}

\begin{lemma}\label{lem:I-3-4}
Let \(\mathcal A\) be an involutive algebra of operators on
\(\mathfrak H\). Let \(\mathfrak X\) be the closed vector subspace of
\(\mathfrak H\) spanned by the \(Tx\), where
\(T\in\mathcal A\) and \(x\in\mathfrak H\). Let
\(\mathfrak X'\) be the closed vector subspace of \(\mathfrak H\)
formed by the \(x\in\mathfrak H\) such that \(Tx=0\) for every
\(T\in\mathcal A\). Then \(\mathfrak X\) and \(\mathfrak X'\) are
complementary orthogonal subspaces of \(\mathfrak H\), and, for every
\(T\in\mathcal A\),
\[
  TP_{\mathfrak X}=P_{\mathfrak X}T=T.
\]
\end{lemma}

\begin{proof}
For every \(T\in\mathcal L(\mathfrak H)\), denote by
\(\mathfrak N_T\) the set of \(x\in\mathfrak H\) such that \(Tx=0\).
We know that the closure of \(T(\mathfrak H)\) is the orthogonal
complement of \(\mathfrak N_{T^*}\). Thus the closed vector subspace
spanned by the \(T(\mathfrak H)\), \(T\in\mathcal A\), has as its
orthogonal complement the intersection of the
\(\mathfrak N_{T^*}\), \(T\in\mathcal A\). On the other hand, for
\(T\in\mathcal A\), we have \(TP_{\mathfrak X'}=0\), and therefore
\[
  T=T(I_{\mathfrak H}-P_{\mathfrak X'})
    =TP_{\mathfrak X}.
\]
Likewise \(T^*=T^*P_{\mathfrak X}\), whence, on taking adjoints,
\[
  T=P_{\mathfrak X}T.
\]
\end{proof}

\begin{lemma}\label{lem:I-3-5}
Let \(\mathcal A\) be an involutive algebra of operators on
\(\mathfrak H\) such that
\[
  \mathfrak X_{\mathfrak H}^{\mathcal A}=\mathfrak H.
\]
Then, for every \(y\in\mathfrak H\), one has
\[
  y\in\mathfrak X_y^{\mathcal A}.
\]
\end{lemma}

\begin{proof}
Put
\[
  \mathfrak X=\mathfrak X_y^{\mathcal A},
  \qquad
  y'=P_{\mathfrak X}y,
  \qquad
  y''=y-y'.
\]
For every \(T\in\mathcal A\), we have \(Ty\in\mathfrak X\) by the
definition of \(\mathfrak X\), and \(Ty'\in\mathfrak X\) because
\(\mathfrak X\) is stable under \(\mathcal A\). Hence
\(Ty''\in\mathfrak X\), and therefore
\[
  (T^*Ty''\mid y'')=0,
\]
that is, \(Ty''=0\). Thus \(y''=0\) by \cref{lem:I-3-4}, which means
that \(y\in\mathfrak X\).
\end{proof}

The following lemma constitutes the essential step.

\begin{lemma}\label{lem:I-3-6}
Let \(\mathcal A\) be an involutive algebra of operators on
\(\mathfrak H\) such that
\[
  \mathfrak X_{\mathfrak H}^{\mathcal A}=\mathfrak H.
\]
Then every operator in \(\mathcal A''\) is in the ultra-strong closure
of \(\mathcal A\).
\end{lemma}

\begin{proof}
Let \(S\in\mathcal A''\) and \(x\in\mathfrak H\). We have
\(E_x^{\mathcal A}\in\mathcal A'\), so that
\(\mathfrak X_x^{\mathcal A}\) is stable under \(S\), and
\(x\in\mathfrak X_x^{\mathcal A}\) by \cref{lem:I-3-5}. Hence
\[
  Sx\in\mathfrak X_x^{\mathcal A}.
\]

This being so, let \((x_i)_{i\in I}\) be a family of elements of
\(\mathfrak H\) such that
\[
  \sum_{i\in I}\lVert x_i\rVert^2<+\infty,
\]
and let us prove that there is a \(T\in\mathcal A\) such that
\[
  \sum_{i\in I}\lVert(S-T)x_i\rVert^2
\]
is arbitrarily small. For this purpose, let \(\mathfrak K\) be a
Hilbert space which is the Hilbert direct sum of a family
\((\mathfrak K_i)_{i\in I}\) of pairwise orthogonal closed vector
subspaces, and suppose there is a linear isometry
% Source PDF page 048 (printed p. 41)
\phantomsection\label{srcpage:41}
\(U_i\) of \(\mathfrak H\) onto \(\mathfrak K_i\). Put
\[
  y_i=U_ix_i.
\]
Since \(\sum_{i\in I}\lVert y_i\rVert^2<+\infty\), we may form
\[
  y=\sum_{i\in I}y_i\in\mathfrak K.
\]
For every \(R\in\mathcal L(\mathfrak H)\), denote by
\(\widetilde R\) the operator in \(\mathcal L(\mathfrak K)\)
represented by the matrix \((R_{i\kappa})\), where
\[
  R_{i\kappa}=0\quad(i\ne\kappa),
  \qquad
  R_{ii}=R\quad(i\in I)
\]
(\hyperref[no:I-2-3]{\S~2, no.~3}). As \(T\) ranges over
\(\mathcal A\), \(\widetilde T\) ranges over an involutive algebra of
operators \(\widetilde{\mathcal A}\) on \(\mathfrak K\). By
\cref{lem:I-2-2}, we have
\[
  \widetilde S\in\widetilde{\mathcal A}''.
\]
Since the \(\widetilde Tz\), \(T\in\mathcal A\),
\(z\in\mathfrak K\), plainly span \(\mathfrak K\), the beginning of
the proof shows that \(\widetilde Sy\) is a strong limit of elements
\(\widetilde Ty\), \(T\in\mathcal A\); that is,
\[
  \lVert\widetilde Sy-\widetilde Ty\rVert^2
    =\sum_{i\in I}\lVert(S-T)x_i\rVert^2
\]
can be made arbitrarily small.
\end{proof}

\begin{theorem}\label{thm:I-3-2}
Let \(\mathcal A\) be an involutive algebra of operators on
\(\mathfrak H\), and let \(\mathcal A_1\) be the unit ball of
\(\mathcal A\).

\begin{enumerate}[label=(\roman*)]
\item The following eight conditions are equivalent:
\begin{enumerate}[label=(i \arabic*),leftmargin=3.2em]
\item \(\mathcal A\) is weakly closed;
\item \(\mathcal A_1\) is weakly closed;
\item \(\mathcal A\) is strongly closed;
\item \(\mathcal A_1\) is strongly closed;
\item \(\mathcal A\) is ultra-weakly closed;
\item \(\mathcal A_1\) is ultra-weakly closed;
\item \(\mathcal A\) is ultra-strongly closed;
\item \(\mathcal A_1\) is ultra-strongly closed.
\end{enumerate}
\item Suppose that the conditions in (i) are satisfied. Put
\[
  \mathfrak X=\mathfrak X_{\mathfrak H}^{\mathcal A}.
\]
Then \(P_{\mathfrak X}\) is the greatest projection in
\(\mathcal A\). For every \(T\in\mathcal A\),
\[
  P_{\mathfrak X}T=TP_{\mathfrak X}=T.
\]
The operators in \(\mathcal A''\) are the operators
\[
  T+\lambda I_{\mathfrak H}
  \qquad(T\in\mathcal A,\ \lambda\in\mathbb C).
\]
\end{enumerate}
\end{theorem}

\begin{proof}
We already have the equivalences
\[
  \text{(i 1)}\Longleftrightarrow\text{(i 3)},\qquad
  \text{(i 2)}\Longleftrightarrow\text{(i 4)}
    \Longleftrightarrow\text{(i 5)}
    \Longleftrightarrow\text{(i 6)}
    \Longleftrightarrow\text{(i 7)}
    \Longleftrightarrow\text{(i 8)}
\]
by \cref{thm:I-3-1}(i),(iv), as well as the implication
\(\text{(i 1)}\Longrightarrow\text{(i 2)}\). Suppose henceforth that
\(\mathcal A\) is ultra-strongly closed. Put
\[
  \mathfrak X'=\mathfrak H\ominus\mathfrak X.
\]
By \cref{lem:I-3-4}, every operator in \(\mathcal A\) is reduced by
\(\mathfrak X\) and \(\mathfrak X'\), and induces zero on
\(\mathfrak X'\). The involutive algebra of operators
\[
  \mathcal B=\mathcal A_{\mathfrak X}
\]
on \(\mathfrak X\) is ultra-strongly closed. Applying
\cref{lem:I-3-6} to \(\mathcal B\) first shows that
\(I_{\mathfrak X}\in\mathcal B\), and hence that
\(P_{\mathfrak X}\) is indeed the greatest projection in
\(\mathcal A\).

The operators in \(\mathcal A'\) are reduced by
\(\mathfrak X\) and \(\mathfrak X'\), induce on \(\mathfrak X'\) an
arbitrary operator in \(\mathcal L(\mathfrak X')\), and induce on
\(\mathfrak X\) an operator in \(\mathcal B'\). Thus the operators in
\(\mathcal A''\) are reduced by \(\mathfrak X\) and
\(\mathfrak X'\), induce on \(\mathfrak X'\) a scalar operator, and
induce on \(\mathfrak X\) an operator in \(\mathcal B''\). By
\cref{lem:I-3-6}, \(\mathcal B''=\mathcal B\). The operators in
\(\mathcal A''\) are therefore precisely the operators
\[
  T+\lambda I_{\mathfrak H},
  \qquad T\in\mathcal A,\quad\lambda\in\mathbb C.
\]
% Source PDF page 049 (printed p. 42)
\phantomsection\label{srcpage:42}
Finally, if \(S\in\mathcal L(\mathfrak H)\) is in the weak closure of
\(\mathcal A\), then \(S\in\mathcal A''\), and hence
\[
  S=T+\lambda I_{\mathfrak H}
\]
for some \(T\in\mathcal A\); on the other hand,
\[
  S=SP_{\mathfrak X}
    =TP_{\mathfrak X}+\lambda P_{\mathfrak X}
    \in\mathcal A.
\]
Thus \(\mathcal A\) is weakly closed, which proves the implication
\(\text{(i 7)}\Longrightarrow\text{(i 1)}\).
\end{proof}

\begin{corollary}\label{cor:I-3-1}
Let \(\mathcal A\) be an involutive algebra of operators on
\(\mathfrak H\) such that
\[
  \mathfrak X_{\mathfrak H}^{\mathcal A}=\mathfrak H
\]
(for example, such that \(I_{\mathfrak H}\in\mathcal A\)). Then
\(\mathcal A''\), that is, the von Neumann algebra generated by
\(\mathcal A\), is the closure of \(\mathcal A\) for any one of the
weak, strong, ultra-weak, and ultra-strong topologies.
\end{corollary}

\begin{proof}
Let \(\mathcal A_1\) (respectively, \(\mathcal A_2\)) be the closure
of \(\mathcal A\) for the weak (respectively, ultra-weak) topology;
it is an involutive algebra of operators. We have
\[
  \mathcal A\subset\mathcal A_1\subset\mathcal A'',
\]
and therefore
\[
  \mathcal A''\subset\mathcal A_1''\subset\mathcal A'',
\]
that is, \(\mathcal A_1''=\mathcal A''\). On the other hand,
\(\mathcal A_1=\mathcal A_1''\) by \cref{thm:I-3-2}. Hence
\(\mathcal A_1=\mathcal A''\). We similarly see that
\(\mathcal A_2=\mathcal A''\). Finally,
\(\mathcal A_1\) (respectively, \(\mathcal A_2\)) is also the strong
(respectively, ultra-strong) closure of \(\mathcal A\), by
\cref{thm:I-3-1}.
\end{proof}

\begin{corollary}\label{cor:I-3-2}
Let \(\mathcal A\) be an involutive algebra of operators on
\(\mathfrak H\) such that
\(\mathfrak X_{\mathfrak H}^{\mathcal A}=\mathfrak H\), and let
\(\mathcal A_1\) be its unit ball. For \(\mathcal A\) to be a von
Neumann algebra, it is necessary and sufficient that
\(\mathcal A\), or \(\mathcal A_1\), be closed for any one of the weak,
strong, ultra-weak, and ultra-strong topologies.
\end{corollary}

\begin{proof}
For every set \(\mathfrak M\subset\mathcal L(\mathfrak H)\), its
commutant \(\mathfrak M'\) is weakly closed. The conditions are
therefore necessary. They are sufficient by \cref{thm:I-3-2}.
\end{proof}

A norm-closed involutive subalgebra of \(\mathcal L(\mathfrak H)\) is
called a \(C^*\)-algebra of operators on \(\mathfrak H\), or a
\(C^*\)-subalgebra of \(\mathcal L(\mathfrak H)\). Von Neumann
algebras are special cases of \(C^*\)-algebras.

\begin{corollary}\label{cor:I-3-3}
Let \(\mathcal A\) be a von Neumann algebra, and let \(\mathfrak m\)
be an ultra-weakly closed left ideal of \(\mathcal A\). Then
\(\mathfrak m\) is weakly closed. There exists a unique projection
\(E\in\mathcal A\) such that \(\mathfrak m\) is the set of
\(T\in\mathcal A\) satisfying \(T=TE\). If \(\mathfrak m\) is a
two-sided ideal, \(E\) belongs to the centre of \(\mathcal A\).
\end{corollary}

\begin{proof}
Let
\[
  \mathfrak n=\mathfrak m\cap\mathfrak m^*,
\]
which is an ultra-weakly closed involutive algebra of operators. Let
\(E\) be the greatest projection in \(\mathfrak n\), and let
\(\mathfrak m'\) be the weakly closed set of \(T\in\mathcal A\) such
that \(T=TE\). We have \(E\in\mathfrak m\), so the relation
\(T\in\mathfrak m'\) implies \(T\in\mathfrak m\); hence
\(\mathfrak m'\subset\mathfrak m\). Conversely, let
\(T\in\mathfrak m\), and let \(T=W\lvert T\rvert\) be its polar
decomposition. We have \(\lvert T\rvert\in\mathfrak m\), hence
\(\lvert T\rvert\in\mathfrak n\), and therefore
\[
  \lvert T\rvert E=\lvert T\rvert.
\]
Consequently
\[
  TE=W\lvert T\rvert E=W\lvert T\rvert=T,
\]
so that \(T\in\mathfrak m'\). Thus
\(\mathfrak m=\mathfrak m'\). If \(F\) is a projection in
\(\mathcal A\) such that \(\mathfrak m=\mathcal AF\), then
\(E=EF\), whence \(E\leq F\), and similarly \(F\leq E\); hence
\(E=F\).
% Source PDF page 050 (printed p. 43)
\phantomsection\label{srcpage:43}
Finally, if \(\mathfrak m\) is a two-sided ideal, then
\(\mathfrak m\), and hence \(E\), is invariant under the unitary
operators in \(\mathcal A\). Thus \(E\in\mathcal A'\), so \(E\)
belongs to the centre of \(\mathcal A\).
\end{proof}

\begin{corollary}\label{cor:I-3-4}
If \(\mathcal A\) is a factor, every nonzero two-sided ideal of
\(\mathcal A\) is ultra-strongly dense in \(\mathcal A\).
\end{corollary}

\begin{proof}
If \(\mathfrak m\) is a nonzero two-sided ideal of \(\mathcal A\), its
ultra-strong closure is a nonzero ultra-weakly closed two-sided ideal,
and the central projection \(E\) of \cref{cor:I-3-3} is necessarily
equal to \(I_{\mathfrak H}\), since \(\mathcal A\) is a factor.
\end{proof}

\begin{corollary}\label{cor:I-3-5}
Let \(\mathcal A\) be a von Neumann algebra, let \(\mathfrak m\) be a
two-sided ideal of \(\mathcal A\), and let
\(\overline{\mathfrak m}\) be its weak closure. For every
\(T\in(\overline{\mathfrak m})^+\), there is an upward directed set
\(\mathfrak F\subset\mathfrak m^+\) whose supremum is \(T\).
\end{corollary}

\begin{proof}
Let \((T_i)_{i\in I}\) be a maximal family of nonzero operators in
\(\mathfrak m^+\) such that
\[
  \sum_{i\in J}T_i\leq T
\]
for every finite subset \(J\) of \(I\). The operators
\(\sum_{i\in J}T_i\) form an upward directed set
\(\mathfrak F\subset\mathfrak m^+\), whose supremum \(S\) is an
element of \((\overline{\mathfrak m})^+\) majorized by \(T\). Put
\[
  R=T-S\in(\overline{\mathfrak m})^+.
\]
As \(A\in\mathfrak m\) tends weakly to the greatest projection \(E\)
in \(\overline{\mathfrak m}\),
\[
  R^{1/2}AR^{1/2}
    \ \longrightarrow\ R^{1/2}ER^{1/2}=R
\]
weakly. Hence, if \(R\ne0\), then
\(R^{1/2}AR^{1/2}\ne0\) for some \(A\in\mathfrak m\), and therefore
for some \(A\in\mathfrak m^+\) such that
\[
  0\leq A\leq I_{\mathfrak H}.
\]
But then
\[
  R^{1/2}AR^{1/2}\leq R,
  \qquad
  R^{1/2}AR^{1/2}\in\mathfrak m^+,
\]
which contradicts the maximality of the family \((T_i)_{i\in I}\).
Therefore \(R=0\).
\end{proof}

The procedure used in the proof of \cref{lem:I-3-6} is used
analogously in the proof of Jacobson's algebraic density theorem. The
equivalence (i 1) \(\Longleftrightarrow\) (i 5) in
\cref{thm:I-3-2} is not valid if \(\mathcal A\) is an arbitrary vector
subspace of \(\mathcal L(\mathfrak H)\). Indeed, it is easy to
construct ultra-weakly continuous linear functionals on
\(\mathcal L(\mathfrak H)\) which are not weakly continuous (when
\(\mathfrak H\) is infinite-dimensional).

\noindent\textbf{Bibliography:} \ArticleRef{7}, \ArticleRef{12},
\ArticleRef{67}, \ArticleRef{73}, \ArticleRef{74}.

\NumberedParagraph{5}{Kaplansky's density theorem}
\label{no:I-3-5}

\begin{theorem}\label{thm:I-3-3}
Let \(\mathcal A\) and \(\mathcal B\) be involutive algebras of
operators on \(\mathfrak H\), with
\(\mathcal A\subset\mathcal B\), and suppose that \(\mathcal A\) is
strongly dense in \(\mathcal B\). Let \(\mathfrak M\) (respectively,
\(\mathfrak N\)) be the set of Hermitian elements of \(\mathcal A\)
(respectively, \(\mathcal B\)).
% Source PDF page 051 (printed p. 44)
\phantomsection\label{srcpage:44}
Then the unit ball of \(\mathcal A\) (respectively, \(\mathfrak M\)) is
strongly dense in the unit ball of \(\mathcal B\) (respectively,
\(\mathfrak N\)).
\end{theorem}

\begin{proof}
We may suppose that \(\mathcal A\) and \(\mathcal B\) are norm closed,
since the unit ball of \(\mathcal A\) (respectively, \(\mathfrak M\))
is norm dense in the unit ball of the norm closure of \(\mathcal A\)
(respectively, \(\mathfrak M\)).

Let \(S\in\mathfrak N\). Then \(S\) is in the weak closure of
\(\mathcal A\). But the mapping
\[
  A\longmapsto\frac12(A+A^*)
\]
is weakly continuous, leaves \(S\) fixed, and maps \(\mathcal A\) into
\(\mathfrak M\). Thus \(S\) is in the weak closure of \(\mathfrak M\),
and consequently in the strong closure of \(\mathfrak M\), since
\(\mathfrak M\) is convex (\cref{thm:I-3-1}).

Suppose, in addition, that \(\lVert S\rVert\leq1\). The function
\[
  x\longmapsto2x(1+x^2)^{-1}
\]
is strictly increasing from \(-1\) to \(+1\) on \([-1,+1]\). Since
\(\mathcal B\) is norm closed, there is therefore an
\(S'\in\mathfrak N\) such that
\[
  S=2S'(I_{\mathfrak H}+S'^2)^{-1}.
\]
Let \(T'\) be any element of \(\mathfrak M\), and put
\[
  T=2T'(I_{\mathfrak H}+T'^2)^{-1}.
\]
Since
\[
  \left|2x(1+x^2)^{-1}\right|\leq1
  \qquad(x\in\mathbb R),
\]
we have \(\lVert T\rVert\leq1\), and
\begin{align}
  \frac12(T-S)
    &=(I_{\mathfrak H}+T'^2)^{-1}
      \bigl[T'(I_{\mathfrak H}+S'^2)
        -(I_{\mathfrak H}+T'^2)S'\bigr]
      (I_{\mathfrak H}+S'^2)^{-1}
      \notag\\
    &=(I_{\mathfrak H}+T'^2)^{-1}(T'-S')
      (I_{\mathfrak H}+S'^2)^{-1}
      \notag\\
    &\quad+
      (I_{\mathfrak H}+T'^2)^{-1}T'(S'-T')S'
      (I_{\mathfrak H}+S'^2)^{-1}
      \notag\\
    &=(I_{\mathfrak H}+T'^2)^{-1}(T'-S')
      (I_{\mathfrak H}+S'^2)^{-1}
      +\frac14T(S'-T')S.
  \tag{1}\label{eq:I-3-kaplansky-1}
\end{align}
As \(T'\) tends strongly to \(S'\), which is possible by what precedes,
we see that \(T\) tends strongly to \(S\), in view of the inequalities
\[
  \lVert(I_{\mathfrak H}+T'^2)^{-1}\rVert\leq1,
  \qquad
  \lVert T\rVert\leq1.
\]
We have thus shown that the unit ball of \(\mathfrak M\) is strongly
dense in that of \(\mathfrak N\).

To complete the proof, form the Hilbert space
\[
  \widetilde{\mathfrak H}
    =\mathfrak H_1\oplus\mathfrak H_2,
  \qquad
  \mathfrak H_1=\mathfrak H_2=\mathfrak H,
\]
and represent every operator
\(T\in\mathcal L(\widetilde{\mathfrak H})\) by a matrix
\((T_{ij})_{i,j=1,2}\), where
\(T_{ij}\in\mathcal L(\mathfrak H)\)
(\hyperref[no:I-2-3]{\S~2, no.~3}). Let
\(\widetilde{\mathcal A}\) (respectively,
\(\widetilde{\mathcal B}\)) be the set of
\(T\in\mathcal L(\widetilde{\mathfrak H})\) such that
\[
  T_{ij}\in\mathcal A
  \quad\text{(respectively, \(T_{ij}\in\mathcal B\))}
  \qquad(i,j=1,2).
\]
It is immediate that \(\widetilde{\mathcal A}\) and
\(\widetilde{\mathcal B}\) are involutive algebras of operators on
\(\widetilde{\mathfrak H}\), and that
\(\widetilde{\mathcal A}\) is strongly dense in
\(\widetilde{\mathcal B}\). This being so, let \(S\in\mathcal B\),
with \(\lVert S\rVert\leq1\), and define
\(\widetilde S\in\mathcal L(\widetilde{\mathfrak H})\) by the matrix
\[
  \widetilde S=
  \begin{pmatrix}
    0&S^*\\
    S&0
  \end{pmatrix}.
\]
We see at once that \(\lVert\widetilde S\rVert\leq1\) and that
\(\widetilde S\) is Hermitian. Hence there are Hermitian operators
\(R=(R_{ij})\) in \(\widetilde{\mathcal A}\), with
\(\lVert R\rVert\leq1\), which tend strongly to \(\widetilde S\).
Then \(R_{21}\) tends strongly to \(S\), and
\(\lVert R_{21}\rVert\leq1\).
\end{proof}

% Source PDF page 052 (printed p. 45)
\phantomsection\label{srcpage:45}
\begin{corollary*}
Let \(\mathfrak H\) be a Hilbert space, let \(\mathcal A\) be a
\(C^*\)-algebra of operators on \(\mathfrak H\), let
\(\mathcal A_1^+\) be the set of elements of norm at most \(1\) in
\(\mathcal A^+\), and let \(\varphi\) be a linear functional on
\(\mathcal A\). If
\(\varphi\vert_{\mathcal A_1^+}\) is strongly continuous at \(0\),
then \(\varphi\) is ultra-weakly continuous.
\end{corollary*}

\begin{proof}
Let \(\mathcal A^h\) be the set of Hermitian elements of
\(\mathcal A\), and let \(\mathcal A_1\) and
\(\mathcal A_1^h\) be the unit balls of \(\mathcal A\) and
\(\mathcal A^h\). The mappings
\[
  T\longmapsto T^+,
  \qquad
  T\longmapsto T^-
\]
from \(\mathcal A_1^h\) into \(\mathcal A_1^+\) are strongly
continuous at \(0\), since
\[
  \lVert T^+x\rVert\leq\lVert Tx\rVert,
  \qquad
  \lVert T^-x\rVert\leq\lVert Tx\rVert
  \qquad(x\in\mathfrak H).
\]
Thus \(\varphi\vert_{\mathcal A_1^h}\) is strongly continuous at
\(0\). By linearity,
\(\varphi\vert_{2\mathcal A_1^h}\) is strongly continuous at \(0\),
and hence \(\varphi\vert_{\mathcal A^h}\) is strongly continuous.

Let \(D\) be a closed convex subset of \(\mathbb C\). Then
\[
  \varphi^{-1}(D)\cap\mathcal A_1^h
\]
is a strongly closed convex subset of \(\mathcal A_1^h\). Every point
of \(\mathcal L(\mathfrak H)\) in its weak closure is also in its
strong closure by \cref{thm:I-3-1}(iv). Therefore
\(\varphi^{-1}(D)\cap\mathcal A_1^h\) is weakly closed in
\(\mathcal A_1^h\). The complement of every open square
\(D'\subset\mathbb C\) is the union of four closed half-planes; hence
\[
  \varphi^{-1}(D')\cap\mathcal A_1^h
\]
is weakly open in \(\mathcal A_1^h\). Thus
\(\varphi\vert_{\mathcal A_1^h}\) is weakly continuous. Since
\[
  T\longmapsto\frac12(T+T^*),
  \qquad
  T\longmapsto\frac1{2i}(T-T^*)
\]
are weakly continuous mappings of \(\mathcal A_1\) into
\(\mathcal A_1^h\), the restriction
\(\varphi\vert_{\mathcal A_1}\) is weakly continuous.

Let \(\mathcal B\) be the weak closure of \(\mathcal A\). For every
\(r\geq0\), let \(\mathcal A_r\) (respectively, \(\mathcal B_r\)) be
the closed ball with centre \(0\) and radius \(r\) in \(\mathcal A\)
(respectively, \(\mathcal B\)). Then \(\mathcal A_r\) is weakly dense
in \(\mathcal B_r\) by \cref{thm:I-3-3}. Thus
\(\varphi\vert_{\mathcal A_r}\) extends to a weakly continuous
function \(\psi_r\) on \(\mathcal B_r\). If \(r\leq r'\), then
\(\psi_{r'}\) extends \(\psi_r\). Hence there is a function \(\psi\)
on \(\mathcal B\) extending all the \(\psi_r\). It is immediate that
\(\psi\) is a linear functional on \(\mathcal B\), whose restriction
to \(\mathcal B_1\) is weakly continuous. Therefore \(\psi\) is
ultra-weakly continuous by \cref{thm:I-3-1}(ii).
\end{proof}

\noindent\textbf{Bibliography:} \ArticleRef{41}, \ArticleRef{67},
\ArticleRef{367}.

\par\medskip
\noindent\textit{Exercises.}---
\begin{enumerate}[label=\arabic*.,leftmargin=*,itemsep=0.5\baselineskip]
\item\label{ex:I-3-1}
Let \(\mathfrak H\) be an infinite-dimensional complex Hilbert space.
\begin{enumerate}[label=\alph*.,leftmargin=2em]
\item Let \((e_1,e_2,\ldots)\) be an orthonormal system in
\(\mathfrak H\), and let \(U_n\) be the partial isometry
\[
  x\longmapsto(x\mid e_n)e_1.
\]
Show that \(U_n^*\) is the operator
\[
  x\longmapsto(x\mid e_1)e_n,
\]
that \(U_n\) tends ultra-strongly to zero, that \(U_n^*\) does not tend
strongly to zero, and that \(U_n\) does not tend in norm to zero.
\item Deduce from (a) that the mapping \(T\mapsto T^*\) is not
continuous for the strong (respectively, ultra-strong) topology.
\item Deduce from (b) that the strong (respectively, ultra-strong)
topology is strictly finer than the weak (respectively, ultra-weak)
topology.
\item Deduce from (a) that the norm topology is strictly finer than the
ultra-strong topology \ArticleRef{74}.
\end{enumerate}

\item\label{ex:I-3-2}
Let \(\mathfrak H\) be an infinite-dimensional complex Hilbert space.
\begin{enumerate}[label=\alph*.,leftmargin=2em]
\item Let \(\mathfrak V\) be an ultra-strong neighbourhood of zero in
\(\mathcal L(\mathfrak H)\). Deduce from
\hyperref[ex:I-3-1]{Exercise~1(d)} that there is a
\(z\in\mathfrak H\) such that
\[
  \sup_{T\in\mathfrak V}\lVert Tz\rVert=+\infty.
\]
(Assuming the contrary, apply Theorem~2 of \TreatiseRef{3},
Chapter~III, \S~3, to \(\mathfrak V\).)
\item Deduce from (a) that, if
\(\mathcal L(\mathfrak H)\times\mathcal L(\mathfrak H)\) is endowed
with the product of the ultra-strong topologies and
\(\mathcal L(\mathfrak H)\) with the weak topology, the mapping
\[
  (S,T)\longmapsto ST
\]
is not continuous.

\(\bigl[\)Let \(x,y\) be nonzero elements, and let
\((x_i)_{i\in I}\) be a family of elements of \(\mathfrak H\) such
that
\[
  \sum_{i\in I}\lVert x_i\rVert^2<+\infty.
\]
Let \(\mathfrak W\) be the set of \(T\in\mathcal L(\mathfrak H)\)
such that
\[
  \sum_{i\in I}\lVert Tx_i\rVert^2\leq1.
\]
Let \(z\in\mathfrak H\) be such that
\[
  \sup_{T\in\mathfrak W}\lVert Tz\rVert=+\infty.
\]
Choose \(S\in\mathfrak W\) such that
\(Sx=\lambda z\), \(\lambda\ne0\); then choose
\(T\in\mathfrak W\) such that
\[
  \lVert TSx\rVert
    =\lvert\lambda\rvert\lVert Tz\rVert
    >\lVert y\rVert^{-1};
\]
finally choose a unitary operator \(U\) such that
\[
  \lvert(UTSx\mid y)\rvert>1.
\]
We have \(S\in\mathfrak W\) and \(UT\in\mathfrak W\).\(\bigr]\)
\ArticleRef{74}.
\end{enumerate}

% Source PDF page 053 (printed p. 46)
\phantomsection\label{srcpage:46}
\item\label{ex:I-3-3}
Let \((e_1,e_2,\ldots)\) be an orthonormal basis of the complex
Hilbert space \(\mathfrak H\). Let \(P_n\) be the projection onto the
line \(\mathbb Ce_n\), and put
\[
  T_{m,n}=P_m+mP_n.
\]
Show that zero is in the ultra-strong closure of the set
\(\mathfrak M\) of the \(T_{m,n}\), but that no sequence extracted
from \(\mathfrak M\) converges weakly to zero. (If
\(T_{m_\nu,n_\nu}\) tends weakly to zero, then
\(\lVert T_{m_\nu,n_\nu}\rVert\) is bounded, so \(m_\nu\) is bounded
and assumes the same value infinitely often.) Deduce that
\(\mathcal L(\mathfrak H)\) is not metrizable for any of the
ultra-strong, ultra-weak, strong, and weak topologies. (Since there is
an ultra-strongly dense sequence in \(\mathcal L(\mathfrak H)\), if
\(\mathcal L(\mathfrak H)\) were metrizable, it would have a
countable base.) \ArticleRef{73}, \ArticleRef{74}.

\item\label{ex:I-3-4}
Use the notation of \cref{thm:I-3-1}. Show that, if
\(\mathfrak H\) has a countable basis, the Banach space
\(\mathfrak M_*\), endowed with the norm induced by that of
\(\mathfrak M^*\), has a countable basis. (If \(\mathfrak D\) is a
dense sequence in \(\mathfrak H\), the linear combinations, with
rational complex coefficients, of the
\(\omega_{x,y}\), where \(x,y\in\mathfrak D\), are dense in
\(\mathfrak M_*\).)

\item\label{ex:I-3-5}
Let \(\mathcal A\) be a von Neumann algebra, let
\(\mathfrak m\) and \(\mathfrak m_1\) be two-sided ideals of
\(\mathcal A\), and let \(\mathfrak n\) be their product ideal in the
usual algebraic sense. Let
\(\overline{\mathfrak m}\), \(\overline{\mathfrak m_1}\), and
\(\overline{\mathfrak n}\) be the strong closures of
\(\mathfrak m\), \(\mathfrak m_1\), and \(\mathfrak n\). Show that
\[
  \overline{\mathfrak n}
    =\overline{\mathfrak m}\cap\overline{\mathfrak m_1}.
\]
\(\bigl[\)Let \(E,E_1\) be the greatest projections in
\(\overline{\mathfrak m}\), \(\overline{\mathfrak m_1}\). There are
elements \(S\) (respectively, \(S_1\)) of \(\mathfrak m\)
(respectively, \(\mathfrak m_1\)) which tend strongly to \(E\)
(respectively, \(E_1\)) while remaining bounded. Then
\(SS_1\in\mathfrak n\) tends strongly to \(EE_1\), and therefore
\[
  EE_1\in\overline{\mathfrak n},
  \qquad
  \overline{\mathfrak m}\cap\overline{\mathfrak m_1}
    \subset\overline{\mathfrak n}.
\]
\(\bigr]\) \ArticleRef{12}.

\item\label{ex:I-3-6}
\begin{enumerate}[label=\alph*.,leftmargin=2em]
\item Let \(\mathfrak H\) be a complex Hilbert space, let
\(\mathcal A\) be an involutive algebra of operators on
\(\mathfrak H\), and let \(\mathcal B\) be the weak closure of
\(\mathcal A\). Suppose that \(I_{\mathfrak H}\in\mathcal B\).
Suppose also that every linear functional \(\varphi\) on
\(\mathcal A\) which is continuous for the norm topology is
ultra-strongly continuous. Let \(\psi\) be the ultra-strongly
continuous extension of \(\varphi\) to \(\mathcal B\). Show that
\[
  \varphi\longmapsto\psi
\]
is an isometric linear mapping of the dual \(\mathcal A^*\) of
\(\mathcal A\) onto the space \(\mathcal B^*\) of ultra-strongly
continuous linear functionals on \(\mathcal B\). (Use
\cref{thm:I-3-3}.) Deduce that the bidual of \(\mathcal A\) is
canonically isomorphic to \(\mathcal B\), the canonical mapping of
\(\mathcal A\) into its bidual being identified with the identity
mapping of \(\mathcal A\) into \(\mathcal B\). (Use
\cref{thm:I-3-1}.)
\item Show that, if \(\mathcal A\) is taken to be the set of compact
operators on \(\mathfrak H\), the conditions in (a) are satisfied and
\(\mathcal B=\mathcal L(\mathfrak H)\).

\(\bigl[\)Let \(\varphi\) be a norm-continuous linear functional on
\(\mathcal A\). To show that \(\varphi\) is ultra-strongly continuous,
reduce to the case in which
\[
  \varphi(S^*)=\overline{\varphi(S)}
  \qquad(S\in\mathcal A).
\]
For \(x,y\in\mathfrak H\), let \(\{x,y\}\) be the operator
\[
  z\longmapsto(z\mid x)y.
\]
Show that
\[
  \varphi(\{x,y\})=(A_\varphi y\mid x),
\]
where \(A_\varphi\in\mathcal L(\mathfrak H)\) depends only on
\(\varphi\) and is Hermitian. Show that, for every orthonormal system
\((e_1,e_2,\ldots)\) in \(\mathfrak H\),
\[
  \sum_i\lvert(A_\varphi e_i\mid e_i)\rvert<+\infty,
\]
% Source PDF page 054 (printed p. 47)
\phantomsection\label{srcpage:47}
by observing that
\[
  \sum_i\lambda_i\varphi(\{e_i,e_i\})
\]
converges for every sequence \((\lambda_i)\) of real numbers tending
to zero. Deduce that \(A_\varphi\) is compact and that, if
\((\mu_i)\) is the family of eigenvalues of \(A_\varphi\)
corresponding to an orthonormal basis \((e_i)\) formed of
eigenvectors, then
\[
  \sum_i\lvert\mu_i\rvert<+\infty.
\]
Conclude that
\[
  \varphi=\sum_i\mu_i\omega_{e_i,e_i}
\]
is ultra-strongly continuous.\(\bigr]\)
\item In the dual \(\mathcal L^*\) of the Banach space
\(\mathcal L=\mathcal L(\mathfrak H)\), let \(\mathcal L_*\) be the
subspace of ultra-strongly continuous linear functionals, and let
\(\mathcal A^\perp\) be the subspace orthogonal to the set
\(\mathcal A\) of compact operators. Deduce from (b) that
\(\mathcal L^*\) is the direct sum of the subspaces
\(\mathcal L_*\) and \(\mathcal A^\perp\). Show that, if
\(\varphi_1\in\mathcal L_*\) and
\(\varphi_2\in\mathcal A^\perp\), then
\[
  \lVert\varphi_1+\varphi_2\rVert
    =\lVert\varphi_1\rVert+\lVert\varphi_2\rVert.
\]

\(\bigl[\)Construct \(S,T\in\mathcal L\), with \(S\) of finite rank,
such that
\[
  \lVert S\rVert\leq1,
  \qquad
  \lvert\varphi_1(S)\rvert\geq\lVert\varphi_1\rVert-\varepsilon,
\]
and
\[
  \lVert T\rVert\leq1,
  \qquad
  \lvert\varphi_2(T)\rvert\geq\lVert\varphi_2\rVert-\varepsilon;
\]
then, by adding a suitable finite-rank operator to \(T\), arrange in
addition that
\[
  \lVert S+T\rVert\leq1.
\]
\(\bigr]\) \ArticleRef{74}, \ArticleRef{111}, \ArticleRef{112},
\ArticleRef{113}, \TreatiseRef{13}.
\end{enumerate}

\item\label{ex:I-3-7}
Let \(\mathcal A\) be a von Neumann algebra on \(\mathfrak H\), let
\(\mathfrak M\) be a strongly dense subset of \(\mathcal A\), and let
\(T'\) be a closed operator on \(\mathfrak H\).
\begin{enumerate}[label=\alph*.,leftmargin=2em]
\item For \(T'\,\eta\,\mathcal A'\)
(\hyperref[ex:I-1-10]{\S~1, Exercise~10}), it is necessary and
sufficient that \(T'\) commute with every \(S\in\mathfrak M\), that
is, that \(T'S\) extend \(ST'\).
\item Suppose that \(T'\) is self-adjoint and \(T'\geq0\), and let
\(\mathfrak D\) be its domain. For
\(T'\,\eta\,\mathcal A'\), it is necessary and sufficient that, for
every \(S\in\mathfrak M\), \(\mathfrak D\) be stable under \(S\) and
\(S^*\), and that, for \(x,y\in\mathfrak D\),
\[
  (T'S^*x\mid T'y)=(T'x\mid T'Sy).
\]

\(\bigl[\)Taking \(x\) and \(y\) in the domain \(\mathfrak N\) of
\(T'^2\), we have
\[
  (T'^2x\mid Sy)=(S^*x\mid T'^2y)=(x\mid ST'^2y),
\]
whence \(Sy\in\mathfrak N\) and
\[
  T'^2Sy=ST'^2y.
\]
Thus \(T'^2\,\eta\,\mathcal A'\), by (a).\(\bigr]\)
\end{enumerate}
\end{enumerate}

\section{Positive Linear Forms}\label{sec:I-4}

\NumberedParagraph{1}{Positive linear forms on an involutive operator
algebra}\label{no:I-4-1}
Let \(\mathcal A\) be an involutive algebra of operators on
\(\mathfrak H\), containing \(I_{\mathfrak H}\). A linear functional
\(\varphi\) on \(\mathcal A\) is said to be \emph{positive} if
\(\varphi(T)\geq0\) for every \(T\in\mathcal A^+\). We then have
\[
  \varphi(T^*)=\overline{\varphi(T)},
  \qquad
  \varphi(T^*T)\geq0
  \qquad(T\in\mathcal A).
\]
Consequently,
\[
  \lvert\varphi(S^*T)\rvert^2
    \leq\varphi(S^*S)\varphi(T^*T)
  \qquad(S,T\in\mathcal A)
\]
(the Cauchy-Schwarz inequality). In particular,
\[
  \lvert\varphi(T)\rvert^2
    \leq\varphi(I_{\mathfrak H})\varphi(T^*T)
    \leq\varphi(I_{\mathfrak H})^2\lVert T^*T\rVert
    =\varphi(I_{\mathfrak H})^2\lVert T\rVert^2,
\]
so that \(\varphi\) is continuous and has norm
\(\varphi(I_{\mathfrak H})\). For every \(T_0\in\mathcal A\), the
linear functional
\[
  T\longmapsto\varphi(T_0TT_0^*)
\]
is positive, since \(T\geq0\) implies \(T_0TT_0^*\geq0\).

% Source PDF page 055 (printed p. 48)
\phantomsection\label{srcpage:48}
A positive linear functional \(\varphi\) is said to be
\emph{faithful} if the conditions
\[
  \varphi(T)=0,\qquad T\in\mathcal A^+
\]
imply \(T=0\).

Let \(\varphi\) and \(\psi\) be two linear functionals on
\(\mathcal A\). We say that \(\varphi\) \emph{majorizes} \(\psi\), and
write \(\varphi\geq\psi\), if \(\varphi-\psi\) is positive.

If \(x\in\mathfrak H\), the linear functional
\(\omega_{x,x}\) on \(\mathcal A\) is positive; we shall denote it
simply by \(\omega_x\).

\begin{lemma}\label{lem:I-4-1}
Let \(x\in\mathfrak H\), and let \(\psi\) be a positive linear
functional on \(\mathcal A\) majorized by \(\omega_x\). There is an
operator \(T'\in\mathcal A'\) such that
\[
  0\leq T'\leq I_{\mathfrak H},
  \qquad
  \psi=\omega_{T'x}.
\]
\end{lemma}

\begin{proof}
For \(S,T\in\mathcal A\), we have
\[
  \lvert\psi(S^*T)\rvert^2
    \leq\psi(S^*S)\psi(T^*T)
    \leq\lVert Sx\rVert^2\lVert Tx\rVert^2.
\]
Thus, on putting
\[
  ((Tx\mid Sx))=\psi(S^*T),
\]
we uniquely define on the subspace \(\mathcal Ax\) of \(\mathfrak H\)
a sesquilinear form which is plainly Hermitian, positive, and of norm
at most \(1\). Hence there is a Hermitian positive operator \(T_0\) on
the space
\[
  \overline{\mathcal Ax}=\mathfrak X
\]
such that
\[
  \psi(S^*T)=(Tx\mid T_0Sx),
  \qquad
  \lVert T_0\rVert\leq1.
\]
For \(R,S,T\in\mathcal A\), we have
\begin{align*}
  (Rx\mid T_0TSx)
    &=\psi((TS)^*R)
      =\psi(S^*(T^*R))\\
    &=(T^*Rx\mid T_0Sx)
      =(Rx\mid TT_0Sx),
\end{align*}
so that \(T_0T=TT_0\) on \(\mathfrak X\). Thus
\(T_0P_{\mathfrak X}\) is a positive Hermitian operator in
\(\mathcal A'\), of norm at most \(1\). Let \(T'\) be its square root.
For \(T\in\mathcal A\), we have
\[
  \psi(T)=(Tx\mid T_0x)
    =(Tx\mid T'^2x)
    =(TT'x\mid T'x).
\]
\end{proof}

\begin{lemma}\label{lem:I-4-2}
If \(\omega_{x,y}\) is a positive linear functional on
\(\mathcal A\), there is a \(z\in\mathfrak H\) such that
\[
  \omega_{x,y}=\omega_z
\]
on \(\mathcal A\).
\end{lemma}

\begin{proof}
For every \(T\in\mathcal A^+\), we have
\begin{align*}
  4\omega_{x,y}(T)
    &=2\omega_{x,y}(T)+2\omega_{x,y}(T^*)\\
    &=2(Tx\mid y)+2(Ty\mid x)\\
    &=(T(x+y)\mid x+y)-(T(x-y)\mid x-y)\\
    &\leq(T(x+y)\mid x+y).
\end{align*}
Thus
\[
  4\omega_{x,y}\leq\omega_{x+y}.
\]
The result now follows from \cref{lem:I-4-1}.
\end{proof}

\begin{lemma}\label{lem:I-4-3}
Let \(\mathcal B\) (respectively, \(\mathcal B_1\)) be an involutive
algebra of operators containing the identity on a Hilbert space
\(\mathfrak K\) (respectively, \(\mathfrak K_1\)). Let \(\Phi\)
(respectively, \(\Phi_1\)) be an involutive-algebra homomorphism of
\(\mathcal A\) onto \(\mathcal B\) (respectively,
\(\mathcal B_1\)). Suppose there is an \(x\in\mathfrak K\)
(respectively, \(x_1\in\mathfrak K_1\)), totalizing for
\(\mathcal B\) (respectively, \(\mathcal B_1\)), such that
\[
  (\Phi(T)x\mid x)=(\Phi_1(T)x_1\mid x_1)
  \qquad(T\in\mathcal A).
\]
Then there is an isomorphism of the Hilbert space \(\mathfrak K\) onto
the Hilbert space \(\mathfrak K_1\) which transforms
\(\mathcal B\) into \(\mathcal B_1\), and \(\Phi\) into \(\Phi_1\).
\end{lemma}

\begin{proof}
For every \(T\in\mathcal A\), we have
\[
  \lVert\Phi(T)x\rVert^2
    =(\Phi(T)^*\Phi(T)x\mid x)
    =(\Phi(T^*T)x\mid x)
    =(\Phi_1(T^*T)x_1\mid x_1)
    =\lVert\Phi_1(T)x_1\rVert^2.
\]
% Source PDF page 056 (printed p. 49)
\phantomsection\label{srcpage:49}
Hence there is a linear isometry \(U\) of \(\mathcal Bx\) onto
\(\mathcal B_1x_1\) such that
\[
  \Phi_1(T)x_1=U\Phi(T)x
  \qquad(T\in\mathcal A).
\]
The mapping \(U\) extends to an isomorphism \(V\) of
\(\mathfrak K\) onto \(\mathfrak K_1\). For \(S,T\in\mathcal A\),
\begin{align*}
  \Phi_1(S)V\Phi(T)x
    &=\Phi_1(S)\Phi_1(T)x_1
      =\Phi_1(ST)x_1\\
    &=U\Phi(ST)x
      =V\Phi(S)\Phi(T)x.
\end{align*}
Consequently
\[
  \Phi_1(S)V=V\Phi(S).
\]
\end{proof}

The preceding lemma is a uniqueness result. The following lemma will
be an existence result. The meaning of the adverb
\emph{canonically} in the statement of the lemma will be made precise
in the proof.

\begin{lemma}\label{lem:I-4-4}
Every positive linear functional \(\varphi\) on \(\mathcal A\)
canonically defines a Hilbert space \(\mathfrak K\), a linear mapping
\(\Gamma\) of \(\mathcal A\) onto a dense vector subspace of
\(\mathfrak K\), and a norm-decreasing homomorphism
\[
  \Phi:\mathcal A\longrightarrow\mathcal L(\mathfrak K),
\]
such that, on putting
\[
  x=\Gamma(I_{\mathfrak H})\in\mathfrak K,
\]
we have
\[
  \Gamma(T)=\Phi(T)x,
  \qquad
  \varphi(T)=(\Phi(T)x\mid x)
  \qquad(T\in\mathcal A).
\]
Moreover,
\[
  \Phi(I_{\mathfrak H})=I_{\mathfrak K}.
\]
If \(\varphi\) is faithful, \(\Phi\) is an isomorphism of
\(\mathcal A\) onto \(\Phi(\mathcal A)\), and \(x\) is separating for
\(\Phi(\mathcal A)\).
\end{lemma}

\begin{proof}
For \(S,T\in\mathcal A\), put
\[
  (S\mid T)=\varphi(T^*S).
\]
Then \(\mathcal A\) becomes a pre-Hilbert space. By the inequality
\[
  \lvert\varphi(T^*S)\rvert^2
    \leq\varphi(S^*S)\varphi(T^*T),
\]
the set \(\mathfrak m\) of \(S\in\mathcal A\) such that
\(\varphi(S^*S)=0\) is also the set of \(S\in\mathcal A\) such that
\(\varphi(T^*S)=0\) for every \(T\in\mathcal A\); it is therefore a
left ideal of \(\mathcal A\). The quotient space
\(\mathcal A/\mathfrak m\) is a separated pre-Hilbert space. Let
\(\mathfrak K\) be its completion. The canonical mapping
\(\Gamma\) of \(\mathcal A\) onto \(\mathcal A/\mathfrak m\) is
therefore a linear mapping of \(\mathcal A\) onto a dense vector
subspace of \(\mathfrak K\).

On the other hand, for \(S\in\mathcal A\), left multiplication by
\(S\) in \(\mathcal A\) induces, on passage to the quotient, a linear
operator \(\widetilde S\) on \(\mathcal A/\mathfrak m\). This operator
has norm at most \(\lVert S\rVert\), since, for \(T\in\mathcal A\),
\[
  (ST\mid ST)
    =\varphi(T^*S^*ST)
    \leq\lVert S^*S\rVert\varphi(T^*T)
    =\lVert S\rVert^2(T\mid T).
\]
Thus \(\widetilde S\) extends to a continuous linear operator
\(\Phi(S)\) on \(\mathfrak K\). We readily see that \(\Phi\) is a
homomorphism such that
\(\Phi(I_{\mathfrak H})=I_{\mathfrak K}\). For example, for
\(R,S,T\in\mathcal A\),
\begin{align*}
  (\Phi(R)^*\Gamma(S)\mid\Gamma(T))
    &=(\Gamma(S)\mid\Phi(R)\Gamma(T))
      =(S\mid RT)\\
    &=\varphi(T^*R^*S)
      =(R^*S\mid T)\\
    &=(\Phi(R^*)\Gamma(S)\mid\Gamma(T)),
\end{align*}
and hence
\[
  \Phi(R)^*=\Phi(R^*).
\]
Finally,
\[
  \Phi(T)x
    =\Phi(T)\Gamma(I_{\mathfrak H})
    =\Gamma(T),
\]
and
\[
  (\Phi(T)x\mid x)
    =(\Phi(T)\Gamma(I_{\mathfrak H})
       \mid\Gamma(I_{\mathfrak H}))
    =(T\mid I_{\mathfrak H})
    =\varphi(T).
\]
% Source PDF page 057 (printed p. 50)
\phantomsection\label{srcpage:50}
If \(\varphi\) is faithful, then \(\mathfrak m=0\). The condition
\(\Phi(T)x=0\) implies \(\Gamma(T)=0\), hence
\(T\in\mathfrak m\), and therefore \(T=0\).
\end{proof}

\begin{lemma}\label{lem:I-4-5}
Let \(\varphi\) be a norm-continuous linear functional on
\(\mathcal A\). If
\[
  \varphi(I_{\mathfrak H})=\lVert\varphi\rVert,
\]
then \(\varphi\) is positive.
\end{lemma}

\begin{proof}
We may suppose that
\(\varphi(I_{\mathfrak H})=\lVert\varphi\rVert=1\). Let
\(T\in\mathcal A^+\), and suppose that the number \(\varphi(T)\) is
not nonnegative. There is then a closed disk
\[
  \lvert z-z_0\rvert\leq\rho
\]
in \(\mathbb C\) which contains the spectrum of \(T\) without
containing \(\varphi(T)\). The spectrum of the normal operator
\(T-z_0I_{\mathfrak H}\) is contained in the disk
\(\lvert z\rvert\leq\rho\), and hence
\[
  \lVert T-z_0I_{\mathfrak H}\rVert\leq\rho.
\]
Therefore
\begin{align*}
  \lvert\varphi(T)-z_0\rvert
    &=\lvert\varphi(T)-z_0\varphi(I_{\mathfrak H})\rvert\\
    &=\lvert\varphi(T-z_0I_{\mathfrak H})\rvert\\
    &\leq\lVert\varphi\rVert
      \,\lVert T-z_0I_{\mathfrak H}\rVert
    \leq\rho,
\end{align*}
which is absurd.
\end{proof}

The argument in \cref{lem:I-4-5} is due to Phelps
[\emph{The range of \(Tf\) for certain linear operators \(T\)}
(Proc. Amer. Math. Soc., vol.~16, 1965, pp.~381--382)].

\noindent\textbf{Bibliography:} \ArticleRef{19}, \ArticleRef{25},
\ArticleRef{28}, \ArticleRef{31}, \ArticleRef{96}.

\NumberedParagraph{2}{Normal positive linear functionals on a von
Neumann algebra}\label{no:I-4-2}

\begin{definition}\label{def:I-4-1}
Let \(\mathcal A\) be a von Neumann algebra. A positive linear
functional \(\varphi\) on \(\mathcal A\) is said to be \emph{normal}
if, for every upward directed set
\(\mathfrak F\subset\mathcal A^+\) with supremum
\(T\in\mathcal A^+\), \(\varphi(T)\) is the supremum of
\(\varphi(\mathfrak F)\).
\end{definition}

Every positive linear functional majorized by a normal positive linear
functional is normal. If \(\varphi\) is normal and
\(T_0\in\mathcal A\), then the positive linear functional
\[
  T\longmapsto\varphi(T_0TT_0^*)
\]
is normal.

\begin{lemma}\label{lem:I-4-6}
Let \(\varphi\) be a linear functional on \(\mathcal A\) defined by
\[
  \varphi=\sum_{i=1}^{n}\omega_{x_i,y_i}
\]
(respectively,
\[
  \varphi=\sum_{i=1}^{\infty}\omega_{x_i,y_i};
  \qquad
  \sum_{i=1}^{\infty}\lVert x_i\rVert^2<+\infty,
  \qquad
  \sum_{i=1}^{\infty}\lVert y_i\rVert^2<+\infty).
\]
Let \(\mathfrak K\) be a Hilbert space of dimension \(n\)
(respectively, of infinite dimension and with a countable basis). Let
\[
  \Phi:T\longmapsto T\otimes I_{\mathfrak K}
\]
be the amplification of \(\mathcal A\) onto
\(\mathcal A\otimes\mathbb C_{\mathfrak K}\). There are
\(x,y\in\mathfrak H\otimes\mathfrak K\) such that
\[
  \varphi(T)=(\Phi(T)x\mid y).
\]

If \(y_i=x_i\) for every \(i\), we may suppose that \(y=x\).
\end{lemma}

\begin{proof}
Suppose, for example, that
\[
  \varphi=\sum_{i=1}^{\infty}\omega_{x_i,y_i},
  \qquad
  \sum_{i=1}^{\infty}\lVert x_i\rVert^2<+\infty,
  \qquad
  \sum_{i=1}^{\infty}\lVert y_i\rVert^2<+\infty.
\]
% Source PDF page 058 (printed p. 51)
\phantomsection\label{srcpage:51}
Let \((e_i)\) be an orthonormal basis of \(\mathfrak K\). Put
\[
  x=\sum_{i=1}^{\infty}x_i\otimes e_i
    \in\mathfrak H\otimes\mathfrak K,
  \qquad
  y=\sum_{i=1}^{\infty}y_i\otimes e_i
    \in\mathfrak H\otimes\mathfrak K.
\]
Then
\[
  (\Phi(T)x\mid y)
    =\left(
      \sum_{i=1}^{\infty}Tx_i\otimes e_i
      \,\middle|\,
      \sum_{i=1}^{\infty}y_i\otimes e_i
     \right)
    =\sum_{i=1}^{\infty}(Tx_i\mid y_i).
\]
\end{proof}

\begin{lemma}\label{lem:I-4-7}
Let \(\mathcal A\) be a von Neumann algebra, let
\(R\in\mathcal A^+\), and let \(\varphi\) and \(\psi\) be two normal
positive linear functionals on \(\mathcal A\) such that
\[
  \varphi(R)<\psi(R).
\]
There is a nonzero element \(S\in\mathcal A^+\), majorized by \(R\),
such that
\[
  \varphi(T)<\psi(T)
\]
for every nonzero \(T\in\mathcal A^+\) majorized by \(S\).
\end{lemma}

\begin{proof}
Let \((R_i)_{i\in I}\) be a totally ordered family of operators in
\(\mathcal A^+\) such that
\[
  R_i\leq R,
  \qquad
  \varphi(R_i)\geq\psi(R_i),
\]
and let \(R_0\) be its supremum. Then
\[
  R_0\in\mathcal A^+,\qquad R_0\leq R,
\]
and
\[
  \varphi(R_0)
    =\sup_i\varphi(R_i)
    \geq\sup_i\psi(R_i)
    =\psi(R_0).
\]
Zorn's theorem therefore proves the existence of a maximal operator
\(R_1\in\mathcal A^+\) such that
\[
  R_1\leq R,
  \qquad
  \varphi(R_1)\geq\psi(R_1).
\]
Put \(S=R-R_1\). Then \(S\in\mathcal A^+\), \(S\leq R\), and
\(S\ne0\). If \(T\in\mathcal A^+\) is nonzero and majorized by \(S\),
then \(\varphi(T)<\psi(T)\), for otherwise \(R_1\) would not be
maximal.
\end{proof}

\begin{theorem}\label{thm:I-4-1}
Let \(\mathcal A\) be a von Neumann algebra, and let \(\varphi\) be a
positive linear functional on \(\mathcal A\). The following conditions
are equivalent:
\begin{enumerate}[label=(\roman*)]
\item \(\varphi\) is normal;
\item \(\varphi\) is ultra-weakly continuous;
\item
\[
  \varphi=\sum_{i=1}^{\infty}\omega_{x_i},
  \qquad
  \sum_{i=1}^{\infty}\lVert x_i\rVert^2<+\infty.
\]
\end{enumerate}
Every ultra-weakly continuous linear functional on \(\mathcal A\) is
a linear combination of normal positive linear functionals.
\end{theorem}

\begin{proof}
The implications
\[
  \text{(iii)}\Longrightarrow\text{(ii)}
  \Longrightarrow\text{(i)}
\]
are immediate. On the other hand, the identity
\begin{align}
  4(Tx\mid y)
    &=(T(x+y)\mid x+y)-(T(x-y)\mid x-y)\notag\\
    &\quad+i(T(x+iy)\mid x+iy)
      -i(T(x-iy)\mid x-iy)
  \tag{1}\label{eq:I-4-polarization}
\end{align}
shows that every ultra-weakly continuous linear functional on
\(\mathcal A\) is a linear combination of ultra-weakly continuous
positive linear functionals.

Let us prove (ii) \(\Longrightarrow\) (iii). Let \(\varphi\) be an
ultra-weakly continuous positive linear functional on \(\mathcal A\).
With the notation of \cref{lem:I-4-6}, we have
\[
  \varphi(T)=(\Phi(T)x\mid y).
\]
% Source PDF page 059 (printed p. 52)
\phantomsection\label{srcpage:52}
By \cref{lem:I-4-2}, there is a
\(z\in\mathfrak H\otimes\mathfrak K\) such that
\[
  \varphi(T)=(\Phi(T)z\mid z).
\]
It follows that there are elements \(z_i\in\mathfrak H\) such that
\[
  \varphi(T)=\sum_{i=1}^{\infty}(Tz_i\mid z_i).
\]

Let us prove (i) \(\Longrightarrow\) (ii). Let \(\varphi\) be a normal
positive linear functional on \(\mathcal A\). Let
\((R_i)_{i\in I}\) be a totally ordered family of operators in
\(\mathcal A^+\) such that \(R_i\leq I_{\mathfrak H}\), and such that
the linear functionals
\[
  T\longmapsto\varphi(TR_i)
\]
on \(\mathcal A\) are ultra-strongly continuous; let \(R\) be the
supremum of the \(R_i\). Then
\[
  R\in\mathcal A^+,\qquad R\leq I_{\mathfrak H}.
\]
Moreover, for every \(T\) in the unit ball \(\mathcal A_1\) of
\(\mathcal A\),
\begin{align*}
  \lvert\varphi(T(R-R_i))\rvert^2
    &\leq\varphi(T(R-R_i)T^*)\varphi(R-R_i)\\
    &\leq\varphi(I_{\mathfrak H})\varphi(R-R_i).
\end{align*}
Since \(\varphi(R-R_i)\) can be made arbitrarily small,
\(\varphi(TR)\) is the uniform limit of \(\varphi(TR_i)\) on
\(\mathcal A_1\), and is therefore ultra-strongly continuous on
\(\mathcal A_1\), and consequently ultra-strongly continuous on
\(\mathcal A\) by \cref{thm:I-3-1}. This proves the existence of a
maximal operator \(S\in\mathcal A^+\) such that
\[
  S\leq I_{\mathfrak H}
\]
and such that the linear functional
\[
  T\longmapsto\varphi(TS)
\]
is ultra-strongly continuous.

Suppose that \(S\ne I_{\mathfrak H}\). We shall obtain a contradiction,
which will complete the proof of the theorem. Put
\[
  S'=I_{\mathfrak H}-S,
\]
and choose \(z\in\mathfrak H\) such that
\[
  \varphi(S')<(S'z\mid z).
\]
By \cref{lem:I-4-7}, there is an \(S_1\in\mathcal A^+\) such that
\[
  S_1\leq S',\qquad S_1\ne0,
\]
and such that
\[
  \varphi(T)<(Tz\mid z)
\]
for every nonzero \(T\in\mathcal A^+\) majorized by \(S_1\). Then, by
the Cauchy-Schwarz inequality, for every \(T\in\mathcal A\),
\begin{align*}
  \lvert\varphi(TS_1)\rvert^2
    &\leq\varphi(I_{\mathfrak H})
      \varphi(S_1T^*TS_1)\\
    &\leq\varphi(I_{\mathfrak H})(S_1T^*TS_1z\mid z)\\
    &=\varphi(I_{\mathfrak H})\lVert TS_1z\rVert^2,
\end{align*}
since
\[
  S_1T^*TS_1
    \leq\lVert T\rVert^2S_1^2
    \leq\lVert T\rVert^2\lVert S_1\rVert S_1.
\]
Thus the linear functional \(T\mapsto\varphi(TS_1)\) is strongly
continuous, and consequently
\[
  T\longmapsto\varphi(T(S+S_1))
\]
is ultra-strongly continuous. This contradicts the maximality of
\(S\).
\end{proof}

Instead of \emph{normal}, one also says \emph{completely additive}.
For another definition of normal positive linear functionals (the
classical definition), see Exercise~9. \Cref{lem:I-4-7} is the analogue
of a familiar lemma in integration theory used in one of the proofs
of the Lebesgue-Nikodym theorem.

Let \(\mathcal A\) be a von Neumann algebra on \(\mathfrak H\). Let
\[
  T\longmapsto\varphi(T)
\]
be a function on \(\mathcal A^+\) having the following properties:
\begin{enumerate}[label=\arabic*\textsuperscript{\(\circ\)},leftmargin=3em]
\item \(0\leq\varphi(T)\leq+\infty\);
\item
\[
  \varphi(T_1+T_2)=\varphi(T_1)+\varphi(T_2);
\]
\item
\[
  \varphi(\lambda T)=\lambda\varphi(T)
  \qquad(\lambda\geq0);
\]
\item If \(\mathfrak F\) is an upward directed set in
\(\mathcal A^+\), with supremum \(T\in\mathcal A^+\), then
\[
  \varphi(T)=\sup\varphi(\mathfrak F).
\]
\end{enumerate}

% Source PDF page 060 (printed p. 53)
\phantomsection\label{srcpage:53}
\noindent\textit{Problem.} Does there exist a family
\((x_i)_{i\in I}\) in \(\mathfrak H\) such that
\[
  \varphi(T)=\sum_{i\in I}(Tx_i\mid x_i)?
\]
We shall see
(\hyperref[prop:I-6-2]{\S~6, Corollary to Proposition~2}) that this is
indeed so in a particular case.

Let \(\mathcal A\) be a complete normed involutive algebra with an
identity element, such that
\[
  \lVert x^*x\rVert=\lVert x\rVert^2,
\]
that is, a normed involutive algebra isomorphic to a \(C^*\)-algebra
of operators containing the identity on a complex Hilbert space. The
elements of the form \(x^*x\) may be taken as the nonnegative elements
of \(\mathcal A\). Suppose that every majorized upward directed family
in \(\mathcal A^+\) has a supremum. This does not imply that
\(\mathcal A\) is isomorphic to a von Neumann algebra, even if
\(\mathcal A\) is abelian. Suppose in addition that, for every nonzero
\(T\in\mathcal A^+\), there is a normal positive linear functional
\(\varphi\) on \(\mathcal A\) such that \(\varphi(T)\ne0\). Then, by
\ArticleRef{166}, \(\mathcal A\) is isomorphic to a von Neumann
algebra.

\noindent\textbf{Bibliography:} \ArticleRef{9}, \ArticleRef{15},
\ArticleRef{17}, \ArticleRef{19}, \ArticleRef{42}, \ArticleRef{43},
\ArticleRef{111}, \ArticleRef{112}, \ArticleRef{113},
\ArticleRef{166}.

\NumberedParagraph{3}{Normal positive linear mappings}
\label{no:I-4-3}

\begin{definition}\label{def:I-4-2}
Let \(\mathcal A\) and \(\mathcal B\) be von Neumann algebras. A linear
mapping \(\Phi\) of \(\mathcal A\) into \(\mathcal B\) is said to be
\emph{positive} if
\[
  \Phi(\mathcal A^+)\subset\mathcal B^+
\]
(which implies \(\Phi(T^*)=\Phi(T)^*\)). The mapping \(\Phi\) is said
to be \emph{normal positive} if, in addition, for every upward directed
set \(\mathfrak F\subset\mathcal A^+\) with supremum
\(T\in\mathcal A^+\), the set \(\Phi(\mathfrak F)\) has
\(\Phi(T)\) as its supremum.
\end{definition}

\begin{theorem}\label{thm:I-4-2}
Let \(\mathcal A\) and \(\mathcal B\) be two von Neumann algebras, and
let \(\Phi\) be a normal positive linear mapping of \(\mathcal A\) into
\(\mathcal B\). Then \(\Phi\) is ultra-weakly continuous, and hence
the restriction of \(\Phi\) to the bounded subsets of \(\mathcal A\)
is weakly continuous. If, in addition, there is a constant \(k\geq0\)
such that
\[
  \Phi(T)^*\Phi(T)\leq k\Phi(T^*T)
\]
(a condition which is always satisfied if \(\mathcal B\) is abelian),
then \(\Phi\) is ultra-strongly continuous, and hence the restriction
of \(\Phi\) to the bounded subsets of \(\mathcal A\) is strongly
continuous.
\end{theorem}

\begin{proof}
For every normal positive linear functional \(\varphi\) on
\(\mathcal B\), \(\varphi\circ\Phi\) is a normal positive linear
functional on \(\mathcal A\). Thus, for every ultra-weakly continuous
linear functional \(\varphi'\) on \(\mathcal B\),
\(\varphi'\circ\Phi\) is ultra-weakly continuous by
\cref{thm:I-4-1}. Therefore \(\Phi\) is ultra-weakly continuous.

Now, if \(T\) tends ultra-strongly to zero, then \(T^*T\) tends
ultra-weakly to zero, and therefore so does \(\Phi(T^*T)\), and hence
so does \(\Phi(T)^*\Phi(T)\) if
\[
  \Phi(T)^*\Phi(T)\leq k\Phi(T^*T).
\]
Thus \(\Phi(T)\) tends ultra-strongly to zero, which proves that
\(\Phi\) is ultra-strongly continuous.

Finally, if \(\mathcal B\) is abelian, we shall show that, for
\(S,T\in\mathcal A\),
\[
  \Phi(T^*S)\Phi(T^*S)^*
    \leq\Phi(T^*T)\Phi(S^*S).
\]
It is enough to prove that, for every character \(\chi\) of
\(\mathcal B\),
\[
  \chi(\Phi(T^*S))\chi(\Phi(T^*S)^*)
    \leq\chi(\Phi(T^*T))\chi(\Phi(S^*S)).
\]
But this is precisely the Cauchy-Schwarz inequality, since
\(\chi\circ\Phi\) is a positive linear functional on \(\mathcal A\).
\end{proof}
% Source PDF page 061 (printed p. 54)
\phantomsection\label{srcpage:54}
\begin{corollary}\label{cor:I-4-2-1}\label{cor:I-4-1}
Let \(\mathcal A\) and \(\mathcal B\) be two von Neumann algebras, and
let \(\Phi\) be an isomorphism of \(\mathcal A\) onto \(\mathcal B\).
Then \(\Phi\) is a homeomorphism for the ultra-weak and ultra-strong
topologies. The restriction of \(\Phi\) to bounded subsets is a
homeomorphism for the weak and strong topologies.
\end{corollary}

\begin{proof}
Since \(\Phi\) is an isomorphism for the order structures, \(\Phi\) is
positive and normal. Moreover,
\(\Phi(T)^*\Phi(T)=\Phi(T^*T)\) for every \(T\in\mathcal A\).
\end{proof}

\begin{corollary}\label{cor:I-4-2}\label{cor:I-4-2-2}
Let \(\mathcal A\) and \(\mathcal B\) be von Neumann algebras, and let
\(\Phi\) be a normal homomorphism or antihomomorphism of \(\mathcal A\)
into \(\mathcal B\) such that \(\Phi(1)=1\). Then \(\Phi(\mathcal A)\)
is a von Neumann algebra. The algebra \(\mathcal A\) can be identified
with a product \(\mathcal A_1\times\mathcal A_2\) of two von Neumann
algebras, with \(\Phi\) vanishing on \(\mathcal A_1\) and being
injective on \(\mathcal A_2\). If \(\mathfrak M\) is a subset of
\(\mathcal A\), and if \(\mathcal C\) is the von Neumann algebra
generated by \(\mathfrak M\), then \(\Phi(\mathcal C)\) is the von
Neumann algebra generated by \(\Phi(\mathfrak M)\).
\end{corollary}

\begin{proof}
By \cref{thm:I-4-2}, \(\Phi\) is ultra-weakly continuous. The kernel
of \(\Phi\) is an ultra-weakly closed two-sided ideal of \(\mathcal A\).
Then, by \hyperref[sec:I-3]{\S~3}, \hyperref[thm:I-3-2]{Theorem~2},
\hyperref[cor:I-3-3]{Corollary~3}, \(\mathcal A\) can be identified
with a product \(\mathcal A_1\times\mathcal A_2\) of von Neumann
algebras, with \(\Phi\) vanishing on \(\mathcal A_1\) and being
injective on \(\mathcal A_2\). We are therefore reduced to the case
where \(\Phi\) is injective and hence isometric
[\hyperref[prop:I-1-8]{\S~1, Proposition~8~(iii)}].
Now let \(\mathcal D\) be a von Neumann algebra contained in
\(\mathcal A\). The unit ball of \(\Phi(\mathcal D)\), the image under
\(\Phi\) of the unit ball of \(\mathcal D\), is ultra-weakly compact.
Hence, by \hyperref[sec:I-3]{\S~3},
\hyperref[thm:I-3-2]{Theorem~2}, \(\Phi(\mathcal D)\) is a von Neumann
algebra. This proves the other assertions of the corollary.
\end{proof}

We saw in \hyperref[no:I-4-1]{no.~1} the close relations that exist
between positive linear functionals and homomorphisms. Let now
\(\mathcal A\) and \(\mathcal B\) be von Neumann algebras on the
complex Hilbert spaces \(\mathfrak H\) and \(\mathfrak K\), let
\(\Phi\) be a normal homomorphism of \(\mathcal A\) onto
\(\mathcal B\), and let \(x\in\mathfrak K\). Then the positive linear
functional
\[
  T\longmapsto(\Phi(T)x\mid x)
\]
on \(\mathcal A\) is normal. Conversely:

\begin{proposition}\label{prop:I-4-1}
Let \(\mathcal A\) be a von Neumann algebra, let \(\varphi\) be a
normal positive linear functional on \(\mathcal A\), and let \(\Phi\)
be the canonical homomorphism defined by \(\varphi\)
(\hyperref[lem:I-4-4]{Lemma~4}). Then \(\Phi\) is normal, and
\(\Phi(\mathcal A)\) is a von Neumann algebra.
\end{proposition}

\begin{proof}
Let \(\mathfrak K\) be the Hilbert space on which
\(\Phi(\mathcal A)\) acts, and let \(\Gamma\) be the canonical map of
\(\mathcal A\) into \(\mathfrak K\). Let \(\mathcal F\) be an upward
directed subset of \(\mathcal A^+\), with supremum
\(T\in\mathcal A^+\). Then \(\Phi(\mathcal F)\) is an upward directed
subset of \(\Phi(\mathcal A)\), bounded above by \(\Phi(T)\).
Moreover, for every \(S\in\mathcal A\),
\[
  \bigl(\Phi(T)\Gamma(S)\mid\Gamma(S)\bigr)
  =\varphi(S^*TS)
\]
is the supremum of
\[
  \bigl(\Phi(\mathcal F)\Gamma(S)\mid\Gamma(S)\bigr)
  =\varphi(S^*\mathcal F S),
\]
by the normality of \(\varphi\). Since the \(\Gamma(S)\) are dense
throughout \(\mathfrak K\), \(\Phi(T)\) is the supremum of
\(\Phi(\mathcal F)\)
% Source PDF page 062 (printed p. 55)
\phantomsection\label{srcpage:55}
in \(\mathcal L(\mathfrak K)\). Thus \(\Phi\) is normal, and
\(\Phi(\mathcal A)\) is a von Neumann algebra by
\cref{cor:I-4-2}.
\end{proof}

A homomorphism of a von Neumann algebra onto a von Neumann algebra is
not always normal (cf. \hyperref[sec:I-8]{\S~8}, Exercise~5).

\noindent\textbf{Bibliography:} \ArticleRef{7}, \ArticleRef{15},
\ArticleRef{19}, \ArticleRef{31}, \ArticleRef{79}, \ArticleRef{89},
\ArticleRef{100}, \ArticleRef{136}.

\NumberedParagraph{4}{Structure of normal homomorphisms}\label{no:I-4-4}
We already know three particular types of normal homomorphisms between
von Neumann algebras: spatial isomorphisms
(\hyperref[no:I-1-5]{\S~1, no.~5}), inductions
(\hyperref[no:I-2-1]{\S~2, no.~1}), and ampliations
(\hyperref[no:I-2-4]{\S~2, no.~4}). This being so, the following
theorem gives the general form of normal homomorphisms, and therefore,
in particular, of isomorphisms:

\begin{theorem}\label{thm:I-4-3}
Let \(\mathcal A\) and \(\mathcal B\) be two von Neumann algebras, and
let \(\Phi\) be a normal homomorphism of \(\mathcal A\) onto
\(\mathcal B\). There exist an ampliation \(\Phi_1\) of \(\mathcal A\)
onto a von Neumann algebra \(\mathcal C\), an induction \(\Phi_2\) of
\(\mathcal C\) onto a von Neumann algebra \(\mathcal D\), and a spatial
isomorphism \(\Phi_3\) of \(\mathcal D\) onto \(\mathcal B\), such that
\[
  \Phi=\Phi_3\circ\Phi_2\circ\Phi_1.
\]
\end{theorem}

\begin{proof}
Let \(\mathfrak H\) and \(\mathfrak K\) be the Hilbert spaces on which
\(\mathcal A\) and \(\mathcal B\) act. Suppose first that there exists
a vector \(y\in\mathfrak K\) totalizing for \(\mathcal B\). Put, for
\(T\in\mathcal A\),
\[
  \varphi(T)=(\Phi(T)y\mid y).
\]
{\emergencystretch=1em
The positive linear functional \(\varphi\) on \(\mathcal A\) is
normal. Hence, by \hyperref[thm:I-4-1]{Theorem~1} and
\hyperref[lem:I-4-6]{Lemma~6}, there exist a Hilbert space
\(\mathfrak H'\), a von Neumann algebra \(\mathcal C\) on
\(\mathfrak H'\), an ampliation \(\Phi_1\) of \(\mathcal A\) onto
\(\mathcal C\), and a vector \(x\in\mathfrak H'\), such that\par}
\[
  \varphi(T)=(\Phi_1(T)x\mid x).
\]
Let \(\mathfrak H''=\mathfrak X_x^{\mathcal C}\). The map
\(S\mapsto S_{\mathfrak H''}\) is an induction \(\Phi_2\) of
\(\mathcal C\) onto the von Neumann algebra
\(\mathcal D=\mathcal C_{\mathfrak H''}\), and
\[
  (\Phi(T)y\mid y)
  =\varphi(T)
  =\bigl((\Phi_2\circ\Phi_1)(T)x\mid x\bigr).
\]
Thus \(\Phi_2\circ\Phi_1\) is a homomorphism of \(\mathcal A\) onto
\(\mathcal D\), \(\Phi\) is a homomorphism of \(\mathcal A\) onto
\(\mathcal B\),
\[
  \mathfrak X_x^{\mathcal D}=\mathfrak H'',\qquad
  \mathfrak X_y^{\mathcal B}=\mathfrak K,
  \qquad
  (\Phi(T)y\mid y)
   =\bigl((\Phi_2\circ\Phi_1)(T)x\mid x\bigr).
\]
By \hyperref[lem:I-4-3]{Lemma~3}, there is an isomorphism of
\(\mathfrak H''\) onto \(\mathfrak K\) which transforms
\(\Phi_2\circ\Phi_1\) into \(\Phi\). This isomorphism defines a spatial
isomorphism \(\Phi_3\) of \(\mathcal D\) onto \(\mathcal B\), and
\(\Phi=\Phi_3\circ\Phi_2\circ\Phi_1\).

Pass to the general case. Consider a maximal family
\((y_i)_{i\in I}\) of nonzero vectors of \(\mathfrak K\) such that the
closed vector subspaces
\(\mathfrak K_i=\mathfrak X_{y_i}^{\mathcal B}\) are pairwise
orthogonal. By the maximality of the family \((y_i)\), one has
\[
  \mathfrak K=\bigoplus_{i\in I}\mathfrak K_i.
\]
Let \(\Phi^i\) be the normal homomorphism
\[
  T\longmapsto\Phi(T)_{\mathfrak K_i}
\]
of \(\mathcal A\) onto \(\mathcal B_{\mathfrak K_i}\). By the first
part of the proof, for every \(i\in I\) one may define the following:
a family \((\mathfrak H_\lambda^i)_{\lambda\in I_i}\) of Hilbert
spaces identifiable with \(\mathfrak H\), with Hilbert
% Source PDF page 063 (printed p. 56)
\phantomsection\label{srcpage:56}
sum \(\mathfrak H_i'\); an ampliation \(\Phi_1^i\) of
\(\mathcal A\) onto a von Neumann algebra \(\mathcal C_i\) on
\(\mathfrak H_i'\); a closed vector subspace \(\mathfrak H_i''\) of
\(\mathfrak H_i'\), invariant under \(\mathcal C_i\); the induction
\(\Phi_2^i\) of \(\mathcal C_i\) onto
\(\mathcal D_i=(\mathcal C_i)_{\mathfrak H_i''}\); and a spatial
isomorphism \(\Phi_3^i\) of \(\mathcal D_i\) onto
\(\mathcal B_{\mathfrak K_i}\), such that
\[
  \Phi^i=\Phi_3^i\circ\Phi_2^i\circ\Phi_1^i.
\]
Let \(\mathfrak H'\) be the Hilbert sum of the
\(\mathfrak H_i'\), for \(i\in I\), that is, of the
\(\mathfrak H_\lambda^i\), for \(\lambda\in I_i\) and \(i\in I\);
the Hilbert sum \(\mathfrak H''\) of the \(\mathfrak H_i''\) is a
closed vector subspace of \(\mathfrak H'\). For every
\(T\in\mathcal A\), let \(\Phi_1(T)\) be the element of
\(\mathcal L(\mathfrak H')\) which induces \(\Phi_1^i(T)\) on
\(\mathfrak H_i'\), for every \(i\in I\); it is clear that \(\Phi_1\)
is an ampliation of \(\mathcal A\) onto a von Neumann algebra
\(\mathcal C\) on \(\mathfrak H'\). The subspace \(\mathfrak H''\) is
invariant under \(\mathcal C\); let \(\Phi_2\) be the induction of
\(\mathcal C\) onto \(\mathcal D=\mathcal C_{\mathfrak H''}\). The
isomorphisms of \(\mathfrak H_i''\) onto \(\mathfrak K_i\) define an
isomorphism of
\[
  \mathfrak H''=\bigoplus_{i\in I}\mathfrak H_i''
  \quad\hbox{onto}\quad
  \mathfrak K=\bigoplus_{i\in I}\mathfrak K_i,
\]
and hence a spatial isomorphism \(\Phi_3\) of \(\mathcal D\) onto
\(\mathcal B\). We have
\(\Phi=\Phi_3\circ\Phi_2\circ\Phi_1\).
\end{proof}

\begin{corollary*}\phantomsection\label{cor:I-4-3}
Let \(\mathcal A\) and \(\mathcal B\) be two von Neumann algebras, and
let \(\Phi\) be an isomorphism of \(\mathcal A\) onto \(\mathcal B\).
There exist a von Neumann algebra \(\mathcal C\) and two projections
\(E',F'\) in \(\mathcal C'\), of central support \(1\), such that one
may identify \(\mathcal A\) with \(\mathcal C_{E'}\),
\(\mathcal B\) with \(\mathcal C_{F'}\), and \(\Phi\) with the
isomorphism
\[
  T_{E'}\longmapsto T_{F'}\qquad(T\in\mathcal C).
\]
\end{corollary*}

\begin{proof}
With the preceding notation, \(\Phi_2\) is an injective induction
\(T\mapsto T_{F'}\) of \(\mathcal C\) onto
\(\mathcal C_{F'}\), and \(\Phi_1^{-1}\) can be identified with an
injective induction \(T\mapsto T_{E'}\) of \(\mathcal C\) onto
\(\mathcal C_{E'}\).
\end{proof}

\begin{remark*}
Suppose that \(\mathcal B'\) is of countable type. Then the ampliation
of \(\mathcal A\) onto \(\mathcal C\) may be supposed to be of the
form
\[
  T\longmapsto T\otimes I_{\widetilde{\mathfrak H}},
\]
where \(\widetilde{\mathfrak H}\) is a Hilbert space with a countable
base. Indeed, \(I\) is countable, since the
\(P_{\mathfrak K_i}\) are nonzero orthogonal projections in
\(\mathcal B'\). And the \(I_i\), for \(i\in I\), may be supposed
countable by \hyperref[lem:I-4-6]{Lemma~6}.
\end{remark*}

\noindent\textbf{Bibliography:} \ArticleRef{17}, \ArticleRef{31},
\ArticleRef{89}.

\NumberedParagraph{5}{Application: Isomorphisms of tensor products}
\label{no:I-4-5}

\begin{proposition}\label{prop:I-4-2}
Let \(\mathcal A_1,\mathcal A_2,\mathcal B_1,\mathcal B_2\) be von
Neumann algebras, let \(\Phi_1\) be a normal homomorphism of
\(\mathcal A_1\) onto \(\mathcal B_1\), and let \(\Phi_2\) be a
normal homomorphism of \(\mathcal A_2\) onto \(\mathcal B_2\). There
exists one and only one normal homomorphism \(\Phi\) of
\(\mathcal A_1\otimes\mathcal A_2\) onto
\(\mathcal B_1\otimes\mathcal B_2\) such that
\[
  \Phi(T_1\otimes T_2)=\Phi_1(T_1)\otimes\Phi_2(T_2)
  \quad
  (T_1\in\mathcal A_1,\ T_2\in\mathcal A_2).
\]
If \(\Phi_1\) and \(\Phi_2\) are isomorphisms (respectively, spatial
isomorphisms), \(\Phi\) is an isomorphism (respectively, a spatial
isomorphism).
\end{proposition}

\begin{proof}
The uniqueness of \(\Phi\) is immediate, since \(\Phi\) must be
ultra-weakly continuous and the involutive algebra generated by the
operators \(T_1\otimes T_2\), where \(T_1\in\mathcal A_1\) and
\(T_2\in\mathcal A_2\), is ultra-weakly dense throughout
\(\mathcal A_1\otimes\mathcal A_2\). To prove the existence of
\(\Phi\), it suffices, by \cref{thm:I-4-3}, to consider the following
two cases: \textit{a}. \(\Phi_1\) is the identity map of
\(\mathcal A_1\), and \(\Phi_2\) is the ampliation
\[
  T_2\longmapsto T_2\otimes I_{\mathfrak K}
\]
of \(\mathcal A_2\) onto
\(\mathcal A_2\otimes\mathbb C_{\mathfrak K}\), with
% Source PDF page 064 (printed p. 57)
\phantomsection\label{srcpage:57}
\(\mathfrak K\) a complex Hilbert space; \textit{b}. \(\Phi_1\) is
the identity map of \(\mathcal A_1\), and \(\Phi_2\) is the induction
\[
  T_2\longmapsto(T_2)_{E'},
\]
where \(E'\) is a projection in \(\mathcal A_2'\).

\textit{Case a}. Let \(\Phi\) be the ampliation of
\(\mathcal A_1\otimes\mathcal A_2\) onto
\[
  (\mathcal A_1\otimes\mathcal A_2)\otimes
      \mathbb C_{\mathfrak K}
  =\mathcal A_1\otimes
      (\mathcal A_2\otimes\mathbb C_{\mathfrak K})
\]
(associativity of the tensor product). For
\(T_1\in\mathcal A_1\) and \(T_2\in\mathcal A_2\), one has
\[
  \Phi(T_1\otimes T_2)
  =(T_1\otimes T_2)\otimes I_{\mathfrak K}
  =T_1\otimes T_2\otimes I_{\mathfrak K}
  =T_1\otimes\Phi_2(T_2).
\]

\textit{Case b}. Let \(\Phi\) be the induction of
\(\mathcal A_1\otimes\mathcal A_2\) onto
\[
  (\mathcal A_1\otimes\mathcal A_2)_{I\otimes E'}
  =\mathcal A_1\otimes(\mathcal A_2)_{E'}.
\]
For \(T_1\in\mathcal A_1\) and \(T_2\in\mathcal A_2\), one has
\[
  \Phi(T_1\otimes T_2)
  =(T_1\otimes T_2)_{I\otimes E'}
  =T_1\otimes(T_2)_{E'}
  =T_1\otimes\Phi_2(T_2).
\]

Suppose that \(\Phi_1\) and \(\Phi_2\) are isomorphisms, and let us
prove that \(\Phi\) is an isomorphism. Here again, it suffices to
consider cases \textit{a} and \textit{b}. In case \textit{a},
\(\Phi\) is an ampliation, and hence is an isomorphism. In case
\textit{b}, \(E'\) has central support \(1\). Let
\(\mathfrak H_1,\mathfrak H_2\) be the Hilbert spaces on which
\(\mathcal A_1,\mathcal A_2\) act, and let
\(\mathfrak H_2'=E'(\mathfrak H_2)\). One has
\[
  \mathfrak X_{\mathfrak H_2'}^{\mathcal A_2'}
  =\mathfrak H_2
\]
(\hyperref[prop:I-1-7]{\S~1, Proposition~7, Corollary~1}), and hence
\[
  \mathfrak X_{\mathfrak H_1\otimes\mathfrak H_2'}
    ^{\mathbb C_{\mathfrak H_1}\otimes\mathcal A_2'}
  =\mathfrak H_1\otimes\mathfrak H_2,
\]
so that \(I\otimes E'\), regarded as a projection of
\((\mathcal A_1\otimes\mathcal A_2)'\), has central support \(1\).
Consequently \(\Phi\) is an isomorphism. Finally, it is clear that if
\(\Phi_1\) and \(\Phi_2\) are spatial isomorphisms, \(\Phi\) is a
spatial isomorphism.
\end{proof}

\noindent\textbf{Bibliography:} \ArticleRef{67}, \ArticleRef{113},
\ArticleRef{128}, \ArticleRef{133}.

\NumberedParagraph{6}{Support of a normal positive linear functional}
\label{no:I-4-6}

\begin{proposition}\label{prop:I-4-3}
Let \(\mathcal A\) be a von Neumann algebra, and let \(\varphi\) be a
normal positive linear functional on \(\mathcal A\). Among the
projections \(G\in\mathcal A\) such that \(\varphi(G)=0\), there is
one, say \(F\), greater than all the others. One has
\[
  \varphi(TF)=\varphi(FT)=0
  \qquad\text{for every }T\in\mathcal A.
\]
\end{proposition}

\begin{proof}
Let \(\mathfrak m\) be the set of \(T\in\mathcal A\) such that
\(\varphi(T^*T)=0\). By the Cauchy--Schwarz inequality,
\(\mathfrak m\) is also the set of \(T\in\mathcal A\) such that
\(\varphi(S^*T)=0\) for every \(S\in\mathcal A\). Thus
\(\mathfrak m\) is an ultra-weakly closed left ideal, by
\hyperref[thm:I-4-1]{Theorem~1}. The existence of \(F\) then follows
from \hyperref[sec:I-3]{\S~3},
\hyperref[thm:I-3-2]{Theorem~2, Corollary~3}. Finally, for
\(T\in\mathcal A\),
\[
  |\varphi(TF)|^2\leq\varphi(TT^*)\varphi(F)=0,
  \qquad
  |\varphi(FT)|^2\leq\varphi(T^*T)\varphi(F)=0.
\]
\end{proof}

\begin{definition}\label{def:I-4-3}
Let \(F\) be the projection defined by
\hyperref[prop:I-4-3]{Proposition~2}. The projection \(1-F\) is called
the \emph{support} of \(\varphi\). Two normal positive functionals on
\(\mathcal A\) are said to be \emph{disjoint} if their supports are
orthogonal.
\end{definition}

% Source PDF page 065 (printed p. 58)
\phantomsection\label{srcpage:58}
Denote this support by \(E_\varphi\). One has
\[
  \varphi(T)=\varphi(E_\varphi T E_\varphi)
  \qquad(T\in\mathcal A),
\]
and \(\varphi\) defines a faithful normal positive linear functional
on \(\mathcal A_{E_\varphi}\). In particular, to say that
\(\varphi\) is faithful amounts to saying that \(E_\varphi=1\).

Let
\[
  \varphi=\sum_{i=1}^{\infty}\omega_{x_i},
  \qquad
  \sum_{i=1}^{\infty}\lVert x_i\rVert^2<+\infty.
\]
A projection \(G\) of \(\mathcal A\) satisfies \(\varphi(G)=0\) if
and only if \(Gx_i=0\) for every \(i\); hence the support of
\(\varphi\) is \(E_{\mathfrak M}^{\mathcal A'}\), where
\(\mathfrak M\) denotes the set of the \(x_i\). In particular,
\(\varphi\) is faithful if and only if \(\mathfrak M\) is separating
for \(\mathcal A\).

For a von Neumann algebra \(\mathcal A\) to be of countable type, it
is necessary and sufficient that there exist on \(\mathcal A\) a
faithful normal positive linear functional \(\varphi\). The condition
is plainly sufficient. Conversely, if \(\mathcal A\) is of countable
type, there exists
(\hyperref[prop:I-1-6]{\S~1, Proposition~6}) a separating sequence
\((x_i)\) for \(\mathcal A\); one may suppose that
\[
  \sum_{i=1}^{\infty}\lVert x_i\rVert^2<+\infty,
\]
and then take
\[
  \varphi=\sum_{i=1}^{\infty}\omega_{x_i}.
\]

\begin{proposition}\label{prop:I-4-4}
Let \(\mathcal A\) be a von Neumann algebra, let
\(\mathcal A_1\) be its unit ball, let \(\varphi\) be a normal positive
linear functional on \(\mathcal A\), and let \(E_\varphi\) be its
support. On \(\mathcal A_1\), convergence of \(\varphi(T^*T)\) to zero
is equivalent to strong convergence of \(TE_\varphi\) to zero.
\end{proposition}

\begin{proof}
The functional \(\varphi\) canonically defines a Hilbert space
\(\mathfrak K\), a map \(\Gamma\) of \(\mathcal A\) into
\(\mathfrak K\), and a homomorphism \(\Phi\) of \(\mathcal A\) onto a
von Neumann algebra \(\mathcal B=\Phi(\mathcal A)\) on
\(\mathfrak K\) (\hyperref[lem:I-4-4]{Lemma~4},
\cref{prop:I-4-1}). Put \(x=\Gamma(1)\). Suppose first that \(\Phi\)
is an isomorphism. Let
\[
  F=E_x^{\mathcal B'},\qquad F_1=1-F,
\]
and let \(E_1\) be the projection of \(\mathcal A\) such that
\(F_1=\Phi(E_1)\). Then \(F_1\) is the greatest projection of
\(\mathcal B\) such that \(F_1x=0\), or such that
\((F_1x\mid x)=0\). Hence \(E_1\) is the greatest projection of
\(\mathcal A\) such that
\[
  0=(\Phi(E_1)x\mid x)=\varphi(E_1).
\]
Thus \(E_\varphi=1-E_1\), and \(\Phi(E_\varphi)=F\). This being so,
for \(TE_\varphi\) to tend strongly to zero, it is necessary and
sufficient, for \(T\in\mathcal A_1\), that
\(\Phi(TE_\varphi)=\Phi(T)F\) tend strongly to zero
(\cref{cor:I-4-1}), and hence that
\[
  \lVert\Phi(T)T'x\rVert\longrightarrow0
  \qquad\text{for every }T'\in\mathcal B',
\]
and hence that
\[
  \lVert\Phi(T)x\rVert^2
  =(\Phi(T^*T)x\mid x)
  =\varphi(T^*T)
\]
tend to zero.

Pass to the general case. The kernel of \(\Phi\) is an ultra-weakly
closed two-sided ideal of \(\mathcal A\); by
\hyperref[sec:I-3]{\S~3},
\hyperref[thm:I-3-2]{Theorem~2, Corollary~3}, \(\mathcal A\) can be
identified with the product \(\mathcal C\times\mathcal D\) of two von
Neumann algebras, with \(\varphi\) and \(\Phi\) inducing zero on
\(\mathcal C\), while the restriction of \(\Phi\) to \(\mathcal D\)
is an isomorphism. It then suffices to apply the result of the first
part of the proof.
\end{proof}

% Source PDF page 066 (printed p. 59)
\phantomsection\label{srcpage:59}
\begin{proposition}\label{prop:I-4-5}
Let \(\mathcal A\) be a von Neumann algebra, and let \(\varphi\) and
\(\psi\) be two normal positive linear functionals on \(\mathcal A\),
with supports \(E_\varphi\) and \(E_\psi\). The following conditions
are equivalent:
\begin{enumerate}[label=\textup{(\roman*)}]
\item \(\varphi(T)=0\) and \(T\in\mathcal A^+\) imply
      \(\psi(T)=0\);
\item \(E_\psi\leq E_\varphi\);
\item on the unit ball \(\mathcal A_1\) of \(\mathcal A\), the
      topology defined by the seminorm
      \(\bigl[\varphi(T^*T)\bigr]^{1/2}\) is finer than the topology
      defined by the seminorm
      \(\bigl[\psi(T^*T)\bigr]^{1/2}\).
\end{enumerate}
\end{proposition}

\begin{proof}
The implication \textup{(ii)} \(\Rightarrow\) \textup{(iii)} follows
from \cref{prop:I-4-4}. The implication
\textup{(iii)} \(\Rightarrow\) \textup{(i)} is immediate. Finally,
suppose that condition \textup{(i)} holds. One has
\(\varphi(1-E_\varphi)=0\), hence
\(\psi(1-E_\varphi)=0\), and therefore
\[
  1-E_\varphi\leq1-E_\psi,
\]
whence \(E_\psi\leq E_\varphi\).
\end{proof}

\noindent\textbf{Bibliography:} \ArticleRef{19}.

\NumberedParagraph{7}{Polar decomposition of a linear functional}
\label{no:I-4-7}
Let \(A\) be an algebra. If \(f\) is a linear functional on \(A\) and
if \(x\in A\), the linear functionals \(x\mathbin{\cdot}f\) and
\(f\mathbin{\cdot}x\) on \(A\) are defined by
\[
  (x\mathbin{\cdot}f)(y)=f(yx),
  \qquad
  (f\mathbin{\cdot}x)(y)=f(xy),
\]
for every \(y\in A\). One has
\[
  x_1\mathbin{\cdot}(x_2\mathbin{\cdot}f)
    =(x_1x_2)\mathbin{\cdot}f,
  \qquad
  (f\mathbin{\cdot}x_1)\mathbin{\cdot}x_2
    =f\mathbin{\cdot}(x_1x_2).
\]
If \(A\) is an involutive algebra and \(f\) is a linear functional on
\(A\), define the linear functional \(f^*\) by
\[
  f^*(y)=\overline{f(y^*)}.
\]
One has
\[
  (x\mathbin{\cdot}f)^*
  =f^*\mathbin{\cdot}x^*.
\]
The functional \(f\) is said to be Hermitian if \(f=f^*\).

If \(A\) is a normed algebra and \(f\) is continuous, then
\(x\mathbin{\cdot}f\) and \(f\mathbin{\cdot}x\) are continuous, and
\[
  \lVert x\mathbin{\cdot}f\rVert
     \leq\lVert x\rVert\,\lVert f\rVert,
  \qquad
  \lVert f\mathbin{\cdot}x\rVert
     \leq\lVert f\rVert\,\lVert x\rVert.
\]
If \(A\) is a von Neumann algebra and \(f\) is ultra-weakly
continuous, then \(x\mathbin{\cdot}f\),
\(f\mathbin{\cdot}x\), and \(f^*\) are ultra-weakly continuous.

\begin{lemma}\label{lem:I-4-8}
Let \(\mathfrak H\) be a Hilbert space, let \(\mathcal A\) be a
\(C^*\)-algebra of operators on \(\mathfrak H\) containing \(1\), let
\(\mathcal A_1\) be the unit ball of \(\mathcal A\), and let \(T\) be
an extreme point of \(\mathcal A_1\). Then \(T^*T\) is a projection.
\end{lemma}

\begin{proof}
One has \(\lVert T\rVert=1\), and hence
\(\lVert T^*T\rVert=1\). Let \(\mathcal B\) be the \(C^*\)-algebra of
operators generated by \(1\) and \(T^*T\); it is commutative. Let
\(\Omega\) be its spectrum, which is compact (Appendix~I). Let
\(\mathcal B'\) be the involutive Banach algebra formed by the
continuous complex-valued functions on \(\Omega\), and let
\(\sigma\) be the Gelfand isomorphism of \(\mathcal B'\) onto
\(\mathcal B\). Put
\[
  f=\sigma^{-1}\bigl((T^*T)^{1/2}\bigr),
  \qquad g=f^2.
\]
One has \(f=\overline f\) and \(0\leq g\leq1\). Suppose that there is
an \(\omega\in\Omega\) such that \(0<g(\omega)<1\). There then exists
a continuous function \(h\colon\Omega\to\mathbb R\), with \(h\geq0\),
such that
\[
  gh\neq0,\qquad
  \sup g(1+h)^2=\sup g(1-h)^2=1
\]
(it suffices to take the support of \(h\) in a sufficiently small
neighborhood of \(\omega\), and \(h\) sufficiently small). Let
\(U=\sigma(h)\in\mathcal B\). One has \(U=U^*\) and
% Source PDF page 067 (printed p. 60)
\phantomsection\label{srcpage:60}
\[
  \lVert T^*T(1+U)^2\rVert
  =\lVert T^*T(1-U)^2\rVert=1.
\]
Thus \(T(1+U)\in\mathcal A_1\) and
\(T(1-U)\in\mathcal A_1\). The equality
\[
  T=\tfrac12\bigl(T(1+U)+T(1-U)\bigr)
\]
implies, since \(T\) is extreme, that
\[
  T=T(1+U)=T(1-U),
\]
whence \(TU=0\) and \(gh=0\), which is absurd. Thus \(g\) can take
only the values \(0\) and \(1\), so that \(T^*T\) is a projection.
\end{proof}

\begin{lemma}\label{lem:I-4-9}
Let \(\mathcal A\) be a von Neumann algebra, let \(E\) be a projection
of \(\mathcal A\), and let \(f\) be an element of the predual of
\(\mathcal A\).
\begin{enumerate}[label=\textup{(\roman*)}]
\item
\[
  \lVert f\rVert^2
  \geq
  \lVert f\mathbin{\cdot}E\rVert^2
  +\lVert f\mathbin{\cdot}(1-E)\rVert^2.
\]
\item If \(\lVert f\rVert=\lVert f\mathbin{\cdot}E\rVert\), then
      \(f=f\mathbin{\cdot}E\).
\end{enumerate}
\end{lemma}

\begin{proof}
Assertion \textup{(ii)} follows from \textup{(i)}. Let us prove
\textup{(i)}. Let \(\mathfrak H\) be the Hilbert space on which
\(\mathcal A\) acts. Let \(\mathcal L_*\) be the predual of
\(\mathcal L(\mathfrak H)\), let \(\mathcal A^\perp\) be the
annihilator of \(\mathcal A\) in \(\mathcal L_*\), and let
\(\varepsilon>0\). As was seen in the proof of
\hyperref[thm:I-3-1]{\S~3, Theorem~1}, the predual of \(\mathcal A\)
can be identified with
\(\mathcal L_*/\mathcal A^\perp\). Hence there exists
\(g\in\mathcal L_*\), extending \(f\), such that
\(\lVert g\rVert\leq\lVert f\rVert+\varepsilon\). If we prove that
\[
  \lVert g\rVert^2
  \geq
  \lVert g\mathbin{\cdot}E\rVert^2
  +\lVert g\mathbin{\cdot}(1-E)\rVert^2,
\]
it will follow that
\[
  (\lVert f\rVert+\varepsilon)^2
  \geq
  \lVert f\mathbin{\cdot}E\rVert^2
  +\lVert f\mathbin{\cdot}(1-E)\rVert^2,
\]
and hence \textup{(i)}, by the arbitrariness of \(\varepsilon\). We
are therefore reduced to the case
\(\mathcal A=\mathcal L(\mathfrak H)\). Since the weakly continuous
linear functionals on \(\mathcal L(\mathfrak H)\) are norm-dense in
\(\mathcal L_*\), we may suppose that \(f\) is weakly continuous.
There then exist two orthonormal systems
\[
  (e_1,\ldots,e_n),\qquad(e_1',\ldots,e_n')
\]
in \(\mathfrak H\), and numbers
\(\lambda_1\geq0,\ldots,\lambda_n\geq0\), such that
\[
  f=\lambda_1\omega_{e_1,e_1'}+\cdots+
      \lambda_n\omega_{e_n,e_n'},
  \qquad
  \lVert f\rVert=\lambda_1+\cdots+\lambda_n
\]
(\hyperref[lem:I-3-3]{\S~3, proof of Lemma~3}). Thus, for every
\(T\in\mathcal L(\mathfrak H)\),
\[
  (f\mathbin{\cdot}E)(T)
    =\sum_i\lambda_i(Te_i\mid Ee_i'),
  \qquad
  (f\mathbin{\cdot}(1-E))(T)
    =\sum_i\lambda_i(Te_i\mid(1-E)e_i'),
\]
whence
\begin{align*}
 &\lVert f\mathbin{\cdot}E\rVert^2
  +\lVert f\mathbin{\cdot}(1-E)\rVert^2\\
 &\quad\leq
  \left(\sum_i\lambda_i\lVert Ee_i'\rVert\right)^2
  +\left(\sum_i\lambda_i\lVert(1-E)e_i'\rVert\right)^2\\
 &\quad=
  \sum_{i=1}^n\lambda_i^2
       \bigl(\lVert Ee_i'\rVert^2+
             \lVert(1-E)e_i'\rVert^2\bigr)\\
 &\qquad\quad
  +2\sum_{1\leq i<j\leq n}\lambda_i\lambda_j
       \bigl(\lVert Ee_i'\rVert\,\lVert Ee_j'\rVert+
             \lVert(1-E)e_i'\rVert\,
             \lVert(1-E)e_j'\rVert\bigr)\\
 &\quad\leq
  \sum_{i=1}^n\lambda_i^2+
  2\sum_{1\leq i<j\leq n}\lambda_i\lambda_j
  =\left(\sum_{i=1}^n\lambda_i\right)^2
  =\lVert f\rVert^2.
\end{align*}
\end{proof}

% Source PDF page 068 (printed p. 61)
\phantomsection\label{srcpage:61}
\begin{lemma}\label{lem:I-4-10}
Let \(\mathfrak H\) be a Hilbert space, let \(\mathcal A\) be an
involutive algebra of operators on \(\mathfrak H\) containing \(1\),
let \(f\) be a continuous linear functional on \(\mathcal A\), let
\(\varphi\) be a positive functional on \(\mathcal A\), and let
\(S\in\mathcal A\) be such that
\[
  \lVert S\rVert\leq1,\qquad
  \varphi\mathbin{\cdot}S=f,\qquad
  \lVert\varphi\rVert=\lVert f\rVert.
\]
If \(T\in\mathcal A\) satisfies
\[
  \lVert T\rVert\leq1,\qquad
  f\mathbin{\cdot}T\geq0,\qquad
  \lVert f\mathbin{\cdot}T\rVert=\lVert f\rVert,
\]
then \(f\mathbin{\cdot}T=\varphi\).
\end{lemma}

\begin{proof}
We may suppose that
\(\lVert\varphi\rVert=\lVert f\rVert=1\). Apply
\hyperref[lem:I-4-4]{Lemma~4} to
\(\mathcal A\) and \(\varphi\), obtaining a Hilbert space
\(\mathfrak K\), a vector \(x\in\mathfrak K\), and a homomorphism
\(\Phi\) of \(\mathcal A\) into \(\mathcal L(\mathfrak K)\). One has
\[
  (x\mid\Phi(T^*S^*)x)
  =(\Phi(ST)x\mid x)
  =\varphi(ST)
  =f(T)
  =(f\mathbin{\cdot}T)(1)
  =\lVert f\mathbin{\cdot}T\rVert=1.
\]
Since \((x\mid x)=\varphi(1)=1\) and
\(\lVert T^*S^*\rVert\leq1\), it follows that
\(\Phi(T^*S^*)x=x\), by the equality case of the Cauchy--Schwarz
inequality. Therefore, for every \(R\in\mathcal A\),
\begin{align*}
  (f\mathbin{\cdot}T)(R)
   &=f(TR)=\varphi(STR)
     =(\Phi(STR)x\mid x)\\
   &=(\Phi(R)x\mid\Phi(T^*S^*)x)
     =(\Phi(R)x\mid x)=\varphi(R),
\end{align*}
whence \(f\mathbin{\cdot}T=\varphi\).
\end{proof}

\begin{theorem}\label{thm:I-4-4}
Let \(\mathcal A\) be a von Neumann algebra, and let \(f\) be an
element of the predual of \(\mathcal A\).
\begin{enumerate}[label=\textup{(\roman*)}]
\item There exists a pair \((\varphi,U)\) with the following
      properties:
  \begin{enumerate}[label=\textit{\alph*.}]
  \item \(\varphi\) is a normal positive functional on \(\mathcal A\),
        and \(\lVert\varphi\rVert=\lVert f\rVert\);
  \item \(U\) is a partially isometric element of \(\mathcal A\)
        whose final projection is the support of \(\varphi\);
  \item
  \[
    f=\varphi\mathbin{\cdot}U,
    \qquad
    \varphi=f\mathbin{\cdot}U^*.
  \]
  \end{enumerate}
\item Let \(\varphi'\) be a normal positive functional on
      \(\mathcal A\), and \(U'\) a partially isometric element of
      \(\mathcal A\) whose final projection is bounded above by the
      support of \(\varphi'\), with
      \[
        f=\varphi'\mathbin{\cdot}U',
        \qquad
        \varphi'=f\mathbin{\cdot}U'^*.
      \]
      Then \(\varphi'=\varphi\) and \(U'=U\).
\end{enumerate}
\end{theorem}

\begin{proof}
We may suppose that \(\lVert f\rVert=1\). Let \(\mathcal A_1\) be the
unit ball of \(\mathcal A\), and let \(\mathcal A_2\) be the set of
\(T\in\mathcal A_1\) such that \(f(T)=1\). Since
\(\mathcal A_1\) is ultra-weakly compact and \(f\) is ultra-weakly
continuous, there are \(T\in\mathcal A_1\) such that
\(\lvert f(T)\rvert=1\); multiplying \(T\) by a suitable scalar shows
that \(\mathcal A_2\neq\varnothing\). Moreover, \(\mathcal A_2\) is
convex and ultra-weakly compact; it therefore contains an extreme
point \(V\). This point is also extreme in \(\mathcal A_1\). Indeed,
if \(V=\tfrac12(S+T)\), with \(S,T\in\mathcal A_1\), then
\[
  |f(S)|\leq1,\qquad |f(T)|\leq1,
  \qquad
  1=f(V)=\tfrac12(f(S)+f(T)),
\]
whence \(f(S)=f(T)=1\), so that \(S,T\in\mathcal A_2\), and
\(S=T=V\), since \(V\) is extreme in \(\mathcal A_2\). By
\cref{lem:I-4-8}, \(V^*V\) is a projection, and hence \(V\) is
partially isometric.

Put \(\varphi=f\mathbin{\cdot}V\). One has
\(\lVert\varphi\rVert\leq\lVert f\rVert\lVert V\rVert\leq1\).
Since \(\varphi(1)=f(V)=1\), \(\varphi\) is positive
(\cref{lem:I-4-5}), has
norm \(1\), and is normal because it is ultra-weakly continuous. Let
\(E\) be the support of \(\varphi\). One has
\(\varphi(V^*V)=f(VV^*V)=f(V)=1\), and hence
% Source PDF page 069 (printed p. 62)
\phantomsection\label{srcpage:62}
\(V^*V\geq E\). Put \(U=EV^*\). Then
\[
  UU^*=E(V^*V)E=E,
\]
so \(U\) is partially isometric with final projection \(E\). For
every \(T\in\mathcal A\),
\[
  (f\mathbin{\cdot}U^*)(T)
  =f(U^*T)=f(VET)=\varphi(ET)=\varphi(T),
\]
and hence \(\varphi=f\mathbin{\cdot}U^*\). Put \(U^*U=F\). One has
\(\lVert f\mathbin{\cdot}F\rVert\leq1\), and
\[
  (f\mathbin{\cdot}F)(U^*)
  =f(U^*UU^*)=f(U^*)=f(VE)=\varphi(E)=1.
\]
Consequently
\(\lVert f\mathbin{\cdot}F\rVert=1=\lVert f\rVert\), and hence
\(f\mathbin{\cdot}F=f\), by \cref{lem:I-4-9}\textup{(ii)}. Thus, for every
\(T\in\mathcal A\),
\[
  (\varphi\mathbin{\cdot}U)(T)
  =\varphi(UT)=f(U^*UT)=f(FT)
  =(f\mathbin{\cdot}F)(T)=f(T),
\]
so that \(f=\varphi\mathbin{\cdot}U\). This proves \textup{(i)}.

Let \(\varphi',U'\) have the properties in \textup{(ii)}. One has
\(\lVert U'^*\rVert\leq1\), \(f\mathbin{\cdot}U'^*\geq0\), and
\[
  \lVert\varphi'\rVert
  =\lVert f\mathbin{\cdot}U'^*\rVert
  \leq\lVert f\rVert
  =\lVert\varphi'\mathbin{\cdot}U'\rVert
  \leq\lVert\varphi'\rVert.
\]
Thus
\(\lVert f\mathbin{\cdot}U'^*\rVert=\lVert f\rVert\). By
\cref{lem:I-4-10}, applied with \(S=U\) and \(T=U'^*\), one has
\[
  \varphi=f\mathbin{\cdot}U^*
  =f\mathbin{\cdot}U'^*=\varphi'.
\]
Let us prove that \(U=U'\). Put \(X=UU'^*\). Then
\[
  \varphi(X)
  =(\varphi\mathbin{\cdot}U)(U'^*)
  =f(U'^*)=(f\mathbin{\cdot}U'^*)(1)=\varphi(1)=1.
\]
Hence \(\varphi(X^*)=1\), and
\[
  1=\varphi(X)^2\leq\varphi(XX^*)\leq1;
\]
thus \(\varphi(XX^*)=1\), and
\[
  \varphi\bigl((E-X)(E-X)^*\bigr)=1-1-1+1=0.
\]
Now \(EXE=X\), since the final projections of \(U\) and \(U'\) are
bounded above by \(E\). Since \(\varphi\) is faithful on
\(E\mathcal A E\), it follows that \(X=E\). Let \(\mathfrak H\) be
the Hilbert space on which \(\mathcal A\) acts. For
\(x\in E(\mathfrak H)\), one has
\[
  \lVert U(U'^*x)\rVert
  =\lVert Ex\rVert=\lVert x\rVert,
  \qquad
  \lVert U(U'^*x)\rVert\geq\lVert U'^*x\rVert.
\]
Thus \(U'^*x\in F(\mathfrak H)\). Since \(U\) maps
\(F(\mathfrak H)\) isometrically onto \(E(\mathfrak H)\), the equality
\[
  UU'^*x=x=UU^*x
\]
implies \(U'^*x=U^*x\). Moreover, \(U^*\) and \(U'^*\) vanish on
\(\mathfrak H\ominus E(\mathfrak H)\), so \(U'^*=U^*\) and \(U'=U\).
\end{proof}

\begin{definition}\label{def:I-4-4}
With the notation of \cref{thm:I-4-4}, \(\varphi\) is called the
\emph{absolute value} of \(f\), and is denoted by \(\lvert f\rvert\).
The equality
\[
  f=\lvert f\rvert\mathbin{\cdot}U
\]
is called the \emph{polar decomposition} of \(f\).
\end{definition}

\begin{proposition}\label{prop:I-4-6}
Retain the notation of \cref{thm:I-4-4}. Then
\[
  \lvert f^*\rvert
  =U^*\mathbin{\cdot}\lvert f\rvert\mathbin{\cdot}U,
\]
and the polar decomposition of \(f^*\) is
\[
  f^*=\lvert f^*\rvert\mathbin{\cdot}U^*.
\]
\end{proposition}

\begin{proof}
Since the support of \(\lvert f\rvert\) is \(UU^*\), one has
\begin{equation}
  f^*=(\lvert f\rvert\mathbin{\cdot}U)^*
  =U^*\mathbin{\cdot}\lvert f\rvert
  =U^*\mathbin{\cdot}\lvert f\rvert\mathbin{\cdot}UU^*.
  \tag{1}\label{eq:I-4-prop6-1}
\end{equation}
It is clear that
\(U^*\mathbin{\cdot}\lvert f\rvert\mathbin{\cdot}U\) is a normal
positive functional on \(\mathcal A\). If \(G\) is a projection of
\(\mathcal A\), one has
\[
  (U^*\mathbin{\cdot}\lvert f\rvert\mathbin{\cdot}U)(G)=0
  \quad\Longleftrightarrow\quad
  \lvert f\rvert(UGU^*)=0.
\]
% Source PDF page 070 (printed p. 63)
\phantomsection\label{srcpage:63}
Since the support of
\[
  UGU^*=UG(UG)^*
\]
is bounded above by the support of \(\lvert f\rvert\), the preceding
condition is equivalent to \(UG=0\), and hence to \(U^*UG=0\).
It follows that
\begin{equation}
  \operatorname{supp}
    \bigl(U^*\mathbin{\cdot}\lvert f\rvert\mathbin{\cdot}U\bigr)
  =U^*U.
  \tag{2}\label{eq:I-4-prop6-2}
\end{equation}
Finally, by \eqref{eq:I-4-prop6-1},
\begin{equation}
  f^*\mathbin{\cdot}U
  =U^*\mathbin{\cdot}\lvert f\rvert\mathbin{\cdot}U.
  \tag{3}\label{eq:I-4-prop6-3}
\end{equation}
The proposition follows from \eqref{eq:I-4-prop6-1},
\eqref{eq:I-4-prop6-2}, \eqref{eq:I-4-prop6-3}, and
\cref{thm:I-4-4}\textup{(ii)}.
\end{proof}

\begin{lemma}\label{lem:I-4-11}
Let \(\mathfrak H\) be a Hilbert space, let \(\mathcal A\) be a
\(C^*\)-algebra of operators on \(\mathfrak H\), let \(\varphi\) be a
continuous positive functional on \(\mathcal A\), and let
\(S\in\mathcal A\) be such that the functional
\(\varphi\mathbin{\cdot}S\) is Hermitian. Then
\[
  \bigl|(\varphi\mathbin{\cdot}S)(T)\bigr|
  \leq\lVert S\rVert\,\varphi(T)
\]
for every \(T\in\mathcal A^+\).
\end{lemma}

\begin{proof}
By hypothesis,
\[
  (\varphi\mathbin{\cdot}S)(T^*)
  =\overline{(\varphi\mathbin{\cdot}S)(T)}
\]
for every \(T\in\mathcal A\), that is,
\[
  \varphi(ST^*)=\overline{\varphi(ST)}=\varphi(T^*S^*).
\]
Consequently, if \(T\in\mathcal A^+\),
\[
  |\varphi(ST)|
  =\bigl|\varphi(ST^{1/2}T^{1/2})\bigr|
  \leq\varphi(STS^*)^{1/2}\varphi(T)^{1/2}
  =\varphi(S^2T)^{1/2}\varphi(T)^{1/2}.
\]
Similarly,
\[
  \varphi(S^2T)
  \leq\varphi(S^4T)^{1/2}\varphi(T)^{1/2},
\]
and, more generally,
\[
  \varphi(S^{2^n}T)
  \leq\varphi(S^{2^{n+1}}T)^{1/2}\varphi(T)^{1/2}.
\]
Step by step, therefore,
\begin{align*}
  |\varphi(ST)|
  &\leq
  \varphi(S^{2^n}T)^{1/2^n}
  \varphi(T)^{1/2+1/4+\cdots+1/2^n}\\
  &\leq
  \lVert\varphi\rVert^{1/2^n}
  \lVert T\rVert^{1/2^n}
  \lVert S\rVert\,
  \varphi(T)^{1/2+1/4+\cdots+1/2^n}.
\end{align*}
Letting \(n\to+\infty\) gives
\[
  |\varphi(ST)|\leq\lVert S\rVert\,\varphi(T).
\]
\end{proof}

\begin{theorem}\label{thm:I-4-5}
Let \(\mathcal A\) be a von Neumann algebra, and let \(\varphi\) and
\(\psi\) be normal positive functionals on \(\mathcal A\) such that
\(\varphi\leq\psi\). There exists an element \(T_0\in\mathcal A\)
such that \(0\leq T_0\leq1\) and
\[
  \varphi(T)=\psi(T_0TT_0)
  \qquad\text{for every }T\in\mathcal A.
\]
\end{theorem}

\begin{proof}
Suppose first that \(\psi\) is faithful. Replacing \(\mathcal A\) by
an isomorphic von Neumann algebra, we may suppose that there exists a
vector \(x\), separating for \(\mathcal A\), such that
\[
  \psi(T)=(Tx\mid x)
  \qquad(T\in\mathcal A)
\]
(\hyperref[lem:I-4-4]{Lemma~4} and \cref{prop:I-4-1}). By
\cref{lem:I-4-1},
there exists \(T_0'\in\mathcal A'\) such that \(0\leq T_0'\leq1\) and
\begin{equation}
  \varphi(T)=(TT_0'x\mid T_0'x)
  \qquad\text{for every }T\in\mathcal A.
  \tag{4}\label{eq:I-4-rn-4}
\end{equation}
For every \(T'\in\mathcal A'\), put
\begin{align}
  \psi'(T')&=(T'x\mid x), \tag{5}\label{eq:I-4-rn-5}\\
  f'(T')&=(T'x\mid T_0'x)
          =(T_0'^*T'x\mid x)
          =(\psi'\mathbin{\cdot}T_0'^*)(T').
          \tag{6}\label{eq:I-4-rn-6}
\end{align}
Let
\begin{equation}
  f'=\lvert f'\rvert\mathbin{\cdot}U_0'
  \tag{7}\label{eq:I-4-rn-7}
\end{equation}
% Source PDF page 071 (printed p. 64)
\phantomsection\label{srcpage:64}
be the polar decomposition of \(f'\), so that
\begin{equation}
  \lvert f'\rvert
  =f'\mathbin{\cdot}U_0'^*
  =\psi'\mathbin{\cdot}T_0'^*U_0'^*.
  \tag{8}\label{eq:I-4-rn-8}
\end{equation}

By \cref{lem:I-4-11}, applied to
\(\mathcal A'\), \(\psi'\), and \(T_0'^*U_0'^*\), one has
\[
  \lvert f'\rvert
  \leq\lVert T_0'^*U_0'^*\rVert\psi'
  \leq\psi'.
\]
By \cref{lem:I-4-1} applied to \(\mathcal A'\), there therefore exists
\(T_0\in\mathcal A\) such that \(0\leq T_0\leq1\) and such that, for
every \(T'\in\mathcal A'\),
\begin{equation}
  \lvert f'\rvert(T')
  =\bigl(T'T_0^{1/2}x\mid T_0^{1/2}x\bigr).
  \tag{9}\label{eq:I-4-rn-9}
\end{equation}
By \eqref{eq:I-4-rn-5}, \eqref{eq:I-4-rn-8}, and
\eqref{eq:I-4-rn-9},
\begin{align}
  (T'x\mid U_0'T_0'x)
   &=(T_0'^*U_0'^*T'x\mid x)
     =\lvert f'\rvert(T') \notag\\
   &=\bigl(T'T_0^{1/2}x\mid T_0^{1/2}x\bigr)
     =(T'x\mid T_0x).
  \tag{10}\label{eq:I-4-rn-10}
\end{align}
By \eqref{eq:I-4-rn-5}, \eqref{eq:I-4-rn-8},
\eqref{eq:I-4-rn-7}, and \eqref{eq:I-4-rn-6},
\begin{align}
  (T'x\mid U_0'^*U_0'T_0'x)
   &=(T_0'^*U_0'^*U_0'T'x\mid x)\notag\\
   &=\lvert f'\rvert(U_0'T')
     =(\lvert f'\rvert\mathbin{\cdot}U_0')(T')\notag\\
   &=f'(T')=(T'x\mid T_0'x).
  \tag{11}\label{eq:I-4-rn-11}
\end{align}
Since \(x\) is totalizing for \(\mathcal A'\),
\eqref{eq:I-4-rn-10} and \eqref{eq:I-4-rn-11} imply
\begin{align}
  U_0'T_0'x&=T_0x, \tag{12}\label{eq:I-4-rn-12}\\
  U_0'^*U_0'T_0'x&=T_0'x. \tag{13}\label{eq:I-4-rn-13}
\end{align}
Therefore, for every \(T\in\mathcal A\), by
\eqref{eq:I-4-rn-4}, \eqref{eq:I-4-rn-13}, and
\eqref{eq:I-4-rn-12},
\begin{align*}
  \varphi(T)
   &=(TT_0'x\mid T_0'x)
     =(TU_0'^*U_0'T_0'x\mid T_0'x)\\
   &=(TU_0'T_0'x\mid U_0'T_0'x)
     =(TT_0x\mid T_0x)\\
   &=(T_0TT_0x\mid x)
     =\psi(T_0TT_0).
\end{align*}

If \(\psi\) is not necessarily faithful, let \(E\) be its support.
By what precedes, there exists \(T_0\in\mathcal A\) such that
\[
  0\leq T_0\leq1,\qquad
  T_0E=ET_0=T_0,
\]
and
\[
  \varphi(T)=\psi(T_0TT_0)
  \qquad(T\in E\mathcal A E).
\]
Since
\[
  0\leq\varphi(1-E)\leq\psi(1-E)=0,
\]
the support of \(\varphi\) is bounded above by \(E\). Thus, for every
\(T\in\mathcal A\),
\[
  \varphi(T)
  =\varphi(ETE)
  =\psi\bigl(T_0(ETE)T_0\bigr)
  =\psi(T_0TT_0).
\]
\end{proof}

\noindent\textbf{Bibliography:} \ArticleRef{34}, \ArticleRef{214},
\ArticleRef{239}, \ArticleRef{258}, \ArticleRef{289}, \ArticleRef{301},
\ArticleRef{343}.

\NumberedParagraph{8}{Decomposition of a Hermitian functional into
positive and negative parts}\label{no:I-4-8}

\begin{theorem}\label{thm:I-4-6}
Let \(\mathcal A\) be a von Neumann algebra, and let \(f\) be an
ultra-weakly continuous Hermitian linear functional on \(\mathcal A\).
\begin{enumerate}[label=\textup{(\roman*)}]
\item There exists one and only one pair \((\varphi,\varphi')\) of
      disjoint normal positive functionals on \(\mathcal A\) such
      that \(f=\varphi-\varphi'\).
\item One has \(\lvert f\rvert=\varphi+\varphi'\).
\item If \(E,E'\) are the supports of \(\varphi,\varphi'\), then the
      support of \(\lvert f\rvert\) is \(E+E'\); one has
      \[
        \varphi=f\mathbin{\cdot}E,\qquad
        \varphi'=-f\mathbin{\cdot}E',
      \]
      and
      \[
        f=\lvert f\rvert\mathbin{\cdot}(E-E')
      \]
      is the polar decomposition of \(f\).
\end{enumerate}
\end{theorem}

% Source PDF page 072 (printed p. 65)
\phantomsection\label{srcpage:65}
\begin{proof}
Let
\[
  f=\lvert f\rvert\mathbin{\cdot}U
\]
be the polar decomposition of \(f\). Since \(f=f^*\),
\cref{prop:I-4-6} first shows that \(U=U^*\). Now \(U^*U\) is a
projection, and therefore \(U^2\) is a projection; hence the spectrum
of \(U\) is contained in \(\{-1,0,1\}\), so that
\[
  U=E_1-E_2,
\]
where \(E_1,E_2\) are two orthogonal projections of \(\mathcal A\).
\Cref{prop:I-4-6} also gives
\begin{equation}
  \lvert f\rvert
  =U\mathbin{\cdot}\lvert f\rvert\mathbin{\cdot}U
  =(E_1-E_2)\mathbin{\cdot}\lvert f\rvert
       \mathbin{\cdot}(E_1-E_2).
  \tag{1}\label{eq:I-4-jordan-1}
\end{equation}
But the support of \(\lvert f\rvert\) is
\(U^2=E_1+E_2\), and hence
\begin{equation}
  \lvert f\rvert
  =(E_1+E_2)\mathbin{\cdot}\lvert f\rvert
       \mathbin{\cdot}(E_1+E_2).
  \tag{2}\label{eq:I-4-jordan-2}
\end{equation}
Adding \eqref{eq:I-4-jordan-1} and
\eqref{eq:I-4-jordan-2} gives
\[
  \lvert f\rvert
  =E_1\mathbin{\cdot}\lvert f\rvert\mathbin{\cdot}E_1
   +E_2\mathbin{\cdot}\lvert f\rvert\mathbin{\cdot}E_2,
\]
whence
\begin{align*}
  E_1\mathbin{\cdot}\lvert f\rvert
   &=\lvert f\rvert\mathbin{\cdot}E_1
    =E_1\mathbin{\cdot}\lvert f\rvert\mathbin{\cdot}E_1,\\
  E_2\mathbin{\cdot}\lvert f\rvert
   &=\lvert f\rvert\mathbin{\cdot}E_2
    =E_2\mathbin{\cdot}\lvert f\rvert\mathbin{\cdot}E_2,
\end{align*}
and
\[
  f=\lvert f\rvert\mathbin{\cdot}U
   =\lvert f\rvert\mathbin{\cdot}(E_1-E_2)
   =E_1\mathbin{\cdot}\lvert f\rvert\mathbin{\cdot}E_1
    -E_2\mathbin{\cdot}\lvert f\rvert\mathbin{\cdot}E_2.
\]
Now
\[
  E_1\mathbin{\cdot}\lvert f\rvert\mathbin{\cdot}E_1,
  \qquad
  E_2\mathbin{\cdot}\lvert f\rvert\mathbin{\cdot}E_2
\]
are disjoint normal positive functionals on \(\mathcal A\).

Let now \(\varphi,\varphi'\) be disjoint normal positive functionals
on \(\mathcal A\) such that \(f=\varphi-\varphi'\), and let
\(E,E'\) be their supports. For a projection \(D\) of \(\mathcal A\)
to satisfy
\[
  (\varphi+\varphi')(D)=0,
\]
it is necessary and sufficient that
\(\varphi(D)=\varphi'(D)=0\), that is, that \(DE=DE'=0\). Thus
\(\varphi+\varphi'\) has support \(E+E'\). Moreover, \(E-E'\) is
partially isometric, with initial and final projection \(E+E'\). For
every \(T\in\mathcal A\), one has
\begin{align*}
  \bigl((\varphi+\varphi')\mathbin{\cdot}(E-E')\bigr)(T)
   &=(\varphi+\varphi')(ET-E'T)\\
   &=\varphi(ET)-\varphi'(E'T)
     =\varphi(T)-\varphi'(T)=f(T),\\
  \bigl(f\mathbin{\cdot}(E-E')^*\bigr)(T)
   &=(\varphi-\varphi')(ET-E'T)\\
   &=\varphi(ET)+\varphi'(E'T)
     =\varphi(T)+\varphi'(T)
     =(\varphi+\varphi')(T).
\end{align*}
By \cref{thm:I-4-4}\textup{(ii)}, therefore,
\[
  \lvert f\rvert=\varphi+\varphi',
  \qquad
  U=E-E'.
\]
The equalities \(f=\varphi-\varphi'\) and
\(\lvert f\rvert=\varphi+\varphi'\) prove that the pair
\((\varphi,\varphi')\) is uniquely determined by \(f\). Hence
\[
  \varphi=E_1\mathbin{\cdot}\lvert f\rvert\mathbin{\cdot}E_1,
  \qquad
  \varphi'=E_2\mathbin{\cdot}\lvert f\rvert\mathbin{\cdot}E_2,
  \qquad
  E=E_1,\quad E'=E_2,
\]
and all the assertions of the theorem are now evident.
\end{proof}

\begin{proposition}\label{prop:I-4-7}
Let \(\mathcal A\) be a von Neumann algebra, and let
\(\varphi,\varphi'\) be two normal positive functionals on
\(\mathcal A\). For \(\varphi\) and \(\varphi'\) to be disjoint, it
is necessary and sufficient that
\[
  \lVert\varphi-\varphi'\rVert
  =\lVert\varphi\rVert+\lVert\varphi'\rVert.
\]
\end{proposition}

\begin{proof}
If \(\varphi\) and \(\varphi'\) are disjoint, then
\(\lvert\varphi-\varphi'\rvert=\varphi+\varphi'\) by
\cref{thm:I-4-6}, and hence
\[
  \lVert\varphi-\varphi'\rVert
  =\lVert\varphi+\varphi'\rVert
  =(\varphi+\varphi')(1)
  =\varphi(1)+\varphi'(1)
  =\lVert\varphi\rVert+\lVert\varphi'\rVert.
\]

% Source PDF page 073 (printed p. 66)
\phantomsection\label{srcpage:66}
Suppose that
\[
  \lVert\varphi-\varphi'\rVert
  =\lVert\varphi\rVert+\lVert\varphi'\rVert.
\]
Since \(\varphi-\varphi'\) is weakly continuous on the unit ball of
\(\mathcal A\), and this unit ball is weakly compact, there exists
\(T\in\mathcal A\) such that
\[
  \lVert T\rVert\leq1,\qquad
  (\varphi-\varphi')(T)
  =\lVert\varphi\rVert+\lVert\varphi'\rVert.
\]
Replacing \(T\) by \(\tfrac12(T+T^*)\), we may suppose that \(T\) is
Hermitian. Then
\[
  \varphi(T^+)+\varphi'(T^-)-\varphi(T^-)-\varphi'(T^+)
  =(\varphi-\varphi')(T)
  =\lVert\varphi\rVert+\lVert\varphi'\rVert.
\]
But
\[
  \varphi(T^+)\leq\lVert\varphi\rVert,\qquad
  \varphi'(T^-)\leq\lVert\varphi'\rVert,\qquad
  -\varphi(T^-)-\varphi'(T^+)\leq0,
\]
whence
\[
  \varphi(T^+)=\lVert\varphi\rVert,
  \qquad
  \varphi'(T^-)=\lVert\varphi'\rVert.
\]
Let \(E,E'\) be the supports of \(T^+,T^-\). One has
\[
  \lVert\varphi\rVert
  =\varphi(T^+)\leq\varphi(E)\leq\lVert\varphi\rVert.
\]
Thus \(\varphi(1-E)=0\), and \(E\) majorizes the support of
\(\varphi\). Similarly, \(E'\) majorizes the support of
\(\varphi'\). Since \(EE'=0\), \(\varphi\) and \(\varphi'\) are
disjoint.
\end{proof}

The comparison of positive functionals, begun in
\hyperref[no:I-4-2]{nos.~2}, \hyperref[no:I-4-6]{6},
\hyperref[no:I-4-7]{7}, and \hyperref[no:I-4-8]{8}, can be developed
further (Exercise~8; Chapter~III, \S~1, Exercise~14). The analogy
with the theory of integration should be observed.

\noindent\textbf{Bibliography:} \ArticleRef{92}, \ArticleRef{111},
\ArticleRef{112}, \ArticleRef{113}, \ArticleRef{188}, \ArticleRef{190},
\ArticleRef{239}, \ArticleRef{258}, \ArticleRef{289}.

\par\medskip
\noindent\textit{Exercises.}---
\begin{enumerate}[label=\arabic*.,leftmargin=*,itemsep=0.5\baselineskip]
\item\label{ex:I-4-1}
Let \(\mathcal A\) be a von Neumann algebra, and let \(E\) be a
projection of \(\mathcal A\). Deduce from
\hyperref[cor:I-4-2]{Corollary~2 of Theorem~2} a
new proof of the fact that \((\mathcal A')_E\) is a von Neumann
algebra.

\item\label{ex:I-4-2}
Let \(\mathcal A\) be a von Neumann algebra, let
\(E_1,E_2,\ldots,E_n\) be pairwise orthogonal projections of
\(\mathcal A\), let \(E=E_1+E_2+\cdots+E_n\), and let \(\varphi\) be
a positive linear functional on \(\mathcal A\). The positive linear
functional
\[
  T\longmapsto\varphi(ETE)
\]
is majorized by a scalar multiple of the positive linear functional
\[
  T\longmapsto\sum_{i=1}^n\varphi(E_iTE_i).
\]
[When \(E=1\), for \(T\in\mathcal A^+\) one has
\[
  \varphi(T)
  =\sum_{i=1}^n
      \varphi\bigl(T^{1/2}T^{1/2}E_i\bigr)
  \leq\varphi(T)^{1/2}
      \sum_{i=1}^n\bigl[\varphi(E_iTE_i)\bigr]^{1/2},
\]
whence
\[
  \varphi(T)
  \leq
  \left[\sum_{i=1}^n
      \bigl[\varphi(E_iTE_i)\bigr]^{1/2}\right]^2
  \leq k\sum_{i=1}^n\varphi(E_iTE_i).
\]]
\ArticleRef{19}

\item\label{ex:I-4-3}
Let \(\mathcal A\) and \(\mathcal B\) be two von Neumann algebras, let
\(\mathcal A_u\) (respectively, \(\mathcal B_u\)) be the group of
unitary operators of \(\mathcal A\) (respectively, \(\mathcal B\)),
let \(\varphi\) (respectively, \(\psi\)) be a faithful positive linear
functional on \(\mathcal A\) (respectively, \(\mathcal B\)), and let
\(\Phi\) be an isomorphism of the group \(\mathcal A_u\) onto the
group \(\mathcal B_u\) such that
\[
  \psi(\Phi(U))=\varphi(U)
  \qquad(U\in\mathcal A_u).
\]
Show that \(\Phi\) extends uniquely to an isomorphism of
\(\mathcal A\) onto \(\mathcal B\). (Endow \(\mathcal A\) and
\(\mathcal B\) with the pre-Hilbert structures defined by
\(\varphi\) and \(\psi\). Then \(\Phi\) is isometric. Use
\hyperref[prop:I-1-3]{\S~1, Proposition~3}.)
\ArticleRef{20}, \ArticleRef{150}, \ArticleRef{170},
\ArticleRef{182}

\item\label{ex:I-4-4}
Let \(\mathcal A\) be an involutive operator algebra, and let
\(\varphi\) be a linear functional on \(\mathcal A\), continuous in
the norm topology. Show that \(\varphi\) is a finite linear
combination of
% Source PDF page 074 (printed p. 67)
\phantomsection\label{srcpage:67}
positive functionals on \(\mathcal A\) continuous in the norm
topology. [Let \(\mathcal A_0\) be the set of Hermitian elements of
\(\mathcal A\), which is a real normed space. Its dual
\(\mathcal A_0^*\) is the set of continuous linear functionals
\(\psi\) on \(\mathcal A\) such that
\(\psi(T^*)=\overline{\psi(T)}\). We may suppose that
\(\varphi\in\mathcal A_0^*\). The set \(\mathcal K\) of positive
functionals of norm at most \(1\) in \(\mathcal A_0^*\) is convex and
weakly compact; the set of \(T\in\mathcal A_0\) such that
\[
  |\psi(T)|\leq1\qquad\text{for every }\psi\in\mathcal K
\]
is the unit ball of \(\mathcal A_0\). Hence the unit ball of
\(\mathcal A_0^*\) is the set of
\[
  \lambda_1\psi_1-\lambda_2\psi_2,
  \qquad
  \lambda_1,\lambda_2\geq0,\quad
  \lambda_1+\lambda_2=1,\quad
  \psi_1,\psi_2\in\mathcal K.
\]]
\ArticleRef{111}, \ArticleRef{188}

\item\label{ex:I-4-5}
Let \(\mathcal A\) be an involutive operator algebra on the complex
Hilbert space \(\mathfrak H\), such that
\(\mathfrak X_{\mathfrak H}^{\mathcal A}=\mathfrak H\); let
\(\overline{\mathcal A}\) be the weak closure of \(\mathcal A\),
which is a von Neumann algebra; and let
\((\varphi_i)_{i\in I}\) be the family of positive linear functionals
on \(\mathcal A\) continuous in the norm topology. For every
\(i\in I\), let \(\Phi_i\) be the homomorphism of \(\mathcal A\)
defined by \(\varphi_i\), and let \(\mathfrak H_i\) be the Hilbert
space on which \(\Phi_i(\mathcal A)\) acts. Let \(\mathfrak H'\) be
the Hilbert sum of the \(\mathfrak H_i\), and let \(\Phi\) be the
homomorphism of \(\mathcal A\) into \(\mathcal L(\mathfrak H')\)
defined by the \(\Phi_i\).
\begin{enumerate}[label=\textit{\alph*.}]
\item Show that \(\Phi\) is an isomorphism of \(\mathcal A\) onto
      \(\mathcal B=\Phi(\mathcal A)\), and that every linear
      functional on \(\mathcal B\) continuous in the norm topology
      is weakly continuous. (The
      \(\varphi_i\circ\Phi^{-1}\) are weakly continuous by the
      construction of \(\Phi\). Then apply Exercise~4.)
\item There exists a canonical isomorphism of the bidual of the
      normed space \(\mathcal A\) onto the weak closure
      \(\overline{\mathcal B}\) of \(\mathcal B\), which transforms
      the canonical map of \(\mathcal A\) into its bidual into
      \(\Phi\). (Use \hyperref[ex:I-3-6]{\S~3, Exercise~6a}.)
\item The isomorphism \(\Phi^{-1}\) of \(\mathcal B\) onto
      \(\mathcal A\) is weakly continuous and extends to a weakly
      continuous homomorphism of \(\overline{\mathcal B}\) onto
      \(\overline{\mathcal A}\).
      \ArticleRef{111}, \ArticleRef{112},
      \ArticleRef{113}, \ArticleRef{188}
\end{enumerate}

\item\label{ex:I-4-6}
Let \(\mathcal A_1\) (respectively, \(\mathcal A_2\)) be a von
Neumann algebra, and let \(\varphi_1\) (respectively, \(\varphi_2\))
be an ultra-weakly continuous linear functional on \(\mathcal A_1\)
(respectively, \(\mathcal A_2\)). There exists one and only one
ultra-weakly continuous linear functional \(\varphi\) on
\(\mathcal A=\mathcal A_1\otimes\mathcal A_2\) such that
\[
  \varphi(T_1\otimes T_2)=\varphi_1(T_1)\varphi_2(T_2)
  \qquad
  (T_1\in\mathcal A_1,\ T_2\in\mathcal A_2).
\]
If \(\varphi_1\) and \(\varphi_2\) are positive, \(\varphi\) is
positive; if \(\varphi_1\) and \(\varphi_2\) are faithful,
\(\varphi\) is faithful.

\(\bigl[\)If
\[
  \varphi_1=\sum_{i=1}^{\infty}\omega_{x_i,y_i},
  \qquad
  \varphi_2=\sum_{j=1}^{\infty}\omega_{z_j,t_j},
\]
show that one may take
\[
  \varphi
  =\sum_{i,j=1}^{\infty}
     \omega_{x_i\otimes z_j,\,y_i\otimes t_j}.
\]
If
\[
  \varphi_1=\sum_{i=1}^{\infty}\omega_{x_i},
  \qquad
  \varphi_2=\sum_{j=1}^{\infty}\omega_{z_j}
\]
are faithful, \((x_i)\) is separating for \(\mathcal A_1\) and
\((z_j)\) is separating for \(\mathcal A_2\), and hence
\((x_i\otimes z_j)\) is separating for
\(\mathcal A_1\otimes\mathcal A_2\).\(\bigr]\)

\item\label{ex:I-4-7}
Let \(\mathcal A\) be a von Neumann algebra, let \(\mathcal Z\) be its
center, let \(\varphi\) be a normal positive linear functional on
\(\mathcal A\), let \(\Phi\) be the canonical homomorphism defined by
\(\varphi\), let \(\mathfrak m\) be the kernel of \(\Phi\), which is
a strongly closed two-sided ideal of \(\mathcal A\), let \(F\) be the
projection of \(\mathcal Z\) such that
\(\mathfrak m=\mathcal A F\), and let \(E_\varphi\) be the support of
\(\varphi\). Show that \(1-F\) is the central support of
\(E_\varphi\). [Let \(F'\) be this central support; one has
\(\Phi(1-F')=0\), hence \(1-F'\leq F\); and
\(\varphi(F)=0\), hence \(FE_\varphi=0\), and therefore
\(FF'=0\).]

\item\label{ex:I-4-8}
Let \(\mathcal A\) be a von Neumann algebra on \(\mathfrak H\), and
let \(\varphi\) be a normal positive linear functional on
\(\mathcal A\). For \(S,T\in\mathcal A\), put
\[
  (S\mid T)_\varphi=\varphi(T^*S).
\]
Let \(\mathcal A_\varphi\) be the separated pre-Hilbert space obtained
by passing to the quotient, and let \(\widehat{\mathcal A}_\varphi\)
be the completed Hilbert space. Let \(\psi\) be another normal
positive linear functional on \(\mathcal A\).
\begin{enumerate}[label=\textit{\alph*.}]
\item For the identity map of \(\mathcal A\) to define, on passing to
      the quotients, a linear map of \(\mathcal A_\psi\) onto
      \(\mathcal A_\varphi\), it is necessary and sufficient that
      \(E_\varphi\leq E_\psi\).

% Source PDF page 075 (printed p. 68)
\phantomsection\label{srcpage:68}
\item For the identity map of \(\mathcal A\) to define, on passing to
      the quotients, a continuous linear map of
      \(\mathcal A_\psi\) onto \(\mathcal A_\varphi\), it is necessary
      and sufficient that \(\varphi\) be majorized by a scalar
      multiple of \(\psi\).
\item For the identity map of \(\mathcal A\) to define, on passing to
      the quotients, a linear map of \(\mathcal A_\psi\) onto
      \(\mathcal A_\varphi\) admitting a closed extension, it is
      necessary and sufficient that the conditions
      \[
        \psi(T_n^*T_n)\longrightarrow0,\qquad
        \varphi\bigl((T_n-T_m)^*(T_n-T_m)\bigr)\longrightarrow0
      \]
      imply
      \[
        \varphi(T_n^*T_n)\longrightarrow0.
      \]
      One then says that \(\varphi\) is \emph{almost dominated} by
      \(\psi\).
\item Let \(\varphi_1,\varphi_2,\ldots\) be normal positive linear
      functionals on \(\mathcal A\), almost dominated by \(\psi\),
      and such that
      \[
        \sum_{i=1}^{\infty}\varphi_i(1)<+\infty.
      \]
      Show that
      \[
        \varphi=\sum_{i=1}^{\infty}\varphi_i
      \]
      is normal and almost dominated by \(\psi\). (Show that
      \(\widehat{\mathcal A}_\varphi\) can be identified with a
      subspace of the Hilbert sum of the
      \(\widehat{\mathcal A}_{\varphi_i}\).)
\item Let \(\varphi\) be a normal positive linear functional on
      \(\mathcal A\), let \(z\in\mathfrak H\), and let
      \(\psi=\omega_z\). For \(\varphi\) to be almost dominated by
      \(\psi\), it is necessary and sufficient that there exist a
      closed operator \(T'\mathrel{\eta}\mathcal A'\)
      (\hyperref[ex:I-1-10]{\S~1, Exercise~10}), with domain dense
      throughout, such that
      \[
        \varphi=\omega_{T'z}.
      \]
      [To show that the condition is necessary, reduce to the case
      where \(z\) is totalizing and separating for \(\mathcal A\).
      Then \(\mathcal A_\psi\) can be identified with
      \(\mathcal A z\subset\mathfrak H\). The identity map of
      \(\mathcal A\), regarded as a map of \(\mathcal A z\) into
      \(\mathcal A_\varphi\), has a smallest closed linear extension
      of the form \(WT'\), where \(T'\) is a self-adjoint operator
      \(\geq0\) on \(\mathfrak H\), and where
      \[
        \lVert WT'x\rVert=\lVert T'x\rVert
      \]
      whenever \(x\) belongs to the domain \(\mathfrak M\) of \(T'\).
      Show that \(\mathfrak M\) is invariant under every operator of
      \(\mathcal A\). For \(S,T,T_1\in\mathcal A\), one has
      \[
        (T'ST_1z\mid T'Tz)
        =\varphi(T^*ST_1)
        =(T'T_1z\mid T'S^*Tz);
      \]
      deduce that
      \[
        (T'Sx\mid T'y)=(T'x\mid T'S^*y)
        \qquad(x,y\in\mathfrak M).
      \]
      Then use \hyperref[ex:I-3-7]{\S~3, Exercise~7b}.]
      \ArticleRef{19}
\end{enumerate}

\item\label{ex:I-4-9}
Let \(\mathcal A\) be a von Neumann algebra, and let \(\varphi\) be a
positive linear functional on \(\mathcal A\). For \(\varphi\) to be
normal, it is necessary and sufficient that, for every family
\((E_i)_{i\in I}\) of pairwise orthogonal projections of
\(\mathcal A\), one have
\[
  \varphi\left(\sum_{i\in I}E_i\right)
  =\sum_{i\in I}\varphi(E_i).
\]
[Adapt the proof of the implication
\textup{(i)} \(\Rightarrow\) \textup{(ii)} of
\hyperref[thm:I-4-1]{Theorem~1}.] \ArticleRef{15}

\item\label{ex:I-4-10}
Let \(\mathfrak H\) be a Hilbert space, let \(\mathcal A\) be a von
Neumann algebra on \(\mathfrak H\), and let \(f\) be an ultra-weakly
continuous linear functional on \(\mathcal A\). There exists an
ultra-weakly continuous linear functional \(g\) on
\(\mathcal L(\mathfrak H)\), extending \(f\), with
\[
  \lVert f\rVert=\lVert g\rVert.
\]
[If \(f\geq0\), there exist \(x_i\in\mathfrak H\) such that
\[
  \sum_i\lVert x_i\rVert^2<+\infty,
  \qquad
  f(T)=\sum_i(Tx_i\mid x_i)
  \quad(T\in\mathcal A).
\]
Put
\[
  g(T)=\sum_i(Tx_i\mid x_i)
  \qquad(T\in\mathcal L(\mathfrak H));
\]
then
\[
  \lVert g\rVert=g(1)=f(1)=\lVert f\rVert.
\]
In the general case, use the polar decomposition of \(f\).]
\ArticleRef{214}, \ArticleRef{289}

\item\label{ex:I-4-11}
Let \(\mathcal A\) be a von Neumann algebra, let
\(\varphi,\varphi'\) be two disjoint normal positive functionals on
\(\mathcal A\), and let \(f=\varphi-\varphi'\). One has
\[
  \lVert\varphi\rVert
  =\sup_{\substack{T\in\mathcal A^+\\\lVert T\rVert\leq1}}f(T),
  \qquad
  \lVert\varphi'\rVert
  =\sup_{\substack{T\in\mathcal A^+\\\lVert T\rVert\leq1}}-f(T).
\]
[If \(T\in\mathcal A^+\) and \(\lVert T\rVert\leq1\), then
\[
  f(T)\leq\varphi(T)\leq\lVert\varphi\rVert;
\]
on the other hand, if \(E\) denotes the support of \(\varphi\), then
\[
  f(E)=\varphi(E)=\lVert\varphi\rVert.
\]]
\end{enumerate}
% Source PDF page 076 (printed p. 69)
\phantomsection\label{srcpage:69}
\section{Hilbert Algebras}\label{sec:I-5}

\NumberedParagraph{1}{Definition of Hilbert Algebras}
\label{no:I-5-1}
Let \(\mathfrak A\) be an associative algebra over the field
\(\mathbb C\), endowed with an inner product \((x\mid y)\) which
makes it a separated pre-Hilbert space. Let \(\mathfrak H\) be the
Hilbert-space completion of \(\mathfrak A\). We suppose given:
\begin{enumerate}[label=\textup{\arabic*\textdegree},leftmargin=*]
\item A bijective linear mapping \(x\mapsto x^{\wedge}\) of
\(\mathfrak A\) onto \(\mathfrak A\); the inverse mapping will be
denoted by \(x\mapsto x^{\vee}\).

\item An involutive antiautomorphism \(x\mapsto x^*\) of
\(\mathfrak A\), that is, a bijective mapping of \(\mathfrak A\) onto
\(\mathfrak A\) such that
\[
  (\lambda x+\mu y)^*=\overline\lambda x^*+\overline\mu y^*,
  \qquad
  (xy)^*=y^*x^*,
  \qquad
  x^{**}=x
\]
(this antiautomorphism makes \(\mathfrak A\) an involutive algebra).
\end{enumerate}

\begin{definition}\label{def:I-5-1}
The algebra \(\mathfrak A\) is called a \emph{quasi-Hilbert algebra}
if the following axioms are satisfied:
\begin{enumerate}[label=\textup{(\roman*)},leftmargin=*]
\item \((x\mid y)=(y^*\mid x^*)\) for
\(x,y\in\mathfrak A\);

\item \((xy\mid z)=(y\mid x^{*\wedge}z)\) for
\(x,y,z\in\mathfrak A\);

\item for every \(x\in\mathfrak A\), the mapping
\(y\mapsto xy\) is continuous;

\item the set of elements \(xy\), where
\(x,y\in\mathfrak A\), is total in \(\mathfrak A\);

\item if \(a\) and \(b\) are two elements of \(\mathfrak H\) such
that
\[
  (a\mid xy)=(b\mid x^{\wedge}y^{\wedge})
\]
for every \(x,y\in\mathfrak A\), there is a sequence \((x_n)\) in
\(\mathfrak A\) such that \(x_n\to b\) and
\(x_n^{\wedge}\to a\).
\end{enumerate}
The algebra \(\mathfrak A\) is called a \emph{Hilbert algebra} if, in
addition, \(x^{\wedge}=x\) for every \(x\in\mathfrak A\).
\end{definition}

Axioms~\textup{(i)} and~\textup{(iii)} imply that the mapping
\[
  y\longmapsto yx=(x^*y^*)^*
\]
is continuous. Axioms~\textup{(i)} and~\textup{(ii)} imply
\begin{equation}
  (xy\mid z)
   =(z^*\mid y^*x^*)
   =(y^{\wedge}z^*\mid x^*)
   =(x\mid z y^{\wedge *}),
  \tag{1}\label{eq:I-5-1}
\end{equation}
which restores the symmetry between left and right multiplication.
Moreover,
\[
  (xy\mid z)
   =(y\mid x^{*\wedge}z)
   =(x^{*\wedge *\wedge}y\mid z),
\]
and hence
\[
  (x\mid z y^{\wedge *})
   =(x^{*\wedge *\wedge}\mid z y^{\wedge *}).
\]
Since \(y^{\wedge *}\) is an arbitrary element of \(\mathfrak A\),
axiom~\textup{(iv)} gives \(x=x^{*\wedge *\wedge}\), whence
\begin{align}
  x^{\vee *}&=x^{*\wedge}, \tag{2}\\
  x^{\wedge *}&=x^{*\vee}. \tag{3}
\end{align}
Relation~\eqref{eq:I-5-1} and axiom~\textup{(ii)} may therefore also be
written
\begin{align}
  (xy\mid z)&=(x\mid z y^{*\vee}), \tag{4}\\
  (xy\mid z)&=(y\mid x^{\vee *}z). \tag{5}
\end{align}
% Source PDF page 077 (printed p. 70)
\phantomsection\label{srcpage:70}
For Hilbert algebras, axiom~\textup{(v)} is satisfied automatically.

By axiom~\textup{(i)}, the mapping \(x\mapsto x^*\) extends uniquely
to an involution \(J\) of \(\mathfrak H\), that is, we recall, to a
bijective mapping \(J\) of \(\mathfrak H\) onto itself such that
\[
  J^2=I_{\mathfrak H},
  \qquad
  (Ja\mid Jb)=(b\mid a),
  \qquad
  J(\lambda a+\mu b)=\overline\lambda Ja+\overline\mu Jb.
\]
The mapping \(J\) is called the involution of \(\mathfrak H\)
canonically defined by \(\mathfrak A\).

The mappings \(y\mapsto xy\) and \(y\mapsto yx\) extend uniquely to
elements \(U_x,V_x\) of \(\mathcal L(\mathfrak H)\). We immediately
have
\begin{align}
  U_{\lambda x+\mu y}&=\lambda U_x+\mu U_y,
  &U_{xy}&=U_xU_y,
  &U_{x^{*\wedge}}&=U_x^*, \tag{6}\\
  V_{\lambda x+\mu y}&=\lambda V_x+\mu V_y,
  &V_{xy}&=V_yV_x,
  &V_{x^{\wedge *}}&=V_x^*, \tag{7}\\
  U_xV_y&=V_yU_x, \tag{8}\label{eq:I-5-8}\\
  JU_xJ&=V_{x^*},
  &JV_xJ&=U_{x^*}. \tag{9}\label{eq:I-5-9}
\end{align}
The \(U_x\) (respectively, the \(V_x\)) form an involutive algebra of
operators on \(\mathfrak H\). By
\hyperref[thm:I-3-2]{\S~3, Corollary~1 of Theorem~2} and
axiom~\textup{(iv)}, the weak closure of this algebra is a von Neumann
algebra \(\mathcal U(\mathfrak A)\) (respectively,
\(\mathcal V(\mathfrak A)\)), called the von Neumann algebra
associated on the left (respectively, on the right) with
\(\mathfrak A\). Consequently, the elements \(xy\), where
\(x,y\in\mathfrak A\), are dense in \(\mathfrak A\). The algebras
\(\mathcal U(\mathfrak A)\) and \(\mathcal V(\mathfrak A)\) commute
by~\eqref{eq:I-5-8}, and
\[
  J\mathcal U(\mathfrak A)J=\mathcal V(\mathfrak A),
  \qquad
  J\mathcal V(\mathfrak A)J=\mathcal U(\mathfrak A)
\]
by~\eqref{eq:I-5-9}. The mappings \(x\mapsto U_x\) and
\(x\mapsto V_x\) are called the canonical mappings of
\(\mathfrak A\) into \(\mathcal U(\mathfrak A)\) and
\(\mathcal V(\mathfrak A)\).

Hilbert algebras are also sometimes called unitary algebras. As we
shall see, they constitute a powerful means of studying von Neumann
algebras.

\noindent\textbf{Bibliography:}
\ArticleRef{1}, \ArticleRef{2}, \ArticleRef{11}, \ArticleRef{13},
\ArticleRef{19}, \ArticleRef{27}, \ArticleRef{29}, \ArticleRef{30},
\ArticleRef{66}, \ArticleRef{72}, \ArticleRef{87}, \ArticleRef{90},
\ArticleRef{95}, \ArticleRef{101}, \ArticleRef{115}, \ArticleRef{117},
\ArticleRef{127}, \ArticleRef{189}, \ArticleRef{261}, \TreatiseRef{9}.

\NumberedParagraph{2}{The Commutation Theorem}
\label{no:I-5-2}

\begin{definition}\label{def:I-5-2}
An element \(a\in\mathfrak H\) is said to be \emph{left-bounded} if
there exists a continuous operator \(U_a\) on \(\mathfrak H\); it is
said to be \emph{right-bounded} if there exists a continuous operator
\(V_a\) on \(\mathfrak H\). These operators are required to satisfy,
respectively,
\[
  U_ax=V_xa,
  \qquad
  V_ax=U_xa
\]
for \(x\in\mathfrak A\).
\end{definition}

The elements of \(\mathfrak A\) are left- and right-bounded, and the
notations \(U_a,V_a\) agree with the earlier notations \(U_x,V_x\)
when \(a\in\mathfrak A\). Moreover, the equality
\(U_ax=V_xa\) (respectively, \(V_ax=U_xa\)), on letting \(V_x\)
(respectively, \(U_x\)) tend weakly to \(I_{\mathfrak H}\), proves
that
\[
  a\in\overline{U_a(\mathfrak H)}
  \qquad
  \text{(respectively, \(a\in\overline{V_a(\mathfrak H)}\)).}
\]
It follows in particular that the mappings \(a\mapsto U_a\) and
\(a\mapsto V_a\) are injective.

\begin{lemma}\label{lem:I-5-1}
If \(a\) is left-bounded and
\(T\in\mathcal V(\mathfrak A)'\), then \(Ta\) is left-bounded and
\(TU_a=U_{Ta}\); the operators \(U_a\) form a left ideal
\(\mathfrak m\) of \(\mathcal V(\mathfrak A)'\). If \(a\) is
right-bounded and \(T\in\mathcal U(\mathfrak A)'\), then \(Ta\) is
right-bounded and \(TV_a=V_{Ta}\); the operators \(V_a\) form a left
ideal \(\mathfrak n\) of \(\mathcal U(\mathfrak A)'\).
\end{lemma}

% Source PDF page 078 (printed p. 71)
\phantomsection\label{srcpage:71}
\begin{proof}
Let \(x,y\in\mathfrak A\), and suppose that \(a\) is left-bounded.
We have
\[
  U_aV_xy=U_a(yx)=V_{yx}a=V_xV_ya=V_xU_ay;
\]
thus \(U_a\) commutes with the \(V_x\), and
\(U_a\in\mathcal V(\mathfrak A)'\). If
\(T\in\mathcal V(\mathfrak A)'\), then
\[
  TU_ax=TV_xa=V_xTa,
\]
so \(Ta\) is left-bounded and \(U_{Ta}=TU_a\). The argument is
analogous when \(a\) is right-bounded.
\end{proof}

Let \(\mathfrak m^*\) (respectively, \(\mathfrak n^*\)) denote the
image of \(\mathfrak m\) (respectively, \(\mathfrak n\)) under the
mapping \(T\mapsto T^*\).

\begin{lemma}\label{lem:I-5-2}
Put
\[
  \mathfrak m_1=\mathfrak m\cap\mathfrak m^*,
  \qquad
  \mathfrak n_1=\mathfrak n\cap\mathfrak n^*.
\]
Then
\[
  \mathfrak m_1''=\mathcal V(\mathfrak A)',
  \qquad
  \mathfrak n_1''=\mathcal U(\mathfrak A)'.
\]
\end{lemma}

\begin{proof}
By \cref{lem:I-5-1},
\[
  \mathfrak m_1''\subset\mathcal V(\mathfrak A)',
  \qquad
  \mathfrak n_1''\subset\mathcal U(\mathfrak A)'.
\]
We shall prove that
\(\mathcal V(\mathfrak A)'\subset\mathfrak m_1''\); the proof of
\(\mathcal U(\mathfrak A)'\subset\mathfrak n_1''\) is analogous.
Let, then, \(T\in\mathcal V(\mathfrak A)'\) and
\(T_1\in\mathfrak m_1'\), and let us prove that \(TT_1=T_1T\).
Lemma~1 immediately implies that, for \(x,x'\in\mathfrak A\),
\[
  U_{x'^*}TU_x\in\mathfrak m_1.
\]
Consequently,
\[
  U_{x'^*}TU_xT_1=T_1U_{x'^*}TU_x,
\]
and it remains only to let \(U_x\) tend weakly to
\(I_{\mathfrak H}\), and then \(U_{x'}\) tend weakly to
\(I_{\mathfrak H}\).
\end{proof}

\begin{lemma}\label{lem:I-5-3}
The algebras \(\mathfrak m_1\) and \(\mathfrak n_1\) commute.
\end{lemma}

\begin{proof}
Let \(U_a\in\mathfrak m_1\) and \(V_c\in\mathfrak n_1\). We have
\[
  U_a^*=U_b,
  \qquad
  V_c^*=V_d,
\]
where \(a,b\) are left-bounded and \(c,d\) are right-bounded. For
every \(x,y\in\mathfrak A\),
\begin{align*}
  (a\mid xy)
   &=(a\mid V_yx)
    =(V_{y^{\wedge *}}a\mid x)
    =(U_ay^{\wedge *}\mid x)\\
   &=(y^{\wedge *}\mid U_bx)
    =(y^{\wedge *}\mid V_xb)
    =(V_{x^{\wedge *}}y^{\wedge *}\mid b)\\
   &=((x^{\wedge}y^{\wedge})^*\mid b)
    =(Jb\mid x^{\wedge}y^{\wedge}).
\end{align*}
By axiom~\textup{(v)}, there is a sequence \((x_n)\) in
\(\mathfrak A\) such that
\[
  x_n^*\longrightarrow b,
  \qquad
  x_n^{\wedge}\longrightarrow a.
\]
An analogous calculation proves the existence of a sequence \((y_n)\)
in \(\mathfrak A\) such that
\[
  y_n^*\longrightarrow c,
  \qquad
  y_n^{\vee}\longrightarrow d.
\]
Then
\begin{align*}
  (U_aV_cx\mid y)
   &=(V_cx\mid U_by)
    =(U_xc\mid V_yb)\\
   &=\lim_n(U_xy_n^*\mid V_yx_n^*)
    =\lim_n(xy_n^*\mid x_n^*y)\\
   &=\lim_n(x_n^{\wedge}x\mid y y_n^{\vee})
    =\lim_n(V_xx_n^{\wedge}\mid U_yy_n^{\vee})\\
   &=(V_xa\mid U_yd)
    =(U_ax\mid V_dy)
    =(V_cU_ax\mid y),
\end{align*}
and therefore \(U_aV_c=V_cU_a\).
\end{proof}

\begin{theorem}[Commutation theorem]\label{thm:I-5-1}
\[
  \mathcal U(\mathfrak A)'=\mathcal V(\mathfrak A),
  \qquad
  \mathcal V(\mathfrak A)'=\mathcal U(\mathfrak A).
\]
\end{theorem}

\begin{proof}
We already know that
\(\mathcal U(\mathfrak A)\subset\mathcal V(\mathfrak A)'\). On the
other hand,
\[
  \mathcal V(\mathfrak A)'
   =\mathfrak m_1''
   \subset\mathfrak n_1'
   =\mathcal U(\mathfrak A)''
   =\mathcal U(\mathfrak A).
\]
Thus \(\mathcal U(\mathfrak A)=\mathcal V(\mathfrak A)'\), and,
on taking commutants, \(\mathcal U(\mathfrak A)'=
\mathcal V(\mathfrak A)\).
\end{proof}

% Source PDF page 079 (printed p. 72)
\phantomsection\label{srcpage:72}
\begin{proposition}\label{prop:I-5-1}
Let \(\mathfrak A\) be a Hilbert algebra and \(\mathfrak A'\) an
involutive subalgebra of \(\mathfrak A\).
\begin{enumerate}[label=\textup{(\roman*)},leftmargin=*]
\item \(\mathfrak A'\) is a Hilbert algebra.
\item If the operators \(U_x\), where \(x\in\mathfrak A'\), are
strongly dense in \(\mathcal U(\mathfrak A)\), then
\(\mathfrak A'\) is dense in \(\mathfrak A\).
\item If \(\mathfrak A'\) is dense in \(\mathfrak A\), then
\[
  \mathcal U(\mathfrak A')=\mathcal U(\mathfrak A),
  \qquad
  \mathcal V(\mathfrak A')=\mathcal V(\mathfrak A).
\]
\end{enumerate}
\end{proposition}

\begin{proof}
\textup{(i)} It is clear that \(\mathfrak A'\) satisfies
axioms~\textup{(i)}, \textup{(ii)}, and~\textup{(iii)}. If
\(z\in\overline{\mathfrak A'}\) is orthogonal to the elements \(xy\)
with \(x,y\in\mathfrak A'\), then
\[
  (U_{x^*}z\mid y)=(z\mid xy)=0
\]
for \(x,y\in\mathfrak A'\), and therefore \(U_{x^*}z=0\), since
\(U_{x^*}z\) is a limit of elements of \(\mathfrak A'\). For every
\(t\in\mathfrak A\), we then have
\[
  (V_tz\mid x)=(z\mid xt^*)=(U_{x^*}z\mid t^*)=0;
\]
thus \(V_tz\) is orthogonal to \(\mathfrak A'\). But \(z\) is a
limit of elements of the form \(V_tz\), and hence is orthogonal to
\(\mathfrak A'\); finally, \(z=0\).

\textup{(ii)} Suppose the condition in~\textup{(ii)} holds. Let
\(a,b\in\mathfrak A\). There are operators \(U_x\), with
\(x\in\mathfrak A'\), which tend strongly to \(U_a\). Thus one can
find \(x\in\mathfrak A'\) such that \(xb\) is arbitrarily close to
\(ab\). By an analogous argument, one can then find
\(y\in\mathfrak A'\) such that \(xy\) is arbitrarily close to
\(xb\). Hence \(\mathfrak A'\) is dense in the set of the \(ab\),
with \(a,b\in\mathfrak A\), and therefore in \(\mathfrak A\).

\textup{(iii)} If \(\mathfrak A'\) is dense in \(\mathfrak A\), the
two algebras have the same Hilbert-space completion. Plainly,
\[
  \mathcal U(\mathfrak A')\subset\mathcal U(\mathfrak A),
  \qquad
  \mathcal V(\mathfrak A')\subset\mathcal V(\mathfrak A).
\]
Consequently,
\[
  \mathcal U(\mathfrak A')
   =\mathcal V(\mathfrak A')'
   \supset\mathcal V(\mathfrak A)'
   =\mathcal U(\mathfrak A),
\]
which gives
\(\mathcal U(\mathfrak A')=\mathcal U(\mathfrak A)\), and similarly
\(\mathcal V(\mathfrak A')=\mathcal V(\mathfrak A)\).
\end{proof}

Throughout the remainder of the book (except in Exercise~5), we shall
deal only with Hilbert algebras. We have nevertheless proved
\cref{thm:I-5-1} for quasi-Hilbert algebras, because it is useful in
certain applications and because restricting to Hilbert algebras does
not shorten the proofs.

\noindent\textbf{Bibliography:}
\ArticleRef{13}, \ArticleRef{19}, \ArticleRef{27}, \ArticleRef{29},
\ArticleRef{30}, \ArticleRef{53}, \ArticleRef{67}, \ArticleRef{90},
\ArticleRef{98}, \ArticleRef{101}, \ArticleRef{115}, \ArticleRef{123},
\ArticleRef{255}, \ArticleRef{277}.

\NumberedParagraph{3}{Bounded Elements in Hilbert Algebras}
\label{no:I-5-3}
Throughout the rest of this section, \(\mathfrak A\) is assumed to be
a Hilbert algebra.

\begin{proposition}\label{prop:I-5-2}
For an element \(a\in\mathfrak H\) to be left-bounded, it is necessary
and sufficient that it be right-bounded. The element \(Ja\) then has
the same properties, and
\[
  U_{Ja}=U_a^*=JV_aJ,
  \qquad
  V_{Ja}=V_a^*=JU_aJ.
\]
\end{proposition}

\begin{proof}
For \(x\in\mathfrak A\) and \(a\in\mathfrak H\),
\[
  U_{Jx}a=JV_xJa.
\]
If \(a\) is right-bounded, it follows that
\[
  JV_aJx=V_xJa;
\]
thus \(Ja\) is left-bounded and \(U_{Ja}=JV_aJ\). If \(Ja\) is
left-bounded, the same equality gives
\[
  U_{Jx}a=JU_{Ja}x,
\]
or, equivalently,
\[
  U_ya=JU_{Ja}Jy\qquad(y\in\mathfrak A);
\]
thus \(a\) is right-bounded and \(V_a=JU_{Ja}J\).

Suppose again that \(a\) is right-bounded. For
\(x,y\in\mathfrak A\),
\begin{align*}
  (U_xJa\mid y)
   &=(Ja\mid x^*y)
    =(y^*x\mid a)
    =(x\mid U_ya)\\
   &=(x\mid V_ay)
    =(V_a^*x\mid y),
\end{align*}
and hence \(U_xJa=V_a^*x\). Thus \(Ja\) is right-bounded and
\(V_{Ja}=V_a^*\).

% Source PDF page 080 (printed p. 73)
\phantomsection\label{srcpage:73}
We have therefore
\[
 \begin{aligned}
  a\text{ right-bounded}
  &\Longleftrightarrow Ja\text{ left-bounded}\\
  &\Longrightarrow Ja\text{ right-bounded}
   \Longrightarrow a\text{ left-bounded}.
 \end{aligned}
\]
Interchanging the roles of \(a\) and \(Ja\), one sees that all these
conditions are equivalent. The formula
\[
  V_{Ja}=V_a^*=JU_aJ
\]
follows from the preceding argument. It then gives
\[
  U_{Ja}=JV_aJ=JV_{Ja}^*J=U_a^*.
\]
\end{proof}

\begin{corollary*}
For every element \(C\) of the common centre of
\(\mathcal U(\mathfrak A)\) and \(\mathcal V(\mathfrak A)\),
\[
  JCJ=C^*.
\]
\end{corollary*}

\begin{proof}
For \(x\in\mathfrak A\), \cref{lem:I-5-1} gives
\[
  U_{Cx^*}=CU_{x^*}=(C^*U_x)^*=U_{C^*x}.
\]
Consequently, by \cref{prop:I-5-2},
\[
  U_{CJx}=U_{JC^*x},
  \qquad
  CJx=JC^*x,
  \qquad
  CJ=JC^*.
\]
\end{proof}

\begin{definition}\label{def:I-5-3}
An element \(a\in\mathfrak H\) having the properties of
\cref{prop:I-5-2} is said to be \emph{bounded relative to}
\(\mathfrak A\).
\end{definition}

\textbf{Remark.} --- Let \(a\in\mathfrak H\). For
\(x,y\in\mathfrak A\),
\[
  (xy\mid a)=(x\mid V_{y^*}a).
\]
Thus, to say that \(a\) is bounded is to say that the bilinear form
\((x,y)\mapsto(xy\mid a)\) is continuous in both variables \(x,y\).
It follows that it is equivalent to say that \(a\) is bounded relative
to \(\mathfrak A\), or relative to any Hilbert subalgebra of
\(\mathfrak A\) which is dense in \(\mathfrak A\).

\begin{proposition}\label{prop:I-5-3}
The operators \(U_a\) (respectively, \(V_a\)), where \(a\) is bounded,
form a two-sided ideal of \(\mathcal U(\mathfrak A)\) (respectively,
\(\mathcal V(\mathfrak A)\)). For \(T\in\mathcal U(\mathfrak A)\),
\[
  TU_a=U_{Ta},
  \qquad
  U_aT=U_{JT^*Ja};
\]
for \(T'\in\mathcal V(\mathfrak A)\),
\[
  T'V_a=V_{T'a},
  \qquad
  V_aT'=V_{JT'^*Ja}.
\]
\end{proposition}

\begin{proof}
The formulas \(TU_a=U_{Ta}\) and \(T'V_a=V_{T'a}\) were established
in \cref{lem:I-5-1}. Moreover,
\begin{align*}
  U_aT
   &=(T^*U_a^*)^*
    =(T^*U_{Ja})^*
    =(U_{T^*Ja})^*
    =U_{JT^*Ja},\\
  V_aT'
   &=(T'^*V_a^*)^*
    =(T'^*V_{Ja})^*
    =(V_{T'^*Ja})^*
    =V_{JT'^*Ja}.
\end{align*}
\end{proof}

\begin{proposition}\label{prop:I-5-4}
For an element \(a\in\mathfrak H\) to be bounded, it is necessary and
sufficient that there exist a sequence \((x_n)\) of elements of
\(\mathfrak A\) such that
\[
  \lVert x_n-a\rVert\longrightarrow0,
  \qquad
  \sup_n\lVert U_{x_n}\rVert<+\infty.
\]
Then \(U_{x_n}\) tends strongly to \(U_a\).
\end{proposition}

\begin{proof}
Suppose there is a sequence \((x_n)\) in \(\mathfrak A\) such that
\[
  \lVert x_n-a\rVert\longrightarrow0,
  \qquad
  \sup_n\lVert U_{x_n}\rVert=M<+\infty.
\]
Then, for every \(y\in\mathfrak A\),
\[
  \lVert V_yx_n\rVert
   =\lVert U_{x_n}y\rVert
   \leq M\lVert y\rVert,
\]
and hence \(\lVert V_ya\rVert\leq M\lVert y\rVert\).
% Source PDF page 081 (printed p. 74)
\phantomsection\label{srcpage:74}
Thus \(a\) is bounded. Moreover,
\[
  U_{x_n}y=V_yx_n\longrightarrow V_ya=U_ay,
\]
so \(U_{x_n}\) tends strongly to \(U_a\).

Conversely, suppose that \(a\) is bounded, and let \(\varepsilon>0\).
There is an \(x\in\mathfrak A\) such that
\[
  \lVert V_xa-a\rVert\leq\varepsilon,
  \qquad
  \lVert V_x\rVert\leq1
\]
(\cref{thm:I-3-3}). There is then a \(y\in\mathfrak A\) such that
\[
  \lVert U_yx-U_ax\rVert\leq\varepsilon,
  \qquad
  \lVert U_y\rVert\leq\lVert U_a\rVert
\]
(\cref{thm:I-3-3}). Since \(U_ax=V_xa\), it follows that
\(\lVert yx-a\rVert\leq2\varepsilon\). Furthermore,
\[
  \lVert U_{yx}\rVert
   \leq\lVert U_y\rVert\lVert U_x\rVert
   =\lVert U_y\rVert\lVert V_x\rVert
   \leq\lVert U_a\rVert.
\]
This proves the existence of a sequence having the properties in the
proposition.
\end{proof}

\begin{definition}\label{def:I-5-4}
The Hilbert algebra \(\mathfrak A\) is called \emph{completed} if
every bounded element of \(\mathfrak H\) belongs to \(\mathfrak A\).
\end{definition}

Let \(\mathfrak A\) be a Hilbert algebra, and let \(\mathfrak B\) be
the vector space of bounded elements of \(\mathfrak H\). As a subspace
of \(\mathfrak H\), \(\mathfrak B\) is endowed with a pre-Hilbert
space structure. For \(a,b\in\mathfrak B\),
\[
  U_ab=V_ba.
\]
Indeed, let \((x_n)\) be a sequence of elements of \(\mathfrak A\)
tending strongly to \(b\). For every \(y\in\mathfrak A\),
\[
  (U_ax_n\mid y)
   =(V_{x_n}a\mid y)
   =(a\mid yx_n^*)
   =(a\mid U_yx_n^*),
\]
and hence, on passing to the limit,
\[
  (U_ab\mid y)
   =(a\mid U_yJb)
   =(a\mid V_{Jb}y)
   =(V_ba\mid y),
\]
which proves the assertion. For \(a,b\in\mathfrak B\), put
\[
  ab=U_ab=V_ba.
\]
This defines on \(\mathfrak B\) a multiplication extending that of
\(\mathfrak A\), and makes \(\mathfrak B\) an associative algebra,
since, for \(a,b,c\in\mathfrak B\),
\[
  a(bc)=U_aV_cb=V_cU_ab=(ab)c.
\]
Finally, for \(a\in\mathfrak B\), put \(a^*=Ja\). Then
\[
  a^*b^*=U_{Ja}Jb=JV_ab=J(ba)=(ba)^*;
\]
thus \(\mathfrak B\) becomes an involutive algebra, and
\(\mathfrak A\) is an involutive subalgebra of \(\mathfrak B\). One
immediately verifies that \(\mathfrak B\) is a Hilbert algebra, called
the \emph{Hilbert algebra of bounded elements}. Every element of
\(\mathfrak H\) bounded relative to \(\mathfrak B\) belongs to
\(\mathfrak B\). Thus the Hilbert algebra of bounded elements is
completed.

By \cref{prop:I-5-1}\textup{(iii)},
\[
  \mathcal U(\mathfrak A)=\mathcal U(\mathfrak B),
  \qquad
  \mathcal V(\mathfrak A)=\mathcal V(\mathfrak B).
\]
The preceding construction makes it possible, in almost every case,
to reduce problems concerning Hilbert algebras to problems concerning
completed Hilbert algebras.

Instead of ``completed,'' one also says ``maximal.''

\noindent\textbf{Bibliography:}
\ArticleRef{13}, \ArticleRef{27}, \ArticleRef{29}, \ArticleRef{30},
\ArticleRef{90}, \ArticleRef{101}, \ArticleRef{115}.

% Source PDF page 082 (printed p. 75)
\phantomsection\label{srcpage:75}
\NumberedParagraph{4}{Central Elements in Hilbert Algebras}
\label{no:I-5-4}

\begin{proposition}\label{prop:I-5-5}
For an element \(a\in\mathfrak H\), the following conditions are
equivalent:
\begin{enumerate}[label=\textup{(\roman*)},leftmargin=*]
\item \((a\mid xy)=(a\mid yx)\) for \(x,y\in\mathfrak A\);
\item \(U_xa=V_xa\) for \(x\in\mathfrak A\);
\item \(Ta=JT^*Ja\) for \(T\in\mathcal U(\mathfrak A)\), and hence
for \(T\in\mathcal V(\mathfrak A)\).
\end{enumerate}
\end{proposition}

\begin{proof}
For \(x,y\in\mathfrak A\),
\[
  (a\mid xy)=(U_{x^*}a\mid y),
  \qquad
  (a\mid yx)=(V_{x^*}a\mid y),
\]
which proves the equivalence of conditions~\textup{(i)} and
\textup{(ii)}. Moreover,
\[
  V_x=JU_x^*J,
\]
so condition~\textup{(iii)} implies condition~\textup{(ii)}. Finally,
condition~\textup{(ii)} implies condition~\textup{(iii)}, since every
\(T\in\mathcal U(\mathfrak A)\) is a weak limit of operators \(U_x\).
\end{proof}

\begin{definition}\label{def:I-5-5}
If \(a\in\mathfrak H\) satisfies the conditions of
\cref{prop:I-5-5}, it is said to be \emph{central relative to}
\(\mathfrak A\).
\end{definition}

For \(a\in\mathfrak A\), condition~\textup{(ii)} of
\cref{prop:I-5-5} recovers the usual algebraic notion. The set
\(\mathfrak Z\) of central elements of \(\mathfrak H\) is a closed
vector subspace of \(\mathfrak H\). It is unchanged if
\(\mathfrak A\) is replaced by an involutive subalgebra which is dense
in \(\mathfrak A\), since \(\mathcal U(\mathfrak A)\) is unchanged.
In particular, the central elements are the same relative to
\(\mathfrak A\) and relative to the Hilbert algebra of bounded
elements.

\begin{proposition}\label{prop:I-5-6}
Let
\[
  \mathcal Z=\mathcal U(\mathfrak A)\cap\mathcal V(\mathfrak A)
\]
be the common centre of \(\mathcal U(\mathfrak A)\) and
\(\mathcal V(\mathfrak A)\). Then:
\begin{enumerate}[label=\textup{(\roman*)},leftmargin=*]
\item \(J(\mathfrak Z)=\mathfrak Z\);
\item \(P_{\mathfrak Z}\in\mathcal Z\);
\item for \(a\in\mathfrak Z\),
\[
  E_a^{\mathcal U(\mathfrak A)}
   =E_a^{\mathcal V(\mathfrak A)}\in\mathcal Z;
\]
\item
\[
  E_{\mathfrak Z}^{\mathcal U(\mathfrak A)}
   =E_{\mathfrak Z}^{\mathcal V(\mathfrak A)}\in\mathcal Z.
\]
\end{enumerate}
\end{proposition}

\begin{proof}
If \(a\in\mathfrak Z\) and \(T\in\mathcal U(\mathfrak A)\), then
\[
  T(Ja)=J(JTJa)=JT^*a=(JT^*J)(Ja),
\]
and hence \(Ja\in\mathfrak Z\), proving~\textup{(i)}. If
\(R\in\mathcal Z\), then
\[
  T(Ra)=R(Ta)=R(JT^*Ja)=(JT^*J)(Ra),
\]
so \(Ra\in\mathfrak Z\), proving~\textup{(ii)}. The equality
\[
  E_a^{\mathcal U(\mathfrak A)}
   =E_a^{\mathcal V(\mathfrak A)}
\]
is immediate and implies
\[
  E_a^{\mathcal U(\mathfrak A)}
   \in\mathcal U(\mathfrak A)\cap\mathcal V(\mathfrak A)
   =\mathcal Z.
\]
Finally, \textup{(iii)} implies~\textup{(iv)}.
\end{proof}

% Source PDF page 083 (printed p. 76)
\phantomsection\label{srcpage:76}
\begin{definition}\label{def:I-5-6}
The projection
\[
  E_{\mathfrak Z}^{\mathcal U(\mathfrak A)}
   =E_{\mathfrak Z}^{\mathcal V(\mathfrak A)}
\]
of \cref{prop:I-5-6} is called the \emph{characteristic projection}
of \(\mathfrak A\).
\end{definition}

In particular, if \(\mathfrak A\) has an identity element, its
characteristic projection is \(I_{\mathfrak H}\).

\begin{proposition}\label{prop:I-5-7}
For a bounded element \(a\in\mathfrak H\) to be central, it is
necessary and sufficient that \(U_a\in\mathcal Z\), or that
\(V_a\in\mathcal Z\). One then has \(U_a=V_a\).
\end{proposition}

\begin{proof}
If the bounded element \(a\) is central, then, for
\(x\in\mathfrak A\),
\[
  U_ax=V_xa=U_xa=V_ax,
\]
and hence \(U_a=V_a\), so that \(U_a\in\mathcal Z\). Conversely, if
\(U_a\in\mathcal Z\), then, for \(x\in\mathfrak A\),
\[
  U_{U_ax}=U_aU_x=U_xU_a=U_{U_xa};
\]
thus \(U_xa=U_ax=V_xa\), and \(a\) is central. The argument is
analogous if \(V_a\in\mathcal Z\).
\end{proof}

Important supplements are given in Exercise~4.

\noindent\textbf{Bibliography:}
\ArticleRef{13}, \ArticleRef{29}, \ArticleRef{90}.

\NumberedParagraph{5}{Elementary Operations on Hilbert Algebras}
\label{no:I-5-5}
Let \(\mathfrak A\) be a Hilbert algebra. Without changing the
pre-Hilbert structure or the involution of \(\mathfrak A\), replace
the multiplication \((x,y)\mapsto xy\) by the multiplication
\((x,y)\mapsto yx\). One immediately verifies that the result is a
Hilbert algebra \(\mathfrak A'\), called the \emph{Hilbert algebra
opposite to} \(\mathfrak A\). One has
\[
  \mathcal U(\mathfrak A')=\mathcal V(\mathfrak A),
  \qquad
  \mathcal V(\mathfrak A')=\mathcal U(\mathfrak A).
\]

Let \((\mathfrak A_i)_{i\in I}\) be a family of Hilbert algebras, and
let \(\mathfrak H_i\) be the Hilbert-space completion of
\(\mathfrak A_i\). Let \(\mathfrak A\) be the direct sum of the
\(\mathfrak A_i\): an element of \(\mathfrak A\) is a family
\((x_i)_{i\in I}\), with \(x_i\in\mathfrak A_i\), all but finitely
many of the \(x_i\) being zero. Endow \(\mathfrak A\) with the usual
algebra and pre-Hilbert-space structures, and put
\[
  (x_i)^*=(x_i^*).
\]
One immediately verifies that \(\mathfrak A\) is then a Hilbert
algebra, called the \emph{direct-sum Hilbert algebra} of the
\(\mathfrak A_i\). Its Hilbert-space completion \(\mathfrak H\) is
the Hilbert direct sum of the \(\mathfrak H_i\). One has
\[
  \mathcal U(\mathfrak A)=\prod_{i\in I}\mathcal U(\mathfrak A_i),
  \qquad
  \mathcal V(\mathfrak A)=\prod_{i\in I}\mathcal V(\mathfrak A_i).
\]
Indeed, it is clear that, for every \(x\in\mathfrak A\),
\[
  U_x\in\prod_{i\in I}\mathcal U(\mathfrak A_i),
  \qquad
  \mathcal U(\mathfrak A)
   \subset\prod_{i\in I}\mathcal U(\mathfrak A_i);
\]
% Source PDF page 084 (printed p. 77)
\phantomsection\label{srcpage:77}
similarly,
\[
  \mathcal V(\mathfrak A)
   \subset\prod_{i\in I}\mathcal V(\mathfrak A_i),
\]
and hence
\[
  \mathcal U(\mathfrak A)
   =\mathcal V(\mathfrak A)'
   \supset\prod_{i\in I}\mathcal V(\mathfrak A_i)'
   =\prod_{i\in I}\mathcal U(\mathfrak A_i),
\]
which proves the assertion. If \(E_i=P_{\mathfrak H_i}\), then the
\(E_i\) are projections of
\(\mathcal U(\mathfrak A)\cap\mathcal V(\mathfrak A)\), and
\[
  \mathcal U(\mathfrak A_i)
   =\bigl(\mathcal U(\mathfrak A)\bigr)_{E_i},
  \qquad
  \mathcal V(\mathfrak A_i)
   =\bigl(\mathcal V(\mathfrak A)\bigr)_{E_i}.
\]

Conversely, let \(\mathfrak A\) be a Hilbert algebra, and let
\((E_i)_{i\in I}\) be a family of pairwise orthogonal projections in
\(\mathcal U(\mathfrak A)\cap\mathcal V(\mathfrak A)\), with sum
\(I_{\mathfrak H}\). Put
\[
  \mathfrak H_i=E_i(\mathfrak H),
  \qquad
  \mathfrak A_i=\mathfrak A\cap\mathfrak H_i.
\]
Since \(\mathfrak H_i\) is invariant under
\(\mathcal U(\mathfrak A)\) and \(\mathcal V(\mathfrak A)\),
\[
  \mathfrak A\mathfrak A_i\subset\mathfrak A_i,
  \qquad
  \mathfrak A_i\mathfrak A\subset\mathfrak A_i.
\]
If \(x\in\mathfrak A_i\), then \(x^*\in\mathfrak A_i\), since
\[
  E_iJx=JE_ix=Jx
\]
(the corollary to \cref{prop:I-5-2}). Suppose that \(\mathfrak A\)
is completed. For \(z\in\mathfrak A\), \(E_iz\) is bounded, and
hence \(E_iz\in\mathfrak A\); it follows that
\(\mathfrak A_i=E_i(\mathfrak A)\) is dense in \(\mathfrak H_i\).
Thus \(\mathfrak A_i\) is endowed with a Hilbert-algebra structure,
its Hilbert-space completion being \(\mathfrak H_i\). It is immediate
that the \(\mathfrak A_i\) are completed. We also have
\[
  \mathcal U(\mathfrak A_i)
   \subset\bigl(\mathcal U(\mathfrak A)\bigr)_{E_i},
  \qquad
  \mathcal V(\mathfrak A_i)
   \subset\bigl(\mathcal V(\mathfrak A)\bigr)_{E_i},
\]
and therefore
\[
  \mathcal U(\mathfrak A_i)
   =\bigl(\mathcal U(\mathfrak A)\bigr)_{E_i},
  \qquad
  \mathcal V(\mathfrak A_i)
   =\bigl(\mathcal V(\mathfrak A)\bigr)_{E_i}.
\]
The direct-sum Hilbert algebra of the \(\mathfrak A_i\) is an
involutive subalgebra of \(\mathfrak A\), dense in \(\mathfrak A\).

\begin{proposition}\label{prop:I-5-8}
Let \(\mathfrak A\) be a Hilbert algebra and \(\mathfrak H\) its
Hilbert-space completion.
\begin{enumerate}[label=\textup{(\roman*)},leftmargin=*]
\item Let \(\mathfrak I\) be a two-sided ideal of \(\mathfrak A\),
and let \(\overline{\mathfrak I}\) be its closure in \(\mathfrak H\).
Then
\[
  P_{\overline{\mathfrak I}}
   \in\mathcal U(\mathfrak A)\cap\mathcal V(\mathfrak A).
\]
\item Let \(E\) be a projection of
\(\mathcal U(\mathfrak A)\cap\mathcal V(\mathfrak A)\). Then
\[
  \mathfrak K=E(\mathfrak H)\cap\mathfrak A
\]
is a self-adjoint two-sided ideal of \(\mathfrak A\), which is dense
in \(E(\mathfrak H)\) if \(\mathfrak A\) is completed.
\item Let \(E'\) be another projection of
\(\mathcal U(\mathfrak A)\cap\mathcal V(\mathfrak A)\), and put
\[
  \mathfrak K'=E'(\mathfrak H)\cap\mathfrak A.
\]
Suppose that \(\mathfrak A\) is completed. Then
\(\mathfrak K\mathfrak K'=0\) if and only if \(E\) and \(E'\) are
orthogonal.
\end{enumerate}
\end{proposition}

\begin{proof}
\textup{(i)} The spaces \(\mathfrak I\), and hence
\(\overline{\mathfrak I}\), are invariant under left and right
multiplication by the elements of \(\mathfrak A\). Therefore
\[
  P_{\overline{\mathfrak I}}
   \in\mathcal U(\mathfrak A)'
      \cap\mathcal V(\mathfrak A)'
   =\mathcal U(\mathfrak A)\cap\mathcal V(\mathfrak A).
\]

\textup{(ii)} Assertion~\textup{(ii)} was established above.

\textup{(iii)} If \(EE'=0\), then
\(\mathfrak K\mathfrak K'\subset\mathfrak K\cap\mathfrak K'=0\).
Conversely, suppose that \(\mathfrak K\mathfrak K'=0\). Let
\(x\in\mathfrak K\) and \(y\in\mathfrak K'\); then
\(x^*\in\mathfrak K\), and hence, for every
\(z\in\mathfrak A\),
% Source PDF page 085 (printed p. 78)
\phantomsection\label{srcpage:78}
\[
  (zx\mid y)=(z\mid yx^*)=0.
\]
It follows that \((x\mid y)=0\); therefore \(\mathfrak K\) and
\(\mathfrak K'\), and hence \(E\) and \(E'\), are orthogonal.
\end{proof}

\textbf{Remark.} --- On putting, for \(\lambda,\mu\in\mathbb C\),
\[
  (\lambda\mid\mu)=\lambda\overline\mu,
  \qquad
  \lambda^*=\overline\lambda,
\]
one plainly endows \(\mathbb C\) with the structure of a
one-dimensional Hilbert algebra. By taking direct sums of such
algebras, one sees that every Hilbert space contains Hilbert algebras
which are dense in it.

Let \(\mathfrak A_1,\mathfrak A_2\) be Hilbert algebras. The algebra
\(\mathfrak A_1\otimes\mathfrak A_2\) becomes an involutive algebra
on putting
\[
  \left(\sum_{i=1}^{n}x_1^i\otimes x_2^i\right)^*
   =\sum_{i=1}^{n}x_1^{i*}\otimes x_2^{i*}.
\]
It is also known that there is on
\(\mathfrak A_1\otimes\mathfrak A_2\) a unique pre-Hilbert-space
structure such that
\[
  (x_1\otimes x_2\mid y_1\otimes y_2)
   =(x_1\mid y_1)(x_2\mid y_2).
\]
One immediately verifies that
\(\mathfrak A=\mathfrak A_1\otimes\mathfrak A_2\) is then a Hilbert
algebra, called the \emph{tensor-product Hilbert algebra} of
\(\mathfrak A_1\) and \(\mathfrak A_2\). Let
\(\mathfrak H,\mathfrak H_1,\mathfrak H_2\) be the Hilbert-space
completions of
\(\mathfrak A_1\otimes\mathfrak A_2,\mathfrak A_1,\mathfrak A_2\),
respectively. Then
\[
  \mathfrak H=\mathfrak H_1\otimes\mathfrak H_2.
\]

\begin{proposition}\label{prop:I-5-9}
\begin{enumerate}[label=\textup{(\roman*)},leftmargin=*]
\item
\[
  \mathcal U(\mathfrak A_1\otimes\mathfrak A_2)
   =\mathcal U(\mathfrak A_1)\otimes\mathcal U(\mathfrak A_2),
  \qquad
  \mathcal V(\mathfrak A_1\otimes\mathfrak A_2)
   =\mathcal V(\mathfrak A_1)\otimes\mathcal V(\mathfrak A_2).
\]
\item If the characteristic projections of \(\mathfrak A_1\) and
\(\mathfrak A_2\) are \(I_{\mathfrak H_1}\) and
\(I_{\mathfrak H_2}\), respectively, then the characteristic
projection of \(\mathfrak A\) is \(I_{\mathfrak H}\).
\end{enumerate}
\end{proposition}

\begin{proof}
For \(x_1,y_1\in\mathfrak A_1\) and
\(x_2,y_2\in\mathfrak A_2\),
\begin{align*}
  (U_{x_1}\otimes U_{x_2})(y_1\otimes y_2)
   &=(U_{x_1}y_1)\otimes(U_{x_2}y_2)\\
   &=x_1y_1\otimes x_2y_2
    =(x_1\otimes x_2)(y_1\otimes y_2)\\
   &=U_{x_1\otimes x_2}(y_1\otimes y_2).
\end{align*}
Thus \(U_{x_1}\otimes U_{x_2}=U_{x_1\otimes x_2}\). Since
\(\mathcal U(\mathfrak A_1\otimes\mathfrak A_2)\) is generated by
the operators \(U_{x_1\otimes x_2}\), it follows that
\[
  \mathcal U(\mathfrak A_1\otimes\mathfrak A_2)
   \subset\mathcal U(\mathfrak A_1)\otimes\mathcal U(\mathfrak A_2).
\]
Similarly,
\[
  \mathcal V(\mathfrak A_1\otimes\mathfrak A_2)
   \subset\mathcal V(\mathfrak A_1)\otimes\mathcal V(\mathfrak A_2).
\]
Consequently,
\begin{align*}
  \mathcal U(\mathfrak A_1\otimes\mathfrak A_2)
   &=\mathcal V(\mathfrak A_1\otimes\mathfrak A_2)'\\
   &\supset
     \mathcal V(\mathfrak A_1)'
       \otimes\mathcal V(\mathfrak A_2)'\\
   &=\mathcal U(\mathfrak A_1)\otimes\mathcal U(\mathfrak A_2),
\end{align*}
which proves~\textup{(i)}.

Now let \(a_1\) be a central element of \(\mathfrak H_1\), and
\(a_2\) a central element of \(\mathfrak H_2\). For
\(x_1,y_1\in\mathfrak A_1\) and
\(x_2,y_2\in\mathfrak A_2\),
\begin{align*}
 &(a_1\otimes a_2\mid
       (x_1\otimes x_2)(y_1\otimes y_2))\\
 &\quad=(a_1\otimes a_2\mid x_1y_1\otimes x_2y_2)
       =(a_1\mid x_1y_1)(a_2\mid x_2y_2)\\
 &\quad=(a_1\mid y_1x_1)(a_2\mid y_2x_2)
       =(a_1\otimes a_2\mid
          (y_1\otimes y_2)(x_1\otimes x_2)),
\end{align*}
% Source PDF page 086 (printed p. 79)
\phantomsection\label{srcpage:79}
so that \(a_1\otimes a_2\) is a central element of \(\mathfrak H\).
This being so, if the characteristic projections of
\(\mathfrak A_1\) and \(\mathfrak A_2\) are equal to \(1\), the
elements \(T_1a_1\) (respectively, \(T_2a_2\)), where
\(T_1\in\mathcal U(\mathfrak A_1)\) (respectively,
\(T_2\in\mathcal U(\mathfrak A_2)\)) and where \(a_1\)
(respectively, \(a_2\)) is a central element of \(\mathfrak H_1\)
(respectively, \(\mathfrak H_2\)), form a total set in
\(\mathfrak H_1\) (respectively, \(\mathfrak H_2\)). Hence the
elements
\[
  T_1a_1\otimes T_2a_2
   =(T_1\otimes T_2)(a_1\otimes a_2)
\]
form a total set in \(\mathfrak H\). Since
\(T_1\otimes T_2\in
\mathcal U(\mathfrak A_1\otimes\mathfrak A_2)\), the
characteristic projection of \(\mathfrak A\) is equal to \(1\).
\end{proof}

\begin{definition}\label{def:I-5-7}
Let \(\mathfrak H\) be a complex Hilbert space, and let
\(\mathcal A\) be a von Neumann algebra on \(\mathfrak H\). The
algebra \(\mathcal A\) is said to be \emph{standard} if there exists
a Hilbert algebra \(\mathfrak A\), dense in \(\mathfrak H\), such
that
\[
  \mathcal A=\mathcal U(\mathfrak A).
\]
\end{definition}

The results of this numbered paragraph then imply the following
proposition.

\begin{proposition}\label{prop:I-5-10}
\begin{enumerate}[label=\textup{(\roman*)},leftmargin=*]
\item If \(\mathcal A\) is standard, then \(\mathcal A'\) is
standard.
\item If \(\mathcal A\) and \(\mathcal B\) are standard, then
\(\mathcal A\otimes\mathcal B\) is standard.
\item Let
\[
  \mathcal A=\prod_{i\in I}\mathcal A_i.
\]
For \(\mathcal A\) to be standard, it is necessary and sufficient
that the \(\mathcal A_i\) be standard.
\end{enumerate}
\end{proposition}

Later (\hyperref[sec:I-6]{\S~6} and Chapter~III,
\S~1), we shall be able to characterize standard von Neumann
algebras completely.

For another ``elementary operation,'' cf. Exercise~3 below.

\noindent\textbf{Bibliography:} \ArticleRef{90}, \ArticleRef{101},
\ArticleRef{116}.

\noindent\textit{Exercises.}---
\begin{enumerate}[label=\arabic*.,leftmargin=*,itemsep=0.5\baselineskip]
\item\label{ex:I-5-1}
Let \(\mathfrak A\) be a Hilbert algebra. If
\[
  \mathcal Z
   =\mathcal U(\mathfrak A)\cap\mathcal V(\mathfrak A)
\]
is of countable type, there exists a bounded separating element for
\(\mathcal Z\). [Let \((x_i)_{i\in I}\) be a maximal family of
nonzero bounded elements such that the
\[
  E_i=E_{x_i}^{\mathcal Z'}
\]
are pairwise orthogonal. Show that
\[
  \sum_{i\in I}E_i=1,
\]
that \(I\) is countable, that one may suppose
\[
  \sum_{i\in I}\lVert x_i\rVert^2<+\infty,
  \qquad
  \sup_{i\in I}\lVert U_{x_i}\rVert<+\infty,
\]
and that
\[
  x=\sum_{i\in I}x_i
\]
is then a bounded separating element for \(\mathcal Z\).]

\item\label{ex:I-5-2}
Let \(\mathfrak A_1\) and \(\mathfrak A_2\) be two Hilbert
algebras, and let \(a_1\) (respectively, \(a_2\)) be a bounded
element relative to \(\mathfrak A_1\) (respectively,
\(\mathfrak A_2\)). Then \(a_1\otimes a_2\) is bounded relative to
\(\mathfrak A_1\otimes\mathfrak A_2\), and
\[
  U_{a_1\otimes a_2}=U_{a_1}\otimes U_{a_2},
  \qquad
  V_{a_1\otimes a_2}=V_{a_1}\otimes V_{a_2}.
\]

\item\label{ex:I-5-3}
Let \(\mathfrak A\) be a completed Hilbert algebra, let
\(\mathfrak H\) be its completed Hilbert space, and let \(J\) be
the involution of \(\mathfrak H\) canonically defined by
\(\mathfrak A\).
\begin{enumerate}[label=\textit{\alph*.},leftmargin=2em]
\item For \(T\in\mathcal U(\mathfrak A)\) (respectively,
\(T'\in\mathcal V(\mathfrak A)\)), the space
\(T(\mathfrak A)\) (respectively, \(T'(\mathfrak A)\)) is a right
(respectively, left) ideal of \(\mathfrak A\), and
\[
  T(\mathfrak A)\cap JTJ(\mathfrak A)
\]
is an involutive subalgebra of \(\mathfrak A\).
\item Taking \(T\) to be a projection
\(E\in\mathcal U(\mathfrak A)\), and putting \(E'=JEJ\),
\[
  E(\mathfrak A)\cap E'(\mathfrak A)=EE'(\mathfrak A)
\]
% Source PDF page 087 (printed p. 80)
\phantomsection\label{srcpage:80}
is a completed Hilbert algebra \(\mathfrak A_1\), dense in
\[
  E(\mathfrak H)\cap E'(\mathfrak H)=EE'(\mathfrak H),
\]
and
\[
  \mathcal U(\mathfrak A_1)
   =\bigl((\mathcal U(\mathfrak A))_E\bigr)_{E'},
  \qquad
  \mathcal V(\mathfrak A_1)
   =\bigl((\mathcal V(\mathfrak A))_{E'}\bigr)_E.
\]
\ArticleRef{116}
\end{enumerate}

\item\label{ex:I-5-4}
Let \(\mathfrak A\) be a Hilbert algebra, let \(\mathfrak H\) be
its completed Hilbert space, let \(J\) be the involution of
\(\mathfrak H\) canonically defined by \(\mathfrak A\), and let
\(\mathcal H\) be the group of automorphisms of
\(\mathfrak A\).
\begin{enumerate}[label=\textit{\alph*.},leftmargin=2em]
\item Let \(\mathcal G\) be the group of unitary operators of
\(\mathcal U(\mathfrak A)\). As \(U\) runs through
\(\mathcal G\), \(UJUJ\) runs through a group
\(\mathcal G_1\subset\mathcal H\) of unitary operators on
\(\mathfrak H\), and
\[
  U\longmapsto UJUJ
\]
is a homomorphism of the group \(\mathcal G\) (regarded as a
nontopological group) onto \(\mathcal G_1\). The kernel of this
homomorphism is the group of unitary operators in the centre of
\(\mathcal U(\mathfrak A)\).

\item For an element \(a\in\mathfrak H\) to be central, it is
necessary and sufficient that
\[
  UJUJa=a
  \qquad\text{for every }U\in\mathcal G.
\]

\item For every \(a\in\mathfrak H\), let \(\mathcal K_a\) be the
convex set generated by the \(UJUJa\), \(U\in\mathcal G\), and let
\(\mathcal K'_a\) be the closure of \(\mathcal K_a\) in
\(\mathfrak H\). Show that \(\mathcal K'_a\) is invariant under
\(\mathcal G_1\), and hence that the point \(b\in\mathcal K'_a\)
at minimum distance from zero is invariant under
\(\mathcal G_1\), so that \(b\in\mathfrak Z\), the set of central
elements. On the other hand, put \(c=P_{\mathfrak Z}a\); show that
\[
  \mathcal K'_a
   \subset c+(\mathfrak H\ominus\mathfrak Z),
\]
and hence that \(b=c\).

\item Let \(a\) be a bounded element of \(\mathfrak H\). Let
\(U_1,U_2,\ldots,U_n\) be elements of \(\mathcal G\), and let
\(\lambda_1,\lambda_2,\ldots,\lambda_n\) be nonnegative real
numbers such that
\[
  \sum_{i=1}^{n}\lambda_i=1,
  \qquad
  b=\sum_{i=1}^{n}\lambda_iU_iJU_iJa.
\]
Show that
\[
  U_b=\sum_{i=1}^{n}\lambda_iU_iU_aU_i^{-1}.
\]
Deduce that
\[
  \lVert V_xc\rVert
   \leq\lVert U_a\rVert\,\lVert x\rVert
\]
for every \(x\in\mathfrak A\) and every \(c\in\mathcal K_a\), and
therefore for every \(c\in\mathcal K'_a\). Conclude that all the
elements of \(\mathcal K'_a\), and in particular
\(P_{\mathfrak Z}a\), are bounded.

\item Show that the set of bounded central elements is dense in
\(\mathfrak Z\). \ArticleRef{29}, \ArticleRef{106}
\end{enumerate}

\item\label{ex:I-5-5}
Let \(G\) be a locally compact group, and let \(\mathfrak A\) be
the algebra of complex-valued continuous functions of compact
support on \(G\), multiplication being defined by convolution of
functions. Let \(dg\) be the element of left-invariant Haar measure
on \(G\). Let \(\Delta(g)\) be the modulus of \(G\), so that
\[
  d(g^{-1})=\Delta(g)^{-1}\,dg.
\]
For \(x\in\mathfrak A\), put
\[
  x^*(g)=\overline{x(g^{-1})},
  \qquad
  x^\wedge(g)=\Delta(g)^{-1/2}x(g).
\]
For \(x,y\in\mathfrak A\), put
\[
  (x\mid y)
   =\int_Gx(g)\overline{y(g)}\Delta(g)^{-1/2}\,dg.
\]
\begin{enumerate}[label=\textit{\alph*.},leftmargin=2em]
\item Show that \(\mathfrak A\) is a quasi-Hilbert algebra. [To
prove axiom~\textup{(iii)}, show that, for
\(x,y,z\in\mathfrak A\),
\[
  |(xy\mid z)|
   \leq \lVert y\rVert\,\lVert z\rVert
          \int_G|x(g)|\Delta(g)^{-1/4}\,dg,
\]
and hence
\[
  \lVert xy\rVert
   \leq \lVert y\rVert
          \int_G|x(g)|\Delta(g)^{-1/4}\,dg.
\]
To prove axiom~\textup{(v)}, show that, if \(a\) and \(b\) are
functions on \(G\) which are square-integrable with respect to the
measure \(\Delta(g)^{-1/2}dg\), the equality
\[
  (a\mid xy)=(b\mid x^\wedge y^\wedge)
  \qquad(x,y\in\mathfrak A)
\]
implies equality of the measures
\[
  a(g)\,dg
  \qquad\text{and}\qquad
  \Delta(g)^{-1/2}b(g)\,dg.
\]]

% Source PDF page 088 (printed p. 81)
\phantomsection\label{srcpage:81}
\item For \(g\in G\) and \(x\in\mathfrak H\), where
\(\mathfrak H\) is the completed Hilbert space of
\(\mathfrak A\), put
\[
  (U_gx)(g')=x(g^{-1}g')\Delta(g)^{1/4},
  \qquad
  (V_gx)(g')=x(g'g)\Delta(g)^{1/4}.
\]
Show that the \(U_g\) (respectively, the \(V_g\)) are unitary
operators on \(\mathfrak H\) which generate
\(\mathcal U(\mathfrak A)\) (respectively,
\(\mathcal V(\mathfrak A)\)).
\end{enumerate}

We shall return to this example in
\hyperref[no:III-7-6]{Chapter~III, \S~7, no.~6}.

\noindent\textbf{Bibliography:} \ArticleRef{1}, \ArticleRef{2},
\ArticleRef{13}, \ArticleRef{27}, \ArticleRef{29}, \ArticleRef{53},
\ArticleRef{56}, \ArticleRef{67}, \ArticleRef{90}, \ArticleRef{98},
\ArticleRef{99}, \TreatiseRef{9}.

\item\label{ex:I-5-6}
Let \(\mathfrak A\) be a Hilbert algebra and \(\mathfrak H\) its
completed Hilbert space. For \(a\in\mathfrak H\), define linear maps
\(u_a,v_a\) from \(\mathfrak A\) into \(\mathfrak H\) by
\[
  u_ax=V_xa,
  \qquad
  v_ax=U_xa.
\]
\begin{enumerate}[label=\textit{\alph*.},leftmargin=2em]
\item The operators \(u_a^*,v_a^*\) extend \(u_{Ja},v_{Ja}\),
respectively. Thus \(u_a,v_a\) admit the smallest closed extensions
\[
  U'_a=u_a^{**},
  \qquad
  V'_a=v_a^{**}.
\]
On the other hand, put
\[
  U_a=u_{Ja}^*,
  \qquad
  V_a=v_{Ja}^*.
\]
The operators \(U_a,V_a\) are closed, extend \(U'_a,V'_a\),
respectively, and
\[
  U'_a=U_{Ja}^*,
  \qquad
  U_a=(U'_{Ja})^*,
  \qquad
  V'_a=V_{Ja}^*,
  \qquad
  V_a=(V'_{Ja})^*.
\]

\item
\[
  U'_a\mathrel{\eta}\mathcal U(\mathfrak A),
  \qquad
  U_a\mathrel{\eta}\mathcal U(\mathfrak A),
  \qquad
  V'_a\mathrel{\eta}\mathcal V(\mathfrak A),
  \qquad
  V_a\mathrel{\eta}\mathcal V(\mathfrak A)
\]
(\hyperref[ex:I-1-10]{\S~1, Exercise~10}). [Show that
\[
  u_aV_xy=V_xu_ay
  \qquad(x,y\in\mathfrak A),
\]
and deduce that
\[
  U'_aV_xb=V_xU'_ab
\]
for \(x\in\mathfrak A\) and \(b\) in the domain of \(U'_a\).]

\item If \(y\in\mathfrak H\), there exists an increasing family of
projections \((E(r))_{r\geq0}\) of
\(\mathcal V(\mathfrak A)\), tending strongly to \(1\) as
\(r\to+\infty\), such that \(E(r)y\) is bounded for every \(r\).
[Let
\[
  U'_{Jy}=WT
\]
be the polar decomposition of \(U'_{Jy}\), where \(T\) is
self-adjoint and nonnegative. Take for \(E(r)\) the spectral
projections of \(T\).]

\item Let \(\mathfrak B\) be the set of bounded elements. Then
\(\mathfrak B\) is contained in the domain of \(U'_a\). Let \(y\)
belong to the domain of \(U'_a\), and let \((E(r))\) have the
properties in~\textit{c}. Then
\[
  U'_aE(r)y=E(r)U'_ay.
\]

\item One has
\[
  U_a=U'_a,
  \qquad
  V_a=V'_a.
\]
(Use~\textit{d}.) \ArticleRef{90}, \ArticleRef{99},
\ArticleRef{459}
\end{enumerate}
\end{enumerate}

\section{Traces}\label{sec:I-6}

\NumberedParagraph{1}{Definition of Traces}\label{no:I-6-1}

\begin{definition}\label{def:I-6-1}
Let \(\mathcal A\) be a von Neumann algebra. A \emph{trace} on
\(\mathcal A^+\) is a function \(\varphi\) defined on
\(\mathcal A^+\), with finite or infinite nonnegative values, and
having the following properties:
\begin{enumerate}[label=\textup{(\roman*)},leftmargin=*]
\item If \(S,T\in\mathcal A^+\), then
\[
  \varphi(S+T)=\varphi(S)+\varphi(T).
\]
\item If \(S\in\mathcal A^+\) and \(\lambda\geq0\), then
\[
  \varphi(\lambda S)=\lambda\varphi(S)
\]
(with the convention \(0\cdot(+\infty)=0\)).
\item If \(S\in\mathcal A^+\) and \(U\) is a unitary operator of
\(\mathcal A\), then
\[
  \varphi(USU^{-1})=\varphi(S).
\]
\end{enumerate}

% Source PDF page 089 (printed p. 82)
\phantomsection\label{srcpage:82}
The trace \(\varphi\) is said to be \emph{faithful} if
\[
  S\in\mathcal A^+,
  \qquad
  \varphi(S)=0
\]
imply \(S=0\).

The trace \(\varphi\) is said to be \emph{finite} if
\[
  \varphi(S)<+\infty
  \qquad\text{for every }S\in\mathcal A^+.
\]

The trace \(\varphi\) is said to be \emph{semifinite} if, for
every \(S\in\mathcal A^+\), \(\varphi(S)\) is the supremum of the
numbers \(\varphi(T)\), where \(T\in\mathcal A^+\), \(T\leq S\),
and \(\varphi(T)<+\infty\).

The trace \(\varphi\) is said to be \emph{normal} if, for every
increasing directed set \(\mathcal F\subset\mathcal A^+\) with
supremum \(S\in\mathcal A^+\), \(\varphi(S)\) is the supremum of
\(\varphi(\mathcal F)\).
\end{definition}

\begin{proposition}\label{prop:I-6-1}
Let \(\mathcal A\) be a von Neumann algebra, and let \(\varphi\) be
a trace on \(\mathcal A^+\). The set of
\(T\in\mathcal A^+\) such that \(\varphi(T)<+\infty\) is the
positive part of a two-sided ideal \(\mathfrak m\) of
\(\mathcal A\). There exists a unique linear functional
\(\dot\varphi\) on \(\mathfrak m\) which agrees with
\(\varphi\) on \(\mathfrak m^+\), and
\[
  \dot\varphi(ST)=\dot\varphi(TS)
  \qquad(S\in\mathfrak m,\ T\in\mathcal A).
\]
Finally, let \(S\in\mathfrak m\). If \(\varphi\) is normal, then
the linear functional
\[
  T\longmapsto\dot\varphi(ST)
\]
on \(\mathcal A\) is ultra-weakly continuous.
\end{proposition}

\begin{proof}
The set of \(T\in\mathcal A^+\) such that
\(\varphi(T)<+\infty\) satisfies the conditions of
\hyperref[prop:I-1-10]{\S~1, Proposition~10}, and is therefore the
positive part of a unique two-sided ideal \(\mathfrak m\) of
\(\mathcal A\). Every element of \(\mathfrak m\) is a linear
combination of elements of \(\mathfrak m^+\)
(\hyperref[prop:I-1-9]{\S~1, Proposition~9}), and the properties of
\(\varphi\) immediately imply that there exists on
\(\mathfrak m\) a unique linear functional \(\dot\varphi\) which
agrees with \(\varphi\) on \(\mathfrak m^+\). If
\(S\in\mathfrak m\) and \(U\) is a unitary operator of
\(\mathcal A\), then
\[
  \dot\varphi(USU^{-1})=\dot\varphi(S),
\]
by the properties of \(\varphi\) when \(S\in\mathfrak m^+\), and
by linearity in the general case. Hence
\[
  \dot\varphi(SU)
   =\dot\varphi(USUU^{-1})
   =\dot\varphi(US).
\]
By \hyperref[prop:I-1-3]{\S~1, Proposition~3}, it follows that
\[
  \dot\varphi(ST)=\dot\varphi(TS)
  \qquad(S\in\mathfrak m,\ T\in\mathcal A).
\]
Finally, let \(S\in\mathfrak m\), and put
\[
  \psi(T)=\dot\varphi(ST)
  \qquad(T\in\mathcal A).
\]
Let us show that \(\psi\) is ultra-weakly continuous if
\(\varphi\) is normal. We may restrict ourselves to the case where
\(S\in\mathfrak m^+\), and we shall show that \(\psi\) is then
positive and normal; this, together with
\hyperref[thm:I-4-1]{\S~4, Theorem~1}, will complete the proof.
First, if \(T\in\mathcal A^+\), then
\[
  \psi(T)
   =\dot\varphi(ST^{1/2}T^{1/2})
   =\dot\varphi(T^{1/2}ST^{1/2})
   \geq0
  \qquad\bigl(T^{1/2}ST^{1/2}\geq0\bigr),
\]
so \(\psi\) is positive. Now let
\(\mathcal F\subset\mathcal A^+\) be an increasing directed set
with supremum \(T\in\mathcal A^+\). Then
\(S^{1/2}\mathcal F S^{1/2}\) is increasing and directed, is
majorized by \(S^{1/2}TS^{1/2}\), and
\(S^{1/2}TS^{1/2}\) belongs to the strong closure of
\(S^{1/2}\mathcal F S^{1/2}\); hence
\(S^{1/2}TS^{1/2}\) is the supremum of
\(S^{1/2}\mathcal F S^{1/2}\). Observe that
\[
  S^{1/2}(S^{1/2})^*=S\in\mathfrak m,
\]
and that, for \(R\in\mathcal A\),
\[
  (RS^{1/2})(RS^{1/2})^*=RSR^*\in\mathfrak m.
\]
Thus, by \hyperref[prop:I-1-11]{\S~1, Proposition~11},
\[
  S^{1/2}\mathcal F S^{1/2}\subset\mathfrak m,
  \qquad
  S^{1/2}TS^{1/2}\in\mathfrak m,
\]
% Source PDF page 090 (printed p. 83)
\phantomsection\label{srcpage:83}
and
\[
  \psi(T)
   =\dot\varphi(S^{1/2}S^{1/2}T)
   =\dot\varphi(S^{1/2}TS^{1/2})
\]
is the supremum of
\[
  \dot\varphi(S^{1/2}\mathcal F S^{1/2})
   =\dot\varphi(S^{1/2}S^{1/2}\mathcal F)
   =\psi(\mathcal F),
\]
by the normality of \(\varphi\).
\end{proof}

By an abuse of language, the linear functional \(\dot\varphi\) on
\(\mathfrak m\) is sometimes itself called a trace. If
\(\varphi\) is finite, \(\dot\varphi\) is a positive linear
functional on \(\mathcal A\), and the results of \S~4 are
applicable. But it would not be sufficient for what follows to
consider only finite traces.

\begin{corollary}\label{cor:I-6-1}
Let \(\varphi\) be a function on \(\mathcal A^+\), with finite or
infinite nonnegative values, such that
\[
  \varphi(S+T)=\varphi(S)+\varphi(T),
  \qquad
  \varphi(\lambda S)=\lambda\varphi(S)
\]
for \(S,T\in\mathcal A^+\) and \(\lambda\geq0\). For
\(\varphi\) to be a trace, it is necessary and sufficient that
\[
  \varphi(R^*R)=\varphi(RR^*)
  \qquad\text{for every }R\in\mathcal A.
\]
\end{corollary}

\begin{proof}
Suppose first that \(\varphi\) is a trace. Let
\(R\in\mathcal A\), and let
\[
  R=U|R|
\]
be its polar decomposition. Then
\[
  R^*R=|R|^2=U^*U|R|^2,
  \qquad
  RR^*=U|R|^2U^*.
\]
If \(\mathfrak m\) denotes the two-sided ideal on which
\(\dot\varphi\) is defined, the condition
\(R^*R\in\mathfrak m\) implies \(RR^*\in\mathfrak m\), and is
therefore equivalent to it; moreover, when it is satisfied,
\[
  \varphi(RR^*)
   =\dot\varphi(U|R|^2U^*)
   =\dot\varphi(U^*U|R|^2)
   =\varphi(R^*R).
\]
Thus
\[
  \varphi(RR^*)=\varphi(R^*R)
  \qquad\text{for every }R\in\mathcal A.
\]
Conversely, suppose that
\[
  \varphi(RR^*)=\varphi(R^*R)
  \qquad\text{for every }R\in\mathcal A.
\]
Let \(S\in\mathcal A^+\), and let \(V\) be a unitary operator of
\(\mathcal A\). Then
\[
  \varphi(VSV^{-1})
   =\varphi\bigl((VS^{1/2})(VS^{1/2})^*\bigr)
   =\varphi\bigl((VS^{1/2})^*(VS^{1/2})\bigr)
   =\varphi(S),
\]
so \(\varphi\) is a trace.
\end{proof}

\begin{corollary}\label{cor:I-6-2}
Let \(\varphi\) be a normal trace on \(\mathcal A^+\). Then
\(\mathcal A\) is canonically identified with the product of three
von Neumann algebras \(\mathcal A_1,\mathcal A_2,\mathcal A_3\),
in such a way that \(\varphi\) induces on
\(\mathcal A_1^+,\mathcal A_2^+,\mathcal A_3^+\) normal traces
\(\varphi_1,\varphi_2,\varphi_3\) having the following properties:
\begin{enumerate}[label=\arabic*\textsuperscript{\(\circ\)}.,leftmargin=*]
\item \(\varphi_1\) is faithful and semifinite;
\item \(\varphi_2=0\);
\item \(\varphi_3(S)=+\infty\) for every nonzero
\(S\in\mathcal A_3^+\).
\end{enumerate}
\end{corollary}

\begin{proof}
By \hyperref[prop:I-1-10]{\S~1, Proposition~10}, the set of
\(T\in\mathcal A^+\) such that \(\varphi(T)=0\) is the positive
part of a two-sided ideal \(\mathfrak n\) of \(\mathcal A\), which
is evidently contained in the ideal \(\mathfrak m\) on which
\(\dot\varphi\) is defined. Let \(E\) (respectively, \(F\)) be the
greatest projection in the strong closure of \(\mathfrak m\)
(respectively, \(\mathfrak n\)). Then \(E\geq F\), and \(E,F\) are
projections of the centre of \(\mathcal A\)
(\hyperref[thm:I-3-2]{\S~3, Theorem~2},
\hyperref[cor:I-3-3]{Corollary~3}). Put
\[
  \mathcal A_1=\mathcal A_{E-F},
  \qquad
  \mathcal A_2=\mathcal A_F,
  \qquad
  \mathcal A_3=\mathcal A_{1-E}.
\]
% Source PDF page 091 (printed p. 84)
\phantomsection\label{srcpage:84}
Then \(\mathcal A\) is identified with the product of
\(\mathcal A_1,\mathcal A_2,\mathcal A_3\), and the asserted properties
of \(\varphi_1,\varphi_2,\varphi_3\) follow at once from
\hyperref[cor:I-3-5]{\S~3, Theorem~2, Corollary~5} and from the
normality of \(\varphi\).
\end{proof}

This corollary reduces the study of normal traces to that of faithful,
semifinite normal traces. The projection \(1-F\) is called the
\emph{support} of \(\varphi\). When \(\varphi\) is finite, the support
of \(\varphi\) is identical with the support of \(\dot\varphi\), in the
sense of \hyperref[def:I-4-3]{\S~4, Definition~3}.

\begin{corollary}\label{cor:I-6-3}
Let \(\varphi\) be a normal trace on \(\mathcal A^+\). For
\(\varphi\) to be semifinite, it is necessary and sufficient that
every nonzero element of \(\mathcal A^+\) majorize a nonzero element
\(T\) of \(\mathcal A^+\) such that \(\varphi(T)<+\infty\).
\end{corollary}

\begin{proof}
The condition is clearly necessary. On the other hand, with the
notation of \cref{cor:I-6-2}, it implies that
\(\mathcal A_3=0\), and hence that \(\varphi\) is semifinite.
\end{proof}

\begin{proposition}\label{prop:I-6-2}
Let \(\mathcal A\) be a von Neumann algebra, let \(\varphi\) be a
normal trace on \(\mathcal A^+\), and let \((T_i)_{i\in I}\) be a
family of elements of \(\mathcal A^+\) such that
\[
  \sum_{i\in I}T_i=1
\]
in the weak topology. For every \(T\in\mathcal A^+\), one has
\[
  \varphi(T)
    =\sum_{i\in I}\varphi\bigl(T_i^{1/2}TT_i^{1/2}\bigr).
\]
\end{proposition}

\begin{proof}
For every subset \(J\) of \(I\), put
\[
  T_J=\sum_{i\in J}T_i.
\]
Let \(\mathcal F\) be the set of finite subsets of \(I\), and let
\(\mathfrak m\) be the ideal of definition of \(\dot\varphi\). If
\(T\in\mathfrak m^+\), then, for \(J\in\mathcal F\),
\[
  \dot\varphi(TT_J)
    =\sum_{i\in J}\dot\varphi(TT_i)
    =\sum_{i\in J}\varphi\bigl(T_i^{1/2}TT_i^{1/2}\bigr).
\]
On the other hand, by Appendix~II, \(1\) is the ultra-strong limit of
the \(T_J\); hence, by \cref{prop:I-6-1}, \(\varphi(T)\) is the
limit of the \(\dot\varphi(TT_J)\). Consequently,
\[
  \varphi(T)
    =\sum_{i\in I}\varphi\bigl(T_i^{1/2}TT_i^{1/2}\bigr).
\]
Let us pass to the general case. By \cref{cor:I-6-2}, it is enough
to consider the case in which \(\varphi\) is semifinite. Let
\(\mathcal G\) be an increasing directed subset of
\(\mathfrak m^+\) with supremum \(T\)
(\hyperref[cor:I-3-5]{\S~3, Theorem~2, Corollary~5}). Then
\begin{align*}
  \varphi(T)
    &=\sup_{S\in\mathcal G}\varphi(S) \\
    &=\sup_{S\in\mathcal G}
       \left(\sum_{i\in I}
       \varphi\bigl(T_i^{1/2}ST_i^{1/2}\bigr)\right) \\
    &=\sup_{\substack{S\in\mathcal G\\J\in\mathcal F}}
       \left(\sum_{i\in J}
       \varphi\bigl(T_i^{1/2}ST_i^{1/2}\bigr)\right) \\
    &=\sup_{J\in\mathcal F}
       \left(\sum_{i\in J}
       \varphi\bigl(T_i^{1/2}TT_i^{1/2}\bigr)\right) \\
    &=\sum_{i\in I}\varphi\bigl(T_i^{1/2}TT_i^{1/2}\bigr).
\end{align*}
\end{proof}

% Source PDF page 092 (printed p. 85)
\phantomsection\label{srcpage:85}
\begin{corollary*}
Let \(\mathcal A\) be a von Neumann algebra on \(\mathfrak H\), and
let \(\varphi\) be a normal trace on \(\mathcal A^+\). There exists a
family \((x_i)_{i\in I}\) of vectors of \(\mathfrak H\) such that
\[
  \varphi=\sum_{i\in I}\omega_{x_i}
  \qquad\text{on }\mathcal A^+.
\]
\end{corollary*}

\begin{proof}
By \cref{cor:I-6-2}, it is enough to consider the following two cases:
\(1^\circ\) \(\varphi(T)=+\infty\) for every nonzero
\(T\in\mathcal A^+\); \(2^\circ\) \(\varphi\) is semifinite. In the
first case one may take for \((x_i)_{i\in I}\) the family of all the
vectors of \(\mathfrak H\). Consider the second case. In
\cref{prop:I-6-2}, one may suppose that
\(T_i\in\mathfrak m\) for every \(i\)
(\hyperref[cor:I-3-5]{\S~3, Theorem~2, Corollary~5}). Then each linear
functional
\[
  T\longmapsto\dot\varphi(TT_i)
    =\dot\varphi\bigl(T_i^{1/2}TT_i^{1/2}\bigr)
\]
is positive and normal on \(\mathcal A\), and it remains only to apply
\hyperref[thm:I-4-1]{\S~4, Theorem~1}.
\end{proof}

In some papers, normal traces, or faithful traces, are called traces.
Instead of ``semifinite,'' one sometimes says ``essential.'' Instead
of ``trace,'' one sometimes says ``pseudo-trace.''

Let \(\mathcal A\) be a von Neumann algebra, let \(\mathfrak m\) be a
two-sided ideal of \(\mathcal A\), and let \(\varphi\) be a positive
linear functional on \(\mathfrak m\) such that
\[
  \varphi(ST)=\varphi(TS)
  \qquad(S\in\mathfrak m,\ T\in\mathcal A).
\]
In certain cases, the restriction of \(\varphi\) to
\(\mathfrak m^+\) can then be extended to a normal trace on
\(\mathcal A^+\) (cf. Chapter~III, \S~1, Exercise~11).

On an abelian von Neumann algebra there are, in general, finite
nonnormal traces, as is easily seen. For the case of factors, see
Exercise~6 and Chapter~III, \S~2, the corollary to Proposition~15.

\noindent\textit{Bibliography.}\quad
\ArticleRef{12}, \ArticleRef{14}, \ArticleRef{19}, \ArticleRef{65},
\ArticleRef{66}, \ArticleRef{70}, \ArticleRef{78}, \ArticleRef{80},
\ArticleRef{89}, \ArticleRef{101}, \ArticleRef{117}.

\NumberedParagraph{2}{Traces and Hilbert Algebras}
\label{no:I-6-2}

\begin{theorem}\label{thm:I-6-1}
Let \(\mathfrak A\) be a Hilbert algebra, and let \(\mathfrak H\) be
the Hilbert-space completion of \(\mathfrak A\). For
\(S\in\mathcal U(\mathfrak A)^+\)
[respectively, \(S\in\mathcal V(\mathfrak A)^+\)], put
\[
  \varphi(S)=
  \begin{cases}
    (a\mid a),
      &\text{if }S^{1/2}=U_a
       \text{ (respectively, }S^{1/2}=V_a\text{) for a bounded }
       a\in\mathfrak H,\\
    +\infty,&\text{otherwise.}
  \end{cases}
\]
Then \(\varphi\) is a faithful, semifinite, normal trace on
\(\mathcal U(\mathfrak A)^+\)
[respectively, on \(\mathcal V(\mathfrak A)^+\)]. The two-sided ideal
of those \(T\in\mathcal U(\mathfrak A)\)
[respectively, \(T\in\mathcal V(\mathfrak A)\)] which can be written
in the form \(U_a\) (respectively, \(V_a\)) with \(a\) bounded is
identical with the two-sided ideal of those
\(T\in\mathcal U(\mathfrak A)\)
[respectively, \(T\in\mathcal V(\mathfrak A)\)] such that
\(\varphi(T^*T)<+\infty\). If \(a\) and \(b\) are bounded elements,
then
\[
  \dot\varphi(U_b^*U_a)=(a\mid b)
  \qquad
  [\text{respectively, }\dot\varphi(V_b^*V_a)=(a\mid b)].
\]
\end{theorem}

\begin{proof}
We shall consider only the case of \(\mathcal U(\mathfrak A)\). It is
clear that \(\varphi\) is a function with nonnegative values on
\(\mathcal U(\mathfrak A)^+\) and has property~\textup{(ii)} of
\cref{def:I-6-1}. On the other hand, for
\(T\in\mathcal U(\mathfrak A)\) to be of the form \(U_a\), with
\(a\) bounded, it is necessary and sufficient that \(|T|\) be of the
form \(U_b\), with \(b\) bounded (since the \(U_a\) form a two-sided
ideal of \(\mathcal U(\mathfrak A)\)), and hence that
\[
  \varphi(T^*T)=\varphi(|T|^2)<+\infty.
\]

% Source PDF page 093 (printed p. 86)
\phantomsection\label{srcpage:86}
Let us prove property~\textup{(i)} of \cref{def:I-6-1}. Let
\[
  R\in\mathcal U(\mathfrak A)^+,
  \qquad
  S\in\mathcal U(\mathfrak A)^+,
  \qquad
  T=R+S.
\]
One has
\[
  R^{1/2}=AT^{1/2},
  \qquad
  S^{1/2}=BT^{1/2},
\]
where \(A\) and \(B\) are elements of \(\mathcal U(\mathfrak A)\)
which vanish on
\(\mathfrak H\ominus\overline{T(\mathfrak H)}\)
(\hyperref[lem:I-1-2]{\S~1, Lemma~2}). Put
\[
  C=A^*A+B^*B.
\]
Then
\begin{align*}
  T
    &=R^{1/2}R^{1/2}+S^{1/2}S^{1/2} \\
    &=T^{1/2}A^*AT^{1/2}+T^{1/2}B^*BT^{1/2} \\
    &=T^{1/2}CT^{1/2}.
\end{align*}
Consequently,
\[
  (CT^{1/2}x\mid T^{1/2}x)
    =(T^{1/2}x\mid T^{1/2}x)
  \qquad(x\in\mathfrak H),
\]
and hence
\[
  (Cu\mid u)=(u\mid u)
  \qquad\bigl(u\in\overline{T(\mathfrak H)}\bigr).
\]
Since \(C\) vanishes on
\(\mathfrak H\ominus\overline{T(\mathfrak H)}\), it follows that
\(C\) is the projection onto \(\overline{T(\mathfrak H)}\), and thus
\[
  T^{1/2}=CT^{1/2}.
\]
This being so, if \(T^{1/2}=U_a\), with \(a\) bounded, then
\[
  R^{1/2}=AU_a=U_{Aa},
  \qquad
  S^{1/2}=BU_a=U_{Ba}.
\]
Moreover, \(a\in\overline{T(\mathfrak H)}\), so that \(Ca=a\);
hence
\begin{align*}
  \varphi(R)+\varphi(S)
    &=(Aa\mid Aa)+(Ba\mid Ba) \\
    &=(Ca\mid a)=(a\mid a)=\varphi(T).
\end{align*}
If, on the contrary, \(\varphi(T)=+\infty\), then
\(\varphi(R)=+\infty\) or \(\varphi(S)=+\infty\). Indeed, if one had
\(R^{1/2}=U_b\) and \(S^{1/2}=U_c\), with \(b,c\) bounded, then
\begin{align*}
  T^{1/2}
    &=A^*AT^{1/2}+B^*BT^{1/2} \\
    &=A^*R^{1/2}+B^*S^{1/2}
      =U_{A^*b+B^*c},
\end{align*}
and therefore \(\varphi(T)<+\infty\). Thus
\[
  \varphi(R+S)=\varphi(R)+\varphi(S)
\]
in all cases.

Now let \(T\) be an arbitrary element of
\(\mathcal U(\mathfrak A)\), and let us prove that
\(\varphi(TT^*)=\varphi(T^*T)\), which will establish that
\(\varphi\) is a trace on \(\mathcal U(\mathfrak A)^+\). Let
\[
  T=W|T|
\]
be the polar decomposition of \(T\), and suppose that
\(\varphi(T^*T)<+\infty\). Then \(|T|=U_a\), with \(a\) bounded, and
\[
  \varphi(T^*T)=\lVert a\rVert^2.
\]
Moreover, \(a\in\overline{|T|(\mathfrak H)}\), and hence
\(\lVert Wa\rVert=\lVert a\rVert\). On the other hand,
\[
  T^*=(WU_a)^*=U_{Wa}^*=U_{JWa},
  \qquad
  JWa\in\overline{T^*(\mathfrak H)},
\]
so that
\[
  \lVert WJWa\rVert
    =\lVert JWa\rVert
    =\lVert Wa\rVert
    =\lVert a\rVert.
\]
Since
\[
  |T^*|=WT^*=U_{WJWa},
\]
it follows that
\[
  \varphi(TT^*)
    =\lVert WJWa\rVert^2
    =\lVert a\rVert^2
    =\varphi(T^*T).
\]
The condition \(\varphi(T^*T)<+\infty\) therefore implies, and is
equivalent to, \(\varphi(TT^*)<+\infty\); and
\[
  \varphi(T^*T)=\varphi(TT^*)
  \qquad\text{for every }T\in\mathcal U(\mathfrak A).
\]
% Source PDF page 094 (printed p. 87)
\phantomsection\label{srcpage:87}
At the same time, if we put \(T=U_b=U_{Wa}\), we have proved that
\[
  \dot\varphi(U_b^*U_b)
    =\varphi(T^*T)
    =\lVert a\rVert^2
    =(b\mid b).
\]
By polarization, it follows that
\[
  \dot\varphi(U_b^*U_a)=(a\mid b)
\]
for arbitrary bounded elements \(a\) and \(b\).

The trace \(\varphi\) is faithful. Indeed, if \(\varphi(T)=0\) for
some \(T\in\mathcal U(\mathfrak A)^+\), then
\(T^{1/2}=U_a\) and \((a\mid a)=0\), whence
\(a=0\), \(T^{1/2}=0\), and \(T=0\).

Let us prove that \(\varphi\) is normal. Let \((T_\alpha)\) be an
increasing directed set in \(\mathcal U(\mathfrak A)^+\), with
supremum \(T\in\mathcal U(\mathfrak A)^+\). Everything reduces to
showing that
\[
  \varphi(T)\leq\sup_\alpha\varphi(T_\alpha).
\]
This is clear if \(\sup_\alpha\varphi(T_\alpha)=+\infty\). Suppose,
then, that
\[
  \varphi(T_\alpha)\leq\mu<+\infty.
\]
Thus
\[
  T_\alpha^{1/2}=U_{a_\alpha},
  \qquad
  (a_\alpha\mid a_\alpha)=\varphi(T_\alpha)\leq\mu.
\]
Since \(T_\alpha\) converges strongly to \(T\) while remaining
majorized by \(T\), \(T_\alpha^{1/2}\) converges strongly to
\(T^{1/2}\). [Indeed, if \(p(x)\) is a polynomial very close to
\(x^{1/2}\) on the interval \([0,\lVert T\rVert]\), then
\(p(T)\) and \(p(T_\alpha)\) are very close, in norm, to
\(T^{1/2}\) and \(T_\alpha^{1/2}\), respectively; and, on the other
hand, \(p(T_\alpha)\) converges strongly to \(p(T)\).] Hence, for
\(x,y\in\mathfrak A\),
\[
  (T^{1/2}y\mid x)
    =\lim_\alpha(T_\alpha^{1/2}y\mid x)
    =\lim_\alpha(U_{a_\alpha}y\mid x)
    =\lim_\alpha(a_\alpha\mid xy^*).
\]
Since the elements \(xy^*\) are dense in \(\mathfrak H\), and the
\(a_\alpha\) remain in a fixed ball, \(a_\alpha\) converges weakly to
an element \(a\in\mathfrak H\) satisfying
\[
  (T^{1/2}y\mid x)=(a\mid xy^*)=(V_ya\mid x).
\]
Thus \(a\) is bounded, \(T^{1/2}=U_a\), and
\[
  \varphi(T)
    =(a\mid a)
    \leq\sup_\alpha(a_\alpha\mid a_\alpha)
    =\sup_\alpha\varphi(T_\alpha).
\]

The trace \(\varphi\) is semifinite. Indeed, let
\(T\in\mathcal U(\mathfrak A)^+\), \(T\neq0\). For every bounded
\(a\), one has
\[
  U_aT^{1/2}=U_b
\]
with \(b\) bounded, and hence
\[
  \varphi\bigl(T^{1/2}U_a^*U_aT^{1/2}\bigr)<+\infty.
\]
Moreover,
\[
  T^{1/2}U_a^*U_aT^{1/2}\leq\lVert U_a\rVert^2T.
\]
On the other hand, one cannot have \(U_aT^{1/2}=0\) for every
bounded \(a\), since, on making \(U_a\) converge weakly to \(1\), one
would obtain \(T=0\). One may therefore find \(U_a\) such that
\[
  T^{1/2}U_a^*U_aT^{1/2}\neq0,
  \qquad
  T^{1/2}U_a^*U_aT^{1/2}\leq T,
\]
which proves the assertion.
\end{proof}


% Source PDF page 095 (printed p. 88)
\phantomsection\label{srcpage:88}
\begin{definition}\label{def:I-6-2}
Given a Hilbert algebra \(\mathfrak A\), the traces defined in
\cref{thm:I-6-1} are called the \emph{natural traces} on
\(\mathcal U(\mathfrak A)^+\) and \(\mathcal V(\mathfrak A)^+\).
\end{definition}

These traces do not change if \(\mathfrak A\) is replaced by the
Hilbert algebra of bounded elements.

Let \(\mathcal A\) be a von Neumann algebra, let \(\omega\) be a trace
on \(\mathcal A^+\), and let \(\mathfrak m\) be the ideal of definition
of \(\dot\omega\). If
\(S,T\in\mathfrak m^{1/2}\), then \(T^*S\in\mathfrak m\); put
\[
  (S\mid T)=\dot\omega(T^*S).
\]
It is clear that this defines on \(\mathfrak m^{1/2}\) a pre-Hilbert
space structure. We shall always use the notation \((S\mid T)\) in
the preceding sense (when no confusion concerning \(\omega\) is
possible), and shall denote by \(\lVert S\rVert_2\) the corresponding
seminorm \((S\mid S)^{1/2}\), to distinguish it from the usual norm
\(\lVert S\rVert\) of the operator \(S\).

\begin{theorem}\label{thm:I-6-2}
Let \(\mathcal A\) be a von Neumann algebra, let \(\omega\) be a
faithful, semifinite, normal trace on \(\mathcal A^+\), and let
\(\mathfrak m\) be the ideal of definition of \(\dot\omega\). Endowed
with the scalar product
\[
  (S\mid T)=\dot\omega(T^*S),
\]
\(\mathfrak m^{1/2}\) is a complete Hilbert algebra. Let
\(\mathfrak K\) be the Hilbert-space completion of
\(\mathfrak m^{1/2}\), and let \(J\) be the involution on
\(\mathfrak K\) canonically defined by \(\mathfrak m^{1/2}\). For
\(R\in\mathcal A\), the mapping
\[
  S\longmapsto RS
  \qquad
  (\text{respectively, }S\longmapsto SR)
\]
of \(\mathfrak m^{1/2}\) into itself extends continuously to an
operator \(\Phi(R)\) [respectively, \(\Psi(R)\)] in
\(\mathcal L(\mathfrak K)\). The mapping \(\Phi\) (respectively,
\(\Psi\)) is an isomorphism (respectively, an antiisomorphism) of
\(\mathcal A\) onto \(\mathcal U(\mathfrak m^{1/2})\)
[respectively, \(\mathcal V(\mathfrak m^{1/2})\)] which extends the
canonical mapping of the Hilbert algebra \(\mathfrak m^{1/2}\) into
\(\mathcal U(\mathfrak m^{1/2})\)
[respectively, into \(\mathcal V(\mathfrak m^{1/2})\)], and
\[
  \Psi(R)=J\Phi(R^*)J.
\]
Finally, let \(\varphi\) and \(\psi\) be the natural traces on
\(\mathcal U(\mathfrak m^{1/2})^+\) and
\(\mathcal V(\mathfrak m^{1/2})^+\), respectively. For
\(R\in\mathcal A^+\), one has
\[
  \omega(R)=\varphi(\Phi(R))=\psi(\Psi(R)).
\]
\end{theorem}

\begin{proof}
For \(R,S,T\in\mathfrak m^{1/2}\), one has, in view of
\hyperref[prop:I-1-11]{\S~1, Proposition~11},
\begin{align}
  (T^*\mid S^*)
    &=\dot\omega(ST^*)
      =\dot\omega(T^*S)
      =(S\mid T), \tag{1}\label{eq:I-6-thm2-1}\\
  (RS\mid T)
    &=\dot\omega(T^*RS)
      =\dot\omega((R^*T)^*S)
      =(S\mid R^*T), \tag{2}\label{eq:I-6-thm2-2}\\
  (RS\mid RS)
    &=\dot\omega(S^*R^*RS)
      \leq\lVert R^*R\rVert\,\dot\omega(S^*S)
      =\lVert R\rVert^2(S\mid S). \tag{3}\label{eq:I-6-thm2-3}
\end{align}
Moreover,
\[
  \bigl(S(1-T^*)\mid S(1-T^*)\bigr)
   =\dot\omega\bigl((1-T)S^*S(1-T^*)\bigr)
   =\dot\omega\bigl(S^*S(1-T^*)(1-T)\bigr).
\]
% Source PDF page 096 (printed p. 89)
\phantomsection\label{srcpage:89}
When \(T\) tends ultra-strongly to \(1\),
\((1-T^*)(1-T)\) tends ultra-weakly to zero; hence, by
\cref{prop:I-6-1},
\[
  (S-ST^*\mid S-ST^*)
\]
tends to zero, which proves axiom~\textup{(iv)} for Hilbert algebras.

For \(R\in\mathcal A\) and \(S\in\mathfrak m^{1/2}\), the inequality
\eqref{eq:I-6-thm2-3} remains valid and proves the existence of the
continuous extension \(\Phi(R)\). It is immediate that \(\Phi\)
(respectively, \(\Psi\)) is a homomorphism (respectively, an
antihomomorphism) of \(\mathcal A\) into \(\mathcal L(\mathfrak K)\)
which extends the canonical mapping of \(\mathfrak m^{1/2}\) into
\(\mathcal U(\mathfrak m^{1/2})\)
[respectively, into \(\mathcal V(\mathfrak m^{1/2})\)], and that
\[
  \Psi(R)=J\Phi(R^*)J.
\]
Moreover, \(\Phi\) and \(\Psi\) are injective. Indeed, if, for
example, \(\Phi(R)=0\), then \(RS=0\) for every
\(S\in\mathfrak m^{1/2}\), and hence \(R=0\), on making \(S\) tend
strongly to \(1\). It is clear that \(\Phi(\mathcal A)\) and
\(\Psi(\mathcal A)\) commute, and therefore
\[
  \Phi(\mathcal A)
    \subset\mathcal V(\mathfrak m^{1/2})'
    =\mathcal U(\mathfrak m^{1/2}),
  \qquad
  \Psi(\mathcal A)
    \subset\mathcal U(\mathfrak m^{1/2})'
    =\mathcal V(\mathfrak m^{1/2}).
\]
Since
\[
  \Phi(\mathcal A)\supset\Phi(\mathfrak m^{1/2}),
  \qquad
  \Psi(\mathcal A)\supset\Psi(\mathfrak m^{1/2}),
\]
the von Neumann algebras generated by \(\Phi(\mathcal A)\) and
\(\Psi(\mathcal A)\), respectively, are
\[
  \mathcal U(\mathfrak m^{1/2})
  \qquad\text{and}\qquad
  \mathcal V(\mathfrak m^{1/2}).
\]
Finally, as in the proof of
\hyperref[prop:I-4-1]{\S~4, Proposition~1}, one sees that \(\Phi\) and
\(\Psi\) are normal. Hence
\[
  \Phi(\mathcal A)=\mathcal U(\mathfrak m^{1/2}),
  \qquad
  \Psi(\mathcal A)=\mathcal V(\mathfrak m^{1/2})
\]
(\hyperref[cor:I-4-2]{\S~4, Theorem~2, Corollary~2}).

For \(S\in\mathfrak m^+\), one has
\[
  \omega(S)
    =(S^{1/2}\mid S^{1/2})
    =\varphi(\Phi(S))
    =\psi(\Psi(S)).
\]
If now \(S\) is arbitrary in \(\mathcal A^+\), let
\(\mathcal F\subset\mathfrak m^+\) be an increasing directed set with
supremum \(S\). Then
\[
  \varphi(\Phi(S))
    =\sup_{F\in\mathcal F}\varphi(\Phi(F))
    =\sup_{F\in\mathcal F}\omega(F)
    =\omega(S),
\]
and likewise \(\psi(\Psi(S))=\omega(S)\). Finally, if \(a\) is a
bounded element of \(\mathfrak K\), then, by \cref{thm:I-6-1},
\[
  \varphi(U_a^*U_a)<+\infty.
\]
Thus \(U_a=\Phi(S)\), with \(\omega(S^*S)<+\infty\), and hence
\(S\in\mathfrak m^{1/2}\). Therefore \(U_a=U_S\), so that
\(a=S\in\mathfrak m^{1/2}\): the Hilbert algebra
\(\mathfrak m^{1/2}\) is complete.
\end{proof}

\noindent\textit{Bibliography.}\quad
\ArticleRef{13}, \ArticleRef{19}, \ArticleRef{29}, \ArticleRef{30},
\ArticleRef{66}, \ArticleRef{99}, \ArticleRef{101}, \ArticleRef{116},
\ArticleRef{117}, \ArticleRef{120}.

\NumberedParagraph{3}{Trace Elements}
\label{no:I-6-3}

\begin{definition}\label{def:I-6-3}
Let \(\mathcal A\) be a von Neumann algebra on \(\mathfrak H\). An
element \(a\in\mathfrak H\) is called a \emph{trace element} for
\(\mathcal A\) if \(\omega_a\) is a trace on \(\mathcal A\), that is,
if
\[
  (T_1T_2a\mid a)=(T_2T_1a\mid a)
  \qquad(T_1,T_2\in\mathcal A).
\]
\end{definition}


% Source PDF page 097 (printed p. 90)
\phantomsection\label{srcpage:90}
\begin{proposition}\label{prop:I-6-3}
Let \(\mathfrak A\) be a Hilbert algebra, and let \(\mathfrak H\) be
its Hilbert-space completion. Every central element of
\(\mathfrak H\) is a trace element for \(\mathcal U(\mathfrak A)\)
and \(\mathcal V(\mathfrak A)\).
\end{proposition}

\begin{proof}
Let \(T_1,T_2\in\mathcal U(\mathfrak A)\), and put
\[
  T_1'=JT_1^*J\in\mathcal V(\mathfrak A),
\]
where \(J\) is the involution of \(\mathfrak H\) canonically defined
by \(\mathfrak A\). Since \(a\) is central,
\[
  T_1a=T_1'a,
  \qquad
  T_1^*a=T_1'^*a.
\]
Consequently,
\begin{align*}
  (T_2T_1a\mid a)
    &=(T_2T_1'a\mid a)
      =(T_1'T_2a\mid a) \\
    &=(T_2a\mid T_1'^*a)
      =(T_2a\mid T_1^*a)
      =(T_1T_2a\mid a),
\end{align*}
so \(a\) is a trace element for \(\mathcal U(\mathfrak A)\). The
proof for \(\mathcal V(\mathfrak A)\) is analogous.
\end{proof}

Let \(\mathcal A\) be a von Neumann algebra, and let \(a\) be a trace
element for \(\mathcal A\). Then
\(E_a^{\mathcal A'}\), which is the support of \(\omega_a\)
(\S~4, no.~6), is a projection of the centre \(\mathcal Z\) of
\(\mathcal A\). Since \(E_a^{\mathcal A}\) and
\(E_a^{\mathcal A'}\) have the same central support
(\hyperref[cor:I-1-7-2]{\S~1, Proposition~7, Corollary~2}),
\(E_a^{\mathcal A}\) has \(E_a^{\mathcal A'}\) as its central
support. In particular, if \(a\) is totalizing for \(\mathcal A\),
then \(a\) is separating for \(\mathcal A\). Put, for
\(T\in\mathcal A\),
\[
  \Phi(T)=Ta.
\]
Then \(\Phi\) is a bijective mapping of \(\mathcal A\) onto
\(\mathfrak A=\Phi(\mathcal A)\), and one has the following
proposition.

\begin{proposition}\label{prop:I-6-4}
\begin{enumerate}[label=\textup{(\roman*)},leftmargin=*]
\item If the involutive-algebra structure of \(\mathcal A\) is
transported to \(\mathfrak A\) by \(\Phi\), then \(\mathfrak A\)
becomes a complete Hilbert algebra.
\item One has
\[
  \mathcal A=\mathcal U(\mathfrak A),
  \qquad
  \mathcal A'=\mathcal V(\mathfrak A).
\]
\item The element \(a\) is an identity element for \(\mathfrak A\)
(and hence is central).
\end{enumerate}
\end{proposition}

\begin{proof}
Transport to \(\mathcal A\), by \(\Phi^{-1}\), the scalar product of
\(\mathfrak A\subset\mathfrak H\). For
\(T_1,T_2\in\mathcal A\), one has
\[
  (T_1\mid T_2)
    =(T_1a\mid T_2a)
    =\omega_a(T_2^*T_1).
\]
Then \(\mathcal A\) becomes a complete Hilbert algebra by
\cref{thm:I-6-2}, which proves~\textup{(i)}. Likewise, on transporting
the results of \cref{thm:I-6-2} by \(\Phi\), one sees that
\(\mathcal A=\mathcal U(\mathfrak A)\). Hence
\(\mathcal A'=\mathcal V(\mathfrak A)\). Finally, \(a\) is the image
under \(\Phi\) of the identity element of \(\mathcal A\), which
proves~\textup{(iii)}.
\end{proof}

\setcounter{corollary}{0}
\begin{corollary}\label{cor:I-6-4}
If \(a\) is a trace element and is totalizing for \(\mathcal A\),
then \(a\) is a trace element and is totalizing for
\(\mathcal A'\).
\end{corollary}

\begin{proof}
This follows at once from \cref{prop:I-6-3,prop:I-6-4}.
\end{proof}

\begin{corollary}\label{cor:I-6-5}
Let \(\mathcal A\) be an abelian von Neumann algebra. If there exists
a totalizing element for \(\mathcal A\), then
\(\mathcal A'=\mathcal A\).
\end{corollary}

\begin{proof}
Every element of \(\mathfrak H\) is a trace element for
\(\mathcal A\). By \cref{prop:I-6-4}, \(\mathcal A'\) is
antiisomorphic to \(\mathcal A\), and is therefore abelian. Thus
\[
  \mathcal A\subset\mathcal A'\subset\mathcal A''=\mathcal A.
\]
\end{proof}

For a direct proof of \cref{cor:I-6-5}, see Exercise~5.

\noindent\textit{Bibliography.}\quad
\ArticleRef{29}, \ArticleRef{66}, \ArticleRef{89}, \ArticleRef{100}.

% Source PDF page 098 (printed p. 91)
\phantomsection\label{srcpage:91}
\NumberedParagraph{4}{Order Relation in the Set of Traces}
\label{no:I-6-4}
The present subsection will clarify the uniqueness questions relating
to traces. Subsections~6 and~7 will be devoted to existence questions.

\begin{definition}\label{def:I-6-4}
Let \(\mathcal A\) be a von Neumann algebra, and let \(\varphi\) and
\(\varphi'\) be two traces on \(\mathcal A^+\). One says that
\(\varphi\) \emph{majorizes} \(\varphi'\), and writes
\(\varphi\geq\varphi'\), if
\[
  \varphi(T)\geq\varphi'(T)
  \qquad\text{for every }T\in\mathcal A^+.
\]
\end{definition}

\begin{theorem}\label{thm:I-6-3}
Let \(\mathcal Z\) be the centre of \(\mathcal A\), and let
\(\varphi\) be a semifinite normal trace on \(\mathcal A^+\). For
every \(S\in\mathcal Z\) such that \(0\leq S\leq1\), the function
\[
  T\longmapsto\varphi(ST)
\]
on \(\mathcal A^+\) is a normal trace \(\varphi_S\) majorized by
\(\varphi\), and every normal trace majorized by \(\varphi\) is of
this form. If \(\varphi\) is faithful, the mapping
\(S\longmapsto\varphi_S\) is injective.
\end{theorem}

\begin{proof}
For every \(S\in\mathcal Z\) such that \(0\leq S\leq1\), it is clear
that the function \(T\longmapsto\varphi(ST)\) on \(\mathcal A^+\) is
a normal trace majorized by \(\varphi\). Conversely, let
\(\varphi'\) be a normal trace majorized by \(\varphi\), and suppose
first that \(\varphi\) is faithful. Let \(\mathfrak m\) and
\(\mathfrak m'\) be the ideals of definition of \(\dot\varphi\) and
\(\dot\varphi'\), respectively; one has
\(\mathfrak m\subset\mathfrak m'\). By \cref{thm:I-6-2},
\(\dot\varphi\) canonically defines a Hilbert-algebra structure on
\(\mathfrak m^{1/2}\), two von Neumann algebras
\(\mathcal U(\mathfrak m^{1/2})\) and
\(\mathcal V(\mathfrak m^{1/2})\) on the Hilbert space
\(\mathfrak K\) which is the completion of
\(\mathfrak m^{1/2}\), an isomorphism \(\Phi\) of \(\mathcal A\)
onto \(\mathcal U(\mathfrak m^{1/2})\), and an antiisomorphism
\(\Psi\) of \(\mathcal A\) onto
\(\mathcal V(\mathfrak m^{1/2})\). For
\(T,T_1\in\mathfrak m^{1/2}\), put
\[
  ((T\mid T_1))=\dot\varphi'(T_1^*T).
\]
Since \(\varphi'\leq\varphi\), there is a unique Hermitian operator
\(A\in\mathcal L(\mathfrak K)\) such that \(0\leq A\leq1\) and
\[
  ((T\mid T_1))=(T\mid A(T_1)).
\]
For \(T,T_1\in\mathfrak m^{1/2}\) and \(R\in\mathcal A\), one has
\begin{align}
  (T\mid\Phi(R)A(T_1))
    &=(R^*T\mid A(T_1))
      =((R^*T\mid T_1))
      =\dot\varphi'(T_1^*R^*T),
      \tag{4}\label{eq:I-6-thm3-4}\\
  (T\mid A\Phi(R)(T_1))
    &=(T\mid A(RT_1))
      =((T\mid RT_1))
      =\dot\varphi'(T_1^*R^*T),
      \tag{5}\label{eq:I-6-thm3-5}\\
  (T\mid\Psi(R)A(T_1))
    &=(TR^*\mid A(T_1))
      =((TR^*\mid T_1))
      =\dot\varphi'(T_1^*TR^*),
      \tag{6}\label{eq:I-6-thm3-6}\\
  (T\mid A\Psi(R)(T_1))
    &=(T\mid A(T_1R))
      =((T\mid T_1R))
      =\dot\varphi'(R^*T_1^*T)
      =\dot\varphi'(T_1^*TR^*).
      \tag{7}\label{eq:I-6-thm3-7}
\end{align}
Thus
\[
  A\in\mathcal U(\mathfrak m^{1/2})'
       \cap\mathcal V(\mathfrak m^{1/2})',
\]
and hence \(A=\Psi(S)\), where \(S\in\mathcal Z\) and
\(0\leq S\leq1\). Consequently, for
\(T,T_1\in\mathfrak m^{1/2}\),
\[
  \dot\varphi'(T_1^*T)
    =(T\mid T_1S)
    =\dot\varphi(ST_1^*T).
\]
Thus \(\varphi'(T)=\varphi(ST)\) for \(T\in\mathfrak m^+\). Now
\(T\longmapsto\varphi'(T)\) and \(T\longmapsto\varphi(ST)\) are two
normal traces on \(\mathcal A^+\) which agree on
\(\mathfrak m^+\), and therefore are identical
(\hyperref[cor:I-3-5]{\S~3, Theorem~2, Corollary~5}). The uniqueness
of \(S\) follows from the preceding argument and from the fact that
\(\Psi\) is injective.

Finally, if \(\varphi\) is no longer assumed faithful,
\(\mathcal A\) is identified with the product of two von Neumann
algebras \(\mathcal A_1,\mathcal A_2\), on whose positive parts
\(\varphi\) induces normal traces
% Source PDF page 099 (printed p. 92)
\phantomsection\label{srcpage:92}
\(\varphi_1,\varphi_2\) with the following properties:
\(1^\circ\) \(\varphi_1\) is faithful and semifinite;
\(2^\circ\) \(\varphi_2=0\). Let \(\varphi_1',\varphi_2'\) be the
traces induced by \(\varphi'\) on
\(\mathcal A_1^+,\mathcal A_2^+\). One has
\[
  \varphi_1'(T)=\varphi_1(S_1T),
\]
where \(S_1\) is an element of the centre of \(\mathcal A_1\) such
that \(0\leq S_1\leq1\). Let \(S\) be the element of
\(\mathcal Z\) defined by the element \(S_1\) of \(\mathcal A_1\)
and the zero element of \(\mathcal A_2\). Evidently,
\[
  \varphi'(T)=\varphi(ST)
  \qquad(T\in\mathcal A^+),
  \qquad
  0\leq S\leq1.
\]
\end{proof}

\begin{corollary*}
On a factor, two faithful, semifinite, normal traces are proportional.
\end{corollary*}

\begin{proof}
Let \(\varphi\) and \(\varphi'\) be these two traces. The trace
\[
  \varphi''=\varphi+\varphi'
\]
is normal, faithful, and semifinite, and majorizes \(\varphi\) and
\(\varphi'\). Hence
\[
  \varphi=\lambda\varphi'',
  \qquad
  \varphi'=\lambda'\varphi'',
\]
where \(\lambda>0\) and \(\lambda'>0\) are constants. Consequently,
\[
  \varphi=\lambda\lambda'^{-1}\varphi'.
\]
\end{proof}

Let \(\mathcal A\) be a von Neumann algebra, let \(\varphi\) be a
trace on \(\mathcal A^+\), and let \(\varphi'\) be a positive linear
functional on \(\mathcal A\). We say that \(\varphi\) majorizes
\(\varphi'\) if
\[
  \varphi(T)\geq\varphi'(T)
  \qquad(T\in\mathcal A^+).
\]
The method used in the proof of \cref{thm:I-6-3} also gives the
following result.

\begin{proposition}\label{prop:I-6-5}
Let \(\varphi\) be a semifinite normal trace on \(\mathcal A^+\), and
let \(\mathfrak m\) be the ideal of definition of \(\dot\varphi\).
For every \(S\in\mathfrak m\) such that \(0\leq S\leq1\), the
function
\[
  T\longmapsto
    \dot\varphi(S^{1/2}TS^{1/2})
    =\dot\varphi(ST)
\]
on \(\mathcal A\) is a normal positive linear functional
\(\varphi_S\) majorized by \(\varphi\), and every normal positive
linear functional majorized by \(\varphi\) is of this form. If
\(\varphi\) is faithful, the mapping
\(S\longmapsto\varphi_S\) is injective.
\end{proposition}

\begin{proof}
Let \(S\in\mathfrak m\), with \(0\leq S\leq1\). It is clear that
\(\varphi_S\) is a normal positive linear functional. Moreover, for
\(T\in\mathcal A^+\),
\[
  \dot\varphi(ST)
    =\dot\varphi(T^{1/2}ST^{1/2})
    \leq\dot\varphi(T^{1/2}T^{1/2})
    =\varphi(T),
\]
so \(\varphi_S\) is majorized by \(\varphi\). Conversely, let
\(\varphi'\) be a positive linear functional majorized by
\(\varphi\). As in the proof of \cref{thm:I-6-3}, introduce
\(\mathfrak K\) and \(A\). Equations
\eqref{eq:I-6-thm3-4} and \eqref{eq:I-6-thm3-5}, with
\(\dot\varphi'\) replaced by \(\varphi'\), remain valid. Hence
\[
  A\in\mathcal U(\mathfrak m^{1/2})',
\]
so \(A=\Psi(S)\) for some \(S\in\mathcal A\), with
\(0\leq S\leq1\). The proof then ends exactly as that of
\cref{thm:I-6-3}.
\end{proof}

\noindent\textit{Bibliography.}\quad
\ArticleRef{19}, \ArticleRef{65}, \ArticleRef{66}, \ArticleRef{92},
\ArticleRef{101}.

\NumberedParagraph{5}{Application: Isomorphisms of Standard von
Neumann Algebras}
\label{no:I-6-5}

\begin{lemma}\label{lem:I-6-1}
Let \(\mathcal A\) be a standard von Neumann algebra on
\(\mathfrak H\), and let \(\varphi\) be a faithful, semifinite,
normal trace on \(\mathcal A^+\). There exists a Hilbert algebra
\(\mathfrak A\), dense in \(\mathfrak H\), such that
\begin{enumerate}[label=\arabic*\textsuperscript{\(\circ\)}.,leftmargin=*]
\item \(\mathcal A=\mathcal U(\mathfrak A)\);
\item \(\varphi\) is the corresponding natural trace on
\(\mathcal A^+\).
\end{enumerate}
\end{lemma}

% Source PDF page 100 (printed p. 93)
\phantomsection\label{srcpage:93}
\begin{proof}
There exists a complete Hilbert algebra \(\mathfrak B\), dense in
\(\mathfrak H\), such that
\[
  \mathcal A=\mathcal U(\mathfrak B).
\]
Let \(\psi\) be the corresponding natural trace on
\(\mathcal A^+\), and put
\[
  \chi=\varphi+\psi.
\]
The traces \(\varphi,\psi,\chi\) are normal, faithful, and semifinite,
and \(\psi\leq\chi\), \(\varphi\leq\chi\). Hence there are operators
\(S,S'\) in the centre \(\mathcal Z\) of \(\mathcal A\) such that
\[
  0\leq S\leq1,
  \qquad
  0\leq S'\leq1,
  \qquad
  \psi(T)=\chi(ST),
  \qquad
  \varphi(T)=\chi(S'T)
\]
for \(T\in\mathcal A^+\), by \cref{thm:I-6-3}. The supports of
\(S\) and \(S'\) are \(1\), since \(\varphi\) and \(\psi\) are
faithful. Using the Gelfand isomorphism of an algebra of continuous
functions on \(\mathcal Z\), one sees that there are projections
\(E_1,E_2,\ldots\) of \(\mathcal Z\), pairwise orthogonal and with
sum \(1\), and numbers \(\lambda_i>0\), such that
\[
  SE_i\geq\lambda_iE_i,
  \qquad
  S'E_i\geq\lambda_iE_i.
\]
It is enough to prove the lemma for each algebra
\(\mathcal A_{E_i}\). We may therefore suppose henceforth that
\[
  S\geq\lambda>0,
  \qquad
  S'\geq\lambda>0.
\]
Put
\[
  Z=S'^{-1/2}S^{1/2}\in\mathcal Z.
\]
Then
\[
  Z\geq\mu>0,
  \qquad
  \varphi(T)=\psi(Z^{-2}T)
  \quad(T\in\mathcal A^+).
\]

We now define a Hilbert algebra \(\mathfrak A\). Its underlying
pre-Hilbert space is the same as that of \(\mathfrak B\), its
involution is that of \(\mathfrak B\), and its multiplication is
defined by
\begin{align*}
  (x,y)\longmapsto Z(xy)
    &=ZU_xy=U_xZy=x(Zy) \\
    &=ZV_yx=V_yZx=(Zx)y.
\end{align*}
It is immediate that this defines a Hilbert-algebra structure (in
particular, by \hyperref[prop:I-5-2]{\S~5, the corollary to
Proposition~2}). Let
\[
  x\longmapsto\widetilde U_x,
  \qquad
  x\longmapsto\widetilde V_x
\]
be the canonical mappings of \(\mathfrak A\) into
\(\mathcal U(\mathfrak A)\) and \(\mathcal V(\mathfrak A)\). One has
\[
  \widetilde U_x=ZU_x,
  \qquad
  \widetilde V_x=ZV_x.
\]
This shows that
\[
  \mathcal U(\mathfrak A)\subset\mathcal U(\mathfrak B),
  \qquad
  \mathcal V(\mathfrak A)\subset\mathcal V(\mathfrak B),
\]
and therefore
\[
  \mathcal U(\mathfrak A)=\mathcal U(\mathfrak B)=\mathcal A.
\]
On the other hand, an element of \(\mathfrak H\) which is bounded
relative to \(\mathfrak A\) is bounded relative to
\(\mathfrak B\), and therefore belongs to
\(\mathfrak B=\mathfrak A\); thus \(\mathfrak A\) is complete. Let
\(\varphi'\) be the natural trace on
\(\mathcal U(\mathfrak A)^+\). If
\(T\in\mathcal U(\mathfrak A)^+\), then
\(\varphi'(T)=+\infty\) if and only if \(\psi(T)=+\infty\), that is,
if and only if \(\varphi(T)=+\infty\). On the other hand, if
\(T^{1/2}=\widetilde U_x\), with \(x\in\mathfrak A\), then
\[
  Z^{-1}T^{1/2}=Z^{-1}\widetilde U_x=U_x,
\]
and hence
\[
  \varphi(T)=\psi(Z^{-2}T)=\lVert x\rVert^2=\varphi'(T).
\]
Therefore \(\varphi=\varphi'\).
\end{proof}

\begin{theorem}\label{thm:I-6-4}
Let \(\mathcal A\) and \(\mathcal A_1\) be standard von Neumann
algebras. Every isomorphism of \(\mathcal A\) onto
\(\mathcal A_1\) is spatial.
\end{theorem}


% Source PDF page 101 (printed p. 94)
\phantomsection\label{srcpage:94}
\begin{proof}
Let \(\Phi\) be an isomorphism of \(\mathcal A\) onto
\(\mathcal A_1\). Let \(\mathfrak A\) be a complete Hilbert algebra
such that \(\mathcal A=\mathcal U(\mathfrak A)\), and let
\(\varphi\) be the corresponding natural trace on \(\mathcal A^+\).
Transporting \(\varphi\) to \(\mathcal A_1^+\) by \(\Phi\), one
obtains a faithful, semifinite, normal trace \(\varphi_1\) on
\(\mathcal A_1^+\). By \cref{lem:I-6-1}, there exists a complete
Hilbert algebra \(\mathfrak A_1\) such that
\(\mathcal A_1=\mathcal U(\mathfrak A_1)\) and such that
\(\varphi_1\) is the corresponding natural trace. Let
\[
  \Omega:x\longmapsto U_x,
  \qquad
  \Omega_1:x\longmapsto U_x'
\]
be the canonical mappings of \(\mathfrak A\) into \(\mathcal A\) and
of \(\mathfrak A_1\) into \(\mathcal A_1\). Then
\[
  \Omega_1^{-1}\circ\Phi\circ\Omega
\]
is an isomorphism of the Hilbert algebra \(\mathfrak A\) onto the
Hilbert algebra \(\mathfrak A_1\), and extends to an isomorphism
\(W\) of the Hilbert space \(\mathfrak H\), the completion of
\(\mathfrak A\), onto the Hilbert space \(\mathfrak H_1\), the
completion of \(\mathfrak A_1\). For \(x,y\in\mathfrak A\), one has
\begin{align*}
  W^{-1}\Phi(U_x)Wy
    &=W^{-1}U_{Wx}'Wy
      =W^{-1}\bigl((Wx)(Wy)\bigr) \\
    &=W^{-1}W(xy)=xy=U_xy.
\end{align*}
Thus
\[
  \Phi(U_x)=WU_xW^{-1}.
\]
Consequently, the isomorphisms \(\Phi\) and
\(T\longmapsto WTW^{-1}\) agree on a two-sided ideal of
\(\mathcal A\) which is ultra-strongly dense in \(\mathcal A\).
Since these isomorphisms are ultra-strongly continuous
(\hyperref[cor:I-4-1]{\S~4, Theorem~2, Corollary~1}), they are
identical.
\end{proof}

\noindent\textit{Bibliography.}\quad \ArticleRef{101}.

\NumberedParagraph{6}{Normal Traces on \(\mathcal L(\mathfrak H)\)}
\label{no:I-6-6}

\begin{theorem}\label{thm:I-6-5}
Let \(\mathfrak H\) be a complex Hilbert space, and let
\((e_i)_{i\in I}\) be an orthonormal basis of \(\mathfrak H\). For
\(T\in\mathcal L(\mathfrak H)^+\), put
\[
  \varphi(T)=\sum_{i\in I}(Te_i\mid e_i).
\]
Then \(\varphi\) is a faithful, semifinite, normal trace on
\(\mathcal L(\mathfrak H)^+\), independent of the choice of the
orthonormal basis \((e_i)_{i\in I}\). If \(E\) is a projection, then
\(\varphi(E)\) is the Hilbert-space dimension of
\(E(\mathfrak H)\).

In the last assertion, all infinite cardinals are understood to be
identified with \(+\infty\).
\end{theorem}

\begin{proof}
It is clear that \(\varphi\) has properties~\textup{(i)} and
\textup{(ii)} of \cref{def:I-6-1}. Now let
\(T\in\mathcal L(\mathfrak H)\), and let
\((e_\kappa')_{\kappa\in K}\) be any orthonormal basis of
\(\mathfrak H\). Then
\begin{align*}
  \sum_{i\in I}(T^*Te_i\mid e_i)
    &=\sum_{i\in I}\lVert Te_i\rVert^2
      =\sum_{i\in I}\sum_{\kappa\in K}
        |(Te_i\mid e_\kappa')|^2 \\
    &=\sum_{\kappa\in K}\sum_{i\in I}
        |(Te_i\mid e_\kappa')|^2,
\end{align*}
whereas
\begin{align*}
  \sum_{\kappa\in K}(TT^*e_\kappa'\mid e_\kappa')
    &=\sum_{\kappa\in K}\lVert T^*e_\kappa'\rVert^2
      =\sum_{\kappa\in K}\sum_{i\in I}
        |(T^*e_\kappa'\mid e_i)|^2 \\
    &=\sum_{\kappa\in K}\sum_{i\in I}
        |(Te_i\mid e_\kappa')|^2.
\end{align*}
This proves both that \(\varphi(T^*T)=\varphi(TT^*)\), and hence that
\(\varphi\) is a trace, and that this trace is independent of the
chosen orthonormal basis.
% Source PDF page 102 (printed p. 95)
\phantomsection\label{srcpage:95}
Since each positive linear functional \(\omega_{e_i}\) is normal,
\(\varphi\) is normal. If \(T\in\mathcal L(\mathfrak H)^+\) is such
that
\[
  0=\varphi(T)=\sum_{i\in I}\lVert T^{1/2}e_i\rVert^2,
\]
then \(T^{1/2}=0\), and hence \(T=0\): the trace \(\varphi\) is
faithful. Now let \(E\) be a projection. The \(e_i\) may be chosen so
that \(Ee_i=0\) or \(Ee_i=e_i\) for every \(i\). Then
\(\varphi(E)\) is the number of the \(e_i\) for which \(Ee_i=e_i\),
that is, the Hilbert-space dimension of \(E(\mathfrak H)\), with the
stated convention concerning infinite cardinals. This proves, in
particular, that \(\varphi(T)<+\infty\) for some nonzero
\(T\in\mathcal L(\mathfrak H)^+\). Since
\(\mathcal L(\mathfrak H)\) is a factor, \(\varphi\) is therefore
semifinite by the second corollary to \cref{prop:I-6-1}.
\end{proof}

By the corollary to \cref{thm:I-6-3}, every normal trace on
\(\mathcal L(\mathfrak H)^+\) is proportional to the preceding trace
\(\varphi\), or is identically infinite on the nonzero operators of
\(\mathcal L(\mathfrak H)^+\).

\begin{corollary*}
The set \(\mathfrak n\) of those \(T\in\mathcal L(\mathfrak H)\) for
which
\[
  \sum_{i\in I}\lVert Te_i\rVert^2
    =\sum_{i,\kappa\in I}|(Te_i\mid e_\kappa)|^2
    <+\infty
\]
is a two-sided ideal of \(\mathcal L(\mathfrak H)\), independent of
the orthonormal basis \((e_i)_{i\in I}\), and consists of compact
operators. Every matrix \((t_{i\kappa})\) such that
\[
  \sum_{i,\kappa\in I}|t_{i\kappa}|^2<+\infty
\]
represents, relative to \((e_i)_{i\in I}\), an operator in
\(\mathfrak n\). If, for \(T,T'\in\mathfrak n\), one puts
\[
  (T\mid T')
    =\sum_{i\in I}(Te_i\mid T'e_i)
    =\sum_{i,\kappa\in I}t_{i\kappa}\overline{t'_{i\kappa}},
\]
then \(\mathfrak n\) becomes a Hilbert space, independently of the
orthonormal basis \((e_i)\). The two-sided ideal \(\mathfrak f\) of
finite-rank operators is dense in \(\mathfrak n\) for this Hilbert
space structure.
\end{corollary*}

\begin{proof}
By \cref{thm:I-6-5}, \(\mathfrak n\) is a two-sided ideal of
\(\mathcal L(\mathfrak H)\), independent of the basis
\((e_i)_{i\in I}\). The inequality
\[
  \lVert T\rVert^2\leq\sum_{i\in I}\lVert Te_i\rVert^2,
\]
which follows at once from the fact that
\(\sum_{i\in I}\lVert Te_i\rVert^2\) is independent of
\((e_i)\), proves that every element of \(\mathfrak n\) is a limit,
in the norm topology, of operators represented by matrices having
only finitely many nonzero entries, and therefore, a fortiori, of
finite-rank operators. Thus every \(T\in\mathfrak n\) is compact.
Now consider a matrix \((t_{i\kappa})\) such that
\[
  \sum_{i,\kappa\in I}|t_{i\kappa}|^2<+\infty.
\]
% Source PDF page 103 (printed p. 96)
\phantomsection\label{srcpage:96}
By putting
\[
  T\left(\sum_{\kappa\in I}\lambda_\kappa e_\kappa\right)
    =\sum_{i\in I}
       \left(\sum_{\kappa\in I}t_{i\kappa}\lambda_\kappa\right)e_i,
\]
one defines an operator \(T\in\mathcal L(\mathfrak H)\), by virtue of
the inequality
\[
  \sum_{i\in I}
    \left|\sum_{\kappa\in I}t_{i\kappa}\lambda_\kappa\right|^2
  \leq
  \left(\sum_{i,\kappa\in I}|t_{i\kappa}|^2\right)
  \left(\sum_{\kappa\in I}|\lambda_\kappa|^2\right),
\]
and one has \(T\in\mathfrak n\). The pre-Hilbert structure of
\(\mathfrak n\) is obtained either directly or by applying
\cref{thm:I-6-2}. It is now immediate that \(\mathfrak n\) is
complete and that \(\mathfrak f\) is dense in \(\mathfrak n\).
\end{proof}

The operators in \(\mathfrak n\) are called
\emph{Hilbert-Schmidt operators}.

Let \(\mathfrak H'\) be the conjugate Hilbert space of
\(\mathfrak H\), that is, we recall, the space \(\mathfrak H\)
endowed with the operations
\[
  (\lambda,x)\longmapsto\overline\lambda x,
  \qquad
  (x,y)\longmapsto x+y,
\]
and with the scalar product
\[
  (x,y)\longmapsto(y\mid x).
\]
For the bilinear form \((x,y)\longmapsto(x\mid y)\) on
\(\mathfrak H\times\mathfrak H'\), \(\mathfrak H'\) is identified
with the dual of \(\mathfrak H\). We shall not identify
\(\mathfrak H\) with its dual in the remainder of this subsection.
Every element \(T\) of \(\mathfrak f\) is of the form
\[
  y\longmapsto\sum_{j=1}^{n}(y\mid y_j)x_j,
\]
and the mapping
\[
  T\longmapsto\sum_{j=1}^{n}y_j\otimes x_j
\]
identifies the vector space \(\mathfrak f\) with the algebraic tensor
product of \(\mathfrak H'\) and \(\mathfrak H\). The \(y_j\) may,
moreover, be assumed orthonormal and the \(x_j\) orthogonal, in which
case
\[
  (T\mid T)
    =\sum_{j=1}^{n}\lVert x_j\rVert^2
    =\left\lVert\sum_{j=1}^{n}y_j\otimes x_j\right\rVert^2.
\]
This proves that the preceding mapping is isometric for the
pre-Hilbert structures of the vector spaces in question. It therefore
extends to a canonical isomorphism of the Hilbert space
\(\mathfrak n\) onto the Hilbert space
\(\mathfrak H'\otimes\mathfrak H\). Let \(S\in\mathfrak n\). Since
\(S\) is compact, one has
\[
  Sy=\sum_{j=1}^{\infty}(y\mid u_j)v_j,
\]
where the \(u_j\) are orthonormal and the \(v_j\) are orthogonal.
Then
\[
  (S\mid S)=\sum_{j=1}^{\infty}\lVert v_j\rVert^2.
\]
Thus \(S\) is the limit, in the Hilbert-space structure of
\(\mathfrak n\), of the operator
\[
  y\longmapsto\sum_{j=1}^{n}(y\mid u_j)v_j,
\]
% Source PDF page 104 (printed p. 97)
\phantomsection\label{srcpage:97}
which was identified with \(\sum_{j=1}^{n}u_j\otimes v_j\). Finally,
\(S\) is identified with the element
\[
  \sum_{j=1}^{\infty}u_j\otimes v_j
\]
of \(\mathfrak H'\otimes\mathfrak H\).

\begin{proposition}\label{prop:I-6-6}
\begin{enumerate}[label=\textup{(\roman*)},leftmargin=*]
\item The mapping \(T\longmapsto T^*\) of \(\mathfrak n\) onto
\(\mathfrak n\) is identified with the isometric linear mapping of
\(\mathfrak H'\otimes\mathfrak H\) onto itself which sends
\(u\otimes v\) to \(v\otimes u\).
\item If \(S\in\mathcal L(\mathfrak H)\), the mapping
\(T\longmapsto ST\) of \(\mathfrak n\) into itself is identified with
\(I_{\mathfrak H'}\otimes S\), and the mapping
\(T\longmapsto TS\) of \(\mathfrak n\) into itself is identified with
\(S^*\otimes I_{\mathfrak H}\).
\end{enumerate}
\end{proposition}

\begin{proof}
If \(T\) is the operator
\[
  y\longmapsto(y\mid u)v,
\]
then
\[
  (T^*y\mid y')
    =(y\mid Ty')
    =\overline{(y'\mid u)}(y\mid v)
    =((y\mid v)u\mid y').
\]
Thus \(T^*\) is the operator
\[
  y\longmapsto(y\mid v)u,
\]
which proves~\textup{(i)}. On the other hand, for
\(S\in\mathcal L(\mathfrak H)\), one has
\[
  STy=(y\mid u)Sv,
\]
so \(ST\) is identified with
\[
  u\otimes Sv
    =(I_{\mathfrak H'}\otimes S)(u\otimes v).
\]
This proves the first assertion of~\textup{(ii)}. The second follows
from the first and from~\textup{(i)}.
\end{proof}

\begin{corollary*}
If \(\mathfrak H\) is a Hilbert space, the von Neumann algebras
\[
  \mathcal L(\mathfrak H')\otimes\mathbb C_{\mathfrak H}
  \qquad\text{and}\qquad
  \mathbb C_{\mathfrak H'}\otimes\mathcal L(\mathfrak H)
\]
on \(\mathfrak H'\otimes\mathfrak H\) are standard von Neumann
algebras.
\end{corollary*}

\begin{proof}
This follows from \cref{thm:I-6-2,prop:I-6-6} and from the fact that
the Hilbert spaces \(\mathfrak H\) and \(\mathfrak H'\), having the
same Hilbert-space dimension, are isomorphic.
\end{proof}

\noindent\textit{Bibliography.}\quad \TreatiseRef{13}.

\NumberedParagraph{7}{First Classification of von Neumann Algebras}
\label{no:I-6-7}

\begin{definition}\label{def:I-6-5}
A von Neumann algebra \(\mathcal A\) is said to be \emph{finite}
(respectively, \emph{semifinite}) if, for every nonzero
\(T\in\mathcal A^+\), there exists a finite (respectively,
semifinite) normal trace \(\varphi\) on \(\mathcal A^+\) such that
\(\varphi(T)\neq0\). It is said to be \emph{properly infinite}
(respectively, \emph{purely infinite}) if there is no finite
(respectively, semifinite) normal trace on \(\mathcal A^+\) other
than zero.
\end{definition}

A von Neumann algebra which is not finite is also said to be
\emph{infinite}. This classification is invariant under isomorphisms
and antiisomorphisms.

% Source PDF page 105 (printed p. 98)
\phantomsection\label{srcpage:98}
The logical relations among these various notions are represented by
the following diagram (where \(A\mathrel{\not\leftrightarrow}B\)
means that properties \(A\) and \(B\) are incompatible, except when
\(\mathfrak H=0\)):
\[
\begin{array}{r@{\qquad}c@{\qquad}l}
  \text{Purely infinite}
    &\not\longleftrightarrow&
  \text{Semifinite}\\[-0.1em]
  \Downarrow\quad\not\searrow
    &&\Uparrow\\[-0.1em]
  \text{Properly infinite}
    &\not\longleftrightarrow&
  \text{Finite}.
\end{array}
\]

We shall see in \S~9 that purely infinite von Neumann algebras exist.
By \cref{thm:I-6-5}, \(\mathcal L(\mathfrak H)\) is semifinite, and
is properly infinite (respectively, finite) if the Hilbert-space
dimension of \(\mathfrak H\) is infinite (respectively, finite). An
abelian von Neumann algebra \(\mathcal A\) is finite; indeed, for every
\(x\in\mathfrak H\), \(\omega_x\) is a finite normal trace on
\(\mathcal A^+\).

Every von Neumann subalgebra of a finite von Neumann algebra is
finite. On the other hand, nothing can be said about the finiteness of
the von Neumann subalgebras of a semifinite von Neumann algebra, since
\(\mathcal L(\mathfrak H)\) is semifinite.

\begin{proposition}\label{prop:I-6-7}
Let \((\mathcal A_i)_{i\in I}\) be a family of von Neumann algebras,
and let \(\mathcal A\) be their product von Neumann algebra. For
\(\mathcal A\) to be finite (respectively, semifinite, properly
infinite, purely infinite), it is necessary and sufficient that each
\(\mathcal A_i\) be finite (respectively, semifinite, properly
infinite, purely infinite).
\end{proposition}

\begin{proof}
Every finite (respectively, semifinite) normal trace \(\varphi\) on
\(\mathcal A^+\) induces on each \(\mathcal A_i^+\) a finite
(respectively, semifinite) normal trace \(\varphi_i\); and if
\(\varphi\) is nonzero, at least one of the \(\varphi_i\) is nonzero.
Conversely, let \(\alpha\in I\), and let \(\varphi_\alpha\) be a
nonzero finite (respectively, semifinite) normal trace on
\(\mathcal A_\alpha^+\). For
\(T=(T_i)\in\mathcal A^+\), put
\[
  \varphi(T)=\varphi_\alpha(T_\alpha).
\]
Then \(\varphi\) is a finite (respectively, semifinite) normal trace
on \(\mathcal A^+\). The proposition follows at once from these
remarks.
\end{proof}

\begin{proposition}\label{prop:I-6-8}
Let \(\mathcal A\) be a von Neumann algebra and \(\mathcal Z\) its
centre. Among the projections of \(\mathcal Z\), there is a greatest
projection, denoted by \(E_1\) (respectively, \(E_2,F_1,F_2\)), such
that \(\mathcal A_{E_1}\) (respectively,
\(\mathcal A_{E_2},\mathcal A_{F_1},\mathcal A_{F_2}\)) is finite
(respectively, semifinite, properly infinite, purely infinite). One
has
\[
  E_1F_1=0,
  \qquad
  E_1+F_1=1;
  \qquad
  E_2F_2=0,
  \qquad
  E_2+F_2=1;
  \qquad
  E_1\leq E_2,
  \qquad
  F_1\geq F_2.
\]
\end{proposition}

\begin{proof}
Let \((G_i)_{i\in I}\) be a maximal family of nonzero projections of
\(\mathcal Z\), pairwise orthogonal, such that the
\(\mathcal A_{G_i}\) are finite (respectively, semifinite, properly
infinite, purely infinite). Put
\[
  G=\sum_{i\in I}G_i.
\]
By \cref{prop:I-6-7}, \(G\) is the greatest projection of
\(\mathcal Z\) such that \(\mathcal A_G\) is finite (respectively,
semifinite, properly infinite, purely infinite). This proves the
existence of \(E_1,E_2,F_1,F_2\). It is clear that
\[
  E_1F_1=0,
  \qquad
  E_2F_2=0,
  \qquad
  E_1\leq E_2,
  \qquad
  F_1\geq F_2.
\]
% Source PDF page 106 (printed p. 99)
\phantomsection\label{srcpage:99}
To prove that \(E_1+F_1=1\) (respectively,
\(E_2+F_2=1\)), it is enough to prove that
\(\mathcal A_{1-E_1}\) (respectively,
\(\mathcal A_{1-E_2}\)) is properly infinite (respectively, purely
infinite). But if there existed on
\(\mathcal A_{1-E_1}^+\) (respectively,
\(\mathcal A_{1-E_2}^+\)) a nonzero finite (respectively,
semifinite) normal trace, its support \(E\) would be a nonzero
projection of \(\mathcal Z\) such that \(\mathcal A_E\) was finite
(respectively, semifinite); hence \(E_1\) (respectively, \(E_2\))
would not be the greatest projection of \(\mathcal Z\) such that
\(\mathcal A_{E_1}\) (respectively, \(\mathcal A_{E_2}\)) was finite
(respectively, semifinite).
\end{proof}

\setcounter{corollary}{0}
\begin{corollary}\label{cor:I-6-8-1}
A von Neumann algebra is canonically isomorphic to the product of a
finite von Neumann algebra, a purely infinite von Neumann algebra,
and a properly infinite semifinite von Neumann algebra.
\end{corollary}

\begin{proof}
One has
\[
  \mathcal A
   =\mathcal A_{E_1}\times\mathcal A_{F_2}
      \times\mathcal A_{E_2F_1}.
\]
\end{proof}

\begin{corollary}\label{cor:I-6-8-2}
A factor is either finite, purely infinite, or infinite and
semifinite.
\end{corollary}

\begin{proof}
If \(\mathcal A\) is a factor, then \(E_1=1\), or \(F_2=1\), or
\(E_2F_1=1\).
\end{proof}

\begin{proposition}\label{prop:I-6-9}
Let \(\mathcal A\) be a von Neumann algebra and \(\mathcal Z\) its
centre.
\begin{enumerate}[label=\textup{(\roman*)},leftmargin=*]
\item The following conditions are equivalent:
  \begin{enumerate}[label=\textup{(i\arabic*)},leftmargin=2.6em]
  \item There exists on \(\mathcal A^+\) a faithful, semifinite,
        normal trace;
  \item \(\mathcal A\) is semifinite.
  \end{enumerate}
\item The following conditions are equivalent:
  \begin{enumerate}[label=\textup{(ii\arabic*)},leftmargin=2.9em]
  \item There exists on \(\mathcal A^+\) a faithful finite normal
        trace;
  \item \(\mathcal A\) is finite and of countable type;
  \item \(\mathcal A\) is finite and \(\mathcal Z\) is of countable
        type.
  \end{enumerate}
\item Every finite von Neumann algebra is a product of finite von
Neumann algebras of countable type.
\end{enumerate}
\end{proposition}

\begin{proof}
The implications \(\textup{(i1)}\Rightarrow\textup{(i2)}\) and
\(\textup{(ii1)}\Rightarrow\textup{(ii2)}\Rightarrow
\textup{(ii3)}\) are immediate. Now suppose that \(\mathcal A\) is
semifinite (respectively, finite). Let
\((\varphi_i)_{i\in I}\) be a maximal family of nonzero semifinite
(respectively, finite) normal traces on \(\mathcal A^+\), whose
supports \(E_i\), which are nonzero projections of \(\mathcal Z\),
are pairwise orthogonal. Put \(E=\sum_{i\in I}E_i\), and let us show
that \(E=1\). If \(E\ne1\), there exists a semifinite
(respectively, finite) normal trace \(\varphi\) on \(\mathcal A^+\)
such that \(\varphi(1-E)\ne0\). The trace
\[
  T\longmapsto\varphi\bigl(T(1-E)\bigr)
\]
is normal, semifinite (respectively, finite), nonzero, and its
support is dominated by \(1-E\), which contradicts the maximality
of the family \((\varphi_i)_{i\in I}\). Hence
\(\sum_{i\in I}E_i=1\). This being so,
\(\psi=\sum_{i\in I}\varphi_i\) is a normal trace on
\(\mathcal A^+\), whose support
% Source PDF page 107 (printed p. 100)
\phantomsection\label{srcpage:100}
is \(\sum_{i\in I}E_i=1\), and is therefore faithful; moreover,
\(\psi\) is semifinite. Indeed, let \(T\in\mathcal A^+\), \(T\ne0\).
There exists an \(E_i\) such that \(TE_i\ne0\); let
\(S\in\mathcal A^+\), \(S\ne0\), be such that
\(S\leq TE_i\) and \(\varphi_i(S)<+\infty\). Then \(S\leq T\) and
\(\psi(S)=\varphi_i(S)<+\infty\), which proves our assertion. Thus
the implication \(\textup{(i2)}\Rightarrow\textup{(i1)}\) is
established. If now the \(\varphi_i\) are finite, the algebras
\(\mathcal A_{E_i}\) are finite and of countable type, which proves
\textup{(iii)}. If, in addition, \(\mathcal Z\) is of countable
type, \(I\) is countable; multiplying the \(\varphi_i\) by suitable
scalars, we may suppose that
\(\sum_{i\in I}\varphi_i(1)<+\infty\). Then \(\psi\) is finite,
which proves the implication
\(\textup{(ii3)}\Rightarrow\textup{(ii1)}\).
\end{proof}

\begin{corollary*}\phantomsection\label{cor:I-6-9}
For a von Neumann algebra to be isomorphic to a standard von Neumann
algebra, it is necessary and sufficient that it be semifinite.
\end{corollary*}

\begin{proof}
This follows from \(\cref{thm:I-6-1,thm:I-6-2}\) and part
\textup{(i)} of \(\cref{prop:I-6-9}\).
\end{proof}

The search for conditions under which a von Neumann algebra is
spatially isomorphic to a standard von Neumann algebra is more
difficult. It will be taken up several times in Chapter~III
(\hyperref[no:III-1-5]{\S~1, no.~5};
\hyperref[no:III-6-2]{\S~6, no.~2}).

\begin{proposition}\label{prop:I-6-10}
Let \(\mathcal A\) be a finite von Neumann algebra, \(\mathcal Z\) its
centre, and \(\varphi\) a semifinite normal trace on
\(\mathcal A^+\). For every nonzero \(T\in\mathcal Z^+\), there
exists a nonzero \(T_1\in\mathcal Z^+\), dominated by \(T\), such
that \(\varphi(T_1)<+\infty\).
\end{proposition}

\begin{proof}
Let \(\psi\) be a finite normal trace on \(\mathcal A^+\) such that
\(\psi(T)\ne0\). Replacing \(\varphi\) by \(\varphi+\psi\), we may
suppose that \(\psi\leq\varphi\). There then exists
(\(\cref{thm:I-6-3}\)) an \(S\in\mathcal Z^+\) such that
\(\psi(R)=\varphi(SR)\) for every \(R\in\mathcal A^+\). Since
\(ST\ne0\), there exist a nonzero \(T_1\in\mathcal Z^+\), dominated
by \(T\), and a number \(\lambda>0\), such that
\(ST_1\geq\lambda T_1\). Then
\[
  \varphi(T_1)
   \leq\lambda^{-1}\varphi(ST_1)
   =\lambda^{-1}\psi(T_1)<+\infty.
\]
\end{proof}

\begin{theorem}\label{thm:I-6-6}
Let \(\mathfrak A\) be a Hilbert algebra, \(F\) its characteristic
projection, \(\mathcal Z=\mathcal U(\mathfrak A)\cap
\mathcal V(\mathfrak A)\), and \(E\) the greatest projection of
\(\mathcal Z\) such that \(\mathcal U(\mathfrak A)_E\)
[respectively, \(\mathcal V(\mathfrak A)_E\)] is finite. Then
\(E=F\).
\end{theorem}

\begin{proof}
First let us prove that \(F\leq E\), that is, that
\(\mathcal U(\mathfrak A)_F\) is finite. To this end, let \(T\) be a
nonzero element of \(\mathcal U(\mathfrak A)^+\) such that \(T=TF\).
There exists a central element \(a\in\mathfrak H\) such that
\[
  \lVert T^{1/2}a\rVert^2=(Ta\mid a)\ne0;
\]
for the equality \(Ta=0\) for every central \(a\) would imply
\(TT'a=T'Ta=0\) for every \(T'\in\mathcal V(\mathfrak A)\) and every
central \(a\), and hence \(TF=0\). Thus \(\omega_a\) is a
% Source PDF page 108 (printed p. 101)
\phantomsection\label{srcpage:101}
finite normal trace on \(\mathcal U(\mathfrak A)^+\)
(\(\cref{prop:I-6-3}\)), nonzero on \(T\). Consequently,
\(\mathcal U(\mathfrak A)_F\) is finite.

Let us now prove that \(E\leq F\). For this it will be enough to
establish the following point: for every nonzero projection \(G\) of
\(\mathcal Z\) dominated by \(E\), there exists a nonzero central
element \(a\) such that \(a\in G(\mathfrak H)\). Let \(\varphi\) be
the natural trace on \(\mathcal U(\mathfrak A)^+\). By
\(\cref{prop:I-6-10}\), there exists a nonzero projection \(G'\) of
\(\mathcal Z\), dominated by \(G\), such that
\(\varphi(G')<+\infty\). One has \(G'=U_a\), with \(a\) a nonzero
bounded central element
(\(\cref{thm:I-6-1}\) and
\hyperref[prop:I-5-7]{\S~5, Proposition~7}), and
\(a\in G'(\mathfrak H)\).
\end{proof}

On a semifinite factor \(\mathcal F\), there exists a faithful,
semifinite, normal trace \(\varphi\); let \(\mathfrak m_\varphi\) be
the ideal of definition of \(\dot\varphi\). Every normal trace on
\(\mathcal F^+\) is of the form \(\lambda\varphi\), where
\(0\leq\lambda\leq+\infty\), with the convention
\(0\cdot(+\infty)=0\)
(\hyperref[cor:I-6-2]{Proposition~1, Corollary~2}, and the corollary
to \(\cref{thm:I-6-3}\)). In particular, \(\mathfrak m_\varphi\)
depends only on \(\mathcal F\), and not on the choice of
\(\varphi\). The elements of \(\mathfrak m_\varphi\) (respectively,
\(\mathfrak m_\varphi^{1/2}\)) are called the \emph{trace-class
elements} (respectively, the \emph{Hilbert--Schmidt elements})
relative to \(\mathcal F\).

\begin{theorem}\label{thm:I-6-7}
Let \(\mathcal F\) be a semifinite factor on a Hilbert space
\(\mathfrak H\). Let \(G\) be a group of unitary operators on
\(\mathfrak H\) such that
\[
  U\mathcal F U^{-1}=\mathcal F
  \qquad(U\in G).
\]
Suppose that there exists a vector \(\xi\in\mathfrak H\) having the
following properties:
\begin{enumerate}[label=\textup{(\roman*)},leftmargin=*]
\item \(\xi\) is separating for \(\mathcal F\);
\item \(\mathbb C\xi\) is the set of vectors of \(\mathfrak H\)
invariant under \(G\).
\end{enumerate}
Then \(\mathcal F\) is a finite factor, and \(\xi\) is a trace element
for \(\mathcal F\).
\end{theorem}

\begin{proof}
\textit{a.} Each \(U\in G\) defines the automorphism
\(T\mapsto UTU^{-1}\) of \(\mathcal F\). For every
\(T\in\mathcal F\), put
\[
  \omega(T)=(T\xi\mid\xi).
\]
By \textup{(ii)}, \(\omega\) is invariant under \(G\).

\textit{b.} Choose once and for all a faithful, semifinite, normal
trace \(\psi\) on \(\mathcal F^+\), and a projection \(E\) of
\(\mathcal F\) such that \(0<\psi(E)<+\infty\). One has
\(E\xi\ne0\), since \(\xi\) is separating for \(\mathcal F\). We may
therefore suppose that
\[
  \lVert E\xi\rVert^2=(E\xi\mid\xi)=1.
\]

\textit{c.} For every \(U\in G\), there exists a number
\(\lambda(U)>0\) such that
\[
  \psi(UTU^{-1})=\lambda(U)\psi(T)
\]
for every \(T\in\mathcal F^+\) (corollary to
\(\cref{thm:I-6-3}\)). It is clear that \(\lambda\) is a
homomorphism of \(G\) into the group of positive numbers. Let us
prove that this homomorphism is trivial. Otherwise, there exists a
\(U\in G\) such that \(\lambda(U)<1\). Put
\(E_n=U^nEU^{-n}\), for \(n=1,2,3,\ldots\). For every projection
\(P\) of \(\mathcal F\) such that \(\psi(P)<+\infty\), the linear
functional \(T\mapsto\psi(TP)=\psi(PTP)\) on \(\mathcal F\) is
positive and normal (\(\cref{prop:I-6-1}\)), and its support is
\(P\). Since
\[
  \psi(E_n^*E_nP)=\psi(E_nPE_n)
   \leq\psi(E_n)=\lambda(U)^n\psi(E)\longrightarrow0,
\]
% Source PDF page 109 (printed p. 102)
\phantomsection\label{srcpage:102}
it follows from \hyperref[prop:I-4-4]{\S~4, Proposition~4} that
\(E_nP\) tends strongly to zero. But the union of the \(P(\mathfrak
H)\), for all projections \(P\) of \(\mathcal F\) such that
\(\psi(P)<+\infty\), is total in \(\mathfrak H\); hence \(E_n\) tends
strongly to zero. Yet
\[
  \lVert E_n\xi\rVert=\lVert U^nE\xi\rVert
   =\lVert E\xi\rVert,
  \qquad\text{and hence}\qquad E\xi=0,
\]
which is a contradiction. Thus we have proved that \(\psi\) is
invariant under \(G\).

\textit{d.} Let \(T\in\mathcal F^+\). Let
\(K_T\subset\mathcal F^+\) be the convex hull of the set of
\(UTU^{-1}\), for \(U\in G\), and let
\(K_T'\subset\mathcal F^+\) be the weak closure of \(K_T\). Then
\(K_T'\) is weakly compact, and \(UK_T'U^{-1}=K_T'\) for every
\(U\in G\). Hence the function
\[
  S\longmapsto\lVert S\xi\rVert
\]
on \(K_T'\), which is lower semicontinuous for the weak topology,
attains its minimum at a point \(S_0\) of \(K_T'\). If \(U\in G\),
then
\[
  \lVert US_0U^{-1}\xi\rVert=\lVert S_0\xi\rVert,
  \qquad
  \tfrac12(S_0+US_0U^{-1})\in K_T',
\]
and therefore
\[
  \lVert S_0\xi\rVert
   \leq\left\lVert\tfrac12(S_0+US_0U^{-1})\xi\right\rVert
   \leq\lVert S_0\xi\rVert.
\]
Thus
\[
  \lVert S_0\xi\rVert
   =\left\lVert\tfrac12(S_0+US_0U^{-1})\xi\right\rVert.
\]
Since the norm in \(\mathfrak H\) is strictly convex, it follows that
\[
  S_0\xi=US_0U^{-1}\xi=US_0\xi.
\]
By hypothesis \textup{(ii)}, \(S_0\xi\in\mathbb C\xi\). Since
\(\xi\) is separating for \(\mathcal F\), \(S_0\) is a scalar.
Thus \(K_T'\) contains a nonnegative scalar.

Furthermore, the function \(\omega\) is constant on \(K_T\), and
hence on \(K_T'\). In particular, it has value \(1\) on \(K_E'\).
Thus \(K_E'\) contains a scalar \(\lambda>0\). The function \(\psi\)
is lower semicontinuous for the weak topology
(\hyperref[prop:I-6-2]{corollary to Proposition~2}), and is finite
and constant on \(K_E\) by \textit{c}; it is therefore finite on
\(K_E'\). In particular, \(\psi(\lambda)<+\infty\), and hence
\(\mathcal F\) is finite.

\textit{e.} We may now suppose that \(\psi(1)=1\). Let
\(T\in\mathcal F^+\). Then \(\psi\) and \(\omega\) are constant on
\(K_T\), and hence on \(K_T'\). By \textit{d}, there exists a scalar
\(\mu\in K_T'\). Then
\[
  \omega(T)=\omega(\mu)=\mu\omega(1)
   =\psi(\mu)\omega(1)=\psi(T)\omega(1),
\]
so that \(\omega\) is a trace.
\end{proof}

In Chapter~III, \S~8, we shall be able to formulate
\(\cref{def:I-6-5}\) in an entirely different way.

If \(\mathfrak A\) is a quasi-Hilbert algebra,
\(\mathcal U(\mathfrak A)\) may be purely infinite.

\noindent\textbf{Bibliography:}
\ArticleRef{10}, \ArticleRef{13}, \ArticleRef{19}, \ArticleRef{29},
\ArticleRef{42}, \ArticleRef{65}, \ArticleRef{66}, \ArticleRef{90},
\ArticleRef{117}, \ArticleRef{123}, \ArticleRef{400},
\ArticleRef{416}.

% Source PDF page 110 (printed p. 103)
\phantomsection\label{srcpage:103}
\NumberedParagraph{8}{Classification and Elementary Operations}
\label{no:I-6-8}

\begin{proposition}\label{prop:I-6-11}
Let \(\mathcal A\) be a semifinite (respectively, finite) von Neumann
algebra, and let \(E\) be a projection of \(\mathcal A\) or of
\(\mathcal A'\). Then \(\mathcal A_E\) is semifinite (respectively,
finite).
\end{proposition}

\begin{proof}
Suppose first that \(E\in\mathcal A\). For every
\(T\in\mathcal A_E^+\), one has \(T\circ E\in\mathcal A^+\). For
every normal trace \(\varphi\) on \(\mathcal A^+\), put
\[
  \varphi_E(T)=\varphi(T\circ E).
\]
It is immediate that \(\varphi_E\) is a normal trace on
\(\mathcal A_E^+\), semifinite if \(\varphi\) is semifinite and
finite if \(\varphi\) is finite. This being so, for every nonzero
\(T\in\mathcal A_E^+\), there exists a semifinite (respectively,
finite) normal trace \(\varphi\) on \(\mathcal A^+\) such that
\(\varphi_E(T)=\varphi(T\circ E)\ne0\). Hence \(\mathcal A_E\) is
semifinite (respectively, finite). If now \(E\in\mathcal A'\), then
\(\mathcal A_E\) is isomorphic to an algebra \(\mathcal A_F\), where
\(F\) is a projection of the centre of \(\mathcal A\)
(\hyperref[prop:I-2-2]{\S~2, Proposition~2}); hence
\(\mathcal A_E\) is semifinite (respectively, finite), by the first
part of the proof.
\end{proof}

\begin{proposition}\label{prop:I-6-12}
The tensor product of two semifinite (respectively, finite) von
Neumann algebras is a semifinite (respectively, finite) von Neumann
algebra.
\end{proposition}

\begin{proof}
If two von Neumann algebras are replaced by two isomorphic von
Neumann algebras, their tensor product is replaced by an isomorphic
von Neumann algebra
(\hyperref[prop:I-4-2]{\S~4, Proposition~2}). By the corollary to
\(\cref{prop:I-6-9}\), we may therefore suppose that the given von
Neumann algebras are
\[
  \mathcal A_1=\mathcal U(\mathfrak A_1),
  \qquad
  \mathcal A_2=\mathcal U(\mathfrak A_2),
\]
where \(\mathfrak A_1,\mathfrak A_2\) are Hilbert algebras. Then
\[
  \mathcal U(\mathfrak A_1)\otimes\mathcal U(\mathfrak A_2)
   =\mathcal U(\mathfrak A_1\otimes\mathfrak A_2)
\]
(\hyperref[prop:I-5-9]{\S~5, Proposition~9(i)}), and hence
\(\mathcal U(\mathfrak A_1)\otimes\mathcal U(\mathfrak A_2)\) is
semifinite. If, in addition, \(\mathcal A_1\) and \(\mathcal A_2\)
are finite, the characteristic projections of \(\mathfrak A_1\)
and \(\mathfrak A_2\) are equal to \(1\) (\(\cref{thm:I-6-6}\));
hence the characteristic projection of
\(\mathfrak A_1\otimes\mathfrak A_2\) is equal to \(1\)
(\hyperref[prop:I-5-9]{\S~5, Proposition~9(ii)}), and therefore
\(\mathcal U(\mathfrak A_1)\otimes\mathcal U(\mathfrak A_2)\) is
finite (\(\cref{thm:I-6-6}\)).
\end{proof}

\begin{proposition}\label{prop:I-6-13}
Let \(\mathcal A\) be a semifinite von Neumann algebra. There exist a
von Neumann algebra \(\mathcal A_1\) antiisomorphic to \(\mathcal A\),
a Hilbert space \(\mathfrak K\), and a projection \(E\) of
\(\mathcal A_1\otimes\mathcal L(\mathfrak K)\), such that
\(\mathcal A'\) is spatially isomorphic to
\[
  \bigl(\mathcal A_1\otimes\mathcal L(\mathfrak K)\bigr)_E.
\]
\end{proposition}

\begin{proof}
There exists an isomorphism \(\Phi\) of \(\mathcal A\) onto a
standard von Neumann algebra \(\mathcal B\) (corollary to
\(\cref{prop:I-6-9}\)). Now
\(\Phi^{-1}=\Phi_3\circ\Phi_2\circ\Phi_1\), where \(\Phi_1\) is the
ampliation of \(\mathcal B\) onto
\[
  \mathcal C=\mathcal B\otimes\mathbb C_{\mathfrak K}
\]
(\(\mathfrak K\) being a suitable Hilbert space), where \(\Phi_2\)
is the induction of \(\mathcal C\) on
\(\mathcal D=\mathcal C_E\) (\(E\) being a suitable projection of
\(\mathcal C'\)), and where \(\Phi_3\) is a spatial isomorphism of
\(\mathcal D\) onto \(\mathcal A\). Then \(\mathcal A'\) is
spatially isomorphic to
\[
  \mathcal D'=\mathcal C_E'
    =\bigl(\mathcal B'\otimes\mathcal L(\mathfrak K)\bigr)_E,
\]
and it is enough to put \(\mathcal B'=\mathcal A_1\).
\end{proof}

% Source PDF page 111 (printed p. 104)
\phantomsection\label{srcpage:104}
\setcounter{corollary}{0}
\begin{corollary}\label{cor:I-6-13-1}
The commutant of a semifinite von Neumann algebra is a semifinite von
Neumann algebra.
\end{corollary}

\begin{proof}
This follows from \hyperref[prop:I-6-11]{Propositions~11},
\hyperref[prop:I-6-12]{12}, and \hyperref[prop:I-6-13]{13}.
\end{proof}

\begin{corollary}\label{cor:I-6-13-2}
Let \(\mathcal A\) be a von Neumann algebra, and let \(\mathcal Z\)
be its centre, which is also the centre of \(\mathcal A'\). The
greatest projection \(E\) of \(\mathcal Z\) such that
\(\mathcal A_E\) is semifinite is equal to the greatest projection
\(F\) of \(\mathcal Z\) such that \(\mathcal A'_F\) is semifinite.
\end{corollary}

\begin{proof}
By \(\cref{cor:I-6-13-1}\), \(\mathcal A'_E\) is semifinite, and
hence \(E\leq F\). Interchanging the roles of \(\mathcal A\) and
\(\mathcal A'\), one has \(F\leq E\).
\end{proof}

\begin{corollary}\label{cor:I-6-13-3}
Let \(\mathcal A\) be a purely infinite von Neumann algebra. Then
\(\mathcal A'\) is purely infinite.
\end{corollary}

\begin{proof}
The projection \(E\) of \(\cref{cor:I-6-13-2}\) is equal to zero.
\end{proof}

\begin{corollary}\label{cor:I-6-13-4}
Let \(\mathcal A\) be a purely infinite von Neumann algebra, and let
\(E\) be a projection of \(\mathcal A\) or of \(\mathcal A'\). Then
\(\mathcal A_E\) is purely infinite.
\end{corollary}

\begin{proof}
Since \(\mathcal A'\) is purely infinite by
\(\cref{cor:I-6-13-3}\), it is enough to consider the case in which
\(E\in\mathcal A'\). There then exists a projection \(F\) of the
centre of \(\mathcal A\) such that \(\mathcal A_E\) is isomorphic to
\(\mathcal A_F\). Since
\[
  \mathcal A=\mathcal A_F\times\mathcal A_{1-F},
\]
\(\mathcal A_F\) is purely infinite (\(\cref{prop:I-6-7}\)).
\end{proof}

The commutant of a finite (respectively, properly infinite) von
Neumann algebra is not necessarily finite (respectively, properly
infinite), as is shown by the case of
\(\mathcal L(\mathfrak H)\) and \(\mathbb C_{\mathfrak H}\).
(Cf., however, Chapter~III, \S~2, Exercise~2.)

For some supplements, cf. Exercises~10 and~11, and Chapter~III,
\S~8, Theorem~2.

\noindent\textit{Problem.}\quad Is \(\cref{prop:I-6-13}\) valid for
an arbitrary von Neumann algebra?

\noindent\textbf{Bibliography:}
\ArticleRef{40}, \ArticleRef{65}, \ArticleRef{67}, \ArticleRef{89}.

\Needspace{4\baselineskip}
\NumberedParagraph{9}{Commutant of the Tensor Product of Two
Semifinite von Neumann Algebras}
\label{no:I-6-9}

\begin{proposition}\label{prop:I-6-14}
Let \(\mathcal A_1\) and \(\mathcal A_2\) be semifinite von Neumann
algebras. Then
\[
  (\mathcal A_1\otimes\mathcal A_2)'
    =\mathcal A_1'\otimes\mathcal A_2'.
\]
\end{proposition}

\begin{proof}
There exist Hilbert algebras \(\mathfrak A_1,\mathfrak A_2\) such
that \(\mathcal A_1\) is isomorphic to
\(\mathcal U(\mathfrak A_1)\), and \(\mathcal A_2\) to
\(\mathcal U(\mathfrak A_2)\). Therefore
(\hyperref[thm:I-4-3]{\S~4, Theorem~3}), we may suppose
\[
  \mathcal A_1
   =\bigl(\mathcal U(\mathfrak A_1)
          \otimes\mathbb C_{\mathfrak K_1}\bigr)_{E_1'},
  \qquad
  \mathcal A_2
   =\bigl(\mathcal U(\mathfrak A_2)
          \otimes\mathbb C_{\mathfrak K_2}\bigr)_{E_2'},
\]
where
\[
  E_1'\in
   \bigl(\mathcal U(\mathfrak A_1)
          \otimes\mathbb C_{\mathfrak K_1}\bigr)',
  \qquad
  E_2'\in
   \bigl(\mathcal U(\mathfrak A_2)
          \otimes\mathbb C_{\mathfrak K_2}\bigr)'.
\]
% Source PDF page 112 (printed p. 105)
\phantomsection\label{srcpage:105}
Putting \(E'=E_1'\otimes E_2'\), one then has
\begin{align*}
  \mathcal A_1\otimes\mathcal A_2
   &=
    \bigl(\mathcal U(\mathfrak A_1)
      \otimes\mathbb C_{\mathfrak K_1}
      \otimes\mathcal U(\mathfrak A_2)
      \otimes\mathbb C_{\mathfrak K_2}\bigr)_{E'}\\
   &=
    \bigl(\mathcal U(\mathfrak A_1\otimes\mathfrak A_2)
      \otimes\mathbb C_{\mathfrak K_1\otimes\mathfrak K_2}
    \bigr)_{E'}.
\end{align*}
On the other hand,
\begin{align*}
  \mathcal A_1'
   &=\bigl(\mathcal U(\mathfrak A_1)
      \otimes\mathbb C_{\mathfrak K_1}\bigr)_{E_1'}'
     =\bigl(\mathcal V(\mathfrak A_1)
      \otimes\mathcal L(\mathfrak K_1)\bigr)_{E_1'},\\
  \mathcal A_2'
   &=\bigl(\mathcal U(\mathfrak A_2)
      \otimes\mathbb C_{\mathfrak K_2}\bigr)_{E_2'}'
     =\bigl(\mathcal V(\mathfrak A_2)
      \otimes\mathcal L(\mathfrak K_2)\bigr)_{E_2'},
\end{align*}
and hence
\begin{align*}
  \mathcal A_1'\otimes\mathcal A_2'
   &=
    \bigl(\mathcal V(\mathfrak A_1)\otimes\mathcal L(\mathfrak K_1)
      \otimes\mathcal V(\mathfrak A_2)
      \otimes\mathcal L(\mathfrak K_2)\bigr)_{E'}\\
   &=
    \bigl(\mathcal V(\mathfrak A_1\otimes\mathfrak A_2)
      \otimes\mathcal L(\mathfrak K_1\otimes\mathfrak K_2)
    \bigr)_{E'}\\
   &=
    \bigl(\mathcal U(\mathfrak A_1\otimes\mathfrak A_2)
      \otimes\mathbb C_{\mathfrak K_1\otimes\mathfrak K_2}
    \bigr)_{E'}'
    =(\mathcal A_1\otimes\mathcal A_2)'.
\end{align*}
\end{proof}

\begin{corollary*}\phantomsection\label{cor:I-6-14}
Let \(\mathcal A_1\) and \(\mathcal A_2\) be arbitrary von Neumann
algebras, let \(\mathcal Z_1\) and \(\mathcal Z_2\) be their centres,
and put \(\mathcal A=\mathcal A_1\otimes\mathcal A_2\). The centre
\(\mathcal Z\) of \(\mathcal A\) is equal to
\(\mathcal Z_1\otimes\mathcal Z_2\). In particular, if
\(\mathcal A_1\) and \(\mathcal A_2\) are factors, then
\(\mathcal A_1\otimes\mathcal A_2\) is a factor.
\end{corollary*}

\begin{proof}
Let \(\mathfrak H_1\) and \(\mathfrak H_2\) be the spaces on which
\(\mathcal A_1\) and \(\mathcal A_2\) act. One has
\[
  \mathcal Z_1\otimes\mathbb C_{\mathfrak H_2}\subset\mathcal Z,
  \qquad
  \mathbb C_{\mathfrak H_1}\otimes\mathcal Z_2\subset\mathcal Z,
  \qquad\text{hence}\qquad
  \mathcal Z_1\otimes\mathcal Z_2\subset\mathcal Z.
\]
On the other hand,
\[
  \mathcal A_1\otimes\mathbb C_{\mathfrak H_2}
      \subset\mathcal A\subset\mathcal Z',
  \qquad
  \mathcal A_1'\otimes\mathbb C_{\mathfrak H_2}
      \subset\mathcal A'\subset\mathcal Z'.
\]
But the von Neumann algebra generated by \(\mathcal A_1\) and
\(\mathcal A_1'\) is \(\mathcal Z_1'\); hence
(\hyperref[prop:I-2-6]{\S~2, Proposition~6})
\(\mathcal Z_1'\otimes\mathbb C_{\mathfrak H_2}\subset\mathcal Z'\).
Similarly,
\(\mathbb C_{\mathfrak H_1}\otimes\mathcal Z_2'\subset\mathcal Z'\).
Thus \(\mathcal Z_1'\otimes\mathcal Z_2'\subset\mathcal Z'\). Now,
by \(\cref{prop:I-6-14}\),
\[
  \mathcal Z_1'\otimes\mathcal Z_2'
    =(\mathcal Z_1\otimes\mathcal Z_2)'.
\]
Consequently, \(\mathcal Z_1\otimes\mathcal Z_2\supset\mathcal Z\).
\end{proof}

The last assertion of the corollary is very elementary
(\hyperref[ex:I-2-4]{\S~2, Exercise~4}). Recall, on the other hand,
the problem posed at the end of \hyperref[no:I-2-4]{\S~2, no.~4}.

\noindent\textbf{Bibliography:}
\ArticleRef{65}, \ArticleRef{128}, \ArticleRef{432}.

\NumberedParagraph{10}{The Space \(L^1\) Defined by a Trace}
\label{no:I-6-10}
Let \(\mathcal A\) be a von Neumann algebra, and let \(\varphi\) be a
trace defined on the two-sided ideal \(\mathfrak m\) of
\(\mathcal A\).

\begin{lemma}\label{lem:I-6-2}
Let \(S\in\mathcal A\) and \(T\in\mathfrak m\). Then
\[
  |\varphi(ST)|
   \leq
   \bigl[
     \varphi(|S^*|\,|T|)
     \varphi(|S|\,|T^*|)
   \bigr]^{1/2}.
\]
\end{lemma}

\begin{proof}
Let \(S=U|S|\) and \(T=V|T|\) be the polar decompositions of \(S\)
and \(T\). One has
\(|T|\in\mathfrak m\) and
\(|T|^{1/2}\in\mathfrak m^{1/2}\); hence, by
\hyperref[prop:I-1-11]{\S~1, Proposition~11},
\[
  \varphi(ST)
   =\varphi\bigl(|T|^{1/2}U|S|V|T|^{1/2}\bigr)
   =\varphi\!\left[
      \bigl(|T|^{1/2}U|S|^{1/2}\bigr)
      \bigl(|T|^{1/2}V^*|S|^{1/2}\bigr)^*
    \right].
\]
% Source PDF page 113 (printed p. 106)
\phantomsection\label{srcpage:106}
Since
\[
  |T|^{1/2}U|S|^{1/2}\in\mathfrak m^{1/2},
  \qquad
  |T|^{1/2}V^*|S|^{1/2}\in\mathfrak m^{1/2},
\]
the Cauchy--Schwarz inequality gives
\begin{align*}
  |\varphi(ST)|^2
  &\leq
   \varphi\!\left[
      \bigl(|T|^{1/2}U|S|^{1/2}\bigr)
      \bigl(|T|^{1/2}U|S|^{1/2}\bigr)^*
   \right]\\
  &\quad{}\times
   \varphi\!\left[
      \bigl(|T|^{1/2}V^*|S|^{1/2}\bigr)
      \bigl(|T|^{1/2}V^*|S|^{1/2}\bigr)^*
   \right]\\
  &=\varphi(U|S|U^*|T|)\,
    \varphi(|S|V|T|V^*).
\end{align*}
But \(|S^*|=U|S|U^*\) and \(|T^*|=V|T|V^*\).
\end{proof}

\begin{theorem}\label{thm:I-6-8}
Let \(\mathcal A\) be a von Neumann algebra, let \(\varphi\) be a
trace defined on the two-sided ideal \(\mathfrak m\) of
\(\mathcal A\), let \(S\in\mathcal A\), and let
\(T\in\mathfrak m\). Then
\[
  |\varphi(ST)|\leq\varphi(|ST|)
    \leq\lVert S\rVert\,\varphi(|T|).
\]
\end{theorem}

\begin{proof}
Suppose first that \(S\geq0\) and \(T\geq0\). One has
\[
  T^{1/2}ST^{1/2}\leq\lVert S\rVert T,
  \qquad\text{and hence}\qquad
  \varphi(ST)=\varphi(T^{1/2}ST^{1/2})
     \leq\lVert S\rVert\varphi(T).
\]
If now \(S\in\mathcal A\) and \(T\in\mathfrak m\) are arbitrary,
then
\[
  \varphi(|S^*|\,|T|)
    \leq\lVert|S^*|\rVert\varphi(|T|)
    =\lVert S\rVert\varphi(|T|),
\]
\[
  \varphi(|S|\,|T^*|)
    \leq\lVert|S|\rVert\varphi(|T^*|)
    =\lVert S\rVert\varphi(|T|);
\]
hence, by \(\cref{lem:I-6-2}\),
\[
  |\varphi(ST)|\leq\lVert S\rVert\varphi(|T|).
\]
Consequently,
\[
  |\varphi(ST)|
   =|\varphi(ST\cdot1)|
   \leq\lVert1\rVert\varphi(|ST|)
   =\varphi(|ST|),
\]
which is the first inequality of the theorem.

On the other hand, let \(ST=W|ST|\) be the polar decomposition of
\(ST\). One has
\[
  \varphi(|ST|)
    =\varphi(W^*ST)
    \leq\lVert W^*S\rVert\varphi(|T|)
    \leq\lVert S\rVert\varphi(|T|),
\]
which is the second inequality of the theorem.
\end{proof}

\setcounter{corollary}{0}
\begin{corollary}\label{cor:I-6-thm8-1}
Let \(T\in\mathfrak m\). The number \(\varphi(|T|)\) is the supremum
of the numbers \(|\varphi(ST)|\), where \(S\) ranges over the set of
operators of \(\mathcal A\) such that \(\lVert S\rVert\leq1\). This
supremum is attained.
\end{corollary}

\begin{proof}
Let \(T=V|T|\) still be the polar decomposition of \(T\). One has
\[
  \varphi(|T|)=|\varphi(V^*T)|,
  \qquad
  \lVert V^*\rVert\leq1.
\]
\end{proof}

\begin{corollary}\label{cor:I-6-thm8-2}
The function \(T\mapsto\varphi(|T|)\), defined on
\(\mathfrak m\), is a seminorm.
\end{corollary}

\begin{proof}
By \(\cref{cor:I-6-thm8-1}\), this function is the upper envelope of
seminorms.
\end{proof}

Henceforth suppose that \(\varphi\) is normal, faithful, and
semifinite. The space \(\mathfrak m\), endowed with the norm
\(T\mapsto\varphi(|T|)\), is a complex normed vector space.

% Source PDF page 114 (printed p. 107)
\phantomsection\label{srcpage:107}
We shall denote by \(L^1(\varphi)\), or more briefly by \(L^1\), the
complex Banach space obtained by completing this normed space, and
\(\mathfrak m\) will be identified with an everywhere dense vector
subspace of \(L^1\). The norm of an element \(f\in L^1\) will be
denoted by \(\lVert f\rVert_{1,\varphi}\), or by
\(\lVert f\rVert_1\).

Let \(T\in\mathfrak m\). The linear functional
\(S\mapsto\varphi(ST)\) on \(\mathcal A\) is an ultra-weakly
continuous linear functional \(\varphi_T\)
(\(\cref{prop:I-6-1}\)); hence, with the notation of
\hyperref[no:I-3-3]{\S~3, no.~3},
\(\varphi_T\in\mathcal A_*\). Moreover,
\[
  \lVert\varphi_T\rVert=\varphi(|T|)=\lVert T\rVert_1,
\]
by \(\cref{cor:I-6-thm8-1}\). Thus the mapping
\(T\mapsto\varphi_T\) is a linear isometry of \(\mathfrak m\) onto a
vector subspace of \(\mathcal A_*\). We shall prove that this
subspace is everywhere dense in \(\mathcal A_*\) in norm. Since
\(\mathcal A\) is identified with the dual of \(\mathcal A_*\)
(\hyperref[thm:I-3-1]{\S~3, Theorem~1(iii)}), it is enough to show
that, for every \(S\in\mathcal A\), there exists a
\(T\in\mathfrak m\) such that \(\varphi_T(S)\ne0\). Let
\(S=U|S|\) be the polar decomposition of \(S\). There exists an
\(R\in\mathfrak m^+\) such that \(R|S|R\ne0\) (for \(R\) may tend
strongly to \(1\) in \(\mathfrak m^+\) while remaining dominated by
\(1\)); then
\[
  \varphi_{R^2U^*}(S)
   =\varphi(R^2U^*S)
   =\varphi(R|S|R)\ne0,
\]
which proves our assertion. It follows that the mapping
\(T\mapsto\varphi_T\) extends uniquely to an isomorphism of the
Banach space \(L^1(\varphi)\) onto the Banach space
\(\mathcal A_*\) of ultra-weakly continuous linear functionals on
\(\mathcal A\). Using again
\hyperref[thm:I-3-1]{\S~3, Theorem~1(iii)}, we obtain the following
result.

\begin{theorem}\label{thm:I-6-9}
Let \(\mathcal A\) be a von Neumann algebra, and let \(\varphi\) be
a faithful, semifinite, normal trace defined on the two-sided ideal
\(\mathfrak m\) of \(\mathcal A\). When the normed spaces
\(\mathcal A\) (with norm \(\lVert\,\cdot\,\rVert\)) and
\(\mathfrak m\) (with norm \(\lVert\,\cdot\,\rVert_1\)) are put in
duality by the bilinear form
\[
  (S,T)\longmapsto\varphi(ST),
\]
\(\mathcal A\) is the dual of \(\mathfrak m\).
\end{theorem}

\noindent\textbf{Bibliography:}
\ArticleRef{15}, \ArticleRef{32}, \ArticleRef{76}, \ArticleRef{101},
\TreatiseRef{13}.

\NumberedParagraph{11}{Trace and Determinant}
\label{no:I-6-11}
Let \(\mathcal A\) be a von Neumann algebra, and let \(T\) be an
invertible element of \(\mathcal A\), with \(T=U|T|\) its polar
decomposition. Then \(U\) is unitary and \(|T|\) is invertible. We
may therefore form \(\log|T|\), which is a Hermitian element of
\(\mathcal A\).

\begin{definition}\label{def:I-6-6}
Let \(\varphi\) be a finite trace on \(\mathcal A\) such that
\(\varphi(1)=1\). The \emph{determinant associated with
\(\varphi\)} is the function \(\Delta\), defined on the set of
invertible elements of \(\mathcal A\) by the formula
\[
  \Delta(T)=\exp\bigl[\varphi(\log|T|)\bigr].
\]
\end{definition}

It is immediate that
\[
  \Delta(T)=\Delta(|T|)
   =\bigl[\Delta(|T|^2)\bigr]^{1/2}
   =\bigl[\Delta(T^*T)\bigr]^{1/2}.
\]

\begin{lemma}\label{lem:I-6-3}
Let \(f(\lambda)\) be an analytic function in an open subset of the
complex plane, and let \(\Gamma\) be a Jordan curve in this subset.
Let \(t\mapsto S_t\) be
% Source PDF page 115 (printed p. 108)
\phantomsection\label{srcpage:108}
a differentiable mapping (in the sense of the norm in
\(\mathcal A\)) of an open interval \(I\) of the real line into
\(\mathcal A\). Suppose that the spectrum of \(S_t\) is in the
interior of \(\Gamma\), for \(t\in I\). Then
\(f(S_t)\in\mathcal A\) is a differentiable function of \(t\), and
\[
  \varphi\!\left[\frac{d}{dt}f(S_t)\right]
    =\varphi\bigl[g(S_t)S_t'\bigr],
  \qquad\text{where}\qquad
  g(\lambda)=\frac{df(\lambda)}{d\lambda},
  \quad
  S_t'=\frac{dS_t}{dt}.
\]
\end{lemma}

\begin{proof}
For \(t\in I\), one has
\[
  f(S_t)=\frac{1}{2\pi i}
    \int_\Gamma f(\lambda)(\lambda-S_t)^{-1}\,d\lambda.
\]
Taking account of the equality
\[
  (\lambda-S_{t'})^{-1}-(\lambda-S_t)^{-1}
   =(\lambda-S_{t'})^{-1}(S_{t'}-S_t)(\lambda-S_t)^{-1},
\]
one has
\[
  \frac{f(S_{t'})-f(S_t)}{t'-t}
   =\frac{1}{2\pi i}\int_\Gamma
      f(\lambda)(\lambda-S_{t'})^{-1}
      \frac{S_{t'}-S_t}{t'-t}
      (\lambda-S_t)^{-1}\,d\lambda.
\]
As \(t'\to t\), \((\lambda-S_{t'})^{-1}\) tends, in the norm of
\(\mathcal A\), to \((\lambda-S_t)^{-1}\), uniformly with respect to
\(\lambda\) on \(\Gamma\), and
\((t'-t)^{-1}(S_{t'}-S_t)\) tends to \(S_t'\). Hence \(f(S_t)\) is
differentiable with respect to \(t\), and
\[
  \frac{d}{dt}f(S_t)
   =\frac{1}{2\pi i}\int_\Gamma
      f(\lambda)(\lambda-S_t)^{-1}S_t'
      (\lambda-S_t)^{-1}\,d\lambda.
\]
On the other hand, the equality
\[
  (\lambda'-S_t)^{-1}-(\lambda-S_t)^{-1}
   =(\lambda'-S_t)^{-1}(\lambda-\lambda')
      (\lambda-S_t)^{-1}
\]
gives, as \(\lambda'\to\lambda\),
\[
  \frac{d}{d\lambda}(\lambda-S_t)^{-1}
    =-(\lambda-S_t)^{-2}.
\]
Thus
\[
  \frac{d}{d\lambda}
    \bigl[f(\lambda)(\lambda-S_t)^{-1}\bigr]
   =g(\lambda)(\lambda-S_t)^{-1}
      -f(\lambda)(\lambda-S_t)^{-2},
\]
and consequently
\[
  g(S_t)
   =\frac{1}{2\pi i}\int_\Gamma
      g(\lambda)(\lambda-S_t)^{-1}\,d\lambda
   =\frac{1}{2\pi i}\int_\Gamma
      f(\lambda)(\lambda-S_t)^{-2}\,d\lambda.
\]
Therefore
\begin{align*}
  \varphi\!\left[\frac{d}{dt}f(S_t)\right]
   &=\frac{1}{2\pi i}\int_\Gamma
      f(\lambda)\,
      \varphi\bigl[(\lambda-S_t)^{-1}S_t'
                       (\lambda-S_t)^{-1}\bigr]\,d\lambda\\
   &=\frac{1}{2\pi i}\int_\Gamma
      f(\lambda)\,
      \varphi\bigl[(\lambda-S_t)^{-2}S_t'\bigr]\,d\lambda
    =\varphi\bigl[g(S_t)S_t'\bigr].
\end{align*}
\end{proof}

% Source PDF page 116 (printed p. 109)
\phantomsection\label{srcpage:109}
\begin{lemma}\label{lem:I-6-4}
Let \(S\) and \(T\) be elements of \(\mathcal A\), with \(T\)
Hermitian. Then
\[
  \Delta(\exp S^*\exp T\exp S)
    =\exp\varphi(S^*+T+S).
\]
\end{lemma}

\begin{proof}
For real \(t\), put
\[
  S_t=\exp(tS^*)\exp T\exp(tS)\in\mathcal A.
\]
Each \(S_t\) is positive Hermitian and invertible. Moreover,
\(S_t\) is a differentiable function of \(t\). There exist numbers
\(\alpha>0\), \(\alpha'<+\infty\), such that
\[
  \alpha<\lVert S_t^{-1}\rVert^{-1}
    \leq\lVert S_t\rVert<\alpha'
  \qquad(t\in[0,1]).
\]
Apply \(\cref{lem:I-6-3}\), taking for \(f(\lambda)\) the principal
determination of \(\log\lambda\) in the complex plane with the
negative part of the real axis removed, and for \(\Gamma\) a
suitable rectangle containing \(\alpha\) and \(\alpha'\) in its
interior. Since \(S_t'=S^*S_t+S_tS\), one has
\[
  \frac{d}{dt}\varphi(\log S_t)
   =\varphi(S_t^{-1}S_t')
   =\varphi(S_t^{-1}S^*S_t+S)
   =\varphi(S^*+S),
\]
and hence
\[
  \varphi(\log S_1)-\varphi(\log S_0)
    =\varphi(S^*+S),
\]
that is,
\[
  \log\Delta(\exp S^*\exp T\exp S)-\varphi(T)
    =\varphi(S^*+S).
\]
\end{proof}

\begin{theorem}\label{thm:I-6-10}
The function \(\Delta\) has the following properties:
\begin{enumerate}[label=\textup{(\roman*)},leftmargin=*]
\item \(\Delta(\lambda1)=|\lambda|\) for \(\lambda\ne0\);
\item
  \[
    \Delta(T^*)=\Delta(T)=\Delta(T^*T)^{1/2}
  \]
  for \(T\in\mathcal A\), \(T\) invertible;
\item
  \[
    \Delta(ST)=\Delta(S)\Delta(T)
  \]
  for \(S,T\in\mathcal A\), \(S\) and \(T\) invertible;
\item \(\Delta(\exp T)=|\exp\varphi(T)|\)\enspace for \(T\in\mathcal A\);
\item
  \[
    \Delta(T)\leq\lim_{n\to+\infty}\lVert T^n\rVert^{1/n}
  \]
  for \(T\in\mathcal A\), \(T\) invertible.
\end{enumerate}
\end{theorem}

\begin{proof}
\textup{(i)} For \(\lambda\ne0\),
\[
  \Delta(\lambda1)
    =\exp\varphi(\log|\lambda|1)
    =\exp\log|\lambda|=|\lambda|.
\]

\textup{(ii)} For \(T\in\mathcal A\), \(T\) invertible, we have
already observed that \(\Delta(T^*T)=\Delta(T)^2\). On the other
hand, \(TT^*=U(T^*T)U^{-1}\), with a unitary operator
\(U\in\mathcal A\), and hence
\[
  \log(TT^*)=U\log(T^*T)U^{-1},
  \qquad\text{so that}\qquad
  \Delta(T^*T)=\Delta(TT^*).
\]

\textup{(iii)} Let \(S,T\) be invertible elements of \(\mathcal A\),
and let \(S=U|S|\), \(T=V|T|\) be their polar decompositions. Taking
account of \(\cref{lem:I-6-4}\) and of \textup{(ii)}, one has
\begin{align*}
  \Delta(ST)
   &=\Delta(U|S|V|T|)\\
   &=\Delta\bigl(
       |T|V^*|S|U^*U|S|V|T|
     \bigr)^{1/2}\\
   &=\Delta\bigl(|T|V^*|S|^2V|T|\bigr)^{1/2}\\
   &=\bigl[
      \exp\varphi\{
        \log|T|+\log(V^*|S|^2V)+\log|T|
      \}
     \bigr]^{1/2}\\
   &=\exp\varphi(\log|S|+\log|T|)
    =\Delta(S)\Delta(T).
\end{align*}
% Source PDF page 117 (printed p. 110)
\phantomsection\label{srcpage:110}
\textup{(iv)} Taking account of \textup{(ii)}, and writing
\(\Re\lambda\) for the real part of a complex number \(\lambda\),
\(\cref{lem:I-6-4}\) gives
\[
  \Delta(\exp T)
   =\bigl[\Delta(\exp T^*\exp T)\bigr]^{1/2}
   =\bigl[\exp\varphi(T^*+T)\bigr]^{1/2}
   =\exp\Re\varphi(T)
   =|\exp\varphi(T)|.
\]

\textup{(v)} For \(T\in\mathcal A\), \(T\) invertible, one has
\(|T|\leq\lVert T\rVert\), and hence
\[
  \log|T|\leq\log\lVert T\rVert,
  \qquad\text{so that}\qquad
  \Delta(T)\leq\exp\log\lVert T\rVert=\lVert T\rVert.
\]
Consequently, taking account of \textup{(iii)},
\[
  \Delta(T)=\Delta(T^n)^{1/n}
     \leq\lVert T^n\rVert^{1/n}
\]
for every positive integer \(n\).
\end{proof}

\begin{corollary*}\phantomsection\label{cor:I-6-10}
For every \(T\in\mathcal A\), \(\varphi(T)\) belongs to the convex
hull of the spectrum of \(T\).
\end{corollary*}

\begin{proof}
It is enough to prove that \(\varphi(T)\) belongs to every closed
half-plane containing the spectrum \(\mathfrak S(T)\) of \(T\).
[Recall that, since \(\mathfrak S(T)\) is a compact subset of the
complex plane, its convex hull is closed.] On the other hand, if
\(T\) is replaced by \(\alpha T+\alpha'\), where \(\alpha\) and
\(\alpha'\) are two complex numbers, \(\varphi(T)\) is replaced by
\(\alpha\varphi(T)+\alpha'\), and \(\mathfrak S(T)\) by
\(\alpha\mathfrak S(T)+\alpha'\). We may therefore suppose that the
half-plane under consideration is the half-plane
\(\Re\lambda\leq0\). Then \(\mathfrak S(\exp T)\), which is
identical with \(\exp\mathfrak S(T)\), is contained inside the
circle \(|\lambda|=1\). Hence
\[
  |\exp\varphi(T)|=\Delta(\exp T)\leq1,
\]
by \(\cref{thm:I-6-10}\textup{(iv),(v)}\). Consequently,
\(\Re\varphi(T)\leq0\).
\end{proof}

Cf. Chapter~III, \S~2, Exercise~6.

In proving the results of this number (results which will not be
used in the sequel), we have appealed to theories not appearing in
the list drawn up in the Introduction. The reader may, however,
regard analytic functions of an operator as defined by Cauchy's
integral, and the necessary knowledge is then concentrated in
Theorem~24~D of \TreatiseRef{9}.

\noindent\textbf{Bibliography:}
\ArticleRef{23}, \ArticleRef{24}.

\par\medskip
\noindent\textit{Exercises.}---
\begin{enumerate}[label=\arabic*.,leftmargin=*,itemsep=0.5\baselineskip]
\item\label{ex:I-6-1}
Let \(\mathcal A\) be a von Neumann algebra, and let \(\varphi\) be a
trace on \(\mathcal A^+\). For \(\varphi\) to be normal, it is
necessary and sufficient that \(\varphi\) be lower semicontinuous
for the weak topology (use the corollary to
\(\cref{prop:I-6-2}\)). If this is so, and if
\(T_0\in\mathcal A^+\) is such that \(\varphi(T_0)<+\infty\), the
restriction of \(\varphi\) to the set of elements of
\(\mathcal A^+\) dominated by \(T_0\) is weakly continuous.

\item\label{ex:I-6-2}
Let \(\mathcal A\) be a von Neumann algebra of countable type on
\(\mathfrak H\), and let \(\varphi\) be a normal trace on
\(\mathcal A^+\). Show that there exists a sequence
\((x_1,x_2,\ldots)\) of elements of \(\mathfrak H\) such that
\[
  \varphi=\sum_{i=1}^{\infty}\omega_{x_i}.
\]
[If \(\varphi(T)=+\infty\) for every nonzero
\(T\in\mathcal A^+\), use a countable separating set for
\(\mathcal A\). If \(\varphi\) is semifinite, form a sequence
\((E_i)\) of pairwise orthogonal projections of \(\mathcal A\) such
that
\[
  \sum_{i=1}^{\infty}E_i=1,
  \qquad
  \varphi(E_i)<+\infty
  \quad\text{for every }i;
\]
then use \(\cref{prop:I-6-2}\).]

% Source PDF page 118 (printed p. 111)
\phantomsection\label{srcpage:111}
\item\label{ex:I-6-3}
Let \(\mathcal A\) be a von Neumann algebra, let \(\varphi\) be a
faithful trace on \(\mathcal A^+\), let \(\mathfrak m\) be the ideal
of definition of \(\dot\varphi\), and let \(S,T\) be Hermitian
operators in \(\mathfrak m^{1/2}\). If
\[
  \dot\varphi(ST)\leq
  \dot\varphi(SUTU^{-1})
\]
for every unitary operator \(U\) of \(\mathcal A\), then \(S\) and
\(T\) commute. [Let \(R\) be a Hermitian operator of
\(\mathcal A\). Show that the real-variable function
\[
  t\longmapsto
  \dot\varphi\bigl(Se^{itR}Te^{-itR}\bigr)
\]
is differentiable, with derivative
\[
  \dot\varphi\bigl(
    iSRe^{itR}Te^{-itR}
    -iSe^{itR}TR e^{-itR}
  \bigr).
\]
Write that this derivative is zero for \(t=0\), and take
\(R=i(ST-TS)\).] The result is identical if
\[
  \dot\varphi(ST)\geq
  \dot\varphi(SUTU^{-1})
\]
for every unitary operator \(U\) of \(\mathcal A\).
\ArticleRef{15}, \ArticleRef{76}

\item\label{ex:I-6-4}
Let \(\mathfrak A\) be a Hilbert algebra having an identity element
\(e\). The trace \(\omega_e\) on
\(\mathcal U(\mathfrak A)^+\) is the natural trace.

\item\label{ex:I-6-5}
Prove directly that, if an abelian von Neumann algebra
\(\mathcal A\) on \(\mathfrak H\) possesses a totalizing element
\(x\), then \(\mathcal A=\mathcal A'\). [Let
\(T'\in\mathcal A'\); it is a matter of proving that \(T'\) is
normal. There exists a sequence \((T_n)\) in \(\mathcal A\) such
that \(T_nx\to T'x\). Since \(T_m-T_n\) is normal, the \(T_n^*x\)
form a Cauchy sequence and have a limit \(y\). For every
\(T\in\mathcal A\),
\[
  (y\mid Tx)
   =\lim_n(T_n^*x\mid Tx)
   =\lim_n(x\mid TT_nx)
   =(x\mid TT'x)
   =(T'^*x\mid Tx),
\]
and hence \(y=T'^*x\). Thus, for every \(T\in\mathcal A\),
\[
  T_n(Tx)\longrightarrow T(T'x)=T'(Tx),
  \qquad
  T_n^*(Tx)\longrightarrow T(T'^*x)=T'^*(Tx).
\]
The equality
\(\lVert T_n(Tx)\rVert=\lVert T_n^*(Tx)\rVert\) gives, on passing
to the limit,
\[
  \lVert T'(Tx)\rVert=\lVert T'^*(Tx)\rVert;
\]
hence \(\lVert T'z\rVert=\lVert T'^*z\rVert\) for every
\(z\in\mathfrak H\).]

\item\label{ex:I-6-6}
Let \(\mathfrak H\) be an infinite-dimensional complex Hilbert
space, let \(\mathfrak m\) be a two-sided ideal of
\(\mathcal L(\mathfrak H)\), and let \(\varphi\) be a trace on
\(\mathfrak m\) [that is, a linear form on \(\mathfrak m\), positive
on \(\mathfrak m^+\), such that
\(\varphi(ST)=\varphi(TS)\) for \(T\in\mathfrak m\) and
\(S\in\mathcal L(\mathfrak H)\)].
\begin{enumerate}[label=\textit{\alph*.},leftmargin=*]
\item If \(\mathfrak m=\mathcal L(\mathfrak H)\), then
      \(\varphi=0\). [Show that, if
      \(\mathfrak H_1,\mathfrak H_2\) are two complementary
      orthogonal subspaces of \(\mathfrak H\) having the same
      dimension, then
      \[
        \varphi(P_{\mathfrak H_1})
         =\varphi(P_{\mathfrak H_2})
         =\varphi(1)
         =\varphi(P_{\mathfrak H_1})
           +\varphi(P_{\mathfrak H_2}),
      \]
      and hence \(\varphi(1)=0\).]

\item Let \(\mathfrak f\) be the two-sided ideal of finite-rank
      operators. Suppose that
      \(\mathfrak m\supset\mathfrak f\) and
      \(\varphi(\mathfrak f)=0\). Let \(S,S'\) be compact operators
      of \(\mathfrak m^+\), and let
      \((\lambda_i(S))\) be the sequence of positive eigenvalues of
      \(S\), arranged in decreasing order, each eigenvalue being
      written a number of times equal to its multiplicity (the
      sequence being completed by \(0,0,0,\ldots\) if \(S\) has
      finite rank). If
      \[
        \lambda_i(S)=o\bigl(\lambda_i(S')\bigr)
        \qquad(i\to+\infty),
      \]
      then \(\varphi(S)=0\). [Reduce to the case in which \(S\) and
      \(S'\) have the same eigenvectors. Then, by adding to \(S\)
      and \(S'\) operators of \(\mathfrak f\), which changes neither
      \(\varphi(S)\) nor \(\varphi(S')\), prove that
      \[
        0\leq\varphi(S)\leq\varepsilon\varphi(S'),
      \]
      where \(\varepsilon>0\) is arbitrarily small.]

\item Let \(\psi\) be the trace on
      \(\mathcal L(\mathfrak H)^+\) defined in
      \(\cref{thm:I-6-5}\), and let \(\mathfrak n\) be the ideal of
      definition of \(\dot\psi\). Let \(\varphi\) be a trace on
      \(\mathfrak n\). Show that \(\varphi\) is proportional to
      \(\dot\psi\). [First show that \(\varphi\) is proportional to
      \(\dot\psi\) on \(\mathfrak f\), observing that \(\varphi\)
      takes the same value on all rank-one projections. Replacing
      \(\varphi\) by \(\varphi-\lambda\dot\psi\), reduce to the case
      \(\varphi(\mathfrak f)=0\). Then prove that
      \(\varphi(\mathfrak n^+)=0\), using \textit{b}.]

\item Let \(\mathfrak p\) be the set of compact operators
      \(S\geq0\) such that
      \[
        \lambda_1(S)+\cdots+\lambda_i(S)=o(\log i)
        \qquad(i\to+\infty).
      \]
      There exists a two-sided ideal \(\mathfrak m\) such that
      \(\mathfrak m\supset\mathfrak n\) and
      \(\mathfrak m^+=\mathfrak p\), and a nonzero trace
      \(\varphi\) on \(\mathfrak m\) such that
      \(\varphi(\mathfrak n)=0\). \ArticleRef{359}
\end{enumerate}

% Source PDF page 119 (printed p. 112)
\phantomsection\label{srcpage:112}
\item\label{ex:I-6-7}
Let \(\mathfrak H\) and \(\mathfrak K\) be complex Hilbert spaces.
The choice of an orthonormal basis \((e_i)_{i\in I}\) of
\(\mathfrak K\) canonically defines isometries \(U_i\) of
\(\mathfrak H\) onto subspaces of
\(\mathfrak H\otimes\mathfrak K\). Let \(\mathcal A\) be a von
Neumann algebra on \(\mathfrak H\), let \(\varphi\) be a normal trace
on \(\mathcal A^+\), and put
\[
  \mathcal B=\mathcal A\otimes\mathcal L(\mathfrak K).
\]
For \(T\in\mathcal B^+\), put
\[
  \psi(T)=\sum_{i\in I}\varphi(U_i^*TU_i).
\]
Show that \(\psi\) is a normal trace on \(\mathcal B^+\), faithful
(respectively, semifinite) if \(\varphi\) is faithful
(respectively, semifinite). (Argue as in \(\cref{thm:I-6-5}\).)

\item\label{ex:I-6-8}
Let \(\mathfrak H\) be a complex Hilbert space, let \(\varphi\) be
the trace of \hyperref[thm:I-6-5]{Theorem~5} on
\(\mathcal L(\mathfrak H)^+\), and let \(\mathfrak m\) be the ideal
of definition of \(\dot\varphi\).
\begin{enumerate}[label=\textit{\alph*.},leftmargin=*]
\item If \(T\in\mathfrak m\), then, for every orthonormal basis
      \((e_i)_{i\in I}\) of \(\mathfrak H\),
      \[
        \sum_{i\in I}|(Te_i\mid e_i)|<+\infty,
        \qquad
        \dot\varphi(T)=\sum_{i\in I}(Te_i\mid e_i).
      \]
      (Write \(T\) as a linear combination of elements of
      \(\mathfrak m^+\).)

\item If \(T\in\mathcal L(\mathfrak H)\) and
      \[
        \sum_{i\in I}|(Te_i\mid e_i)|<+\infty
      \]
      for every orthonormal basis \((e_i)_{i\in I}\) of
      \(\mathfrak H\), then \(T\in\mathfrak m\). [Writing
      \(T=T_1+iT_2\), with
      \[
        T_1=\tfrac12(T+T^*),
        \qquad
        T_2=\tfrac1{2i}(T-T^*),
      \]
      reduce to the case in which \(T\) is Hermitian. By choosing a
      suitable orthonormal basis, show that
      \(T^+\in\mathfrak m^+\) and \(T^-\in\mathfrak m^+\).]

\item Show that it is possible to have \(T\notin\mathfrak m\) and
      \[
        \sum_{i\in I}|(Te_i\mid e_i)|<+\infty
      \]
      for one orthonormal basis \((e_i)_{i\in I}\) of
      \(\mathfrak H\). [Let \(\mathfrak H_1,\mathfrak H_2\) be two
      complementary orthogonal subspaces of \(\mathfrak H\) having
      the same dimension; take \(T\) unitary such that
      \[
        T(\mathfrak H_1)=\mathfrak H_2,
        \qquad
        T(\mathfrak H_2)=\mathfrak H_1,
      \]
      and \(e_i\in\mathfrak H_1\cup\mathfrak H_2\) for every \(i\).]

\item If \(\mathfrak H\) is infinite-dimensional and if
      \(T\in\mathcal L(\mathfrak H)\) is such that
      \[
        \sum_{i,\lambda}|(Te_i\mid e_\lambda)|<+\infty
      \]
      for every orthonormal basis \((e_i)_{i\in I}\) of
      \(\mathfrak H\), show that \(T=0\). (First show that, for every
      \(x\in\mathfrak H\), \(Tx\) is proportional to \(x\).)

\item If \(T\in\mathcal L(\mathfrak H)\) is such that
      \[
        \sum_{i,\lambda}|(Te_i\mid e_\lambda)|<+\infty
      \]
      for one orthonormal basis \((e_i)_{i\in I}\) of
      \(\mathfrak H\), then \(T\in\mathfrak m\). [Let
      \((f_\lambda)_{\lambda\in L}\) be another orthonormal basis of
      \(\mathfrak H\). Putting
      \[
        f_\lambda=\sum_{i\in I}u_{\lambda,i}e_i,
      \]
      one has
      \[
        |(Tf_\lambda\mid f_\lambda)|
          \leq\sum_{i,\chi\in I}
            |u_{\lambda,i}|\,|u_{\lambda,\chi}|\,
            |(Te_i\mid e_\chi)|,
      \]
      and hence
      \[
        \sum_{\lambda\in L}|(Tf_\lambda\mid f_\lambda)|
        \leq
        \sum_{i,\chi\in I}
          \left[
            |(Te_i\mid e_\chi)|
            \left(\sum_{\lambda\in L}|u_{\lambda,i}|^2\right)^{1/2}
            \left(\sum_{\lambda\in L}|u_{\lambda,\chi}|^2\right)^{1/2}
          \right]
        <+\infty.
      \]
      Then use \textit{b}.]
\end{enumerate}

% Source PDF page 120 (printed p. 113)
\phantomsection\label{srcpage:113}
\noindent\textit{Problem.}\quad If \(T\in\mathfrak m\), does there
exist an orthonormal basis \((e_i)_{i\in I}\) of \(\mathfrak H\)
such that
\[
  \sum_{i,\lambda\in I}|(Te_i\mid e_\lambda)|<+\infty?
\]

\item\label{ex:I-6-9}
If a finite von Neumann algebra \(\mathcal A\) possesses a countable
totalizing set \(\mathfrak M\), then \(\mathcal A\) is of countable
type. [\(\mathfrak M\) separates the centre of
\(\mathcal A\); use
\(\hyperref[prop:I-6-9]{Proposition~9}\textup{(ii)}\).]

\item\label{ex:I-6-10}
Let \(\mathcal A\) and \(\mathcal B\) be two von Neumann algebras on
the Hilbert spaces \(\mathfrak H\) and \(\mathfrak K\).
\begin{enumerate}[label=\textit{\alph*.},leftmargin=*]
\item If \(\mathcal A\) is properly infinite, then
      \(\mathcal A\otimes\mathcal B\) is properly infinite. [Let
      \(E\) be a projection of the centre of
      \(\mathcal A\otimes\mathcal B\) such that
      \((\mathcal A\otimes\mathcal B)_E\) is finite. Then
      \[
        (\mathcal A\otimes\mathbb C_{\mathfrak K})_E
      \]
      is finite, and hence \(E=0\).]

\item If \(E\) (respectively, \(F\)) is the greatest projection of
      the centre of \(\mathcal A\) (respectively, of
      \(\mathcal B\)) such that \(\mathcal A_E\)
      (respectively, \(\mathcal B_F\)) is finite, then
      \(G=E\otimes F\) is the greatest projection of the centre of
      \(\mathcal A\otimes\mathcal B\) such that
      \((\mathcal A\otimes\mathcal B)_G\) is finite. [One has
      \begin{align*}
        \mathcal A\otimes\mathcal B
        ={}&(\mathcal A_E\otimes\mathcal B_F)
          \times(\mathcal A_E\otimes\mathcal B_{1-F})\\
         &{}\times(\mathcal A_{1-E}\otimes\mathcal B_F)
          \times(\mathcal A_{1-E}\otimes\mathcal B_{1-F});
      \end{align*}
      apply \textit{a} and \(\cref{prop:I-6-12}\).]

\item If \(\mathcal A\otimes\mathcal B\) is properly infinite, one
      of the algebras \(\mathcal A,\mathcal B\) is properly infinite
      (use \textit{b}).
\end{enumerate}

\item\label{ex:I-6-11}
Let \(\mathcal A\) and \(\mathcal B\) be two von Neumann algebras. If
\(\mathcal A\otimes\mathcal B\) is purely infinite, one of the
algebras \(\mathcal A,\mathcal B\) is purely infinite. [If there
exists a projection \(E\) (respectively, \(F\)) in the centre of
\(\mathcal A\) (respectively, \(\mathcal B\)) such that
\(\mathcal A_E\) (respectively, \(\mathcal B_F\)) is semifinite,
then \(\mathcal A_E\otimes\mathcal B_F\) is semifinite.]

The converse is true (Chapter~III, \S~8, Theorem~2).

\item\label{ex:I-6-12}
Let \(\mathfrak A_1,\mathfrak A_2\) be two Hilbert algebras, and let
\(E_1,E_2\) be their characteristic projections. The
characteristic projection of
\(\mathfrak A_1\otimes\mathfrak A_2\) is \(E_1\otimes E_2\). (Use
\(\cref{thm:I-6-6}\) and \hyperref[ex:I-6-10]{Exercise~10b}.)

\item\label{ex:I-6-13}
Let \(\mathcal A\) be a von Neumann algebra, let \(\varphi\) be a
finite trace on \(\mathcal A\) such that \(\varphi(1)=1\), and let
\(\Delta\) be the associated determinant.
\begin{enumerate}[label=\textit{\alph*.},leftmargin=*]
\item Show that \(\Delta\) is continuous for the norm topology.
\item Show that, if \(T_1,T_2\) are positive invertible Hermitian
      operators of \(\mathcal A\) such that \(T_1\leq T_2\), then
      \(\Delta(T_2)\geq\Delta(T_1)\). [Observe that
      \[
        T_1^{-1/2}T_2T_1^{-1/2}\geq1,
      \]
      and hence that
      \[
        \Delta(T_1^{-1/2}T_2T_1^{-1/2})\geq1.
      \]]
      \ArticleRef{23}, \ArticleRef{24}
\end{enumerate}
\end{enumerate}

\section{Abelian von Neumann Algebras}\label{sec:I-7}

\NumberedParagraph{1}{Basic Measures}\label{no:I-7-1}
Let \(\mathfrak H\) be a complex Hilbert space, and let
\(\mathcal Y\subset\mathcal L(\mathfrak H)\) be an abelian
\(C^*\)-algebra. Let \(Z\) be the spectrum of \(\mathcal Y\),
\(\mathcal C_\infty(Z)\) the set of complex-valued numerical
functions on \(Z\) which are continuous and tend to zero at
infinity, \(f\mapsto T_f\) the Gelfand isomorphism of
\(\mathcal C_\infty(Z)\) onto \(\mathcal Y\), and
\(\nu_{x,y}\) the spectral measure on \(Z\) defined by the pair
\((x,y)\) of elements of \(\mathfrak H\). (On all this, cf.
Appendix~I, nos.~1 and~2.)
% Source PDF page 121 (printed p. 114)
\phantomsection\label{srcpage:114}
\begin{definition}\label{def:I-7-1}
A positive measure \(\nu\) on \(Z\) is called \emph{basic} if it has
the following property: for a subset of \(Z\) to be locally
\(\nu\)-negligible, it is necessary and sufficient that it be locally
\(\nu_{x,x}\)-negligible for every \(x\in\mathfrak H\).
\end{definition}

(Throughout this book, when we speak of a measure on a locally compact
space, we shall mean a Radon measure; cf. \(\TreatiseRef{4}\).)

If there exists a basic measure \(\nu\) on \(Z\), every basic measure is
a positive measure equivalent to \(\nu\), and conversely.

If \(\nu\) is a basic measure, every measure \(\nu_{x,x}\) has a
density with respect to \(\nu\), and this density is
\(\nu\)-integrable. By
\hyperref[eq:App-I-6]{formula~(6) of Appendix~I}, for every measure
\(\nu_{x,y}\) there is a \(\nu\)-integrable function \(h_{x,y}\) on
\(Z\) such that \(\nu_{x,y}=h_{x,y}\nu\). The \(h_{x,y}\) have the
following properties, which follow from the properties of the
\(\nu_{x,y}\):
\begin{align}
 h_{\lambda x+\lambda'x',y}
   &=\lambda h_{x,y}+\lambda'h_{x',y}
     &&(\lambda,\lambda'\in\mathbb C), \tag{1}\label{eq:I-7-1}\\
 h_{y,x}&=\overline{h_{x,y}}, \tag{2}\label{eq:I-7-2}\\
 h_{x,x}&\geq0, \tag{3}\label{eq:I-7-3}\\
 h_{T_gx,T_{g'}y}
   &=g\overline{g'}\,h_{x,y}
     &&(g,g'\in\mathcal C_\infty(Z)), \tag{4}\label{eq:I-7-4}
\end{align}
except on \(\nu\)-negligible sets which depend on
\(x,x',y,\lambda,\lambda',g,g'\).

Since the union of the supports of the \(\nu_{x,x}\) is everywhere
dense in \(Z\), the support of a basic measure \(\nu\) is all of
\(Z\). Thus \(\mathcal C_\infty(Z)\), with its structures as an
involutive algebra and as a Banach space, is identified with a subspace
of \(L_{\mathbb C}^{\infty}(Z,\nu)\) (the space of complex-valued
\(\nu\)-measurable functions essentially bounded on \(Z\), in which
two functions are identified when they are equal locally almost
everywhere).

\begin{proposition}\label{prop:I-7-1}
There exists a homomorphism of the involutive algebra
\(L_{\mathbb C}^{\infty}(Z,\nu)\) onto an involutive subalgebra of
\(\mathcal L(\mathfrak H)\), extending the Gelfand isomorphism and
continuous when \(L_{\mathbb C}^{\infty}(Z,\nu)\) is endowed with its
weak topology as the dual of \(L_{\mathbb C}^{1}(Z,\nu)\), and
\(\mathcal L(\mathfrak H)\) with the weak operator topology. This
homomorphism is unique. It is isometric. Its image is the weak closure
of \(\mathcal Y\). If the extended isomorphism is still denoted by
\(f\mapsto T_f\), then, for
\(f\in L_{\mathbb C}^{\infty}(Z,\nu)\) and
\(x,y\in\mathfrak H\),
\[
 (T_fx\mid y)=\int_Z f(\zeta)\,d\nu_{x,y}(\zeta).
\]
\end{proposition}

\begin{proof}
The equality
\[
 (T_fx\mid y)
   =\int_Z f(\zeta)h_{x,y}(\zeta)\,d\nu(\zeta),
\]
valid for \(f\in\mathcal C_\infty(Z)\), shows that the Gelfand
isomorphism is continuous for the weak topologies occurring in the
statement. Now the unit ball of \(\mathcal L(\mathfrak H)\) is
% Source PDF page 122 (printed p. 115)
\phantomsection\label{srcpage:115}
weakly complete; on the other hand, every point of
\(L_{\mathbb C}^{\infty}(Z,\nu)\) is weakly adherent to a bounded
subset of \(\mathcal C_\infty(Z)\). [Indeed, if
\(f\in L_{\mathbb C}^{\infty}(Z,\nu)\) and \(\lVert f\rVert\leq1\),
then \(f\) is the weak limit of functions
\(g\in L_{\mathbb C}^{\infty}(Z,\nu)\) of compact support such that
\(\lVert g\rVert\leq1\); next, \(g\) is a mean limit of continuous
functions \(h\), supported in a fixed compact set, such that
\(\lVert h\rVert\leq1\); and \(g\) is, \emph{a fortiori}, the weak
limit of the \(h\).] Hence, by \(\TreatiseRef{3}\), Chapter~III,
\S~2, Proposition~8, the Gelfand isomorphism extends uniquely to a
weakly continuous linear map from
\(L_{\mathbb C}^{\infty}(Z,\nu)\) into
\(\mathcal L(\mathfrak H)\), whose image \(\mathcal A\) is contained
in the weak closure of \(\mathcal Y\). We continue to denote this map
by \(f\mapsto T_f\). The formulas
\[
 (T_fx\mid y)=\int_Z f(\zeta)h_{x,y}(\zeta)\,d\nu(\zeta),
 \qquad T_{\overline f}=T_f^*,
 \qquad T_{fg}=T_fT_g,
\]
which are true for \(f,g\in\mathcal C_\infty(Z)\), remain true by
continuity for \(f,g\in L_{\mathbb C}^{\infty}(Z,\nu)\) [for the
last formula, one must let \(f\), and then \(g\), converge separately
to elements of \(L_{\mathbb C}^{\infty}(Z,\nu)\)].

Let us prove that \(\lVert T_f\rVert=\lVert f\rVert\). First,
\(T_{g\overline g}=T_gT_g^*\), and therefore the homomorphism
\(f\mapsto T_f\) preserves the natural order relations. The constant
function \(1\) corresponds to a self-adjoint idempotent operator,
that is, to a projection \(E\). If \(0\leq f\leq\lambda\), then
\(0\leq T_f\leq\lambda E\), and consequently
\(\lVert T_f\rVert\leq\lambda\); this proves that
\(\lVert T_f\rVert\leq\lVert f\rVert\) for \(f\geq0\). Suppose now
that the function \(f\geq0\) majorizes a number \(\mu>0\) on a
measurable set \(X\) which is not locally \(\nu\)-negligible. Then
\(T_f\) majorizes \(\mu E'\), where \(E'\) is the projection
corresponding to the characteristic function of \(X\). Now \(X\) is
not negligible for at least one spectral measure \(\nu_{x,x}\), and
then
\[
 (E'x\mid x)=\int_Xd\nu_{x,x}(\zeta)\neq0,
 \qquad\text{so that}\qquad E'\neq0;
\]
this proves that \(\lVert T_f\rVert\geq\mu\). Thus
\(\lVert T_f\rVert=\lVert f\rVert\) for \(f\geq0\), and consequently
for arbitrary \(f\) (by considering \(f\overline f\)).

The unit ball \(\mathcal A_1\) of \(\mathcal A\) is therefore the
image of the unit ball of \(L_{\mathbb C}^{\infty}(Z,\nu)\). The
latter is weakly compact. Hence \(\mathcal A_1\) is weakly compact,
and consequently \(\mathcal A\) is weakly closed
(\hyperref[thm:I-3-2]{\S~3, Theorem~2}).
\end{proof}

This extension of the Gelfand isomorphism may be compared with the
extension defined in Appendix~I, no.~3. See also Exercise~2 of
Appendix~I.

\noindent\textit{Bibliography.}\quad
\ArticleRef{9}, \ArticleRef{18}, \ArticleRef{28}, \ArticleRef{45},
\ArticleRef{70}, \ArticleRef{100}, \TreatiseRef{9}.

\NumberedParagraph{2}{Existence of Basic Measures}
\label{no:I-7-2}

\begin{proposition}\label{prop:I-7-2}
If \(x\in\mathfrak H\) is totalizing for \(\mathcal Y'\), then
\(\nu_{x,x}\) is a basic measure on \(Z\).
\end{proposition}

\begin{proof}
Let \(y\in\mathfrak H\). We must show that \(\nu_{y,y}\) is based on
\(\nu_{x,x}\). First, if \(y=T'x\), with \(T'\in\mathcal Y'\), then,
% Source PDF page 123 (printed p. 116)
\phantomsection\label{srcpage:116}
for every positive function \(f\in\mathcal C_\infty(Z)\),
\[
 \nu_{y,y}(f)
   =(T_fT'x\mid T'x)
   =\lVert T'T_f^{1/2}x\rVert^2
   \leq\lVert T'\rVert^2(T_fx\mid x)
   =\lVert T'\rVert^2\nu_{x,x}(f);
\]
hence \(\nu_{y,y}\leq\lVert T'\rVert^2\nu_{x,x}\), which proves our
assertion in this case. Let us pass to the general case. We shall use
criterion~5 of \(\TreatiseRef{4}\), Chapter~V, \S~5, Corollary~5 to
Theorem~2. Fix a function \(f\geq0\) in
\(\mathcal C_\infty(Z)\), and \(\varepsilon>0\). We may first find
\(y'\in\mathcal Y'x\) such that
\[
 \lVert y-y'\rVert
   \leq\frac{\varepsilon}{4}\lVert f\rVert^{-1}\lVert y\rVert^{-1},
 \qquad
 \lVert y'\rVert\leq\lVert y\rVert.
\]
Then, for every \(h\geq0\) in \(\mathcal C_\infty(Z)\) majorized by
\(f\), one has
\begin{align*}
 \bigl|\nu_{y,y}(h)-\nu_{y',y'}(h)\bigr|
 &\leq\bigl|(T_hy\mid y)-(T_hy\mid y')\bigr|
      +\bigl|(T_hy\mid y')-(T_hy'\mid y')\bigr| \\
 &\leq\lVert T_h\rVert\,\lVert y\rVert\,\lVert y-y'\rVert
      +\lVert T_h\rVert\,\lVert y-y'\rVert\,\lVert y'\rVert \\
 &\leq2\lVert f\rVert\,\lVert y\rVert\,\lVert y-y'\rVert
 \leq\frac{\varepsilon}{2}.
\end{align*}
Next, by the first part of the proof, there exists \(\delta>0\) such
that \(\nu_{x,x}(h)\leq\delta\) implies
\(\nu_{y',y'}(h)\leq\varepsilon/2\). Hence, for
\(h\in\mathcal C_\infty(Z)\), the conditions
\(0\leq h\leq f\) and \(\nu_{x,x}(h)\leq\delta\) imply
\(\nu_{y,y}(h)\leq\varepsilon\).
\end{proof}

\begin{proposition}\label{prop:I-7-3}
For there to exist a bounded basic measure on \(Z\), it is necessary
and sufficient that there exist an element totalizing for
\(\mathcal Y'\).
\end{proposition}

\begin{proof}
The condition is sufficient by \cref{prop:I-7-2}. Conversely, suppose
that there exists a bounded basic measure \(\nu\) on \(Z\). We shall
show that \(\mathcal Y''\) is of countable type; together with
\hyperref[prop:I-2-3]{\S~2, the corollary to Proposition~3}, this will
prove the proposition.

We may restrict ourselves to the case in which \(\mathcal Y''\) is the
weak closure of \(\mathcal Y\). Let \((E_i)_{i\in I}\) be a family of
nonzero pairwise orthogonal projections of \(\mathcal Y''\). Let
\(Z_i\) be a \(\nu\)-measurable subset of \(Z\) whose characteristic
function corresponds to \(E_i\) under the isomorphism of
\cref{prop:I-7-1}. The intersection of any two distinct \(Z_i\) is
negligible; hence, for every finite subset \(J\) of \(I\),
\[
 \sum_{i\in J}\nu(Z_i)\leq\nu(Z)<+\infty.
\]
Thus, for every positive integer \(n\), the inequality
\(\nu(Z_i)\geq1/n\) holds for only finitely many indices \(i\); on the
other hand, \(\nu(Z_i)>0\) for every \(i\), since \(E_i\neq0\).
Therefore \(I\) is countable.
\end{proof}

\begin{proposition}\label{prop:I-7-4}
Let \(\mathcal A\) be an abelian von Neumann algebra on
\(\mathfrak H\).
\begin{enumerate}[label=\textup{(\roman*)},leftmargin=*]
\item There exists a \(C^*\)-subalgebra \(\mathcal Y\) of
\(\mathcal A\), weakly dense in \(\mathcal A\), whose spectrum \(Z\)
supports a basic measure.
\item If \(\mathfrak H\) has a countable basis, \(\mathcal Y\) may be
chosen so that, in addition, \(Z\) is compact and has a countable
basis.

% Source PDF page 124 (printed p. 117)
\phantomsection\label{srcpage:117}
\item Suppose that \(\mathcal A\) is of countable type (as is the case
if \(\mathfrak H\) has a countable basis). For every
\(C^*\)-subalgebra \(\mathcal Y\) of \(\mathcal A\), the spectrum of
\(\mathcal Y\) supports a bounded basic measure.
\end{enumerate}
\end{proposition}

\begin{proof}
If \(\mathcal A\) is of countable type, there exists an element \(x\)
totalizing for \(\mathcal A'\). Then, whatever the
\(C^*\)-subalgebra \(\mathcal Y\) of \(\mathcal A\), one has
\(\mathcal Y'\supset\mathcal A'\), and therefore \(x\) is totalizing
for \(\mathcal Y'\). Together with \cref{prop:I-7-2}, this proves
\textup{(iii)}.

If \(\mathfrak H\) has a countable basis, \(\mathcal A\) is the von
Neumann algebra generated by a sequence \((T_i)\). Taking for
\(\mathcal Y\) the \(C^*\)-subalgebra of \(\mathcal A\) generated by
\(1\) and the \(T_i\), the spectrum of \(\mathcal Y\) is compact and
has a countable basis (Appendix~I, no.~1), and supports a basic measure
by \textup{(iii)}. This proves \textup{(ii)}.

Finally consider the general case. There exist
\begin{enumerate}[label=\arabic*\textsuperscript{\(\circ\)}.,leftmargin=*]
\item a family \((\mathfrak H_i)_{i\in I}\) of Hilbert spaces;
\item on each \(\mathfrak H_i\), a von Neumann algebra
\(\mathcal A_i\) of countable type,
\end{enumerate}
such that \(\mathfrak H\) is identified with the Hilbert direct sum of
the \(\mathfrak H_i\), and \(\mathcal A\) with the product of the
\(\mathcal A_i\) [\hyperref[prop:I-6-9]{\S~6, Proposition~9,
\textup{(iii)}}]. Let \(E_i=P_{\mathfrak H_i}\), and let \(Z'\) be the
spectrum of \(\mathcal A\). The projection \(E_i\) corresponds to a
continuous function on \(Z'\) equal to its square, that is, to the
characteristic function of a set \(Z_i\) which is both open and closed.
Since \(E_iE_k=0\) for \(i\neq k\), the \(Z_i\) are pairwise disjoint.
Consider the \(C^*\)-subalgebra \(\mathcal Y\) of \(\mathcal A\)
formed by those \(T=(T_i)_{i\in I}\) having the following property:
for every \(\varepsilon>0\), the set of \(i\in I\) such that
\(\lVert T_i\rVert\geq\varepsilon\) is finite. The
\(C^*\)-algebra \(\mathcal Y\) is weakly dense in \(\mathcal A\).
Every nonzero character of \(\mathcal A\) situated in one of the
\(Z_i\) defines, by restriction to \(\mathcal Y\), a nonzero character
of \(\mathcal Y\). Hence there is a continuous map \(\theta\) from
\(\bigcup_{i\in I}Z_i\) into the spectrum \(Z\) of \(\mathcal Y\).
Conversely, let \(\chi\) be a nonzero character of \(\mathcal Y\); one
has
\[
 \chi(E_i)\chi(E_k)=\chi(E_iE_k)=0
 \qquad(i\neq k).
\]
Thus \(\chi(E_i)=0\) for every \(i\) in a subset \(J\) of \(I\) whose
complement has at most one element. If \(J=I\), norm continuity would
give \(\chi=0\); hence \(I\setminus J=\{i_0\}\) and
\(\chi(E_{i_0})=1\). For every \(T\in\mathcal A\), one has
\(TE_{i_0}\in\mathcal Y\); put
\[
 \chi'(T)=\chi(TE_{i_0}).
\]
We see at once that the character \(\chi\) of \(\mathcal Y\) is thus
extended to a nonzero character \(\chi'\) of \(\mathcal A\) belonging
to \(Z_{i_0}\). Hence there is a continuous map \(\theta'\) from
\(Z\) into \(\bigcup_{i\in I}Z_i\). It is easily checked that
\(\theta'\circ\theta\) and \(\theta\circ\theta'\) are the identity
maps on
% Source PDF page 125 (printed p. 118)
\phantomsection\label{srcpage:118}
\(\bigcup_{i\in I}Z_i\) and \(Z\), respectively. In short, \(Z\) is
identified (with its topology) with the union of the \(Z_i\). Likewise,
\(Z_i\) is identified with the spectrum of \(\mathcal A_i\). By
\textup{(iii)}, there is a basic measure \(\nu_i\) on \(Z_i\). The
\(\nu_i\) define on \(Z\) a positive measure \(\nu\) whose restriction
to each \(Z_i\) is \(\nu_i\).

On the other hand, if
\(x=\sum_{i\in I}x_i\in\mathfrak H\), with
\(x_i\in\mathfrak H_i\), and \(T=(T_i)_{i\in I}\in\mathcal Y\), then
\[
 (Tx\mid x)=\sum_{i\in I}(T_ix_i\mid x_i),
\]
and therefore \(\nu_{x,x}\) induces on \(Z_i\) the measure
\(\nu_{x_i,x_i}\). This being so, for a subset \(Y\) of \(Z\) to be
locally negligible for every \(\nu_{x,x}\), it is necessary and
sufficient that each \(Y\cap Z_i\) be negligible for every
\(\nu_{x_i,x_i}\), that is, that each \(Y\cap Z_i\) be negligible for
\(\nu_i\), that is, that \(Y\) be locally negligible for \(\nu\).
Thus \(\nu\) is a basic measure on \(Z\).
\end{proof}

It may happen that the spectrum \(Z\) of an abelian
\(C^*\)-algebra of operators \(\mathcal Y\) supports no basic measure,
even when \(Z\) is compact and has a countable basis (Exercise~3e).

\noindent\textit{Bibliography.}\quad
\ArticleRef{9}, \ArticleRef{18}, \ArticleRef{45}, \ArticleRef{71},
\ArticleRef{89}, \ArticleRef{100}, \ArticleRef{101}.

\NumberedParagraph{3}{Structure of Abelian von Neumann Algebras}
\label{no:I-7-3}

\begin{theorem}\label{thm:I-7-1}
Let \(\mathfrak H\) be a complex Hilbert space and \(\mathcal A\) an
abelian von Neumann algebra on \(\mathfrak H\). There exist a locally
compact space \(Z\), a positive measure \(\nu\) on \(Z\) with support
\(Z\), and an isometric isomorphism of the normed involutive algebra
\(\mathcal A\) onto the normed involutive algebra
\(L_{\mathbb C}^{\infty}(Z,\nu)\). If \(\mathfrak H\) has a countable
basis, one may require \(Z\) to be compact and to have a countable
basis.
\end{theorem}

\begin{proof}
Let \(\mathcal Y\) be a \(C^*\)-subalgebra of \(\mathcal A\),
weakly dense in \(\mathcal A\), whose spectrum \(Z\) supports a
basic measure \(\nu\) [\cref{prop:I-7-4}\textup{(i)}]. By
\cref{prop:I-7-1}, there is an isometric isomorphism of
\(\mathcal A\) onto \(L_{\mathbb C}^{\infty}(Z,\nu)\). If
\(\mathfrak H\) has a countable basis, one may suppose that \(Z\)
is compact and has a countable basis
[\cref{prop:I-7-4}\textup{(ii)}].
\end{proof}

Conversely:

\begin{theorem}\label{thm:I-7-2}
Let \(Z\) be a locally compact space, let \(\nu\) be a positive
measure on \(Z\), and put
\(\mathfrak H=L_{\mathbb C}^{2}(Z,\nu)\). For
\(f\in L_{\mathbb C}^{\infty}(Z,\nu)\), let
\(T_f\in\mathcal L(\mathfrak H)\) be the operator defined by
\(T_fg=fg\) \((g\in\mathfrak H)\), and let \(\mathcal A\) be the
set of the \(T_f\).
\begin{enumerate}[label=\textup{(\roman*)},leftmargin=*]
\item \(\mathcal A\) is a von Neumann algebra and
      \(\mathcal A=\mathcal A'\).
\item The map \(f\mapsto T_f\) is an isometric isomorphism of the
      normed involutive algebra
      \(L_{\mathbb C}^{\infty}(Z,\nu)\) onto the normed involutive
      algebra \(\mathcal A\).
\end{enumerate}
\end{theorem}

% Source PDF page 126 (printed p. 119)
\phantomsection\label{srcpage:119}
\begin{proof}
It is clear that the map \(f\mapsto T_f\) is a homomorphism, for the
involutive-algebra structure, of
\(L_{\mathbb C}^{\infty}(Z,\nu)\) onto \(\mathcal A\), and that
\(\lVert T_f\rVert\leq\lVert f\rVert\). If
\(\lvert f(\zeta)\rvert\geq\lambda\) on a non-\(\nu\)-negligible
integrable set \(Y\), then, writing \(\chi_Y\) for the characteristic
function of \(Y\),
\[
 \lvert T_f\chi_Y\rvert\geq\lambda\chi_Y,
\]
and hence \(\lVert T_f\rVert\geq\lambda\). Thus
\(\lVert T_f\rVert=\lVert f\rVert\). For \(g,h\in\mathfrak H\) and
\(f\in L_{\mathbb C}^{\infty}(Z,\nu)\), one has
\begin{equation}
 (T_fg\mid h)
   =\int_Z f(\zeta)g(\zeta)\overline{h(\zeta)}\,d\nu(\zeta),
 \tag{5}\label{eq:I-7-5}
\end{equation}
and \(g\overline h\in L_{\mathbb C}^{1}(Z,\nu)\). Hence the map
\(f\mapsto T_f\) is weakly continuous (for the weak topologies already
considered). The unit ball of \(\mathcal A\), being the image of the
unit ball of \(L_{\mathbb C}^{\infty}(Z,\nu)\), is weakly compact;
therefore \(\mathcal A\) is an abelian von Neumann algebra. Put
\[
 \mathfrak A
   =\mathfrak H\cap L_{\mathbb C}^{\infty}(Z,\nu).
\]
It is immediate that \(\mathfrak A\) is an abelian Hilbert algebra
dense in \(\mathfrak H\). Hence
\(\mathcal U(\mathfrak A)\subset
\mathcal U(\mathfrak A)'=\mathcal V(\mathfrak A)\)
(\hyperref[thm:I-5-1]{\S~5, Theorem~1}); by symmetry,
\[
 \mathcal U(\mathfrak A)=\mathcal V(\mathfrak A)
  =\mathcal U(\mathfrak A)'=\mathcal V(\mathfrak A)'.
\]
For \(f\in\mathfrak A\),
\(U_f=V_f=T_f\in\mathcal A\), and therefore
\[
 \mathcal U(\mathfrak A)\subset\mathcal A\subset\mathcal A',
 \qquad\text{hence}\qquad
 \mathcal A\subset\mathcal U(\mathfrak A)'=\mathcal U(\mathfrak A),
\]
and finally \(\mathcal A=\mathcal U(\mathfrak A)=\mathcal A'\).
\end{proof}

\noindent\textbf{Remark.}\quad
If \(\nu\) has support \(Z\), \(\mathcal C_\infty(Z)\), with its norm,
is identified with an involutive subalgebra of
\(L_{\mathbb C}^{\infty}(Z,\nu)\). The image \(\mathcal Y\) of
\(\mathcal C_\infty(Z)\) under the isomorphism \(f\mapsto T_f\) is a
\(C^*\)-subalgebra of \(\mathcal A\). The spectrum of \(\mathcal Y\)
is identified with \(Z\), in such a way that the Gelfand isomorphism is
identified with the isomorphism \(f\mapsto T_f\) of
\(\mathcal C_\infty(Z)\) onto \(\mathcal Y\). If
\(g,h\in\mathfrak H\), it follows from \eqref{eq:I-7-5} that
\(\nu_{g,h}\) is the measure having density
\(g\overline h\) with respect to \(\nu\). Now, as \(g\) and \(h\)
range over \(L_{\mathbb C}^{2}(Z,\nu)\), \(g\overline h\) ranges over
all of \(L_{\mathbb C}^{1}(Z,\nu)\); hence \(\nu\) is a basic measure
on \(Z\). The map \(f\mapsto T_f\) of
\(L_{\mathbb C}^{\infty}(Z,\nu)\) onto \(\mathcal A\) is the weakly
continuous extension of the Gelfand isomorphism defined by
\cref{prop:I-7-1}.

In \cref{thm:I-7-1}, the datum of \(\mathcal A\) does not determine
\(Z\) and \(\nu\) uniquely. See, however, Appendix~IV.

\noindent\textit{Bibliography.}\quad
\ArticleRef{9}, \ArticleRef{45}, \ArticleRef{71}, \ArticleRef{89},
\ArticleRef{100}, \ArticleRef{101}.

\noindent\textit{Exercises.}---
\begin{enumerate}[label=\arabic*.,leftmargin=*,itemsep=0.5\baselineskip]
\item\label{ex:I-7-1}
Let \(\mathfrak H\) be a complex Hilbert space, let \(\mathcal Y\) be
a commutative \(C^*\)-algebra of operators on \(\mathfrak H\), let
\(Z\) be the spectrum of \(\mathcal Y\), and let \(\nu\) be a positive
measure on \(Z\). For the Gelfand isomorphism to be continuous
[\(\mathcal C_\infty(Z)\) being endowed with the topology induced by
the weak topology of \(L_{\mathbb C}^{\infty}(Z,\nu)\), and
\(\mathcal L(\mathfrak H)\) with the weak operator topology], it is
necessary and sufficient that every measure \(\nu_{x,x}\) be based on
\(\nu\). If this is so, the Gelfand isomorphism admits a weakly
continuous extension to \(L_{\mathbb C}^{\infty}(Z,\nu)\); and for
\(\nu\) to be basic, it is necessary and sufficient that this
extension be injective.

% Source PDF page 127 (printed p. 120)
\phantomsection\label{srcpage:120}
\item\label{ex:I-7-2}
Let \(Z\) be a locally compact space, \(\nu\) a positive measure on
\(Z\), and let \(f\mapsto T_f\) be an isomorphism of the involutive
algebra \(L_{\mathbb C}^{\infty}(Z,\nu)\) onto an abelian von Neumann
algebra \(\mathcal A\). Let \(X\) be a measurable subset of \(Z\), and
let \(E\) be the projection of \(\mathcal A\) corresponding to the
characteristic function of \(X\). For \(\mathcal A_E\) to be of
countable type, it is necessary and sufficient that \(X\), up to a
locally negligible set, be a countable union of integrable sets. (For
sufficiency, argue as in \cref{prop:I-7-3}. For necessity, consider
families of nonzero pairwise orthogonal projections of \(\mathcal A\),
majorized by \(E\), which correspond to integrable subsets of \(Z\),
and apply Zorn's theorem to them.)

\item\label{ex:I-7-3}
\begin{enumerate}[label=\textit{\alph*.},leftmargin=2em]
\item The property of a von Neumann algebra of possessing a countable
family of generators is invariant under isomorphism (use
\hyperref[cor:I-4-2]{\S~4, Theorem~2, Corollary~2}).

\item If a von Neumann algebra \(\mathcal A\) is isomorphic to a von
Neumann algebra acting on a Hilbert space with a countable basis, then
\(\mathcal A\) is of countable type and is generated by a countable
family of elements.

\item Let \(\mathcal A\) be a von Neumann algebra and \(\mathcal Z\)
its centre. If \(\mathcal Z\) is of countable type, and if
\(\mathcal A\) is generated by a countable family of elements, then
\(\mathcal A\) is isomorphic to a von Neumann algebra acting on a
Hilbert space with a countable basis. (Let \(x\) be a separating
element for \(\mathcal Z\); \(\mathcal A\) is isomorphic to the algebra
induced by \(\mathcal A\) on \(\mathfrak X_x^{\mathcal A}\); show that
the Hilbert space \(\mathfrak X_x^{\mathcal A}\) has a countable
basis.)

\item Adopt the notation of \cref{thm:I-7-2}, and suppose \(\nu\)
bounded. Show that the constant function equal to \(1\) on \(Z\) is a
separating and totalizing element for \(\mathcal A\). If, moreover,
there exists a countable family \(\mathcal M\) of elements of
\(\mathcal L(\mathfrak H)\) generating the von Neumann algebra
\(\mathcal A\), then \(\mathfrak H\) has a countable basis. (The
transforms of the function \(1\) by the operators of \(\mathcal M\)
form a total set in \(\mathfrak H\).) Deduce an example of an abelian
von Neumann algebra of countable type which is not generated by a
countable family of elements. (Take for \(Z\) a suitable product of
spaces each identical with the interval \([0,1]\) of the real line,
and for \(\nu\) the product of the Lebesgue measures, in such a way
that \(L_{\mathbb C}^{2}(Z,\nu)\) does not have a countable basis.)

\item Adopt the notation of \cref{thm:I-7-2}; take for \(Z\) the
interval \([0,1]\) with the discrete topology, and for \(\nu\) the
measure defined by placing mass \(+1\) at every point. Show that
\(\mathcal A\) is not of countable type, but is generated by one
element (for example, by the operator \(T_f\), where \(f\) is the
function \(\zeta\mapsto\zeta\)). Let \(\mathcal Y\) be the
\(C^*\)-subalgebra of \(\mathcal A\) generated by this operator
\(T_f\) and by \(1\). Show that the spectrum of \(\mathcal Y\) is
compact and has a countable basis, but supports no basic measure.

\item Let \(\mathcal A\) be an abelian von Neumann algebra on a Hilbert
space with a countable basis. Show that \(\mathcal A\) is generated by
a single Hermitian element. [Form a sequence
\((E_1,E_2,\ldots)\) of projections generating \(\mathcal A\). Deduce
from it a family \((F_r)\) of projections generating \(\mathcal A\),
where \(r\) ranges over the rational numbers in \([0,1]\), with the
following properties:
\[
 F_0=0,\qquad F_1=1,\qquad
 r\leq r'\ \Longrightarrow\ F_r\leq F_{r'},\qquad
 \lim_{\substack{r\to r_0\\r>r_0}}F_r=F_{r_0}.
\]
Construct a Hermitian operator \(T\) such that \(F_r\) is the greatest
spectral projection for which \(F_rT\leq rF_r\). Show that \(T\)
generates \(\mathcal A\).] \ArticleRef{73}

\item Let \(\mathfrak H\) be a Hilbert space with a countable basis,
let \(\mathcal A\) and \(\mathcal B\) be abelian von Neumann algebras
on \(\mathfrak H\), and let \(\mathcal C\) be the von Neumann algebra
generated by \(\mathcal A\) and \(\mathcal B\). Then \(\mathcal C\)
is generated by one element. (Use \textit{f}.)

\item Let \((\mathcal A_i)\) be a family of pairwise commuting, weakly
closed involutive algebras on \(\mathfrak H\), and let \(\mathcal A\)
be the von Neumann algebra generated by the \(\mathcal A_i\). If each
\(\mathcal A_i\) is generated by one element, then \(\mathcal A\) is
generated by one element. (Use \textit{f} and \textit{g}.)
\ArticleRef{313}

% Source PDF page 128 (printed p. 121)
\phantomsection\label{srcpage:121}
\item Let \(\mathcal A\) be a von Neumann algebra on \(\mathfrak H\),
the product of a family \((\mathcal A_i)\) of von Neumann algebras. For
\(\mathcal A\) to be generated by one element, it is necessary and
sufficient that every \(\mathcal A_i\) be generated by one element.
(Use \textit{h}.)
\end{enumerate}

It is not known whether every von Neumann algebra on a Hilbert space
with a countable basis can be generated by one element, or even by a
finite number of elements. (See, however, Chapter~III, \S~3,
Exercise~8, and \S~8, Exercise~12.)

\item\label{ex:I-7-4}
Adopt the notation of \cref{thm:I-7-2}. Every
\(h\in L_{\mathbb C}^{1}(Z,\nu)\) defines a linear form
\(\varphi_h\) on \(\mathcal A\) by
\[
 \varphi_h(T_f)=\int_Z f(\zeta)h(\zeta)\,d\nu(\zeta).
\]
Show that the forms \(\varphi_h\) are precisely the forms
\(\omega_{x,y}\) on \(\mathcal A\), and also precisely the ultraweakly
continuous linear forms on \(\mathcal A\). (Use
\hyperref[thm:I-3-1]{\S~3, Theorem~1}.) Compare
\hyperref[thm:I-6-8]{\S~6, Theorem~8}, Chapter~III, \S~1, the
corollary to Theorem~4, and \S~6, no.~3, Remark~1.
\ArticleRef{19}, \ArticleRef{89}, \ArticleRef{100}, \ArticleRef{109}

\item\label{ex:I-7-5}
Let \(\mathcal A\) be an abelian von Neumann algebra. There exist a
locally compact space \(Z\) and a positive measure \(\nu\) on \(Z\)
with support \(Z\), having the following properties:
\begin{enumerate}[label=\textit{\alph*.},leftmargin=2em]
\item the normed involutive algebra \(\mathcal A\) is isomorphic to the
normed involutive algebra \(L_{\mathbb C}^{\infty}(Z,\nu)\);
\item every function in \(L_{\mathbb C}^{\infty}(Z,\nu)\) agrees
locally almost everywhere with a continuous function.
\end{enumerate}
[Show that the space \(Z\) and the measure \(\nu\) constructed in the
proof of \cref{prop:I-7-4}(i) have property~\textit{b}.]
\ArticleRef{9}, \ArticleRef{45}, \ArticleRef{100}, \ArticleRef{110}

\item\label{ex:I-7-6}
Let \(\mathfrak A\) be a completed abelian Hilbert algebra. There exist
a locally compact space \(Z\) and a positive measure \(\nu\) on \(Z\)
with support \(Z\), such that \(\mathfrak A\) is isomorphic to the
Hilbert algebra
\[
 L_{\mathbb C}^{2}(Z,\nu)\cap L_{\mathbb C}^{\infty}(Z,\nu).
\]
(Use \cref{thm:I-7-1}, and \S~6, Theorems~1 and~3.)
\end{enumerate}

\section{Discrete von Neumann Algebras}\label{sec:I-8}

\NumberedParagraph{1}{Second Classification of von Neumann Algebras}
\label{no:I-8-1}

\begin{definition}\label{def:I-8-1}
Let \(\mathcal A\) be a von Neumann algebra and \(\mathcal Z\) its
centre. The algebra \(\mathcal A\) is called \emph{discrete} if it is
isomorphic to a von Neumann algebra whose commutant is abelian. It is
called \emph{continuous} if there is no nonzero projection
\(E\in\mathcal Z\) such that the von Neumann algebra \(\mathcal A_E\)
is discrete.
\end{definition}

The preceding classification is invariant under isomorphisms. It is
also invariant under anti-isomorphisms.

A factor is either discrete or continuous.

A discrete factor \(\mathcal A\) is isomorphic to the algebra of all
operators on a suitable Hilbert space. Indeed, after transforming
\(\mathcal A\) by a suitably chosen isomorphism, one may suppose that
\(\mathcal A'\), which, like \(\mathcal A\), is a factor, is abelian.
Then \(\mathcal A'\) consists only of the scalars, which proves our
assertion.

\begin{proposition}\label{prop:I-8-1}
Let \((\mathcal A_i)_{i\in I}\) be a family of von Neumann algebras,
and let \(\mathcal A\) be their product von Neumann algebra. For
\(\mathcal A\) to be discrete (respectively, continuous), it is
necessary and sufficient that every \(\mathcal A_i\) be discrete
(respectively, continuous).
\end{proposition}

% Source PDF page 129 (printed p. 122)
\phantomsection\label{srcpage:122}
\begin{proof}
If the \(\mathcal A_i'\) are abelian, then
\[
 \mathcal A'=\prod_{i\in I}\mathcal A_i'
\]
is abelian. On the other hand, replacing the \(\mathcal A_i\) by
isomorphic von Neumann algebras replaces \(\mathcal A\) by an
isomorphic von Neumann algebra. This proves that if the
\(\mathcal A_i\) are discrete, then \(\mathcal A\) is discrete.
Conversely, if \(\mathcal A\) is discrete, an isomorphism of
\(\mathcal A\) onto a von Neumann algebra \(\mathcal B\) such that
\(\mathcal B'\) is abelian defines an isomorphism of each
\(\mathcal A_i\) onto a von Neumann algebra \(\mathcal B_i\) such that
\(\mathcal B_i'\) is abelian; hence the \(\mathcal A_i\) are discrete.

If some algebra \(\mathcal A_i\) is not continuous, it is the product
of a nonzero discrete von Neumann algebra and another von Neumann
algebra; hence \(\mathcal A\) is a product of von Neumann algebras one
of which is nonzero and discrete. Consequently \(\mathcal A\) is not
continuous. Conversely, if \(\mathcal A\) is not continuous, there is
a nonzero projection \(E\) of its centre such that \(\mathcal A_E\) is
discrete. Write \(E=(E_i)_{i\in I}\), where \(E_i\) is a projection of
the centre of \(\mathcal A_i\). Then
\[
 \mathcal A_E=\prod_{i\in I}(\mathcal A_i)_{E_i};
\]
one of the \((\mathcal A_i)_{E_i}\) is nonzero and discrete, and
therefore one of the \(\mathcal A_i\) is not continuous.
\end{proof}

Just as Proposition~8 and its Corollary~1 were deduced from
\hyperref[prop:I-6-7]{\S~6, Proposition~7}, the following results are
deduced from \cref{prop:I-8-1}:

\setcounter{corollary}{0}
\begin{corollary}\label{cor:I-8-1-1}
Let \(\mathcal A\) be a von Neumann algebra and \(\mathcal Z\) its
centre. Among the projections of \(\mathcal Z\), there is a greatest
projection \(E\) (respectively, \(F\)) such that \(\mathcal A_E\)
(respectively, \(\mathcal A_F\)) is discrete (respectively,
continuous). One has \(EF=0\) and \(E+F=1\).
\end{corollary}

\begin{corollary}\label{cor:I-8-1-2}
Every von Neumann algebra is canonically isomorphic to the product of
a discrete von Neumann algebra and a continuous von Neumann algebra.
\end{corollary}

\begin{proposition}\label{prop:I-8-2}
Every discrete von Neumann algebra is semifinite.
\end{proposition}

\begin{proof}
Let \(\mathcal A\) be a discrete von Neumann algebra. Transforming
\(\mathcal A\) by an isomorphism, we may suppose \(\mathcal A'\)
abelian. Then \(\mathcal A'\) is finite, and therefore \(\mathcal A\)
is semifinite (\hyperref[prop:I-6-13]{\S~6, Proposition~13,
Corollary~1}).
\end{proof}

\setcounter{corollary}{0}
\begin{corollary}\label{cor:I-8-2-1}
Let \(\mathcal A\) be a von Neumann algebra and \(\mathcal Z\) its
centre. Let \(E\) (respectively, \(F\)) be the greatest projection of
\(\mathcal Z\) such that \(\mathcal A_E\) (respectively,
\(\mathcal A_F\)) is semifinite (respectively, discrete). Then
\(E\geq F\).
\end{corollary}

\begin{corollary}\label{cor:I-8-2-2}
Every purely infinite von Neumann algebra is continuous.
\end{corollary}

\noindent\textbf{Remark.}\quad
Since a discrete factor is isomorphic to an algebra
\(\mathcal L(\mathfrak H)\), we already knew very elementarily
(\hyperref[thm:I-6-5]{\S~6, Theorem~5}) that a discrete factor is
semifinite.

% Source PDF page 130 (printed p. 123)
\phantomsection\label{srcpage:123}
In \S~9 we shall use the following result.

\begin{proposition}\label{prop:I-8-3}
Let \(\mathcal A\) be a factor and \(\varphi\) a faithful normal trace
on \(\mathcal A^+\). If there exists an infinite strictly decreasing
sequence \((E_n)\) of projections of \(\mathcal A\) such that
\(\varphi(E_n)<+\infty\) for every \(n\), then \(\mathcal A\) is
continuous.
\end{proposition}

\begin{proof}
Suppose that \(\mathcal A\) is discrete. Replacing \(\mathcal A\) by
an isomorphic algebra, we may suppose
\(\mathcal A=\mathcal L(\mathfrak H)\), where \(\mathfrak H\) is a
suitable Hilbert space. By \hyperref[thm:I-6-5]{\S~6, Theorem~5}, the
Hilbert-space dimension of \(E_n(\mathfrak H)\) is finite; hence the
infinite sequence \((E_n)\) cannot be strictly decreasing.
\end{proof}

\noindent\textit{Bibliography.}\quad
\ArticleRef{10}, \ArticleRef{42}, \ArticleRef{43}, \ArticleRef{65},
\ArticleRef{100}, \ArticleRef{117}.

\NumberedParagraph{2}{Abelian Projections}\label{no:I-8-2}

Let \(\mathcal A\) be a von Neumann algebra, and let \(E\) be a
nonzero projection of \(\mathcal A\). To say that a von Neumann
algebra consists only of the scalars is equivalent to saying that it
has no projections other than \(0\) and \(1\); hence, to say that
\(\mathcal A_E\) consists only of the scalars is equivalent to saying
that the only projections of \(\mathcal A\) majorized by \(E\) are
\(0\) and \(E\).

\begin{definition}\label{def:I-8-2}
Let \(\mathcal A\) be a von Neumann algebra and \(E\) a projection of
\(\mathcal A\). The projection \(E\) is called \emph{minimal}
(relative to \(\mathcal A\)) if \(E\neq0\) and every projection of
\(\mathcal A\) majorized by \(E\) is equal to \(0\) or to \(E\).
\end{definition}

We generalize \cref{def:I-8-2} as follows.

\begin{definition}\label{def:I-8-3}
Let \(\mathcal A\) be a von Neumann algebra and \(E\) a projection of
\(\mathcal A\). The projection \(E\) is called \emph{abelian}
(relative to \(\mathcal A\)) if the algebra \(\mathcal A_E\) is
abelian.
\end{definition}

Let \(\mathcal Z\) be the centre of \(\mathcal A\). For every
projection \(E\) of \(\mathcal A\), the centre of \(\mathcal A_E\) is
\(\mathcal Z_E\)
(\hyperref[cor:I-2-2]{\S~2, corollary to Proposition~2}). Since a von
Neumann algebra is generated by its projections, we see that a
projection \(E\) of \(\mathcal A\) is abelian if and only if every
projection \(F\) of \(\mathcal A\) majorized by \(E\) is of the form
\(EG\), where \(G\) is a projection of \(\mathcal Z\) (which may then
be taken equal to the central support of \(F\)). It follows that an
abelian projection \(E\) is minimal among the projections of
\(\mathcal A\) having the same central support as \(E\).

In particular, if \(\mathcal A\) is a factor, every nonzero abelian
projection of \(\mathcal A\) is minimal.

\begin{theorem}\label{thm:I-8-1}
Let \(\mathcal A\) be a von Neumann algebra on \(\mathfrak H\). The
following six conditions are equivalent:
\begin{enumerate}[label=\textup{(\roman*)},leftmargin=*]
\item \(\mathcal A\) is discrete;
\item \(\mathcal A'\) is discrete;
\item every nonzero projection of \(\mathcal A\) majorizes a nonzero
abelian projection of \(\mathcal A\);
\item every nonzero projection of \(\mathcal A'\) majorizes a nonzero
abelian projection of \(\mathcal A'\);
\item there exists an abelian projection of \(\mathcal A\) whose
central support is \(1\);
\item there exists an abelian projection of \(\mathcal A'\) whose
central support is \(1\).
\end{enumerate}
\end{theorem}

% Source PDF page 131 (printed p. 124)
\phantomsection\label{srcpage:124}
\begin{proof}
We shall establish the implications
\[
 \textup{(i)}\Longrightarrow\textup{(iii)}
 \Longrightarrow\textup{(v)}\Longrightarrow\textup{(ii)}
 \Longrightarrow\textup{(i)}.
\]
The implications from \textup{(i)} to \textup{(iii)} and from
\textup{(ii)} to \textup{(iv)} [respectively, those from
\textup{(iii)} to \textup{(v)} and from \textup{(iv)} to
\textup{(vi)}, and those from \textup{(v)} to \textup{(ii)} and from
\textup{(vi)} to \textup{(i)}] are equivalent, as is seen by
interchanging the roles of \(\mathcal A\) and \(\mathcal A'\). It
therefore suffices to establish the implications from \textup{(i)} to
\textup{(iii)}, from \textup{(iii)} to \textup{(v)}, and from
\textup{(v)} to \textup{(ii)}.

\textup{(i)}\(\Rightarrow\)\textup{(iii)}: We may suppose
\(\mathcal A'\) abelian, and we must show that every nonzero projection
\(E\) of \(\mathcal A\) majorizes a nonzero abelian projection of
\(\mathcal A\). Let \(x\) be a nonzero element of
\(E(\mathfrak H)\), and let \(F=E_x^{\mathcal A'}\in\mathcal A\); then
\(F\leq E\). The algebra \(\mathcal A'_F\) is abelian and has \(x\) as
a totalizing element, and hence
\(\mathcal A_F\) is abelian
(\hyperref[cor:I-6-5]{\S~6, Proposition~4, Corollary~2}). Thus \(F\)
is an abelian projection of \(\mathcal A\) majorized by \(E\).

\textup{(iii)}\(\Rightarrow\)\textup{(v)}: Let \(\mathcal Z\) be the
centre of \(\mathcal A\), and let \((E_i)_{i\in I}\) be a maximal
family of nonzero pairwise orthogonal projections of \(\mathcal Z\)
having the following property: for every \(i\in I\), there is an
abelian projection \(F_i\) of \(\mathcal A\) with central support
\(E_i\). If every nonzero projection of \(\mathcal Z\) majorizes a
nonzero abelian projection of \(\mathcal A\), then
\[
 \sum_{i\in I}E_i=1.
\]
Consequently \(F=\sum_{i\in I}F_i\) is a projection of
\(\mathcal A\) with central support \(1\). On the other hand,
\[
 \mathcal A_F=\prod_{i\in I}\mathcal A_{F_i},
\]
and therefore \(F\) is abelian.

\textup{(v)}\(\Rightarrow\)\textup{(ii)}: Let \(F\) be an abelian
projection of \(\mathcal A\) with central support \(1\). Then
\(\mathcal A_F\) is abelian, so \(\mathcal A'_F\) is discrete; hence
\(\mathcal A'\), which is isomorphic to \(\mathcal A'_F\)
(\hyperref[prop:I-2-2]{\S~2, Proposition~2}), is discrete.
\end{proof}

\setcounter{corollary}{0}
\begin{corollary}\label{cor:I-8-thm1-1}
If \(\mathcal A\) is continuous, then \(\mathcal A'\) is continuous.
\end{corollary}

\begin{proof}
If \(\mathcal A'\) is not continuous, there exists a nonzero
projection \(E\) in \(\mathcal A\cap\mathcal A'\) such that
\(\mathcal A'_E\) is discrete. Then \(\mathcal A_E\) is discrete, and
therefore \(\mathcal A\) is not continuous.
\end{proof}

\begin{corollary}\label{cor:I-8-thm1-2}
Every abelian von Neumann algebra is discrete.
\end{corollary}

\begin{proof}
If \(\mathcal A\) is abelian, \(\mathcal A'\) is discrete, and hence
\(\mathcal A\) is discrete.
\end{proof}

\begin{corollary}\label{cor:I-8-thm1-3}
Let \(\mathcal A\) be a factor on \(\mathfrak H\). The following five
conditions are equivalent:
\begin{enumerate}[label=\textup{(\roman*)},leftmargin=*]
\item \(\mathcal A\) is discrete;
\item \(\mathcal A'\) is discrete;
\item \(\mathcal A\) has minimal projections;
\item \(\mathcal A'\) has minimal projections;
\item there exist Hilbert spaces \(\mathfrak H_1,\mathfrak H_2\) and a
spatial isomorphism of \(\mathfrak H\) onto
\(\mathfrak H_1\otimes\mathfrak H_2\) which transforms \(\mathcal A\)
into \(\mathcal L(\mathfrak H_1)\otimes\mathbb C_{\mathfrak H_2}\),
and \(\mathcal A'\) into
\(\mathbb C_{\mathfrak H_1}\otimes\mathcal L(\mathfrak H_2)\).
\end{enumerate}
\end{corollary}

\begin{proof}
Conditions \textup{(i)}, \textup{(ii)}, \textup{(iii)}, and
\textup{(iv)} are equivalent by \cref{thm:I-8-1}. Condition
\textup{(v)} plainly implies that \(\mathcal A\) and \(\mathcal A'\)
are discrete. Finally, suppose \(\mathcal A\) is isomorphic to an
algebra \(\mathcal L(\mathfrak H_1)\). There exist
(\hyperref[thm:I-4-3]{\S~4, Theorem~3}) a Hilbert space \(\mathfrak K\)
and a projection
\[
 E'\in
 \bigl(\mathcal L(\mathfrak H_1)\otimes\mathbb C_{\mathfrak K}\bigr)'
 =\mathbb C_{\mathfrak H_1}\otimes\mathcal L(\mathfrak K)
\]
such that \(\mathcal A\) is spatially isomorphic to
\[
 \bigl(\mathcal L(\mathfrak H_1)
       \otimes\mathbb C_{\mathfrak K}\bigr)_{E'}.
\]
One has \(E'=1\otimes E\), where \(E\) is a projection of
\(\mathcal L(\mathfrak K)\). On putting
\(\mathfrak H_2=E(\mathfrak K)\), we see that \(\mathcal A\) is
% Source PDF page 132 (printed p. 125)
\phantomsection\label{srcpage:125}
spatially isomorphic to
\(\mathcal L(\mathfrak H_1)\otimes\mathbb C_{\mathfrak H_2}\).
\end{proof}

Condition~\textup{(iii)} of \cref{cor:I-8-thm1-3} therefore gives an
``internal'' characterization of von Neumann algebras isomorphic to
algebras \(\mathcal L(\mathfrak H)\).

\Cref{cor:I-8-thm1-2} says that an abelian von Neumann algebra
\(\mathcal A\) is isomorphic to an abelian von Neumann algebra
\(\mathcal B\) such that \(\mathcal B'\) is abelian, that is, such that
\(\mathcal B'=\mathcal B\), which was also known from
\cref{thm:I-7-1,thm:I-7-2}.

Instead of an ``abelian'' projection, one sometimes says an
``irreducible'' projection.

\noindent\textit{Bibliography.}\quad
\ArticleRef{6}, \ArticleRef{10}, \ArticleRef{42}, \ArticleRef{43},
\ArticleRef{65}, \ArticleRef{100}, \ArticleRef{117}.

\Needspace{4\baselineskip}
\NumberedParagraph{3}{Discrete Algebras and Elementary Operations}
\label{no:I-8-3}

We already know that the commutant of a discrete (respectively,
continuous) von Neumann algebra is a discrete (respectively,
continuous) von Neumann algebra, and that a product of discrete
(respectively, continuous) von Neumann algebras is a discrete
(respectively, continuous) von Neumann algebra.

\begin{proposition}\label{prop:I-8-4}
Let \(\mathcal A\) be a discrete (respectively, continuous) von
Neumann algebra, and let \(E\) be a projection of \(\mathcal A\) or of
\(\mathcal A'\). Then \(\mathcal A_E\) is discrete (respectively,
continuous).
\end{proposition}

\begin{proof}
Since \(\mathcal A'\) is discrete (respectively, continuous), it is
enough to consider the case \(E\in\mathcal A'\). There is then a
projection \(F\) of the centre of \(\mathcal A\) such that
\(\mathcal A_E\) is isomorphic to \(\mathcal A_F\). Now
\(\mathcal A=\mathcal A_F\times\mathcal A_{1-F}\), and therefore
\(\mathcal A_F\) is discrete (respectively, continuous) by
\cref{prop:I-8-1}.
\end{proof}

\begin{proposition}\label{prop:I-8-5}
Let \(\mathcal A_1\) and \(\mathcal A_2\) be discrete von Neumann
algebras. Then \(\mathcal A_1\otimes\mathcal A_2\)
is discrete.
\end{proposition}

\begin{proof}
If \(\mathcal A_1\) and \(\mathcal A_2\) are replaced by isomorphic
von Neumann algebras, then
\(\mathcal A_1\otimes\mathcal A_2\) is replaced by
an isomorphic von Neumann algebra
(\hyperref[prop:I-4-2]{\S~4, Proposition~2}). We may therefore suppose
that \(\mathcal A_1'\) and \(\mathcal A_2'\) are abelian. Then
\(\mathcal A_1'\otimes\mathcal A_2'\) is abelian.
But
\[
 \bigl(\mathcal A_1\otimes\mathcal A_2\bigr)'
  =\mathcal A_1'\otimes\mathcal A_2'
  \qquad(\hyperref[prop:I-6-14]{\S~6, Proposition~14}).
\]
Thus \(\mathcal A_1\otimes\mathcal A_2\) is
discrete.
\end{proof}

In Chapter~III, \S~3, we shall study the structure of discrete von
Neumann algebras more thoroughly (cf. also Chapter~III, \S~8, the
corollary to Theorem~2).

\noindent\textit{Bibliography.}\quad \ArticleRef{65}.

% Source PDF page 133 (printed p. 126)
\phantomsection\label{srcpage:126}
\NumberedParagraph{4}{Definition of the Types}\label{no:I-8-4}

Purely infinite von Neumann algebras are called algebras of type~III.
Continuous semifinite von Neumann algebras are also called algebras of
type~II, and, more precisely, of type~\(\mathrm{II}_1\) if they are
finite and of type~\(\mathrm{II}_\infty\) if they are properly
infinite. Discrete algebras are also called algebras of type~I. For a
von Neumann algebra \(\mathcal A\) to be of type~I (respectively,
II, III), it is necessary and sufficient that \(\mathcal A'\) be of
type~I (respectively, II, III).

Let \(\mathcal A\) be a discrete factor, and let \(\mathfrak K\) be a
complex Hilbert space such that \(\mathcal A\) is isomorphic to
\(\mathcal L(\mathfrak K)\). Let \(n\) be the Hilbert-space dimension
of \(\mathfrak K\). Then \(n\) is the cardinality of every family of
pairwise orthogonal minimal projections of \(\mathcal A\) whose sum is
\(1\). Thus \(n\) is an algebraic invariant of \(\mathcal A\). We say
that \(\mathcal A\) is of type \(\mathrm I_n\). If \(n\) is an
infinite cardinal, we also say that \(\mathcal A\) is of type
\(\mathrm I_\infty\). We shall generalize these notations in
Chapter~III, \S~3.

\noindent\textit{Bibliography.}\quad
\ArticleRef{42}, \ArticleRef{43}, \ArticleRef{65}.

\NumberedParagraph{5}{Complete Hilbert Algebras and Type I Factors}
\label{no:I-8-5}

\begin{lemma}\label{lem:I-8-1}
Let \(\mathfrak A\) be a nonzero complete Hilbert algebra. Then
\(\mathfrak A\), \(\mathcal U(\mathfrak A)\), and
\(\mathcal V(\mathfrak A)\) possess nonzero minimal Hermitian
idempotents.
\end{lemma}

\begin{proof}
\textit{a.} Let \(\mathcal E\) be the set of nonzero Hermitian
idempotents of \(\mathfrak A\). Let us show that \(\mathcal E\) is
nonempty. There exists a nonzero Hermitian element \(a\) in
\(\mathfrak A\). Then \(U_a\) is nonzero and Hermitian. There exists
\(T\in\mathcal U(\mathfrak A)\) such that \(TU_a\) is a nonzero
spectral projection of \(U_a\). Since \(TU_a=U_{Ta}\), \(Ta\) is a
nonzero Hermitian idempotent of \(\mathfrak A\).

\textit{b.} The map \((x,y)\mapsto xy\) from
\(\mathfrak A\times\mathfrak A\) into \(\mathfrak A\) is separately
continuous. Since \(\mathfrak A\) is complete, this map is continuous
(\TreatiseRef{3}, Chapter~III, \S~4, Proposition~2). Hence there is a
constant \(M>0\) such that
\[
 \lVert xy\rVert\leq M\lVert x\rVert\,\lVert y\rVert
 \qquad(x,y\in\mathfrak A).
\]
If \(e\in\mathfrak A\) is a nonzero idempotent, then
\[
 \lVert e\rVert=\lVert e^2\rVert\leq M\lVert e\rVert^2,
 \qquad\text{so that}\qquad \lVert e\rVert\geq M^{-1}.
\]

\textit{c.} Put
\[
 \alpha=\inf_{e\in\mathcal E}\lVert e\rVert^2.
\]
Then \(\alpha\geq M^{-2}\). There exists \(e\in\mathcal E\) such that
\(\lVert e\rVert^2<2\alpha\). If
\(e=e_1+e_2\), with \(e_1,e_2\in\mathcal E\) and
\(e_1e_2=e_2e_1=0\), then
\[
 (e_1\mid e_2)=(e_1^2\mid e_2)=(e_1\mid e_1e_2)=0,
\]
and hence
\[
 2\alpha\leq\lVert e_1\rVert^2+\lVert e_2\rVert^2
  =\lVert e\rVert^2<2\alpha,
\]
which is absurd. Therefore \(e\) is a minimal element of
\(\mathcal E\). Every nonzero projection of
\(\mathcal U(\mathfrak A)\) majorized by \(U_e\) is of the form
\(U_{e'}\), with \(e'\in\mathfrak A\); then \(e'\) is an element of
\(\mathcal E\) majorized by \(e\), whence \(e'=e\) and
\(U_{e'}=U_e\). Thus \(U_e\) is a minimal projection of
\(\mathcal U(\mathfrak A)\). The assertion for
\(\mathcal V(\mathfrak A)\) follows symmetrically.
\end{proof}

% Source PDF page 134 (printed p. 127)
\phantomsection\label{srcpage:127}
\begin{proposition}\label{prop:I-8-6}
Let \(\mathfrak A\) be a completed Hilbert algebra such that
\(\mathcal U(\mathfrak A)\) is a factor. The following conditions are
equivalent:
\begin{enumerate}[label=\textup{(\roman*)},leftmargin=*]
\item \(\mathfrak A\) is complete;
\item \(\mathcal U(\mathfrak A)\) is of type~I;
\item after multiplication of the norm of \(\mathfrak A\) by a
constant, \(\mathfrak A\) is isomorphic to the Hilbert algebra of
Hilbert--Schmidt operators on a suitable Hilbert space.
\end{enumerate}
\end{proposition}

\begin{proof}
\textup{(iii)}\(\Rightarrow\)\textup{(i)} is obvious.

\textup{(i)}\(\Rightarrow\)\textup{(ii)} follows from
\cref{lem:I-8-1} and \cref{cor:I-8-thm1-3}.

\textup{(ii)}\(\Rightarrow\)\textup{(iii)}: since \(\mathfrak A\) is
completed, \(a\mapsto U_a\) is an isomorphism of the Hilbert algebra
\(\mathfrak A\) onto the Hilbert algebra of the Hilbert--Schmidt
elements of \(\mathcal U(\mathfrak A)\), endowed with the natural
trace (\hyperref[thm:I-6-1]{\S~6, Theorem~1}). If
\(\mathcal U(\mathfrak A)\) is of type~I, there exist a Hilbert space
\(\mathfrak K\) and an isomorphism of \(\mathcal U(\mathfrak A)\) onto
\(\mathcal L(\mathfrak K)\). This isomorphism transforms the natural
trace on \(\mathcal U(\mathfrak A)^+\) into a trace on
\(\mathcal L(\mathfrak K)^+\), which is a multiple of the trace
defined in \hyperref[thm:I-6-5]{Theorem~5}. Consequently, after
multiplication of the norm of \(\mathfrak A\) by a constant, the
Hilbert algebra \(\mathfrak A\) is isomorphic to the Hilbert algebra of
Hilbert--Schmidt operators on \(\mathfrak K\).
\end{proof}

\noindent\textbf{Remark.}\quad
By dimension considerations, \(\mathfrak A\) determines, up to
isomorphism, the Hilbert space \(\mathfrak K\) in
\hyperref[prop:I-8-6]{Proposition~6}\textup{(iii)}.

\begin{proposition}\label{prop:I-8-7}
Let \(\mathfrak A\) be a complete Hilbert algebra. Let \((E_i)\) be
the family of minimal projections of
\(\mathcal U(\mathfrak A)\cap\mathcal V(\mathfrak A)\).
\begin{enumerate}[label=\textup{(\roman*)},leftmargin=*]
\item The Hilbert space \(\mathfrak A\) is the Hilbert direct sum of
the \(\mathfrak A_i=E_i(\mathfrak A)\).
\item The \(\mathfrak A_i\) are self-adjoint two-sided ideals of
\(\mathfrak A\) which annihilate one another pairwise.
\item After multiplication of the norm of \(\mathfrak A_i\) by a
constant, \(\mathfrak A_i\) is isomorphic to the Hilbert algebra of
Hilbert--Schmidt operators on a suitable Hilbert space.
\item If \(\mathcal U(\mathfrak A)\) and
\(\mathcal V(\mathfrak A)\) are finite von Neumann algebras, every
\(\mathfrak A_i\) is finite-dimensional.
\end{enumerate}
\end{proposition}

\begin{proof}
It is clear that the \(E_i\) are pairwise orthogonal. Let
\[
 E=\sum_iE_i,
 \qquad F=1-E.
\]
If \(F\neq0\), then \(\mathfrak B=F(\mathfrak A)\) is a nonzero
complete Hilbert algebra. Hence \(\mathcal U(\mathfrak B)\) has a
minimal projection \(G\) (\cref{lem:I-8-1}). The central support
\(G'\) of \(G\) is a minimal projection of
\(\mathcal U(\mathfrak B)\cap\mathcal V(\mathfrak B)\). Thus there is
in \(\mathcal U(\mathfrak A)\cap\mathcal V(\mathfrak A)\) a minimal
projection orthogonal to the \(E_i\), which is absurd. Hence \(F=0\),
which proves \textup{(i)}. Assertion~\textup{(ii)} follows from
\hyperref[prop:I-5-8]{\S~5, Proposition~8}. Assertion~\textup{(iii)}
follows from \cref{prop:I-8-6}. Suppose
\(\mathcal U(\mathfrak A)\) finite. For every \(i\),
\(\mathcal U(\mathfrak A_i)\)
% Source PDF page 135 (printed p. 128)
\phantomsection\label{srcpage:128}
is a finite factor, of type~I (\cref{prop:I-8-6}), and hence of type
\(\mathrm I_n\), with \(n<+\infty\). Therefore \(\mathfrak A_i\) is a
finite-dimensional vector space; this proves \textup{(iv)}.
\end{proof}

\begin{corollary*}\phantomsection\label{cor:I-8-7}
The closed two-sided ideals of \(\mathfrak A\) are the Hilbert direct
sums of some of the \(\mathfrak A_i\). The minimal nonzero closed
two-sided ideals (respectively, the maximal proper closed two-sided
ideals) are the \(\mathfrak A_i\) (respectively, the
\(\mathfrak A\ominus\mathfrak A_i\)).
\end{corollary*}

\begin{proof}
The Hilbert direct sum of some of the \(\mathfrak A_i\) is plainly a
closed two-sided ideal of \(\mathfrak A\). Conversely, if \(I\) is a
closed two-sided ideal of \(\mathfrak A\), then
\(P_I\in\mathcal U(\mathfrak A)\cap\mathcal V(\mathfrak A)\)
(\hyperref[prop:I-5-8]{\S~5, Proposition~8}); hence, for every \(i\),
\(P_I\) majorizes \(E_i\) or is orthogonal to it. Thus \(I\) is the
Hilbert direct sum of those \(\mathfrak A_i\) which it contains. The
rest of the corollary is then obvious.
\end{proof}

\noindent\textit{Bibliography.}\quad
\ArticleRef{1}, \ArticleRef{148}, \TreatiseRef{9}.

\noindent\textit{Exercises.}---
\begin{enumerate}[label=\arabic*.,leftmargin=*,itemsep=0.5\baselineskip]
\item\label{ex:I-8-1}
Let \(\mathcal A\) be a continuous von Neumann algebra on
\(\mathfrak K\), and let \(E\) be a nonzero projection of
\(\mathcal A\). Show that the Hilbert-space dimension of
\(E(\mathfrak K)\) is infinite. [If
\(\dim E(\mathfrak K)<+\infty\), then \(\mathcal A_E\) is discrete.]

\item\label{ex:I-8-2}
Give another proof of the implication
\textup{(i)}\(\Rightarrow\)\textup{(v)} in
\cref{cor:I-8-thm1-3}. [If \(\mathcal A\) is isomorphic to
\(\mathcal L(\mathfrak K)\), there exists in \(\mathcal A\) a family
\((E_i)_{i\in I}\) of pairwise orthogonal minimal projections such
that
\[
 \sum_{i\in I}E_i=1,
\]
and operators \(U_{ik}\) such that
\[
 U_{ik}^*U_{ik}=E_i,
 \qquad
 U_{ik}U_{ik}^*=E_k.
\]
Use \hyperref[prop:I-2-5]{\S~2, Proposition~5}.]

\item\label{ex:I-8-3}
Let \(\mathcal A\) be a discrete factor on \(\mathfrak H\), and let
\(\mathcal B\) be a von Neumann algebra contained in \(\mathcal A'\).
If the von Neumann algebra \(\mathcal C\) generated by \(\mathcal A\)
and \(\mathcal B\) is \(\mathcal L(\mathfrak H)\), then
\(\mathcal B=\mathcal A'\). [If \(\mathcal B\neq\mathcal A'\), there
exist in \(\mathcal A'\), which is isomorphic to an algebra
\(\mathcal L(\mathfrak K)\), nonscalar operators commuting with
\(\mathcal B\), and hence with \(\mathcal C\).] \ArticleRef{65}
(Cf. Chapter~III, \S~7, Exercise~13d.)

\item\label{ex:I-8-4}
Let \(\mathcal A\) be a von Neumann algebra, \(\mathcal Z\) its
centre, \(E\) a minimal projection of \(\mathcal A\), and \(F\) its
central support. Show that \(F\) is a minimal projection of
\(\mathcal Z\), and that \(\mathcal A_F\) is a discrete factor. Deduce
that the supremum of the minimal projections of \(\mathcal A\) is a
projection \(G\) of \(\mathcal Z\), and that \(\mathcal A_G\) is a
product of discrete factors.

\item\label{ex:I-8-5}
\begin{enumerate}[label=\textit{\alph*.},leftmargin=2em]
\item Let \(\mathcal Z\) be an abelian von Neumann algebra and
\(\chi\) a normal character of \(\mathcal Z\). Show that the support
of \(\chi\) is a minimal projection of \(\mathcal Z\).
\item Let \(Z\) be the interval \([0,1]\), let \(\nu\) be Lebesgue
measure on \(Z\), let
\(\mathfrak H=L_{\mathbb C}^{2}(Z,\nu)\), and let \(\mathcal A\) be
the von Neumann algebra of \cref{thm:I-7-2}. Show that no character of
\(\mathcal A\) is normal (use \textit{a}). \ArticleRef{70}
\end{enumerate}

\item\label{ex:I-8-6}
Let \(\mathcal A\) be a von Neumann algebra having the following
property:

\smallskip
\noindent\textup{(M)} Every decreasing sequence of projections of
\(\mathcal A\) is stationary.
\begin{enumerate}[label=\textit{\alph*.},leftmargin=2em]
\item Show that every nonzero projection of \(\mathcal A\) majorizes
a minimal projection.
\item Deduce from \textit{a} that \(1\) is a finite sum of pairwise
orthogonal minimal projections.
\item Show that \(\mathcal A\) is the product of a finite number of
factors having property~\textup{(M)}. (Apply \textit{b} to the centre
of \(\mathcal A\).)
\end{enumerate}
\end{enumerate}
% Source PDF page 136 (printed p. 129)
\phantomsection\label{srcpage:129}
\begin{enumerate}[label=\textit{\alph*.},leftmargin=2em,start=4]
\item Show that a factor having property~\textup{(M)} is a
finite-dimensional factor. (Use \textit{b}.)
\item Show that \textup{(M)} is equivalent to the following property:
the vector space \(\mathcal A\) is finite-dimensional. (Use
\textit{c} and \textit{d}.)
\end{enumerate}

\begin{enumerate}[label=\arabic*.,leftmargin=*,itemsep=0.5\baselineskip,start=7]
\item\label{ex:I-8-7}
\begin{enumerate}[label=\textit{\alph*.},leftmargin=2em]
\item Let \(\mathcal A\) be a von Neumann algebra on
\(\mathfrak H\), and let
\((E_n)=(P_{\mathfrak X_n})\) be an infinite sequence of nonzero,
pairwise orthogonal projections of \(\mathcal A\). Let
\(x_n\in\mathfrak X_n\), with \(\lVert x_n\rVert=1\). Let
\(\mathcal U\) be a nonprincipal ultrafilter on the set of integers
\(\geq0\). For \(T\in\mathcal A\), put
\[
 \varphi(T)=\lim_{\mathcal U}(Tx_n\mid x_n).
\]
Show that \(\varphi\) is a linear form on \(\mathcal A\), continuous
for the norm topology but not for the ultra-strong topology.
[Consider the projection \(F_n\), the supremum of
\(E_{n+1},E_{n+2},\ldots\); \(F_n\) tends ultra-strongly to zero and
\(\varphi(F_n)=1\).]

\item Show that a von Neumann algebra which, as a vector space, has
infinite dimension is a nonreflexive Banach space. (Use \textit{a},
Exercise~6, and \hyperref[thm:I-3-1]{\S~3, Theorem~1}.)
\end{enumerate}
\end{enumerate}

\section{Existence of the Different Types of Factors}\label{sec:I-9}

\NumberedParagraph{1}{A Lemma}\label{no:I-9-1}

\begin{lemma}\label{lem:I-9-1}
Let \(\mathcal A\) be a von Neumann algebra, \(\mathcal B\) a von
Neumann subalgebra of \(\mathcal A\), maximal abelian in
\(\mathcal A\), \(\varphi\) a faithful normal trace on
\(\mathcal A^+\), and \(\mathfrak m\) the ideal of definition of
\(\dot\varphi\). For \(T\in\mathfrak m^{1/2}\), let
\(\mathcal K_T\) be the convex hull in \(\mathcal A\) of the set of
the \(UTU^{-1}\), where \(U\) runs through the unitary operators of
\(\mathcal B\). Let \(\mathcal K_T'\) be the weak closure of
\(\mathcal K_T\). Then
\[
 \mathcal K_T'\subset\mathfrak m^{1/2},
 \qquad
 \mathcal K_T'\cap\mathcal B\ne\varnothing.
\]
\end{lemma}

\begin{proof}
Recall that \(\varphi\) defines on \(\mathfrak m^{1/2}\) a
pre-Hilbert-space structure, and that the corresponding norm is
denoted by \(\lVert\cdot\rVert_2\). We may clearly restrict ourselves
to the case in which \(\lVert T\rVert\leq1\) and
\(\lVert T\rVert_2\leq1\). Then, for every unitary operator \(U\) of
\(\mathcal B\),
\[
 \lVert UTU^{-1}\rVert\leq1,
 \qquad
 \lVert UTU^{-1}\rVert_2\leq1.
\]
If \(\mathcal A_1\) denotes the unit ball of \(\mathcal A\) for the
norm \(\lVert\cdot\rVert\), and \(\mathcal K\) the unit ball of the
pre-Hilbert space \(\mathfrak m^{1/2}\), then
\(\mathcal K_T\subset\mathcal A_1\cap\mathcal K\). On the other hand,
\[
 \varphi=\sum_{i\in I}\omega_{x_i}
\]
(\hyperref[srcpage:85]{\S~6, corollary to Proposition~2}). Hence
\[
 \varphi(S^*S)=\sum_{i\in I}\lVert Sx_i\rVert^2
\]
is a lower semicontinuous function of \(S\) for the weak topology on
\(\mathcal A\). Thus \(\mathcal K\) is weakly closed; consequently,
\(\mathcal K_T'\subset\mathcal A_1\cap\mathcal K\), and, in
particular, \(\mathcal K_T'\subset\mathfrak m^{1/2}\).

The set \(\mathcal K_T\) is invariant under the group \(\mathcal G\)
of transformations \(S\mapsto USU^{-1}\), where \(U\) runs through
the unitary operators of \(\mathcal B\); the same is therefore true
of \(\mathcal K_T'\). The set \(\mathcal K_T'\) is weakly compact,
% Source PDF page 137 (printed p. 130)
\phantomsection\label{srcpage:130}
and hence the function \(S\mapsto\lVert S\rVert_2\) attains its
minimum on \(\mathcal K_T'\) at an element
\(T_0\in\mathcal K_T'\). This element is unique because, since
\(\varphi\) is faithful, \(\mathfrak m^{1/2}\) is strictly convex
for the norm \(\lVert\cdot\rVert_2\). Thus \(T_0\) is a fixed point
of \(\mathcal G\), that is, \(T_0\) commutes with the unitary
operators of \(\mathcal B\); hence
\[
 T_0\in\mathcal B'\cap\mathcal A=\mathcal B
\]
(because \(\mathcal B\) is maximal abelian), and finally
\(T_0\in\mathcal K_T'\cap\mathcal B\).
\end{proof}

\noindent\textit{Bibliography.}\quad
\ArticleRef{15}, \ArticleRef{78}.

\NumberedParagraph{2}{Construction of Certain von Neumann Algebras}
\label{no:I-9-2}
Let \(\mathfrak H\) be a complex Hilbert space and \(\mathcal A\) an
abelian von Neumann algebra on \(\mathfrak H\). Let \(G\) be a
discrete group with identity element \(e\), and let
\(s\mapsto U_s\) be a unitary representation of \(G\) on
\(\mathfrak H\). Suppose that
\[
 U_s^{-1}\mathcal A U_s=\mathcal A
 \qquad(s\in G).
\]
It follows that, for every \(s\in G\), the map
\[
 T\longmapsto U_s^{-1}TU_s=T^s
\]
is an automorphism of the von Neumann algebra \(\mathcal A\).

For every \(s\in G\), let \(\mathfrak H_s\) be a Hilbert space of
the same Hilbert-space dimension as \(\mathfrak H\), and let
\(J_s\) be an isometric linear map from \(\mathfrak H\) onto
\(\mathfrak H_s\). Let \(\widetilde{\mathfrak H}\) be the Hilbert
direct sum of the \(\mathfrak H_s\), for \(s\in G\). As in \S~2,
no.~3, we shall represent every element \(R\) of
\(\mathcal L(\widetilde{\mathfrak H})\) by the matrix
\((R_{s,t})_{s\in G,t\in G}\), where
\[
 R_{s,t}=J_s^*RJ_t\in\mathcal L(\mathfrak H).
\]

For every \(T\in\mathcal A\), let \(\Phi(T)\) be the element of
\(\mathcal L(\widetilde{\mathfrak H})\) defined by the following
matrix \((R_{s,t})\): \(R_{s,t}=0\) if \(s\ne t\), and
\(R_{s,s}=T\) for every \(s\in G\). In other words,
\(\Phi(T)\) is the element of
\(\mathcal L(\widetilde{\mathfrak H})\) reduced by the
\(\mathfrak H_s\) such that
\[
 J_s^*\Phi(T)J_s=T
 \qquad(s\in G).
\]
The map \(\Phi\) is an isomorphism of the von Neumann algebra
\(\mathcal A\) onto an abelian von Neumann algebra
\(\widetilde{\mathcal A}\subset
\mathcal L(\widetilde{\mathfrak H})\).

For every \(y\in G\), let \(\widetilde U_y\) be the element of
\(\mathcal L(\widetilde{\mathfrak H})\) defined by the following
matrix \((R_{s,t})\): \(R_{s,t}=0\) if \(st^{-1}\ne y\), and
\(R_{yt,t}=U_y\) for every \(t\in G\). In other words,
\(\widetilde U_y\) is a unitary operator which permutes the
\(\mathfrak H_s\) among themselves and, more precisely, carries
\(\mathfrak H_s\) onto \(\mathfrak H_{ys}\), with
\[
 J_{ys}^*\widetilde U_yJ_s=U_y.
\]
Thus \(\widetilde U_z\widetilde U_y\) carries
\(\mathfrak H_s\) onto \(\mathfrak H_{zys}\), with
\[
 J_{zys}^*\widetilde U_z\widetilde U_yJ_s
  =J_{zys}^*\widetilde U_zJ_{ys}
   J_{ys}^*\widetilde U_yJ_s
  =U_zU_y=U_{zy}.
\]
Consequently,
\begin{equation}
 \widetilde U_z\widetilde U_y=\widetilde U_{zy}.
 \tag{1}\label{eq:I-9-1}
\end{equation}
On the other hand, \(\widetilde U_y^{-1}\Phi(T)\widetilde U_y\)
is reduced by every \(\mathfrak H_s\), and
\begin{align*}
 J_s^*\bigl(\widetilde U_y^{-1}\Phi(T)\widetilde U_y\bigr)J_s
 &=\bigl(J_s^*\widetilde U_{y^{-1}}J_{ys}\bigr)
   \bigl(J_{ys}^*\Phi(T)J_{ys}\bigr)
   \bigl(J_{ys}^*\widetilde U_yJ_s\bigr)\\
 &=U_{y^{-1}}TU_y=T^y;
\end{align*}
thus
\begin{equation}
 \widetilde U_y^{-1}\Phi(T)\widetilde U_y=\Phi(T^y).
 \tag{2}\label{eq:I-9-2}
\end{equation}

% Source PDF page 138 (printed p. 131)
\phantomsection\label{srcpage:131}
The operators of the form
\[
 \Phi(T_1)\widetilde U_{y_1}
 +\Phi(T_2)\widetilde U_{y_2}
 +\cdots+
 \Phi(T_n)\widetilde U_{y_n}
\]
constitute an involutive subalgebra \(\mathcal B_0\) of
\(\mathcal L(\widetilde{\mathfrak H})\). Indeed,
\[
 \bigl(\Phi(T)\widetilde U_y\bigr)^*
 =\widetilde U_{y^{-1}}\Phi(T^*)
 =\Phi(T^{*y})\widetilde U_{y^{-1}},
\]
and
\begin{align*}
 \Phi(T_1)\widetilde U_{y_1}\Phi(T_2)\widetilde U_{y_2}
 &=\Phi(T_1)\Phi\bigl(T_2^{y_1^{-1}}\bigr)
   \widetilde U_{y_1}\widetilde U_{y_2}\\
 &=\Phi\bigl(T_1T_2^{y_1^{-1}}\bigr)
   \widetilde U_{y_1y_2}.
\end{align*}

Let \(\mathcal B\) be the von Neumann algebra generated by
\(\mathcal B_0\), that is, the weak closure of \(\mathcal B_0\).
The algebra \(\widetilde{\mathcal A}\) is an abelian von Neumann
subalgebra of \(\mathcal B\).

Let us calculate the matrix \((R_{s,t})\) of
\(\Phi(T)\widetilde U_y\). One has \(R_{s,t}=0\) if
\(st^{-1}\ne y\), and
\[
 R_{yt,t}
 =J_{yt}^*\Phi(T)\widetilde U_yJ_t
 =\bigl(J_{yt}^*\Phi(T)J_{yt}\bigr)
  \bigl(J_{yt}^*\widetilde U_yJ_t\bigr)
 =TU_y.
\]
There is therefore a family \((T_y)_{y\in G}\) in \(\mathcal A\)
such that
\[
 R_{s,t}=T_{st^{-1}}U_{st^{-1}}.
\]
This property is consequently also true for the operators of
\(\mathcal B_0\), and it is preserved on passing to the weak limit:
every element \(S\in\mathcal B\) is represented by a matrix of the
form
\[
 \bigl(T_{st^{-1}}U_{st^{-1}}\bigr),
 \qquad T_y\in\mathcal A\quad(y\in G).
\]
In particular, for \(S\in\mathcal B^+\),
\[
 T_e=J_e^*SJ_e\in\mathcal A^+.
\]

\begin{proposition}\label{prop:I-9-1}
Let \(\varphi\) be a faithful, semifinite normal trace on
\(\mathcal A^+\), invariant under \(G\). For every
\[
 S=\bigl(T_{st^{-1}}U_{st^{-1}}\bigr)\in\mathcal B^+,
\]
put \(\psi(S)=\varphi(T_e)\). Then \(\psi\) is a faithful,
semifinite normal trace on \(\mathcal B^+\). This trace is finite if
and only if \(\varphi\) is finite. For \(T\in\mathcal A^+\),
\[
 \psi\bigl(\Phi(T)\bigr)=\varphi(T).
\]
\end{proposition}

\begin{proof}
For \(S,S_1\in\mathcal B^+\) and \(\lambda\geq0\), one immediately
has
\[
 \psi(S+S_1)=\psi(S)+\psi(S_1),
 \qquad
 \psi(\lambda S)=\lambda\psi(S).
\]
Let us show that \(\psi(SS^*)=\psi(S^*S)\) for
\(S\in\mathcal B\). Let
\(S=(T_{st^{-1}}U_{st^{-1}})\). One has
\[
 S^*=\bigl(U_{ts^{-1}}^*T_{ts^{-1}}^*\bigr).
\]
Put
\[
 SS^*=\bigl(R_{st^{-1}}U_{st^{-1}}\bigr),
 \qquad
 S^*S=\bigl(R'_{st^{-1}}U_{st^{-1}}\bigr).
\]
Then
\begin{align*}
 R_e
 &=\sum_{t\in G}T_{t^{-1}}U_{t^{-1}}U_{t^{-1}}^*T_{t^{-1}}^*
   =\sum_{t\in G}T_tT_t^*,\\
 R'_e
 &=\sum_{t\in G}U_t^*T_t^*T_tU_t,
\end{align*}
% Source PDF page 139 (printed p. 132)
\phantomsection\label{srcpage:132}
the series converging in the strong topology. All the terms of these
series belong to \(\mathcal A^+\). Hence
\begin{align*}
 \psi(S^*S)
 &=\varphi(R'_e)
  =\sum_{t\in G}\varphi(U_t^*T_t^*T_tU_t)
  =\sum_{t\in G}\varphi(T_tT_t^*)\\
 &=\varphi(R_e)=\psi(SS^*).
\end{align*}
Thus \(\psi\) is a trace, and the normality of \(\varphi\) at once
implies that of \(\psi\).

Let us show that \(\psi\) is faithful. Let
\(S=(T_{st^{-1}}U_{st^{-1}})\in\mathcal B^+\), and suppose that
\(\psi(S)=\varphi(T_e)=0\). Since \(\varphi\) is faithful, it follows
that \(T_e=0\). Hence, for every \(x\in\mathfrak H_s\),
\begin{align*}
 \lVert S^{1/2}x\rVert^2
 &=(Sx\mid x)
  =(SJ_sJ_s^*x\mid J_sJ_s^*x)\\
 &=(T_eJ_s^*x\mid J_s^*x)=0.
\end{align*}
Thus \(S^{1/2}x=0\); hence \(S^{1/2}=0\), and consequently
\(S=0\).

Let us show that \(\psi\) is semifinite. The operator
\(I_{\mathfrak H}\in\mathcal A\) is the supremum of an increasing
directed family of operators \(T_i\in\mathcal A^+\) such that
\(\varphi(T_i)<+\infty\). The formula
\(\psi(\Phi(T))=\varphi(T)\) being evident, the operator
\(I_{\widetilde{\mathfrak H}}\in\mathcal B\) is the supremum of the
increasing directed family \((\Phi(T_i))\), and
\(\psi(\Phi(T_i))<+\infty\). Finally, it is clear that \(\psi\) is
finite if and only if \(\varphi\) is finite.
\end{proof}

\begin{proposition}\label{prop:I-9-2}
Suppose that \(\widetilde{\mathcal A}\) is a maximal abelian von
Neumann subalgebra of \(\mathcal B\). For \(\mathcal B\) to be
semifinite, it is necessary and sufficient that there exist on
\(\mathcal A^+\) a faithful, semifinite normal trace invariant under
\(G\).
\end{proposition}

\begin{proof}
The condition is sufficient by \cref{prop:I-9-1}. Conversely, suppose
that \(\mathcal B\) is semifinite. Let \(\psi\) be a faithful,
semifinite normal trace on \(\mathcal B^+\). This trace induces on
\(\widetilde{\mathcal A}^+\) a faithful normal trace. There exists on
\(\mathcal A^+\) a faithful normal trace \(\varphi\) such that
\[
 \varphi(T)=\psi\bigl(\Phi(T)\bigr)
 \qquad(T\in\mathcal A^+).
\]
This trace is invariant under \(G\). Indeed, for
\(T\in\mathcal A^+\) and \(y\in G\), formula~\eqref{eq:I-9-2}
gives
\begin{align*}
 \varphi(U_y^{-1}TU_y)
 &=\varphi(T^y)
  =\psi\bigl(\Phi(T^y)\bigr)
  =\psi\bigl(\widetilde U_y^{-1}\Phi(T)\widetilde U_y\bigr)\\
 &=\psi\bigl(\Phi(T)\bigr)=\varphi(T).
\end{align*}

The whole difficulty is to show that \(\varphi\) is semifinite. To
this end we first establish the following result: if
\[
 S=\bigl(T_{st^{-1}}U_{st^{-1}}\bigr)\in\mathcal B
 \quad\text{and}\quad
 \psi(S^*S)<+\infty,
\]
then \(\varphi(T_e^*T_e)<+\infty\). Indeed, for every
\(T_1=\Phi(T)\in\widetilde{\mathcal A}\) and every \(s\in G\),
\[
 J_s^*T_1SJ_s
 =T J_s^*SJ_s
 =TT_e
 =T_eT
 =J_s^*SJ_sT
 =J_s^*ST_1J_s,
\]
where we have used, for the first time in an essential way, the fact
that \(\mathcal A\) is abelian. In particular, if \(T_1\) is a
unitary operator of \(\widetilde{\mathcal A}\), then, on replacing
\(S\) by \(ST_1^{-1}\) in the preceding equality, one obtains
\[
 J_s^*(T_1ST_1^{-1})J_s=J_s^*SJ_s.
\]
% Source PDF page 140 (printed p. 133)
\phantomsection\label{srcpage:133}
The operator \(J_s^*RJ_s\) therefore remains fixed as \(R\) runs
through the convex set \(\mathcal K_S\) generated by the operators
\(T_1ST_1^{-1}\), where \(T_1\) is a variable unitary operator of
\(\widetilde{\mathcal A}\), and hence as \(R\) runs through the weak
closure \(\mathcal K_S'\) of \(\mathcal K_S\). Now
\(\widetilde{\mathcal A}\) is a maximal abelian von Neumann
subalgebra of \(\mathcal B\). By \cref{lem:I-9-1},
\(\mathcal K_S'\) contains an element
\(S_0\in\widetilde{\mathcal A}\) such that
\(\psi(S_0^*S_0)<+\infty\). The equality
\[
 J_s^*S_0J_s=J_s^*SJ_s=T_e
\]
implies \(S_0=\Phi(T_e)\), which proves our assertion.

The operator \(I_{\widetilde{\mathfrak H}}\) is the strong limit of
operators \(S\in\mathcal B\) such that \(\psi(S^*S)<+\infty\). If
\(S=(T_{st^{-1}}U_{st^{-1}})\), then
\(T_e=J_e^*SJ_e\) tends strongly to \(I_{\mathfrak H}\), and
\(\varphi(T_e^*T_e)<+\infty\). By
\hyperref[cor:I-6-2]{\S~6, Proposition~1, Corollary~2}, this proves
that \(\varphi\) is semifinite.
\end{proof}

\begin{lemma}\label{lem:I-9-2}
Suppose that \(\mathcal A\) is a maximal abelian von Neumann
subalgebra of \(\mathcal L(\mathfrak H)\). Suppose, moreover, that
\[
 \mathcal A\cap\mathcal A U_y=0
 \qquad(y\ne e).
\]
Then \(\widetilde{\mathcal A}\) is a maximal abelian von Neumann
subalgebra of \(\mathcal B\).
\end{lemma}

\begin{proof}
Let \(S=(T_{st^{-1}}U_{st^{-1}})\) be an element of \(\mathcal B\)
which commutes with \(\widetilde{\mathcal A}\). Thus, for every
\(T\in\mathcal A\), one has \(S\Phi(T)=\Phi(T)S\), whence, by
calculating matrices,
\[
 TT_{st^{-1}}U_{st^{-1}}
 =T_{st^{-1}}U_{st^{-1}}T.
\]
Since \(\mathcal A\) is maximal abelian, this implies
\(T_yU_y\in\mathcal A\), that is,
\(T_y\in\mathcal A U_y^{-1}\), for every \(y\in G\). On the other
hand, \(T_y\in\mathcal A\); the second hypothesis of the lemma
therefore implies \(T_y=0\) for \(y\ne e\). Hence
\(S=\Phi(T_e)\in\widetilde{\mathcal A}\).
\end{proof}

\begin{lemma}\label{lem:I-9-3}
To the hypotheses of \cref{lem:I-9-2}, add the hypothesis that the
only elements of \(\mathcal A\) invariant under \(G\) are scalar
operators. Then \(\mathcal B\) is a factor.
\end{lemma}

\begin{proof}
Let \(S\) be an element of the centre of \(\mathcal B\). By
\cref{lem:I-9-2}, \(S\in\widetilde{\mathcal A}\), and hence
\(S=\Phi(T)\) for some \(T\in\mathcal A\). On the other hand, for
every \(y\in G\),
\[
 \Phi(T)
 =\widetilde U_y^{-1}\Phi(T)\widetilde U_y
 =\Phi(T^y),
 \qquad\text{so that}\qquad T=T^y.
\]
By the hypothesis of the lemma, \(T\) is a scalar operator, and
therefore \(S=\Phi(T)\) is a scalar operator.
\end{proof}

\noindent\textit{Bibliography.}\quad
\ArticleRef{65}, \ArticleRef{78}, \ArticleRef{210}, \ArticleRef{211},
\ArticleRef{224}, \ArticleRef{231}, \ArticleRef{232}, \ArticleRef{235},
\ArticleRef{236}, \ArticleRef{237}, \ArticleRef{246}, \ArticleRef{247},
\ArticleRef{275}, \ArticleRef{276}, \ArticleRef{290}, \ArticleRef{392},
\ArticleRef{441}.

\NumberedParagraph{3}{Examples from Measure Theory}\label{no:I-9-3}
Retain the notation of no.~2. We shall now place ourselves in a
situation in which the conditions of \cref{lem:I-9-2,lem:I-9-3} are
satisfied.

Let \(Z\) be a locally compact space countable at infinity, and let
\(\nu\) be a positive measure on \(Z\). Take for \(\mathfrak H\) the
Hilbert space \(L_{\mathbb C}^{2}(Z,\nu)\).

% Source PDF page 141 (printed p. 134)
\phantomsection\label{srcpage:134}
For every \(f\in L_{\mathbb C}^{\infty}(Z,\nu)\), let \(T_f\) be the
element of \(\mathcal L(\mathfrak H)\) defined by
\[
 T_f(g)=fg\qquad(g\in\mathfrak H).
\]
The set of the \(T_f\) is a maximal abelian von Neumann subalgebra
\(\mathcal A\) of \(\mathcal L(\mathfrak H)\)
(\hyperref[thm:I-7-2]{\S~7, Theorem~2}).

Let, on the other hand, \(G\) be a discrete group acting on the left
on \(Z\). For every \(s\in G\), suppose that the map
\(\zeta\mapsto s\zeta\) of \(Z\) onto \(Z\) is a homeomorphism which
transforms \(\nu\) into an equivalent measure. (Thus \(\nu\) is
``quasi-invariant'' under \(G\).) For every \(s\in G\), there is
therefore on \(Z\) a finite, strictly positive, locally
\(\nu\)-integrable function \(\zeta\mapsto r_s(\zeta)\) such that
\[
 d\nu(s\zeta)=r_s(\zeta)\,d\nu(\zeta).
\]
For \(s,t\in G\), one has
\begin{align*}
 r_{st}(\zeta)\,d\nu(\zeta)
 &=d\nu(st\zeta)
  =r_s(t\zeta)\,d\nu(t\zeta)\\
 &=r_s(t\zeta)r_t(\zeta)\,d\nu(\zeta),
\end{align*}
and hence
\[
 r_{st}(\zeta)=r_s(t\zeta)r_t(\zeta)
\]
almost everywhere on \(Z\). Taking \(t=s^{-1}\), one deduces
\[
 1=r_s(s^{-1}\zeta)r_{s^{-1}}(\zeta)
\]
almost everywhere on \(Z\).

For \(g\in\mathfrak H\) and \(s\in G\), put
\[
 (U_sg)(\zeta)
 =r_{s^{-1}}(\zeta)^{1/2}g(s^{-1}\zeta).
\]
Then
\begin{align*}
 \int_Z |(U_sg)(\zeta)|^2\,d\nu(\zeta)
 &=\int_Z r_{s^{-1}}(\zeta)|g(s^{-1}\zeta)|^2\,d\nu(\zeta)\\
 &=\int_Z|g(s^{-1}\zeta)|^2\,d\nu(s^{-1}\zeta)
  =\int_Z|g(\zeta)|^2\,d\nu(\zeta),
\end{align*}
so \(U_s\) is isometric. Moreover,
\begin{align*}
 (U_sU_tg)(\zeta)
 &=r_{s^{-1}}(\zeta)^{1/2}(U_tg)(s^{-1}\zeta)\\
 &=r_{s^{-1}}(\zeta)^{1/2}
   r_{t^{-1}}(s^{-1}\zeta)^{1/2}g(t^{-1}s^{-1}\zeta)\\
 &=r_{t^{-1}s^{-1}}(\zeta)^{1/2}g(t^{-1}s^{-1}\zeta)
  =(U_{st}g)(\zeta);
\end{align*}
thus \(s\mapsto U_s\) is a unitary representation of \(G\) on
\(\mathfrak H\). For \(f\in L_{\mathbb C}^{\infty}(Z,\nu)\), one
has, on the other hand,
\begin{align*}
 (U_s^{-1}T_fU_sg)(\zeta)
 &=r_s(\zeta)^{1/2}(T_fU_sg)(s\zeta)\\
 &=r_s(\zeta)^{1/2}f(s\zeta)(U_sg)(s\zeta)\\
 &=r_s(\zeta)^{1/2}f(s\zeta)
   r_{s^{-1}}(s\zeta)^{1/2}g(\zeta)\\
 &=f(s\zeta)g(\zeta).
\end{align*}
Thus, putting \(f_s(\zeta)=f(s\zeta)\),
\[
 U_s^{-1}T_fU_s=T_{f_s}.
\]
The conditions at the beginning of no.~2 are therefore satisfied, as
is the first hypothesis of \cref{lem:I-9-2}. We now make the
following additional hypothesis:

\begin{quote}
\textup{(\(\star\))} For every \(s\ne e\) in \(G\), and every
non-negligible measurable subset \(Z'\) of \(Z\), there exists a
non-negligible measurable subset \(Z''\subset Z'\) such that
\(Z''\cap sZ''=\varnothing\).
\end{quote}

% Source PDF page 142 (printed p. 135)
\phantomsection\label{srcpage:135}
We shall show that then
\[
 \mathcal A\cap\mathcal A U_y=0
 \qquad(y\in G,\ y\ne e).
\]
An element of \(\mathcal A\cap\mathcal A U_y\) has the form \(T_f\),
with \(f\in L_{\mathbb C}^{\infty}(Z,\nu)\), and the form
\(T_gU_y\), with \(g\in L_{\mathbb C}^{\infty}(Z,\nu)\). In
general, denote by \(\chi_Y\) the characteristic function of a subset
\(Y\) of \(Z\). Let \(X\) be the set of the \(\zeta\in Z\) such that
\(f(\zeta)\ne0\). If \(X\) is non-negligible, there exists a
non-negligible integrable subset \(X'\) of \(X\) such that
\(X'\cap yX'=\varnothing\). One has, almost everywhere,
\[
 f\chi_{X'}
 =T_f\chi_{X'}
 =T_gU_y\chi_{X'}
 =g(U_y\chi_{X'}).
\]
But \(U_y\chi_{X'}\) vanishes almost everywhere on \(X'\); hence
\(f\chi_{X'}=0\) almost everywhere, which is a contradiction. Thus
\(X\), and consequently \(f\), is negligible, so that \(T_f=0\).

Suppose finally that:
\begin{quote}
\textup{(\(\star\star\))} \(G\) is ergodic in \(Z\).
\end{quote}
Let \(T_f\) be an element of \(\mathcal A\) invariant under \(G\).
For every \(s\in G\), one has
\[
 T_{f_s}=U_s^{-1}T_fU_s=T_f,
\]
and therefore \(f=f_s\) almost everywhere. By the ergodicity of
\(G\), there is a constant equal to \(f\) almost everywhere. Hence
\(T_f\) is a scalar.

Thus, when hypotheses \textup{(\(\star\))} and
\textup{(\(\star\star\))} are satisfied,
\cref{lem:I-9-2,lem:I-9-3} apply, and \(\mathcal B\) is a factor.

The hypothesis that \(Z\) is countable at infinity serves only to
simplify the language a little.

\noindent\textit{Bibliography.}\quad
\ArticleRef{65}, \ArticleRef{78}.

\NumberedParagraph{4}{Existence of the Different Types of Factors}
\label{no:I-9-4}
We retain the notation of nos.~2 and~3, and suppose that hypotheses
\textup{(\(\star\))} and \textup{(\(\star\star\))} are satisfied.

\begin{proposition}\label{prop:I-9-3}
Suppose that \(\nu\) is invariant under \(G\), that
\(\nu(\{\zeta\})=0\) for every \(\zeta\in Z\), and that
\[
 0<\nu(Z)<+\infty
 \qquad\textup{[respectively, }\nu(Z)=+\infty\textup{]}.
\]
Then \(\mathcal B\) is a factor of type \(\mathrm{II}_1\)
[respectively, \(\mathrm{II}_\infty\)].
\end{proposition}

\begin{proof}
For every \(T_f\in\mathcal A^+\), put
\[
 \varphi(T_f)=\int_Z f(\zeta)\,d\nu(\zeta).
\]
It is clear that \(\varphi\) is a faithful trace on
\(\mathcal A^+\). The invariance of \(\nu\) under \(G\) implies the
invariance of \(\varphi\) under \(G\). The trace \(\varphi\) is
normal, because the isomorphism \(f\mapsto T_f\) is compatible with
the natural order relations in
\(L_{\mathbb C}^{\infty}(Z,\nu)\) and \(\mathcal A\); and, if
\((f_\alpha)\) is an increasing directed family of positive functions
in \(L_{\mathbb C}^{\infty}(Z,\nu)\), with supremum
\(f\in L_{\mathbb C}^{\infty}(Z,\nu)\) in the natural order of that
space, then
\[
 \int_Z f(\zeta)\,d\nu(\zeta)
\]
is the supremum of the
\(\int_Z f_\alpha(\zeta)\,d\nu(\zeta)\). If
\(\nu(Z)<+\infty\), it is clear
% Source PDF page 143 (printed p. 136)
\phantomsection\label{srcpage:136}
that \(\varphi\) is finite. If \(\nu(Z)=+\infty\), the function
\(1\) on \(Z\) is the supremum in
\(L_{\mathbb C}^{\infty}(Z,\nu)\) of an increasing directed family
of positive functions \(f\in L_{\mathbb C}^{\infty}(Z,\nu)\) such
that
\[
 \int_Z f(\zeta)\,d\nu(\zeta)<+\infty;
\]
hence \(\varphi\) is semifinite but not finite. It follows from
\cref{prop:I-9-1} that on \(\mathcal B^+\) there exists a faithful,
semifinite normal trace, finite if and only if
\(\nu(Z)<+\infty\). Since \(\mathcal B\) is a factor, we see that it
is semifinite, and finite if and only if \(\nu(Z)<+\infty\).

To finish proving the proposition, it is therefore enough to show
that \(\mathcal B\) is not discrete. Now the hypotheses
\(\nu(Z)>0\) and \(\nu(\{\zeta\})=0\) for every \(\zeta\in Z\)
imply that there exists a decreasing sequence \((Z_n)\) of measurable
subsets of \(Z\) such that
\[
 +\infty>\nu(Z_1)>\nu(Z_2)>\cdots
\]
(consider a point of the support of \(\nu\), and open neighbourhoods
of that point whose measures tend to zero). The \(\chi_{Z_n}\) define
strictly decreasing projections \(E_n\) of \(\mathcal A\) such that
\(\varphi(E_n)<+\infty\). The \(\Phi(E_n)\) form a strictly decreasing
sequence of projections of \(\mathcal B\) such that
\[
 \psi\bigl(\Phi(E_n)\bigr)<+\infty.
\]
Thus \(\mathcal B\) is not discrete
(\hyperref[prop:I-8-3]{\S~8, Proposition~3}).
\end{proof}

\begin{theorem}\label{thm:I-9-1}
There exist factors of type \(\mathrm{II}_1\), and factors of type
\(\mathrm{II}_\infty\).
\end{theorem}

\begin{proof}
Take for \(Z\) the one-dimensional torus (respectively, the real
line) with its usual topology, for \(\nu\) Haar measure on \(Z\), and
for \(G\) an everywhere dense subgroup of \(Z\), regarded as a
discrete group. Let \(G\) act on \(Z\) by translation. The measure
\(\nu\) is invariant under \(G\), one has
\(\nu(\{\zeta\})=0\) for every \(\zeta\in G\), and
\[
 0<\nu(Z)<+\infty
 \qquad\textup{[respectively, }\nu(Z)=+\infty\textup{]}.
\]
It is classical that \(G\) is ergodic in \(Z\). Finally, let
\(y\ne e\) be an element of \(G\), and let \(Z'\) be a
non-negligible measurable subset of \(Z\). There exists a
neighbourhood \(V\) of \(e\) in \(Z\), small enough that
\(V\cap yV=\varnothing\). Then, if \(s\in G\) is such that
\(Z''=Z'\cap sV\) is non-negligible, one has
\(Z''\cap yZ''=\varnothing\). The factor \(\mathcal B\) is therefore
of type \(\mathrm{II}_1\) (respectively, \(\mathrm{II}_\infty\)).
\end{proof}

We return to the general situation of no.~3, with hypotheses
\textup{(\(\star\))} and \textup{(\(\star\star\))}.

\begin{proposition}\label{prop:I-9-4}
If the subgroup \(G_0\) of \(G\), formed by the elements which
preserve \(\nu\), is distinct from \(G\) and is ergodic in \(Z\),
then \(\mathcal B\) is a factor of type \(\mathrm{III}\).
\end{proposition}

\begin{proof}
For every \(T_f\in\mathcal A^+\), put
\[
 \varphi(T_f)=\int_Z f(\zeta)\,d\nu(\zeta).
\]
As in the proof of \cref{prop:I-9-3}, one sees that \(\varphi\) is a
faithful, semifinite normal trace on \(\mathcal A^+\), invariant under
\(G_0\). Arguing
% Source PDF page 144 (printed p. 137)
\phantomsection\label{srcpage:137}
by contradiction, suppose that \(\mathcal B\) is semifinite. There is
then on \(\mathcal A^+\) a faithful, semifinite normal trace
\(\varphi'\) invariant under \(G\) (\cref{prop:I-9-2}). Put
\(\omega=\varphi+\varphi'\). There exist
\(T_g,T_{g'}\in\mathcal A^+\) such that
\[
 \varphi(T_f)=\omega(T_fT_g),
 \qquad
 \varphi'(T_f)=\omega(T_fT_{g'})
\]
for every \(T_f\in\mathcal A^+\)
(\hyperref[thm:I-6-3]{\S~6, Theorem~3}). For \(s\in G_0\), one has
\begin{align*}
 \omega(T_fT_{g_s})
 &=\omega(T_fU_s^{-1}T_gU_s)
  =\omega(U_sT_fU_s^{-1}T_g)\\
 &=\varphi(T_{f_{s^{-1}}})
  =\varphi(T_f)
  =\omega(T_fT_g)
\end{align*}
for every \(T_f\in\mathcal A^+\); hence \(g_s=g\) almost everywhere.
Since \(G_0\) is ergodic, it follows that \(g\) is almost everywhere
equal to a constant \(>0\); the same is true of \(g'\). Hence
\(\varphi\) is proportional to \(\varphi'\), and consequently is
invariant under \(G\), contrary to the hypotheses.
\end{proof}

\begin{theorem}\label{thm:I-9-2}
There exist factors of type \(\mathrm{III}\).
\end{theorem}

\begin{proof}
Take for \(Z\) the real line, for \(\nu\) Lebesgue measure, and for
\(G\) the group of transformations \(\zeta\mapsto a\zeta+b\) of
\(Z\), where \(a\) and \(b\) are rational and \(a>0\). The subgroup
\(G_0\) is obtained by taking \(a=1\); it is ergodic in \(Z\). It is
easy to see that hypothesis \textup{(\(\star\))} is satisfied. It is
then enough to apply \cref{prop:I-9-4}.
\end{proof}

\begin{remark*}
The factors just constructed, of types \(\mathrm{II}_1\),
\(\mathrm{II}_\infty\), and \(\mathrm{III}\), act on Hilbert spaces with a
countable basis, provided that in the proof of \cref{thm:I-9-1} the
group \(G\) is chosen countable.
\end{remark*}

There are other methods of constructing factors. (Cf. Chapter~III,
\S~5, Exercise~8; Chapter~III, \S~7, no.~6; and \ArticleRef{77},
\ArticleRef{78}, \ArticleRef{296}, \ArticleRef{315}, \ArticleRef{316},
\ArticleRef{317}.) The theory of \(C^*\)-algebras also supplies
examples.

\noindent\textit{Bibliography.}\quad
\ArticleRef{13}, \ArticleRef{65}, \ArticleRef{78}.

We shall see (Chapter~II, \S~6, no.~2) that the study of arbitrary von
Neumann algebras reduces, to a certain extent, to that of factors.
Factors of type \(\mathrm I\) are well known
(\hyperref[cor:I-8-thm1-3]{\S~8, Theorem~1, Corollary~3}). Factors of
type \(\mathrm{II}_\infty\) are tensor products of factors of type
\(\mathrm{II}_1\) and factors of type \(\mathrm I\) (Chapter~III,
\S~8, Exercise~11). Unfortunately, we do not have a classification of
factors of type \(\mathrm{II}_1\) or \(\mathrm{III}\).
\emph{This is one of the principal problems of the theory.} In
Hilbert spaces with a countable basis, nine pairwise nonisomorphic
factors of type \(\mathrm{II}_1\) are known (Chapter~III, \S~7,
no.~7, Theorem~4; and \ArticleRef{309}, \ArticleRef{441},
\ArticleRef{443}, \ArticleRef{463}, \ArticleRef{474}), as are
continuous families of pairwise nonisomorphic factors of type
\(\mathrm{III}\) (\ArticleRef{171}, \ArticleRef{310},
\ArticleRef{411}, \ArticleRef{435}).

\noindent\textit{Exercises.}---
\begin{enumerate}[label=\arabic*.,leftmargin=*,itemsep=0.5\baselineskip]
\item\label{ex:I-9-1}
Adopt the notation of no.~2. Let
\(\Phi'(T),\widetilde U'_y,W\) be the elements of
\(\mathcal L(\widetilde{\mathfrak H})\) defined by the matrices
\[
 \bigl(\delta_{s,t}T^{s^{-1}}\bigr),
 \qquad
 \bigl(\delta_{sy^{-1},t}I_{\mathfrak H}\bigr),
 \qquad
 \bigl(\delta_{s,t^{-1}}U_s\bigr),
\]
where \(\delta_{s,t}\) is the Kronecker symbol.
\begin{enumerate}[label=\textit{\alph*.},leftmargin=2em]
\item Show that \(y\mapsto\widetilde U_y'{}^*\) is a unitary
representation of \(G\) on \(\widetilde{\mathfrak H}\), and that
\(W\) is a unitary operator such that
\[
 W^2=1,
 \qquad
 W\Phi(T)W=\Phi'(T),
 \qquad
 W\widetilde U_yW=\widetilde U_y'{}^*;
\]
show also that
\[
 \Phi'(T)\in\mathcal B',
 \qquad
 \widetilde U'_y\in\mathcal B'.
\]

% Source PDF page 145 (printed p. 138)
\phantomsection\label{srcpage:138}
\item Henceforth suppose that \(\mathcal A'=\mathcal A\). Let
\(S\in\mathcal L(\widetilde{\mathfrak H})\) commute with the
\(\widetilde U'_y\) and the \(\Phi'(T)\). Show that \(S\) and
\(WSW\) are represented by the matrices
\[
 \bigl(T_{st^{-1}}U_{st^{-1}}\bigr)
 \quad\text{and}\quad
 \bigl(U_sT_{s^{-1}t}U_s^{-1}\bigr),
\]
where \(T_y\in\mathcal A\) for every \(y\in G\).

\item Show that two elements of
\(\mathcal L(\widetilde{\mathfrak H})\), one commuting with the
\(\Phi(T)\) and the \(\widetilde U_y\), the other commuting with the
\(\Phi'(T)\) and the \(\widetilde U'_y\), commute with each other
(use \textit{b}). Deduce that \(\mathcal B'\) is the von Neumann
algebra generated by the \(\Phi'(T)\) and the
\(\widetilde U'_y\), and that
\[
 \mathcal B'=W\mathcal BW.
\]
\ArticleRef{13}, \ArticleRef{65}, \ArticleRef{78}.
\end{enumerate}

\item\label{ex:I-9-2}
Adopt the notation of no.~2, and suppose that the conditions of
\cref{lem:I-9-2} are satisfied. Show that, if \(T\) is a minimal
projection of \(\mathcal A\), then \(\Phi(T)\) is a minimal
projection of \(\mathcal B\). [A projection of \(\mathcal B\)
majorized by \(\Phi(T)\) commutes with
\(\widetilde{\mathcal A}\).] Deduce examples in which the construction
of no.~3 leads to factors of type \(\mathrm I\). \ArticleRef{65}

\item\label{ex:I-9-3}
Show that there exist two involutive algebras of operators
\(\mathcal B\) and \(\mathcal C\) on a complex Hilbert space, isomorphic as
involutive algebras, whose weak closures are respectively finite and
properly infinite. (Take for \(\mathcal B\) and \(\mathcal C\) the
algebras of \hyperref[ex:I-2-6]{\S~2, Exercise~6b}, with
\(\mathcal A\) and \(\mathcal A'\) finite factors on an
infinite-dimensional Hilbert space.)

\item\label{ex:I-9-4}
Show that there exists a factor of type \(\mathrm{II}_1\) having a
totalizing element and acting on a Hilbert space which does not have a
countable basis. (In the proof of \cref{thm:I-9-1}, take \(G\)
uncountable. The totalizing element is the image under one of the
\(J_s\) of the function \(1\in\mathfrak H\).) Show that this factor is
of countable type (as is every finite factor), but is not generated by
a countable family of elements. (A von Neumann algebra generated by a
countable family of elements and possessing a countable totalizing
set acts on a Hilbert space with a countable basis.)

\item\label{ex:I-9-5}
In \hyperref[ex:I-7-3]{\S~7, Exercise~3}, we studied the following
properties which a von Neumann algebra may possess: (i) to be of
countable type; (ii) to be generated by a countable family of
elements. We now propose to study neighbouring properties.

Let \(\mathcal A\) be a von Neumann algebra on \(\mathfrak H\), and
let \(\mathcal Z\) be its centre.
\begin{enumerate}[label=\textit{\alph*.},leftmargin=2em]
\item The following conditions are equivalent:
\begin{enumerate}[label=\textup{(\roman*)},leftmargin=2.8em]
\item For every projection \(E\) of \(\mathcal Z\) such that
\(\mathcal Z_E\) is of countable type, \(\mathcal A_E\) is of
countable type;
\item \(\mathcal A\) is a product of von Neumann algebras of
countable type;
\item For every decomposition
\[
 \mathcal Z=\prod_{i\in I}\mathcal Z_{E_i}
\]
such that the \(\mathcal Z_{E_i}\) are of countable type, the
\(\mathcal A_{E_i}\) are of countable type.
\end{enumerate}
[Show that \textup{(i)}\(\Rightarrow\)\textup{(iii)}
\(\Rightarrow\)\textup{(ii)}\(\Rightarrow\)\textup{(i)}.] One then
says that \(\mathcal A\) is \emph{of countable type over its centre}.

\item If \(\mathcal A\) is the von Neumann algebra generated by
\(\mathcal Z\) and a countable family of elements, then
\(\mathcal A\) is of countable type over its centre. (Reduce to the
case in which \(\mathcal Z\) has a separating element \(x\). Then
\(\mathcal A\) is isomorphic to the algebra induced on
\(\mathfrak X_x^{\mathcal A}\), and we may therefore suppose that
\(x\) is totalizing for \(\mathcal A\). By the hypothesis on
\(\mathcal A\), there exists a sequence \(T_1,T_2,\ldots\) of
elements of \(\mathcal A\) such that the vectors
\(TT_ix\), \(T\in\mathcal Z\), form a total set in \(\mathfrak H\).
Then the \(T_ix\) form a separating set for \(\mathcal Z'\).
\emph{A fortiori}, \(\mathcal A\subset\mathcal Z'\) is of countable
type.)

\item Show that there exist von Neumann algebras of countable type
over their centre which are not generated by their centre and a
countable family of elements (use Exercise~4). (Cf., however,
Chapter~III, \S~3, Exercise~1.)
\end{enumerate}
\end{enumerate}
% Source PDF page 146 (printed p. 139)
\chapter{Reduction of von Neumann Algebras}\label{ch:II}\label{srcpage:139}

\section{Fields of Hilbert Spaces}\label{sec:II-1}

\NumberedParagraph{1}{Borel spaces, measures}\label{no:II-1-1}
A \emph{Borel space} is a set \(E\) furnished with a set \(\mathcal B\) of
subsets of \(E\) having the following properties:
\(\varnothing\in\mathcal B\), and \(\mathcal B\) is closed under countable
unions and passage to the complement (and hence under countable
intersections). The elements of \(\mathcal B\) are called the \emph{Borel
subsets} of \(E\).

Let \(E,F\) be two Borel spaces. A mapping \(f\) of \(E\) into \(F\) is said
to be \emph{Borel} if the inverse image under \(f\) of every Borel subset of
\(F\) is a Borel subset of \(E\).

Let \(E\) be a Borel space and \(E'\) a subset of \(E\). The intersections
with \(E'\) of the Borel subsets of \(E\) define on \(E'\) a Borel
structure, said to be \emph{induced} by the Borel structure of \(E\).

Let \(E\) be a topological space. The subsets of \(E\) which are Borel for
the topology define on \(E\) a Borel structure said to \emph{underlie} the
topology.

A Borel space \(E\) is said to be \emph{discrete} if every subset of \(E\)
is Borel.

A Borel space is said to be \emph{standard} if its Borel structure
underlies a Polish-space topology (for example, the topology of a
second-countable locally compact space). It is clear that, if \(E\) is a
countable Borel space, to say that \(E\) is standard is equivalent to saying
that \(E\) is discrete. If \(E\) is an uncountable standard Borel space,
\(E\) is isomorphic to the Borel space underlying the topological space
\([0,1]\). A Borel subset of a standard Borel space is a standard Borel space
(cf. C.~Kuratowski, \emph{Topologie}, I, \emph{Espaces m\'etrisables,
espaces complets}, 2nd ed., p.~358, Remark~1).

Let \(Z\) be a Borel space and \(\mathcal B\) the set of Borel subsets of
\(Z\). In this book a \emph{positive measure} on \(Z\) means a mapping
\(\nu\) of \(\mathcal B\) into \([0,+\infty]\) such that: \(1^\circ\), if
\(Z_1,Z_2,\ldots\) are pairwise disjoint elements of \(\mathcal B\), then
\[
  \nu(Z_1\cup Z_2\cup\cdots)=\nu(Z_1)+\nu(Z_2)+\cdots;
\]
\(2^\circ\), \(Z\) is the union of a sequence of Borel subsets
\(Y_1,Y_2,\ldots\) such that \(\nu(Y_i)<+\infty\) for every \(i\). A subset
of \(Z\) is said to be \(\nu\)-\emph{negligible} if it is contained in a
set \(Y\in\mathcal B\) such that \(\nu(Y)=0\). A subset of \(Z\) is said to
be \(\nu\)-\emph{measurable} if it is of the form \(X\cup N\), where
\(X\in\mathcal B\) and \(N\) is
% Source PDF page 147 (printed p. 140)
\phantomsection\label{srcpage:140}
\(\nu\)-negligible. The set \(\mathcal M\) of \(\nu\)-measurable subsets is
closed under countable unions, countable intersections, and passage to the
complement; putting \(\nu(X\cup N)=\nu(X)\) gives an extension of \(\nu\)
to \(\mathcal M\) which still satisfies the axioms of a measure. We shall
assume the theory of measurable functions, integrable functions, and so on,
relative to \(\nu\), to be known.

A positive measure \(\nu\) on \(Z\) is said to be \emph{standard} if there
exists a \(\nu\)-negligible subset \(N\) of \(Z\) such that the Borel space
\(Z\setminus N\) is standard.

If \(Z\) is a locally compact, countable-at-infinity space, a positive
(Radon) measure on \(Z\), regarded as a function on the set of Borel subsets
of \(Z\), is a positive measure in the preceding sense. When \(Z\) is
second-countable, this measure is standard.

\NumberedParagraph{2}{Fields of vectors}\label{no:II-1-2}
Let \(Z\) be a Borel space and \(\nu\) a positive measure on \(Z\). Whenever
this can be done without risk of confusion, we shall say ``measurable,''
``integrable,'' and so forth, instead of ``\(\nu\)-measurable,''
``\(\nu\)-integrable,'' and so forth.

A \emph{field of complex Hilbert spaces} on \(Z\) is a mapping
\(\zeta\mapsto\mathfrak H(\zeta)\), defined on \(Z\), such that
\(\mathfrak H(\zeta)\) is a complex Hilbert space for every \(\zeta\in Z\).
Then
\[
  \mathfrak F=\prod_{\zeta\in Z}\mathfrak H(\zeta)
\]
is a complex vector space; an element \(x\) of \(\mathfrak F\) is a mapping
\(\zeta\mapsto x(\zeta)\), defined on \(Z\), such that
\(x(\zeta)\in\mathfrak H(\zeta)\) for every \(\zeta\in Z\); such a mapping
is called a \emph{field of vectors} on \(Z\). If \(Y\) is a subset of \(Z\),
an element of \(\prod_{\zeta\in Y}\mathfrak H(\zeta)\) is called a field of
vectors on \(Y\).

\begin{lemma}\label{lem:II-1-1}
Let \(x_1,x_2,\ldots\) be a sequence of fields of vectors such that the
functions
\[
  \zeta\longmapsto (x_m(\zeta)\mid x_n(\zeta))
\]
are measurable. For \(\zeta\in Z\), let \(\mathfrak X(\zeta)\) be the
complex vector space algebraically spanned by the \(x_n(\zeta)\), and let
\(d(\zeta)\) be the (algebraic) dimension of \(\mathfrak X(\zeta)\). Then
the set \(Z_p\) of \(\zeta\in Z\) such that
\(d(\zeta)=p\), \(p=0,1,2,\ldots,\aleph_0\), is measurable. There exists a
sequence \(y_1,y_2,\ldots\) of fields of vectors such that:
\begin{enumerate}[label=\textup{(\roman*)}]
\item for every \(\zeta\in Z\), the vectors
  \(y_1(\zeta),y_2(\zeta),\ldots\) algebraically span
  \(\mathfrak X(\zeta)\);
\item if \(d(\zeta)=\aleph_0\), then
  \(y_1(\zeta),y_2(\zeta),\ldots\) form an orthonormal system; if
  \(d(\zeta)<\aleph_0\), then
  \(y_1(\zeta),\ldots,y_{d(\zeta)}(\zeta)\) form an orthonormal system and
  \(y_n(\zeta)=0\) for \(n>d(\zeta)\);
\item for each field \(y_n\), there is a partition of \(Z\) into
  measurable subsets \(Z_1,Z_2,\ldots\) such that, on each \(Z_k\),
  \(y_n\) has the form
  \[
    \zeta\longmapsto\sum_i f_i(\zeta)x_i(\zeta),
  \]
  where the \(f_i\) are measurable complex-valued functions which vanish
  identically for all sufficiently large \(i\).
\end{enumerate}
\end{lemma}

% Source PDF page 148 (printed p. 141)
\phantomsection\label{srcpage:141}
\begin{proof}
By induction, first define a sequence \(z_1,z_2,\ldots\) of fields of
vectors as follows. Let \(\zeta\in Z\); take for \(z_k(\zeta)\) the first
of the vectors \(x_1(\zeta),x_2(\zeta),\ldots\) which is linearly
independent of the \(z_i(\zeta)\) with \(i<k\), if such a vector exists;
otherwise put \(z_k(\zeta)=0\). Then: \(1^\circ\), for every
\(\zeta\in Z\), the vectors \(z_1(\zeta),z_2(\zeta),\ldots\) span the
vector space \(\mathfrak X(\zeta)\); \(2^\circ\), if
\(d(\zeta)=\aleph_0\), the vectors \(z_1(\zeta),z_2(\zeta),\ldots\) are
linearly independent; \(3^\circ\), if \(d(\zeta)<\aleph_0\), then
\(z_i(\zeta)=0\) for \(i>d(\zeta)\), and the \(z_i(\zeta)\) with
\(i\leq d(\zeta)\) are linearly independent.

We prove the following assertion, where \(j\) denotes an integer \(>0\):

\smallskip
\noindent\((A_j)\) There exists a sequence \(Y_1,Y_2,\ldots\) of disjoint
measurable subsets of \(Z\), covering \(Z\), such that, for every \(i\leq j\)
and every \(k=1,2,\ldots\), either \(z_i(\zeta)=0\) for all
\(\zeta\in Y_k\), or \(z_i(\zeta)=x_{n_i}(\zeta)\neq0\) for all
\(\zeta\in Y_k\), with an index \(n_i\) independent of \(\zeta\).

Suppose \((A_j)\) established for \(j<l\), and prove \((A_l)\). Let \(Z'\)
be a measurable subset of \(Z\) such that, for \(i<l\), on \(Z'\) one has
\(z_i(\zeta)=x_{n_i}(\zeta)\neq0\), with an index \(n_i\) independent of
\(\zeta\). Then the set \(X_q\) of \(\zeta\in Z'\) such that
\(z_l(\zeta)=x_q(\zeta)\neq0\) is defined by the vanishing of the Gram
determinants
\[
  \Delta\bigl(x_{n_1}(\zeta),x_{n_2}(\zeta),\ldots,
  x_{n_{l-1}}(\zeta),x_m(\zeta)\bigr)
\]
of
\[
  x_{n_1}(\zeta)=z_1(\zeta),\quad
  x_{n_2}(\zeta)=z_2(\zeta),\quad\ldots,\quad
  x_{n_{l-1}}(\zeta)=z_{l-1}(\zeta),\quad x_m(\zeta),
  \qquad m<q,
\]
and by the condition
\[
  \Delta\bigl(x_{n_1}(\zeta),\ldots,x_{n_{l-1}}(\zeta),
  x_q(\zeta)\bigr)\neq0.
\]
Thus \(X_q\) is measurable, whence \((A_l)\) follows at once.

On each of the subsets \(Y_1,Y_2,\ldots\) in \((A_j)\),
\(d(\zeta)\) is either constant or at least \(j\). The set of \(\zeta\) for
which \(d(\zeta)<j\) is therefore measurable, proving the first assertion
of the lemma. Now, at each point \(\zeta\in Z\), orthonormalize the
sequence of nonzero \(z_i(\zeta)\) by the Gram--Schmidt process. This gives
vectors \(y_i(\zeta)\). If \(p<\aleph_0\), put \(y_i(\zeta)=0\) for
\(i>p\). The fields \(y_i\) plainly have properties (i) and (ii) of the
lemma. Suppose property (iii) established for indices \(n<j\), and prove
it for \(y_j\), using \((A_j)\). If \(z_j(\zeta)=0\) on \(Y_k\), then
\(y_j(\zeta)=0\) on \(Y_k\). If
\(z_j(\zeta)=x_{n_j}(\zeta)\neq0\) on \(Y_k\), then
\[
  y_j(\zeta)=\lVert u(\zeta)\rVert^{-1}u(\zeta),
  \qquad
  u(\zeta)=z_j(\zeta)-\sum_{n<j}
    (z_j(\zeta)\mid y_n(\zeta))y_n(\zeta),
\]
and our induction hypothesis, together with the measurability of the
functions
\(\zeta\mapsto(x_p(\zeta)\mid x_q(\zeta))\), gives property (iii) for the
index \(j\).
\end{proof}

\begin{remark*}
The statement and proof of the lemma remain valid if the word
``measurable'' is replaced throughout by the word ``Borel.''
\end{remark*}

Bibliography: \ArticleRef{29}, \ArticleRef{80}.

% Source PDF page 149 (printed p. 142)
\phantomsection\label{srcpage:142}
\NumberedParagraph{3}{Measurable fields of Hilbert spaces}
\label{no:II-1-3}
Let \(\zeta\mapsto\mathfrak H(\zeta)\) be a field of complex Hilbert spaces
on \(Z\), and let \(\mathfrak F\) be the vector space of fields of vectors.

\begin{definition}\label{def:II-1-1}
The \(\mathfrak H(\zeta)\) are said to form a \(\nu\)-\emph{measurable field
of complex Hilbert spaces} if a vector subspace \(\mathfrak G\) of
\(\mathfrak F\) is specified having the following properties:
\begin{enumerate}[label=\textup{(\roman*)}]
\item for every \(x\in\mathfrak G\), the function
  \(\zeta\mapsto\lVert x(\zeta)\rVert\) is \(\nu\)-measurable;
\item if \(y\in\mathfrak F\) is such that, for every
  \(x\in\mathfrak G\), the complex-valued function
  \(\zeta\mapsto(x(\zeta)\mid y(\zeta))\) is \(\nu\)-measurable, then
  \(y\in\mathfrak G\);
\item there exists a sequence \(x_1,x_2,\ldots\) of elements of
  \(\mathfrak G\) such that, for every \(\zeta\in Z\), the
  \(x_n(\zeta)\) form a total sequence in \(\mathfrak H(\zeta)\).
\end{enumerate}
\end{definition}

The fields of vectors belonging to \(\mathfrak G\) are then called
\(\nu\)-\emph{measurable fields of vectors}. A sequence
\(x_1,x_2,\ldots\) of \(\nu\)-measurable fields of vectors having property
(iii) is called a \emph{fundamental sequence of \(\nu\)-measurable fields
of vectors}.

Property (iii) implies that the \(\mathfrak H(\zeta)\) have countable
bases.

If \(x\) and \(y\) are measurable fields of vectors, then
\((x(\zeta)\mid y(\zeta))\), which is a linear combination with constant
coefficients of
\[
  \lVert(x+y)(\zeta)\rVert^2,
  \quad \lVert(x-y)(\zeta)\rVert^2,
  \quad \lVert(x+iy)(\zeta)\rVert^2,
  \quad \lVert(x-iy)(\zeta)\rVert^2,
\]
is a measurable function of \(\zeta\). Hence, by property (ii), the product
of a measurable field of vectors by a measurable complex-valued function
is a measurable field; the limit of a sequence of measurable fields of
vectors which converges weakly at every point of \(Z\) is a measurable
field.

\begin{remarks*}
1. By multiplying the fields \(x_i\) of Definition~1 by measurable
functions, one may suppose that each function
\(\zeta\mapsto\lVert x_i(\zeta)\rVert\) is bounded and vanishes outside a
set of finite measure. By adjoining to the \(x_i\) their linear combinations
with rational complex coefficients, one obtains a sequence
\(y_1,y_2,\ldots\) of measurable fields of vectors such that, in addition,
for every \(\zeta\in Z\), the \(y_n(\zeta)\) are dense in
\(\mathfrak H(\zeta)\).

2. If \(\nu_1\) is a positive measure on \(Z\) equivalent to \(\nu\), the
field of the \(\mathfrak H(\zeta)\), furnished with the space
\(\mathfrak G\), is a \(\nu_1\)-measurable field of Hilbert spaces. In
other words, the notion of a \(\nu\)-measurable field of Hilbert spaces
involves only the measure class of \(\nu\).
\end{remarks*}

\noindent\emph{Examples.}---1. Suppose that \(Z\) is discrete. Every
complex-valued function on \(Z\) is measurable. The only subspace
\(\mathfrak G\) of \(\mathfrak F\) having the properties of Definition~1
is \(\mathfrak F\) itself.

% Source PDF page 150 (printed p. 143)
\phantomsection\label{srcpage:143}
2. Let \(\mathfrak H_0\) be a complex Hilbert space with a countable base.
The \emph{constant field corresponding to \(\mathfrak H_0\) on \(Z\)} is
the \(\nu\)-measurable field defined as follows: (a)
\(\mathfrak H(\zeta)=\mathfrak H_0\) for every \(\zeta\in Z\); (b) the
\(\nu\)-measurable fields of vectors are the \(\nu\)-measurable mappings
of \(Z\) into \(\mathfrak H_0\). (Since \(\mathfrak H_0\) has a countable
base, the measurable mappings of \(Z\) into \(\mathfrak H_0\) are the same
whether \(\mathfrak H_0\) is given the strong or the weak topology.)
Property (i) of Definition~1 is satisfied. Moreover, the constant mappings
of \(Z\) into \(\mathfrak H_0\) are measurable; hence, if a field of vectors
\(\zeta\mapsto x(\zeta)\) satisfies property (ii) of Definition~1, then
\((x(\zeta)\mid a)\) depends measurably on \(\zeta\) for every
\(a\in\mathfrak H_0\), and therefore \(x\) is measurable. Finally, if
\((a_n)\) is a total sequence in \(\mathfrak H_0\), the constant fields
\(\zeta\mapsto a_n\) form a fundamental sequence of measurable fields of
vectors.

Let \(Y\) be a measurable subset of \(Z\). A field of vectors \(x\) on \(Y\)
is said to be measurable if it can be extended to a measurable field of
vectors on \(Z\). For this it is necessary and sufficient that, for every
measurable field of vectors \(y\) on \(Z\), the function
\(\zeta\mapsto(x(\zeta)\mid y(\zeta))\) (defined on \(Y\)) be measurable.
The condition is plainly necessary. Conversely, suppose it satisfied, and
extend the field \(x\) to \(Z\) by putting \(x(\zeta)=0\) for
\(\zeta\notin Y\). The field so extended is measurable by
Definition~1(ii). The measurable fields of vectors on \(Y\) define on the
field \((\mathfrak H(\zeta))_{\zeta\in Y}\) a measurable-field structure;
this field is denoted by \(\mathcal E\vert Y\) (if \(\mathcal E\) denotes
the given field on \(Z\)) and is called the field \emph{induced} by
\(\mathcal E\) on \(Y\).

If \(Z\) is the union of a sequence \(Z_1,Z_2,\ldots\) of measurable
subsets, a field of vectors on \(Z\) is measurable if and only if its
restriction to each \(Z_i\) is measurable.

Let \(Z,Z'\) be Borel spaces; let \(\nu,\nu'\) be positive measures on
\(Z,Z'\); and let
\[
  \mathcal E=((\mathfrak H(\zeta)),\mathfrak G),
  \qquad
  \mathcal E'=((\mathfrak H'(\zeta')),\mathfrak G')
\]
be \(\nu\)-measurable and \(\nu'\)-measurable fields of Hilbert spaces.
Let \(\eta:Z\to Z'\) be an isomorphism of Borel spaces which transforms
\(\nu\) into \(\nu'\). An \(\eta\)-\emph{isomorphism} of \(\mathcal E\)
onto \(\mathcal E'\) is a family \((V(\zeta))_{\zeta\in Z}\) with the
following properties: (i) for every \(\zeta\in Z\), \(V(\zeta)\) is an
isomorphism of the Hilbert space \(\mathfrak H(\zeta)\) onto the Hilbert
space \(\mathfrak H'(\eta(\zeta))\); (ii) a field
\(\zeta\mapsto x(\zeta)\in\mathfrak H(\zeta)\) belongs to
\(\mathfrak G\) if and only if the field
\(\eta(\zeta)\mapsto V(\zeta)x(\zeta)\in
\mathfrak H'(\eta(\zeta))\) belongs to \(\mathfrak G'\).

If \(Z=Z'\) and \(\nu=\nu'\), an \emph{isomorphism} of \(\mathcal E\)
onto \(\mathcal E'\) means an \(\eta\)-isomorphism, where \(\eta\) is the
identity mapping. A field isomorphic to a constant field is called
\emph{trivial}.

\par\addvspace{6pt}\noindent\textbf{Remark 3.}\ ---\enspace
If in Definition~1 ``\(\nu\)-measurable'' is replaced throughout by
``Borel,'' one obtains the definition of a \emph{Borel field of Hilbert
spaces on \(Z\)}. (This definition involves only the
% Source PDF page 151 (printed p. 144)
\phantomsection\label{srcpage:144}
Borel structure of \(Z\), and not \(\nu\).) The definitions and properties
of this numbered paragraph carry over upon replacing ``measurable''
throughout by ``Borel.'' Let
\(\mathcal E=((\mathfrak H(\zeta)),\mathfrak G)\) be a Borel field of
Hilbert spaces on \(Z\), and let \(\nu\) be a measure on \(Z\). A field of
vectors \(\zeta\mapsto x(\zeta)\in\mathfrak H(\zeta)\) on \(Z\) is said to
be \(\nu\)-measurable \emph{relative to \(\mathcal E\)} if, for every
\(y\in\mathfrak G\), the function
\(\zeta\mapsto(x(\zeta)\mid y(\zeta))\) is \(\nu\)-measurable. Let
\(\mathfrak G'\) be the set of \(\nu\)-measurable fields of vectors on
\(Z\). It is readily seen that \(\mathfrak G'\) has properties (i), (ii),
and (iii) of Definition~1. Thus
\(((\mathfrak H(\zeta)),\mathfrak G')\) is a \(\nu\)-measurable field of
Hilbert spaces, said to be \emph{derived} from \(\mathcal E\). Every element
of \(\mathfrak G'\) is equal almost everywhere to an element of
\(\mathfrak G\) (this follows at once from the analogous property for
measurable numerical functions).
\par\addvspace{6pt}

If condition (iii) is omitted from Definition~1, one comes to a halt as
soon as one tries to prove the nontrivial results of the theory.

Bibliography: \ArticleRef{26}, \ArticleRef{28}, \ArticleRef{48},
\ArticleRef{52}, \ArticleRef{80}, \ArticleRef{100}, \ArticleRef{117},
\ArticleRef{118}, \ArticleRef{145}, \ArticleRef{193}, \ArticleRef{194},
\ArticleRef{205}, \ArticleRef{273}, \ArticleRef{329}, \ArticleRef{375}.

\NumberedParagraph{4}{First properties of measurable fields of Hilbert
spaces}\label{no:II-1-4}
Let \(Z\) be a Borel space, \(\nu\) a positive measure on \(Z\), and
\(\zeta\mapsto\mathfrak H(\zeta)\) a \(\nu\)-measurable field of Hilbert
spaces on \(Z\).

\begin{proposition}\label{prop:II-1-1}
\begin{enumerate}[label=\textup{(\roman*)}]
\item The set \(Z_p\) of \(\zeta\in Z\) for which the Hilbert dimension
  \(d(\zeta)\) of \(\mathfrak H(\zeta)\) is equal to \(p\) is measurable.
\item There exists a sequence \(y_1,y_2,\ldots\) of measurable fields of
  vectors with the following properties:
  \begin{enumerate}[label=\textup{\alph*.}]
  \item if \(d(\zeta)=\aleph_0\), then
    \(y_1(\zeta),y_2(\zeta),\ldots\) is an orthonormal basis of
    \(\mathfrak H(\zeta)\);
  \item if \(d(\zeta)<\aleph_0\), then
    \(y_1(\zeta),y_2(\zeta),\ldots,y_{d(\zeta)}(\zeta)\) is an
    orthonormal basis of \(\mathfrak H(\zeta)\), and
    \(y_i(\zeta)=0\) for \(i>d(\zeta)\).
  \end{enumerate}
\end{enumerate}
\end{proposition}

\begin{proof}
Let \(x_1,x_2,\ldots\) be a fundamental sequence of measurable fields of
vectors. Observe that \(d(\zeta)\) is also the algebraic dimension of the
vector space spanned by the \(x_i(\zeta)\). It is then enough to apply
Lemma~1 to the \(x_i\). The resulting \(y_i\) are measurable by part (iii)
of Lemma~1.
\end{proof}

\begin{definition}\label{def:II-1-2}
A sequence \(y_1,y_2,\ldots\) of measurable fields of vectors having the
properties of Proposition~1(ii) is called a \emph{measurable field of
orthonormal bases}.
\end{definition}

\begin{proposition}\label{prop:II-1-2}
Let \(x_1,x_2,\ldots\) be a fundamental sequence of measurable fields of
vectors. A field of vectors \(x\) on \(Z\) is measurable if and only if the
functions \(\zeta\mapsto(x(\zeta)\mid x_i(\zeta))\) are measurable.
\end{proposition}

\begin{proof}
The condition is plainly necessary. Now suppose that the functions
\(\zeta\mapsto(x(\zeta)\mid x_i(\zeta))\) are measurable. Let
\(y_1,y_2,\ldots\) be the sequence of fields of vectors derived from the
sequence \((x_i)\) by application
% Source PDF page 152 (printed p. 145)
\phantomsection\label{srcpage:145}
of Lemma~1. The sequence \((y_i)\) is a measurable field of orthonormal
bases. By part (iii) of Lemma~1, the functions
\(\zeta\mapsto(x(\zeta)\mid y_i(\zeta))\) are measurable. For every
measurable field of vectors \(y\) on \(Z\), the number
\((x(\zeta)\mid y(\zeta))\), which is equal to
\[
  \sum_{i=1}^{\infty}
  (x(\zeta)\mid y_i(\zeta))
  \overline{(y(\zeta)\mid y_i(\zeta))},
\]
depends measurably on \(\zeta\). Hence \(x\) is measurable.
\end{proof}

\begin{proposition}\label{prop:II-1-3}
For \(p=1,2,\ldots,\aleph_0\), let \(Z_p\) be the set of \(\zeta\in Z\)
such that \(d(\zeta)=p\), and let \(\mathfrak H_p\) be a complex Hilbert
space of Hilbert dimension \(p\). For every \(p\), the field induced by
\((\mathfrak H(\zeta))_{\zeta\in Z}\) on \(Z_p\) is isomorphic to the
constant field defined by \(\mathfrak H_p\).
\end{proposition}

\begin{proof}
Let \(y_1,y_2,\ldots\) be a measurable field of orthonormal bases and let
\((e_i^p)\) be an orthonormal basis of \(\mathfrak H_p\). For
\(\zeta\in Z_p\), let \(U(\zeta)\) be the isomorphism of
\(\mathfrak H(\zeta)\) onto \(\mathfrak H_p\) which transforms
\(y_1(\zeta)\) into \(e_1^p\), \(y_2(\zeta)\) into \(e_2^p\), and so on.
By Proposition~2, to say that the field of vectors \(x\) on \(Z\) is
measurable amounts to saying that, for every \(p\), the functions
\[
  \zeta\longmapsto(x(\zeta)\mid y_i(\zeta))
  =(U(\zeta)x(\zeta)\mid e_i^p)
\]
on \(Z_p\) are measurable, that is, that for every \(p\) the mapping
\(\zeta\mapsto U(\zeta)x(\zeta)\) of \(Z_p\) into \(\mathfrak H_p\) is
measurable.
\end{proof}

Proposition~3 makes it possible, in most cases, to reduce the study of
measurable fields of complex Hilbert spaces to the study of constant
fields.

\begin{proposition}\label{prop:II-1-4}
Let \(Z\) be a Borel space, \(\nu\) a positive measure on \(Z\), and
\(\zeta\mapsto\mathfrak H(\zeta)\) a field of Hilbert spaces on \(Z\). Let
\(x_1,x_2,\ldots\) be a sequence of fields of vectors having the following
properties: \(1^\circ\), the functions
\(\zeta\mapsto(x_i(\zeta)\mid x_j(\zeta))\) are measurable;
\(2^\circ\), for every \(\zeta\in Z\), the \(x_i(\zeta)\) form a total
sequence in \(\mathfrak H(\zeta)\). Then there exists on the
\(\mathfrak H(\zeta)\) a unique measurable-field structure for which the
fields \(x_i\) are measurable.
\end{proposition}

\begin{proof}
Uniqueness follows at once from Proposition~2. We prove existence. Apply
Lemma~1 to the fields \(x_i\), and let \(y_1,y_2,\ldots\) be the resulting
fields. Let \(\mathfrak G\) be the set of fields of vectors \(x\) for which
the functions \(\zeta\mapsto(x(\zeta)\mid y_i(\zeta))\) are measurable. We
show that \(\mathfrak G\) satisfies conditions (i), (ii), and (iii) of
Definition~1. If \(x\in\mathfrak G\),
\[
  \lVert x(\zeta)\rVert^2
  =\sum_{i=1}^{\infty}|(x(\zeta)\mid y_i(\zeta))|^2
\]
depends measurably on \(\zeta\). Moreover,
\(y_1,y_2,\ldots\in\mathfrak G\); hence, if a field of vectors \(y\) is such
that \((x(\zeta)\mid y(\zeta))\) depends measurably
% Source PDF page 153 (printed p. 146)
\phantomsection\label{srcpage:146}
on \(\zeta\) for every \(x\in\mathfrak G\), then
\(y\in\mathfrak G\). Finally, for every \(\zeta\in Z\), the
\(y_n(\zeta)\) form a total system in \(\mathfrak H(\zeta)\). Thus the
specification of \(\mathfrak G\) furnishes the \(\mathfrak H(\zeta)\) with
a measurable-field structure. Finally, the functions
\(\zeta\mapsto(x_i(\zeta)\mid y_j(\zeta))\) are measurable by part (iii) of
Lemma~1. Hence \(x_i\in\mathfrak G\) for every \(i\).
\end{proof}

\begin{remark*}
All the results of this numbered paragraph remain valid if ``measurable''
is replaced throughout by ``Borel.'' This remark is also valid for many
(but not all) of the results which follow. We shall henceforth refrain from
mentioning it.
\end{remark*}

Bibliography: \ArticleRef{28}, \ArticleRef{48}, \ArticleRef{80}.

\NumberedParagraph{5}{Fields of square-integrable vectors}
\label{no:II-1-5}
Let \(\zeta\mapsto\mathfrak H(\zeta)\) be a \(\nu\)-measurable field of
complex Hilbert spaces on \(Z\). A field of vectors \(x\) on \(Z\) is said
to be \emph{square-integrable} if it is measurable and if
\[
  \int\lVert x(\zeta)\rVert^2\,d\nu(\zeta)<+\infty.
\]
The set of square-integrable fields is a complex vector space
\(\mathfrak R\). For \(x,y\in\mathfrak R\),
\((x(\zeta)\mid y(\zeta))\) is an integrable function of \(\zeta\). Put
\[
  (x\mid y)=\int(x(\zeta)\mid y(\zeta))\,d\nu(\zeta).
\]
The space \(\mathfrak R\) is thus furnished with the structure of a complex
pre-Hilbert space. For \(x\in\mathfrak R\),
\[
  \lVert x\rVert^2=\int\lVert x(\zeta)\rVert^2\,d\nu(\zeta).
\]
The elements \(x\in\mathfrak R\) such that \(\lVert x\rVert=0\) are
precisely those \(x\) which vanish almost everywhere. We shall identify two
elements of \(\mathfrak R\) which are equal almost everywhere; in other
words, we consider the separated pre-Hilbert space \(\mathfrak H\)
associated with \(\mathfrak R\).

In what follows, the elements of \(\mathfrak H\) will be treated as fields
of vectors whenever this causes no confusion. Thus, for
\(x\in\mathfrak H\), we shall speak of the values
\(x(\zeta)\in\mathfrak H(\zeta)\) of \(x\). It must be noted, however, that
the \(x(\zeta)\) are determined only ``up to negligible sets.'' This is the
origin of the negligible exceptional sets which enter into most later
statements. Often these sets create no difficulty, and we shall sometimes
even refrain from mentioning them, preferring statements which are
slightly incorrect but easy to understand. At other times, on the contrary,
such an omission would be a source of serious errors.

% Source PDF page 154 (printed p. 147)
\phantomsection\label{srcpage:147}
\begin{proposition}\label{prop:II-1-5}
\begin{enumerate}[label=\textup{(\roman*)}]
\item The space \(\mathfrak H\) is a Hilbert space.
\item If a sequence \((x_n)\) of elements of \(\mathfrak H\) converges to
an element \(x\) of \(\mathfrak H\) in the metric of \(\mathfrak H\), then
some subsequence of \((x_n)\) converges to \(x\) almost everywhere.
\end{enumerate}
\end{proposition}

\begin{proof}
Let \(x_1,x_2,\ldots\) be a Cauchy sequence in \(\mathfrak H\). It is
enough to prove that a subsequence \((x_{n_k})\) converges almost everywhere
to an element \(x\) of \(\mathfrak H\), and that
\(\lVert x_{n_k}-x\rVert\to0\). By passing to a subsequence, we may suppose
that
\[
  \sum_{n=1}^{\infty}\lVert x_{n+1}-x_n\rVert<+\infty.
\]
Then
\[
  \sum_{n=1}^{\infty}
  \lVert x_{n+1}(\zeta)-x_n(\zeta)\rVert<+\infty
\]
for every \(\zeta\) outside a certain negligible set \(N\)
(\TreatiseRef{4}, Chapter~IV, \S~3, Theorem~1). For \(\zeta\notin N\), the
series
\[
  x_1(\zeta)+\sum_{n=1}^{\infty}
  \bigl(x_{n+1}(\zeta)-x_n(\zeta)\bigr)
\]
converges in the metric of \(\mathfrak H(\zeta)\) to an element
\(x(\zeta)\) of \(\mathfrak H(\zeta)\) such that
\[
  \lVert x(\zeta)\rVert\leq \lVert x_1(\zeta)\rVert
  +\sum_{n=1}^{\infty}
   \lVert x_{n+1}(\zeta)-x_n(\zeta)\rVert.
\]
Put \(x(\zeta)=0\) for \(\zeta\in N\). The field of vectors
\(\zeta\mapsto x(\zeta)\) is \(\nu\)-measurable as the almost-everywhere
limit of the measurable fields \(x_p\). By \TreatiseRef{4}, Chapter~IV,
\S~3, Theorem~1,
\(\int\lVert x(\zeta)\rVert^2d\nu(\zeta)<+\infty\) (so
\(x\in\mathfrak H\)), and
\[
  \lVert x-x_p\rVert\leq
  \sum_{n=p}^{\infty}\lVert x_{n+1}-x_n\rVert,
\]
which proves that \(x\) is the limit of the \(x_p\) in the metric of
\(\mathfrak H\).
\end{proof}

\begin{definition}\label{def:II-1-3}
The space \(\mathfrak H\) is called the \emph{Hilbert integral} of the
\(\mathfrak H(\zeta)\), and is denoted by
\[
  \int^{\oplus}\mathfrak H(\zeta)\,d\nu(\zeta).
\]
\end{definition}

The space \(\mathfrak H\) in fact depends on the choice of the space
\(\mathfrak G\) of measurable fields of vectors. When necessary, one may
use the notation
\[
  \int_{\mathfrak G}^{\oplus}\mathfrak H(\zeta)\,d\nu(\zeta).
\]
If \(x\) is a square-integrable field of vectors, by analogy with the
preceding notation we shall also sometimes denote the corresponding
element of \(\mathfrak H\) by
\[
  \int^{\oplus}x(\zeta)\,d\nu(\zeta).
\]

% Source PDF page 155 (printed p. 148)
\phantomsection\label{srcpage:148}
Let \(Y\) be a measurable subset of \(Z\). The square-integrable fields of
vectors which vanish outside \(Y\) form a closed vector subspace of
\(\mathfrak H\), denoted by
\[
  \int_Y^{\oplus}\mathfrak H(\zeta)\,d\nu(\zeta),
\]
and canonically identified with the Hilbert integral of the induced
measurable field \(\zeta\mapsto\mathfrak H(\zeta)\), \(\zeta\in Y\).

\noindent\emph{Examples.}---1. Suppose that \(Z\) is discrete and that
each point of \(Z\) has measure \(1\) (cf. no.~2, Example~1). Then
\(\mathfrak R\) can only be the space of fields of vectors \(x\) such that
\[
  \sum_{\zeta\in Z}\lVert x(\zeta)\rVert^2<+\infty,
\]
and \(\mathfrak H=\mathfrak R\) is the Hilbert sum of the
\(\mathfrak H(\zeta)\). In this case the \(\mathfrak H(\zeta)\) may be
canonically identified with subspaces of \(\mathfrak H\). This is not so in
general.

2. Let \(\mathfrak H_0\) be a complex Hilbert space with a countable base,
and let \(\zeta\mapsto\mathfrak H(\zeta)\) be the corresponding constant
field on \(Z\). Then the square-integrable fields of vectors are the
square-integrable mappings of \(Z\) into \(\mathfrak H_0\), and
\[
  \int^{\oplus}\mathfrak H(\zeta)\,d\nu(\zeta)
  =L_{\mathfrak H_0}^{2}(Z,\nu).
\]
Thus Proposition~5 appears as a generalization of the Fischer--Riesz
theorem.

\begin{remark*}
Let \(\nu_1\) be a positive measure equivalent to \(\nu\). Put
\(\nu_1=\rho\nu\), where \(\zeta\mapsto\rho(\zeta)\) is a
\(\nu\)-measurable function on \(Z\) such that
\(0<\rho(\zeta)<+\infty\) for every \(\zeta\). Then the mapping which sends
a square-\(\nu\)-integrable field of vectors \(x\) to the field of vectors
\[
  \zeta\longmapsto\rho(\zeta)^{-1/2}x(\zeta)
\]
is an isomorphism of
\(\int^{\oplus}\mathfrak H(\zeta)d\nu(\zeta)\) onto
\(\int^{\oplus}\mathfrak H(\zeta)d\nu_1(\zeta)\), since
\begin{align*}
  \int\bigl\lVert\rho(\zeta)^{-1/2}x(\zeta)\bigr\rVert^2d\nu_1(\zeta)
  &=\int\lVert x(\zeta)\rVert^2\rho(\zeta)^{-1}
    \rho(\zeta)d\nu(\zeta)\\
  &=\int\lVert x(\zeta)\rVert^2d\nu(\zeta).
\end{align*}
For fixed \(\nu\) and \(\nu_1\), this isomorphism is independent of the
choice of \(\rho\). It is called the \emph{canonical isomorphism} of
\[
  \int^{\oplus}\mathfrak H(\zeta)\,d\nu(\zeta)
  \quad\text{onto}\quad
  \int^{\oplus}\mathfrak H(\zeta)\,d\nu_1(\zeta).
\]
\end{remark*}

Bibliography: \ArticleRef{26}, \ArticleRef{28}, \ArticleRef{48},
\ArticleRef{52}, \ArticleRef{80}, \ArticleRef{100}, \ArticleRef{117},
\ArticleRef{118}.

\NumberedParagraph{6}{First properties of Hilbert integrals}
\label{no:II-1-6}
\begin{proposition}\label{prop:II-1-6}
Let \(y_1,y_2,\ldots\) be a measurable field of orthonormal bases.

% Source PDF page 156 (printed p. 149)
\phantomsection\label{srcpage:149}
\begin{enumerate}[label=\textup{(\roman*)}]
\item Let \(x\) be a field of vectors. In order that
  \(x\in\mathfrak H\), it is necessary and sufficient that the functions
  \(\zeta\mapsto(x(\zeta)\mid y_i(\zeta))\) be square-integrable and that
  \[
    \sum_{i=1}^{\infty}\int
      |(x(\zeta)\mid y_i(\zeta))|^2\,d\nu(\zeta)<+\infty.
  \]
\item If \(x,y\in\mathfrak H\), then
  \[
    (x\mid y)=\sum_{i=1}^{\infty}\int
    (x(\zeta)\mid y_i(\zeta))
    \overline{(y(\zeta)\mid y_i(\zeta))}\,d\nu(\zeta).
  \]
\item If \(x\in\mathfrak H\), then, in the metric of \(\mathfrak H\),
  \(x\) is the limit, as \(n\to+\infty\), of the fields of vectors
  \[
    \zeta\longmapsto\sum_{i=1}^{n}
    (x(\zeta)\mid y_i(\zeta))y_i(\zeta).
  \]
\end{enumerate}
\end{proposition}

\begin{proof}
Let \(x\) be a field of vectors. For every \(\zeta\),
\[
  \lVert x(\zeta)\rVert^2
  =\sum_{i=1}^{\infty}|(x(\zeta)\mid y_i(\zeta))|^2.
\]
Part (i) follows at once, in view of Proposition~2. If
\(x,y\in\mathfrak H\), then, for every \(\zeta\),
\[
  (x(\zeta)\mid y(\zeta))
  =\sum_{i=1}^{\infty}(x(\zeta)\mid y_i(\zeta))
   \overline{(y(\zeta)\mid y_i(\zeta))},
\]
and
\[
  \sum_{i=1}^{\infty}\int
  |(x(\zeta)\mid y_i(\zeta))|\,
  |(y(\zeta)\mid y_i(\zeta))|\,d\nu(\zeta)
  \leq\int\lVert x(\zeta)\rVert\lVert y(\zeta)\rVert\,d\nu(\zeta)
  <+\infty,
\]
which gives (ii). In particular,
\[
  \lVert x\rVert^2=\sum_{i=1}^{\infty}\int
  |(x(\zeta)\mid y_i(\zeta))|^2\,d\nu(\zeta).
\]
Thus, if
\[
  x_n(\zeta)=\sum_{i=1}^{n}
  (x(\zeta)\mid y_i(\zeta))y_i(\zeta),
\]
then
\[
  \lVert x-x_n\rVert^2=\sum_{i=n+1}^{\infty}\int
  |(x(\zeta)\mid y_i(\zeta))|^2\,d\nu(\zeta),
\]
so \(\lVert x-x_n\rVert\to0\) as \(n\to+\infty\), proving (iii).
\end{proof}

\begin{corollary*}
If \(\nu\) is standard, \(\mathfrak H\) has a countable base.
\end{corollary*}

\begin{proof}
By the hypothesis on \(\nu\), there exists a sequence
\(f_1,f_2,\ldots\) of complex-valued functions which is dense in
% Source PDF page 157 (printed p. 150)
\phantomsection\label{srcpage:150}
\(L_{\mathbb C}^{2}(Z,\nu)\). Hence the fields of vectors of the form
\[
  \zeta\longmapsto f_{n_1}(\zeta)y_1(\zeta)
  +f_{n_2}(\zeta)y_2(\zeta)+\cdots+f_{n_p}(\zeta)y_p(\zeta)
\]
are dense in \(\mathfrak H\) by Proposition~6(iii), and the set of these
fields is countable.
\end{proof}

\begin{proposition}\label{prop:II-1-7}
Let \(x_1,x_2,\ldots\) be a fundamental sequence of measurable fields of
vectors such that the functions
\(\zeta\mapsto\lVert x_i(\zeta)\rVert^2\) are bounded. Let \(\mathfrak M\)
be the set of fields of vectors obtained by multiplying the \(x_i\) by
bounded measurable numerical functions which vanish outside a set of
finite measure. Then \(\mathfrak M\) is a total set in \(\mathfrak H\).
\end{proposition}

\begin{proof}
It is clear that \(\mathfrak M\subset\mathfrak H\). Let
\(x\in\mathfrak H\) be orthogonal to \(\mathfrak M\). For every bounded
measurable function \(f\) on \(Z\) which vanishes outside a set of finite
measure, and for every \(i=1,2,\ldots\),
\[
  \int f(\zeta)(x_i(\zeta)\mid x(\zeta))\,d\nu(\zeta)=0.
\]
By the arbitrariness of \(f\), it follows that
\((x_i(\zeta)\mid x(\zeta))=0\) except on a negligible set \(N_i\). Let
\(N\) be the negligible union of the \(N_i\). For \(\zeta\notin N\), one
has \((x_i(\zeta)\mid x(\zeta))=0\) for every \(i\), and hence
\(x(\zeta)=0\). Thus \(\lVert x\rVert=0\).
\end{proof}

\begin{proposition}\label{prop:II-1-8}
Let \(x_1,x_2,\ldots\) be a sequence of measurable fields of vectors. Let
\(\mathfrak M\) be the set of square-integrable fields of vectors obtained
by multiplying the \(x_i\) by measurable numerical functions. If
\(\mathfrak M\) is total in \(\mathfrak H\), then, almost everywhere on
\(Z\), the \(x_i(\zeta)\) form a total sequence in
\(\mathfrak H(\zeta)\).
\end{proposition}

\begin{proof}
Let \(\mathfrak X(\zeta)\) be the closed vector subspace of
\(\mathfrak H(\zeta)\) spanned by the \(x_i(\zeta)\). We must show that
\(\mathfrak X(\zeta)=\mathfrak H(\zeta)\) almost everywhere on \(Z\). Let
\(y_1,y_2,\ldots\) be a fundamental sequence of square-integrable fields of
vectors (no.~3, Remark~1). Then \(y_i\) is, in the metric of
\(\mathfrak H\), the limit of a sequence \(z_1,z_2,\ldots\) of finite
linear combinations of fields in \(\mathfrak M\). By passing to a
subsequence, we may suppose, by Proposition~5, that the
\(z_n(\zeta)\) converge to \(y_i(\zeta)\) except on a negligible set
\(N_i\). For fixed \(\zeta\), every \(z_n(\zeta)\) is a finite linear
combination of the \(x_p(\zeta)\). Hence, for \(\zeta\notin N_i\),
\(y_i(\zeta)\in\mathfrak X(\zeta)\). The union \(N\) of the \(N_i\) is
negligible. For \(\zeta\notin N\), one has
\(y_i(\zeta)\in\mathfrak X(\zeta)\) for every \(i\), and therefore
\(\mathfrak X(\zeta)=\mathfrak H(\zeta)\).
\end{proof}

Bibliography: \ArticleRef{28}, \ArticleRef{80}.

\NumberedParagraph{7}{Measurable fields of subspaces}
\label{no:II-1-7}
\begin{proposition}\label{prop:II-1-9}
Let \(\zeta\mapsto\mathfrak H(\zeta)\) be a \(\nu\)-measurable field of
complex Hilbert spaces on \(Z\). For every \(\zeta\in Z\), let
\(\mathfrak K(\zeta)\) be a closed vector subspace
% Source PDF page 158 (printed p. 151)
\phantomsection\label{srcpage:151}
of \(\mathfrak H(\zeta)\), and let
\(E(\zeta)=P_{\mathfrak K(\zeta)}\). Let \(\mathfrak G\) be the set of
measurable fields of vectors \(x\) such that
\(x(\zeta)\in\mathfrak K(\zeta)\) for every \(\zeta\). The following
conditions are equivalent:
\begin{enumerate}[label=\textup{(\roman*)}]
\item the field of Hilbert spaces
  \(\zeta\mapsto\mathfrak K(\zeta)\), furnished with \(\mathfrak G\), is a
  \(\nu\)-measurable field;
\item there exists a sequence \((x_i)\) of measurable fields of vectors
  for the field \(\zeta\mapsto\mathfrak H(\zeta)\) such that, for every
  \(\zeta\), \(x_1(\zeta),x_2(\zeta),\ldots\) is a total sequence in
  \(\mathfrak K(\zeta)\);
\item for every measurable field of vectors \(x\) for the field
  \(\zeta\mapsto\mathfrak H(\zeta)\), the field
  \(\zeta\mapsto E(\zeta)x(\zeta)\) is measurable.
\end{enumerate}
\end{proposition}

\begin{proof}
\((i)\Rightarrow(ii)\) is clear.

\((ii)\Rightarrow(iii)\). Suppose the sequence \((x_i)\) exists. By
applying Lemma~1, the \(x_i\) may be replaced by measurable fields \(y_i\)
such that, for every \(\zeta\in Z\), the nonzero \(y_i(\zeta)\) form an
orthonormal basis of \(\mathfrak K(\zeta)\). Then, for every measurable
field of vectors \(x\),
\[
  \zeta\longmapsto E(\zeta)x(\zeta)
  =\sum_i(x(\zeta)\mid y_i(\zeta))y_i(\zeta)
\]
is a measurable field.

\((iii)\Rightarrow(i)\). Suppose condition (iii) satisfied, and show that
\(\mathfrak G\) satisfies properties (i), (ii), and (iii) of Definition~1.
Property (i) is clear. Now let \(x\) be a field of vectors such that
\(x(\zeta)\in\mathfrak K(\zeta)\) for every \(\zeta\in Z\), and such that,
for every \(y\in\mathfrak G\), \((x(\zeta)\mid y(\zeta))\) depends
measurably on \(\zeta\). For every measurable field of vectors \(z\) for
the field \(\zeta\mapsto\mathfrak H(\zeta)\), the field
\(\zeta\mapsto E(\zeta)z(\zeta)\) belongs to \(\mathfrak G\), and hence
\[
  (x(\zeta)\mid z(\zeta))
  =(x(\zeta)\mid E(\zeta)z(\zeta))
\]
depends measurably on \(\zeta\); thus \(x\in\mathfrak G\). Finally, if
\((x_i)\) is a fundamental sequence of measurable fields of vectors for
the field \(\zeta\mapsto\mathfrak H(\zeta)\), then the fields
\(\zeta\mapsto E(\zeta)x_i(\zeta)\) belong to \(\mathfrak G\), and, for
every \(\zeta\in Z\), the \(E(\zeta)x_i(\zeta)\) form a total sequence in
\(\mathfrak K(\zeta)\).
\end{proof}

Under the conditions of Proposition~9, the \(\mathfrak K(\zeta)\) are said
to form a \emph{measurable field of subspaces}. Then
\[
  \int^{\oplus}\mathfrak K(\zeta)\,d\nu(\zeta)
\]
is a complete, hence closed, subspace of
\(\int^{\oplus}\mathfrak H(\zeta)\,d\nu(\zeta)\).

Bibliography: \ArticleRef{59}.

\NumberedParagraph{8}{Measurable fields of tensor products}
\label{no:II-1-8}
\begin{proposition}\label{prop:II-1-10}
Let \(\zeta\mapsto\mathfrak H(\zeta)\) and
\(\zeta\mapsto\mathfrak K(\zeta)\) be two \(\nu\)-measurable fields of
complex Hilbert spaces on \(Z\). There exists on the field
\(\zeta\mapsto\mathfrak H(\zeta)\otimes\mathfrak K(\zeta)\) a unique
\(\nu\)-measurable-field structure
% Source PDF page 159 (printed p. 152)
\phantomsection\label{srcpage:152}
having the following property: whenever
\(\zeta\mapsto x(\zeta)\in\mathfrak H(\zeta)\) and
\(\zeta\mapsto y(\zeta)\in\mathfrak K(\zeta)\) are measurable fields of
vectors, the field \(\zeta\mapsto x(\zeta)\otimes y(\zeta)\) is
measurable.
\end{proposition}

\begin{proof}
Let \((x_i)\), respectively \((y_j)\), be a fundamental sequence of
measurable fields of vectors for the \(\mathfrak H(\zeta)\), respectively
the \(\mathfrak K(\zeta)\). The fields
\(\zeta\mapsto x_i(\zeta)\otimes y_j(\zeta)\) have the properties which
permit Proposition~4 to be applied. Thus there exists on the
\(\mathfrak H(\zeta)\otimes\mathfrak K(\zeta)\) a measurable-field
structure such that the fields
\(\zeta\mapsto x_i(\zeta)\otimes y_j(\zeta)\) form a fundamental sequence
of measurable fields of vectors. If
\(\zeta\mapsto x(\zeta)\in\mathfrak H(\zeta)\) and
\(\zeta\mapsto y(\zeta)\in\mathfrak K(\zeta)\) are measurable fields of
vectors, then
\begin{align*}
 &(x(\zeta)\otimes y(\zeta)\mid
   x_i(\zeta)\otimes y_j(\zeta))\\
 &\qquad=(x(\zeta)\mid x_i(\zeta))
   (y(\zeta)\mid y_j(\zeta))
\end{align*}
depends measurably on \(\zeta\), and therefore the field
\(\zeta\mapsto x(\zeta)\otimes y(\zeta)\) is measurable. This proves the
existence of the structure asserted in the proposition. On the other hand,
for every measurable-field structure on the
\(\mathfrak H(\zeta)\otimes\mathfrak K(\zeta)\) having the stated
property, the fields \(\zeta\mapsto x_i(\zeta)\otimes y_j(\zeta)\) form a
fundamental sequence of measurable fields of vectors. This proves
uniqueness.
\end{proof}

\begin{proposition}\label{prop:II-1-11}
Let \(\zeta\mapsto\mathfrak H(\zeta)\) be a \(\nu\)-measurable field of
complex Hilbert spaces on \(Z\), and let
\(\zeta\mapsto\mathfrak K(\zeta)\) be the constant field on \(Z\)
corresponding to a Hilbert space \(\mathfrak K_0\) with a countable base.
There exists a unique isomorphism of the Hilbert space
\[
  \left(\int^{\oplus}\mathfrak H(\zeta)\,d\nu(\zeta)\right)
  \otimes\mathfrak K_0
\]
onto the Hilbert space
\[
  \int^{\oplus}\bigl(\mathfrak H(\zeta)\otimes
  \mathfrak K(\zeta)\bigr)\,d\nu(\zeta)
\]
which, for every
\(x\in\int^{\oplus}\mathfrak H(\zeta)d\nu(\zeta)\) and every
\(y\in\mathfrak K_0\), transforms \(x\otimes y\) into the field
\(\zeta\mapsto x(\zeta)\otimes y\).
\end{proposition}

\begin{proof}
If \(x_1,x_2,\ldots,x_n\) are elements of
\(\int^{\oplus}\mathfrak H(\zeta)d\nu(\zeta)\), and
\(y_1,y_2,\ldots,y_n\) are elements of \(\mathfrak K_0\), then
\begin{align*}
 \int\left\lVert\sum_{i=1}^{n}x_i(\zeta)\otimes y_i\right\rVert^2
 d\nu(\zeta)
 &=\int\sum_{i,j=1}^{n}
   (x_i(\zeta)\mid x_j(\zeta))(y_i\mid y_j)\,d\nu(\zeta)\\
 &=\sum_{i,j=1}^{n}\left(\int
   (x_i(\zeta)\mid x_j(\zeta))\,d\nu(\zeta)\right)(y_i\mid y_j)\\
 &=\sum_{i,j=1}^{n}(x_i\mid x_j)(y_i\mid y_j)
 =\left\lVert\sum_{i=1}^{n}x_i\otimes y_i\right\rVert^2.
\end{align*}
% Source PDF page 160 (printed p. 153)
\phantomsection\label{srcpage:153}
Thus there exists an isometric linear mapping of the algebraic tensor
product of \(\int^{\oplus}\mathfrak H(\zeta)d\nu(\zeta)\) and
\(\mathfrak K_0\) into
\(\int^{\oplus}(\mathfrak H(\zeta)\otimes\mathfrak K(\zeta))
d\nu(\zeta)\) which transforms \(\sum_{i=1}^{n}x_i\otimes y_i\) into the
field
\[
  \zeta\longmapsto\sum_{i=1}^{n}x_i(\zeta)\otimes y_i.
\]
The isometric extension \(\theta\) of this mapping to
\[
  \left(\int^{\oplus}\mathfrak H(\zeta)d\nu(\zeta)\right)
  \otimes\mathfrak K_0
\]
has the whole of
\[
  \int^{\oplus}\bigl(\mathfrak H(\zeta)\otimes
  \mathfrak K(\zeta)\bigr)d\nu(\zeta)
\]
as its image. Indeed, let \((z_i)\) be a fundamental sequence of measurable
fields of vectors for the \(\mathfrak H(\zeta)\), such that the
\(\lVert z_i(\zeta)\rVert\) are bounded, and let \((z_j')\) be a total
sequence in \(\mathfrak K_0\). Then the fields
\(\zeta\mapsto z_i(\zeta)\otimes z_j'\) form a fundamental sequence of
measurable fields of vectors for the
\(\mathfrak H(\zeta)\otimes\mathfrak K(\zeta)\); hence the fields
\(\zeta\mapsto f(\zeta)z_i(\zeta)\otimes z_j'\), where \(f\) is a bounded
measurable numerical function which vanishes outside a set of finite
measure, form a total set in
\(\int^{\oplus}(\mathfrak H(\zeta)\otimes\mathfrak K(\zeta))d\nu(\zeta)\)
by Proposition~7; but these fields lie in the image of \(\theta\). This
proves the assertion. Finally, uniqueness is clear.
\end{proof}

Henceforth, by the isomorphism of Proposition~11, we shall identify
\[
  \int^{\oplus}\bigl(\mathfrak H(\zeta)\otimes\mathfrak K(\zeta)\bigr)
  d\nu(\zeta)
  \quad\text{and}\quad
  \left(\int^{\oplus}\mathfrak H(\zeta)d\nu(\zeta)\right)
  \otimes\mathfrak K_0.
\]

\begin{corollary*}
Let \(\mathfrak K_0\) be a Hilbert space with a countable base, and let
\(\zeta\mapsto\mathfrak K(\zeta)\) be the corresponding constant field on
\(Z\). There exists a unique isomorphism of the Hilbert space
\(L_{\mathbb C}^{2}(Z,\nu)\otimes\mathfrak K_0\) onto the Hilbert space
\(\int^{\oplus}\mathfrak K(\zeta)d\nu(\zeta)\) which, for every
\(f\in L_{\mathbb C}^{2}(Z,\nu)\) and every \(y\in\mathfrak K_0\),
transforms \(f\otimes y\) into the field
\(\zeta\mapsto f(\zeta)y\).
\end{corollary*}

\begin{proof}
Take for the field \(\zeta\mapsto\mathfrak H(\zeta)\) the constant field
corresponding to the one-dimensional Hilbert space \(\mathbb C\). Then
\[
  \int^{\oplus}\mathfrak H(\zeta)\,d\nu(\zeta)
  =L_{\mathbb C}^{2}(Z,\nu)
\]
% Source PDF page 161 (printed p. 154)
\phantomsection\label{srcpage:154}
(no.~5, Example~2). Moreover,
\(\mathfrak H(\zeta)\otimes\mathfrak K(\zeta)\) is canonically identified
with \(\mathfrak K(\zeta)\).
\end{proof}

Henceforth we shall identify
\[
  \int^{\oplus}\mathfrak K(\zeta)\,d\nu(\zeta)
  =L_{\mathfrak K_0}^{2}(Z,\nu)
  \quad\text{with}\quad
  L_{\mathbb C}^{2}(Z,\nu)\otimes\mathfrak K_0
\]
by the isomorphism of the corollary.

Bibliography: \ArticleRef{53}.

\noindent\emph{Exercises.}---
\begin{enumerate}[label=\arabic*.]
\item Let \(\zeta\mapsto\mathfrak H(\zeta)\) be a \(\nu\)-measurable
field of complex Hilbert spaces on \(Z\), and put
\(\mathfrak H=\int^{\oplus}\mathfrak H(\zeta)d\nu(\zeta)\). Let
\(x_1,x_2,\ldots\) be a sequence of measurable fields of vectors. Let
\(\mathfrak M\) be the set of square-integrable fields of vectors obtained
by multiplying the \(x_i\) by characteristic functions of measurable
sets. If \(\mathfrak M\) is strongly dense in \(\mathfrak H\), then the
\(x_i(\zeta)\) are strongly dense in \(\mathfrak H(\zeta)\) almost
everywhere. [Argue as for Proposition~8, replacing \(\mathfrak X(\zeta)\)
by the strong closure of the set of the \(x_i(\zeta)\), and \((y_i)\) by a
sequence of measurable fields of vectors such that, for every
\(\zeta\in Z\), the \(y_i(\zeta)\) are dense in
\(\mathfrak H(\zeta)\), and such that
\(\sup_{\zeta\in Z}\lVert y_i(\zeta)\rVert<+\infty\).]

\item Let \(\zeta\mapsto\mathfrak H(\zeta)\) be the constant field
corresponding to the one-dimensional Hilbert space \(\mathbb C\), and put
\(\mathfrak H=\int^{\oplus}\mathfrak H(\zeta)d\nu(\zeta)\). Show that the
weak convergence to zero of a sequence \((x_n)\) of elements of
\(\mathfrak H\) does not necessarily imply the weak convergence of
\(x_n(\zeta)\) to zero almost everywhere.

\item Let \(\zeta\mapsto\mathfrak H(\zeta)\) be a \(\nu\)-measurable
field of complex Hilbert spaces on \(Z\), and let \((x_i)\) be a fundamental
sequence of measurable fields of vectors. For every measurable field of
vectors \(\zeta\mapsto x(\zeta)\), there exists a sequence of fields of
vectors of the form
\[
  \zeta\longmapsto\sum_{i=1}^{n}f_i(\zeta)x_i(\zeta),
\]
where the \(f_i\) are measurable complex-valued functions on \(Z\), which
converges to \(x(\zeta)\) almost everywhere on \(Z\). [Reduce to the case
where \(\nu(Z)<+\infty\), and where \(\lVert x(\zeta)\rVert\) and the
\(\lVert x_i(\zeta)\rVert\) are bounded. Then use Proposition~7 and
Proposition~5(ii).] \ArticleRef{28}

\item Let \(Z\) be a Borel space, \(\nu\) a positive measure on \(Z\),
\(\zeta\mapsto\mathfrak H(\zeta)\) a field of complex Hilbert spaces on
\(Z\), \(\mathfrak F=\prod_{\zeta\in Z}\mathfrak H(\zeta)\), and
\(\Omega\) the set of vector subspaces \(\mathfrak G\) of \(\mathfrak F\)
having the following property: for every \(x\in\mathfrak G\), the function
\(\zeta\mapsto\lVert x(\zeta)\rVert\) is \(\nu\)-measurable.
\begin{enumerate}[label=\alph*.]
\item If the field \(\zeta\mapsto\mathfrak H(\zeta)\) is a
  \(\nu\)-measurable field, the set of \(\nu\)-measurable fields of vectors
  is a maximal element of \(\Omega\).
\item Let \(\mathfrak G\) be a maximal element of \(\Omega\). Suppose
  that there exists a sequence \((x_i)\) of elements of \(\mathfrak G\)
  such that, for every \(\zeta\in Z\), the \(x_i(\zeta)\) form a total
  sequence in \(\mathfrak H(\zeta)\). Then the field
  \(\zeta\mapsto\mathfrak H(\zeta)\), furnished with \(\mathfrak G\), is
  \(\nu\)-measurable \ArticleRef{48}, \ArticleRef{52}.
\end{enumerate}

\item\label{ex:II-1-5} Let \(Z\) be a locally compact space,
\(\zeta\mapsto\mathfrak H(\zeta)\) a field of complex Hilbert spaces on
\(Z\), and \(\mathfrak F=\prod_{\zeta\in Z}\mathfrak H(\zeta)\). The
field \(\zeta\mapsto\mathfrak H(\zeta)\) is said to be a \emph{continuous
field of Hilbert spaces} if a vector subspace \(\mathfrak G\) of
\(\mathfrak F\) is specified having the following properties:

% Source PDF page 162 (printed p. 155)
\phantomsection\label{srcpage:155}
\begin{enumerate}[label=\textup{(\roman*)}]
\item for every \(x\in\mathfrak G\), \(\lVert x(\zeta)\rVert\) depends
  continuously on \(\zeta\);
\item if \(y\in\mathfrak F\) is such that, for every
  \(x\in\mathfrak G\), \(\lVert y(\zeta)-x(\zeta)\rVert\) depends
  continuously on \(\zeta\), then \(y\in\mathfrak G\);
\item for every \(\zeta\in Z\), the \(x(\zeta)\),
  \(x\in\mathfrak G\), form a total set in \(\mathfrak H(\zeta)\).
\end{enumerate}
The fields in \(\mathfrak G\) are then called \emph{continuous fields of
vectors}.
\begin{enumerate}[label=\alph*.]
\item If \(x,y\in\mathfrak G\), then
  \((x(\zeta)\mid y(\zeta))\) depends continuously on \(\zeta\).
\item Property (ii) may be replaced by:
  \begin{quote}
  (ii') if \(y\in\mathfrak F\) is such that
  \(\lVert y(\zeta)\rVert\) depends continuously on \(\zeta\), and, for
  every \(x\in\mathfrak G\), \((x(\zeta)\mid y(\zeta))\) depends
  continuously on \(\zeta\), then \(y\in\mathfrak G\).
  \end{quote}
  (Cf. j.)
\item If \(x\in\mathfrak G\) and \(f\) is a continuous complex-valued
  function on \(Z\), the field \(\zeta\mapsto f(\zeta)x(\zeta)\) is
  continuous. (Use b.)
\item Let \((x_n)\) be a sequence of continuous fields of vectors and let
  \(x\in\mathfrak F\). If
  \(\lVert x_n(\zeta)-x(\zeta)\rVert\to0\) uniformly on \(Z\), then
  \(x\in\mathfrak G\).
\item In order that \(x\in\mathfrak F\) be continuous, it is necessary and
  sufficient that, for every \(\zeta\in Z\) and every \(\varepsilon>0\),
  there exist a continuous field \(y\) such that
  \(\lVert x(\zeta')-y(\zeta')\rVert\leq\varepsilon\) in a neighborhood
  of \(\zeta\).
\item Let \((x_i)_{i\in I}\) be a family of continuous fields of vectors
  such that, for every \(\zeta\in Z\), the \(x_i(\zeta)\) form a total set
  in \(\mathfrak H(\zeta)\). Let \(x\in\mathfrak F\). If the functions
  \(\zeta\mapsto\lVert x(\zeta)\rVert\) and
  \(\zeta\mapsto(x(\zeta)\mid x_i(\zeta))\) are continuous, then
  \(x\in\mathfrak G\). [Use both (ii) and (ii').]
\item For every \(\zeta\in Z\) and every
  \(a\in\mathfrak H(\zeta)\), there exists \(x\in\mathfrak G\) such that
  \(x(\zeta)=a\). [Construct a sequence of fields
  \(x_n\in\mathfrak G\) such that
  \(\lVert x_{n+1}(\zeta')-x_n(\zeta')\rVert\leq2^{-n}\) on \(Z\), and
  such that \(x_n(\zeta)\) converges strongly to \(a\).]
\item Let \(n\) be an integer \(\geq0\). The set of \(\zeta\in Z\) such
  that \(\dim\mathfrak H(\zeta)>n\) is open.
\item Let \(\mathfrak F_0\) be the set of \(x\in\mathfrak F\) such that
  \(\lVert x(\zeta)\rVert\) depends continuously on \(\zeta\). A vector
  subspace of \(\mathfrak F\) satisfies (i) and (ii) if and only if it is
  maximal among the vector subspaces of \(\mathfrak F\) contained in
  \(\mathfrak F_0\).
\item If \(\mathfrak H(\zeta)\), for every \(\zeta\in Z\), is identical
  with a fixed Hilbert space \(\mathfrak H_0\), the set of strongly
  continuous mappings of \(Z\) into \(\mathfrak H_0\) satisfies (i), (ii),
  and (iii). But a mapping \(y\) of \(Z\) into \(\mathfrak H_0\) such that
  \((y(\zeta)\mid x(\zeta))\) depends continuously on \(\zeta\) for every
  strongly continuous mapping \(x\) of \(Z\) into \(\mathfrak H_0\) is not
  always strongly continuous.
\item If \(\mathfrak G\) satisfies the property
  \begin{quote}
  (iv) there exists a sequence \((x_n)\) of elements of
  \(\mathfrak G\) such that, for every \(\zeta\in Z\), the
  \(x_n(\zeta)\) form a total sequence in \(\mathfrak H(\zeta)\),
  \end{quote}
  then on the field \(\zeta\mapsto\mathfrak H(\zeta)\) there exists a
  unique Borel-field structure for which the \(x_n\) are Borel fields of
  vectors \ArticleRef{28}, \ArticleRef{48}.
\end{enumerate}
\end{enumerate}

\section{Fields of Operators}\label{sec:II-2}

\NumberedParagraph{1}{Measurable fields of linear mappings}
\label{no:II-2-1}
Let \(Z\) be a Borel space, \(\nu\) a positive measure on \(Z\), and let
\(\zeta\mapsto\mathfrak H(\zeta)\) and
\(\zeta\mapsto\mathfrak H'(\zeta)\) be two \(\nu\)-measurable fields of
complex Hilbert spaces on \(Z\). For every \(\zeta\in Z\), let
\(T(\zeta)\in\mathcal L(\mathfrak H(\zeta),\mathfrak H'(\zeta))\); that
is, recall, \(T(\zeta)\) is a continuous linear mapping of
\(\mathfrak H(\zeta)\) into \(\mathfrak H'(\zeta)\). The mapping
\(\zeta\mapsto T(\zeta)\) is called a \emph{field of continuous linear
mappings} on \(Z\). Most often the measurable field
\(\zeta\mapsto\mathfrak H'(\zeta)\) will be
% Source PDF page 163 (printed p. 156)
\phantomsection\label{srcpage:156}
identical with the field \(\zeta\mapsto\mathfrak H(\zeta)\), and one then
speaks of a \emph{field of operators}.

\begin{definition}\label{def:II-2-1}
The field of continuous linear mappings \(\zeta\mapsto T(\zeta)\) is said
to be \emph{measurable} if, for every measurable field of vectors
\(\zeta\mapsto x(\zeta)\in\mathfrak H(\zeta)\), the field of vectors
\(\zeta\mapsto T(\zeta)x(\zeta)\in\mathfrak H'(\zeta)\) is measurable.
\end{definition}

If the field \(\zeta\mapsto T(\zeta)\) is measurable, then
\(\lVert T(\zeta)\rVert\) depends measurably on \(\zeta\). Indeed, let
\(x_1,x_2,\ldots\) be a sequence of measurable fields of vectors, with
values in the \(\mathfrak H(\zeta)\), such that, for every \(\zeta\in Z\),
the \(x_i(\zeta)\) are dense in \(\mathfrak H(\zeta)\) (\S~1, no.~2,
Remark~1). Then \(\lVert T(\zeta)\rVert\) is the supremum of the numbers
\[
  \lVert T(\zeta)x_i(\zeta)\rVert\,
  \lVert x_i(\zeta)\rVert^{-1}
\]
(with the convention \(0\cdot+\infty=0\)), and these numbers depend
measurably on \(\zeta\), proving our assertion.

\begin{proposition}\label{prop:II-2-1}
Let \(x_1,x_2,\ldots\) and \(x_1',x_2',\ldots\) be two fundamental
sequences of measurable fields of vectors, with values respectively in the
\(\mathfrak H(\zeta)\) and the \(\mathfrak H'(\zeta)\). The field
\(\zeta\mapsto T(\zeta)\) is measurable if and only if the functions
\[
  \zeta\longmapsto(T(\zeta)x_i(\zeta)\mid x_j'(\zeta))
\]
are measurable.
\end{proposition}

\begin{proof}
The condition is plainly necessary. Conversely, suppose that the
functions
\[
  \zeta\longmapsto(T(\zeta)x_i(\zeta)\mid x_j'(\zeta))
  =(x_i(\zeta)\mid T(\zeta)^*x_j'(\zeta))
\]
are measurable. Then the fields of vectors
\(\zeta\mapsto T(\zeta)^*x_j'(\zeta)\) are measurable (\S~1,
Proposition~2). Hence, for every measurable field of vectors \(x\) with
values in the \(\mathfrak H(\zeta)\), the functions
\[
  \zeta\longmapsto(T(\zeta)x(\zeta)\mid x_j'(\zeta))
  =(x(\zeta)\mid T(\zeta)^*x_j'(\zeta))
\]
are measurable. Thus, by \S~1, Proposition~2, the field of vectors
\(\zeta\mapsto T(\zeta)x(\zeta)\) is measurable.
\end{proof}

If \(\zeta\mapsto T(\zeta)\) and \(\zeta\mapsto T'(\zeta)\) are
measurable fields of operators, the fields
\(\zeta\mapsto T(\zeta)+T'(\zeta)\) and
\(\zeta\mapsto T(\zeta)T'(\zeta)\) are plainly measurable. The field
\(\zeta\mapsto T(\zeta)^*\) is measurable by Proposition~1. There are
analogous properties for measurable fields of linear mappings.

Let \(Y\) be a measurable subset of \(Z\). A field of linear mappings
\(\zeta\mapsto T(\zeta)\) on \(Y\) is said to be measurable if it extends
to a measurable field on \(Z\). This is equivalent to saying that, for
every measurable field of vectors \(x\) on \(Y\), the field of vectors
\(\zeta\mapsto T(\zeta)x(\zeta)\) on \(Y\) is measurable. If \(Z\) is the
union of a sequence \(Z_1,Z_2,\ldots\) of measurable subsets, a field of
linear mappings on \(Z\) is measurable if and only if its restriction to
each \(Z_i\) is measurable.

Let \(\zeta\mapsto T(\zeta)\in\mathcal L(\mathfrak H(\zeta))\) be a
measurable field of operators on \(Z\), and let
\(\zeta\mapsto\mathfrak K(\zeta)\) be a measurable field of subspaces such
that
% Source PDF page 164 (printed p. 157)
\phantomsection\label{srcpage:157}
\[
  P_{\mathfrak K(\zeta)}T(\zeta)P_{\mathfrak K(\zeta)}=T(\zeta)
  \qquad(\zeta\in Z).
\]
Then the field of operators
\(\zeta\mapsto T(\zeta)_{\mathfrak K(\zeta)}\) is measurable, relative to
the field \(\zeta\mapsto\mathfrak K(\zeta)\), as follows immediately from
\S~1, Proposition~9.

If \(\zeta\mapsto\mathfrak H_1(\zeta)\) is another \(\nu\)-measurable
field of Hilbert spaces on \(Z\), and if
\(\zeta\mapsto T_1(\zeta)\in\mathcal L(\mathfrak H_1(\zeta))\) is a
measurable field of operators, the field
\(\zeta\mapsto T(\zeta)\otimes T_1(\zeta)\) is measurable, relative to
the field
\(\zeta\mapsto\mathfrak H(\zeta)\otimes\mathfrak H_1(\zeta)\). Indeed,
let
\[
  \zeta\longmapsto x(\zeta)\in\mathfrak H(\zeta),
  \qquad
  \zeta\longmapsto x_1(\zeta)\in\mathfrak H_1(\zeta)
\]
be measurable fields of vectors. Then the field
\[
  \zeta\longmapsto
  \bigl(T(\zeta)\otimes T_1(\zeta)\bigr)
  \bigl(x(\zeta)\otimes x_1(\zeta)\bigr)
  =T(\zeta)x(\zeta)\otimes T_1(\zeta)x_1(\zeta)
\]
is measurable, and our assertion then follows easily from Proposition~1.

Bibliography: \ArticleRef{28}, \ArticleRef{48}, \ArticleRef{52},
\ArticleRef{80}, \ArticleRef{100}, \ArticleRef{117}, \ArticleRef{118}.

\NumberedParagraph{2}{Examples}\label{no:II-2-2}
\begin{enumerate}
\item Suppose that \(Z\) is discrete. A measurable field of continuous
  linear mappings is any mapping
  \(\zeta\mapsto T(\zeta)\) such that
  \(T(\zeta)\in\mathcal L(\mathfrak H(\zeta),\mathfrak H'(\zeta))\).
\item Let
  \(\zeta\mapsto T(\zeta)\in
  \mathcal L(\mathfrak H(\zeta),\mathfrak H'(\zeta))\) be a measurable
  field of continuous linear mappings. Suppose that, for every
  \(\zeta\in Z\), \(T(\zeta)\) is an isomorphism of the Hilbert space
  \(\mathfrak H(\zeta)\) onto the Hilbert space
  \(\mathfrak H'(\zeta)\). Then the field
  \(\zeta\mapsto T(\zeta)^{-1}=T(\zeta)^*\) is measurable, so that
  \((T(\zeta))_{\zeta\in Z}\) is an isomorphism of the measurable field
  \((\mathfrak H(\zeta))_{\zeta\in Z}\) onto the measurable field
  \((\mathfrak H'(\zeta))_{\zeta\in Z}\).
\item Let \(Z\) be a locally compact space countable at infinity,
  \(\nu\) a positive measure on \(Z\), \(\mathfrak H_0\) a Hilbert space
  with a countable base, and \(\zeta\mapsto\mathfrak H(\zeta)\) the
  corresponding constant field on \(Z\). Let
  \(\zeta\mapsto T(\zeta)\) be a field of operators, that is, a mapping
  of \(Z\) into \(\mathcal L(\mathfrak H_0)\). The following hypotheses
  are equivalent:
  \begin{enumerate}
  \item the field \(\zeta\mapsto T(\zeta)\) is measurable;
  \item the functions
    \(\zeta\mapsto(T(\zeta)x_i\mid x_j)\), where \((x_i)\) is a total
    family in \(\mathfrak H_0\), are measurable;
  \item the mapping \(\zeta\mapsto T(\zeta)\) of \(Z\) into
    \(\mathcal L(\mathfrak H_0)\), furnished with the strong topology,
    is measurable.
  \end{enumerate}
  Indeed, (a) and (b) are equivalent by Proposition~1. It is clear that
  (c) implies (b). Finally, suppose (a) holds, and let \(Y\) be a compact
  subset of \(Z\). Since the function
  \(\zeta\mapsto\lVert T(\zeta)\rVert\) is measurable, there is a compact
  subset \(Y_1\subset Y\) on which \(\lVert T(\zeta)\rVert\) is bounded
  and such that \(\nu(Y-Y_1)\) is arbitrarily small. Moreover, if
  \((y_1,y_2,\ldots)\) is an everywhere dense sequence in
  \(\mathfrak H_0\), there is a compact subset \(Y_2\subset Y_1\) such
  that the mappings \(\zeta\mapsto T(\zeta)y_i\) are continuous on
  \(Y_2\), where \(\mathfrak H_0\) has its strong topology, and such that
  \(\nu(Y_1-Y_2)\) is arbitrarily small. Thus \(\nu(Y-Y_2)\)
% Source PDF page 165 (printed p. 158)
\phantomsection\label{srcpage:158}
  is arbitrarily small, and on \(Y_2\) the mapping
  \(\zeta\mapsto T(\zeta)y\) is continuous for every
  \(y\in\mathfrak H_0\), with \(\mathfrak H_0\) furnished with the strong
  topology. Hence (c) holds.
\end{enumerate}

\NumberedParagraph{3}{Decomposable linear mappings}
\label{no:II-2-3}
Let \(\zeta\mapsto\mathfrak H(\zeta)\) and
\(\zeta\mapsto\mathfrak H'(\zeta)\) be two \(\nu\)-measurable fields of
complex Hilbert spaces on \(Z\). Put
\[
  \mathfrak H=\int^{\oplus}\mathfrak H(\zeta)\,d\nu(\zeta),
  \qquad
  \mathfrak H'=\int^{\oplus}\mathfrak H'(\zeta)\,d\nu(\zeta).
\]
A measurable field
\[
  \zeta\longmapsto T(\zeta)\in
  \mathcal L(\mathfrak H(\zeta),\mathfrak H'(\zeta))
\]
is said to be \emph{essentially bounded} if the essential supremum
\(\lambda\) of the function \(\zeta\mapsto\lVert T(\zeta)\rVert\) is
finite. If this is so, then, for every square-integrable field of vectors
\(x\), the field \(x'\) defined by
\(x'(\zeta)=T(\zeta)x(\zeta)\) is square-integrable, and
\(\lVert x'\rVert\leq\lambda\lVert x\rVert\). Thus the mapping
\(x\mapsto x'\) is an element \(T\) of
\(\mathcal L(\mathfrak H,\mathfrak H')\) such that
\(\lVert T\rVert\leq\lambda\).

\begin{proposition}\label{prop:II-2-2}
We have \(\lVert T\rVert=\lambda\).
\end{proposition}

\begin{proof}
Let \(x\) be a square-integrable field of vectors, and let \(f\) be
an essentially bounded measurable complex-valued function on \(Z\). Then
\begin{align*}
  \int |f(\zeta)|^2\lVert T(\zeta)x(\zeta)\rVert^2\,d\nu(\zeta)
  &=\lVert T(fx)\rVert^2 \\
  &\leq\lVert T\rVert^2\lVert fx\rVert^2 \\
  &=\lVert T\rVert^2\int |f(\zeta)|^2
       \lVert x(\zeta)\rVert^2\,d\nu(\zeta).
\end{align*}
Consequently, since \(f\) is arbitrary,
\begin{equation}\tag{1}\label{eq:II-2-1}
  \lVert T(\zeta)x(\zeta)\rVert
  \leq\lVert T\rVert\,\lVert x(\zeta)\rVert
  \quad\text{almost everywhere.}
\end{equation}
Now let \((x_n)\) be a sequence of square-integrable fields of vectors
such that, for every \(\zeta\in Z\), the \(x_n(\zeta)\) are everywhere
dense in \(\mathfrak H(\zeta)\). Almost everywhere,
\[
  \lVert T(\zeta)x_n(\zeta)\rVert
  \leq\lVert T\rVert\,\lVert x_n(\zeta)\rVert
  \qquad\text{for every }n,
\]
and hence \(\lVert T(\zeta)\rVert\leq\lVert T\rVert\). Therefore
\(\lambda\leq\lVert T\rVert\).
\end{proof}

\begin{corollary*}
If two essentially bounded measurable fields of linear mappings define
the same element of \(\mathcal L(\mathfrak H,\mathfrak H')\), they are
equal almost everywhere.
\end{corollary*}

\begin{proof}
The difference \(\zeta\mapsto T(\zeta)\) of these two fields defines the
zero mapping; hence \(\lVert T(\zeta)\rVert=0\) almost everywhere.
\end{proof}

% Source PDF page 166 (printed p. 159)
\phantomsection\label{srcpage:159}
\begin{definition}\label{def:II-2-2}
A mapping \(T\in\mathcal L(\mathfrak H,\mathfrak H')\) is said to be
\emph{decomposable} if it is defined by an essentially bounded measurable
field \(\zeta\mapsto T(\zeta)\). We then write
\[
  T=\int^{\oplus}T(\zeta)\,d\nu(\zeta).
\]
If the fields \(\zeta\mapsto\mathfrak H(\zeta)\) and
\(\zeta\mapsto\mathfrak H'(\zeta)\) are identical, one speaks of
\emph{decomposable operators}.
\end{definition}

By the corollary to Proposition~2, the \(T(\zeta)\) are defined ``up to
null sets.'' In particular, at a given point \(\zeta\in Z\) of measure
zero, \(T(\zeta)\) may be chosen arbitrarily.

\emph{Example.} Let \(\zeta\mapsto T(\zeta)\) be an isomorphism of the
measurable field \((\mathfrak H(\zeta))_{\zeta\in Z}\) onto the measurable
field \((\mathfrak H'(\zeta))_{\zeta\in Z}\). Then
\[
  T=\int^{\oplus}T(\zeta)\,d\nu(\zeta),
  \qquad
  T^{-1}=\int^{\oplus}T(\zeta)^{-1}\,d\nu(\zeta)
\]
are reciprocal isomorphisms of \(\mathfrak H\) onto \(\mathfrak H'\) and
of \(\mathfrak H'\) onto \(\mathfrak H\).

\begin{proposition}\label{prop:II-2-3}
Let \(T_1,T_2\) be decomposable operators. If
\[
  T_1=\int^{\oplus}T_1(\zeta)\,d\nu(\zeta),
  \qquad
  T_2=\int^{\oplus}T_2(\zeta)\,d\nu(\zeta),
\]
then
\begin{align*}
  T_1+T_2
    &=\int^{\oplus}\bigl(T_1(\zeta)+T_2(\zeta)\bigr)\,d\nu(\zeta),
  &
  T_1T_2
    &=\int^{\oplus}T_1(\zeta)T_2(\zeta)\,d\nu(\zeta),\\
  \lambda T_1
    &=\int^{\oplus}\lambda T_1(\zeta)\,d\nu(\zeta),
  &
  T_1^*
    &=\int^{\oplus}T_1(\zeta)^*\,d\nu(\zeta).
\end{align*}
\end{proposition}

\begin{proof}
We prove, for example, the last assertion. Put
\[
  T_1'=\int^{\oplus}T_1(\zeta)^*\,d\nu(\zeta).
\]
For \(x,y\in\mathfrak H\),
\begin{align*}
  (T_1'x\mid y)
  &=\int(T_1(\zeta)^*x(\zeta)\mid y(\zeta))\,d\nu(\zeta)\\
  &=\int(x(\zeta)\mid T_1(\zeta)y(\zeta))\,d\nu(\zeta)
   =(x\mid T_1y),
\end{align*}
and therefore \(T_1'=T_1^*\).
\end{proof}

There are analogous properties for decomposable linear mappings.

Let
\[
  T=\int^{\oplus}T(\zeta)\,d\nu(\zeta)
\]
% Source PDF page 167 (printed p. 160)
\phantomsection\label{srcpage:160}
be a decomposable operator. By Propositions~2 and~3, \(T\) is unitary
(respectively, Hermitian, partially isometric, and so forth) if and only
if \(T(\zeta)\) is unitary (respectively, Hermitian, partially isometric,
and so forth) almost everywhere.

Let \(\zeta\mapsto E(\zeta)\) be a measurable field of projections, and
put
\[
  E=\int^{\oplus}E(\zeta)\,d\nu(\zeta),
\]
which is a projection in \(\mathcal L(\mathfrak H)\). Let
\[
  \mathfrak K(\zeta)=E(\zeta)\mathfrak H(\zeta),
  \qquad
  \mathfrak K=E\mathfrak H.
\]
A field \(x\in\mathfrak H\) belongs to \(\mathfrak K\) if and only if
\(x=Ex\), that is, if and only if
\(x(\zeta)=E(\zeta)x(\zeta)\) almost everywhere. Thus
\(\mathfrak K\) is the Hilbert integral of the measurable field of
subspaces \(\zeta\mapsto\mathfrak K(\zeta)\) (\S~1, Proposition~9).

\begin{proposition}\label{prop:II-2-4}
Let
\[
  T_i=\int^{\oplus}T_i(\zeta)\,d\nu(\zeta)
  \quad(i=1,2,\ldots),
  \qquad
  T=\int^{\oplus}T(\zeta)\,d\nu(\zeta)
\]
be decomposable operators.
\begin{enumerate}
\item If \(T_i\) converges strongly to \(T\), there exists a subsequence
  \((T_{n_k})\) such that, almost everywhere,
  \(T_{n_k}(\zeta)\) converges strongly to \(T(\zeta)\).
\item If, almost everywhere, \(T_i(\zeta)\) converges strongly to
  \(T(\zeta)\), and if \(\sup_i\lVert T_i\rVert<+\infty\), then
  \(T_i\) converges strongly to \(T\).
\end{enumerate}
\end{proposition}

\begin{proof}
Suppose that \(T_i\) converges strongly to \(T\). Put
\[
  a=\sup_i\lVert T_i\rVert<+\infty.
\]
Except on a null set \(N\), we have
\(\lVert T_i(\zeta)\rVert\leq a\) for every \(i\). Let \((x_j)\) be a
fundamental sequence of measurable fields of vectors such that
\(x_j\in\mathfrak H\) for every \(j\). We have
\(\lVert T_i x_j-Tx_j\rVert\to0\) as \(i\to+\infty\). Hence (\S~1,
Proposition~5) there exist a null set \(N_j\) and a subsequence
\((T_{n_k})\) such that
\[
  \lVert T_{n_k}(\zeta)x_j(\zeta)-T(\zeta)x_j(\zeta)\rVert\longrightarrow0
  \quad(k\to+\infty)
\]
for \(\zeta\notin N_j\). Let \(N'\) be the union of the \(N_j\), which is
a null set. By the diagonal procedure, the sequence
\((n_1,n_2,\ldots)\) may be assumed independent of \(j\). Then, for
\(\zeta\notin N\cup N'\), the preceding convergence holds for every
\(j\), and
\(\lVert T_{n_k}(\zeta)\rVert\leq a\); hence \(T_{n_k}(\zeta)\) converges
strongly to \(T(\zeta)\).

% Source PDF page 168 (printed p. 161)
\phantomsection\label{srcpage:161}
Now suppose that, almost everywhere, \(T_i(\zeta)\) converges strongly to
\(T(\zeta)\), and that \(a=\sup_i\lVert T_i\rVert<+\infty\). Except on a
null set, \(\lVert T_i(\zeta)\rVert\leq a\) for every \(i\), and therefore
\(\lVert T(\zeta)\rVert\leq a\). Retain the preceding notation. For every
\(j\),
\[
  \lVert T_i(\zeta)x_j(\zeta)-T(\zeta)x_j(\zeta)\rVert^2
\]
converges to zero almost everywhere as \(i\to+\infty\); moreover,
\[
  \lVert T_i(\zeta)x_j(\zeta)-T(\zeta)x_j(\zeta)\rVert^2
  \leq(2a\lVert x_j(\zeta)\rVert)^2,
\]
and the function \(\zeta\mapsto\lVert x_j(\zeta)\rVert^2\) is integrable.
Consequently,
\[
  \lVert T_i x_j-Tx_j\rVert\longrightarrow0
  \quad(i\to+\infty).
\]
For every measurable complex-valued function \(f\) on \(Z\) such that
\[
  \sup_{\zeta\in Z}|f(\zeta)|=b<+\infty,
\]
we have
\[
  \lVert(T_i-T)(fx_j)\rVert
  \leq b\lVert(T_i-T)x_j\rVert\longrightarrow0
  \quad(i\to+\infty).
\]
Now, if each function \(\zeta\mapsto\lVert x_j(\zeta)\rVert\) is assumed
bounded, the \(fx_j\) form a total set in \(\mathfrak H\) (\S~1,
Proposition~7). Thus \(T_i\) converges strongly to \(T\).
\end{proof}

\begin{proposition}\label{prop:II-2-5}
There exists in \(\mathfrak H\) a sequence
\[
  T_1=\int^{\oplus}T_1(\zeta)\,d\nu(\zeta),\quad
  T_2=\int^{\oplus}T_2(\zeta)\,d\nu(\zeta),\quad\ldots
\]
of decomposable operators such that, for every \(\zeta\in Z\),
\(\mathcal L(\mathfrak H(\zeta))\) is the von Neumann algebra generated
by the \(T_i(\zeta)\).
\end{proposition}

\begin{proof}
It is enough to prove the proposition when the field
\(\zeta\mapsto\mathfrak H(\zeta)\) on \(Z\) is the constant field
corresponding to a Hilbert space \(\mathfrak H_0\) with a countable base.
Then \(\mathcal L(\mathfrak H_0)\) is the von Neumann algebra generated
by a sequence \(S_1,S_2,\ldots\) of elements, and the constant fields
\(\zeta\mapsto S_i\) may be taken for the fields
\(\zeta\mapsto T_i(\zeta)\).
\end{proof}

Bibliography: \ArticleRef{28}, \ArticleRef{48}, \ArticleRef{52},
\ArticleRef{57}, \ArticleRef{80}, \ArticleRef{100}, \ArticleRef{119},
\ArticleRef{123}.

\NumberedParagraph{4}{Diagonalizable operators}
\label{no:II-2-4}
Let \(L_{\mathbb C}^{\infty}(Z,\nu)\) be the set of essentially bounded
measurable complex-valued functions on \(Z\), in which two functions
equal almost everywhere are identified. If \(Z\) is a locally compact
space countable at infinity, let \(L_{\infty}(Z)\) be the set of
complex-valued functions on \(Z\) which are continuous and tend to zero
at infinity. We furnish \(L_{\mathbb C}^{\infty}(Z,\nu)\) and
\(L_{\infty}(Z)\) with their usual structures
% Source PDF page 169 (printed p. 162)
\phantomsection\label{srcpage:162}
of involutive algebras and Banach spaces. If
\(f\in L_{\mathbb C}^{\infty}(Z,\nu)\), or if \(f\in L_{\infty}(Z)\), the
field of operators
\(\zeta\mapsto f(\zeta)I\in\mathcal L(\mathfrak H(\zeta))\) is measurable
and essentially bounded. We denote by \(T_f\) the corresponding operator
in \(\mathcal L(\mathfrak H)\).

\begin{definition}\label{def:II-2-3}
Operators of the form \(T_f\), where
\(f\in L_{\mathbb C}^{\infty}(Z,\nu)\) (respectively,
\(f\in L_{\infty}(Z)\)), are called \emph{diagonalizable}
(respectively, \emph{continuously diagonalizable}).
\end{definition}

By Proposition~3, the mapping \(f\mapsto T_f\) is a homomorphism of the
involutive algebra \(L_{\mathbb C}^{\infty}(Z,\nu)\) (respectively,
\(L_{\infty}(Z)\)) onto the involutive algebra \(\mathfrak Z\)
(respectively, \(\mathfrak Y\)) of diagonalizable (respectively,
continuously diagonalizable) operators. These homomorphisms are called
\emph{canonical}.

\textbf{Remark.} Let \(Z'\) be the measurable set of \(\zeta\in Z\) such
that \(\mathfrak H(\zeta)\neq0\), let \(\chi_{Z'}\) be its characteristic
function, and put \(\nu'=\chi_{Z'}\nu\). It is immediate that the field
\(\zeta\mapsto\mathfrak H(\zeta)\), with the same measurable fields of
vectors, is also \(\nu'\)-measurable, and that
\[
  \int^{\oplus}\mathfrak H(\zeta)\,d\nu(\zeta)
  =\int^{\oplus}\mathfrak H(\zeta)\,d\nu'(\zeta).
\]
The measurable fields of operators, the decomposable operators, and the
diagonalizable operators are the same for \(\nu\) and \(\nu'\). The
canonical homomorphism of \(L_{\mathbb C}^{\infty}(Z,\nu)\) onto
\(\mathfrak Z\) is the composite of the evident canonical homomorphism
of \(L_{\mathbb C}^{\infty}(Z,\nu)\) onto
\(L_{\mathbb C}^{\infty}(Z,\nu')\), and the canonical homomorphism of
\(L_{\mathbb C}^{\infty}(Z,\nu')\) onto \(\mathfrak Z\), the latter being
an isomorphism by Proposition~2. Thus the canonical homomorphism of
\(L_{\mathbb C}^{\infty}(Z,\nu)\) onto \(\mathfrak Z\) is an isomorphism
if and only if \(\mathfrak H(\zeta)\neq0\) almost everywhere (in which
case \(\nu=\nu'\)). On the other hand, the canonical homomorphism of
\(L_{\infty}(Z)\) into \(L_{\mathbb C}^{\infty}(Z,\nu)\) is injective if
and only if the support of \(\nu\) is \(Z\). Thus, when
\(\mathfrak H(\zeta)\neq0\) almost everywhere and the support of \(\nu\)
is \(Z\), \(Z\) is identified with the spectrum of \(\mathfrak Y\) in
such a way that \(f(\zeta)=\zeta(T_f)\) for \(\zeta\in Z\) and
\(f\in L_{\infty}(Z)\).

\begin{proposition}\label{prop:II-2-6}
Suppose that \(Z\) is locally compact and countable at infinity, that
\(\mathfrak H(\zeta)\neq0\) almost everywhere, and that \(\nu\) has
support \(Z\). Then \(\nu\) is a basic measure on the spectrum
\(Z\) of \(\mathfrak Y\).
\end{proposition}

\begin{proof}
If \(x\in\mathfrak H\) and \(f\in L_{\infty}(Z)\), then
\[
  (T_fx\mid x)=\int f(\zeta)(x(\zeta)\mid x(\zeta))\,d\nu(\zeta),
\]
which proves that, with the notation of Chapter~I, \S~7,
\[
  d\nu_{x,x}(\zeta)=\lVert x(\zeta)\rVert^2\,d\nu(\zeta).
\]
Thus the measure \(\nu_{x,x}\) has base \(\nu\). Conversely, let \(N\)
be a subset of \(Z\) which is null for all the measures \(\nu_{x,x}\).
Then, for every \(x\in\mathfrak H\),
% Source PDF page 170 (printed p. 163)
\phantomsection\label{srcpage:163}
\(x(\zeta)=0\) \(\nu\)-almost everywhere on \(N\). Since
\(\mathfrak H(\zeta)\neq0\) almost everywhere, there exists a field
\(x\in\mathfrak H\) such that \(x(\zeta)\neq0\) \(\nu\)-almost everywhere
on \(Z\). Hence \(N\) is \(\nu\)-null.
\end{proof}

\begin{proposition}\label{prop:II-2-7}
\begin{enumerate}
\item The algebra \(\mathfrak Z\) is an abelian von Neumann algebra. If
  \(Z\) is locally compact and countable at infinity, \(\mathfrak Z\) is the
  weak closure of \(\mathfrak Y\).
\item The canonical homomorphism of
  \(L_{\mathbb C}^{\infty}(Z,\nu)\) onto \(\mathfrak Z\) is continuous
  when \(L_{\mathbb C}^{\infty}(Z,\nu)\) is furnished with its weak
  topology as the dual of \(L_{\mathbb C}^{1}(Z,\nu)\), and
  \(\mathfrak Z\) with the weak operator topology.
\item \(\mathfrak Z'\) is of countable type.
\end{enumerate}
\end{proposition}

\begin{proof}
We prove (ii). Let
\[
  x,y\in\mathfrak H,
  \qquad f\in L_{\mathbb C}^{\infty}(Z,\nu).
\]
Then
\[
  (T_fx\mid y)=\int f(\zeta)(x(\zeta)\mid y(\zeta))\,d\nu(\zeta)
\]
depends continuously on \(f\) when
\(L_{\mathbb C}^{\infty}(Z,\nu)\) is furnished with its weak topology as
the dual of \(L_{\mathbb C}^{1}(Z,\nu)\), since the function
\(\zeta\mapsto(x(\zeta)\mid y(\zeta))\) is integrable.

We prove (i). By the remarks preceding Proposition~6, we may suppose
that \(\mathfrak H(\zeta)\neq0\) almost everywhere. The unit ball of
\(L_{\mathbb C}^{\infty}(Z,\nu)\) is weakly compact. Hence the unit ball
of \(\mathfrak Z\) is weakly compact, and consequently \(\mathfrak Z\)
is a von Neumann algebra (Chapter~I, \S~3, Theorem~2). Suppose that
\(Z\) is locally compact and countable at infinity. Since \(L_{\infty}(Z)\)
is weakly dense in \(L_{\mathbb C}^{\infty}(Z,\nu)\), \(\mathfrak Y\) is
weakly dense in \(\mathfrak Z\).

We prove (iii). The fields \(x_i\) of \S~1, Proposition~7 may be assumed
to belong to \(\mathfrak H\). By that proposition, the \(x_i\) form a
totalizing set for \(\mathfrak Z\), and hence a separating set for
\(\mathfrak Z'\).
\end{proof}

In \S~6 we shall prove the converses of the preceding results.

Bibliography: \ArticleRef{28}, \ArticleRef{80}, \ArticleRef{100}.

\NumberedParagraph{5}{Characterization of decomposable mappings}
\label{no:II-2-5}
Let, as before, \(\zeta\mapsto\mathfrak H(\zeta)\) and
\(\zeta\mapsto\mathfrak H'(\zeta)\) be two \(\nu\)-measurable fields of
complex Hilbert spaces on \(Z\), and put
\[
  \mathfrak H=\int^{\oplus}\mathfrak H(\zeta)\,d\nu(\zeta),
  \qquad
  \mathfrak H'=\int^{\oplus}\mathfrak H'(\zeta)\,d\nu(\zeta).
\]
Every function \(f\in L_{\mathbb C}^{\infty}(Z,\nu)\) defines a
diagonalizable operator \(T_f\in\mathcal L(\mathfrak H)\) and a
diagonalizable operator \(T_f'\in\mathcal L(\mathfrak H')\). If
\(T\in\mathcal L(\mathfrak H,\mathfrak H')\) is decomposable, it is clear
that \(T_f'T=TT_f\) for every
\(f\in L_{\mathbb C}^{\infty}(Z,\nu)\). Conversely:

% Source PDF page 171 (printed p. 164)
\phantomsection\label{srcpage:164}
\begin{theorem}\label{thm:II-2-1}
Let \(T\) be a continuous linear mapping from \(\mathfrak H\) into
\(\mathfrak H'\) such that, for every
\(f\in L_{\mathbb C}^{\infty}(Z,\nu)\),
\(T_f'T=TT_f\). Then \(T\) is decomposable.
\end{theorem}

\begin{proof}
Let \((x_1,x_2,\ldots)\) be a fundamental sequence of measurable fields
of vectors on \(Z\), relative to
\(\zeta\mapsto\mathfrak H(\zeta)\), such that
\(\sup_{\zeta\in Z}\lVert x_i(\zeta)\rVert<+\infty\) and
\(x_i\in\mathfrak H\) for every \(i\). Put
\(y_i=Tx_i\in\mathfrak H'\). For every finite sequence
\(\rho=(\rho_1,\rho_2,\ldots,\rho_n)\) of rational complex numbers, put
\[
  x_{\rho}=\sum_{i=1}^{n}\rho_i x_i,
  \qquad
  y_{\rho}=\sum_{i=1}^{n}\rho_i y_i.
\]
For every \(f\in L_{\mathbb C}^{\infty}(Z,\nu)\), we have
\(TT_f x_{\rho}=T_f'Tx_{\rho}=T_f'y_{\rho}\), and therefore
\[
  \int |f(\zeta)|^2\lVert y_{\rho}(\zeta)\rVert^2\,d\nu(\zeta)
  \leq\lVert T\rVert^2
  \int |f(\zeta)|^2\lVert x_{\rho}(\zeta)\rVert^2\,d\nu(\zeta).
\]
Since \(f\) is arbitrary,
\[
  \lVert y_{\rho}(\zeta)\rVert
  \leq\lVert T\rVert\,\lVert x_{\rho}(\zeta)\rVert
\]
except on a null set \(N_{\rho}\). Let \(N\) be the null set which is
the union of the \(N_{\rho}\) as \(\rho\) ranges over the set of finite
sequences of rational complex numbers. For \(\zeta\notin N\), and for
all rational complex numbers \(\rho_1,\rho_2,\ldots,\rho_n\),
\[
  \left\lVert\sum_{i=1}^{n}\rho_i y_i(\zeta)\right\rVert
  \leq\lVert T\rVert
  \left\lVert\sum_{i=1}^{n}\rho_i x_i(\zeta)\right\rVert.
\]
Thus there exists a continuous linear mapping \(T(\zeta)\) of
\(\mathfrak H(\zeta)\) into \(\mathfrak H'(\zeta)\) such that
\(T(\zeta)x_i(\zeta)=y_i(\zeta)\) for every \(i\), and
\(\lVert T(\zeta)\rVert\leq\lVert T\rVert\). Put \(T(\zeta)=0\) for
\(\zeta\in N\). The field \(\zeta\mapsto T(\zeta)\) is measurable and
therefore defines a decomposable mapping \(T_1\) from \(\mathfrak H\)
into \(\mathfrak H'\). We have \(T_1x_i=y_i=Tx_i\) for every \(i\), and
so
\[
  T_1(T_fx_i)=T_f'T_1x_i=T_f'Tx_i=T(T_fx_i)
  \qquad\text{for every }f\in L_{\mathbb C}^{\infty}(Z,\nu).
\]
But the vectors \(T_fx_i\) form a total system in \(\mathfrak H\)
(\S~1, Proposition~7). Therefore
\[
  T=T_1=\int^{\oplus}T(\zeta)\,d\nu(\zeta).
\]
\end{proof}

\begin{corollary*}
An operator \(T\in\mathcal L(\mathfrak H)\) is decomposable if and only
if it commutes with the diagonalizable operators.
\end{corollary*}

% Source PDF page 172 (printed p. 165)
\phantomsection\label{srcpage:165}
\textbf{Remark.} Suppose \(Z\) is locally compact and countable at infinity.
Let \(T\in\mathcal L(\mathfrak H,\mathfrak H')\) satisfy
\(T_f'T=TT_f\) for every \(f\in L_{\infty}(Z)\). Then \(T\) is
decomposable. This follows from Theorem~1 and Proposition~7.

Bibliography: \ArticleRef{28}, \ArticleRef{80}, \ArticleRef{117}.

\NumberedParagraph{6}{Constant fields of operators}
\label{no:II-2-6}
\begin{proposition}\label{prop:II-2-8}
Let \(\zeta\mapsto\mathfrak H(\zeta)\) be a \(\nu\)-measurable field of
complex Hilbert spaces on \(Z\), put
\[
  \mathfrak H=\int^{\oplus}\mathfrak H(\zeta)\,d\nu(\zeta),
\]
and let \(\zeta\mapsto\mathfrak K(\zeta)\) be the constant field on
\(Z\) corresponding to a Hilbert space \(\mathfrak K_0\) with a
countable base. Let
\[
  S=\int^{\oplus}S(\zeta)\,d\nu(\zeta)\in\mathcal L(\mathfrak H),
  \qquad T\in\mathcal L(\mathfrak K_0).
\]
Then, identifying
\[
  \left(\int^{\oplus}\mathfrak H(\zeta)\,d\nu(\zeta)\right)
  \otimes\mathfrak K_0
  \quad\text{and}\quad
  \int^{\oplus}\bigl(\mathfrak H(\zeta)\otimes
  \mathfrak K(\zeta)\bigr)\,d\nu(\zeta),
\]
we have
\[
  S\otimes T=\int^{\oplus}\bigl(S(\zeta)\otimes T\bigr)\,d\nu(\zeta).
\]
\end{proposition}

\begin{proof}
Let
\(x\in\int^{\oplus}\mathfrak H(\zeta)\,d\nu(\zeta)\) and
\(y\in\mathfrak K_0\). Then \(x\otimes y\) is identified with the field
\(\zeta\mapsto x(\zeta)\otimes y\). We have
\[
  (S\otimes T)(x\otimes y)=(Sx)\otimes(Ty),
\]
which is identified with the field
\[
  \zeta\longmapsto(Sx)(\zeta)\otimes Ty
  =S(\zeta)x(\zeta)\otimes Ty
  =\bigl(S(\zeta)\otimes T\bigr)(x(\zeta)\otimes y).
\]
Consequently,
\[
  (S\otimes T)(x\otimes y)
  =\left(\int^{\oplus}\bigl(S(\zeta)\otimes T\bigr)
  \,d\nu(\zeta)\right)(x\otimes y).
\]
\end{proof}

\begin{corollary*}
Let \(\mathfrak K_0\) be a complex Hilbert space with a countable base,
let \(\zeta\mapsto\mathfrak K(\zeta)\) be the corresponding constant
field on \(Z\), and put
\[
  \mathfrak K=\int^{\oplus}\mathfrak K(\zeta)\,d\nu(\zeta)
  =L_{\mathbb C}^{2}(Z,\nu)\otimes\mathfrak K_0.
\]
For every \(T\in\mathcal L(\mathfrak K_0)\),
\[
  I\otimes T=\int^{\oplus}T(\zeta)\,d\nu(\zeta),
\]
where \(T(\zeta)=T\) for every \(\zeta\in Z\).
\end{corollary*}

% Source PDF page 173 (printed p. 166)
\phantomsection\label{srcpage:166}
\begin{proof}
Apply Proposition~8 when \(\zeta\mapsto\mathfrak H(\zeta)\) is the
constant field corresponding to the one-dimensional Hilbert space
\(\mathbb C\), and \(S(\zeta)=I\) for every \(\zeta\).
\end{proof}

\begin{lemma}\label{lem:II-2-1}
Let \(\mathfrak K_0\) be a complex Hilbert space with a countable base,
and let \(\mathcal U\) be the set of unitary operators in
\(\mathcal L(\mathfrak K_0)\). For the weak topology, and a fortiori for
the strong topology, \(\mathcal U\) is a Borel subset of
\(\mathcal L(\mathfrak K_0)\).
\end{lemma}

\begin{proof}
For every \(x\in\mathfrak K_0\) and every \(\lambda\geq0\), the set of
\(T\in\mathcal L(\mathfrak K_0)\) such that
\(\lVert Tx\rVert\leq\lambda\) is weakly closed. Hence the set of
\(T\in\mathcal L(\mathfrak K_0)\) such that
\(\lVert Tx\rVert=\lambda\) is Borel for the weak topology, since it is
defined by the conditions
\[
  \lVert Tx\rVert\leq\lambda,
  \qquad
  \lVert Tx\rVert>\lambda-\frac1n
  \quad(n=1,2,\ldots).
\]
Now, if \((x_i)\) is an everywhere dense sequence in \(\mathfrak K_0\),
\(\mathcal U\) is the set of \(T\in\mathcal L(\mathfrak K_0)\) such that
\[
  \lVert Tx_i\rVert=\lVert x_i\rVert,
  \qquad
  \lVert T^*x_i\rVert=\lVert x_i\rVert
  \quad\text{for every }i,
\]
and is therefore Borel for the weak topology.
\end{proof}

Before stating Lemma~2, observe that, if \(Y\) is any subset of \(Z\),
there exists a \(\nu\)-measurable subset \(X\) of \(Z\), containing
\(Y\), with the following property: if \(X'\) is another
\(\nu\)-measurable subset of \(Z\) containing \(Y\), then \(X'\)
contains \(X\) up to a null set. (It is enough to take a Borel set
\(X\) containing \(Y\) and having measure equal to the outer measure of
\(Y\).) The set \(X\) is determined by \(Y\), up to a null set; we shall
say that \(X\) is the \emph{measurable envelope} of \(Y\).

\begin{lemma}\label{lem:II-2-2}
Let \(Y\) be a subset of \(Z\), and \(X\) its measurable envelope. Let
\(\mathscr E:\zeta\mapsto\mathfrak H(\zeta)\) and
\(\mathscr E':\zeta\mapsto\mathfrak H'(\zeta)\) be
\(\nu\)-measurable fields of complex Hilbert spaces on \(Z\), and let
\(J\) be a countable set. For every \(j\in J\), let
\(\zeta\mapsto T_j(\zeta)\in\mathcal L(\mathfrak H(\zeta))\)
(respectively,
\(\zeta\mapsto T_j'(\zeta)\in\mathcal L(\mathfrak H'(\zeta))\)) be a
measurable field of operators relative to \(\mathscr E\) (respectively,
\(\mathscr E'\)). Suppose that, for every \(\zeta\in Y\), there exists
an isomorphism \(U(\zeta)\) of \(\mathfrak H(\zeta)\) onto
\(\mathfrak H'(\zeta)\) such that
\[
  U(\zeta)T_j(\zeta)U(\zeta)^{-1}=T_j'(\zeta)
  \qquad(j\in J).
\]
Suppose \(\nu\) is standard. Then, after modification of \(\mathscr E\)
and \(\mathscr E'\) on null sets, there exists an isomorphism
\(\zeta\mapsto V(\zeta)\) of \(\mathscr E|X\) onto
\(\mathscr E'|X\) such that
\[
  V(\zeta)T_j(\zeta)V(\zeta)^{-1}=T_j'(\zeta)
  \qquad(j\in J,\ \zeta\in X).
\]
\end{lemma}

\begin{proof}
We may suppose that \(X=Z\), and that the given set of fields of
operators is stable under taking adjoints. We have
\(\dim\mathfrak H(\zeta)=\dim\mathfrak H'(\zeta)\) everywhere on \(Y\),
and hence almost everywhere on \(Z\). We are thus reduced to the case in
which \(\mathscr E\) and \(\mathscr E'\) are both the constant field on
\(Z\) defined by a Hilbert space \(\mathfrak K_0\) with a countable
base. Let \((x_i)\) be a total sequence in \(\mathfrak K_0\). The
mappings
\[
  \zeta\longmapsto T_j(\zeta)x_i,
  \qquad
  \zeta\longmapsto T_j'(\zeta)x_i
  \quad(j\in J;\ i=1,2,\ldots)
\]
are measurable. Moreover, we may suppose that \(Z\) is compact and
metrizable. Thus \(Z\) is, up to a null set, the union
% Source PDF page 174 (printed p. 167)
\phantomsection\label{srcpage:167}
of a sequence of compact subsets such that the restrictions to each of
these subsets of the mappings
\[
  \zeta\longmapsto T_j(\zeta)x_i,
  \qquad
  \zeta\longmapsto T_j'(\zeta)x_i
  \quad(j\in J;\ i=1,2,\ldots)
\]
are continuous for the strong topology of \(\mathfrak K_0\). It is
enough to construct the \(V(\zeta)\) on each of these compact subsets.
We may therefore suppose from now on that all the preceding mappings are
continuous on \(Z\) for the strong topology of \(\mathfrak K_0\).

Let \(\mathcal L_1\) be the unit ball of
\(\mathcal L(\mathfrak K_0)\); furnished with the weak topology,
\(\mathcal L_1\) is a compact metrizable space. For every
\(\zeta\in Z\), let \(\mathcal M_{\zeta}\) be the set of
\(T\in\mathcal L_1\) satisfying the following conditions:
\begin{enumerate}
\item \(T_j'(\zeta)T=TT_j(\zeta)\) for every \(j\in J\);
\item \(T\) belongs to the set \(\mathcal U\) of unitary operators on
  \(\mathfrak K_0\).
\end{enumerate}
For \(\zeta\in Y\), we have \(U(\zeta)\in\mathcal M_{\zeta}\), and so
\(\mathcal M_{\zeta}\neq\varnothing\). Let \(\mathcal A_1\)
(respectively, \(\mathcal A_2\)) be the set of pairs
\((\zeta,T)\in Z\times\mathcal L_1\) satisfying condition 1
(respectively, condition 2). Then \(\mathcal A_1\) is the set of pairs
\((\zeta,T)\in Z\times\mathcal L_1\) such that
\[
  (Tx_i\mid T_j'(\zeta)^*x_k)
  =(T_j(\zeta)x_i\mid T^*x_k)
  \qquad(j\in J;\ i,k=1,2,\ldots).
\]
Since the functions
\[
  (\zeta,T)\longmapsto(Tx_i\mid T_j'(\zeta)^*x_k),
  \qquad
  (\zeta,T)\longmapsto(T_j(\zeta)x_i\mid T^*x_k)
\]
are continuous on \(Z\times\mathcal L_1\), \(\mathcal A_1\) is closed
in \(Z\times\mathcal L_1\). By Lemma~1,
\(\mathcal A_2=Z\times\mathcal U\) is Borel in
\(Z\times\mathcal L_1\); hence \(\mathcal A_1\cap\mathcal A_2\) is
Borel in \(Z\times\mathcal L_1\). It follows from Appendix~V that there
exists a \(\nu\)-measurable mapping
\(\zeta\mapsto V(\zeta)\), defined on \(X\), with values in
\(\mathcal L_1\), such that \(V(\zeta)\in\mathcal M_{\zeta}\) for
\(\zeta\in X\). Then \(V(\zeta)\) is unitary, and
\[
  T_j'(\zeta)V(\zeta)=V(\zeta)T_j(\zeta)
  \qquad(j\in J).
\]
\end{proof}

\begin{theorem}\label{thm:II-2-2}
Let \(\mathfrak K_0\) be a complex Hilbert space with a countable base,
and let \((S_{\iota})_{\iota\in I}\) be a family of operators in
\(\mathcal L(\mathfrak K_0)\). Let
\(\zeta\mapsto\mathfrak H(\zeta)\) be a \(\nu\)-measurable field of
complex Hilbert spaces on \(Z\), put
\[
  \mathfrak H=\int^{\oplus}\mathfrak H(\zeta)\,d\nu(\zeta),
\]
and, for every \(\iota\in I\), let
\(\zeta\mapsto T_{\iota}(\zeta)\) be an essentially bounded measurable
field of operators on \(Z\); put
\[
  T_{\iota}=\int^{\oplus}T_{\iota}(\zeta)\,d\nu(\zeta)
  \in\mathcal L(\mathfrak H).
\]
Suppose that \(\nu\) is standard and that, for every \(\zeta\in Z\),
there exists an isomorphism \(U(\zeta)\) of the Hilbert space
\(\mathfrak H(\zeta)\) onto the Hilbert space \(\mathfrak K_0\) such
that
\[
  U(\zeta)^{-1}S_{\iota}U(\zeta)=T_{\iota}(\zeta)
  \qquad(\iota\in I).
\]

% Source PDF page 175 (printed p. 168)
\phantomsection\label{srcpage:168}
Then there exists an isomorphism of \(\mathfrak H\) onto
\(L_{\mathbb C}^{2}(Z,\nu)\otimes\mathfrak K_0\) which transforms, for
every \(\iota\in I\), \(T_{\iota}\) into \(I\otimes S_{\iota}\).
\end{theorem}

\begin{proof}
The spaces \(\mathfrak H(\zeta)\) all have the same Hilbert dimension as
\(\mathfrak K_0\). We are therefore reduced to the case in which
\(\zeta\mapsto\mathfrak H(\zeta)\) is the constant field on \(Z\)
corresponding to \(\mathfrak K_0\). Let \(J\) be a countable subset of
\(I\) such that the family \((S_{\iota})_{\iota\in J}\) is weakly
everywhere dense in the family \((S_{\iota})_{\iota\in I}\). By Lemma~2,
for every \(\zeta\in Z\) there exists a unitary operator
\(V(\zeta)\) on \(\mathfrak K_0\), depending measurably on \(\zeta\),
such that
\[
  V(\zeta)^{-1}S_{\iota}V(\zeta)
  =U(\zeta)^{-1}S_{\iota}U(\zeta)
\]
for every \(\iota\in J\), and hence for every \(\iota\in I\). Then
\[
  \int^{\oplus}V(\zeta)\,d\nu(\zeta)=V
\]
is a unitary operator on
\(\mathfrak H=L_{\mathbb C}^{2}(Z,\nu)\otimes\mathfrak K_0\) which
transforms \(T_{\iota}\) into the operator defined by the constant field
\(\zeta\mapsto S_{\iota}\), namely \(I\otimes S_{\iota}\) (corollary to
Proposition~8).
\end{proof}

\begin{theorem}\label{thm:II-2-3}
For \(i=1,2,\ldots\), let \(\mathfrak K_i\) be a complex Hilbert space
with a countable base, and let
\((S_{\iota}^{i})_{\iota\in I}\) be a family of operators in
\(\mathcal L(\mathfrak K_i)\). Let
\(\zeta\mapsto\mathfrak H(\zeta)\) be a \(\nu\)-measurable field of
complex Hilbert spaces on \(Z\), put
\[
  \mathfrak H=\int^{\oplus}\mathfrak H(\zeta)\,d\nu(\zeta),
\]
and, for every \(\iota\in I\), let
\(\zeta\mapsto T_{\iota}(\zeta)\) be an essentially bounded measurable
field of operators on \(Z\); put
\[
  T_{\iota}=\int^{\oplus}T_{\iota}(\zeta)\,d\nu(\zeta)
  \in\mathcal L(\mathfrak H).
\]
Let \(Z_1,Z_2,\ldots\) be a sequence of subsets of \(Z\) whose union is
\(Z\). Suppose that \(\nu\) is standard and that, for every
\(\zeta\in Z_i\), there exists an isomorphism \(U_i(\zeta)\) of
\(\mathfrak H(\zeta)\) onto \(\mathfrak K_i\) which transforms each
\(T_{\iota}(\zeta)\) into \(S_{\iota}^{i}\).

Then, if \(\nu\neq0\), there exists an isomorphism of one of the
\(\mathfrak K_i\) onto a closed vector subspace \(\mathfrak H'\) of
\(\mathfrak H\), reducing the \(T_{\iota}\), which transforms, for every
\(\iota\in I\), \(S_{\iota}^{i}\) into
\((T_{\iota})_{\mathfrak H'}\).
\end{theorem}

\begin{proof}
By restricting to a non-null measurable subset of \(Z\) on which the
Hilbert dimension of \(\mathfrak H(\zeta)\) is constant, we are first
reduced to the case in which \(\zeta\mapsto\mathfrak H(\zeta)\) is a
constant field on \(Z\), corresponding to a Hilbert space
\(\mathfrak K_0\) which may then be assumed identical with all the
\(\mathfrak K_i\). Arguing as for Theorem~2, we may suppose that \(I\) is
countable. By Lemma~2, we are next reduced to the case in which the
\(Z_i\) are measurable and \(U_i(\zeta)\) depends measurably on
\(\zeta\) for \(\zeta\in Z_i\). Then one of the \(Z_i\), say \(Z_1\),
% Source PDF page 176 (printed p. 169)
\phantomsection\label{srcpage:169}
is non-null. Thus, by restricting to a non-null measurable subset of
\(Z_1\), we are reduced to proving the conclusion of the theorem under
the hypotheses of Theorem~2, whose notation we now adopt, with in
addition \(\nu\neq0\). Let
\(f\in L_{\mathbb C}^{2}(Z,\nu)\) satisfy \(\lVert f\rVert=1\). Then the
mapping \(x\mapsto f\otimes x\) is an isomorphism \(W\) of
\(\mathfrak K_0\) onto a subspace \(\mathfrak H'\) of
\(L_{\mathbb C}^{2}(Z,\nu)\otimes\mathfrak K_0\). Since
\[
  (I\otimes S_{\iota})(f\otimes x)=f\otimes S_{\iota}x
\]
and
\[
  (I\otimes S_{\iota})^*(f\otimes x)=f\otimes S_{\iota}^*x,
\]
we see that \(\mathfrak H'\) reduces the \(I\otimes S_{\iota}\), and that
\(W\) transforms \(S_{\iota}\) into
\((I\otimes S_{\iota})_{\mathfrak H'}\). Theorem~3 is therefore a
consequence of Theorem~2.
\end{proof}

Bibliography: \ArticleRef{57}, \ArticleRef{58}, \ArticleRef{91}. By
\ArticleRef{466}, the hypothesis that \(\nu\) is standard may be omitted
from Theorems~2 and~3.

\medskip
\noindent\textit{Exercises.}---
\begin{enumerate}
\item Let \(\mathfrak Z\) be the algebra of diagonalizable operators in
  \(\int^{\oplus}\mathfrak H(\zeta)\,d\nu(\zeta)\). We have
  \(\mathfrak Z'=\mathfrak Z\) if and only if the Hilbert dimension of
  \(\mathfrak H(\zeta)\) is zero or one almost everywhere.
\item
  \begin{enumerate}
  \item Let
    \(T=\int^{\oplus}T(\zeta)\,d\nu(\zeta)\) be a decomposable operator
    on \(\int^{\oplus}\mathfrak H(\zeta)\,d\nu(\zeta)\). Let \(E\) be
    the support of \(T\), which is decomposable, and write
    \(E=\int^{\oplus}E(\zeta)\,d\nu(\zeta)\). Show that, almost
    everywhere, \(E(\zeta)\) is the support of \(T(\zeta)\).
    [Let \((x_i)\) be a fundamental sequence of square-integrable fields
    of vectors and put
    \(\mathfrak K(\zeta)=E(\zeta)\mathfrak H(\zeta)\). The vectors
    \(T'T^*x_i\), where \(T'\) ranges over the algebra of diagonalizable
    operators, form a total system in \(E\mathfrak H\); hence, almost
    everywhere, the \(T(\zeta)^*x_i(\zeta)\) form a total system in
    \(\mathfrak K(\zeta)\).]
  \item Let \(T=WS\) be the polar decomposition of \(T\), where
    \[
      W=\int^{\oplus}W(\zeta)\,d\nu(\zeta),
      \qquad
      S=\int^{\oplus}S(\zeta)\,d\nu(\zeta).
    \]
    Using (a), show that
    \(T(\zeta)=W(\zeta)S(\zeta)\) is almost everywhere the polar
    decomposition of \(T(\zeta)\).
  \item Let \(\zeta\mapsto T_1(\zeta)\) be a measurable field of
    operators, and let
    \[
      T_1(\zeta)=W_1(\zeta)S_1(\zeta)
    \]
    be the polar decomposition of \(T_1(\zeta)\). Show that the fields
    \(\zeta\mapsto W_1(\zeta)\) and
    \(\zeta\mapsto S_1(\zeta)\) are measurable. [Reduce to the case in
    which \(\zeta\mapsto T_1(\zeta)\) is essentially bounded and use
    (b).]
  \end{enumerate}
\item Let \(Z\) be locally compact and countable at infinity, let \(\nu\)
  be a positive measure on \(Z\), and let
  \(\zeta\mapsto\mathfrak H(\zeta)\) be a \(\nu\)-measurable field of
  complex Hilbert spaces on \(Z\). Let \(Z'\) be the set of
  \(\zeta\in Z\) for which \(\mathfrak H(\zeta)\neq0\), and put
  \(\nu'=\chi_{Z'}\nu\), where \(\chi_{Z'}\) is the characteristic
  function of \(Z'\). If the support of \(\nu'\) is \(Z\), the canonical
  homomorphism of \(L_{\infty}(Z)\) onto the algebra \(\mathfrak Y\) of
  continuously diagonalizable operators is an isomorphism, so that
  \(Z\) is identified with the spectrum of \(\mathfrak Y\). But if
  \(Z-Z'\) is not \(\nu\)-null, \(\nu\) is not a basic measure on \(Z\).

% Source PDF page 177 (printed p. 170)
\phantomsection\label{srcpage:170}
\item Let \(\mathfrak Z\) be the algebra of diagonalizable operators.
  Prove again that every \(T\in\mathfrak Z'\) is decomposable by the
  following procedure. Reduce to the case in which
  \(\zeta\mapsto\mathfrak H(\zeta)\) is the constant field corresponding
  to a Hilbert space \(\mathfrak H_0\). Then
  \[
    \mathfrak H=L_{\mathbb C}^{2}(Z,\nu)\otimes\mathfrak H_0,
    \qquad
    \mathfrak Z=\mathfrak Z_1\otimes\mathbb C I_{\mathfrak H_0},
  \]
  where \(\mathfrak Z_1\) is the algebra of multiplication operators by
  functions in \(L_{\mathbb C}^{\infty}(Z,\nu)\) on
  \(L_{\mathbb C}^{2}(Z,\nu)\). Then
  \(\mathfrak Z'=\mathfrak Z_1\otimes\mathcal L(\mathfrak H_0)\).
  Moreover, \(\mathfrak Z'\) is of countable type, so \(T\) is the strong
  limit of a sequence of operators of the form
  \[
    T_1\otimes S_1+T_2\otimes S_2+\cdots+T_n\otimes S_n,
  \]
  where \(T_i\in\mathfrak Z_1\) and
  \(S_i\in\mathcal L(\mathfrak H_0)\). Show that an operator
  \(T_i\otimes S_i\) is decomposable, and deduce that \(T\) is
  decomposable.
\item Let \(\zeta\mapsto\mathfrak H(\zeta)\) be a
  \(\nu\)-measurable field of complex Hilbert spaces on \(Z\), and let
  \(Y\) be a metrizable topological space. For every \(y\in Y\), let
  \(\zeta\mapsto T_y(\zeta)\) be a \(\nu\)-measurable field of operators
  such that
  \[
    \lVert T_y(\zeta)\rVert\leq M<+\infty
    \qquad(\zeta\in Z,\ y\in Y),
  \]
  and put
  \[
    T_y=\int^{\oplus}T_y(\zeta)\,d\nu(\zeta).
  \]
  If, for every \(\zeta\in Z\), \(T_y(\zeta)\) is a strongly continuous
  function of \(y\), then \(T_y\) is a strongly continuous function of
  \(y\). [Use Proposition~4(ii).]
\item Let \(Z=[0,1]\), and let \(\nu\) be Lebesgue measure on \(Z\).
  \begin{enumerate}
  \item For \(\zeta\in Z\), let \(I_{\zeta}\) be the set of
    \(f\in L_{\mathbb C}^{\infty}(Z,\nu)\) having the following
    property: there exists a \(\nu\)-measurable subset \(Y\) of \(Z\),
    of density one at \(\zeta\), such that \(f(\zeta')\) tends to zero
    as \(\zeta'\) tends to \(\zeta\) while remaining in \(Y\). Show that
    \(I_{\zeta}\) is a self-adjoint ideal of the involutive algebra
    \(L_{\mathbb C}^{\infty}(Z,\nu)\).
  \item Let \(J_{\zeta}\) be a maximal ideal of
    \(L_{\mathbb C}^{\infty}(Z,\nu)\) containing \(I_{\zeta}\), and
    let \(\varphi_{\zeta}\) be the corresponding character of
    \(L_{\mathbb C}^{\infty}(Z,\nu)\). For
    \(f\in L_{\mathbb C}^{\infty}(Z,\nu)\), put
    \(f'(\zeta)=\varphi_{\zeta}(f)\). Show that the function
    \(\zeta\mapsto f'(\zeta)\) agrees almost everywhere with \(f\).
    [Observe that, for almost every \(\zeta\in Z\), \(f(\zeta')\) tends
    to \(f(\zeta)\) as \(\zeta'\) tends to \(\zeta\) on a suitably
    chosen set of density one at \(\zeta\).]
  \item Let \(\zeta\mapsto\mathfrak H(\zeta)\) be a
    \(\nu\)-measurable field of finite-dimensional complex Hilbert
    spaces on \(Z\), put
    \(\mathfrak H=\int^{\oplus}\mathfrak H(\zeta)\,d\nu(\zeta)\), and
    let \(\mathfrak Z\) be the algebra of diagonalizable operators on
    \(\mathfrak H\). Show that, for every \(T\in\mathfrak Z'\), one can
    choose an essentially bounded \(\nu\)-measurable field of operators
    \(\zeta\mapsto T(\zeta)\) such that:
    \begin{enumerate}
    \item \(T=\int^{\oplus}T(\zeta)\,d\nu(\zeta)\);
    \item the mapping \(T\mapsto T(\zeta)\) is a homomorphism of
      \(\mathfrak Z'\) onto \(\mathcal L(\mathfrak H(\zeta))\), for
      every \(\zeta\).
    \end{enumerate}
    [Use (b).] \ArticleRef{57}. See also J.~Dieudonn\'e, ``Sur le
    th\'eor\`eme de Lebesgue--Nikodym (IV),'' \emph{J. Indian Math.
    Soc.}, vol.~15 (1951), pp.~77--86.
  \end{enumerate}
\item Let \(Z=[0,1]\), let \(\nu\) be Lebesgue measure on \(Z\), and
  put
  \[
    \mathfrak H=L_{\mathbb C}^{2}(Z,\nu)
    =\int^{\oplus}\mathfrak H(\zeta)\,d\nu(\zeta),
  \]
  where \(\zeta\mapsto\mathfrak H(\zeta)\) is the constant field
  corresponding to the one-dimensional Hilbert space \(\mathbb C\).
  Let \(T\) be the identity mapping of \(\mathfrak H\) onto itself.
  Show that there is no field
% Source PDF page 178 (printed p. 171)
\phantomsection\label{srcpage:171}
  \(\zeta\mapsto T(\zeta)\in
  \mathcal L(\mathfrak H,\mathfrak H(\zeta))\) such that, for every
  \(x\in\mathfrak H\),
  \[
    Tx=\int^{\oplus}\bigl(T(\zeta)x\bigr)\,d\nu(\zeta).
  \]
  [Let \(x_n\in\mathfrak H\) be the function
  \(\zeta\mapsto e^{2\pi i n\zeta}\). The operator \(T(\zeta)\) would
  be identified with a continuous linear functional \(f_{\zeta}\) on
  \(\mathfrak H\) such that, almost everywhere on \(Z\),
  \(f_{\zeta}(x_n)=e^{2\pi i n\zeta}\) for every \(n\). But
  \(\sum_n|e^{2\pi i n\zeta}|^2=+\infty\).]
\end{enumerate}

\section{Fields of von Neumann Algebras}\label{sec:II-3}

\NumberedParagraph{1}{Preliminary theorem}\label{no:II-3-1}
Let, as before, \(Z\) be a Borel space, \(\nu\) a positive measure on
\(Z\), and \(\zeta\mapsto\mathfrak H(\zeta)\) a \(\nu\)-measurable
field of complex Hilbert spaces on \(Z\). Put
\[
  \mathfrak H=\int^{\oplus}\mathfrak H(\zeta)\,d\nu(\zeta),
\]
and let \(\mathfrak Z\) be the algebra of diagonalizable operators.

\begin{theorem}\label{thm:II-3-1}
For every \(\iota\in I\), let
\(\zeta\mapsto T_{\iota}(\zeta)\) be an essentially bounded measurable
field of operators, and put
\[
  T_{\iota}=\int^{\oplus}T_{\iota}(\zeta)\,d\nu(\zeta).
\]
Let \(\mathcal A(\zeta)\) be the von Neumann algebra on
\(\mathfrak H(\zeta)\) generated by the \(T_{\iota}(\zeta)\), and let
\(\mathcal A\) be the von Neumann algebra on \(\mathfrak H\) generated
by the \(T_{\iota}\) and by \(\mathfrak Z\). Let
\[
  S=\int^{\oplus}S(\zeta)\,d\nu(\zeta)
\]
be a decomposable operator.
\begin{enumerate}
\item If \(S\in\mathcal A\), then
  \(S(\zeta)\in\mathcal A(\zeta)\) almost everywhere.
\item If \(S(\zeta)\in\mathcal A(\zeta)\) almost everywhere and \(I\)
  is countable, then \(S\in\mathcal A\).
\end{enumerate}
\end{theorem}

\begin{proof}
Assertion (i) is evident when \(S\) is one of the \(T_{\iota}\), or
when \(S\in\mathfrak Z\), and therefore when \(S\) belongs to the
involutive algebra \(\mathcal B\) generated by \(\mathfrak Z\) and the
\(T_{\iota}\). Moreover, \(\mathfrak Z'\) is of countable type
(\S~2, Proposition~7(iii)); hence \(\mathcal A\) is of countable type.
Thus, if \(S\in\mathcal A\), \(S\) is the strong limit of a sequence of
elements of \(\mathcal B\) (Chapter~I, \S~3, Proposition~1 and
Theorem~3). Using \S~2, Proposition~4, we see that
\(S(\zeta)\in\mathcal A(\zeta)\) almost everywhere.

Suppose that \(S(\zeta)\in\mathcal A(\zeta)\) almost everywhere and
that \(I\) is countable. Let \(T'\in\mathcal A'\). Since
\(\mathfrak Z\subset\mathcal A\), we have
\(T'\in\mathfrak Z'\), and hence
\[
  T'=\int^{\oplus}T'(\zeta)\,d\nu(\zeta)
  \qquad\text{(\S~2, Theorem~1).}
\]
% Source PDF page 179 (printed p. 172)
\phantomsection\label{srcpage:172}
Since \(T'\) commutes with the \(T_{\iota}\) and the
\(T_{\iota}^*\), \(T'(\zeta)\) commutes with
\(T_{\iota}(\zeta)\) and \(T_{\iota}(\zeta)^*\), for every \(\iota\),
almost everywhere. Thus \(T'(\zeta)\) commutes with
\(S(\zeta)\) and \(S(\zeta)^*\) almost everywhere, so that \(T'\)
commutes with \(S\) and \(S^*\). Hence
\(S\in\mathcal A''=\mathcal A\).
\end{proof}

\begin{corollary}\label{cor:II-3-1}
For every \(\iota\in I\), let
\(\zeta\mapsto T_{\iota}(\zeta)\) be an essentially bounded measurable
field of operators, and put
\[
  T_{\iota}=\int^{\oplus}T_{\iota}(\zeta)\,d\nu(\zeta).
\]
Let \(\mathcal A(\zeta)\) be the von Neumann algebra on
\(\mathfrak H(\zeta)\) generated by the \(T_{\iota}(\zeta)\), and let
\(\mathcal B\) be the von Neumann algebra on \(\mathfrak H\) generated
by the \(T_{\iota}\).
\begin{enumerate}
\item If \(\mathfrak Z\) is a maximal abelian von Neumann subalgebra of
  \(\mathcal B'\), then
  \(\mathcal A(\zeta)=\mathcal L(\mathfrak H(\zeta))\) almost
  everywhere.
\item If \(\mathcal A(\zeta)=\mathcal L(\mathfrak H(\zeta))\) almost
  everywhere and \(I\) is countable, then \(\mathfrak Z\) is a maximal
  abelian von Neumann subalgebra of \(\mathcal B'\).
\end{enumerate}
\end{corollary}

\begin{proof}
If \(\mathfrak Z\) is a maximal abelian von Neumann subalgebra of
\(\mathcal B'\), the von Neumann algebra \(\mathcal A\) generated by
\(\mathfrak Z\) and the \(T_{\iota}\), or by \(\mathfrak Z\) and
\(\mathcal B\), is \(\mathfrak Z'\) (Chapter~I, \S~1,
Proposition~13). Let
\[
  S_1=\int^{\oplus}S_1(\zeta)\,d\nu(\zeta),\quad
  S_2=\int^{\oplus}S_2(\zeta)\,d\nu(\zeta),\quad\ldots
\]
be a sequence of operators in \(\mathfrak Z'\) such that, almost
everywhere, \(\mathcal L(\mathfrak H(\zeta))\) is the von Neumann algebra
generated by the \(S_i(\zeta)\) (\S~2, Proposition~5). By Theorem~1(i),
\(S_i(\zeta)\in\mathcal A(\zeta)\) almost everywhere; hence
\(\mathcal A(\zeta)=\mathcal L(\mathfrak H(\zeta))\) almost everywhere.

Conversely, suppose that \(I\) is countable and that
\(\mathcal A(\zeta)=\mathcal L(\mathfrak H(\zeta))\) almost everywhere.
Then, by Theorem~1(ii), every decomposable operator belongs to
\(\mathcal A\). Thus \(\mathfrak Z'=\mathcal A\), and \(\mathfrak Z\)
is a maximal abelian von Neumann subalgebra of \(\mathcal B'\)
(Chapter~I, \S~1, Proposition~13).
\end{proof}

\begin{corollary}\label{cor:II-3-2}
The von Neumann algebra \(\mathfrak Z'\) is generated by
\(\mathfrak Z\) and a countable family of elements.
\end{corollary}

\begin{proof}
This follows from Theorem~1(ii) and \S~2, Proposition~5.
\end{proof}

Without the countability hypothesis on \(I\), Theorem~1(ii) and
Corollary~1(ii) are false (Exercise~2).

Bibliography: \ArticleRef{28}, \ArticleRef{55}, \ArticleRef{57},
\ArticleRef{100}, \ArticleRef{118}.

\NumberedParagraph{2}{Measurable fields of von Neumann algebras}
\label{no:II-3-2}
For every \(\zeta\in Z\), let \(\mathcal A(\zeta)\) be a von Neumann
algebra on \(\mathfrak H(\zeta)\). The mapping
\(\zeta\mapsto\mathcal A(\zeta)\) is called a \emph{field of von
Neumann algebras} on \(Z\).

% Source PDF page 180 (printed p. 173)
\phantomsection\label{srcpage:173}
\begin{definition}\label{def:II-3-1}
A field of von Neumann algebras
\(\zeta\mapsto\mathcal A(\zeta)\) on \(Z\) is said to be
\emph{measurable} if there exists a sequence
\(\zeta\mapsto T_1(\zeta),\zeta\mapsto T_2(\zeta),\ldots\) of
measurable fields of operators such that, almost everywhere,
\(\mathcal A(\zeta)\) is the von Neumann algebra generated by the
\(T_i(\zeta)\).
\end{definition}

Replacing \(T_i(\zeta)\) by
\(\lVert T_i(\zeta)\rVert^{-1}T_i(\zeta)\) at the points \(\zeta\) for
which \(\lVert T_i(\zeta)\rVert\neq0\), we may suppose the fields
\(\zeta\mapsto T_i(\zeta)\) essentially bounded; they then define
operators in \(\mathfrak Z'\).

\emph{Examples.} The field
\(\zeta\mapsto\mathbb C I_{\mathfrak H(\zeta)}\) is measurable. The
field \(\zeta\mapsto\mathcal L(\mathfrak H(\zeta))\) is measurable
(\S~2, Proposition~5).

Let \(Y\) be a measurable subset of \(Z\), and let
\(\zeta\mapsto\mathcal A(\zeta)\) be a field of von Neumann algebras on
\(Y\). This field is said to be measurable if it extends to a measurable
field of von Neumann algebras on \(Z\). Equivalently, there must exist a
sequence
\[
  \zeta\longmapsto T_1(\zeta),\quad
  \zeta\longmapsto T_2(\zeta),\quad\ldots
\]
of measurable fields of operators on \(Y\) such that, almost everywhere
on \(Y\), \(\mathcal A(\zeta)\) is the von Neumann algebra generated by
the \(T_i(\zeta)\). If \(Z\) is the union of a sequence
\((Z_1,Z_2,\ldots)\) of measurable subsets, a field of von Neumann
algebras on \(Z\) is measurable if and only if its restriction to each
\(Z_i\) is measurable.

\begin{proposition}\label{prop:II-3-1}
Let \(\zeta\mapsto\mathcal A(\zeta)\) be a measurable field of von
Neumann algebras.
\begin{enumerate}
\item The set \(\mathcal A\) of decomposable operators
  \[
    T=\int^{\oplus}T(\zeta)\,d\nu(\zeta)
  \]
  such that \(T(\zeta)\in\mathcal A(\zeta)\) almost everywhere is a von
  Neumann algebra on \(\mathfrak H\) such that
  \(\mathfrak Z\subset\mathcal A\subset\mathfrak Z'\), generated by
  \(\mathfrak Z\) and a countable family of elements.
\item If \(\zeta\mapsto\mathcal B(\zeta)\) is a measurable field of von
  Neumann algebras defining the same von Neumann algebra \(\mathcal A\)
  on \(\mathfrak H\), then
  \(\mathcal A(\zeta)=\mathcal B(\zeta)\) almost everywhere.
\end{enumerate}
\end{proposition}

\begin{proof}
Assertion (i) is an immediate consequence of Theorem~1. We prove (ii).
Let
\[
  T_1=\int^{\oplus}T_1(\zeta)\,d\nu(\zeta),\quad
  T_2=\int^{\oplus}T_2(\zeta)\,d\nu(\zeta),\quad\ldots
\]
be a sequence of decomposable operators such that, almost everywhere,
\(\mathcal A(\zeta)\) is the von Neumann algebra generated by the
\(T_i(\zeta)\). We have \(T_i\in\mathcal A\), and hence
\(T_i(\zeta)\in\mathcal B(\zeta)\) almost everywhere. Thus
\(\mathcal A(\zeta)\subset\mathcal B(\zeta)\) almost everywhere.
Similarly, \(\mathcal B(\zeta)\subset\mathcal A(\zeta)\) almost
everywhere.
\end{proof}

% Source PDF page 181 (printed p. 174)
\phantomsection\label{srcpage:174}
\begin{definition}\label{def:II-3-2}
A von Neumann algebra \(\mathcal A\) on \(\mathfrak H\) is said to be
\emph{decomposable} if it is defined by a measurable field
\(\zeta\mapsto\mathcal A(\zeta)\) of von Neumann algebras. We then write
\[
  \mathcal A=\int^{\oplus}\mathcal A(\zeta)\,d\nu(\zeta).
\]
\end{definition}

By Proposition~1(ii), the \(\mathcal A(\zeta)\) are determined by
\(\mathcal A\) up to null sets.

\begin{theorem}\label{thm:II-3-2}
A von Neumann algebra \(\mathcal A\) on \(\mathfrak H\) is decomposable
if and only if it is the von Neumann algebra generated by
\(\mathfrak Z\) and a countable family of decomposable operators.
\end{theorem}

\begin{proof}
The condition is necessary by Proposition~1(i). Conversely, suppose
that \(\mathcal A\) is generated by \(\mathfrak Z\) and a sequence
\[
  T_1=\int^{\oplus}T_1(\zeta)\,d\nu(\zeta),\quad
  T_2=\int^{\oplus}T_2(\zeta)\,d\nu(\zeta),\quad\ldots
\]
of decomposable operators, and let
\(\mathcal A(\zeta)\subset\mathcal L(\mathfrak H(\zeta))\) be the von
Neumann algebra generated by the \(T_i(\zeta)\). The field
\(\zeta\mapsto\mathcal A(\zeta)\) is measurable, and, by Theorem~1,
\[
  \mathcal A=\int^{\oplus}\mathcal A(\zeta)\,d\nu(\zeta).
\]
\end{proof}

\textbf{Remark.} If \(\nu\) is standard, \(\mathfrak H\) has a countable
base; to say that \(\mathcal A\) is generated by \(\mathfrak Z\) and a
countable family of decomposable operators then amounts to saying that
\[
  \mathfrak Z\subset\mathcal A\subset\mathfrak Z'.
\]

\emph{Problem.} Is the same true in general?

Bibliography: \ArticleRef{28}, \ArticleRef{48}, \ArticleRef{49},
\ArticleRef{80}, \ArticleRef{100}, \ArticleRef{338}.

\NumberedParagraph{3}{Relations between a decomposable von Neumann
algebra and its components}\label{no:II-3-3}
\begin{theorem}\label{thm:II-3-3}
Suppose \(\mathcal A\) and \(\mathcal A'\) are decomposable. Write
\[
  \mathcal A=\int^{\oplus}\mathcal A(\zeta)\,d\nu(\zeta),
  \qquad
  \mathcal A'=\int^{\oplus}\mathcal A'(\zeta)\,d\nu(\zeta).
\]
\begin{enumerate}
\item The algebras \(\mathcal A(\zeta)\) and
  \(\mathcal A'(\zeta)\) commute almost everywhere.
\item If \(\mathfrak Z\) is the centre of \(\mathcal A\), and hence of
  \(\mathcal A'\), then \(\mathcal A(\zeta)\) and
  \(\mathcal A'(\zeta)\) are almost everywhere factors which generate
  the von Neumann algebra \(\mathcal L(\mathfrak H(\zeta))\).
\item Conversely, if \(\mathcal A(\zeta)\) and
  \(\mathcal A'(\zeta)\) generate
  \(\mathcal L(\mathfrak H(\zeta))\) almost everywhere, then
  \(\mathfrak Z\) is the centre of \(\mathcal A\) and
  \(\mathcal A'\).
\end{enumerate}
\end{theorem}

% Source PDF page 182 (printed p. 175)
\phantomsection\label{srcpage:175}
\begin{proof}
Let
\[
  T_1=\int^{\oplus}T_1(\zeta)\,d\nu(\zeta),\quad
  T_2=\int^{\oplus}T_2(\zeta)\,d\nu(\zeta),\quad\ldots
\]
(respectively,
\[
  T_1'=\int^{\oplus}T_1'(\zeta)\,d\nu(\zeta),\quad
  T_2'=\int^{\oplus}T_2'(\zeta)\,d\nu(\zeta),\quad\ldots)
\]
be decomposable operators which, together with \(\mathfrak Z\),
generate \(\mathcal A\) (respectively, \(\mathcal A'\)). Almost
everywhere, \(\mathcal A(\zeta)\) is generated by the
\(T_i(\zeta)\), and \(\mathcal A'(\zeta)\) by the
\(T_j'(\zeta)\). Since \(T_i\) commutes with \(T_j'\) and
\(T_j'^*\), \(T_i(\zeta)\) commutes with \(T_j'(\zeta)\) and
\(T_j'(\zeta)^*\) almost everywhere. Thus
\(\mathcal A(\zeta)\) and \(\mathcal A'(\zeta)\) commute almost
everywhere, proving (i).

We prove (ii) and (iii) together. To say that
\(\mathcal A\cap\mathcal A'=\mathfrak Z\) amounts to saying that
\(\mathcal A\) and \(\mathcal A'\) generate the von Neumann algebra
\(\mathfrak Z'\), hence that \(\mathfrak Z\), the \(T_i\), and the
\(T_j'\) generate \(\mathfrak Z'\), hence, by Theorem~1, that almost
everywhere the \(T_i(\zeta)\) and \(T_j'(\zeta)\) generate
\(\mathcal L(\mathfrak H(\zeta))\), and hence that almost everywhere
\(\mathcal A(\zeta)\) and \(\mathcal A'(\zeta)\) generate
\(\mathcal L(\mathfrak H(\zeta))\). Finally, if
\(\mathcal A(\zeta)\) and \(\mathcal A'(\zeta)\) generate
\(\mathcal L(\mathfrak H(\zeta))\), then
\[
  \mathcal A(\zeta)\cap\mathcal A'(\zeta)
  \subset
  \bigl(\mathcal A'(\zeta)\cup\mathcal A(\zeta)\bigr)'
  =\mathcal L(\mathfrak H(\zeta))',
\]
so \(\mathcal A(\zeta)\), and likewise
\(\mathcal A'(\zeta)\), is a factor.
\end{proof}

\begin{proposition}\label{prop:II-3-2}
Let
\[
  \mathcal A_1=\int^{\oplus}\mathcal A_1(\zeta)\,d\nu(\zeta),\quad
  \mathcal A_2=\int^{\oplus}\mathcal A_2(\zeta)\,d\nu(\zeta),\quad\ldots
\]
be a sequence of decomposable von Neumann algebras, and let
\(\mathcal A\) be the von Neumann algebra which they generate. Then
\(\mathcal A\) is decomposable; and if
\[
  \mathcal A=\int^{\oplus}\mathcal A(\zeta)\,d\nu(\zeta),
\]
then \(\mathcal A(\zeta)\) is almost everywhere the von Neumann algebra
generated by the \(\mathcal A_i(\zeta)\).
\end{proposition}

\begin{proof}
Let
\[
  T_1^i=\int^{\oplus}T_1^i(\zeta)\,d\nu(\zeta),\quad
  T_2^i=\int^{\oplus}T_2^i(\zeta)\,d\nu(\zeta),\quad\ldots
\]
be operators which, together with \(\mathfrak Z\), generate
\(\mathcal A_i\). Almost everywhere, \(\mathcal A_i(\zeta)\) is
generated by \(T_1^i(\zeta),T_2^i(\zeta),\ldots\). Thus
\(\mathcal A\) is generated by \(\mathfrak Z\) and the
\(T_j^i\) (\(i=1,2,\ldots; j=1,2,\ldots\)); hence \(\mathcal A\) is
decomposable and, almost everywhere, \(\mathcal A(\zeta)\) is generated
by the \(T_j^i(\zeta)\), that is, by the
\(\mathcal A_i(\zeta)\).
\end{proof}

\begin{lemma}\label{lem:II-3-1}
Suppose \(\nu\) is standard. If the field of von Neumann algebras
\(\zeta\mapsto\mathcal A(\zeta)\) is measurable, then the field
\(\zeta\mapsto\mathcal A(\zeta)'\) is also measurable.
\end{lemma}

% Source PDF page 183 (printed p. 176)
\phantomsection\label{srcpage:176}
\begin{proof}
It is enough to prove the lemma when
\(\zeta\mapsto\mathfrak H(\zeta)\) is the constant field on \(Z\)
corresponding to a Hilbert space \(\mathfrak H_0\) with a countable base,
and when \(Z\) is compact and metrizable. Let \(\mathcal L_1\) be the
unit ball of \(\mathcal L(\mathfrak H_0)\), furnished with the strong
topology. For \(i=1,2,\ldots\), let \(T_i(\zeta)\) be an operator in
\(\mathcal L_1\), depending measurably on \(\zeta\), such that, for
every \(\zeta\in Z\), \(\mathcal A(\zeta)\) is the von Neumann algebra
generated by the \(T_i(\zeta)\). Up to a null set, \(Z\) is the union of
a sequence of compact subsets \(Y_j\) such that the restrictions to each
\(Y_j\) of the mappings
\(\zeta\mapsto T_i(\zeta)\) and
\(\zeta\mapsto T_i(\zeta)^*\) are continuous. Since it is enough to
prove that the restriction of the field
\(\zeta\mapsto\mathcal A(\zeta)'\) to each \(Y_j\) is measurable, we
shall henceforth restrict ourselves to the case in which the mappings
\(\zeta\mapsto T_i(\zeta)\) and
\(\zeta\mapsto T_i(\zeta)^*\) are continuous.

Let \(\mathcal M\) be the subset of \(Z\times\mathcal L_1\) consisting
of the pairs \((\zeta,T')\) such that
\(T'\in\mathcal A(\zeta)'\), that is,
\[
  T'T_i(\zeta)=T_i(\zeta)T',
  \qquad
  T'T_i(\zeta)^*=T_i(\zeta)^*T'.
\]
Since the mappings
\[
  (\zeta,T')\longmapsto T'T_i(\zeta)-T_i(\zeta)T',
  \qquad
  (\zeta,T')\longmapsto T'T_i(\zeta)^*-T_i(\zeta)^*T'
\]
from \(Z\times\mathcal L_1\) into \(\mathcal L(\mathfrak H_0)\) are
continuous, \(\mathcal M\) is closed. Let
\((\mathcal M_1,\mathcal M_2,\ldots)\) be a base for the topology of
\(\mathcal M\). The \(\mathcal M_n\) are Borel subsets of
\(Z\times\mathcal L_1\), which is a complete metric space with a
countable base. Let \(Z_n\) be the \(\nu\)-measurable projection of
\(\mathcal M_n\) on \(Z\). By Appendix~V, there exists a
\(\nu\)-measurable mapping \(\zeta\mapsto T_n'(\zeta)\) from \(Z_n\)
into \(\mathcal L_1\) such that
\((\zeta,T_n'(\zeta))\in\mathcal M_n\) for every \(\zeta\in Z_n\).
Extend this field to \(Z\) by assigning it the value zero on
\(Z-Z_n\), and continue to denote the resulting \(\nu\)-measurable
field by \(\zeta\mapsto T_n'(\zeta)\). It remains to show that, for
every \(\zeta\in Z\), the \(T_n'(\zeta)\) generate
\(\mathcal A(\zeta)'\). Let \(T'\in\mathcal A(\zeta)'\) with
\(\lVert T'\rVert\leq1\). Then \((\zeta,T')\in\mathcal M\). If
\((\mathcal M_{n_1},\mathcal M_{n_2},\ldots)\) is a fundamental system
of neighborhoods of \((\zeta,T')\) in \(\mathcal M\), then
\((\zeta,T')\) is the limit of
\((\zeta,T_{n_p}'(\zeta))\) in \(Z\times\mathcal L_1\); that is,
\(T'\) is the strong limit of the \(T_{n_p}'(\zeta)\).
\end{proof}

\begin{theorem}\label{thm:II-3-4}
Suppose \(\nu\) is standard.
\begin{enumerate}
\item If
  \(\mathcal A=\int^{\oplus}\mathcal A(\zeta)\,d\nu(\zeta)\) is a
  decomposable von Neumann algebra, then
  \[
    \mathcal A'=\int^{\oplus}\mathcal A(\zeta)'\,d\nu(\zeta).
  \]
\item Let
  \[
    \mathcal B_1=\int^{\oplus}\mathcal B_1(\zeta)\,d\nu(\zeta),\quad
    \mathcal B_2=\int^{\oplus}\mathcal B_2(\zeta)\,d\nu(\zeta),\quad
    \ldots
  \]
  be decomposable von Neumann algebras, and let \(\mathcal B\) be their
  intersection. Then \(\mathcal B\) is decomposable; and if
  \(\mathcal B=\int^{\oplus}\mathcal B(\zeta)\,d\nu(\zeta)\), then
  \(\mathcal B(\zeta)\) is almost everywhere the intersection of the
  \(\mathcal B_i(\zeta)\).
\end{enumerate}
\end{theorem}

% Source PDF page 184 (printed p. 177)
\phantomsection\label{srcpage:177}
\begin{proof}
We prove (i). Let
\[
  T_1=\int^{\oplus}T_1(\zeta)\,d\nu(\zeta),\quad
  T_2=\int^{\oplus}T_2(\zeta)\,d\nu(\zeta),\quad\ldots
\]
be decomposable operators which, together with \(\mathfrak Z\),
generate \(\mathcal A\). Almost everywhere,
\(\mathcal A(\zeta)\) is generated by the \(T_i(\zeta)\). A
decomposable operator
\[
  T'=\int^{\oplus}T'(\zeta)\,d\nu(\zeta)
\]
belongs to \(\mathcal A'\) if and only if it commutes with the \(T_i\)
and the \(T_i^*\), hence if and only if, almost everywhere,
\(T'(\zeta)\) commutes with \(T_i(\zeta)\) and
\(T_i(\zeta)^*\), hence if and only if
\(T'(\zeta)\in\mathcal A(\zeta)'\) almost everywhere. By Lemma~1, the
field \(\zeta\mapsto\mathcal A(\zeta)'\) is measurable. Therefore
\[
  \mathcal A'=\int^{\oplus}\mathcal A(\zeta)'\,d\nu(\zeta).
\]

We prove (ii). A decomposable operator
\[
  T=\int^{\oplus}T(\zeta)\,d\nu(\zeta)
\]
belongs to every \(\mathcal B_i\) if and only if, almost everywhere,
\(T(\zeta)\) belongs to every \(\mathcal B_i(\zeta)\), that is, to their
intersection \(\mathcal B(\zeta)\). To prove (ii), it is therefore enough
to prove that the field \(\zeta\mapsto\mathcal B(\zeta)\) is measurable.
By Lemma~1, it is enough to prove that the field
\(\zeta\mapsto\mathcal B(\zeta)'\) is measurable. But
\(\mathcal B(\zeta)'\) is the von Neumann algebra generated by
\(\mathcal B_1(\zeta)',\mathcal B_2(\zeta)',\ldots\) (Chapter~I,
\S~1, Proposition~1). It is therefore enough to apply Lemma~1 and
Proposition~2 once again.
\end{proof}

Bibliography: \ArticleRef{28}, \ArticleRef{80}, \ArticleRef{100},
\ArticleRef{117}, \ArticleRef{118}. By \ArticleRef{338}, the hypothesis
that \(\nu\) is standard may be omitted from Lemma~1 and Theorem~4.

\Needspace{4\baselineskip}
\NumberedParagraph{4}{Constant fields of von Neumann algebras}
\label{no:II-3-4}
\begin{proposition}\label{prop:II-3-3}
Let \(\zeta\mapsto\mathfrak H(\zeta)\) be a \(\nu\)-measurable field of
complex Hilbert spaces on \(Z\), put
\[
  \mathfrak H=\int^{\oplus}\mathfrak H(\zeta)\,d\nu(\zeta),
\]
let \(\mathfrak K_0\) be a complex Hilbert space with a countable base,
and let \(\zeta\mapsto\mathfrak K(\zeta)\) be the corresponding constant
field on \(Z\). Let
\[
  \mathcal A=\int^{\oplus}\mathcal A(\zeta)\,d\nu(\zeta)
\]
be a decomposable von Neumann algebra on \(\mathfrak H\), and let
\(\mathcal B\) be a von Neumann algebra on \(\mathfrak K_0\). Then,
identifying \(\mathfrak H\otimes\mathfrak K_0\) with
\[
  \int^{\oplus}\bigl(\mathfrak H(\zeta)\otimes
  \mathfrak K(\zeta)\bigr)\,d\nu(\zeta),
\]
we have
\[
  \mathcal A\otimes\mathcal B
  =\int^{\oplus}\bigl(\mathcal A(\zeta)\otimes\mathcal B\bigr)
  \,d\nu(\zeta).
\]
\end{proposition}

\begin{proof}
Let
\[
  S_1=\int^{\oplus}S_1(\zeta)\,d\nu(\zeta),\quad
  S_2=\int^{\oplus}S_2(\zeta)\,d\nu(\zeta),\quad\ldots
\]
% Source PDF page 185 (printed p. 178)
\phantomsection\label{srcpage:178}
be decomposable operators on \(\mathfrak H\) which, together with
\(\mathfrak Z\), generate \(\mathcal A\). Let
the operators \(T_1,T_2,\ldots\) on \(\mathfrak K_0\) generate
\(\mathcal B\). Then the \(S_i(\zeta)\otimes T_j\) generate
\(\mathcal A(\zeta)\otimes\mathcal B\) almost everywhere, and hence the
field \(\zeta\mapsto\mathcal A(\zeta)\otimes\mathcal B\) is measurable.
Moreover, since
\[
  S_i\otimes T_j
  =\int^{\oplus}\bigl(S_i(\zeta)\otimes T_j\bigr)\,d\nu(\zeta)
  \qquad\text{(\S~2, Proposition~8),}
\]
the algebra
\(\int^{\oplus}(\mathcal A(\zeta)\otimes\mathcal B)\,d\nu(\zeta)\)
is the von Neumann algebra generated by the \(S_i\otimes T_j\) and by
the diagonalizable operators on
\(\mathfrak H\otimes\mathfrak K_0\), which form the algebra
\(\mathfrak Z\otimes\mathbb C I_{\mathfrak K_0}\). Therefore this
integral algebra, which is generated by
\(\mathcal A\otimes\mathbb C I_{\mathfrak K_0}\) and
\(\mathbb C I_{\mathfrak H}\otimes\mathcal B\), is equal to
\(\mathcal A\otimes\mathcal B\).
\end{proof}

\begin{corollary*}
Let \(\mathfrak K_0\) be a complex Hilbert space with a countable base,
let \(\zeta\mapsto\mathfrak K(\zeta)\) be the corresponding constant
field on \(Z\), and put
\[
  \mathfrak K=\int^{\oplus}\mathfrak K(\zeta)\,d\nu(\zeta)
  =L_{\mathbb C}^{2}(Z,\nu)\otimes\mathfrak K_0.
\]
For every von Neumann algebra \(\mathcal B\) on \(\mathfrak K_0\),
\[
  \mathfrak Z\otimes\mathcal B
  =\int^{\oplus}\mathcal B(\zeta)\,d\nu(\zeta),
\]
where \(\mathcal B(\zeta)=\mathcal B\) for every \(\zeta\in Z\).
\end{corollary*}

\begin{proof}
Apply Proposition~3 when \(\zeta\mapsto\mathfrak H(\zeta)\) is the
constant field \(\zeta\mapsto\mathbb C\).
\end{proof}

\begin{lemma}\label{lem:II-3-2}
Let \(Y\) be a subset of \(Z\), let \(X\) be its measurable envelope,
let \(\mathfrak K_0\) be a complex Hilbert space with a countable base,
let \(\mathcal A_0\) be a von Neumann algebra on \(\mathfrak K_0\), and
let \(\mathcal A(\zeta)\) be a von Neumann algebra on
\(\mathfrak K_0\), generated by a sequence of operators depending
measurably on \(\zeta\). Suppose that, for every \(\zeta\in Y\), there
exists a unitary operator \(U(\zeta)\) on \(\mathfrak K_0\) such that
\[
  U(\zeta)^{-1}\mathcal A_0U(\zeta)=\mathcal A(\zeta).
\]
Suppose also that \(\nu\) is standard. Then, for almost every
\(\zeta\in X\), there exists a unitary operator \(V(\zeta)\) on
\(\mathfrak K_0\), depending measurably on \(\zeta\), such that
\[
  V(\zeta)^{-1}\mathcal A_0V(\zeta)=\mathcal A(\zeta).
\]
\end{lemma}

\begin{proof}
Let \((x_i)\) be a total sequence in \(\mathfrak K_0\), let
\((S_i)\) (respectively, \((S_i')\)) be a sequence of operators on
\(\mathfrak K_0\) generating \(\mathcal A_0\) (respectively,
\(\mathcal A_0'\)), and let
\(\zeta\mapsto T_i(\zeta)\) (respectively,
\(\zeta\mapsto T_i'(\zeta)\)) be a sequence of measurable mappings of
\(Z\) into \(\mathcal L(\mathfrak K_0)\) such that, for every
\(\zeta\in Z\), the \(T_i(\zeta)\) (respectively,
\(T_i'(\zeta)\)) generate the von Neumann algebra
% Source PDF page 186 (printed p. 179)
\phantomsection\label{srcpage:179}
\(\mathcal A(\zeta)\) (respectively, \(\mathcal A(\zeta)'\)). We may
suppose that, for every \(i\), the mapping
\(\zeta\mapsto T_i(\zeta)^*\) is one of the mappings
\(\zeta\mapsto T_k(\zeta)\), and similarly for the \(T_i'(\zeta)\).
The mappings \(\zeta\mapsto T_i(\zeta)x_j\) and
\(\zeta\mapsto T_i'(\zeta)x_j\) are measurable. We may suppose that
\(Z\) is compact and metrizable. Then, up to a null set, \(Z\) is the
union of a sequence of compact subsets such that the restrictions to
each of them of the mappings
\(\zeta\mapsto T_i(\zeta)x_j\) and
\(\zeta\mapsto T_i'(\zeta)x_j\) are continuous for the strong topology
of \(\mathfrak K_0\). We may therefore suppose henceforth that these
mappings are continuous on \(Z\) for the strong topology of
\(\mathfrak K_0\).

Let \(\mathcal L_1\) be the unit ball of
\(\mathcal L(\mathfrak K_0)\), furnished with the strong topology, and
let \(\mathcal M\subset Z\times\mathcal L_1\) be the set of pairs
\((\zeta,T)\) such that \(T\) is unitary and
\(T^{-1}\mathcal A_0T=\mathcal A(\zeta)\). To say that
\((\zeta,T)\in\mathcal M\) amounts to saying that \(T\) is unitary,
that \(T^*S_iT\) commutes with \(T_j'(\zeta)\), and that
\(T^*S_i'T\) commutes with \(T_j(\zeta)\), for all \(i,j\); or,
equivalently, that \(T\) is unitary and
\begin{align*}
  (S_iTT_j'(\zeta)x_k\mid Tx_l)
  &=(S_iTx_k\mid TT_j'(\zeta)^*x_l),\\
  (S_i'TT_j(\zeta)x_k\mid Tx_l)
  &=(S_i'Tx_k\mid TT_j(\zeta)^*x_l)
\end{align*}
for all \(i,j,k,l\). Thus \(\mathcal M\) is Borel in
\(Z\times\mathcal L_1\) (compare \S~2, Lemma~1). Moreover, for every
\(\zeta\in Y\), \((\zeta,U(\zeta))\in\mathcal M\), so the projection of
\(\mathcal M\) on \(Z\) contains \(Y\). By Appendix~V, there is a
measurable mapping \(\zeta\mapsto V(\zeta)\) of \(X\) into
\(\mathcal L(\mathfrak K_0)\) such that
\((\zeta,V(\zeta))\in\mathcal M\) for almost every \(\zeta\in X\).
\end{proof}

\begin{proposition}\label{prop:II-3-4}
Let \(\mathfrak K_0\) be a complex Hilbert space with a countable base,
let \(\mathcal A_0\) be a von Neumann algebra on \(\mathfrak K_0\), and
let \(\zeta\mapsto\mathfrak H(\zeta)\) be a \(\nu\)-measurable field of
complex Hilbert spaces on \(Z\). Put
\[
  \mathfrak H=\int^{\oplus}\mathfrak H(\zeta)\,d\nu(\zeta),
\]
let \(\mathfrak Z\) be the algebra of diagonalizable operators, and let
\(\zeta\mapsto\mathcal A(\zeta)\subset
\mathcal L(\mathfrak H(\zeta))\) be a measurable field of von Neumann
algebras on \(Z\), with
\[
  \mathcal A=\int^{\oplus}\mathcal A(\zeta)\,d\nu(\zeta).
\]
Suppose that, for every \(\zeta\in Z\), there exists an isomorphism
\(U(\zeta)\) of \(\mathfrak H(\zeta)\) onto \(\mathfrak K_0\) such that
\[
  U(\zeta)^{-1}\mathcal A_0U(\zeta)=\mathcal A(\zeta).
\]
Suppose also that \(\nu\) is standard. Then there exists an isomorphism
of \(\mathfrak H\) onto
\(L_{\mathbb C}^{2}(Z,\nu)\otimes\mathfrak K_0\) which transforms
\(\mathcal A\) into \(\mathfrak Z\otimes\mathcal A_0\).
\end{proposition}

\begin{proof}
We are immediately reduced to the case in which
\(\zeta\mapsto\mathfrak H(\zeta)\) is the constant field corresponding
to \(\mathfrak K_0\). By Lemma~2, for
% Source PDF page 187 (printed p. 180)
\phantomsection\label{srcpage:180}
every \(\zeta\in Z\) there exists a unitary operator \(V(\zeta)\) on
\(\mathfrak K_0\), depending measurably on \(\zeta\), such that
\(V(\zeta)^{-1}\mathcal A_0V(\zeta)=\mathcal A(\zeta)\). Then
\[
  \int^{\oplus}V(\zeta)\,d\nu(\zeta)
\]
is a unitary operator on
\(\mathfrak H=L_{\mathbb C}^{2}(Z,\nu)\otimes\mathfrak K_0\) which
transforms \(\mathcal A\) into the von Neumann algebra defined by the
constant field \(\zeta\mapsto\mathcal A_0\), that is,
\(\mathfrak Z\otimes\mathcal A_0\) (corollary to Proposition~3).
\end{proof}

\begin{proposition}\label{prop:II-3-5}
Let \(\mathfrak K_0\) be a complex Hilbert space, let \(\mathcal A_0\)
be a von Neumann algebra on \(\mathfrak K_0\), let
\(\zeta\mapsto\mathfrak H(\zeta)\) be a \(\nu\)-measurable field of
complex Hilbert spaces on \(Z\), and let
\(\zeta\mapsto\mathcal A(\zeta)\subset
\mathcal L(\mathfrak H(\zeta))\) be a measurable field of von Neumann
algebras on \(Z\). If \(\nu\) is standard, the set of
\(\zeta\in Z\) for which \(\mathcal A(\zeta)\) is spatially isomorphic
to \(\mathcal A_0\) is measurable.
\end{proposition}

\begin{proof}
This follows immediately from Lemma~2.
\end{proof}

Bibliography: \ArticleRef{80}. By \ArticleRef{338} and
\ArticleRef{466}, the hypothesis that \(\nu\) is standard may be omitted
from Propositions~4 and~5.

\NumberedParagraph{5}{Reduction of discrete or continuous von Neumann
algebras}\label{no:II-3-5}
\begin{proposition}\label{prop:II-3-6}
Let
\[
  \mathcal A=\int^{\oplus}\mathcal A(\zeta)\,d\nu(\zeta)
\]
be a decomposable von Neumann algebra, and let
\[
  E=\int^{\oplus}E(\zeta)\,d\nu(\zeta)
\]
be a projection in \(\mathcal A\) (respectively, \(\mathcal A'\)), where
for every \(\zeta\in Z\), \(E(\zeta)\) is a projection in
\(\mathcal A(\zeta)\) (respectively, \(\mathcal A(\zeta)'\)). Let
\(\zeta\mapsto\mathfrak K(\zeta)\) be the measurable field of subspaces
corresponding to the \(E(\zeta)\), and put
\[
  \mathfrak K=E\mathfrak H
  =\int^{\oplus}\mathfrak K(\zeta)\,d\nu(\zeta).
\]
Then
\[
  \mathcal A_{\mathfrak K}
  =\int^{\oplus}\mathcal A(\zeta)_{\mathfrak K(\zeta)}
  \,d\nu(\zeta).
\]
\end{proposition}

\begin{proof}
Let
\[
  T_1=\int^{\oplus}T_1(\zeta)\,d\nu(\zeta),\quad
  T_2=\int^{\oplus}T_2(\zeta)\,d\nu(\zeta),\quad\ldots
\]
be operators in \(\mathcal A\) such that the \(T_i(\zeta)\) generate
\(\mathcal A(\zeta)\) almost everywhere. We may suppose the sequence
\((T_i)\) stable under taking adjoints and multiplication. Then the
\((ET_iE)_{\mathfrak K}\), together with
\(\mathfrak Z_{\mathfrak K}\), generate
\(\mathcal A_{\mathfrak K}\) (Chapter~I, \S~2, Proposition~1). But
\[
  ET_iE=\int^{\oplus}E(\zeta)T_i(\zeta)E(\zeta)\,d\nu(\zeta),
\]
and hence
\[
  (ET_iE)_{\mathfrak K}
  =\int^{\oplus}
  \bigl(E(\zeta)T_i(\zeta)E(\zeta)\bigr)_{\mathfrak K(\zeta)}
  \,d\nu(\zeta).
\]
% Source PDF page 188 (printed p. 181)
\phantomsection\label{srcpage:181}
Since the
\(\bigl(E(\zeta)T_i(\zeta)E(\zeta)\bigr)_{\mathfrak K(\zeta)}\)
generate \(\mathcal A(\zeta)_{\mathfrak K(\zeta)}\) almost everywhere,
we have
\[
  \mathcal A_{\mathfrak K}
  =\int^{\oplus}\mathcal A(\zeta)_{\mathfrak K(\zeta)}
  \,d\nu(\zeta).
\]
\end{proof}

\begin{corollary*}
The projection
\[
  E=\int^{\oplus}E(\zeta)\,d\nu(\zeta)
\]
is an abelian projection of
\[
  \mathcal A=\int^{\oplus}\mathcal A(\zeta)\,d\nu(\zeta)
\]
if and only if \(E(\zeta)\) is almost everywhere an abelian projection
of \(\mathcal A(\zeta)\).
\end{corollary*}

\begin{proof}
To say that \(E\) is abelian is to say that
\(\mathcal A_{\mathfrak K}\) is abelian, hence that the
\(\mathcal A(\zeta)_{\mathfrak K(\zeta)}\) are abelian almost
everywhere, and hence that \(E(\zeta)\) is abelian almost everywhere.
\end{proof}

\begin{lemma}\label{lem:II-3-3}
Let
\(\mathcal A=\int^{\oplus}\mathcal A(\zeta)\,d\nu(\zeta)\) be a
decomposable von Neumann algebra, and let
\(\zeta\mapsto\mathfrak K(\zeta)\) be a measurable field of subspaces.
Put
\[
  \mathfrak K=\int^{\oplus}\mathfrak K(\zeta)\,d\nu(\zeta),
  \qquad
  E(\zeta)=E_{\mathfrak K(\zeta)}^{\mathcal A(\zeta)'},
  \qquad
  E=E_{\mathfrak K}^{\mathcal A'}.
\]
Then
\[
  E=\int^{\oplus}E(\zeta)\,d\nu(\zeta).
\]
\end{lemma}

\begin{proof}
Let \(\mathfrak K_1=E\mathfrak H\), and let
\(\zeta\mapsto\mathfrak K_1(\zeta)\) be the measurable field of
subspaces such that
\(\mathfrak K_1=\int^{\oplus}\mathfrak K_1(\zeta)\,d\nu(\zeta)\). Let
\[
  T_1=\int^{\oplus}T_1(\zeta)\,d\nu(\zeta),\quad
  T_2=\int^{\oplus}T_2(\zeta)\,d\nu(\zeta),\quad\ldots
\]
be operators in \(\mathcal A\) such that (1) the \(T_i(\zeta)\)
generate \(\mathcal A(\zeta)\) almost everywhere, and (2) the \(T_i\)
are strongly everywhere dense in \(\mathcal A\). Let \((x_j)\) be a
fundamental sequence of square-integrable fields of vectors. Then the
vectors
\[
  TT_iP_{\mathfrak K}x_j
  \quad(T\in\mathfrak Z;\ i=1,2,\ldots;\ j=1,2,\ldots)
\]
form a total sequence in \(\mathfrak K_1\). Hence the vectors
\(T_i(\zeta)P_{\mathfrak K(\zeta)}x_j(\zeta)\) form almost everywhere a
total sequence in \(\mathfrak K_1(\zeta)\). Therefore, almost
everywhere, \(P_{\mathfrak K_1(\zeta)}=E(\zeta)\).
\end{proof}

\begin{proposition}\label{prop:II-3-7}
Let
\(\mathcal A=\int^{\oplus}\mathcal A(\zeta)\,d\nu(\zeta)\) be a
decomposable von Neumann algebra.
\begin{enumerate}
\item If \(\mathcal A\) is discrete, then \(\mathcal A(\zeta)\) is
  discrete almost everywhere.
\item If \(\mathcal A\) is continuous and \(\nu\) is standard, then
  \(\mathcal A(\zeta)\) is continuous almost everywhere.
\end{enumerate}
\end{proposition}

% Source PDF page 189 (printed p. 182)
\phantomsection\label{srcpage:182}
\begin{proof}
Suppose \(\mathcal A\) is discrete. Let
\[
  G=\int^{\oplus}G(\zeta)\,d\nu(\zeta)
\]
be an abelian projection of \(\mathcal A\) with central support \(I\)
(Chapter~I, \S~8, Theorem~1). Almost everywhere on \(Z\),
\(G(\zeta)\) is an abelian projection of \(\mathcal A(\zeta)\)
(corollary to Proposition~6) with central support \(I\) (Lemma~3).
Thus \(\mathcal A(\zeta)\) is discrete almost everywhere.

Suppose \(\nu\) is standard, and let \(Y\) be the set of
\(\zeta\in Z\) for which \(\mathcal A(\zeta)\) is not continuous. We
shall show that there exists an abelian projection
\[
  G=\int^{\oplus}G(\zeta)\,d\nu(\zeta)
\]
of \(\mathcal A\), with \(G(\zeta)\neq0\) for \(\zeta\in Y\). If
\(\mathcal A\) is continuous, it will follow that \(Y\) is null, and
hence that \(\mathcal A(\zeta)\) is continuous almost everywhere.

To prove our assertion, the usual technique reduces us to the case in
which: (a) \(\zeta\mapsto\mathfrak H(\zeta)\) is the constant field
corresponding to a Hilbert space \(\mathfrak H_0\); and (b) \(Z\) is
compact and metrizable, and there exist continuous mappings
\[
  \zeta\mapsto T_1(\zeta),\quad
  \zeta\mapsto T_2(\zeta),\quad\ldots,\qquad
  \zeta\mapsto T_1'(\zeta),\quad
  \zeta\mapsto T_2'(\zeta),\quad\ldots
\]
of \(Z\) into the unit ball \(\mathcal L_1\) of
\(\mathcal L(\mathfrak H_0)\), furnished with the strong topology, such
that, for every \(\zeta\in Z\), the \(T_i(\zeta)\) generate
\(\mathcal A(\zeta)\), and the \(T_i'(\zeta)\) generate
\(\mathcal A(\zeta)'\). We may further suppose that, for every
\(i=1,2,\ldots\), there is a \(j=1,2,\ldots\) such that
\[
  T_i(\zeta)^*=T_j(\zeta),
  \qquad
  T_i'(\zeta)^*=T_j'(\zeta)
\]
for every \(\zeta\in Z\).

Let \(\mathcal M\) be the set of pairs
\((\zeta,T)\in Z\times\mathcal L_1\) such that:
\begin{enumerate}
\item \(TT_i'(\zeta)=T_i'(\zeta)T\) for \(i=1,2,\ldots\);
\item \(T\) is a projection;
\item
  \[
    (TT_i(\zeta)T)(TT_j(\zeta)T)
    =(TT_j(\zeta)T)(TT_i(\zeta)T)
    \quad(i,j=1,2,\ldots);
  \]
\item \(T\neq0\).
\end{enumerate}
Conditions 1, 2, and 3 define closed subsets of
\(Z\times\mathcal L_1\), while condition 4 defines an open subset; thus
\(\mathcal M\) is Borel. Moreover, for \(\zeta\in Y\), the set of
\(T\) such that \((\zeta,T)\in\mathcal M\) is nonempty. By Appendix~V,
there is a \(\nu\)-measurable mapping
\(\zeta\mapsto G(\zeta)\), defined on a \(\nu\)-measurable subset
\(X\) of \(Z\) containing \(Y\), such that
\((\zeta,G(\zeta))\in\mathcal M\) for every \(\zeta\in X\). The
\(G(\zeta)\) are projections by condition 2, belong to
\(\mathcal A(\zeta)\) by condition 1, are abelian by condition 3, and
are nonzero by condition 4. Put \(G(\zeta)=0\) for \(\zeta\in Z-X\),
and let
% Source PDF page 190 (printed p. 183)
\phantomsection\label{srcpage:183}
\[
  G=\int^{\oplus}G(\zeta)\,d\nu(\zeta).
\]
Then \(G\) is an abelian projection of \(\mathcal A\), and
\(G(\zeta)\neq0\) on \(X\), hence on \(Y\).
\end{proof}

\setcounter{corollary}{0}
\begin{corollary}\label{cor:II-3-1-discrete-continuous}
Suppose \(\nu\) is standard. Let \(E\) (respectively, \(F\)) be the
greatest projection in the centre of \(\mathcal A\) such that
\(\mathcal A_E\) (respectively, \(\mathcal A_F\)) is discrete
(respectively, continuous). Write
\[
  E=\int^{\oplus}E(\zeta)\,d\nu(\zeta),
  \qquad
  F=\int^{\oplus}F(\zeta)\,d\nu(\zeta).
\]
Then, almost everywhere, \(E(\zeta)\) (respectively, \(F(\zeta)\)) is
the greatest projection in the centre of \(\mathcal A(\zeta)\) such
that \(\mathcal A(\zeta)_{E(\zeta)}\) (respectively,
\(\mathcal A(\zeta)_{F(\zeta)}\)) is discrete (respectively,
continuous).
\end{corollary}

\begin{proof}
Almost everywhere, \(E(\zeta)\) and \(F(\zeta)\) are orthogonal
projections with sum \(I\), in the centre of \(\mathcal A(\zeta)\)
(Theorem~4), while \(\mathcal A(\zeta)_{E(\zeta)}\) is discrete and
\(\mathcal A(\zeta)_{F(\zeta)}\) is continuous (Propositions~6 and~7).
\end{proof}

\begin{corollary}\label{cor:II-3-2-discrete-continuous}
Suppose \(\nu\) is standard. The algebra \(\mathcal A\) is discrete
(respectively, continuous) if and only if \(\mathcal A(\zeta)\) is
discrete (respectively, continuous) almost everywhere.
\end{corollary}

\begin{proof}
This follows from Corollary~1.
\end{proof}

Bibliography: \ArticleRef{10}, \ArticleRef{80}, \ArticleRef{117}.

\NumberedParagraph{6}{Measurable fields of homomorphisms}
\label{no:II-3-6}
Let \(\zeta\mapsto\mathfrak H(\zeta)\) and
\(\zeta\mapsto\mathfrak K(\zeta)\) be two \(\nu\)-measurable fields of
complex Hilbert spaces on \(Z\), and put
\[
  \mathfrak H=\int^{\oplus}\mathfrak H(\zeta)\,d\nu(\zeta),
  \qquad
  \mathfrak K=\int^{\oplus}\mathfrak K(\zeta)\,d\nu(\zeta).
\]
Let
\[
  \zeta\longmapsto\mathcal A(\zeta)
  \subset\mathcal L(\mathfrak H(\zeta)),
  \qquad
  \zeta\longmapsto\mathcal B(\zeta)
  \subset\mathcal L(\mathfrak K(\zeta))
\]
be measurable fields of von Neumann algebras on \(Z\), and put
\[
  \mathcal A=\int^{\oplus}\mathcal A(\zeta)\,d\nu(\zeta),
  \qquad
  \mathcal B=\int^{\oplus}\mathcal B(\zeta)\,d\nu(\zeta).
\]
For every \(\zeta\in Z\), let \(\Phi_{\zeta}\) be a homomorphism of
\(\mathcal A(\zeta)\) into \(\mathcal B(\zeta)\). The mapping
\(\zeta\mapsto\Phi_{\zeta}\) is called a \emph{field of
homomorphisms}. This field is said to be \emph{measurable} if, for every
measurable field of operators
\[
  \zeta\longmapsto T(\zeta)\in\mathcal A(\zeta),
\]
the field
\[
  \zeta\longmapsto\Phi_{\zeta}(T(\zeta))
  \in\mathcal B(\zeta)
\]
is measurable. Then, for
\(T=\int^{\oplus}T(\zeta)\,d\nu(\zeta)\in\mathcal A\), put
\[
  \Phi(T)=\int^{\oplus}\Phi_{\zeta}(T(\zeta))\,d\nu(\zeta)
  \in\mathcal B.
\]
% Source PDF page 191 (printed p. 184)
\phantomsection\label{srcpage:184}
It is immediate that \(\Phi\) is a homomorphism of the involutive
algebra \(\mathcal A\) into the involutive algebra \(\mathcal B\).

\begin{proposition}\label{prop:II-3-8}
If the \(\Phi_{\zeta}\) are normal, then \(\Phi\) is normal.
\end{proposition}

\begin{proof}
Let \((T_{\iota})\) be an increasing directed set in
\(\mathcal A_+\), with supremum \(T\in\mathcal A_+\). Since
\(\mathcal A\) is of countable type, there is an increasing sequence
\((T_1,T_2,\ldots)\) of operators \(T_{\iota}\), with supremum \(T\)
(Chapter~I, \S~3, corollary to Proposition~1). Put
\[
  T_i=\int^{\oplus}T_i(\zeta)\,d\nu(\zeta),
  \qquad
  T=\int^{\oplus}T(\zeta)\,d\nu(\zeta),
\]
where
\[
  T_i(\zeta)\in\mathcal A(\zeta)_+,
  \qquad
  T(\zeta)\in\mathcal A(\zeta)_+
  \quad(\zeta\in Z).
\]
By passing to a subsequence, we may suppose that, almost everywhere,
the \(T_i(\zeta)\) increase and converge strongly to \(T(\zeta)\)
(\S~2, Proposition~4(i)). Hence, almost everywhere,
\(\Phi_{\zeta}(T(\zeta))\) is the strong limit of the increasing
sequence \(\Phi_{\zeta}(T_i(\zeta))\). Thus \(\Phi(T)\) is the strong
limit of the increasing sequence \(\Phi(T_i)\) (\S~2,
Proposition~4(ii)), and consequently is its supremum.
\end{proof}

\begin{proposition}\label{prop:II-3-9}
Let \(\zeta\mapsto\Phi_{\zeta}\) and
\(\zeta\mapsto\Psi_{\zeta}\) be two measurable fields of normal
homomorphisms of \(\mathcal A(\zeta)\) into \(\mathcal B(\zeta)\),
defining the same homomorphism of \(\mathcal A\) into \(\mathcal B\).
Then \(\Phi_{\zeta}=\Psi_{\zeta}\) almost everywhere.
\end{proposition}

\begin{proof}
Let
\[
  T_1=\int^{\oplus}T_1(\zeta)\,d\nu(\zeta),
  \qquad
  T_2=\int^{\oplus}T_2(\zeta)\,d\nu(\zeta),
  \qquad\ldots
\]
be operators of \(\mathcal A\) such that, almost everywhere, the
\(T_i(\zeta)\) generate \(\mathcal A(\zeta)\). We have
\[
  \int^{\oplus}\Phi_{\zeta}(T_i(\zeta))\,d\nu(\zeta)
  =\int^{\oplus}\Psi_{\zeta}(T_i(\zeta))\,d\nu(\zeta),
\]
and therefore, almost everywhere,
\[
  \Phi_{\zeta}(T_i(\zeta))
  =\Psi_{\zeta}(T_i(\zeta))
  \qquad\text{for every }i.
\]
Thus, almost everywhere, \(\Phi_{\zeta}\) and \(\Psi_{\zeta}\)
coincide on an involutive subalgebra of \(\mathcal A(\zeta)\) which is
everywhere ultraweakly dense in \(\mathcal A(\zeta)\). Since
\(\Phi_{\zeta}\) and \(\Psi_{\zeta}\) are ultraweakly continuous
(Chapter~I, \S~4, Theorem~2), we have
\(\Phi_{\zeta}=\Psi_{\zeta}\) almost everywhere.
\end{proof}

\begin{definition}\label{def:II-3-3}
A normal homomorphism \(\Phi\) of \(\mathcal A\) into
\(\mathcal B\) is said to be \emph{decomposable} if it is defined by
a measurable field \(\zeta\mapsto\Phi_{\zeta}\) of normal
homomorphisms of \(\mathcal A(\zeta)\) into \(\mathcal B(\zeta)\).
% Source PDF page 192 (printed p. 185)
\phantomsection\label{srcpage:185}
We then write
\[
  \Phi=\int^{\oplus}\Phi_{\zeta}\,d\nu(\zeta).
\]
\end{definition}

By Proposition~9, the \(\Phi_{\zeta}\) are determined by \(\Phi\),
up to null sets.

\begin{proposition}\label{prop:II-3-10}
If \(\Phi_{\zeta}\) is, almost everywhere, an isomorphism of
\(\mathcal A(\zeta)\) onto \(\mathcal B(\zeta)\), then \(\Phi\) is an
isomorphism of \(\mathcal A\) onto \(\mathcal B\).
\end{proposition}

\begin{proof}
Let
\[
  T_1=\int^{\oplus}T_1(\zeta)\,d\nu(\zeta),
  \qquad
  T_2=\int^{\oplus}T_2(\zeta)\,d\nu(\zeta),
  \qquad\ldots
\]
be operators of \(\mathcal A\) such that, almost everywhere, the
\(T_i(\zeta)\) generate \(\mathcal A(\zeta)\). Then, almost
everywhere, the \(\Phi_{\zeta}(T_i(\zeta))\) generate
\(\mathcal B(\zeta)\). Hence \(\mathcal B\) is generated by the
\(\Phi(T_i)\) and by the algebra of diagonalizable operators on
\(\mathfrak K\). But \(\Phi(T_i)\in\Phi(\mathcal A)\), and every
diagonalizable operator on \(\mathfrak K\) is the image under \(\Phi\)
of a diagonalizable operator on \(\mathfrak H\). Therefore
\(\Phi(\mathcal A)\), which is a von Neumann algebra (Chapter~I,
\S~4, Theorem~2, Corollary~2), is equal to \(\mathcal B\). Moreover,
if
\[
  T=\int^{\oplus}T(\zeta)\,d\nu(\zeta)\in\mathcal A,
\]
the relation \(\Phi(T)=0\) implies
\(\Phi_{\zeta}(T(\zeta))=0\) almost everywhere, hence
\(T(\zeta)=0\) almost everywhere, and therefore \(T=0\).
\end{proof}

\begin{proposition}\label{prop:II-3-11}
Let \(\zeta\mapsto\mathfrak H(\zeta)\) be a \(\nu\)-measurable field
of complex Hilbert spaces on \(Z\), let
\[
  \mathfrak H=\int^{\oplus}\mathfrak H(\zeta)\,d\nu(\zeta),
\]
let
\(\zeta\mapsto\mathcal A(\zeta)\subset
\mathcal L(\mathfrak H(\zeta))\) be a \(\nu\)-measurable field of von
Neumann algebras, and let
\[
  \mathcal A=\int^{\oplus}\mathcal A(\zeta)\,d\nu(\zeta).
\]
Let \(\mathfrak K\) be a Hilbert space, let \(\mathcal B\) be a von
Neumann algebra on \(\mathfrak K\) such that \(\mathcal B'\) is of
countable type, and let \(\Phi\) be an isomorphism of \(\mathcal A\)
onto \(\mathcal B\). Suppose that \(\nu\) is standard. Then there
exist a \(\nu\)-measurable field
\(\zeta\mapsto\mathfrak K(\zeta)\) of Hilbert spaces on \(Z\), a
\(\nu\)-measurable field
\(\zeta\mapsto\mathcal B(\zeta)\subset
\mathcal L(\mathfrak K(\zeta))\) of von Neumann algebras, a
\(\nu\)-measurable field \(\zeta\mapsto\Phi_{\zeta}\) of
isomorphisms of \(\mathcal A(\zeta)\) onto \(\mathcal B(\zeta)\),
and an isomorphism of
\(\int^{\oplus}\mathfrak K(\zeta)\,d\nu(\zeta)\) onto
\(\mathfrak K\), which carries
\(\int^{\oplus}\mathcal B(\zeta)\,d\nu(\zeta)\) onto
\(\mathcal B\), and carries
\(\int^{\oplus}\Phi_{\zeta}\,d\nu(\zeta)\) onto \(\Phi\).
\end{proposition}

% Source PDF page 193 (printed p. 186)
\phantomsection\label{srcpage:186}
\begin{proof}
By Chapter~I, \S~4, Theorem~3, it is enough to consider the following
two cases.
\begin{enumerate}[label=\textit{\alph*.},leftmargin=*]
\item Let \(\mathfrak H_0\) be a Hilbert space, and suppose that
\[
  \mathfrak K=\mathfrak H\otimes\mathfrak H_0,
  \qquad
  \mathcal B=\mathcal A\otimes\mathbb C I_{\mathfrak H_0},
\]
and that \(\Phi\) is the isomorphism
\(T\mapsto T\otimes I_{\mathfrak H_0}\). Moreover (Chapter~I,
\S~4, no.~4, Remark), we may suppose that \(\mathfrak H_0\) has a
countable base. Then \(\mathfrak K\) is identified with
\[
  \int^{\oplus}
    \bigl(\mathfrak H(\zeta)\otimes\mathfrak H_0\bigr)
    \,d\nu(\zeta)
\]
(\S~1, Proposition~11),
\(\mathcal A\otimes\mathbb C I_{\mathfrak H_0}\) is identified with
\[
  \int^{\oplus}
    \bigl(\mathcal A(\zeta)\otimes
          \mathbb C I_{\mathfrak H_0}\bigr)
    \,d\nu(\zeta)
\]
(Proposition~3), and \(\Phi\) is identified with
\(\int^{\oplus}\Phi_{\zeta}\,d\nu(\zeta)\), where, for every
\(\zeta\), \(\Phi_{\zeta}\) is the isomorphism
\[
  S\longmapsto S\otimes I_{\mathfrak H_0}
\]
of \(\mathcal A(\zeta)\) onto
\(\mathcal A(\zeta)\otimes\mathbb C I_{\mathfrak H_0}\)
(\S~2, Proposition~9).

\item Let \(E'\) be a projection of \(\mathcal A'\) with central
support \(I\), and suppose that
\[
  \mathfrak K=E'(\mathfrak H),
  \qquad
  \mathcal B=\mathcal A_{E'},
\]
and that \(\Phi\) is the isomorphism \(T\mapsto T_{E'}\). Put
\[
  E'=\int^{\oplus}E'(\zeta)\,d\nu(\zeta),
  \qquad
  \mathfrak K(\zeta)=E'(\zeta)\mathfrak H(\zeta),
  \qquad
  \mathcal B(\zeta)=\mathcal A(\zeta)_{E'(\zeta)}.
\]
Almost everywhere, \(E'(\zeta)\) is a projection of
\(\mathcal A(\zeta)'\) with central support \(I\), by Lemma~3
[taking account of the fact that the field
\(\zeta\mapsto\mathcal A(\zeta)'\) is measurable and that
\(\mathcal A'=\int^{\oplus}\mathcal A(\zeta)'\,d\nu(\zeta)\), since
\(\nu\) is standard]. Hence the mapping
\(S\mapsto S_{E'(\zeta)}\) is an isomorphism \(\Phi_{\zeta}\) of
\(\mathcal A(\zeta)\) onto \(\mathcal B(\zeta)\). We have
\[
  \mathfrak K
  =\int^{\oplus}\mathfrak K(\zeta)\,d\nu(\zeta)
  \quad(\S~2,\text{ no.~3}),
\]
\[
  \mathcal B
  =\int^{\oplus}\mathcal B(\zeta)\,d\nu(\zeta)
  \quad(\text{Proposition~6}),
\]
and
\[
  \Phi=\int^{\oplus}\Phi_{\zeta}\,d\nu(\zeta).
\]
\end{enumerate}
\end{proof}

\par\medskip
\noindent\textit{Exercises.}---
\begin{enumerate}[label=\arabic*.,leftmargin=*,itemsep=0.5\baselineskip]
\item\label{ex:II-3-1} Show that Theorem~1(ii) is false if the hypothesis that \(I\)
is countable is omitted. [Choose the fields
\(\zeta\mapsto T_{\iota}(\zeta)\) so that the \(T_{\iota}\) are
diagonalizable, but, for every \(\zeta\in Z\), the
\(T_{\iota}(\zeta)\) generate
\(\mathcal L(\mathfrak H(\zeta))\).]

\item\label{ex:II-3-2} Let
\[
  \zeta\mapsto\mathcal A(\zeta)
  \subset\mathcal L(\mathfrak H(\zeta)),
  \qquad
  \zeta\mapsto\mathcal B(\zeta)
  \subset\mathcal L(\mathfrak K(\zeta))
\]
be two measurable fields of von Neumann algebras on \(Z\).
\begin{enumerate}[label=\textit{\alph*.},leftmargin=*]
\item The set \(Y\) of \(\zeta\in Z\) such that
\(\mathcal A(\zeta)\) and \(\mathcal B(\zeta)\) are spatially
isomorphic is measurable.
\item If \(\nu\) is standard, the set of \(\zeta\in Z\) such that
\(\mathcal A(\zeta)=\mathcal B(\zeta)\), or such that
\(\mathcal A(\zeta)\subset\mathcal B(\zeta)\), is measurable.
[Use Lemma~1 and part~\textit{a}.]

% Source PDF page 194 (printed p. 187)
\phantomsection\label{srcpage:187}
\item If \(\nu\) is standard, the set of \(\zeta\in Z\) such that
\(\mathcal A(\zeta)\) is a factor is measurable. [To say that
\(\mathcal A(\zeta)\) is a factor amounts to saying that
\(\mathcal A(\zeta)\cap\mathcal A(\zeta)'\) is spatially isomorphic
to \(\mathcal L(\mathfrak H(\zeta))\).]
\item Suppose that
\(\mathcal A(\zeta)\subset\mathcal B(\zeta)\) and
\(\mathcal A(\zeta)\neq\mathcal B(\zeta)\) for every \(\zeta\).
Show that there exists a measurable field
\(\zeta\mapsto E(\zeta)\) of projections on \(Z\) such that
\(E(\zeta)\in\mathcal B(\zeta)\) and
\(E(\zeta)\notin\mathcal A(\zeta)\) for every \(\zeta\). [Use
Chapter~I, \S~2, Exercise~2(b).] \ArticleRef{80},
\ArticleRef{338}, \ArticleRef{360}.
\end{enumerate}
\end{enumerate}

\section{Fields of Hilbert Algebras}\label{sec:II-4}

\NumberedParagraph{1}{Measurable Fields of Hilbert Algebras}
\label{no:II-4-1}
Let \(\zeta\mapsto\mathfrak H(\zeta)\) be a \(\nu\)-measurable
field of complex Hilbert spaces on \(Z\). For every \(\zeta\in Z\),
let \(\mathfrak A(\zeta)\) be a Hilbert algebra having
\(\mathfrak H(\zeta)\) as its completed Hilbert space. The mapping
\(\zeta\mapsto\mathfrak A(\zeta)\) is called a \emph{field of Hilbert
algebras}.

\begin{definition}\label{def:II-4-1}
A field \(\zeta\mapsto\mathfrak A(\zeta)\) of Hilbert algebras on
\(Z\) is said to be \emph{measurable} if the following conditions are
satisfied:
\begin{enumerate}[label=\textup{(\roman*)},leftmargin=*]
\item If \(\zeta\mapsto x(\zeta)\) is a measurable field of vectors
such that \(x(\zeta)\in\mathfrak A(\zeta)\) for every \(\zeta\), then
the field \(\zeta\mapsto x(\zeta)^*\) is measurable.
\item If \(\zeta\mapsto x(\zeta)\) and
\(\zeta\mapsto y(\zeta)\) are measurable fields of vectors such that
\(x(\zeta),y(\zeta)\in\mathfrak A(\zeta)\) for every \(\zeta\), then
the field \(\zeta\mapsto x(\zeta)y(\zeta)\) is measurable.
\item There exists a fundamental sequence \((x_i)\) of measurable
fields of vectors such that
\(x_i(\zeta)\in\mathfrak A(\zeta)\) for every \(i\) and every
\(\zeta\).
\end{enumerate}
\end{definition}

The measurable fields of vectors \(x\) such that
\(x(\zeta)\in\mathfrak A(\zeta)\) for every \(\zeta\) form an
involutive algebra \(\mathfrak A_0\). If
\(x\in\mathfrak A_0\), the field of operators
\(\zeta\mapsto U_{x(\zeta)}\) is measurable; indeed, the fields of
vectors
\[
  \zeta\longmapsto U_{x(\zeta)}x_i(\zeta)
  =x(\zeta)x_i(\zeta)
\]
are measurable, and it is enough to apply \S~2, Proposition~1.
Similarly, the field of operators
\(\zeta\mapsto V_{x(\zeta)}\) is measurable.

Whenever convenient, we may suppose that the fields \(x_i\) in
Definition~1(iii) also have the following properties:
\begin{enumerate}[label=\textup{\arabic*\textdegree},leftmargin=*]
\item the functions \(\zeta\mapsto\lVert x_i(\zeta)\rVert\) are
essentially bounded;
\item the functions
\(\zeta\mapsto\lVert U_{x_i(\zeta)}\rVert\) are essentially bounded;
\item the set of these fields is stable under multiplication and
involution.
\end{enumerate}
Properties~1\textdegree{} and~2\textdegree{} are ensured by multiplying
the \(x_i\) by suitable numerical functions. Property~3\textdegree{}
is then obtained by enlarging the sequence \((x_i)\).

% Source PDF page 197 (printed p. 188)
\phantomsection\label{srcpage:188}
\begin{proposition}\label{prop:II-4-1}
Let \(\zeta\mapsto\mathfrak A(\zeta)\) be a field of Hilbert algebras
on \(Z\). Suppose that there exists a fundamental sequence \((x_i)\)
of measurable fields of vectors having the following properties:
\begin{enumerate}[label=\textup{\arabic*\textdegree},leftmargin=*]
\item \(x_i(\zeta)\in\mathfrak A(\zeta)\) for every \(\zeta\);
\item the fields \(\zeta\mapsto x_i(\zeta)^*\) are measurable;
\item the fields
\(\zeta\mapsto x_i(\zeta)x_j(\zeta)^*\) are measurable.
\end{enumerate}
Then the field \(\zeta\mapsto\mathfrak A(\zeta)\) is measurable.
\end{proposition}

\begin{proof}
If \(\zeta\mapsto x(\zeta)\) is a measurable field of vectors such
that \(x(\zeta)\in\mathfrak A(\zeta)\) for every \(\zeta\in Z\), then
\[
  \bigl(x(\zeta)^*\mid x_i(\zeta)\bigr)
  =\bigl(x_i(\zeta)^*\mid x(\zeta)\bigr)
\]
depends measurably on \(\zeta\); hence the field
\(\zeta\mapsto x(\zeta)^*\) is measurable. This is condition~(i) of
Definition~1. Moreover,
\[
  \bigl(x(\zeta)^*x_i(\zeta)\mid x_j(\zeta)\bigr)
  =\bigl(x(\zeta)^*\mid x_j(\zeta)x_i(\zeta)^*\bigr)
\]
depends measurably on \(\zeta\), and so the field
\(\zeta\mapsto x(\zeta)^*x_i(\zeta)\) is measurable. If
\(\zeta\mapsto y(\zeta)\) is another measurable field of vectors such
that \(y(\zeta)\in\mathfrak A(\zeta)\) for every \(\zeta\in Z\), then
\[
  \bigl(x(\zeta)y(\zeta)\mid x_i(\zeta)\bigr)
  =\bigl(y(\zeta)\mid x(\zeta)^*x_i(\zeta)\bigr)
\]
depends measurably on \(\zeta\). Hence the field
\(\zeta\mapsto x(\zeta)y(\zeta)\) is measurable. This is
condition~(ii) of Definition~1. Finally, condition~(iii) of
Definition~1 is plainly satisfied.
\end{proof}

Bibliography: \ArticleRef{13}, \ArticleRef{29}, \ArticleRef{86},
\ArticleRef{106}, \ArticleRef{116}, \ArticleRef{117}.

\NumberedParagraph{2}{Decomposable Hilbert Algebras}
\label{no:II-4-2}
Put
\[
  \mathfrak H
  =\int^{\oplus}\mathfrak H(\zeta)\,d\nu(\zeta),
\]
and let \(\mathfrak A\) be the set of fields
\(\zeta\mapsto x(\zeta)\) belonging to
\(\mathfrak A_0\cap\mathfrak H\) for which the field of operators
\(\zeta\mapsto U_{x(\zeta)}\) is essentially bounded (and hence so is
the field \(\zeta\mapsto V_{x(\zeta)}\)). Then \(\mathfrak A\) is an
involutive subalgebra of \(\mathfrak A_0\). Indeed, if
\(x,y\in\mathfrak A\), then, almost everywhere,
\[
  \lVert(xy)(\zeta)\rVert
  =\lVert U_{x(\zeta)}y(\zeta)\rVert
  \leq M\lVert y(\zeta)\rVert,
\]
where \(M\) denotes the essential supremum of
\(\zeta\mapsto\lVert U_{x(\zeta)}\rVert\); hence
\(xy\in\mathfrak H\). Also,
\[
  U_{(xy)(\zeta)}=U_{x(\zeta)}U_{y(\zeta)},
\]
so that the field of operators
\(\zeta\mapsto U_{(xy)(\zeta)}\) is essentially bounded; therefore
\(xy\in\mathfrak A\). On the other hand,
\[
  \lVert x^*(\zeta)\rVert=\lVert x(\zeta)\rVert,
  \qquad
  \lVert U_{x^*(\zeta)}\rVert=\lVert U_{x(\zeta)}\rVert,
\]
so \(x^*\in\mathfrak A\), which proves our assertion.

As a subspace of \(\mathfrak H\), \(\mathfrak A\) is a pre-Hilbert
space. We then have:

% Source PDF page 198 (printed p. 189)
\phantomsection\label{srcpage:189}
\begin{proposition}\label{prop:II-4-2}
The algebra \(\mathfrak A\) is a Hilbert algebra dense in
\(\mathfrak H\).
\end{proposition}

\begin{proof}
If \(x,y,z\in\mathfrak A\), then
\begin{align*}
  (x\mid y)
    &=\int\bigl(x(\zeta)\mid y(\zeta)\bigr)\,d\nu(\zeta)
      =\int\bigl(y^*(\zeta)\mid x^*(\zeta)\bigr)\,d\nu(\zeta)
      =(y^*\mid x^*),\\
  (xy\mid z)
    &=\int\bigl(x(\zeta)y(\zeta)\mid z(\zeta)\bigr)\,d\nu(\zeta)
      =\int\bigl(y(\zeta)\mid x^*(\zeta)z(\zeta)\bigr)\,d\nu(\zeta)
      =(y\mid x^*z),\\
  \lVert xy\rVert^2
    &=\int\lVert x(\zeta)y(\zeta)\rVert^2\,d\nu(\zeta)
      \leq M^2\int\lVert y(\zeta)\rVert^2\,d\nu(\zeta)
      =M^2\lVert y\rVert^2,
\end{align*}
where \(M\) is the essential supremum of the function
\(\zeta\mapsto\lVert U_{x(\zeta)}\rVert\). Finally, let \((x_i)\) be
a fundamental sequence of measurable fields of vectors such that
\(x_i(\zeta)\in\mathfrak A(\zeta)\) for every \(\zeta\), and such
that the functions
\(\zeta\mapsto\lVert x_i(\zeta)\rVert\) and
\(\zeta\mapsto\lVert U_{x_i(\zeta)}\rVert\) are bounded. For every
\(\zeta\in Z\), the \(x_i(\zeta)x_j(\zeta)\) form a total system in
\(\mathfrak H(\zeta)\). Hence the fields
\[
  \zeta\longmapsto
  f(\zeta)g(\zeta)x_i(\zeta)x_j(\zeta),
\]
where \(f\) and \(g\) are bounded measurable numerical functions
vanishing outside a set of finite measure, form a total system in
\(\mathfrak H\) (\S~1, Proposition~7). But
\(fx_i\in\mathfrak A\) and \(gx_j\in\mathfrak A\). This proves both
that \(\mathfrak A\) is dense in \(\mathfrak H\) and that it is a
Hilbert algebra.
\end{proof}

\begin{definition}\label{def:II-4-2}
A Hilbert algebra \(\mathfrak A\) in \(\mathfrak H\) is said to be
\emph{decomposable} if it is defined by a measurable field
\(\zeta\mapsto\mathfrak A(\zeta)\) of Hilbert algebras. We then write
\[
  \mathfrak A
  =\int^{\oplus}\mathfrak A(\zeta)\,d\nu(\zeta).
\]
\end{definition}

Bibliography: \ArticleRef{13}, \ArticleRef{29}, \ArticleRef{86},
\ArticleRef{106}, \ArticleRef{116}, \ArticleRef{117}.

\NumberedParagraph{3}{Involution and the von Neumann Algebras
Associated with \(\mathfrak A\)}
\label{no:II-4-3}
Let \(J\) (respectively, \(J(\zeta)\)) be the involution of
\(\mathfrak H\) (respectively, \(\mathfrak H(\zeta)\)) canonically
defined by \(\mathfrak A\) (respectively,
\(\mathfrak A(\zeta)\)). Let \(\zeta\mapsto x(\zeta)\) be a
measurable field of vectors. With the notation of Definition~1,
\[
  \bigl(J(\zeta)x(\zeta)\mid x_i(\zeta)\bigr)
  =\bigl(x_i(\zeta)^*\mid x(\zeta)\bigr)
\]
depends measurably on \(\zeta\); hence the field
\(\zeta\mapsto J(\zeta)x(\zeta)\) is measurable. If the field
\(\zeta\mapsto x(\zeta)\) belongs to \(\mathfrak H\), then the field
\(\zeta\mapsto J(\zeta)x(\zeta)\) also belongs to \(\mathfrak H\).
% Source PDF page 199 (printed p. 190)
\phantomsection\label{srcpage:190}
Denote by
\(\int^{\oplus}J(\zeta)\,d\nu(\zeta)\) the mapping of
\(\mathfrak H\) into \(\mathfrak H\) so defined. It is an isometric
antilinear mapping of \(\mathfrak H\) into itself, which coincides
with \(J\) on \(\mathfrak A\). Consequently,
\[
  J=\int^{\oplus}J(\zeta)\,d\nu(\zeta).
\]

For \(x\in\mathfrak A\), the field of operators
\(\zeta\mapsto U_{x(\zeta)}\) is measurable and essentially bounded.
We have
\[
  U_x=\int^{\oplus}U_{x(\zeta)}\,d\nu(\zeta);
\]
indeed, for \(y\in\mathfrak A\),
\[
  (U_xy)(\zeta)=(xy)(\zeta)=U_{x(\zeta)}y(\zeta),
\]
so \(\int^{\oplus}U_{x(\zeta)}\,d\nu(\zeta)\) coincides with \(U_x\)
on \(\mathfrak A\), and hence on \(\mathfrak H\). Similarly,
\[
  V_x=\int^{\oplus}V_{x(\zeta)}\,d\nu(\zeta).
\]

\begin{proposition}\label{prop:II-4-3}
The von Neumann algebras \(\mathcal U(\mathfrak A)\) and
\(\mathcal V(\mathfrak A)\) are decomposable, and
\[
  \mathcal U(\mathfrak A)
   =\int^{\oplus}\mathcal U(\mathfrak A(\zeta))\,d\nu(\zeta),
  \qquad
  \mathcal V(\mathfrak A)
   =\int^{\oplus}\mathcal V(\mathfrak A(\zeta))\,d\nu(\zeta).
\]
\end{proposition}

\begin{proof}
For every \(x\in\mathfrak A\), we have
\(U_x,V_x\in\mathfrak Z'\), where \(\mathfrak Z\) denotes the
algebra of diagonalizable operators. Hence
\[
  \mathcal U(\mathfrak A)\subset\mathfrak Z',
  \qquad
  \mathcal V(\mathfrak A)\subset\mathfrak Z'.
\]
Consequently,
\[
  \mathfrak Z\subset\mathcal U(\mathfrak A)'
  =\mathcal V(\mathfrak A),
  \qquad
  \mathfrak Z\subset\mathcal V(\mathfrak A)'
  =\mathcal U(\mathfrak A).
\]
Let \((x_i)\) be a fundamental sequence of measurable fields of
vectors. Suppose that \(x_i(\zeta)\in\mathfrak A(\zeta)\) for every
\(\zeta\) and every \(i\), and that the functions
\[
  \zeta\longmapsto\lVert x_i(\zeta)\rVert,
  \qquad
  \zeta\longmapsto\lVert U_{x_i(\zeta)}\rVert
\]
are bounded. Put
\[
  T_i=\int^{\oplus}U_{x_i(\zeta)}\,d\nu(\zeta).
\]
By Chapter~I, \S~5, Proposition~1,
\(\mathcal U(\mathfrak A(\zeta))\) is the von Neumann algebra
generated by the \(U_{x_i(\zeta)}\). Using \S~3, Theorem~1, we see
that, in order to prove the equality
\[
  \mathcal U(\mathfrak A)
  =\int^{\oplus}\mathcal U(\mathfrak A(\zeta))\,d\nu(\zeta),
\]
it is enough to show that \(\mathcal U(\mathfrak A)\) is equal to the
von Neumann algebra \(\mathcal C\) generated by \(\mathfrak Z\) and
the \(T_i\). First, \(\mathfrak Z\subset\mathcal U(\mathfrak A)\),
and, since the \(T_i\) commute with the \(V_x\),
% Source PDF page 200 (printed p. 191)
\phantomsection\label{srcpage:191}
for \(x\in\mathfrak A\), they belong to
\(\mathcal U(\mathfrak A)\); hence
\(\mathcal C\subset\mathcal U(\mathfrak A)\). Next, among the
operators \(T_iT_f\), where \(T_f\) is the diagonalizable operator
defined by a bounded measurable numerical function \(f\), there are
operators \(U_y\) for which the \(y\)'s form a total subset of
\(\mathfrak A\) (by \S~1, Proposition~7); these operators generate
\(\mathcal U(\mathfrak A)\) (Chapter~I, \S~5, Proposition~1).
Therefore \(\mathcal C=\mathcal U(\mathfrak A)\), and consequently
\[
  \mathcal U(\mathfrak A)
  =\int^{\oplus}\mathcal U(\mathfrak A(\zeta))\,d\nu(\zeta).
\]
Similarly,
\[
  \mathcal V(\mathfrak A)
  =\int^{\oplus}\mathcal V(\mathfrak A(\zeta))\,d\nu(\zeta).
\]
\end{proof}

\begin{corollary*}
The algebra \(\mathfrak Z\) is the centre of both
\(\mathcal U(\mathfrak A)\) and \(\mathcal V(\mathfrak A)\) if and
only if, almost everywhere,
\(\mathcal U(\mathfrak A(\zeta))\) and
\(\mathcal V(\mathfrak A(\zeta))\) are factors.
\end{corollary*}

\begin{proof}
This follows from Proposition~3 and \S~3, Theorem~3.
\end{proof}

Bibliography: \ArticleRef{13}, \ArticleRef{29}, \ArticleRef{86},
\ArticleRef{106}, \ArticleRef{116}, \ArticleRef{117}.

\NumberedParagraph{4}{Elements Bounded Relative to \(\mathfrak A\)}
\label{no:II-4-4}

\begin{proposition}\label{prop:II-4-4}
If \(\zeta\mapsto a(\zeta)\) is a measurable field of vectors such
that, for every \(\zeta\in Z\), \(a(\zeta)\) is bounded relative to
\(\mathfrak A(\zeta)\), then the fields
\(\zeta\mapsto U_{a(\zeta)}\) and
\(\zeta\mapsto V_{a(\zeta)}\) are measurable.
\end{proposition}

\begin{proof}
Let \((x_i)\) be a fundamental sequence of measurable fields of
vectors such that \(x_i(\zeta)\in\mathfrak A(\zeta)\) for every
\(\zeta\). The fields
\[
  \zeta\longmapsto U_{a(\zeta)}x_i(\zeta)
  =V_{x_i(\zeta)}a(\zeta)
\]
are measurable; hence (\S~2, Proposition~1) the field
\(\zeta\mapsto U_{a(\zeta)}\) is measurable. Similarly, the field
\(\zeta\mapsto V_{a(\zeta)}\) is measurable.
\end{proof}

\begin{proposition}\label{prop:II-4-5}
Let \(a\in\mathfrak H\). The element \(a\) is bounded if and only if
\(a(\zeta)\) is bounded almost everywhere and the fields
\(\zeta\mapsto U_{a(\zeta)}\) and
\(\zeta\mapsto V_{a(\zeta)}\) are essentially bounded. In that case,
\[
  U_a=\int^{\oplus}U_{a(\zeta)}\,d\nu(\zeta),
  \qquad
  V_a=\int^{\oplus}V_{a(\zeta)}\,d\nu(\zeta).
\]
\end{proposition}

\begin{proof}
Suppose that \(a(\zeta)\) is bounded almost everywhere and that
\(\lVert U_{a(\zeta)}\rVert\leq M\) almost everywhere. For every
\(x\in\mathfrak A\), we have, almost everywhere,
\[
  \lVert V_{x(\zeta)}a(\zeta)\rVert
  =\lVert U_{a(\zeta)}x(\zeta)\rVert
  \leq M\lVert x(\zeta)\rVert;
\]
hence \(\lVert V_xa\rVert\leq M\lVert x\rVert\), and \(a\) is
bounded.

% Source PDF page 201 (printed p. 192)
\phantomsection\label{srcpage:192}
Conversely, suppose that \(a\) is bounded. Since
\(U_a\in\mathfrak Z'\), the operator \(U_a\) is decomposable. Write
\[
  U_a=\int^{\oplus}T(\zeta)\,d\nu(\zeta).
\]
For every \(x\in\mathfrak A\),
\[
  U_ax=V_xa,
  \qquad\text{so}\qquad
  T(\zeta)x(\zeta)=V_{x(\zeta)}a(\zeta)
\]
almost everywhere. Let \((x_i)\) be a fundamental sequence of
measurable fields of vectors such that:
\begin{enumerate}[label=\textup{\arabic*\textdegree},leftmargin=*]
\item \(x_i(\zeta)\in\mathfrak A(\zeta)\) for every \(\zeta\);
\item the functions
\(\zeta\mapsto\lVert x_i(\zeta)\rVert\) and
\(\zeta\mapsto\lVert U_{x_i(\zeta)}\rVert\) are bounded;
\item the set of these fields is stable under multiplication and
involution.
\end{enumerate}
Let \(Y\) be a measurable subset of \(Z\) of finite measure, and let
\(y_i\) be the field equal to \(x_i\) on \(Y\) and to \(0\) on
\(Z-Y\). Then \(y_i\in\mathfrak A\), so
\[
  T(\zeta)y_i(\zeta)=V_{y_i(\zeta)}a(\zeta)
\]
for every \(i\), almost everywhere. Thus, almost everywhere on \(Y\),
\(a(\zeta)\) is bounded relative to an involutive subalgebra dense in
\(\mathfrak A(\zeta)\), and consequently is bounded relative to
\(\mathfrak A(\zeta)\). Hence \(a(\zeta)\) is bounded relative to
\(\mathfrak A(\zeta)\) almost everywhere on \(Z\). Moreover,
\(T(\zeta)y_i(\zeta)=U_{a(\zeta)}y_i(\zeta)\) almost everywhere;
therefore \(T(\zeta)=U_{a(\zeta)}\) almost everywhere on \(Y\), and
hence almost everywhere on \(Z\). Thus the field
\(\zeta\mapsto U_{a(\zeta)}\) is essentially bounded, and
\[
  U_a=\int^{\oplus}U_{a(\zeta)}\,d\nu(\zeta).
\]
Similarly,
\[
  V_a=\int^{\oplus}V_{a(\zeta)}\,d\nu(\zeta).
\]
\end{proof}

\begin{proposition}\label{prop:II-4-6}
For every \(\zeta\in Z\), let \(\mathfrak B(\zeta)\) be the completed
Hilbert algebra corresponding to \(\mathfrak A(\zeta)\). Then the
field \(\zeta\mapsto\mathfrak B(\zeta)\) is measurable, and
\[
  \mathfrak B
  =\int^{\oplus}\mathfrak B(\zeta)\,d\nu(\zeta)
\]
is the completed Hilbert algebra corresponding to \(\mathfrak A\).
\end{proposition}

\begin{proof}
Let \(\zeta\mapsto a(\zeta)\) and
\(\zeta\mapsto b(\zeta)\) be two measurable fields of vectors such
that \(a(\zeta),b(\zeta)\in\mathfrak B(\zeta)\) for every
\(\zeta\), that is, such that \(a(\zeta)\) and \(b(\zeta)\) are
bounded for every \(\zeta\). Then the field
\[
  \zeta\longmapsto a(\zeta)b(\zeta)
  =U_{a(\zeta)}b(\zeta)
\]
is measurable. On the other hand, the field
\[
  \zeta\longmapsto a(\zeta)^*=J(\zeta)a(\zeta)
\]
% Source PDF page 202 (printed p. 193)
\phantomsection\label{srcpage:193}
is measurable. Since the field
\(\zeta\mapsto\mathfrak B(\zeta)\) plainly satisfies
condition~(iii) of Definition~1, this field is measurable. Thus
\[
  \mathfrak B
  =\int^{\oplus}\mathfrak B(\zeta)\,d\nu(\zeta)
\]
is a Hilbert algebra dense in \(\mathfrak H\), of which
\(\mathfrak A\) is an involutive subalgebra. By the definition of
\(\mathfrak B\) and Proposition~5, every element bounded relative to
\(\mathfrak A\) belongs to \(\mathfrak B\).
\end{proof}

Bibliography: \ArticleRef{116}.

\NumberedParagraph{5}{Elements Central Relative to \(\mathfrak A\)}
\label{no:II-4-5}
Let \(\mathfrak C\) (respectively, \(\mathfrak C(\zeta)\)) be the
subspace of \(\mathfrak H\) (respectively,
\(\mathfrak H(\zeta)\)) formed by the elements central relative to
\(\mathfrak A\) (respectively, \(\mathfrak A(\zeta)\)).

\begin{proposition}\label{prop:II-4-7}
The \(\mathfrak C(\zeta)\) form a measurable field of subspaces, and
\[
  \mathfrak C
  =\int^{\oplus}\mathfrak C(\zeta)\,d\nu(\zeta).
\]
\end{proposition}

\begin{proof}
Let \((x_i)\) be a fundamental sequence of measurable fields of
vectors such that \(x_i(\zeta)\in\mathfrak A(\zeta)\) for every
\(i\) and every \(\zeta\). For every \(\zeta\in Z\), the vectors
\[
  x_i(\zeta)x_j(\zeta)-x_j(\zeta)x_i(\zeta)
\]
generate the same closed vector subspace \(\mathfrak C'(\zeta)\) of
\(\mathfrak H(\zeta)\) as the vectors \(xy-yx\), with
\(x,y\in\mathfrak A(\zeta)\). The field of subspaces
\(\mathfrak C'(\zeta)\) is measurable (\S~1, Proposition~9(ii)).
Since
\[
  P_{\mathfrak C(\zeta)}=I-P_{\mathfrak C'(\zeta)},
\]
the field \(\zeta\mapsto P_{\mathfrak C(\zeta)}\), and consequently
the field \(\zeta\mapsto\mathfrak C(\zeta)\), are measurable
(\S~1, Proposition~9(iii)). Suppose the fields \(x_i\) have been
chosen so that the functions
\[
  \zeta\longmapsto\lVert x_i(\zeta)\rVert,
  \qquad
  \zeta\longmapsto\lVert U_{x_i(\zeta)}\rVert
\]
are bounded, and so that the set of the \(x_i\) is also stable under
involution and multiplication. Let \(a\in\mathfrak H\). If
\(a(\zeta)\in\mathfrak C(\zeta)\) almost everywhere, then
\[
  U_{x_i(\zeta)}a(\zeta)
  =J(\zeta)U_{x_i(\zeta)}^*J(\zeta)a(\zeta)
\]
almost everywhere. Put
\[
  T_i=\int^{\oplus}U_{x_i(\zeta)}\,d\nu(\zeta).
\]
Then \(T_ia=JT_i^*Ja\), and consequently, for every bounded
measurable numerical function \(f\) vanishing outside a set of finite
measure,
\[
  T_fT_ia
  =T_fJT_i^*Ja
  =JT_{\overline f}T_i^*Ja
  =J(T_fT_i)^*Ja.
\]
The linear combinations of the operators \(T_fT_i\) form an
involutive subalgebra of \(\mathcal U(\mathfrak A)\) which, as was
observed in proving Proposition~3, generates
\(\mathcal U(\mathfrak A)\). Passing to the limit, we therefore have
\[
  Ta=JT^*Ja
  \qquad\bigl(T\in\mathcal U(\mathfrak A)\bigr),
\]
and consequently \(a\in\mathfrak C\). Conversely, if
\(a\in\mathfrak C\), then, for every \(i\),
\[
  U_{x_i(\zeta)}a(\zeta)
  =J(\zeta)U_{x_i(\zeta)}^*J(\zeta)a(\zeta)
\]
almost everywhere. Hence, almost everywhere, \(a(\zeta)\) is central
relative to an involutive subalgebra dense in
\(\mathfrak A(\zeta)\), so that
\(a(\zeta)\in\mathfrak C(\zeta)\).
\end{proof}

% Source PDF page 203 (printed p. 194)
\phantomsection\label{srcpage:194}
\begin{corollary*}
Let \(E\) be the characteristic projection of \(\mathfrak A\), and
let \(E(\zeta)\) be the characteristic projection of
\(\mathfrak A(\zeta)\). Then
\[
  E=\int^{\oplus}E(\zeta)\,d\nu(\zeta).
\]
\end{corollary*}

\begin{proof}
We have
\[
  E=E_{\mathfrak C}^{\mathcal U(\mathfrak A)},
  \qquad
  E(\zeta)=E_{\mathfrak C(\zeta)}^{\mathcal U(\mathfrak A(\zeta))}.
\]
It is therefore enough to apply Proposition~7 and \S~3, Lemma~3.
\end{proof}

Bibliography: \ArticleRef{29}, \ArticleRef{117}.

\NumberedParagraph{6}{Uniqueness and Existence of the Decomposition}
\label{no:II-4-6}

\begin{proposition}\label{prop:II-4-8}
Suppose that
\[
  \mathfrak A
  =\int^{\oplus}\mathfrak A(\zeta)\,d\nu(\zeta)
  =\int^{\oplus}\mathfrak A'(\zeta)\,d\nu(\zeta).
\]
If the \(\mathfrak A(\zeta)\) and the \(\mathfrak A'(\zeta)\) are
completed almost everywhere, then
\(\mathfrak A(\zeta)=\mathfrak A'(\zeta)\) almost everywhere.
\end{proposition}

\begin{proof}
Let \((x_i)\) be a fundamental sequence of measurable fields of
vectors such that:
\begin{enumerate}[label=\textup{\arabic*\textdegree},leftmargin=*]
\item \(x_i(\zeta)\in\mathfrak A(\zeta)\) for every \(\zeta\);
\item the functions
\(\zeta\mapsto\lVert x_i(\zeta)\rVert\) and
\(\zeta\mapsto\lVert U_{x_i(\zeta)}\rVert\) are bounded;
\item the set of these fields is stable under multiplication and
involution, the multiplication and involution being those defined by
the \(\mathfrak A(\zeta)\).
\end{enumerate}
Let \(Y\) be a measurable subset of \(Z\) of finite measure, and let
\(y_i\) be the field equal to \(x_i\) on \(Y\) and to \(0\) on
\(Z-Y\). We have \(y_i\in\mathfrak A\). Hence
\(y_i(\zeta)\in\mathfrak A'(\zeta)\) almost everywhere, and thus
\(x_i(\zeta)\in\mathfrak A'(\zeta)\) almost everywhere on \(Y\).
Moreover, almost everywhere on \(Y\), the products
\(x_i(\zeta)x_j(\zeta)\) and the adjoints \(x_i(\zeta)^*\), computed
in \(\mathfrak A(\zeta)\) and in \(\mathfrak A'(\zeta)\), are the
same. This proves that, almost everywhere, there exists a common
involutive subalgebra of \(\mathfrak A(\zeta)\) and
\(\mathfrak A'(\zeta)\), dense in \(\mathfrak H(\zeta)\). Since
\(\mathfrak A(\zeta)\) and \(\mathfrak A'(\zeta)\) are completed
almost everywhere, we have
\(\mathfrak A(\zeta)=\mathfrak A'(\zeta)\) almost everywhere
(Chapter~I, \S~5, no.~3, Remark).
\end{proof}

\begin{theorem}\label{thm:II-4-1}
For a completed Hilbert algebra \(\mathfrak A\) in
\[
  \mathfrak H
  =\int^{\oplus}\mathfrak H(\zeta)\,d\nu(\zeta)
\]
to be decomposable, it is necessary and sufficient that the von
Neumann algebra \(\mathcal U(\mathfrak A)\) be decomposable.
\end{theorem}

% Source PDF page 204 (printed p. 195)
\phantomsection\label{srcpage:195}
\begin{proof}
We already know that the condition is necessary (Proposition~3).
Conversely, suppose that \(\mathcal U(\mathfrak A)\) is decomposable.
Then the algebra \(\mathfrak Z\) of diagonalizable operators is
contained in \(\mathcal U(\mathfrak A)\), and
\(\mathcal U(\mathfrak A)\subset\mathfrak Z'\). Let \(J\) be the
involution of \(\mathfrak H\) associated with \(\mathfrak A\).

There exists a countable subset of \(\mathfrak H\) that is totalizing for
\(\mathfrak Z\). Since each element of this set is the limit of a sequence
of elements of \(\mathfrak A\), there exists a sequence of elements of
\(\mathfrak A\) that is totalizing for \(\mathfrak Z\). By suitably enlarging
this sequence, we finally obtain a sequence \((y_i)\) of distinct elements
of \(\mathfrak A\), totalizing for \(\mathfrak Z\), and stable under
addition, multiplication, involution, and multiplication by rational
complex numbers.

For every \(i\), choose a representative field
\(\zeta\mapsto y_i(\zeta)\) of \(y_i\). For every \(\zeta\in Z\), let
\(\mathfrak B(\zeta)\) be the set of the \(y_i(\zeta)\),
\(i=1,2,\ldots\). Almost everywhere,
\[
  r_1y_i(\zeta)+r_2y_j(\zeta)
  =(r_1y_i+r_2y_j)(\zeta)
\]
for all \(i,j\) and all rational complex numbers \(r_1,r_2\). We
shall endow the \(\mathfrak B(\zeta)\) with structures of algebras over
the field of rational complex numbers by setting
\[
  y_i(\zeta)y_j(\zeta)=(y_iy_j)(\zeta).
\]
This definition of multiplication is not legitimate if \(\zeta\)
belongs to one of the sets \(N_{ijkl}\) defined by the conditions
\[
  y_i(\zeta)=y_k(\zeta),
  \qquad
  y_j(\zeta)=y_l(\zeta),
  \qquad
  (y_iy_j)(\zeta)\neq(y_ky_l)(\zeta).
\]
But this set is null. Indeed, let \(Y\) be the measurable set of the
\(\zeta\in Z\) such that
\[
  y_i(\zeta)=y_k(\zeta),
  \qquad
  y_j(\zeta)=y_l(\zeta),
\]
and let \(E\) be the corresponding projection of \(\mathfrak Z\). We
have
\[
  Ey_i=Ey_k,
  \qquad
  Ey_j=Ey_l,
\]
and hence, taking account of the fact that \(E\in\mathfrak Z\),
\[
  E(y_iy_j)=EU_{y_i}y_j
  =U_{Ey_i}Ey_j
  =U_{Ey_k}Ey_l
  =E(y_ky_l).
\]
Thus \((y_iy_j)(\zeta)=(y_ky_l)(\zeta)\) almost everywhere on \(Y\),
which proves our assertion. Outside a null set, therefore, we have
indeed defined a multiplication on \(\mathfrak B(\zeta)\) such that
\(y_i(\zeta)y_j(\zeta)=(y_iy_j)(\zeta)\); and, after removing another
null set, the fact that the \(y_i\) form an algebra over the rational
complex field implies that the \(\mathfrak B(\zeta)\) are algebras
over this field. Similarly, using the fact that
\(J(Ey_i)=E(Jy_i)\) for every projection \(E\) of \(\mathfrak Z\), we
see that, on setting
\(y_i(\zeta)^*=y_{i'}(\zeta)\) whenever \(y_i^*=y_{i'}\), we
% Source PDF page 205 (printed p. 196)
\phantomsection\label{srcpage:196}
define on the \(\mathfrak B(\zeta)\) a structure of involutive algebra
over the rational complex field, with
\[
  \bigl(y_i(\zeta)\mid y_j(\zeta)\bigr)
  =\bigl(y_j(\zeta)^*\mid y_i(\zeta)^*\bigr)
\]
(after removing one more null set). Consider the complex vector
subspace \(\mathfrak A_1(\zeta)\) of \(\mathfrak H(\zeta)\) spanned
by \(\mathfrak B(\zeta)\). There exists on
\(\mathfrak A_1(\zeta)\) one and only one structure of involutive
algebra over \(\mathbb C\) inducing on \(\mathfrak B(\zeta)\) the
structure of involutive algebra just constructed. We shall show that
\(\mathfrak A_1(\zeta)\) is a Hilbert algebra dense in
\(\mathfrak H(\zeta)\). The equality
\[
  \bigl(y_i(\zeta)\mid y_j(\zeta)\bigr)
  =\bigl(y_j(\zeta)^*\mid y_i(\zeta)^*\bigr)
\]
for all \(i,j\) implies \((x\mid y)=(y^*\mid x^*)\) for all
\(x,y\in\mathfrak A_1(\zeta)\). Put
\[
  U_{y_i}=\int^{\oplus}U_{y_i}(\zeta)\,d\nu(\zeta).
\]
We have
\begin{align*}
  \int^{\oplus}y_i(\zeta)y_j(\zeta)\,d\nu(\zeta)
    &=\int^{\oplus}(y_iy_j)(\zeta)\,d\nu(\zeta)\\
    &=y_iy_j
      =U_{y_i}y_j
      =\int^{\oplus}U_{y_i}(\zeta)y_j(\zeta)\,d\nu(\zeta),
\end{align*}
and therefore, almost everywhere,
\[
  U_{y_i}(\zeta)y_j(\zeta)=y_i(\zeta)y_j(\zeta).
\]
Similarly, almost everywhere,
\[
  U_{y_i}(\zeta)^*y_j(\zeta)
  =U_{y_i^*}(\zeta)y_j(\zeta)
  =y_i(\zeta)^*y_j(\zeta).
\]
Hence, almost everywhere,
\[
  \bigl(y_i(\zeta)y_j(\zeta)\mid y_k(\zeta)\bigr)
  =\bigl(y_j(\zeta)\mid y_i(\zeta)^*y_k(\zeta)\bigr)
\]
for all \(i,j,k\). Thus \((xy\mid z)=(y\mid x^*z)\) for all
\(x,y,z\in\mathfrak A_1(\zeta)\), almost everywhere. At the same time
we see that, almost everywhere,
\[
  \lVert y_i(\zeta)y_j(\zeta)\rVert
  \leq\lVert U_{y_i}\rVert\,\lVert y_j(\zeta)\rVert,
\]
and hence the mapping \(y\mapsto xy\) on
\(\mathfrak A_1(\zeta)\) is continuous. Finally, the vectors
\(Ty_i\), with \(T\in\mathfrak Z\), form a total set in
\(\mathfrak H\) and belong to \(\mathfrak A\). Every element of
\(\mathfrak H\) can be approximated arbitrarily closely by products
of two elements of \(\mathfrak A\), and consequently by linear
combinations of elements of the form
\[
  (Ty_i)(T_1y_j)
  \qquad(T,T_1\in\mathfrak Z).
\]
But
\[
  (Ty_i)(T_1y_j)
  =U_{Ty_i}T_1y_j
  =TT_1U_{y_i}y_j
  =(TT_1)(y_iy_j),
\]
% Source PDF page 206 (printed p. 197)
\phantomsection\label{srcpage:197}
so the family \((y_iy_j)_{i,j\geq 1}\) is totalizing for
\(\mathfrak Z\). Hence,
almost everywhere, the vectors
\[
  y_i(\zeta)y_j(\zeta)
\]
form a total set in \(\mathfrak H(\zeta)\) (\S~1, Proposition~8).
We have therefore shown
that, except on a null set \(N\), \(\mathfrak A_1(\zeta)\) is a
Hilbert algebra dense in \(\mathfrak H(\zeta)\). On \(N\), redefine
\(\mathfrak A_1(\zeta)\) arbitrarily, subject only to the condition
that it be a Hilbert algebra dense in \(\mathfrak H(\zeta)\) (cf.
Chapter~I, \S~5, no.~5, Remark). By Proposition~1, the field
\(\zeta\mapsto\mathfrak A_1(\zeta)\) is measurable. Put
\[
  \mathfrak A_1
  =\int^{\oplus}\mathfrak A_1(\zeta)\,d\nu(\zeta).
\]
Since
\[
  U_{y_i(\zeta)}y_j(\zeta)=U_{y_i}(\zeta)y_j(\zeta)
\]
almost everywhere, it follows that
\(U_{y_i(\zeta)}=U_{y_i}(\zeta)\) almost everywhere. Let \(T_f\)
denote the diagonalizable operator corresponding to the function
\(f\in L_{\mathbb C}^{\infty}(Z,\nu)\). We have
\[
  (T_fy_i)(\zeta)=f(\zeta)y_i(\zeta),
  \qquad
  U_{(T_fy_i)(\zeta)}
  =f(\zeta)U_{y_i(\zeta)}
  =f(\zeta)U_{y_i}(\zeta).
\]
Thus \(\lVert U_{(T_fy_i)(\zeta)}\rVert\) is essentially bounded,
and \(T_fy_i\in\mathfrak A_1\). If
\(f_1\in L_{\mathbb C}^{\infty}(Z,\nu)\), the product
\((T_fy_i)(T_{f_1}y_j)\), computed in \(\mathfrak A_1\), is the field
\[
  \zeta\longmapsto
  f(\zeta)y_i(\zeta)f_1(\zeta)y_j(\zeta),
\]
that is,
\[
  (T_fT_{f_1})(y_iy_j)
  =(T_fy_i)(T_{f_1}y_j),
\]
computed in \(\mathfrak A\). Therefore there exists a common
involutive subalgebra of \(\mathfrak A\) and \(\mathfrak A_1\), dense
in \(\mathfrak H\). Since \(\mathfrak A\) is completed, it is the
completed Hilbert algebra corresponding to \(\mathfrak A_1\). Let
\(\mathfrak A(\zeta)\) be the completed Hilbert algebra corresponding
to \(\mathfrak A_1(\zeta)\). Then the field
\(\zeta\mapsto\mathfrak A(\zeta)\) is measurable, and
\[
  \mathfrak A
  =\int^{\oplus}\mathfrak A(\zeta)\,d\nu(\zeta)
\]
(Proposition~6).
\end{proof}

It is unlikely that Proposition~8 remains valid if the hypothesis that
\(\mathfrak A(\zeta)\) and \(\mathfrak A'(\zeta)\) are completed is
omitted.

Bibliography: \ArticleRef{13}, \ArticleRef{29}, \ArticleRef{86},
\ArticleRef{106}, \ArticleRef{116}, \ArticleRef{117}.

\par\medskip
\noindent\textit{Exercises.}---
\begin{enumerate}[label=\arabic*.,leftmargin=*,itemsep=0.5\baselineskip]
\item\label{ex:II-4-1} Let \(Z\) be a Borel space and \(\nu\) a positive measure on
\(Z\). For every \(\zeta\in Z\), let \(\mathfrak A(\zeta)\) be a
Hilbert algebra. Then
\[
  \mathfrak B=\prod_{\zeta\in Z}\mathfrak A(\zeta)
\]
is an involutive algebra. Suppose given an involutive subalgebra
\(\mathfrak C\) of \(\mathfrak B\) having the following properties:
\begin{enumerate}[label=\textup{(\roman*)},leftmargin=*]
\item for every \(x\in\mathfrak C\), the function
\(\zeta\mapsto\lVert x(\zeta)\rVert\) is measurable;
\item if \(y\in\mathfrak B\) is such that, for every
\(x\in\mathfrak C\), the function
\(\zeta\mapsto(x(\zeta)\mid y(\zeta))\) is measurable, then
\(y\in\mathfrak C\);
\item there exists a sequence \((x_i)\) of elements of
\(\mathfrak C\) such that, for every \(\zeta\in Z\), the
\(x_i(\zeta)\) form a total sequence in \(\mathfrak A(\zeta)\).
\end{enumerate}
Let \(\mathfrak H(\zeta)\) be the completed Hilbert space of
\(\mathfrak A(\zeta)\). Show that there exists on the field
\(\zeta\mapsto\mathfrak H(\zeta)\) one and only one structure of a
\(\nu\)-measurable field having the following property: for a field of
vectors
% Source PDF page 207 (printed p. 198)
\phantomsection\label{srcpage:198}
\(\zeta\mapsto x(\zeta)\) to belong to \(\mathfrak C\), it is
necessary and sufficient that it be measurable and that
\(x(\zeta)\in\mathfrak A(\zeta)\) for every \(\zeta\in Z\). With the
field \(\zeta\mapsto\mathfrak H(\zeta)\) endowed with this structure,
show that the field \(\zeta\mapsto\mathfrak A(\zeta)\) is measurable.

\item\label{ex:II-4-2} Let \(\zeta\mapsto\mathfrak H(\zeta)\) be a
\(\nu\)-measurable field of complex Hilbert spaces on \(Z\), and let
\(\zeta\mapsto\mathfrak A(\zeta)\subset\mathfrak H(\zeta)\) be a
\(\nu\)-measurable field of Hilbert algebras. If, for every
\(\zeta\in Z\), \(\mathfrak A(\zeta)\) has a unit element
\(1(\zeta)\), then the field \(\zeta\mapsto1(\zeta)\) is measurable.
[Let \((x_i)\) be a fundamental sequence of measurable fields of
vectors such that \(x_i(\zeta)\in\mathfrak A(\zeta)\) for every
\(i\) and every \(\zeta\). Then
\[
  \bigl(1(\zeta)\mid x_i(\zeta)x_j(\zeta)\bigr)
  =\bigl(x_i(\zeta)^*\mid x_j(\zeta)\bigr)
\]
depends measurably on \(\zeta\).]
\end{enumerate}

\section{Fields of Traces}\label{sec:II-5}

\NumberedParagraph{1}{Measurable Fields of Traces}
\label{no:II-5-1}
Let \(\zeta\mapsto\mathfrak H(\zeta)\) be a \(\nu\)-measurable
field of complex Hilbert spaces on \(Z\), and let
\[
  \mathfrak H
  =\int^{\oplus}\mathfrak H(\zeta)\,d\nu(\zeta),
  \qquad
  \mathcal A
  =\int^{\oplus}\mathcal A(\zeta)\,d\nu(\zeta)
\]
be a decomposable von Neumann algebra. For every \(\zeta\in Z\), let
\(\varphi_{\zeta}\) be a trace on \(\mathcal A(\zeta)_+\). The
mapping \(\zeta\mapsto\varphi_{\zeta}\) is called a \emph{field of
traces} on \(Z\).

\begin{definition}\label{def:II-5-1}
The field of traces \(\zeta\mapsto\varphi_{\zeta}\) is said to be
\emph{measurable} if, for every measurable field of operators
\(\zeta\mapsto T(\zeta)\) such that
\(T(\zeta)\in\mathcal A(\zeta)_+\) for every \(\zeta\), the function
\(\zeta\mapsto\varphi_{\zeta}(T(\zeta))\) is measurable.
\end{definition}

Now suppose that the field \(\zeta\mapsto T(\zeta)\) is essentially
bounded, and put
\[
  T=\int^{\oplus}T(\zeta)\,d\nu(\zeta)\in\mathcal A_+.
\]
The number
\(\int\varphi_{\zeta}(T(\zeta))\,d\nu(\zeta)\) depends only on
\(T\); denote it by \(\varphi(T)\). Plainly, \(\varphi\) is a trace
on \(\mathcal A_+\), denoted by
\(\int^{\oplus}\varphi_{\zeta}\,d\nu(\zeta)\).

\begin{proposition}\label{prop:II-5-1}
\begin{enumerate}[label=\textup{(\roman*)},leftmargin=*]
\item The trace \(\varphi\) is finite if and only if
\[
  \int\varphi_{\zeta}(I)\,d\nu(\zeta)<+\infty
\]
(which implies that \(\varphi_{\zeta}\) is finite almost everywhere).
\item If the \(\varphi_{\zeta}\) are normal almost everywhere, then
\(\varphi\) is normal. If, in addition, \(\varphi\) is semifinite,
then the \(\varphi_{\zeta}\) are semifinite almost everywhere.
\item If the \(\varphi_{\zeta}\) are faithful almost everywhere,
then \(\varphi\) is faithful.
\end{enumerate}
\end{proposition}

\begin{proof}
For \(\varphi\) to be finite, it is necessary and sufficient that
\(\varphi(I)<+\infty\). But
\[
  \varphi(I)=\int\varphi_{\zeta}(I)\,d\nu(\zeta),
\]
which proves~\textup{(i)}.
% Source PDF page 208 (printed p. 199)
\phantomsection\label{srcpage:199}
Suppose that the \(\varphi_{\zeta}\) are faithful almost everywhere.
Then, if
\[
  T=\int^{\oplus}T(\zeta)\,d\nu(\zeta)\in\mathcal A_+,
  \qquad \varphi(T)=0,
\]
we have \(\varphi_{\zeta}(T(\zeta))=0\) almost everywhere, hence
\(T(\zeta)=0\) almost everywhere, and therefore \(T=0\). This
proves~\textup{(iii)}.

Suppose that the \(\varphi_{\zeta}\) are normal almost everywhere. Let
\((T_{\iota})\) be an increasing directed family in \(\mathcal A_+\)
with supremum \(T\in\mathcal A_+\). Since \(\mathcal A\) is of
countable type [\S~2, Proposition~7\textup{(iii)}], one can extract
from the family \((T_{\iota})\) an increasing sequence
\(T_1,T_2,\ldots\) with supremum \(T\) (Chapter~I, \S~3, corollary
to Proposition~1). Put
\[
  T_i=\int^{\oplus}T_i(\zeta)\,d\nu(\zeta),
  \qquad
  T=\int^{\oplus}T(\zeta)\,d\nu(\zeta).
\]
There is an increasing sequence of integers \((n_1,n_2,\ldots)\)
such that \(T_{n_k}(\zeta)\) converges strongly and increasingly to
\(T(\zeta)\) almost everywhere (\S~2, Proposition~4). Hence, almost
everywhere,
\(\varphi_{\zeta}(T_{n_k}(\zeta))\) increases to
\(\varphi_{\zeta}(T(\zeta))\). Consequently
\[
  \varphi(T_{n_k})
  =\int\varphi_{\zeta}(T_{n_k}(\zeta))\,d\nu(\zeta)
\]
increases to
\[
  \varphi(T)
  =\int\varphi_{\zeta}(T(\zeta))\,d\nu(\zeta).
\]

If, in addition, \(\varphi\) is semifinite, let us show that
\(\varphi_{\zeta}\) is semifinite almost everywhere. Since
\(\mathcal A\) is of countable type, \(I\) is the supremum of an
increasing sequence \((S_i)\) of operators of \(\mathcal A_+\) such
that \(\varphi(S_i)<+\infty\). Put
\[
  S_i=\int^{\oplus}S_i(\zeta)\,d\nu(\zeta).
\]
We have \(\varphi_{\zeta}(S_i(\zeta))<+\infty\) almost everywhere,
and, after passing to a subsequence, \(S_i(\zeta)\) converges strongly
and increasingly to \(I\), almost everywhere. Thus
\(\varphi_{\zeta}\) is semifinite almost everywhere.
\end{proof}

\noindent\textit{Problems.} If \(\varphi\) is normal, are the
\(\varphi_{\zeta}\) normal? If \(\varphi\) is faithful, are the
\(\varphi_{\zeta}\) faithful? If the \(\varphi_{\zeta}\) are
semifinite, is \(\varphi\) semifinite?

Bibliography: \ArticleRef{80}, \ArticleRef{117}.

\NumberedParagraph{2}{Decomposition of Traces}
\label{no:II-5-2}

\begin{theorem}\label{thm:II-5-1}
Let
\[
  \mathfrak A
  =\int^{\oplus}\mathfrak A(\zeta)\,d\nu(\zeta)
\]
be a decomposable Hilbert algebra in \(\mathfrak H\), let
\(\varphi\) be the natural trace on
\(\mathcal U(\mathfrak A)_+\), and let
\(\varphi_{\zeta}\) be the natural trace on
\(\mathcal U(\mathfrak A(\zeta))_+\). Then the field of traces
\(\zeta\mapsto\varphi_{\zeta}\) is measurable, and
\[
  \varphi=\int^{\oplus}\varphi_{\zeta}\,d\nu(\zeta).
\]
\end{theorem}

% Source PDF page 209 (printed p. 200)
\phantomsection\label{srcpage:200}
\begin{proof}
Let
\[
  a=\int^{\oplus}a(\zeta)\,d\nu(\zeta)
\]
be an element of \(\mathfrak H\) which is bounded for
\(\mathfrak A\). We have
\[
  \varphi_{\zeta}
    \bigl(U_{a(\zeta)}^*U_{a(\zeta)}\bigr)
  =\bigl(a(\zeta)\mid a(\zeta)\bigr)
\]
and
\begin{align*}
  \varphi(U_a^*U_a)
    &=(a\mid a)
      =\int\bigl(a(\zeta)\mid a(\zeta)\bigr)\,d\nu(\zeta)\\
    &=\int\varphi_{\zeta}
      \bigl(U_{a(\zeta)}^*U_{a(\zeta)}\bigr)\,d\nu(\zeta).
\end{align*}

Now let \(T\) be an arbitrary element of
\(\mathcal U(\mathfrak A)_+\), and put
\[
  T=\int^{\oplus}T(\zeta)\,d\nu(\zeta),
  \qquad
  T(\zeta)\in\mathcal U(\mathfrak A(\zeta))_+
  \quad(\zeta\in Z).
\]
The algebra \(\mathcal U(\mathfrak A)\) is of countable type. By
Chapter~I, \S~3, the corollary to Proposition~1 and Corollary~5 to
Theorem~2, there is an increasing sequence \((T_i)\) of operators of
\(\mathcal U(\mathfrak A)_+\), with supremum \(T\), such that
\(\varphi(T_i)<+\infty\) for every \(i\). Put
\[
  T_i=\int^{\oplus}T_i(\zeta)\,d\nu(\zeta).
\]
We already know that the functions
\(\zeta\mapsto\varphi_{\zeta}(T_i(\zeta))\) are integrable and that
\[
  \varphi(T_i)
  =\int\varphi_{\zeta}(T_i(\zeta))\,d\nu(\zeta).
\]
On the other hand, after passing to a subsequence, the functions
\(\zeta\mapsto\varphi_{\zeta}(T_i(\zeta))\) form, almost
everywhere, an increasing sequence converging to the function
\(\zeta\mapsto\varphi_{\zeta}(T(\zeta))\). Hence this function is
measurable, and
\[
  \int\varphi_{\zeta}(T(\zeta))\,d\nu(\zeta)
  =\lim_i\varphi(T_i)=\varphi(T).
\]
\end{proof}

\begin{theorem}\label{thm:II-5-2}
Let
\[
  \mathcal A
  =\int^{\oplus}\mathcal A(\zeta)\,d\nu(\zeta)
\]
be a decomposable semifinite von Neumann algebra. Suppose that
\(\nu\) is standard.
\begin{enumerate}[label=\textup{(\roman*)},leftmargin=*]
\item The \(\mathcal A(\zeta)\) are semifinite almost everywhere.
\item The algebra \(\mathcal A\) is finite if and only if
\(\mathcal A(\zeta)\) is finite almost everywhere.
\item Let \(\varphi\) be a faithful semifinite normal trace on
\(\mathcal A_+\). Then there is a measurable field
\(\zeta\mapsto\varphi_{\zeta}\) of faithful semifinite normal traces
on the \(\mathcal A(\zeta)_+\) such that
\[
  \varphi=\int^{\oplus}\varphi_{\zeta}\,d\nu(\zeta).
\]
\end{enumerate}
\end{theorem}

% Source PDF page 210 (printed p. 201)
\phantomsection\label{srcpage:201}
\begin{proof}
There exist a Hilbert space \(\mathfrak K\), a standard von Neumann
algebra \(\mathcal B\) on \(\mathfrak K\), and an isomorphism
\(\Phi\) of \(\mathcal A\) onto \(\mathcal B\). Since
\(\mathcal A\) is of countable type, \(\mathcal B\), and hence
\(\mathcal B'\), are of countable type. Consequently (\S~3,
Proposition~11) there exist a \(\nu\)-measurable field
\(\zeta\mapsto\mathfrak K(\zeta)\) of Hilbert spaces on \(Z\), a
\(\nu\)-measurable field
\[
  \zeta\longmapsto
  \mathcal B(\zeta)\subset\mathcal L(\mathfrak K(\zeta))
\]
of von Neumann algebras, and a \(\nu\)-measurable field
\(\zeta\mapsto\Phi_{\zeta}\) of isomorphisms of
\(\mathcal A(\zeta)\) onto \(\mathcal B(\zeta)\), such that we may
identify
\[
  \mathfrak K\quad\text{with}\quad
  \int^{\oplus}\mathfrak K(\zeta)\,d\nu(\zeta),
  \qquad
  \mathcal B\quad\text{with}\quad
  \int^{\oplus}\mathcal B(\zeta)\,d\nu(\zeta),
\]
and \(\Phi\) with
\(\int^{\oplus}\Phi_{\zeta}\,d\nu(\zeta)\). Let \(\psi\) be the
normal trace on \(\mathcal B_+\) transported from \(\varphi\) by
\(\Phi\). There exists (Chapter~I, \S~6, Lemma~1) a Hilbert algebra
\(\mathfrak A\), dense in \(\mathfrak K\), such that
\(\mathcal B=\mathcal U(\mathfrak A)\) and \(\psi\) is the natural
trace on \(\mathcal U(\mathfrak A)_+\). Since
\(\mathcal U(\mathfrak A)=\mathcal B\) is decomposable,
\(\mathfrak A\) is decomposable (\S~4, Theorem~1). Let
\[
  \mathfrak A
  =\int^{\oplus}\mathfrak A(\zeta)\,d\nu(\zeta).
\]
We have
\[
  \mathcal B=\mathcal U(\mathfrak A)
  =\int^{\oplus}\mathcal U(\mathfrak A(\zeta))\,d\nu(\zeta)
  \qquad(\S~4,\text{ Proposition~3}),
\]
and hence
\(\mathcal B(\zeta)=\mathcal U(\mathfrak A(\zeta))\) almost
everywhere. This already proves~\textup{(i)}. Moreover, for
\(\mathcal A\), or equivalently \(\mathcal B\), to be finite, it is
necessary and sufficient that the characteristic projection of
\(\mathfrak A\) be \(I\) (Chapter~I, \S~6, Theorem~6), hence that
the characteristic projection of \(\mathfrak A(\zeta)\) be \(I\)
almost everywhere (\S~4, corollary to Proposition~7), and therefore
that \(\mathcal B(\zeta)\), or equivalently
\(\mathcal A(\zeta)\), be finite almost everywhere. This
proves~\textup{(ii)}.

Let \(\psi_{\zeta}\) be the natural trace on
\(\mathcal U(\mathfrak A(\zeta))_+\). We have
\[
  \psi=\int^{\oplus}\psi_{\zeta}\,d\nu(\zeta)
  \qquad\text{(Theorem~1)}.
\]
Finally, let \(\varphi_{\zeta}\) be the normal trace on
\(\mathcal A(\zeta)_+\) transported from \(\psi_{\zeta}\) by
\(\Phi_{\zeta}^{-1}\). For every measurable field of operators
\[
  \zeta\longmapsto T(\zeta)\in\mathcal A(\zeta)_+,
\]
the number
\[
  \varphi_{\zeta}(T(\zeta))
  =\psi_{\zeta}\bigl(\Phi_{\zeta}(T(\zeta))\bigr)
\]
depends measurably on \(\zeta\), so the field
\(\zeta\mapsto\varphi_{\zeta}\) is measurable. Furthermore, if
\[
  T=\int^{\oplus}T(\zeta)\,d\nu(\zeta)\in\mathcal A_+,
\]
then
\[
  \Phi(T)
  =\int^{\oplus}\Phi_{\zeta}(T(\zeta))\,d\nu(\zeta),
\]
and hence
\[
  \varphi(T)
  =\psi(\Phi(T))
  =\int\psi_{\zeta}\bigl(\Phi_{\zeta}(T(\zeta))\bigr)
     \,d\nu(\zeta)
  =\int\varphi_{\zeta}(T(\zeta))\,d\nu(\zeta);
\]%
% Source PDF page 211 (printed p. 202)
\phantomsection\label{srcpage:202}%
thus
\[
  \varphi=\int^{\oplus}\varphi_{\zeta}\,d\nu(\zeta).
\]
\end{proof}

\begin{corollary*}\phantomsection\label{cor:II-5-2}
Suppose that \(\nu\) is standard. Let
\[
  \mathcal A
  =\int^{\oplus}\mathcal A(\zeta)\,d\nu(\zeta)
\]
be a decomposable von Neumann algebra, and let \(\varphi\) be a
normal (respectively, faithful normal) trace on \(\mathcal A_+\).
There exists a measurable field
\(\zeta\mapsto\varphi_{\zeta}\) of normal (respectively, faithful
normal) traces on the \(\mathcal A(\zeta)_+\) such that
\[
  \varphi=\int^{\oplus}\varphi_{\zeta}\,d\nu(\zeta).
\]
\end{corollary*}

\begin{proof}
There are projections \(E_1,E_2,E_3\) in the centre of
\(\mathcal A\), with sum \(I\), having the following properties:
\textup{(1)} the trace \(\varphi^1\) induced by \(\varphi\) on
\(\mathcal A_{E_1,+}\) is faithful and semifinite;
\textup{(2)} the trace \(\varphi^2\) induced by \(\varphi\) on
\(\mathcal A_{E_2,+}\) is infinite on every nonzero operator;
\textup{(3)} the trace \(\varphi^3\) induced by \(\varphi\) on
\(\mathcal A_{E_3,+}\) is zero. Let
\[
  E_j=\int^{\oplus}E_j(\zeta)\,d\nu(\zeta)
  \qquad(j=1,2,3).
\]
Almost everywhere, \(E_1(\zeta),E_2(\zeta),E_3(\zeta)\) are
pairwise orthogonal projections in the centre of
\(\mathcal A(\zeta)\), with sum \(I\). We have
\[
  \mathcal A_{E_j}
  =\int^{\oplus}\mathcal A(\zeta)_{E_j(\zeta)}\,d\nu(\zeta)
  \qquad(j=1,2,3)
  \quad(\S~3,\text{ Proposition~6}).
\]
Let \(\zeta\mapsto\varphi_{\zeta}^1\) be a measurable field of
traces on \(\mathcal A(\zeta)_{E_1(\zeta),+}\). Suppose that these
traces are faithful, normal, and semifinite, and that
\[
  \varphi^1
  =\int^{\oplus}\varphi_{\zeta}^1\,d\nu(\zeta)
  \qquad\text{(Theorem~2)}.
\]
For every \(\zeta\), let \(\varphi_{\zeta}^2\) be the trace on
\(\mathcal A(\zeta)_{E_2(\zeta),+}\) which is infinite on every
nonzero operator, and let \(\varphi_{\zeta}\) be the trace on
\(\mathcal A(\zeta)_+\) which induces
\(\varphi_{\zeta}^1\) on
\(\mathcal A(\zeta)_{E_1(\zeta),+}\),
\(\varphi_{\zeta}^2\) on
\(\mathcal A(\zeta)_{E_2(\zeta),+}\), and zero on
\(\mathcal A(\zeta)_{E_3(\zeta),+}\). For every measurable field
\(\zeta\mapsto T(\zeta)\in\mathcal A(\zeta)_+\),
\[
  \varphi_{\zeta}(T(\zeta))
  =\varphi_{\zeta}^1(T(\zeta)E_1(\zeta))
   +\varphi_{\zeta}^2(T(\zeta)E_2(\zeta))
\]
depends measurably on \(\zeta\), and therefore the field
\(\zeta\mapsto\varphi_{\zeta}\) is measurable. Furthermore, if
\[
  T=\int^{\oplus}T(\zeta)\,d\nu(\zeta)\in\mathcal A_+,
\]
then%
% Source PDF page 212 (printed p. 203)
\phantomsection\label{srcpage:203}%
\begin{align*}
  \varphi(T)
    &=\varphi^1(TE_1)+\varphi^2(TE_2)\\
    &=\int\varphi_{\zeta}^1(T(\zeta)E_1(\zeta))\,d\nu(\zeta)
      +\int\varphi_{\zeta}^2(T(\zeta)E_2(\zeta))\,d\nu(\zeta)\\
    &=\int\varphi_{\zeta}(T(\zeta))\,d\nu(\zeta),
\end{align*}
and hence
\[
  \varphi=\int^{\oplus}\varphi_{\zeta}\,d\nu(\zeta).
\]
Finally, if \(\varphi\) is faithful, then \(E_3=0\), and hence the
\(\varphi_{\zeta}\) are faithful.
\end{proof}

Bibliography: \ArticleRef{10}, \ArticleRef{80}, \ArticleRef{117},
\ArticleRef{123}.

\NumberedParagraph{3}{Uniqueness of the Decomposition}
\label{no:II-5-3}
The following lemma is a lemma of global theory.

\begin{lemma}\label{lem:II-5-1}
Let \(\mathcal A\) be a von Neumann algebra, let \(\omega\) and
\(\omega'\) be two normal traces on \(\mathcal A_+\), and let
\(\mathcal B\) be an involutive subalgebra of \(\mathcal A\),
strongly dense in \(\mathcal A\). Suppose that
\[
  \omega(T^*T)=\omega'(T^*T)<+\infty
  \qquad\text{for every }T\in\mathcal B.
\]
Then \(\omega=\omega'\).
\end{lemma}

\begin{proof}
The set of \(T\in\mathcal A_+\) such that \(\omega(T)<+\infty\)
is the positive part of a two-sided ideal \(\mathfrak m\) of
\(\mathcal A\). Let \(\dot\omega\) be the linear form on
\(\mathfrak m\) which coincides with \(\omega\) on
\(\mathfrak m_+\). If \(S\in\mathfrak m\), the linear form
\[
  T\longmapsto\dot\omega(ST)=\dot\omega(TS)
  \qquad(T\in\mathcal A)
\]
is ultraweakly continuous (Chapter~I, \S~6, Proposition~1). Define
\(\mathfrak m'\) and \(\dot\omega'\) similarly. Let
\(T_1,T_2\in\mathcal B\). Then \(T_1T_2\) is a linear combination
of elements of the form \(R^*R\), where \(R\in\mathcal B\). Hence
\(T_1T_2\in\mathfrak m\cap\mathfrak m'\) and
\(\dot\omega(T_1T_2)=\dot\omega'(T_1T_2)\). It follows that the
linear forms
\[
  T\longmapsto\dot\omega(T_1T_2T),
  \qquad
  T\longmapsto\dot\omega'(T_1T_2T)
  \qquad(T\in\mathcal A)
\]
coincide on \(\mathcal B\), and therefore on \(\mathcal A\). Let
\(S\in\mathfrak m\cap\mathfrak m'\). From what has just been shown,
the linear forms
\[
  T\longmapsto\dot\omega(TS),
  \qquad
  T\longmapsto\dot\omega'(TS)
  \qquad(T\in\mathcal A)
\]
coincide on \(\mathcal B^2\), hence on \(\mathcal B\), and
therefore on \(\mathcal A\); in particular,
\(\dot\omega(S)=\dot\omega'(S)\).

We have \(\mathfrak m\cap\mathfrak m'\supset\mathcal B^2\), and
hence \(\mathfrak m\cap\mathfrak m'\) is strongly dense in
\(\mathcal A\). If \(R\in\mathcal A_+\), there is an increasing
directed family \((R_{\iota})\) of elements of
\((\mathfrak m\cap\mathfrak m')_+\) with supremum \(R\)
(Chapter~I, \S~3, Corollary~5 to Theorem~2). Thus
\[
  \omega(R)=\sup_{\iota}\omega(R_{\iota})
  =\sup_{\iota}\omega'(R_{\iota})=\omega'(R).
\]
\end{proof}

\begin{theorem}\label{thm:II-5-3}
Let
\[
  \mathcal A
  =\int^{\oplus}\mathcal A(\zeta)\,d\nu(\zeta)
\]
be a decomposable von Neumann algebra. Let
\(\zeta\mapsto\varphi_{\zeta}\) and
\(\zeta\mapsto\psi_{\zeta}\) be two measurable fields of normal
traces on the \(\mathcal A(\zeta)_+\). If
\(\int^{\oplus}\varphi_{\zeta}\,d\nu(\zeta)\) and
\(\int^{\oplus}\psi_{\zeta}\,d\nu(\zeta)\) are identical with the
same semifinite trace \(\varphi\), then
\(\varphi_{\zeta}=\psi_{\zeta}\) almost everywhere.
\end{theorem}

% Source PDF page 213 (printed p. 204)
\phantomsection\label{srcpage:204}
\begin{proof}
By \S~2, Proposition~7\textup{(iii)}, \(\mathcal A\) is of
countable type. Let \(\mathfrak m\) be the ideal of the
\(T\in\mathcal A\) such that \(\varphi(T^*T)<+\infty\). Since
\(\varphi\) is semifinite, \(\mathfrak m\) is strongly dense in
\(\mathcal A\). Let \((T_j)\) be a sequence of elements of
\(\mathfrak m\) such that \(\mathcal A\) is the von Neumann algebra
generated by \(\mathfrak Z\), the algebra of diagonalizable
operators, and the \(T_j\). Since the unit ball of \(\mathfrak m\)
is strongly dense in that of \(\mathcal A\), and the latter is
metrizable for the strong topology, every \(T_j\) is the strong
limit of a sequence of elements of \(\mathfrak m\). We may therefore
suppose that \(T_j\in\mathfrak m\) for every \(j\), that \(I\) is
the strong limit of a subsequence of \((T_j)\), and, moreover, that
the \(T_j\) form an involutive algebra over the rational complex
field. Put
\[
  T_j=\int^{\oplus}T_j(\zeta)\,d\nu(\zeta).
\]
For every measurable subset \(Y\) of \(Z\), denoting by \(E_Y\)
the corresponding projection of \(\mathfrak Z\), we have
\begin{align*}
  \int_Y\varphi_{\zeta}
    \bigl(T_j(\zeta)^*T_j(\zeta)\bigr)\,d\nu(\zeta)
    &=\varphi(T_j^*T_jE_Y)\\
    &=\int_Y\psi_{\zeta}
      \bigl(T_j(\zeta)^*T_j(\zeta)\bigr)\,d\nu(\zeta)
      <+\infty.
\end{align*}
Thus
\[
  \varphi_{\zeta}\bigl(T_j(\zeta)^*T_j(\zeta)\bigr)
  =\psi_{\zeta}\bigl(T_j(\zeta)^*T_j(\zeta)\bigr)<+\infty
\]
for every \(j\), except on a null set \(N\). Enlarging \(N\), we
may suppose that, for every \(\zeta\in Z\setminus N\), the
\(T_j(\zeta)\) form an involutive algebra over the rational complex
field, strongly dense in \(\mathcal A(\zeta)\). For
\(\zeta\in Z\setminus N\), we have
\[
  \varphi_{\zeta}\bigl(T_i(\zeta)^*T_j(\zeta)\bigr)
  =\psi_{\zeta}\bigl(T_i(\zeta)^*T_j(\zeta)\bigr)
\]
for all \(i,j\), and therefore
\(\varphi_{\zeta}(T^*T)=\psi_{\zeta}(T^*T)\) for every \(T\) in
the involutive algebra over \(\mathbb C\) generated by the
\(T_i(\zeta)\). Hence, by Lemma~1,
\(\varphi_{\zeta}=\psi_{\zeta}\).
\end{proof}

Observe that Lemma~1 is unnecessary if the
\(\mathcal A(\zeta)\) are factors.

Theorem~3 is visibly incomplete. Here again, the problems remain
unsolved even when \(\nu\) is standard.

Bibliography: \ArticleRef{80}, \ArticleRef{306}.

\NumberedParagraph{4}{Reduction of Semifinite, Finite, Purely Infinite,
and Properly Infinite von Neumann Algebras}
\label{no:II-5-4}
Exceptionally, in this number we shall use certain results of
Chapter~III.

\begin{theorem}\label{thm:II-5-4}
Let
\[
  \mathcal A
  =\int^{\oplus}\mathcal A(\zeta)\,d\nu(\zeta)
\]
be a decomposable von Neumann algebra. Suppose that \(\nu\) is
standard. If \(\mathcal A\) is purely infinite, then
\(\mathcal A(\zeta)\) is purely infinite almost everywhere.
\end{theorem}

% Source PDF page 214 (printed p. 205)
\phantomsection\label{srcpage:205}
\begin{proof}
Put
\[
  \mathfrak H
  =\int^{\oplus}\mathfrak H(\zeta)\,d\nu(\zeta).
\]
Let \(Y\) be the set of \(\zeta\in Z\) for which
\(\mathcal A(\zeta)\) is not purely infinite. We shall show that
there exist a projection
\[
  G=\int^{\oplus}G(\zeta)\,d\nu(\zeta)
\]
of \(\mathcal A\), and
\[
  x=\int^{\oplus}x(\zeta)\,d\nu(\zeta)\in G(\mathfrak H),
\]
such that \(x(\zeta)\ne0\) for \(\zeta\in Y\) and \(x\) is a trace
element for \(\mathcal A_G\). It will follow that \(x=0\), and hence
that \(Y\) is null.

To prove this assertion, the usual technique reduces us to the case
where: \textup{(a)} \(\zeta\mapsto\mathfrak H(\zeta)\) is the
constant field corresponding to a Hilbert space \(\mathfrak H_0\);
\textup{(b)} \(Z\) is compact metrizable and there exist continuous
maps
\[
  \zeta\mapsto T_1(\zeta),\quad
  \zeta\mapsto T_2(\zeta),\quad\ldots,
  \qquad
  \zeta\mapsto T'_1(\zeta),\quad
  \zeta\mapsto T'_2(\zeta),\quad\ldots
\]
from \(Z\) into the unit ball \(\mathcal L_1\) of
\(\mathcal L(\mathfrak H_0)\), furnished with the strong topology,
such that, for \(\zeta\in Z\), the \(T_i(\zeta)\) generate
\(\mathcal A(\zeta)\) and the \(T'_i(\zeta)\) generate
\(\mathcal A(\zeta)'\). We may also suppose that, for every
\(i=1,2,\ldots\), there is a \(j=1,2,\ldots\) such that
\[
  T_i(\zeta)^*=T_j(\zeta),
  \qquad
  T'_i(\zeta)^*=T'_j(\zeta)
  \qquad(\zeta\in Z).
\]
Let \(\mathfrak M\) be the set of
\((\zeta,T,y)\in Z\times\mathcal L_1\times\mathfrak H_0\) such
that
\begin{enumerate}[label=\textup{\arabic*\textsuperscript{o}},leftmargin=*]
\item \(TT'_i(\zeta)=T'_i(\zeta)T\) for \(i=1,2,\ldots\);
\item \(T\) is a projection;
\item \(Ty=y\);
\item
\[
  \bigl((TT_i(\zeta)T)(TT_j(\zeta)T)y\mid y\bigr)
  =\bigl((TT_j(\zeta)T)(TT_i(\zeta)T)y\mid y\bigr)
\]
for \(i,j=1,2,\ldots\);
\item \(\lVert y\rVert=1\).
\end{enumerate}
Then \(\mathfrak M\) is closed (\(\mathfrak H_0\) being furnished
with the strong topology). For \(\zeta\in Y\), the set of
\((T,y)\) such that \((\zeta,T,y)\in\mathfrak M\) is nonempty
(Chapter~III, \S~2, Corollary~1 to Proposition~7). Hence
(Appendix~V) there exist measurable maps
\(\zeta\mapsto G(\zeta)\) and \(\zeta\mapsto x(\zeta)\), defined
on a measurable subset \(X\) of \(Z\) containing \(Y\), such that
\((\zeta,G(\zeta),x(\zeta))\in\mathfrak M\) for every
\(\zeta\in X\). Put \(G(\zeta)=0\), \(x(\zeta)=0\) for
\(\zeta\in Z\setminus X\). The \(G(\zeta)\) are projections
(condition~\textup{2}), belong to \(\mathcal A(\zeta)\)
(condition~\textup{1}), and \(x(\zeta)\ne0\) for \(\zeta\in Y\)
(condition~\textup{5}). Put
\[
  G=\int^{\oplus}G(\zeta)\,d\nu(\zeta),
  \qquad
  x=\int^{\oplus}x(\zeta)\,d\nu(\zeta).
\]
Then \(G\) is a projection of \(\mathcal A\), and
\(x\in G(\mathfrak H)\) (condition~\textup{3}). Let \(\omega\) be
the normal positive form \(S\mapsto(Sx\mid x)\) on
\(G\mathcal A G\). By condition~\textup{4},
\(\omega(S_1S_2)=\omega(S_2S_1)\) when \(S_1,S_2\) belong to the
involutive algebra generated by \(G\mathfrak ZG\) and the
\[
  G\left(\int^{\oplus}T_i(\zeta)\,d\nu(\zeta)\right)G.
\]%
% Source PDF page 215 (printed p. 206)
\phantomsection\label{srcpage:206}%
Since this algebra is weakly dense in \(G\mathcal A G\)
(\S~3, Theorem~1), it follows
that \(x\) is a trace element for \(\mathcal A_G\).
\end{proof}

\setcounter{corollary}{0}
\begin{corollary}\label{cor:II-5-4-1}
Suppose that \(\nu\) is standard. Let \(E\) (respectively, \(F\))
be the largest projection in the centre of \(\mathcal A\) such that
\(\mathcal A_E\) (respectively, \(\mathcal A_F\)) is semifinite
(respectively, purely infinite). Let
\[
  E=\int^{\oplus}E(\zeta)\,d\nu(\zeta),
  \qquad
  F=\int^{\oplus}F(\zeta)\,d\nu(\zeta).
\]
Then, almost everywhere, \(E(\zeta)\) (respectively,
\(F(\zeta)\)) is the largest projection in the centre of
\(\mathcal A(\zeta)\) such that
\(\mathcal A(\zeta)_{E(\zeta)}\) (respectively,
\(\mathcal A(\zeta)_{F(\zeta)}\)) is semifinite (respectively,
purely infinite).
\end{corollary}

\begin{proof}
Almost everywhere, \(E(\zeta)\) and \(F(\zeta)\) are orthogonal
projections, with sum \(I\), in the centre of
\(\mathcal A(\zeta)\) (\S~3, Theorem~4);
\(\mathcal A(\zeta)_{E(\zeta)}\) is semifinite
[\S~3, Proposition~6, and Theorem~2\textup{(i)}], and
\(\mathcal A(\zeta)_{F(\zeta)}\) is purely infinite
(\S~3, Proposition~6, and Theorem~4).
\end{proof}

\begin{corollary}\label{cor:II-5-4-2}
Suppose that \(\nu\) is standard. In order that \(\mathcal A\) be
semifinite (respectively, purely infinite), it is necessary and
sufficient that \(\mathcal A(\zeta)\) be semifinite (respectively,
purely infinite) almost everywhere.
\end{corollary}

\begin{proof}
This follows from Corollary~1.
\end{proof}

\begin{theorem}\label{thm:II-5-5}
Let
\[
  \mathcal A
  =\int^{\oplus}\mathcal A(\zeta)\,d\nu(\zeta)
\]
be a decomposable von Neumann algebra. If \(\mathcal A\) is properly
infinite, then \(\mathcal A(\zeta)\) is properly infinite almost
everywhere.
\end{theorem}

\begin{proof}
There are orthogonal projections \(E,F\in\mathcal A\) such that
\[
  E\sim F\sim E+F=I
\]
(Chapter~III, \S~8, Corollary~2 to Theorem~1). Let
\[
  E=\int^{\oplus}E(\zeta)\,d\nu(\zeta),
  \qquad
  F=\int^{\oplus}F(\zeta)\,d\nu(\zeta).
\]
We have
\[
  E(\zeta)\sim F(\zeta)\sim E(\zeta)+F(\zeta)=I
\]
except on a null set \(N\). Let \(\zeta\in Z\setminus N\). For
every nonzero projection \(G\) in the centre of
\(\mathcal A(\zeta)\), we have
\(GE(\zeta)\sim GF(\zeta)\sim G\); hence \(G\) is infinite
(Chapter~III, \S~2, Proposition~4). Thus \(\mathcal A(\zeta)\) is
properly infinite (Chapter~III, \S~2, Proposition~9).
\end{proof}

\setcounter{corollary}{0}
\begin{corollary}\label{cor:II-5-5-1}
Suppose that \(\nu\) is standard. Let \(E\) (respectively, \(F\))
be the largest projection in the centre of \(\mathcal A\) such that
\(\mathcal A_E\) (respectively, \(\mathcal A_F\)) is finite
(respectively, properly infinite). Let
\[
  E=\int^{\oplus}E(\zeta)\,d\nu(\zeta),
  \qquad
  F=\int^{\oplus}F(\zeta)\,d\nu(\zeta).
\]
Then, almost everywhere, \(E(\zeta)\) (respectively,
\(F(\zeta)\)) is the largest projection in the centre of
\(\mathcal A(\zeta)\) such that
\(\mathcal A(\zeta)_{E(\zeta)}\) (respectively,
\(\mathcal A(\zeta)_{F(\zeta)}\)) is finite (respectively,
properly infinite).
\end{corollary}

% Source PDF page 216 (printed p. 207)
\phantomsection\label{srcpage:207}
\begin{proof}
The argument is the same as above, using this time
Theorem~2\textup{(ii)} and Theorem~5.
\end{proof}

\begin{corollary}\label{cor:II-5-5-2}
Suppose that \(\nu\) is standard. In order that \(\mathcal A\) be
finite (respectively, properly infinite), it is necessary and
sufficient that \(\mathcal A(\zeta)\) be finite (respectively,
properly infinite) almost everywhere.
\end{corollary}

\begin{proof}
This follows from Corollary~1.
\end{proof}

Bibliography: \ArticleRef{10}, \ArticleRef{80}, \ArticleRef{311}.

\par\medskip
\noindent\textit{Exercises.}---
\begin{enumerate}[label=\arabic*.,leftmargin=*,itemsep=0.5\baselineskip]
\item\label{ex:II-5-1}
Let \(\zeta\mapsto\mathfrak H(\zeta)\) be a
\(\nu\)-measurable field of complex Hilbert spaces on \(Z\). For
every \(\zeta\in Z\), let \(\varphi_{\zeta}\) be the trace on
\(\mathcal L(\mathfrak H(\zeta))_+\) defined in Chapter~I, \S~6,
no.~6. Show that the field \(\zeta\mapsto\varphi_{\zeta}\) is
measurable. (Consider a measurable field of orthonormal bases.)

\item\label{ex:II-5-2}
Let
\[
  \mathcal A
  =\int^{\oplus}\mathcal A(\zeta)\,d\nu(\zeta)
\]
be a decomposable von Neumann algebra, and let
\(\zeta\mapsto\varphi_{\zeta}\) be a measurable field of finite
traces on the \(\mathcal A(\zeta)\), such that
\(\varphi_{\zeta}(I_{\mathfrak H(\zeta)})=1\) for every \(\zeta\).
Suppose that
\[
  \int d\nu(\zeta)=1.
\]
Then \(\varphi=\int^{\oplus}\varphi_{\zeta}\,d\nu(\zeta)\) is a
finite trace on \(\mathcal A\) such that \(\varphi(I)=1\). Let
\(\Delta\) (respectively, \(\Delta_{\zeta}\)) be the determinant
associated with \(\varphi\) (respectively, \(\varphi_{\zeta}\)).
Show that, if
\[
  T=\int^{\oplus}T(\zeta)\,d\nu(\zeta)\in\mathcal A
\]
is invertible, then \(T(\zeta)\) is invertible almost everywhere,
and
\[
  \Delta(T)
  =\exp\int\log\Delta_{\zeta}(T(\zeta))\,d\nu(\zeta).
\]
(Use Exercise~2b of \S~2.)

\item\label{ex:II-5-3}
Let
\[
  \mathcal A
  =\int^{\oplus}\mathcal A(\zeta)\,d\nu(\zeta)
\]
be a decomposable von Neumann algebra. Suppose that \(\nu\) is
standard. The set of \(\zeta\) for which \(\mathcal A(\zeta)\) is
of type~I (respectively, type~\(\mathrm{II}_1\),
type~\(\mathrm{II}_{\infty}\), type~III) is measurable. (Use
\S~3, Corollary~1 to Proposition~7, Corollary~1 to Theorem~4, and
Corollary~1 to Theorem~5.) \ArticleRef{80}, \ArticleRef{311}.
\end{enumerate}

\section{Decomposition of a Hilbert Space as a Hilbertian Integral}
\label{sec:II-6}

\NumberedParagraph{1}{Statement of the Problem}
\label{no:II-6-1}
The data of \S\S~1, 2, and~3 being given, namely:
\begin{enumerate}[label=\arabic*\textsuperscript{o},leftmargin=*]
\item a compact metrizable space \(Z\);
\item a positive measure \(\nu\) on \(Z\), with support \(Z\);
\item a \(\nu\)-measurable field
\(\zeta\mapsto\mathfrak H(\zeta)\) of nonzero Hilbert spaces on
\(Z\),
\end{enumerate}
one can construct canonically:
\begin{enumerate}[label=\arabic*\textsuperscript{o},leftmargin=*]
\item the Hilbert space with a countable basis
\[
  \mathfrak H
  =\int^{\oplus}\mathfrak H(\zeta)\,d\nu(\zeta);
\]
\item the abelian von Neumann algebra \(\mathfrak Z\) of
diagonalizable operators;
\item more precisely, the \(C^*\)-algebra \(\mathcal Y\) of
continuously diagonalizable operators, whose weak closure is
\(\mathfrak Z\).
\end{enumerate}

% Source PDF page 217 (printed p. 208)
\phantomsection\label{srcpage:208}
We shall now show that the course of these constructions can be
reversed in an essentially unique manner. This will greatly increase
the scope of the preceding sections.

\Needspace{4\baselineskip}
\NumberedParagraph{2}{Existence Theorems}
\label{no:II-6-2}

\begin{theorem}\label{thm:II-6-1}
Let \(\mathfrak H\) be a complete Hilbert space with a countable
basis, let \(\mathcal Y\) be an abelian \(C^*\)-algebra of operators
on \(\mathfrak H\), let \(Z\) be the spectrum of \(\mathcal Y\),
and let \(\nu\) be a basic measure on \(Z\). Suppose that \(I\)
belongs to the weak closure of \(\mathcal Y\). Then there exist a
\(\nu\)-measurable field \(\zeta\mapsto\mathfrak H(\zeta)\) of
nonzero complete Hilbert spaces on \(Z\), and an isomorphism of
\(\mathfrak H\) onto
\[
  \int^{\oplus}\mathfrak H(\zeta)\,d\nu(\zeta),
\]
which transforms the Gelfand isomorphism into the canonical
isomorphism of \(\mathcal C_\infty(Z)\) onto the algebra of
continuously diagonalizable operators.
\end{theorem}

\begin{proof}
\textup{(i)} Denote by \(f\mapsto T_f\) the weakly continuous
isomorphism of \(L_{\mathbb C}^{\infty}(Z,\nu)\) onto the weak
closure \(\mathfrak Z\) of \(\mathcal Y\), extending the Gelfand
isomorphism (Chapter~I, \S~7, Proposition~1); this weak closure is,
moreover, \(\mathcal Y''\), since it contains \(I\) (Chapter~I,
\S~3, Theorem~2). Let \((x_1,x_2,\ldots)\) be an everywhere dense
sequence in \(\mathfrak H\). By adjoining to the \(x_i\) their
linear combinations with rational complex coefficients, we may
suppose that the \(x_i\) form a vector subspace \(\mathfrak H'\) of
\(\mathfrak H\) over the rational complex field.

For \(x,y\in\mathfrak H\), let \(h_{x,y}\) be the density, with
respect to \(\nu\), of the spectral measure \(\nu_{x,y}\). By the
formulas of Chapter~I, \S~7, no.~1, and the countability of
\(\mathfrak H'\), there exists a \(\nu\)-null subset \(N\) of
\(Z\) such that, for \(\zeta\notin N\), the function
\[
  (x,y)\longmapsto h_{x,y}(\zeta)
\]
is, on \(\mathfrak H'\), a positive Hermitian sesquilinear form.
Let \(\mathfrak H(\zeta)\) be the Hilbert space obtained, by passage
to the quotient and completion, from the space \(\mathfrak H'\)
endowed with this sesquilinear form, and let \(\phi(\zeta)\) be the
canonical mapping of \(\mathfrak H'\) into
\(\mathfrak H(\zeta)\). Let \(N_1\subset Z\setminus N\) be the set
of those \(\zeta\in Z\setminus N\) for which
\(\mathfrak H(\zeta)=0\). For \(\zeta\in Z\setminus N\), the
condition \(\zeta\in N_1\) is equivalent to
\[
  h_{x_i,x_j}(\zeta)=0
  \qquad\text{for all }i,j;
\]
thus \(N_1\) is \(\nu\)-measurable. Let \(f\) be its
characteristic function. We have
\[
  (T_fx_i\mid T_fx_i)
  =(T_fx_i\mid x_i)
  =\int f(\zeta)\,d\nu_{x_i,x_i}(\zeta)
  =\int f(\zeta)h_{x_i,x_i}(\zeta)\,d\nu(\zeta)=0,
\]
and hence \(T_fx_i=0\) for every \(i\), so that \(T_f=0\). Thus
\(N_1\) is \(\nu\)-null. By choosing arbitrarily new spaces
\(\mathfrak H(\zeta)\) on the \(\nu\)-null set \(N\cup N_1\), we
may arrange that \(\mathfrak H(\zeta)\ne0\) for every
\(\zeta\in Z\).

% Source PDF page 218 (printed p. 209)
\phantomsection\label{srcpage:209}
\textup{(ii)} We now endow the \(\mathfrak H(\zeta)\) with the
structure of a measurable field. For \(\zeta\notin N\cup N_1\), put
\(x_i(\zeta)=\phi(\zeta)x_i\). The number
\[
  \bigl(x_i(\zeta)\mid x_j(\zeta)\bigr)=h_{x_i,x_j}(\zeta)
\]
depends measurably on \(\zeta\); and, for every
\(\zeta\notin N\cup N_1\), the \(x_i(\zeta)\) form a total set in
\(\mathfrak H(\zeta)\). There is therefore on the
\(\mathfrak H(\zeta)\) one and only one measurable-field structure
such that the fields \(\zeta\mapsto x_i(\zeta)\) are measurable
fields of vectors (\S~1, Proposition~4).

\textup{(iii)} We shall define an isomorphism of \(\mathfrak H\)
onto \(\int^{\oplus}\mathfrak H(\zeta)\,d\nu(\zeta)\). Let
\[
  x=\sum_{i=1}^{n}T_{f_i}x_i\in\mathfrak H,
  \qquad
  f_1,f_2,\ldots,f_n\in L_{\mathbb C}^{\infty}(Z,\nu).
\]
We have
\begin{align*}
  \lVert x\rVert^2
    &=\int d\nu_{x,x}(\zeta)
      =\int h_{x,x}(\zeta)\,d\nu(\zeta)\\
    &=\int\sum_{i,j=1}^{n}
      f_i(\zeta)\overline{f_j(\zeta)}
      h_{x_i,x_j}(\zeta)\,d\nu(\zeta)\\
    &=\int\sum_{i,j=1}^{n}
      f_i(\zeta)\overline{f_j(\zeta)}
      \bigl(x_i(\zeta)\mid x_j(\zeta)\bigr)\,d\nu(\zeta)\\
    &=\int\left\lVert
      \sum_{i=1}^{n}f_i(\zeta)x_i(\zeta)
      \right\rVert^2\,d\nu(\zeta).
\end{align*}
This shows first that the field of vectors
\[
  \zeta\longmapsto\sum_{i=1}^{n}f_i(\zeta)x_i(\zeta)
\]
is square-integrable; next, that this field depends (up to null sets)
only on the vector \(x\), and not on its representation in the form
\(\sum_{i=1}^{n}T_{f_i}x_i\); and finally, that the mapping \(U_0\)
which assigns to \(x\in\mathfrak H\) the field
\(\zeta\mapsto\sum_{i=1}^{n}f_i(\zeta)x_i(\zeta)\) in
\(\int^{\oplus}\mathfrak H(\zeta)\,d\nu(\zeta)\) is isometric.
These fields are everywhere dense in
\(\int^{\oplus}\mathfrak H(\zeta)\,d\nu(\zeta)\)
(\S~1, Proposition~7). On the other hand, vectors of the form
\(\sum_{i=1}^{n}T_{f_i}x_i\) are everywhere dense in
\(\mathfrak H\). Thus \(U_0\) extends to an isomorphism \(U\) of
\(\mathfrak H\) onto
\(\int^{\oplus}\mathfrak H(\zeta)\,d\nu(\zeta)\).

A function \(f\in L_{\mathbb C}^{\infty}(Z,\nu)\) defines, on the
one hand, an operator \(T_f\) on \(\mathfrak H\), and, on the other
hand, a diagonalizable operator \(T'_f\) on
\(\int^{\oplus}\mathfrak H(\zeta)\,d\nu(\zeta)\). With the
preceding notation,
\[
  T_fx=\sum_{i=1}^{n}T_{ff_i}x_i.
\]%
% Source PDF page 219 (printed p. 210)
\phantomsection\label{srcpage:210}%
On the other hand, \(T'_fU_0x\) is the field of vectors
\[
  \zeta\longmapsto
  \sum_{i=1}^{n}f(\zeta)f_i(\zeta)x_i(\zeta).
\]
Thus \(T'_fU_0x=U_0T_fx\), and consequently
\(T'_f=UT_fU^{-1}\).
\end{proof}

\begin{theorem}\label{thm:II-6-2}
Let \(\mathfrak H\) be a complete Hilbert space with a countable
basis, and let \(\mathfrak Z\) be an abelian von Neumann algebra on
\(\mathfrak H\). Then there exist a compact metrizable space \(Z\),
a positive measure \(\nu\) on \(Z\) with support \(Z\), a
\(\nu\)-measurable field \(\zeta\mapsto\mathfrak H(\zeta)\) of
nonzero complete Hilbert spaces on \(Z\), and an isomorphism of
\(\mathfrak H\) onto
\[
  \int^{\oplus}\mathfrak H(\zeta)\,d\nu(\zeta)
\]
which transforms \(\mathfrak Z\) into the algebra of diagonalizable
operators.
\end{theorem}

\begin{proof}
Let \(\mathcal Y\) be a weakly dense \(C^*\)-subalgebra of
\(\mathfrak Z\), whose spectrum \(Z\) is compact metrizable and
carries a basic measure \(\nu\) (Chapter~I, \S~7, Proposition~4).
By Theorem~1, there exist a \(\nu\)-measurable field
\(\zeta\mapsto\mathfrak H(\zeta)\) of nonzero Hilbert spaces on
\(Z\), and an isomorphism of \(\mathfrak H\) onto
\(\int^{\oplus}\mathfrak H(\zeta)\,d\nu(\zeta)\), which transforms
\(\mathcal Y\) into the algebra of continuously diagonalizable
operators, and hence \(\mathfrak Z\) into the algebra of
diagonalizable operators [\S~2, Proposition~7\textup{(i)}].
\end{proof}

\begin{corollary*}\phantomsection\label{cor:II-6-2}
Let \(\mathfrak H\) be a complete Hilbert space with a countable
basis, and let \(\mathcal A\) be a von Neumann algebra on
\(\mathfrak H\). There exist a compact metrizable space \(Z\), a
positive measure \(\nu\) on \(Z\) with support \(Z\), a
\(\nu\)-measurable field \(\zeta\mapsto\mathfrak H(\zeta)\) of
nonzero complete Hilbert spaces on \(Z\), a \(\nu\)-measurable
field \(\zeta\mapsto\mathcal A(\zeta)\) of factors on the
\(\mathfrak H(\zeta)\), and an isomorphism of \(\mathfrak H\) onto
\(\int^{\oplus}\mathfrak H(\zeta)\,d\nu(\zeta)\) which transforms
\(\mathcal A\) into
\[
  \int^{\oplus}\mathcal A(\zeta)\,d\nu(\zeta).
\]
\end{corollary*}

\begin{proof}
Apply Theorem~2 to the centre \(\mathfrak Z\) of \(\mathcal A\).
We obtain \(Z\), \(\nu\), \(\zeta\mapsto\mathfrak H(\zeta)\),
and an isomorphism \(U\) of \(\mathfrak H\) onto
\(\int^{\oplus}\mathfrak H(\zeta)\,d\nu(\zeta)\) which transforms
\(\mathfrak Z\) into the algebra of diagonalizable operators. We
have \(\mathfrak Z\subset\mathcal A\subset\mathfrak Z'\), and hence
\(U\mathcal AU^{-1}\) is decomposable (\S~3, Theorem~2).
Consequently there is a \(\nu\)-measurable field
\(\zeta\mapsto\mathcal A(\zeta)\) of von Neumann algebras on the
\(\mathfrak H(\zeta)\) such that
\[
  U\mathcal AU^{-1}
  =\int^{\oplus}\mathcal A(\zeta)\,d\nu(\zeta).
\]
Since the centre of \(U\mathcal AU^{-1}\) is the algebra of
diagonalizable operators, the \(\mathcal A(\zeta)\) are factors
almost everywhere (\S~3, Theorem~3).
\end{proof}

The corollary reduces, in a certain measure, the study of von Neumann
algebras to that of factors; this was one of the principal aims of
``reduction theory.''\space%
% Source PDF page 220 (printed p. 211)
\phantomsection\label{srcpage:211}%
Nevertheless, we saw in Chapter~I that arbitrary von Neumann
algebras can be studied directly by methods
which constitute ``global theory.''

Bibliography: \ArticleRef{28}, \ArticleRef{49}, \ArticleRef{80},
\ArticleRef{100}, \ArticleRef{117}, \ArticleRef{145},
\ArticleRef{193}, \ArticleRef{194}, \ArticleRef{205},
\ArticleRef{206}.

\NumberedParagraph{3}{Uniqueness Theorems}
\label{no:II-6-3}
Given \(\mathfrak H\) and \(\mathcal Y\), we wish to show that
\(Z\), \(\nu\), and the field \(\zeta\mapsto\mathfrak H(\zeta)\)
are essentially unique. We have already observed (\S~2, no.~4) that
\(Z\) is canonically identified with the spectrum of \(\mathcal Y\),
in such a way that the canonical isomorphism of
\(\mathcal C_\infty(Z)\) onto \(\mathcal Y\) is identified with the
Gelfand isomorphism. This proves the uniqueness of \(Z\) (up to a
homeomorphism), and permits the uniqueness theorem to be stated as
follows.

\begin{theorem}\label{thm:II-6-3}
Let \(Z\) be a locally compact space countable at infinity. Let
\(\nu\) be a positive measure on \(Z\), with support \(Z\), let
\(\zeta\mapsto\mathfrak H(\zeta)\) be a \(\nu\)-measurable field of
nonzero Hilbert spaces on \(Z\), let
\[
  \mathfrak H=\int^{\oplus}\mathfrak H(\zeta)\,d\nu(\zeta),
\]
let \(\mathcal Y\) be the algebra of continuously diagonalizable
operators on \(\mathfrak H\), and let \(f\mapsto T_f\) be the
canonical isomorphism of \(\mathcal C_\infty(Z)\) onto
\(\mathcal Y\). Define similarly
\(\nu_1\), \(\zeta\mapsto\mathfrak H_1(\zeta)\),
\(\mathfrak H_1\), \(\mathcal Y_1\), and \(f\mapsto T_f^1\).
Let \(U\) be an isomorphism of \(\mathfrak H\) onto
\(\mathfrak H_1\) which transforms \(T_f\) into \(T_f^1\) for
every \(f\in\mathcal C_\infty(Z)\). Then \(\nu\) and \(\nu_1\)
are equivalent; and, after modifying the \(\mathfrak H(\zeta)\) and
\(\mathfrak H_1(\zeta)\), if necessary, on null sets, there exists
an isomorphism \(\zeta\mapsto V(\zeta)\) of the field
\((\mathfrak H(\zeta))\) onto the field
\((\mathfrak H_1(\zeta))\), such that \(U=WV\), where
\[
  V=\int^{\oplus}V(\zeta)\,d\nu(\zeta)
\]
is the isomorphism of \(\mathfrak H\) onto
\(\int^{\oplus}\mathfrak H_1(\zeta)\,d\nu(\zeta)\), and where
\(W\) is the canonical isomorphism of
\(\int^{\oplus}\mathfrak H_1(\zeta)\,d\nu(\zeta)\) onto
\(\mathfrak H_1\).
\end{theorem}

\begin{proof}
Identify \(Z\) with the spectrum of \(\mathcal Y\) and with that of
\(\mathcal Y_1\), so that, for every
\(f\in\mathcal C_\infty(Z)\), \(T_f\) and
\(T_f^1=UT_fU^{-1}\) are the elements of \(\mathcal Y\) and
\(\mathcal Y_1\) corresponding to \(f\) under the Gelfand
isomorphisms. Then, if \(x,y\in\mathfrak H\), we have
\(\nu_{x,y}=\nu_{Ux,Uy}\), and hence the basic measures on \(Z\)
defined by \(\mathcal Y\) and \(\mathcal Y_1\) are the same. Thus
\(\nu\) and \(\nu_1\) are equivalent.

Let
\[
  \widetilde{\mathfrak H}
  =\int^{\oplus}\mathfrak H_1(\zeta)\,d\nu(\zeta),
\]
let \(W\) be the canonical isomorphism of
\(\widetilde{\mathfrak H}\) onto \(\mathfrak H_1\), and let
\(\widetilde{\mathcal Y}\) be the transform of \(\mathcal Y_1\) by
\(W^{-1}\); this is the algebra of continuously diagonalizable
operators on \(\widetilde{\mathfrak H}\). Every
\(f\in\mathcal C_\infty(Z)\) defines an operator
\(\widetilde T_f\) on \(\widetilde{\mathfrak H}\), and
\[
  \widetilde T_f
  =(W^{-1}U)T_f(W^{-1}U)^{-1}.
\]
% Source PDF page 221 (printed p. 212)
\phantomsection\label{srcpage:212}
Thus \(W^{-1}U=V\) is a decomposable linear mapping from
\(\mathfrak H\) into \(\widetilde{\mathfrak H}\)
(\S~2, Theorem~1), so that there is a measurable field
\(\zeta\mapsto V(\zeta)\) of linear mappings from
\(\mathfrak H(\zeta)\) into \(\mathfrak H_1(\zeta)\), with
\[
  V=\int^{\oplus}V(\zeta)\,d\nu(\zeta).
\]
Finally, since \(V\) is an isomorphism of \(\mathfrak H\) onto
\(\widetilde{\mathfrak H}\), we have
\[
  V^*V=I_{\mathfrak H},
  \qquad
  VV^*=I_{\widetilde{\mathfrak H}},
\]
and therefore
\[
  V(\zeta)^*V(\zeta)=I_{\mathfrak H(\zeta)},
  \qquad
  V(\zeta)V(\zeta)^*=I_{\mathfrak H_1(\zeta)}
\]
almost everywhere; that is, \(V(\zeta)\) is almost everywhere an
isomorphism of \(\mathfrak H(\zeta)\) onto
\(\mathfrak H_1(\zeta)\).
\end{proof}

\begin{theorem}\label{thm:II-6-4}
Let \(Z\) be a Borel space, let \(\nu\) be a standard positive
measure on \(Z\), let
\(\mathcal E=(\mathfrak H(\zeta))\) be a \(\nu\)-measurable field
of nonzero complete Hilbert spaces on \(Z\), let
\[
  \mathfrak H
  =\int^{\oplus}\mathfrak H(\zeta)\,d\nu(\zeta),
\]
and let \(\mathfrak Z\) be the algebra of diagonalizable operators
on \(\mathfrak H\). Define similarly
\(Z_1\), \(\nu_1\),
\(\mathcal E_1=(\mathfrak H_1(\zeta_1))\),
\(\mathfrak H_1\), and \(\mathfrak Z_1\). Let \(U\) be an
isomorphism of \(\mathfrak H\) onto \(\mathfrak H_1\) which
transforms \(\mathfrak Z\) into \(\mathfrak Z_1\). Then there exist:
\begin{enumerate}[label=\arabic*\textsuperscript{o},leftmargin=*]
\item a \(\nu\)-null Borel set \(N\) in \(Z\), and a
\(\nu_1\)-null Borel set \(N_1\) in \(Z_1\);
\item a Borel isomorphism \(\eta\) of \(Z\setminus N\) onto
\(Z_1\setminus N_1\), which transforms \(\nu\) into a measure
\(\widetilde\nu_1\) equivalent to \(\nu_1\);
\item an \(\eta\)-isomorphism \((V(\zeta))\) of
\(\mathcal E\mathbin{|}(Z\setminus N)\) onto
\(\mathcal E_1\mathbin{|}(Z_1\setminus N_1)\), defining an
isomorphism \(V\) of \(\mathfrak H\) onto
\[
  \widetilde{\mathfrak H}_1
  =\int^{\oplus}\mathfrak H_1(\zeta_1)\,d\widetilde\nu_1(\zeta_1)
\]
such that \(U=WV\), where \(W\) is the canonical isomorphism of
\(\widetilde{\mathfrak H}_1\) onto \(\mathfrak H_1\).
\end{enumerate}
\end{theorem}

\begin{proof}
We may suppose that \(Z\) and \(Z_1\) are compact metrizable. The
isomorphism \(U\) defines an isomorphism of \(\mathfrak Z\) onto
\(\mathfrak Z_1\), and hence an isomorphism \(\Phi\) of
\(L_{\mathbb C}^{\infty}(Z,\nu)\) onto
\(L_{\mathbb C}^{\infty}(Z_1,\nu_1)\). Consequently there exist
\(N,N_1,\widetilde\nu_1,\eta\), with properties~\textup{1} and
\textup{2} of the theorem, \(\eta\) also defining the isomorphism
\(\Phi\) (Appendix~IV).

For \(f\in L_{\mathbb C}^{\infty}(Z,\nu)\) [respectively,
\(f_1\in L_{\mathbb C}^{\infty}(Z_1,\nu_1)\)], denote by \(T_f\)
[respectively, \(\widetilde T_{f_1}^1\)] the diagonalizable operator
on \(\mathfrak H\) (respectively, \(\widetilde{\mathfrak H}_1\))
defined by \(f\) (respectively, \(f_1\)). If \(f\) and \(f_1\)
correspond under \(\Phi\), then \(T_f\) and
\(W\widetilde T_{f_1}^1W^{-1}\) correspond under \(U\); that is,
\[
  W\widetilde T_{f_1}^1W^{-1}=UT_fU^{-1},
\]
or equivalently
\[
  \widetilde T_{f_1}^1W^{-1}U=W^{-1}UT_f,
\]
or again
\[
  \widetilde T_{f_1}^1V=VT_f,
  \qquad
  V=W^{-1}U\in\mathcal L(\mathfrak H,\widetilde{\mathfrak H}_1).
\]%
% Source PDF page 222 (printed p. 213)
\phantomsection\label{srcpage:213}%
It remains, therefore, to show that \(V\) is ``decomposable'' in a
generalized sense. We could have presented this generalization in
\S~2, but it would have been cumbersome. We shall briefly indicate
how the arguments of \S~2 extend, leaving to the reader the details,
which present no difficulty. Let \((x_1,x_2,\ldots)\) be a
fundamental sequence of measurable fields of vectors on
\(Z\setminus N\), such that \(x_i\in\mathfrak H\); put
\(y_i=Vx_i\in\widetilde{\mathfrak H}_1\). Using the equality
\(\widetilde T_{f_1}^1V=VT_f\), one proves, exactly as for
Theorem~1 of \S~2, that
\[
  \left\lVert\sum_{i=1}^{n}
    \rho_i y_i(\eta(\zeta))\right\rVert
  \leq
  \left\lVert\sum_{i=1}^{n}\rho_i x_i(\zeta)\right\rVert
\]
almost everywhere on \(Z\setminus N\), for all rational complex
numbers \(\rho_1,\rho_2,\ldots,\rho_n\). Hence there exists a
continuous linear mapping \(V(\zeta)\) from
\(\mathfrak H(\zeta)\) into \(\mathfrak H_1(\eta(\zeta))\), for
\(\zeta\in Z\setminus N\), such that
\[
  \lVert V(\zeta)\rVert\leq1,
  \qquad
  V(\zeta)x_i(\zeta)=y_i(\eta(\zeta))
\]
for every \(i\), almost everywhere. Generalizing Proposition~1 of
\S~2, one concludes that the \(V(\zeta)\) transform every
measurable field of vectors on \(Z\setminus N\) into a measurable
field of vectors on \(Z_1\setminus N_1\).

Reasoning similarly with \(V^{-1}\), there exists, for
\(\zeta_1\in Z_1\setminus N_1\), a continuous linear mapping
\(V'(\zeta_1)\) from \(\mathfrak H_1(\zeta_1)\) into
\(\mathfrak H(\eta^{-1}(\zeta_1))\), such that
\[
  \lVert V'(\zeta_1)\rVert\leq1,
  \qquad
  V'(\zeta_1)y_i(\zeta_1)
  =x_i(\eta^{-1}(\zeta_1))
\]
for every \(i\), almost everywhere; and the \(V'(\zeta_1)\)
transform every measurable field of vectors on
\(Z_1\setminus N_1\) into a measurable field of vectors on
\(Z\setminus N\). Consequently \(V(\zeta)\) and
\(V'(\eta(\zeta))\) are almost everywhere reciprocal isomorphisms
of \(\mathfrak H(\zeta)\) onto
\(\mathfrak H_1(\eta(\zeta))\) and of
\(\mathfrak H_1(\eta(\zeta))\) onto \(\mathfrak H(\zeta)\). After
modification on a null set,
\((V(\zeta))_{\zeta\in Z\setminus N}\) is an
\(\eta\)-isomorphism of
\(\mathcal E\mathbin{|}(Z\setminus N)\) onto
\(\mathcal E_1\mathbin{|}(Z_1\setminus N_1)\). The corresponding
isomorphism of \(\mathfrak H\) onto \(\widetilde{\mathfrak H}_1\)
transforms, as does \(V\), the elements of the form \(T_fx_i\), and
hence is equal to \(V\).
\end{proof}

Bibliography: \ArticleRef{80}.

\par\medskip
\noindent\textit{Exercises.}---
\begin{enumerate}[label=\arabic*.,leftmargin=*,itemsep=0.5\baselineskip]
\item\label{ex:II-6-1}
Let \(\zeta\mapsto\mathfrak H(\zeta)\) be a
\(\nu\)-measurable field of complex Hilbert spaces on \(Z\), let
\[
  \mathfrak H
  =\int^{\oplus}\mathfrak H(\zeta)\,d\nu(\zeta),
\]
let \(\mathfrak Z\) be the algebra of diagonalizable operators, and
let \(\mathcal A\) be a factor contained in \(\mathfrak Z'\).
Suppose that \(\mathfrak Z\) is a maximal abelian von Neumann
subalgebra of \(\mathcal A'\).
\begin{enumerate}[label=\textit{\alph*.},leftmargin=*]
\item If \(\mathcal A\) is discrete, and if
\(\mathfrak H(\zeta)\ne0\) almost everywhere, there exist:
\textup{(1)} a Hilbert space \(\mathfrak H_0\) with a countable
basis; \textup{(2)} almost everywhere on \(Z\), an isomorphism
\(U(\zeta)\) of \(\mathfrak H(\zeta)\) onto
\(\mathfrak H_0\); \textup{(3)} for every \(T\in\mathcal A\), a
unique operator \(T_1\in\mathcal L(\mathfrak H_0)\) such that
\[
  T=\int^{\oplus}
    U(\zeta)^{-1}T_1U(\zeta)\,d\nu(\zeta).
\]
% Source PDF page 223 (printed p. 214)
\phantomsection\label{srcpage:214}
(Use \hyperref[prop:II-2-8]{\S~2, Proposition~8} and Exercise~1, and
\hyperref[thm:II-6-3]{Theorem~3 of the present section}.) As \(T\) ranges
over \(\mathcal A\), \(T_1\) ranges over
\(\mathcal L(\mathfrak H_0)\).

\item In the general case, for every
\[
  T=\int^{\oplus}T(\zeta)\,d\nu(\zeta)\in\mathcal A,
\]
\(T\ne0\) implies \(T(\zeta)\ne0\) almost everywhere. [Let \(Y\) be the
measurable subset of \(Z\) consisting of the \(\zeta\) such that
\(T(\zeta)=0\). Let \(E\) be the diagonalizable projection corresponding
to \(Y\). We have \(T_E=0\), hence \(E=0\).]
\ArticleRef{16}, \ArticleRef{57}, \ArticleRef{58}.

\medskip
\noindent\textit{Problem.} What relations are there among the
\(T(\zeta)\)? \ArticleRef{245}, \ArticleRef{265}.
\end{enumerate}

\item We place ourselves under the hypotheses of
\hyperref[thm:II-6-2]{Theorem~2}. Show that,
with the notation of that theorem, \(Z\) may be taken to be a compact
interval of the real line. (By
\hyperref[ex:I-7-3]{Chapter~I, \S~7, Exercise~3(f)}, take
\(\mathcal Y\) to be generated by a single Hermitian element. The spectrum
of \(\mathcal Y\) is then a compact subset of the real line.)
\ArticleRef{80}.
\end{enumerate}

% Source PDF page 224 (printed p. 215)
\chapter{Further Topics}\label{ch:III}\label{srcpage:215}

\section{Comparison of Projections}\label{sec:III-1}

\NumberedParagraph{1}{Comparison of Projections}
\label{no:III-1-1}

\begin{definition}\label{def:III-1-1}
Let \(\mathcal A\) be a von Neumann algebra and let \(E,F\) be two
projections of \(\mathcal A\). The projections \(E\) and \(F\) are said to
be \emph{equivalent} (relative to \(\mathcal A\)), and one writes
\(E\sim F\), if there is an element \(U\in\mathcal A\) such that
\[
  U^*U=E,\qquad UU^*=F.
\]
One writes \(E\prec F\), or \(F\succ E\), if there is a projection of
\(\mathcal A\) equivalent to \(E\) and majorized by \(F\).
\end{definition}

If \(\mathfrak X=E(\mathfrak H)\) and
\(\mathfrak Y=F(\mathfrak H)\), the notations
\(\mathfrak X\sim\mathfrak Y\), \(\mathfrak X\prec\mathfrak Y\), and
\(\mathfrak Y\succ\mathfrak X\) are also used.

Two projections \(E,F\in\mathcal A\) are equivalent if there is in
\(\mathcal A\) a partially isometric operator having \(E\) as initial
projection and \(F\) as final projection, or, again, if there is in
\(\mathcal A\) an operator whose restriction to \(E(\mathfrak H)\) is an
isometric mapping of \(E(\mathfrak H)\) onto \(F(\mathfrak H)\). This
restriction then defines a spatial isomorphism of \(\mathcal A_E\) onto
\(\mathcal A_F\), and of \(\mathcal A'_E\) onto \(\mathcal A'_F\). It is
immediate that the relation \(E\sim F\) is an equivalence relation, and
that the relation \(E\prec F\) is reflexive and transitive (in other words,
it is a preorder relation).

Let \(E,F\) be two projections of \(\mathcal A\), and let \(E',F'\) be
their central supports. If \(E\sim F\), then \(E'\leq F'\) (by
\hyperref[cor:I-1-7-1]{Chapter~I, \S~1, Proposition~7, Corollary~1}),
and similarly \(F'\leq E'\), hence
\(E'=F'\). Consequently, if \(E\prec F\), then \(E'\leq F'\).

Let \((E_\iota)_{\iota\in I}\) [respectively
\((F_\iota)_{\iota\in I}\)] be a family of pairwise orthogonal
projections of \(\mathcal A\), and let
\[
  E=\sum_{\iota\in I}E_\iota
  \qquad
  \left(\text{respectively }F=\sum_{\iota\in I}F_\iota\right).
\]
If \(E_\iota\sim F_\iota\) for every \(\iota\in I\), then
\(E\sim F\). Indeed, let \(U_\iota\) be a partially isometric operator of
\(\mathcal A\) having \(E_\iota\) as initial projection and \(F_\iota\) as
final projection. There is a unique partially isometric operator \(U\)
having \(E\) as initial projection, \(F\) as final projection, and agreeing
with \(U_\iota\) on \(E_\iota(\mathfrak H)\). This operator \(U\) is
invariant under the unitary operators of \(\mathcal A'\), hence
\(U\in\mathcal A\), and our assertion is proved. It follows that, if
\(E_\iota\prec F_\iota\) for every \(\iota\in I\), then \(E\prec F\).

% Source PDF page 225 (printed p. 216)
\phantomsection\label{srcpage:216}
Let \((\mathcal A_x)_{x\in K}\) be a family of von Neumann algebras, and
let \(\mathcal A\) be their product. For each \(x\in K\), let \(E_x,F_x\)
be two projections of \(\mathcal A_x\). Put \(E=(E_x)\), \(F=(F_x)\),
which are projections of \(\mathcal A\). In order that \(E\sim F\)
(respectively \(E\prec F\)), it is necessary and sufficient that
\(E_x\sim F_x\) (respectively \(E_x\prec F_x\)) for every \(x\in K\).

It is clear that \(E\sim F\) implies \(E\prec F\) and \(F\prec E\).
Conversely:

\begin{proposition}\label{prop:III-1-1}
If \(E\prec F\) and \(F\prec E\), then \(E\sim F\).
\end{proposition}

\begin{proof}
We have \(E\sim F'\leq F\) and \(F\sim E'\leq E\). The relations
\(F'\leq F\) and \(F\sim E'\) imply \(F'\sim E''\), with
\(E''\leq E'\leq E\); moreover, \(E\sim F'\sim E''\). Let \(U\) be a
partially isometric operator of \(\mathcal A\) having \(E\) as initial
projection and \(E''\) as final projection. Put
\[
  \mathfrak X=E(\mathfrak H),\qquad
  \mathfrak X'=E'(\mathfrak H),\qquad
  \mathfrak X''=E''(\mathfrak H),
\]
and
\[
  \mathfrak X^{(2n)}=U^n(\mathfrak X),\qquad
  \mathfrak X^{(2n+1)}=U^n(\mathfrak X')
\]
(so that \(\mathfrak X^{(0)}=\mathfrak X\),
\(\mathfrak X^{(1)}=\mathfrak X'\), and
\(\mathfrak X^{(2)}=\mathfrak X''\)). The relations
\(\mathfrak X\supset\mathfrak X'\supset\mathfrak X''\) imply
\(\mathfrak X^{(2n)}\supset\mathfrak X^{(2n+1)}
\supset\mathfrak X^{(2n+2)}\), so the sequence
\((\mathfrak X^{(i)})\) is decreasing. Let
\[
  \mathfrak Y=\bigcap_{i=0}^{\infty}\mathfrak X^{(i)}.
\]
Then \(\mathfrak X\) is the Hilbert sum of the mutually orthogonal
subspaces
\[
  \mathfrak Y,\quad
  \mathfrak X^{(0)}\ominus\mathfrak X^{(1)},\quad
  \mathfrak X^{(1)}\ominus\mathfrak X^{(2)},\quad
  \mathfrak X^{(2)}\ominus\mathfrak X^{(3)},\quad
  \mathfrak X^{(3)}\ominus\mathfrak X^{(4)},\quad\ldots,
\]
whereas \(\mathfrak X'\) is the Hilbert sum of the mutually orthogonal
subspaces
\[
  \mathfrak Y,\quad
  \mathfrak X^{(2)}\ominus\mathfrak X^{(3)},\quad
  \mathfrak X^{(1)}\ominus\mathfrak X^{(2)},\quad
  \mathfrak X^{(4)}\ominus\mathfrak X^{(5)},\quad
  \mathfrak X^{(3)}\ominus\mathfrak X^{(4)},\quad\ldots.
\]
Now \(U\) maps \(\mathfrak X^{(i)}\ominus\mathfrak X^{(i+1)}\)
isometrically onto
\(\mathfrak X^{(i+2)}\ominus\mathfrak X^{(i+3)}\), and hence
\[
  \mathfrak X^{(i)}\ominus\mathfrak X^{(i+1)}
  \sim
  \mathfrak X^{(i+2)}\ominus\mathfrak X^{(i+3)}.
\]
Thus \(\mathfrak X\sim\mathfrak X'\), and consequently
\(E\sim E'\sim F\).
\end{proof}

\begin{proposition}\label{prop:III-1-2}
Let \(T\in\mathcal A\). The supports of \(T\) and \(T^*\) are equivalent.
\end{proposition}

\begin{proof}
Let \(T=W|T|\) be the polar decomposition of \(T\). The supports of \(T\)
and \(T^*\) are \(W^*W\) and \(WW^*\), respectively. Moreover,
\(W\in\mathcal A\).
\end{proof}

\setcounter{corollary}{0}
\begin{corollary}\label{cor:III-1-prop2-1}
Let \(E=P_{\mathfrak X}\) and \(F=P_{\mathfrak Y}\) be two projections of
\(\mathcal A\), and put
\(\mathfrak X'=\mathfrak H\ominus\mathfrak X\),
\(\mathfrak Y'=\mathfrak H\ominus\mathfrak Y\). Then
\[
  \mathfrak X\ominus(\mathfrak X\cap\mathfrak Y')
  \sim
  \mathfrak Y\ominus(\mathfrak Y\cap\mathfrak X').
\]
\end{corollary}

\begin{proof}
The null space of \(FE\) is generated by \(\mathfrak X'\) and
\(\mathfrak X\cap\mathfrak Y'\). Hence the support of \(FE\) is the
projection onto
\(\mathfrak X\ominus(\mathfrak X\cap\mathfrak Y')\).
% Source PDF page 226 (printed p. 217)
\phantomsection\label{srcpage:217}
Similarly, the support of \(EF\) is the projection onto
\(\mathfrak Y\ominus(\mathfrak Y\cap\mathfrak X')\). But
\(EF=(FE)^*\).
\end{proof}

\begin{corollary}\label{cor:III-1-prop2-2}
Let \(T'\in\mathcal A'\), let \(E'\) be the support of \(T'\), let
\(x\in\mathfrak H\), and put \(y=T'x\). Then
\(\mathfrak X_y^{\mathcal A}\prec\mathfrak X_x^{\mathcal A}\).
Moreover, if \(E'x=x\), then
\(\mathfrak X_y^{\mathcal A}\sim\mathfrak X_x^{\mathcal A}\).
\end{corollary}

\begin{proof}
For every \(T\in\mathcal A\), we have
\(Ty=TT'x=T'Tx\), and therefore
\(\mathfrak X_y^{\mathcal A}\) is the closure of
\(T'(\mathfrak X_x^{\mathcal A})\), that is, of
\(T'E_x^{\mathcal A}(\mathfrak H)\). By
\hyperref[prop:III-1-2]{Proposition~2},
\(\mathfrak X_y^{\mathcal A}\) is consequently equivalent to the closure
of
\[
  (T'E_x^{\mathcal A})^*(\mathfrak H)
  =E_x^{\mathcal A}T'^*(\mathfrak H)
  \subset\mathfrak X_x^{\mathcal A},
\]
whence
\(\mathfrak X_y^{\mathcal A}\prec\mathfrak X_x^{\mathcal A}\).

If \(E'x=x\), then \(\mathfrak X_x^{\mathcal A}\) is contained in
\(E'(\mathfrak H)=\overline{T'^*(\mathfrak H)}\), and hence the closure of
\(E_x^{\mathcal A}T'^*(\mathfrak H)\) equals
\(\mathfrak X_x^{\mathcal A}\). This time we therefore have
\(\mathfrak X_y^{\mathcal A}\sim\mathfrak X_x^{\mathcal A}\).
\end{proof}

\begin{proposition}\label{prop:III-1-3}
If \(\mathcal A\) is finite, every projection of \(\mathcal A\) equivalent
to \(I\) is equal to \(I\).
\end{proposition}

\begin{proof}
If \(E\sim I\), then \(I=U^*U\) and \(E=UU^*\) for some
\(U\in\mathcal A\). Hence, for every finite trace \(\varphi\) on
\(\mathcal A\),
\[
  \varphi(I-E)=\varphi(U^*U)-\varphi(UU^*)=0.
\]
Thus \(I-E=0\).
\end{proof}

The proof of \hyperref[prop:III-1-1]{Proposition~1} is analogous to a
familiar proof in the theory
of cardinals (the role played there by equipotence of sets is played here
by equivalence of projections).

In \S~8 we shall prove, not without difficulty, the converse of
\hyperref[prop:III-1-3]{Proposition~3}.

In some papers the notation \(E\prec F\) means (with the notation of the
text): ``\(E\prec F\) and \(E\) is not equivalent to \(F\).''

\noindent\textit{Bibliography.} \ArticleRef{42}, \ArticleRef{65}.

\Needspace{4\baselineskip}
\NumberedParagraph{2}{Comparability Theorem}
\label{no:III-1-2}

\begin{lemma}\label{lem:III-1-1}
Let \(E_1,E_2\) be projections of \(\mathcal A\), and let \(F_1,F_2\) be
their central supports. If \(F_1F_2\ne0\), then there are two equivalent
nonzero projections \(G_1,G_2\) of \(\mathcal A\), majorized respectively
by \(E_1,E_2\).
\end{lemma}

\begin{proof}
Put
\[
  \mathfrak X_1=E_1(\mathfrak H),\quad
  \mathfrak X_2=E_2(\mathfrak H),\quad
  \mathfrak Y_1=F_1(\mathfrak H),\quad
  \mathfrak Y_2=F_2(\mathfrak H).
\]
By \hyperref[cor:I-1-7-1]{Chapter~I, \S~1, Proposition~7, Corollary~1},
\(\mathfrak Y_i\) is the closed vector subspace generated by the
\(T(\mathfrak X_i)\), as \(T\) ranges over \(\mathcal A\). One may even,
by \hyperref[prop:I-1-3]{Chapter~I, \S~1, Proposition~3}, restrict \(T\)
to the unitary elements
of \(\mathcal A\), in which case
\(T(\mathfrak X_i)\sim\mathfrak X_i\). Since \(F_1F_2\ne0\), by replacing
\(\mathfrak X_1,\mathfrak X_2\), if necessary, by equivalent subspaces,
we may therefore suppose that \(\mathfrak X_1\) and \(\mathfrak X_2\)
are not orthogonal.
\hyperref[cor:III-1-prop2-1]{Corollary~1 of Proposition~2} then supplies two
equivalent nonzero subspaces contained respectively in
\(\mathfrak X_1\) and \(\mathfrak X_2\).
\end{proof}

% Source PDF page 227 (printed p. 218)
\phantomsection\label{srcpage:218}
\begin{theorem}\label{thm:III-1-1}
Let \(\mathcal A\) be a von Neumann algebra, let \(\mathcal Z\) be its
centre, and let \(E,F\) be projections of \(\mathcal A\). There exists a
projection \(G\in\mathcal Z\) such that
\[
  FG\prec EG,
  \qquad
  E(I-G)\prec F(I-G).
\]
\end{theorem}

\begin{proof}
By Zorn's theorem, there is a maximal family
\(((E_\iota,F_\iota))_{\iota\in I}\) of pairs of projections of
\(\mathcal A\) with the following properties: the \(E_\iota\)
(respectively, the \(F_\iota\)) are nonzero, pairwise orthogonal,
majorized by \(E\) (respectively, \(F\)), and
\(E_\iota\sim F_\iota\). Let
\[
  E^0=\sum_{\iota\in I}E_\iota\leq E,
  \qquad
  F^0=\sum_{\iota\in I}F_\iota\leq F.
\]
We have \(E^0\sim F^0\). Put \(E^1=E-E^0\), \(F^1=F-F^0\). By
\hyperref[lem:III-1-1]{Lemma~1} and the maximality of the family
\(((E_\iota,F_\iota))_{\iota\in I}\), the central supports of \(E^1\) and
\(F^1\) are orthogonal. There is therefore a projection
\(G\in\mathcal Z\) such that
\[
  E^1\leq G,\qquad F^1\leq I-G.
\]
With this choice, we have
\[
  E^0G\sim F^0G,
  \qquad
  E^0(I-G)\sim F^0(I-G),
\]
and hence
\[
  FG=F^0G\prec E^0G+E^1G=EG,
\]
and
\[
  E(I-G)=E^0(I-G)
  \prec F^0(I-G)+F^1(I-G)=F(I-G).
\]
\end{proof}

\setcounter{corollary}{0}
\begin{corollary}\label{cor:III-1-thm1-1}
Let \(\mathcal A\) be a factor and let \(E,F\) be two projections of
\(\mathcal A\). Either \(E\prec F\) or \(F\prec E\).
\end{corollary}

\begin{proof}
The projection \(G\) of \hyperref[thm:III-1-1]{Theorem~1} is either
\(0\) or \(I\).
\end{proof}

\begin{corollary}\label{cor:III-1-thm1-2}
Let \(\mathcal A\) be a von Neumann algebra, let \(\mathcal Z\) be its
centre, and let \((E_\iota)_{\iota\in I}\) be a family of nonzero,
pairwise orthogonal, equivalent projections of \(\mathcal A\). There are
a projection \(G\in\mathcal Z\) and a family
\((F_x)_{x\in K}\) of nonzero, pairwise orthogonal, equivalent projections
of \(\mathcal A\), majorized by \(G\), having the following properties:
\begin{enumerate}[label=\arabic*\textsuperscript{o}.,leftmargin=*]
\item \(I\subset K\);
\item for \(x\in I\), \(F_x\sim E_xG\);
\item if
\[
  F_0=G-\sum_{x\in K}F_x,
\]
then \(F_0\prec F_x\), and \(F_0\) is not equivalent to the \(F_x\).
Moreover, if \(I\) is infinite, one may suppose that \(F_0=0\).
\end{enumerate}
\end{corollary}

\begin{proof}
Let \((E_x)_{x\in K}\) be a maximal family of projections of
\(\mathcal A\) extending \((E_\iota)_{\iota\in I}\) (so that
\(I\subset K\)), in which the \(E_x\) are pairwise orthogonal and
equivalent. Put
\[
  E_0=I-\sum_{x\in K}E_x.
\]%
% Source PDF page 228 (printed p. 219)
\phantomsection\label{srcpage:219}%
Let \(G\in\mathcal Z\) be a projection such that
\[
  E_0G\prec E_xG,
  \qquad
  E_x(I-G)\prec E_0(I-G)
\]
(the choice of \(x\in K\) is immaterial, since the \(E_x\) are
equivalent). Put \(F_0=E_0G\), \(F_x=E_xG\). The relation
\(F_x\prec F_0\) would imply \(E_x\prec E_0\), contrary to the maximality
of the family \((E_x)_{x\in K}\). Thus, on the one hand \(F_x\ne0\), and
on the other hand \(F_0\) is not equivalent to the \(F_x\). We do indeed
have
\[
  F_0=E_0G=G-\sum_{x\in K}E_xG=G-\sum_{x\in K}F_x.
\]
Finally, if \(I\) is infinite, then \(K\) is infinite. Let \(K'\) be the
complement in \(K\) of an arbitrarily chosen element of \(K\); \(K'\) is
equipotent to \(K\), and
\[
  G=F_0+\sum_{x\in K}F_x
  \sim F_0+\sum_{x\in K'}F_x
  \prec\sum_{x\in K}F_x\leq G.
\]
Thus, by \hyperref[prop:III-1-1]{Proposition~1},
\(G\sim\sum_{x\in K}F_x\). By replacing the \(F_x\) by suitably chosen
equivalent projections, we see that we may suppose
\(G=\sum_{x\in K}F_x\).
\end{proof}

\begin{corollary}\label{cor:III-1-thm1-3}
Let \(\mathcal A\) be a von Neumann algebra. In order that
\(\mathcal A\) be continuous, it is necessary and sufficient that every
projection of \(\mathcal A\) be the sum of two orthogonal equivalent
projections.
\end{corollary}

\begin{proof}
If \(\mathcal A\) is not continuous, it has a nonzero abelian projection
\(E\). Two equivalent projections of \(\mathcal A\) majorized by \(E\)
have the same central support (\hyperref[no:III-1-1]{no.~1}), and hence
are equal (\hyperref[no:I-8-2]{Chapter~I, \S~8, no.~2}). Thus
\(\mathcal A\) does not have the property in the
corollary.

Now suppose that \(\mathcal A\) is continuous, and let \(E\) be a nonzero
projection of \(\mathcal A\). We shall show that \(E\) majorizes two
nonzero orthogonal equivalent projections \(E_1,E_2\) of \(\mathcal A\).
Applying the same result to \(E-(E_1+E_2)\), and continuing transfinitely,
we see that \(E\) will be the sum of two orthogonal equivalent projections
of \(\mathcal A\), and the proof will be complete. To prove our assertion,
we may replace \(\mathcal A\) by \(\mathcal A_E\), which is continuous.
It therefore suffices to prove this: if \(\mathcal A\) is continuous and
nonzero, there are two nonzero orthogonal equivalent projections of
\(\mathcal A\). Now \(\mathcal A\) is nonabelian. Let \(E_1\) be a
projection of \(\mathcal A\) not belonging to its centre and let
\(E'_1=I-E_1\). Applying \hyperref[thm:III-1-1]{Theorem~1} to
\(E_1,E'_1\), we obtain two
orthogonal projections \(E_2,E'_2\) of \(\mathcal A\) such that
\(0\ne E_2\prec E'_2\). Then \(E_2\sim E_3\leq E'_2\), and \(E_2,E_3\)
are orthogonal.
\end{proof}

Applying \hyperref[cor:III-1-thm1-2]{Corollary~2 of Theorem~1} to the
case of a factor, the reader will
see that this corollary plays the role of ``Euclidean division.''

\noindent\textit{Bibliography.}
\ArticleRef{6}, \ArticleRef{10}, \ArticleRef{42}, \ArticleRef{54},
\ArticleRef{65}, \ArticleRef{105}, \ArticleRef{114}, \ArticleRef{117}.

% Source PDF page 229 (printed p. 220)
\phantomsection\label{srcpage:220}
\NumberedParagraph{3}{Cyclic Projections of \(\mathcal A\) and Cyclic
Projections of \(\mathcal A'\)}
\label{no:III-1-3}

\begin{lemma}\label{lem:III-1-2}
Let \(\mathcal A\) be a von Neumann algebra on \(\mathfrak H\), and let
\((T_\iota)_{\iota\in I}\) be a family of elements of \(\mathcal A\).
There exists a \(T\in\mathcal A\) such that:
\begin{enumerate}[label=\arabic*\textsuperscript{o}.,leftmargin=*]
\item \(\lVert T_\iota T\rVert\leq1\) for every \(\iota\in I\);
\item if \(y\in\mathfrak H\) is such that
\(\sum_{\iota\in I}\lVert T_\iota y\rVert^2<+\infty\), then
\(y\in T(\mathfrak H)\).
\end{enumerate}
\end{lemma}

\begin{proof}
Let \(\mathfrak H'\) be a complex Hilbert space whose Hilbert dimension is
\(\operatorname{Card}I+1\). Let \(\mathcal B\) be the algebra
\(\mathcal A\mathbin{\bar\otimes}\mathcal L(\mathfrak H')\) on
\(\mathfrak H\otimes\mathfrak H'\). By means of an orthonormal basis of
\(\mathfrak H'\), identify
\[
  \mathfrak H\otimes\mathfrak H'
  \quad\text{with}\quad
  \mathfrak H\oplus\Bigl(\bigoplus_{\iota\in I}\mathfrak H_\iota\Bigr),
\]
where \(\mathfrak H\) and the \(\mathfrak H_\iota\) are pairwise
orthogonal subspaces of \(\mathfrak H\otimes\mathfrak H'\). For every
\(\iota\in I\), there is a partially isometric operator
\(U_\iota\in\mathcal B\) having \(\mathfrak H\) and
\(\mathfrak H_\iota\) as initial and final subspaces. The algebra
\(\mathcal A\) is identified with \(\mathcal B_{\mathfrak H}\)
(\hyperref[prop:I-2-5]{Chapter~I, \S~2, Proposition~5}).

Let \(\mathfrak X\) be the set of the
\(x\in\mathfrak H\otimes\mathfrak H'\) such that
\[
  P_{\mathfrak H_\iota}x
  =U_\iota T_\iota P_{\mathfrak H}x
  \qquad(\iota\in I).
\]
Plainly, \(\mathfrak X\) is a closed vector subspace of
\(\mathfrak H\otimes\mathfrak H'\), invariant under every unitary
operator of \(\mathcal B'\). Hence \(P_{\mathfrak X}\in\mathcal B\).
Moreover, the conditions \(x\in\mathfrak X\) and
\(P_{\mathfrak H}x=0\) imply
\(P_{\mathfrak H_\iota}x=0\) for every \(\iota\in I\), hence \(x=0\).
\hyperref[cor:III-1-prop2-1]{Corollary~1 of Proposition~2} therefore
shows that
\(P_{\mathfrak X}\prec P_{\mathfrak H}\), relative to \(\mathcal B\).
There is a partially isometric operator \(U\in\mathcal B\) whose initial
projection is majorized by \(P_{\mathfrak H}\) and whose final projection
is \(P_{\mathfrak X}\). Let \(\widetilde T\) be the restriction to
\(\mathfrak H\) of \(P_{\mathfrak H}U\). Then
\(\widetilde T\in\mathcal B_{\mathfrak H}=\mathcal A\), and
\(\widetilde T(\mathfrak H)=P_{\mathfrak H}(\mathfrak X)\), so that
\(\widetilde T(\mathfrak H)\) is the set of those
\(y\in\mathfrak H\) for which
\[
  \sum_{\iota\in I}\lVert T_\iota y\rVert^2
  =\sum_{\iota\in I}\lVert U_\iota T_\iota y\rVert^2
  <+\infty.
\]
It follows that, for every \(z\in\mathfrak H\), the numbers
\(\lVert T_\iota\widetilde Tz\rVert\) are bounded. The numbers
\(\lVert T_\iota\widetilde T\rVert\) are therefore bounded
(\TreatiseRef{3}, Chapter~III, \S~3, Theorem~2). Multiplying
\(\widetilde T\) by a suitable scalar gives the operator \(T\) of the
lemma.
\end{proof}

\begin{lemma}\label{lem:III-1-3}
Let \(\mathcal A\) be a von Neumann algebra on \(\mathfrak H\), let
\(x\in\mathfrak H\), and let
\(y\in\mathfrak X_x^{\mathcal A}\). There exist \(z\in\mathfrak H\) and
\(S,T\in\mathcal A\) such that
\[
  x=Tz,\qquad y=Sz,\qquad z=Ez,
\]
where \(E\) denotes the support of \(T\).
\end{lemma}

\begin{proof}
Let \((S_1,S_2,\ldots)\) be a sequence of elements of \(\mathcal A\) such
that \(\lVert S_nx-y\rVert\leq4^{-n}\). Put
\(T_n=2^n(S_{n+1}-S_n)\), and let \(T\in\mathcal A\) have the properties
of \hyperref[lem:III-1-2]{Lemma~2} with respect to the family \((T_n)\).
We have
\[
  \lVert(S_{n+1}-S_n)T\rVert\leq2^{-n},
\]
so the operators \(S_nT\) converge in norm to an operator
\(S\in\mathcal A\). On the other hand,
\[
  \sum_{n=1}^{\infty}\lVert T_nx\rVert^2
  =\sum_{n=1}^{\infty}4^n\lVert S_{n+1}x-S_nx\rVert^2
  \leq\sum_{n=1}^{\infty}4^n(2\cdot4^{-n})^2<+\infty.
\]
% Source PDF page 230 (printed p. 221)
\phantomsection\label{srcpage:221}
Thus there is a \(z\in\mathfrak H\) such that \(Tz=x\), and we may plainly
suppose that \(z\) is orthogonal to the null space of \(T\), that is,
\(Ez=z\), where \(E\) is the support of \(T\). Finally, the equality
\(S_nTz=S_nx\) gives \(Sz=y\) on passing to the limit.
\end{proof}

\begin{theorem}\label{thm:III-1-2}
Let \(\mathcal A\) be a von Neumann algebra on \(\mathfrak H\), and let
\(x,x_1\in\mathfrak H\). The relations
\[
  E_{x_1}^{\mathcal A}\prec E_x^{\mathcal A},
  \qquad
  E_{x_1}^{\mathcal A'}\prec E_x^{\mathcal A'}
\]
are equivalent.
\end{theorem}

\begin{proof}
If \(E_{x_1}^{\mathcal A}\prec E_x^{\mathcal A}\), there are a
\(y\in\mathfrak X_x^{\mathcal A}\) and a (partially isometric) operator
\(U'\in\mathcal A'\) such that \(x_1=U'y\). Apply
\hyperref[lem:III-1-3]{Lemma~3} to
\(\mathcal A,x,y\), and let \(z\) be the element of \(\mathfrak H\) thus
obtained. By \hyperref[cor:III-1-prop2-2]{Corollary~2 of Proposition~2},
with the roles of
\(\mathcal A\) and \(\mathcal A'\) interchanged, we have
\[
  E_{x_1}^{\mathcal A'}\leq E_y^{\mathcal A'}
  \prec E_z^{\mathcal A'}\sim E_x^{\mathcal A'}.
\]
Similarly,
\(E_{x_1}^{\mathcal A'}\prec E_x^{\mathcal A'}\) implies
\(E_{x_1}^{\mathcal A}\prec E_x^{\mathcal A}\).
\end{proof}

\begin{corollary*}\phantomsection\label{cor:III-1-thm2}
The relations
\(E_{x_1}^{\mathcal A}\sim E_x^{\mathcal A}\) and
\(E_{x_1}^{\mathcal A'}\sim E_x^{\mathcal A'}\) are equivalent.
\end{corollary*}

\begin{proof}
Apply \hyperref[thm:III-1-2]{Theorem~2} and
\hyperref[prop:III-1-1]{Proposition~1}.
\end{proof}

\noindent\textit{Bibliography.} \ArticleRef{65}.

\NumberedParagraph{4}{Applications: I. Properties of Totalizing and
Separating Elements}
\label{no:III-1-4}

\begin{lemma}\label{lem:III-1-4}
Let \(\mathcal A\) be a von Neumann algebra, let \(E,F\) be projections of
\(\mathcal A\), and let \(E',F'\) be their central supports. Suppose that
\(F'\leq E'\). If there is a totalizing element \(x\) for
\(\mathcal A_E\) and a separating element \(y\) for \(\mathcal A_F\),
then \(F\prec E\).
\end{lemma}

\begin{proof}
By identifying \(\mathcal A\) suitably with a product of two von Neumann
algebras, \hyperref[thm:III-1-1]{Theorem~1} allows us to consider only
the case \(E\prec F\).
Let \(E_1\) be a projection of \(\mathcal A\) such that
\(E\sim E_1\leq F\). There is a totalizing element for
\(\mathcal A_{E_1}\). We may therefore restrict ourselves to the case
\(E\leq F\). We then have \(E'=F'\). The element \(x\) is separating for
\(\mathcal A'_E\), hence separating for \(\mathcal A'_F\) (because the
inductions of \(\mathcal A'_{E'}\) on \(\mathcal A'_E\) and
\(\mathcal A'_F\) are isomorphisms), and hence totalizing for
\(\mathcal A_F\). Thus
\[
  \mathfrak X_x^{\mathcal A_F}=F(\mathfrak H)
  \supset\mathfrak X_y^{\mathcal A_F},
\]
and consequently, by \hyperref[thm:III-1-2]{Theorem~2},
\(\mathfrak X_x^{\mathcal A'_F}\succ
\mathfrak X_y^{\mathcal A'_F}\). But \(y\) is totalizing for
\(\mathcal A'_F\), so
\(\mathfrak X_y^{\mathcal A'_F}=F(\mathfrak H)\). On the other hand,
\(E(\mathfrak H)\supset\mathfrak X_x^{\mathcal A'_F}\). Thus
\(E(\mathfrak H)\succ F(\mathfrak H)\).
\end{proof}

\begin{proposition}\label{prop:III-1-4}
Let \(x\) be a separating element for \(\mathcal A\). If there exists a
totalizing element for \(\mathcal A\), then
\(E_x^{\mathcal A}\sim I\).
\end{proposition}

\begin{proof}
Put \(E'=E_x^{\mathcal A}\in\mathcal A'\). By
\hyperref[cor:I-1-7-2]{Chapter~I, \S~1, Proposition~7, Corollary~2},
\(E'\) has central support \(I\). Moreover,
\(x\) is totalizing for \(\mathcal A_{E'}\).
% Source PDF page 231 (printed p. 222)
\phantomsection\label{srcpage:222}
Hence, by \hyperref[lem:III-1-4]{Lemma~4}, \(I\prec E'\), and
consequently, by \hyperref[prop:III-1-1]{Proposition~1},
\(I\sim E'\).
\end{proof}

\begin{corollary*}\phantomsection\label{cor:III-1-prop4}
If there are a totalizing element and a separating element for
\(\mathcal A\), then there is an element which is both totalizing and
separating for \(\mathcal A\).
\end{corollary*}

\begin{proof}
Let \(x\) be a separating element for \(\mathcal A\). By
\hyperref[prop:III-1-4]{Proposition~4},
there is a partially isometric operator \(U'\in\mathcal A'\) having
\(E'=E_x^{\mathcal A}\) as initial projection and \(I\) as final
projection. Since \(x\) is separating and totalizing for
\(\mathcal A_{E'}\), \(U'x\) is separating and totalizing for
\(\mathcal A\).
\end{proof}

\begin{theorem}\label{thm:III-1-3}
Let \(\mathcal A\) (respectively, \(\mathcal A_1\)) be a von Neumann
algebra, and suppose that there is a totalizing and separating element for
\(\mathcal A\) (respectively, \(\mathcal A_1\)). Then every isomorphism
\(\Phi\) of \(\mathcal A\) onto \(\mathcal A_1\) is spatial.
\end{theorem}

\begin{proof}
There exist, by
\hyperref[cor:I-4-3]{Chapter~I, \S~4, the corollary to Theorem~3}, a complex
Hilbert space \(\mathfrak K\), a von Neumann algebra \(\mathcal C\) on
\(\mathfrak K\), and projections \(E',F'\) of \(\mathcal C'\), with
central support \(I\), such that \(\mathcal A\) is identified with
\(\mathcal C_{E'}\), \(\mathcal A_1\) with \(\mathcal C_{F'}\), and
\(\Phi\) with the isomorphism
\(T_{E'}\mapsto T_{F'}\) (\(T\in\mathcal C\)). By the hypothesis of the
theorem and \hyperref[lem:III-1-4]{Lemma~4}, \(E'\sim F'\). Let \(U'\)
be a partially isometric
operator of \(\mathcal C'\) having \(E'\) as initial projection and \(F'\)
as final projection. For every \(T\in\mathcal C\), the restriction of
\(U'\) to \(E'(\mathfrak K)\) transforms \(T_{E'}\) into \(T_{F'}\).
Hence \(\Phi\) is spatial.
\end{proof}

We shall encounter on several occasions
(\hyperref[thm:III-1-6]{Theorem~6};
\hyperref[prop:III-3-3]{\S~3, Proposition~3};
\hyperref[prop:III-6-10]{\S~6, Proposition~10}; \S~8, Corollaries~8 and~9 of
Theorem~1) conditions ensuring that an isomorphism is spatial.

\begin{theorem}\label{thm:III-1-4}
Let \(\mathcal A\) be a von Neumann algebra on \(\mathfrak H\) possessing
a separating element, and let \(\varphi\) be a normal positive linear
form on \(\mathcal A\). There exists an \(x\in\mathfrak H\) such that
\(\varphi=\omega_x\).
\end{theorem}

\begin{proof}
Let \(z\) be a separating element for \(\mathcal A\), and put
\(E'=E_z^{\mathcal A}\). The central support of \(E'\) is \(I\) (because
\(E_z^{\mathcal A'}=I\)), so induction of \(\mathcal A\) on
\(\mathcal A_{E'}\) is an isomorphism and allows \(\varphi\) to be
transported to \(\mathcal A_{E'}\). It suffices to prove the theorem for
\(\mathcal A_{E'}\) and this transported form. We may therefore suppose
henceforth that \(z\) is totalizing for \(\mathcal A\). If the theorem is
true for \(\varphi+\omega_z\), it is true for \(\varphi\), which is
majorized by \(\varphi+\omega_z\)
(\hyperref[lem:I-4-1]{Chapter~I, \S~4, Lemma~1}). We may
therefore suppose henceforth that \(\varphi\geq\omega_z\), and hence that
\(\varphi\) is faithful. Then, by
\hyperref[lem:I-4-4]{Chapter~I, \S~4, Lemma~4} and
\hyperref[prop:I-4-1]{Proposition~1}, there exist a Hilbert space
\(\mathfrak K\), a von Neumann
algebra \(\mathcal B\) on \(\mathfrak K\), a totalizing and separating
element \(y\) for \(\mathcal B\), and an isomorphism \(\Phi\) of
\(\mathcal A\) onto \(\mathcal B\), such that
\[
  \varphi(T)=(\Phi(T)y\mid y)\qquad(T\in\mathcal A).
\]
By \hyperref[thm:III-1-3]{Theorem~3}, \(\Phi\) is spatial.
\end{proof}

% Source PDF page 232 (printed p. 223)
\phantomsection\label{srcpage:223}
\begin{corollary*}\phantomsection\label{cor:III-1-thm4}
Let \(\mathcal A\) be a von Neumann algebra. The following conditions are
equivalent:
\begin{enumerate}[label=\textup{(\roman*)},leftmargin=*]
\item Every normal positive linear form on \(\mathcal A\) is a form
\(\omega_x\).
\item For every projection \(E\) of \(\mathcal A\) such that
\(\mathcal A_E\) is of countable type, \(\mathcal A_E\) has a separating
element.
\end{enumerate}
These conditions hold, in particular, when \(\mathcal A\) is abelian and
when \(\mathcal A\) is standard.
\end{corollary*}

\begin{proof}
If \(\mathcal A\) has property~\textup{(i)}, it is clear that, for every
projection \(E\) of \(\mathcal A\), \(\mathcal A_E\) has
property~\textup{(i)}. If, in addition, \(\mathcal A_E\) is of countable
type, there exists on \(\mathcal A_E\) a faithful normal positive linear
form. This form is an \(\omega_x\), and \(x\) is separating for
\(\mathcal A_E\).

Suppose condition~\textup{(ii)} holds. Let \(\varphi\) be a normal
positive linear form on \(\mathcal A\), and let \(E\) be its support.
Then \(\mathcal A_E\) is of countable type, and hence has a separating
element. Applying \hyperref[thm:III-1-4]{Theorem~4} to \(\mathcal A_E\),
there is an
\(x\in E(\mathfrak H)\) such that
\(\varphi(T)=(Tx\mid x)\) for every \(T\in E\mathcal A E\). Thus, for
every \(T\in\mathcal A\),
\[
  \varphi(T)=\varphi(ETE)=(ETEx\mid x)=(Tx\mid x).
\]

If \(\mathcal A\) is abelian, it satisfies condition~\textup{(ii)} by
Chapter~I, \S~2, the corollary to Proposition~3. Suppose that
\(\mathcal A\) is standard, and let \(E\) be a projection of
\(\mathcal A\) such that \(\mathcal A_E\) is of countable type. We show
that \(\mathcal A_E\) has a separating element. By
\hyperref[prop:I-2-3]{Chapter~I, \S~2, Proposition~3}, applied to
\(\mathcal A_E\), it suffices to consider the
case in which \(\mathcal A_E\) has a totalizing element \(x\). Let \(F\)
be the central support of \(E\). The element \(x\) is separating for
\(\mathcal A'_E\), hence for \(\mathcal A'_F\). Since
\(\mathcal A'_F\) is standard
(\hyperref[no:I-5-5]{Chapter~I, \S~5, no.~5}),
\(\mathcal A_F\) has a separating element \(y\). Then \(Ey\) is
separating for \(\mathcal A_E\).
\end{proof}

See, on this subject, \hyperref[no:III-6-3]{\S~6, no.~3}, and \S~8,
Corollaries~10 and~11 of
Theorem~1.

When a von Neumann algebra \(\mathcal A\) has a totalizing and separating
element, \(\mathcal A\) and \(\mathcal A'\) are, in certain respects,
``equally large''; this idea will be made precise in
\hyperref[sec:III-6]{\S~6} (and already to
some extent in the next number).

\noindent\textit{Bibliography.}
\ArticleRef{17}, \ArticleRef{19}, \ArticleRef{31}, \ArticleRef{40},
\ArticleRef{62}, \ArticleRef{65}, \ArticleRef{66}, \ArticleRef{67},
\ArticleRef{70}, \ArticleRef{89}, \ArticleRef{100}, \ArticleRef{119}.

\NumberedParagraph{5}{Applications: II. Characterization of Standard von
Neumann Algebras}
\label{no:III-1-5}

\begin{lemma}\label{lem:III-1-5}
Let \(\mathcal A\) be a von Neumann algebra of countable type, and let
\(\mathcal Z\) be its centre. If there is an involution \(J\) of
\(\mathfrak H\), commuting with the projections of \(\mathcal Z\), such
that \(J\mathcal A J=\mathcal A'\), then \(\mathcal A\) has a totalizing
and separating element.
\end{lemma}

\begin{proof}
By \hyperref[prop:I-2-3]{Chapter~I, \S~2, Proposition~3}, we have
\(\mathcal A=\mathcal A_1\times\mathcal A_2\), where
\(\mathcal A_1\) has a totalizing element and \(\mathcal A_2\) has a
separating element. But \(J\),
% Source PDF page 233 (printed p. 224)
\phantomsection\label{srcpage:224}
commuting with the projections of \(\mathcal Z\), defines involutions
which interchange \(\mathcal A_1\) and
\(\mathcal A'_1\), and \(\mathcal A_2\) and \(\mathcal A'_2\).
Consequently, \(\mathcal A_1\) has a separating element and
\(\mathcal A_2\) has a totalizing element. Thus \(\mathcal A\) has a
totalizing element and a separating element, and hence, by the
\hyperref[cor:III-1-prop4]{corollary to Proposition~4}, a totalizing and separating
element.
\end{proof}

\begin{theorem}\label{thm:III-1-5}
Let \(\mathcal A\) be a von Neumann algebra of countable type. In order
that \(\mathcal A\) be standard, it is necessary and sufficient that
\(\mathcal A\) be semifinite and have a totalizing and separating element.
\end{theorem}

\begin{proof}
The conditions are necessary by \hyperref[lem:III-1-5]{Lemma~5} and
\hyperref[cor:I-6-9]{Chapter~I, \S~6, the corollary to Proposition~9}.
Conversely, suppose that \(\mathcal A\) is
semifinite and has a totalizing and separating element. By
\hyperref[cor:I-6-9]{Chapter~I, \S~6, the corollary to Proposition~9},
there is an isomorphism \(\Phi\) of
\(\mathcal A\) onto a standard von Neumann algebra \(\mathcal B\). Since
\(\mathcal A\) is of countable type, so is \(\mathcal B\), and hence
\(\mathcal B\) has a totalizing and separating element by
\hyperref[lem:III-1-5]{Lemma~5}. Thus \(\Phi\) is spatial by
\hyperref[thm:III-1-3]{Theorem~3}, and \(\mathcal A\) is standard.
\end{proof}

\begin{corollary*}\phantomsection\label{cor:III-1-thm5}
Let \(\mathcal A\) be a finite von Neumann algebra having a totalizing and
separating element. Then every separating element for \(\mathcal A\) is
totalizing for \(\mathcal A\), and every totalizing element for
\(\mathcal A\) is separating for \(\mathcal A\).
\end{corollary*}

\begin{proof}
By \hyperref[thm:III-1-5]{Theorem~5}, \(\mathcal A'\) is anti-isomorphic
to \(\mathcal A\), and
hence is finite. The two assertions of the corollary are therefore
equivalent, and follow at once from
\hyperref[prop:III-1-3]{Proposition~3} and
\hyperref[prop:III-1-4]{Proposition~4}.
\end{proof}

We shall now seek to free ourselves from the countability hypothesis in
\hyperref[thm:III-1-5]{Theorem~5}.

\begin{lemma}\label{lem:III-1-6}
Let \(\mathcal A\) be a von Neumann algebra, let
\((E_\iota)_{\iota\in I}\) be an infinite family of projections of
\(\mathcal A\), with supremum \(I\), such that the
\(\mathcal A_{E_\iota}\) are of countable type, and let
\((F_x)_{x\in K}\) be a family of nonzero pairwise orthogonal projections
of \(\mathcal A\). Then
\[
  \operatorname{Card}K\leq\operatorname{Card}I.
\]
\end{lemma}

\begin{proof}
Since \(\mathcal A_{E_\iota}\) is of countable type, there is a
countable set \(\mathfrak M_\iota\) such that
\(E_\iota=E_{\mathfrak M_\iota}^{\mathcal A}\). Let \(K_\iota\) be the
(countable) subset of \(K\) consisting of those \(x\in K\) for which
\(F_x(\mathfrak M_\iota)\ne0\). We have
\(K=\bigcup_{\iota\in I}K_\iota\), for if some \(x\in K\) satisfied
\(F_x(\mathfrak M_\iota)=0\) for every \(\iota\in I\), then \(F_x\)
would be orthogonal to all the \(E_\iota\), and hence would be zero. Thus
\[
  \operatorname{Card}K
  \leq(\operatorname{Card}I)\aleph_0
  =\operatorname{Card}I.
\]
\end{proof}

\begin{lemma}\label{lem:III-1-7}
Every von Neumann algebra \(\mathcal A\) is a product of von Neumann
algebras \(\mathcal B\) having one or the other of the following two
properties:
\begin{enumerate}[label=\textup{(\roman*)},leftmargin=*]
\item \(\mathcal B\) is of countable type;
\item there is an uncountable family
% Source PDF page 234 (printed p. 225)
\phantomsection\label{srcpage:225}
\((E_\iota)_{\iota\in I}\) of equivalent, pairwise orthogonal projections
of \(\mathcal B\), with sum \(I\), such that the
\(\mathcal B_{E_\iota}\) are of countable type.
\end{enumerate}
\end{lemma}
\begin{proof}
Let \(E\) be a nonzero projection of \(\mathcal A\) such that
\(\mathcal A_E\) is of countable type (for example, a nonzero cyclic
projection). By decreasing \(E\), if necessary, we may suppose, by
\hyperref[cor:III-1-thm1-2]{Theorem~1, Corollary~2}, that there are a family
\((E_x)_{x\in K}\) of pairwise orthogonal projections of
\(\mathcal A\), equivalent to \(E\), and a projection \(G\) of the centre
of \(\mathcal A\), majorizing the \(E_x\), such that
\[
  F=G-\sum_{x\in K}E_x\prec E.
\]
The algebra \(\mathcal A_F\) is of countable type. If \(K\) is countable,
there is a countable totalizing set for \(\mathcal A'_G\), and hence
\(\mathcal A_G\) is of countable type. If \(K\) is uncountable, we may,
by \hyperref[cor:III-1-thm1-2]{Theorem~1, Corollary~2}, suppose that
\(F=0\). In both cases,
\(\mathcal A_G\) has property~\textup{(i)} or property~\textup{(ii)}
stated in the lemma for \(\mathcal B\). It remains only to repeat the
argument for \(\mathcal A_{I-G}\), and so on transfinitely.
\end{proof}

\begin{theorem}\label{thm:III-1-6}
Let \(\mathcal A\) be a von Neumann algebra on \(\mathfrak H\), and let
\(\mathcal Z\) be its centre. Suppose that there is an involution \(J\) of
\(\mathfrak H\) such that
\[
  J\mathcal A J=\mathcal A',
  \qquad
  JCJ=C^*\quad(C\in\mathcal Z).
\]
Let \(\mathcal A_1,\mathcal Z_1,J_1\), on another Hilbert space
\(\mathfrak H_1\), have the same properties. Then every isomorphism
\(\Phi\) of \(\mathcal A\) onto \(\mathcal A_1\) is spatial.
\end{theorem}

\begin{proof}
If \(\mathcal A=\prod_{\iota\in I}\mathcal A_\iota\), the
\(\mathcal A_\iota\) satisfy the same hypotheses as \(\mathcal A\)
(because \(J\) commutes with the projections of \(\mathcal Z\)). On the
other hand, if \(\mathcal A\) is of countable type,
\hyperref[thm:III-1-6]{Theorem~6} follows from
\hyperref[lem:III-1-5]{Lemma~5} and \hyperref[thm:III-1-3]{Theorem~3}.
By \hyperref[lem:III-1-7]{Lemma~7}, we may therefore restrict
ourselves to the case in which \(\mathcal A\) contains an infinite family
of pairwise orthogonal equivalent projections, with sum \(I\), for which
the corresponding reduced algebras are of countable type. There are then
analogous families in \(\mathcal A',\mathcal A_1,\mathcal A'_1\), these
four families being equipotent.

By \hyperref[cor:I-4-3]{Chapter~I, \S~4, Theorem~3 and its corollary},
there are a von Neumann
algebra \(\mathcal B\) and projections
\(P_{\mathfrak X},P_{\mathfrak Y}\) of \(\mathcal B'\), with central
support \(I\), such that \(\mathcal A\) is identified with
\(\mathcal B_{\mathfrak X}\), \(\mathcal A_1\) with
\(\mathcal B_{\mathfrak Y}\), and \(\Phi\) with the isomorphism
\(T_{\mathfrak X}\mapsto T_{\mathfrak Y}\) (\(T\in\mathcal B\)).
As in the proof of \hyperref[thm:III-1-3]{Theorem~3}, it is enough to show
that \(\mathfrak X\sim\mathfrak Y\).
\hyperref[thm:III-1-1]{Theorem~1} allows us to restrict ourselves
to the cases \(\mathfrak X\prec\mathfrak Y\) and
\(\mathfrak Y\prec\mathfrak X\). Suppose, for example, that
\(\mathfrak X\prec\mathfrak Y\). We are immediately reduced to the case
\(\mathfrak X\subset\mathfrak Y\). There is an infinite family
\((P_{\mathfrak X_\iota})_{\iota\in I}\) [respectively,
\((P_{\mathfrak Y_\iota})_{\iota\in I}\)] of pairwise orthogonal
equivalent projections of \(\mathcal B'\), with sum
\(P_{\mathfrak X}\) (respectively, \(P_{\mathfrak Y}\)), such that the
\(\mathcal A'_{\mathfrak X_\iota}\) (respectively, the
\(\mathcal A'_{\mathfrak Y_\iota}\)) are of countable type. Moreover, we
may add to the family \((P_{\mathfrak X_\iota})_{\iota\in I}\) nonzero
projections \(\prec P_{\mathfrak X}\) so that the resulting family
\((P_{\mathfrak X_x})_{x\in K}\) consists of pairwise
% Source PDF page 235 (printed p. 226)
\phantomsection\label{srcpage:226}
orthogonal projections of \(\mathcal B'\) with sum \(P_{\mathfrak Y}\)
(as follows
at once by applying \hyperref[thm:III-1-1]{Theorem~1} transfinitely).
\hyperref[lem:III-1-6]{Lemma~6}, applied to
\(\mathcal B'_{\mathfrak Y}\) and the families
\((\mathfrak X_x)_{x\in K}\), \((\mathfrak Y_\iota)_{\iota\in I}\),
shows that \(I\) and \(K\) are equipotent. Hence
\[
  P_{\mathfrak Y}=\sum_{x\in K}P_{\mathfrak X_x}
  \prec\sum_{\iota\in I}P_{\mathfrak X_\iota}
  =P_{\mathfrak X},
\]
and finally \(\mathfrak Y\sim\mathfrak X\).
\end{proof}

\begin{corollary*}\phantomsection\label{cor:III-1-thm6}
Let \(\mathcal A\) be a von Neumann algebra on \(\mathfrak H\), and let
\(\mathcal Z\) be its centre. In order that \(\mathcal A\) be standard,
it is necessary and sufficient that \(\mathcal A\) be semifinite and that
there exist on \(\mathfrak H\) an involution \(J\) with the following
properties:
\begin{enumerate}[label=\textup{(\roman*)},leftmargin=*]
\item \(J\mathcal A J=\mathcal A'\);
\item \(JCJ=C^*\) for every \(C\in\mathcal Z\).
\end{enumerate}
\end{corollary*}

\begin{proof}
The conditions are necessary by Chapter~I, \S~5, the corollary to
Proposition~2, and
\hyperref[cor:I-6-9]{\S~6, the corollary to Proposition~9}. Now suppose that
\(\mathcal A\) is semifinite. By
\hyperref[cor:I-6-9]{Chapter~I, \S~6, the corollary to Proposition~9},
there is an isomorphism \(\Phi\) of \(\mathcal A\) onto a
standard von Neumann algebra \(\mathcal A_1\), which satisfies the
conditions of the corollary. If \(\mathcal A\) also satisfies them, then
\(\Phi\) is spatial by \hyperref[thm:III-1-6]{Theorem~6}, and therefore
\(\mathcal A\) is
standard.
\end{proof}

The conclusion of the \hyperref[cor:III-1-thm5]{corollary to Theorem~5}
does not necessarily remain
valid if the hypothesis that \(\mathcal A\) is finite is suppressed
(\hyperref[ex:III-1-6]{Exercise~6}).

\hyperref[lem:III-1-6]{Lemma~6} plainly makes it possible to attach to
every von Neumann algebra
a cardinal which is invariant under algebraic isomorphism.

The conditions of \hyperref[thm:III-1-6]{Theorem~6} are satisfied by
certain von Neumann algebras
which are not necessarily standard; see \ArticleRef{13}.

\hyperref[thm:III-1-6]{Theorem~6} generalizes
\hyperref[thm:I-6-4]{Theorem~4 of Chapter~I, \S~6}.

\noindent\textit{Bibliography.}
\ArticleRef{17}, \ArticleRef{66}, \ArticleRef{101}.

\medskip
\noindent\textit{Exercises.}
\begin{enumerate}[label=\arabic*.,leftmargin=*]
\item\label{ex:III-1-1} Let \(\mathcal A\) be a von Neumann algebra and
let \(E,F\) be two projections of \(\mathcal A\). In order that
\(E\prec F\), it is necessary and sufficient that there exist a
projection \(G\in\mathcal A\), with \(G\geq E\), such that \(G\sim F\).
(Let \(E\sim E_1\leq F\). Apply the comparability theorem to \(F-E_1\)
and \(I-E\). If \(F-E_1\prec I-E\), the proposition is evident. If
\(I-E\prec F-E_1\), then \(F\sim I\), and one may take \(G=I\).)
\ArticleRef{65}.

\item\label{ex:III-1-2} Let \(\mathcal A\) be a von Neumann algebra, and
let \(P_{\mathfrak X},P_{\mathfrak Y}\) be two projections of
\(\mathcal A\).
\begin{enumerate}[label=\textit{\alph*.},leftmargin=*]
\item Show that
\[
  P_{\overline{\mathfrak X+\mathfrak Y}}-P_{\mathfrak Y}
  \sim
  P_{\mathfrak X}-P_{\mathfrak X\cap\mathfrak Y}.
\]
(Use \hyperref[cor:III-1-prop2-1]{Corollary~1 of Proposition~2}.)
\item Deduce that, if \(\varphi\) is a trace on \(\mathcal A_+\), then
\[
  \varphi(P_{\mathfrak X})+\varphi(P_{\mathfrak Y})
  =\varphi(P_{\overline{\mathfrak X+\mathfrak Y}})
   +\varphi(P_{\mathfrak X\cap\mathfrak Y}).
\]
\ArticleRef{65}.
\end{enumerate}

\item\label{ex:III-1-3} Let \(\mathcal A\) be a finite von Neumann
algebra, and let \(S,T\in\mathcal A\) satisfy \(ST=I\). Then \(TS=I\).
(The support of \(T\) is \(I\), and hence, by
\hyperref[prop:III-1-2]{Proposition~2} and
\hyperref[prop:III-1-3]{Proposition~3}, so is
the support of \(T^*\). Show that \(T\) is invertible.)

This property characterizes finite algebras; see \S~8, Theorem~1.

\item\label{ex:III-1-4} Give all the details of the following argument,
which is a new proof of \hyperref[lem:III-1-1]{Lemma~1}. Let
\(\mathcal A\) be a von Neumann
algebra, and let \(E_1,E_2\) be two projections of \(\mathcal A\) with
central support \(I\). Since the mapping
\(T'_{E_1}\mapsto T'_{E_2}\) (\(T'\in\mathcal A'\)) is an isomorphism of
\(\mathcal A'_{E_1}\) onto \(\mathcal A'_{E_2}\), normal positive linear
forms \(\varphi_1\) on \(\mathcal A'_{E_1}\) and \(\varphi_2\) on
\(\mathcal A'_{E_2}\) correspond under this isomorphism;
% Source PDF page 236 (printed p. 227)
\phantomsection\label{srcpage:227}
by decreasing \(\varphi_1,\varphi_2\), if necessary, one may suppose
\(\varphi_1=\omega_{x_1}\), \(\varphi_2=\omega_{x_2}\), with
\(x_1\ne0\), \(x_2\ne0\). Put
\(\mathfrak X_1=\mathfrak X_{x_1}^{\mathcal A'}\) and
\(\mathfrak X_2=\mathfrak X_{x_2}^{\mathcal A'}\). By
\hyperref[lem:I-4-3]{Chapter~I, \S~4, Lemma~3}, there is an isomorphism
of the Hilbert space \(\mathfrak X_1\)
onto the Hilbert space \(\mathfrak X_2\) which transforms
\(T'_{\mathfrak X_1}\) into \(T'_{\mathfrak X_2}\) for every
\(T'\in\mathcal A'\). Conclude that
\(P_{\mathfrak X_1}\sim P_{\mathfrak X_2}\).

\item\label{ex:III-1-5} Let \(\mathcal A\) be a von Neumann algebra,
let \(\varphi\) be a normal positive linear form on \(\mathcal A\), and
let \(E\) be a nonzero projection of \(\mathcal A\). Show that there is a
nonzero projection \(F\in\mathcal A\), majorized by \(E\), such that the
form \(T\mapsto\varphi(FTF)\), \(T\in\mathcal A\), is a form
\(\omega_x\). (Choose \(F\) of the form \(EE_y^{\mathcal A'}\), and apply
\hyperref[thm:III-1-4]{Theorem~4}.)

\item\label{ex:III-1-6} Let \(\mathfrak K\) be a complex Hilbert space
which is the Hilbert sum of subspaces
\(\mathfrak K_1,\allowbreak \mathfrak K_2,\allowbreak \ldots\), let
\(\mathfrak H\) be another
complex Hilbert space, and let \(U_n\) be an isometry of
\(\mathfrak H\) onto \(\mathfrak K_n\) (\(n=1,2,\ldots\)). These data
define an amplification of \(\mathcal L(\mathfrak H)\) onto a von Neumann
algebra \(\mathcal A\) on \(\mathfrak K\). Let
\((e_1,e_2,\ldots)\) be an orthonormal basis of \(\mathfrak H\), and put
\(e_{np}=U_ne_p\). Show that
\[
  x=e_{11}+\frac12e_{22}+\frac13e_{33}+\cdots
\]
is totalizing and separating for \(\mathcal A\), whereas
\[
  y=e_{21}+\frac12e_{32}+\frac13e_{43}+\cdots
\]
is separating without being totalizing for \(\mathcal A\).

\item\label{ex:III-1-7} Let \(\mathfrak H_1,\mathfrak H_2\) be two
two-dimensional complex Hilbert spaces, and put
\[
  \mathfrak H=\mathfrak H_1\oplus\mathfrak H_2,
  \qquad
  \mathcal A=\mathcal L(\mathfrak H_1)\times\mathbb C_{\mathfrak H_2}.
\]
Then
\[
  \mathcal A'=\mathbb C_{\mathfrak H_1}\times\mathcal L(\mathfrak H_2).
\]
There is an involution of \(\mathfrak H\) which transforms
\(\mathcal A\) into \(\mathcal A'\), but \(\mathcal A\) is not standard.

\item\label{ex:III-1-8} Let \(\mathcal A\) be a von Neumann algebra and
let \(\mathfrak m\) be a two-sided ideal of \(\mathcal A\).
\begin{enumerate}[label=\textit{\alph*.},leftmargin=*]
\item Let \(E\) be a projection of \(\mathfrak m\). If \(F\prec E\), then
\(F\in\mathfrak m\). (Reduce to the case \(F\sim E\); then
\(F=UU^*\), \(E=U^*U\), with \(U\in\mathcal A\), and
\(U=UE\in\mathfrak m\), hence \(F\in\mathfrak m\).)
\item Let \(E_1,E_2\) be projections of \(\mathfrak m\), and let \(E\) be
their supremum. Then \(E\in\mathfrak m\). (Use part~a and
\hyperref[ex:III-1-2]{Exercise~2(a)}.) \ArticleRef{12}.
\end{enumerate}

\item\label{ex:III-1-9} Let \(\mathcal A\) be a von Neumann algebra and
let \(\mathscr P\) be the set of projections of \(\mathcal A\). An
\emph{ideal of \(\mathscr P\)} is a subset \(\mathscr I\) of
\(\mathscr P\) having the following properties:
\begin{enumerate}[label=\textup{(\roman*)},leftmargin=*]
\item if \(E\in\mathscr I\) and \(F\prec E\), then \(F\in\mathscr I\);
\item if \(E_1,E_2\in\mathscr I\), their supremum belongs to
\(\mathscr I\).
\end{enumerate}
\begin{enumerate}[label=\textit{\alph*.},leftmargin=*]
\item Let \(\mathscr I\) be an ideal of \(\mathscr P\), and let
\(\mathfrak m\) be the two-sided ideal of \(\mathcal A\) generated by
\(\mathscr I\). Show that \(\mathfrak m\) is identical with the set
\(\mathfrak m'\) of those \(T\in\mathcal A\) for which
\(ET=TE=T\), for some (variable) \(E\in\mathscr I\). (Show that
\(\mathfrak m'\) is a two-sided ideal of \(\mathcal A\).) Deduce that the
set of projections of \(\mathfrak m\) is \(\mathscr I\).
\item Deduce a bijection between the ideals of
\(\mathscr P\) and the restricted two-sided ideals
(\hyperref[ex:I-1-6]{Chapter~I, \S~1, Exercise~6}) of \(\mathcal A\).
(Use part~a and \hyperref[ex:III-1-8]{Exercise~8}.)
\ArticleRef{12},\allowbreak \ArticleRef{124}.
\end{enumerate}

\item\label{ex:III-1-10} Let \(\mathcal A\) be a von Neumann algebra on
\(\mathfrak H\), and let \(\mathfrak m\) be a restricted two-sided ideal
of \(\mathcal A\)
(\hyperref[ex:I-1-6]{Chapter~I, \S~1, Exercise~6}).
\begin{enumerate}[label=\textit{\alph*.},leftmargin=*]
\item Let \(T\in\mathcal A_+\). In order that an element \(S\) of
\(\mathcal L(\mathfrak H)\) belong to \(\mathfrak m_+\) and be majorized
by \(T\), it is necessary and sufficient that
\[
  S=T^{1/2}T_1T^{1/2},
  \qquad
  T_1\in\mathfrak m_+,\quad T_1\leq I.
\]
(If \(S\in\mathfrak m_+\) and \(S\leq T\), show, using
\hyperref[ex:III-1-9]{Exercise~9(a)} and
\hyperref[lem:I-1-2]{Chapter~I, \S~1, Lemma~2}, that
\(S^{1/2}=RT^{1/2}\), where \(R\in\mathfrak m\) and
\(\lVert R\rVert\leq1\); then
\(S=(S^{1/2})^*S^{1/2}=T^{1/2}R^*RT^{1/2}\).)

% Source PDF page 237 (printed p. 228)
\phantomsection\label{srcpage:228}
\item Show that the set \(\mathscr F\) of the operators of
\(\mathfrak m_+\) majorized by \(T\) is upward directed. (Use part~a and
\hyperref[ex:III-1-9]{Exercise~9(a)}.)

\noindent\textit{Problem.} Is this result still true if
\(\mathfrak m\) is not restricted?

\item If the strong closure of \(\mathfrak m\) is \(\mathcal A\), then
\(T\) is the supremum of \(\mathscr F\). \ArticleRef{12}.
\end{enumerate}

\item\label{ex:III-1-11} Let \(\mathcal A\) be a von Neumann algebra,
let \(\mathfrak m\) be a two-sided ideal of \(\mathcal A\) strongly
everywhere dense in \(\mathcal A\), and let \(\varphi\) be a linear form
on \(\mathfrak m\), positive on \(\mathfrak m_+\), such that
\[
  \varphi(ST)=\varphi(TS)
  \qquad(S\in\mathfrak m,\ T\in\mathcal A),
\]
and having the following property: if \(\mathscr F\subset\mathfrak m_+\)
is an upward directed set with supremum \(T\in\mathfrak m_+\), then
\(\varphi(T)\) is the supremum of \(\varphi(\mathscr F)\). Then there is a
unique normal trace \(\psi\) on \(\mathcal A_+\) which agrees with
\(\varphi\) on \(\mathfrak m_+\), and \(\psi\) is semifinite.

[Let \(\mathfrak n\) be the restricted ideal generated by the projections
of \(\mathfrak m\)
(\hyperref[ex:I-1-6]{Chapter~I, \S~1, Exercise~6}). For
\(R\in\mathcal A_+\), put
\[
  \psi(R)=\sup\varphi(\mathscr F_R),
\]
where \(\mathscr F_R\) is the set (upward directed by
\hyperref[ex:III-1-10]{Exercise~10}) of the
\(R_1\in\mathfrak n_+\) majorized by \(R\). It remains only to prove the
following: if \(\mathscr F\subset\mathcal A_+\) is an upward directed set
with supremum \(R\), and if \(\lambda=\sup\psi(\mathscr F)\), then
\(\psi(R)=\lambda\). Let \(\mathscr G\) be the upward directed set of the
\(S\in\mathfrak n_+\) which are majorized by some (variable) operator of
\(\mathscr F\). The set \(\mathscr G\) is majorized by \(R\), and every
upper bound of \(\mathscr G\) majorizes \(\mathscr F\), by
\hyperref[ex:III-1-10]{Exercise~10(c)}, hence majorizes \(R\). Thus
\(R\) is the supremum of
\(\mathscr G\). Put \(\mu=\sup\varphi(\mathscr G)\). Plainly,
\(\mu\leq\lambda\leq\psi(R)\). We shall show that
\(\varphi(R_1)\leq\mu\) for \(R_1\in\mathscr F_R\), which will imply
\(\psi(R)\leq\mu\). There is a projection \(E\in\mathfrak m\) such that
\(R_1=ER_1E\). The set \(E\mathscr GE\) is upward directed, is majorized
by \(ERE\), and \(ERE\) is strongly adherent to it; hence \(ERE\) is the
supremum of \(E\mathscr GE\). Therefore
\(\varphi(ERE)=\sup\varphi(E\mathscr GE)\). The number \(\mu\) majorizes
\(\varphi(\mathscr G)\), hence \(\varphi(E\mathscr GE)\), and consequently
\(\mu\geq\varphi(ERE)\geq\varphi(ER_1E)=\varphi(R_1)\).]
\ArticleRef{12}.

\item\label{ex:III-1-12} Let \(\mathcal A\) be a von Neumann algebra and
let \(T\) be a closed operator with everywhere dense domain, affiliated
with \(\mathcal A\)
(\hyperref[ex:I-1-10]{Chapter~I, \S~1, Exercise~10}). The supports of \(T\)
and \(T^*\) are equivalent. (The support \(E\) of \(T\) is defined as when
\(T\) is continuous: \(I-E\) is the projection onto the null space of
\(T\).) \ArticleRef{65}.

\item\label{ex:III-1-13} Let \(\mathcal A\) be a finite von Neumann
algebra on \(\mathfrak H\), and let \(\mathfrak M\) be a vector subspace
of \(\mathfrak H\). The subspace \(\mathfrak M\) is said to be
\emph{essentially dense} in \(\mathfrak H\) (relative to
\(\mathcal A\)) if there is an increasing sequence
\(\mathfrak X_1,\mathfrak X_2,\ldots\) of closed vector subspaces of
\(\mathfrak H\), contained in \(\mathfrak M\) and total in
\(\mathfrak H\), such that \(P_{\mathfrak X_i}\in\mathcal A\) for every
\(i\).
\begin{enumerate}[label=\textit{\alph*.},leftmargin=*]
\item The intersection of two essentially dense vector subspaces
\(\mathfrak M,\mathfrak M'\) is essentially dense. [Let
\((\mathfrak X_1,\mathfrak X_2,\ldots)\) [respectively,
\((\mathfrak X'_1,\mathfrak X'_2,\ldots)\)] be a sequence having the
properties in the definition relative to \(\mathfrak M\) (respectively,
\(\mathfrak M'\)). Then the \(\mathfrak X_i\cap\mathfrak X'_i\) form an
increasing sequence of closed vector subspaces, and
\(P_{\mathfrak X_i\cap\mathfrak X'_i}\in\mathcal A\). For every finite
normal trace \(\varphi\) on \(\mathcal A\),
\[
  \varphi\bigl(P_{\mathfrak H\ominus
    (\mathfrak X_i\cap\mathfrak X'_i)}\bigr)
  \leq
  \varphi(P_{\mathfrak H\ominus\mathfrak X_i})
  +\varphi(P_{\mathfrak H\ominus\mathfrak X'_i})
\]
by \hyperref[ex:III-1-2]{Exercise~2(b)}; hence
\(\varphi(P_{\mathfrak H\ominus
(\mathfrak X_i\cap\mathfrak X'_i)})\to0\), and therefore the
\(\mathfrak X_i\cap\mathfrak X'_i\) are total in \(\mathfrak H\).]

\item If \(\mathfrak M\) is essentially dense and \(T\) is a closed
operator with everywhere dense domain, affiliated with \(\mathcal A\)
(\hyperref[ex:I-1-10]{Chapter~I, \S~1, Exercise~10}), then
\(T^{-1}(\mathfrak M)\) is essentially dense. [Let \(T=W|T|\) be the
polar decomposition of \(T\), and let \(\mathfrak X(\lambda)\)
\((0\leq\lambda\leq+\infty)\) be the largest spectral subspace of
\(|T|\) such that
\[
  |T|P_{\mathfrak X(\lambda)}
  \leq\lambda P_{\mathfrak X(\lambda)}.
\]
Let \(\mathfrak X^i\) be the set of the \(x\in\mathfrak H\) such that
\(TP_{\mathfrak X(i)}x\in\mathfrak X_i\). The subspaces
\(\mathfrak X(i)\cap\mathfrak X^i\) are closed, increasing, and contained
in \(T^{-1}(\mathfrak M)\); their projections belong to \(\mathcal A\).
For every finite normal trace \(\varphi\) on \(\mathcal A\),
\(\varphi(P_{\mathfrak H\ominus\mathfrak X(i)})\to0\).
Furthermore, \(P_{\mathfrak H\ominus\mathfrak X^i}\) is the support of
\(P_{\mathfrak H\ominus\mathfrak X_i}TP_{\mathfrak X(i)}\), and the
support of
\((P_{\mathfrak H\ominus\mathfrak X_i}TP_{\mathfrak X(i)})^*\) is
majorized by \(P_{\mathfrak H\ominus\mathfrak X_i}\). Conclude that
\(\varphi(P_{\mathfrak H\ominus\mathfrak X^i})\to0\), and that the
\(\mathfrak X(i)\cap\mathfrak X^i\) are total in \(\mathfrak H\).]
% Source PDF page 238 (printed p. 229)
\phantomsection\label{srcpage:229}
\item Let \(T_1,T_2\) be closed operators, with everywhere dense
domains, affiliated with \(\mathcal A\). Show that \(T_1+T_2\) and
\(T_1T_2\) have smallest closed extensions, with everywhere dense
domains, affiliated with \(\mathcal A\). (By \textit{a} and \textit{b},
\(T_1+T_2\), \(T_1T_2\), \(T_1^*+T_2^*\), and \(T_2^*T_1^*\) have
everywhere dense domains.)

\item Let \(x\in\mathfrak H\). If
\(y\in\mathfrak X_x^{\mathcal A}\), there is a closed operator \(T\),
with everywhere dense domain, affiliated with \(\mathcal A\), such that
\(y=Tx\) (use \textit{c} and \hyperref[lem:III-1-3]{Lemma~3}).
\ArticleRef{65}, \ArticleRef{101}.
\end{enumerate}

\item\label{ex:III-1-14}
Let \(\mathcal A\) be a finite von Neumann algebra on \(\mathfrak H\),
and let \(\varphi\) and \(\psi\) be two normal positive linear forms on
\(\mathcal A\), with supports \(E_\varphi,E_\psi\), such that
\(E_\varphi\leq E_\psi\). Show that there is a closed operator
\(T\mathrel{\eta}\mathcal A\)
(\hyperref[ex:I-1-10]{Chapter~I, \S~1, Exercise~10}), with
everywhere dense domain, such that, for every \(S\in\mathcal A\),
\[
  \varphi(S)=\lim_{n\to+\infty}\psi(T_n^*ST_n),
\]
where \(T_n\) is defined as follows: if \(T=W\lvert T\rvert\) is the
polar decomposition of \(T\), and \(E(\lambda)\) is the largest spectral
projection of \(\lvert T\rvert\) such that
\(\lvert T\rvert E(\lambda)\leq\lambda E(\lambda)\), put
\[
  T_n=W\lvert T\rvert E(n)\in\mathcal A.
\]
[First reduce, by considering \(\mathcal A_{E_\psi}\), to the case in
which \(\psi\) is faithful; then, by an isomorphism onto \(\mathcal A\),
to the case in which \(\psi=\omega_x\), where \(x\) is totalizing and
separating for \(\mathcal A\). Then \(\varphi=\omega_y\), by
\hyperref[thm:III-1-4]{Theorem~4}; apply
\hyperref[ex:III-1-13]{Exercise~13\textit{d}} to \(x\) and \(y\).]
\ArticleRef{19}, \ArticleRef{92}, \ArticleRef{101}.

\item\label{ex:III-1-15}
Let \(Z\) be a Borel space, let \(\nu\) be a standard positive measure
on \(Z\), let \(\zeta\mapsto\mathfrak H(\zeta)\) be a
\(\nu\)-measurable field of complex Hilbert spaces on \(Z\), and let
\(\zeta\mapsto\mathcal A(\zeta)\subset
\mathcal L(\mathfrak H(\zeta))\) be a \(\nu\)-measurable field of von
Neumann algebras. Let
\[
  E=\int^{\oplus}E(\zeta)\,d\nu(\zeta),
  \qquad
  F=\int^{\oplus}F(\zeta)\,d\nu(\zeta)
\]
be two projections of
\[
  \mathcal A=\int^{\oplus}\mathcal A(\zeta)\,d\nu(\zeta).
\]
In order that \(E\prec F\), it is necessary and sufficient that
\(E(\zeta)\prec F(\zeta)\) almost everywhere. [The condition is plainly
necessary. To prove that it is sufficient, argue, for example, as in
\hyperref[lem:II-2-2]{Chapter~II, \S~2, Lemma~2}.] \ArticleRef{80}
\end{enumerate}

\section{Classification of Projections}\label{sec:III-2}

\NumberedParagraph{1}{Definitions}
\label{no:III-2-1}

\begin{definition}\label{def:III-2-1}
Let \(\mathcal A\) be a von Neumann algebra and let \(E\) be a projection
of \(\mathcal A\). The projection \(E\) is said to be \emph{finite}
(respectively, \emph{semifinite}, \emph{properly infinite},
\emph{purely infinite}, \emph{infinite}) if the algebra
\(\mathcal A_E\) is finite (respectively, semifinite, properly infinite,
purely infinite, infinite).
\end{definition}

When precision is desired, one says that \(E\) is finite (respectively,
semifinite, and so on) relative to \(\mathcal A\).

If \(E\) is finite (respectively, semifinite, and so on) and
\(F\sim E\), then \(F\) is finite (respectively, semifinite, and so on).

A finite projection is semifinite. A purely infinite projection is
properly infinite. An abelian projection is finite.

\begin{proposition}\label{prop:III-2-1}
Let \((E_\iota)_{\iota\in I}\) be a family of projections of
\(\mathcal A\) whose central supports are pairwise orthogonal, and put
\[
  E=\sum_{\iota\in I}E_\iota.
\]%
% Source PDF page 239 (printed p. 230)
\phantomsection\label{srcpage:230}%
In order that \(E\) be finite (respectively, semifinite, properly
infinite, purely infinite), it is necessary and sufficient that every
\(E_\iota\) be finite (respectively, semifinite, properly infinite,
purely infinite).
\end{proposition}

\begin{proof}
We have
\[
  \mathcal A_E=\prod_{\iota\in I}\mathcal A_{E_\iota},
\]
so the proposition follows from
\hyperref[prop:I-6-7]{Chapter~I, \S~6, Proposition~7}.
\end{proof}

\begin{proposition}\label{prop:III-2-2}
Let \(E\) be a finite (respectively, semifinite, purely infinite)
projection of \(\mathcal A\), and let \(F\) be a projection such that
\(F\prec E\). Then \(F\) is finite (respectively, semifinite, purely
infinite).
\end{proposition}

\begin{proof}
It is enough to study the case in which \(F\leq E\). Then
\(\mathcal A_F=(\mathcal A_E)_F\), so the proposition follows from
\hyperref[prop:I-6-11]{Chapter~I, \S~6, Proposition~11} and
\hyperref[cor:I-6-13-4]{Proposition~13, Corollary~4}.
\end{proof}

Observe that the zero projection is at once finite, semifinite, properly
infinite, and purely infinite. This is somewhat more convenient than the
usual convention, according to which zero is only a finite and
semifinite projection.

\noindent\textit{Bibliography.}
\ArticleRef{10}, \ArticleRef{42}, \ArticleRef{65}.

\NumberedParagraph{2}{Cyclic projections of \(\mathcal A\) and cyclic
projections of \(\mathcal A'\)}
\label{no:III-2-2}

\begin{lemma}\label{lem:III-2-1}
Let \(P_{\mathfrak X}\) be a projection of \(\mathcal A\), let
\(P_{\mathfrak X'}\) be a projection of \(\mathcal A'\), and put
\(\mathfrak Y=\mathfrak X\cap\mathfrak X'\). Set
\[
  \mathcal B=\mathcal A_{\mathfrak X},
  \qquad
  \mathcal C=\mathcal A'_{\mathfrak X'}.
\]
Then \(\mathcal B_{\mathfrak Y}\) and
\(\mathcal C'_{\mathfrak Y}\) are identical with one and the same von
Neumann algebra \(\mathcal W\), while
\(\mathcal C_{\mathfrak Y}\) and
\(\mathcal B'_{\mathfrak Y}\) are identical with \(\mathcal W'\).
If \(P_{\mathfrak X}\) and \(P_{\mathfrak X'}\) have the same central
support, \(\mathcal B\) is isomorphic to \(\mathcal W\), and
\(\mathcal C\) to \(\mathcal W'\). If
\[
  \mathfrak X=\mathfrak X_x^{\mathcal A'},
  \qquad
  \mathfrak X'=\mathfrak X_x^{\mathcal A}
\]
(in which case \(P_{\mathfrak X}\) and \(P_{\mathfrak X'}\) have the
same central support), then \(x\) is totalizing and separating for
\(\mathcal W\) and \(\mathcal W'\).
\end{lemma}

\begin{proof}
We have
\[
  \mathcal B_{\mathfrak Y}
  =(\mathcal A_{\mathfrak X})_{\mathfrak Y}
  =\mathcal A_{\mathfrak Y}
  =(\mathcal A'_{\mathfrak X'})'_{\mathfrak Y}
  =\mathcal C'_{\mathfrak Y},
\]
which defines \(\mathcal W\); and
\[
  \mathcal C_{\mathfrak Y}
  =(\mathcal C'_{\mathfrak Y})'=\mathcal W',
  \qquad
  \mathcal B'_{\mathfrak Y}
  =(\mathcal B_{\mathfrak Y})'=\mathcal W'.
\]
If \(P_{\mathfrak X}\) and \(P_{\mathfrak X'}\) have the same central
support \(P_3\), the induction of \(\mathcal A'_3\) on
\(\mathcal A'_{\mathfrak X}=\mathcal B'\) is an isomorphism. Hence the
central support of \(P_{\mathfrak Y}\), considered as an operator of
\(\mathcal B'\), is \(I_{\mathfrak X}\); consequently \(\mathcal B\)
is isomorphic to \(\mathcal B_{\mathfrak Y}=\mathcal W\). Similarly,
\(\mathcal C\) is isomorphic to
\(\mathcal C_{\mathfrak Y}=\mathcal W'\). Finally, if
\(\mathfrak X=\mathfrak X_x^{\mathcal A'}\) and
\(\mathfrak X'=\mathfrak X_x^{\mathcal A}\), then \(x\) is totalizing
for \(\mathcal A'_{\mathfrak X}=\mathcal B'\), and hence for
\(\mathcal B'_{\mathfrak Y}=\mathcal W'\); it is also totalizing for
\(\mathcal A_{\mathfrak X'}=\mathcal C'\), and hence for
\(\mathcal C'_{\mathfrak Y}=\mathcal W\). Thus it is separating as
well for each of these two mutually commutant algebras.
\end{proof}

% Source PDF page 240 (printed p. 231)
\phantomsection\label{srcpage:231}
\begin{proposition}\label{prop:III-2-3}
Let \(x\in\mathfrak H\). Then \(E_x^{\mathcal A}\) is abelian
(respectively, finite, semifinite, properly infinite, purely infinite) if
and only if \(E_x^{\mathcal A'}\) is abelian (respectively, finite,
semifinite, properly infinite, purely infinite).
\end{proposition}

\begin{proof}
We have \(\mathcal A=\mathcal A_1\times\mathcal A_2\) and
\(\mathcal A'=\mathcal A'_1\times\mathcal A'_2\), where
\(\mathcal A_1\) and \(\mathcal A'_1\) are semifinite, while
\(\mathcal A_2\) and \(\mathcal A'_2\) are purely infinite. Moreover,
every projection of \(\mathcal A_2\) or \(\mathcal A'_2\) is purely
infinite
(\hyperref[cor:I-6-13-4]{Chapter~I, \S~6, Proposition~13, Corollary~4}).
It is therefore
enough to study the case in which \(\mathcal A\) is semifinite. Adopt the
notation of \hyperref[lem:III-2-1]{Lemma~1}. All the algebras occurring
there are semifinite. The
algebras \(\mathcal W\) and \(\mathcal W'\) are standard
(\hyperref[thm:III-1-5]{\S~1, Theorem~5}), so
\(\mathcal A_{\mathfrak X}=\mathcal B\) is anti-isomorphic to
\(\mathcal A'_{\mathfrak X'}=\mathcal C\). Hence
\(E_x^{\mathcal A}\) is abelian (respectively, finite, semifinite,
properly infinite) if and only if \(E_x^{\mathcal A'}\) is abelian
(respectively, finite, semifinite, properly infinite).
\end{proof}

\noindent\textit{Bibliography.}
\ArticleRef{17}, \ArticleRef{19}, \ArticleRef{31}, \ArticleRef{40},
\ArticleRef{65}, \ArticleRef{100}.

\NumberedParagraph{3}{Finite projections}
\label{no:III-2-3}

\begin{proposition}\label{prop:III-2-4}
Let \(E\) be a finite projection of \(\mathcal A\). Every projection of
\(\mathcal A\) majorized by \(E\) and equivalent to \(E\) is equal to
\(E\).
\end{proposition}

\begin{proof}
This follows from \hyperref[prop:III-1-3]{\S~1, Proposition~3}.
\end{proof}

We shall prove later (\S~8, Theorem~1, Corollary~1) a converse of
\hyperref[prop:III-2-4]{Proposition~4}.

\begin{proposition}\label{prop:III-2-5}
Let \(E_1,E_2,\ldots,E_n\) be projections of \(\mathcal A\), and let
\(E\) be their supremum. If the \(E_i\) are finite, then \(E\) is finite.
\end{proposition}

\begin{proof}
It is enough to consider the case \(n=2\).

\smallskip
\noindent\textup{(i)} Suppose first that \(E_1\) and \(E_2\) are
equivalent and orthogonal. Then
\[
  \mathcal A_E\cong
  \mathcal A_{E_1}\otimes\mathcal L(\mathfrak K),
\]
where \(\mathfrak K\) is a complex Hilbert space of dimension \(2\)
(\hyperref[prop:I-2-5]{Chapter~I, \S~2, Proposition~5}). Since the algebras
\(\mathcal A_{E_1}\) and \(\mathcal L(\mathfrak K)\) are finite,
\(\mathcal A_E\) is finite
(\hyperref[prop:I-6-12]{Chapter~I, \S~6, Proposition~12}).

\smallskip
\noindent\textup{(ii)} Suppose only that \(E_1\) and \(E_2\) are
orthogonal. By \hyperref[thm:III-1-1]{\S~1, Theorem~1} and
\cref{prop:III-2-1}, it is enough to study
the case \(E_1\prec E_2\) and the case \(E_2\prec E_1\). Suppose, for
example, that \(E_2\prec E_1\). Let \(\mathfrak K\) be a complex
Hilbert space, let \(F\) be a rank-one projection on \(\mathfrak K\),
and identify each \(T\in\mathcal A\) with
\(T\otimes F\in\mathcal B=\mathcal A\otimes\mathcal L(\mathfrak K)\).
Instead of proving that \(E\) is finite relative to \(\mathcal A\), we
shall prove the equivalent assertion that it is finite relative to
\(\mathcal B\). By taking \(\dim\mathfrak K\geq2\), we may suppose that
there is in \(\mathcal B\) a projection \(E'_1\) equivalent to \(E_1\)
and orthogonal to \(E_1\). Then
\[
  E_1+E_2\prec E_1+E'_1,
\]
and \(E_1+E'_1\) is finite by part~\textup{(i)} of the proof.

% Source PDF page 241 (printed p. 232)
\phantomsection\label{srcpage:232}
\noindent\textup{(iii)} Finally consider the general case. Put
\[
  \mathfrak X_1=E_1(\mathfrak H),\qquad
  \mathfrak X_2=E_2(\mathfrak H),
\]
let \(\mathfrak X\) be the closure of
\(\mathfrak X_1+\mathfrak X_2\), that is, \(E(\mathfrak H)\), and put
\(\mathfrak Y=\mathfrak X\ominus\mathfrak X_1\). An element of
\(\mathfrak Y\) orthogonal to \(\mathfrak X_2\) is orthogonal to
\(\mathfrak X_1+\mathfrak X_2\), hence to \(\mathfrak Y\), and is
therefore zero. Thus \(\mathfrak Y\prec\mathfrak X_2\)
(\hyperref[cor:III-1-prop2-1]{\S~1, Proposition~2, Corollary~1}).
Consequently \(P_{\mathfrak Y}\) is finite,
and hence
\[
  P_{\mathfrak X}=P_{\mathfrak X_1}+P_{\mathfrak Y}
\]
is finite by part~\textup{(ii)} of the proof.
\end{proof}

\begin{proposition}\label{prop:III-2-6}
Let \(E,F\) be two equivalent finite projections of \(\mathcal A\), and
let \(G\) be a projection of \(\mathcal A\) majorizing \(E\) and \(F\).
There is a unitary operator \(U\in\mathcal A\) such that
\[
  UEU^{-1}=F,
  \qquad
  UGU^{-1}=G.
\]
In particular, \(G-E\sim G-F\).
\end{proposition}

\begin{proof}
Let \(G_1\) be the supremum of \(E\) and \(F\). We have \(G_1\leq G\),
and \(G_1\) is finite (\hyperref[prop:III-2-5]{Proposition~5}). By
\hyperref[thm:III-1-1]{\S~1, Theorem~1}, it is enough
to study the case \(G_1-E\prec G_1-F\) and the case
\(G_1-F\prec G_1-E\). If, for example,
\(G_1-E\sim E_1\), where \(E_1\leq G_1-F\), then
\[
  G_1\geq F+E_1
  \sim E+(G_1-E)=G_1;
\]
hence \(F+E_1=G_1\) by \hyperref[prop:III-2-4]{Proposition~4}, and consequently
\(G_1-F=E_1\sim G_1-E\). Let \(V\) (respectively, \(W\)) be a partially
isometric operator of \(\mathcal A\) having \(E\) (respectively,
\(G_1-E\)) as initial projection and \(F\) (respectively,
\(G_1-F\)) as final projection. Let \(U\) be the unitary operator that
agrees with \(V\) on \(E(\mathfrak H)\), with \(W\) on
\((G_1-E)(\mathfrak H)\), and with the identity on
\((I-G_1)(\mathfrak H)\). Then
\[
  U\in\mathcal A,
  \qquad UEU^{-1}=F,
  \qquad UGU^{-1}=G.
\]
\end{proof}

\begin{remark*}
The proof of the proposition and its conclusions remain valid under the
following hypotheses: \(E\) and \(F\) are equivalent projections of
\(\mathcal A\) with supremum \(G_1\); every projection of \(\mathcal A\)
majorized by \(G_1\) and equivalent to \(G_1\) is equal to \(G_1\); and
\(G\) is a projection of \(\mathcal A\) majorizing \(E\) and \(F\). We
shall use this remark in \S~8.
\end{remark*}

\noindent\textit{Bibliography.}
\ArticleRef{10}, \ArticleRef{42}, \ArticleRef{65}.

\Needspace{4\baselineskip}
\NumberedParagraph{4}{Semifinite projections}
\label{no:III-2-4}

\begin{proposition}\label{prop:III-2-7}
Let \(E\) be a nonzero semifinite projection of \(\mathcal A\). There is
a nonzero finite projection of \(\mathcal A\) majorized by \(E\).
\end{proposition}

\begin{proof}
There is a faithful normal semifinite trace \(\varphi\) on
\(\mathcal A_E^+\). Let \(T\) be a nonzero operator of
\(\mathcal A_E^+\) such that \(\varphi(T)<+\infty\), and let \(F\) be a
nonzero spectral projection of \(T\) such that \(F\leq\lambda T\). Then
\(\varphi(F)<+\infty\). The restriction of \(\varphi\) to
\(F\mathcal A_E^+F\) defines a faithful finite normal trace on
\(\mathcal A_F^+\), so \(\mathcal A_F\) is finite.
\end{proof}

\setcounter{corollary}{0}
\begin{corollary}\label{cor:III-2-prop7-1}
Let \(E\) be a semifinite projection of \(\mathcal A\).
\begin{enumerate}[label=\textup{(\roman*)},leftmargin=*]
\item There is a family \((E_\iota)_{\iota\in I}\) of pairwise
orthogonal finite projections of \(\mathcal A\) such that
\[
  E=\sum_{\iota\in I}E_\iota.
\]

% Source PDF page 242 (printed p. 233)
\phantomsection\label{srcpage:233}
\item If \(E\ne0\), there is a finite projection \(F\) of
\(\mathcal A\), majorized by \(E\), such that \(\mathcal A_F\) has a
nonzero trace vector.
\end{enumerate}
\end{corollary}

\begin{proof}
Let \((E_\iota)_{\iota\in I}\) be a maximal family of nonzero,
pairwise orthogonal finite projections of \(\mathcal A\), majorized by
\(E\). Put
\[
  E'=E-\sum_{\iota\in I}E_\iota.
\]
Then \(E'\) is semifinite and majorizes no nonzero finite projection;
hence \(E'=0\) by \hyperref[prop:III-2-7]{Proposition~7}.

Let us prove \textup{(ii)}. By \hyperref[prop:III-2-7]{Proposition~7},
we may suppose that \(E\)
is finite and nonzero. Let \(x\) be a nonzero vector in the range of
\(E\), and put \(F=E_x^{\mathcal A'}\). There is a nonzero finite normal
trace \(\varphi\) on \(\mathcal A_F\). Since \(x\) is separating for
\(\mathcal A_F\), \(\varphi\) has the form \(\omega_y\)
(\hyperref[thm:III-1-4]{\S~1, Theorem~4}), and \(y\) is a trace vector
for \(\mathcal A_F\).
\end{proof}

\begin{corollary}\label{cor:III-2-prop7-2}
Let \(E\) be a semifinite projection of \(\mathcal A\). There is a
finite projection \(F\) of \(\mathcal A\), majorized by \(E\), having
the same central support as \(E\).
\end{corollary}

\begin{proof}
Let \((E_\iota)_{\iota\in I}\) be a maximal family of nonzero finite
projections of \(\mathcal A\), majorized by \(E\), whose central supports
\(E'_\iota\) are pairwise orthogonal. Put
\[
  F=\sum_{\iota\in I}E_\iota\leq E;
\]
this projection is finite by \hyperref[prop:III-2-1]{Proposition~1}.
Let \(F'\) (respectively,
\(E'\)) be the central support of \(F\) (respectively, \(E\)). We have
\(F'\leq E'\). If \(E'-F'\ne0\), then \(E(E'-F')\ne0\), and hence, by
\hyperref[prop:III-2-7]{Proposition~7}, there is a nonzero finite
projection majorized by \(E\)
and orthogonal to \(F'\), and therefore to all the \(E_\iota\). This
contradicts the maximality of the family \((E_\iota)_{\iota\in I}\).
Thus \(F'=E'\).
\end{proof}

\begin{corollary}\label{cor:III-2-prop7-3}
In order that \(\mathcal A\) be semifinite, it is necessary and sufficient
that \(\mathcal A\) be isomorphic to a von Neumann algebra \(\mathcal B\)
such that \(\mathcal B'\) is finite.
\end{corollary}

\begin{proof}
The condition is sufficient by
\hyperref[cor:I-6-13-1]{Chapter~I, \S~6, Proposition~13, Corollary~1}.
Conversely, if \(\mathcal A\) is semifinite, then
\(\mathcal A'\) is semifinite, so by
\hyperref[cor:III-2-prop7-2]{Corollary~2} there is a finite
projection \(E'\) of \(\mathcal A'\) with central support \(I\). Then
\(\mathcal A\) is isomorphic to
\(\mathcal B=\mathcal A_{E'}\), and
\(\mathcal B'=\mathcal A'_{E'}\) is finite.
\end{proof}

\begin{corollary}\label{cor:III-2-prop7-4}
Let \(\mathcal A\) be a semifinite von Neumann algebra. In order that
\(\mathcal A\) be continuous, it is necessary and sufficient that there
exist a decreasing sequence \((E_1,E_2,\ldots)\) of finite projections
of \(\mathcal A\), with central support \(I\), such that
\[
  E_n-E_{n+1}\sim E_{n+1}
  \qquad(n=1,2,\ldots).
\]
\end{corollary}

\begin{proof}
If \(\mathcal A\) is continuous, the sequence \((E_i)\) exists by
\hyperref[cor:III-2-prop7-2]{Corollary~2} and
\hyperref[cor:III-1-thm1-3]{\S~1, Theorem~1, Corollary~3}. Now suppose
that there is a
sequence \((E_i)\) with the properties in the statement. It is enough to
show that \(\mathcal A_{E_1}\) is continuous (for then
\(\mathcal A'_{E_1}\) is continuous, hence \(\mathcal A'\), which is
isomorphic to \(\mathcal A'_{E_1}\), is continuous, and therefore so is
\(\mathcal A\)). We shall consequently suppose from now on that
\(\mathcal A\) is finite. Using
\hyperref[prop:I-6-9]{Chapter~I, \S~6, Proposition~9(iii)}, we
next reduce to the case in which \(\mathcal A\) is finite and of
countable type. There is then\space%
% Source PDF page 243 (printed p. 234)
\phantomsection\label{srcpage:234}%
a faithful finite trace \(\varphi\) on \(\mathcal A\)
[\hyperref[prop:I-6-9]{Chapter~I, \S~6, Proposition~9(ii)}]. For every
abelian projection \(E\) of \(\mathcal A\), we have \(E\prec E_n\) for
every \(n\) (by \hyperref[thm:III-1-1]{\S~1, Theorem~1} and the
minimality property of abelian projections), and therefore
\[
  \varphi(E)\leq\varphi(E_n)=2^{-n}\varphi(E_1).
\]
Thus \(\varphi(E)=0\), and hence \(E=0\). Therefore \(\mathcal A\) is
continuous.
\end{proof}

\begin{proposition}\label{prop:III-2-8}
Let \(E\) be the largest projection in the centre of \(\mathcal A\) that
is semifinite relative to \(\mathcal A\)
(\hyperref[prop:I-6-8]{Chapter~I, \S~6, Proposition~8}). In order that
a projection \(F\) of \(\mathcal A\) be
semifinite, it is necessary and sufficient that \(F\leq E\).
\end{proposition}

\begin{proof}
If \(F\leq E\), then \(F\) is semifinite by
\hyperref[prop:III-2-2]{Proposition~2}. If \(F\) is
semifinite, let \(G\) be its central support. Then \(\mathcal A_F\) is
semifinite, hence \(\mathcal A'_F\) is semifinite, hence
\(\mathcal A'_G\) is semifinite, hence \(\mathcal A_G\) is semifinite,
and consequently \(G\leq E\).
\end{proof}

\hyperref[cor:III-2-prop7-3]{Corollary~3 of Proposition~7} gives a
definition of semifinite algebras
analogous to the definition of discrete algebras.

\noindent\textit{Bibliography.}
\ArticleRef{10}, \ArticleRef{42}, \ArticleRef{65}.

\NumberedParagraph{5}{Properly infinite projections}
\label{no:III-2-5}

\begin{proposition}\label{prop:III-2-9}
Let \(\mathcal Z\) be the centre of \(\mathcal A\). In order that a
projection \(E\) of \(\mathcal A\) be properly infinite, it is necessary
and sufficient that, for every projection \(F\) of \(\mathcal Z\) such
that \(FE\ne0\), the projection \(FE\) be infinite.
\end{proposition}

\begin{proof}
The condition is necessary by \hyperref[prop:III-2-1]{Proposition~1}.
On the other hand, if \(E\)
is not properly infinite, there is a nonzero projection \(E_1\) in the
centre of \(\mathcal A_E\) such that \(\mathcal A_{E_1}\) is finite
(\hyperref[prop:I-6-8]{Chapter~I, \S~6, Proposition~8}). We have
\(E_1=FE\) for some
\(F\in\mathcal Z\). Thus \(FE\ne0\), and \(FE\) is finite. The condition
in the proposition is therefore sufficient.
\end{proof}

\setcounter{corollary}{0}
\begin{corollary}\label{cor:III-2-prop9-1}
If a projection \(E\) of \(\mathcal A\) is properly infinite, its central
support \(E'\) is properly infinite.
\end{corollary}

\begin{proof}
If \(F\) is a projection of the centre of \(\mathcal A\) such that
\(FE'\ne0\), then \(FE\ne0\), so \(FE\), and \emph{a fortiori} \(FE'\),
are infinite.
\end{proof}

\begin{corollary}\label{cor:III-2-prop9-2}
The largest projection in the centre of \(\mathcal A\) that is properly
infinite relative to \(\mathcal A\) is also the largest properly infinite
projection of \(\mathcal A\).
\end{corollary}

\begin{proof}
This follows at once from
\hyperref[cor:III-2-prop9-1]{Corollary~1}.
\end{proof}

\begin{corollary}\label{cor:III-2-prop9-3}
Let \((E_\iota)_{\iota\in I}\) be a family of projections of
\(\mathcal A\), and let \(E\) be their supremum. If the \(E_\iota\) are
properly infinite, then \(E\) is properly infinite.
\end{corollary}

\begin{proof}
This follows at once from \hyperref[prop:III-2-9]{Proposition~9}.
\end{proof}

% Source PDF page 244 (printed p. 235)
\phantomsection\label{srcpage:235}
\begin{proposition}\label{prop:III-2-10}
Let \((E_\iota)_{\iota\in I}\) be an infinite family of nonzero,
equivalent, pairwise orthogonal projections of \(\mathcal A\), and put
\[
  E=\sum_{\iota\in I}E_\iota.
\]
Then \(E\) is properly infinite.
\end{proposition}

\begin{proof}
Let \(J\) be the complement in \(I\) of an arbitrarily chosen element of
\(I\). Then \(I\) and \(J\) are equipotent. For every projection \(F\)
of \(\mathcal Z\) such that \(EF\ne0\), we have
\[
  EF\sim\sum_{\iota\in J}E_\iota F\ne EF;
\]
hence \(EF\) is infinite by \hyperref[prop:III-2-4]{Proposition~4}, and
consequently \(E\) is properly infinite by
\hyperref[prop:III-2-9]{Proposition~9}.
\end{proof}

We shall prove later (\S~8, Theorem~1, Corollary~2) a converse of
\hyperref[prop:III-2-10]{Proposition~10}.

\noindent\textit{Bibliography.}
\ArticleRef{10}, \ArticleRef{42}, \ArticleRef{65}.

\NumberedParagraph{6}{Purely infinite projections}
\label{no:III-2-6}

\begin{proposition}\label{prop:III-2-11}
In order that a projection \(E\) of \(\mathcal A\) be purely infinite,
it is necessary and sufficient that \(E\) majorize no nonzero finite
projection.
\end{proposition}

\begin{proof}
The condition is necessary by \hyperref[prop:III-2-2]{Proposition~2}.
It is sufficient, for if
\(E\) is not purely infinite, then \(E\) majorizes a nonzero semifinite
projection, and hence a nonzero finite projection by
\hyperref[prop:III-2-7]{Proposition~7}.
\end{proof}

\begin{proposition}\label{prop:III-2-12}
Let \(E\) be the largest projection in the centre of \(\mathcal A\) that
is purely infinite relative to \(\mathcal A\)
(\hyperref[prop:I-6-8]{Chapter~I, \S~6, Proposition~8}). In order that
a projection \(F\) of \(\mathcal A\) be
purely infinite, it is necessary and sufficient that \(F\leq E\).
\end{proposition}

\begin{proof}
If \(F\leq E\), then \(F\) is purely infinite by
\hyperref[prop:III-2-2]{Proposition~2}. If \(F\)
is purely infinite, let \(G\) be its central support. Then
\(\mathcal A_F\) is purely infinite, hence \(\mathcal A'_F\) is purely
infinite, hence \(\mathcal A'_G\) is purely infinite, hence
\(\mathcal A_G\) is purely infinite, and consequently \(G\leq E\).
\end{proof}

Cf. also \hyperref[ex:III-2-1]{Exercise~1} and
\S~8, Theorem~1, Corollary~5.

\noindent\textit{Bibliography.}
\ArticleRef{10}, \ArticleRef{42}, \ArticleRef{65}.

\NumberedParagraph{7}{Comparison of projections and dimension}
\label{no:III-2-7}

Let \(\mathcal A\) be a factor, and let \(\varphi\) be a faithful normal
trace on \(\mathcal A^+\). If \(\mathcal A\) is purely infinite, then
\(\varphi\) is identically infinite on the set of nonzero operators of
\(\mathcal A^+\). If \(\mathcal A\) is semifinite, suppose that
\(\varphi\) is semifinite; then \(\varphi\) is determined up to
multiplication by a constant \(\lambda\) such that
\(0<\lambda<+\infty\).

\begin{definition}\label{def:III-2-2}
Let \(\mathcal A\) be a factor, and let \(\mathfrak M\) be the set of its
projections. A \emph{relative dimension} on \(\mathfrak M\) is the
restriction to \(\mathfrak M\) of a faithful normal trace on
\(\mathcal A^+\), assumed semifinite if \(\mathcal A\) is semifinite.
\end{definition}

This relative dimension \(D\) is determined up to multiplication by a
constant \(\lambda\), \(0<\lambda<+\infty\).

Let \(E\) and \(F\) be projections of \(\mathcal A\). If \(E\sim F\),
then \(D(E)=D(F)\)
(\hyperref[cor:I-6-1]{Chapter~I, \S~6, Proposition~1, Corollary~1}).
Consequently, if \(E\prec F\), then \(D(E)\leq D(F)\).

% Source PDF page 245 (printed p. 236)
\phantomsection\label{srcpage:236}
\begin{proposition}\label{prop:III-2-13}
Let \(\mathcal A\) be a factor, let \(\mathfrak M\) be the set of its
projections, and let \(D\) be a relative dimension on \(\mathfrak M\).
\begin{enumerate}[label=\textup{(\roman*)},leftmargin=*]
\item If \(E\) is a finite projection of \(\mathcal A\), then
\(D(E)<+\infty\).
\item If \(E\) is an infinite projection of \(\mathcal A\), then
\(D(E)=+\infty\).
\item If \(E\) and \(F\) are finite projections such that
\(D(E)=D(F)\) [respectively, \(D(E)\leq D(F)\)], then
\(E\sim F\) (respectively, \(E\prec F\)).
\end{enumerate}
\end{proposition}

\begin{proof}
Let \(\varphi\) be a faithful normal trace on \(\mathcal A^+\),
semifinite if \(\mathcal A\) is semifinite, whose restriction to
\(\mathfrak M\) is \(D\), and let \(E\) be a projection of
\(\mathcal A\). The algebra \(\mathcal A_E\) is a factor
(\hyperref[cor:I-2-2]{Chapter~I, \S~2, Corollary to Proposition~2}).
The trace \(\psi\) induced by
\(\varphi\) on \(\mathcal A_E^+\) is faithful and normal, and is
semifinite if \(\varphi\) is semifinite. Since two faithful normal traces
on a factor are proportional, we see that \(E\) is finite (respectively,
infinite) if and only if
\[
  D(E)=\varphi(E)=\psi(I_{E(\mathfrak H)})
\]
is finite (respectively, infinite), which proves \textup{(i)} and
\textup{(ii)}. Now let \(E\) and \(F\) be finite projections of
\(\mathcal A\) such that \(D(E)=D(F)\). Suppose, for example, that
\(E\sim E_1\leq F\)
(\hyperref[cor:III-1-thm1-1]{\S~1, Theorem~1, Corollary~1}). We have
\(D(E_1)=D(F)<+\infty\), hence \(D(F-E_1)=0\), and therefore
\(F-E_1=0\), since \(\varphi\) is faithful. Thus \(E\sim F\). Finally,
if two projections \(E,F\) of \(\mathcal A\) satisfy \(D(E)<D(F)\), then
we cannot have \(F\prec E\), and hence \(E\prec F\).
\end{proof}

\begin{proposition}\label{prop:III-2-14}
Let \(\mathcal A\) be a factor, let \(\mathfrak M\) be the set of its
projections, let \(D\) be a relative dimension on \(\mathfrak M\), and
let \(\mathfrak D\) be the set of values taken by \(D\).
\begin{enumerate}[label=\textup{(\roman*)},leftmargin=*]
\item If \(\mathcal A\) is of type \(\mathrm I_n\), where \(n\) is
finite (respectively, \(n=+\infty\)), then, by multiplying \(D\) by a
suitable scalar, one can arrange that
\[
  \mathfrak D=\{0,1,\ldots,n\}
  \quad\text{(respectively, }\mathfrak D=
  \{0,1,2,\ldots,+\infty\}\text{)}.
\]
If \(\mathcal A=\mathcal L(\mathfrak H)\), then \(D(E)\) is the
Hilbert-space dimension of \(E(\mathfrak H)\), with all infinite
cardinals identified with \(+\infty\).
\item If \(\mathcal A\) is of type \(\mathrm{II}_1\), then, by
multiplying \(D\) by a suitable scalar, one can arrange that
\(\mathfrak D=[0,1]\).
\item If \(\mathcal A\) is of type \(\mathrm{II}_\infty\), then
\(\mathfrak D=[0,+\infty]\).
\item If \(\mathcal A\) is of type \(\mathrm{III}\), then
\(\mathfrak D=\{0,+\infty\}\).
\end{enumerate}
\end{proposition}

\begin{proof}
If \(\mathcal A\) is of type \(\mathrm I_n\), we may suppose that
\(\mathcal A=\mathcal L(\mathfrak H)\), where \(\mathfrak H\) has
Hilbert-space dimension \(n\). With a suitable choice of \(D\),
\(D(E)\) is then the Hilbert-space dimension of \(E(\mathfrak H)\)
(\hyperref[thm:I-6-5]{Chapter~I, \S~6, Theorem~5}), which proves
\textup{(i)}. If
\(\mathcal A\) is of type \(\mathrm{II}_1\), then \(D(I)=1\) for a
suitable choice of \(D\). By
\hyperref[cor:III-1-thm1-3]{\S~1, Theorem~1, Corollary~3},
\(\mathfrak D\) then contains all numbers of the form
\(p2^{-n}\), where \(p,n\) are nonnegative integers and \(p\leq2^n\),
so it is dense in \([0,1]\). On the other\space%
% Source PDF page 246 (printed p. 237)
\phantomsection\label{srcpage:237}%
hand, let \((\lambda_1,\lambda_2,\ldots)\) be an increasing sequence of
numbers in \(\mathfrak D\), with limit \(\lambda\), and let \(E_i\) be
a projection of \(\mathcal A\) such that \(D(E_i)=\lambda_i\). By
induction, using \hyperref[prop:III-2-6]{Proposition~6} and
\hyperref[prop:III-2-13]{Proposition~13}, one constructs an increasing
sequence of projections \(F_i\) of \(\mathcal A\) such that
\(F_i\sim E_i\). Let \(F\) be the supremum of the \(F_i\). Then
\(D(F)=\lambda\), and hence \(\lambda\in\mathfrak D\), which proves
\textup{(ii)}. If \(\mathcal A\) is of type \(\mathrm{II}_\infty\),
there is an increasing directed family of finite projections \(E_i\) of
\(\mathcal A\) with supremum \(I\)
(\hyperref[cor:III-2-prop7-1]{Proposition~7, Corollary~1}). The
numbers \(D(E_i)\) are then arbitrarily large, and
\([0,D(E_i)]\subset\mathfrak D\) by \textup{(ii)}; therefore
\(\mathfrak D=[0,+\infty]\). Finally, \textup{(iv)} is immediate.
\end{proof}

\begin{proposition}\label{prop:III-2-15}
Let \(\mathcal A\) be a finite factor, let \(\mathfrak M\) be the set of
projections of \(\mathcal A\), and let \(D\) be a finite nonnegative
function on \(\mathfrak M\) having the following properties:
\begin{enumerate}[label=\textup{(\roman*)},leftmargin=*]
\item if \(E,F\in\mathfrak M\) are orthogonal, then
\(D(E+F)=D(E)+D(F)\);
\item if \(U\) is a unitary operator of \(\mathcal A\) and
\(E\in\mathfrak M\), then \(D(UEU^{-1})=D(E)\).
\end{enumerate}
If \(D\ne0\), then \(D\) is a relative dimension.
\end{proposition}

\begin{proof}
If \(\mathcal A\) is of type \(\mathrm I_n\), with \(n\) finite, we may
suppose that \(\mathcal A=\mathcal L(\mathfrak K)\), where
\(\mathfrak K\) is a complex Hilbert space of dimension \(n\). The
function \(D\) takes the same value \(\lambda\) on every minimal
projection. Hence, for every projection \(E\) of \(\mathcal A\),
\[
  D(E)=\lambda\dim E(\mathfrak K),
\]
which proves the proposition in this case. Suppose henceforth that
\(\mathcal A\) is continuous. Suppose that \(D(I)\ne0\), and let \(D'\)
be a relative dimension such that \(D'(I)=D(I)\). For every projection
\(E\) such that \(D'(E)\) has the form \(p2^{-n}\), where \(p,n\) are
nonnegative integers, there are equivalent, pairwise orthogonal
projections \(E_1,E_2,\ldots,E_{2^n}\) of \(\mathcal A\) such that
\[
  I=\sum_{i=1}^{2^n}E_i,
  \qquad
  E=\sum_{i=1}^{p}E_i.
\]
We have
\[
  2^nD(E_1)=D(I)=D'(I)=2^nD'(E_1),
\]
and therefore \(D(E_1)=D'(E_1)\), whence \(D(E)=D'(E)\). Now let \(F\)
be any projection of \(\mathcal A\). There is a projection
\(F_1\leq F\) such that \(D'(F_1)\) has the form \(p2^{-n}\) and
\(D'(F)-D'(F_1)\) is arbitrarily small. Since
\(D(F)\geq D(F_1)=D'(F_1)\), it follows that \(D(F)\geq D'(F)\).
Similarly, \(D(I-F)\geq D'(I-F)\); hence \(D(F)\leq D'(F)\), and finally
\(D(F)=D'(F)\).
\end{proof}

\phantomsection\label{cor:III-2-prop15}
\begin{corollary*}
Let \(\mathcal A\) be a finite factor, and let \(\psi\) be a numerical
function defined on the set \(\mathfrak N\) of Hermitian elements of
\(\mathcal A\), having the following properties:
\begin{enumerate}[label=\textup{(\roman*)},leftmargin=*]
\item \(\psi(\lambda T)=\lambda\psi(T)\) for
\(\lambda\in\mathbb R\), \(T\in\mathfrak N\);
\item \(\psi(T_1+T_2)=\psi(T_1)+\psi(T_2)\) if \(T_1,T_2\) are
commuting elements of \(\mathfrak N\);\TranslatorNote{In the French text the
second summand is printed as \(\psi(T_1)\); it has been corrected here to
\(\psi(T_2)\).}
\item \(\psi(T)\geq0\) for \(T\in\mathcal A^+\);
\item \(\psi(UTU^{-1})=\psi(T)\) if \(T\in\mathfrak N\) and \(U\) is
a unitary element of \(\mathcal A\).
\end{enumerate}%
% Source PDF page 247 (printed p. 238)
\phantomsection\label{srcpage:238}%
Let \(\varphi\) be a faithful finite normal trace on \(\mathcal A\).
Then the restriction of \(\varphi\) to \(\mathfrak N\) is proportional
to \(\psi\).
\end{corollary*}

\begin{proof}
By properties \textup{(ii)}, \textup{(iii)}, and \textup{(iv)}, there is
a constant \(\lambda\geq0\) such that \(\psi=\lambda\varphi\) on the set
of projections \(\mathfrak M\)
(\hyperref[prop:III-2-15]{Proposition~15}). Let
\(T\in\mathfrak N\), and let \(\varepsilon>0\). There are pairwise
orthogonal projections \(E_1,E_2,\ldots,E_n\) of \(\mathcal A\),
commuting with \(T\), and real numbers
\(\mu_1,\mu_2,\ldots,\mu_n\), such that
\[
  \sum_{i=1}^{n}\mu_iE_i-\varepsilon I
  \leq T\leq
  \sum_{i=1}^{n}\mu_iE_i+\varepsilon I.
\]
Taking properties \textup{(i)}, \textup{(ii)}, and \textup{(iii)} into
account, we see that \(\psi(T)\) and \(\lambda\varphi(T)\) both lie
between
\[
  \sum_{i=1}^{n}\mu_i\psi(E_i)-\varepsilon\psi(I)
  \quad\text{and}\quad
  \sum_{i=1}^{n}\mu_i\psi(E_i)+\varepsilon\psi(I).
\]
Therefore \(\psi(T)=\lambda\varphi(T)\).
\end{proof}

\hyperref[prop:III-2-14]{Proposition~14(i)} justifies the term
``relative dimension.''

\hyperref[prop:III-2-14]{Proposition~14(i) and (ii)} make it possible
to define a normalized
relative dimension when \(\mathcal A\) is discrete or finite. These two
normalizations, however, do not agree when \(\mathcal A\) is of type
\(\mathrm I_n\), \(1<n<+\infty\).

For factors of countable type,
\hyperref[prop:III-2-13]{Proposition~13(iii)} may be freed from the
finiteness hypothesis (cf. \S~8, Theorem~1, Corollary~5); for an extension
of \hyperref[prop:III-2-15]{Proposition~15}, cf.
\S~8, Theorem~1, Corollary~7.

\noindent\textit{Bibliography.} \ArticleRef{65}.

\medskip
\noindent\textit{Exercises.}
\begin{enumerate}[label=\arabic*.,leftmargin=*]
\item\label{ex:III-2-1}
Let \(\mathcal A\) be a von Neumann algebra, let \(E_1,E_2\) be
projections of \(\mathcal A\), and let \(F_1,F_2\) be their central
supports. If \(E_1\) is finite, \(E_2\) is properly infinite, and
\(F_1\leq F_2\), then \(E_1\prec E_2\). (Apply the
\hyperref[thm:III-1-1]{comparability theorem} to \(E_1\) and \(E_2\).)
\ArticleRef{65}

\item\label{ex:III-2-2}
Let \(\mathcal A\) be a finite von Neumann algebra. If there is a finite
totalizing set
\(\{x_1,x_2,\allowbreak\ldots,\allowbreak x_n\}\) for \(\mathcal A\), then
\(\mathcal A'\) is finite. (The projections \(E_{x_i}^{\mathcal A}\)
are finite by \hyperref[prop:III-2-3]{Proposition~3} and have supremum
\(I\).) \ArticleRef{65}

\item\label{ex:III-2-3}
Let \(\mathcal A\) be a von Neumann algebra. A projection \(E\) of
\(\mathcal A\) is said to be \emph{discrete} (respectively,
\emph{continuous}) if \(\mathcal A_E\) is discrete (respectively,
continuous).
\begin{enumerate}[label=\textit{\alph*.},leftmargin=2em]
\item Let \(\mathcal Z\) be the centre of \(\mathcal A\), and let \(F\)
be the largest projection of \(\mathcal Z\) that is discrete
(respectively, continuous) relative to \(\mathcal A\). In order that a
projection \(E\) of \(\mathcal A\) be discrete (respectively,
continuous), it is necessary and sufficient that \(E\leq F\).
\item If \(E\) is discrete, there is an abelian projection of
\(\mathcal A\), majorized by \(E\), having the same central support as
\(E\).
\end{enumerate}

\item\label{ex:III-2-4}
Let \(\mathcal A\) be a finite factor, let \(\mathfrak G\) be the set of
invertible elements of \(\mathcal A\), and let \(\Delta\) be a finite
strictly positive function on \(\mathfrak G\) having the following properties:
\begin{enumerate}[label=\textup{(\roman*)},leftmargin=*]
\item if \(S,T\in\mathfrak G\), then
\(\Delta(ST)=\Delta(S)\Delta(T)\);
\item if \(S\in\mathfrak G\), then \(\Delta(S^*)=\Delta(S)\);
\item \(\Delta(\lambda I)=\lambda\) for \(\lambda>0\);
\item \(\Delta(S)\leq1\) if \(0\leq S\leq I\).
\end{enumerate}%
% Source PDF page 248 (printed p. 239)
\phantomsection\label{srcpage:239}%
Show that \(\Delta\) is the determinant defined in
\hyperref[no:I-6-11]{Chapter~I, \S~6, no.~11}. [First show that
\(\Delta(U)=1\) for every unitary \(U\). For
every Hermitian element \(S\) of \(\mathcal A\), put
\[
  \omega(S)=\log\Delta(\exp S),
\]
and show, with the aid of the
\hyperref[cor:III-2-prop15]{corollary to Proposition~15}, that
\(\omega\), extended by linearity to \(\mathcal A\), is a faithful finite
normal trace on \(\mathcal A\). To show that
\(\omega(\lambda S)=\lambda\omega(S)\), where \(\lambda\) is real and
\(S\geq0\), observe that the function
\(\lambda\mapsto\omega(\lambda S)\) is additive and nonnegative for
\(\lambda\geq0\), and is therefore linear.] \ArticleRef{24}
\end{enumerate}

\section{Complements on Discrete von Neumann Algebras}
\label{sec:III-3}

\NumberedParagraph{1}{Structure of discrete von Neumann algebras}
\label{no:III-3-1}

\begin{definition}\label{def:III-3-1}
A von Neumann algebra \(\mathcal A\) is said to be \emph{homogeneous} if
there is in \(\mathcal A\) a family of equivalent, pairwise orthogonal
abelian projections with sum \(I\).
\end{definition}

Such an algebra is discrete. Every von Neumann algebra isomorphic or
anti-isomorphic to a homogeneous algebra is homogeneous. An abelian von
Neumann algebra is homogeneous. A discrete factor is homogeneous.

A homogeneous algebra is spatially isomorphic to an algebra of the form
\(\mathcal B\otimes\mathcal L(\mathfrak K)\), with \(\mathcal B\)
abelian; conversely, an algebra
\(\mathcal B\otimes\mathcal L(\mathfrak K)\), with \(\mathcal B\)
abelian, is homogeneous
(\hyperref[prop:I-2-5]{Chapter~I, \S~2, Proposition~5}). It follows that
the tensor product of two homogeneous algebras is homogeneous.

\begin{lemma}\label{lem:III-3-1}
Let \(\mathcal A\) be a von Neumann algebra, let \(E,F\) be projections
of \(\mathcal A\), and let \(E',F'\) be their central supports. If
\(E'\leq F'\) and \(E\) is abelian, then \(E\prec F\).
\end{lemma}

\begin{proof}
By \hyperref[thm:III-1-1]{\S~1, Theorem~1}, it is enough to study the
case in which
\(F\prec E\), and hence the case in which \(F\leq E\). We then have
\(F=EF'\), since \(E\) is abelian
(\hyperref[no:I-8-2]{Chapter~I, \S~8, no.~2}). Since
\(E\leq E'\leq F'\), it follows that \(F=E\).
\end{proof}

\begin{proposition}\label{prop:III-3-1}
Let \(\mathcal A\) be a homogeneous von Neumann algebra on
\(\mathfrak H\). Let \((E_\iota)_{\iota\in I}\) [respectively,
\((F_x)_{x\in K}\)] be a family of equivalent, pairwise orthogonal
abelian projections of \(\mathcal A\), with sum \(I\). Then
\(\operatorname{Card}I=\operatorname{Card}K\), provided
\(\mathfrak H\ne0\).
\end{proposition}

\begin{proof}
If \(K\) is finite, then \(\mathcal A\) is finite
(\hyperref[prop:III-2-5]{\S~2, Proposition~5}), and hence \(I\) is
finite (\hyperref[prop:III-2-10]{\S~2, Proposition~10}).
Thus \(I\) and \(K\) are either both finite or both infinite. Suppose
first that \(I\) and \(K\) are finite. The \(E_\iota\) and the \(F_x\)
have central support \(I\), and therefore \(E_\iota\sim F_x\) by
\hyperref[lem:III-3-1]{Lemma~1}. Suppose, for example, that
\(\operatorname{Card}I=\operatorname{Card}K'\), where \(K'\subset K\).
Then
\[
  I=\sum_{\iota\in I}E_\iota
  \sim\sum_{x\in K'}F_x;
\]%
% Source PDF page 249 (printed p. 240)
\phantomsection\label{srcpage:240}%
by \hyperref[prop:III-1-3]{\S~1, Proposition~3}, it follows that
\(K'=K\), and therefore
\(\operatorname{Card}I=\operatorname{Card}K\). If \(I\) and \(K\) are
infinite, then, by decomposing \(\mathcal A\) suitably as a product, we
may suppose that the algebras \(\mathcal A_{E_\iota}\) and
\(\mathcal A_{F_x}\) are of countable type (for these algebras are
abelian). \hyperref[lem:III-1-6]{Lemma~6 of \S~1} then gives
\(\operatorname{Card}I\leq\operatorname{Card}K\) and
\(\operatorname{Card}K\leq\operatorname{Card}I\).
\end{proof}

\hyperref[prop:III-3-1]{Proposition~1} proves that
\(n=\operatorname{Card}I\) is an algebraic
invariant of \(\mathcal A\). One says that \(\mathcal A\) is of
\emph{type \(\mathrm I_n\)} (which generalizes the terminology adopted
in \hyperref[no:I-8-4]{Chapter~I, \S~8, no.~4} when \(\mathcal A\) is
a factor). If
\(\mathcal A=\mathcal B\otimes\mathcal L(\mathfrak K)\), with
\(\mathcal B\) abelian, then \(n\) is the Hilbert-space dimension of
\(\mathfrak K\). In order that \(\mathcal A\) be finite, it is necessary
and sufficient that \(n\) be finite. If \(\mathcal A\) is of type
\(\mathrm I_n\), one also says that \(n\) is the \emph{multiplicity} of
\(\mathcal A'\). However, \(n\) is not an algebraic invariant of
\(\mathcal A'\).

\begin{proposition}\label{prop:III-3-2}
Every discrete von Neumann algebra \(\mathcal A\) has the form
\[
  \prod_{\iota\in I}\mathcal A_\iota,
\]
where the \(\mathcal A_\iota\) are of types
\(\mathrm I_{n_\iota}\), with the \(n_\iota\) pairwise distinct.
\end{proposition}

\begin{proof}
Apply \hyperref[cor:III-1-thm1-2]{\S~1, Theorem~1, Corollary~2}, taking
for the family
\((E_\iota)_{\iota\in I}\) a family consisting of one nonzero abelian
projection of \(\mathcal A\). With the notation of that corollary, the
\(F_x\) are nonzero, abelian, pairwise orthogonal, and equivalent, and
\(F_0\) is abelian. If the central support of \(F_0\) were equal to that
of the \(F_x\), then \(F_0\sim F_x\) by
\hyperref[lem:III-3-1]{Lemma~1}, contrary to the
corollary being used. Hence there is a projection \(E\) in the centre of
\(\mathcal A\), majorized by \(G\), orthogonal to \(F_0\), and such that
the \(F_xE\) are nonzero. The equality
\[
  G=F_0+\sum_{x\in K}F_x
\]
gives
\[
  E=\sum_{x\in K}F_xE.
\]
Thus \(\mathcal A_E\) is homogeneous and nonzero. It is now enough to
repeat the argument on \(\mathcal A_{I-E}\), which is discrete, and to
continue in this way transfinitely.
\end{proof}

\phantomsection\label{cor:III-3-prop2}
\begin{corollary*}
Every discrete von Neumann algebra has the form
\[
  \prod_{\iota\in I}
  \bigl(\mathcal B_\iota\otimes\mathbb C_{\mathfrak H_\iota}\bigr),
\]
where the \(\mathcal B_\iota\) are abelian.
\end{corollary*}

\begin{proof}
If \(\mathcal A\) is discrete, then \(\mathcal A'\) is discrete, and
hence
\[
  \mathcal A'
  =\prod_{\iota\in I}
    \bigl(\mathcal C_\iota\otimes\mathcal L(\mathfrak H_\iota)\bigr)
\]
with the \(\mathcal C_\iota\) abelian. Therefore
\[
  \mathcal A
  =\prod_{\iota\in I}
    \bigl(\mathcal C'_\iota\otimes\mathbb C_{\mathfrak H_\iota}\bigr),
\]
and it is enough to put \(\mathcal C'_\iota=\mathcal B_\iota\).
\end{proof}

\noindent\textit{Bibliography.}
\ArticleRef{10}, \ArticleRef{43}, \ArticleRef{61}, \ArticleRef{100}.

\NumberedParagraph{2}{Isomorphisms of discrete von Neumann algebras}
\label{no:III-3-2}

\begin{proposition}\label{prop:III-3-3}
Let \(\mathcal A,\mathcal B\) be von Neumann algebras such that
\(\mathcal A'\) and \(\mathcal B'\) are homogeneous. Suppose that
\(\mathcal A\) and \(\mathcal B\) have the same multiplicity. Then every
isomorphism \(\Phi\) of \(\mathcal A\) onto \(\mathcal B\) is spatial.
\end{proposition}

% Source PDF page 250 (printed p. 241)
\phantomsection\label{srcpage:241}
\begin{proof}
There are
(\hyperref[cor:I-4-3]{Chapter~I, \S~4, Corollary to Theorem~3}) a
von Neumann algebra
\(\mathcal C\) and projections \(E',F'\) of \(\mathcal C'\), with
central support \(I\), such that \(\mathcal A\) is identified with
\(\mathcal C_{E'}\), \(\mathcal B\) with \(\mathcal C_{F'}\), and
\(\Phi\) with the isomorphism
\[
  T_{E'}\longmapsto T_{F'}
  \qquad(T\in\mathcal C).
\]
Let \((E'_\iota)_{\iota\in I}\) [respectively,
\((F'_x)_{x\in K}\)] be a family of equivalent, pairwise orthogonal,
abelian projections of \(\mathcal C'\), with sum \(E'\) (respectively,
\(F'\)). We have \(E'_\iota\sim F'_x\) by
\hyperref[lem:III-3-1]{Lemma~1}. Since
\(\mathcal A=\mathcal C_{E'}\) and
\(\mathcal B=\mathcal C_{F'}\) have the same multiplicity, \(I\) and
\(K\) are equipotent. Therefore \(E'\sim F'\).
\end{proof}

\setcounter{corollary}{0}
\begin{corollary}\label{cor:III-3-prop3-1}
Let \(\mathcal A\) and \(\mathcal B\) be two von Neumann algebras such
that \(\mathcal A'\) and \(\mathcal B'\) are abelian. Every isomorphism
of \(\mathcal A\) onto \(\mathcal B\) is spatial.
\end{corollary}

\begin{proof}
The algebras \(\mathcal A\) and \(\mathcal B\) have multiplicity \(1\).
\end{proof}

\begin{corollary}\label{cor:III-3-prop3-2}
Let \(\mathcal A\) be a discrete von Neumann algebra, let \(\mathcal Z\)
be its centre, and let \(\Phi\) be an automorphism of \(\mathcal A\) that
fixes the elements of \(\mathcal Z\). There is a unitary operator
\(U\in\mathcal A\) such that
\[
  \Phi(T)=UTU^{-1}
  \qquad(T\in\mathcal A).
\]
\end{corollary}

\begin{proof}
By replacing \(\mathcal A\) by an isomorphic von Neumann algebra, we may
suppose that \(\mathcal A'\) is commutative, and hence equal to
\(\mathcal Z\). By
\hyperref[cor:III-3-prop3-1]{Corollary~1}, there is a unitary operator
\(U\) such
that \(\Phi(T)=UTU^{-1}\) for every \(T\in\mathcal A\). If
\(T\in\mathcal Z\), then
\(UTU^{-1}=\Phi(T)=T\); hence
\[
  U\in\mathcal Z'=\mathcal A''=\mathcal A.
\]
\end{proof}

\noindent\textit{Bibliography.}
\ArticleRef{43}, \ArticleRef{44}, \ArticleRef{61}, \ArticleRef{125}.

\medskip
\noindent\textit{Exercises.}
\begin{enumerate}[label=\arabic*.,leftmargin=*]
\item\label{ex:III-3-1}
Let \(\mathcal A\) be a discrete von Neumann algebra, and let
\(\mathcal Z\) be its centre. If \(\mathcal A\) is of countable type over
\(\mathcal Z\)
(\hyperref[ex:I-9-5]{Chapter~I, \S~9, Exercise~5}), then \(\mathcal A\)
is
generated by \(\mathcal Z\) and a countable family of elements. [By
\hyperref[prop:III-3-2]{Proposition~2}, reduce to the case in which
\(\mathcal A\) is homogeneous,
and then to the case in which \(\mathcal Z\) is also of countable type.
Then
\(\mathcal A=\mathcal B\otimes\mathcal L(\mathfrak K)\), with
\(\mathcal B\) abelian, and \(\mathcal A\) is of countable type. Deduce
from \hyperref[ex:I-2-3]{Chapter~I, \S~2, Exercise~3} that
\(\mathfrak K\) has a countable
basis.]

(Cf. \hyperref[cor:II-3-2]{Chapter~II, \S~3, Theorem~1, Corollary~2}.)

\item\label{ex:III-3-2}
Let \(\mathcal A\) be a von Neumann algebra. In order that \(\mathcal A\)
and \(\mathcal A'\) be homogeneous of types \(\mathrm I_n\) and
\(\mathrm I_{n'}\), it is necessary and sufficient that
\[
  \mathcal A
    =\mathcal B\otimes\mathcal L(\mathfrak K)
       \otimes\mathbb C_{\mathfrak K'},
  \qquad
  \mathcal A'
    =\mathcal B\otimes\mathbb C_{\mathfrak K}
       \otimes\mathcal L(\mathfrak K'),
\]
where \(\mathcal B\) is an abelian von Neumann algebra such that
\(\mathcal B=\mathcal B'\), and \(\mathfrak K\) (respectively,
\(\mathfrak K'\)) is a complex Hilbert space of Hilbert-space dimension
\(n\) (respectively, \(n'\)). Deduce that the von Neumann algebra
generated by \(\mathcal A\) and \(\mathcal A'\) is of type
\(\mathrm I_{nn'}\). \ArticleRef{100}

\item\label{ex:III-3-3}
A ring is said to be \emph{binormal} if, for every four elements
\(x_1,x_2,x_3,x_4\) of the ring,
\[
  \sum_{\sigma}\varepsilon_\sigma
    x_{\sigma(1)}x_{\sigma(2)}x_{\sigma(3)}x_{\sigma(4)}=0,
\]
where \(\sigma\) ranges over the permutations of \(\{1,2,3,4\}\), and
\(\varepsilon_\sigma\in\{+1,-1\}\) is the signature of \(\sigma\).
\begin{enumerate}[label=\textit{\alph*.},leftmargin=2em]
\item The ring of square matrices of degree \(2\) over a commutative ring
is binormal. By linearity, it is enough to consider the case
\[
  x_1=\begin{pmatrix}1&0\\0&0\end{pmatrix},\quad
  x_2=\begin{pmatrix}0&0\\1&0\end{pmatrix},\quad
  x_3=\begin{pmatrix}0&1\\0&0\end{pmatrix},\quad
  x_4=\begin{pmatrix}0&0\\0&1\end{pmatrix}.
\]

% Source PDF page 251 (printed p. 242)
\phantomsection\label{srcpage:242}
\item The ring of square matrices of degree \(3\) over \(\mathbb C\) is
not binormal. Consider the matrices
\[
  x_1=\begin{pmatrix}1&0&0\\0&0&0\\0&0&0\end{pmatrix},\quad
  x_2=\begin{pmatrix}0&0&0\\1&0&0\\0&0&0\end{pmatrix},
\]
\[
  x_3=\begin{pmatrix}0&0&1\\0&0&0\\0&0&0\end{pmatrix},\quad
  x_4=\begin{pmatrix}0&0&0\\0&0&0\\0&0&1\end{pmatrix}.
\]
\item Deduce from \textit{b} that a von Neumann algebra containing three
equivalent, nonzero, pairwise orthogonal projections is not binormal.
\item Deduce from \textit{a} and \textit{c} that a von Neumann algebra is
binormal if and only if it is a product of an abelian von Neumann algebra
and a von Neumann algebra of type \(\mathrm I_2\). \ArticleRef{3}
\end{enumerate}

\item\label{ex:III-3-4}
A von Neumann algebra \(\mathcal A\) such that \(\mathcal A=\mathcal A'\)
is spatially isomorphic to a von Neumann algebra of the type defined in
\hyperref[thm:I-7-2]{Chapter~I, \S~7, Theorem~2}. (Use
\hyperref[thm:I-7-1]{Theorem~1} and
\hyperref[thm:I-7-2]{Theorem~2 of Chapter~I, \S~7}, and
\hyperref[cor:III-3-prop3-1]{Corollary~1 of Proposition~3 of the
present section}.)
\ArticleRef{100}

\item\label{ex:III-3-5}
A von Neumann algebra \(\mathcal A\) on \(\mathfrak H\) is said to be
\emph{uniform} if there is in \(\mathcal A\) a family
\((E_\iota)_{\iota\in I}\) of equivalent, pairwise orthogonal, finite
projections with sum \(I\).
\begin{enumerate}[label=\textit{\alph*.},leftmargin=2em]
\item In order that a von Neumann algebra be uniform, it is necessary and
sufficient that it be spatially isomorphic to an algebra
\(\mathcal B\otimes\mathcal L(\mathfrak K)\), with \(\mathcal B\)
finite.
\item Let \(\mathcal A\) be a uniform von Neumann algebra.
Let \((E_\iota)_{\iota\in I}\) and, respectively,
\((F_x)_{x\in K}\) be families of equivalent, pairwise orthogonal,
finite projections of \(\mathcal A\), each with sum \(I\). If \(I\) is
infinite, then \(K\) is infinite and
\(\operatorname{Card}I=\operatorname{Card}K\), provided
\(\mathfrak H\ne0\). (Argue much as in
\hyperref[prop:III-3-1]{Proposition~1}.) The number
\(\operatorname{Card}I\) is then called the \emph{order} of
\(\mathcal A\).
\item Let \(\mathcal A\) be a properly infinite von Neumann algebra,
let \(\mathcal Z\) be its centre, and let \(E\) be a nonzero finite
projection of \(\mathcal A\). There are a nonzero projection
\(G\in\mathcal Z\) and a family \((E_\iota)\) of projections equivalent
to \(EG\), pairwise orthogonal, with sum \(G\). (Use
\hyperref[cor:III-1-thm1-2]{Corollary~2 of Theorem~1, \S~1}.)
\item Every properly infinite semifinite von Neumann algebra is a product
of uniform von Neumann algebras of pairwise distinct orders. (Use
\textit{a}, \textit{b}, and \textit{c}.)
\item Every semifinite von Neumann algebra \(\mathcal A\) has the form
\[
  \prod_{\iota\in I}
  \bigl(\mathcal B_\iota\otimes\mathbb C_{\mathfrak K_\iota}\bigr),
\]
where the \(\mathcal B_\iota\) are finite. (Apply \textit{d} to
\(\mathcal A'\).)
\item Let \(\mathcal A\) be a properly infinite uniform von Neumann
algebra, let \(\mathcal Z\) be its centre, and let \(E\) be a finite
projection of \(\mathcal A\). There is a family \((E_\iota)\) of
projections equivalent to \(E\), pairwise orthogonal, whose sum belongs
to \(\mathcal Z\). (Use \textit{c}.)
\item Let \(\mathcal A,\mathcal B\) be von Neumann algebras such that
\(\mathcal A'\) and \(\mathcal B'\) are properly infinite and uniform,
of the same order. Every isomorphism of \(\mathcal A\) onto
\(\mathcal B\) is spatial. (Use a proof analogous to that of
\hyperref[prop:III-3-3]{Proposition~3}, making use of \textit{b} and
\textit{f}.)
\item Let \(\mathcal A\) be a properly infinite uniform von Neumann
algebra. In order that \(\mathcal A\) be standard, it is necessary and
sufficient that \(\mathcal A'\) be properly infinite and uniform, of the
same order as \(\mathcal A\). (To prove that the condition is sufficient,
observe that \(\mathcal A\) is isomorphic to a standard von Neumann
algebra \(\mathcal B\), and use \textit{g}.) \ArticleRef{31}
\ArticleRef{89}
\end{enumerate}

% Source PDF page 252 (printed p. 243)
\phantomsection\label{srcpage:243}
\item\label{ex:III-3-6}
Let \(Z\) be a Borel space, let \(\nu\) be a positive measure on \(Z\),
let \(\zeta\mapsto\mathfrak H(\zeta)\) be a \(\nu\)-measurable field
of complex Hilbert spaces on \(Z\), and put
\[
  \mathfrak H
   =\int^{\oplus}\mathfrak H(\zeta)\,d\nu(\zeta).
\]
Let \(\zeta\mapsto\mathcal A(\zeta)\) be a \(\nu\)-measurable field of
von Neumann algebras on \(Z\), and put
\[
  \mathcal A
   =\int^{\oplus}\mathcal A(\zeta)\,d\nu(\zeta).
\]
If \(\mathcal A\) is of type \(\mathrm I_n\), where
\(n=1,2,\ldots,\aleph_0\), then \(\mathcal A(\zeta)\) is of type
\(\mathrm I_n\) almost everywhere. (Decompose \(I\) into equivalent,
orthogonal, abelian projections of \(\mathcal A\).) The converse is true
if \(\nu\) is standard. \ArticleRef{10}, \ArticleRef{80}

\item\label{ex:III-3-7}
Prove \hyperref[thm:II-6-2]{Chapter~II, \S~6, Theorem~2} by the
following method. Reduce to the
case in which \(\mathcal Z'\) is homogeneous; then
\[
  \mathcal Z'=\mathcal B\otimes\mathcal L(\mathfrak K),
  \qquad \mathcal B=\mathcal B'.
\]
Thus, by \hyperref[ex:III-3-4]{Exercise~4}, there are a compact
metrizable space \(Z\) and a
positive measure \(\nu\) on \(Z\) such that \(\mathcal B\) is the algebra
of multiplication operators by the functions in
\(L_{\mathbb C}^{\infty}(Z,\nu)\) on
\(\mathfrak H_0=L_{\mathbb C}^{2}(Z,\nu)\). Then
\[
  \mathfrak H=\mathfrak H_0\otimes\mathfrak K
  =\int^{\oplus}\mathfrak H(\zeta)\,d\nu(\zeta),
\]
where \(\zeta\mapsto\mathfrak H(\zeta)\) is the constant field
corresponding to \(\mathfrak K\) on \(Z\), and
\(\mathcal Z=\mathcal B\otimes\mathbb C_{\mathfrak K}\) is the algebra
of diagonalizable operators on \(\mathfrak H\).

\item\label{ex:III-3-8}
Let \(\mathcal A\) be a discrete von Neumann algebra on a Hilbert space
\(\mathfrak H\) with a countable basis. Show that \(\mathcal A\) is
generated by a single element. [Using
\hyperref[ex:I-7-3]{Chapter~I, \S~7, Exercises~3\textit{f},
\textit{h}, and \textit{i}}, reduce to the case in
which \(\mathcal A\) is homogeneous, and then to the case
\(\mathcal A=\mathcal L(\mathfrak H)\). Observe that
\(\mathcal L(\mathfrak H)\) can be generated by a compact Hermitian
operator whose eigenvalues have multiplicity \(1\), and by a rank-one
projection.] \ArticleRef{288}
\end{enumerate}

\section{Operator-Valued Traces}\label{sec:III-4}

\NumberedParagraph{1}{Definition}
\label{no:III-4-1}

Let \(\mathcal A\) be a von Neumann algebra, and let \(\mathcal Z\) be a
von Neumann algebra contained in the centre of \(\mathcal A\). Throughout
this section we choose, once and for all, a locally compact space \(Z\),
a positive measure \(\nu\) on \(Z\), and an isomorphism of the normed
involutive algebra \(L_{\mathbb C}^{\infty}(Z,\nu)\) onto \(\mathcal Z\)
(\hyperref[thm:I-7-1]{Chapter~I, \S~7, Theorem~1}); we identify
\(L_{\mathbb C}^{\infty}(Z,\nu)\) with \(\mathcal Z\) by this
isomorphism. This identification is compatible with the natural order
relations. Thus \(\mathcal Z^+\) is embedded in the set
\(\widehat{\mathcal Z}_+\) of measurable nonnegative functions, finite
or not, on \(Z\). (In \(\widehat{\mathcal Z}_+\), as in
\(L_{\mathbb C}^{\infty}(Z,\nu)\), two functions equal locally almost
everywhere are identified.) Every increasing directed set that is
bounded above in \(\mathcal Z^+\) has a supremum; it follows readily that
every increasing directed set in \(\widehat{\mathcal Z}_+\) has a
supremum.
% Source PDF page 253 (printed p. 244)
\phantomsection\label{srcpage:244}
\begin{definition}\label{def:III-4-1}
A trace associated with \(\mathcal Z\) and defined on \(\mathcal A_+\)
(or a \(\mathcal Z\)-trace on \(\mathcal A_+\)) is a mapping \(\Phi\),
defined on \(\mathcal A_+\), with values in
\(\widehat{\mathcal Z}_+\), having the following properties:
\begin{enumerate}[label=\textup{(\roman*)},leftmargin=*]
\item if \(S\in\mathcal A_+\) and \(T\in\mathcal A_+\), then
\(\Phi(S+T)=\Phi(S)+\Phi(T)\);
\item if \(S\in\mathcal A_+\) and \(T\in\mathcal Z_+\), then
\(\Phi(TS)=T\Phi(S)\);
\item if \(S\in\mathcal A_+\) and \(U\) is a unitary operator in
\(\mathcal A\), then
\(\Phi(USU^{-1})=\Phi(S)\).
\end{enumerate}

The trace \(\Phi\) is said to be \emph{faithful} if
\(S\in\mathcal A_+\) and \(\Phi(S)=0\) imply \(S=0\).

It is said to be \emph{finite} if
\(\Phi(S)\in\mathcal Z_+\) for every \(S\in\mathcal A_+\).

It is said to be \emph{semifinite} if, for every nonzero
\(S\in\mathcal A_+\), there is a nonzero \(T\in\mathcal A_+\),
majorized by \(S\), such that \(\Phi(T)\in\mathcal Z_+\).

It is said to be \emph{normal} if, for every increasing directed set
\(\mathcal F\subset\mathcal A_+\) with supremum
\(S\in\mathcal A_+\), \(\Phi(S)\) is the supremum of
\(\Phi(\mathcal F)\).

If \(\mathcal Z\) is the set of scalar operators, a
\(\mathcal Z\)-trace on \(\mathcal A_+\) is of the form
\(S\mapsto\varphi(S)I\), where \(\varphi\) is a trace on
\(\mathcal A_+\).
\end{definition}

\begin{proposition}\label{prop:III-4-1}
Let \(\Phi\) be a \(\mathcal Z\)-trace on \(\mathcal A_+\). The set
of \(S\in\mathcal A_+\) such that \(\Phi(S)\in\mathcal Z_+\) is the
positive part of an ideal \(\mathfrak m\) of \(\mathcal A\). There is a
unique linear mapping \(\dot\Phi\) from \(\mathfrak m\) into
\(\mathcal Z\) which agrees with \(\Phi\) on \(\mathfrak m_+\), and
\[
  \dot\Phi(ST)=\dot\Phi(TS)
  \qquad(S\in\mathfrak m,\ T\in\mathcal A).
\]
Finally, let \(S\in\mathfrak m\). If \(\Phi\) is normal, the linear
mapping \(T\mapsto\dot\Phi(ST)\) from \(\mathcal A\) into
\(\mathcal Z\) is ultra-strongly and ultra-weakly continuous, and its
restriction to bounded subsets of \(\mathcal A\) is strongly and weakly
continuous.
\end{proposition}

\begin{proof}
Proceed as in Chapter~I, \S~6, Proposition~1. Instead of using
Chapter~I, \S~4, Theorem~1, use
\hyperref[thm:I-4-2]{Chapter~I, \S~4, Theorem~2}. We leave the details
to the reader.
\end{proof}

The mapping \(\dot\Phi\) is sometimes called a
\(\mathcal Z\)-trace on \(\mathfrak m\), when no confusion can result.

Bibliography: \ArticleRef{6}, \ArticleRef{12}, \ArticleRef{63},
\ArticleRef{101}, \ArticleRef{104}, \ArticleRef{105}.

\NumberedParagraph{2}{Traces on \(\widehat{\mathcal Z}_+\)}
\label{no:III-4-2}
A finite or infinite nonnegative function \(\omega\) on
\(\widehat{\mathcal Z}_+\) will be called a \emph{trace on
\(\widehat{\mathcal Z}_+\)} if
\[
  \omega(S+T)=\omega(S)+\omega(T),
  \qquad \omega(\lambda S)=\lambda\omega(S)
\]
for \(S,T\in\widehat{\mathcal Z}_+\) and \(\lambda\geq0\) (we agree
once and for all that \(0\cdot(+\infty)=0\)). By analogy with the
definition in Chapter~I, \S~6, faithful, semifinite, and normal traces on
\(\widehat{\mathcal Z}_+\) are defined in the evident way. Let
\(\varphi\) be a normal trace on \(\mathcal Z_+\); for every
\(S\in\widehat{\mathcal Z}_+\), put
\[
  \omega(S)=\sup_i\varphi(S_i),
\]
where \((S_i)\) is the family of functions in \(\mathcal Z_+\)
majorized by \(S\). It is clear that \(\omega\) is a normal trace on
\(\widehat{\mathcal Z}_+\), faithful (respectively, semifinite) if and
only if \(\varphi\) is faithful (respectively, semifinite); and\space%
% Source PDF page 254 (printed p. 245)
\phantomsection\label{srcpage:245}%
every normal trace \(\omega\) on \(\widehat{\mathcal Z}_+\) is obtained
in this way, \(\varphi\) being simply the restriction of \(\omega\) to
\(\mathcal Z_+\). Henceforth we identify \(\omega\) and \(\varphi\).

Let \(S\in\widehat{\mathcal Z}_+\). The number
\(\int S(\zeta)\,d\nu(\zeta)\) does not change if \(S\) is modified on
a locally negligible set. Put
\[
  \omega(S)=\int S(\zeta)\,d\nu(\zeta).
\]
Then \(\omega\) is a faithful semifinite normal trace on
\(\widehat{\mathcal Z}_+\).

\begin{lemma}\label{lem:III-4-1}
Let \(\omega\) and \(\omega_1\) be two normal traces on
\(\widehat{\mathcal Z}_+\), with \(\omega\) faithful and semifinite.
There is a unique \(Q\in\widehat{\mathcal Z}_+\) such that
\[
  \omega_1(S)=\omega(SQ)
  \qquad\text{for every }S\in\widehat{\mathcal Z}_+.
\]
If \(\omega_1\) is faithful (respectively, semifinite), then
\(Q(\zeta)>0\) (respectively, \(Q(\zeta)<+\infty\)) locally almost
everywhere on \(Z\).
\end{lemma}

\begin{proof}
By Chapter~I, \S~6, Proposition~1, Corollary~2, there is a projection
\(E\in\mathcal Z\), identified with the characteristic function of a
measurable subset \(Y\) of \(Z\), having the following properties:
\begin{enumerate}[label=\textup{(\roman*)},leftmargin=*]
\item if \(S\in\widehat{\mathcal Z}_+\) is not locally almost everywhere
zero on \(Z-Y\), then \(\omega_1(S)=+\infty\);
\item the trace \(\omega'_1\) induced by \(\omega_1\) on
\(\widehat{\mathcal Z}_E^+\) is semifinite.
\end{enumerate}
Let \(\omega'\) be the trace induced by \(\omega\) on
\(\widehat{\mathcal Z}_E^+\), and put
\(\psi=\omega'+\omega'_1\), which is faithful and semifinite. There are
(Chapter~I, \S~6, Theorem~3) elements
\(Q_0,Q_1\in\widehat{\mathcal Z}_E^+\) such that
\[
  \omega'(S)=\psi(SQ_0),
  \qquad
  \omega'_1(S)=\psi(SQ_1)
  \qquad(S\in\widehat{\mathcal Z}_E^+).
\]
Identify \(\widehat{\mathcal Z}_E^+\) with
\(\widehat{\mathcal Z}_+E\), that is, with the functions in
\(\widehat{\mathcal Z}_+\) which vanish on \(Z-Y\). Then
\(Q'=Q_1Q_0^{-1}\) is a function in \(\widehat{\mathcal Z}_+\) which
vanishes on \(Z-Y\). For \(S\in\widehat{\mathcal Z}_E^+\),
\[
  \omega'_1(S)=\psi(SQ_1)
  =\psi(SQ_1Q_0^{-1}Q_0).
\]
Since \(\omega\) is faithful, \(Q_0(\zeta)>0\) locally almost
everywhere on \(Y\), and therefore
\(Q_0(\zeta)^{-1}Q_0(\zeta)=1\); consequently
\[
  \omega'_1(S)=\psi(SQ'Q_0)=\omega'(SQ').
\]
Put \(Q(\zeta)=Q'(\zeta)\) for \(\zeta\in Y\), and
\(Q(\zeta)=+\infty\) for \(\zeta\in Z-Y\). Then
\[
  \omega_1(S)=\omega(SQ)
  \qquad\text{for every }S\in\widehat{\mathcal Z}_+.
\]

Let us prove the uniqueness of \(Q\). Suppose that
\(Q'\in\widehat{\mathcal Z}_+\) and
\(\omega(SQ)=\omega(SQ')\) for every
\(S\in\widehat{\mathcal Z}_+\). If \(Q\neq Q'\), there is a measurable
set \(X\subset Z\), not locally negligible, on which, for example,
\(Q(\zeta)<Q'(\zeta)\). Hence there is an
\(S\in\widehat{\mathcal Z}_+\) such that
\[
  SQ\leq SQ',\qquad SQ\neq SQ',\qquad SQ\in\mathcal Z_+,
  \qquad \omega(SQ)<+\infty,
\]
the last two conditions being possible because \(\omega\) is semifinite.
Then
\[
  \omega(SQ'-SQ)=\omega(SQ')-\omega(SQ)=0,
\]
contrary to the faithfulness of \(\omega\).

Finally, if \(Q(\zeta)=0\) (respectively, \(Q(\zeta)=+\infty\)) on a
set which is not locally negligible, then \(\omega_1\) is not faithful
(respectively, is not semifinite), which proves the last assertion of the
lemma.
\end{proof}

% Source PDF page 255 (printed p. 246)
\phantomsection\label{srcpage:246}
\NumberedParagraph{3}{Relations Between Scalar Traces and
Operator-Valued Traces}
\label{no:III-4-3}

\begin{proposition}\label{prop:III-4-2}
Let \(\mathcal A\) be a von Neumann algebra, let \(\mathcal Z\) be a
von Neumann algebra contained in the centre of \(\mathcal A\), let
\(\Phi\) be a normal \(\mathcal Z\)-trace on \(\mathcal A_+\), and let
\(\omega\) be a normal trace on \(\widehat{\mathcal Z}_+\).
\begin{enumerate}[label=\textup{(\roman*)},leftmargin=*]
\item \(\varphi=\omega\circ\Phi\) is a normal trace on
\(\mathcal A_+\).
\item The trace \(\varphi\) is faithful if and only if \(\omega\) and
\(\Phi\) are faithful.
\item If \(\omega\) and \(\Phi\) are semifinite, then \(\varphi\) is
semifinite.
\item If \(\varphi\) is semifinite and \(\omega\) faithful, then
\(\Phi\) is semifinite.
\item If \(\varphi\) is semifinite and \(\Phi\) faithful, then
\(\omega\) is semifinite.
\end{enumerate}
\end{proposition}

\begin{proof}
Assertion \textup{(i)} is immediate.

\textup{(ii)} If \(\omega\) and \(\Phi\) are faithful, then
\(\varphi\) is faithful. Conversely, suppose that \(\varphi\) is
faithful. It is clear that \(\Phi\) is faithful. Now let \(T\) be a
nonzero element of \(\mathcal Z_+\). We have
\[
  \omega\bigl(T\Phi(I)\bigr)=\omega\bigl(\Phi(T)\bigr)
  =\varphi(T)>0.
\]
But \(\Phi(I)\) is the supremum of an increasing directed family of
elements of \(\mathcal Z_+\); hence \(\omega(TS)>0\) for some
\(S\in\mathcal Z_+\), and therefore \(\omega(T)>0\). Thus \(\omega\)
is faithful.

\textup{(iii)} Suppose that \(\omega\) and \(\Phi\) are semifinite, and
let \(S\) be a nonzero element of \(\mathcal A_+\). There is a nonzero
\(S_1\in\mathcal A_+\) such that
\(S_1\leq S\) and \(\Phi(S_1)\in\mathcal Z_+\). If
\(\Phi(S_1)=0\), then \(\varphi(S_1)=0\), and \textup{(iii)} is proved.
If \(\Phi(S_1)\neq0\), there is a nonzero \(T\in\mathcal Z_+\),
majorized by \(\Phi(S_1)\), such that \(\omega(T)<+\infty\). We may
write \(T=T_1\Phi(S_1)\), where \(T_1\in\mathcal Z_+\) and
\(T_1\leq I\). Put
\[
  S_2=S_1T_1\in\mathcal A_+.
\]
Since \(\Phi(S_2)=T_1\Phi(S_1)=T\), we see that
\(S_2\neq0\), \(S_2\leq S_1\leq S\), and
\(\varphi(S_2)=\omega(T)<+\infty\). Thus \(\varphi\) is semifinite.

\textup{(iv)} Suppose that \(\varphi\) is semifinite and \(\omega\)
faithful. Let \(S\) be a nonzero element of \(\mathcal A_+\). There is a
nonzero \(S_1\in\mathcal A_+\) such that \(S_1\leq S\) and
\[
  \varphi(S_1)=\omega\bigl(\Phi(S_1)\bigr)<+\infty.
\]
Since \(\omega\) is faithful, this requires
\(\Phi(S_1)<+\infty\) locally almost everywhere. The proof is complete
if \(\Phi(S_1)=0\); hence we may suppose that
\(0<\Phi(S_1)\leq\lambda<+\infty\) on a measurable set which is not
locally negligible. There is then a projection \(E\in\mathcal Z\) such
that \(E\Phi(S_1)\neq0\) and \(E\Phi(S_1)\in\mathcal Z\). Put
\(S_2=ES_1\). Since \(\Phi(S_2)=E\Phi(S_1)\), we have
\(S_2\neq0\), \(S_2\leq S\), and \(\Phi(S_2)\in\mathcal Z\). Thus
\(\Phi\) is semifinite.

\textup{(v)} Suppose that \(\Phi\) is faithful and \(\omega\) is not
semifinite. There is a nonzero \(T\in\mathcal Z_+\) such that
\(\omega(T_1)=+\infty\) for every nonzero \(T_1\in\mathcal Z_+\)
majorized by \(T\). Since
\(\Phi(T_1)=T_1\Phi(I)\) vanishes, apart from a locally negligible set,
on the same set as \(T_1\), we have
\[
  \varphi(T_1)=\omega\bigl(\Phi(T_1)\bigr)=+\infty.
\]
Thus \(\varphi\) is not semifinite.
\end{proof}

% Source PDF page 256 (printed p. 247)
\phantomsection\label{srcpage:247}
\begin{proposition}\label{prop:III-4-3}
Let \(\mathcal A\) be a von Neumann algebra, let \(\mathcal Z\) be a
von Neumann algebra contained in the centre of \(\mathcal A\), let
\(\varphi\) be a normal trace on \(\mathcal A_+\), and let \(\omega\)
be a faithful semifinite normal trace on \(\widehat{\mathcal Z}_+\).
There is a unique normal \(\mathcal Z\)-trace \(\Phi\) on
\(\mathcal A_+\) such that \(\varphi=\omega\circ\Phi\).
\end{proposition}

\begin{proof}
Let \(S\in\mathcal A_+\). The mapping
\(T\mapsto\varphi(ST)\), where \(T\) ranges over \(\mathcal Z_+\), is
a normal trace on \(\mathcal Z_+\), hence on
\(\widehat{\mathcal Z}_+\). By Lemma~1, there is a unique
\(\Phi(S)\in\widehat{\mathcal Z}_+\) such that
\[
  \varphi(ST)=\omega\bigl(\Phi(S)T\bigr)
  \qquad(T\in\mathcal Z_+).
\]
Let \(S,S_1\in\mathcal A_+\), let \(U\) be a unitary element of
\(\mathcal A\), and let \(S_2,T\in\mathcal Z_+\). We have
\begin{align*}
\omega\bigl(\Phi(S+S_1)T\bigr)
  &=\varphi((S+S_1)T)
   =\varphi(ST)+\varphi(S_1T)\\
  &=\omega\bigl((\Phi(S)+\Phi(S_1))T\bigr),\\
\omega\bigl(\Phi(SS_2)T\bigr)
  &=\varphi(SS_2T)
   =\omega\bigl(\Phi(S)S_2T\bigr),\\
\omega\bigl(\Phi(USU^{-1})T\bigr)
  &=\varphi(USU^{-1}T)
   =\varphi(SU^{-1}TU)\\
  &=\varphi(ST)=\omega\bigl(\Phi(S)T\bigr).
\end{align*}
Thus \(\Phi\) is a \(\mathcal Z\)-trace. Let us show that \(\Phi\) is
normal. If \((S_i)\) is an increasing directed family in
\(\mathcal A_+\), with supremum \(S\in\mathcal A_+\), then the
\(\Phi(S_i)\) form an increasing directed family in
\(\widehat{\mathcal Z}_+\), with supremum \(Q\). For
\(T\in\mathcal Z_+\), the equality
\(\varphi(S_iT)=\omega(\Phi(S_i)T)\), on passing to the limit, gives
\(\varphi(ST)=\omega(QT)\); hence \(Q=\Phi(S)\), proving the assertion.
Finally, uniqueness is immediate, since \(\varphi=\omega\circ\Phi\)
implies, for \(S\in\mathcal A_+\),
\(\omega(\Phi(S)T)=\varphi(ST)\) for every \(T\in\mathcal Z_+\).
\end{proof}

\begin{proposition}\label{prop:III-4-4}
Let \(\mathcal A\) be a von Neumann algebra, let \(\mathcal Z\) be its
centre, let \(\Phi\) be a faithful semifinite normal
\(\mathcal Z\)-trace on \(\mathcal A_+\), and let \(\varphi\) be a
normal trace on \(\mathcal A_+\). There is a unique normal trace
\(\omega\) on \(\widehat{\mathcal Z}_+\) such that
\(\varphi=\omega\circ\Phi\).
\end{proposition}

\begin{proof}
Let \(\omega_1\) be any faithful semifinite normal trace on
\(\widehat{\mathcal Z}_+\). Put
\(\varphi_1=\omega_1\circ\Phi\); this is a faithful semifinite normal
trace on \(\mathcal A_+\) by Proposition~2. By Chapter~I, \S~6,
Proposition~1, Corollary~2, there is a projection \(E\in\mathcal Z\),
identified with the characteristic function of a measurable subset
\(Y\) of \(Z\), having the following properties:
\begin{enumerate}[label=\textup{(\roman*)},leftmargin=*]
\item if \(S\in\mathcal A_+\) does not belong to \(\mathcal A_E\),
then \(\varphi(S)=+\infty\);
\item the trace \(\varphi'\) induced by \(\varphi\) on
\((\mathcal A_E)_+\) is semifinite.
\end{enumerate}
Let \(\varphi'_1\) be the trace induced by \(\varphi_1\) on
\((\mathcal A_E)_+\), and put \(\psi=\varphi'+\varphi'_1\), which is
faithful and semifinite. There are (Chapter~I, \S~6, Theorem~3)
elements \(Q_0,Q_1\in\widehat{\mathcal Z}_E^+\) such that
\[
  \varphi'(S)=\psi(SQ_0),
  \qquad
  \varphi'_1(S)=\psi(SQ_1)
  \qquad(S\in(\mathcal A_E)_+).
\]
Let \(Q\) be the function in \(\widehat{\mathcal Z}_+\) equal to
\(+\infty\) on \(Z-Y\), and to \(Q_0Q_1^{-1}\) on \(Y\), and let
\(\omega\) be the normal trace on \(\widehat{\mathcal Z}_+\) defined by
\[
  \omega(T)=\omega_1(TQ)
  \qquad(T\in\widehat{\mathcal Z}_+).
\]
We shall prove that \(\varphi(S)=\omega(\Phi(S))\) for every
\(S\in\mathcal A_+\). If \(S\notin\mathcal A_E\), then
\(\Phi(S)\) is not locally almost everywhere zero on \(Z-Y\), and
therefore \(\Phi(S)Q\) is infinite on a set which is not locally
negligible; hence
\(\omega(\Phi(S))=\omega_1(\Phi(S)Q)=+\infty=\varphi(S)\). It remains
only\space%
% Source PDF page 257 (printed p. 248)
\phantomsection\label{srcpage:248}%
to consider the case \(S\in\mathcal A_E\). Identify
\(\widehat{\mathcal Z}_E^+\) with \(\widehat{\mathcal Z}_+E\), that is,
with the set of functions in \(\widehat{\mathcal Z}_+\) which vanish on
\(Z-Y\). Let \((Q_i)_{i\in I}\) be the increasing directed family of
functions in \(\mathcal Z_+\) majorized by \(Q_1^{-1}\). Then
\(\Phi(S)Q_0Q_1^{-1}\) is the supremum of the
\(\Phi(S)Q_0Q_i\), and therefore
\begin{align*}
\omega(\Phi(S))
  &=\omega_1(\Phi(S)Q)
   =\omega_1(\Phi(S)Q_0Q_1^{-1})\\
  &=\sup_i\omega_1(\Phi(S)Q_0Q_i)
   =\sup_i\omega_1(\Phi(SQ_0Q_i))\\
  &=\sup_i\varphi_1(SQ_0Q_i)
   =\sup_i\psi(SQ_0Q_iQ_1)
   =\sup_i\varphi(SQ_iQ_1).
\end{align*}
Now \(Q_1(\zeta)\) vanishes only on a locally negligible subset of
\(Y\), since \(\varphi'_1\) is faithful; hence
\(Q_1(\zeta)^{-1}Q_1(\zeta)=1\) locally almost everywhere on \(Y\).
Thus the supremum of the \(Q_iQ_1\) is \(E\), and consequently
\[
  \omega(\Phi(S))=\varphi(SE)=\varphi(S).
\]

Let us prove that the trace \(\omega\), whose existence has just been
established, is unique. The elements of \(\mathcal Z_+\) of the form
\(\Phi(S)\), with \(S\in\mathcal A_+\), for which the value of
\(\omega\) is prescribed, form the positive part of an ideal
\(\mathfrak m\) of \(\mathcal Z\). Let \(F\) be the greatest projection
in the strong closure of \(\mathfrak m\). Suppose \(I-F\neq0\). Since
\(\Phi\) is semifinite, there is a nonzero \(S\in\mathcal A_+\),
majorized by \(I-F\), such that \(\Phi(S)\in\mathcal Z_+\). Then
\[
  \Phi(S)=\Phi(S(I-F))=\Phi(S)(I-F)=0,
\]
contrary to the faithfulness of \(\Phi\). Hence \(F=I\). By normality,
\(\omega\) is therefore unique on \(\mathcal Z_+\), and consequently on
\(\widehat{\mathcal Z}_+\).
\end{proof}

It is easy to see that Proposition~4 is false if \(\mathcal Z\) is merely
contained in the centre of \(\mathcal A\) (for example, if
\(\mathcal Z\) consists only of the scalars). Likewise, the faithfulness
and semifiniteness hypotheses in Propositions~3 and~4 are indispensable.

Bibliography: \ArticleRef{6}, \ArticleRef{12}, \ArticleRef{101}.

\NumberedParagraph{4}{Existence and Uniqueness Theorems for
Operator-Valued Traces}
\label{no:III-4-4}

\begin{theorem}\label{thm:III-4-1}
Let \(\mathcal A\) be a von Neumann algebra and let \(\mathcal Z\) be a
von Neumann algebra contained in the centre of \(\mathcal A\). The
algebra \(\mathcal A\) is semifinite if and only if there is a faithful
semifinite normal \(\mathcal Z\)-trace on \(\mathcal A_+\).
\end{theorem}

\begin{proof}
Let \(\omega\) be a faithful semifinite normal trace on
\(\widehat{\mathcal Z}_+\). If there is a faithful semifinite normal
\(\mathcal Z\)-trace \(\Phi\) on \(\mathcal A_+\), then
\(\omega\circ\Phi\) is a faithful semifinite normal trace on
\(\mathcal A_+\) by Proposition~2, and therefore \(\mathcal A\) is
semifinite. Conversely, if \(\mathcal A\) is semifinite, there is a
faithful semifinite normal trace \(\varphi\) on \(\mathcal A_+\), and
therefore, by Proposition~3, there is a normal \(\mathcal Z\)-trace
\(\Phi\) on \(\mathcal A_+\) such that
\(\varphi=\omega\circ\Phi\). By Proposition~2, \(\Phi\) is faithful and
semifinite.
\end{proof}

\begin{theorem}\label{thm:III-4-2}
Let \(\mathcal A\) be a von Neumann algebra, let \(\mathcal Z\) be its
centre, and let \(\Phi_1,\Phi_2\) be two faithful semifinite normal
\(\mathcal Z\)-traces on \(\mathcal A_+\). There is a unique element
\(Q\in\widehat{\mathcal Z}_+\) such that
\[
  \Phi_1(S)=\Phi_2(S)Q
  \qquad(S\in\mathcal A_+).
\]
Moreover, \(0<Q(\zeta)<+\infty\) locally almost everywhere on \(Z\).
\end{theorem}

% Source PDF page 258 (printed p. 249)
\phantomsection\label{srcpage:249}
\begin{proof}
Let \(\omega_1\) be a faithful semifinite normal trace on
\(\widehat{\mathcal Z}_+\). Then
\(\varphi=\omega_1\circ\Phi_1\) is a faithful semifinite normal trace on
\(\mathcal A_+\) by Proposition~2. By Proposition~4, there is a normal
trace \(\omega_2\) on \(\widehat{\mathcal Z}_+\) such that
\(\varphi=\omega_2\circ\Phi_2\), and \(\omega_2\) is faithful and
semifinite by Proposition~2. By Lemma~1, there is an element
\(Q\in\widehat{\mathcal Z}_+\) such that
\[
  \omega_2(T)=\omega_1(TQ)
  \qquad(T\in\widehat{\mathcal Z}_+),
\]
and \(0<Q(\zeta)<+\infty\) locally almost everywhere on \(Z\). Hence,
for every \(S\in\mathcal A_+\),
\[
  \omega_1(\Phi_1(S))
   =\omega_2(\Phi_2(S))
   =\omega_1(\Phi_2(S)Q).
\]
By Proposition~3, the mappings
\(S\mapsto\Phi_1(S)\) and \(S\mapsto\Phi_2(S)Q\) are identical.

Now suppose that \(\Phi_1(S)=\Phi_2(S)Q'\), where
\(Q'\in\widehat{\mathcal Z}_+\), and put
\(\omega'_2(T)=\omega_1(TQ')\) for every
\(T\in\widehat{\mathcal Z}_+\). For every \(S\in\mathcal A_+\),
\[
  \omega_2(\Phi_2(S))
   =\omega_1(\Phi_1(S))
   =\omega_1(\Phi_2(S)Q')
   =\omega'_2(\Phi_2(S)).
\]
By Proposition~4, \(\omega_2=\omega'_2\), and by Lemma~1,
\(Q=Q'\).
\end{proof}

\begin{theorem}\label{thm:III-4-3}
Let \(\mathcal A\) be a finite von Neumann algebra and let
\(\mathcal Z\) be its centre. There is a unique normal
\(\mathcal Z\)-trace \(\Phi\) on \(\mathcal A_+\) such that
\(\Phi(T)=T\) for \(T\in\mathcal Z_+\). This \(\mathcal Z\)-trace is
finite and faithful.
\end{theorem}

\begin{proof}
Let \(\Psi\) be a faithful semifinite normal \(\mathcal Z\)-trace on
\(\mathcal A_+\), whose existence follows from Theorem~1. Put
\(Q=\Psi(I)\in\widehat{\mathcal Z}_+\). Since \(\Psi\) is faithful,
\(Q(\zeta)>0\) locally almost everywhere. Suppose that
\(Q(\zeta)=+\infty\) on a measurable subset of \(Z\) which is not
locally negligible. Let \(\omega\) be a faithful semifinite normal trace
on \(\widehat{\mathcal Z}_+\). There would then be a nonzero
\(T\in\mathcal Z_+\) such that, for every nonzero
\(T'\in\mathcal Z_+\) majorized by \(T\),
\[
  \omega(\Psi(T'))=\omega(QT')=+\infty.
\]
Since the trace \(\omega\circ\Psi\) is semifinite, this would contradict
\hyperref[prop:I-6-10]{Chapter~I, \S~6, Proposition~10}. Hence
\(0<Q(\zeta)<+\infty\) locally almost everywhere. The mapping
\[
  S\longmapsto\Psi(S)Q^{-1}
\]
on \(\mathcal A_+\) is therefore a normal \(\mathcal Z\)-trace
\(\Phi\) such that \(\Phi(I)=I\). Thus
\(\Phi(T)=T\Phi(I)=T\) for \(T\in\mathcal Z_+\). This property implies
that \(\Phi\) is finite. It also implies that \(\Phi\) is faithful: the
set of \(S\in\mathcal A_+\) such that \(\Phi(S)=0\) is the positive
part of a two-sided ideal \(\mathfrak m\) of \(\mathcal A\), and the
greatest projection \(E\) in the strong closure of \(\mathfrak m\)
satisfies \(E\in\mathcal Z\) and \(\Phi(E)=0\). Finally, the uniqueness
of \(\Phi\) follows from Theorem~2.
\end{proof}

\begin{definition}\label{def:III-4-2}
Let \(\mathcal A\) be a finite von Neumann algebra. The
\(\mathcal Z\)-trace on \(\mathcal A\) whose existence is proved by
Theorem~3 is called the \emph{canonical \(\mathcal Z\)-trace} of
\(\mathcal A\).
\end{definition}

This mapping, often denoted by \(T\mapsto T^\natural\), is therefore a
positive normal linear mapping from \(\mathcal A\) onto its centre
\(\mathcal Z\), such that \(T^\natural=T\) for \(T\in\mathcal Z\). We
have
\[
  (S^*)^\natural=(S^\natural)^*,
  \qquad
  (ST)^\natural=(TS)^\natural
  \quad(S,T\in\mathcal A),
\]
and%
% Source PDF page 259 (printed p. 250)
\phantomsection\label{srcpage:250}%
\[
  (ST)^\natural=S^\natural T
  \qquad(S\in\mathcal A,\ T\in\mathcal Z).
\]
By \hyperref[thm:I-4-2]{Chapter~I, \S~4, Theorem~2}, the mapping
\(T\mapsto T^\natural\) is ultra-weakly and ultra-strongly continuous.

If \(\mathcal A\) is a finite factor, the canonical mapping is of the
form \(T\mapsto\varphi(T)I\), where \(\varphi\) is the unique normal
trace on \(\mathcal A\) such that \(\varphi(I)=1\). This trace
\(\varphi\) is called the \emph{canonical trace} on \(\mathcal A\).

Bibliography: \ArticleRef{6}, \ArticleRef{12}, \ArticleRef{101},
\ArticleRef{104}, \ArticleRef{105}, and T.~Iwamura, ``On continuous
geometries I,'' \emph{Japanese Journal of Mathematics}, vol.~9 (1944),
pp.~57--71.

\noindent\textit{Exercises.}
\begin{enumerate}[label=\arabic*.,leftmargin=*]
\item Let \(\mathcal A\) be a continuous finite von Neumann algebra,
let \(\mathcal Z\) be its centre, let \(S\mapsto S^\natural\) be the
canonical \(\mathcal Z\)-trace of \(\mathcal A\), and let
\(T\in\mathcal Z\) satisfy \(0\leq T\leq I\). Show that there is a
projection \(E\in\mathcal A\) such that \(E^\natural=T\).

[First show that the result is true if \(T\) is a step function taking
only finitely many values of the form \(p2^{-n}\), where \(p,n\) are
nonnegative integers. For this use
\hyperref[cor:III-1-thm1-3]{\S~1, Theorem~1, Corollary~3}. Then establish the
general result by arguing as in \S~2, Proposition~14.]
\ArticleRef{89}; and T.~Iwamura, ``On continuous geometries I,''
\emph{Japanese Journal of Mathematics}, vol.~9 (1944), pp.~57--71.

\item\label{ex:III-4-2} Let \(\mathfrak A\) be a Hilbert algebra, let \(\mathfrak H\) be
its completed Hilbert space, and let \(\mathfrak Z\) be the set of
central elements of \(\mathfrak H\). Suppose that
\(\mathcal A=\mathcal U(\mathfrak A)\) is finite. Let
\(T\mapsto T^\natural\) be the canonical \(\mathcal Z\)-trace of
\(\mathcal A\), where \(\mathcal Z\) is the centre of \(\mathcal A\),
and let \(a\) be a bounded element of \(\mathfrak H\). Adopt the notation
of Chapter~I, \S~5, Exercise~4.
\begin{enumerate}[label=\textit{\alph*.},leftmargin=2em]
\item Show that, for every \(b\in\mathcal K'_a\),
\(U_b^\natural=U_a^\natural\). First prove this for
\(b\in\mathcal K_a\). Then, if \(b_n\to b\), with
\(b_n\in\mathcal K_a\) and \(b\in\mathcal K'_a\), the sequence
\(\lVert U_{b_n}\rVert\) remains bounded and \(U_{b_n}\) converges
weakly to \(U_b\). Use the ultra-weak continuity of the canonical
\(\mathcal Z\)-trace.
\item Show that \(U_a^\natural=U_{P_{\mathfrak Z}a}\).
\item Show that \(\mathcal A\) is a factor if and only if
\(\mathfrak Z\) has dimension~1. If \(\dim\mathfrak Z>1\), then
\(\mathcal A\) is not a factor by Chapter~I, \S~5, Exercise~4e. If
\(\mathcal A\) is not a factor, there are bounded elements
\(a,b\in\mathfrak H\) such that \(U_a^\natural,U_b^\natural\) are
linearly independent; then use part~b. \ArticleRef{29},
\ArticleRef{106}
\end{enumerate}

\item Let \(\mathcal A\) be a semifinite von Neumann algebra, and let
its centre \(\mathcal Z\) be identified with
\(L_{\mathbb C}^{\infty}(Z,\nu)\). Let \(\Phi\) be a faithful semifinite
normal \(\mathcal Z\)-trace on \(\mathcal A_+\). Show that the set of
\(T\in\mathcal A_+\) for which \(\Phi(T)\) is finite locally almost
everywhere is independent of \(\Phi\) and of the identification of
\(\mathcal Z\) with \(L_{\mathbb C}^{\infty}(Z,\nu)\), and is the
positive part of a two-sided ideal \(\mathfrak m\) of \(\mathcal A\).
Show that \(\mathfrak m\) is strongly dense everywhere in
\(\mathcal A\), and is the union of the definition ideals of \(\Psi\),
as \(\Psi\) ranges over the faithful normal \(\mathcal Z\)-traces on
\(\mathcal A_+\). Show that \(\mathfrak m=\mathcal A\) if and only if
\(\mathcal A\) is finite. \ArticleRef{12}

\item Let \(Z\) be a Borel space, let \(\nu\) be a positive measure on
\(Z\), let \(\zeta\mapsto\mathfrak H(\zeta)\) be a
\(\nu\)-measurable field of nonzero complex Hilbert spaces on \(Z\),
and let
\[
  \mathfrak H=\int^{\oplus}\mathfrak H(\zeta)\,d\nu(\zeta),
  \qquad
  \mathcal A=\int^{\oplus}\mathcal A(\zeta)\,d\nu(\zeta)
\]
be a decomposable von Neumann algebra on \(\mathfrak H\). Let
\(\zeta\mapsto\varphi_\zeta\) be a measurable field of traces on the
\(\mathcal A(\zeta)_+\), and let \(\mathcal Z\) be the algebra of
diagonalizable operators, identified with
\(L_{\mathbb C}^{\infty}(Z,\nu)\).
\begin{enumerate}[label=\textit{\alph*.},leftmargin=2em]
\item For every
\[
  T=\int^{\oplus}T(\zeta)\,d\nu(\zeta)\in\mathcal A_+,
\]
let \(\Phi(T)\) be the function
\(\zeta\mapsto\varphi_\zeta(T(\zeta))\), which belongs to
\(\widehat{\mathcal Z}_+\). Show that \(\Phi\) is a
\(\mathcal Z\)-trace on \(\mathcal A_+\), and is normal if the
\(\varphi_\zeta\) are normal. Argue as in Chapter~II, \S~5,
Proposition~1.

% Source PDF page 260 (printed p. 251)
\phantomsection\label{srcpage:251}
\item Conversely, every faithful semifinite normal
\(\mathcal Z\)-trace \(\Phi\) on \(\mathcal A_+\) is of this form.
[Let \(\omega\) be the trace
\(S\mapsto\int S(\zeta)\,d\nu(\zeta)\) on
\(\widehat{\mathcal Z}_+\). Form the trace \(\omega\circ\Phi\). Apply
Chapter~II, \S~5, Theorem~2 to obtain the field
\(\zeta\mapsto\varphi_\zeta\). Then use part~a and the uniqueness result
of Proposition~3.]
\end{enumerate}
\end{enumerate}

\section{Approximation Theorem}\label{sec:III-5}

\NumberedParagraph{1}{The Approximation Theorem}
\label{no:III-5-1}
In the complex Hilbert space \(\mathfrak H\), let \(T\) be a bounded
Hermitian operator and let \(E\) be a nonzero projection commuting with
\(T\). In this numbered paragraph put
\begin{align*}
  M_E(T)&=\sup_{\lVert x\rVert=1,\ Ex=x}(Tx\mid x),
  &m_E(T)&=\inf_{\lVert x\rVert=1,\ Ex=x}(Tx\mid x),\\
  \omega_E(T)&=M_E(T)-m_E(T).
\end{align*}
When \(E=I\), we write simply \(M(T),m(T),\omega(T)\). We agree,
moreover, that \(\omega_0(T)=0\). If \(\mathcal F\) is a family of
projections commuting with \(T\), put
\[
  \omega_{\mathcal F}(T)=\sup_{E\in\mathcal F}\omega_E(T).
\]

\begin{lemma}\label{lem:III-5-1}
Let \(\mathcal A\) be a von Neumann algebra, let \(\mathcal Z\) be its
centre, and let \(T\) be a Hermitian operator in \(\mathcal A\). There
are projections \(G,G'\in\mathcal Z\), orthogonal and with sum \(I\),
and a unitary operator \(U\in\mathcal A\), such that
\begin{align}
  \omega_G\!\left(\frac12(T+UTU^{-1})\right)
    &\leq\frac34\omega(T),\tag{1}\label{eq:III-5-1}\\
  \omega_{G'}\!\left(\frac12(T+UTU^{-1})\right)
    &\leq\frac34\omega(T).\tag{2}\label{eq:III-5-2}
\end{align}
\end{lemma}

\begin{proof}
Put \(n(T)=\frac12(M(T)+m(T))\). There are orthogonal spectral
projections \(E,F\) of \(T\), with sum \(I\), such that
\[
  M_E(T)\leq n(T),
  \qquad
  m_F(T)\geq n(T).
\]
There are projections \(G,G'\in\mathcal Z\) such that
\[
  GG'=0,\qquad G+G'=I,\qquad EG\prec FG,\qquad FG'\prec EG'
\]
by \hyperref[thm:III-1-1]{\S~1, Theorem~1}. There is a partial isometry
\(V\) (respectively, \(W\)) in \(\mathcal A\), with initial projection
\(EG\) (respectively, \(FG'\)) and final projection
\(G_1\leq FG\) (respectively, \(G'_1\leq EG'\)). Consequently there is
a unitary operator \(U\in\mathcal A\) which maps
\[
  EG(\mathfrak H)\ \text{onto}\ G_1(\mathfrak H),\quad
  G_1(\mathfrak H)\ \text{onto}\ EG(\mathfrak H),
\]
maps \(FG'(\mathfrak H)\) onto \(G'_1(\mathfrak H)\) and
\(G'_1(\mathfrak H)\) onto \(FG'(\mathfrak H)\), and is the identity
on every vector orthogonal to these subspaces. We show, for example,
that \(G,G',U\) satisfy inequality~\eqref{eq:III-5-1}. We have
\[
  TG\geq m(T)EG+n(T)FG
   =m(T)EG+n(T)G_1+n(T)(FG-G_1),
\]
and therefore
\[
  (UTU^{-1})G
   \geq m(T)G_1+n(T)EG+n(T)(FG-G_1).
\]%
% Source PDF page 261 (printed p. 252)
\phantomsection\label{srcpage:252}%
Adding these inequalities term by term gives
\begin{align*}
  \frac12(T+UTU^{-1})G
  &\geq \frac12\bigl(m(T)+n(T)\bigr)(EG+G_1)
        +n(T)(FG-G_1)\\
  &\geq \frac12\bigl(m(T)+n(T)\bigr)
        (EG+G_1+FG-G_1)\\
  &=\left(M(T)-\frac34\omega(T)\right)G.
\end{align*}
Since it is clear that
\[
  \frac12(T+UTU^{-1})G\leq M(T)G,
\]
inequality~\eqref{eq:III-5-1} is proved.
\end{proof}

\begin{lemma}\label{lem:III-5-2}
Let \(T\) be a Hermitian operator in \(\mathcal A\), and let
\(\mathcal F\) be a finite family of pairwise orthogonal projections in
\(\mathcal Z\) whose sum is \(I\). There are a finite family
\(\mathcal F'\) of pairwise orthogonal projections in \(\mathcal Z\)
whose sum is \(I\), and a unitary operator \(U\in\mathcal A\), such that
\[
  \omega_{\mathcal F'}\!\left(\frac12(T+UTU^{-1})\right)
  \leq\frac34\omega_{\mathcal F}(T).
\]
\end{lemma}

\begin{proof}
If the lemma is proved for a finite family of von Neumann algebras, it
follows at once for their product von Neumann algebra. We are therefore
reduced to the case where \(\mathcal F\) consists of the single
projection \(I\). Lemma~1 then constructs a two-element family
\(\mathcal F'\) and a unitary \(U\) having the required properties.
\end{proof}

For \(T\in\mathcal A\), denote by \(\mathcal K_T\) the convex subset of
\(\mathcal A\) generated by the elements \(UTU^{-1}\), where \(U\)
ranges over the group \(\mathcal G\) of unitary operators in
\(\mathcal A\). Denote by \(\mathcal K'_T\) the closure of
\(\mathcal K_T\) in \(\mathcal A\) for the norm topology. On the other
hand, let \(\mathcal E\) be the set of functions \(U\mapsto f(U)\) on
\(\mathcal G\), with nonnegative real values, vanishing except at
finitely many points, and such that
\[
  \sum_{U\in\mathcal G}f(U)=1.
\]
For \(f\in\mathcal E\) and \(T\in\mathcal A\), put
\[
  f\mathbin{\cdot}T
  =\sum_{U\in\mathcal G}f(U)UTU^{-1}.
\]
As \(f\) ranges over \(\mathcal E\), \(f\mathbin{\cdot}T\) ranges over
\(\mathcal K_T\). It is immediate that, for \(g\in\mathcal E\),
\[
  g\mathbin{\cdot}(f\mathbin{\cdot}T)=(gf)\mathbin{\cdot}T,
\]
where \(gf\) denotes the function on \(\mathcal G\) given by
\[
  W\longmapsto
  \sum_{\substack{U,V\in\mathcal G\\VU=W}}f(U)g(V),
\]
that is, the convolution product of \(g\) and \(f\). This function is an
element of \(\mathcal E\), where \(\mathcal G\) is regarded as a
discrete group.

\begin{lemma}\label{lem:III-5-3}
Let \(T\) be a Hermitian operator in \(\mathcal A\), and let
\(\varepsilon>0\). There are \(f\in\mathcal E\) and
\(S\in\mathcal Z\) such that
\[
  \lVert f\mathbin{\cdot}T-S\rVert\leq\varepsilon.
\]
\end{lemma}

\begin{proof}
For every integer \(p>0\), there are a family
\(\mathcal F=(E_1,E_2,\ldots,E_n)\) of pairwise orthogonal projections
in \(\mathcal Z\), whose sum is \(I\), and an \(f\in\mathcal E\), such
that
\[
  \omega_{\mathcal F}(f\mathbin{\cdot}T)
  \leq\left(\frac34\right)^p\omega(T).
\]
This follows at once from Lemma~2 by induction on \(p\). Choosing
suitable real numbers \(\alpha_i\), we therefore have%
% Source PDF page 262 (printed p. 253)
\phantomsection\label{srcpage:253}%
\[
  \left\lVert f\mathbin{\cdot}T-
    \sum_{i=1}^{n}\alpha_iE_i\right\rVert
  \leq\left(\frac34\right)^p\omega(T).
\]
But \(S=\sum_{i=1}^{n}\alpha_iE_i\) belongs to \(\mathcal Z\), and
\((3/4)^p\omega(T)\) is arbitrarily small when \(p\) is sufficiently
large.
\end{proof}

\begin{lemma}\label{lem:III-5-4}
Let \(T_1,T_2,\ldots,T_n\) be elements of \(\mathcal A\), and let
\(\varepsilon>0\). There are \(f\in\mathcal E\) and elements
\(S_1,S_2,\ldots,S_n\in\mathcal Z\) such that
\[
  \lVert f\mathbin{\cdot}T_h-S_h\rVert\leq\varepsilon
  \qquad(h=1,2,\ldots,n).
\]
\end{lemma}

\begin{proof}
Since every element of \(\mathcal A\) is a linear combination of two
Hermitian elements of \(\mathcal A\), we may suppose that the \(T_h\)
are Hermitian. Arguing by induction on \(n\), suppose that we have found
\(g\in\mathcal E\) and \(S_1,S_2,\ldots,S_{n-1}\in\mathcal Z\) such
that
\[
  \lVert g\mathbin{\cdot}T_h-S_h\rVert\leq\varepsilon
  \qquad(h=1,2,\ldots,n-1).
\]
By Lemma~3 there are \(g'\in\mathcal E\) and \(S_n\in\mathcal Z\)
such that
\[
  \lVert g'\mathbin{\cdot}(g\mathbin{\cdot}T_n)-S_n\rVert
  \leq\varepsilon.
\]
For \(h=1,2,\ldots,n-1\),
\begin{align*}
  \lVert g'\mathbin{\cdot}(g\mathbin{\cdot}T_h)-S_h\rVert
  &=\lVert g'\mathbin{\cdot}(g\mathbin{\cdot}T_h-S_h)\rVert\\
  &\leq\lVert g\mathbin{\cdot}T_h-S_h\rVert
  \leq\varepsilon.
\end{align*}
It is therefore enough to put \(f=g'g\).
\end{proof}

\begin{lemma}\label{lem:III-5-5}
Let \(T_1,T_2,\ldots,T_n\) be elements of \(\mathcal A\). There are
elements \(S_1,S_2,\ldots,S_n\in\mathcal Z\) and a sequence
\((f_1,f_2,\ldots)\) of elements of \(\mathcal E\) such that, for
\(h=1,2,\ldots,n\),
\[
  \lVert f_i\mathbin{\cdot}T_h-S_h\rVert\longrightarrow0
  \qquad\text{as }i\longrightarrow+\infty.
\]
\end{lemma}

\begin{proof}
Using Lemma~4 and arguing by induction on \(i\), for every integer
\(i>0\) construct \(g_i\in\mathcal E\) and
\(S_1^i,S_2^i,\ldots,S_n^i\in\mathcal Z\) such that
\[
  \lVert(g_ig_{i-1}\cdots g_1)\mathbin{\cdot}T_h-S_h^i\rVert
  \leq2^{-i}
  \qquad(h=1,2,\ldots,n).
\]
Put \(f_i=g_ig_{i-1}\cdots g_1\). We have
\begin{align*}
  \lVert f_{i+1}\mathbin{\cdot}T_h-S_h^i\rVert
  &=\lVert g_{i+1}\mathbin{\cdot}
       (f_i\mathbin{\cdot}T_h)-S_h^i\rVert\\
  &=\lVert g_{i+1}\mathbin{\cdot}
       (f_i\mathbin{\cdot}T_h-S_h^i)\rVert\\
  &\leq\lVert f_i\mathbin{\cdot}T_h-S_h^i\rVert
  \leq2^{-i},
\end{align*}
and therefore
\[
  \lVert f_{i+1}\mathbin{\cdot}T_h-f_i\mathbin{\cdot}T_h\rVert
  \leq2^{-i+1}.
\]
Thus, for each \(h=1,2,\ldots,n\), the sequence
\((f_1\mathbin{\cdot}T_h,f_2\mathbin{\cdot}T_h,\ldots)\) is a Cauchy
sequence. Its limit \(S_h\) is also the limit of the \(S_h^i\), and
therefore belongs to \(\mathcal Z\).
\end{proof}

Notice that \(S_h\in\mathcal K'_{T_h}\). As a special case of Lemma~5,
we therefore have the following result.

\begin{theorem}\label{thm:III-5-1}
For every \(T\in\mathcal A\),
\[
  \mathcal K''_T=\mathcal K'_T\cap\mathcal Z
\]
is nonempty.
\end{theorem}

% Source PDF page 263 (printed p. 254)
\phantomsection\label{srcpage:254}
\phantomsection
\begin{corollary*}\phantomsection\label{cor:III-5-1}
Let \(\mathcal A\) be a finite von Neumann algebra, let
\(\mathcal Z\) be its centre, and let \(\Phi\) be the canonical
\(\mathcal Z\)-trace of \(\mathcal A\).
\begin{enumerate}[label=\textup{(\roman*)},leftmargin=*]
\item For every \(T\in\mathcal A\), \(\mathcal K''_T\) consists of the
single point \(\Phi(T)\).
\item Let \(\Psi:\mathcal A\to\mathcal Z\) be a norm-continuous linear
mapping such that \(\Psi(T_1T_2)=\Psi(T_2T_1)\) for
\(T_1,T_2\in\mathcal A\), and such that \(\Psi(T)=T\) for
\(T\in\mathcal Z\). Then \(\Psi=\Phi\).
\end{enumerate}
\end{corollary*}

\begin{proof}
We shall show that \(\mathcal K''_T\) consists of the single point
\(\Psi(T)\). Since \(\Phi\) has the properties assumed for \(\Psi\),
this will prove both \textup{(i)} and \textup{(ii)}. For every
\(T\in\mathcal A\) and every unitary operator \(U\in\mathcal A\),
\[
  \Psi(UTU^{-1})=\Psi(U^{-1}UT)=\Psi(T).
\]
Thus \(\Psi\) is constant on \(\mathcal K_T\), and consequently on
\(\mathcal K'_T\). Hence, for \(S\in\mathcal K''_T\),
\(S=\Psi(S)=\Psi(T)\).
\end{proof}

We shall prove later some converses to this corollary
(\hyperref[sec:III-8]{\S~8, Theorem~1, Corollaries~3 and~4}).

Bibliography: \ArticleRef{6}.

\NumberedParagraph{2}{Application: Two-Sided Ideals of
\(\mathcal A\) and Ideals of \(\mathcal Z\)}
\label{no:III-5-2}
The algebra \(\mathcal A\) is no longer assumed finite.

\begin{lemma}\label{lem:III-5-6}
Let \(T_1,T_2\in\mathcal A\). Then \(\mathcal K''_{T_1+T_2}\) is
contained in the norm closure of
\(\mathcal K''_{T_1}+\mathcal K''_{T_2}\).
\end{lemma}

\begin{proof}
Let \(S\in\mathcal K''_{T_1+T_2}\), and let \(\varepsilon>0\). There is
an \(f\in\mathcal E\) such that
\[
  \lVert f\mathbin{\cdot}(T_1+T_2)-S\rVert\leq\varepsilon.
\]
Then, by Lemma~5, there are \(g\in\mathcal E\) and elements
\[
  S_1\in\mathcal K''_{f\mathbin{\cdot}T_1}
       \subset\mathcal K''_{T_1},
  \qquad
  S_2\in\mathcal K''_{f\mathbin{\cdot}T_2}
       \subset\mathcal K''_{T_2}
\]
such that
\[
  \lVert g\mathbin{\cdot}(f\mathbin{\cdot}T_1)-S_1\rVert
  \leq\varepsilon,
  \qquad
  \lVert g\mathbin{\cdot}(f\mathbin{\cdot}T_2)-S_2\rVert
  \leq\varepsilon.
\]
Since
\(\lVert g\mathbin{\cdot}(f\mathbin{\cdot}(T_1+T_2))-S\rVert
\leq\varepsilon\), it follows that
\[
  \lVert S-(S_1+S_2)\rVert\leq3\varepsilon,
\]
which proves the lemma.
\end{proof}

\begin{lemma}\label{lem:III-5-7}
Let \(T\in\mathcal A\) and \(S\in\mathcal Z\). Then
\(\mathcal K''_{ST}\) is contained in the norm closure of
\(S\mathcal K''_T\).
\end{lemma}

\begin{proof}
Let \(R\in\mathcal K''_{ST}\), and let \(\varepsilon>0\). There is an
\(f\in\mathcal E\) such that
\(\lVert f\mathbin{\cdot}(ST)-R\rVert\leq\varepsilon\). By Theorem~1,
there are \(g\in\mathcal E\) and
\(R_1\in\mathcal K''_{f\mathbin{\cdot}T}\subset\mathcal K''_T\) such
that
\(\lVert g\mathbin{\cdot}(f\mathbin{\cdot}T)-R_1\rVert
\leq\varepsilon\). The preceding inequalities imply
\begin{align*}
  \lVert S(g\mathbin{\cdot}(f\mathbin{\cdot}T))-R\rVert
   &=\lVert g\mathbin{\cdot}(f\mathbin{\cdot}(ST)-R)\rVert\\
   &\leq\lVert f\mathbin{\cdot}(ST)-R\rVert\leq\varepsilon,\\
  \lVert S(g\mathbin{\cdot}(f\mathbin{\cdot}T))-SR_1\rVert
   &\leq\varepsilon\lVert S\rVert.
\end{align*}
Thus \(\lVert R-SR_1\rVert\leq\varepsilon(1+\lVert S\rVert)\), which
proves the lemma.
\end{proof}

\begin{lemma}\label{lem:III-5-8}
Let \(\mathfrak m\) be a norm-closed two-sided ideal of \(\mathcal A\).
Then \(\mathfrak m\cap\mathcal Z\) is the union of the sets
\(\mathcal K''_T\) as \(T\) ranges over \(\mathfrak m\).
\end{lemma}

% Source PDF page 264 (printed p. 255)
\phantomsection\label{srcpage:255}
\begin{proof}
Since \(\mathcal K''_T=\{T\}\) for \(T\in\mathcal Z\), it is clear
that \(\mathfrak m\cap\mathcal Z\) is contained in the union of the
\(\mathcal K''_T\), \(T\in\mathfrak m\). Conversely, if
\(T\in\mathfrak m\), then \(\mathcal K_T\subset\mathfrak m\), hence
\(\mathcal K'_T\subset\mathfrak m\), and therefore
\(\mathcal K''_T\subset\mathfrak m\cap\mathcal Z\).
\end{proof}

\begin{lemma}\label{lem:III-5-9}
Let \(\mathfrak m\) be a two-sided ideal of \(\mathcal A\), and let
\(\overline{\mathfrak m}\) be its norm closure. Then
\(\overline{\mathfrak m}\cap\mathcal Z\) is the norm closure of
\(\mathfrak m\cap\mathcal Z\).
\end{lemma}

\begin{proof}
Let \(\mathfrak n\) be the norm closure of
\(\mathfrak m\cap\mathcal Z\). Clearly
\(\mathfrak n\subset\overline{\mathfrak m}\cap\mathcal Z\). Suppose
that there is an \(S\in\overline{\mathfrak m}\cap\mathcal Z\) which
does not belong to \(\mathfrak n\). We shall obtain a contradiction,
which will prove the lemma. Let \(Z\) be the spectrum of
\(\mathcal Z\), let \(L_\infty(Z)\) be the algebra of continuous
complex-valued functions on \(Z\), and let
\(T\mapsto f_T\) be the canonical isomorphism of \(\mathcal Z\) onto
\(L_\infty(Z)\). Then \(\mathfrak n\) is the set of
\(T\in\mathcal Z\) such that \(f_T\) vanishes on a certain closed subset
\(Y\) of \(Z\). There is a \(\zeta\in Y\) such that
\(f_S(\zeta)\neq0\). Hence there is a spectral projection
\(E\in\mathcal Z\) of \(S\) such that \(f_E(\zeta)=1\) and the
restriction of \(S\) to \(E(\mathfrak H)\) is invertible. Since \(S\)
is a norm limit of elements of \(\mathfrak m\), there is a
\(T\in\mathfrak m\) whose restriction to \(E(\mathfrak H)\) is also
invertible. For a suitably chosen \(T_1\in\mathcal A\), we then have
\(TT_1=E\), and hence \(E\in\mathfrak m\). This contradicts
\(f_E(\zeta)=1\).
\end{proof}

\begin{lemma}\label{lem:III-5-10}
Let \(\mathfrak n\) be a norm-closed ideal of \(\mathcal Z\), and let
\(\mathfrak m\) be the set of all \(T\in\mathcal A\) such that
\[
  \mathcal K''_{T_1TT_2}\subset\mathfrak n
  \qquad\text{for all }T_1,T_2\in\mathcal A.
\]
Then \(\mathfrak m\) is the greatest two-sided ideal of
\(\mathcal A\) such that
\(\mathfrak m\cap\mathcal Z\subset\mathfrak n\); it is norm closed,
and \(\mathfrak m\cap\mathcal Z=\mathfrak n\).
\end{lemma}

\begin{proof}
If \(T,T'\in\mathfrak m\), then \(T+T'\in\mathfrak m\) by Lemma~6.
It is therefore clear that \(\mathfrak m\) is a two-sided ideal of
\(\mathcal A\). If \(T\in\mathfrak m\cap\mathcal Z\), then
\(\{T\}=\mathcal K''_T\subset\mathfrak n\), and hence
\(T\in\mathfrak n\); thus
\(\mathfrak m\cap\mathcal Z\subset\mathfrak n\). Conversely, if
\(T\in\mathfrak n\), then for all \(T_1,T_2\in\mathcal A\),
\[
  T\mathcal K''_{T_1T_2}\subset\mathfrak n.
\]
Consequently, by Lemma~7,
\[
  \mathcal K''_{T_1TT_2}
   =\mathcal K''_{TT_1T_2}\subset\mathfrak n,
\]
so \(T\in\mathfrak m\cap\mathcal Z\). Thus
\(\mathfrak n=\mathfrak m\cap\mathcal Z\).

Finally, let \(\mathfrak m'\) be a two-sided ideal of \(\mathcal A\)
such that \(\mathfrak m'\cap\mathcal Z\subset\mathfrak n\), and let
\(\mathfrak m''\) be its norm closure. By Lemma~9,
\(\mathfrak m''\cap\mathcal Z\subset\mathfrak n\). Hence, for every
\(T\in\mathfrak m''\), Lemma~8 gives
\(\mathcal K''_T\subset\mathfrak n\). Applying the same fact to
\(T_1TT_2\in\mathfrak m''\), for arbitrary
\(T_1,T_2\in\mathcal A\), shows that
\(\mathfrak m''\subset\mathfrak m\). This proves at once that
\(\mathfrak m\) is the greatest such ideal and that it is norm closed.
\end{proof}

\begin{proposition}\label{prop:III-5-1}
Let \(\mathfrak n\) be a norm-closed ideal of \(\mathcal Z\). There is
a greatest two-sided ideal \(\mathfrak m\) of \(\mathcal A\) such that
\(\mathfrak m\cap\mathcal Z\subset\mathfrak n\), and in fact
\(\mathfrak m\cap\mathcal Z=\mathfrak n\). There is also a least
two-sided ideal \(\mathfrak m'\) of \(\mathcal A\) such that
\(\mathfrak m'\cap\mathcal Z=\mathfrak n\).
\end{proposition}

\begin{proof}
The existence of \(\mathfrak m\), and the equality
\(\mathfrak m\cap\mathcal Z=\mathfrak n\), follow from Lemma~10. On the
other hand, let \(\mathfrak m'\) be the two-sided ideal of
\(\mathcal A\) generated by \(\mathfrak n\). We have
\(\mathfrak m'\cap\mathcal Z\supset\mathfrak n\) and
\(\mathfrak m'\subset\mathfrak m\), hence
\(\mathfrak m'\cap\mathcal Z=\mathfrak n\). Thus \(\mathfrak m'\) is
the least two-sided ideal of \(\mathcal A\) having this intersection
with \(\mathcal Z\).
\end{proof}

% Source PDF page 265 (printed p. 256)
\phantomsection\label{srcpage:256}
\setcounter{corollary}{0}
\begin{corollary}\label{cor:III-5-1-1}
The mapping
\(\mathfrak M\mapsto\mathfrak M\cap\mathcal Z\) is a bijection from
the set of maximal two-sided ideals of \(\mathcal A\) onto the set of
maximal ideals of \(\mathcal Z\).
\end{corollary}

\begin{proof}
Let \(\mathfrak m\) be a maximal two-sided ideal of \(\mathcal A\),
which is therefore norm closed. Let \(\mathfrak n'\) be a norm-closed
ideal of \(\mathcal Z\) which contains
\(\mathfrak m\cap\mathcal Z\) and is distinct from it. The greatest
two-sided ideal \(\mathfrak m'\) of \(\mathcal A\) such that
\(\mathfrak m'\cap\mathcal Z=\mathfrak n'\) contains
\(\mathfrak m\) and is distinct from it, hence equals \(\mathcal A\).
Thus \(\mathfrak n'=\mathcal Z\), proving that
\(\mathfrak m\cap\mathcal Z\) is a maximal ideal of \(\mathcal Z\).
Conversely, if \(\mathfrak n\) is a maximal ideal of \(\mathcal Z\),
it is clear that the greatest two-sided ideal \(\mathfrak m\) of
\(\mathcal A\) such that
\(\mathfrak m\cap\mathcal Z=\mathfrak n\) is maximal. If
\(\mathfrak m_1\) is another two-sided ideal of \(\mathcal A\) such
that \(\mathfrak m_1\cap\mathcal Z=\mathfrak n\), then
\(\mathfrak m_1\subset\mathfrak m\), and therefore
\(\mathfrak m_1=\mathfrak m\) when \(\mathfrak m_1\) is maximal.
\end{proof}

\begin{corollary}\label{cor:III-5-1-2}
If \(\mathfrak m\) is the intersection of the maximal two-sided ideals
of \(\mathcal A\), then \(\mathfrak m\cap\mathcal Z=0\).
\end{corollary}

\begin{proof}
This follows from Corollary~1 and from the fact that the intersection of
the maximal ideals of \(\mathcal Z\) is zero.
\end{proof}

\begin{corollary}\label{cor:III-5-1-3}
If \(\mathcal A\) is a factor, the set of two-sided ideals of
\(\mathcal A\) distinct from \(\mathcal A\) has a greatest element.
\end{corollary}

\begin{proof}
To say that a two-sided ideal \(\mathfrak m\) of \(\mathcal A\) is
distinct from \(\mathcal A\) is equivalent to saying that
\(\mathfrak m\cap\mathcal Z=0\).
\end{proof}

When \(\mathcal A\) is finite, Lemmas~8 and~10 may be stated, by the
corollary to Theorem~1, as follows.

\begin{proposition}\label{prop:III-5-2}
Let \(\mathcal A\) be a finite von Neumann algebra.
\begin{enumerate}[label=\textup{(\roman*)},leftmargin=*]
\item If \(\mathfrak m\) is a norm-closed two-sided ideal of
\(\mathcal A\), then \(\mathfrak m\cap\mathcal Z\) is the image of
\(\mathfrak m\) under the canonical \(\mathcal Z\)-trace of
\(\mathcal A\).
\item Let \(\mathfrak n\) be a norm-closed ideal of \(\mathcal Z\).
The greatest two-sided ideal \(\mathfrak m\) of \(\mathcal A\) such
that \(\mathfrak m\cap\mathcal Z=\mathfrak n\) is the set of all
\(T\in\mathcal A\) such that
\[
  (TT_1)^\natural\in\mathfrak n
  \qquad\text{for every }T_1\in\mathcal A.
\]
\end{enumerate}
\end{proposition}

\setcounter{corollary}{0}
\begin{corollary}\label{cor:III-5-2-1}
Let \(\mathfrak m\) be a two-sided ideal of a finite von Neumann
algebra \(\mathcal A\). If \(\mathfrak m\cap\mathcal Z=0\), then
\(\mathfrak m=0\).
\end{corollary}

\begin{proof}
Let \(\mathfrak m'\) be the greatest two-sided ideal of \(\mathcal A\)
such that \(\mathfrak m'\cap\mathcal Z=0\). For every
\(T\in\mathfrak m'\), we have \(TT^*\in\mathfrak m'\), and therefore
\((TT^*)^\natural\in\mathfrak m'\cap\mathcal Z=0\); hence \(T=0\).
\end{proof}

\begin{corollary}\label{cor:III-5-2-2}
The intersection of the maximal two-sided ideals of a finite von Neumann
algebra \(\mathcal A\) is zero.
\end{corollary}

\begin{proof}
This follows from Corollary~1 and from Corollary~2 to Proposition~1.
\end{proof}

% Source PDF page 266 (printed p. 257)
\phantomsection\label{srcpage:257}
\begin{corollary}\label{cor:III-5-2-3}
If \(\mathcal A\) is a finite factor, its only two-sided ideals are
\(0\) and \(\mathcal A\).
\end{corollary}

\begin{proof}
This follows from Corollary~1.
\end{proof}

Bibliography: \ArticleRef{29}, \ArticleRef{60}, \ArticleRef{63},
\ArticleRef{124}.

\NumberedParagraph{3}{Characters of Finite von Neumann Algebras}
\label{no:III-5-3}
In this numbered paragraph, \(\mathcal A\) denotes a finite von Neumann
algebra, \(\mathcal Z\) its centre, and \(T\mapsto T^\natural\) its
canonical \(\mathcal Z\)-trace. Let \(\mathcal L\) be the set of
norm-continuous central linear forms on \(\mathcal A\), that is, forms
such that
\[
  \varphi(T_1T_2)=\varphi(T_2T_1)
  \qquad(T_1,T_2\in\mathcal A),
\]
and let \(\mathcal L'\) be the set of norm-continuous linear forms on
\(\mathcal Z\). The following proposition generalizes, for a finite von
Neumann algebra and its canonical \(\mathcal Z\)-trace, Propositions~2
and~4 of \S~4.

\begin{proposition}\label{prop:III-5-3}
Let \(\varphi\in\mathcal L\), and let \(\varphi'\in\mathcal L'\) be
the restriction of \(\varphi\) to \(\mathcal Z\). For every
\(T\in\mathcal A\),
\[
  \varphi(T)=\varphi'(T^\natural).
\]
The mapping \(\varphi\mapsto\varphi'\) is a bijection from
\(\mathcal L\) onto \(\mathcal L'\).
\end{proposition}

\begin{proof}
Since \(\varphi\) is central, it is constant on \(\mathcal K_T\).
Since it is norm continuous, it is constant on \(\mathcal K'_T\).
Thus, by the corollary to Theorem~1,
\[
  \varphi(T)=\varphi(T^\natural)=\varphi'(T^\natural).
\]
This proves that \(\varphi\mapsto\varphi'\) is injective. To prove
surjectivity, let \(\varphi'\in\mathcal L'\) and define a linear form
\(\varphi\) on \(\mathcal A\) by
\(\varphi(T)=\varphi'(T^\natural)\). The properties of the canonical
\(\mathcal Z\)-trace show that \(\varphi\in\mathcal L\), and the
restriction of \(\varphi\) to \(\mathcal Z\) is \(\varphi'\).
\end{proof}

\begin{definition}\label{def:III-5-1}
A \emph{character} of the finite von Neumann algebra \(\mathcal A\) is
a norm-continuous central linear form \(\chi\) on \(\mathcal A\) whose
restriction to the centre \(\mathcal Z\) of \(\mathcal A\) is a
character of \(\mathcal Z\).
\end{definition}

Recall that the characters of \(\mathcal Z\) are the homomorphisms from
\(\mathcal Z\) onto the complex field. They are norm continuous. Their
set, which is in bijective correspondence with the set of maximal ideals
of \(\mathcal Z\), carries a natural topology making it a compact space,
the spectrum \(Z\) of \(\mathcal Z\); and \(\mathcal Z\) is identified
with the algebra of continuous complex-valued functions on \(Z\). By
Proposition~3, the characters of \(\mathcal A\) are identified with the
points of \(Z\).

The characters of \(\mathcal A\) are finite traces on \(\mathcal A\)
(in general nonnormal; cf. Chapter~I, \S~8, Exercise~5). An element
\(\chi\in\mathcal L\) is a character if and only if, for all
\(T\in\mathcal A\) and \(S\in\mathcal Z\),%
% Source PDF page 267 (printed p. 258)
\phantomsection\label{srcpage:258}%
\[
  \chi(ST)=\chi(S)\chi(T).
\]
The condition is plainly sufficient. Conversely, if \(\chi\) is a
character, then
\[
  \chi(ST)=\chi((ST)^\natural)
  =\chi(ST^\natural)=\chi(S)\chi(T^\natural)
  =\chi(S)\chi(T).
\]

Since the elements of \(\mathcal L'\) are precisely the complex Radon
measures on \(Z\), Proposition~3 yields the following results.

\begin{proposition}\label{prop:III-5-4}
There is a canonical bijection \(\varphi\mapsto\nu\) from
\(\mathcal L\) onto the set of complex Radon measures on \(Z\). It is
defined by
\[
  \varphi(T)=\int_Z\chi(T)\,d\nu(\chi)
  \qquad(T\in\mathcal A).
\]
The form \(\varphi\) is a trace if and only if \(\nu\) is positive.
\end{proposition}

\phantomsection
\begin{corollary*}\phantomsection\label{cor:III-5-4}
Let \(\mathcal C\) be the convex subset of \(\mathcal L\) consisting
of the finite traces on \(\mathcal A\) such that \(\varphi(I)\leq1\).
Then the nonzero extreme points of \(\mathcal C\) are the characters of
\(\mathcal A\).
\end{corollary*}

\begin{proof}
Indeed, the point masses of mass~1 on \(Z\) are the nonzero extreme
points of the convex set of positive measures \(\nu\) such that
\(\nu(1)\leq1\).
\end{proof}

\begin{proposition}\label{prop:III-5-5}
Let \(\chi\) be a character of \(\mathcal A\). The set
\[
  \mathfrak m=\{T\in\mathcal A:\chi(T^*T)=0\}
\]
is a maximal two-sided ideal of \(\mathcal A\). The mapping
\(\chi\mapsto\mathfrak m\) is a bijection from the set of characters of
\(\mathcal A\) onto the set of maximal two-sided ideals of
\(\mathcal A\).
\end{proposition}

\begin{proof}
First observe that the conditions \(\chi(T^*T)=0\), on the one hand,
and \(\chi(T_1T)=0\) for every \(T_1\in\mathcal A\), on the other, are
equivalent by the Cauchy--Schwarz inequality. Let \(\chi'\) be the
restriction of \(\chi\) to \(\mathcal Z\). This is a character of
\(\mathcal Z\), and hence determines a maximal ideal \(\mathfrak n\)
of \(\mathcal Z\). The condition \(\chi(T_1T)=0\) for every
\(T_1\in\mathcal A\) is equivalent to
\(\chi'((T_1T)^\natural)=0\) for every \(T_1\in\mathcal A\), hence to
\[
  (T_1T)^\natural\in\mathfrak n
  \qquad\text{for every }T_1\in\mathcal A.
\]
By Proposition~2 this condition defines a maximal two-sided ideal
\(\mathfrak m\) of \(\mathcal A\). Moreover, the mappings
\(\chi\mapsto\chi'\), \(\chi'\mapsto\mathfrak n\), and
\(\mathfrak n\mapsto\mathfrak m\) are bijections among, respectively,
the characters of \(\mathcal A\), the characters of \(\mathcal Z\),
the maximal ideals of \(\mathcal Z\), and the maximal two-sided ideals
of \(\mathcal A\).
\end{proof}

Bibliography: \ArticleRef{29}, \ArticleRef{63}, \ArticleRef{69},
\ArticleRef{70}, \ArticleRef{107}, \ArticleRef{121}, \ArticleRef{123},
\ArticleRef{124}.

\noindent\textit{Exercises.}
\begin{enumerate}[label=\arabic*.,leftmargin=*]
\item Let \(\mathcal A\) be a finite von Neumann algebra, let
\(\mathcal Z\) be its centre, let \(T\in\mathcal A\), and let
\(\mathcal K'''_T\) be the weak closure of \(\mathcal K_T\). Show that
\[
  \mathcal K'''_T\cap\mathcal Z=\{T^\natural\},
\]
where \(T\mapsto T^\natural\) denotes the canonical
\(\mathcal Z\)-trace of \(\mathcal A\). Argue as in the corollary to
Theorem~1, using the ultra-weak continuity of
\(T\mapsto T^\natural\). \ArticleRef{6}
% Source PDF page 268 (printed p. 259)
\phantomsection\label{srcpage:259}
\item\label{ex:III-5-2}
\begin{enumerate}[label=\textit{\alph*.},leftmargin=2em]
\item Let \(\mathcal A\) be a von Neumann algebra, let \(\mathcal Z\)
be its centre, let \(T_1,T_2,\ldots,T_n\) be elements of \(\mathcal A\),
and let \(S_1\in\mathcal K''_{T_1}\). There is a sequence
\((f_1,f_2,\ldots)\) of elements of \(\mathcal E\) such that
\(f_j\mathbin{\cdot}T_1\) converges (in the norm topology) to \(S_1\),
and \(f_j\mathbin{\cdot}T_h\) converges to an element of
\(\mathcal K''_{T_h}\), for \(h=2,3,\ldots,n\). (Choose
\(g_q\in\mathcal E\) so that
\(\lVert g_q\mathbin{\cdot}T_1-S_1\rVert\leq2^{-q}\); then apply
\hyperref[lem:III-5-4]{Lemma~4} to
\(g_q\mathbin{\cdot}T_1,g_q\mathbin{\cdot}T_2,\ldots,
g_q\mathbin{\cdot}T_n\) in order to determine \(f_q\).)

\item Let \(\mathfrak n\) be an ideal of \(\mathcal Z\), closed in the
norm topology. Let \(S\in\mathcal A_+\) and \(T\in\mathcal A_+\).
If \(\mathcal K''_T\subset\mathfrak n\) and \(S\leq T\), then
\(\mathcal K''_S\subset\mathfrak n\). (Using \textit{a}, show that every
element of \(\mathcal K''_S\) is majorized by an element of
\(\mathcal K''_T\).)
\end{enumerate}

\item\label{ex:III-5-3}
Let \(\mathcal A\) be a von Neumann algebra, let \(\mathcal Z\) be its
centre, let \(\mathfrak n\) be an ideal of \(\mathcal Z\), closed in the
norm topology, and let \(\mathfrak m\) be the largest two-sided ideal of
\(\mathcal A\) such that \(\mathfrak m\cap\mathcal Z=\mathfrak n\).
Show that \(\mathfrak m\) is equal to the set \(\mathfrak m'\) of the
\(T\in\mathcal A\) such that
\(\mathcal K''_{T^*T}\subset\mathfrak n\). [First show that
\(\mathfrak m\subset\mathfrak m'\) and that
\(\mathfrak m'\cap\mathcal Z=\mathfrak n\); it then suffices to prove
that \(\mathfrak m'\) is a two-sided ideal of \(\mathcal A\). For every
unitary operator \(U\in\mathcal A\) and every \(T\in\mathfrak m'\), one
has \(UT\in\mathfrak m'\) and \(TU\in\mathfrak m'\). Finally, let
\(T_1,T_2\in\mathfrak m'\). One has
\[
 (T_1+T_2)^*(T_1+T_2)
 \leq 2T_1^*T_1+2T_2^*T_2,
\]
and \(\mathcal K''_{T_1^*T_1+T_2^*T_2}\) is contained in the closure of
\(\mathcal K''_{T_1^*T_1}+\mathcal K''_{T_2^*T_2}\), hence in
 \(\mathfrak n\); thus \(T_1+T_2\in\mathfrak m'\), by
 \hyperref[ex:III-5-2]{Exercise~2\textit{b}}.]
If \(\mathcal A\) is finite, \(\mathfrak m\) is the set of the
\(T\in\mathcal A\) such that \((T^*T)^\natural\in\mathfrak n\).
Recover this result by using \hyperref[prop:III-5-5]{Proposition~5}.
\ArticleRef{29},
\ArticleRef{124}

\item\label{ex:III-5-4}
Let \(\mathcal A\) be a finite von Neumann algebra and let \(\Phi\) be a
finite \(\mathcal Z\)-trace on \(\mathcal A\). Show that
\[
 \Phi(T)=T^\natural\Phi(I)\qquad(T\in\mathcal A),
\]
where \(T\mapsto T^\natural\) denotes the canonical
\(\mathcal Z\)-trace of \(\mathcal A\). (Use
\hyperref[thm:III-5-1]{Theorem~1}.) Deduce that
\(\Phi\) is normal. \ArticleRef{12}

\item\label{ex:III-5-5}
Let \(\mathfrak H\) be a complex Hilbert space, and let
\(\mathfrak H_1,\mathfrak H_2\) be two supplementary orthogonal
subspaces of \(\mathfrak H\), of infinite dimensions. Let \(\mathcal A\)
be the set of the operators of the form
\[
 S=\lambda_1P_{\mathfrak H_1}+\lambda_2P_{\mathfrak H_2}+T,
\]
where \(T\in\mathcal L(\mathfrak H)\) is compact and
\(\lambda_1,\lambda_2\in\mathbb C\). Show that \(\mathcal A\) is a
unital C*-algebra whose centre consists only of the scalar operators.
Show that, if \(S\) is unitary, then
\(\lvert\lambda_1\rvert=\lvert\lambda_2\rvert=1\). Deduce that, if
\(U_1,U_2,\ldots,U_n\) are unitary operators of \(\mathcal A\) and
\(\mu_1,\mu_2,\ldots,\mu_n\) are nonnegative numbers with sum \(1\),
then
\[
 \sum_{i=1}^n
 \mu_iU_i(P_{\mathfrak H_1}-P_{\mathfrak H_2})U_i^{-1}
\]
has the form
\(P_{\mathfrak H_1}-P_{\mathfrak H_2}+R\), with \(R\) compact, and
therefore cannot converge in norm to an operator of the centre of
\(\mathcal A\). \ArticleRef{6}

\item\label{ex:III-5-6}
For every \(n\in\mathbb N\) (the set of integers \(\geq0\)), let
\(\mathfrak H_n\) be a Hilbert space of dimension \(n\), let
\(\mathcal A_n=\mathcal L(\mathfrak H_n)\), and let \(\mathcal Z_n\)
be the centre of \(\mathcal A_n\) (identified with \(\mathbb C\)). Let
\[
 \mathcal A=\prod_{n\in\mathbb N}\mathcal A_n
\]
be the finite von Neumann algebra with centre
\[
 \mathcal Z=\prod_{n\in\mathbb N}\mathcal Z_n;
\]
thus \(\mathcal Z\) is identified with the algebra
\(L_{\mathbb C}^{\infty}(\mathbb N)\) of bounded sequences
\((\lambda_0,\lambda_1,\ldots)\) of complex numbers. Let \(\mathcal U\)
be a nontrivial ultrafilter on \(\mathbb N\). Let \(\mathfrak n\) be the
maximal ideal of \(\mathcal Z\) formed by the sequences
\((\lambda_i)\) such that
\(\lim_{\mathcal U}\lambda_i=0\). Let \(\mathfrak m\) be the largest
two-sided ideal of \(\mathcal A\) such that
\(\mathfrak m\cap\mathcal Z=\mathfrak n\), and let \(\mathfrak m'\) be
the two-sided ideal of \(\mathcal A\) generated by \(\mathfrak n\).
Show that \(\mathfrak m\) is distinct from the norm closure of
\(\mathfrak m'\). [Let
\(T=(T_n)_{n\in\mathbb N}\in\mathcal A\), where \(T_n\) is a rank-one
projection in \(\mathfrak H_n\). One has
\((TS)^\natural\in\mathfrak n\) for every \(S\in\mathcal A\), since the
sequence with which \((TS)^\natural\) is identified tends to zero. But,
for every \(R\in\mathfrak m'\), one has
\(\lVert T-R\rVert\geq1\).]

\item\label{ex:III-5-7}
Let \(\mathcal A\) be a von Neumann algebra, let \(\mathcal Z\) be its
centre, let \(\mathfrak m\) be a two-sided ideal of \(\mathcal A\),
closed in the norm topology, let \(\mathcal B\) be the involutive
algebra \(\mathcal A/\mathfrak m\), let \(\mathcal Y\) be the centre of
\(\mathcal B\), and let\space%
% Source PDF page 269 (printed p. 260)
\phantomsection\label{srcpage:260}%
\(\Phi\) be the canonical mapping of \(\mathcal A\) onto \(\mathcal B\).
Show that \(\Phi(\mathcal Z)=\mathcal Y\). [If \(T\in\mathcal A\) is
such that \(\Phi(T)\in\mathcal Y\), then
\(T-UTU^{-1}\in\mathfrak m\) for every unitary operator
\(U\in\mathcal A\), and hence \(T-T_1\in\mathfrak m\) for every
\(T_1\in\mathcal K'_T\), therefore, in particular, for every
\(T_1\in\mathcal K''_T\).]

\item\label{ex:III-5-8}
Let \(\mathcal A\) be a finite von Neumann algebra, let \(\mathcal Z\)
be its centre, let \(\varphi\) be a trace on \(\mathcal A\), and let
\(\mathfrak m\) be the two-sided ideal of \(\mathcal A\) formed by the
\(T\in\mathcal A\) such that \(\varphi(T^*T)=0\). The sesquilinear form
\[
 (S,T)\longmapsto\varphi(ST^*)
\]
on \(\mathcal A\) defines, on passing to the quotient, an inner product
on \(\mathfrak A=\mathcal A/\mathfrak m\) which makes
\(\mathfrak A\) a Hilbert algebra with an identity element. Let
\(\mathfrak H\) be the completed Hilbert space of \(\mathfrak A\), let
\(\mathfrak Z\) be the set of the central elements of \(\mathfrak A\),
and let \(\Phi\) be the canonical mapping of \(\mathcal A\) onto
\(\mathfrak A\).
\begin{enumerate}[label=\textit{\alph*.},leftmargin=2em]
\item \(\Phi(\mathcal Z)\) is everywhere dense in \(\mathfrak Z\).
[Let \(T\in\mathcal A\), and let \(U\) be a unitary element of
\(\mathcal A\). Show that
\[
 \Phi(T)-\Phi(U)\Phi(T)\Phi(U)^{-1}
\]
is orthogonal to \(\mathfrak Z\). Deduce that \(\Phi(T^\natural)\) is
the orthogonal projection of \(\Phi(T)\) onto \(\mathfrak Z\).]

\noindent\textit{Problem.} Is \(\mathfrak A\) complete?

\item In order that \(\mathcal U(\mathfrak A)\) and
\(\mathcal V(\mathfrak A)\) be factors, it is necessary and sufficient
that \(\varphi\) be a character of \(\mathcal A\). [Using \textit{a}
and \hyperref[ex:III-4-2]{\S~4, Exercise~2\textit{c}}, show that it is
necessary and sufficient
that \(\mathcal Z/(\mathfrak m\cap\mathcal Z)\) have dimension \(1\).]

\item Take for \(\mathcal A\) the example in
\hyperref[ex:III-5-6]{Exercise~6}. Show that
\(\varphi\) can be chosen so that \(\mathfrak m\) is the ideal considered
in \hyperref[ex:III-5-6]{Exercise~6} and
\(\mathcal U(\mathfrak A)\) is then a factor of type
\(\mathrm{II}_1\). (By \textit{b}, it is enough to prove that
\(\mathcal A/\mathfrak m\) has infinite dimension.) This gives a new
proof of the existence of factors of type \(\mathrm{II}_1\).
\ArticleRef{124}
\end{enumerate}

\item\label{ex:III-5-9}
\begin{enumerate}[label=\textit{\alph*.},leftmargin=2em]
\item Let \(\mathcal A\) be a unital C*-algebra, let \(\mathcal B\) be
a unital C*-subalgebra, and let \(\pi\) be a linear projection of
\(\mathcal A\) onto \(\mathcal B\) such that \(\lVert\pi\rVert\leq1\).
Then \(\pi\) is positive,
\[
 \pi(ATB)=A\pi(T)B
 \quad(A,B\in\mathcal B,\ T\in\mathcal A),
\]
and
\[
 \pi(T)^*\pi(T)\leq\pi(T^*T)
 \quad(T\in\mathcal A).
\]
\ArticleRef{195}

\item Let \(\mathcal A\) be a von Neumann algebra on
\(\mathfrak H\). For every \(T\in\mathcal L(\mathfrak H)\), let
\(\mathcal K'''_T\) be the weakly closed convex hull of the set of the
\(UTU^{-1}\), where \(U\) ranges over the unitary elements of
\(\mathcal A\). One says that \(\mathcal A\) has the \emph{Schwartz
property} if, for every \(T\in\mathcal L(\mathfrak H)\),
\(\mathcal K'''_T\) contains an element of \(\mathcal A'\). If this is so,
there is a linear projection \(\pi\) of \(\mathcal L(\mathfrak H)\) onto
\(\mathcal A'\) such that \(\lVert\pi\rVert\leq1\) and
\(\pi(T)\in\mathcal K'''_T\) for every
\(T\in\mathcal L(\mathfrak H)\).
\ArticleRef{309}

\item There are factors of type \(\mathrm{II}_1\) which have the
Schwartz property, and others which do not. \ArticleRef{309}

\item Suppose that the unitary group of \(\mathcal A\) contains an
amenable subgroup \(G\) such that \(\mathcal A\) is the von Neumann
algebra generated by \(G\). Then \(\mathcal A\) has the Schwartz
property; this is so, in particular, if \(\mathcal A\) is abelian.
\ArticleRef{296}, \ArticleRef{319}, \ArticleRef{336}, \ArticleRef{399}

If \(\mathcal A\) is the algebra of multiplication operators by bounded
measurable functions on \(L^2([0,1])\), there are at least two distinct
mappings
\[
 \pi:\mathcal L(\mathfrak H)\longrightarrow\mathcal A'=\mathcal A
\]
having the properties in \textit{b}. \ArticleRef{230}

\item Let \(\mathcal A\) be an abelian von Neumann algebra, and let
\(\pi\) be a positive linear mapping of \(\mathcal L(\mathfrak H)\)
into \(\mathcal A'\) such that \(\pi(T)=T\) for \(T\in\mathcal A\).
Then
\[
 \pi(ST)=\pi(S)T,
 \qquad
 \pi(TS)=T\pi(S)
\]
for every \(S\in\mathcal L(\mathfrak H)\) and every
\(T\in\mathcal A\). \ArticleRef{230}

\item Let \(\mathcal A\) be a von Neumann algebra, let \(\varphi\) be
a faithful semifinite normal trace on \(\mathcal A_+\), let
\(\mathfrak m\) be the definition ideal of \(\varphi\), and let
\(\mathcal B\) be a von Neumann subalgebra of \(\mathcal A\). Put
\(\mathfrak n=\mathfrak m\cap\mathcal B\), and suppose that
\(\mathfrak n\) is strongly dense in \(\mathcal B\). Let
\(\mathcal B^{\perp}\) be the set of the \(T\in\mathcal A\) such that
\(\varphi(TS)=0\) for \(S\in\mathfrak n\). The sum
\(\mathcal B+\mathcal B^{\perp}\) is direct and equal to
\(\mathcal A\). Let \(\pi\) be the corresponding projection of
\(\mathcal A\) onto \(\mathcal B\). Then \(\lVert\pi\rVert\leq1\),
and \(\pi\) is positive, faithful, ultra-weakly continuous, and
ultra-strongly continuous; if \(S\in\mathcal A\) and if
\(T\in\mathcal A\cap\mathcal B'=\mathcal C\), then%
% Source PDF page 270 (printed p. 261)
\phantomsection\label{srcpage:261}%
\[
 \pi(ST)=\pi(TS).
\]
If \(\mathcal B\) is the set of the elements of \(\mathcal A\) which
commute with \(\mathcal C\), then, for every \(T\in\mathfrak m\), the
weakly closed convex hull of the set
\[
 \{UTU^{-1}\mid U\text{ is a unitary element of }\mathcal C\}
\]
meets \(\mathcal B\) in a unique point, namely \(\pi(T)\), and this point
belongs to \(\mathfrak n\). \ArticleRef{15}, \ArticleRef{130},
\ArticleRef{134}, \ArticleRef{179}, \ArticleRef{241}, \ArticleRef{278},
\ArticleRef{293}

\item Let \(\mathcal A\) be a von Neumann algebra on \(\mathfrak H\),
let \(G\) be a group, and let \(g\mapsto U_g\) be a unitary
representation of \(G\) on \(\mathfrak H\) such that
\[
 U_g\mathcal A U_g^{-1}=\mathcal A\qquad(g\in G).
\]
Let \(U_G\) be the set of the \(U_g\), for \(g\in G\), put
\(\mathcal B=\mathcal A\cap U_G'\), and let \(E_0\) be the projection of
\(\mathfrak H\) onto the subspace of the vectors invariant under
\(U_G\). One has
\[
 E_0\in U_G'\cap U_G''\subset U_G'\subset\mathcal B'.
\]
Suppose that the central support of \(E_0\), relative to
\(\mathcal B'\), is \(I\); in other words, suppose that
\(\mathcal A'E_0(\mathfrak H)\) is total in \(\mathfrak H\). Then there
is a unique normal \(G\)-invariant linear projection \(\pi\) of
\(\mathcal A\) onto \(\mathcal B\), and \(\pi\) is positive and
faithful. For \(T\in\mathcal A\), \(\pi(T)\) is the unique point of
\(\mathcal B\) in the weakly closed convex hull of the
\(U_gTU_g^{-1}\), \(g\in G\), and \(\pi(T)\) is the unique element of
\(\mathcal A\) such that
\[
 \pi(T)E_0=E_0TE_0.
\]
A normal positive linear functional \(\omega\) on \(\mathcal A\) is
\(G\)-invariant if and only if
\[
 \omega=(\omega\vert_{\mathcal B})\circ\pi.
\]
\ArticleRef{370}, \ArticleRef{371}, \ArticleRef{444}
\end{enumerate}
\end{enumerate}

% This section is defined here only to defer its expansion until after
% the intervening Section 6 material, which follows below in source order.
\long\gdef\DixmierDeferredSectionSeven{
% Source PDF page 279 (printed p. 270)
\phantomsection\label{srcpage:270}
\section{Hyperfinite Factors}\label{sec:III-7}

\NumberedParagraph{1}{Factors Contained in a Finite von Neumann Algebra}
\label{no:III-7-1}

\begin{lemma}\label{lem:III-7-1}
Let \(\mathcal A\) be a von Neumann algebra, let \(\varphi\) be a
faithful finite normal trace on \(\mathcal A\), let \(\mathcal B\) be
an involutive subalgebra of \(\mathcal A\) containing \(I\), let
\(T\in\mathcal A\), and let \(\lVert\cdot\rVert_2\) be the norm defined
by \(\varphi\) (\hyperref[no:I-6-2]{Chapter~I, \S~6, no.~2}). The
following conditions are
equivalent:
\begin{enumerate}[label=\textup{(\roman*)},leftmargin=*]
\item \(T\) is a limit, for the norm \(\lVert\cdot\rVert_2\), of
elements of \(\mathcal B\).
\item \(T\) is a limit, for the norm \(\lVert\cdot\rVert_2\), of
elements \(S\in\mathcal B\) whose norms \(\lVert S\rVert\) are bounded.
\item \(T\) is an ultra-strong limit of elements of \(\mathcal B\).
\end{enumerate}
\end{lemma}

\begin{proof}
The implication \textup{(ii)}\(\Rightarrow\)\textup{(iii)} follows
from \hyperref[prop:I-4-4]{Chapter~I, \S~4, Proposition~4}. One has
\textup{(iii)}\(\Rightarrow\)\textup{(i)} because, if \(S\) converges
ultra-strongly to zero, then \(S^*S\) converges ultra-weakly to zero,
and hence
\[
 \lVert S\rVert_2=\varphi(S^*S)^{1/2}\longrightarrow0.
\]
Finally, suppose that condition~\textup{(i)} holds, and let us show that
condition~\textup{(ii)} holds. The closure \(\overline{\mathcal B}\) of
\(\mathcal B\) for the norm \(\lVert\cdot\rVert_2\) is an involutive
subalgebra of \(\mathcal A\) containing \(I\), and, by what precedes,
is ultra-strongly closed; it is therefore a von Neumann algebra. By
\hyperref[thm:I-6-2]{Chapter~I, \S~6, Theorem~2},
\(\overline{\mathcal B}\) is canonically
equipped with a Hilbert-algebra structure. The element \(T\) is bounded
relative to \(\overline{\mathcal B}\); hence
(\hyperref[prop:I-5-4]{Chapter~I, \S~5, Proposition~4}) \(T\) is a
limit, for the norm
\(\lVert\cdot\rVert_2\), of elements \(T_n\in\mathcal B\) such that the
\(\lVert U_{T_n}\rVert\) are bounded. The \(\lVert T_n\rVert\) are then
bounded, since the mapping
\[
 S\longmapsto U_S\qquad(S\in\overline{\mathcal B})
\]
is an isomorphism of von Neumann algebras.
\end{proof}

\begin{lemma}\label{lem:III-7-2}
Let \(\mathcal A\) be a finite von Neumann algebra, let \(\mathcal Z\)
be its centre, let \(\varphi\) be a faithful finite normal trace on
\(\mathcal A\), let \(A\mapsto A^\natural\) be the canonical
\(\mathcal Z\)-trace of \(\mathcal A\), and let
\(\lVert\cdot\rVert_2\) be the norm defined by \(\varphi\). Let
\(T\in\mathcal A\) and \(\varepsilon>0\) be such that
\[
 \lVert TT'-T'T\rVert_2\leq\varepsilon\lVert T'\rVert
 \qquad(T'\in\mathcal A).
\]
Then
\[
 \lVert T-T^\natural\rVert_2\leq\varepsilon.
\]
\end{lemma}

\begin{proof}
Let \(\eta>0\). There are unitary operators
\(U_1,U_2,\ldots,U_n\in\mathcal A\) and nonnegative numbers
\(\lambda_1,\lambda_2,\ldots,\lambda_n\) with sum \(1\), such that
\[
 \left\lVert
   \sum_{i=1}^n\lambda_iU_i^{-1}TU_i-T^\natural
 \right\rVert\leq\eta
\]
(\hyperref[cor:III-5-1]{\S~5, the corollary to Theorem~1}). One has
\begin{align*}
 \left\lVert
   \sum_{i=1}^n\lambda_iU_i^{-1}TU_i-T
 \right\rVert_2
 &\leq\sum_{i=1}^n\lambda_i
       \lVert U_i^{-1}TU_i-T\rVert_2\\
 &=\sum_{i=1}^n\lambda_i
       \lVert TU_i-U_iT\rVert_2\\
 &\leq\sum_{i=1}^n\lambda_i\varepsilon
  =\varepsilon.
\end{align*}%
% Source PDF page 280 (printed p. 271)
\phantomsection\label{srcpage:271}%
Consequently
\[
 \lVert T-T^\natural\rVert_2
 \leq\varepsilon+\eta\varphi(I)^{1/2},
\]
which, in view of the arbitrariness of \(\eta\), proves the lemma.
\end{proof}

\begin{proposition}\label{prop:III-7-1}
Let \(\mathcal A\) be a finite von Neumann algebra, let
\((\mathcal A_\iota)_{\iota\in I}\) be a family of factors contained in
\(\mathcal A\), totally ordered by inclusion, and let \(\mathcal B\) be
the von Neumann algebra generated by the \(\mathcal A_\iota\). Then
\(\mathcal B\) is a finite factor.
\end{proposition}

\begin{proof}
Let \(S_0\) be an element of the centre of \(\mathcal B\). It is enough
to show that, for every projection \(E\) of the centre \(\mathcal Z\) of
\(\mathcal A\) such that \(\mathcal Z_E\) is of countable type,
\((S_0)_E\) is a scalar. We may therefore reduce to the case in which
the centre of \(\mathcal A\) is of countable type. Let \(\varphi\) be a
faithful finite normal trace on \(\mathcal A\). For every
\(\varepsilon>0\), there are (\hyperref[lem:III-7-1]{Lemma~1}) an
\(\iota\in I\) and an
\(S\in\mathcal A_\iota\) such that
\[
 \lVert S-S_0\rVert_2\leq\varepsilon,
\]
where \(\lVert\cdot\rVert_2\) denotes the norm defined by \(\varphi\)
on \(\mathcal A\). For every \(T\in\mathcal A_\iota\), one has
\[
 \lVert ST-TS\rVert_2
 =\lVert(S-S_0)T-T(S-S_0)\rVert_2
 \leq2\varepsilon\lVert T\rVert.
\]
Thus (\hyperref[lem:III-7-2]{Lemma~2})
\[
 \lVert S-S^\natural\rVert_2\leq2\varepsilon,
\]
where \(S\mapsto S^\natural\) denotes the canonical trace of
\(\mathcal A_\iota\). Hence
\[
 \lVert S_0-S^\natural\rVert_2\leq3\varepsilon.
\]
Since \(S^\natural\) is scalar, \(S_0\), which is a limit of scalar
operators, is scalar.
\end{proof}

Concerning \hyperref[prop:III-7-1]{Proposition~1}, cf.
\hyperref[ex:III-7-1]{Exercise~1}.

Bibliography: \ArticleRef{22}, \ArticleRef{67}.

\NumberedParagraph{2}{Existence and Uniqueness of Continuous
Hyperfinite Factors}
\label{no:III-7-2}

\begin{definition}\label{def:III-7-1}
A factor \(\mathcal A\) is said to be \emph{hyperfinite} if it satisfies
one of the following conditions:
\begin{enumerate}[label=\textup{(\roman*)},leftmargin=*]
\item \(\mathcal A\) is of type \(\mathrm I_p\), with \(p<+\infty\).
\item \(\mathcal A\) is finite and is the von Neumann algebra generated
by an increasing sequence of factors of types
\[
 \mathrm I_1,\ \mathrm I_2,\ \mathrm I_4,\ \mathrm I_8,\ldots,
 \mathrm I_{2^n},\ldots.
\]
\end{enumerate}
\end{definition}

Every factor isomorphic or anti-isomorphic to a hyperfinite factor is
hyperfinite.

\begin{lemma}\label{lem:III-7-3}
Let \(\mathcal A\) be a factor of type \(\mathrm I_n\), with \(n\)
finite, on \(\mathfrak H\), and let \(\mathcal B\) be a factor
containing \(\mathcal A\) which is either finite continuous or of type
\(\mathrm I_{2pn}\), where \(p\geq1\). There is a factor
\(\mathcal A_1\) of type \(\mathrm I_{2n}\) such that
\[
 \mathcal A\subset\mathcal A_1\subset\mathcal B.
\]
\end{lemma}

\begin{proof}
Since \(\mathcal A\) is isomorphic to \(\mathcal L(\mathfrak H_n)\),
where \(\mathfrak H_n\) is a Hilbert space of dimension \(n\), there
are \(n\) pairwise orthogonal, equivalent minimal projections
\(E_1,E_2,\ldots,E_n\) of \(\mathcal A\), with sum \(I\) (equivalent
relative to \(\mathcal A\), hence relative to \(\mathcal B\)). Put
\[
 \mathfrak K=E_1(\mathfrak H),
 \qquad
 \mathcal C=\mathcal B_{\mathfrak K}.
\]
There is a spatial isomorphism of \(\mathcal B\) onto
\(\mathcal C\otimes\mathcal L(\mathfrak H_n)\) which transforms
\(\mathcal A\) into\space%
% Source PDF page 281 (printed p. 272)
\phantomsection\label{srcpage:272}%
\(\mathbb C_{\mathfrak K}\otimes\mathcal L(\mathfrak H_n)\). Now, if
\(\mathcal B\) is finite continuous (respectively, of type
\(\mathrm I_{2pn}\)), then \(\mathcal C\) is finite continuous
(respectively, of type \(\mathrm I_{2p}\)). In both cases,
\(\mathcal C\) has two orthogonal, equivalent projections
\(F_1,F_2\) with sum \(I\) (in the first case, this follows from
\hyperref[cor:III-1-thm1-3]{\S~1, Theorem~1, Corollary~3}). Let
\(V\in\mathcal C\) be a partially
isometric operator such that
\[
 V^*V=F_1,
 \qquad
 VV^*=F_2.
\]
Then the linear combinations of \(F_1,F_2,V,V^*\) form a factor
\(\mathcal C_1\) of type \(\mathrm I_2\) contained in \(\mathcal C\).
Putting
\[
 \mathcal A_1=\mathcal C_1\otimes\mathcal L(\mathfrak H_n),
\]
one has \(\mathcal A\subset\mathcal A_1\subset\mathcal B\), and
\(\mathcal A_1\) is of type \(\mathrm I_{2n}\).
\end{proof}

\begin{theorem}\label{thm:III-7-1}
Every finite continuous factor contains a continuous hyperfinite factor.
\end{theorem}

\begin{proof}
Let \(\mathcal A\) be a finite continuous factor.
\hyperref[lem:III-7-3]{Lemma~3}, applied
recursively, proves that there is an increasing sequence
\(\mathcal A_0,\mathcal A_1,\ldots\) of factors contained in
\(\mathcal A\), such that \(\mathcal A_i\) is of type
\(\mathrm I_{2^i}\). Let \(\mathcal B\) be the von Neumann algebra
generated by the \(\mathcal A_i\). By
\hyperref[prop:III-7-1]{Proposition~1},
\(\mathcal B\) is a finite factor. One has
\(\dim\mathcal B=+\infty\), and hence \(\mathcal B\) is continuous.
Finally, \(\mathcal B\) is plainly hyperfinite.
\end{proof}

\begin{corollary}\label{cor:III-7-1}
There are continuous hyperfinite factors acting on a Hilbert space with
a countable basis.
\end{corollary}

\begin{proof}
This follows from \hyperref[thm:III-7-1]{Theorem~1} and from the
existence of finite continuous
factors acting on a Hilbert space with a countable basis (Chapter~I,
\hyperref[thm:I-9-1]{\S~9, no.~4, Theorem~1} and the remark).
\end{proof}

Given a finite factor \(\mathcal A\), in what follows we shall always
denote by \(\lVert\cdot\rVert_2\) the norm defined by the canonical
trace of \(\mathcal A\), and shall call the corresponding metric the
\emph{canonical metric} of \(\mathcal A\).

\begin{theorem}\label{thm:III-7-2}
Any two continuous hyperfinite factors are isomorphic.
\end{theorem}

\begin{proof}
Let \(\mathcal A,\mathcal B\) be two continuous hyperfinite factors,
and let \((\mathcal A_i)\) [respectively, \((\mathcal B_i)\)] be an
increasing sequence of factors contained in \(\mathcal A\)
[respectively, in \(\mathcal B\)], of types \(\mathrm I_{2^i}\),
generating \(\mathcal A\) [respectively, \(\mathcal B\)]. Suppose that,
for \(1\leq i\leq n\), isomorphisms \(\Phi_i\) of \(\mathcal A_i\)
onto \(\mathcal B_i\) have been constructed which extend one another.
Since \(\mathcal A_{n+1}\) is isomorphic to
\(\mathcal A_n\otimes\mathcal L(\mathfrak H_2)\), and
\(\mathcal B_{n+1}\) to
\(\mathcal B_n\otimes\mathcal L(\mathfrak H_2)\), where
\(\mathfrak H_2\) is a Hilbert space of dimension \(2\), there is an
isomorphism of \(\mathcal A_{n+1}\) onto \(\mathcal B_{n+1}\) which
extends \(\Phi_n\). Recursively, there is therefore, for every \(i\),
an isomorphism \(\Phi_i\) of \(\mathcal A_i\) onto \(\mathcal B_i\)
such that \(\Phi_j\) extends \(\Phi_i\) for \(j\geq i\). Put
\[
 \mathcal C=\bigcup_{i=1}^{\infty}\mathcal A_i,
 \qquad
 \mathcal D=\bigcup_{i=1}^{\infty}\mathcal B_i.
\]
Then \(\mathcal C\) (respectively, \(\mathcal D\)) is an involutive
subalgebra of \(\mathcal A\) (respectively, of \(\mathcal B\)),
everywhere dense in \(\mathcal A\) (respectively, in \(\mathcal B\))
for the ultra-strong topology; and the \(\Phi_i\) define\space%
% Source PDF page 282 (printed p. 273)
\phantomsection\label{srcpage:273}%
an isomorphism \(\Phi\) of \(\mathcal C\) onto \(\mathcal D\). The
canonical metric of \(\mathcal A\) (respectively, of \(\mathcal B\))
induces on \(\mathcal A_i\) (respectively, on \(\mathcal B_i\)) the
canonical metric of \(\mathcal A_i\) (respectively, of
\(\mathcal B_i\)). Since \(\Phi_i\) transforms the canonical trace of
\(\mathcal A_i\) into that of \(\mathcal B_i\), \(\Phi_i\) is
isometric, and ultimately \(\Phi\) is isometric for the canonical
metrics. Now \(\mathcal C\) (respectively, \(\mathcal D\)) is
everywhere dense in \(\mathcal A\) (respectively, in \(\mathcal B\))
for the canonical metric (\hyperref[lem:III-7-1]{Lemma~1}). The
isomorphism \(\Phi\) of the
Hilbert algebra \(\mathcal C\) onto the Hilbert algebra \(\mathcal D\)
extends to an isomorphism of the completed Hilbert algebra extending
\(\mathcal C\) onto the completed Hilbert algebra extending
\(\mathcal D\). Finally, these completed Hilbert algebras are
\(\mathcal A\) and \(\mathcal B\)
(\hyperref[thm:I-6-2]{Chapter~I, \S~6, Theorem~2}).
\end{proof}

Continuous hyperfinite factors are also called ``approximately finite.''
This terminology accords poorly with the now commonly accepted
terminology ``finite factor.''

Bibliography: \ArticleRef{67}.

\NumberedParagraph{3}{Some Inequalities}
\label{no:III-7-3}

\begin{lemma}\label{lem:III-7-4}
Let \(\mathcal A\) be a finite factor, let \(E\) be a projection of
\(\mathcal A\), and let \(T\) be a Hermitian element of \(\mathcal A\)
such that \(\lVert T\rVert\leq1\).
\begin{enumerate}[label=\textup{(\roman*)},leftmargin=*]
\item There is a spectral projection \(F\) of \(T\) such that
\[
 \lVert F-E\rVert_2\leq9\lVert T-E\rVert_2^{1/2}.
\]
\item If \(T\geq0\), then
\[
 \lVert T^{1/2}-E\rVert_2
 \leq13\lVert T-E\rVert_2^{1/4}.
\]
\end{enumerate}
\end{lemma}

\begin{proof}
One has
\[
 T^2-T=(T^2-ET)+(ET-E^2)+(E-T),
\]
and hence
\begin{align*}
 \lVert T^2-T\rVert_2
 &\leq\lVert T\rVert\,\lVert T-E\rVert_2
     +\lVert E\rVert\,\lVert T-E\rVert_2
     +\lVert T-E\rVert_2\\
 &\leq3\lVert T-E\rVert_2.
\end{align*}
Let \(\varepsilon\) be a number such that
\(0<\varepsilon<\tfrac12\). Let \(F,F_1\) be the largest spectral
projections of \(T\) such that
\[
 (1-\varepsilon)F\leq TF\leq F,
 \qquad
 -\varepsilon F_1\leq TF_1\leq\varepsilon F_1.
\]
Put \(F_2=I-F-F_1\). Since
\[
 \lvert\lambda^2-\lambda\rvert
 \geq\varepsilon-\varepsilon^2
 \geq\tfrac12\varepsilon
\]
for \(\lambda\) outside the intervals
\([ -\varepsilon,+\varepsilon]\) and
\([1-\varepsilon,1+\varepsilon]\), one has
\begin{align*}
 \tfrac12\varepsilon\lVert F_2\rVert_2
 &\leq\lVert(T^2-T)F_2\rVert_2\\
 &\leq\lVert T^2-T\rVert_2
 \leq3\lVert T-E\rVert_2,
\end{align*}
and therefore
\[
 \lVert F_2\rVert_2
 \leq\frac6\varepsilon\lVert T-E\rVert_2.
\]

}

\section{Coupling Operator}\label{sec:III-6}

\NumberedParagraph{1}{Coupling Operator}
\label{no:III-6-1}

\begin{lemma}\label{lem:III-6-1}
Let \(\mathcal A\) be a von Neumann algebra on \(\mathfrak H\), let
\(\mathcal Z\) be its centre, and let \(J\) be an involution of
\(\mathfrak H\) which commutes with the projections of \(\mathcal Z\)
and is such that \(J\mathcal A J=\mathcal A'\). For every
\(x\in\mathfrak H\), one has
\[
 J(\mathfrak X_x^{\mathcal A'})=\mathfrak X_{Jx}^{\mathcal A},
 \qquad
 \mathfrak X_x^{\mathcal A}\sim\mathfrak X_{Jx}^{\mathcal A}.
\]
\end{lemma}

\begin{proof}
As \(T'\) ranges over \(\mathcal A'\), \(T=JT'J\) ranges over
\(\mathcal A\), and
\[
 J(T'x)=T(Jx),
 \qquad\text{hence}\qquad
 J(\mathfrak X_x^{\mathcal A'})=\mathfrak X_{Jx}^{\mathcal A}.
\]
To prove the second assertion of the lemma, it is enough to consider the
case
\(\mathfrak X_x^{\mathcal A}\prec\mathfrak X_{Jx}^{\mathcal A}\)
(by \hyperref[thm:III-1-1]{Theorem~1 of \S~1} and because \(J\)
commutes with the projections of \(\mathcal Z\)). One then has
\(\mathfrak X_x^{\mathcal A'}\prec\mathfrak X_{Jx}^{\mathcal A'}\)
(\hyperref[thm:III-1-2]{\S~1, Theorem~2}); since \(J\) defines an
anti-isomorphism of
\(\mathcal A\) onto \(\mathcal A'\), it follows that
\[
 \mathfrak X_{Jx}^{\mathcal A}
 =J(\mathfrak X_x^{\mathcal A'})
 \prec J(\mathfrak X_{Jx}^{\mathcal A'})
 =\mathfrak X_x^{\mathcal A};
\]
hence \(\mathfrak X_x^{\mathcal A}\sim\mathfrak X_{Jx}^{\mathcal A}\)
(\hyperref[prop:III-1-1]{\S~1, Proposition~1}).
\end{proof}

\begin{corollary}\label{cor:III-6-1}
Let \(\mathfrak A\) be a Hilbert algebra, let \(\mathfrak H\) be the
completed Hilbert space of \(\mathfrak A\), let \(J\) be the involution
of \(\mathfrak H\) canonically defined by \(\mathfrak A\), and put
\[
 \mathcal Z=\mathcal U(\mathfrak A)\cap\mathcal V(\mathfrak A).
\]
Let \(\Phi\) be a normal \(\mathcal Z\)-trace on
\(\mathcal U(\mathfrak A)_+\), and let \(\Phi'\) be the
\(\mathcal Z\)-trace on \(\mathcal V(\mathfrak A)_+\) obtained by
transporting \(\Phi\) by \(J\). For every \(x\in\mathfrak H\), one has
\[
 \Phi\bigl(E_x^{\mathcal V(\mathfrak A)}\bigr)
 =\Phi'\bigl(E_x^{\mathcal U(\mathfrak A)}\bigr).
\]
\end{corollary}

\begin{proof}
By \hyperref[lem:III-6-1]{Lemma~1},
\(E_x^{\mathcal U(\mathfrak A)}\sim
JE_x^{\mathcal V(\mathfrak A)}J\); hence
\[
 \Phi\bigl(E_x^{\mathcal V(\mathfrak A)}\bigr)
 =\Phi'\bigl(JE_x^{\mathcal V(\mathfrak A)}J\bigr)
 =\Phi'\bigl(E_x^{\mathcal U(\mathfrak A)}\bigr).
\]
\end{proof}

% Source PDF page 271 (printed p. 262)
\phantomsection\label{srcpage:262}
Let \(\mathcal A\) be a von Neumann algebra on \(\mathfrak H\), and
let \(\mathcal Z\) be its centre. As in \hyperref[sec:III-4]{\S~4},
we choose, once and for
all, a locally compact space \(Z\), a positive measure \(\nu\) on
\(Z\), and an identification of \(\mathcal Z\) with
\(L_{\mathbb C}^{\infty}(Z,\nu)\), which embeds
\(\mathcal Z_+\) in the set \(\widehat{\mathcal Z}_+\) of measurable
nonnegative functions, finite or not, on \(Z\).

\begin{proposition}\label{prop:III-6-1}
Suppose that \(\mathcal A\) is semifinite. Let \(\Phi\) (respectively,
\(\Phi'\)) be a faithful semifinite normal \(\mathcal Z\)-trace on
\(\mathcal A_+\) (respectively, on \(\mathcal A'_+\)). There is a
unique element \(C\in\widehat{\mathcal Z}_+\) such that, for every
\(x\in\mathfrak H\),
\[
 \Phi(E_x^{\mathcal A'})=C\Phi'(E_x^{\mathcal A}).
\]
One has \(0<C(\zeta)<+\infty\) locally almost everywhere.
\end{proposition}

\begin{proof}
\noindent\textup{(i)} Let us prove the uniqueness of \(C\). Let
\(C_1\in\widehat{\mathcal Z}_+\) be such that
\[
 \Phi(E_x^{\mathcal A'})=C_1\Phi'(E_x^{\mathcal A})
 \qquad(x\in\mathfrak H).
\]
Suppose, for example, that \(C(\zeta)<C_1(\zeta)\) on a measurable
subset \(Z'\) of \(Z\) which is not locally negligible; we shall obtain
a contradiction. Let \(G\) be the projection of \(\mathcal Z\)
corresponding to the characteristic function of \(Z'\). Since \(\Phi'\)
is semifinite, there is a nonzero projection \(E'\in\mathcal A'\) such
that \(E'\leq G\) and \(\Phi'(E')\in\mathcal Z_+\). By decreasing
\(E'\), if necessary, we may suppose that \(E'\) has the form
\(E_x^{\mathcal A}\). Since \(\Phi'\) is faithful, one has
\[
 0<\Phi'(E_x^{\mathcal A})<+\infty
\]
on a subset of \(Z'\) which is not locally negligible. Then
\[
 C\Phi'(E_x^{\mathcal A})
 \ne C_1\Phi'(E_x^{\mathcal A}),
\]
which is absurd.

\noindent\textup{(ii)} When \(\mathcal A\) is given, it is enough to
verify the existence of \(C\) and the last assertion of the proposition
for one particular choice of \(\Phi\) and \(\Phi'\). This follows from
the relations between two faithful semifinite normal
\(\mathcal Z\)-traces (\hyperref[thm:III-4-2]{\S~4, Theorem~2}).

\noindent\textup{(iii)} Suppose that the existence of \(C\) and the
last assertion of the proposition have been established for
\(\mathcal A\) and \(\mathcal A'\). Let
\(P_{\mathfrak X}\) (respectively, \(P_{\mathfrak X'}\)) be a
projection of \(\mathcal A\) (respectively, of \(\mathcal A'\)) with
central support \(I\). Adopt the notation of
\hyperref[lem:III-2-1]{\S~2, Lemma~1}. The
induction
\[
 T'\longmapsto T'_{\mathfrak X}
\]
is an isomorphism \(\Theta\) of \(\mathcal A'\) onto
\(\mathcal B'\); in particular, this isomorphism allows
\(\mathcal Z\) to be identified with \(\mathcal Z_{\mathfrak X}\),
the centre of \(\mathcal B\) and \(\mathcal B'\). Transporting
\(\Phi'\) to \(\mathcal B'_+\) by \(\Theta\), one obtains a faithful
semifinite normal \(\mathcal Z\)-trace \(\Phi'_1\) on
\(\mathcal B'_+\). On the other hand, let \(\mathcal M\) be the set of
the \(T\in\mathcal A\) such that
\[
 TP_{\mathfrak X}=P_{\mathfrak X}T=T;
\]
the restriction of \(\Phi\) to \(\mathcal M_+\) defines a faithful
semifinite normal \(\mathcal Z\)-trace \(\Phi_1\) on
\(\mathcal B_+\). For \(x\in\mathfrak X\), one has
\[
 \mathfrak X_x^{\mathcal A'}
 =\mathfrak X_x^{\mathcal B'}\subset\mathfrak X,
 \qquad
 \mathfrak X_x^{\mathcal A}\cap\mathfrak X
 =P_{\mathfrak X}(\mathfrak X_x^{\mathcal A})
 =\mathfrak X_x^{\mathcal B};
\]
hence
\[
 \Phi(E_x^{\mathcal A'})=\Phi_1(E_x^{\mathcal B'}),
 \qquad
 \Phi'(E_x^{\mathcal A})=\Phi'_1(E_x^{\mathcal B}).
\]%
% Source PDF page 272 (printed p. 263)
\phantomsection\label{srcpage:263}%
The proposition is therefore true for \(\mathcal B\) and
\(\mathcal B'\). A further application of the same argument shows that
the proposition is also true for \(\mathcal W\) and \(\mathcal W'\).

\noindent\textup{(iv)} We shall show that every semifinite von Neumann
algebra is spatially isomorphic to a von Neumann algebra such that the
algebra and its commutant are finite, the algebra being, moreover,
standard. Together with the \hyperref[cor:III-6-1]{corollary to
Lemma~1}, this will complete the
proof. Every semifinite von Neumann algebra is spatially isomorphic to
a von Neumann algebra of the form
\[
 (\mathcal A_1\otimes\mathbb C_{\mathfrak K})_{E'_1},
\]
where \(\mathcal A_1\) is a standard von Neumann algebra on
\(\mathfrak H_1\), and
\(E'_1\in(\mathcal A_1\otimes\mathbb C_{\mathfrak K})'\)
(\hyperref[thm:I-4-3]{Chapter~I, \S~4, Theorem~3}, and
\hyperref[cor:I-6-9]{\S~6, the corollary to Proposition~9}). On the
Hilbert space
\(\mathfrak H_1\otimes(\mathfrak K\otimes\mathfrak K)\), consider the
von Neumann algebra
\[
 \mathcal A_1\otimes
 \bigl(\mathbb C_{\mathfrak K}\otimes
       \mathcal L(\mathfrak K)\bigr).
\]
It is the tensor product of two standard von Neumann algebras
(Chapter~I, \S~6, the corollary to Proposition~6), and is therefore
standard
[\hyperref[prop:I-5-9]{Chapter~I, \S~5, Proposition~9\textup{(ii)}}];
furthermore,
\(\mathcal A_1\otimes\mathbb C_{\mathfrak K}\) is spatially
isomorphic to
\[
 \bigl(\mathcal A_1\otimes\mathbb C_{\mathfrak K}
       \otimes\mathcal L(\mathfrak K)\bigr)_{E_1},
\]
where \(E_1\) is a suitably chosen projection of
\(\mathcal A_1\otimes\mathbb C_{\mathfrak K}
 \otimes\mathcal L(\mathfrak K)\).
\end{proof}

\begin{definition}\label{def:III-6-1}
Let \(\mathcal A\) be a von Neumann algebra such that \(\mathcal A\)
and \(\mathcal A'\) are finite. In
\hyperref[prop:III-6-1]{Proposition~1}, take for \(\Phi\)
and \(\Phi'\) the canonical \(\mathcal Z\)-traces of \(\mathcal A\)
and \(\mathcal A'\). The corresponding function in
\(\widehat{\mathcal Z}_+\) is called the \emph{coupling operator}
between \(\mathcal A\) and \(\mathcal A'\), and is denoted by
\(C_{\mathcal A}\).
\end{definition}

When \(\mathcal A\) is a factor, \(C_{\mathcal A}\) is identified
with a finite number \(>0\), the \emph{coupling number} between
\(\mathcal A\) and \(\mathcal A'\).

In English, \(C_{\mathcal A}\) is called sometimes ``coupling
operator,'' sometimes ``invariant.'' Roughly speaking,
\(C_{\mathcal A}\) measures the ``relative size'' of \(\mathcal A\)
and \(\mathcal A'\) (cf. \hyperref[ex:III-6-2]{Exercise~2}). One can
define
\(C_{\mathcal A}\) for arbitrary \(\mathcal A\).

Bibliography: \ArticleRef{20}, \ArticleRef{31}, \ArticleRef{40},
\ArticleRef{65}, \ArticleRef{66}, \ArticleRef{67}, \ArticleRef{89}.

\NumberedParagraph{2}{Properties of the Coupling Operator}
\label{no:III-6-2}

One plainly has
\[
 C_{\mathcal A'}=C_{\mathcal A}^{-1}.
\]

Let \(Z'\) be a locally compact subset of \(Z\), let \(\nu'\) be the
measure induced by \(\nu\) on \(Z'\), let \(G\) be the projection of
\(\mathcal Z\) associated with the characteristic function of \(Z'\),
and put \(\mathcal B=\mathcal A_G\). The centre of \(\mathcal A_G\) is
\(\mathcal Z_G\), which is identified with
\(L_{\mathbb C}^{\infty}(Z',\nu')\). One verifies at once that
\(C_{\mathcal B}\) is the restriction of \(C_{\mathcal A}\) to
\(Z'\).

Now let \(E'\) be a projection of \(\mathcal A'\) with central support
\(I\). Induction \(T\mapsto T_{E'}\) allows, in particular, the centre
of \(\mathcal A_{E'}\) to be identified with \(\mathcal Z\). Let
\(\Phi\) (respectively, \(\Phi'\)) be the canonical
\(\mathcal Z\)-trace of \(\mathcal A\) (respectively, of
\(\mathcal A'\)). Then:

\begin{proposition}\label{prop:III-6-2}
\(C_{\mathcal A_{E'}}\) is the product of \(C_{\mathcal A}\) and
\(\Phi'(E')\).
\end{proposition}

\begin{proof}
For every operator \(T'\in E'\mathcal A'E'\), there is a unique
function (up to locally negligible sets)\space%
% Source PDF page 273 (printed p. 264)
\phantomsection\label{srcpage:264}%
\(\Omega(T')\in L_{\mathbb C}^{\infty}(Z,\nu)\) such that
\[
 \Phi'(T')=\Omega(T')\Phi'(E').
\]
Uniqueness is immediate, because, by the faithfulness of \(\Phi'\),
the function \(\Phi'(E')\) vanishes only on a locally negligible set.
When \(T'\) is Hermitian, one has
\(-kE'\leq T'\leq kE'\) for some \(k>0\), and hence
\[
 -k\Phi'(E')\leq\Phi'(T')\leq k\Phi'(E');
\]
this proves the existence of \(\Omega(T')\) in this case, and linearity
proves it in the general case. The mapping
\(T'\mapsto\Omega(T')\) is a linear mapping of
\(E'\mathcal A'E'\) into \(\mathcal Z\), and therefore is identified
with a linear mapping of \(\mathcal A'_{E'}\) into
\(\mathcal Z_{E'}\). This mapping is plainly positive. One has
\[
 \Omega(T'_1T'_2)=\Omega(T'_2T'_1).
\]
If \(Q\in\mathcal Z\), then
\[
 \Phi'(E'QE')=Q\Phi'(E'),
 \qquad\text{hence}\qquad
 \Omega(E'QE')=Q.
\]
Thus (\hyperref[cor:III-5-1]{\S~5, the corollary to Theorem~1}),
\(\Omega\) is identified with
the canonical \(\mathcal Z_{E'}\)-trace \(\Psi'\) of
\(\mathcal A'_{E'}\).
For every \(x\in E'(\mathfrak H)\), the projections
\(E_x^{\mathcal A'}\) and \(E_x^{\mathcal A'_{E'}}\) correspond under
the isomorphism \(T\mapsto T_{E'}\) of \(\mathcal A\) onto
\(\mathcal A_{E'}\). Thus, denoting by \(\Psi\) the canonical
\(\mathcal Z_{E'}\)-trace of \(\mathcal A_{E'}\), one has
\[
 \Psi(E_x^{\mathcal A'_{E'}})=\Phi(E_x^{\mathcal A'}).
\]
On the other hand,
\[
 \Phi'(E_x^{\mathcal A})
 =\Omega(E_x^{\mathcal A})\Phi'(E')
 =\Psi'(E_x^{\mathcal A_{E'}})\Phi'(E'),
\]
by the first part of the proof. Consequently
\[
 \Psi(E_x^{\mathcal A'_{E'}})
 =C_{\mathcal A}\Phi'(E')
  \Psi'(E_x^{\mathcal A_{E'}}).
\]
\end{proof}

\begin{proposition}\label{prop:III-6-3}
Suppose that \(\mathcal A\) and \(\mathcal A'\) are finite and of
countable type. In order that \(\mathcal A\) have a separating
(respectively, totalizing) element, it is necessary and sufficient that
\(C_{\mathcal A}\geq1\) (respectively,
\(C_{\mathcal A}\leq1\)).
\end{proposition}

\begin{proof}
If there is an element \(x\) separating for \(\mathcal A\), then
\[
 E_x^{\mathcal A'}=I,
 \qquad E_x^{\mathcal A}\leq I,
 \qquad\text{hence}\qquad
 I=\Phi(E_x^{\mathcal A'})\geq\Phi'(E_x^{\mathcal A}),
\]
and consequently \(C_{\mathcal A}\geq1\). Conversely, suppose that
\(C_{\mathcal A}\geq1\), and let us show that \(\mathcal A\) has a
separating element. It is enough
(\hyperref[prop:I-2-3]{Chapter~I, \S~2, Proposition~3}) to
consider the case in which there is an element \(x\) separating for
\(\mathcal A'\). Then \(E_x^{\mathcal A}=I\), and hence
\(\Phi(E_x^{\mathcal A'})\geq I\), so that
\(E_x^{\mathcal A'}=I\): the element \(x\) is separating for
\(\mathcal A\). Interchanging the roles of \(\mathcal A\) and
\(\mathcal A'\), one sees that the existence of a totalizing element
for \(\mathcal A\) is equivalent to the condition
\(C_{\mathcal A}\leq1\).
\end{proof}

\begin{proposition}\label{prop:III-6-4}
Suppose that \(\mathcal A\) and \(\mathcal A'\) are finite. In order
that \(\mathcal A\) be standard, it is necessary and sufficient that
\(C_{\mathcal A}=1\).
\end{proposition}

% Source PDF page 274 (printed p. 265)
\phantomsection\label{srcpage:265}
\begin{proof}
The condition is necessary by the
\hyperref[cor:III-6-1]{corollary to Lemma~1}. Conversely,
suppose that \(C_{\mathcal A}=1\), and let us show that \(\mathcal A\)
is standard. Since a product of standard von Neumann algebras is
standard, it is enough to study the case in which \(\mathcal A\) and
\(\mathcal A'\) are of countable type. Then
\hyperref[prop:III-6-3]{Proposition~3} gives a
separating element for \(\mathcal A\) and a separating element for
\(\mathcal A'\). Therefore \(\mathcal A\) is standard
(\hyperref[thm:III-1-5]{\S~1, Theorem~5}).
\end{proof}

Bibliography: \ArticleRef{20}, \ArticleRef{31}, \ArticleRef{40},
\ArticleRef{65}, \ArticleRef{66}, \ArticleRef{67}, \ArticleRef{89}.

\NumberedParagraph{3}{Applications: I. Comparison of the Strong and
Ultra-Strong Topologies, and of the Weak and Ultra-Weak Topologies}
\label{no:III-6-3}

\begin{proposition}\label{prop:III-6-5}
Let \(\mathcal A\) be a von Neumann algebra on \(\mathfrak H\), and
let \(n\) be an integer \(>0\). The following properties are
equivalent:
\begin{enumerate}[label=\textup{(\roman*)},leftmargin=*]
\item Every ultra-weakly continuous linear functional on \(\mathcal A\)
is a sum of \(n\) functionals \(\omega_{x,y}\).
\item Every normal positive linear functional on \(\mathcal A\) is a
sum of \(n\) functionals \(\omega_x\).
\item The subsets of \(\mathcal A\) defined by a condition of the form
\[
 \sum_{i=1}^{n}\lVert Tx_i\rVert^2\leq1
\]
form a fundamental system of ultra-strong neighborhoods of \(0\) in
\(\mathcal A\).
\item Every projection \(E\in\mathcal A\) such that \(\mathcal A_E\)
is of countable type is the supremum of \(n\) cyclic projections.
\end{enumerate}
\end{proposition}

\begin{proof}
Let \(\mathfrak K\) be a Hilbert space of dimension \(n\). To say that
\(\mathcal A\) and the integer \(n\) satisfy condition~\textup{(i)}
[respectively, \textup{(ii)}, \textup{(iii)}, \textup{(iv)}] amounts,
as is seen at once, to saying that, on
\(\mathfrak H\otimes\mathfrak K\),
\(\mathcal A\otimes\mathbb C_{\mathfrak K}\) and the integer \(1\)
satisfy \textup{(i)} [respectively, \textup{(ii)}, \textup{(iii)},
\textup{(iv)}]. It is therefore enough to study the case \(n=1\).

\noindent\textup{(i)\(\Rightarrow\)(ii).} Suppose that
condition~\textup{(i)} holds. Let \(\varphi\) be a normal positive
linear functional. Since \(\varphi\) is ultra-weakly continuous, it is
a functional \(\omega_{x,y}\), and therefore
(\hyperref[lem:I-4-2]{Chapter~I, \S~4, Lemma~2}) is a functional
\(\omega_z\).

\noindent\textup{(ii)\(\Rightarrow\)(iii).} For every sequence
\((y_i)\) of elements of \(\mathfrak H\) such that
\[
 \sum_{i=1}^{\infty}\lVert y_i\rVert^2<+\infty,
\]
let \(\mathcal V_{(y_i)}\) be the ultra-strong neighborhood of \(0\) in
\(\mathcal A\) defined by the condition
\[
 \sum_{i=1}^{\infty}\lVert Ty_i\rVert^2\leq1.
\]
One has
\[
 \sum_{i=1}^{\infty}\lVert Ty_i\rVert^2=\varphi(T^*T),
\]
where \(\varphi\) is the normal positive linear functional
\(\sum_{i=1}^{\infty}\omega_{y_i}\). If condition~\textup{(ii)} holds,
then \(\varphi=\omega_x\), and hence
\(\mathcal V_{(y_i)}=\mathcal V_{(x)}\), which proves
condition~\textup{(iii)}.

% Source PDF page 275 (printed p. 266)
\phantomsection\label{srcpage:266}
\noindent\textup{(iii)\(\Rightarrow\)(i).} Let \(\varphi\) be an
ultra-weakly continuous linear functional on \(\mathcal A\). If
condition~\textup{(iii)} holds, there is an \(x\in\mathfrak H\) such
that \(\lVert Tx\rVert\leq1\) implies
\(\lvert\varphi(T)\rvert\leq1\). Thus \(\varphi(T)\) is a continuous
linear functional of \(Tx\), so there is a \(y\in\mathfrak H\) such
that
\[
 \varphi(T)=(Tx\mid y).
\]

The equivalence \textup{(ii)}\(\Longleftrightarrow\)\textup{(iv)} is
the \hyperref[cor:III-1-thm4]{corollary to Theorem~4 of \S~1}.
\end{proof}

\par\addvspace{6pt}\noindent\textbf{Remark 1.}\ ---\enspace
The conditions of \hyperref[prop:III-6-5]{Proposition~5} are satisfied,
with \(n=1\), when \(\mathcal A\) is abelian and when \(\mathcal A\)
is standard
(\hyperref[cor:III-1-thm4]{\S~1, the corollary to Theorem~4}). Another
very extensive case will be seen in
\hyperref[cor:III-8-thm1-10]{\S~8, Theorem~1, Corollary~10}.

\begin{proposition}\label{prop:III-6-6}
Suppose that \(\mathcal A\) and \(\mathcal A'\) are finite. In order
that the conditions of \hyperref[prop:III-6-5]{Proposition~5} be
satisfied, it is necessary and
sufficient that
\[
 C_{\mathcal A}\geq\frac1n.
\]
\end{proposition}

\begin{proof}
Let \(\mathfrak K\) be a Hilbert space of dimension \(n\). Put
\[
 \mathcal B=\mathcal A\otimes\mathbb C_{\mathfrak K}
\]
on \(\mathfrak H\otimes\mathfrak K\). One has
\(C_{\mathcal B}=nC_{\mathcal A}\)
(\hyperref[prop:III-6-2]{Proposition~2}). As in the proof of
\hyperref[prop:III-6-5]{Proposition~5}, we are thus reduced to the case
\(n=1\). It is then enough to apply
\hyperref[prop:III-6-3]{Proposition~3}.
\end{proof}

\begin{proposition}\label{prop:III-6-7}
Let \(\mathcal A\) be a von Neumann algebra. The following properties
are equivalent:
\begin{enumerate}[label=\textup{(\roman*)},leftmargin=*]
\item Every ultra-weakly continuous linear functional on \(\mathcal A\)
is a sum of finitely many functionals \(\omega_{x,y}\).
\item Every normal positive linear functional on \(\mathcal A\) is a
sum of finitely many functionals \(\omega_x\).
\item The weak and ultra-weak topologies coincide on \(\mathcal A\).
\item The strong and ultra-strong topologies coincide on \(\mathcal A\).
\end{enumerate}
\end{proposition}

\begin{proof}
The equivalence \textup{(i)}\(\Longleftrightarrow\)\textup{(iii)}
follows from the definition of the weak and ultra-weak topologies
(\hyperref[no:I-3-1]{Chapter~I, \S~3, no.~1}). Since the weak topology
(respectively, the
ultra-weak topology) is the weakened topology of the strong topology
(respectively, of the ultra-strong topology)
(\hyperref[thm:I-3-1]{Chapter~I, \S~3, Theorem~1}), one has
\textup{(iv)}\(\Rightarrow\)\textup{(iii)}. The implication
\textup{(iii)}\(\Rightarrow\)\textup{(iv)} follows from the relations
between the weak and strong topologies (respectively, the ultra-weak and
ultra-strong topologies) given in
\hyperref[no:I-3-2]{Chapter~I, \S~3, no.~2}. Since every
ultra-weakly continuous linear functional is a linear combination of
normal positive linear functionals, one has
\textup{(ii)}\(\Rightarrow\)\textup{(i)}. Finally, the implication
\textup{(i)}\(\Rightarrow\)\textup{(ii)} follows from
\hyperref[lem:I-4-2]{Lemma~2} and \hyperref[lem:I-4-6]{Lemma~6} of
Chapter~I, \S~4.
\end{proof}

\par\addvspace{6pt}\noindent\textbf{Remark 2.}\ ---\enspace
Let \(\mathcal A_1,\mathcal A_2,\ldots,\mathcal A_n\) be von Neumann
algebras, and put
\[
 \mathcal A=\mathcal A_1\times\mathcal A_2\times\cdots\times\mathcal A_n.
\]
% Source PDF page 276 (printed p. 267)
\phantomsection\label{srcpage:267}
In order that the strong and ultra-strong topologies coincide on
\(\mathcal A\), it is necessary and sufficient that they coincide on
\(\mathcal A_1,\mathcal A_2,\ldots,\mathcal A_n\). Thus, in studying
whether the conditions of \hyperref[prop:III-6-7]{Proposition~7} hold,
it is enough to consider
the following cases: \(1^\circ\) \(\mathcal A'\) is properly infinite;
\(2^\circ\) \(\mathcal A'\) is finite and \(\mathcal A\) is properly
infinite; \(3^\circ\) \(\mathcal A'\) and \(\mathcal A\) are finite.
The first case will be studied in \hyperref[sec:III-8]{\S~8}. The other
two cases will now be
studied.

\begin{lemma}\label{lem:III-6-2}
If the strong and ultra-strong topologies coincide on \(\mathcal A\),
every projection of \(\mathcal A\) which is the supremum of a countable
family of cyclic projections is majorized by the supremum of a finite
family of cyclic projections.
\end{lemma}

\begin{proof}
Let \(E\) be the supremum of the projections
\(E_{x_1}^{\mathcal A'},E_{x_2}^{\mathcal A'},\ldots\). We may suppose
that
\[
 \sum_{r=1}^{\infty}\lVert x_r\rVert^2<+\infty.
\]
Let \(\mathcal V_{(x_r)}\) be the ultra-strong neighborhood of \(0\)
in \(\mathcal A\) formed by the \(T\in\mathcal A\) such that
\[
 \sum_{r=1}^{\infty}\lVert Tx_r\rVert^2\leq1.
\]
There are vectors \(y_1,y_2,\ldots,y_n\) such that
\(\mathcal V_{(y_i)}\subset\mathcal V_{(x_r)}\). Let \(F\) be the
supremum of
\(E_{y_1}^{\mathcal A'},E_{y_2}^{\mathcal A'},\ldots,
E_{y_n}^{\mathcal A'}\). For every number \(\alpha>0\), one has
\[
 \alpha(I-F)\in\mathcal V_{(y_i)}\subset\mathcal V_{(x_r)}.
\]
Thus \((I-F)x_r=0\) for every \(r\). Consequently
\(I-F\leq I-E\), and hence \(F\geq E\).
\end{proof}

\begin{proposition}\label{prop:III-6-8}
If \(\mathcal A'\) is finite and \(\mathcal A\) is properly infinite,
the ultra-strong topology on \(\mathcal A\) is strictly finer than the
strong topology.
\end{proposition}

\begin{proof}
For every \(x\in\mathfrak H\), \(E_x^{\mathcal A}\) is finite, and
hence \(E_x^{\mathcal A'}\) is finite
(\hyperref[prop:III-2-3]{\S~2, Proposition~3}). Therefore
the supremum of a finite family of cyclic projections of \(\mathcal A\)
is finite (\S~3, Proposition~5). But there is an infinite family (which
may be assumed countable) of nonzero, pairwise orthogonal, equivalent
projections of \(\mathcal A\) [such a family is furnished by
\hyperref[cor:III-1-thm1-2]{\S~1, Corollary~2 of Theorem~1}, applied to
a family \((E_i)\) consisting of a
single finite projection]. By decreasing these projections, they may
be assumed cyclic; their supremum is infinite
(\hyperref[prop:III-2-10]{\S~2, Proposition~10}). The proposition now
follows from \hyperref[lem:III-6-2]{Lemma~2}.
\end{proof}

\begin{proposition}\label{prop:III-6-9}
Suppose that \(\mathcal A\) and \(\mathcal A'\) are finite. In order
that the strong and ultra-strong topologies coincide on \(\mathcal A\),
it is necessary and sufficient that \(C_{\mathcal A}^{-1}\) be
essentially bounded.
\end{proposition}

\begin{proof}
If \(C_{\mathcal A}^{-1}\) is essentially bounded, the strong and
ultra-strong topologies coincide on \(\mathcal A\)
(\hyperref[prop:III-6-6]{Proposition~6}). If
\(C_{\mathcal A}^{-1}\) is not essentially bounded, then, for every
integer \(p>0\), there is a nonzero projection \(E_p\) of the centre
\(\mathcal Z\) of \(\mathcal A\) such that
\[
 C_{\mathcal A_{E_p}}\leq\frac1p.
\]
By decreasing the \(E_p\), we may suppose that\space%
% Source PDF page 277 (printed p. 268)
\phantomsection\label{srcpage:268}%
the \(\mathcal Z_{E_p}\) are of countable type. Let \(E\) be the
supremum of the \(E_p\). The projection \(E\) is the supremum of a
countable family of cyclic projections. If the strong and ultra-strong
topologies coincide on \(\mathcal A\), then \(E\) is the supremum of a
finite number \(q\) of cyclic projections
(\hyperref[lem:III-6-2]{Lemma~2}); every projection
of \(\mathcal A\) majorized by \(E\) is then the supremum of \(q\)
cyclic projections, and therefore
(\hyperref[prop:III-6-6]{Proposition~6})
\[
 C_{\mathcal A_{E_p}}\geq\frac1q,
\]
which is a contradiction.
\end{proof}

The converse of \hyperref[lem:III-6-2]{Lemma~2} is true
(\S~8, Exercise~7).

Bibliography: \ArticleRef{15}, \ArticleRef{19}, \ArticleRef{31},
\ArticleRef{40}, \ArticleRef{70}, \ArticleRef{89}, \ArticleRef{117}.

\NumberedParagraph{4}{Applications: II. Conditions for an Isomorphism
to Be Spatial}
\label{no:III-6-4}

Let \(\mathcal A\) be a von Neumann algebra, and let \(\mathcal Z\) be
its centre. We identify
\[
 \mathcal Z=L_{\mathbb C}^{\infty}(Z,\nu).
\]
Let \(\mathcal B\) be another von Neumann algebra, and let
\(\mathcal Y\) be its centre. If \(\Omega\) is an isomorphism of
\(\mathcal A\) onto
\(\mathcal B\), then \(\Omega\) defines an isomorphism of
\(\mathcal Z\) onto \(\mathcal Y\), and so allows \(\mathcal Y\) to be
identified with \(L_{\mathbb C}^{\infty}(Z,\nu)\). We make this
convention in the following statement.

\begin{proposition}\label{prop:III-6-10}
Suppose that \(\mathcal A,\mathcal A',\mathcal B,\mathcal B'\) are
finite. Let \(\Omega\) be an isomorphism of \(\mathcal A\) onto
\(\mathcal B\). If \(C_{\mathcal A}=C_{\mathcal B}\), then
\(\Omega\) is spatial.
\end{proposition}

\begin{proof}
There are a von Neumann algebra \(\mathcal C\) and projections
\(E',F'\in\mathcal C'\), with central support \(I\), such that
\(\mathcal A\) may be identified with \(\mathcal C_{E'}\),
\(\mathcal B\) with \(\mathcal C_{F'}\), and \(\Omega\) with the
mapping
\[
 T_{E'}\longmapsto T_{F'},
 \qquad T\in\mathcal C
\]
(\hyperref[cor:I-4-3]{Chapter~I, \S~4, the corollary to Theorem~3}).
Let \(G'\) be the
supremum of \(E'\) and \(F'\), which is a projection of
\(\mathcal C'\). Since
\(\mathcal A'=\mathcal C'_{E'}\) and
\(\mathcal B'=\mathcal C'_{F'}\) are finite, \(E'\) and \(F'\) are
finite projections of \(\mathcal C'\), and hence \(G'\) is finite.
Replacing \(\mathcal C\) by \(\mathcal C_{G'}\), one sees that
\(\mathcal C'\) may be assumed finite. On the other hand,
\(\mathcal C\), being isomorphic to \(\mathcal A\) and
\(\mathcal B\), is finite. Let \(\Phi'\) be the canonical
\(\mathcal Z\)-trace of \(\mathcal C'\). By
\hyperref[prop:III-6-2]{Proposition~2},
\[
 C_{\mathcal C}\Phi'(E')=C_{\mathcal A},
 \qquad
 C_{\mathcal C}\Phi'(F')=C_{\mathcal B}.
\]
If \(C_{\mathcal A}=C_{\mathcal B}\), it follows that
\(\Phi'(E')=\Phi'(F')\), and hence \(E'\sim F'\) (by the comparison
theorem for projections and the faithfulness of \(\Phi'\)). Therefore
\(\Omega\) is spatial.
\end{proof}

\begin{corollary}\label{cor:III-6-2}
Let \(\mathcal A\) be a finite von Neumann algebra, let \(\mathcal Z\)
be its centre, and let \(\Omega\) be an automorphism of \(\mathcal A\)
which fixes the elements of \(\mathcal Z\). Then \(\Omega\) is spatial.
\end{corollary}

\begin{proof}
There is a family \((E'_i)\) of pairwise orthogonal finite projections
of \(\mathcal A'\), with sum \(I\)
(\hyperref[cor:III-2-prop7-1]{\S~2, Proposition~7, Corollary~1}). Let
\(F_i\) be the central support of \(E'_i\). Then
\(\Omega\) induces an automorphism \(\Psi_i\) of
\(\mathcal A_{F_i}\), and induction
\(T_{F_i}\mapsto T_{E'_i}\) transforms \(\Psi_i\) into an
automorphism \(\Phi_i\) of \(\mathcal A_{E'_i}\). Since\space%
% Source PDF page 278 (printed p. 269)
\phantomsection\label{srcpage:269}%
\(\mathcal A_{E'_i}\) and \(\mathcal A'_{E'_i}\) are finite, and since
\(\Phi_i\) fixes the elements of the centre of \(\mathcal A_{E'_i}\),
\(\Phi_i\) is defined by a unitary operator \(U_i\) on
\(E'_i(\mathfrak H)\) (\hyperref[prop:III-6-10]{Proposition~10}). Let
\(U\) be the unitary
operator on \(\mathfrak H\) which, for every \(i\), coincides with
\(U_i\) on \(E'_i(\mathfrak H)\). Plainly, \(U\) defines the
automorphism \(\Omega\).
\end{proof}

Bibliography: \ArticleRef{20}, \ArticleRef{31}, \ArticleRef{66},
\ArticleRef{67}, \ArticleRef{89}.

\medskip
\noindent\textit{Exercises.}
\begin{enumerate}[label=\arabic*.,leftmargin=*]
\item\label{ex:III-6-1}
Let \(\mathfrak A\) be a Hilbert algebra, let \(\mathfrak H\) be the
completed Hilbert space of \(\mathfrak A\), put
\(\mathcal Z=\mathcal U(\mathfrak A)\cap\mathcal V(\mathfrak A)\),
let \(J\) be the involution of \(\mathfrak H\) canonically defined by
\(\mathfrak A\), let \(\Phi\) be a normal \(\mathcal Z\)-trace on
\(\mathcal U(\mathfrak A)_+\), and let \(\Phi'\) be the normal
\(\mathcal Z\)-trace on \(\mathcal V(\mathfrak A)_+\) obtained by
transporting \(\Phi\) by \(J\).
\begin{enumerate}[label=\textit{\alph*.},leftmargin=2em]
\item If \(x\in\mathfrak H\) is bounded,
\(E_x^{\mathcal U(\mathfrak A)}\) is the support of \(V_x^*\), and
\(JE_x^{\mathcal V(\mathfrak A)}J\) is the support of
\(JU_x^*J=V_x\); therefore
\[
 \Phi(E_x^{\mathcal V(\mathfrak A)})
 =\Phi'(E_x^{\mathcal U(\mathfrak A)}).
\]

\item If \(x\) is arbitrary in \(\mathfrak H\), let \((x_n)\) be a
sequence of bounded elements converging to \(x\). Put
\[
 E=E_x^{\mathcal V(\mathfrak A)},\qquad
 E'=E_x^{\mathcal U(\mathfrak A)},\qquad
 E_n=E_{x_n}^{\mathcal V(\mathfrak A)},\qquad
 E'_n=E_{x_n}^{\mathcal U(\mathfrak A)}.
\]
Replacing \(x_n\) by \(EE'x_n\), we may suppose that \(E_n\)
(respectively, \(E'_n\)) converges strongly to \(E\) (respectively,
to \(E'\)), with \(E_n\leq E\) (respectively, \(E'_n\leq E'\)). For
every trace \(\omega\) on \(\widehat{\mathcal Z}_+\), one has
\[
 \omega(\Phi(E_n))\longrightarrow\omega(\Phi(E)),
 \qquad
 \omega(\Phi'(E'_n))\longrightarrow\omega(\Phi'(E')).
\]
(Use Chapter~I, \S~6, the corollary to Proposition~2.) With the aid of
\textit{a}, conclude that
\(\omega(\Phi(E))=\omega(\Phi'(E'))\). This gives another proof of the
\hyperref[cor:III-6-1]{corollary to Lemma~1}.

\item For arbitrary \(x\in\mathfrak H\), prove directly that
\[
 \Phi(E_x^{\mathcal V(\mathfrak A)})
 =\Phi'(E_x^{\mathcal U(\mathfrak A)})
\]
by arguing as in \textit{a}, but also using
\hyperref[ex:I-1-10]{Chapter~I, \S~1, Exercise~10}, and
\hyperref[ex:I-5-6]{\S~5, Exercise~6}.
\end{enumerate}

\item\label{ex:III-6-2}
Let \(\mathcal A\) be a von Neumann algebra. Suppose that
\(\mathcal A\) and \(\mathcal A'\) are homogeneous and finite, of
multiplicities \(n\) and \(n'\), respectively. Show that the coupling
operator is constant and equal to \(n(n')^{-1}\).

\item\label{ex:III-6-3}
Let \(\mathcal A\) be a von Neumann algebra, and let \(\mathcal Z\) be
its centre.
\begin{enumerate}[label=\textit{\alph*.},leftmargin=2em]
\item Let \(x\) be a trace element for \(\mathcal A\). If
\(T\in\mathcal Z\), then \(Tx\) is a trace element for
\(\mathcal A\). Every element of \(\mathfrak X_x^{\mathcal Z}\) is a
trace element for \(\mathcal A\).

\item Put \(E=E_x^{\mathcal Z}\in\mathcal Z'\). The mapping
\[
 T\longmapsto ETE
\]
of \(\mathcal A\) into \(\mathcal Z'_E\) is linear and positive, and
\[
 ESTE=ETSE\qquad(S,T\in\mathcal A).
\]

\item Suppose, in addition, that \(x\) is separating for
\(\mathcal Z\). Then, for every \(T\in\mathcal A\), there is a unique
\(T'\in\mathcal Z\) such that
\[
 ETE=T'E.
\]
(Observe that \(\mathcal Z_E\) has a totalizing element, and hence that
\(\mathcal Z_E=\mathcal Z'_E\).) Show that \(T\mapsto T'\) is the
canonical \(\mathcal Z\)-trace of \(\mathcal A\), which is finite.

\item Suppose that \(\mathcal A\) and \(\mathcal A'\) are finite and
that \(C_{\mathcal A}\geq1\). There is a projection
\(F\in\mathcal Z'\) having the following property: for every
\(T\in\mathcal A\), there is a unique \(T'\in\mathcal Z\) such that
\[
 FTF=T'F,
\]
and \(T\mapsto T'\) is the canonical \(\mathcal Z\)-trace of
\(\mathcal A\). (Reduce to the case in which \(\mathcal A\) is of
countable type. Then use \hyperref[prop:III-6-3]{Proposition~3} and
\textit{c}.)
\ArticleRef{66}, \ArticleRef{89}
\end{enumerate}
\end{enumerate}

\DixmierDeferredSectionSeven
\global\let\DixmierDeferredSectionSeven\relax
% Source PDF page 283 (printed p. 274)
\phantomsection\label{srcpage:274}
On the other hand, \(\|TF-F\|\leq\varepsilon\) and
\(\|TF_1\|\leq\varepsilon\); hence
\[
  \|T-F\|_2
  \leq \|TF-F\|_2+\|TF_1\|_2+\|TF_2\|_2
  \leq \varepsilon+\varepsilon+\frac{6}{\varepsilon}\|T-E\|_2.
\]
If \(\|T-E\|_2^{1/2}\geq\tfrac12\), the inequalities of the lemma
are plainly satisfied. If \(\|T-E\|_2^{1/2}<\tfrac12\), we may take
\(\varepsilon=\|T-E\|_2^{1/2}\). We then obtain
\[
  \|T-F\|_2\leq 8\|T-E\|_2^{1/2}.
\]
Consequently,
\[
  \|F-E\|_2
  \leq \|T-E\|_2+\|T-F\|_2
  \leq \|T-E\|_2+8\|T-E\|_2^{1/2}
  \leq 9\|T-E\|_2^{1/2},
\]
which proves (i).

Now suppose that \(T\geq0\). With the same notation as above, we have
\[
  \|T^{1/2}F-F\|\leq\varepsilon,
  \qquad
  \|T^{1/2}F_1\|\leq\varepsilon^{1/2},
\]
and therefore
\[
  \|T^{1/2}-F\|_2
  \leq \|T^{1/2}F-F\|_2+\|T^{1/2}F_1\|_2
       +\|T^{1/2}F_2\|_2
  \leq \varepsilon+\varepsilon^{1/2}
       +\frac{6}{\varepsilon}\|T-E\|_2.
\]
It follows that
\[
  \|T^{1/2}-E\|_2
  \leq \|T^{1/2}-F\|_2+\|F-E\|_2
  \leq \varepsilon+\varepsilon^{1/2}
       +\frac{6}{\varepsilon}\|T-E\|_2
       +9\|T-E\|_2^{1/2}.
\]
Taking \(\varepsilon=\|T-E\|_2^{1/2}\), if
\(\|T-E\|_2^{1/2}<\tfrac12\), gives
\[
  \|T^{1/2}-E\|_2\leq13\|T-E\|_2^{1/4};
\]
if \(\|T-E\|_2^{1/2}\geq\tfrac12\), the inequality is immediate.
\end{proof}

\begin{lemma}\label{lem:III-7-5}
Let \(\mathcal A\) be a finite factor and let \(E,F\) be projections
of \(\mathcal A\). There is a partial isometry \(W\in\mathcal A\)
whose initial projection is dominated by \(E\), whose final projection
is dominated by \(F\), and such that
\[
  \|W-E\|_2\leq14\|F-E\|_2^{1/4}.
\]
\end{lemma}

\begin{proof}
Let \(FE=WB\) be the polar decomposition of \(FE\). We have
\(0\leq B\leq I\). The operator \(W\) is partially isometric and its
final projection is dominated by \(F\); since \((FE)^*=EF\), its
initial projection is dominated by \(E\). We have
\[
  B^2=(FE)^*(FE)=EFE,
  \qquad
  \|B^2-E\|_2=\|E(F-E)E\|_2\leq\|F-E\|_2.
\]
By Lemma~4(ii),
\[
  \|B-E\|_2\leq13\|F-E\|_2^{1/4}.
\]%
% Source PDF page 284 (printed p. 275)
\phantomsection\label{srcpage:275}%
Hence
\begin{align*}
  \|W-E\|_2
  &\leq \|W-FE\|_2+\|FE-E\|_2 \\
  &=\|W(E-B)\|_2+\|(F-E)E\|_2 \\
  &\leq \|E-B\|_2+\|F-E\|_2 \\
  &\leq13\|F-E\|_2^{1/4}+\|F-E\|_2
   \leq14\|F-E\|_2^{1/4}.
\end{align*}
\end{proof}

\begin{lemma}\label{lem:III-7-6}
Let \(\mathcal A\) be a finite factor, let \(U\in\mathcal A\) be
unitary, and let \(W\in\mathcal A\) be a partial isometry that agrees
with \(U\) on its initial subspace. Then
\[
  \|U-W\|_2^2\leq2\|W-I\|_2.
\]
\end{lemma}

\begin{proof}
The initial subspace of \(W\) is the final subspace of \(W^*\).
Thus \(UW^*=WW^*\), and, on taking adjoints, \(WU^*=WW^*\).
Consequently,
\[
  (U-W)(U-W)^*=UU^*-UW^*-WU^*+WW^*=I-WW^*.
\]
If \(\varphi\) denotes the canonical trace on \(\mathcal A\), it
follows that
\begin{align*}
  \|U-W\|_2^2
  &=\varphi(I-WW^*) \\
  &\leq |\varphi(I-W)|+|\varphi(W-WW^*)| \\
  &\leq \|W-I\|_2+\|(I-W)W^*\|_2
   \leq2\|W-I\|_2.
\end{align*}
\end{proof}

\begin{lemma}\label{lem:III-7-7}
Let \(\mathcal A\) be a finite factor, and let \(E,F\) be equivalent
projections of \(\mathcal A\). There is a unitary operator
\(U\in\mathcal A\) such that
\[
  F=UEU^{-1},
  \qquad
  \|U-I\|_2\leq36\|F-E\|_2^{1/8}.
\]
\end{lemma}

\begin{proof}
Let \(W\) be the partial isometry whose existence is asserted by
Lemma~5, and let \(E'\leq E\) and \(F'\leq F\) be its initial and
final projections. Since \(E'\sim F'\),
\hyperref[prop:III-2-6]{Proposition~6 of \S~2} gives
\(E-E'\sim F-F'\). Let \(V\in\mathcal A\) be a partial isometry with
initial projection \(E-E'\) and final projection \(F-F'\). By arguing
similarly with \(I-E\) and \(I-F\), we construct partial isometries
\(W_1,V_1\in\mathcal A\); the sum
\(W+V+W_1+V_1\) is a unitary operator \(U\in\mathcal A\) such that
\(F=UEU^{-1}\). Now
\[
  \|W-E\|_2\leq14\|F-E\|_2^{1/4},
\]
and likewise
\[
  \|W_1-(I-E)\|_2
  \leq14\|(I-F)-(I-E)\|_2^{1/4}
  =14\|F-E\|_2^{1/4}.
\]
Thus
\[
  \|W+W_1-I\|_2\leq28\|F-E\|_2^{1/4}.
\]
Since \(U\) agrees with the partial isometry \(W+W_1\) on the
initial subspace of the latter, Lemma~6 gives
\[
  \|U-(W+W_1)\|_2
  \leq2^{1/2}\|W+W_1-I\|_2^{1/2}
  \leq8\|F-E\|_2^{1/8}.
\]%
% Source PDF page 285 (printed p. 276)
\phantomsection\label{srcpage:276}%
Therefore
\[
  \|U-I\|_2\leq36\|F-E\|_2^{1/8}.
\]
\end{proof}

Bibliography: \ArticleRef{67}.

\NumberedParagraph{4}{A New Definition of Hyperfinite Factors}
\label{no:III-7-4}
In Lemmas~8, 9, 10, and~11, \(\mathcal A\) denotes a continuous
finite factor having the following property:
\begin{equation*}
\tag{\(\star\)}\label{eq:III-7-star}
\begin{minipage}{0.88\linewidth}
For arbitrary \(T_1,T_2,\ldots,T_m\in\mathcal A\) and
\(\varepsilon>0\), there are a finite-dimensional involutive
subalgebra \(\mathcal B\) of \(\mathcal A\), over \(\mathbb C\), and
elements \(S_1,S_2,\ldots,S_m\in\mathcal B\) such that
\(\|S_h-T_h\|_2\leq\varepsilon\) for \(h=1,2,\ldots,m\).
\end{minipage}
\end{equation*}
We continue to denote the canonical trace on \(\mathcal A\) by
\(\varphi\).

\begin{lemma}\label{lem:III-7-8}
Given \(T_1,T_2,\ldots,T_m\in\mathcal A\) and
\(\varepsilon_1>0\), there are a factor
\(\mathcal B\subset\mathcal A\) of type \(\mathrm{I}_{2^n}\) and elements
\(S_1,S_2,\ldots,S_m\in\mathcal B\) such that
\[
  \|S_h-T_h\|_2\leq\varepsilon_1
  \qquad(h=1,2,\ldots,m).
\]
\end{lemma}

\begin{proof}
Let \(\mathcal C\) be an involutive subalgebra of finite dimension in
\(\mathcal A\), and let
\[
  R_1,R_2,\ldots,R_m\in\mathcal C,
\]
whose existence is asserted by \eqref{eq:III-7-star}. Let
\(E_1,E_2,\ldots,E_p\) be the mutually orthogonal minimal projections
of the centre of \(\mathcal C\), with sum \(I\) (adjoining \(I\) to
\(\mathcal C\), if necessary). Each \(\mathcal C_{E_i}\) is a
factor of finite dimension. Let
\[
  E_i^1,E_i^2,\ldots,E_i^{p_i}
\]
be mutually orthogonal minimal
projections of \(\mathcal C E_i\), with sum \(E_i\), and let
\(W_i^j\in\mathcal C E_i\) be a partial isometry with initial
projection \(E_i^1\) and final projection \(E_i^j\). Each \(R_h\) is
a linear combination of the operators
\(W_i^kW_i^{j*}\), where \(j,k=1,2,\ldots,p_i\) and
\(i=1,2,\ldots,p\). It is therefore enough to prove the existence of a
factor \(\mathcal B\subset\mathcal A\) of type \(\mathrm{I}_{2^n}\) containing
elements \(V_i^j\), with \(\|V_i^j\|\leq1\), arbitrarily close to
the \(W_i^j\) in the norm \(\|\cdot\|_2\). Indeed,
\begin{align*}
 \|W_i^kW_i^{j*}-V_i^kV_i^{j*}\|_2
 &\leq \|W_i^k(W_i^{j*}-V_i^{j*})\|_2
       +\|(W_i^k-V_i^k)V_i^{j*}\|_2 \\
 &\leq \|W_i^k\|\,\|W_i^{j*}-V_i^{j*}\|_2
       +\|V_i^{j*}\|\,\|W_i^k-V_i^k\|_2 \\
 &\leq \|W_i^j-V_i^j\|_2+\|W_i^k-V_i^k\|_2.
\end{align*}
Choose mutually orthogonal equivalent projections
\(F_i^1,F_i^2,\ldots,F_i^{q_i}\leq E_i^1\) such that
\[
  \varphi(F_i^1)=\varphi(F_i^2)=\cdots=
  \varphi(F_i^{q_i})=2^{-n}.
\]
For the moment keep \(n\) fixed and choose \(q_i\) as large as
possible, so that
\[
  \varphi\!\left(E_i^1-
    (F_i^1+F_i^2+\cdots+F_i^{q_i})\right)\leq2^{-n}.
\]%
% Source PDF page 286 (printed p. 277)
\phantomsection\label{srcpage:277}%
Transporting \(F_i^1,F_i^2,\ldots,F_i^{q_i}\) by \(W_i^j\) gives
equivalent projections, dominated by \(E_i^j\), that are mutually
orthogonal. Carrying this construction out for
\(j=1,2,\ldots,p_i\) gives what we shall call the projections \(F\)
of index \(i\). As \(i\) ranges from \(1\) to \(p\), we have obtained
a number of mutually orthogonal projections \(F\), each of trace
\(2^{-n}\). Add arbitrarily enough further projections \(F\) in
\(\mathcal A\), also of trace \(2^{-n}\), so that all the projections
\(F\) are mutually orthogonal and have sum \(I\). We may regard them
as the minimal projections of a factor
\(\mathcal B\subset\mathcal A\) of type \(\mathrm{I}_{2^n}\). Moreover,
\(\mathcal B\) may be chosen so that every partial isometry in
\(\mathcal B\) whose initial projection is one of
\(F_i^1,F_i^2,\ldots,F_i^{q_i}\) and whose final projection is one of
the projections \(F\) of index \(i\), agrees on its initial subspace
with one of the \(W_i^j\). Thus there are partial isometries
\(V_i^j\in\mathcal B\) such that \(W_i^j-V_i^j\) is a partial
isometry with initial projection
\[
  E_i^1-(F_i^1+F_i^2+\cdots+F_i^{q_i}).
\]
Consequently,
\[
  \|W_i^j-V_i^j\|_2^2
  =\varphi\!\left(E_i^1-
    (F_i^1+F_i^2+\cdots+F_i^{q_i})\right)
  \leq2^{-n}.
\]
Taking \(n\) arbitrarily large therefore makes
\(\|W_i^j-V_i^j\|_2\) arbitrarily small.
\end{proof}

\begin{lemma}\label{lem:III-7-9}
Given \(T_1,T_2,\ldots,T_m\in\mathcal A\), a projection
\(E\in\mathcal A\) such that \(\varphi(E)=2^{-n}\), and
\(\varepsilon_2>0\), there are a factor
\(\mathcal B\subset\mathcal A\) of type \(\mathrm{I}_{2^p}\), with
\(p\geq n\), elements \(S_1,S_2,\ldots,S_m\in\mathcal B\), and a
projection \(F\in\mathcal B\) such that
\[
  \|S_h-T_h\|_2\leq\varepsilon_2\quad(h=1,2,\ldots,m),
  \qquad
  \|F-E\|_2\leq\varepsilon_2,
  \qquad
  \varphi(F)=2^{-n}.
\]
\end{lemma}

\begin{proof}
By Lemma~8, there is a factor \(\mathcal B\subset\mathcal A\) of
type \(\mathrm{I}_{2^p}\). It contains elements
\[
  S_1,S_2,\ldots,S_m,S_{m+1},
\]
such that
\[
  \|S_h-T_h\|_2\leq\varepsilon_1\quad(h=1,2,\ldots,m),
  \qquad
  \|S_{m+1}-E\|_2\leq\varepsilon_1.
\]
By Lemma~3 we may, after enlarging \(\mathcal B\) if necessary,
suppose that \(p\geq n\). Replacing \(S_{m+1}\) by
\(\tfrac12(S_{m+1}+S_{m+1}^*)\), we may suppose that it is
self-adjoint. Put
\[
  R=2S_{m+1}(I+S_{m+1}^2)^{-1}.
\]
We have \(\|R\|\leq1\), and, observing that
\(E=2E(I+E^2)^{-1}\),
\begin{align*}
 \frac12(R-E)
 &=(I+S_{m+1}^2)^{-1}(S_{m+1}-E)(I+E)^{-1} \\
 &\quad+\frac14R(E-S_{m+1})E
\end{align*}
[cf. \hyperref[eq:I-3-kaplansky-1]{Chapter~I, \S~3, formula~(1)}]; hence
\[
  \|R-E\|_2\leq\frac52\|S_{m+1}-E\|_2
  \leq\frac52\varepsilon_1.
\]%
% Source PDF page 287 (printed p. 278)
\phantomsection\label{srcpage:278}%
By Lemma~4(i), there is a projection \(F_1\in\mathcal B\) such that
\[
  \|F_1-E\|_2\leq15\varepsilon_1^{1/2}.
\]
Thus
\[
  |\varphi(F_1)-2^{-n}|
  =|\varphi(F_1)-\varphi(E)|
  =|\varphi(F_1-E)|
  \leq\|F_1-E\|_2
  \leq15\varepsilon_1^{1/2}.
\]
There is a projection \(F\in\mathcal B\), either dominated by or
dominating \(F_1\), such that \(\varphi(F)=2^{-n}\), because
\(\mathcal B\) is of type \(\mathrm{I}_{2^p}\) and \(p\geq n\). We have
\[
  \|F_1-F\|_2^2
  =\varphi((F_1-F)^2)
  =|\varphi(F_1)-\varphi(F)|
  =|\varphi(F_1)-2^{-n}|
  \leq15\varepsilon_1^{1/2}.
\]
Hence
\[
  \|F-E\|_2
  \leq\|F-F_1\|_2+\|F_1-E\|_2
  \leq15^{1/2}\varepsilon_1^{1/4}
       +15\varepsilon_1^{1/2}
  \leq\varepsilon_2
\]
when \(\varepsilon_1\) is sufficiently small.
\end{proof}

\begin{lemma}\label{lem:III-7-10}
Given \(T_1,T_2,\ldots,T_m\in\mathcal A\), a projection
\(E\in\mathcal A\) such that
\[
  \varphi(E)=2^{-n},
  \qquad
  ET_h=T_hE=T_h\quad(h=1,2,\ldots,m),
\]
and \(\varepsilon_3>0\), there are a factor
\(\mathcal B\subset\mathcal A\) of type \(\mathrm{I}_{2^p}\), with
\(p\geq n\), such that \(E\in\mathcal B\), and elements
\(S_1,S_2,\ldots,S_m\in\mathcal B\) such that
\[
  ES_h=S_hE=S_h,
  \qquad
  \|S_h-T_h\|_2\leq\varepsilon_3
  \quad(h=1,2,\ldots,m).
\]
Moreover, if \(E\) is a minimal projection of a factor
\(\mathcal C\subset\mathcal A\) of type \(\mathrm{I}_{2^n}\), then we may
suppose that \(\mathcal B\supset\mathcal C\).
\end{lemma}

\begin{proof}
By Lemma~9, there are a factor
\(\mathcal B_1\subset\mathcal A\) of type \(\mathrm{I}_{2^p}\), with
\(p\geq n\), elements \(R_1,R_2,\ldots,R_m\in\mathcal B_1\), and a
projection \(F\in\mathcal B_1\), such that
\[
  \|R_h-T_h\|_2\leq\varepsilon_2\quad(h=1,2,\ldots,m),
  \qquad
  \|F-E\|_2\leq\varepsilon_2,
  \qquad
  \varphi(F)=2^{-n}.
\]
By Lemma~7 there is a unitary \(U\in\mathcal A\) such that
\[
  E=U^{-1}FU,
  \qquad
  \|U-I\|_2\leq36\|F-E\|_2^{1/8}
  \leq36\varepsilon_2^{1/8}.
\]
Then \(\mathcal B=U^{-1}\mathcal B_1U\) is a factor of type
\(I_{2^p}\), contained in \(\mathcal A\), and containing \(E\). Put
\[
  R_h'=U^{-1}R_hU\in\mathcal B.
\]
We have
\begin{align*}
 \|R_h'-T_h\|_2
 &=\|R_h-U T_hU^{-1}\|_2 \\
 &\leq\|R_h-T_h\|_2+\|T_h-UT_hU^{-1}\|_2 \\
 &\leq\varepsilon_2+\|T_h-UT_h\|_2
       +\|UT_h-UT_hU^{-1}\|_2 \\
 &\leq\varepsilon_2+\|T_h\|\,\|I-U\|_2
       +\|T_h\|\,\|I-U^{-1}\|_2 \\
 &\leq\varepsilon_2+72\|T_h\|\varepsilon_2^{1/8}.
\end{align*}%
% Source PDF page 288 (printed p. 279)
\phantomsection\label{srcpage:279}%
Choose the number \(\varepsilon_2\) in Lemma~9 so that
\[
  \varepsilon_2+72\|T_h\|\varepsilon_2^{1/8}
  \leq\varepsilon_3
  \qquad(h=1,2,\ldots,m).
\]
Then \(\|R_h'-T_h\|_2\leq\varepsilon_3\), and therefore
\[
  \|ER_h'E-T_h\|_2
  =\|E(R_h'-T_h)E\|_2
  \leq\|R_h'-T_h\|_2
  \leq\varepsilon_3.
\]
It is enough to put \(S_h=ER_h'E\).

Now let \(E_1,E_2,\ldots,E_{2^n}\) be mutually orthogonal minimal
projections of \(\mathcal C\), with sum \(I\), and let \(E_1=E\).
These projections identify \(\mathcal A\) with
\(\mathcal A_E\otimes\mathcal L(\mathfrak H_{2^n})\),
where \(\mathfrak H_{2^n}\) is a Hilbert space of dimension \(2^n\),
and identify \(\mathcal C\) itself with
\(\mathcal C_E\otimes\mathcal L(\mathfrak H_{2^n})\).
Then
\[
  \mathcal B_E\otimes
  \mathcal L(\mathfrak H_{2^n})
\]
is a factor of type \(\mathrm{I}_{2^p}\), lying between \(\mathcal C\) and
\(\mathcal A\), and containing \(E\) and the \(S_h\). Replacing
\(\mathcal B\) by this factor proves the final assertion.
\end{proof}

\begin{lemma}\label{lem:III-7-11}
Given \(T_1,T_2,\ldots,T_m\in\mathcal A\), a factor
\(\mathcal B\subset\mathcal A\) of type \(\mathrm{I}_{2^n}\), and
\(\varepsilon_4>0\), there are a factor
\(\mathcal C\subset\mathcal A\) of type \(\mathrm{I}_{2^p}\), with
\(p\geq n\), such that
\(\mathcal B\subset\mathcal C\subset\mathcal A\), and elements
\(S_1,S_2,\ldots,S_m\in\mathcal C\) such that
\[
  \|S_h-T_h\|_2\leq\varepsilon_4
  \qquad(h=1,2,\ldots,m).
\]
\end{lemma}

\begin{proof}
Let \(E_1,E_2,\ldots,E_{2^n}\) be mutually orthogonal minimal
projections of \(\mathcal B\), with sum \(I\), and let
\(W_i\in\mathcal B\) be a partial isometry with initial projection
\(E_1\) and final projection \(E_i\). Put
\[
  T_{ijh}=W_i^*T_hW_j.
\]
Then
\[
  E_1T_{ijh}=T_{ijh}E_1=T_{ijh}.
\]
By Lemma~10 there are a factor \(\mathcal C\) of type
\(I_{2^p}\), with \(p\geq n\) and
\(\mathcal B\subset\mathcal C\subset\mathcal A\), and elements
\(S_{ijh}\in\mathcal C\) such that
\(\|S_{ijh}-T_{ijh}\|_2\leq\varepsilon_3\) for all \(i,j,h\).
We have
\[
  T_h=\sum_{i=1}^{2^n}\sum_{j=1}^{2^n}E_iT_hE_j
     =\sum_{i,j}W_iW_i^*T_hW_jW_j^*
     =\sum_{i,j}W_iT_{ijh}W_j^*.
\]
Put
\[
  S_h=\sum_{i,j}W_iS_{ijh}W_j^*\in\mathcal C.
\]
Then
\[
  \|S_h-T_h\|_2
  \leq\sum_{i,j}\|W_i(S_{ijh}-T_{ijh})W_j^*\|_2
  \leq\sum_{i,j}\|S_{ijh}-T_{ijh}\|_2
  \leq2^{2n}\varepsilon_3
  \leq\varepsilon_4
\]
if \(\varepsilon_3\) is sufficiently small.
\end{proof}
% Source PDF page 289 (printed p. 280)
\phantomsection\label{srcpage:280}
\begin{theorem}\label{thm:III-7-3}
Let \(\mathcal A\) be a finite factor. The following three conditions
are equivalent:
\begin{enumerate}[label=\textup{(\roman*)},leftmargin=*]
\item \(\mathcal A\) is hyperfinite;
\item \(\mathcal A\) is the von Neumann algebra generated by an
increasing sequence of involutive subalgebras of finite dimension;
\item
  \begin{enumerate}[label=\textup{(\alph*)},leftmargin=2em]
  \item \(\mathcal A\) is the von Neumann algebra generated by a
  countable family of elements;
  \item given \(T_1,T_2,\ldots,T_m\in\mathcal A\) and
  \(\varepsilon>0\), there are a finite-dimensional involutive
  subalgebra \(\mathcal B\) of \(\mathcal A\) and elements
  \(S_1,S_2,\ldots,S_m\in\mathcal B\) such that
  \[
    \|S_h-T_h\|_2\leq\varepsilon
    \qquad(h=1,2,\ldots,m).
  \]
  \end{enumerate}
\end{enumerate}
\end{theorem}

\begin{proof}
It is clear that (i) implies (ii). Suppose that \(\mathcal A\) is the
von Neumann algebra generated by an increasing sequence
\(\mathcal A_0,\mathcal A_1,\ldots\) of finite-dimensional
involutive subalgebras. Condition (iii)(a) is then plainly satisfied,
and condition (iii)(b) follows from Lemma~1. Finally, suppose that
(iii)(a) and (iii)(b) hold. Let \(R_1,R_2,\ldots\) be a sequence of
elements generating \(\mathcal A\). Condition (iii)(b), which is just
condition \eqref{eq:III-7-star} at the beginning of this numbered paragraph,
allows us to apply Lemma~11. By applying that lemma inductively, we can
construct an increasing sequence
\(\mathcal B_0,\mathcal B_1,\ldots\) of factors contained in
\(\mathcal A\), with the following properties:
\begin{enumerate}[label=\arabic*\textsuperscript{o},leftmargin=*]
\item \(\mathcal B_i\) is of type \(\mathrm{I}_{2^{p_i}}\);
\item \(\mathcal B_i\) contains elements
\(S_1^i,S_2^i,\ldots,S_i^i\) such that
\[
  \|S_h^i-R_h\|_2\leq\frac1i
  \qquad(h=1,2,\ldots,i).
\]
\end{enumerate}
By Lemma~1, the von Neumann algebra generated by the
\(\mathcal B_i\) contains every \(R_h\), and hence is
\(\mathcal A\). If
\(\mathcal B_i=\mathcal B_{i+1}=\cdots\), then \(\mathcal A\) is of
type \(\mathrm{I}_{2^{p_i}}\). If \(p_i\to+\infty\), factors of the intervening
types \(\mathrm{I}_{2^n}\) may be inserted between the \(\mathcal B_i\), by
Lemma~3; after renumbering, we may therefore suppose that
\(\mathcal B_i\) is of type \(\mathrm{I}_{2^i}\). Thus \(\mathcal A\) is
hyperfinite in this case as well.
\end{proof}

Bibliography: \ArticleRef{67}, \ArticleRef{144}, \ArticleRef{181},
\ArticleRef{197}.

\NumberedParagraph{5}{Hyperfinite Factors and Elementary Operations}
\label{no:III-7-5}

\begin{proposition}\label{prop:III-7-2}
The tensor product of two hyperfinite factors is a hyperfinite factor.
\end{proposition}

\begin{proof}
We already know that, if \(\mathcal A\) and \(\mathcal B\) are finite
  factors, then \(\mathcal A\otimes\mathcal B\) is a
finite factor (\hyperref[prop:I-6-12]{Chapter~I, \S~6, Proposition~12}
and the \hyperref[cor:I-6-14]{corollary to Proposition~14}). Suppose,
moreover, that \(\mathcal A\) (respectively
\(\mathcal B\)) is generated as a von Neumann algebra by an increasing
sequence \((\mathcal A_i)\) (respectively \((\mathcal B_i)\)) of
finite-dimensional involutive subalgebras. Then
\(\mathcal A\otimes\mathcal B\) is generated by the
increasing sequence \((\mathcal A_i\otimes\mathcal B_i)\)
(\hyperref[prop:I-2-6]{Chapter~I, \S~2, Proposition~6}), and is
therefore hyperfinite.
\end{proof}

% Source PDF page 290 (printed p. 281)
\phantomsection\label{srcpage:281}
\begin{proposition}\label{prop:III-7-3}
Let \(\mathcal A\) be a hyperfinite factor and let \(E\) be a
projection of \(\mathcal A\). Then \(\mathcal A_E\) is hyperfinite.
\end{proposition}

\begin{proof}
Let \(\varphi\) be the canonical trace of \(\mathcal A\).

\emph{(i)} Let \(\mathcal A_1,\mathcal A_2,\ldots\) be an increasing
sequence of factors of types \(I_2,I_{2^2},\ldots\) generating
\(\mathcal A\). Suppose that \(E\in\mathcal A_i\). Then the
\((\mathcal A_j)_E\), for \(j\geq i\), are increasing
finite-dimensional involutive algebras generating
\(\mathcal A_E\); hence \(\mathcal A_E\) is hyperfinite. More
generally, if \(\varphi(E)=p2^{-n}\), where \(p,n\) are integers, then
\(E\) is equivalent to a projection belonging to one of the
\(\mathcal A_i\), and \(\mathcal A_E\) is again hyperfinite.

\emph{(ii)} We pass to the general case. The factor
\(\mathcal A_E\) plainly satisfies condition (iii)(a) of Theorem~3.
We show that it also satisfies condition (iii)(b). Let
\(T_1,T_2,\ldots,T_m\in E\mathcal A E\), and let
\(\varepsilon>0\). There is a projection \(F\in\mathcal A\), with
\(F\leq E\), such that \(\varphi(F)\) is of the form \(p2^{-n}\) and
\(\|E-F\|_2\leq\varepsilon\). We have
\[
  \|T_i-FT_iF\|_2
  \leq\|ET_i-FT_i\|_2+\|FT_iE-FT_iF\|_2
  \leq2\varepsilon\|T_i\|.
\]
By part (i) of the proof, there are a finite-dimensional involutive
subalgebra \(\mathcal B\) of \(F\mathcal A F\), and elements
\(S_1,S_2,\ldots,S_m\in\mathcal B\), such that
\[
  \|FT_iF-S_i\|_2\leq\varepsilon
  \qquad(i=1,2,\ldots,m).
\]
Hence
\[
  \|T_i-S_i\|_2\leq\varepsilon(1+2\|T_i\|)
  \qquad(i=1,2,\ldots,m).
\]
Moreover, \(\mathcal B\subset E\mathcal A E\). Finally, the canonical
trace of \(\mathcal A\) induces on \(\mathcal A_E\) a trace
proportional to its canonical trace. Thus the norm \(\|\cdot\|_2\) of
\(\mathcal A\) induces on \(\mathcal A_E\) a norm proportional to its
own \(2\)-norm.
\end{proof}

\begin{proposition}\label{prop:III-7-4}
Let \(\mathcal A\) be a hyperfinite factor. If \(\mathcal A'\) is a
finite factor, then \(\mathcal A'\) is hyperfinite.
\end{proposition}

\begin{proof}
There exist (\hyperref[prop:I-6-13]{Chapter~I, \S~6, Proposition~13})
Hilbert spaces
\(\mathfrak H_1\) and \(\mathfrak K_1\), a factor
\(\mathcal A_1\subset\mathcal L(\mathfrak H_1)\) anti-isomorphic to
\(\mathcal A\), and a projection
\(E\in\mathcal A_1\otimes
\mathcal L(\mathfrak K_1)\), such that \(\mathcal A'\) is isomorphic
to
\[
  \bigl(E(\mathcal A_1\otimes
  \mathcal L(\mathfrak K_1))E\bigr).
\]
The space \(\mathfrak H_1\otimes\mathfrak K_1\) is the Hilbert direct
sum of a family \((\mathfrak K_i)_{i\in I}\) of subspaces having the
following properties:
\begin{enumerate}[label=\arabic*\textsuperscript{o},leftmargin=*]
\item \(P_i\in\mathcal A_1\otimes
\mathcal L(\mathfrak K_1)\), where \(P_i\) is the projection onto
\(\mathfrak K_i\);
\item the algebras
  \(P_i(\mathcal A_1\otimes
\mathcal L(\mathfrak K_1))P_i\) are isomorphic to \(\mathcal A_1\).
\end{enumerate}
Let \(\psi\) be a faithful normal semifinite trace on
\((\mathcal A_1\otimes
\mathcal L(\mathfrak K_1))^+\). Since \(\mathcal A'\) is finite and
\(\mathcal A_1\otimes\mathcal L(\mathfrak K_1)\)
is a factor, \(\psi(E)<+\infty\). Hence, for some sufficiently large
finite subset \(J\subset I\),
\[
  \psi(E)\leq\psi\!\left(\sum_{i\in J}P_i\right).
\]
Thus, by \hyperref[prop:III-2-13]{Proposition~13 of \S~2},
\[
  E\precsim\sum_{i\in J}P_i.
\]%
% Source PDF page 291 (printed p. 282)
\phantomsection\label{srcpage:282}%
After replacing \(E\) by an equivalent projection, we may suppose that
\(E\leq\sum_{i\in J}P_i\). Replacing
\(\mathfrak H_1\otimes\mathfrak K_1\) by
\(\sum_{i\in J}\mathfrak K_i\), we see that we may ultimately suppose
that \(\mathfrak K_1\) is finite-dimensional. Since \(\mathcal A_1\)
is hyperfinite, Proposition~2 shows that
\(\mathcal A_1\otimes
\mathcal L(\mathfrak K_1)\) is hyperfinite; Proposition~3 then shows
that its corner by \(E\) is hyperfinite.
\end{proof}

Bibliography: \ArticleRef{67}.

\NumberedParagraph{6}{New Examples of Finite Factors}
\label{no:III-7-6}
Let \(G\) be a discrete group with identity element \(e\), and let
\(\mathfrak A\) be the algebra of complex-valued functions on \(G\)
that vanish except at finitely many points, multiplication in
\(\mathfrak A\) being convolution:
\[
  (x\star y)(g)=\sum_{h\in G}x(h)y(h^{-1}g)
\]
(the sum has only finitely many nonzero terms). For
\(x\in\mathfrak A\), define
\[
  x^*(g)=\overline{x(g^{-1})}.
\]
For \(x,y\in\mathfrak A\), put
\[
  (x\mid y)=\sum_{g\in G}x(g)\overline{y(g)}.
\]
Then \(\mathfrak A\) becomes a pre-Hilbert space; its Hilbert
completion is the space \(\mathfrak H\) of complex-valued functions
\(x\) on \(G\) such that
\[
  \sum_{g\in G}|x(g)|^2<+\infty.
\]
Let \(\varepsilon_g\in\mathfrak A\) be the characteristic function of
the element \(g\in G\).

\begin{proposition}\label{prop:III-7-5}
\begin{enumerate}[label=\textup{(\roman*)},leftmargin=*]
\item \(\mathfrak A\) is a Hilbert algebra with identity
\(\varepsilon_e\).
\item \(\mathcal U(\mathfrak A)\) and
\(\mathcal V(\mathfrak A)\) are finite; every element of
\(\mathcal U(\mathfrak A)\) (respectively
\(\mathcal V(\mathfrak A)\)) is of the form \(U_a\) (respectively
\(V_a\)), where \(a\) is a bounded element of \(\mathfrak H\).
\item An element \(b\in\mathfrak H\) is central if and only if the
function \(b\) is constant on the conjugacy classes of \(G\).
\item The algebras \(\mathcal U(\mathfrak A)\) and
\(\mathcal V(\mathfrak A)\) are factors if and only if every conjugacy
class of \(G\) distinct from \(\{e\}\) is infinite.
\item If \(a,b\) are bounded elements of \(\mathfrak H\), then
\(U_ab\) (respectively \(V_ab\)) is the convolution product
\(a\star b\) (respectively \(b\star a\)).
\end{enumerate}
\end{proposition}

\begin{proof}
A direct calculation shows that, for
\(x,y,z\in\mathfrak A\),
\[
  (x\star y)^*=y^*\star x^*,
  \qquad
  (x\star y\mid z)=(y\mid x^*\star z).
\]
It is clear that \(\varepsilon_e\) is an identity element of
\(\mathfrak A\), and hence the elements of the form
\(x\star y\), \(x,y\in\mathfrak A\), are all the elements of
\(\mathfrak A\). If
\[
  \|x\|_1=\sum_{g\in G}|x(g)|,
\]%
% Source PDF page 292 (printed p. 283)
\phantomsection\label{srcpage:283}%
then
\[
  \|x\star y\|_2\leq\|x\|_1\|y\|_2.
\]
Thus \(\mathfrak A\) is a Hilbert algebra. Since \(\mathfrak A\) has
an identity, its characteristic projection is \(I\), and therefore
\(\mathcal U(\mathfrak A)\) and \(\mathcal V(\mathfrak A)\) are
finite. Moreover, the two-sided ideal of \(\mathcal U(\mathfrak A)\)
(respectively \(\mathcal V(\mathfrak A)\)) formed by the elements
\(U_a\) (respectively \(V_a\)), with \(a\) bounded, contains
\(U_{\varepsilon_e}=I\), and hence is the whole of
\(\mathcal U(\mathfrak A)\) (respectively
\(\mathcal V(\mathfrak A)\)). This proves (i) and (ii).

We prove (v). For \(a,b\in\mathfrak A\) and \(g\in G\),
\[
  (U_ab)(g)=(V_ba)(g)=(a\star b)(g).
\]
For fixed \(g\in G\), each of \((U_ab)(g)\), \((V_ba)(g)\), and
\((a\star b)(g)\) is separately continuous in \(a\) and \(b\) as
\(a,b\) range over the bounded elements of \(\mathfrak H\). Hence
\(U_ab=V_ba=a\star b\) for bounded \(a,b\in\mathfrak H\).

We prove (iii). An element \(b\in\mathfrak H\) is central if and only
if \(U_{\varepsilon_g}b=V_{\varepsilon_g}b\) for every \(g\in G\),
that is, if and only if
\(\varepsilon_g\star b=b\star\varepsilon_g\) for every \(g\in G\).
This says precisely that the functions
\(h\mapsto b(h)\) and \(h\mapsto b(g^{-1}hg)\) are identical, or,
equivalently, that \(b\) is constant on conjugacy classes.

If \(G\) has a finite conjugacy class distinct from \(\{e\}\), then
\(\mathfrak A\) has a central element not proportional to
\(\varepsilon_e\), so \(\mathcal U(\mathfrak A)\) has a nonscalar
central element and is not a factor. Conversely, if
\(\mathcal U(\mathfrak A)\) has a nonscalar central element \(U_a\),
with \(a\) bounded, then \(a\) is central and not proportional to
\(\varepsilon_e\). Hence it is constant and nonzero on a conjugacy
class distinct from \(\{e\}\); since
\(\sum_{g\in G}|a(g)|^2<+\infty\), that class is finite. This proves
(iv).
\end{proof}

Bibliography: \ArticleRef{1}, \ArticleRef{2}, \ArticleRef{13},
\ArticleRef{29}, \ArticleRef{30}, \ArticleRef{53}, \ArticleRef{56},
\ArticleRef{67}, \ArticleRef{90}, \ArticleRef{98}, \ArticleRef{99},
\TreatiseRef{9}.

\NumberedParagraph{7}{Existence of Nonhyperfinite Finite Factors}
\label{no:III-7-7}

\begin{definition}\label{def:III-7-2}
Let \(\mathcal A\) be a finite factor, and let \(\varphi\) be its
canonical trace. The factor \(\mathcal A\) is said to have
\emph{property \(\Gamma\)} if, given
\(T_1,T_2,\ldots,T_m\in\mathcal A\) and \(\varepsilon>0\), there is a
unitary operator \(U\in\mathcal A\) such that \(\varphi(U)=0\) and
\[
  \|U^{-1}T_hU-T_h\|_2\leq\varepsilon
  \qquad(h=1,2,\ldots,m).
\]
\end{definition}

Property \(\Gamma\) is invariant under isomorphism.

\begin{proposition}\label{prop:III-7-6}
A continuous hyperfinite factor has property \(\Gamma\).
\end{proposition}

\begin{proof}
Let \(\mathcal A\) be a continuous hyperfinite factor, let
\(T_1,T_2,\ldots,T_m\in\mathcal A\), and let \(\varepsilon>0\). Let
\(\mathcal A_0,\mathcal A_1,\mathcal A_2,\ldots\) be an increasing
sequence of factors of types \(I_2,I_{2^2},I_{2^3},\ldots\) generating
\(\mathcal A\). There are an integer \(i\) and elements
\(S_1,S_2,\ldots,S_m\in\mathcal A_i\) such that
\[
  \|T_h-S_h\|_2\leq\frac{\varepsilon}{2}
  \qquad(h=1,2,\ldots,m).
\]
The factor \(\mathcal A_{i+1}\) is isomorphic to the tensor product of
\(\mathcal A_i\) and a factor of type \(\mathrm{I}_2\). Hence there is a factor
\(\mathcal B\) of type \(\mathrm{I}_2\), commuting with \(\mathcal A_i\) and
contained in \(\mathcal A_{i+1}\). Let \(U\) be a unitary operator of
\(\mathcal B\) whose canonical trace, relative to \(\mathcal B\) and
therefore relative to \(\mathcal A\), is zero; a symmetry may be used.
Then%
% Source PDF page 293 (printed p. 284)
\phantomsection\label{srcpage:284}%
\[
  \|U^{-1}(T_h-S_h)U\|_2=\|T_h-S_h\|_2
  \leq\frac{\varepsilon}{2},
  \qquad
  U^{-1}S_hU=S_h.
\]
Consequently,
\[
  \|U^{-1}T_hU-T_h\|_2\leq\varepsilon.
\]
\end{proof}

\begin{lemma}\label{lem:III-7-12}
Suppose that, in the group \(G\) of Proposition~5, every conjugacy
class distinct from \(\{e\}\) is infinite, and that there is a subset
\(F\subset G\) with the following properties:
\begin{enumerate}[label=\textup{(\roman*)},leftmargin=*]
\item there is an element \(g_1\in G\) such that
\[
  F\cup g_1Fg_1^{-1}\cup\{e\}=G;
\]
\item there are elements \(g_2,g_3\in G\) such that
\[
  F,\qquad g_2Fg_2^{-1},\qquad g_3Fg_3^{-1}
\]
are pairwise disjoint.
\end{enumerate}
Then \(\mathcal U(\mathfrak A)\) is a finite factor that does not have
property \(\Gamma\).
\end{lemma}

\begin{proof}
Suppose that \(\mathcal U(\mathfrak A)\) has property \(\Gamma\).
Let \(\varepsilon>0\), and let \(U_y\), where \(y\) is a bounded
element of \(\mathfrak H\), be a unitary operator of canonical trace
zero such that
\[
  \|U_{g_i}-U_y^*U_{g_i}U_y\|_2\leq\varepsilon
  \qquad(i=1,2,3),
\]
where \(U_g=U_{\varepsilon_g}\). We have
\[
  0=\varphi(U_y)=(U_yI\mid I)=(y\mid\varepsilon_e)=y(e).
\]
Furthermore,
\[
  \|U_{g_i}-U_y^*U_{g_i}U_y\|_2
  =\|U_y-U_{g_i}^{-1}U_yU_{g_i}\|_2
  =\|y-\varepsilon_{g_i^{-1}}\star y\star\varepsilon_{g_i}\|_2
  \leq\varepsilon,
\]
and
\[
  \|y\|_2=\|U_y\|_2=1.
\]
It follows from (i) that
\begin{align*}
  1
  &\leq \sum_{g\in F}|y(g)|^2
       +\sum_{g\in g_1Fg_1^{-1}}|y(g)|^2 \\
  &=\sum_{g\in F}|y(g)|^2
       +\sum_{g\in F}
       |(\varepsilon_{g_1^{-1}}\star y\star\varepsilon_{g_1})(g)|^2
   \leq2\sum_{g\in F}|y(g)|^2+2\varepsilon.
\end{align*}
On the other hand, (ii) gives
\begin{align*}
  1
  &\geq \sum_{g\in F}|y(g)|^2
       +\sum_{g\in g_2Fg_2^{-1}}|y(g)|^2
       +\sum_{g\in g_3Fg_3^{-1}}|y(g)|^2 \\
  &=\sum_{g\in F}|y(g)|^2
       +\sum_{g\in F}
       |(\varepsilon_{g_2^{-1}}\star y\star\varepsilon_{g_2})(g)|^2
       +\sum_{g\in F}
       |(\varepsilon_{g_3^{-1}}\star y\star\varepsilon_{g_3})(g)|^2 \\
  &\geq3\sum_{g\in F}|y(g)|^2-4\varepsilon.
\end{align*}
Thus
\[
  \frac12-\varepsilon
  \leq\sum_{g\in F}|y(g)|^2
  \leq\frac13+\frac43\varepsilon,
\]
which is impossible when \(\varepsilon\) is sufficiently small.
\end{proof}

% Source PDF page 294 (printed p. 285)
\phantomsection\label{srcpage:285}
\begin{theorem}\label{thm:III-7-4}
There exist nonhyperfinite finite factors acting on a Hilbert space
with a countable basis.
\end{theorem}

\begin{proof}
Let \(G\) be the free group on two generators \(g_1,g_2\). Every
element of \(G\) distinct from \(e\) has a canonical expression of the
form
\[
  \cdots g_1^{n}g_2^{p}g_1^{q}g_2^{r}\cdots,
\]
with all exponents nonzero. Let \(a\in G\). The elements
\(g_1^nag_1^{-n}\) are all distinct unless \(a\) is a power of
\(g_1\); in the latter case, the elements \(g_2^nag_2^{-n}\) are all
distinct unless \(a=e\). Thus every conjugacy class distinct from
\(\{e\}\) is infinite. Let \(F\) be the set of all \(a\neq e\) whose
canonical expression ends in \(g_1^n\), where
\(n=\pm1,\pm2,\ldots\). Conditions (i) and (ii) of Lemma~12 hold, with
\(g_3=g_2^{-1}\). Hence \(\mathcal U(\mathfrak A)\) is a continuous
finite factor without property \(\Gamma\), and is therefore not
hyperfinite.
\end{proof}

For the classification of factors of type \(\mathrm{II}_1\), see the
\hyperref[no:I-9-4]{remarks at the end of Chapter~I, \S~9, no.~4}.

Bibliography: \ArticleRef{67}.

\noindent\textit{Exercises.}
\begin{enumerate}[label=\arabic*.,leftmargin=*]
\item\label{ex:III-7-1} Let \(\nu\) be Lebesgue measure on \([0,2]\), and put
\[
  \mathfrak H=L^2_{\mathbb C}([0,2],\nu).
\]
Every \(\nu\)-measurable
set \(Z\subset[0,2]\) defines a projection of \(\mathfrak H\), namely
multiplication by the characteristic function of \(Z\). In
particular, let \(E\) be the projection corresponding to \([0,1]\).

For \(n=1,2,\ldots\), let
\(Z_1^n,Z_2^n,\ldots,Z_{2^n}^n\) be a partition of \([0,1]\) into
measurable sets with the following properties:
\begin{enumerate}[label=\textit{\alph*.},leftmargin=2em]
\item \(Z_i^{n-1}=Z_{2i-1}^n\cup Z_{2i}^n\);
\item \(\nu(Z_{2i-1}^n)=\tfrac34\nu(Z_i^{n-1})\);
\item \(\nu(Z_{2i}^n)=\tfrac14\nu(Z_i^{n-1})\).
\end{enumerate}
Let \(Y_i^n\subset[1,2]\) be defined as follows:
\(Y_{2i-1}^n\) (respectively \(Y_{2i}^n\)) is obtained from
\(Z_{2i}^n\) (respectively \(Z_{2i-1}^n\)) by the translation
\(\zeta\mapsto\zeta+1\). Put
\(X_i^n=Z_i^n\cup Y_i^n\), and let \(E_i^n\) be the projection of
\(\mathcal L(\mathfrak H)\) corresponding to \(X_i^n\). For each
\(n\), the projections \(E_1^n,E_2^n,\ldots,E_{2^n}^n\) may be
regarded as the minimal projections of a factor \(\mathcal A_n\) of
type \(\mathrm{I}_{2^n}\) commuting with \(E\). The \(\mathcal A_n\) may also
be chosen to form an increasing sequence.

Show that the von Neumann algebra \(\mathcal A\) generated by the
\(\mathcal A_n\) is not a factor. Show that
\(E_1^n+E_3^n+\cdots+E_{2^n-1}^n\) converges strongly to \(E\) as
\(n\to+\infty\), so that \(E\) belongs to the centre of
\(\mathcal A\).

\item Show that, if \(\mathcal A\) is a finite von Neumann algebra,
the set of factors contained in \(\mathcal A\), ordered by inclusion,
is inductive. Use Proposition~1. \ArticleRef{22}

\item Let \(\mathcal A\) be a finite von Neumann algebra and let
\((\mathcal A_i)_{i\in I}\) be a family, totally ordered by inclusion,
of von Neumann subalgebras of \(\mathcal A\), all having the same
centre \(\mathcal Z\) as \(\mathcal A\). Show that the von Neumann
algebra generated by the \(\mathcal A_i\) has centre \(\mathcal Z\).
Argue as in Proposition~1.

\item
\begin{enumerate}[label=\textit{\alph*.},leftmargin=2em]
\item Let \(\mathcal A\) be a von Neumann algebra, let \(\varphi\) be
a faithful finite normal trace on \(\mathcal A\), and let
\(\mathcal B\) be a maximal abelian von Neumann subalgebra of
\(\mathcal A\). Suppose that \(T\in\mathcal A\) and
\(\varepsilon>0\) satisfy
\[
  \|TT'-T'T\|_2\leq\varepsilon\|T'\|
  \qquad(T'\in\mathcal B).
\]%
% Source PDF page 295 (printed p. 286)
\phantomsection\label{srcpage:286}%
Then the distance from \(T\) to \(\mathcal B\), for the norm
\(\|\cdot\|_2\), is at most \(\varepsilon\). By
\hyperref[lem:I-9-1]{Chapter~I, \S~9, Lemma~1}, there is an element
\(S\in\mathcal B\) strongly adherent to
the convex set \(\mathcal K\) generated by the
\(UTU^{-1}\), where \(U\) ranges over the unitary operators of
\(\mathcal B\). For every \(R\in\mathcal K\),
\(\|R-T\|_2\leq\varepsilon\), and hence
\(\|S-T\|_2\leq\varepsilon\).

\item Let \(\mathcal A\) be a finite von Neumann algebra, let
\((\mathcal B_i)_{i\in I}\) be a family, totally ordered by inclusion,
of von Neumann subalgebras of \(\mathcal A\), and let
\(\mathcal C_i\) be a maximal abelian von Neumann subalgebra of
\(\mathcal B_i\), with \(\mathcal C_i\subset\mathcal C_j\) whenever
\(\mathcal B_i\subset\mathcal B_j\). Let \(\mathcal B\)
(respectively \(\mathcal C\)) be the von Neumann algebra generated by
the \(\mathcal B_i\) (respectively \(\mathcal C_i\)). Show that
\(\mathcal C\) is a maximal abelian von Neumann subalgebra of
\(\mathcal B\). Argue as in Proposition~1, using (a).
\end{enumerate}

\item Let \(\mathcal A\) be a factor of type \(\mathrm{I}_n\), with \(n\)
finite, and let \(\mathcal B\) be a continuous finite factor
(respectively a factor of type \(\mathrm{I}_{pn}\)) containing \(\mathcal A\).
For every integer \(q\) (respectively every divisor \(q\) of \(p\)),
there is a factor \(\mathcal A_1\) of type \(\mathrm{I}_{qn}\) such that
\(\mathcal A\subset\mathcal A_1\subset\mathcal B\). Argue as in
Lemma~3. \ArticleRef{67}

\item Let \(\mathcal A\) be a continuous hyperfinite factor and let
\((p_1,p_2,\ldots)\) be a sequence of integers tending to
\(+\infty\), with \(p_i\) dividing \(p_{i+1}\). Show that there is an
increasing sequence of factors \(\mathcal A_i\subset\mathcal A\), of
type \(\mathrm{I}_{p_i}\), that generates \(\mathcal A\) as a von Neumann
algebra. By arguing as in Theorem~1, first show that the assertion holds
for some continuous hyperfinite factor; then use Theorem~2.
\ArticleRef{67}

\item A hyperfinite factor is isomorphic to a factor acting on a
Hilbert space with a countable basis; see
\hyperref[ex:I-7-3]{Chapter~I, \S~7, Exercise~3(c)}.

\item Let \(I\) be an uncountable set, and let \(G\) be the discrete
group of permutations of \(I\) that fix all but finitely many elements
of \(I\).
\begin{enumerate}[label=\textit{\alph*.},leftmargin=2em]
\item Show that every conjugacy class of \(G\) distinct from
\(\{e\}\) is infinite.
\item Show that the factor \(\mathcal U(\mathfrak A)\) of
Proposition~5 has property (iii)(b) of Theorem~3. If
\(T_1,T_2,\ldots,T_m\in\mathcal U(\mathfrak A)\), the \(T_i\) can be
approximated arbitrarily closely in the norm \(\|\cdot\|_2\) by
operators \(S_i=U_{x_i}\), with \(x_i\in\mathfrak A\). The \(x_i\)
vanish outside a finite subgroup of \(G\), and consequently the
\(S_i\) generate a finite-dimensional involutive subalgebra of
\(\mathcal U(\mathfrak A)\).
\item Show that \(\mathcal U(\mathfrak A)\) does not have property
(iii)(a) of Theorem~3. If a von Neumann algebra is generated by a
countable family of elements and has a totalizing element, it acts on
a Hilbert space with a countable basis; but \(G\) is uncountable.
\end{enumerate}

\item Use the notation of no.~6. If \(a\in\mathfrak H\) is such that,
for every \(b\in\mathfrak H\), the bounded function \(a\star b\) on
\(G\) belongs to \(\mathfrak H\), then \(a\) is bounded. Show that the
mapping \(b\mapsto a\star b\) is continuous by applying the closed
graph theorem. \ArticleRef{90}

\item Let \(G\neq\{e\}\) be a discrete group every conjugacy class of
which, except \(\{e\}\), is infinite.
\begin{enumerate}[label=\textit{\alph*.},leftmargin=2em]
\item If \(G\) is the union of an increasing sequence of finite
subgroups, the factor \(\mathcal U(\mathfrak A)\) of Proposition~5 is
continuous and hyperfinite. Use Theorem~3(iii).
\item If, for every finite family \(g_1,g_2,\ldots,g_n\in G\), there
is an element \(g\neq e\) of \(G\) commuting with all the \(g_i\),
then \(\mathcal U(\mathfrak A)\) has property \(\Gamma\). Observe that
\(U_{\varepsilon_g}\) is unitary and has canonical trace zero.

% Source PDF page 296 (printed p. 287)
\phantomsection\label{srcpage:287}
\item Deduce from (a) and (b) a new proof of Proposition~6. Consider
the group of permutations of a countably infinite set \(I\) that fix
all but finitely many elements of \(I\). \ArticleRef{67}
\end{enumerate}

\item Let \(\mathcal A\) be a continuous finite factor on a Hilbert
space with a countable basis.
\begin{enumerate}[label=\textit{\alph*.},leftmargin=2em]
\item Let \((\mathcal A_i)_{i\in I}\) be a family, totally ordered by
inclusion, of hyperfinite factors contained in \(\mathcal A\). Show
that the von Neumann algebra generated by the \(\mathcal A_i\) is a
hyperfinite factor. Use Theorem~3(iii).
\item For every hyperfinite factor
\(\mathcal B\subset\mathcal A\), there is a maximal hyperfinite factor
\(\mathcal C\) such that
\(\mathcal B\subset\mathcal C\subset\mathcal A\); use (a). Moreover,
\(\mathcal C\) is continuous; use Lemma~3. \ArticleRef{22}
\end{enumerate}

\item Let \(G\) be a discrete group and let \(G_1\) be a subgroup.
Use the notation of no.~6.
\begin{enumerate}[label=\textit{\alph*.},leftmargin=2em]
\item Show that the \(U_a\), where \(a\in\mathfrak H\) is bounded
and vanishes off \(G_1\), form a von Neumann subalgebra
\(\mathcal A\) of \(\mathcal U(\mathfrak A)\). Use Lemma~1.
\item The algebra \(\mathcal A\) is a maximal abelian von Neumann
subalgebra of \(\mathcal U(\mathfrak A)\) if and only if \(G_1\) is
abelian and, for every \(g\in G\setminus G_1\), the set
\(\{g_1gg_1^{-1}:g_1\in G_1\}\) is infinite.
\item The algebra \(\mathcal A'\cap\mathcal U(\mathfrak A)\) consists
only of scalars if and only if, for every \(g\in G\), the set
\(\{g_1gg_1^{-1}:g_1\in G_1\}\) is infinite.
\item Let \(K\) be an infinite commutative field that is the union of
an increasing sequence of finite subfields, for example the algebraic
closure of a finite field. For \(\alpha\in K\), \(\alpha\neq0\), and
\(\beta\in K\), let \((\alpha,\beta)\) be the bijection of \(K\) onto
itself defined by
\[
  (\alpha,\beta)\xi=\alpha\xi+\beta.
\]
The \((\alpha,\beta)\) form a group \(G\), which is the union of an
increasing sequence of finite subgroups. Let \(G_1\) be the subgroup
of \(G\) formed by those \((\alpha,\beta)\) for which
\(\alpha,\beta\in K'\), where \(K'\) is an infinite proper subfield of
\(K\). Show that \(\mathcal U(\mathfrak A)\) is a continuous
hyperfinite factor, that \(\mathcal A\neq\mathcal U(\mathfrak A)\),
and that \(\mathcal A'\cap\mathcal U(\mathfrak A)\) consists only of
scalars. Use (c) and Exercise~10(a). \ArticleRef{16}, \ArticleRef{67}
\end{enumerate}

\item\label{ex:III-7-13} Let \(\mathcal A\) be a von Neumann algebra with centre
\(\mathcal Z\). For every subset \(\mathcal M\subset\mathcal A\),
put \(\mathcal M'_{\mathcal A}=\mathcal M'\cap\mathcal A\).
\begin{enumerate}[label=\textit{\alph*.},leftmargin=2em]
\item If \(\mathcal M\) is stable under adjunction, then
\((\mathcal M'_{\mathcal A})'_{\mathcal A}\) is a von Neumann
algebra containing \(\mathcal Z\) and \(\mathcal M\).
\item The algebra \(\mathcal A\) is called \emph{normal} if, for every
von Neumann subalgebra \(\mathcal B\subset\mathcal A\) containing
\(\mathcal Z\),
\[
  (\mathcal B'_{\mathcal A})'_{\mathcal A}=\mathcal B.
\]
Show that a discrete \(\mathcal A\) is normal. One may suppose that
\(\mathcal A'\) is abelian, and hence \(\mathcal A'=\mathcal Z\).
Then \(\mathcal B'_{\mathcal A}=\mathcal B'\) and
\((\mathcal B'_{\mathcal A})'_{\mathcal A}=
\mathcal B''=\mathcal B\).
\item A continuous finite factor \(\mathcal A\) on a Hilbert space
with a countable basis is not normal. If \(\mathcal A\) is
hyperfinite, use Exercise~12(d). If it is not hyperfinite, use
Exercise~11(b) to choose a maximal hyperfinite factor
\(\mathcal B\subset\mathcal A\), and put
\(\mathcal C=\mathcal B'_{\mathcal A}\). Distinguish three cases:
\(1^\circ\) \(\mathcal C\) is not a factor;
\(2^\circ\) \(\mathcal C\) consists only of scalars;
\(3^\circ\) \(\mathcal C\) is a nontrivial factor. Show that the third
case contradicts the maximality of \(\mathcal B\), and in the first
two cases show that
\((\mathcal B'_{\mathcal A})'_{\mathcal A}\neq\mathcal B\).

\emph{Problem (Fuglede--Kadison).} Can the first case occur?
\item A continuous semifinite factor \(\mathcal A\) on a Hilbert
space with a countable basis is not normal. Write
\(\mathcal A=\mathcal A_1\otimes
\mathcal L(\mathfrak K)\), where \(\mathcal A_1\) is a continuous
finite factor, and use (c).

\emph{Problem.} Is every continuous von Neumann algebra nonnormal?
\ArticleRef{16}, \ArticleRef{22}, \ArticleRef{65}, \ArticleRef{78},
\ArticleRef{283}
\end{enumerate}

% Source PDF page 297 (printed p. 288)
\phantomsection\label{srcpage:288}
\item Use the notation of no.~6. Let \(G_0\) be the subgroup of
\(G\) formed by its finite conjugacy classes. If \(G/G_0\) is
infinite, then \(\mathcal U(\mathfrak A)\) is continuous.

Let \(\mathcal Z=\mathcal U(\mathfrak A)\cap
\mathcal V(\mathfrak A)\), and let \((G_i)\) be the family of cosets
\(xG_0\) of \(G\) modulo \(G_0\). Let \(E_i\) be multiplication by
the characteristic function of \(G_i\) on \(\mathfrak H\). Show that
the \(E_i\) form an infinite family of mutually orthogonal equivalent
projections of \(\mathcal Z'\). Then apply Exercise~2 of \S~3.
\ArticleRef{56}

\item Use the notation of no.~6, taking \(G\) to be the discrete free
group generated by two elements \(a,b\). The automorphism of \(G\)
interchanging \(a\) and \(b\) defines an automorphism \(\Phi\) of
\(\mathcal U(\mathfrak A)\) that is not inner (Kadison). Suppose that
\(\Phi\) is implemented by a unitary \(U_x\in\mathcal U(\mathfrak A)\),
where \(x\) is a bounded element of \(\mathfrak H\). For every integer
\(n\),
\[
  U_xU_{\varepsilon_{a^n}}U_x^*=U_{\varepsilon_{b^n}},
  \qquad
  x\star\varepsilon_{a^n}\star x^*=\varepsilon_{b^n}.
\]
But if \(y\in\mathfrak H\) has finite support and is very close to
\(x\), then
\[
  (y\star\varepsilon_{a^n}\star y^*)(b^n)=0
\]
for all sufficiently large \(n\).

\item Let \(\mathcal A\) be a continuous finite von Neumann algebra
on a complex Hilbert space \(\mathfrak H\) with a countable basis, and
let \(\mathcal Z\) be its centre. Suppose there is an involutive
subalgebra \(\mathcal B\subset\mathcal A\) with the following
properties:
\begin{enumerate}[label=\textup{(\roman*)},leftmargin=2em]
\item every finite subset of \(\mathcal B\) generates a
finite-dimensional involutive algebra;
\item \(\mathcal Z\) and \(\mathcal B\) generate \(\mathcal A\).
\end{enumerate}
\begin{enumerate}[label=\textit{\alph*.},leftmargin=2em]
\item Suppose, in addition, that \(\mathcal A\) is standard. Show
that \(\mathcal A\) is spatially isomorphic to
\(\mathcal Z\otimes\mathcal A_0\), where
\(\mathcal A_0\) is a continuous hyperfinite factor acting on a
Hilbert space with a countable basis. Write
\[
  \mathcal A=\int_Z^\oplus\mathcal A(\zeta)\,d\nu(\zeta),
\]
where \(Z\) is compact metrizable and the \(\mathcal A(\zeta)\) are
factors. Show that the \(\mathcal A(\zeta)\) may be supposed standard
and hyperfinite, and hence all spatially isomorphic to
\(\mathcal A_0\). Then use
\hyperref[prop:II-3-4]{Chapter~II, \S~3, Proposition~4}.
\item In the general case, show that \(\mathcal A\) is isomorphic to
\(\mathcal Z\otimes\mathcal A_0\). Show first that
\(\mathcal A\) is isomorphic to a standard algebra acting on a Hilbert
space with a countable basis. \ArticleRef{40}, \ArticleRef{360}
\end{enumerate}
\end{enumerate}

\section{A New Definition of Finite von Neumann Algebras}
\label{sec:III-8}

\NumberedParagraph{1}{Statement of the Theorem}
\label{no:III-8-1}

\begin{theorem}\label{thm:III-8-1}
Let \(\mathcal A\) be a von Neumann algebra. It is finite if and only
if every projection of \(\mathcal A\) equivalent to \(I\) is equal to
\(I\); equivalently, the conditions
\[
  T\in\mathcal A,
  \qquad T^*T=I
\]
imply \(TT^*=I\).
\end{theorem}

The condition is necessary by \S~1, Proposition~3. Throughout the
remainder of this section, except in numbered paragraphs 6 and~7,
\(\mathcal A\) denotes a von Neumann algebra such that every
projection of \(\mathcal A\) equivalent to \(I\) is equal to \(I\).
We shall prove that \(\mathcal A\) is finite.

\begin{remark}\label{rem:III-8-1}
Suppose that the following assertion has been proved:
\begin{equation*}
\tag{\(\star\)}\label{eq:III-8-star}
\begin{minipage}{0.86\linewidth}
If \(\mathcal Z\neq0\), there is a nonzero projection \(F\) in the
centre \(\mathcal Z\) of \(\mathcal A\) such that
\(\mathcal A_F\) is finite.
\end{minipage}
\end{equation*}
% Source PDF page 298 (printed p. 289)
\phantomsection\label{srcpage:289}
Then the theorem will follow. Indeed, let \((F_i)_{i\in I}\) be a
maximal family of nonzero, pairwise orthogonal projections of
\(\mathcal Z\) such that the \(\mathcal A_{F_i}\) are finite. Put
\[
 G=\sum_{i\in I}F_i,\qquad G'=I-G.
\]
Every projection of \(\mathcal A_{G'}\) equivalent to \(I_{G'}\) is
equal to \(I_{G'}\); hence, if \(G'\ne0\), there is a nonzero
projection \(G''\) of \(\mathcal Z\), dominated by \(G'\), such that
\(\mathcal A_{G''}\) is finite, contrary to the maximality of the
family \((F_i)_{i\in I}\). Thus \(G'=0\),
\(\sum_{i\in I}F_i=I\), and consequently
\[
 \mathcal A=\prod_{i\in I}\mathcal A_{F_i}
\]
is finite.

It is assertion \eqref{eq:III-8-star} that we shall prove below.
\end{remark}

\begin{lemma}\label{lem:III-8-1}
\begin{enumerate}[label=\textup{(\roman*)},leftmargin=2em]
\item Let \(E,F\) be two equivalent projections of \(\mathcal A\). If
\(E\leq F\), then \(E=F\).
\item Let \((E_i)_{i\in I}\) be a family of nonzero equivalent,
pairwise orthogonal projections. Then \(I\) is finite.
\end{enumerate}
\end{lemma}

\begin{proof}
For (i), one has
\[
 I=F+(I-F)\sim E+(I-F),
\]
and therefore \(I=E+(I-F)\), whence \(E=F\).

For (ii), if \(I\) were infinite, there would be a proper subset
\(J\) of \(I\) equipotent to \(I\). Then
\[
 \sum_{i\in I}E_i\sim\sum_{i\in J}E_i,
\]
contrary to (i).
\end{proof}

\begin{remark}\label{rem:III-8-2}
Suppose that \(\mathcal A\) is not continuous. There is
(\hyperref[prop:III-3-2]{\S~3, Proposition~2}) a nonzero projection
\(F\) of \(\mathcal Z\) such that \(\mathcal A_F\) is homogeneous of
type \(\mathrm I_n\), where \(n\) is finite by
\hyperref[lem:III-8-1]{Lemma~1(ii)}. Thus \(\mathcal A_F\) is finite,
which proves \eqref{eq:III-8-star}. We may therefore make the following
assumption throughout the rest of this section:
\begin{equation*}
\tag{\(\star\star\)}
 \mathcal A\text{ is continuous.}
\end{equation*}
\end{remark}

\noindent Bibliography: \ArticleRef{6}, \ArticleRef{65},
\ArticleRef{66}.

\NumberedParagraph{2}{Fundamental Projections}
\label{no:III-8-2}

\begin{definition}\label{def:III-8-1}
A projection \(E\) of \(\mathcal A\) is called \emph{fundamental} if
there exist a projection \(F\) of \(\mathcal Z\) and equivalent,
pairwise orthogonal projections
\(E_1,E_2,E_3,\ldots,E_{2^n}\) of \(\mathcal A\), with sum \(F\),
such that \(E_1=E\).
\end{definition}

Every projection equivalent to a fundamental projection is
fundamental, as follows at once from the remark in \S~2, no.~3.

\begin{lemma}\label{lem:III-8-2}
Every nonzero projection \(E\) of \(\mathcal A\) majorizes a nonzero
fundamental projection.
\end{lemma}

% Source PDF page 299 (printed p. 290)
\phantomsection\label{srcpage:290}
\begin{proof}
There exist (\hyperref[cor:III-1-thm1-2]{\S~1, Theorem~1,
Corollary~2}, and \hyperref[lem:III-8-1]{Lemma~1}) a projection \(F\)
of \(\mathcal Z\) and pairwise orthogonal projections
\(E_1,E_2,\ldots,E_n\) of \(\mathcal A\), having the following
properties:
\begin{enumerate}[label=\textup{\arabic*\textdegree},leftmargin=2.3em]
\item \(E_1\sim EF\ne0\);
\item \(E_1\sim E_2\sim\cdots\sim E_{n-1}\succ E_n\);
\item \(F=E_1+E_2+\cdots+E_n\).
\end{enumerate}
Let \(p\) be an integer such that \(2^p>n\). By repeated application
of \hyperref[cor:III-1-thm1-3]{Corollary~3 of Theorem~1 in \S~1},
there is a family \((G_1,G_2,\ldots,G_{2^p})\) of equivalent,
pairwise orthogonal projections of \(\mathcal A\), with sum \(F\).
Apply \hyperref[thm:III-1-1]{Theorem~1 of \S~1} to \(E_1\) and
\(G_1\). We are reduced successively to the case \(G_1\prec E_1\)
and the case \(E_1\prec G_1\). In the first case the lemma is proved
(for \(E_1\sim EF\ne0\), hence \(F\ne0\) and \(G_1\ne0\)). In the
second case,
\[
 F=\sum_{i=1}^{n}E_i\prec\sum_{i=1}^{n}G_i\ne F,
\]
which is absurd by \hyperref[lem:III-8-1]{Lemma~1}.
\end{proof}

\begin{corollary*}\phantomsection\label{cor:III-8-lem2}
Every projection \(E\) of \(\mathcal A\) is the sum of a family of
pairwise orthogonal fundamental projections.
\end{corollary*}

\begin{proof}
Let \((E_i)_{i\in I}\) be a maximal family of nonzero, pairwise
orthogonal fundamental projections dominated by \(E\). By
\hyperref[lem:III-8-2]{Lemma~2},
\(E-\sum_{i\in I}E_i=0\), whence \(E=\sum_{i\in I}E_i\).
\end{proof}

\begin{remark*}
Let \(x\in\mathfrak H\), \(x\ne0\), and let \(F\) be a nonzero
fundamental projection dominated by \(E_x^{\mathcal A'}\). Let
\(F_1,F_2,\ldots,F_{2^n}\) be pairwise orthogonal equivalent
projections of \(\mathcal A\), whose sum \(G\) is a projection of
\(\mathcal Z\), with \(F_1=F\). Then
\[
 \mathcal A_G\cong\mathcal A_F\otimes
 \mathcal L(\mathfrak K),
\]
where \(\mathfrak K\) is a Hilbert space of dimension \(2^n\)
(Chapter~I, \S~2, Proposition~5). If \(\mathcal A_F\) is finite,
then \(\mathcal A_G\) is finite and the proof is complete
(no.~1, Remark~1). It therefore suffices to prove that
\(\mathcal A_F\) is finite. Now \(\mathcal A_F\) satisfies the same
hypothesis as \(\mathcal A\) (Lemma~1), is continuous, and \(Fx\) is
separating for \(\mathcal A_F\). In short, we may add to the preceding
assumptions on \(\mathcal A\) the following one:
\begin{equation*}
\tag{\(\star\star\star\)}\label{eq:III-8-triple-star}
 \text{There exists a separating element }x\text{ for }\mathcal A.
\end{equation*}
\end{remark*}

\noindent Bibliography: \ArticleRef{6}, \ArticleRef{65}.

\NumberedParagraph{3}{Weights on the Set of Fundamental Projections}
\label{no:III-8-3}

\begin{lemma}\label{lem:III-8-3}
Let \(E\) be a nonzero fundamental projection, let \(F\) be a
projection of \(\mathcal Z\), and let
\(E_1,E_2,\ldots,E_{2^n}\) be pairwise orthogonal projections of
\(\mathcal A\), equivalent to \(E\), with sum \(F\). Then \(F\) and
the integer \(n\) depend only on \(E\).
\end{lemma}

\begin{proof}
Every projection of \(\mathcal Z\) majorizing \(E\) majorizes the
\(E_i\), and hence \(F\), so that \(F\) is the central support of
\(E\). On the other hand,
% Source PDF page 300 (printed p. 291)
\phantomsection\label{srcpage:291}
let \(E'_1,E'_2,\ldots,E'_p\) be pairwise orthogonal projections of
\(\mathcal A\), equivalent to \(E\), with sum \(F\). If \(p<2^n\),
then
\[
 F=\sum_{i=1}^{p}E'_i\sim\sum_{i=1}^{p}E_i\ne F,
\]
which is absurd by \hyperref[lem:III-8-1]{Lemma~1}; the hypothesis
\(2^n<p\) is shown absurd in the same way.
\end{proof}

\begin{definition}\label{def:III-8-2}
Let \(E\) be a fundamental projection. If \(E\ne0\), put
\[
 \Phi(E)=2^{-n}F,
\]
where \(n\) and \(F\) are defined by \hyperref[lem:III-8-3]{Lemma~3}.
If \(E=0\), put \(\Phi(E)=0\) (an element of \(\mathcal Z\)).
\end{definition}

\begin{lemma}\label{lem:III-8-4}
Let \(E,E'\) be fundamental projections. The relations
\(\Phi(E)=\Phi(E')\) and \(E\sim E'\) are equivalent.
\end{lemma}

\begin{proof}
Lemma~3 immediately proves that \(E\sim E'\) implies
\(\Phi(E)=\Phi(E')\). Conversely, suppose that
\(\Phi(E)=\Phi(E')\), and let us prove that \(E\sim E'\). Using
\hyperref[thm:III-1-1]{Theorem~1 of \S~1}, we are reduced to one of
the cases \(E\prec E'\) and \(E'\prec E\). Suppose, for example, that
\(E\prec E'\). Let \((E_1,E_2,\ldots,E_{2^n})\) and
\((E'_1,E'_2,\ldots,E'_{2^n})\) be families of orthogonal projections,
equivalent respectively to \(E\) and \(E'\), each with sum
\(F\in\mathcal Z\). One has \(E_i\sim E\prec E'\sim E'_i\) for
\(1\leq i\leq2^n\); hence there are projections \(E''_i\leq E'_i\)
with \(E''_i\sim E_i\). Consequently
\[
 F=\sum_{i=1}^{2^n}E_i\sim\sum_{i=1}^{2^n}E''_i\leq F.
\]
Thus \(\sum E''_i=F\), by Lemma~1; hence \(E''_i=E'_i\), and finally
\(E_i\sim E'_i\).
\end{proof}

\begin{lemma}\label{lem:III-8-5}
Let \(E,E'\) be fundamental projections of \(\mathcal A\) such that
\[
 \Phi(E)=2^{-p}F,\qquad \Phi(E')=2^{-q}F,\qquad q\geq p.
\]
There exist \(2^{q-p}\) mutually orthogonal projections of
\(\mathcal A\), equivalent to \(E'\), with sum \(E\).
\end{lemma}

\begin{proof}
Let \((E'_i)_{1\leq i\leq2^q}\) be equivalent, orthogonal
projections, equivalent to \(E'\), with sum \(F\). Plainly
\[
 E''=\sum_{i=1}^{2^{q-p}}E'_i
\]
is fundamental and \(\Phi(E'')=2^{-p}F\). Hence
\(E''\sim E\), by Lemma~4. The equality above then gives the
projections asserted in the lemma.
\end{proof}

\begin{lemma}\label{lem:III-8-6}
Let \(E_1,E_2,\ldots,E_n,F_1,F_2,\ldots,F_p,F^{\sim}\) be
fundamental projections, the \(E_i\) (respectively the \(F_j\)) being
pairwise orthogonal. Put
% Source PDF page 301 (printed p. 292)
\phantomsection\label{srcpage:292}
\[
 E=\sum_{i=1}^{n}E_i,\qquad F=\sum_{j=1}^{p}F_j.
\]
Suppose that \(F\leq E\) and
\[
 \Phi(F^{\sim})+\sum_{j=1}^{p}\Phi(F_j)
 \leq\sum_{i=1}^{n}\Phi(E_i).
\]
Then there is a projection \(F_{p+1}\), equivalent to
\(F^{\sim}\), orthogonal to \(F\), and dominated by \(E\).
\end{lemma}

\begin{proof}
Consider the central supports of the \(E_i\), the \(F_j\), and
\(F^{\sim}\). By decomposing \(\mathcal A\) suitably into a product
of von Neumann algebras, we are reduced to the case in which all these
supports are equal to either \(0\) or \(I\). Since zero projections
play no role in the question, we may suppose that all these supports
are \(I\). Write
\[
 \Phi(E_i)=2^{-r_i},\qquad
 \Phi(F_j)=2^{-s_j},\qquad
 \Phi(F^{\sim})=2^{-s}.
\]
Let \(q\) be an integer larger than the \(r_i\), the \(s_j\), and
\(s\). Let \(G\) be a fundamental projection (which will serve as a
``unit of measure'') such that \(\Phi(G)=2^{-q}\)
(\hyperref[cor:III-1-thm1-3]{\S~1, Theorem~1, Corollary~3}). By
Lemma~5, let
\[
 E_i^1,E_i^2,\ldots,E_i^{2^{q-r_i}},\qquad
 F_j^1,F_j^2,\ldots,F_j^{2^{q-s_j}},\qquad
 \widetilde F^1,\widetilde F^2,\ldots,
 \widetilde F^{2^{q-s}}
\]
be projections equivalent to \(G\), orthogonal within each family,
and with sums \(E_i\), \(F_j\), and \(F^{\sim}\), respectively. The
\(E_i^k\) are all pairwise orthogonal and have sum \(E\); the
\(F_j^k\) are all pairwise orthogonal and have sum \(F\). The
hypothesis of the lemma may be written
\begin{equation}
 2^{q-s}+\sum_{j=1}^{p}2^{q-s_j}
 \leq\sum_{i=1}^{n}2^{q-r_i}.
 \tag{1}
\end{equation}
This permits us to choose, in the set \(M\) of all the \(E_i^k\), two
disjoint subsets \(M_1,M_2\), having respectively
\(2^{q-s}\) and \(\sum_j2^{q-s_j}\) elements. The sum
\(\widetilde F'\) of the elements of \(M_1\) is equivalent to
\(F^{\sim}\), while the sum \(F'\) of the elements of \(M_2\) is
equivalent to \(F\); and \(\widetilde F',F'\) are orthogonal and
dominated by \(E\). We have
\[
 F^{\sim}\sim\widetilde F'
 \leq E-F'\sim E-F
\]
(\S~2, no.~3, Remark), and the lemma follows.
\end{proof}

\begin{definition}\label{def:III-8-3}
A function \(E\mapsto\varphi(E)\), with finite nonnegative values, on
the set of fundamental projections of \(\mathcal A\), is called a
\emph{weight} if
% Source PDF page 302 (printed p. 293)
\phantomsection\label{srcpage:293}
it has the following properties:
\begin{enumerate}[label=\textup{(\roman*)},leftmargin=2em]
\item if \(E\) is fundamental and \((E_i)_{i\in I}\) is a family of
orthogonal fundamental projections with sum \(E\), then
\[
 \varphi(E)=\sum_{i\in I}\varphi(E_i);
\]
\item if \(\varphi(E)=0\), then \(E=0\).
\end{enumerate}
The weight is called \emph{central} if \(E_1\sim E_2\) implies
\(\varphi(E_1)=\varphi(E_2)\).
\end{definition}

\begin{lemma}\label{lem:III-8-7}
Let \(x\) be a separating element for \(\mathcal A\). For every
fundamental projection \(E\), put
\[
 \varphi(E)=(\Phi(E)x\mid x).
\]
Then \(\varphi\) is a central weight.
\end{lemma}

\begin{proof}
By Lemma~4, \(E_1\sim E_2\) implies
\(\varphi(E_1)=\varphi(E_2)\). Since \(x\) is separating for
\(\mathcal A\), the relation \(\varphi(E)=0\) implies
\(\Phi(E)=0\), and hence \(E=0\). Finally, let \(E\) be a fundamental
projection and let \((E_i)_{i\in I}\) be a family of orthogonal
fundamental projections with sum \(E\). We shall prove, thereby
completing the proof of the lemma, that
\(\varphi(E)=\sum_i\varphi(E_i)\); for this it is enough to prove that
\(\sum_i\Phi(E_i)\) exists in \(\mathcal Z\) and is equal to
\(\Phi(E)\).

First, for every finite subset \(J\) of \(I\),
\[
 \sum_{i\in J}\Phi(E_i)\leq\Phi(E).
\]
Indeed, otherwise, by suitably decomposing \(\mathcal A\) into a
product of von Neumann algebras, we could reduce to the case
\[
 \Phi(E)+2^{-p}\leq\sum_{i\in J}\Phi(E_i).
\]
Let \(G\) be a fundamental projection such that
\(\Phi(G)=2^{-p}\). Applying Lemma~6 twice, there would be orthogonal
projections \(E',G'\), equivalent respectively to \(E,G\), and
dominated by \(\sum_{i\in J}E_i\). Hence
\[
 E\sim E'\geq\sum_{i\in I}E_i\geq E'+G',
\]
which is absurd by Lemma~1. Thus
\(\sum_{i\in J}\Phi(E_i)\leq\Phi(E)\). Consequently
\(\sum_{i\in I}\Phi(E_i)\) exists and is dominated by \(\Phi(E)\).

We shall now prove that
\(\Phi(E)\leq\sum_{i\in I}\Phi(E_i)\). If this were not so,
% Source PDF page 303 (printed p. 294)
\phantomsection\label{srcpage:294}
by suitably decomposing \(\mathcal A\) into a product of von Neumann
algebras, we could reduce to the case
\[
 2^{-p}+\sum_{i\in I}\Phi(E_i)\leq\Phi(E).
\]
Let \(E_0\) be a fundamental projection such that
\(\Phi(E_0)=2^{-p}\). Since zero projections play no role in the
question, we may suppose that the \(E_i\) are nonzero, and hence form
a sequence \((E_1,E_2,\ldots)\), because \(\mathcal A\) has a
separating element. By repeated application of Lemma~6, there would
then be a sequence \((E'_0,E'_1,E'_2,\ldots)\) of orthogonal
projections dominated by \(E\), with \(E'_i\sim E_i\) for
\(i=0,1,2,\ldots\). We would therefore have
\[
 E=\sum_{i=1}^{\infty}E_i
 \sim\sum_{i=1}^{\infty}E'_i\leq E,
 \qquad
 E-\sum_{i=1}^{\infty}E'_i\geq E'_0\ne0,
\]
which is absurd by Lemma~1.
\end{proof}

\noindent Bibliography: \ArticleRef{6}, \ArticleRef{65}.

\NumberedParagraph{4}{Construction of a Trace to Within
\(\boldsymbol{\varepsilon}\)}
\label{no:III-8-4}

\begin{lemma}\label{lem:III-8-8}
Let \(\varphi\) and \(\psi\) be two weights, and let
\(\varepsilon>0\). There exist a number \(\theta>0\) and a nonzero
fundamental projection \(E\) such that, for every fundamental
projection \(F\leq E\),
\[
 \theta\varphi(F)\leq\psi(F)
 \leq\theta(1+\varepsilon)\varphi(F).
\]
\end{lemma}

\begin{proof}
We may suppose that \(\varphi(I)=\psi(I)\). Suppose that, for every
nonzero fundamental projection \(G_1\), there were a fundamental
projection \(G_2\leq G_1\) such that
\(\varphi(G_2)>\psi(G_2)\). There would then be (by Zorn's theorem and
Lemma~2) a family \((G_i)\) of orthogonal fundamental projections with
sum \(I\), such that \(\varphi(G_i)>\psi(G_i)\). Hence
\[
 \varphi(I)=\sum_i\varphi(G_i)>
 \sum_i\psi(G_i)=\psi(I),
\]
which is absurd. Thus there is a nonzero fundamental projection \(G\)
such that \(\varphi(G_1)\leq\psi(G_1)\) for every fundamental
\(G_1\leq G\).

Let \(\theta\) be the supremum of the real numbers \(\eta\) such that
\(\eta\varphi(G_1)\leq\psi(G_1)\) for every fundamental
\(G_1\leq G\). Then \(1\leq\theta<+\infty\), and
\(\theta\varphi(G_1)\leq\psi(G_1)\) for every fundamental
\(G_1\leq G\). Suppose that every nonzero fundamental projection
\(G_1\leq G\) majorized a nonzero fundamental projection \(G_2\)
such that
\(\theta(1+\varepsilon)\varphi(G_2)\leq\psi(G_2)\). Reasoning as
above, we would deduce that
\(\theta(1+\varepsilon)\varphi(G_1)\leq\psi(G_1)\),
% Source PDF page 304 (printed p. 295)
\phantomsection\label{srcpage:295}
contrary to the definition of \(\theta\). Hence there is a nonzero
fundamental projection \(E\leq G\) such that, for every fundamental
projection \(F\leq E\),
\(\psi(F)\leq\theta(1+\varepsilon)\varphi(F)\). On the other hand,
\(\theta\varphi(F)\leq\psi(F)\), since \(F\leq G\).
\end{proof}

\begin{definition}\label{def:III-8-4}
Let \(\varepsilon>0\). A positive linear form \(\varphi\) on
\(\mathcal A\) is called a \emph{trace to within \(\varepsilon\)} if,
for every \(T\in\mathcal A_+\) and every unitary operator
\(U\in\mathcal A\),
\[
 \varphi(UTU^{-1})\leq(1+\varepsilon)\varphi(T).
\]
\end{definition}

If \(S\) is any element of \(\mathcal A\), it follows that
\[
 \varphi(SS^*)\leq(1+\varepsilon)\varphi(S^*S).
\]
Indeed, one has \(S=US_1\), with \(U\) unitary and \(S_1\) Hermitian,
\(U,S_1\in\mathcal A\) (by the remark in \S~2, no.~3), and therefore
\(SS^*=US_1^2U^{-1}\), \(S^*S=S_1^2\). Conversely, if
\(\varphi(SS^*)\leq(1+\varepsilon)\varphi(S^*S)\) for every
\(S\in\mathcal A\), then, for every \(T\in\mathcal A_+\) and every
unitary operator \(U\in\mathcal A\),
\begin{align*}
 \varphi(UTU^{-1})
 &=\varphi\bigl((UT^{1/2})(UT^{1/2})^*\bigr)\\
 &\leq(1+\varepsilon)
 \varphi\bigl((UT^{1/2})^*(UT^{1/2})\bigr)
 =(1+\varepsilon)\varphi(T).
\end{align*}

\begin{lemma}\label{lem:III-8-9}
For every \(\varepsilon>0\), there exist a nonzero fundamental
projection \(E\in\mathcal A\) and a nonzero normal positive linear
form \(\omega\) on \(\mathcal A_E\) which is a trace to within
\(\varepsilon\).
\end{lemma}

\begin{proof}
Let \(\varphi\) be a faithful normal positive linear form on
\(\mathcal A\) [one exists by assumption
\eqref{eq:III-8-triple-star}], and
let \(\psi\) be a central weight on the set of fundamental projections
(Lemma~7). By Lemma~8 there is a nonzero fundamental projection \(E\)
such that, for every fundamental projection \(F\leq E\),
\[
 \psi(F)\leq\varphi(F)\leq(1+\varepsilon)\psi(F)
\]
(after multiplying \(\varphi\), if necessary, by a suitable scalar).
The restriction of \(\varphi\) to \(E\mathcal A E\) defines a nonzero
normal positive linear form \(\omega\) on \(\mathcal A_E\). Let
\(T\in\mathcal A_+\), with \(ET=TE=T\), and let \(U\) be a unitary
operator of \(\mathcal A\) commuting with \(E\). We shall prove that
\begin{equation}
 \varphi(UTU^{-1})\leq(1+\varepsilon)\varphi(T),
 \tag{2}\label{eq:III-8-trace-approximation}
\end{equation}
which will establish that \(\omega\) is a trace to within
\(\varepsilon\). If \(T\) is a fundamental projection dominated by
\(E\), then \(UTU^{-1}\) is a fundamental projection dominated by
\(E\) and equivalent to \(T\), and hence
\[
 \varphi(UTU^{-1})
 \leq(1+\varepsilon)\psi(UTU^{-1})
 =(1+\varepsilon)\psi(T)
 \leq(1+\varepsilon)\varphi(T).
\]
If \(T\) is an arbitrary projection of \(\mathcal A\) dominated by
\(E\), write \(T=\sum_iT_i\),
% Source PDF page 305 (printed p. 296)
\phantomsection\label{srcpage:296}
where the \(T_i\) are orthogonal fundamental projections (Corollary
to Lemma~2); then
\[
 \varphi(T)=\sum_i\varphi(T_i),\qquad
 \varphi(UTU^{-1})=\sum_i\varphi(UT_iU^{-1}),
\]
so that \eqref{eq:III-8-trace-approximation} is still true in this
case. Finally, if \(T\) is an
arbitrary element of \(\mathcal A_+\) such that \(ET=TE=T\), then
\(T\) is the norm limit of operators
\(\sum_{i=1}^{n}\lambda_iT_i\), where the \(T_i\) are projections of
\(\mathcal A\) dominated by \(E\) and the \(\lambda_i\) are
nonnegative scalars; thus \eqref{eq:III-8-trace-approximation} follows
once more by passage to the
limit.
\end{proof}

\begin{lemma}\label{lem:III-8-10}
For every \(\varepsilon>0\), there exists a nonzero normal positive
linear form \(\varphi\) on \(\mathcal A\) which is a trace to within
\(\varepsilon\).
\end{lemma}

\begin{proof}
Let \(E\) and \(\omega\) be the projection and the linear form whose
existence is assured by Lemma~9. Let
\(E_1,E_2,\ldots,E_p\) be equivalent orthogonal projections, with sum
\(F\in\mathcal Z\), and with \(E_1=E\). Let \(W_i\) be a partially
isometric operator having \(E\) as initial projection and \(E_i\) as
final projection. The form \(\omega\) is identified with a linear form
on the set \(\mathcal B\) of all \(T\in\mathcal A\) such that
\(ET=TE=T\). For \(T\in\mathcal A\), put
\[
 \varphi(T)=\sum_{i=1}^{p}\omega(W_i^*TW_i).
\]
Since
\[
 \sum_{j=1}^{p}W_jW_j^*=\sum_{j=1}^{p}E_j=F,
\]
putting \(T_{ij}=W_j^*TW_i\in\mathcal B\), we obtain
\begin{align*}
 \varphi(T^*T)
 &=\sum_{i,j}\omega(W_i^*T^*W_jW_j^*TW_i)
   =\sum_{i,j}\omega(T_{ji}^*T_{ji}),\\
 \varphi(TT^*)
 &=\sum_{i,j}\omega(W_i^*TW_jW_j^*T^*W_i)
   =\sum_{i,j}\omega(T_{ij}T_{ij}^*)
   =\sum_{i,j}\omega(T_{ji}T_{ji}^*).
\end{align*}
It follows that
\[
 \varphi(T^*T)\leq(1+\varepsilon)\varphi(TT^*).
\]
Moreover, \(\varphi(E)=\omega(E)\ne0\).
\end{proof}

\noindent Bibliography: \ArticleRef{6}, \ArticleRef{66},
\ArticleRef{140}, \ArticleRef{262}, \ArticleRef{303}.

\NumberedParagraph{5}{End of the Proof of the Theorem}
\label{no:III-8-5}
We use the notation \(\mathcal E,\mathcal K_T,\mathcal K'_T,
\mathcal K''_T\) of \S~5.

\begin{lemma}\label{lem:III-8-11}
For every \(T\in\mathcal A\), \(\mathcal K''_T\) consists of one
point.
\end{lemma}

% Source PDF page 306 (printed p. 297)
\phantomsection\label{srcpage:297}
\begin{proof}
Suppose first that \(T\) is Hermitian. We may plainly restrict
ourselves to the case \(0\leq T\leq I\). Let \(S_1,S_2\) be elements
of \(\mathcal K''_T\), and suppose that \(S_1\ne S_2\). There exist a
nonzero projection \(E\) of \(\mathcal Z\) and a scalar \(\alpha>0\)
such that, for example,
\(S_1E-S_2E\geq\alpha E\). We shall deduce a contradiction. Replacing
\(\mathcal A\) by \(\mathcal A_E\), it is enough to deduce a
contradiction from the hypothesis
\(S_1-S_2\geq\alpha I\). Let \(\varphi\) be a nonzero positive linear
form on \(\mathcal A\) which is a trace to within \(\varepsilon\). For
every \(f\in\mathcal E\),
\[
 \lvert\varphi(f\mathbin{\cdot}T)-\varphi(T)\rvert
 \leq\varepsilon\varphi(T).
\]
Consequently
\[
 \lvert\varphi(S_1)-\varphi(T)\rvert\leq\varepsilon\varphi(T),
 \qquad
 \lvert\varphi(S_2)-\varphi(T)\rvert\leq\varepsilon\varphi(T),
\]
and therefore
\[
 2\varepsilon\varphi(T)
 \geq\lvert\varphi(S_1-S_2)\rvert
 \geq\varphi(\alpha I),
\]
which is absurd when \(\varepsilon\) is sufficiently small.

Now let \(T\in\mathcal A\) be arbitrary, and write
\(T=T_1+iT_2\), with \(T_1,T_2\) Hermitian. Put
\(\mathcal K''_{T_1}=\{S_1\}\) and
\(\mathcal K''_{T_2}=\{S_2\}\). Let \(S\in\mathcal K''_T\). For
every \(\varepsilon>0\), there is an \(f\in\mathcal E\) such that
\(\lVert f\mathbin{\cdot}T-S\rVert\leq\varepsilon\); hence
\(\lVert f\mathbin{\cdot}T^*-S^*\rVert\leq\varepsilon\), and therefore
\[
 \left\lVert f\mathbin{\cdot}T_1-\frac12(S+S^*)\right\rVert
 \leq\varepsilon,
 \qquad
 \left\lVert f\mathbin{\cdot}T_2-\frac1{2i}(S-S^*)\right\rVert
 \leq\varepsilon.
\]
Since \(\varepsilon>0\) is arbitrary, it follows that
\[
 \frac12(S+S^*)=S_1,
 \qquad
 \frac1{2i}(S-S^*)=S_2,
\]
and hence \(S=S_1+iS_2\).
\end{proof}

We shall henceforth denote by \(T^\natural\) the point to which
\(\mathcal K''_T\) reduces.

\begin{lemma}\label{lem:III-8-12}
The mapping \(T\mapsto T^\natural\) of \(\mathcal A\) into
\(\mathcal Z\) is positive and linear. For all
\(T_1,T_2\in\mathcal A\),
\[
 (T_1T_2)^\natural=(T_2T_1)^\natural.
\]
\end{lemma}

\begin{proof}
Linearity follows from \hyperref[lem:III-5-6]{Lemma~6 of \S~5}. If
\(T\geq0\), then \(T_1\geq0\) for every \(T_1\in\mathcal K'_T\), and
hence \(T^\natural\geq0\). If \(T\in\mathcal A\) and \(U\) is a
unitary operator of \(\mathcal A\), then
\[
 \mathcal K''_{UT}
 =\mathcal K''_{U^{-1}(UT)U}=\mathcal K''_{TU},
\]
and therefore \((UT)^\natural=(TU)^\natural\). Since every element of
\(\mathcal A\) is a linear combination of unitary elements of
\(\mathcal A\), the asserted relation follows.
\end{proof}

We can now prove Theorem~1. Let \(\varphi\) be a nonzero normal
positive linear form on \(\mathcal A\) which is a trace to within
\(\varepsilon\). For \(T\in\mathcal A\), put
\(\varphi'(T)=\varphi(T^\natural)\). By Lemma~12, \(\varphi'\) is a
trace on \(\mathcal A\), and
\(\varphi'(I)=\varphi(I)\ne0\). Finally, let \(T\in\mathcal A_+\).
For every \(T_1\in\mathcal K_T\),
\(\varphi(T_1)\leq(1+\varepsilon)\varphi(T)\), and hence
\(\varphi'(T)=\varphi(T^\natural)\leq
(1+\varepsilon)\varphi(T)\). Since \(\varphi\) is normal, it follows
that \(\varphi'\) is normal. The support \(F\)
% Source PDF page 307 (printed p. 298)
\phantomsection\label{srcpage:298}
of \(\varphi'\) is a nonzero projection of \(\mathcal Z\), and
\(\varphi'\) defines a faithful normal trace on \(\mathcal A_F\), so
that \(\mathcal A_F\) is finite. This completes the proof.

There is probably a shorter proof of Theorem~1. It remains to be
discovered.

\noindent Bibliography: \ArticleRef{6}, \ArticleRef{66},
\ArticleRef{140}, \ArticleRef{262}, \ArticleRef{303}.

\NumberedParagraph{6}{Consequences of the Theorem}
\label{no:III-8-6}

\begin{corollary}\label{cor:III-8-thm1-1}
Let \(\mathcal A\) be a von Neumann algebra and let \(E\) be a
projection of \(\mathcal A\). In order that \(E\) be finite, it is
necessary and sufficient that every projection of \(\mathcal A\)
dominated by \(E\) and equivalent to \(E\) be equal to \(E\).
\end{corollary}

\begin{proof}
It is enough to apply Theorem~1 to \(\mathcal A_E\).
\end{proof}

This gives a new definition of finite projections. Propositions~9
and~11 of \S~2 then furnish a new definition of properly infinite
projections and of purely infinite projections.

\begin{corollary}\label{cor:III-8-thm1-2}
Let \(E\) be a properly infinite projection of \(\mathcal A\). There
exist pairwise orthogonal projections
\(E^1,E^2,E^3,\ldots\), equivalent to \(E\), with sum \(E\).
\end{corollary}

\begin{proof}
Suppose first only that \(E\) is infinite. Let
\(E_2\leq E=E_1\), with \(E_2\ne E_1\) and \(E_2\sim E_1\). Let
\(U\) be a partially isometric operator of \(\mathcal A\) having
\(E_1\) and \(E_2\) as initial and final projections. By induction,
\(U^i\) is a partially isometric operator with initial projection
\(E_1\) and final projection \(E_{i+1}\), where
\(E_1\geq E_2\geq E_3\geq\cdots\). Put \(F_i=E_i-E_{i+1}\). The
restriction of \(U^i\) to \(F_1(\mathfrak H)\) maps
\(F_1(\mathfrak H)\) isometrically onto \(F_{i+1}(\mathfrak H)\), so
\[
 0\ne F_1\sim F_2\sim F_3\sim\cdots.
\]
By \hyperref[cor:III-1-thm1-2]{Corollary~2 of Theorem~1 in \S~1},
applied in \(\mathcal A_E\) to the family
\((F_1,F_2,\ldots)\), there exist a projection \(F\) of
\(\mathcal Z\) and an infinite family \((G_i)_{i\in I}\) of nonzero,
equivalent, pairwise orthogonal projections of \(\mathcal A\), with
sum \(EF\). We may write
\(I=I_1\mathbin{\dot\cup}I_2\mathbin{\dot\cup}\cdots\), where the
\(I_j\) are disjoint and equipotent to \(I\). Then
\[
 \sum_{i\in I_1}G_i,\quad
 \sum_{i\in I_2}G_i,\quad\ldots
\]
are nonzero, pairwise orthogonal projections, equivalent to
\(\sum_{i\in I}G_i\). Thus
\(EF=G^1+G^2+\cdots\), where the \(G^i\) are projections of
\(\mathcal A\), orthogonal, nonzero, and equivalent to \(EF\).

Finally, suppose that \(E\) is properly infinite. Applying
\hyperref[prop:III-2-9]{Proposition~9 of \S~2} and Zorn's theorem, we
construct a family \((F_k)_{k\in K}\) of orthogonal projections of
\(\mathcal Z\) such that \(E=\sum_{k\in K}EF_k\) and, for every
\(k\in K\), there are orthogonal projections
\(G_k^1,G_k^2,\ldots\), equivalent to \(EF_k\), with
% Source PDF page 308 (printed p. 299)
\phantomsection\label{srcpage:299}
sum \(EF_k\). The projections
\[
 E^i=\sum_{k\in K}G_k^i
\]
are the projections asserted in the corollary.
\end{proof}

\begin{corollary}\label{cor:III-8-thm1-3}
Let \(\mathcal A\) be a von Neumann algebra with centre
\(\mathcal Z\). If there is a linear mapping \(\Phi\) of
\(\mathcal A\) into \(\mathcal Z\) such that
\[
 \Phi(T)=T\quad(T\in\mathcal Z),
\]
and such that \(\Phi(S_1S_2)=\Phi(S_2S_1)\) for
\(S_1,S_2\in\mathcal A\), then \(\mathcal A\) is finite.
\end{corollary}

\begin{proof}
Let \(E\) be the largest projection of \(\mathcal Z\) properly
infinite relative to \(\mathcal A\). By Corollary~2,
\(E=E_1+E_2\), where \(E_1,E_2\) are orthogonal projections of
\(\mathcal A\), and \(E\sim E_1\sim E_2\). Thus
\(E=U^*U\), \(E_1=UU^*\), for some \(U\in\mathcal A\), and hence
\(\Phi(E)=\Phi(E_1)\). Similarly \(\Phi(E)=\Phi(E_2)\). Consequently
\[
 2E=2\Phi(E)=\Phi(E_1)+\Phi(E_2)=\Phi(E)=E,
\]
and hence \(E=0\).
\end{proof}

\begin{corollary}\label{cor:III-8-thm1-4}
Let \(\mathcal A\) be a von Neumann algebra with centre
\(\mathcal Z\). If, for every \(T\in\mathcal A\), the set
\(\mathcal K'_T\) of \S~5 meets \(\mathcal Z\) in exactly one point,
then \(\mathcal A\) is finite.
\end{corollary}

\begin{proof}
Let \(\Phi(T)\) be the point at which \(\mathcal K'_T\) meets
\(\mathcal Z\). The mapping \(\Phi\) of \(\mathcal A\) into
\(\mathcal Z\) is linear (\hyperref[lem:III-5-6]{\S~5, Lemma~6}).
Clearly \(\Phi(T)=T\) for \(T\in\mathcal Z\). On the other hand, if
\(S_1,S_2\in\mathcal A\) and \(S_1\) is unitary, then
\[
 \mathcal K'_{S_1S_2}
 =\mathcal K'_{S_1(S_2S_1)S_1^{-1}}
 =\mathcal K'_{S_2S_1},
\]
and therefore \(\Phi(S_1S_2)=\Phi(S_2S_1)\). By linearity, this
equality extends to the case of arbitrary \(S_1\in\mathcal A\).
The result now follows from Corollary~3.
\end{proof}

\begin{corollary}\label{cor:III-8-thm1-5}
Let \(\mathcal A\) be a von Neumann algebra, let \(E,F\) be
projections of \(\mathcal A\), and let \(E',F'\) be their central
supports. Suppose that \(\mathcal A_E\) is of countable type, that
\(F\) is properly infinite, and that \(E'\leq F'\). Then
\(E\prec F\).

In particular, in a factor of countable type, any two infinite
projections are equivalent.
\end{corollary}

\begin{proof}
Consider a maximal family \((E_i)_{i\in I}\) of nonzero orthogonal
projections of \(\mathcal A\), dominated by \(E\), and such that
\(E_i\prec F\). The central support of
\(E-\sum_{i\in I}E_i\) is dominated by that of \(F\), and
\hyperref[lem:III-1-1]{Lemma~1 of \S~1} shows that
\(E-\sum_{i\in I}E_i=0\). Thus \(E=\sum_{i\in I}E_i\). Since
\(\mathcal A_E\) is of
% Source PDF page 309 (printed p. 300)
\phantomsection\label{srcpage:300}
countable type, \(I\) is countable. On the
other hand, by Corollary~2 there is an infinite family
\((F_k)_{k\in K}\) of equivalent orthogonal projections of
\(\mathcal A\), with sum \(F\). Since \(I\) is equipotent to a subset
of \(K\),
\[
 E=\sum_{i\in I}E_i\prec\sum_{k\in K}F_k=F.
\]
\end{proof}

\begin{corollary}\label{cor:III-8-thm1-6}
Let \(\mathcal A\) be a purely infinite von Neumann algebra of
countable type, and let \(T\) be a nonzero element of \(\mathcal A\).
Then the set \(\mathcal K'_T\) of \S~5 meets the centre
\(\mathcal Z\) of \(\mathcal A\) in at least one nonzero point.
\end{corollary}

\begin{proof}
We reduce immediately to the case in which \(T\) is Hermitian and
\(\lVert T\rVert\leq1\). Replacing \(T\) by \(-T\), if necessary, we
may suppose that there are a nonzero spectral projection \(E\) of
\(T\) and an integer \(n>0\) such that \(TE\geq n^{-1}E\); hence
\[
 T\geq n^{-1}E-(I-E).
\]
Let \(F\) be the central support of \(E\). Writing
\(\mathcal A=\mathcal A_F\times\mathcal A_{I-F}\), we reduce to the
case \(F=I\). If \(E\) majorizes a central projection of
\(\mathcal A\), the corollary is immediate. We may therefore suppose
that \(E\) majorizes no central projection of \(\mathcal A\), so that
the central support of \(I-E\) is \(I\). By Corollary~5,
\(E\sim I-E\). By Corollary~2, there are pairwise orthogonal
projections \(E_0,E_1,\ldots,E_n\) of \(\mathcal A\), with sum \(E\),
equivalent to \(E\). Put \(I-E=E_{n+1}\). There are unitary operators
\(U_0,U_1,\ldots,U_{n+1}\) of \(\mathcal A\) such that
\[
 U_jE_iU_j^{-1}=E_{\sigma_j(i)},
\]
where the \(\sigma_j\) are the \(n+2\) cyclic permutations of
\(0,1,\ldots,n+1\). For every \(i\), one has
\[
 \sum_jU_jE_iU_j^{-1}=I.
\]
Put
\[
 S=(n+2)^{-1}\sum_{j=0}^{n+1}U_jTU_j^{-1}.
\]
Then
\begin{align*}
 (n+2)S
 &\geq\sum_{j=0}^{n+1}U_j
 \bigl(n^{-1}(E_0+E_1+\cdots+E_n)-E_{n+1}\bigr)U_j^{-1}\\
 &=(n+1)n^{-1}I-I=n^{-1}I.
\end{align*}
Thus \(S_1\geq[n(n+2)]^{-1}I\) for every
\(S_1\in\mathcal K_S\), and therefore for every
\(S_1\in\mathcal K'_S\). Hence
\(\mathcal K'_S\cap\mathcal Z\) contains a nonzero point. But
\(S\in\mathcal K_T\), so \(\mathcal K'_S\subset\mathcal K'_T\).
\end{proof}

\begin{corollary}\label{cor:III-8-thm1-7}
Let \(\mathcal A\) be a factor, let \(M\) be the set of projections of
\(\mathcal A\), and let \(D\) be a function on \(M\) with the
following properties:
\begin{enumerate}[label=\textup{(\roman*)},leftmargin=2em]
\item \(0=D(0)<D(E)\leq+\infty\) for every nonzero
\(E\in M\);
\item if \(E,F\in M\) are orthogonal, then
\(D(E+F)=D(E)+D(F)\);
\item if \(U\) is a unitary operator of \(\mathcal A\) and
\(E\in M\), then \(D(UEU^{-1})=D(E)\).
\end{enumerate}%
% Source PDF page 310 (printed p. 301)
\phantomsection\label{srcpage:301}%
Then \(D\) is a relative dimension on \(\mathcal A\), or is
identically infinite on the set of nonzero projections of
\(\mathcal A\).
\end{corollary}

\begin{proof}
First let us prove that, if \(E\in M\) is infinite, then
\(D(E)=+\infty\). By Corollary~2, there are orthogonal projections
\(E_1,E_2,E_3\) of \(\mathcal A\) such that
\[
 E=E_1+E_2+E_3,
 \qquad E\sim E_1\sim E_2\sim E_3.
\]
One has
\(E_1\prec E_1+E_2\prec E\sim E_1\), so
\(E_1+E_2\sim E_1\); similarly \(E_2+E_3\sim E_3\), and hence
\(I-E_1\sim I-(E_1+E_2)\). Consequently there is a unitary operator
\(U\in\mathcal A\) such that \(UE_1U^{-1}=E_1+E_2\). Thus
\[
 D(E_1)=D(E_1+E_2)=D(E_1)+D(E_2),
\]
and consequently \(D(E_1)=+\infty\). Since \(D(E)\geq D(E_1)\),
we have \(D(E)=+\infty\).

We now prove that if \(D\) is finite for one finite projection
\(E_0\), then \(D\) is finite for every finite projection \(E\). By
\hyperref[cor:III-1-thm1-2]{\S~1, Theorem~1, Corollary~2}, there are
pairwise orthogonal projections \(E_1,E_2,\ldots,E_n\) of
\(\mathcal A\), with sum \(E\), such that
\[
 E_0\sim E_1\sim E_2\sim\cdots\sim E_{n-1}\succ E_n.
\]
Taking account of \hyperref[prop:III-2-6]{Proposition~6 of \S~2}, we
then have
\[
 D(E)=\sum_{i=1}^{n}D(E_i)\leq nD(E_0)<+\infty.
\]

This being established, the assertion of the corollary is immediate
if \(\mathcal A\) is purely infinite. If \(\mathcal A\) is not purely
infinite, let \(E,F\) be finite projections of \(\mathcal A\), and
let \(D'\) be a relative dimension on \(\mathcal A\). We shall prove,
thereby completing the proof, that
\[
 \frac{D(E)}{D(F)}=\frac{D'(E)}{D'(F)}.
\]
Let \(G\) be the supremum of \(E\) and \(F\); it is finite. The
restriction of \(D\) to the projections of \(G\mathcal A G\) defines,
on the set of projections of \(\mathcal A_G\), a function to which
\hyperref[prop:III-2-15]{Proposition~15 of \S~2} applies. Thus \(D\)
and \(D'\) are proportional on the projections of
\(G\mathcal A G\).
\end{proof}

\begin{corollary}\label{cor:III-8-thm1-8}
Let \(\mathcal A\) and \(\mathcal B\) be two von Neumann algebras.
Suppose that \(\mathcal A'\) (respectively \(\mathcal B'\)) contains
an infinite family \((E'_i)_{i\in I}\) (respectively
\((F'_i)_{i\in I}\)) of equivalent, pairwise orthogonal projections,
with sum \(I\), such that the \(\mathcal A'_{E'_i}\) (respectively
the \(\mathcal B'_{F'_i}\)) are of countable type (as is the case if
\(\mathcal A'\) and \(\mathcal B'\) are of countable type and
properly infinite). Then every isomorphism \(\Phi\) of \(\mathcal A\)
onto \(\mathcal B\) is spatial.
\end{corollary}

% Source PDF page 311 (printed p. 302)
\phantomsection\label{srcpage:302}
\begin{proof}
Let \((I_j)_{j\in K}\) be a partition of \(I\), the sets \(I_j\)
being countably infinite. Put \(E'_j=\sum_{i\in I_j}E'_i\). The
\(E'_j\) are properly infinite (\hyperref[prop:III-2-10]{\S~2,
Proposition~10}), equivalent, pairwise orthogonal, with sum \(I\), and
the \(\mathcal A'_{E'_j}\) are of countable type. Replacing the
\(E'_i\) by the \(E'_j\), we may henceforth suppose that the \(E'_i\)
are properly infinite; and similarly that the \(F'_i\) are properly
infinite.

There exist a von Neumann algebra \(\mathcal C\) and projections
\(E',F'\in\mathcal C'\), with central support \(I\), such that
\(\mathcal A\) is identified with \(\mathcal C_{E'}\),
\(\mathcal B\) with \(\mathcal C_{F'}\), and \(\Phi\) with the
isomorphism \(T_{E'}\mapsto T_{F'}\) \((T\in\mathcal C)\). By
Corollary~5, \(E'_i\sim F'_i\), and hence
\[
 E'=\sum_{i\in I}E'_i\sim\sum_{i\in I}F'_i=F'.
\]
Thus \(\Phi\) is implemented by a partially isometric operator of
\(\mathcal C'\).
\end{proof}

\begin{corollary}\label{cor:III-8-thm1-9}
Let \(\mathcal A\) be a von Neumann algebra, let \(\mathcal Z\) be
its centre, and let \(\Phi\) be an automorphism of \(\mathcal A\)
which fixes the elements of \(\mathcal Z\). If \(\mathcal A'\) is
properly infinite, then \(\Phi\) is spatial.
\end{corollary}

\begin{proof}
If \(\mathcal A=\prod_{i\in I}\mathcal A_i\), then \(\Phi\) induces
automorphisms \(\Phi_i\) of the \(\mathcal A_i\), and it is enough to
prove the corollary for the \(\mathcal A_i\) and \(\Phi_i\). By
\hyperref[lem:III-1-7]{Lemma~7 of \S~1}, we may then confine ourselves
to the following two cases: \(1^\circ\) \(\mathcal A'\) is of
countable type (and properly infinite); \(2^\circ\) there is in
\(\mathcal A'\) an infinite family \((E'_i)_{i\in I}\) of equivalent,
pairwise orthogonal projections, with sum \(I\), such that the
\(\mathcal A'_{E'_i}\) are of countable type. In both cases,
Corollary~8 proves that \(\Phi\) is spatial.
\end{proof}

\begin{corollary}\label{cor:III-8-thm1-10}
Let \(\mathcal A\) be a von Neumann algebra. If \(\mathcal A'\) is
properly infinite, every ultra-weakly continuous linear form on
\(\mathcal A\) is a form \(\omega_{x,y}\).
\end{corollary}

\begin{proof}
Let \((E'_1,E'_2,\ldots)\) be a sequence of pairwise orthogonal
projections of \(\mathcal A'\), equivalent to \(I\), with sum \(I\)
(Corollary~2). The induction \(T\mapsto T_{E'_1}\) of
\(\mathcal A\) onto \(\mathcal A_{E'_1}\) is a spatial isomorphism
\(\Phi\). Put \(\mathfrak H_1=E'_1(\mathfrak H)\). We may identify
\(\mathfrak H\) with \(\mathfrak H_1\otimes\mathfrak K\), where
\(\mathfrak K\) is an infinite-dimensional Hilbert space with a
countable basis, and \(\Phi^{-1}\) with the amplification
\[
 T_{E'_1}\longmapsto T_{E'_1}\otimes I_{\mathfrak K}
\]
of \(\mathcal A_{E'_1}\) onto
\(\mathcal A=\mathcal A_{E'_1}\otimes
\mathbb C I_{\mathfrak K}\). The corollary then follows from
\hyperref[lem:I-4-5]{Chapter~I, \S~4, Lemma~5}.
\end{proof}

\begin{corollary}\label{cor:III-8-thm1-11}
Suppose that \(\mathcal A'\) is properly infinite. In order that
\(\mathcal A\) have a separating element, it is necessary and
sufficient that \(\mathcal A\) be of countable type.
\end{corollary}

% Source PDF page 312 (printed p. 303)
\phantomsection\label{srcpage:303}
\begin{proof}
The condition is plainly necessary. Conversely, if \(\mathcal A\) is
of countable type, there is a faithful normal positive linear form on
\(\mathcal A\). By Corollary~10 this form is a form \(\omega_x\).
Then \(x\) is separating for \(\mathcal A\).
\end{proof}

\noindent Bibliography: \ArticleRef{6}, \ArticleRef{10},
\ArticleRef{15}, \ArticleRef{31}, \ArticleRef{42}, \ArticleRef{62},
\ArticleRef{65}, \ArticleRef{66}, \ArticleRef{89}, \ArticleRef{404},
\ArticleRef{414}.

\NumberedParagraph{7}{Further Results on Tensor Products}
\label{no:III-8-7}

\begin{lemma}\label{lem:III-8-13}
Let \(\mathcal A\) be a von Neumann algebra, let \(\varphi\) be a
faithful normal trace on \(\mathcal A_+\), and let
\(S\in\mathcal A\) satisfy \(\varphi(S^*S)<+\infty\). The mapping
\[
 T\longmapsto ST^*
\]
of \(\mathcal A\) into \(\mathcal A\) is strongly continuous on the
bounded subsets of \(\mathcal A\).
\end{lemma}

\begin{proof}
For \(T\in\mathcal A\), put
\(\psi(T)=\varphi(S^*ST)\). Then \(\psi\) is ultra-weakly continuous
(\hyperref[prop:I-6-1]{Chapter~I, \S~6, Proposition~1}). There is a
family \((x_i)\) of vectors such that
\(\varphi=\sum_i\omega_{x_i}\) on \(\mathcal A_+\) (Chapter~I,
\S~6, corollary to Proposition~2). For every \(T\in\mathcal A\),
\begin{align*}
 \sum_i\lVert ST^*x_i\rVert^2
 &=\sum_i\omega_{x_i}(TS^*ST^*)
 =\varphi(TS^*ST^*)\\
 &=\varphi(S^*ST^*T)=\psi(T^*T).
\end{align*}
When \(T\) tends strongly to zero while remaining bounded,
\(T^*T\) tends ultra-weakly to zero; hence \(\psi(T^*T)\to0\), and
therefore \(ST^*x_i\to0\) for every \(i\). Now the \(x_i\) form a
separating set for \(\mathcal A\), since \(\varphi\) is faithful.
Since \(ST^*\) remains bounded, it follows that \(ST^*\) tends
strongly to zero.
\end{proof}

\begin{lemma}\label{lem:III-8-14}
Let \(\mathcal A\) be a von Neumann algebra and let \(E\) be a
nonzero properly infinite projection of \(\mathcal A\). The mapping
\(T\mapsto T^*\) of \(E\mathcal A E\) into \(E\mathcal A E\) is not
strongly continuous on the bounded subsets of \(E\mathcal A E\).
\end{lemma}

\begin{proof}
There are equivalent pairwise orthogonal projections
\(E_1,E_2,\ldots\) of \(\mathcal A\), with sum \(E\) (Corollary~2 of
Theorem~1). By
\hyperref[prop:I-2-5]{Chapter~I, \S~2, Proposition~5(ii)},
\(\mathcal A_E\) is spatially isomorphic to
\[
 \mathcal A_{E_1}\otimes\mathcal L(\mathfrak K),
\]
where \(\mathfrak K\) is an infinite-dimensional Hilbert space. Thus
\(\mathcal A_E\) contains a von Neumann subalgebra isomorphic to
\(\mathcal L(\mathfrak K)\). But the mapping \(T\mapsto T^*\) of
\(\mathcal L(\mathfrak K)\) into itself is not strongly continuous on
bounded sets (\hyperref[no:I-3-1]{Chapter~I, \S~3, no.~1}).
\end{proof}

\begin{theorem}\label{thm:III-8-2}
Let \(\mathcal A_1,\mathcal A_2\) be two von Neumann algebras. If one
of them is purely infinite, then
\(\mathcal A_1\otimes\mathcal A_2\) is purely
infinite.
\end{theorem}

\begin{proof}
Let \(\mathfrak H_1,\mathfrak H_2\) be the Hilbert spaces on which
\(\mathcal A_1,\mathcal A_2\) act. We shall suppose that
\(\mathcal A_1\) is purely infinite. Choose an orthonormal basis
\((e_i)_{i\in I}\) of \(\mathfrak H_2\); this permits the canonical
identification of \(\mathfrak H_1\otimes\mathfrak H_2\) with
\(\bigoplus_{i\in I}\mathfrak H_i\), where
\(\mathfrak H_i=\mathfrak H_1\) for every \(i\in I\). The elements of
\(\mathcal A_1\otimes\mathcal L(\mathfrak H_2)\)
are
% Source PDF page 313 (printed p. 304)
\phantomsection\label{srcpage:304}
represented by matrices \((T_{\iota\chi})\), where
\(T_{\iota\chi}\in\mathcal A_1\), for all \(\iota,\chi\in I\).
If \(X\in\mathcal A_1\), the element \(X\otimes I\) of
\(\mathcal A_1\mathbin{\overline\otimes}\mathcal L(\mathfrak H_2)\)
is represented by the matrix \((\delta_{\iota\chi}X)\)
(\hyperref[prop:I-2-4]{Chapter~I, \S~2, Proposition~4}). If
\(T=(T_{\iota\chi})\in
\mathcal A_1\mathbin{\overline\otimes}\mathcal L(\mathfrak H_2)\)
and \(X\in\mathcal A_1\), then \(T(X\otimes I)\) is represented by
the matrix \((T_{\iota\chi}X)\)
(\hyperref[no:I-2-3]{Chapter~I, \S~2, no.~3}).

If \(\mathcal A_1\mathbin{\overline\otimes}\mathcal A_2\) is not
purely infinite, there exist a faithful normal trace \(\varphi\) on
\((\mathcal A_1\mathbin{\overline\otimes}\mathcal A_2)_+\) and an
\(S\in(\mathcal A_1\mathbin{\overline\otimes}\mathcal A_2)_+\) such
that \(S\ne0\) and \(\varphi(S)<+\infty\). Write
\(S=(S_{\iota\chi})\). If \(S_{\iota\iota}\) were zero for every
\(\iota\in I\), then \((Sx\mid x)=0\) for every
\(x\in\mathfrak H_\iota\) and every \(\iota\), hence
\(S^{1/2}(\mathfrak H_\iota)=0\) for every \(\iota\), and consequently
\(S^{1/2}=0\) and \(S=0\), which is absurd. There is therefore an
index \(\iota_0\) such that \(S_{\iota_0\iota_0}\ne0\). For every
\(T=(T_{\iota\chi})\in
\mathcal A_1\mathbin{\overline\otimes}\mathcal L(\mathfrak H_2)\),
put \(\widetilde T=T_{\iota_0\iota_0}\). We have
\(\widetilde S\geq0\) and \(\widetilde S\ne0\). There are therefore
a nonzero projection \(E\) of \(\mathcal A_1\) and a number
\(\lambda\geq0\) such that
\[
  \lVert Ex\rVert\leq\lambda\lVert\widetilde Sx\rVert
  \qquad(x\in\mathfrak H_1).
\]
Consider the mappings
\[
 \mathcal A_1\xrightarrow{\alpha}
 \mathcal A_1\mathbin{\overline\otimes}\mathcal A_2
 \xrightarrow{\beta}
 \mathcal A_1\mathbin{\overline\otimes}\mathcal A_2
 \xrightarrow{\gamma}\mathcal A_1,
\]
where
\[
 \alpha(X)=X\otimes I,
 \qquad \beta(T)=ST^*,
 \qquad \gamma(T)=\widetilde T.
\]
These mappings are strongly continuous on bounded sets: the first by
\hyperref[cor:I-4-2-1]{Chapter~I, \S~4, Corollary~1 of Theorem~2},
the second by \cref{lem:III-8-13}, and the third plainly. But
\[
 \gamma(\beta(\alpha(X)))
   =\bigl(S(X\otimes I)^*\bigr)^{\sim}
   =\widetilde S X^*.
\]
Thus the mapping \(X\mapsto\widetilde S X^*\) of \(\mathcal A_1\)
into itself, and hence, a fortiori, the mapping \(X\mapsto EX^*\), is
strongly continuous on bounded sets. On taking
\(X\in E\mathcal A_1E\), this contradicts \cref{lem:III-8-14}, since
\(E\) is a purely infinite projection of \(\mathcal A_1\).
\end{proof}

\begin{corollary}\label{cor:III-8-thm2}
Let \(\mathcal A_1,\mathcal A_2\) be two von Neumann algebras. If one
of them is continuous, then
\(\mathcal A_1\mathbin{\overline\otimes}\mathcal A_2\) is continuous.
\end{corollary}

\begin{proof}
Write
\[
 \mathcal A_1=\mathcal B_1\times\mathcal C_1,
 \qquad
 \mathcal A_2=\mathcal B_2\times\mathcal C_2,
\]
where \(\mathcal B_1,\mathcal B_2\) are semifinite and
\(\mathcal C_1,\mathcal C_2\) are purely infinite. Then
\(\mathcal C_1\mathbin{\overline\otimes}\mathcal B_2\),
\(\mathcal B_1\mathbin{\overline\otimes}\mathcal C_2\), and
\(\mathcal C_1\mathbin{\overline\otimes}\mathcal C_2\) are purely
infinite by \cref{thm:III-8-2}, and hence continuous. It remains to
prove that
\(\mathcal B_1\mathbin{\overline\otimes}\mathcal B_2\) is continuous,
assuming, for example, that \(\mathcal B_1\) is continuous. There is a
decreasing sequence \((E_n)\) of finite projections of \(\mathcal B_1\),
with central support \(I\), such that
\[
 E_n-E_{n+1}\sim E_{n+1}\qquad(n\geq1)
\]
(\hyperref[cor:III-2-prop7-4]{\S~2, Corollary~4 of Proposition~7}).
Put \(F_n=E_n\otimes I\in
\mathcal B_1\mathbin{\overline\otimes}\mathcal B_2\). The \(F_n\) are
decreasing projections such that
\(F_n-F_{n+1}\sim F_{n+1}\) for \(n\geq1\). By
\hyperref[cor:I-1-7-1]{Chapter~I, \S~1, Corollary~1 of Proposition~7},
the central support of \(F_n\) is \(I\). Consequently
\(\mathcal B_1\mathbin{\overline\otimes}\mathcal B_2\) is continuous
(\hyperref[cor:III-2-prop7-4]{\S~2, Corollary~4 of Proposition~7}).
\end{proof}

\noindent Bibliography: \ArticleRef{192}.

\noindent\textit{Exercises.}
\begin{enumerate}[label=\arabic*.,leftmargin=*]
\item\label{ex:III-8-1}
Let \(\mathcal A\) be a purely infinite von Neumann algebra of
countable type, and let \(\mathcal Z\) be its centre.
\begin{enumerate}[label=\textit{\alph*.},leftmargin=*]
\item Let \(\mathfrak m\) be a two-sided ideal of \(\mathcal A\). If
\(\mathfrak m\ne0\), then \(\mathfrak m\cap\mathcal Z\ne0\). [Let
\(T\in\mathfrak m_+\), \(T\ne0\). Exercise~6 of Chapter~I, \S~1,
supplies a projection \(E\ne0\) in \(\mathfrak m\); \(E\) is
equivalent to its central support \(F\); use \S~1, Exercise~8.]

% Source PDF page 314 (printed p. 305)
\phantomsection\label{srcpage:305}
\item The intersection of the maximal two-sided ideals of
\(\mathcal A\) is zero
(\hyperref[cor:III-5-1-2]{\S~5, Proposition~1, Corollary~2}).

\item A purely infinite factor of countable type has no nontrivial
two-sided ideals. (Use (a).)

\item Let \(\mathcal A\) be a purely infinite factor of countable type
(such factors exist: Chapter~I, \S~9, no.~4, Remark), and let
\(\mathfrak K\) be a Hilbert space that does not have a countable basis.
Then \(\mathcal A\mathbin{\overline\otimes}\mathcal L(\mathfrak K)\)
is a purely infinite factor having nontrivial two-sided ideals.
[Its commutant
\((\mathcal A\mathbin{\overline\otimes}\mathcal L(\mathfrak K))'
=\mathcal A'\mathbin{\overline\otimes}\mathbb C I_{\mathfrak K}\)
is purely infinite; \(\mathcal L(\mathfrak K)\) is not of countable
type, and hence
\(\mathcal A\mathbin{\overline\otimes}\mathcal L(\mathfrak K)\) is
not of countable type; use Chapter~I, \S~1, Exercise~7.]
\ArticleRef{25}.
\end{enumerate}

\item\label{ex:III-8-2}
Let \(\mathcal A\) be a von Neumann algebra, and let \(\varphi\) be a
faithful positive linear form on \(\mathcal A\). Suppose that there is
a number \(\lambda\geq1\) such that
\[
 \varphi(T^*T)\leq\lambda\varphi(TT^*)
 \qquad(T\in\mathcal A).
\]
Show that \(\mathcal A\) is finite. [Let \(E\) be the greatest properly
infinite projection of \(\mathcal A\); write
\(E=E_1+E_2+\cdots\), where the \(E_i\) are mutually orthogonal
projections equivalent to \(E\). One has
\(\varphi(E_i)\geq\lambda^{-i}\varphi(E)\) for every \(i\), and hence
\(\varphi(E_i)=0\), \(\varphi(E)=0\), and \(E=0\).]

\item\label{ex:III-8-3}
Let \(\mathcal A\) be a von Neumann algebra on \(\mathfrak H\).
\begin{enumerate}[label=\textit{\alph*.},leftmargin=*]
\item Let
\[
 E_0=P_{\mathfrak X_0},\quad E_1=P_{\mathfrak X_1},\quad
 E_2=P_{\mathfrak X_2},\ldots
\]
be equivalent, pairwise orthogonal projections of \(\mathcal A\) with
sum \(I\), and let \(U_i\in\mathcal A\) be a partially isometric
operator having \(E_{i+1}\) as initial projection and \(E_i\) as final
projection \((i=0,1,2,\ldots)\). Put
\[
 T=I-\sum_{i=0}^{\infty}U_i
 \quad\text{(the series converging strongly)},
 \qquad
 \mathfrak Y=T(\mathfrak X_1\oplus\mathfrak X_2\oplus\cdots).
\]
Show that \(Tx=0\) implies \(x=0\). Deduce that
\(\mathfrak Y\cap\mathfrak X_0=0\). Show that
\(T^*(\mathfrak H)\cap\mathfrak X_0=0\). Deduce that \(\mathfrak Y\)
is dense in \(\mathfrak H\).

\item Let \(E,F\) be two equivalent, orthogonal, properly infinite
projections of \(\mathcal A\), with sum \(I\). Show that there is a
\(T\in\mathcal A\) such that: \textup{1\textdegree} \(Tx=0\) implies
\(x=0\); \textup{2\textdegree} \(TF(\mathfrak H)\) is dense in
\(\mathfrak H\); \textup{3\textdegree}
\(TF(\mathfrak H)\cap E(\mathfrak H)=0\). (Apply Corollary~2 of
Theorem~1 to \(F\), and use (a).)

\item If \(\mathcal A'\) is properly infinite, there are two elements
\(S_1,S_2\in\mathcal A\) such that
\[
 S_1(\mathfrak H)\cap S_2(\mathfrak H)=0,
\]
while both \(S_1(\mathfrak H)\) and \(S_2(\mathfrak H)\) are dense in
\(\mathfrak H\). (Apply (b) twice, interchanging the roles of \(E\)
and \(F\).) Let \(x\) be a separating element for \(\mathcal A\). Show
that \(y=S_1x\) and \(z=S_2x\) are separating for \(\mathcal A\), and
that an equality \(T'y=T'z\), with \(T'\in\mathcal A'\), implies
\(T'y=T'z=0\). [One has
\(T'y=T'z=S_1T'x=S_2T'x\in
S_1(\mathfrak H)\cap S_2(\mathfrak H)\).] Deduce that there is no closed
operator \(T'_1\mathrel{\eta}\mathcal A'\)
(Chapter~I, \S~1, Exercise~10), with everywhere dense domain, such that
\(T'_1y=z\).

\item Suppose that \(\mathcal A\) is properly infinite and of countable
type. Show that there are two faithful normal positive linear forms
\(\varphi_1,\varphi_2\) on \(\mathcal A\) such that every positive
linear form majorized by both \(\varphi_1\) and \(\varphi_2\) is zero.
[By making an isomorphic representation of \(\mathcal A\), reduce to
the case in which \(\mathcal A'\) is properly infinite. Then
\(\mathcal A\) has a separating element \(x\); with the notation of
(c), take \(\varphi_1=\omega_y\), \(\varphi_2=\omega_z\), and apply
\hyperref[lem:I-4-1]{Chapter~I, \S~4, Lemma~1}.] \ArticleRef{19}.
\end{enumerate}

\item\label{ex:III-8-4}
Let \(\mathcal A\) be a von Neumann algebra, and let \(E,F\) be
equivalent projections of \(\mathcal A\). There are orthogonal
projections \(E_1,E_2\) (respectively, \(F_1,F_2\)) of \(\mathcal A\),
with sum \(E\) (respectively, \(F\)), and unitary operators
\(U_1,U_2\in\mathcal A\), such that
\[
 U_1E_1U_1^{-1}=F_1,
 \qquad
 U_2E_2U_2^{-1}=F_2.
\]
[If \(E\) is finite, take \(E_1=F_1=0\) and apply \S~2,
Proposition~6. If \(E\) is properly infinite, choose
\(E_1,E_2,F_1,F_2\) so that
\(E\sim E_1\sim E_2\sim F\sim F_1\sim F_2\), and prove that
\(I-E_1\sim I\sim I-F_1\) and
\(I-E_2\sim I\sim I-F_2\).] \ArticleRef{10}.

% Source PDF page 315 (printed p. 306)
\phantomsection\label{srcpage:306}
\item\label{ex:III-8-5}
Let \(\mathcal A\) be a von Neumann algebra, and let \(E,F\) be two
projections of \(\mathcal A\), with \(\mathcal A_F\) of countable
type. In order that \(E\prec F\), it is necessary and sufficient that
\(\varphi(E)\leq\varphi(F)\) for every normal trace \(\varphi\) on
\(\mathcal A_+\). [To prove sufficiency, reduce first to the case
\(F\prec E\), and then to the case \(F\leq E\). If \(F\) is finite,
reduce further to the case in which there is on \(\mathcal A_+\) a
faithful normal semifinite trace \(\varphi\) such that
\(\varphi(F)<+\infty\); the relations
\[
 \varphi(E)\leq\varphi(F),
 \qquad
 \varphi(E)=\varphi(F)+\varphi(E-F)
\]
then give \(E-F=0\). If \(F\) is properly infinite, apply Corollary~5
of Theorem~1.]

\item\label{ex:III-8-6}
Let \(\mathcal A\) be a von Neumann algebra such that every normal
positive linear form on \(\mathcal A\) is the sum of finitely many forms
\(\omega_x\). Then there is an integer \(n\) such that every normal
positive linear form on \(\mathcal A\) is the sum of \(n\) forms
\(\omega_x\). (Use Corollary~10 of Theorem~1 and Propositions~6--9 of
\S~6.) \ArticleRef{19}.

\item\label{ex:III-8-7}
Prove the converse of \S~6, Lemma~2. [If the strong and ultra-strong
topologies do not coincide on \(\mathcal A\), either \(\mathcal A\) is
properly infinite and \(\mathcal A'\) is finite, or both \(\mathcal A\)
 and \(\mathcal A'\) are finite and the inverse coupling operator
 \(C_{\mathcal A}^{-1}\) is not essentially bounded. In the first
(respectively, second) case, use the proof of Proposition~8
(respectively, Proposition~9) of \S~6.] \ArticleRef{19}.

\item\label{ex:III-8-8}
Let \(\mathcal A\) be a von Neumann algebra, let \(\mathcal Z\) be its
centre, and let \(\Phi\) be a linear mapping of \(\mathcal A\) into a
complex vector space \(E\), such that \(\Phi(ST)=\Phi(TS)\) for
\(S,T\in\mathcal A\), and such that \(\Phi(F)\ne0\) for every nonzero
projection \(F\) of \(\mathcal Z\). Then \(\mathcal A\) is finite.
(Adapt the proof of Corollary~3 of Theorem~1.)

\item\label{ex:III-8-9}
Let \(\mathcal A\) be a von Neumann algebra, and let \(\mathcal Z\) be
its centre.
\begin{enumerate}[label=\textit{\alph*.},leftmargin=*]
\item Let \(\Phi\) be a finite \(\mathcal Z\)-trace on \(\mathcal A\).
Show that \(\Phi\) is normal. (Reduce to the case in which
\(\mathcal A\) is either properly infinite or finite. In the first
case, show that \(\Phi=0\), using Corollary~2 of Theorem~1. In the
second case, use \S~5, Exercise~4.)

\item Let \(\mathfrak m\) be a restricted ideal of \(\mathcal A\)
(Chapter~I, \S~1, Exercise~6). Show that a \(\mathcal Z\)-trace on
\(\mathfrak m\) is normal. (Use (a).) \ArticleRef{12}.
\end{enumerate}

\item\label{ex:III-8-10}
Let \(\mathcal A\) be a von Neumann algebra, and let \(\mathcal Z\) be
its centre. Suppose that \(\mathcal A\) is finite and \(\mathcal A'\)
is properly infinite. There is a projection \(F\in\mathcal Z'\) having
the following property: for every \(T\in\mathcal A\), there is a unique
\(T'\in\mathcal Z\) such that
\[
 FTF=T'F,
\]
and \(T\mapsto T'\) is the canonical \(\mathcal Z\)-trace of
\(\mathcal A\). (Reduce to the case in which \(\mathcal A\) is of
countable type. Then use Corollary~11 of Theorem~1 and \S~6,
Exercise~3.) \ArticleRef{89}.

\item\label{ex:III-8-11}
Let \(\mathcal A\) be a factor of type \(\mathrm{II}_\infty\). Show
that \(\mathcal A\) is spatially isomorphic to an algebra
\[
 \mathcal B\mathbin{\overline\otimes}\mathcal L(\mathfrak K),
\]
where \(\mathcal B\) is a factor of type \(\mathrm{II}_1\) and
\(\mathfrak K\) is an infinite-dimensional Hilbert space. [By
Corollary~2 of Theorem~1 there is a sequence
\(E_1,E_2,\ldots\) of equivalent, pairwise orthogonal, nonzero
projections of \(\mathcal A\). By \S~2, Proposition~7, one may suppose
that the \(E_i\) are finite. By \S~1, Corollary~2 of Theorem~1, one may
suppose \(\sum_iE_i=I\). Then use
\hyperref[prop:I-2-5]{Chapter~I, \S~2, Proposition~5(ii)}.]

\item\label{ex:III-8-12}
Let \(\mathfrak H\) (respectively, \(\mathfrak H_2\)) be a Hilbert
space of countably infinite dimension (respectively, of dimension
two).
\begin{enumerate}[label=\textit{\alph*.},leftmargin=*]
\item Let \(\mathcal A\) be a von Neumann algebra on \(\mathfrak H\),
let \((x_1,x_2,\ldots)\) be a sequence generating \(\mathcal A\), and
let \(y\) be a generator of \(\mathcal L(\mathfrak H)\) (cf. \S~3,
Exercise~9). Represent every element of
\(\mathcal A\mathbin{\overline\otimes}\mathcal L(\mathfrak H)\) by a
matrix \((x_{ij})\) of elements of \(\mathcal A\). For
\(k=1,2,\ldots\), let
\(z_k\in\mathcal A\mathbin{\overline\otimes}\mathcal L(\mathfrak H)\)
be represented by the matrix \((x_{ij})\) such that
\(x_{11}=x_k\) and \(x_{ij}=0\) if \(i\ne1\) or \(j\ne1\). Then the
\(z_k\) and \(I\otimes y\) generate
\(\mathcal A\mathbin{\overline\otimes}\mathcal L(\mathfrak H)\).

\item Deduce from (a) that
\(\mathcal A\mathbin{\overline\otimes}\mathcal L(\mathfrak H)\) is
generated by two elements.

% Source PDF page 316 (printed p. 307)
\phantomsection\label{srcpage:307}
\item Deduce from (b) that
\(\mathcal A\mathbin{\overline\otimes}\mathcal L(\mathfrak H)
\mathbin{\overline\otimes}\mathcal L(\mathfrak H_2)\) can be generated
by two elements, one of which is unitary.

\item Deduce from (c) that
\[
 \mathcal A\mathbin{\overline\otimes}\mathcal L(\mathfrak H)
 \mathbin{\overline\otimes}\mathcal L(\mathfrak H_2)
 \mathbin{\overline\otimes}\mathcal L(\mathfrak H_2)
\]
can be generated by two unitary elements.

\item Deduce from (d) that, if \(\mathcal B\) is a properly infinite
von Neumann algebra on \(\mathfrak H\), then \(\mathcal B\) can be
generated by one element. \ArticleRef{473}.
\end{enumerate}
\end{enumerate}

\section{Derivations and Automorphisms of von Neumann Algebras}
\label{sec:III-9}

\NumberedParagraph{1}{Derivations of Algebras}
\label{no:III-9-1}
Let \(A\) be an algebra. For \(a,b\in A\), put
\[
 [a,b]=ab-ba.
\]
A \emph{derivation} of \(A\) is a linear mapping \(\delta\) of \(A\)
into \(A\) such that
\begin{equation}
 \delta(ab)=(\delta a)b+a(\delta b)
 \tag{1}\label{eq:III-9-1}
\end{equation}
for all \(a,b\in A\). One then has
\begin{equation}
 \delta([a,b])=[\delta a,b]+[a,\delta b]
 \tag{2}\label{eq:III-9-2}
\end{equation}
for all \(a,b\in A\). Indeed,
\begin{align*}
 \delta([a,b])
 &=\delta(ab)-\delta(ba)\\
 &=(\delta a)b+a(\delta b)-(\delta b)a-b(\delta a)
   =[\delta a,b]+[a,\delta b].
\end{align*}

If \(\delta\) leaves a subset \(M\) of \(A\) invariant, it leaves the
commutant \(M'\) of \(M\) invariant; for, if \(a\in M'\) and
\(b\in M\), then
\[
 [\delta a,b]=\delta([a,b])-[a,\delta b]=0.
\]
In particular, \(\delta\) leaves the centre of \(A\) invariant. If
\(A\) has an identity \(I\), then
\[
 \delta(I)=\delta(I^2)=\delta(I)I+I\delta(I)=2\delta(I),
\]
and hence \(\delta(I)=0\).

If \(x\in A\), the mapping \(a\mapsto[x,a]\) of \(A\) into \(A\) is
a derivation, since
\[
 x(ab)-(ab)x=(xa-ax)b+a(xb-bx).
\]
This derivation is called the \emph{inner derivation defined by \(x\)}.

Suppose that \(A\) is equipped with an involution \(x\mapsto x^*\).
Let \(\delta\) be a derivation of \(A\), and put
\[
 \delta^*(a)=\bigl(\delta(a^*)\bigr)^*
 \qquad(a\in A).
\]
Then \(\delta^*\) is a derivation, since
\begin{align*}
 \delta^*(ab)
 &=\bigl(\delta(b^*a^*)\bigr)^*
  =\bigl((\delta b^*)a^*+b^*(\delta a^*)\bigr)^*\\
 &=a(\delta^*b)+(\delta^*a)b.
\end{align*}
If \(\delta=\delta^*\), the derivation \(\delta\) is said to be
\emph{Hermitian}; in that case \(\delta a\) is Hermitian for every
Hermitian \(a\in A\). If \(\delta\) is an arbitrary derivation of
\(A\), one may write
\[
 \delta=\frac12(\delta+\delta^*)
       +i\frac1{2i}(\delta-\delta^*),
\]
and both \(\frac12(\delta+\delta^*)\) and
\(\frac1{2i}(\delta-\delta^*)\) are Hermitian derivations of \(A\).

% Source PDF page 317 (printed p. 308)
\phantomsection\label{srcpage:308}
\NumberedParagraph{2}{Derivations of \(\boldsymbol{C^*}\)-Algebras:
Continuity and Extension}
\label{no:III-9-2}

\begin{lemma}\label{lem:III-9-1}
Let \(\mathcal A\) be a \(C^*\)-algebra of operators, let
\(\mathcal B\) be a commutative \(C^*\)-subalgebra of \(\mathcal A\)
containing \(I\), let \(\varphi\) be a positive form on \(\mathcal A\)
whose restriction to \(\mathcal B\) is a character of \(\mathcal B\),
and let \(\delta\) be a derivation of \(\mathcal A\). Then
\[
 \varphi(\delta\mathcal B)=0.
\]
\end{lemma}

\begin{proof}
Let \(T\) be a Hermitian element of \(\mathcal B\), and let us show
that \(\varphi(\delta T)=0\). We may suppose that \(\varphi(T)=0\), by
adding a scalar operator to \(T\). On considering the Gelfand transform
of \(\mathcal B\), one sees that \(T\) may be written in the form
\(U^2-V^2\), where \(U,V\) are Hermitian elements of \(\mathcal B\)
such that \(\varphi(U)=\varphi(V)=0\). We have
\[
 \delta T=U(\delta U)+(\delta U)U-V(\delta V)-(\delta V)V.
\]
Now
\[
 \lvert\varphi(U(\delta U))\rvert^2
 \leq\varphi(U^2)\varphi((\delta U)^*(\delta U))=0,
\]
and hence \(\varphi(U(\delta U))=0\). Similarly,
\[
 \varphi((\delta U)U)=\varphi(V(\delta V))
 =\varphi((\delta V)V)=0.
\]
\end{proof}

\begin{lemma}\label{lem:III-9-2}
Let \(\mathcal A\) be a \(C^*\)-algebra of operators containing \(I\),
let \(\mathcal C\) be its centre, and let \(\delta\) be a derivation of
\(\mathcal A\). Then \(\delta(\mathcal C)=0\).
\end{lemma}

\begin{proof}
Since \(\delta(\mathcal C)\subset\mathcal C\) by no.~1, it is enough
to consider the case in which \(\mathcal A\) is commutative. Every
character of \(\mathcal A\) then vanishes on \(\delta(\mathcal A)\) by
\cref{lem:III-9-1}, whence \(\delta(\mathcal A)=0\).
\end{proof}

\begin{lemma}\label{lem:III-9-3}
Let \(\mathcal A\) be a \(C^*\)-algebra of operators, and let
\(\delta\) be a derivation of \(\mathcal A\). Then \(\delta\) is
continuous in norm.
\end{lemma}

\begin{proof}
If \(I\notin\mathcal A\), then \(\delta\) extends to a derivation of
\(\mathcal A+\mathbb C I\) (which is a \(C^*\)-algebra) that vanishes
at \(I\). We may therefore suppose that \(I\in\mathcal A\). By no.~1,
we reduce to the case in which \(\delta\) is Hermitian. Suppose that
\(\delta\) is not continuous. Its restriction to the real vector
subspace of Hermitian elements of \(\mathcal A\) is not continuous.
By the closed graph theorem
(\TreatiseRef{3}, Chapter~I, \S~3, Corollary~5 of Theorem~1), there is
a sequence \((S_1,S_2,\ldots)\) of Hermitian elements of \(\mathcal A\)
tending to zero such that \(\delta S_n\to A\), where \(A\) is a
nonzero Hermitian element of \(\mathcal A\). We may suppose that
\(S_n\ne0\) for every \(n\). Multiplying the \(S_n\) by a suitable
constant, we may suppose that the spectrum of \(A\) contains a number
greater than \(3\). Replacing \(S_n\) by
\(S_n+2\lVert S_n\rVert I\), which does not change \(\delta S_n\), we
may suppose that
\[
 S_n\geq\frac13\lVert S_n\rVert I.
\]
There is a Hermitian element \(H\) of \(\mathcal A\) such that
\[
 \lVert H\rVert=1,
 \qquad
 HAH\geq3H^2.
\]
Put \(T_n=HS_nH\) and \(B=HAH\). We have
\[
 \delta T_n=(\delta H)S_nH+H(\delta S_n)H+HS_n(\delta H),
\]
% Source PDF page 318 (printed p. 309)
\phantomsection\label{srcpage:309}
and therefore
\begin{equation}
 \delta T_n\longrightarrow B.
 \tag{1}\label{eq:III-9-3}
\end{equation}
On the other hand,
\[
 \frac13\lVert S_n\rVert H^2
 \leq T_n\leq\lVert S_n\rVert H^2,
\]
whence \(\frac13\lVert S_n\rVert\leq\lVert T_n\rVert\), and hence
\[
 \lVert T_n\rVert^{-1}T_n
 \leq3\lVert S_n\rVert^{-1}\lVert S_n\rVert H^2
 =3H^2.
\]
Consequently,
\begin{equation}
 \lVert T_n\rVert^{-1}T_n\leq B.
 \tag{2}\label{eq:III-9-4}
\end{equation}
For every \(n\), there is a character \(\psi_n\) of the commutative
\(C^*\)-algebra generated by \(I\) and \(T_n\) such that
\[
 \psi_n(\lVert T_n\rVert^{-1}T_n)=1.
\]
Since \(I\) is an interior point of the cone of positive elements of
\(\mathcal A\), \(\psi_n\) extends to a positive form \(\varphi_n\) on
\(\mathcal A\)
(\TreatiseRef{3}, Chapter~II, \S~3, Proposition~6). By
\cref{lem:III-9-1}, \(\varphi_n(\delta T_n)=0\). By
\eqref{eq:III-9-3}, \(\lvert\varphi_n(B)\rvert<1\) for sufficiently
large \(n\). But \eqref{eq:III-9-4} gives
\[
 \varphi_n(B)\geq
 \varphi_n(\lVert T_n\rVert^{-1}T_n)=1,
\]
a contradiction.
\end{proof}

\begin{lemma}\label{lem:III-9-4}
Let \(\mathcal A\) be a \(C^*\)-algebra of operators, let
\(\mathcal B\) be its weak closure, and let \(\delta\) be a derivation
of \(\mathcal A\). Then \(\delta\) is ultra-weakly continuous and
extends to a derivation of \(\mathcal B\) having the same norm.
\end{lemma}

\begin{proof}
Let \(\mathfrak H\) be the Hilbert space on which \(\mathcal A\) acts,
and let \(\mathcal A_1^+\) be the set of elements of \(\mathcal A_+\)
of norm at most \(1\). If \(T\in\mathcal A_1^+\) and
\(x,y\in\mathfrak H\), then
\begin{align}
 \bigl\lvert(\delta(T)x\mid y)\bigr\rvert
 &=\bigl\lvert
   (\delta(T^{1/2})T^{1/2}x\mid y)
   +(T^{1/2}\delta(T^{1/2})x\mid y)
   \bigr\rvert \notag\\
 &\leq\lVert\delta\rVert\lVert T^{1/2}x\rVert\lVert y\rVert
   +\lVert\delta\rVert\lVert T^{1/2}y\rVert\lVert x\rVert \notag\\
 &\leq\lVert\delta\rVert
   (\lVert x\rVert^2+\lVert y\rVert^2)^{1/2}
   (\lVert T^{1/2}x\rVert^2+\lVert T^{1/2}y\rVert^2)^{1/2} \notag\\
 &=\lVert\delta\rVert
   (\lVert x\rVert^2+\lVert y\rVert^2)^{1/2}
   \bigl((Tx\mid x)+(Ty\mid y)\bigr)^{1/2}.
 \tag{3}\label{eq:III-9-5}
\end{align}
Here \(\lVert\delta\rVert<+\infty\) by \cref{lem:III-9-3}. Let
\((x_1,x_2,\ldots)\) and \((y_1,y_2,\ldots)\) be sequences in
\(\mathfrak H\) such that
\[
 \sum_i\lVert x_i\rVert^2<+\infty,
 \qquad
 \sum_i\lVert y_i\rVert^2<+\infty,
\]
and let \(\varepsilon>0\). There is an integer \(N\) such that, for
\(\lVert T\rVert\leq1\),
\[
 \left\lvert\sum_{i=N}^{\infty}(\delta(T)x_i\mid y_i)\right\rvert
 \leq\varepsilon.
\]
In view of \eqref{eq:III-9-5}, there is therefore a weak neighborhood
\(V\) of zero in \(\mathcal A_1^+\) such that
\[
 T\in V\quad\Longrightarrow\quad
 \left\lvert\sum_{i=1}^{\infty}(\delta(T)x_i\mid y_i)\right\rvert
 \leq2\varepsilon.
\]
By \hyperref[srcpage:45]{Chapter~I, \S~3, the corollary to
Theorem~3}, the linear form
\[
 T\longmapsto\sum_i(\delta(T)x_i\mid y_i)
\]
on \(\mathcal A\) is ultra-weakly continuous. Thus \(\delta\) is
ultra-weakly continuous.

% Source PDF page 319 (printed p. 310)
\phantomsection\label{srcpage:310}
Every point of \(\mathcal B\) is in the ultra-weak closure of a bounded
subset of \(\mathcal A\)
(\hyperref[thm:I-3-3]{Chapter~I, \S~3, Theorem~3}). On the other hand,
the bounded ultra-weakly closed subsets of \(\mathcal B\) are
ultra-weakly compact, and hence ultra-weakly complete. Consequently
(\TreatiseRef{3}, Chapter~III, \S~2, Proposition~8), \(\delta\) extends
to an ultra-weakly continuous linear mapping
\(\overline\delta\colon\mathcal B\to\mathcal B\).

The relation
\[
 \overline\delta(ST)=(\overline\delta S)T+S(\overline\delta T)
\]
holds for \(S,T\in\mathcal A\), and both sides are separately
ultra-weakly continuous functions of \(S\) and \(T\). Hence
\(\overline\delta\) is a derivation of \(\mathcal B\). The equality
\(\lVert\overline\delta\rVert=\lVert\delta\rVert\) follows from
\hyperref[thm:I-3-3]{Chapter~I, \S~3, Theorem~3}.
\end{proof}

\noindent Bibliography: \ArticleRef{252}, \ArticleRef{367}.

\NumberedParagraph{3}{Derivations of von Neumann Algebras}
\label{no:III-9-3}

\begin{lemma}\label{lem:III-9-5}
Let \(\mathcal A\) be a von Neumann algebra, and let \(\delta\) be a
derivation of \(\mathcal A\).
\begin{enumerate}[label=\textup{(\roman*)},leftmargin=2em]
\item For every projection \(E\) of \(\mathcal A\), the mapping
\[
 T\longmapsto E\delta(T)E
\]
is a derivation of \(E\mathcal A E\) with norm at most
\(\lVert\delta\rVert\).

\item Let \(\mathcal F\) be an upward directed set of projections of
\(\mathcal A\), with supremum \(I\). Suppose that, for every
\(E\in\mathcal F\), there is a \(T_E\in E\mathcal A E\) such that
\(\lVert T_E\rVert\leq\lVert\delta\rVert\) and
\[
 E\delta(T)E=[T_E,T]
 \qquad(T\in E\mathcal A E).
\]
Then there is a \(T_0\in\mathcal A\) such that
\(\lVert T_0\rVert\leq\lVert\delta\rVert\) and
\[
 \delta(T)=[T_0,T]
 \qquad(T\in\mathcal A).
\]
\end{enumerate}
\end{lemma}

\begin{proof}
Part (i) is immediate. We prove (ii). For every \(E\in\mathcal F\),
let \(\mathcal F_E\) be the set of the \(E'\in\mathcal F\) such that
\(E'\geq E\). The sets \(\mathcal F_E\) form a filter base
\(\mathscr B\) on \(\mathcal F\). Since the closed ball with centre
zero and radius \(\lVert\delta\rVert\) in \(\mathcal A\) is weakly
compact, the mapping \(E\mapsto T_E\) has along \(\mathscr B\) a weak
cluster value \(T_0\) such that
\(\lVert T_0\rVert\leq\lVert\delta\rVert\). On the other hand, if
\(T\in\mathcal A\), the mapping \(E\mapsto E\delta(T)E\) tends strongly
to \(\delta(T)\) along \(\mathscr B\). Hence
\[
 \delta(S)=[T_0,S]
 \qquad
 \left(S\in\mathcal A_0=
   \bigcup_{E\in\mathcal F}E\mathcal A E\right).
\]
Thus the derivation
\[
 \delta'\colon T\longmapsto\delta(T)-[T_0,T]
\]
of \(\mathcal A\) vanishes on \(\mathcal A_0\). For every
\(T\in\mathcal A\), we therefore have
\[
 E\delta'(T)E
 =\delta'(ETE)-\delta'(E)TE-ET\delta'(E)=0
\]
for every \(E\in\mathcal F\), and consequently
\(\delta'(T)=0\), on passing to the strong limit.
\end{proof}

\begin{lemma}\label{lem:III-9-6}
Let \(\mathcal A\) be a von Neumann algebra, and let \(\mathcal F\) be
the set of projections \(E\) of \(\mathcal A\) such that
\(\mathcal A_E\) is of countable type. Then \(\mathcal F\) is upward
directed and has supremum \(I\).
\end{lemma}

\begin{proof}
This follows from the fact that the elements of \(\mathcal F\) are the
projections of the form \(E_{\mathfrak M}^{\mathcal A'}\), where
\(\mathfrak M\) is a countable set of vectors
(\hyperref[prop:I-1-6]{Chapter~I, \S~1, Proposition~6}).
\end{proof}

% Source PDF page 320 (printed p. 311)
\phantomsection\label{srcpage:311}
\begin{theorem}\label{thm:III-9-1}
Let \(\mathcal A\) be a von Neumann algebra, and let \(\delta\) be a
derivation of \(\mathcal A\). There is a \(T_0\in\mathcal A\) such
that
\[
 \lVert T_0\rVert\leq\lVert\delta\rVert,
 \qquad
 \delta(T)=[T_0,T]\quad(T\in\mathcal A).
\]
\end{theorem}

\begin{proof}
\noindent\textup{1\textdegree.}
Let \(\mathcal U\) be the unitary group of \(\mathcal A\). For
\(U\in\mathcal U\), define a mapping \(A_U\) of \(\mathcal A\) into
itself by
\[
 A_U(T)=(UT+\delta(U))U^{-1}.
\]
If \(U,V\in\mathcal U\), then
\begin{align*}
 A_UA_V(T)
 &=\bigl[U(VT+\delta(V))V^{-1}+\delta(U)\bigr]U^{-1}\\
 &=UVTV^{-1}U^{-1}+U\delta(V)V^{-1}U^{-1}
   +\delta(U)U^{-1}\\
 &=\bigl(UVT+\delta(UV)\bigr)V^{-1}U^{-1},
\end{align*}
and hence
\begin{equation}
 A_UA_V=A_{UV}.
 \tag{1}\label{eq:III-9-6}
\end{equation}

Let \(\mathscr E\) be the set of the subsets
\(\mathscr K\subset\mathcal A\) having the following properties:
\begin{enumerate}[label=\arabic*\textsuperscript{o}.,leftmargin=*]
\item \(\mathscr K\) is nonempty, convex, and weakly compact;
\item \(A_U(\mathscr K)\subset\mathscr K\) for every
\(U\in\mathcal U\);
\item every element of \(\mathscr K\) has norm at most
\(\lVert\delta\rVert\).
\end{enumerate}
First, \(\mathscr E\ne\varnothing\). Indeed, let
\(\mathscr M\subset\mathcal A\) be the set of the
\(\delta(U)U^{-1}\), for \(U\in\mathcal U\). If
\(U,V\in\mathcal U\), then
\[
 A_V(\delta(U)U^{-1})
 =\bigl(V\delta(U)U^{-1}+\delta(V)\bigr)V^{-1}
 =\delta(VU)U^{-1}V^{-1},
\]
so that \(A_V(\mathscr M)\subset\mathscr M\). Consequently the weakly
closed convex hull of \(\mathscr M\) is an element of \(\mathscr E\).
If \((\mathscr K_i)\) is a totally ordered family of elements of
\(\mathscr E\), then \(\bigcap_i\mathscr K_i\in\mathscr E\). By
Zorn's theorem, \(\mathscr E\) has a minimal element
\(\mathscr K_0\).

Suppose it has been proved that \(\mathscr K_0\) consists of a single
point \(T_0\), or, equivalently, that
\(\mathscr K_0-\mathscr K_0=\{0\}\). Then, for every
\(U\in\mathcal U\),
\[
 T_0=A_U(T_0)=UT_0U^{-1}+\delta(U)U^{-1},
\]
and hence \(\delta(U)=[T_0,U]\). By linearity,
\(\delta(T)=[T_0,T]\) for every \(T\in\mathcal A\).

Observe that \(\mathscr K_0-\mathscr K_0\) is convex and weakly
compact, and that
\[
 U(\mathscr K_0-\mathscr K_0)U^{-1}
   =\mathscr K_0-\mathscr K_0
 \qquad(U\in\mathcal U),
\]
because
\[
 U(S-T)U^{-1}=A_U(S)-A_U(T)
 \qquad(S,T\in\mathscr K_0).
\]

\noindent\textup{2\textdegree.}
If \(\mathcal A\) is a product of von Neumann algebras
\(\mathcal A_1,\mathcal A_2\), then
\(\delta(\mathcal A_1)\subset\mathcal A_1\) and
\(\delta(\mathcal A_2)\subset\mathcal A_2\) by
\cref{lem:III-9-2}, and it is enough to prove the theorem for
\(\mathcal A_1\) and \(\mathcal A_2\). We may therefore suppose that
\(\mathcal A\) is semifinite or purely infinite. If \(\mathcal A\) is
semifinite, the set of projections \(E\) of \(\mathcal A\) such that
\(\mathcal A_E\) is finite is upward directed with supremum \(I\)
(\hyperref[prop:III-2-5]{\S~2, Proposition~5} and
\hyperref[cor:III-2-prop7-1]{Corollary~1 of Proposition~7}); thus it is
enough to consider the case in which \(\mathcal A\) is finite, by
\cref{lem:III-9-5}. Finally, by \cref{lem:III-9-5,lem:III-9-6}, it is
enough to
% Source PDF page 321 (printed p. 312)
\phantomsection\label{srcpage:312}
prove the theorem when \(\mathcal A\) is finite and of countable type,
and when \(\mathcal A\) is purely infinite and of countable type.

\noindent\textup{3\textdegree.}
Suppose that \(\mathcal A\) is finite and of countable type, and let
\(S,T\in\mathscr K_0\). There is a faithful finite normal trace on
\(\mathcal A\); this trace defines a pre-Hilbert norm
\(\lVert\,\cdot\,\rVert_2\) on \(\mathcal A\). Put
\[
 \alpha=\sup_{R\in\mathscr K_0}\lVert R\rVert_2<+\infty.
\]
Let \(\mathscr K\subset\mathscr K_0\) be the set of the
\[
 A_U\left(\frac12(S+T)\right),
 \qquad U\in\mathcal U,
\]
and let \(\mathscr K'\subset\mathscr K_0\) be the weakly closed convex
hull of \(\mathscr K\). By \eqref{eq:III-9-6}, \(\mathscr K\) is
invariant under the mappings \(A_U\), and hence so is
\(\mathscr K'\). Thus \(\mathscr K'=\mathscr K_0\), by the minimality
of \(\mathscr K_0\). For every \(\varepsilon>0\), there is an
\(R\in\mathscr K_0\) such that
\(\alpha-\varepsilon<\lVert R\rVert_2\). Since the norm
\(\lVert\,\cdot\,\rVert_2\) is lower semicontinuous for the weak
topology
(\hyperref[prop:I-6-2]{Chapter~I, \S~6, corollary to Proposition~2}),
there is a \(U\in\mathcal U\) such that
\[
 \alpha-\varepsilon
 <\left\lVert A_U\left(\frac12(S+T)\right)\right\rVert_2
 =\left\lVert\frac12\bigl(A_U(S)+A_U(T)\bigr)\right\rVert_2.
\]
Since \(\lVert A_U(S)\rVert_2\leq\alpha\) and
\(\lVert A_U(T)\rVert_2\leq\alpha\), and since
\(\lVert\,\cdot\,\rVert_2\) is a pre-Hilbert norm, one can choose
\(U\) so that \(\lVert A_U(S)-A_U(T)\rVert_2\) is arbitrarily small.
But
\[
 \lVert A_U(S)-A_U(T)\rVert_2
 =\lVert U(S-T)U^{-1}\rVert_2
 =\lVert S-T\rVert_2.
\]
Thus \(S=T\).

\noindent\textup{4\textdegree.}
Suppose that \(\mathcal A\) is purely infinite and of countable type,
and let \(S,T\in\mathscr K_0\). Let \(f\) be a weakly continuous
linear form on \(\mathcal A\), and put
\[
 \alpha=\sup_{R\in\mathscr K_0}\lvert f(R)\rvert<+\infty.
\]
Reasoning exactly as in \textup{3\textdegree}, for every
\(\varepsilon>0\) there is a \(U\in\mathcal U\) such that
\[
 \alpha-\varepsilon
 \leq\left\lvert\frac12
 \bigl(f(A_U(S))+f(A_U(T))\bigr)\right\rvert.
\]
Since \(\lvert f(A_U(S))\rvert\leq\alpha\) and
\(\lvert f(A_U(T))\rvert\leq\alpha\), one can choose \(U\) so that
\[
 \lvert f(A_U(S)-A_U(T))\rvert
 =\lvert f(U(S-T)U^{-1})\rvert
\]
is arbitrarily small. If
\(\mathscr K_0-\mathscr K_0\ne\{0\}\), then by
\hyperref[cor:III-8-thm1-6]{\S~8, Corollary~6 of Theorem~1} one can
choose \(S,T\in\mathscr K_0\) such that \(S-T\) is a nonzero element
of the centre of \(\mathcal A\). In that event,
\[
 U(S-T)U^{-1}=S-T
 \qquad(U\in\mathcal U),
\]
and the preceding argument proves that \(f(S-T)=0\). Hence
\(S-T=0\), a contradiction.
\end{proof}

\begin{corollary}\label{cor:III-9-thm1}
Let \(\mathfrak H\) be a Hilbert space, let \(\mathcal A\) be a
\(C^*\)-algebra of operators on \(\mathfrak H\), let \(\mathcal B\)
be the weak closure of \(\mathcal A\), and let \(\delta\) be a
derivation of \(\mathcal A\). There is a \(T_0\in\mathcal B\) such
that
\[
 \lVert T_0\rVert\leq\lVert\delta\rVert,
 \qquad
 \delta(T)=[T_0,T]\quad(T\in\mathcal A).
\]
\end{corollary}

\begin{proof}
This follows from \cref{lem:III-9-4,thm:III-9-1}.
\end{proof}

\noindent Bibliography: \ArticleRef{367}, \ArticleRef{368},
\ArticleRef{376}, \ArticleRef{404}, \ArticleRef{452}.

% Source PDF page 322 (printed p. 313)
\phantomsection\label{srcpage:313}
\NumberedParagraph{4}{Automorphisms of von Neumann Algebras}
\label{no:III-9-4}
In this number we shall make extensive use of the properties of the
holomorphic functional calculus (\TreatiseRef{9}, and N.~Bourbaki,
\emph{Theories spectrales}, Chapters~I and II, Paris, Hermann, 1967).

Let \(E\) be a complex Banach space. Denote by \(\mathcal L(E)\) the
set of continuous linear operators on \(E\). The usual norm on
\(\mathcal L(E)\) makes \(\mathcal L(E)\) a Banach space; the
corresponding topology on \(\mathcal L(E)\) is called the norm
topology. The set \(\operatorname{GL}(E)\) of invertible elements of
\(\mathcal L(E)\) is open in norm in \(\mathcal L(E)\).

Let \(\Delta\) be the set of \(z\in\mathbb C\) such that
\(-\pi<\operatorname{Im}z<\pi\), and let \(\Delta'\) be the set of
\(z\in\mathbb C\) that are not nonpositive real numbers. Let
\[
 \log\colon\Delta'\longrightarrow\Delta
\]
be the principal determination of the logarithm. It is a bijection of
\(\Delta'\) onto \(\Delta\), whose inverse is the exponential mapping
of \(\Delta\) onto \(\Delta'\).

Let \(A\) (respectively, \(A'\)) be the set of the
\(x\in\mathcal L(E)\) (respectively, \(x\in\operatorname{GL}(E)\))
such that
\[
 \operatorname{Sp}x\subset\Delta
 \quad\text{(respectively, }\operatorname{Sp}x\subset\Delta'\text{)}.
\]
Then \(A\) (respectively, \(A'\)) is an open subset in norm of
\(\mathcal L(E)\) (respectively, of \(\operatorname{GL}(E)\)), and the
mappings
\[
 x\longmapsto\exp x,
 \qquad
 y\longmapsto\log y
\]
are reciprocal homeomorphisms of \(A\) onto \(A'\), and of \(A'\)
onto \(A\).

\noindent[The operators \(\exp x\) and \(\log y\) are defined by
Cauchy integrals; for every \(x\in\mathcal L(E)\), \(\exp x\) may also
be defined by the series
\(\sum_{n\geq0}\frac1{n!}x^n\).]

\begin{lemma}\label{lem:III-9-7}
Let \(x\in\mathcal L(E)\), \(y=\exp x\), and \(\xi\in E\).
\begin{enumerate}[label=\textup{(\roman*)},leftmargin=2em]
\item If \(x\xi=0\), then \(y\xi=\xi\).
\item If \(y\xi=\xi\), and if \(\operatorname{Sp}x\) contains none
of the points
\[
 \pm2i\pi,\ \pm4i\pi,\ \pm6i\pi,\ldots,
\]
then \(x\xi=0\).
\end{enumerate}
\end{lemma}

\begin{proof}
Part (i) is immediate. Suppose that \(\operatorname{Sp}x\) satisfies
the condition in (ii). Let \(f\) be the entire function such that
\[
 \exp z-1=zf(z).
\]
We have \(y-I=(\exp x)-I=xf(x)\), and
\(0\notin f(\operatorname{Sp}x)=\operatorname{Sp}f(x)\), so that
\(f(x)\) is invertible. Thus the kernel of \(y-I\) is equal to the
kernel of \(x\), which proves (ii).
\end{proof}

\begin{lemma}\label{lem:III-9-8}
Let \(x\in\mathcal L(E)\) and \(y=\exp x\). Let \(F\) be the Banach
space of continuous bilinear mappings of \(E^2\) into \(E\). Define
\(x'\in\mathcal L(F)\) and \(y'\in\operatorname{GL}(F)\) by
\begin{align*}
 (x'f)(\xi_1,\xi_2)
 &=x(f(\xi_1,\xi_2))-f(x\xi_1,\xi_2)-f(\xi_1,x\xi_2),\\
 (y'f)(\xi_1,\xi_2)
 &=y\bigl(f(y^{-1}\xi_1,y^{-1}\xi_2)\bigr)
\end{align*}
for \(\xi_1,\xi_2\in E\) and \(f\in F\). Let \(f\in F\).
\begin{enumerate}[label=\textup{(\roman*)},leftmargin=2em]
\item If \(x'f=0\), then \(y'f=f\).
\item If \(y'f=f\), and if every \(z\in\operatorname{Sp}x\)
satisfies
\[
 \lvert\operatorname{Im}z\rvert<\frac{2\pi}{3},
\]
then \(x'f=0\).
\end{enumerate}
\end{lemma}

% Source PDF page 323 (printed p. 314)
\phantomsection\label{srcpage:314}
\begin{proof}
The mapping \(y\mapsto y'\) is a homomorphism of the Lie group
\(\operatorname{GL}(E)\) into the Lie group \(\operatorname{GL}(F)\).
It is classical that the corresponding homomorphism of the Lie algebra
\(\mathcal L(E)\) of \(\operatorname{GL}(E)\) into the Lie algebra
\(\mathcal L(F)\) of \(\operatorname{GL}(F)\) is \(x\mapsto x'\).
Since \(y=\exp x\), it follows that \(y'=\exp x'\). By
\cref{lem:III-9-7}, it now suffices to show that, if
\(\operatorname{Sp}x\) satisfies the condition in (ii), then
\[
 \pm2i\pi,\ \pm4i\pi,\ldots\notin\operatorname{Sp}x'.
\]
Define \(x_0,x_1,x_2\in\mathcal L(F)\) by
\begin{align*}
 (x_0f)(\xi_1,\xi_2)&=x(f(\xi_1,\xi_2)),\\
 (x_1f)(\xi_1,\xi_2)&=-f(x\xi_1,\xi_2),\\
 (x_2f)(\xi_1,\xi_2)&=-f(\xi_1,x\xi_2).
\end{align*}
Then \(x'=x_0+x_1+x_2\), and the \(x_i\) commute; hence
\[
 \operatorname{Sp}x'
 \subset\operatorname{Sp}x_0+
          \operatorname{Sp}x_1+
          \operatorname{Sp}x_2.
\]
The spectra of the \(x_i\) are contained in
\(\operatorname{Sp}x\cup(-\operatorname{Sp}x)\). Consequently, for
every \(z\in\operatorname{Sp}x'\),
\[
 \lvert\operatorname{Im}z\rvert
 <3\frac{2\pi}{3}=2\pi.
\]
\end{proof}

\begin{lemma}\label{lem:III-9-9}
Let \(E\) be a complex Banach algebra, let \(x\in\mathcal L(E)\), and
put \(y=\exp x\).
\begin{enumerate}[label=\textup{(\roman*)},leftmargin=2em]
\item If \(x\) is a derivation of \(E\), then \(y\) is an
automorphism of the algebra \(E\).
\item If \(y\) is an automorphism of the algebra \(E\), and if
\[
 \lvert\operatorname{Im}z\rvert<\frac{2\pi}{3}
 \qquad(z\in\operatorname{Sp}x),
\]
then \(x\) is a derivation of \(E\).
\end{enumerate}
\end{lemma}

\begin{proof}
Apply \cref{lem:III-9-8}, taking the multiplication of \(E\) for
\(f\).
\end{proof}

\begin{lemma}\label{lem:III-9-10}
Let \(E\) be a complex Banach algebra, and let \(y\) be an isometric
automorphism of the algebra \(E\) such that
\[
 \lVert y-I\rVert<\sqrt3.
\]
\begin{enumerate}[label=\textup{(\roman*)},leftmargin=2em]
\item Every \(z\in\operatorname{Sp}y\) belongs to the open angle
\[
 -\frac{2\pi}{3}<\arg z<\frac{2\pi}{3},
\]
which makes it possible to form \(x=\log y\).
\item The operator \(x\) is a derivation of \(E\).
\end{enumerate}
\end{lemma}

\begin{proof}
The conditions on \(y\) imply that every
\(z\in\operatorname{Sp}y\) satisfies
\(\lvert z\rvert=1\) and \(\lvert z-1\rvert<\sqrt3\), which proves
(i). We have
\(\operatorname{Sp}x=\log(\operatorname{Sp}y)\), and hence (ii)
follows from \cref{lem:III-9-9}.
\end{proof}

\begin{theorem}\label{thm:III-9-2}
Let \(\mathcal A\) be a von Neumann algebra, and let \(\Phi\) be an
automorphism of \(\mathcal A\). Suppose that
\[
 \lVert\Phi-I\rVert<\sqrt3.
\]
There is a unitary element \(U_0\) of \(\mathcal A\) such that
\[
 \Phi(T)=U_0TU_0^{-1}
 \qquad(T\in\mathcal A).
\]
\end{theorem}

% Source PDF page 324 (printed p. 315)
\phantomsection\label{srcpage:315}
\begin{proof}
Put \(\delta=\log\Phi\), which is a derivation of \(\mathcal A\) by
\cref{lem:III-9-10}. By \cref{thm:III-9-1}, there is a
\(T_0\in\mathcal A\) such that
\[
 \delta(T)=L_{T_0}(T)-R_{T_0}(T)
 \qquad(T\in\mathcal A),
\]
where \(L_{T_0}\) (respectively, \(R_{T_0}\)) denotes left
(respectively, right) multiplication by \(T_0\) in \(\mathcal A\).
Put \(S_0=\exp T_0\), which is an invertible element of
\(\mathcal A\). We have
\begin{align*}
 \Phi(T)
 &=\bigl(\exp(L_{T_0}+R_{-T_0})\bigr)(T)
   =\bigl(\exp L_{T_0}\,\exp R_{-T_0}\bigr)(T)\\
 &=\left(\sum_{m,n\geq0}\frac1{m!n!}
          L_{T_0}^{m}R_{-T_0}^{n}\right)(T)\\
 &=\sum_{m,n\geq0}\frac1{m!n!}T_0^mT(-T_0)^n\\
 &=\left(\sum_{m\geq0}\frac1{m!}T_0^m\right)
   T
   \left(\sum_{n\geq0}\frac1{n!}(-T_0)^n\right)\\
 &=S_0TS_0^{-1}.
\end{align*}
On the other hand, for every \(T\in\mathcal A\),
\[
 S_0T^*S_0^{-1}=\Phi(T^*)=\Phi(T)^*
 =S_0^{-1*}T^*S_0^*.
\]
Thus \(S_0^*S_0\) commutes with \(T^*\). Consequently
\(\lvert S_0\rvert\) belongs to the centre of \(\mathcal A\). If
\[
 S_0=U_0\lvert S_0\rvert
\]
is the polar decomposition of \(S_0\), then
\(\Phi(T)=U_0TU_0^{-1}\) for every \(T\in\mathcal A\).
\end{proof}

\noindent Bibliography: \ArticleRef{103}, \ArticleRef{200},
\ArticleRef{216}, \ArticleRef{369}, \ArticleRef{404}, \ArticleRef{405}.

Theorem~2 was first proved in \ArticleRef{404}, with the bound
\(\sqrt3\) replaced by \(2\). The method used in the text (due
essentially to J.-P.~Serre) is a little shorter, but does not yield the
bound \(2\).

There are factors possessing outer automorphisms (cf. \S~7,
Exercise~15).

% Source PDF page 326 (printed p. 317)
\chapter*{Appendix I}\label{srcpage:317}
\addcontentsline{toc}{chapter}{Appendix I}
\label{app:I}

\NumberedParagraph{1}{Spectrum}
\label{no:App-I-1}\label{sec:app-I-1}
For every locally compact space \(Z\), denote by \(L_\infty(Z)\) the
set of continuous complex-valued functions on \(Z\) that tend to zero
at infinity. This set is an abelian involutive algebra (the involution
being defined by passing from a function to its complex conjugate). If,
for every \(f\in L_\infty(Z)\), one puts
\[
 \lVert f\rVert=\sup_{\zeta\in Z}\lvert f(\zeta)\rvert,
\]
then \(L_\infty(Z)\) is at the same time a complete normed algebra.

Let \(\mathfrak H\) be a complex Hilbert space, and let
\(\mathcal Y\) be an abelian \(C^*\)-algebra of operators on
\(\mathfrak H\). Let \(Z\) be the set of characters of \(\mathcal Y\),
that is, the homomorphisms of \(\mathcal Y\) onto \(\mathbb C\).
Equipped with the weak topology of the dual of the Banach space
\(\mathcal Y\), it is a locally compact space. For every
\(f\in L_\infty(Z)\), there is a unique element \(T_f\in\mathcal Y\)
such that
\[
 \zeta(T_f)=f(\zeta)
 \qquad(\zeta\in Z),
\]
and \(f\mapsto T_f\) is an isomorphism of the normed involutive algebra
\(L_\infty(Z)\) onto the normed algebra \(\mathcal Y\), preserving the
natural order relations. The space \(Z\) is called the
\emph{spectrum of \(\mathcal Y\)}, and \(f\mapsto T_f\) the
\emph{Gelfand isomorphism} of \(L_\infty(Z)\) onto \(\mathcal Y\).

If \(\mathcal Y\) has an identity, then \(L_\infty(Z)\) has an
identity, which can only be the function \(1\) on \(Z\). Thus \(Z\) is
compact. Let \(\lambda\mapsto g(\lambda)\) be a continuous
complex-valued function of the complex variable \(\lambda\). If
\(f\in L_\infty(Z)\), then \(g\circ f\in L_\infty(Z)\), and one puts
\[
 g(T_f)=T_{g\circ f};
\]
this agrees with the usual definition when \(g\) is a polynomial.

Suppose that \(\mathcal Y\) has a countable base for the norm topology;
in other words, suppose that \(\mathcal Y\), as a
\(C^*\)-subalgebra of \(\mathcal L(\mathfrak H)\), is generated by a
countable family of elements. Then \(Z\) has a countable base. Indeed,
there is a sequence \((f_i)\) dense in \(L_\infty(Z)\) for the norm.
Let \(Z_i\) be the open subset of \(Z\) consisting of the
\(\zeta\in Z\) such that \(\lvert f_i(\zeta)\rvert>1\). We show that
the \(Z_i\) form a base for the topology of \(Z\). Let
\(\zeta_0\in Z\), and let \(Y\) be an open subset of \(Z\) containing
\(\zeta_0\). Let \(f\in L_\infty(Z)\) be equal to \(2\) at
\(\zeta_0\) and to zero on \(Z\setminus Y\), and choose \(i\) so that
\(\lVert f-f_i\rVert<1\). Then \(\zeta_0\in Z_i\subset Y\), which
proves the assertion.

\NumberedParagraph{2}{Spectral Measures}
\label{no:App-I-2}\label{sec:app-I-2}
Let \(x,y\in\mathfrak H\). The mapping
\[
 f\longmapsto(T_fx\mid y)
\]
is a linear form on \(L_\infty(Z)\), and
\[
 \lvert(T_fx\mid y)\rvert
 \leq\lVert f\rVert\lVert x\rVert\lVert y\rVert.
\]
% Source PDF page 327 (printed p. 318)
\phantomsection\label{srcpage:318}
There is therefore a unique bounded measure \(\nu_{x,y}\) on \(Z\),
called a \emph{spectral measure}, such that
\[
 \nu_{x,y}(f)=(T_fx\mid y)
 \qquad(f\in L_\infty(Z)).
\]
For \(x,x',y\in\mathfrak H\), \(\lambda,\lambda'\in\mathbb C\), and
\(g,g'\in L_\infty(Z)\), one has
\begin{align}
 \lVert\nu_{x,y}\rVert
 &\leq\lVert x\rVert\lVert y\rVert,
 \tag{1}\label{eq:App-I-1}\\
 \nu_{\lambda x+\lambda'x',y}
 &=\lambda\nu_{x,y}+\lambda'\nu_{x',y},
 \tag{2}\label{eq:App-I-2}\\
 \nu_{y,x}&=\overline{\nu_{x,y}},
 \tag{3}\label{eq:App-I-3}\\
 \nu_{x,x}&\geq0,
 \tag{4}\label{eq:App-I-4}\\
 \nu_{T_gx,T_{g'}y}&=g\overline{g'}\,\nu_{x,y}.
 \tag{5}\label{eq:App-I-5}
\end{align}
For example, let us prove formula \eqref{eq:App-I-5}. For
\(f\in L_\infty(Z)\),
\begin{align*}
 \nu_{T_gx,T_{g'}y}(f)
 &=(T_fT_gx\mid T_{g'}y)
  =(T_{g\overline{g'}f}x\mid y)\\
 &=\nu_{x,y}(g\overline{g'}f)
  =(g\overline{g'}\,\nu_{x,y})(f).
\end{align*}
Formulas \eqref{eq:App-I-2} and \eqref{eq:App-I-3} imply
\begin{equation}
 \begin{aligned}
 4\nu_{x,y}
 &=\nu_{x+y,x+y}-\nu_{x-y,x-y}\\
 &\quad+i\nu_{x+iy,x+iy}-i\nu_{x-iy,x-iy}.
 \end{aligned}
 \tag{6}\label{eq:App-I-6}
\end{equation}
Let \(S\in\mathcal L(\mathfrak H)\) and \(f\in L_\infty(Z)\). Then
\[
 (ST_fx\mid y)=(T_fx\mid S^*y)=\nu_{x,S^*y}(f),
 \qquad
 (T_fSx\mid y)=\nu_{Sx,y}(f).
\]
Thus \(S\) commutes with \(\mathcal Y\) if and only if
\[
 \nu_{x,S^*y}=\nu_{Sx,y}
 \qquad(x,y\in\mathfrak H).
\]
Finally, for every nonzero positive function \(f\in L_\infty(Z)\),
there is an \(x\in\mathfrak H\) such that \((T_fx\mid x)>0\). It
follows that the union of the supports of the \(\nu_{x,x}\) is dense
in \(Z\).

\NumberedParagraph{3}{Extension of the Gelfand Isomorphism}
\label{no:App-I-3}\label{sec:app-I-3}
Let \(\mathscr L\) be the set of all bounded complex functions on
\(Z\) measurable with respect to every spectral measure. Then
\(\mathscr L\) is plainly an involutive algebra, and \(L_\infty(Z)\)
is an involutive subalgebra of \(\mathscr L\). Give \(\mathscr L\) the
norm
\[
 \lVert f\rVert=\sup_{\zeta\in Z}\lvert f(\zeta)\rvert.
\]
If \(f\in\mathscr L\), then \(f\) is integrable for all spectral
measures (which are bounded), and
\[
 \left\lvert\int f(\zeta)\,d\nu_{x,y}(\zeta)\right\rvert
 \leq\lVert f\rVert\lVert\nu_{x,y}\rVert
 \leq\lVert f\rVert\lVert x\rVert\lVert y\rVert.
\]
Thus
\[
 (x,y)\longmapsto\int f(\zeta)\,d\nu_{x,y}(\zeta)
\]
is a continuous sesquilinear form on \(\mathfrak H\). Consequently
there is a unique operator \(T_f\in\mathcal L(\mathfrak H)\) such that
\[
 (T_fx\mid y)=\int f(\zeta)\,d\nu_{x,y}(\zeta).
\]
We have \(\lVert T_f\rVert\leq\lVert f\rVert\). When
\(f\in L_\infty(Z)\), the notation \(T_f\) agrees with the preceding
notation, by the very definition of the \(\nu_{x,y}\). In addition,
the following properties hold:
\begin{enumerate}[label=\textup{(\roman*)},leftmargin=2em]
\item \(T_{\lambda f+\lambda'f'}=\lambda T_f+\lambda'T_{f'}\)
\((f,f'\in\mathscr L;\ \lambda,\lambda'\in\mathbb C)\);
% Source PDF page 328 (printed p. 319)
\phantomsection\label{srcpage:319}
\item \(T_{\overline f}=T_f^*\);
\item \(T_{ff'}=T_fT_{f'}\);
\item \(T_f\in\mathcal Y''\);
\item \(\nu_{T_fx,y}=f\nu_{x,y}\).
\end{enumerate}

Property (i) is evident. We have
\begin{align*}
 (T_{\overline f}x\mid y)
 &=\int \overline{f}(\zeta)\,d\nu_{x,y}(\zeta)
  =\overline{\int f(\zeta)\,d\nu_{y,x}(\zeta)}\\
 &=\overline{(T_fy\mid x)}=(T_f^*x\mid y),
\end{align*}
whence (ii). If \(S\in\mathcal Y'\), then
\begin{align*}
 (T_fSx\mid y)
 &=\int f(\zeta)\,d\nu_{Sx,y}(\zeta)
  =\int f(\zeta)\,d\nu_{x,S^*y}(\zeta)\\
 &=(T_fx\mid S^*y)=(ST_fx\mid y),
\end{align*}
so that \(T_fS=ST_f\), whence (iv). If \(g\in L_\infty(Z)\),
\[
 (T_gT_fx\mid y)
 =\int f(\zeta)\,d\nu_{x,T_g^*y}(\zeta)
 =\int f(\zeta)g(\zeta)\,d\nu_{x,y}(\zeta),
\]
whence (v). Finally,
\[
 (T_{ff'}x\mid y)
 =\int f(\zeta)f'(\zeta)\,d\nu_{x,y}(\zeta)
 =\int f(\zeta)\,d\nu_{T_{f'}x,y}(\zeta)
 =(T_fT_{f'}x\mid y),
\]
whence (iii).

A subset \(Y\) of \(Z\) measurable with respect to every spectral
measure has as its characteristic function a function
\(\chi_Y\in\mathscr L\) such that
\(\chi_Y=\overline{\chi_Y}=\chi_Y^2\). It therefore corresponds to a
projection \(E_Y\in\mathcal Y''\). An operator in
\(\mathcal L(\mathfrak H)\) that commutes with all the \(E_Y\) commutes
with \(\mathcal Y\). Indeed, if \(g\in L_\infty(Z)\), then \(g\) is a
norm limit of finite linear combinations of functions \(\chi_Y\), and
therefore \(T_g\) is a norm limit of finite linear combinations of
projections \(E_Y\), which proves our assertion.

When \(\mathcal Y\) is the \(C^*\)-algebra generated by a single
Hermitian operator \(T\in\mathcal L(\mathfrak H)\), the projections
\(E_Y\) are called the \emph{spectral projections of \(T\)}. They
commute with every operator that commutes with \(T\); and every operator
that commutes with the spectral projections of \(T\) commutes with
\(T\).

\noindent\textbf{Bibliography:} \ArticleRef{18}, \ArticleRef{25},
\ArticleRef{28}, \ArticleRef{45}, \ArticleRef{100}, \TreatiseRef{9}.

\noindent\textit{Exercises.}---
\begin{enumerate}[label=\arabic*.,leftmargin=2em]
\item\label{ex:App-I-1} Let \(T\) be a Hermitian operator in
\(\mathcal L(\mathfrak H)\), and let \(\mathcal Y\) be the
\(C^*\)-subalgebra of \(\mathcal L(\mathfrak H)\) generated by \(T\)
and \(I\). Show that the mapping \(\zeta\mapsto\zeta(T)\) is a
homeomorphism of the spectrum \(Z\) of \(\mathcal Y\) onto a compact
subset of \(\mathbb R\), the spectrum of \(T\) in the usual sense, and
that it transforms the function on \(Z\) associated with \(T\) into
the function \(\zeta\mapsto\zeta\).

\item\label{ex:App-I-2} Let
\(\mathfrak H_1=L^2([0,1])\) for Lebesgue measure, and let \(T_1\) be
the operator on \(\mathfrak H_1\) of multiplication by the function
\(x\mapsto x\). Let \(\mathfrak H_2\) be the Hilbert space of
complex-valued functions \(f\) on \([0,1]\) such that
\[
 \sum_{x\in[0,1]}\lvert f(x)\rvert^2<+\infty,
\]
 and let \(T_2\) be the operator on \(\mathfrak H_2\) of multiplication
 by the function \(x\mapsto x\).
% Source PDF page 329 (printed p. 320)
\phantomsection\label{srcpage:320}
Let \(\mathfrak H=\mathfrak H_1\oplus\mathfrak H_2\), let \(T\) be
the continuous linear operator on \(\mathfrak H\) extending \(T_1\)
and \(T_2\), and let \(\mathcal Y\) be the \(C^*\)-algebra of operators
generated by \(T\). Show that the spectrum of \(\mathcal Y\) may be
identified with \([0,1]\) (cf. \hyperref[ex:App-I-1]{Exercise~1}), and
that a function measurable with respect to every spectral measure is
Lebesgue measurable. Deduce that the extension of the Gelfand
isomorphism studied in \hyperref[sec:app-I-3]{no.~3} does not have as
its image the von Neumann algebra generated by \(\mathcal Y\).
\end{enumerate}

\noindent\textbf{Bibliography:} \TreatiseRef{15}, p.~65.
(Cf. \hyperref[sec:I-7]{Chapter~I, \S~7},
\hyperref[prop:I-7-1]{Proposition~1}.)

% Source PDF page 330 (printed p. 321)
\chapter*{Appendix II}\label{srcpage:321}
\addcontentsline{toc}{chapter}{Appendix II}
\label{app:II}

Let \(\mathfrak H\) be a complex Hilbert space, and let \(\mathcal M\)
be the set of Hermitian operators in \(\mathcal L(\mathfrak H)\).
Recall that \(\mathcal M\) is equipped with a natural order relation.
[If \(S,T\in\mathcal M\), one writes \(S\geq T\) when
\((Sx\mid x)\geq(Tx\mid x)\) for every \(x\in\mathfrak H\).] Denote by
\(\mathcal N\) the set of \(T\in\mathcal M\) such that
\(0\leq T\leq I\).

Let \(\mathcal F\subset\mathcal M\). Suppose that \(\mathcal F\) is
increasingly directed; that is, for every pair \(S_1,S_2\in\mathcal F\),
there is an \(S\in\mathcal F\) such that \(S\geq S_1\) and
\(S\geq S_2\). Suppose, moreover, that \(\mathcal F\) is bounded above;
that is, there is a Hermitian operator majorizing all the operators in
\(\mathcal F\). Then:

\begin{quote}
\itshape
The set \(\mathcal F\) has a supremum in \(\mathcal M\), and this
supremum belongs to the strong closure of \(\mathcal F\).
\end{quote}

For every \(S\in\mathcal F\), let \(\mathcal F_S\) be the set of
\(T\in\mathcal F\) majorizing \(S\). It is enough to prove our assertion
for \(\mathcal F_S\) in place of \(\mathcal F\). We may therefore
suppose that \(\mathcal F\) is bounded below as well as above and hence
reduce to the case \(\mathcal F\subset\mathcal N\). On the other hand,
the \(\mathcal F_S\) form a filter base. Since \(\mathcal N\) is weakly
compact, this filter base has a weak cluster point
\(T_0\in\mathcal M\). For every \(S\in\mathcal F\), the set of Hermitian
operators majorizing \(S\) is weakly closed and contains
\(\mathcal F_S\); hence \(T_0\geq S\). Thus \(T_0\) majorizes
\(\mathcal F\) and is a weak cluster point of \(\mathcal F\). If \(T_1\)
is another Hermitian operator majorizing \(\mathcal F\), then \(T_1\)
majorizes the weak closure of \(\mathcal F\), and therefore \(T_0\).
Thus \(T_0\) is the supremum of \(\mathcal F\). Finally, if
\(x_1,x_2,\ldots,x_n\in\mathfrak H\) and \(T\in\mathcal F\), then
\[
 \lVert(T_0-T)x_i\rVert^2
 \leq\lVert(T_0-T)^{1/2}x_i\rVert^2
 =((T_0-T)x_i\mid x_i)
\]
for every \(i\). Hence the quantities \(\lVert(T_0-T)x_i\rVert^2\)
can be made arbitrarily small, which proves that \(T_0\) belongs to
the strong closure of \(\mathcal F\).

In the course of the proof we have also established the following
point: if an upper bound of \(\mathcal F\) is a weak cluster point of
\(\mathcal F\), then that upper bound is the supremum of \(\mathcal F\).

\noindent\textbf{Bibliography:} J.-P.~Vigier,
\emph{Étude sur les suites infinies d'opérateurs hermitiens}, thesis
no.~1089, Geneva, 1946, p.~12.

% Source PDF page 332 (printed p. 323)
\chapter*{Appendix III}\label{srcpage:323}
\addcontentsline{toc}{chapter}{Appendix III}
\label{app:III}

\NumberedParagraph{1}{Support of an Operator}
\label{no:App-III-1}\label{sec:app-III-1}
Let \(\mathfrak H\) be a complex Hilbert space, and let
\(T\in\mathcal L(\mathfrak H)\). The set of \(x\in\mathfrak H\) such
that \(Tx=0\) is a closed vector subspace \(\mathfrak N\) of
\(\mathfrak H\). Let \(\mathfrak X\) be its orthogonal complement. The
projection \(E=P_{\mathfrak X}\) is called the \emph{support of \(T\)}.
For every \(y\in\mathfrak H\), we have
\(y-Ey\in\mathfrak N\), and therefore \(Ty-TEy=0\). Hence \(T=TE\).
Conversely, let \(E_1\) be a projection in
\(\mathcal L(\mathfrak H)\) such that \(TE_1=T\). For every
\(z\in\mathfrak H\), we have \(T(I-E_1)z=0\), and therefore
\((I-E_1)z\in\mathfrak N=\mathfrak X^\perp\); hence
\(I-E_1\leq I-E\), so that \(E_1\geq E\). Thus \(E\) is the smallest
of the projections \(E_1\in\mathcal L(\mathfrak H)\) such that
\(TE_1=T\).

The closure of \(T(\mathfrak H)\) is a closed vector subspace
\(\mathfrak Y\) of \(\mathfrak H\). Let
\(F=P_{\mathfrak Y}\). Clearly, \(F\) is the smallest of the
projections \(F_1\in\mathcal L(\mathfrak H)\) such that \(F_1T=T\), or,
equivalently, such that \(T^*F_1=T^*\). Thus \(F\) is the support of
\(T^*\).

\NumberedParagraph{2}{Partially Isometric Operators}
\label{no:App-III-2}\label{sec:app-III-2}
Let \(U\in\mathcal L(\mathfrak H)\), and let \(E\) be its support. The
operator \(U\) is said to be \emph{partially isometric} if \(U\) is
isometric on \(\mathfrak X=E(\mathfrak H)\). Then
\(U(\mathfrak H)=U(\mathfrak X)\) is a closed vector subspace
\(\mathfrak Y\) of \(\mathfrak H\), and \(U\) maps
\(\mathfrak X\) isometrically onto \(\mathfrak Y\). Let
\(F=P_{\mathfrak Y}\). The projection \(E\) (respectively, the
subspace \(\mathfrak X\)) is called the \emph{initial projection}
(respectively, the \emph{initial subspace}) of \(U\), and \(F\)
(respectively, \(\mathfrak Y\)) is called its \emph{final projection}
(respectively, its \emph{final subspace}).

Let \(x\in\mathfrak X\), and put \(y=Ux\in\mathfrak Y\). For every
\(z\in\mathfrak H\),
\[
 (x\mid z)=(x\mid Ez)=(Ux\mid UEz)=(y\mid Uz),
\]
and therefore \(x=U^*y\). Thus the mapping \(x\mapsto Ux\) from
\(\mathfrak X\) onto \(\mathfrak Y\) has as its inverse the isometric
mapping \(y\mapsto U^*y\) from \(\mathfrak Y\) onto
\(\mathfrak X\). Since, moreover, the support of \(U^*\) is \(F\), it
follows that \(U^*\) is partially isometric, with initial projection
\(F\) and final projection \(E\). We also see that
\[
 U^*U=E,\qquad UU^*=F.
\]

Conversely, let \(V\in\mathcal L(\mathfrak H)\) be such that \(V^*V\)
is a projection \(G\). Then, for every \(x\in\mathfrak H\),
\[
 \lVert Vx\rVert^2=(V^*Vx\mid x)=(Gx\mid x)=\lVert Gx\rVert^2;
\]
% Source PDF page 333 (printed p. 324)
\phantomsection\label{srcpage:324}
hence \(V\) is isometric on \(G(\mathfrak H)\) and vanishes on
\(\mathfrak H\ominus G(\mathfrak H)\), which proves that \(V\) is
partially isometric. Similarly, if \(W\in\mathcal L(\mathfrak H)\) and
\(WW^*\) is a projection, then \(W^*\) is partially isometric, and
therefore so is \(W\).

\NumberedParagraph{3}{Polar Decomposition of an Operator}
\label{no:App-III-3}\label{sec:app-III-3}
Let \(T\in\mathcal L(\mathfrak H)\), let \(E\) be the support of
\(T\), let \(F\) be the support of \(T^*\), and put
\(\mathfrak X=E(\mathfrak H)\), \(\mathfrak Y=F(\mathfrak H)\). Put
\[
 \lvert T\rvert=(T^*T)^{1/2}.
\]
For every \(x\in\mathfrak H\),
\[
 \lVert Tx\rVert^2=(T^*Tx\mid x)=\lVert\lvert T\rvert x\rVert^2.
\]
Consequently, \(\lvert T\rvert\) has support \(E\), and hence the
closure of \(\lvert T\rvert(\mathfrak H)\) is \(\mathfrak X\).
Moreover, the mapping \(\lvert T\rvert x\mapsto Tx\) is a linear
isometry from \(\lvert T\rvert(\mathfrak H)\) onto
\(T(\mathfrak H)\), and therefore extends to a linear isometry \(V\)
from \(\mathfrak X\) onto \(\mathfrak Y\). Let \(U\) be the partially
isometric operator of support \(E\) which agrees with \(V\) on
\(\mathfrak X\); this partially isometric operator has initial
projection \(E\) and final projection \(F\). We have
\[
 T=U\lvert T\rvert,
\]
an equality called the \emph{polar decomposition of \(T\)}.

Conversely, if \(T=U_1T_1\), where \(T_1\) is positive Hermitian and
\(U_1\) is a partially isometric operator whose initial projection is
the support of \(T_1\), then
\[
 T^*T=T_1U_1^*U_1T_1=T_1^2,
\]
so that \(T_1=\lvert T\rvert\), and then \(U_1=U\).

We have
\[
 T^*=\lvert T\rvert U^*=U^*(U\lvert T\rvert U^*).
\]
The operator \(U\lvert T\rvert U^*\) is positive Hermitian with
support \(F\), and \(U^*\) is partially isometric, with initial
projection \(F\) and final projection \(E\). Thus the equality
\(T^*=U^*(U\lvert T\rvert U^*)\) is the polar decomposition of
\(T^*\). Hence
\[
 \lvert T^*\rvert=U\lvert T\rvert U^*,
 \qquad
 \lvert T\rvert=U^*\lvert T^*\rvert U.
\]

\noindent\textbf{Bibliography:} J.~von Neumann,
\emph{Ueber adjungierte Funktionaloperatoren} (\emph{Ann. Math.},
vol.~33, 1932, pp.~294--310).

% Source PDF page 334 (printed p. 325)
\chapter*{Appendix IV}\label{srcpage:325}
\addcontentsline{toc}{chapter}{Appendix IV}
\label{app:IV}

If \(Z\) is a locally compact space and \(\nu\) a positive measure on
\(Z\), we shall denote by \(L_{\mathbb C}^{\infty}(Z,\nu)\) the set of
complex-valued functions on \(Z\) which are measurable and essentially
bounded with respect to \(\nu\), two functions equal locally almost
everywhere being identified; we shall denote by \(L_\infty(Z)\) the set
of continuous complex-valued functions on \(Z\) which tend to zero at
infinity.

The following result is due to J.~von Neumann [\emph{Einige Sätze über
messbare Abbildungen} (\emph{Ann. Math.}, vol.~33, 1932,
pp.~574--586)] in the more general case where \(Z\) and \(Z_1\) are
complete metric spaces with countable bases.

\phantomsection\label{thm:App-IV-main}
\begin{quote}
\itshape
Let \(Z\) and \(Z_1\) be two locally compact spaces with countable
bases, let \(\nu\) be a positive measure on \(Z\), let \(\nu_1\) be a
positive measure on \(Z_1\), and let \(\Phi\) be an isomorphism of the
involutive algebra \(L_{\mathbb C}^{\infty}(Z,\nu)\) onto the involutive
algebra \(L_{\mathbb C}^{\infty}(Z_1,\nu_1)\). Then there exist:

\begin{enumerate}[label=\arabic*\textsuperscript{\(\circ\)},leftmargin=2em]
\item a \(\nu\)-negligible set \(N\) in \(Z\) and a
\(\nu_1\)-negligible set \(N_1\) in \(Z_1\);
\item a Borel isomorphism \(\eta\) of \(Z\setminus N\) onto
\(Z_1\setminus N_1\) which transforms \(\nu\) into a measure equivalent
to \(\nu_1\), and such that
\[
 (\Phi(f))(\eta(\zeta))=f(\zeta)
\]
for every \(f\in L_{\mathbb C}^{\infty}(Z,\nu)\) and every
\(\zeta\in Z\setminus N\).
\end{enumerate}
\end{quote}

We begin with a few remarks.

\begin{enumerate}[label=\arabic*\textsuperscript{\(\circ\)},leftmargin=2em]
\item \(\Phi\) is isometric. Indeed, \(\Phi\) transforms positive
functions, which are of the form \(g\overline g\), into positive
functions; and \(\lVert f\rVert_\infty\) is the least number
\(a\geq0\) such that \(a^2-f\overline f\geq0\).

\item We may replace \(\nu\) and \(\nu_1\) by equivalent positive
measures and hence, in particular (since \(Z\) and \(Z_1\) are
countable at infinity), by bounded measures. Then, adjoining points at
infinity carrying zero mass to \(Z\) and \(Z_1\), we see that it is
enough to consider the case where \(Z\) and \(Z_1\) are compact (with
countable bases), as we shall henceforth suppose.

\item Let \(Y\) be the support of \(\nu\), and let \(Y_1\) be the
support of \(\nu_1\). It is enough to prove the proposition for the
spaces \(Y,Y_1\) and the measures induced on \(Y,Y_1\) by
\(\nu,\nu_1\). We shall therefore henceforth suppose that the supports
of \(\nu\) and \(\nu_1\) are \(Z\) and \(Z_1\), which permits us to
identify \(L_\infty(Z)\) and \(L_\infty(Z_1)\) with involutive
subalgebras of \(L_{\mathbb C}^{\infty}(Z,\nu)\) and
\(L_{\mathbb C}^{\infty}(Z_1,\nu_1)\), respectively.
\end{enumerate}

% Source PDF page 335 (printed p. 326)
\phantomsection\label{srcpage:326}
Let \(A\) be the closed involutive subalgebra (for the topology induced
by the norm) of \(L_{\mathbb C}^{\infty}(Z,\nu)\) generated by
\(L_\infty(Z)\) and \(\Phi^{-1}(L_\infty(Z_1))\). Let
\(A_1=\Phi(A)\), which is the closed involutive subalgebra of
\(L_{\mathbb C}^{\infty}(Z_1,\nu_1)\) generated by \(L_\infty(Z_1)\)
and \(\Phi(L_\infty(Z))\). Then \(A\) and \(A_1\) have countable bases
for the topology induced by the norm. Let \(Z'\) be the spectrum of the
Banach algebra \(A\). It is a compact space with a countable base. Let
\(\Theta\) be the Gelfand isomorphism of \(L_\infty(Z')\) onto \(A\).
The linear form
\[
 f\longmapsto\int f\,d\nu
\]
on \(A\), transported to \(L_\infty(Z')\) by \(\Theta^{-1}\), is a
positive measure \(\nu'\) on \(Z'\) with support \(Z'\). The canonical
injection of \(L_\infty(Z)\) into \(A\) defines, by transposition, a
continuous mapping \(\theta\) from \(Z'\) into \(Z\); for
\(f\in L_\infty(Z)\),
\[
 f\circ\theta=\Theta^{-1}(f),
\]
so that \(\theta\) is necessarily surjective. Clearly
\(\theta(\nu')=\nu\). Every function \(f\in A\) defines, on the one
hand, the function
\(g=f\circ\theta\in L_{\mathbb C}^{\infty}(Z',\nu')\), and, on the
other, the function \(h=\Theta^{-1}(f)\in L_\infty(Z')\). We shall show
that \(g-h\) is \(\nu'\)-negligible. We already know that \(g=h\) if
\(f\in L_\infty(Z)\). In the general case, let \((f_n)\) be a sequence
of functions in \(L_\infty(Z)\) such that
\[
 \int\lvert f-f_n\rvert^2\,d\nu\longrightarrow0,
\]
and put \(g_n=f_n\circ\theta\). We have
\[
 \int\lvert g-g_n\rvert^2\,d\nu'
 =\int\lvert f-f_n\rvert^2\,d\nu,
 \qquad\text{since }\nu=\theta(\nu').
\]
Also,
\[
 \int(h-g_n)(\overline h-\overline{g_n})\,d\nu'
 =\int(f-f_n)(\overline f-\overline{f_n})\,d\nu
\]
by the definition of \(\nu'\). Thus \(g\) and \(h\) are both limits
of \(g_n\) in \(L_{\mathbb C}^{2}(Z',\nu')\), which proves our
assertion.

Since \(Z\) and \(Z'\) are compact spaces with countable bases and
\(\nu=\theta(\nu')\), there is a disintegration
\(\zeta\mapsto\lambda_\zeta\) of \(\nu'\) relative to \(\nu\)
(\TreatiseRef{4}, Chapter~VI, \S~3, Theorem~1). Let \(d\) be a metric
on \(Z'\) compatible with the topology. Let \(p\) be a positive
integer. We shall show that the support of \(\lambda_\zeta\) has
diameter at most \(1/p\), except when \(\zeta\) belongs to a negligible
set \(N_p\). Indeed, let \(h_1,h_2,\ldots,h_n\) be positive functions
in \(L_\infty(Z')\) whose supports \(X_1,X_2,\ldots,X_n\) have
diameter at most \(1/p\), and such that
\[
 \sum_{i=1}^n h_i>0
\]
at every point of \(Z'\). Put
\[
 f_i=\Theta(h_i),\qquad g_i=f_i\circ\theta,
 \qquad k_i=\inf(g_i,h_i).
\]
The functions \(g_i,h_i,k_i\) are equal almost everywhere for
\(\nu'\). We have \(\sum_{i=1}^n g_i>0\) almost everywhere. Hence, if
\(Y_i\) denotes the set of \(\zeta\in Z\) such that
% Source PDF page 336 (printed p. 327)
\phantomsection\label{srcpage:327}
\(f_i(\zeta)>0\), then \(Z\) is the union of the \(Y_i\), up to a
\(\nu\)-negligible set. On the other hand,
\begin{align*}
 \int d\nu(\zeta)\int k_i(\zeta')\,d\lambda_\zeta(\zeta')
 &=\int k_i(\zeta')\,d\nu'(\zeta')\\
 &=\int g_i(\zeta')\,d\nu'(\zeta')\\
 &=\int d\nu(\zeta)\int g_i(\zeta')\,d\lambda_\zeta(\zeta').
\end{align*}
Consequently, taking into account that \(k_i\leq g_i\),
\[
 \int k_i(\zeta')\,d\lambda_\zeta(\zeta')
 =\int g_i(\zeta')\,d\lambda_\zeta(\zeta')
\]
for almost every \(\zeta\in Z\). But, for
\(\zeta'\in\theta^{-1}(Y_i)\cap(Z'\setminus X_i)\), we have
\[
 k_i(\zeta')=0
 \qquad\text{and}\qquad
 g_i(\zeta')>0.
\]
Thus, for almost every \(\zeta\in Z\), \(\lambda_\zeta\) is
concentrated on
\[
 (Z'\setminus\theta^{-1}(Y_i))\cup X_i.
\]
If, in addition, \(\zeta\in Y_i\), then \(\lambda_\zeta\) is
concentrated on \(\theta^{-1}(Y_i)\); consequently, for almost every
\(\zeta\in Y_i\), \(\lambda_\zeta\) is concentrated on \(X_i\), and
its support therefore has diameter at most \(1/p\). Since this is true
for every \(i\), our assertion is proved.

Let \(N\) be the \(\nu\)-negligible set which is the union of the
\(N_p\). For \(\zeta\in Z\setminus N\), \(\lambda_\zeta\) has support
of diameter zero and is therefore a point measure. For every
\(\zeta\in Z\setminus N\), there is a point
\(\theta'(\zeta)\in\theta^{-1}(\{\zeta\})\) such that
\[
 \lambda_\zeta=\varepsilon_{\theta'(\zeta)}.
\]
Since \(\lambda_\zeta\) depends measurably on \(\zeta\), so does
\(\theta'(\zeta)\). The image of \(\nu\) under \(\theta'\) is
\[
 \int\varepsilon_{\theta'(\zeta)}\,d\nu(\zeta)
 =\int\lambda_\zeta\,d\nu(\zeta)=\nu'.
\]
The set \(N'=Z'\setminus\theta'(Z)\), whose inverse image under
\(\theta'\) is empty, is therefore \(\nu'\)-negligible. Thus
\(\theta'\) and the restriction of \(\theta\) to \(Z'\setminus N'\)
are inverse mappings from \(Z\setminus N\) onto \(Z'\setminus N'\)
and from \(Z'\setminus N'\) onto \(Z\setminus N\), transforming
\(\nu\) into \(\nu'\) and \(\nu'\) into \(\nu\); and the isomorphism
of \(L_{\mathbb C}^{\infty}(Z,\nu)\) onto
\(L_{\mathbb C}^{\infty}(Z',\nu')\) which they define reduces to
\(\Theta^{-1}\) on \(A\). Moreover, \(Z\) is the union of a sequence
of pairwise disjoint subsets \(Z_0,Z_1,Z_2,\ldots\), where \(Z_0\) is
\(\nu\)-negligible and \(\theta'|_{Z_i}\) is continuous for
\(i\geq1\). Thus, by enlarging \(N\) and \(N'\), we may suppose that
\(\theta'|_{Z\setminus N}\) is a Borel isomorphism of
\(Z\setminus N\) onto \(Z'\setminus N'\).

Reasoning in the same way with \(A_1\), we may define on the spectrum
\(Z_1'\) of \(A_1\) a positive measure \(\nu_1'\) with support
\(Z_1'\), a \(\nu_1\)-negligible set \(N_1\) in \(Z_1\), a
\(\nu_1'\)-negligible set \(N_1'\) in \(Z_1'\), and reciprocal Borel
isomorphisms \(\theta_1'\), \(\theta_1\), respectively from
\(Z_1\setminus N_1\) onto \(Z_1'\setminus N_1'\) and from
\(Z_1'\setminus N_1'\) onto \(Z_1\setminus N_1\), which transform
\(\nu_1\) into \(\nu_1'\) and \(\nu_1'\) into \(\nu_1\), and which
define an isomorphism of
\(L_{\mathbb C}^{\infty}(Z_1',\nu_1')\) onto
\(L_{\mathbb C}^{\infty}(Z_1,\nu_1)\) reducing on \(A_1\) to
\(\Theta_1^{-1}\), where \(\Theta_1\) denotes the Gelfand isomorphism
of \(L_\infty(Z_1')\) onto \(A_1\).

% Source PDF page 337 (printed p. 328)
\phantomsection\label{srcpage:328}
We have thus reduced the proof of the proposition to the case where
\[
 A=L_\infty(Z),\qquad A_1=L_\infty(Z_1).
\]
The isomorphism of \(L_\infty(Z)\) onto \(L_\infty(Z_1)\) induced by
\(\Phi\) then defines a homeomorphism \(\eta\) of \(Z\) onto \(Z_1\)
such that
\[
 (\Phi(f))(\eta(\zeta))=f(\zeta)
\]
for every \(f\in L_\infty(Z)\) and every \(\zeta\in Z\). Identifying
\(Z\) with \(Z_1\) by \(\eta\), we are reduced to the case where
\(Z=Z_1\) and where \(\Phi\) induces the identity on \(L_\infty(Z)\).
It is then enough to prove the following lemma, which is valid without
any countability hypothesis.

\par\addvspace{.55\baselineskip}
\noindent\textbf{Lemma.} --- \emph{%
Let \(Z\) be a locally compact space, let \(\nu\) and \(\nu_1\) be
positive measures on \(Z\) with support \(Z\), and let \(\Phi\) be an
isomorphism of \(L_{\mathbb C}^{\infty}(Z,\nu)\) onto
\(L_{\mathbb C}^{\infty}(Z,\nu_1)\) which reduces to the identity on
\(L_\infty(Z)\). Then \(\nu\) and \(\nu_1\) are equivalent
[so that
\(L_{\mathbb C}^{\infty}(Z,\nu)=L_{\mathbb C}^{\infty}(Z,\nu_1)\)],
and \(\Phi\) is the identity.}
\par\addvspace{.55\baselineskip}

Let \(f\) be a bounded positive lower semicontinuous function on
\(Z\). Then \(f\) is the supremum, both in
\(L_{\mathbb C}^{\infty}(Z,\nu)\) and in
\(L_{\mathbb C}^{\infty}(Z,\nu_1)\), of the functions in
\(L_\infty(Z)\) majorized by \(f\). Hence \(\Phi(f)=f\). Now let \(g\)
be a positive function in \(L_{\mathbb C}^{\infty}(Z,\nu)\) with
compact support. There is a decreasing sequence of bounded lower
semicontinuous functions \(g_n\geq g\) such that
\[
 \int g_n\,d\nu\longrightarrow\int g\,d\nu,
 \qquad
 \int g_n\,d\nu_1\longrightarrow\int g\,d\nu_1.
\]
Then \(g\) is the infimum of the \(g_n\), both in
\(L_{\mathbb C}^{\infty}(Z,\nu)\) and in
\(L_{\mathbb C}^{\infty}(Z,\nu_1)\). Consequently:
\begin{enumerate}[label=\arabic*\textsuperscript{\(\circ\)},leftmargin=2em]
\item if \(g\) is \(\nu\)-negligible, it is \(\nu_1\)-negligible,
and conversely, which proves that \(\nu\) and \(\nu_1\) are
equivalent;
\item \(\Phi(g)=g\).
\end{enumerate}
Finally, every positive function
\(h\in L_{\mathbb C}^{\infty}(Z,\nu)
=L_{\mathbb C}^{\infty}(Z,\nu_1)\) is the supremum in
\(L_{\mathbb C}^{\infty}(Z,\nu)\) of the positive functions with
compact support in \(L_{\mathbb C}^{\infty}(Z,\nu)\) which it
majorizes; hence \(\Phi(h)=h\). \hfill\rule{0.55ex}{1.5ex}

% Source PDF page 338 (printed p. 329)
\chapter*{Appendix V}\label{srcpage:329}
\addcontentsline{toc}{chapter}{Appendix V}
\label{app:V}

\NumberedParagraph{1}{Borel Sets}
\label{no:App-V-1}\label{sec:app-V-1}
Let \(E\) be a topological space. A collection \(\mathscr T\) of
subsets of \(E\) is called a \emph{sigma-algebra} if it satisfies the
following conditions: (a) the complement of every set in \(\mathscr T\)
belongs to \(\mathscr T\); (b) every countable intersection of sets in
\(\mathscr T\) belongs to \(\mathscr T\). Every intersection of
sigma-algebras on \(E\) is a sigma-algebra on \(E\). There is therefore
a smallest sigma-algebra containing the set of closed subsets of \(E\);
the sets in this sigma-algebra are called the \emph{Borel subsets} of
\(E\). If \(E\) is Hausdorff, the subsets of \(E\) of the form
\[
 A_1\cap A_2\cap A_3\cap\cdots,
\]
where each \(A_i\) is a countable union of compact subsets of \(E\), are
Borel subsets of \(E\), called \(K_{\sigma\delta}\) sets in \(E\).

\NumberedParagraph{2}{Polish Spaces}
\label{no:App-V-2}\label{sec:app-V-2}
A topological space \(E\) is called \emph{Polish} if its topology has a
countable base and there is a metric compatible with the topology of
\(E\) for which \(E\) is complete.

Every countable product of Polish spaces is a Polish space.

Every space \(E\) which is the topological sum of a countable family
\((F_n)\) of Polish spaces is Polish. Indeed, for every \(n\), let
\(d_n\) be a metric on \(F_n\), compatible with the topology of \(F_n\)
and for which \(F_n\) is complete, such that
\(d_n(x,y)\leq1\) on \(F_n\times F_n\). Define on \(E\) a metric
\(d\), compatible with the topology of \(E\), by putting
\(d(x,y)=1\) if \(x\) and \(y\) do not belong to the same set \(F_n\),
and \(d(x,y)=d_n(x,y)\) if \(x\) and \(y\) both belong to \(F_n\).
It is immediately verified that every Cauchy sequence in \(E\)
converges and that \(E\) has a countable dense subset; hence our
assertion.

Let \(E\) be a Polish space and \(F\) a subspace of \(E\). If \(F\) is
closed in \(E\), then \(F\) is Polish. Suppose now that \(F\) is open
in \(E\). Let \(F'=E\setminus F\), and let \(d\) be a metric on \(E\)
compatible with the topology of \(E\) and for which \(E\) is complete.
For \(x,y\in F\), put
\[
 \delta(x,y)=d(x,y)
 +\left\lvert d(x,F')^{-1}-d(y,F')^{-1}\right\rvert.
\]
It is immediate that \(\delta\) is a metric on \(F\) compatible with
the topology of \(F\). On the other hand, if \((x_n)\) is a Cauchy
sequence in \(F\) for \(\delta\), then \(x_n\) converges in \(E\) for
\(d\) to an element \(x\); the numbers \(d(x_n,F')\) are bounded below
by a positive number, and therefore \(x\in F\); moreover, \(x_n\)
converges to \(x\) for \(\delta\). Thus \(F\) is complete for
\(\delta\), and hence \(F\) is Polish.

% Source PDF page 339 (printed p. 330)
\phantomsection\label{srcpage:330}
The set \(\mathbb N\) of natural integers, considered as a subspace of
\(\mathbb R\), is a complete metric space. Consequently
\(\mathbb N^{\mathbb N}\), equipped with the metric
\[
 d\bigl((m_i),(n_i)\bigr)
 =\sum_{i=0}^{\infty}2^{-i}
   \frac{\lvert m_i-n_i\rvert}{1+\lvert m_i-n_i\rvert},
\]
is a complete metric space. If \(E\) is a Polish space, there is a
continuous mapping of \(\mathbb N^{\mathbb N}\) onto \(E\). Indeed,
for every \(\varepsilon>0\), there is a countable covering of \(E\) by
nonempty closed sets of diameter less than \(\varepsilon\), since
\(E\) has a countable base. Hence there is a mapping
\[
 (n_0,\ldots,n_p)\longmapsto F(n_0,\ldots,n_p)
\]
from the set of finite sequences of nonnegative integers into the set
of nonempty closed subsets of \(E\), having the following properties:
\begin{enumerate}[label=\arabic*\textsuperscript{\(\circ\)},leftmargin=2em]
\item \(F(\varnothing)=E\);
\item for every sequence \((n_0,\ldots,n_p)\),
\[
 F(n_0,\ldots,n_p)
 =F(n_0,\ldots,n_p,0)\cup F(n_0,\ldots,n_p,1)\cup\cdots;
\]
\item the diameter of each \(F(n_0,\ldots,n_p)\) is at most
\(2^{-p}\), for every \(p\geq0\).
\end{enumerate}
This being so, for every
\(x=(m_0,m_1,\ldots)\in\mathbb N^{\mathbb N}\), the sequence of the
\(F(m_0,\ldots,m_p)\), for \(p=1,2,\ldots\), is a Cauchy filter base
on \(E\), and therefore converges to a point of \(E\), which we denote
by \(f(x)\). Every point of \(E\), for every \(p\geq0\), belongs to a
set \(F_p=F(m_0,\ldots,m_p)\), and is therefore a limit of the filter
base \((F_0,F_1,\ldots)\); thus \(f(\mathbb N^{\mathbb N})=E\). If
\(x=(m_0,m_1,\ldots)\) and \(y=(n_0,n_1,\ldots)\) are two elements of
\(\mathbb N^{\mathbb N}\) such that
\(d(x,y)<2^{-i-1}\), then
\(m_0=n_0,m_1=n_1,\ldots,m_i=n_i\). Hence \(f(x)\) and \(f(y)\) both
belong to \(F(m_0,\ldots,m_i)\), and are therefore at distance less
than \(2^{-i}\). This proves that \(f\) is continuous.

\NumberedParagraph{3}{Souslin Sets}
\label{no:App-V-3}\label{sec:app-V-3}
Let \(E\) be a Polish space. A subset of \(E\) is said to be
\emph{Souslin} if it is the image of a Polish space \(P\) under a
continuous mapping from \(P\) into \(E\).

The image of a Souslin subset of \(E\) under a continuous mapping from
\(E\) into another Polish space is Souslin.

Let \((A_n)\) be a sequence of Souslin subsets of \(E\). Then the union
\(A\) of the \(A_n\) is Souslin. Indeed, let \(P_n\) be a Polish space
and \(f_n\) a continuous mapping from \(P_n\) into \(E\) such that
\(f_n(P_n)=A_n\); let \(P\) be the topological sum of the \(P_n\),
which is Polish. The mapping \(f\) from \(P\) into \(E\), equal to
\(f_n\) on \(P_n\), is continuous, and \(f(P)=A\); hence \(A\) is
Souslin. Similarly, the intersection \(A'\) of the \(A_n\) is Souslin.
Indeed, let \(P'\) be the product of the \(P_n\), which is Polish, and
let \(Q\) be the subspace of \(P'\) consisting of the \((x_n)\) such
that
\[
 f_m(x_m)=f_n(x_n)
\]
for every pair of indices. Then \(Q\) is a closed subspace of \(P'\),
and is therefore Polish. If \(p_n\) denotes the canonical mapping of
\(P'\) onto \(P_n\), the restriction \(f'\) of \(f_n\circ p_n\) to
\(Q\) is independent of \(n\), is continuous, and satisfies
\(f'(Q)=A'\); hence \(A'\) is Souslin.

Thus the collection of subsets of \(E\) which are Souslin and whose
complements are Souslin is a sigma-algebra which, by
\hyperref[sec:app-V-2]{no.~2}, contains the closed subsets of \(E\).
Consequently every Borel subset of \(E\) is Souslin.

% Source PDF page 340 (printed p. 331)
\phantomsection\label{srcpage:331}
Let \(A\) be a relatively compact Souslin subset of \(E\). We shall
show that there exist a compact space \(F\), a \(K_{\sigma\delta}\) set
\(B\) in \(F\), and a continuous mapping \(f\) from \(F\) into \(E\)
such that \(A=f(B)\). Replacing \(E\) by \(\overline A\), we may suppose
that \(E\) is compact. Let \(g\) be a continuous mapping of
\(\mathbb N^{\mathbb N}\) into \(E\) such that
\(g(\mathbb N^{\mathbb N})=A\). Let \(\overline{\mathbb R}\) be the
real line completed by one point at infinity. The space
\(\mathbb N^{\mathbb N}\) is a subspace of the compact space
\(\overline{\mathbb R}^{\mathbb N}\), and it is easy to see that
\(\mathbb N^{\mathbb N}\) is the intersection of a sequence
\(G_1,G_2,\ldots\) of open subsets of
\(\overline{\mathbb R}^{\mathbb N}\). Put
\[
 F=\overline{\mathbb R}^{\mathbb N}\times E,
\]
which is compact. Let \(B\) be the graph of \(g\), let
\(\overline B\) be its closure in \(F\), and let \(f\) be the canonical
projection of \(F\) onto \(E\). We have \(A=f(B)\). On the other hand,
since \(g\) is continuous, \(B\) is closed in
\(\mathbb N^{\mathbb N}\times E\), and therefore
\[
 B=\overline B\cap(\mathbb N^{\mathbb N}\times E).
\]
Now \(\overline B\) is compact, and
\[
 \mathbb N^{\mathbb N}\times E
 =(G_1\times E)\cap(G_2\times E)\cap\cdots.
\]
Finally, each \(G_n\) is a countable union of closed, hence compact,
subsets of \(\overline{\mathbb R}^{\mathbb N}\); therefore
\(G_n\times E\) is a countable union of compact subsets of \(F\).

\NumberedParagraph{4}{Measurability of Souslin Sets}
\label{no:App-V-4}\label{sec:app-V-4}
Let \(E\) be a locally compact space with a countable base, let
\(\nu\) be a positive measure on \(E\), and let \(A\) be a Souslin
subset of \(E\). We shall show that \(A\) is \(\nu\)-measurable.

It is enough to consider the case where \(A\) is relatively compact.
There then exist a compact space \(F\), a continuous mapping \(f\) from
\(F\) into \(E\), and a family \((B_{n,p})\) of compact subsets of
\(F\) such that, on putting
\[
 B_n=B_{n,1}\cup B_{n,2}\cup\cdots,
 \qquad
 B=B_1\cap B_2\cap\cdots,
\]
we have \(A=f(B)\). We may moreover suppose that
\(B_{n,1}\subset B_{n,2}\subset\cdots\) for every \(n\). We shall
show that, for every \(a<\nu^*(A)\), there is a compact set
\(C\subset B\) such that \(\nu(f(C))\geq a\). Since \(f(C)\) is
compact, this will complete the proof.

First we prove the existence of a sequence \((p_1,p_2,\ldots)\) of
integers such that, if
\[
 C_n=B\cap B_{1,p_1}\cap\cdots\cap B_{n,p_n},
\]
then \(\nu^*(f(C_n))>a\) for every \(n\). Suppose that the
\(B_{i,p_i}\) have been chosen for \(i<n\). Since
\(C_{n-1}\subset B\subset B_n\),
\[
 C_{n-1}=(C_{n-1}\cap B_{n,1})
 \cup(C_{n-1}\cap B_{n,2})\cup\cdots,
\]
and therefore
\[
 \nu^*\bigl(f(C_{n-1}\cap B_{n,i})\bigr)
 \longrightarrow\nu^*(f(C_{n-1}))
 \qquad(i\longrightarrow+\infty).
\]
Thus
\(\nu^*(f(C_{n-1}\cap B_{n,p_n}))>a\) for a suitable choice of
\(p_n\), which proves our assertion.

Put
\[
 C=C_1\cap C_2\cap\cdots
  =B_{1,p_1}\cap B_{2,p_2}\cap\cdots.
\]
Since the \(B_{1,p_1}\cap\cdots\cap B_{n,p_n}\) form a decreasing
sequence of compact sets with intersection \(C\), the set \(f(C)\) is
compact and is equal to the intersection of the
\(f(B_{1,p_1}\cap\cdots\cap B_{n,p_n})\). Hence
\[
 \nu(f(C))
 =\lim_n\nu\bigl(f(B_{1,p_1}\cap\cdots\cap B_{n,p_n})\bigr)
 \geq\lim_n\nu^*(f(C_n))\geq a.
\]

% Source PDF page 341 (printed p. 332)
\phantomsection\label{srcpage:332}
\NumberedParagraph{5}{Existence of Measurable Mappings}
\label{no:App-V-5}\label{sec:app-V-5}
Since \(\mathbb N\) is well ordered, we may equip
\(\mathbb N^{\mathbb N}\) with the lexicographic order. Then
\(\mathbb N^{\mathbb N}\) is totally ordered. We show that every
nonempty closed subset \(P\) of \(\mathbb N^{\mathbb N}\) has a least
element. Let \(n_0\) be the least integer such that \(P\) has an
element of the form \((n_0,\ldots)\). Let \(n_1\) be the least integer
such that \(P\) has an element of the form
\((n_0,n_1,\ldots)\), and so on. Then
\[
 x_0=(n_0,n_1,n_2,\ldots)\leq x
 \qquad(x\in P).
\]
On the other hand, \(x_0\) belongs to the closure of \(P\); hence
\(x_0\in P\), since \(P\) is closed.

Let \(Q\subset\mathbb N^{\mathbb N}\) be the countable set of elements
of \(\mathbb N^{\mathbb N}\) all but finitely many of whose coordinates
are zero. Every open set \(V\) in \(\mathbb N^{\mathbb N}\) is then a
union of intervals \([z_1,z_2)\), where \(z_1,z_2\in Q\). Indeed, let
\(x=(n_0,n_1,\ldots)\in V\). There is an integer \(i\) such that every
element of \(\mathbb N^{\mathbb N}\) whose coordinates with indices
\(0,1,\ldots,i\) are equal to those of \(x\) belongs to \(V\). Put
\[
 \begin{split}
 z_1&=(n_0,n_1,\ldots,n_i,n_{i+1},0,0,\ldots),\\
 z_2&=(n_0,n_1,\ldots,n_i,n_{i+1}+1,0,0,\ldots).
 \end{split}
\]
It is immediate that \(x\in[z_1,z_2)\subset V\), which proves our
assertion.

This being so, we shall establish the following result.

\phantomsection\label{thm:App-V-selection}
\begingroup
\itshape\noindent
Let \(Z\) be a locally compact space with a countable base, let \(\nu\)
be a positive measure on \(Z\), let \(Z'\) be a subset of \(Z\), and
let \(G\) be a Polish space. For every \(\zeta\in Z\), let
\(G_\zeta\) be a subset of \(G\), nonempty for \(\zeta\in Z'\).
Suppose that the set \(E\) of pairs \((\zeta,x)\in Z\times G\) such
that \(x\in G_\zeta\) is Souslin in \(Z\times G\). Then there exist a
\(\nu\)-measurable subset \(Z''\) of \(Z\) containing \(Z'\), and a
\(\nu\)-measurable mapping \(\eta\) from \(Z''\) into \(G\), such that
\(\eta(\zeta)\in G_\zeta\) for every \(\zeta\in Z''\).
\par
\endgroup

We must construct a subset \(Z''\) of \(Z\), measurable with respect
to \(\nu\) and containing \(Z'\), and a mapping \(\eta_1\) from
\(Z''\) into \(Z\times G\), measurable with respect to \(\nu\) and
with values in \(E\), such that
\[
 \pi\eta_1(\zeta)=\zeta\qquad(\zeta\in Z''),
\]
where \(\pi\) denotes the canonical mapping of \(Z\times G\) onto
\(Z\). Regarding \(E\) as a continuous image of
\(\mathbb N^{\mathbb N}\), we are reduced to the following situation:
there is a continuous mapping \(f\) from \(\mathbb N^{\mathbb N}\)
into \(Z\) such that
\[
 A=f(\mathbb N^{\mathbb N})\supset Z',
\]
and we shall prove the existence of a \(\nu\)-measurable mapping \(g\)
from \(A\) into \(\mathbb N^{\mathbb N}\) such that \(f\circ g\) is
the identity mapping of \(A\). This will complete the proof. For every
\(x\in A\), the set \(f^{-1}(\{x\})\) is a closed subset of
\(\mathbb N^{\mathbb N}\). Let \(g(x)\) be its least element, so that
\(f\circ g\) is indeed the identity mapping of \(A\). We show that
\(g\) is \(\nu\)-measurable. By what was proved above, it is enough to
show that, for every \(z\in\mathbb N^{\mathbb N}\),
\[
 g^{-1}\bigl(\{y\in\mathbb N^{\mathbb N}:y<z\}\bigr)
\]
is \(\nu\)-measurable. But this set is plainly equal to
\[
 f\bigl(\{y\in\mathbb N^{\mathbb N}:y<z\}\bigr),
\]
and is therefore Souslin and hence \(\nu\)-measurable.

\noindent\textbf{Bibliography:} \ArticleRef{80}; and N.~Bourbaki,
\emph{Topologie générale}, Chapter~IX, 2nd ed., \emph{Actualités
scientifiques et industrielles}, no.~1045, Paris, Hermann, 1958.

% Source PDF page 342 (printed p. 333)
\chapter*{Remarks on Terminology}\label{srcpage:333}
\addcontentsline{toc}{chapter}{Remarks on Terminology}
\label{back:terminology}

The general notation is customary; it most often conforms to that of
N.~Bourbaki. Certain points are specified below in fairly great
detail. The reader must not, however, imagine that these explanations
are sufficient to enable him to begin reading this book if he does not
possess the knowledge indicated in the Introduction.

\medskip
\noindent\textbf{Set theory.} --- A finite set is regarded as
countable.

Let \(A\) and \(B\) be two sets, and let \(f\) be a mapping from \(A\)
into \(B\). The mapping \(f\) is said to be \emph{injective} if
\(x\neq y\) implies \(f(x)\neq f(y)\). It is said to be
\emph{surjective} if the image \(f(A)\) of \(A\) under \(f\) is equal
to \(B\). It is said to be \emph{bijective} if it is both injective
and surjective. The mapping \(f\) is also denoted by
\(x\mapsto f(x)\).

Let \(A\) be a set, and let \(f\) be a mapping from \(A\) into \(A\).
A subset \(A'\) of \(A\) is said to be \emph{stable under \(f\)}
if \(f(A')\subset A'\).

Let \(A,B,C\) be three sets, let \(f\) be a mapping from \(A\) into
\(B\), and let \(g\) be a mapping from \(B\) into \(C\). The composite
mapping from \(A\) into \(C\) is denoted by \(g\circ f\), or by \(gf\).

The set of nonnegative integers is denoted by \(\mathbb N\), the set of
real numbers by \(\mathbb R\), and the set of complex numbers by
\(\mathbb C\).

Let \(A\) and \(B\) be two sets, and let \(f\) be a mapping from \(A\)
into \(B\). If \(B=\mathbb R\) (respectively, \(\mathbb C\)), then
\(f\) is called a \emph{numerical function} (respectively, a
\emph{complex numerical function}). If \(A=\mathbb R\), then \(f\) is
called a \emph{function of a real variable}.

\medskip
\noindent\textbf{Topology.} --- A topological space \(Z\) is said to
have a \emph{countable base} if there is a sequence
\((V_1,V_2,\ldots)\) of open subsets of \(Z\) such that every open
subset of \(Z\) is a union of some of the \(V_i\). When \(Z\) is
metrizable, this is equivalent to saying that \(Z\) has a countable
dense subset. When \(Z\) is a Hilbert space, this is equivalent to
saying that it has a countable orthonormal basis of vectors, so that
the terminology cannot give rise to any confusion.
% Source PDF page 343 (printed p. 334)
\phantomsection\label{srcpage:334}
\noindent\textbf{Algebra.} --- A vector space \(E\) over the field of complex
numbers is called a \emph{complex vector space}. Let \(E\) and \(F\) be two
complex vector spaces. A mapping \(f\) from \(E\) into \(F\) is said to be
\emph{linear} if
\[
 f(\lambda x+\mu y)=\lambda f(x)+\mu f(y)
\]
for all \(x,y\in E\), \(\lambda,\mu\in\mathbb C\). The \emph{rank} of \(f\)
is the dimension of the vector subspace \(f(E)\) of \(F\).

A linear mapping from \(E\) into \(E\) is also called a \emph{linear
operator} in \(E\), or simply an \emph{operator} in \(E\). An operator of the
form \(x\mapsto\lambda x\) (where \(\lambda\) is a fixed element of
\(\mathbb C\)) is called a \emph{homothety}, or a \emph{scalar operator}; it
is denoted by \(\lambda\) when no confusion results. The operator \(1\) is
called the \emph{identity operator}.

Let \(E^*\) be the algebraic dual of \(E\), and let \(\langle x',x\rangle\)
denote the value at a point \(x\) of \(E\) of the linear form
\(x'\in E^*\). Let \(A'\) be a vector subspace of \(E^*\). The restriction
to \(E\times A'\) of the mapping
\((x,x')\mapsto\langle x,x'\rangle\) is called the \emph{canonical bilinear
form} on \(E\times A'\).

An \emph{algebra over \(\mathbb C\)} is a complex vector space \(A\), endowed
with a bilinear mapping \((x,y)\mapsto xy\) from \(A\times A\) into \(A\),
such that \((xy)z=x(yz)\) for all \(x,y,z\in A\) (we shall consider only
associative algebras). The algebra \(A\) is said to be \emph{Abelian} if any
two elements \(x,y\) of \(A\) commute. In general, the \emph{centre} of \(A\)
is the set of those \(x\in A\) which commute with all the elements of \(A\).
An \emph{involution} in \(A\) is a bijective mapping \(x\mapsto x^*\) of
\(A\) onto \(A\) such that
\[
 x^{**}=x,\qquad (x+y)^*=x^*+y^*,
\]
\[
 (\lambda x)^*=\overline\lambda x^*,\qquad (xy)^*=y^*x^*
 \quad\text{for }x,y\in A,\ \lambda\in\mathbb C.
\]
An algebra over \(\mathbb C\) endowed with an involution is called an
\emph{involutive algebra}.

\medskip
\noindent\textbf{Topological vector spaces.} --- Let \(E\) be a complex
vector space. A \emph{seminorm} on \(E\) is a function \(p\) defined on \(E\),
with values \(\geq0\), such that
\(p(\lambda x)=|\lambda|p(x)\) and
\(p(x+y)\leq p(x)+p(y)\) for all \(x,y\in E\), \(\lambda\in\mathbb C\). If,
in addition, \(p(x)=0\) implies \(x=0\), then \(p\) is called a
\emph{norm}; the notation \(\lVert x\rVert\) is generally used in this case,
and \(E\) is called a \emph{complex normed vector space}. The set of
\(x\in E\) such that \(\lVert x\rVert\leq1\) is called the \emph{unit ball}
of \(E\).

In a complex topological vector space \(E\), a set \(A\) is said to be
\emph{total} if the smallest closed vector subspace of \(E\) containing \(A\)
is \(E\) itself.

A \emph{complex pre-Hilbert space} is a complex vector space \(\mathfrak H\)
endowed with an \emph{inner product} [denoted by \((x\mid y)\) in this book],
that is, with a mapping \((x,y)\mapsto(x\mid y)\) from
\(\mathfrak H\times\mathfrak H\) into \(\mathbb C\) satisfying
\[
 (y\mid x)=\overline{(x\mid y)},\qquad
 (\lambda x+\mu y\mid z)=\lambda(x\mid z)+\mu(y\mid z),\qquad
 (x\mid x)\geq0
\]
% Source PDF page 344 (printed p. 335)
\phantomsection\label{srcpage:335}
for all \(x,y,z\in\mathfrak H\), \(\lambda,\mu\in\mathbb C\). Then the
function \(x\mapsto(x\mid x)^{1/2}\) is a seminorm on \(\mathfrak H\), and
we always endow \(\mathfrak H\) with the topology defined by this seminorm. If
this topology is Hausdorff, that is, if this seminorm is a norm, we have a
\emph{complex Hausdorff pre-Hilbert space}. If, in addition,
\(\mathfrak H\) is complete, \(\mathfrak H\) is called a \emph{complex
Hilbert space}.

Let \(\mathfrak H\) be a complex Hilbert space and \(T\) a continuous operator
in \(\mathfrak H\). By \(\lVert T\rVert\) we denote the supremum of
\(\lVert Tx\rVert\) for \(\lVert x\rVert\leq1\). By \(T^*\) we denote the
unique continuous operator in \(\mathfrak H\) such that
\((T^*x\mid y)=(x\mid Ty)\) for all \(x,y\in\mathfrak H\). If \(T=T^*\),
\(T\) is said to be \emph{Hermitian}. If \(TT^*=T^*T=1\), \(T\) is said to be
\emph{unitary}. If \(T\) maps the unit ball of \(\mathfrak H\) into a strongly
compact subset of \(\mathfrak H\), \(T\) is said to be \emph{compact}. (We
recall that the terminology used by many authors is ``completely
continuous''.) An \emph{involution} in \(\mathfrak H\) is a bijective mapping
\(J\) of \(\mathfrak H\) onto \(\mathfrak H\) such that
\[
 J^2x=x,\qquad J(x+y)=Jx+Jy,\qquad
 J(\lambda x)=\overline\lambda Jx,\qquad
 (Jx\mid Jy)=(y\mid x)
\]
for all \(x,y\in\mathfrak H\), \(\lambda\in\mathbb C\).

\begin{center}
\textsc{End.}
\end{center}

% Source PDF page 345 is blank; intentionally omitted (no srcpage:336 anchor).

% Source PDF page 346 (printed p. 337)
\chapter*{Bibliography}\label{back:bibliography}\label{srcpage:337}
\addcontentsline{toc}{chapter}{Bibliography}

\section*{Journal Articles}\label{bib:journal-articles}
\addcontentsline{toc}{section}{Journal Articles}

The first part of this bibliography extends approximately through 1954. The
articles are arranged there in alphabetical order by author. The second part
extends approximately from 1954 to 1968. The articles are arranged there by
year.

For some time now, it has become increasingly artificial to separate works
devoted to von Neumann algebras from works devoted to \(C^*\)-algebras. The
division made between the present book and the second edition of my book on
\(C^*\)-algebras (\emph{Cahiers scientifiques}, fascicle XXIX) is therefore
rather arbitrary.

Some of the articles in theoretical physics related to von Neumann algebras
have been indicated, but no attempt has been made to give a complete list.

\subsection*{First Part}\label{bib:journal-first-part}

\let\DixmierBibliographyChapter\chapter
\RenewDocumentCommand{\chapter}{s o m}{}
\begin{thebibliography}{999}
\let\chapter\DixmierBibliographyChapter

\bibitem[1]{article:1}\phantomsection\label{bib:article:1}
\textsc{W. Ambrose}, \emph{Structure theorems for a special class of Banach algebras}
(Trans. Amer. Math. Soc., t. 57, 1945, p. 364--386).

\bibitem[2]{article:2}\phantomsection\label{bib:article:2}
\textsc{W. Ambrose}, \emph{The \(L^2\)-system of a unimodular group I}
(Trans. Amer. Math. Soc., t. 65, 1949, p. 27--48).

\bibitem[3]{article:3}\phantomsection\label{bib:article:3}
\textsc{A. Brown}, \emph{The unitary equivalence of binormal operators}
(Amer. J. Math., t. 76, 1954, p. 414--434).

\bibitem[4]{article:4}\phantomsection\label{bib:article:4}
\textsc{J. W. Calkin}, \emph{Two-sided ideals and congruences in the ring of bounded operators in Hilbert space}
(Ann. Math., t. 42, 1941, p. 839--873).

\bibitem[5]{article:5}\phantomsection\label{bib:article:5}
\textsc{J. Dixmier}, \emph{Mesure de Haar et trace d'un opérateur}
(C. R. Acad. Sc., t. 228, 1949, p. 152--154).

\bibitem[6]{article:6}\phantomsection\label{bib:article:6}
\textsc{J. Dixmier}, \emph{Les anneaux d'opérateurs de classe finie}
(Ann. Éc. Norm. Sup., t. 66, 1949, p. 209--261).

\bibitem[7]{article:7}\phantomsection\label{bib:article:7}
\textsc{J. Dixmier}, \emph{Les fonctionnelles linéaires sur l'ensemble des opérateurs bornés d'un espace de Hilbert}
(Ann. Math., t. 51, 1950, p. 387--408).

\bibitem[8]{article:8}\phantomsection\label{bib:article:8}
\textsc{J. Dixmier}, \emph{Application \(\natural\) dans les anneaux d'opérateurs}
(C. R. Acad. Sc., t. 230, 1950, p. 607--608).

\bibitem[9]{article:9}\phantomsection\label{bib:article:9}
\textsc{J. Dixmier}, \emph{Sur certains espaces considérés par M. H. Stone}
(Summa Brasil. Math., t. 2, fasc. 11, 1951, p. 151--182).

\bibitem[10]{article:10}\phantomsection\label{bib:article:10}
\textsc{J. Dixmier}, \emph{Sur la réduction des anneaux d'opérateurs}
(Ann. Éc. Norm. Sup., t. 68, 1951, p. 185--202).

\bibitem[11]{article:11}\phantomsection\label{bib:article:11}
\textsc{J. Dixmier}, \emph{Algèbres quasi-unitaires}
(C. R. Acad. Sc., t. 233, 1951, p. 837--839).

\bibitem[12]{article:12}\phantomsection\label{bib:article:12}
\textsc{J. Dixmier}, \emph{Applications \(\natural\) dans les anneaux d'opérateurs}
(Compos. Math., t. 10, 1952, p. 1--55).

\bibitem[13]{article:13}\phantomsection\label{bib:article:13}
\textsc{J. Dixmier}, \emph{Algèbres quasi-unitaires}
(Comment. Math. Helv., t. 26, 1952, p. 275--322).

% Source PDF page 347 (printed p. 338)
\phantomsection\label{srcpage:338}
\bibitem[14]{article:14}\phantomsection\label{bib:article:14}
\textsc{J. Dixmier}, \emph{Remarques sur les applications \(\natural\)}
(Archiv Math., t. 3, 1952, p. 290--297).

\bibitem[15]{article:15}\phantomsection\label{bib:article:15}
\textsc{J. Dixmier}, \emph{Formes linéaires sur un anneau d'opérateurs}
(Bull. Soc. Math. Fr., t. 81, 1953, p. 9--39).

\bibitem[16]{article:16}\phantomsection\label{bib:article:16}
\textsc{J. Dixmier}, \emph{Sous-anneaux abéliens maximaux dans les facteurs de type fini}
(Ann. Math., t. 59, 1954, p. 279--286).

\bibitem[17]{article:17}\phantomsection\label{bib:article:17}
\textsc{J. Dixmier}, \emph{Sur les anneaux d'opérateurs dans les espaces hilbertiens}
(C. R. Acad. Sc., t. 238, 1954, p. 439--441).

\bibitem[18]{article:18}\phantomsection\label{bib:article:18}
\textsc{N. Dunford}, \emph{Resolution of the identity for commutative B*-algebras of operators}
(Acta Univ. Szeged, t. 12, 1950, p. 51--56).

\bibitem[19]{article:19}\phantomsection\label{bib:article:19}
\textsc{H. A. Dye}, \emph{The Radon-Nikodym theorem for finite rings of operators}
(Trans. Amer. Math. Soc., t. 72, 1952, p. 243--280).

\bibitem[20]{article:20}\phantomsection\label{bib:article:20}
\textsc{H. A. Dye}, \emph{The unitary structure in finite rings of operators}
(Duke Math. J., t. 20, 1953, p. 55--69).

\bibitem[21]{article:21}\phantomsection\label{bib:article:21}
\textsc{H. A. Dye}, \emph{On the geometry of projections in certain operator algebras}
(Ann. Math., t. 61, 1955, p. 73--89).

\bibitem[22]{article:22}\phantomsection\label{bib:article:22}
\textsc{B. Fuglede et R. V. Kadison}, \emph{On a conjecture of Murray and von Neumann}
(Proc. Nat. Acad. Sc. U. S. A., t. 37, 1951, p. 420--425).

\bibitem[23]{article:23}\phantomsection\label{bib:article:23}
\textsc{B. Fuglede et R. V. Kadison}, \emph{On determinants and a property of the trace in finite factors}
(Proc. Nat. Acad. Sc. U. S. A., t. 37, 1951, p. 425--431).

\bibitem[24]{article:24}\phantomsection\label{bib:article:24}
\textsc{B. Fuglede et R. V. Kadison}, \emph{Determinant theory in finite factors}
(Ann. Math., t. 55, 1952, p. 520--530).

\bibitem[25]{article:25}\phantomsection\label{bib:article:25}
\textsc{I. Gelfand et M. Neumark}, \emph{On the imbedding of normed rings into the ring of operators in Hilbert space}
(Mat. Sbornik, t. 12, 1943, p. 197--213).

\bibitem[26]{article:26}\phantomsection\label{bib:article:26}
\textsc{R. Godement}, \emph{Théorie générale des sommes continues d'espaces de Banach}
(C. R. Acad. Sc., t. 228, 1949, p. 1321--1323).

\bibitem[27]{article:27}\phantomsection\label{bib:article:27}
\textsc{R. Godement}, \emph{Sur la théorie des caractères. I. Définition et classification des caractères}
(C. R. Acad. Sc., t. 229, 1949, p. 967--969).

\bibitem[28]{article:28}\phantomsection\label{bib:article:28}
\textsc{R. Godement}, \emph{Sur la théorie des représentations unitaires}
(Ann. Math., t. 53, 1951, p. 68--124).

\bibitem[29]{article:29}\phantomsection\label{bib:article:29}
\textsc{R. Godement}, \emph{Mémoire sur la théorie des caractères dans les groupes localement compacts unimodulaires}
(J. Math. pures et appl., t. 30, 1951, p. 1--110).

\bibitem[30]{article:30}\phantomsection\label{bib:article:30}
\textsc{R. Godement}, \emph{Théorie des caractères. I. Algèbres unitaires}
(Ann. Math., t. 59, 1954, p. 47--62).

\bibitem[31]{article:31}\phantomsection\label{bib:article:31}
\textsc{E. L. Griffin}, \emph{Some contributions to the theory of rings of operators}
(Trans. Amer. Math. Soc., t. 75, 1953, p. 471--504).

\bibitem[32]{article:32}\phantomsection\label{bib:article:32}
\textsc{A. Grothendieck}, \emph{Réarrangements de fonctions et inégalités de convexité dans les algèbres de von Neumann munies d'une trace}, Séminaire Bourbaki, mars 1955, 16 pages.

\bibitem[33]{article:33}\phantomsection\label{bib:article:33}
\textsc{I. Halperin}, \emph{A remark on a preceding paper by J. von Neumann}
(Ann. Math., t. 41, 1940, p. 554--555).

\bibitem[34]{article:34}\phantomsection\label{bib:article:34}
\textsc{R. V. Kadison}, \emph{Isometries of operator algebras}
(Ann. Math., t. 54, 1951, p. 325--338).

\bibitem[35]{article:35}\phantomsection\label{bib:article:35}
\textsc{R. V. Kadison}, \emph{Order properties of bounded self-adjoint operators}
(Proc. Amer. Math. Soc., t. 2, 1951, p. 505--510).

\bibitem[36]{article:36}\phantomsection\label{bib:article:36}
\textsc{R. V. Kadison}, \emph{A generalised Schwarz inequality and algebraic invariants for operator algebras}, t. 56, 1952, p. 494--503.

\bibitem[37]{article:37}\phantomsection\label{bib:article:37}
\textsc{R. V. Kadison}, \emph{Infinite unitary groups}
(Trans. Amer. Math. Soc., t. 72, 1952, p. 386--399).

\bibitem[38]{article:38}\phantomsection\label{bib:article:38}
\textsc{R. V. Kadison}, \emph{Infinite general linear groups}
(Trans. Amer. Math. Soc., t. 76, 1954, p. 66--91).

\bibitem[39]{article:39}\phantomsection\label{bib:article:39}
\textsc{R. V. Kadison}, \emph{The general linear group of infinite factors}
(Duke Math. J., t. 22, 1955, p. 119--122).

% Source PDF page 348 (printed p. 339)
\phantomsection\label{srcpage:339}
\bibitem[40]{article:40}\phantomsection\label{bib:article:40}
\textsc{I. Kaplansky}, \emph{Quelques résultats sur les anneaux d'opérateurs}
(C. R. Acad. Sc., t. 231, 1950, p. 485--486).

\bibitem[41]{article:41}\phantomsection\label{bib:article:41}
\textsc{I. Kaplansky}, \emph{A theorem on rings of operators}
(Pacific J. Math., t. 1, 1951, p. 227--232).

\bibitem[42]{article:42}\phantomsection\label{bib:article:42}
\textsc{I. Kaplansky}, \emph{Projections in Banach algebras}
(Ann. Math., t. 53, 1951, p. 235--249).

\bibitem[43]{article:43}\phantomsection\label{bib:article:43}
\textsc{I. Kaplansky}, \emph{Algebras of type I}
(Ann. Math., t. 56, 1952, p. 460--472).

\bibitem[44]{article:44}\phantomsection\label{bib:article:44}
\textsc{I. Kaplansky}, \emph{Modules over operator algebras}
(Amer. J. Math., t. 75, 1953, p. 839--853).

\bibitem[45]{article:45}\phantomsection\label{bib:article:45}
\textsc{J. L. Kelley}, \emph{Commutative operator algebras}
(Proc. Nat. Acad. Sc. U. S. A., t. 38, 1952, p. 598--605).

\bibitem[46]{article:46}\phantomsection\label{bib:article:46}
\textsc{M. Kondô}, \emph{Sur la notion de dimension}
(Proc. Imp. Acad. Tokyo, t. 19, 1943, p. 215--223).

\bibitem[47]{article:47}\phantomsection\label{bib:article:47}
\textsc{M. Kondô}, \emph{Les anneaux des opérateurs et les dimensions}
(Proc. Imp. Acad. Tokyo, t. 20, 1944, p. 389--398).

\bibitem[48]{article:48}\phantomsection\label{bib:article:48}
\textsc{M. Kondô}, \emph{Sur les sommes directes des espaces linéaires}
(Proc. Imp. Acad. Tokyo, t. 20, 1944, p. 425--431).

\bibitem[49]{article:49}\phantomsection\label{bib:article:49}
\textsc{M. Kondô}, \emph{Sur la réductibilité des anneaux des opérateurs}
(Proc. Imp. Acad. Tokyo, t. 20, 1944, p. 432--438).

\bibitem[50]{article:50}\phantomsection\label{bib:article:50}
\textsc{M. Kondô}, \emph{Le produit kroneckerien infini des espaces linéaires}
(Proc. Imp. Acad. Tokyo, t. 20, 1944, p. 569--579).

\bibitem[51]{article:51}\phantomsection\label{bib:article:51}
\textsc{M. Kondô}, \emph{Les anneaux des opérateurs et les dimensions. II}
(Proc. Imp. Acad. Tokyo, t. 20, 1944, p. 689--693).

\bibitem[52]{article:52}\phantomsection\label{bib:article:52}
\textsc{G. W. Mackey}, \emph{Induced representations of locally compact groups. I}
(Ann. Math., t. 55, 1952, p. 101--139).

\bibitem[53]{article:53}\phantomsection\label{bib:article:53}
\textsc{G. W. Mackey}, \emph{Induced representations of locally compact groups. II}
(Ann. Math., t. 58, 1953, p. 193--221).

\bibitem[54]{article:54}\phantomsection\label{bib:article:54}
\textsc{F. Maeda}, \emph{Relative dimensionality in operator rings}
(J. Science of Hiroshima University, t. 11, 1941, p. 1--6).

\bibitem[55]{article:55}\phantomsection\label{bib:article:55}
\textsc{F. I. Mautner}, \emph{The completeness of the irreducible unitary representations of a locally compact group}
(Proc. Nat. Acad. Sc. U. S. A., t. 34, 1948, p. 52--54).

\bibitem[56]{article:56}\phantomsection\label{bib:article:56}
\textsc{F. I. Mautner}, \emph{The structure of the regular representation of certain discrete groups}
(Duke Math. J., t. 17, 1950, p. 437--441).

\bibitem[57]{article:57}\phantomsection\label{bib:article:57}
\textsc{F. I. Mautner}, \emph{Unitary representations of locally compact groups. I}
(Ann. Math., t. 51, 1950, p. 1--25).

\bibitem[58]{article:58}\phantomsection\label{bib:article:58}
\textsc{F. I. Mautner}, \emph{Unitary representations of locally compact groups. II}
(Ann. Math., t. 52, 1950, p. 528--556).

\bibitem[59]{article:59}\phantomsection\label{bib:article:59}
\textsc{F. I. Mautner}, \emph{Induced representations}
(Amer. J. Math., t. 74, 1952, p. 737--758).

\bibitem[60]{article:60}\phantomsection\label{bib:article:60}
\textsc{Y. Misonou}, \emph{On a weakly central operator algebra}
(Tôhoku Math. J., t. 4, 1952, p. 194--202).

\bibitem[61]{article:61}\phantomsection\label{bib:article:61}
\textsc{Y. Misonou}, \emph{Operator algebras of type I}
(Kodai Math. Sem. Rep., 1953, p. 87--90).

\bibitem[62]{article:62}\phantomsection\label{bib:article:62}
\textsc{Y. Misonou}, \emph{Unitary equivalence of factors of type III}
(Proc. Jap. Acad., t. 29, 1953, p. 482--485).

\bibitem[63]{article:63}\phantomsection\label{bib:article:63}
\textsc{Y. Misonou et M. Nakamura}, \emph{Centering of an operator algebra}
(Tôhoku Math. J., t. 3, 1951, p. 243--248).

\bibitem[64]{article:64}\phantomsection\label{bib:article:64}
\textsc{F. J. Murray}, \emph{Bilinear transformations in Hilbert space}
(Trans. Amer. Math. Soc., t. 45, 1939, p. 474--507).

\bibitem[65]{article:65}\phantomsection\label{bib:article:65}
\textsc{F. J. Murray et J. von Neumann}, \emph{On rings of operators}
(Ann. Math., t. 37, 1936, p. 116--229).

\bibitem[66]{article:66}\phantomsection\label{bib:article:66}
\textsc{F. J. Murray et J. von Neumann}, \emph{On rings of operators II}
(Trans. Amer. Math. Soc., t. 41, 1937, p. 208--248).

\bibitem[67]{article:67}\phantomsection\label{bib:article:67}
\textsc{F. J. Murray et J. von Neumann}, \emph{On rings of operators IV}
(Ann. Math., t. 44, 1943, p. 716--808).

% Source PDF page 349 (printed p. 340)
\phantomsection\label{srcpage:340}
\bibitem[68]{article:68}\phantomsection\label{bib:article:68}
\textsc{M. A. Naimark}, \emph{Anneaux d'opérateurs dans les espaces hilbertiens}
(Uspekhi Mat. Nauk, t. 4, 1949, p. 83--147) (en russe).

\bibitem[69]{article:69}\phantomsection\label{bib:article:69}
\textsc{M. Nakamura}, \emph{The two-sided representations of an operator algebra}
(Proc. Jap. Acad., t. 27, 1951, p. 172--176).

\bibitem[70]{article:70}\phantomsection\label{bib:article:70}
\textsc{M. Nakamura et Z. Takeda}, \emph{The Radon-Nikodym theorem of traces for a certain operator algebra}
(Tôhoku Math. J., t. 4, 1952, p. 275--283).

\bibitem[71]{article:71}\phantomsection\label{bib:article:71}
\textsc{M. Nakamura et Z. Takeda}, \emph{Normal states of commutative operator algebras}
(Tôhoku Math. J., t. 5, 1953, p. 109--121).

\bibitem[72]{article:72}\phantomsection\label{bib:article:72}
\textsc{H. Nakano}, \emph{Hilbert algebras}
(Tôhoku Math. J., t. 2, 1950, p. 4--23).

\bibitem[73]{article:73}\phantomsection\label{bib:article:73}
\textsc{J. von Neumann}, \emph{Zur Algebra der Funktionaloperationen und Theorie der normalen Operatoren}
(Math. Ann., t. 102, 1929, p. 370--427).

\bibitem[74]{article:74}\phantomsection\label{bib:article:74}
\textsc{J. von Neumann}, \emph{On a certain topology for rings of operators}
(Ann. Math., t. 37, 1936, p. 111--115).

\bibitem[75]{article:75}\phantomsection\label{bib:article:75}
\textsc{J. von Neumann}, \emph{On an algebraic generalization of the quantum mechanical formalism I}
(Mat. Sbornik, t. 1, 1936, p. 415--484).

\bibitem[76]{article:76}\phantomsection\label{bib:article:76}
\textsc{J. von Neumann}, \emph{Some matrix-inequalities and metrization of matrix-space}
(Tomsk Univ. Rev., t. 1, 1937, p. 286--300).

\bibitem[77]{article:77}\phantomsection\label{bib:article:77}
\textsc{J. von Neumann}, \emph{On infinite direct products}
(Compos. Math., t. 6, 1938, p. 1--77).

\bibitem[78]{article:78}\phantomsection\label{bib:article:78}
\textsc{J. von Neumann}, \emph{On rings of operators III}
(Ann. Math., t. 41, 1940, p. 94--161).

\bibitem[79]{article:79}\phantomsection\label{bib:article:79}
\textsc{J. von Neumann}, \emph{On some algebraical properties of operator rings}
(Ann. Math., t. 44, 1943, p. 709--715).

\bibitem[80]{article:80}\phantomsection\label{bib:article:80}
\textsc{J. von Neumann}, \emph{On rings of operators. Reduction theory}
(Ann. Math., t. 50, 1949, p. 401--485).

\bibitem[81]{article:81}\phantomsection\label{bib:article:81}
\textsc{M. Orihara}, \emph{Sur les anneaux des opérateurs I}
(Proc. Imp. Acad. Tokyo, t. 20, 1944, p. 399--405).

\bibitem[82]{article:82}\phantomsection\label{bib:article:82}
\textsc{M. Orihara}, \emph{Sur les anneaux des opérateurs II}
(Proc. Imp. Acad. Tokyo, t. 20, 1944, p. 545--553).

\bibitem[83]{article:83}\phantomsection\label{bib:article:83}
\textsc{M. Orihara}, \emph{Rings of operators and their traces}
(Memoirs of the Faculty of Science Kyūsyū Univ., série A, t. V, 1950, p. 107--138).

\bibitem[84]{article:84}\phantomsection\label{bib:article:84}
\textsc{M. Orihara}, \emph{Correction to my paper ``Rings of operators and their traces''}
(Memoirs of the Faculty of Science Kyūsyū Univ., série A, t. VIII, 1953, p. 89--91).

\bibitem[85]{article:85}\phantomsection\label{bib:article:85}
\textsc{R. Pallu de la Barrière}, \emph{Algèbres autoadjointes faiblement fermées et algèbres hilbertiennes de classe finie}
(C. R. Acad. Sc., t. 232, 1951, p. 1994--1995).

\bibitem[86]{article:86}\phantomsection\label{bib:article:86}
\textsc{R. Pallu de la Barrière}, \emph{Décomposition des opérateurs non bornés dans les sommes continues d'espaces de Hilbert}
(C. R. Acad. Sc., t. 232, 1951, p. 2071--2073).

\bibitem[87]{article:87}\phantomsection\label{bib:article:87}
\textsc{R. Pallu de la Barrière}, \emph{Algèbres unitaires et espaces de Ambrose}
(C. R. Acad. Sc., t. 233, 1951, p. 997--999).

\bibitem[88]{article:88}\phantomsection\label{bib:article:88}
\textsc{R. Pallu de la Barrière}, \emph{Isomorphismes des *-algèbres faiblement fermées d'opérateurs}
(C. R. Acad. Sc., t. 232, 1951, p. 2071--2073).

\bibitem[89]{article:89}\phantomsection\label{bib:article:89}
\textsc{R. Pallu de la Barrière}, \emph{Sur les algèbres d'opérateurs dans les espaces hilbertiens}
(Bull. Soc. Math. Fr., t. 82, 1954, p. 1--51).

\bibitem[90]{article:90}\phantomsection\label{bib:article:90}
\textsc{R. Pallu de la Barrière}, \emph{Algèbres unitaires et espaces d'Ambrose}
(Ann. Éc. Norm. Sup., t. 70, 1953, p. 381--401).

\bibitem[91]{article:91}\phantomsection\label{bib:article:91}
\textsc{L. Pukánszky}, \emph{On a theorem of Mautner}
(Acta Sc. Math. Szeged, t. 15, 1954, p. 145--148).

\bibitem[92]{article:92}\phantomsection\label{bib:article:92}
\textsc{L. Pukánszky}, \emph{The theorem of Radon-Nikodym in operator rings}
(Acta Sc. Math. Szeged, t. 15, 1954, p. 149--156).

\bibitem[93]{article:93}\phantomsection\label{bib:article:93}
\textsc{L. Pukánszky}, \emph{On the theory of quasi-unitary algebras}
(Acta Sc. Math. Szeged, t. 16, 1955, p. 103--121).

\bibitem[94]{article:94}\phantomsection\label{bib:article:94}
\textsc{C. E. Rickart}, \emph{Banach algebras with an adjoint operation}
(Ann. Math., t. 47, 1946, p. 528--549).

% Source PDF page 350 (printed p. 341)
\phantomsection\label{srcpage:341}
\bibitem[95]{article:95}\phantomsection\label{bib:article:95}
\textsc{V. Rohlin}, \emph{Anneaux unitaires}
(C. R. Acad. Sc. U. R. S. S., t. 59, 1948, p. 643--646) (en russe).

\bibitem[96]{article:96}\phantomsection\label{bib:article:96}
\textsc{I. E. Segal}, \emph{Irreducible representations of operator algebras}
(Bull. Amer. Math. Soc., t. 53, 1947, p. 73--88).

\bibitem[97]{article:97}\phantomsection\label{bib:article:97}
\textsc{I. E. Segal}, \emph{Two-sided ideals in operator algebras}
(Ann. Math., t. 50, 1949, p. 856--865).

\bibitem[98]{article:98}\phantomsection\label{bib:article:98}
\textsc{I. E. Segal}, \emph{The two-sided regular representation of a unimodular locally compact group}
(Ann. Math., t. 51, 1950, p. 293--298).

\bibitem[99]{article:99}\phantomsection\label{bib:article:99}
\textsc{I. E. Segal}, \emph{An extension of Plancherel's formula to separable unimodular groups}
(Ann. Math., t. 52, 1950, p. 272--292).

\bibitem[100]{article:100}\phantomsection\label{bib:article:100}
\textsc{I. E. Segal}, \emph{Decomposition of operator algebras, I et II}
(Memoirs of the Amer. Math. Soc., t. 9, 1951, p. 1--67 et 1--66).

\bibitem[101]{article:101}\phantomsection\label{bib:article:101}
\textsc{I. E. Segal}, \emph{A non-commutative extension of abstract integration}
(Ann. Math., t. 57, 1953, p. 401--457).

\bibitem[102]{article:102}\phantomsection\label{bib:article:102}
\textsc{I. E. Segal}, \emph{Correction to ``A non-commutative extension of abstract integration''}
(Ann. Math., t. 57, 1953, p. 595--596).

\bibitem[103]{article:103}\phantomsection\label{bib:article:103}
\textsc{I. M. Singer}, \emph{Automorphisms of finite factors}
(Amer. J. Math., t. 77, 1955, p. 117--133).

\bibitem[104]{article:104}\phantomsection\label{bib:article:104}
\textsc{H. Sunouchi}, \emph{On rings of operators of infinite classes}
(Proc. Jap. Acad., t. 28, 1952, p. 9--13).

\bibitem[105]{article:105}\phantomsection\label{bib:article:105}
\textsc{H. Sunouchi}, \emph{On rings of operators of infinite classes II}
(Proc. Jap. Acad., t. 28, 1952, p. 330--335).

\bibitem[106]{article:106}\phantomsection\label{bib:article:106}
\textsc{H. Sunouchi}, \emph{The irreducible decompositions of the maximal Hilbert algebras of the finite class}
(Tôhoku Math. J., t. 4, 1952, p. 207--215).

\bibitem[107]{article:107}\phantomsection\label{bib:article:107}
\textsc{H. Sunouchi}, \emph{An extension of the Plancherel formula to unimodular groups}
(Tôhoku Math. J., t. 4, 1952, p. 216--230).

\bibitem[108]{article:108}\phantomsection\label{bib:article:108}
\textsc{H. Sunouchi}, \emph{A caracterization of the maximal ideal in a factor of the case \(II_\infty\)}
(Kōdai Math. Sem. Rep., 1954, p. 7).

\bibitem[109]{article:109}\phantomsection\label{bib:article:109}
\textsc{Z. Takeda}, \emph{On a theorem of R. Pallu de la Barrière}
(Proc. Jap. Acad., t. 28, 1952, p. 558--563).

\bibitem[110]{article:110}\phantomsection\label{bib:article:110}
\textsc{Z. Takeda}, \emph{Perfection of measure spaces and W*-algebras}
(Kōdai Math. Sem. Rep., 1953, p. 23--26).

\bibitem[111]{article:111}\phantomsection\label{bib:article:111}
\textsc{Z. Takeda}, \emph{Conjugate spaces of operator algebras}
(Proc. Jap. Acad., t. 30, 1954, p. 90--95).

\bibitem[112]{article:112}\phantomsection\label{bib:article:112}
\textsc{Z. Takeda}, \emph{On the representations of operator algebras}
(Proc. Jap. Acad., t. 30, 1954, p. 299--304).

\bibitem[113]{article:113}\phantomsection\label{bib:article:113}
\textsc{Z. Takeda}, \emph{On the representations of operator algebras II}
(Tôhoku Math. J., t. 6, 1954, p. 212--219).

\bibitem[114]{article:114}\phantomsection\label{bib:article:114}
\textsc{Z. Takeda et T. Turumaru}, \emph{On the property ``Position p''}
(Math. Jap., t. 2, 1952, p. 195--197).

\bibitem[115]{article:115}\phantomsection\label{bib:article:115}
\textsc{O. Takenouchi}, \emph{On the maximal Hilbert algebras}
(Tôhoku Math. J., t. 3, 1951, p. 123--131).

\bibitem[116]{article:116}\phantomsection\label{bib:article:116}
\textsc{O. Takenouchi}, \emph{On the structure of maximal Hilbert algebras}
(Math. J. Okayama Univ., t. 1, 1952, p. 1--31).

\bibitem[117]{article:117}\phantomsection\label{bib:article:117}
\textsc{M. Tomita}, \emph{On rings of operators in non separable Hilbert spaces}
(Memoirs of the Faculty of Science Kyūsyū Univ., t. 7, 1953, p. 129--168).

\bibitem[118]{article:118}\phantomsection\label{bib:article:118}
\textsc{M. Tomita}, \emph{Representations of operator algebras}
(Math. J. Okayama Univ., t. 3, 1954, p. 147--173).

\bibitem[119]{article:119}\phantomsection\label{bib:article:119}
\textsc{H. Umegaki}, \emph{On some representation theorems in an operator algebra I}
(Proc. Jap. Acad., t. 27, 1951, p. 328--333).

\bibitem[120]{article:120}\phantomsection\label{bib:article:120}
\textsc{H. Umegaki}, \emph{On some representation theorems in an operator algebra III}
(Proc. Jap. Acad., t. 28, 1952, p. 29--31).

% Source PDF page 351 (printed p. 342)
\phantomsection\label{srcpage:342}
\bibitem[121]{article:121}\phantomsection\label{bib:article:121}
\textsc{H. Umegaki}, \emph{Operator algebra of finite class}
(Kōdai Math. Sem. Rep., 1952, p. 123--129).

\bibitem[122]{article:122}\phantomsection\label{bib:article:122}
\textsc{H. Umegaki}, \emph{Operator algebra of finite class II}
(Kōdai Math. Sem. Rep., 1953, p. 61--63).

\bibitem[123]{article:123}\phantomsection\label{bib:article:123}
\textsc{H. Umegaki}, \emph{Decomposition theorems of operator algebra and their applications}
(Jap. J. Math., t. 22, 1952, p. 27--50).

\bibitem[124]{article:124}\phantomsection\label{bib:article:124}
\textsc{F. B. Wright}, \emph{A reduction for algebras of finite type}
(Ann. Math., t. 60, 1954, p. 560--570).

\bibitem[125]{article:125}\phantomsection\label{bib:article:125}
\textsc{K. Yosida}, \emph{On the unitary equivalence in general Euclid space}
(Proc. Jap. Acad., t. 22, 1946, p. 242--245).

\subsection*{Second Part}\label{bib:journal-second-part}

\begin{center}\bfseries 1954.\end{center}

\bibitem[126]{article:126}\phantomsection\label{bib:article:126}
\textsc{J. Feldman et R. V. Kadison}, \emph{The closure of the regular operators in a ring of operators}
(Proc. Amer. Math. Soc., t. 5, p. 909--916).

\bibitem[127]{article:127}\phantomsection\label{bib:article:127}
\textsc{E. Hongo et M. Orihara}, \emph{A remark on a quasi-unitary algebra}
(Yokohama Math. J., t. 2, p. 69--72).

\bibitem[128]{article:128}\phantomsection\label{bib:article:128}
\textsc{Y. Misonou}, \emph{On the direct product of W*-algebras}
(Tôhoku Math. J., t. 6, p. 189--204).

\bibitem[129]{article:129}\phantomsection\label{bib:article:129}
\textsc{M. Nakamura}, \emph{On the direct product of finite factors}
(Tôhoku Math. J., t. 6, p. 205--207).

\bibitem[130]{article:130}\phantomsection\label{bib:article:130}
\textsc{M. Nakamura et T. Turumaru}, \emph{Expectations in an operator algebra}
(Tôhoku Math. J., t. 6, p. 182--188).

\bibitem[131]{article:131}\phantomsection\label{bib:article:131}
\textsc{M. Nakamura et T. Turumaru}, \emph{On extensions of pure states of an abelian operator algebra}
(Tôhoku Math. J., t. 6, p. 253--257).

\bibitem[132]{article:132}\phantomsection\label{bib:article:132}
\textsc{I. E. Segal}, \emph{Abstract probability spaces and a theorem of Kolmogoroff}
(Amer. J. Math., t. 76, p. 721--732).

\bibitem[133]{article:133}\phantomsection\label{bib:article:133}
\textsc{T. Turumaru}, \emph{On the direct product of operator algebras, III}
(Tôhoku Math. J., t. 6, p. 208--211).

\bibitem[134]{article:134}\phantomsection\label{bib:article:134}
\textsc{H. Umegaki}, \emph{Conditional expectation in an operator algebra, I}
(Tôhoku Math. J., t. 6, p. 177--181).

\begin{center}\bfseries 1955.\end{center}

\bibitem[135]{article:135}\phantomsection\label{bib:article:135}
\textsc{C. Davis}, \emph{Generators of the ring of bounded operators}
(Proc. Amer. Math. Soc., t. 6, p. 970--972).

\bibitem[136]{article:136}\phantomsection\label{bib:article:136}
\textsc{E. L. Griffin}, \emph{Some contributions to the theory of rings of operators, II}
(Trans. Amer. Math. Soc., t. 79, p. 389--400).

\bibitem[137]{article:137}\phantomsection\label{bib:article:137}
\textsc{E. Hongo}, \emph{A note on the commutor of certain operator algebras}
(Bull. Kyushu inst. tech., t. 1, p. 19--22).

\bibitem[138]{article:138}\phantomsection\label{bib:article:138}
\textsc{R. V. Kadison}, \emph{Multiplicity theory for operator algebras}
(Proc. Nat. Acad. Sc. U. S. A., t. 41, p. 169--173).

\bibitem[139]{article:139}\phantomsection\label{bib:article:139}
\textsc{R. V. Kadison}, \emph{Isomorphisms of factors of infinite type}
(Can. J. Math., t. 7, p. 322--327).

\bibitem[140]{article:140}\phantomsection\label{bib:article:140}
\textsc{R. V. Kadison}, \emph{On the additivity of the trace in finite factors}
(Proc. Nat. Acad. Sc. U. S. A., t. 41, p. 385--387).

\bibitem[141]{article:141}\phantomsection\label{bib:article:141}
\textsc{I. Kaplansky}, \emph{Rings of operators} (multigraphié, University of Chicago).

\bibitem[142]{article:142}\phantomsection\label{bib:article:142}
\textsc{L. H. Loomis}, \emph{The lattice theoretic background of the dimension theory of operator algebras}
(Mem. Amer. Math. Soc., n\textsuperscript{o} 18).

\bibitem[143]{article:143}\phantomsection\label{bib:article:143}
\textsc{S. Maeda}, \emph{Dimension functions on certain general lattices}
(J. sci. Hiroshima Univ., t. 19, p. 211--237).

% Source PDF page 352 (printed p. 343)
\phantomsection\label{srcpage:343}
\bibitem[144]{article:144}\phantomsection\label{bib:article:144}
\textsc{Y. Misonou}, \emph{Generalized approximately finite operator algebras}
(Tôhoku Math. J., t. 7, p. 192--205).

\bibitem[145]{article:145}\phantomsection\label{bib:article:145}
\textsc{M. A. Naimark et S. V. Fomin}, \emph{Sommes directes continues d'espaces hilbertiens et quelques applications} (en russe)
(Usp. Mat. nauk, t. 10, p. 111--142).

\bibitem[146]{article:146}\phantomsection\label{bib:article:146}
\textsc{T. Ogasawara}, \emph{A structure theorem for complete quasi-unitary algebras}
(J. sci. Hiroshima Univ., t. 19, p. 79--85).

\bibitem[147]{article:147}\phantomsection\label{bib:article:147}
\textsc{T. Ogasawara}, \emph{Topologies on rings of operators}
(J. sci. Hiroshima Univ., t. 19, p. 255--272).

\bibitem[148]{article:148}\phantomsection\label{bib:article:148}
\textsc{T. Ogasawara et K. Yoshinaga}, \emph{A non-commutative theory of integration for operators}
(J. sci. Hiroshima Univ., t. 18, p. 311--347).

\bibitem[149]{article:149}\phantomsection\label{bib:article:149}
\textsc{T. Ogasawara et K. Yoshinaga}, \emph{Extension of \(\natural\)-application to unbounded operators}
(J. sci. Hiroshima Univ., t. 19, p. 273--299).

\bibitem[150]{article:150}\phantomsection\label{bib:article:150}
\textsc{S. Sakai}, \emph{On the group isomorphism of unitary groups in AW*-algebras}
(Tôhoku Math. J., t. 7, p. 87--95).

\bibitem[151]{article:151}\phantomsection\label{bib:article:151}
\textsc{U. Sasaki}, \emph{Lattices of projections in AW*-algebras}
(J. sci. Hiroshima Univ., t. 19, p. 1--30).

\bibitem[152]{article:152}\phantomsection\label{bib:article:152}
\textsc{H. Sunouchi}, \emph{A characterization of the maximal ideal in a factor, II}
(Kōdai Math. sem. rep., t. 7, p. 65--66).

\bibitem[153]{article:153}\phantomsection\label{bib:article:153}
\textsc{N. Suzuki}, \emph{On the invariants of W*-algebras}
(Tôhoku Math. J., t. 7, p. 177--185).

\bibitem[154]{article:154}\phantomsection\label{bib:article:154}
\textsc{N. Suzuki}, \emph{On automorphisms of W*-algebras leaving the center elementwise invariant}
(Tôhoku Math. J., t. 7, p. 186--191).

\bibitem[155]{article:155}\phantomsection\label{bib:article:155}
\textsc{Z. Takeda}, \emph{Inductive limit and infinite direct product of operator algebras}
(Tôhoku Math. J., t. 7, p. 67--86).

\bibitem[156]{article:156}\phantomsection\label{bib:article:156}
\textsc{C. T. Ionescu Tulcea}, \emph{Deux théorèmes concernant certains espaces de champs de vecteurs}
(Bull. Sc. math., t. 79, p. 1--6).

\bibitem[157]{article:157}\phantomsection\label{bib:article:157}
\textsc{Ti. Yen}, \emph{Trace on finite AW*-algebras}
(Duke Math. J., t. 22, p. 207--222).

\begin{center}\bfseries 1956.\end{center}

\bibitem[158]{article:158}\phantomsection\label{bib:article:158}
\textsc{S. K. Berberian}, \emph{On the projection geometry of a finite AW*-algebra}
(Trans. Amer. Math. Soc., t. 83, p. 493--509).

\bibitem[159]{article:159}\phantomsection\label{bib:article:159}
\textsc{J. Feldman}, \emph{Isomorphisms of finite type II rings of operators}
(Ann. Math., t. 63, p. 565--571).

\bibitem[160]{article:160}\phantomsection\label{bib:article:160}
\textsc{J. Feldman}, \emph{Embedding of AW*-algebras}
(Duke Math. J., t. 23, p. 303--308).

\bibitem[161]{article:161}\phantomsection\label{bib:article:161}
\textsc{J. Feldman}, \emph{Some connections between topological and algebraical properties in rings of operators}
(Duke Math. J., t. 23, p. 365--370).

\bibitem[162]{article:162}\phantomsection\label{bib:article:162}
\textsc{J. Feldman}, \emph{Nonseparability of certain finite factors}
(Proc. Amer. Math. Soc., t. 7, p. 23--26).

\bibitem[163]{article:163}\phantomsection\label{bib:article:163}
\textsc{M. Goldman}, \emph{Structure of AW*-algebras, I}
(Duke Math. J., t. 23, p. 23--34).

\bibitem[164]{article:164}\phantomsection\label{bib:article:164}
\textsc{H. Gonshor}, \emph{Spectral theory for a class of non-normal operators}
(Can. J. Math., t. 8, p. 449--461).

\bibitem[165]{article:165}\phantomsection\label{bib:article:165}
\textsc{E. Hongo}, \emph{On left rings of certain *-algebras}
(Bull. Kyushu inst. tech., t. 2, p. 1--15).

\bibitem[166]{article:166}\phantomsection\label{bib:article:166}
\textsc{R. V. Kadison}, \emph{Operator algebras with a faithful weakly-closed representation}
(Ann. Math., t. 64, p. 175--181).

\bibitem[167]{article:167}\phantomsection\label{bib:article:167}
\textsc{S. Maeda}, \emph{Lenghts of projections in rings of operators}
(J. sci. Hiroshima Univ., t. 20, p. 5--11).

\bibitem[168]{article:168}\phantomsection\label{bib:article:168}
\textsc{Y. Misonou}, \emph{On divisors of factors}
(Tôhoku Math. J., t. 8, p. 63--69).

\bibitem[169]{article:169}\phantomsection\label{bib:article:169}
\textsc{M. Nakamura et H. Umegaki}, \emph{On a proposition of von Neumann}
(Kōdai Math. sem. rep., t. 8, p. 142--144).

\bibitem[170]{article:170}\phantomsection\label{bib:article:170}
\textsc{T. Ogasawara et S. Maeda}, \emph{A generalization of a theorem of Dye}
(J. sci. Hiroshima Univ., t. 20, p. 1--4).

\bibitem[171]{article:171}\phantomsection\label{bib:article:171}
\textsc{L. Pukánszky}, \emph{Some examples of factors}
(Publicationes math., t. 4, p. 135--156).

% Source PDF page 353 (printed p. 344)
\phantomsection\label{srcpage:344}
\bibitem[172]{article:172}\phantomsection\label{bib:article:172}
\textsc{S. Sakai}, \emph{The absolute value of W*-algebras of finite type}
(Tôhoku Math. J., t. 8, p. 70--85).

\bibitem[173]{article:173}\phantomsection\label{bib:article:173}
\textsc{S. Sakai}, \emph{On the \(\sigma\)-weak topology of W*-algebras}
(Proc. Japan Acad., t. 32, p. 329--332).

\bibitem[174]{article:174}\phantomsection\label{bib:article:174}
\textsc{S. Sakai}, \emph{A characterization of W*-algebras}
(Pacific J. Math., t. 6, p. 763--773).

\bibitem[175]{article:175}\phantomsection\label{bib:article:175}
\textsc{I. E. Segal}, \emph{Tensor algebras over Hilbert spaces}
(Trans. Amer. Math. Soc., t. 81, p. 106--134).

\bibitem[176]{article:176}\phantomsection\label{bib:article:176}
\textsc{I. E. Segal}, \emph{Tensor algebras over Hilbert spaces, II}
(Ann. Math., t. 63, p. 160--175).

\bibitem[177]{article:177}\phantomsection\label{bib:article:177}
\textsc{H. Sunouchi}, \emph{Infinite Lie rings}
(Tôhoku Math. J., t. 8, p. 291--307).

\bibitem[178]{article:178}\phantomsection\label{bib:article:178}
\textsc{H. Umegaki}, \emph{Weak compactness in an operator space}
(Kōdai Math. sem. rep., t. 8, p. 145--151).

\bibitem[179]{article:179}\phantomsection\label{bib:article:179}
\textsc{H. Umegaki}, \emph{Conditional expectation in an operator algebra, II}
(Tôhoku Math. J., t. 8, p. 86--100).

\bibitem[180]{article:180}\phantomsection\label{bib:article:180}
\textsc{H. Widom}, \emph{Embedding in algebras of type I}
(Duke Math. J., t. 23, p. 309--324).

\bibitem[181]{article:181}\phantomsection\label{bib:article:181}
\textsc{H. Widom}, \emph{Approximately finite algebras}
(Trans. Amer. Math. Soc., t. 83, p. 170--178).

\bibitem[182]{article:182}\phantomsection\label{bib:article:182}
\textsc{Ti. Yen}, \emph{Isomorphisms of unitary groups in AW*-algebras}
(Tôhoku Math. J., t. 8, p. 275--280).

\bibitem[183]{article:183}\phantomsection\label{bib:article:183}
\textsc{Ti. Yen}, \emph{Quotient algebra of a finite AW*-algebra}
(Pacific J. Math., t. 6, p. 389--395).

\begin{center}\bfseries 1957.\end{center}

\bibitem[184]{article:184}\phantomsection\label{bib:article:184}
\textsc{S. K. Berberian}, \emph{The regular ring of a finite AW*-algebra}
(Ann. Math., t. 65, p. 224--240).

\bibitem[185]{article:185}\phantomsection\label{bib:article:185}
\textsc{J. Feldman et J. M. G. Fell}, \emph{Separable representations of rings of operators}
(Ann. Math., t. 65, p. 241--249).

\bibitem[186]{article:186}\phantomsection\label{bib:article:186}
\textsc{J. M. G. Fell}, \emph{Representations of weakly closed algebras}
(Math. Ann., t. 133, p. 118--126).

\bibitem[187]{article:187}\phantomsection\label{bib:article:187}
\textsc{A. M. Gleason}, \emph{Measures on the closed subspaces of a Hilbert space}
(J. of rational mech. and analysis, t. 6, p. 885--894).

\bibitem[188]{article:188}\phantomsection\label{bib:article:188}
\textsc{A. Grothendieck}, \emph{Un résultat sur le dual d'une \(C^*\)-algèbre}
(J. Math. pures et appl., t. 36, p. 97--108).

\bibitem[189]{article:189}\phantomsection\label{bib:article:189}
\textsc{E. Hongo}, \emph{On quasi-unitary algebras with semi-finite left rings}
(Bull. Kyushu inst. tech., t. 3, p. 1--10).

\bibitem[190]{article:190}\phantomsection\label{bib:article:190}
\textsc{R. V. Kadison}, \emph{Unitary invariants for representations of operator algebras}
(Ann. Math., t. 66, p. 304--379).

\bibitem[191]{article:191}\phantomsection\label{bib:article:191}
\textsc{R. V. Kadison et I. M. Singer}, \emph{Three test problems in operator theory}
(Pacific J. Math., t. 7, p. 1101--1106).

\bibitem[192]{article:192}\phantomsection\label{bib:article:192}
\textsc{S. Sakai}, \emph{On topological properties of W*-algebras}
(Proc. Japan Acad., t. 33, p. 439--444).

\bibitem[193]{article:193}\phantomsection\label{bib:article:193}
\textsc{E. Thoma}, \emph{Zur Reduktionstheorie in separablen Hilbert-Räumen}
(Math. Zeitschr., t. 67, p. 1--9).

\bibitem[194]{article:194}\phantomsection\label{bib:article:194}
\textsc{E. Thoma}, \emph{Zur Reduktionstheorie in allgemeinen Hilbert-Räumen}
(Math. Zeitschr., t. 68, p. 153--188).

\bibitem[195]{article:195}\phantomsection\label{bib:article:195}
\textsc{J. Tomiyama}, \emph{On the projection of norm one in W*-algebras}
(Proc. Japan Acad., t. 33, p. 608--612).

\bibitem[196]{article:196}\phantomsection\label{bib:article:196}
\textsc{K. Tsuji}, \emph{W*-algebras and abstract L-spaces}
(Bull. Kyushu inst. tech., t. 3, p. 11--13).

\bibitem[197]{article:197}\phantomsection\label{bib:article:197}
\textsc{H. Widom}, \emph{Nonisomorphic approximately finite factors}
(Proc. Amer. Math. Soc., t. 8, p. 537--540).

\bibitem[198]{article:198}\phantomsection\label{bib:article:198}
\textsc{Ti Yen}, \emph{Isomorphisms of AW*-algebras}
(Proc. Amer. Math. Soc., t. 8, p. 345--349).

% Source PDF page 354 (printed p. 345)
\phantomsection\label{srcpage:345}
\begin{center}\bfseries 1958.\end{center}

\bibitem[199]{article:199}\phantomsection\label{bib:article:199}
\textsc{S. K. Berberian}, \emph{\(N\times N\) matrices over an AW*-algebra}
(Amer. J. Math., t. 80, p. 37--44).

\bibitem[200]{article:200}\phantomsection\label{bib:article:200}
\textsc{R. J. Blattner}, \emph{Automorphic group representations}
(Pacific J. Math., t. 8, p. 665--677).

\bibitem[201]{article:201}\phantomsection\label{bib:article:201}
\textsc{J. Feldman}, \emph{The uniformly closed ideals in an AW*-algebra}
(Bull. Amer. Math. Soc., t. 62, p. 245--246).

\bibitem[202]{article:202}\phantomsection\label{bib:article:202}
\textsc{H. Gonshor}, \emph{Spectral theory for a class of non-normal operators, II}
(Can. J. Math., t. 10, p. 97--102).

\bibitem[203]{article:203}\phantomsection\label{bib:article:203}
\textsc{E. Hongo}, \emph{On some properties of quasi-unitary algebras}
(Bull. Kyushu inst. tech., t. 4, p. 1--6).

\bibitem[204]{article:204}\phantomsection\label{bib:article:204}
\textsc{R. V. Kadison}, \emph{Theory of operators; Part II: Operator algebras}
(Bull. Amer. Math. Soc., t. 64, p. 61--85).

\bibitem[205]{article:205}\phantomsection\label{bib:article:205}
\textsc{H. Leptin}, \emph{Zur Reduktionstheorie Hilbertscher Räume}
(Math. Zeitschr., t. 69, p. 40--58).

\bibitem[206]{article:206}\phantomsection\label{bib:article:206}
\textsc{H. Leptin}, \emph{Reduktion linearer Funktionale auf Operatorringen}
(Abh. math. Sem. Univ. Hamburg, t. 22, p. 98--113).

\bibitem[207]{article:207}\phantomsection\label{bib:article:207}
\textsc{S. Maeda}, \emph{On the lattice of projection of a Baer *-ring}
(J. sci. Hiroshima Univ., t. 22, p. 75--88).

\bibitem[208]{article:208}\phantomsection\label{bib:article:208}
\textsc{M. Nakai}, \emph{Some expectations in AW*-algebras}
(Proc. Japan Acad., t. 34, p. 411--416).

\bibitem[209]{article:209}\phantomsection\label{bib:article:209}
\textsc{M. Nakamura}, \emph{A proof of a theorem of Takesaki}
(Kōdai Math. sem. rep., t. 10, p. 189--190).

\bibitem[210]{article:210}\phantomsection\label{bib:article:210}
\textsc{M. Nakamura et Z. Takeda}, \emph{On some elementary properties of the crossed products of von Neumann algebras}
(Proc. Japan Acad., t. 34, p. 489--494).

\bibitem[211]{article:211}\phantomsection\label{bib:article:211}
\textsc{M. Nakamura et Z. Takeda}, \emph{On certain examples of the crossed product of finite factors, I--II}
(Proc. Japan Acad., t. 34, p. 495--499 et 500--502).

\bibitem[212]{article:212}\phantomsection\label{bib:article:212}
\textsc{T. Ono}, \emph{Local theory of rings of operators, I--II}
(J. Math. Soc. Japan, t. 10, p. 184--216 et 438--458).

\bibitem[213]{article:213}\phantomsection\label{bib:article:213}
\textsc{T. Saito}, \emph{On incomplete infinite direct product of W*-algebras}
(Tôhoku Math. J., t. 10, p. 165--171).

\bibitem[214]{article:214}\phantomsection\label{bib:article:214}
\textsc{S. Sakai}, \emph{On linear functionals of W*-algebras}
(Proc. Japan Acad., t. 34, p. 571--574).

\bibitem[215]{article:215}\phantomsection\label{bib:article:215}
\textsc{I. E. Segal}, \emph{Distributions in Hilbert space and canonical systems of operators}
(Trans. Amer. Math. Soc., t. 88, p. 12--41).

\bibitem[216]{article:216}\phantomsection\label{bib:article:216}
\textsc{N. Suzuki}, \emph{A linear representation of a countably infinite group}
(Proc. Japan Acad., t. 34, p. 575--579).

\bibitem[217]{article:217}\phantomsection\label{bib:article:217}
\textsc{M. Takesaki}, \emph{On the direct product of W*-factors}
(Tôhoku Math. J., t. 10, p. 116--119).

\bibitem[218]{article:218}\phantomsection\label{bib:article:218}
\textsc{M. Takesaki}, \emph{On the conjugate space of operator algebra}
(Tôhoku Math. J., t. 10, p. 194--203).

\bibitem[219]{article:219}\phantomsection\label{bib:article:219}
\textsc{M. Takesaki}, \emph{A note on the cross norm of the direct product of operator algebra}
(Kōdai Math. sem. rep., t. 10, p. 137--140).

\bibitem[220]{article:220}\phantomsection\label{bib:article:220}
\textsc{S. Teleman}, \emph{Sur les algèbres de J. von Neumann}
(Bull. Sc. math., t. 82, p. 117--126).

\bibitem[221]{article:221}\phantomsection\label{bib:article:221}
\textsc{J. Tomiyama}, \emph{On the projection of norm one in W*-algebras, II}
(Tôhoku Math. J., t. 10, p. 204--209).

\bibitem[222]{article:222}\phantomsection\label{bib:article:222}
\textsc{J. Tomiyama}, \emph{A remark on the invariants of W*-algebras}
(Tôhoku Math. J., t. 10, p. 37--41).

\bibitem[223]{article:223}\phantomsection\label{bib:article:223}
\textsc{J. Tomiyama}, \emph{Generalized dimension function for W*-algebras of infinite type}
(Tôhoku Math. J., t. 10, p. 121--129).

\bibitem[224]{article:224}\phantomsection\label{bib:article:224}
\textsc{T. Turumaru}, \emph{Crossed product of operator algebra}
(Tôhoku Math. J., t. 10, p. 355--365).

\bibitem[225]{article:225}\phantomsection\label{bib:article:225}
\textsc{F. B. Wright}, \emph{The ideals in a factor}
(Ann. Math., t. 68, p. 475--483).

% Source PDF page 355 (printed p. 346)
\phantomsection\label{srcpage:346}
\begin{center}\bfseries 1959.\end{center}

\bibitem[226]{article:226}\phantomsection\label{bib:article:226}
\textsc{C. Davis}, \emph{Various averaging operations onto subalgebras}
(Ill. J. Math., t. 3, p. 538--553).

\bibitem[227]{article:227}\phantomsection\label{bib:article:227}
\textsc{H. A. Dye}, \emph{On groups of measure preserving transformations, I}
(Amer. J. Math., t. 81, p. 119--159).

\bibitem[228]{article:228}\phantomsection\label{bib:article:228}
\textsc{M. Goldman}, \emph{On subfactors of factors of type \(II_1\)}
(Michigan Math. J., t. 6, p. 167--172).

\bibitem[229]{article:229}\phantomsection\label{bib:article:229}
\textsc{E. Hongo}, \emph{On left multiplicative operators on a quasi-unitary algebra}
(Bull. Kyushu inst. tech., t. 5, p. 19--22).

\bibitem[230]{article:230}\phantomsection\label{bib:article:230}
\textsc{R. V. Kadison et I. M. Singer}, \emph{Extensions of pure states}
(Amer. J. Math., t. 81, p. 383--400).

\bibitem[231]{article:231}\phantomsection\label{bib:article:231}
\textsc{M. Nakamura et Z. Takeda}, \emph{On the extensions of finite factors, I}
(Proc. Japan Acad., t. 35, p. 149--154).

\bibitem[232]{article:232}\phantomsection\label{bib:article:232}
\textsc{T. Saito}, \emph{The direct and crossed product of rings of operators}
(Tôhoku Math. J., t. 11, p. 299--304).

\bibitem[233]{article:233}\phantomsection\label{bib:article:233}
\textsc{S. Sankaran}, \emph{The *-algebra of unbounded operators}
(J. London Math. Soc., t. 34, p. 337--344).

\bibitem[234]{article:234}\phantomsection\label{bib:article:234}
\textsc{P. C. Shields}, \emph{A new topology for von Neumann algebras}
(Bull. Amer. Math. Soc., t. 65, p. 267--269).

\bibitem[235]{article:235}\phantomsection\label{bib:article:235}
\textsc{N. Suzuki}, \emph{Crossed products of rings of operators}
(Tôhoku Math. J., t. 11, p. 113--124).

\bibitem[236]{article:236}\phantomsection\label{bib:article:236}
\textsc{N. Suzuki}, \emph{Certain types of groups of automorphisms of a factor}
(Tôhoku Math. J., t. 11, p. 314--320 et t. 12, p. 477--478).

\bibitem[237]{article:237}\phantomsection\label{bib:article:237}
\textsc{Z. Takeda}, \emph{On the extensions of finite factors, II}
(Proc. Japan Acad., t. 35, p. 215--220).

\bibitem[238]{article:238}\phantomsection\label{bib:article:238}
\textsc{M. Takesaki}, \emph{On the singularity of a positive linear functional on operator algebra}
(Proc. Japan Acad., t. 35, p. 365--366).

\bibitem[239]{article:239}\phantomsection\label{bib:article:239}
\textsc{M. Tomita}, \emph{Spectral theory of operator algebras, I}
(Math. J. Okayama Univ., t. 9, p. 63--98).

\bibitem[240]{article:240}\phantomsection\label{bib:article:240}
\textsc{J. Tomiyama}, \emph{On the projection of norm one in W*-algebras, III}
(Tôhoku Math. J., t. 11, p. 125--129).

\bibitem[241]{article:241}\phantomsection\label{bib:article:241}
\textsc{H. Umegaki}, \emph{Conditional expectation in an operator algebra, III}
(Kōdai Math. sem. rep., t. 11, p. 51--64).

\begin{center}\bfseries 1960.\end{center}

\bibitem[242]{article:242}\phantomsection\label{bib:article:242}
\textsc{I. Halperin}, \emph{Introduction to von Neumann algebras and continuous geometry}
(Can. Math. bull., t. 3, p. 273--288 et t. 5, p. 59).

\bibitem[243]{article:243}\phantomsection\label{bib:article:243}
\textsc{R. V. Kadison et I. M. Singer}, \emph{Triangular operator algebras. Fundamentals and hyperreducible theory}
(Amer. J. Math., t. 82, p. 227--259).

\bibitem[244]{article:244}\phantomsection\label{bib:article:244}
\textsc{I. Kovacs}, \emph{Un complément à la théorie de l'« intégration non commutative »}
(Acta Univ. Szeged, t. 21, p. 7--11).

\bibitem[245]{article:245}\phantomsection\label{bib:article:245}
\textsc{M. A. Naimark}, \emph{Représentations factorielles d'un groupe localement compact} (en russe)
(Dokl. Akad. nauk S. S. S. R., t. 134, p. 275--277).

\bibitem[246]{article:246}\phantomsection\label{bib:article:246}
\textsc{M. Nakamura et Z. Takeda}, \emph{A Galois theory for finite factors}
(Proc. Japan Acad., t. 36, p. 258--260).

\bibitem[247]{article:247}\phantomsection\label{bib:article:247}
\textsc{M. Nakamura et Z. Takeda}, \emph{On the fundamental theorem of the Galois theory for finite factors}
(Proc. Japan Acad., t. 36, p. 313--318).

\bibitem[248]{article:248}\phantomsection\label{bib:article:248}
\textsc{M. Nakamura, M. Takesaki et H. Umegaki}, \emph{A remark on the expectation of operator algebras}
(Kōdai Math. sem. rep., t. 12, p. 82--90).

\bibitem[249]{article:249}\phantomsection\label{bib:article:249}
\textsc{L. Pukánszky}, \emph{On maximal abelian subrings of factors of type \(II_1\)}
(Can. J. Math., t. 12, p. 289--296).

% Source PDF page 356 (printed p. 347)
\phantomsection\label{srcpage:347}
\bibitem[250]{article:250}\phantomsection\label{bib:article:250}
\textsc{T. Saito}, \emph{Some remarks on a representation of a group}
(Tôhoku Math. J., t. 12, p. 383--388).

\bibitem[251]{article:251}\phantomsection\label{bib:article:251}
\textsc{T. Saito et J. Tomiyama}, \emph{Some results on the product of W*-algebras}
(Tôhoku Math. J., t. 12, p. 455--458).

\bibitem[252]{article:252}\phantomsection\label{bib:article:252}
\textsc{S. Sakai}, \emph{On a conjecture of Kaplansky}
(Tôhoku Math. J., t. 12, p. 31--33).

\bibitem[253]{article:253}\phantomsection\label{bib:article:253}
\textsc{R. Schatten}, \emph{Norm ideals of completely continuous operators}, Erg. der Math., Springer-Verlag, Berlin-Göttingen-Heidelberg, 1960.

\bibitem[254]{article:254}\phantomsection\label{bib:article:254}
\textsc{I. E. Segal}, \emph{A note on the concept of entropy}
(J. Math. Mech., t. 9, p. 623--630).

\bibitem[255]{article:255}\phantomsection\label{bib:article:255}
\textsc{O. Takenouchi}, \emph{Sur les algèbres de Hilbert}
(C. R. Acad. Sc., t. 250, p. 3436--3437).

\bibitem[256]{article:256}\phantomsection\label{bib:article:256}
\textsc{M. Takesaki}, \emph{On the non-separability of singular representations of operator algebra}
(Kōdai Math. sem. rep., t. 12, p. 102--108).

\bibitem[257]{article:257}\phantomsection\label{bib:article:257}
\textsc{M. Takesaki}, \emph{On the Hahn-Banach type theorem and the Jordan decomposition of module linear mapping over some operator algebra}
(Kōdai Math. sem. rep., t. 12, p. 1--10).

\bibitem[258]{article:258}\phantomsection\label{bib:article:258}
\textsc{M. Tomita}, \emph{Spectral theory of operator algebras, II}
(Math. J. Okayama Univ., t. 10, p. 19--60).

\begin{center}\bfseries 1961.\end{center}

\bibitem[259]{article:259}\phantomsection\label{bib:article:259}
\textsc{C. Davis}, \emph{Operator-valued entropy of a quantum mechanical measurement}
(Proc. Japan Acad., t. 37, p. 533--538).

\bibitem[260]{article:260}\phantomsection\label{bib:article:260}
\textsc{A. Guichardet}, \emph{Une caractérisation des algèbres de von Neumann discrètes}
(Bull. Soc. Math. France, t. 89, p. 77--101).

\bibitem[261]{article:261}\phantomsection\label{bib:article:261}
\textsc{E. Hongo}, \emph{A structure theory for semi-finite quasi-unitary algebras}
(Bull. Kyushu inst. tech., t. 7, p. 1--17).

\bibitem[262]{article:262}\phantomsection\label{bib:article:262}
\textsc{R. V. Kadison}, \emph{The trace in finite operator algebras}
(Proc. Amer. Math. Soc., t. 12, p. 973--977).

\bibitem[263]{article:263}\phantomsection\label{bib:article:263}
\textsc{I. Kovacs}, \emph{Théorèmes ergodiques non commutatifs}
(C. R. Acad. Sc., t. 253, p. 770--771).

\bibitem[264]{article:264}\phantomsection\label{bib:article:264}
\textsc{I. Kovacs}, \emph{Sur certains automorphismes des algèbres hilbertiennes}
(Acta Sc. Math. Szeged, t. 22, p. 234--242).

\bibitem[265]{article:265}\phantomsection\label{bib:article:265}
\textsc{M. A. Naimark}, \emph{Décomposition en représentations factorielles des représentations unitaires des groupes localement compacts} (en russe)
(Sibirsk. mat. Z., t. 2, p. 89--99).

\bibitem[266]{article:266}\phantomsection\label{bib:article:266}
\textsc{M. Nakamura et Z. Takeda}, \emph{On inner automorphisms of certain finite factors}
(Proc. Japan Acad., t. 37, p. 31--32).

\bibitem[267]{article:267}\phantomsection\label{bib:article:267}
\textsc{M. Nakamura et Z. Takeda}, \emph{On outer automorphisms of certain finite factors}
(Proc. Japan Acad., t. 37, p. 215--216).

\bibitem[268]{article:268}\phantomsection\label{bib:article:268}
\textsc{M. Nakamura et H. Umegaki}, \emph{A note on the entropy for operator algebras}
(Proc. Japan Acad., t. 37, p. 149--154).

\bibitem[269]{article:269}\phantomsection\label{bib:article:269}
\textsc{M. Nakamura et H. Umegaki}, \emph{On the Blackwell theorem in operator algebras}
(Proc. Japan Acad., t. 37, p. 312--315).

\bibitem[270]{article:270}\phantomsection\label{bib:article:270}
\textsc{H. Reeh et S. Schlieder}, \emph{Bemerkungen zur Unitarequivalenz von Lorentz-invarianten Feldern}
(Nuovo Cimento, t. 22, p. 1051--1068).

\bibitem[271]{article:271}\phantomsection\label{bib:article:271}
\textsc{T. Saito}, \emph{On a representation of a countably infinite group}
(Tôhoku Math. J., t. 13, p. 268--273).

\bibitem[272]{article:272}\phantomsection\label{bib:article:272}
\textsc{T. Saito}, \emph{On groups of automorphisms of finite factors}
(Tôhoku Math. J., t. 13, p. 427--433).

\bibitem[273]{article:273}\phantomsection\label{bib:article:273}
\textsc{S. Sankaran}, \emph{Ordered decompositions of Hilbert spaces}
(J. London Math. Soc., t. 36, p. 97--107).

\bibitem[274]{article:274}\phantomsection\label{bib:article:274}
\textsc{S. Sankaran}, \emph{Decomposition of von Neumann algebras of type I}
(Math. Ann., t. 142, p. 399--406).

% Source PDF page 357 (printed p. 348)
\phantomsection\label{srcpage:348}
\bibitem[275]{article:275}\phantomsection\label{bib:article:275}
\textsc{Z. Takeda}, \emph{On the extension theorem of the Galois theory for finite factors}
(Proc. Japan Acad., t. 37, p. 78--82).

\bibitem[276]{article:276}\phantomsection\label{bib:article:276}
\textsc{Z. Takeda}, \emph{On the normal basis theorem of the Galois theory for finite factors}
(Proc. Japan Acad., t. 37, p. 144--148).

\bibitem[277]{article:277}\phantomsection\label{bib:article:277}
\textsc{O. Takenouchi}, \emph{Sur les sous-algèbres d'une algèbre de Hilbert}
(Ann. scient. Éc. Norm. Sup., t. 78, p. 211--240).

\bibitem[278]{article:278}\phantomsection\label{bib:article:278}
\textsc{H. Umegaki}, \emph{On information in operator algebras}
(Proc. Japan Acad., t. 37, p. 459--461).

\begin{center}\bfseries 1962.\end{center}

\bibitem[279]{article:279}\phantomsection\label{bib:article:279}
\textsc{H. J. Borchers}, \emph{On the structure of the algebra of field operators I}
(Nuovo Cimento, t. 24, p. 214--236).

\bibitem[280]{article:280}\phantomsection\label{bib:article:280}
\textsc{M. Echigo et M. Nakamura}, \emph{A remark on the concept of channels}
(Proc. Jap. Acad., t. 38, p. 307--309).

\bibitem[281]{article:281}\phantomsection\label{bib:article:281}
\textsc{C. Foias et I. Kovacs}, \emph{Une caractérisation nouvelle des algèbres de von Neumann finies}
(Acta Univ. Szeged, t. 23, p. 274--278).

\bibitem[282]{article:282}\phantomsection\label{bib:article:282}
\textsc{R. Haag et B. Schroer}, \emph{Postulates of quantum field theory}
(J. Math. Phys., t. 3, p. 248--256).

\bibitem[283]{article:283}\phantomsection\label{bib:article:283}
\textsc{R. V. Kadison}, \emph{Normalcy in operator algebras}
(Duke Math. J., t. 29, p. 459--464).

\bibitem[284]{article:284}\phantomsection\label{bib:article:284}
\textsc{L. H. Loomis}, \emph{Unique direct integral decompositions of convex sets}
(Amer. J. Math., t. 84, p. 509--526).

\bibitem[285]{article:285}\phantomsection\label{bib:article:285}
\textsc{E. J. Mc Shane}, \emph{Families of measures and representation of algebras of operators}
(Trans. Amer. Math. Soc., t. 102, p. 328--345).

\bibitem[286]{article:286}\phantomsection\label{bib:article:286}
\textsc{M. Nakamura et H. Umegaki}, \emph{On von Neumann's theory of measurements in quantum statistics}
(Math. Japonicæ, t. 7, p. 151--157).

\bibitem[287]{article:287}\phantomsection\label{bib:article:287}
\textsc{C. Pearcy}, \emph{A complete set of unitary invariants for operators generating finite W*-algebras of type I}
(Pacific J. Math., t. 12, p. 1405--1416).

\bibitem[288]{article:288}\phantomsection\label{bib:article:288}
\textsc{C. Pearcy}, \emph{W*-algebras with a single generator}
(Proc. Amer. Math. Soc., t. 13, p. 831--832).

\bibitem[289]{article:289}\phantomsection\label{bib:article:289}
\textsc{S. Sakai}, \emph{The theory of W*-algebras} (miméographié, Yale University).

\bibitem[290]{article:290}\phantomsection\label{bib:article:290}
\textsc{N. Suzuki}, \emph{Extensions of rings of operators in Hilbert spaces}
(Tôhoku Math. J., t. 14, p. 217--232).

\bibitem[291]{article:291}\phantomsection\label{bib:article:291}
\textsc{J. Tomiyama}, \emph{Topological representations of \(C^*\)-algebras}
(Tôhoku Math. J., t. 14, p. 187--204).

\bibitem[292]{article:292}\phantomsection\label{bib:article:292}
\textsc{H. Umegaki}, \emph{Entropy functional in stationary channels}
(Proc. Japan Acad., t. 38, p. 668--672).

\bibitem[293]{article:293}\phantomsection\label{bib:article:293}
\textsc{H. Umegaki}, \emph{Conditional expectation in an operator algebra; IV. Entropy and information}
(Kōdai Math. sem. rep., t. 14, p. 59--85).

\begin{center}\bfseries 1963.\end{center}

\bibitem[294]{article:294}\phantomsection\label{bib:article:294}
\textsc{H. Araki}, \emph{A lattice of von Neumann algebras associated with the quantum theory of a free Bose field}
(J. Math. Phys., t. 4, p. 1343--1362).

\bibitem[295]{article:295}\phantomsection\label{bib:article:295}
\textsc{H. Araki et E. J. Woods}, \emph{Representations of the canonical commutation relations describing a nonrelativistic infinite free Bose gas}
(J. Math. Phys., t. 4, p. 637--662).

\bibitem[296]{article:296}\phantomsection\label{bib:article:296}
\textsc{D. J. C. Bures}, \emph{Certain factors constructed as infinite tensor product}
(Compos. Math., t. 15, p. 169--191).

\bibitem[297]{article:297}\phantomsection\label{bib:article:297}
\textsc{H. Choda et M. Choda}, \emph{On theorems of Korovkin}
(Proc. Japan Acad., t. 39, p. 107--108).

\bibitem[298]{article:298}\phantomsection\label{bib:article:298}
\textsc{H. Choda et M. Echigo}, \emph{A new algebraical property of certain von Neumann algebras}
(Proc. Japan Acad., t. 39, p. 651--655).
% Source PDF page 358 (printed p. 349)
\phantomsection\label{srcpage:349}
\bibitem[299]{article:299}\phantomsection\label{bib:article:299}
\textsc{Don Deckard and C. Pearcy}, \emph{On matrices over the ring of continuous complex-valued functions on a Stonian space}
(Proc. Amer. Math. Soc., t. 14, p. 322--328).

\bibitem[300]{article:300}\phantomsection\label{bib:article:300}
\textsc{H. A. Dye}, \emph{On groups of measure preserving transformations, II}
(Amer. J. Math., t. 85, p. 551--576).

\bibitem[301]{article:301}\phantomsection\label{bib:article:301}
\textsc{E. G. Effros}, \emph{Order ideals in a \(C^*\)-algebra}
(Duke Math. J., t. 30, p. 391--412).

\bibitem[302]{article:302}\phantomsection\label{bib:article:302}
\textsc{M. Guenin and B. Misra}, \emph{On the von Neumann algebras generated by the field operators}
(Nuovo Cimento, t. 30, p. 1272--1290).

\bibitem[303]{article:303}\phantomsection\label{bib:article:303}
\textsc{R. V. Kadison}, \emph{The trace in finite operator algebras}
(mimeographed, Columbia Univ.).

\bibitem[304]{article:304}\phantomsection\label{bib:article:304}
\textsc{R. V. Kadison}, \emph{Remarks on the type of von Neumann algebras of local observables in quantum field theory}
(J. Math. Phys., t. 4, p. 1511--1516).

\bibitem[305]{article:305}\phantomsection\label{bib:article:305}
\textsc{I. Kovacs}, \emph{Ergodic theorems for gages}
(Acta Univ. Szeged, t. 24, p. 103--118).

\bibitem[306]{article:306}\phantomsection\label{bib:article:306}
\textsc{F. Niiro}, \emph{Sur l'unicit\'e de la d\'ecomposition d'une trace}
(Sci. papers college gen. ed. Univ. Tokyo, t. 13, p. 159--162).

\bibitem[307]{article:307}\phantomsection\label{bib:article:307}
\textsc{T. Ono}, \emph{Local theory in function analysis}
(J. Math. Soc. Japan, t. 15, p. 9--30).

\bibitem[308]{article:308}\phantomsection\label{bib:article:308}
\textsc{C. Pearcy}, \emph{On unitary equivalence of matrices over the ring of continuous complex-valued functions on a Stonian space}
(Can. J. Math., t. 15, p. 323--331).

\bibitem[309]{article:309}\phantomsection\label{bib:article:309}
\textsc{J. Schwartz}, \emph{Two finite, non-hyperfinite, non-isomorphic factors}
(Comm. pure appl. math., t. 16, p. 19--26).

\bibitem[310]{article:310}\phantomsection\label{bib:article:310}
\textsc{J. Schwartz}, \emph{Non-isomorphism of a pair of factors of type III}
(Comm. pure appl. math., t. 16, p. 111--120).

\bibitem[311]{article:311}\phantomsection\label{bib:article:311}
\textsc{J. Schwartz}, \emph{Type II factors in a central decomposition}
(Comm. pure appl. math., t. 16, p. 247--252).

\bibitem[312]{article:312}\phantomsection\label{bib:article:312}
\textsc{N. Suzuki}, \emph{Isometries on Hilbert spaces}
(Proc. Japan Acad., t. 39, p. 435--438).

\bibitem[313]{article:313}\phantomsection\label{bib:article:313}
\textsc{N. Suzuki and T. Saito}, \emph{On the operators which generate continuous von Neumann algebras}
(T\^ohoku Math. J., t. 15, p. 277--280).

\bibitem[314]{article:314}\phantomsection\label{bib:article:314}
\textsc{K. Tsuji}, \emph{Annihilators of von Neumann algebras (annihilating spaces)}
(Bull. Kyushu inst. tech., t. 10, p. 25--39 and t. 13, p. 19).

\begin{center}\bfseries 1964.\end{center}

\bibitem[315]{article:315}\phantomsection\label{bib:article:315}
\textsc{H. Araki}, \emph{Von Neumann algebras of local observables for free scalar field}
(J. Math. Phys., t. 5, p. 1--13).

\bibitem[316]{article:316}\phantomsection\label{bib:article:316}
\textsc{H. Araki}, \emph{On the algebra of all local observables}
(Progr. theoret. phys., t. 32, p. 844--854).

\bibitem[317]{article:317}\phantomsection\label{bib:article:317}
\textsc{H. Araki}, \emph{Type of von Neumann algebra associated with free field}
(Progr. theoret. phys., t. 32, p. 956--965).

\bibitem[318]{article:318}\phantomsection\label{bib:article:318}
\textsc{H. Araki and W. Wyss}, \emph{Representations of canonical anticommutation relations}
(Helv. Phys. Acta, t. 37, p. 136--159).

\bibitem[319]{article:319}\phantomsection\label{bib:article:319}
\textsc{H. Choda and M. Echigo}, \emph{Some remarks on von Neumann algebras with the property Q}
(Mem. Osaka Gakugei Univ., no.~13, p. 13--21).

\bibitem[320]{article:320}\phantomsection\label{bib:article:320}
\textsc{H. Choda and M. Echigo}, \emph{A remark on construction of finite factors, I and II}
(Proc. Japan Acad., t. 40, p. 474--478 and 479--481).

\bibitem[321]{article:321}\phantomsection\label{bib:article:321}
\textsc{Don Deckard and C. Pearcy}, \emph{On continuous matrix-valued functions on a Stonian space}
(Pacific J. Math., t. 14, p. 857--869).

\bibitem[322]{article:322}\phantomsection\label{bib:article:322}
\textsc{P. Miles}, \emph{\(B^*\)-algebra unit ball extremal points}
(Pacific J. Math., t. 14, p. 627--637).

\bibitem[323]{article:323}\phantomsection\label{bib:article:323}
\textsc{C. Pearcy}, \emph{On certain von Neumann algebras which are generated by partial isometries}
(Proc. Amer. Math. Soc., t. 15, p. 393--395).

\bibitem[324]{article:324}\phantomsection\label{bib:article:324}
\textsc{C. Pearcy}, \emph{Entire functions on infinite von Neumann algebras of type I}
(Michigan Math. J., t. 11, p. 1--7).

\bibitem[325]{article:325}\phantomsection\label{bib:article:325}
\textsc{S. Sakai}, \emph{Weakly compact operators on operator algebras}
(Pacific J. Math., t. 14, p. 659--664).

\bibitem[326]{article:326}\phantomsection\label{bib:article:326}
\textsc{S. Sakai}, \emph{On the reduction theory of von Neumann}
(Bull. Amer. Math. Soc., t. 70, p. 393--398).

% Source PDF page 359 (printed p. 350)
\phantomsection\label{srcpage:350}
\bibitem[327]{article:327}\phantomsection\label{bib:article:327}
\textsc{R. F. Streater}, \emph{Intensive observables in quantum theory}
(J. Math. Phys., t. 5, p. 581--590).

\bibitem[328]{article:328}\phantomsection\label{bib:article:328}
\textsc{N. Suzuki}, \emph{On the type of completely continuous operators}
(Proc. Japan Acad., t. 40, p. 683--685).

\bibitem[329]{article:329}\phantomsection\label{bib:article:329}
\textsc{J. L. Taylor}, \emph{The Tomita decomposition of rings of operators}
(Trans. Amer. Math. Soc., t. 113, p. 30--39).

\bibitem[330]{article:330}\phantomsection\label{bib:article:330}
\textsc{A. S. Wightman}, \emph{La th\'eorie quantique locale et la th\'eorie quantique des champs}
(Ann. Inst. H. Poincar\'e, t. 1, p. 403--420).

\begin{center}\bfseries 1965.\end{center}

\bibitem[331]{article:331}\phantomsection\label{bib:article:331}
\textsc{S. Anastasio}, \emph{Maximal abelian subalgebras in hyperfinite factors}
(Amer. J. Math., t. 87, p. 955--971).

\bibitem[332]{article:332}\phantomsection\label{bib:article:332}
\textsc{H. J. Borchers}, \emph{On the structure of the algebra of field operators, II}
(Comm. Math. Phys., t. 1, p. 49--56).

\bibitem[333]{article:333}\phantomsection\label{bib:article:333}
\textsc{H. J. Borchers}, \emph{On the vacuum state in quantum field theory, II}
(Comm. Math. Phys., t. 1, p. 57--79).

\bibitem[334]{article:334}\phantomsection\label{bib:article:334}
\textsc{A. Brown and C. Pearcy}, \emph{Structure of commutators of operators}
(Ann. Math., t. 82, p. 112--127).

\bibitem[335]{article:335}\phantomsection\label{bib:article:335}
\textsc{H. Choda}, \emph{On automorphisms of abelian von Neumann algebras}
(Proc. Japan Acad., t. 41, p. 280--283).

\bibitem[336]{article:336}\phantomsection\label{bib:article:336}
\textsc{H. Choda and M. Choda}, \emph{Some remarks on von Neumann algebras with an algebraic property}
(Proc. Japan Acad., t. 41, p. 71--72).

\bibitem[337]{article:337}\phantomsection\label{bib:article:337}
\textsc{M. Choda}, \emph{On the conditional expectation of a partial isometry in a certain von Neumann algebra}
(Proc. Japan Acad., t. 41, p. 277--279).

\bibitem[338]{article:338}\phantomsection\label{bib:article:338}
\textsc{E. G. Effros}, \emph{The Borel space of von Neumann algebras on a separable Hilbert space}
(Pacific J. Math., t. 15, p. 1153--1164).

\bibitem[339]{article:339}\phantomsection\label{bib:article:339}
\textsc{J. Langerholc and B. Schroer}, \emph{On the structure of the von Neumann algebras generated by local functions of the free Bose field}
(Comm. Math. Phys., t. 1, p. 215--239).

\bibitem[340]{article:340}\phantomsection\label{bib:article:340}
\textsc{C. Pearcy}, \emph{On commutators of operators on Hilbert space}
(Proc. Amer. Math. Soc., t. 16, p. 53--59).

\bibitem[341]{article:341}\phantomsection\label{bib:article:341}
\textsc{G. A. Reid}, \emph{A generalisation of \(W^*\)-algebras}
(Pacific J. Math., t. 15, p. 1019--1026).

\bibitem[342]{article:342}\phantomsection\label{bib:article:342}
\textsc{T. Saito}, \emph{Some remarks on a representation of a group, I--II}
(T\^ohoku Math. J., t. 17, p. 206--209).

\bibitem[343]{article:343}\phantomsection\label{bib:article:343}
\textsc{S. Sakai}, \emph{A Radon-Nikodym theorem in \(W^*\)-algebras}
(Bull. Amer. Math. Soc., t. 71, p. 149--151).

\bibitem[344]{article:344}\phantomsection\label{bib:article:344}
\textsc{S. Sakai}, \emph{On topologies of finite \(W^*\)-algebras}
(Ill. J. Math., t. 9, p. 236--241).

\bibitem[345]{article:345}\phantomsection\label{bib:article:345}
\textsc{I. E. Segal}, \emph{Algebraic integration theory}
(Bull. Amer. Math. Soc., t. 71, p. 419--489).

\bibitem[346]{article:346}\phantomsection\label{bib:article:346}
\textsc{E. St{\o}rmer}, \emph{On the Jordan structure of \(C^*\)-algebras}
(Trans. Amer. Math. Soc., t. 120, p. 438--447).

\bibitem[347]{article:347}\phantomsection\label{bib:article:347}
\textsc{N. Suzuki}, \emph{Algebraic aspects of non self-adjoint operators}
(Proc. Japan Acad., t. 41, p. 706--710).

\bibitem[348]{article:348}\phantomsection\label{bib:article:348}
\textsc{R. J. Tauer}, \emph{Maximal abelian subalgebras in finite factors of type II}
(Trans. Amer. Math. Soc., t. 114, p. 281--308).

\bibitem[349]{article:349}\phantomsection\label{bib:article:349}
\textsc{R. J. Tauer}, \emph{Semi-regular maximal abelian subalgebras in hyperfinite factors}
(Bull. Amer. Math. Soc., t. 71, p. 606--608).

\bibitem[350]{article:350}\phantomsection\label{bib:article:350}
\textsc{D. M. Topping}, \emph{Jordan algebras of self-adjoint operators}
(Bull. Amer. Math. Soc., t. 71, p. 160--164).

\bibitem[351]{article:351}\phantomsection\label{bib:article:351}
\textsc{D. M. Topping}, \emph{Jordan algebras of self-adjoint operators}
(Mem. Amer. Math. Soc., no.~53).

% Source PDF page 360 (printed p. 351)
\phantomsection\label{srcpage:351}
\begin{center}\bfseries 1966.\end{center}

\bibitem[352]{article:352}\phantomsection\label{bib:article:352}
\textsc{J. Aarnes}, \emph{The Vitali-Hahn-Saks theorem for von Neumann algebras}
(Math. Scand., t. 18, p. 87--92).

\bibitem[353]{article:353}\phantomsection\label{bib:article:353}
\textsc{H. Araki and E. J. Woods}, \emph{Complete boolean algebras of type I factors}
(Pub. res. inst. math. sci. Kyoto Univ., t. 2, p. 157--242 and 451--452).

\bibitem[354]{article:354}\phantomsection\label{bib:article:354}
\textsc{M. Broise}, \emph{Sur certaines applications unitaires de l'espace des op\'erateurs de Hilbert-Schmidt}
(C. R. Acad. Sc., t. 263, p. 722--725).

\bibitem[355]{article:355}\phantomsection\label{bib:article:355}
\textsc{M. Broise}, \emph{Sur les isomorphismes de certaines alg\`ebres de von Neumann}
(Ann. scient. \`Ec. Norm. Sup., t. 83, p. 91--111).

\bibitem[356]{article:356}\phantomsection\label{bib:article:356}
\textsc{A. Brown and C. Pearcy}, \emph{Multiplicative commutators of operators}
(Can. J. Math., t. 18, p. 737--749).

\bibitem[357]{article:357}\phantomsection\label{bib:article:357}
\textsc{A. Brown and C. Pearcy}, \emph{Commutators in factors of type III}
(Can. J. Math., t. 18, p. 1152--1160).

\bibitem[358]{article:358}\phantomsection\label{bib:article:358}
\textsc{A. Cordesse and G. Rideau}, \emph{On some representations of anticommutation relations, I--II}
(Nuovo Cimento, t. 45, p. 1--14 and t. 46, p. 624--636).

\bibitem[359]{article:359}\phantomsection\label{bib:article:359}
\textsc{J. Dixmier}, \emph{Existence de traces non normales}
(C. R. Acad. Sc., t. 262, p. 1107--1108).

\bibitem[360]{article:360}\phantomsection\label{bib:article:360}
\textsc{E. G. Effros}, \emph{Global structure in von Neumann algebras}
(Trans. Amer. Math. Soc., t. 121, p. 434--454).

\bibitem[361]{article:361}\phantomsection\label{bib:article:361}
\textsc{P. A. Fillmore}, \emph{On products of symmetries}
(Can. J. Math., t. 18, p. 897--900).

\bibitem[362]{article:362}\phantomsection\label{bib:article:362}
\textsc{E. L. Griffin}, \emph{Everywhere defined linear transformations affiliated with rings of operators}
(Pacific J. Math., t. 18, p. 489--493).

\bibitem[363]{article:363}\phantomsection\label{bib:article:363}
\textsc{A. Guichardet}, \emph{Sur la cat\'egorie des alg\`ebres de von Neumann}
(Bull. Sc. math., t. 90, p. 41--63).

\bibitem[364]{article:364}\phantomsection\label{bib:article:364}
\textsc{H. Halpern}, \emph{The maximal GCR ideal in an \(AW^*\)-algebra}
(Proc. Amer. Math. Soc., t. 17, p. 906--914).

\bibitem[365]{article:365}\phantomsection\label{bib:article:365}
\textsc{H. Halpern}, \emph{An integral representation of a normal functional on a von Neumann algebra}
(Trans. Amer. Math. Soc., t. 125, p. 32--46).

\bibitem[366]{article:366}\phantomsection\label{bib:article:366}
\textsc{R. T. Harris}, \emph{A direct integral construction}
(Duke Math. J., t. 33, p. 535--537).

\bibitem[367]{article:367}\phantomsection\label{bib:article:367}
\textsc{R. V. Kadison}, \emph{Derivation of operator algebras}
(Ann. Math., t. 83, p. 280--293).

\bibitem[368]{article:368}\phantomsection\label{bib:article:368}
\textsc{R. V. Kadison and J. R. Ringrose}, \emph{Derivations of operator group algebras}
(Amer. J. Math., t. 88, p. 562--576).

\bibitem[369]{article:369}\phantomsection\label{bib:article:369}
\textsc{R. V. Kadison and J. R. Ringrose}, \emph{Automorphisms of operator algebras}
(Bull. Amer. Math. Soc., t. 72, p. 1059--1063).

\bibitem[370]{article:370}\phantomsection\label{bib:article:370}
\textsc{I. Kovacs and J. Sz\H{u}cs}, \emph{Th\'eor\`emes de type ergodique dans les alg\`ebres de von Neumann}
(C. R. Acad. Sc., t. 262, p. 341--344).

\bibitem[371]{article:371}\phantomsection\label{bib:article:371}
\textsc{I. Kovacs and J. Sz\H{u}cs}, \emph{Ergodic type theorems in von Neumann algebras}
(Acta sci. Math. Szeged, t. 27, p. 233--246).

\bibitem[372]{article:372}\phantomsection\label{bib:article:372}
\textsc{K. Matoba}, \emph{A remark on a theorem of H. Araki}
(Proc. Japan Acad., t. 42, p. 344--346).

\bibitem[373]{article:373}\phantomsection\label{bib:article:373}
\textsc{A. R. Padmanabhan}, \emph{Some dominated convergence theorems in a von Neumann algebra}
(Proc. Japan Acad., t. 42, p. 347--350).

\bibitem[374]{article:374}\phantomsection\label{bib:article:374}
\textsc{A. R. Padmanabhan}, \emph{Stability and mixing in von Neumann algebras}
(Kodai Math. sem. rep., t. 18, p. 335--342).

\bibitem[375]{article:375}\phantomsection\label{bib:article:375}
\textsc{E. A. Pedersen}, \emph{A decomposition theory for rings of operators}
(Thesis, Louisiana State Univ.).

\bibitem[376]{article:376}\phantomsection\label{bib:article:376}
\textsc{S. Sakai}, \emph{Derivations of \(W^*\)-algebras}
(Ann. Math., t. 83, p. 273--279).

\bibitem[377]{article:377}\phantomsection\label{bib:article:377}
\textsc{E. St{\o}rmer}, \emph{Jordan algebras of type I}
(Acta Math., t. 115, p. 165--184).

\bibitem[378]{article:378}\phantomsection\label{bib:article:378}
\textsc{N. Suzuki}, \emph{The algebraic structure of non-selfadjoint operators}
(Acta sci. Math. Szeged, t. 28, p. 173--184).

\bibitem[379]{article:379}\phantomsection\label{bib:article:379}
\textsc{N. Suzuki}, \emph{On the spectral decomposition of dissipative operators}
(Proc. Japan Acad., t. 42, p. 577--582).

% Source PDF page 361 (printed p. 352)
\phantomsection\label{srcpage:352}
\bibitem[380]{article:380}\phantomsection\label{bib:article:380}
\textsc{D. M. Topping}, \emph{An isomorphism invariant for spin factors}
(J. Math. Mech., t. 15, p. 1055--1064).

\bibitem[381]{article:381}\phantomsection\label{bib:article:381}
\textsc{K. Tsuji}, \emph{The J. von Neumann commutation theorem}
(Bull. Kyushu inst. tech., t. 13, p. 1--3).

\bibitem[382]{article:382}\phantomsection\label{bib:article:382}
\textsc{K. Tsuji}, \emph{Tensor product of annihilating spaces}
(Bull. Kyushu inst. tech., t. 13, p. 15--18).

\begin{center}\bfseries 1967.\end{center}

\bibitem[383]{article:383}\phantomsection\label{bib:article:383}
\textsc{C. A. Akemann}, \emph{The dual space of an operator algebra}
(Trans. Amer. Math. Soc., t. 126, p. 286--302).

\bibitem[384]{article:384}\phantomsection\label{bib:article:384}
\textsc{C. A. Akemann}, \emph{Projections onto separable \(C^*\)-subalgebras of a \(W^*\)-algebra}
(Bull. Amer. Math. Soc., t. 73, p. 925).

\bibitem[385]{article:385}\phantomsection\label{bib:article:385}
\textsc{W. B. Arveson}, \emph{Analyticity in operator algebras}
(Amer. J. Math., t. 89, p. 578--642).

\bibitem[386]{article:386}\phantomsection\label{bib:article:386}
\textsc{H. J. Borchers}, \emph{A remark on a theorem of B. Misra}
(Comm. Math. Phys., t. 4, p. 315--323).

\bibitem[387]{article:387}\phantomsection\label{bib:article:387}
\textsc{M. Breuer}, \emph{A generalization of Kuiper's theorem to factors of type \(\mathrm{II}_\infty\)}
(J. Math. Mech., t. 16, p. 917--925).

\bibitem[388]{article:388}\phantomsection\label{bib:article:388}
\textsc{M. Broise}, \emph{Sur les vecteurs s\'eparateurs des alg\`ebres de von Neumann}
(C. R. Acad. Sc., t. 264, p. 937--940).

\bibitem[389]{article:389}\phantomsection\label{bib:article:389}
\textsc{M. Broise}, \emph{Sur les vecteurs s\'eparateurs des alg\`ebres de von Neumann}
(J. functional anal., t. 1, p. 281--289).

\bibitem[390]{article:390}\phantomsection\label{bib:article:390}
\textsc{M. Broise}, \emph{Commutateurs dans le groupe unitaire d'un facteur}
(J. Math. pures et appl., t. 46, p. 299--312).

\bibitem[391]{article:391}\phantomsection\label{bib:article:391}
\textsc{H. Choda}, \emph{On Maharam subfactors of finite factors, I--II}
(Proc. Japan Acad., t. 43, p. 451--455 and 937--940).

\bibitem[392]{article:392}\phantomsection\label{bib:article:392}
\textsc{H. Choda}, \emph{On the crossed product of abelian von Neumann algebras, I--II}
(Proc. Japan Acad., t. 43, p. 111--116 and 198--201).

\bibitem[393]{article:393}\phantomsection\label{bib:article:393}
\textsc{M. Choda and H. Choda}, \emph{On extensions of automorphisms of abelian von Neumann algebras}
(Proc. Japan Acad., t. 43, p. 295--299).

\bibitem[394]{article:394}\phantomsection\label{bib:article:394}
\textsc{A. Cordesse and G. Rideau}, \emph{On some representations of anticommutation relations, III}
(Nuovo Cimento, t. 50, p. 244--255).

\bibitem[395]{article:395}\phantomsection\label{bib:article:395}
\textsc{G. F. Dell'Antonio}, \emph{On the limits of sequences of normal states}
(Comm. pure appl. math., t. 20, p. 413--429).

\bibitem[396]{article:396}\phantomsection\label{bib:article:396}
\textsc{R. G. Douglas and D. M. Topping}, \emph{Operators whose squares are zero}
(Revue roumaine Math. pures et appl., t. 12, p. 647--652).

\bibitem[397]{article:397}\phantomsection\label{bib:article:397}
\textsc{P. A. Fillmore and D. M. Topping}, \emph{Operator algebras generated by projections}
(Duke math. J., t. 34, p. 333--336).

\bibitem[398]{article:398}\phantomsection\label{bib:article:398}
\textsc{D. J. H. Garling}, \emph{On ideals of operators in Hilbert space}
(Proc. London Math. Soc., t. 17, p. 115--138).

\bibitem[399]{article:399}\phantomsection\label{bib:article:399}
\textsc{J. Hakeda}, \emph{On property P of von Neumann algebras}
(T\^ohoku Math. J., t. 19, p. 238--242).

\bibitem[400]{article:400}\phantomsection\label{bib:article:400}
\textsc{N. M. Hugenholtz}, \emph{On the factor type of equilibrium states in quantum statistical mechanics}
(Comm. Math. Phys., t. 6, p. 189--193).

\bibitem[401]{article:401}\phantomsection\label{bib:article:401}
\textsc{B. E. Johnson}, \emph{\(AW^*\)-algebras are \(QW^*\)-algebras}
(Pacific J. Math., t. 23, p. 97--99).

\bibitem[402]{article:402}\phantomsection\label{bib:article:402}
\textsc{B. E. Johnson}, \emph{Continuity of homomorphisms of algebras of operators}
(J. London Math. Soc., t. 42, p. 537--541).

\bibitem[403]{article:403}\phantomsection\label{bib:article:403}
\textsc{R. V. Kadison}, \emph{Lectures on operator algebras}
(Carg\`ese lectures in theoretical physics, Gordon and Breach, New York).

\bibitem[404]{article:404}\phantomsection\label{bib:article:404}
\textsc{R. V. Kadison and J. R. Ringrose}, \emph{Derivations and automorphisms of operator algebras}
(Comm. Math. Phys., t. 4, p. 32--63).

\bibitem[405]{article:405}\phantomsection\label{bib:article:405}
\textsc{R. V. Kadison, E. C. Lance and J. R. Ringrose}, \emph{Derivations and automorphisms of operator algebras, II}
(J. functional anal., t. 1, p. 204--221).

% Source PDF page 362 (printed p. 353)
\phantomsection\label{srcpage:353}
\bibitem[406]{article:406}\phantomsection\label{bib:article:406}
\textsc{C. C. Moore}, \emph{Invariant measures on product spaces}
(Proceedings of the 5th Berkeley Symposium on mathematical statistics and probability, vol. II, part II, p. 447--459).

\bibitem[407]{article:407}\phantomsection\label{bib:article:407}
\textsc{A. R. Padmanabhan}, \emph{Convergence in measure and related results in finite rings of operators}
(Trans. Amer. Math. Soc., t. 128, p. 359--378 and t. 132, p. 563).

\bibitem[408]{article:408}\phantomsection\label{bib:article:408}
\textsc{C. Pearcy and D. M. Topping}, \emph{Sums of small number of idempotents}
(Michigan Math. J., t. 14, p. 453--465).

\bibitem[409]{article:409}\phantomsection\label{bib:article:409}
\textsc{P. Porcelli and E. A. Pedersen}, \emph{On rings of operators}
(Bull. Amer. Math. Soc., t. 73, p. 142--144).

\bibitem[410]{article:410}\phantomsection\label{bib:article:410}
\textsc{R. T. Powers}, \emph{Representations of uniformly hyperfinite algebras and their associated von Neumann rings}
(Bull. Amer. Math. Soc., t. 73, p. 572--575).

\bibitem[411]{article:411}\phantomsection\label{bib:article:411}
\textsc{R. T. Powers}, \emph{Representations of uniformly hyperfinite algebras and their associated von Neumann rings}
(Ann. Math., t. 86, p. 138--171).

\bibitem[412]{article:412}\phantomsection\label{bib:article:412}
\textsc{K. Saito}, \emph{On the preduals of \(W^*\)-algebras}
(T\^ohoku Math. J., t. 19, p. 324--331).

\bibitem[413]{article:413}\phantomsection\label{bib:article:413}
\textsc{K. Saito}, \emph{Non-commutative extensions of Lusin's theorem}
(T\^ohoku Math. J., t. 19, p. 332--340).

\bibitem[414]{article:414}\phantomsection\label{bib:article:414}
\textsc{J. T. Schwartz}, \emph{\(W^*\)-algebras}
(Notes on mathematics and its applications, Gordon and Breach, New York).

\bibitem[415]{article:415}\phantomsection\label{bib:article:415}
\textsc{E. St{\o}rmer}, \emph{On anti-automorphisms of von Neumann algebras}
(Pacific J. Math., t. 21, p. 349--370).

\bibitem[416]{article:416}\phantomsection\label{bib:article:416}
\textsc{E. St{\o}rmer}, \emph{Types of von Neumann algebras associated with extremal invariant states}
(Comm. Math. Phys., t. 6, p. 194--204).

\bibitem[417]{article:417}\phantomsection\label{bib:article:417}
\textsc{R. T. Tauer}, \emph{M-semiregular subalgebras in hyperfinite factors}
(Trans. Amer. Math. Soc., t. 129, p. 530--541).

\bibitem[418]{article:418}\phantomsection\label{bib:article:418}
\textsc{D. M. Topping}, \emph{Asymptoticity and semimodularity in projection lattices}
(Pacific J. Math., t. 20, p. 317--325).

\begin{center}\bfseries 1968.\end{center}

\bibitem[419]{article:419}\phantomsection\label{bib:article:419}
\textsc{C. A. Akemann}, \emph{Sequential convergence in the dual of a \(W^*\)-algebra}
(Comm. Math. Phys., t. 7, p. 222--224).

\bibitem[420]{article:420}\phantomsection\label{bib:article:420}
\textsc{C. A. Akemann}, \emph{Interpolation in \(W^*\)-algebras}
(Duke Math. J., t. 35, p. 525--533).

\bibitem[421]{article:421}\phantomsection\label{bib:article:421}
\textsc{S. Anastasio}, \emph{Non-normal abelian subalgebras}
(T\^ohoku Math. J., t. 20, p. 26--37).

\bibitem[422]{article:422}\phantomsection\label{bib:article:422}
\textsc{G. F. Dell'Antonio}, \emph{Structure of the algebras of some free systems}
(Comm. Math. Phys., t. 9, p. 81--117).

\bibitem[423]{article:423}\phantomsection\label{bib:article:423}
\textsc{H. Halpern}, \emph{A spectral decomposition for self-adjoint elements in the maximum GCR ideal of a von Neumann algebra with applications to noncommutative integration theory}
(Trans. Amer. Math. Soc., t. 133, p. 281--306).

\bibitem[424]{article:424}\phantomsection\label{bib:article:424}
\textsc{O. Mar\'echal}, \emph{Op\'erateurs d\'ecomposables dans les champs mesurables d'espaces de Hilbert}
(C. R. Acad. Sc., t. 266, s\'erie A, p. 710--713).

\bibitem[425]{article:425}\phantomsection\label{bib:article:425}
\textsc{O. Mar\'echal}, \emph{Op\'erateurs d\'ecomposables dans les espaces \(L_E^p\)}
(C. R. Acad. Sc., t. 266, s\'erie A, p. 1043--1045).

\bibitem[426]{article:426}\phantomsection\label{bib:article:426}
\textsc{T. Okayasu}, \emph{A structure theorem of automorphisms of von Neumann algebras}
(T\^ohoku Math. J., t. 20, p. 199--206).

\bibitem[427]{article:427}\phantomsection\label{bib:article:427}
\textsc{C. Pearcy and J. R. Ringrose}, \emph{Trace-preserving isomorphisms in finite operator algebras}
(Amer. J. Math., t. 90, p. 444--455).

\bibitem[428]{article:428}\phantomsection\label{bib:article:428}
\textsc{B. Russo}, \emph{Unimodular contractions in Hilbert space}
(Pac. J. Math., t. 26, p. 163--169).

\bibitem[429]{article:429}\phantomsection\label{bib:article:429}
\textsc{B. Russo}, \emph{Isometries of \(L^p\)-spaces associated with finite von Neumann algebras}
(Bull. Amer. Math. Soc., t. 74, p. 228--232).

\bibitem[430]{article:430}\phantomsection\label{bib:article:430}
\textsc{T. Saito}, \emph{Generators of certain von Neumann algebras}
(T\^ohoku Math. J., t. 20, p. 101--105).

\bibitem[431]{article:431}\phantomsection\label{bib:article:431}
\textsc{T. Saito}, \emph{On generators of von Neumann algebras}
(Michigan Math. J., t. 15, p. 1--4).

% Source PDF page 363 (printed p. 354)
\phantomsection\label{srcpage:354}
\bibitem[432]{article:432}\phantomsection\label{bib:article:432}
\textsc{S. Sakai}, \emph{On the tensor product of \(W^*\)-algebras}
(Amer. J. Math., t. 90, p. 335--341).

\bibitem[433]{article:433}\phantomsection\label{bib:article:433}
\textsc{E. St{\o}rmer}, \emph{Irreducible Jordan algebras of self-adjoint operators}
(Trans. Amer. Math. Soc., t. 130, p. 153--166).

\bibitem[434]{article:434}\phantomsection\label{bib:article:434}
\textsc{W. I. M. Wils}, \emph{Stone-\v Cech compactification and representations of operator algebras}
(Thesis, Nijmegen).

\begin{center}\bfseries Forthcoming.\end{center}

\bibitem[435]{article:435}\phantomsection\label{bib:article:435}
\textsc{H. Araki and E. J. Woods}, \emph{A classification of factors}.

\bibitem[436]{article:436}\phantomsection\label{bib:article:436}
\textsc{W. B. Arveson}, \emph{A density theorem for operator algebras}.

\bibitem[437]{article:437}\phantomsection\label{bib:article:437}
\textsc{S. K. Berberian}, \emph{Trace and the convex hull of the spectrum in a von Neumann algebra of finite class}.

\bibitem[438]{article:438}\phantomsection\label{bib:article:438}
\textsc{H. J. Borchers}, \emph{On the converse of the Reeh-Schlieder theorem}.

\bibitem[439]{article:439}\phantomsection\label{bib:article:439}
\textsc{M. Breuer}, \emph{Fredholm theories in von Neumann algebras, I--II}.

\bibitem[440]{article:440}\phantomsection\label{bib:article:440}
\textsc{A. Brown, C. Pearcy and D. Topping}, \emph{Commutators and the strong radical}.

\bibitem[441]{article:441}\phantomsection\label{bib:article:441}
\textsc{Wai-Mee Ching}, \emph{Non-isomorphic non-hyperfinite factors}.

\bibitem[442]{article:442}\phantomsection\label{bib:article:442}
\textsc{J. Dixmier}, \emph{Quelques propri\'et\'es des suites centrales dans les facteurs de type \(\mathrm{II}_1\)}.

\bibitem[443]{article:443}\phantomsection\label{bib:article:443}
\textsc{J. Dixmier and E. C. Lance}, \emph{Deux nouveaux facteurs de type \(\mathrm{II}_1\)}.

\bibitem[444]{article:444}\phantomsection\label{bib:article:444}
\textsc{S. Doplicher, D. Kastler and E. St{\o}rmer}, \emph{Invariant states and asymptotic abelianness}.

\bibitem[445]{article:445}\phantomsection\label{bib:article:445}
\textsc{R. G. Douglas}, \emph{On operators similar to normal operators}.

\bibitem[446]{article:446}\phantomsection\label{bib:article:446}
\textsc{R. G. Douglas and C. Pearcy}, \emph{Von Neumann algebras with a single generator}.

\bibitem[447]{article:447}\phantomsection\label{bib:article:447}
\textsc{J. Ernest}, \emph{Hopf-von Neumann algebras}.

\bibitem[448]{article:448}\phantomsection\label{bib:article:448}
\textsc{P. A. Fillmore and D. M. Topping}, \emph{A direct integral decomposition for certain operator algebras}.

\bibitem[449]{article:449}\phantomsection\label{bib:article:449}
\textsc{P. A. Fillmore and D. M. Topping}, \emph{Sums of irreducible operators}.

\bibitem[450]{article:450}\phantomsection\label{bib:article:450}
\textsc{M. Flensted-Jensen}, \emph{A note on desintegration, type and global type of von Neumann algebras}.

\bibitem[451]{article:451}\phantomsection\label{bib:article:451}
\textsc{N. M. Hugenholtz and J. D. Wieringa}, \emph{On locally normal states in quantum statistical mechanics}.

\bibitem[452]{article:452}\phantomsection\label{bib:article:452}
\textsc{B. E. Johnson and J. R. Ringrose}, \emph{Derivations of operator algebras and discrete group algebras}.

\bibitem[453]{article:453}\phantomsection\label{bib:article:453}
\textsc{R. R. Kallman}, \emph{A generalization of free action}.

\bibitem[454]{article:454}\phantomsection\label{bib:article:454}
\textsc{R. R. Kallman}, \emph{Unitary groups and automorphisms of operator algebras}.

\bibitem[455]{article:455}\phantomsection\label{bib:article:455}
\textsc{A. R. Padmanabhan}, \emph{Probabilistic aspects of rings of operators}.

\bibitem[456]{article:456}\phantomsection\label{bib:article:456}
\textsc{A. R. Padmanabhan}, \emph{Some probabilistic implications of automorphisms of finite rings of operators}.

\bibitem[457]{article:457}\phantomsection\label{bib:article:457}
\textsc{C. Pearcy and D. M. Topping}, \emph{Commutators and certain \(\mathrm{II}_1\)-factors}.

\bibitem[458]{article:458}\phantomsection\label{bib:article:458}
\textsc{C. Pearcy and D. M. Topping}, \emph{Sums of small numbers of certain operators}.

\bibitem[459]{article:459}\phantomsection\label{bib:article:459}
\textsc{M. A. Rieffel}, \emph{Square integrable representations of Hilbert algebras}.

\bibitem[460]{article:460}\phantomsection\label{bib:article:460}
\textsc{T. Saito}, \emph{Generators of hyperfinite factors}.

\bibitem[461]{article:461}\phantomsection\label{bib:article:461}
\textsc{S. Sakai}, \emph{On global \(W^*\)-algebras}.

\bibitem[462]{article:462}\phantomsection\label{bib:article:462}
\textsc{S. Sakai}, \emph{On the hyperfinite \(\mathrm{II}_1\)-factor}.

\bibitem[463]{article:463}\phantomsection\label{bib:article:463}
\textsc{S. Sakai}, \emph{Asymptotically abelian type \(\mathrm{II}_1\)-factors}.

\bibitem[464]{article:464}\phantomsection\label{bib:article:464}
\textsc{R. F. Streater and A. Wulfsohn}, \emph{Continuous tensor product of Hilbert spaces}.

\bibitem[465]{article:465}\phantomsection\label{bib:article:465}
\textsc{E. St{\o}rmer}, \emph{States and invariant maps of operator algebras}.

\bibitem[466]{article:466}\phantomsection\label{bib:article:466}
\textsc{M. Takesaki}, \emph{Remarks on the reduction theory of von Neumann algebras}.

\bibitem[467]{article:467}\phantomsection\label{bib:article:467}
\textsc{M. Tomita}, \emph{Quasi-standard von Neumann algebras}.

\bibitem[468]{article:468}\phantomsection\label{bib:article:468}
\textsc{D. M. Topping}, \emph{On linear combinations of special operators}.

\bibitem[469]{article:469}\phantomsection\label{bib:article:469}
\textsc{J. Vesterstr{\o}m and W. Wils}, \emph{Direct integral of Hilbert spaces, II}.

\bibitem[470]{article:470}\phantomsection\label{bib:article:470}
\textsc{W. Wils}, \emph{Direct integrals of Hilbert spaces, I}.

\bibitem[471]{article:471}\phantomsection\label{bib:article:471}
\textsc{W. Wils}, \emph{Two-sided ideals in \(W^*\)-algebras}.

% Source PDF page 364 (printed p. 355)
\phantomsection\label{srcpage:355}
\bibitem[472]{article:472}\phantomsection\label{bib:article:472}
\textsc{W. Wils}, \emph{D\'esint\'egration centrale des formes positives sur les \(C^*\)-alg\`ebres}.

\bibitem[473]{article:473}\phantomsection\label{bib:article:473}
\textsc{W. Wogen}, \emph{On generators for von Neumann algebras}.

\bibitem[474]{article:474}\phantomsection\label{bib:article:474}
\textsc{G. Zeller-Meier}, \emph{Deux autres facteurs de type \(\mathrm{II}_1\)}.

\global\let\DixmierBackMatterChapter\chapter
\end{thebibliography}
\let\chapter\DixmierBackMatterChapter

\section*{General Treatises}\label{bib:general-treatises}
\addcontentsline{toc}{section}{General Treatises}
\begingroup
\makeatletter
\renewcommand{\@biblabel}[1]{[[#1]]}
\makeatother
\let\DixmierTreatiseBibliographyChapter\chapter
\RenewDocumentCommand{\chapter}{s o m}{}
\begin{thebibliography}{99}
\let\chapter\DixmierTreatiseBibliographyChapter

\bibitem{treatise:1}\phantomsection\label{bib:treatise:1}
\textsc{N. I. Ahiezer and I. M. Glazman}, \emph{La th\'eorie des op\'erateurs lin\'eaires dans l'espace hilbertien}
(Gosudarstv. Izdat. Techn.-Teor. Lit., Moscow-Leningrad, 1950) (in Russian). German translation (Akademie-Verlag, Berlin, 1954).

\bibitem{treatise:2}\phantomsection\label{bib:treatise:2}
\textsc{N. Aronszajn}, \emph{Introduction to the Theory of Hilbert Spaces}
(vol. I, Mathematical Monographs, Oklahoma Agricultural and Mechanical College, Stillwater, 1950).

\bibitem{treatise:3}\phantomsection\label{bib:treatise:3}
\textsc{N. Bourbaki}, \emph{Espaces vectoriels topologiques}
(Act. Sc. Ind., nos.~1189 and 1229, Paris, Hermann, 1966 and 1955).

\bibitem{treatise:4}\phantomsection\label{bib:treatise:4}
\textsc{N. Bourbaki}, \emph{Int\'egration}
(Act. Sc. Ind., nos.~1175, 1244 and 1281, Paris, Hermann, 1965, 1967 and 1959).

\bibitem{treatise:5}\phantomsection\label{bib:treatise:5}
\textsc{P. R. Halmos}, \emph{Introduction to Hilbert Space and the Theory of spectral Multiplicity}
(Chelsea Publishing Company, New York, 1951).

\bibitem{treatise:6}\phantomsection\label{bib:treatise:6}
\textsc{E. Hille and R. S. Phillips}, \emph{Functional Analysis and semi-groups}
(Amer. Math. Soc., Colloquium Publications, vol. 31, New York, 1957).

\bibitem{treatise:7}\phantomsection\label{bib:treatise:7}
\textsc{G. Julia}, \emph{Introduction math\'ematique aux th\'eories quantiques}
(Cahiers Scientifiques, fasc. 16, 2nd ed., Paris, Gauthier-Villars, 1953; fasc. 19, 2nd ed., Paris, Gauthier-Villars, 1955).

\bibitem{treatise:8}\phantomsection\label{bib:treatise:8}
\textsc{A. Lichnerowicz}, \emph{Alg\`ebre et Analyse lin\'eaires}
(Masson, Paris, 1947).

\bibitem{treatise:9}\phantomsection\label{bib:treatise:9}
\textsc{L. H. Loomis}, \emph{An introduction to abstract harmonic analysis}
(The University Series in higher Mathematics, New York, 1953).

\bibitem{treatise:10}\phantomsection\label{bib:treatise:10}
\textsc{L. A. Lyusternik and V. I. Sobolev}, \emph{Les \`el\'ements de l'analyse fonctionnelle}
(Gosudarstv. Izdat. Tehn.-Teor. Lit., Moscow-Leningrad, 1951) (in Russian).

\bibitem{treatise:11}\phantomsection\label{bib:treatise:11}
\textsc{J. von Neumann}, \emph{Functional operators. II. The geometry of orthogonal spaces}
(Ann. Math. Studies, no.~22, Princeton University Press, Princeton, 1950).

\bibitem{treatise:12}\phantomsection\label{bib:treatise:12}
\textsc{F. Riesz and B. Sz.-Nagy}, \emph{Le\c cons d'analyse fonctionnelle}, 3rd ed.
(Gauthier-Villars, Paris, and Acad. Sc. Hongrie, Budapest, 1955).

\bibitem{treatise:13}\phantomsection\label{bib:treatise:13}
\textsc{R. Schatten}, \emph{A theory of cross-spaces}
(Ann. Math. Studies, no.~26, Princeton University Press, Princeton, 1950).

\bibitem{treatise:14}\phantomsection\label{bib:treatise:14}
\textsc{M. H. Stone}, \emph{Linear transformations in Hilbert space}
(Amer. Math. Soc., Colloquium Publications, vol. 15, New York, 1932).

\bibitem{treatise:15}\phantomsection\label{bib:treatise:15}
\textsc{B. Sz.-Nagy}, \emph{Spektraldarstellung linearer transformationen des Hilbertschen Raumes}
(Ergebnisse der Mathematik und ihrer Grenzgebiete, t. 5, Berlin, Springer, 1942).

\end{thebibliography}
\let\chapter\DixmierTreatiseBibliographyChapter
\endgroup

% Source PDF page 366 (printed p. 357)
\chapter*{Index of Notation}\label{back:notation-index}\label{srcpage:357}
\addcontentsline{toc}{chapter}{Index of Notation}

The reference numbers indicate successively the chapter, the section, and the
number (or, exceptionally, the exercise).

\begingroup
\newcommand{\NotationEntry}[3]{\par\noindent\(#1\) : \hyperref[#2]{#3}.}
\setlength{\columnsep}{1em}
\begin{multicols}{2}
\raggedcolumns
\raggedright
\NotationEntry{\mathcal L(\mathfrak H)}{no:I-1-1}{I, 1, 1}
\NotationEntry{I_{\mathfrak H}}{no:I-1-1}{I, 1, 1}
\NotationEntry{\mathfrak M'}{no:I-1-1}{I, 1, 1}
\NotationEntry{\mathfrak M''}{no:I-1-1}{I, 1, 1}
\NotationEntry{\mathbb C_{\mathfrak H}}{no:I-1-1}{I, 1, 1}
\NotationEntry{\mathfrak M^+}{no:I-1-2}{I, 1, 2}
\NotationEntry{P_{\mathfrak X}}{no:I-1-4}{I, 1, 4}
\NotationEntry{E_{\mathfrak M}^{\mathcal A}}{no:I-1-4}{I, 1, 4}
\NotationEntry{\mathfrak X_{\mathfrak M}^{\mathcal A}}{no:I-1-4}{I, 1, 4}
\NotationEntry{E_x^{\mathcal A}}{no:I-1-4}{I, 1, 4}
\NotationEntry{\mathfrak X_x^{\mathcal A}}{no:I-1-4}{I, 1, 4}
\NotationEntry{\mathfrak m^2}{no:I-1-6}{I, 1, 6}
\NotationEntry{\mathfrak m^{1/2}}{no:I-1-6}{I, 1, 6}
\NotationEntry{\mathfrak m^0}{ex:I-1-6}{I, 1, Exercise 6}
\NotationEntry{\mathfrak m^\infty}{ex:I-1-6}{I, 1, Exercise 6}
\NotationEntry{T\mathrel{\eta}\mathcal A}{ex:I-1-10}{I, 1, Exercise 10}
\NotationEntry{T_{\mathfrak X}}{no:I-2-1}{I, 2, 1}
\NotationEntry{T_E}{no:I-2-1}{I, 2, 1}
\NotationEntry{\mathfrak M_E}{no:I-2-1}{I, 2, 1}
\NotationEntry{\mathcal A'_E}{no:I-2-1}{I, 2, 1}
\NotationEntry{\prod_{i\in I}\mathcal A_i}{no:I-2-2}{I, 2, 2}
\NotationEntry{\mathcal A_1\times\mathcal A_2\times\cdots\times\mathcal A_n}{no:I-2-2}{I, 2, 2}
\NotationEntry{\mathcal A_1\mathbin{\overline\otimes}\mathcal A_2\mathbin{\overline\otimes}\cdots\mathbin{\overline\otimes}\mathcal A_n}{no:I-2-4}{I, 2, 4}
\NotationEntry{\omega_{x,y}}{no:I-3-3}{I, 3, 3}
\NotationEntry{\omega_x}{no:I-4-1}{I, 4, 1}
\NotationEntry{\varphi\geq\psi}{no:I-4-1}{I, 4, 1}
\NotationEntry{x\mathbin{\cdot}f,\ f\mathbin{\cdot}x,\ f^*,\lvert f\rvert}{no:I-4-7}{I, 4, 7}
\NotationEntry{U_x,V_x}{no:I-5-1}{I, 5, 1}
\NotationEntry{\mathcal U(\mathfrak A),\mathcal V(\mathfrak A)}{no:I-5-1}{I, 5, 1}
\NotationEntry{\widetilde\varphi}{no:I-6-1}{I, 6, 1}
\NotationEntry{(S\mid T)}{no:I-6-2}{I, 6, 2}
\NotationEntry{\lVert S\rVert_2}{no:I-6-2}{I, 6, 2}
\NotationEntry{L^1(\varphi),L^1}{no:I-6-10}{I, 6, 10}
\NotationEntry{\lVert f\rVert_{1,\varphi},\lVert f\rVert_1}{no:I-6-10}{I, 6, 10}
\NotationEntry{L_\infty(Z)}{no:I-7-1}{I, 7, 1}
\NotationEntry{L_{\mathbb C}^\infty(Z,\nu)}{no:I-7-1}{I, 7, 1}
\NotationEntry{\int_{\mathfrak G}^\oplus\mathfrak H(\zeta)\,d\nu(\zeta)}{no:II-1-5}{II, 1, 5}
\NotationEntry{\int^\oplus\mathfrak H(\zeta)\,d\nu(\zeta)}{no:II-1-5}{II, 1, 5}
\NotationEntry{\int^\oplus x(\zeta)\,d\nu(\zeta)}{no:II-1-5}{II, 1, 5}
\NotationEntry{\int_Y^\oplus\mathfrak H(\zeta)\,d\nu(\zeta)}{no:II-1-5}{II, 1, 5}
\NotationEntry{\int^\oplus T(\zeta)\,d\nu(\zeta)}{no:II-2-3}{II, 2, 3}
\NotationEntry{\int^\oplus\mathcal A(\zeta)\,d\nu(\zeta)}{no:II-3-2}{II, 3, 2}
\NotationEntry{\int^\oplus\Phi_\zeta\,d\nu(\zeta)}{no:II-3-6}{II, 3, 6}
\NotationEntry{\int^\oplus\mathfrak A(\zeta)\,d\nu(\zeta)}{no:II-4-2}{II, 4, 2}
\NotationEntry{\int^\oplus\varphi_\zeta\,d\nu(\zeta)}{no:II-5-1}{II, 5, 1}
\NotationEntry{E\sim F}{no:III-1-1}{III, 1, 1}
\NotationEntry{E\prec F}{no:III-1-1}{III, 1, 1}
\NotationEntry{E\succ F}{no:III-1-1}{III, 1, 1}
\NotationEntry{\mathfrak X\sim\mathfrak Y}{no:III-1-1}{III, 1, 1}
\NotationEntry{\mathfrak X\prec\mathfrak Y}{no:III-1-1}{III, 1, 1}
\NotationEntry{\mathfrak X\succ\mathfrak Y}{no:III-1-1}{III, 1, 1}
\NotationEntry{\widehat{\mathcal Z}_+}{no:III-4-1}{III, 4, 1}
\NotationEntry{T^\natural}{no:III-4-4}{III, 4, 4}
\NotationEntry{f\mathbin{\cdot}T}{no:III-5-1}{III, 5, 1}
\NotationEntry{C_{\mathcal A}}{no:III-6-1}{III, 6, 1}
\end{multicols}
\endgroup

% Source PDF page 368 (printed p. 359)
\chapter*{Terminological Index}\label{back:terminological-index}\label{srcpage:359}
\addcontentsline{toc}{chapter}{Terminological Index}

The reference numbers indicate successively the chapter, the section, and the
number (or, exceptionally, the exercise).

\begingroup
\newcommand{\TermEntry}[3]{\par\noindent\hangindent=1em\hangafter=1 #1: \hyperref[#2]{#3}.}
\raggedright
\TermEntry{Abelian (projection)}{no:I-8-2}{I, 8, 2}
\TermEntry{Complete (Hilbert algebra)}{no:I-5-3}{I, 5, 3}
\TermEntry{Affiliated with a von Neumann algebra (operator)}{ex:I-1-10}{I, 1, Exercise 10}
\TermEntry{Von Neumann algebra}{no:I-1-1}{I, 1, 1}
\TermEntry{Von Neumann algebra generated by a set}{no:I-1-1}{I, 1, 1}
\TermEntry{Hilbert algebra}{no:I-5-1}{I, 5, 1}
\TermEntry{Quasi-Hilbert algebra}{no:I-5-1}{I, 5, 1}
\TermEntry{Algebraically anti-isomorphic}{no:I-1-5}{I, 1, 5}
\TermEntry{Algebraically isomorphic}{no:I-1-5}{I, 1, 5}
\TermEntry{Amplification}{no:I-2-4}{I, 2, 4}
\TermEntry{Antihomomorphism}{no:I-1-5}{I, 1, 5}
\TermEntry{Anti-isomorphic}{no:I-1-5}{I, 1, 5}
\TermEntry{Anti-isomorphism}{no:I-1-5}{I, 1, 5}
\TermEntry{Algebraic anti-isomorphism}{no:I-1-5}{I, 1, 5}
\TermEntry{Spatial anti-isomorphism}{no:I-1-5}{I, 1, 5}
\TermEntry{Canonical mapping from \(\mathfrak A\) into \(\mathcal U(\mathfrak A)\) or \(\mathcal V(\mathfrak A)\) (\(\mathfrak A\), a quasi-Hilbert algebra)}{no:I-5-1}{I, 5, 1}
\TermEntry{Left- or right-associated with a quasi-Hilbert algebra (von Neumann algebra)}{no:I-5-1}{I, 5, 1}
\TermEntry{Basic (measure)}{no:I-7-1}{I, 7, 1}
\TermEntry{Bicommutant}{no:I-1-1}{I, 1, 1}
\TermEntry{Binormal}{ex:III-3-3}{III, 3, Exercise 3}
\TermEntry{Borel}{no:III-1-1}{III, 1, 1}
\TermEntry{Bounded relative to a Hilbert algebra}{no:I-5-3}{I, 5, 3}
\TermEntry{Right- or left-bounded relative to a quasi-Hilbert algebra}{no:I-5-2}{I, 5, 2}
\TermEntry{Operator \(C^*\)-algebra}{no:I-3-4}{I, 3, 4}
\TermEntry{Character of a finite von Neumann algebra}{no:III-5-3}{III, 5, 3}
\TermEntry{Central relative to a Hilbert algebra}{no:I-5-4}{I, 5, 4}
\TermEntry{Borel field of Hilbert spaces}{no:II-1-3}{II, 1, 3}
\TermEntry{Constant field of Hilbert spaces}{no:II-1-3}{II, 1, 3}
\TermEntry{Continuous field of Hilbert spaces}{ex:II-1-5}{II, 1, Exercise 5}
\TermEntry{Field of Hilbert spaces}{no:II-1-2}{II, 1, 2}
\TermEntry{Field of vectors}{no:II-1-2}{II, 1, 2}
\TermEntry{Field of square-integrable vectors}{no:II-1-5}{II, 1, 5}
\TermEntry{Measurable field of vectors on a measurable subset}{no:II-1-3}{II, 1, 3}
\TermEntry{Measurable field of von Neumann algebras}{no:II-3-2}{II, 3, 2}
\TermEntry{Measurable field of Hilbert algebras}{no:II-4-1}{II, 4, 1}
\TermEntry{Measurable field of linear mappings}{no:II-2-1}{II, 2, 1}
\TermEntry{Measurable field of orthonormal bases}{no:II-1-4}{II, 1, 4}
\TermEntry{Measurable field of subspaces}{no:II-1-7}{II, 1, 7}
\TermEntry{Measurable field of Hilbert spaces}{no:II-1-3}{II, 1, 3}
\TermEntry{Measurable field of traces}{no:II-5-1}{II, 5, 1}
\TermEntry{Measurable field of vectors}{no:II-1-3}{II, 1, 3}
\TermEntry{Measurable field of homomorphisms}{no:II-3-6}{II, 3, 6}
\TermEntry{Measurable field of operators}{no:II-2-1}{II, 2, 1}
\TermEntry{Induced measurable field}{no:II-1-2}{II, 1, 2}
\TermEntry{Commutant}{no:I-1-1}{I, 1, 1}
\TermEntry{Continuous (von Neumann algebra)}{no:I-8-1}{I, 8, 1}
\TermEntry{Cyclic (projection)}{no:I-1-4}{I, 1, 4}
\TermEntry{Decomposable (von Neumann algebra)}{no:II-3-2}{II, 3, 2}
\TermEntry{Decomposable (Hilbert algebra)}{no:II-4-2}{II, 4, 2}
\TermEntry{Decomposable (linear mapping)}{no:II-2-3}{II, 2, 3}
\endgroup

% Source PDF page 369 (printed p. 360)
\phantomsection\label{srcpage:360}
\begingroup
\newcommand{\TermEntry}[3]{\par\noindent\hangindent=1em\hangafter=1 #1: \hyperref[#2]{#3}.}
\raggedright
\TermEntry{Decomposable (homomorphism)}{no:II-3-6}{II, 3, 6}
\TermEntry{Decomposable (operator)}{no:II-2-3}{II, 2, 3}
\par\noindent\hangindent=1em\hangafter=1 Polar decomposition:
\hyperref[no:I-4-7]{I, 4, 7} and \hyperref[app:III]{Appendix III}.
\TermEntry{Derivation}{no:III-9-1}{III, 9, 1}
\TermEntry{Determinant associated with a trace}{no:I-6-11}{I, 6, 11}
\TermEntry{Diagonalizable (operator)}{no:II-2-4}{II, 2, 4}
\TermEntry{Relative dimension}{no:III-2-7}{III, 2, 7}
\TermEntry{Discrete (von Neumann algebra)}{no:I-8-1}{I, 8, 1}
\TermEntry{Trace-element}{no:I-6-3}{I, 6, 3}
\TermEntry{Measurable envelope}{no:II-2-6}{II, 2, 6}
\TermEntry{Equivalent (projection)}{no:III-1-1}{III, 1, 1}
\TermEntry{Essentially bounded (measurable field of linear mappings)}{no:II-2-3}{II, 2, 3}
\TermEntry{Essentially dense}{ex:III-1-13}{III, 1, Exercise 13}
\TermEntry{Mutually singular (normal positive forms)}{no:I-4-6}{I, 4, 6}
\TermEntry{Factor}{no:I-1-1}{I, 1, 1}
\TermEntry{Weak (topology)}{no:I-3-1}{I, 3, 1}
\TermEntry{Faithful (positive linear form)}{no:I-4-1}{I, 4, 1}
\TermEntry{Faithful (trace)}{no:I-6-1}{I, 6, 1}
\TermEntry{Finite (projection)}{no:III-2-1}{III, 2, 1}
\TermEntry{Finite (von Neumann algebra)}{no:I-6-7}{I, 6, 7}
\TermEntry{Finite (trace)}{no:I-6-1}{I, 6, 1}
\TermEntry{Coupling operator}{no:III-6-1}{III, 6, 1}
\TermEntry{Fundamental (projection)}{no:III-8-2}{III, 8, 2}
\TermEntry{Strong (topology)}{no:I-3-1}{I, 3, 1}
\TermEntry{Of countable type (von Neumann algebra)}{no:I-1-2}{I, 1, 2}
\TermEntry{Of countable type over its centre (von Neumann algebra)}{ex:I-9-5}{I, 9, Exercise 5}
\TermEntry{Homogeneous (von Neumann algebra)}{no:III-3-1}{III, 3, 1}
\TermEntry{Homomorphism}{no:I-1-5}{I, 1, 5}
\TermEntry{Canonical homomorphism of \(L_{\mathbb C}^{\infty}(Z,\nu)\) [respectively, \(L_\infty(Z)\)] into the algebra of diagonalizable (respectively, continuously diagonalizable) operators}{no:II-2-4}{II, 2, 4}
\TermEntry{Hyperfinite}{no:III-7-2}{III, 7, 2}
\TermEntry{Induction}{no:I-2-1}{I, 2, 1}
\TermEntry{Induced (von Neumann algebra)}{no:I-2-1}{I, 2, 1}
\TermEntry{Induced (Borel structure)}{no:II-1-1}{II, 1, 1}
\TermEntry{Infinite (projection)}{no:III-2-1}{III, 2, 1}
\TermEntry{Infinite (von Neumann algebra)}{no:I-6-7}{I, 6, 7}
\TermEntry{Hilbert integral of a measurable field of Hilbert spaces}{no:II-1-5}{II, 1, 5}
\TermEntry{Involution canonically defined by a quasi-Hilbert algebra}{no:I-5-1}{I, 5, 1}
\TermEntry{Isomorphic}{no:I-1-5}{I, 1, 5}
\TermEntry{Isomorphism}{no:I-1-5}{I, 1, 5}
\TermEntry{Algebraic isomorphism}{no:I-1-5}{I, 1, 5}
\TermEntry{Canonical isomorphism of \(\int^\oplus\mathfrak H(\zeta)\,d\nu(\zeta)\) onto \(\int^\oplus\mathfrak H(\zeta)\,d\nu_1(\zeta)\)}{no:II-1-5}{II, 1, 5}
\TermEntry{Isomorphism of measurable fields}{no:II-1-3}{II, 1, 3}
\TermEntry{Gelfand isomorphism}{app:I}{Appendix I}
\TermEntry{Spatial isomorphism}{no:I-1-5}{I, 1, 5}
\TermEntry{Positive measure}{no:II-1-1}{II, 1, 1}
\TermEntry{Spectral measure}{app:I}{Appendix I}
\TermEntry{Minimal (projection)}{no:I-8-2}{I, 8, 2}
\TermEntry{Multiplicity}{no:III-3-1}{III, 3, 1}
\TermEntry{Normal (von Neumann algebra)}{ex:III-7-13}{III, 7, Exercise 13}
\TermEntry{Normal (\(\mathcal Z\)-trace)}{no:III-4-1}{III, 4, 1}
\TermEntry{Normal (positive linear mapping)}{no:I-4-3}{I, 4, 3}
\TermEntry{Normal (positive linear form)}{no:I-4-2}{I, 4, 2}
\TermEntry{Normal (trace)}{no:I-6-1}{I, 6, 1}
\TermEntry{Norm (topology)}{no:I-3-1}{I, 3, 1}
\TermEntry{Hilbert-Schmidt operator}{no:I-6-6}{I, 6, 6}
\TermEntry{Opposite (Hilbert algebra)}{no:I-5-5}{I, 5, 5}
\TermEntry{Order of a von Neumann algebra}{ex:III-3-5}{III, 3, Exercise 5}
\TermEntry{Partially isometric (operator)}{app:III}{Appendix III}
\TermEntry{Weight}{no:III-8-3}{III, 8, 3}
\TermEntry{Central weight}{no:III-8-3}{III, 8, 3}
\TermEntry{Positive (linear form)}{no:I-4-1}{I, 4, 1}
\TermEntry{Predual}{no:I-3-3}{I, 3, 3}
\TermEntry{Almost dominated (positive linear form)}{ex:I-4-8}{I, 4, Exercise 8}
\TermEntry{Product of von Neumann algebras}{no:I-2-2}{I, 2, 2}
\TermEntry{Tensor product of von Neumann algebras}{no:I-2-4}{I, 2, 4}
\TermEntry{Tensor product of Hilbert algebras}{no:I-5-5}{I, 5, 5}
\TermEntry{Characteristic projection}{no:I-5-4}{I, 5, 4}
\TermEntry{Final or initial projection}{app:III}{Appendix III}
\TermEntry{Properly infinite (projection)}{no:III-2-1}{III, 2, 1}
\TermEntry{Properly infinite (von Neumann algebra)}{no:I-6-7}{I, 6, 7}
\TermEntry{Property \(\Gamma\)}{no:III-7-7}{III, 7, 7}
\TermEntry{Purely infinite (projection)}{no:III-2-1}{III, 2, 1}
\TermEntry{Purely infinite (von Neumann algebra)}{no:I-6-7}{I, 6, 7}
\TermEntry{Reduced (homomorphism or antihomomorphism)}{no:I-2-1}{I, 2, 1}
\TermEntry{Reduced (von Neumann algebra)}{no:I-2-1}{I, 2, 1}
\TermEntry{Restricted (two-sided ideal)}{ex:I-1-6}{I, 1, Exercise 6}
\TermEntry{Semifinite (projection)}{no:III-2-1}{III, 2, 1}
\endgroup

% Source PDF page 370 (printed p. 361)
\phantomsection\label{srcpage:361}
\begingroup
\newcommand{\TermEntry}[3]{\par\noindent\hangindent=1em\hangafter=1 #1: \hyperref[#2]{#3}.}
\raggedright
\TermEntry{Semifinite (von Neumann algebra)}{no:I-6-7}{I, 6, 7}
\TermEntry{Semifinite (trace)}{no:I-6-1}{I, 6, 1}
\TermEntry{Separating (element)}{no:I-1-4}{I, 1, 4}
\TermEntry{Separating (set)}{no:I-1-4}{I, 1, 4}
\TermEntry{Direct sum (Hilbert algebra)}{no:I-5-5}{I, 5, 5}
\TermEntry{Final or initial subspace}{app:III}{Appendix III}
\TermEntry{Underlying (Borel structure)}{no:II-1-1}{II, 1, 1}
\TermEntry{Spatially anti-isomorphic}{no:I-1-5}{I, 1, 5}
\TermEntry{Spatially isomorphic}{no:I-1-5}{I, 1, 5}
\TermEntry{Spectrum}{app:I}{Appendix I}
\TermEntry{Standard (von Neumann algebra)}{no:I-5-5}{I, 5, 5}
\TermEntry{Standard (Borel space)}{no:II-1-1}{II, 1, 1}
\TermEntry{Standard (measure)}{no:II-1-1}{II, 1, 1}
\TermEntry{Fundamental sequence of measurable fields of vectors}{no:II-1-3}{II, 1, 3}
\TermEntry{Central support}{no:I-1-4}{I, 1, 4}
\TermEntry{Support of a normal positive linear form}{no:I-4-6}{I, 4, 6}
\TermEntry{Support of a normal trace}{no:I-6-1}{I, 6, 1}
\TermEntry{Support of an operator}{app:III}{Appendix III}
\TermEntry{Totalizing (element)}{no:I-1-4}{I, 1, 4}
\TermEntry{Totalizing (set)}{no:I-1-4}{I, 1, 4}
\TermEntry{Trace within \(\varepsilon\)}{no:III-8-4}{III, 8, 4}
\TermEntry{Canonical trace}{no:III-4-4}{III, 4, 4}
\TermEntry{Natural trace}{no:I-6-2}{I, 6, 2}
\TermEntry{Trace on \(\mathcal A_+\)}{no:I-6-1}{I, 6, 1}
\TermEntry{Trace on \(\widehat{\mathcal Z}_+\)}{no:III-4-2}{III, 4, 2}
\TermEntry{Trace on a two-sided ideal}{no:I-6-1}{I, 6, 1}
\par\noindent\hangindent=1em\hangafter=1 Type of a von Neumann algebra:
\hyperref[no:I-8-4]{I, 8, 4} and \hyperref[no:III-3-1]{III, 3, 1}.
\TermEntry{Ultra-weak (topology)}{no:I-3-1}{I, 3, 1}
\TermEntry{Ultra-strong (topology)}{no:I-3-1}{I, 3, 1}
\TermEntry{Uniform (von Neumann algebra)}{ex:III-3-5}{III, 3, Exercise 5}
\TermEntry{Absolute value of a linear form}{no:I-4-7}{I, 4, 7}
\par\noindent\hangindent=1em\hangafter=1 \(\mathcal Z\)-trace, canonical \(\mathcal Z\)-trace:
\hyperref[no:III-4-1]{III, 4, 1} and \hyperref[no:III-4-4]{III, 4, 4}.
\endgroup


% ===================== END MAIN TRANSLATION =====================

\end{document}
